% Complete English EGA IV4 editable text.
%
% This volume-only reader cannot provide destinations owned by EGA IV, Parts
% One--Three or by the other EGA volumes.  Keep their printed references, but
% suppress broken PDF links here; the cumulative/global reader retains links.
\PassOptionsToPackage{hyperfootnotes=false,draft}{hyperref}
% BEGIN NATIVE INPUT preamble
\documentclass[oneside]{amsart}

\usepackage[all]{xy}
\usepackage{tikz-cd}
\usepackage[T1]{fontenc}
\usepackage{xstring}
\usepackage{xparse}
\usepackage{xcolor}
\definecolor{brightmaroon}{rgb}{0.76, 0.13, 0.28}
\usepackage[linktocpage=true,colorlinks=true,hyperindex,hypertexnames=false,citecolor=blue,linkcolor=brightmaroon]{hyperref}
\usepackage[left=1.25in,right=1.25in,top=0.75in,bottom=0.75in]{geometry}
%\usepackage[charter,greekfamily=didot]{mathdesign}
%\usepackage{Baskervaldx}
\usepackage{amssymb}
\usepackage{mathrsfs}
\usepackage{mathpazo}
\linespread{1.05}

\usepackage[nobottomtitles]{titlesec}
\usepackage{marginnote}
\usepackage[shortlabels]{enumitem}
\usepackage{longtable}
\usepackage{multicol}
\usepackage{aurical}
\usepackage{microtype}

\newtheoremstyle{ega-env-style}%
  {}{}{\rmfamily}{}{\bfseries}{.}{ }{\thmnote{(#3)}}%

\newtheoremstyle{ega-thm-env-style}%
  {}{}{\itshape}{}{\bfseries}{. --- }{ }{\thmname{#1}\thmnote{ (#3)}}%

\newtheoremstyle{ega-defn-env-style}%
  {}{}{\rmfamily}{}{\bfseries}{. --- }{ }{\thmname{#1}\thmnote{ (#3)}}%

\theoremstyle{ega-env-style}
\newtheorem*{env}{---}

\theoremstyle{ega-thm-env-style}
\newtheorem*{theorem}{Theorem}
\newtheorem*{proposition}{Proposition}
\newtheorem*{lemma}{Lemma}
\newtheorem*{corollary}{Corollary}

\theoremstyle{ega-defn-env-style}
\newtheorem*{definition}{Definition}
\newtheorem*{example}{Example}
\newtheorem*{examples}{Examples}
\newtheorem*{remark}{Remark}
\newtheorem*{remarks}{Remarks}
\newtheorem*{scholium}{Scholium}
\newtheorem*{notation}{Notation}

% indent subsections, see https://tex.stackexchange.com/questions/177290/.
% also make section titles bigger.
% also add § to \thesection, https://tex.stackexchange.com/questions/119667/ and https://tex.stackexchange.com/questions/308737/.
\makeatletter
\def\l@subsection{\@tocline{2}{0pt}{2.5pc}{2.2pc}{}}
\def\section{\@startsection{section}{1}%
  \z@{.7\linespacing\@plus\linespacing}{.5\linespacing}%
  {\normalfont\bfseries\Large\scshape\centering}}
\renewcommand{\@seccntformat}[1]{%
  \ifnum\pdfstrcmp{#1}{section}=0\textsection\fi%
  \csname the#1\endcsname.~}
\makeatother

%\allowdisplaybreaks[1]
%\binoppenalty=9999
%\relpenalty=9999

% for Chapter 0, Chapter I, etc.
% credit for ZeroRoman https://tex.stackexchange.com/questions/211414/
% added into scripts/make_book.py
%\newcommand{\ZeroRoman}[1]{\ifcase\value{#1}\relax 0\else\Roman{#1}\fi}
%\renewcommand{\thechapter}{\ZeroRoman{chapter}}

\def\mathcal{\mathscr}
% BEGIN NATIVE INPUT preamble-base
\def\sh{\mathcal}                   % sheaf font
\def\bb{\mathbf}                    % bold font
\def\cat{\mathtt}                   % category font
\def\leq{\leqslant}                 % <=
\def\geq{\geqslant}                 % >=
\def\setmin{-}                      % set minus
\def\rad{\mathfrak{r}}              % radical
\def\nilrad{\mathfrak{N}}           % nilradical
\def\emp{\varnothing}               % empty set
\def\vphi{\phi}                     % for switching \phi and \varphi, change if needed
\def\HH{\mathrm{H}}                 % cohomology H
\def\CHH{\check{\HH}}               % Čech cohomology H
\def\RR{\mathrm{R}}                 % right derived R
\def\LL{\mathrm{L}}                 % left derived L
\def\dual#1{{#1}^\vee}              % dual
\def\kres{k}                        % residue field k
\def\C{\cat{C}}                     % category C
\def\op{^\cat{op}}                  % opposite category
\def\Set{\cat{Set}}                 % category of sets
\def\CHom{\cat{Hom}}                % functor category
\def\supertilde{{\,\widetilde{\,}\,}}   % use \supertilde instead of ^\sim
\def\GL{\bb{GL}}
\def\bbGamma{\boldsymbol\Gamma}
\def\red{\mathrm{red}}
\def\rg{\operatorname{rg}}
\def\gr{\operatorname{gr}}
\def\Gr{\operatorname{Gr}}
\def\Hom{\operatorname{Hom}}
\def\Proj{\operatorname{Proj}}
\def\Tor{\operatorname{Tor}}
\def\Ext{\operatorname{Ext}}
\def\Supp{\operatorname{Supp}}
\def\Ker{\operatorname{Ker}}
\def\Im{\operatorname{Im}}
\def\Coker{\operatorname{Coker}}
\def\Spec{\operatorname{Spec}}
\def\Spm{\operatorname{Spm}}
\def\Spf{\operatorname{Spf}}
\def\grad{\operatorname{grad}}
\def\dimc{\operatorname{dimc}}
\def\codim{\operatorname{codim}}
\def\id{\operatorname{id}}
\def\Der{\operatorname{Der}}
\def\Diff{\operatorname{Diff}}
\def\Hyp{\operatorname{Hyp}}
\def\Tr{\operatorname{Tr}}

\renewcommand{\to}{\mathchoice{\longrightarrow}{\rightarrow}{\rightarrow}{\rightarrow}}
\newcommand{\from}{\mathchoice{\longleftarrow}{\leftarrow}{\leftarrow}{\leftarrow}}
\let\mapstoo\mapsto
\renewcommand{\mapsto}{\mathchoice{\longmapsto}{\mapstoo}{\mapstoo}{\mapstoo}}
\def\isoto{\simeq} % isomorphism
\def\simto{\xrightarrow{\sim}}

\renewcommand{\hat}[1]{\widehat{#1}}
\renewcommand{\tilde}[1]{\widetilde{#1}}
% END NATIVE INPUT preamble-base

\def\shHom{\sh{H}\textup{\kern-2.2pt{\Fontauri\slshape om}}}   % sheaf Hom
\def\shProj{\sh{P}\textup{\kern-2.2pt{\Fontauri\slshape roj}}} % sheaf Proj
\def\shExt{\sh{E}\textup{\kern-2.2pt{\Fontauri\slshape xt}}}   % sheaf Ext
\def\shTor{\sh{T}\textup{\kern-2.2pt{\Fontauri\slshape or}}}   % sheaf Tor
\def\shGr{\sh{G}\textup{\kern-2.2pt{\Fontauri\slshape r}}}     % sheaf Gr
\def\shDer{\sh{D}\,\textup{\kern-2.2pt{\Fontauri\slshape er}}} % sheaf Der
\def\shDiff{\sh{D}\,\textup{\kern-2.2pt{\Fontauri\slshape if{}f}}\,} % sheaf Diff
\def\shHomcont{\sh{H}\textup{\kern-2.2pt{\Fontauri\slshape om.\,cont}}}   % sheaf Hom.cont
\def\shAut{\sh{A}\textup{\kern-2.2pt{\Fontauri\slshape ut}}}   % sheaf Aut

% if unsure of a translation
%\newcommand{\unsure}[2][]{\hl{#2}\marginpar{#1}}
%\newcommand{\completelyunsure}{\unsure{[\ldots]}}
\def\unsure#1{#1 {\color{red}(?)}}
\def\completelyunsure{{\color{red}(???)}}

% use to mark where original page starts
\newcommand{\oldpage}[2][]{{\marginnote{\normalfont{\textbf{#1}~|~#2}}}\ignorespaces}
\def\sectionbreak{\begin{center}***\end{center}}

% Give every source label its own zero-size PDF destination.  The inherited
% unnumbered theorem environments otherwise leave many labels pointing to the
% preceding subsection or display anchor rather than to their exact statement.
\let\EGAOriginalLabel\label
\renewcommand{\label}[1]{\phantomsection\EGAOriginalLabel{#1}}

% for referencing environments.
% use as \sref{chapter-number.x.y.z}, with optional args
% for volume and indices, e.g. \sref[volume]{chapter-number.x.y.z}[i].
\makeatletter
\newcommand{\EGAUnavailableHyperref}[2]{%
  \begingroup
  \let\@latex@warning\@gobble
  \let\@latex@warning@no@line\@gobble
  \let\G@refundefinedtrue\relax
  \hyperref[#1]{#2}%
  \endgroup
}
\NewDocumentCommand{\sref}{o m o}{%
  \@ifundefined{r@#2}%
    {\IfNoValueTF{#1}%
      {\IfNoValueTF{#3}%
        {\EGAUnavailableHyperref{#2}{\normalfont{(\StrBehind{#2}{.})}}}%
        {\EGAUnavailableHyperref{#2}{\normalfont{(\StrBehind{#2}{.},~{#3})}}}}%
      {\IfNoValueTF{#3}%
        {\EGAUnavailableHyperref{#2}{\normalfont{(\textbf{#1},~\StrBehind{#2}{.})}}}%
        {\EGAUnavailableHyperref{#2}{\normalfont{(\textbf{#1},~\StrBehind{#2}{.},~{#3})}}}}}%
    {\IfNoValueTF{#1}%
      {\IfNoValueTF{#3}%
        {\hyperref[#2]{\normalfont{(\StrBehind{#2}{.})}}}%
        {\hyperref[#2]{\normalfont{(\StrBehind{#2}{.},~{#3})}}}}%
      {\IfNoValueTF{#3}%
        {\hyperref[#2]{\normalfont{(\textbf{#1},~\StrBehind{#2}{.})}}}%
        {\hyperref[#2]{\normalfont{(\textbf{#1},~\StrBehind{#2}{.},~{#3})}}}}}%
}
\makeatother

% for marking changes following errata
% use as \erratum[volume]{correction} to say where the erratum is given
\newcommand{\erratum}[2][]{{#2}\marginpar{\footnotesize\textbf{Err}\textsubscript{#1}}}

% for referencing equations, use as \eref{eq:chapter-number.x.y.z}.
\newcommand{\eref}[1]{\hyperref[#1]{\normalfont{(\StrBehind{#1}{.})}}}
% END NATIVE INPUT preamble
% BEGIN NATIVE INPUT EGA4_reference_v2_registry
% Generated stable semantic destinations for the EGA IV reader.
% Visible text and numbering are unchanged.
\makeatletter
\newcommand{\sgaStableTarget}[1]{%
  \Hy@raisedlink{\hyper@anchorstart{#1}\hyper@anchorend}}
\expandafter\def\csname sga@target@IV.1.1.1\endcsname{ega4.c04.def.1-1-1}
\expandafter\def\csname sga@target@IV.1.1.2\endcsname{ega4.c04.prop.1-1-2}
\expandafter\def\csname sga@target@IV.1.1.3\endcsname{ega4.c04.prop.1-1-3}
\expandafter\def\csname sga@target@IV.1.1.4\endcsname{ega4.c04.cor.1-1-4}
\expandafter\def\csname sga@target@IV.1.1.5\endcsname{ega4.c04.prop.1-1-5}
\expandafter\def\csname sga@target@IV.1.1.6\endcsname{ega4.c04.prop.1-1-6}
\expandafter\def\csname sga@target@IV.1.10.1\endcsname{ega4.c04.thm.1-10-1}
\expandafter\def\csname sga@target@IV.1.10.2\endcsname{ega4.c04.cor.1-10-2}
\expandafter\def\csname sga@target@IV.1.10.3\endcsname{ega4.c04.prop.1-10-3}
\expandafter\def\csname sga@target@IV.1.10.4\endcsname{ega4.c04.cor.1-10-4}
\expandafter\def\csname sga@target@IV.1.2.1\endcsname{ega4.c04.def.1-2-1}
\expandafter\def\csname sga@target@IV.1.2.2\endcsname{ega4.c04.prop.1-2-2}
\expandafter\def\csname sga@target@IV.1.2.3\endcsname{ega4.c04.cor.1-2-3}
\expandafter\def\csname sga@target@IV.1.2.4\endcsname{ega4.c04.prop.1-2-4}
\expandafter\def\csname sga@target@IV.1.2.4.1\endcsname{ega4.c04.eq.1-2-4-1}
\expandafter\def\csname sga@target@IV.1.2.5\endcsname{ega4.c04.prop.1-2-5}
\expandafter\def\csname sga@target@IV.1.2.6\endcsname{ega4.c04.prop.1-2-6}
\expandafter\def\csname sga@target@IV.1.2.7\endcsname{ega4.c04.prop.1-2-7}
\expandafter\def\csname sga@target@IV.1.2.8\endcsname{ega4.c04.cor.1-2-8}
\expandafter\def\csname sga@target@IV.1.2.9\endcsname{ega4.c04.prop.1-2-9}
\expandafter\def\csname sga@target@IV.1.3.1\endcsname{ega4.c04.statement.1-3-1}
\expandafter\def\csname sga@target@IV.1.3.10\endcsname{ega4.c04.statement.1-3-10}
\expandafter\def\csname sga@target@IV.1.3.11\endcsname{ega4.c04.prop.1-3-11}
\expandafter\def\csname sga@target@IV.1.3.2\endcsname{ega4.c04.def.1-3-2}
\expandafter\def\csname sga@target@IV.1.3.3\endcsname{ega4.c04.prop.1-3-3}
\expandafter\def\csname sga@target@IV.1.3.4\endcsname{ega4.c04.prop.1-3-4}
\expandafter\def\csname sga@target@IV.1.3.5\endcsname{ega4.c04.cor.1-3-5}
\expandafter\def\csname sga@target@IV.1.3.6\endcsname{ega4.c04.prop.1-3-6}
\expandafter\def\csname sga@target@IV.1.3.7\endcsname{ega4.c04.prop.1-3-7}
\expandafter\def\csname sga@target@IV.1.3.8\endcsname{ega4.c04.statement.1-3-8}
\expandafter\def\csname sga@target@IV.1.3.9\endcsname{ega4.c04.prop.1-3-9}
\expandafter\def\csname sga@target@IV.1.4.1\endcsname{ega4.c04.statement.1-4-1}
\expandafter\def\csname sga@target@IV.1.4.2\endcsname{ega4.c04.def.1-4-2}
\expandafter\def\csname sga@target@IV.1.4.3\endcsname{ega4.c04.prop.1-4-3}
\expandafter\def\csname sga@target@IV.1.4.3.1\endcsname{ega4.c04.cor.1-4-3-1}
\expandafter\def\csname sga@target@IV.1.4.4\endcsname{ega4.c04.prop.1-4-4}
\expandafter\def\csname sga@target@IV.1.4.4.1\endcsname{ega4.c04.lem.1-4-4-1}
\expandafter\def\csname sga@target@IV.1.4.5\endcsname{ega4.c04.cor.1-4-5}
\expandafter\def\csname sga@target@IV.1.4.6\endcsname{ega4.c04.prop.1-4-6}
\expandafter\def\csname sga@target@IV.1.4.7\endcsname{ega4.c04.prop.1-4-7}
\expandafter\def\csname sga@target@IV.1.4.7.1\endcsname{ega4.c04.lem.1-4-7-1}
\expandafter\def\csname sga@target@IV.1.5.1\endcsname{ega4.c04.prop.1-5-1}
\expandafter\def\csname sga@target@IV.1.5.2\endcsname{ega4.c04.def.1-5-2}
\expandafter\def\csname sga@target@IV.1.5.3\endcsname{ega4.c04.prop.1-5-3}
\expandafter\def\csname sga@target@IV.1.5.4\endcsname{ega4.c04.prop.1-5-4}
\expandafter\def\csname sga@target@IV.1.5.5\endcsname{ega4.c04.cor.1-5-5}
\expandafter\def\csname sga@target@IV.1.5.6\endcsname{ega4.c04.cor.1-5-6}
\expandafter\def\csname sga@target@IV.1.5.7\endcsname{ega4.c04.prop.1-5-7}
\expandafter\def\csname sga@target@IV.1.6.1\endcsname{ega4.c04.def.1-6-1}
\expandafter\def\csname sga@target@IV.1.6.2\endcsname{ega4.c04.prop.1-6-2}
\expandafter\def\csname sga@target@IV.1.6.3\endcsname{ega4.c04.cor.1-6-3}
\expandafter\def\csname sga@target@IV.1.6.4\endcsname{ega4.c04.rem.1-6-4}
\expandafter\def\csname sga@target@IV.1.6.5\endcsname{ega4.c04.prop.1-6-5}
\expandafter\def\csname sga@target@IV.1.7.1\endcsname{ega4.c04.statement.1-7-1}
\expandafter\def\csname sga@target@IV.1.7.10\endcsname{ega4.c04.statement.1-7-10}
\expandafter\def\csname sga@target@IV.1.7.11\endcsname{ega4.c04.statement.1-7-11}
\expandafter\def\csname sga@target@IV.1.7.12\endcsname{ega4.c04.statement.1-7-12}
\expandafter\def\csname sga@target@IV.1.7.13\endcsname{ega4.c04.statement.1-7-13}
\expandafter\def\csname sga@target@IV.1.7.14\endcsname{ega4.c04.statement.1-7-14}
\expandafter\def\csname sga@target@IV.1.7.15\endcsname{ega4.c04.statement.1-7-15}
\expandafter\def\csname sga@target@IV.1.7.16\endcsname{ega4.c04.statement.1-7-16}
\expandafter\def\csname sga@target@IV.1.7.17\endcsname{ega4.c04.statement.1-7-17}
\expandafter\def\csname sga@target@IV.1.7.18\endcsname{ega4.c04.statement.1-7-18}
\expandafter\def\csname sga@target@IV.1.7.19\endcsname{ega4.c04.statement.1-7-19}
\expandafter\def\csname sga@target@IV.1.7.2\endcsname{ega4.c04.statement.1-7-2}
\expandafter\def\csname sga@target@IV.1.7.20\endcsname{ega4.c04.statement.1-7-20}
\expandafter\def\csname sga@target@IV.1.7.21\endcsname{ega4.c04.statement.1-7-21}
\expandafter\def\csname sga@target@IV.1.7.22\endcsname{ega4.c04.statement.1-7-22}
\expandafter\def\csname sga@target@IV.1.7.3\endcsname{ega4.c04.statement.1-7-3}
\expandafter\def\csname sga@target@IV.1.7.4\endcsname{ega4.c04.statement.1-7-4}
\expandafter\def\csname sga@target@IV.1.7.5\endcsname{ega4.c04.statement.1-7-5}
\expandafter\def\csname sga@target@IV.1.7.6\endcsname{ega4.c04.statement.1-7-6}
\expandafter\def\csname sga@target@IV.1.7.7\endcsname{ega4.c04.statement.1-7-7}
\expandafter\def\csname sga@target@IV.1.7.8\endcsname{ega4.c04.statement.1-7-8}
\expandafter\def\csname sga@target@IV.1.7.9\endcsname{ega4.c04.statement.1-7-9}
\expandafter\def\csname sga@target@IV.1.8.1\endcsname{ega4.c04.statement.1-8-1}
\expandafter\def\csname sga@target@IV.1.8.2\endcsname{ega4.c04.prop.1-8-2}
\expandafter\def\csname sga@target@IV.1.8.3\endcsname{ega4.c04.lem.1-8-3}
\expandafter\def\csname sga@target@IV.1.8.4\endcsname{ega4.c04.thm.1-8-4}
\expandafter\def\csname sga@target@IV.1.8.4.1\endcsname{ega4.c04.lem.1-8-4-1}
\expandafter\def\csname sga@target@IV.1.8.4.2\endcsname{ega4.c04.lem.1-8-4-2}
\expandafter\def\csname sga@target@IV.1.8.5\endcsname{ega4.c04.cor.1-8-5}
\expandafter\def\csname sga@target@IV.1.8.5.1\endcsname{ega4.c04.lem.1-8-5-1}
\expandafter\def\csname sga@target@IV.1.8.6\endcsname{ega4.c04.cor.1-8-6}
\expandafter\def\csname sga@target@IV.1.8.7\endcsname{ega4.c04.prop.1-8-7}
\expandafter\def\csname sga@target@IV.1.8.7.1\endcsname{ega4.c04.lem.1-8-7-1}
\expandafter\def\csname sga@target@IV.1.9.1\endcsname{ega4.c04.lem.1-9-1}
\expandafter\def\csname sga@target@IV.1.9.10\endcsname{ega4.c04.prop.1-9-10}
\expandafter\def\csname sga@target@IV.1.9.11\endcsname{ega4.c04.prop.1-9-11}
\expandafter\def\csname sga@target@IV.1.9.12\endcsname{ega4.c04.cor.1-9-12}
\expandafter\def\csname sga@target@IV.1.9.13\endcsname{ega4.c04.statement.1-9-13}
\expandafter\def\csname sga@target@IV.1.9.14\endcsname{ega4.c04.prop.1-9-14}
\expandafter\def\csname sga@target@IV.1.9.15\endcsname{ega4.c04.prop.1-9-15}
\expandafter\def\csname sga@target@IV.1.9.16\endcsname{ega4.c04.statement.1-9-16}
\expandafter\def\csname sga@target@IV.1.9.2\endcsname{ega4.c04.prop.1-9-2}
\expandafter\def\csname sga@target@IV.1.9.2.1\endcsname{ega4.c04.statement.1-9-2-1}
\expandafter\def\csname sga@target@IV.1.9.2.i\endcsname{ega4.c04.statement.1-9-2-i}
\expandafter\def\csname sga@target@IV.1.9.2.ii\endcsname{ega4.c04.statement.1-9-2-ii}
\expandafter\def\csname sga@target@IV.1.9.3\endcsname{ega4.c04.prop.1-9-3}
\expandafter\def\csname sga@target@IV.1.9.3.1\endcsname{ega4.c04.lem.1-9-3-1}
\expandafter\def\csname sga@target@IV.1.9.3.2\endcsname{ega4.c04.lem.1-9-3-2}
\expandafter\def\csname sga@target@IV.1.9.4\endcsname{ega4.c04.def.1-9-4}
\expandafter\def\csname sga@target@IV.1.9.5\endcsname{ega4.c04.prop.1-9-5}
\expandafter\def\csname sga@target@IV.1.9.6\endcsname{ega4.c04.statement.1-9-6}
\expandafter\def\csname sga@target@IV.1.9.7\endcsname{ega4.c04.thm.1-9-7}
\expandafter\def\csname sga@target@IV.1.9.8\endcsname{ega4.c04.cor.1-9-8}
\expandafter\def\csname sga@target@IV.1.9.9\endcsname{ega4.c04.cor.1-9-9}
\expandafter\def\csname sga@target@IV.10.1.1\endcsname{ega4.c04.statement.10-1-1}
\expandafter\def\csname sga@target@IV.10.1.2\endcsname{ega4.c04.prop.10-1-2}
\expandafter\def\csname sga@target@IV.10.1.3\endcsname{ega4.c04.def.10-1-3}
\expandafter\def\csname sga@target@IV.10.1.4\endcsname{ega4.c04.cor.10-1-4}
\expandafter\def\csname sga@target@IV.10.10.1\endcsname{ega4.c04.statement.10-10-1}
\expandafter\def\csname sga@target@IV.10.10.2\endcsname{ega4.c04.statement.10-10-2}
\expandafter\def\csname sga@target@IV.10.10.3\endcsname{ega4.c04.statement.10-10-3}
\expandafter\def\csname sga@target@IV.10.10.4\endcsname{ega4.c04.statement.10-10-4}
\expandafter\def\csname sga@target@IV.10.10.5\endcsname{ega4.c04.statement.10-10-5}
\expandafter\def\csname sga@target@IV.10.2.1\endcsname{ega4.c04.prop.10-2-1}
\expandafter\def\csname sga@target@IV.10.2.2\endcsname{ega4.c04.def.10-2-2}
\expandafter\def\csname sga@target@IV.10.2.3\endcsname{ega4.c04.cor.10-2-3}
\expandafter\def\csname sga@target@IV.10.2.4\endcsname{ega4.c04.cor.10-2-4}
\expandafter\def\csname sga@target@IV.10.2.5\endcsname{ega4.c04.cor.10-2-5}
\expandafter\def\csname sga@target@IV.10.2.6\endcsname{ega4.c04.cor.10-2-6}
\expandafter\def\csname sga@target@IV.10.2.7\endcsname{ega4.c04.prop.10-2-7}
\expandafter\def\csname sga@target@IV.10.2.8\endcsname{ega4.c04.item.10-2-8}
\expandafter\def\csname sga@target@IV.10.2.8.1\endcsname{ega4.c04.eq.10-2-8-1}
\expandafter\def\csname sga@target@IV.10.3.1\endcsname{ega4.c04.def.10-3-1}
\expandafter\def\csname sga@target@IV.10.3.2\endcsname{ega4.c04.prop.10-3-2}
\expandafter\def\csname sga@target@IV.10.3.3\endcsname{ega4.c04.prop.10-3-3}
\expandafter\def\csname sga@target@IV.10.4.1\endcsname{ega4.c04.def.10-4-1}
\expandafter\def\csname sga@target@IV.10.4.10\endcsname{ega4.c04.statement.10-4-10}
\expandafter\def\csname sga@target@IV.10.4.11\endcsname{ega4.c04.statement.10-4-11}
\expandafter\def\csname sga@target@IV.10.4.11.1\endcsname{ega4.c04.lem.10-4-11-1}
\expandafter\def\csname sga@target@IV.10.4.12\endcsname{ega4.c04.cor.10-4-12}
\expandafter\def\csname sga@target@IV.10.4.12.1\endcsname{ega4.c04.sec.10-4-12-1}
\expandafter\def\csname sga@target@IV.10.4.12.1.errata\endcsname{ega4.c04.eq.10-4-12-1-errata}
\expandafter\def\csname sga@target@IV.10.4.13\endcsname{ega4.c04.cor.10-4-13}
\expandafter\def\csname sga@target@IV.10.4.14\endcsname{ega4.c04.rem.10-4-14}
\expandafter\def\csname sga@target@IV.10.4.2\endcsname{ega4.c04.prop.10-4-2}
\expandafter\def\csname sga@target@IV.10.4.3\endcsname{ega4.c04.statement.10-4-3}
\expandafter\def\csname sga@target@IV.10.4.4\endcsname{ega4.c04.prop.10-4-4}
\expandafter\def\csname sga@target@IV.10.4.5\endcsname{ega4.c04.prop.10-4-5}
\expandafter\def\csname sga@target@IV.10.4.5.1\endcsname{ega4.c04.lem.10-4-5-1}
\expandafter\def\csname sga@target@IV.10.4.6\endcsname{ega4.c04.cor.10-4-6}
\expandafter\def\csname sga@target@IV.10.4.7\endcsname{ega4.c04.cor.10-4-7}
\expandafter\def\csname sga@target@IV.10.4.8\endcsname{ega4.c04.cor.10-4-8}
\expandafter\def\csname sga@target@IV.10.4.9\endcsname{ega4.c04.statement.10-4-9}
\expandafter\def\csname sga@target@IV.10.5.1\endcsname{ega4.c04.prop.10-5-1}
\expandafter\def\csname sga@target@IV.10.5.10\endcsname{ega4.c04.prop.10-5-10}
\expandafter\def\csname sga@target@IV.10.5.2\endcsname{ega4.c04.cor.10-5-2}
\expandafter\def\csname sga@target@IV.10.5.3\endcsname{ega4.c04.cor.10-5-3}
\expandafter\def\csname sga@target@IV.10.5.4\endcsname{ega4.c04.rem.10-5-4}
\expandafter\def\csname sga@target@IV.10.5.5\endcsname{ega4.c04.prop.10-5-5}
\expandafter\def\csname sga@target@IV.10.5.6\endcsname{ega4.c04.rem.10-5-6}
\expandafter\def\csname sga@target@IV.10.5.7\endcsname{ega4.c04.prop.10-5-7}
\expandafter\def\csname sga@target@IV.10.5.8\endcsname{ega4.c04.cor.10-5-8}
\expandafter\def\csname sga@target@IV.10.5.9\endcsname{ega4.c04.cor.10-5-9}
\expandafter\def\csname sga@target@IV.10.6.1\endcsname{ega4.c04.prop.10-6-1}
\expandafter\def\csname sga@target@IV.10.6.1.1\endcsname{ega4.c04.eq.10-6-1-1}
\expandafter\def\csname sga@target@IV.10.6.1.2\endcsname{ega4.c04.eq.10-6-1-2}
\expandafter\def\csname sga@target@IV.10.6.1.3\endcsname{ega4.c04.eq.10-6-1-3}
\expandafter\def\csname sga@target@IV.10.6.1.4\endcsname{ega4.c04.eq.10-6-1-4}
\expandafter\def\csname sga@target@IV.10.6.2\endcsname{ega4.c04.cor.10-6-2}
\expandafter\def\csname sga@target@IV.10.6.3\endcsname{ega4.c04.prop.10-6-3}
\expandafter\def\csname sga@target@IV.10.6.3.1\endcsname{ega4.c04.eq.10-6-3-1}
\expandafter\def\csname sga@target@IV.10.6.3.2\endcsname{ega4.c04.eq.10-6-3-2}
\expandafter\def\csname sga@target@IV.10.6.3.3\endcsname{ega4.c04.eq.10-6-3-3}
\expandafter\def\csname sga@target@IV.10.6.4\endcsname{ega4.c04.cor.10-6-4}
\expandafter\def\csname sga@target@IV.10.6.4.1\endcsname{ega4.c04.eq.10-6-4-1}
\expandafter\def\csname sga@target@IV.10.7.1\endcsname{ega4.c04.statement.10-7-1}
\expandafter\def\csname sga@target@IV.10.7.2\endcsname{ega4.c04.statement.10-7-2}
\expandafter\def\csname sga@target@IV.10.7.3\endcsname{ega4.c04.statement.10-7-3}
\expandafter\def\csname sga@target@IV.10.8.1\endcsname{ega4.c04.def.10-8-1}
\expandafter\def\csname sga@target@IV.10.8.1.1\endcsname{ega4.c04.eq.10-8-1-1}
\expandafter\def\csname sga@target@IV.10.8.1.2\endcsname{ega4.c04.eq.10-8-1-2}
\expandafter\def\csname sga@target@IV.10.8.2\endcsname{ega4.c04.statement.10-8-2}
\expandafter\def\csname sga@target@IV.10.8.3\endcsname{ega4.c04.lem.10-8-3}
\expandafter\def\csname sga@target@IV.10.8.3.1\endcsname{ega4.c04.eq.10-8-3-1}
\expandafter\def\csname sga@target@IV.10.8.4\endcsname{ega4.c04.cor.10-8-4}
\expandafter\def\csname sga@target@IV.10.8.5\endcsname{ega4.c04.prop.10-8-5}
\expandafter\def\csname sga@target@IV.10.8.6\endcsname{ega4.c04.cor.10-8-6}
\expandafter\def\csname sga@target@IV.10.8.7\endcsname{ega4.c04.prop.10-8-7}
\expandafter\def\csname sga@target@IV.10.8.7.1\endcsname{ega4.c04.eq.10-8-7-1}
\expandafter\def\csname sga@target@IV.10.8.8\endcsname{ega4.c04.cor.10-8-8}
\expandafter\def\csname sga@target@IV.10.9.1\endcsname{ega4.c04.statement.10-9-1}
\expandafter\def\csname sga@target@IV.10.9.2\endcsname{ega4.c04.statement.10-9-2}
\expandafter\def\csname sga@target@IV.10.9.3\endcsname{ega4.c04.statement.10-9-3}
\expandafter\def\csname sga@target@IV.10.9.4\endcsname{ega4.c04.prop.10-9-4}
\expandafter\def\csname sga@target@IV.10.9.5\endcsname{ega4.c04.statement.10-9-5}
\expandafter\def\csname sga@target@IV.10.9.6\endcsname{ega4.c04.prop.10-9-6}
\expandafter\def\csname sga@target@IV.11.1.1\endcsname{ega4.c04.thm.11-1-1}
\expandafter\def\csname sga@target@IV.11.1.1.1\endcsname{ega4.c04.sec.11-1-1-1}
\expandafter\def\csname sga@target@IV.11.1.1.2\endcsname{ega4.c04.lem.11-1-1-2}
\expandafter\def\csname sga@target@IV.11.1.2\endcsname{ega4.c04.cor.11-1-2}
\expandafter\def\csname sga@target@IV.11.1.3\endcsname{ega4.c04.cor.11-1-3}
\expandafter\def\csname sga@target@IV.11.1.4\endcsname{ega4.c04.cor.11-1-4}
\expandafter\def\csname sga@target@IV.11.1.5\endcsname{ega4.c04.cor.11-1-5}
\expandafter\def\csname sga@target@IV.11.1.6\endcsname{ega4.c04.cor.11-1-6}
\expandafter\def\csname sga@target@IV.11.10.1\endcsname{ega4.c04.prop.11-10-1}
\expandafter\def\csname sga@target@IV.11.10.1.1\endcsname{ega4.c04.eq.11-10-1-1}
\expandafter\def\csname sga@target@IV.11.10.1.2\endcsname{ega4.c04.eq.11-10-1-2}
\expandafter\def\csname sga@target@IV.11.10.10\endcsname{ega4.c04.prop.11-10-10}
\expandafter\def\csname sga@target@IV.11.10.2\endcsname{ega4.c04.def.11-10-2}
\expandafter\def\csname sga@target@IV.11.10.3\endcsname{ega4.c04.item.11-10-3}
\expandafter\def\csname sga@target@IV.11.10.4\endcsname{ega4.c04.prop.11-10-4}
\expandafter\def\csname sga@target@IV.11.10.5\endcsname{ega4.c04.thm.11-10-5}
\expandafter\def\csname sga@target@IV.11.10.6\endcsname{ega4.c04.prop.11-10-6}
\expandafter\def\csname sga@target@IV.11.10.7\endcsname{ega4.c04.cor.11-10-7}
\expandafter\def\csname sga@target@IV.11.10.8\endcsname{ega4.c04.def.11-10-8}
\expandafter\def\csname sga@target@IV.11.10.9\endcsname{ega4.c04.thm.11-10-9}
\expandafter\def\csname sga@target@IV.11.2.1\endcsname{ega4.c04.statement.11-2-1}
\expandafter\def\csname sga@target@IV.11.2.1.1\endcsname{ega4.c04.eq.11-2-1-1}
\expandafter\def\csname sga@target@IV.11.2.1.2\endcsname{ega4.c04.eq.11-2-1-2}
\expandafter\def\csname sga@target@IV.11.2.1.3\endcsname{ega4.c04.eq.11-2-1-3}
\expandafter\def\csname sga@target@IV.11.2.1.4\endcsname{ega4.c04.eq.11-2-1-4}
\expandafter\def\csname sga@target@IV.11.2.1.5\endcsname{ega4.c04.eq.11-2-1-5}
\expandafter\def\csname sga@target@IV.11.2.10\endcsname{ega4.c04.rem.11-2-10}
\expandafter\def\csname sga@target@IV.11.2.2\endcsname{ega4.c04.statement.11-2-2}
\expandafter\def\csname sga@target@IV.11.2.2.1\endcsname{ega4.c04.eq.11-2-2-1}
\expandafter\def\csname sga@target@IV.11.2.2.2\endcsname{ega4.c04.lem.11-2-2-2}
\expandafter\def\csname sga@target@IV.11.2.3\endcsname{ega4.c04.statement.11-2-3}
\expandafter\def\csname sga@target@IV.11.2.3.1\endcsname{ega4.c04.eq.11-2-3-1}
\expandafter\def\csname sga@target@IV.11.2.3.2\endcsname{ega4.c04.eq.11-2-3-2}
\expandafter\def\csname sga@target@IV.11.2.3.3\endcsname{ega4.c04.eq.11-2-3-3}
\expandafter\def\csname sga@target@IV.11.2.4\endcsname{ega4.c04.lem.11-2-4}
\expandafter\def\csname sga@target@IV.11.2.5\endcsname{ega4.c04.lem.11-2-5}
\expandafter\def\csname sga@target@IV.11.2.6\endcsname{ega4.c04.thm.11-2-6}
\expandafter\def\csname sga@target@IV.11.2.6.1\endcsname{ega4.c04.cor.11-2-6-1}
\expandafter\def\csname sga@target@IV.11.2.7\endcsname{ega4.c04.cor.11-2-7}
\expandafter\def\csname sga@target@IV.11.2.8\endcsname{ega4.c04.prop.11-2-8}
\expandafter\def\csname sga@target@IV.11.2.9\endcsname{ega4.c04.thm.11-2-9}
\expandafter\def\csname sga@target@IV.11.2.9.1\endcsname{ega4.c04.statement.11-2-9-1}
\expandafter\def\csname sga@target@IV.11.2.9.2\endcsname{ega4.c04.lem.11-2-9-2}
\expandafter\def\csname sga@target@IV.11.2.9.3\endcsname{ega4.c04.lem.11-2-9-3}
\expandafter\def\csname sga@target@IV.11.2.9.4\endcsname{ega4.c04.cor.11-2-9-4}
\expandafter\def\csname sga@target@IV.11.2.9.5\endcsname{ega4.c04.lem.11-2-9-5}
\expandafter\def\csname sga@target@IV.11.2.9.6\endcsname{ega4.c04.lem.11-2-9-6}
\expandafter\def\csname sga@target@IV.11.2.9.7\endcsname{ega4.c04.eq.11-2-9-7}
\expandafter\def\csname sga@target@IV.11.2.9.8\endcsname{ega4.c04.eq.11-2-9-8}
\expandafter\def\csname sga@target@IV.11.2.9.9\endcsname{ega4.c04.eq.11-2-9-9}
\expandafter\def\csname sga@target@IV.11.3\endcsname{ega4.c04.subsec.11-3.key-iv-11-3}
\expandafter\def\csname sga@target@IV.11.3.1\endcsname{ega4.c04.thm.11-3-1}
\expandafter\def\csname sga@target@IV.11.3.10\endcsname{ega4.c04.thm.11-3-10}
\expandafter\def\csname sga@target@IV.11.3.10.1\endcsname{ega4.c04.lem.11-3-10-1}
\expandafter\def\csname sga@target@IV.11.3.10.2\endcsname{ega4.c04.lem.11-3-10-2}
\expandafter\def\csname sga@target@IV.11.3.10.3\endcsname{ega4.c04.eq.11-3-10-3}
\expandafter\def\csname sga@target@IV.11.3.10.4\endcsname{ega4.c04.sec.11-3-10-4}
\expandafter\def\csname sga@target@IV.11.3.11\endcsname{ega4.c04.cor.11-3-11}
\expandafter\def\csname sga@target@IV.11.3.12\endcsname{ega4.c04.rem.11-3-12}
\expandafter\def\csname sga@target@IV.11.3.13\endcsname{ega4.c04.prop.11-3-13}
\expandafter\def\csname sga@target@IV.11.3.14\endcsname{ega4.c04.cor.11-3-14}
\expandafter\def\csname sga@target@IV.11.3.15\endcsname{ega4.c04.prop.11-3-15}
\expandafter\def\csname sga@target@IV.11.3.16\endcsname{ega4.c04.prop.11-3-16}
\expandafter\def\csname sga@target@IV.11.3.17\endcsname{ega4.c04.cor.11-3-17}
\expandafter\def\csname sga@target@IV.11.3.2\endcsname{ega4.c04.cor.11-3-2}
\expandafter\def\csname sga@target@IV.11.3.3\endcsname{ega4.c04.cor.11-3-3}
\expandafter\def\csname sga@target@IV.11.3.4\endcsname{ega4.c04.prop.11-3-4}
\expandafter\def\csname sga@target@IV.11.3.5\endcsname{ega4.c04.cor.11-3-5}
\expandafter\def\csname sga@target@IV.11.3.6\endcsname{ega4.c04.prop.11-3-6}
\expandafter\def\csname sga@target@IV.11.3.7\endcsname{ega4.c04.prop.11-3-7}
\expandafter\def\csname sga@target@IV.11.3.8\endcsname{ega4.c04.thm.11-3-8}
\expandafter\def\csname sga@target@IV.11.3.9\endcsname{ega4.c04.prop.11-3-9}
\expandafter\def\csname sga@target@IV.11.3.9.1\endcsname{ega4.c04.lem.11-3-9-1}
\expandafter\def\csname sga@target@IV.11.4.1\endcsname{ega4.c04.thm.11-4-1}
\expandafter\def\csname sga@target@IV.11.4.1.1\endcsname{ega4.c04.lem.11-4-1-1}
\expandafter\def\csname sga@target@IV.11.4.1.2\endcsname{ega4.c04.lem.11-4-1-2}
\expandafter\def\csname sga@target@IV.11.4.10\endcsname{ega4.c04.cor.11-4-10}
\expandafter\def\csname sga@target@IV.11.4.11\endcsname{ega4.c04.prop.11-4-11}
\expandafter\def\csname sga@target@IV.11.4.12\endcsname{ega4.c04.prop.11-4-12}
\expandafter\def\csname sga@target@IV.11.4.13\endcsname{ega4.c04.cor.11-4-13}
\expandafter\def\csname sga@target@IV.11.4.2\endcsname{ega4.c04.rem.11-4-2}
\expandafter\def\csname sga@target@IV.11.4.3\endcsname{ega4.c04.cor.11-4-3}
\expandafter\def\csname sga@target@IV.11.4.4\endcsname{ega4.c04.cor.11-4-4}
\expandafter\def\csname sga@target@IV.11.4.5\endcsname{ega4.c04.prop.11-4-5}
\expandafter\def\csname sga@target@IV.11.4.6\endcsname{ega4.c04.cor.11-4-6}
\expandafter\def\csname sga@target@IV.11.4.7\endcsname{ega4.c04.prop.11-4-7}
\expandafter\def\csname sga@target@IV.11.4.8\endcsname{ega4.c04.rem.11-4-8}
\expandafter\def\csname sga@target@IV.11.4.9\endcsname{ega4.c04.cor.11-4-9}
\expandafter\def\csname sga@target@IV.11.5.1\endcsname{ega4.c04.thm.11-5-1}
\expandafter\def\csname sga@target@IV.11.5.2\endcsname{ega4.c04.cor.11-5-2}
\expandafter\def\csname sga@target@IV.11.5.3\endcsname{ega4.c04.prop.11-5-3}
\expandafter\def\csname sga@target@IV.11.5.4\endcsname{ega4.c04.cor.11-5-4}
\expandafter\def\csname sga@target@IV.11.5.5\endcsname{ega4.c04.prop.11-5-5}
\expandafter\def\csname sga@target@IV.11.5.5.1\endcsname{ega4.c04.lem.11-5-5-1}
\expandafter\def\csname sga@target@IV.11.5.5.2\endcsname{ega4.c04.lem.11-5-5-2}
\expandafter\def\csname sga@target@IV.11.6.1\endcsname{ega4.c04.thm.11-6-1}
\expandafter\def\csname sga@target@IV.11.6.2\endcsname{ega4.c04.cor.11-6-2}
\expandafter\def\csname sga@target@IV.11.7.1\endcsname{ega4.c04.statement.11-7-1}
\expandafter\def\csname sga@target@IV.11.7.2\endcsname{ega4.c04.statement.11-7-2}
\expandafter\def\csname sga@target@IV.11.7.3\endcsname{ega4.c04.statement.11-7-3}
\expandafter\def\csname sga@target@IV.11.7.4\endcsname{ega4.c04.statement.11-7-4}
\expandafter\def\csname sga@target@IV.11.7.5\endcsname{ega4.c04.statement.11-7-5}
\expandafter\def\csname sga@target@IV.11.8.1\endcsname{ega4.c04.thm.11-8-1}
\expandafter\def\csname sga@target@IV.11.8.2\endcsname{ega4.c04.subsec.11-8-2}
\expandafter\def\csname sga@target@IV.11.9.1\endcsname{ega4.c04.statement.11-9-1}
\expandafter\def\csname sga@target@IV.11.9.1.1\endcsname{ega4.c04.statement.11-9-1-1}
\expandafter\def\csname sga@target@IV.11.9.10\endcsname{ega4.c04.thm.11-9-10}
\expandafter\def\csname sga@target@IV.11.9.11\endcsname{ega4.c04.rem.11-9-11}
\expandafter\def\csname sga@target@IV.11.9.12\endcsname{ega4.c04.lem.11-9-12}
\expandafter\def\csname sga@target@IV.11.9.13\endcsname{ega4.c04.prop.11-9-13}
\expandafter\def\csname sga@target@IV.11.9.14\endcsname{ega4.c04.statement.11-9-14}
\expandafter\def\csname sga@target@IV.11.9.15\endcsname{ega4.c04.prop.11-9-15}
\expandafter\def\csname sga@target@IV.11.9.16\endcsname{ega4.c04.thm.11-9-16}
\expandafter\def\csname sga@target@IV.11.9.17\endcsname{ega4.c04.prop.11-9-17}
\expandafter\def\csname sga@target@IV.11.9.17.1\endcsname{ega4.c04.lem.11-9-17-1}
\expandafter\def\csname sga@target@IV.11.9.18\endcsname{ega4.c04.item.11-9-18}
\expandafter\def\csname sga@target@IV.11.9.2\endcsname{ega4.c04.statement.11-9-2}
\expandafter\def\csname sga@target@IV.11.9.3\endcsname{ega4.c04.statement.11-9-3}
\expandafter\def\csname sga@target@IV.11.9.4\endcsname{ega4.c04.statement.11-9-4}
\expandafter\def\csname sga@target@IV.11.9.5\endcsname{ega4.c04.statement.11-9-5}
\expandafter\def\csname sga@target@IV.11.9.6\endcsname{ega4.c04.statement.11-9-6}
\expandafter\def\csname sga@target@IV.11.9.7\endcsname{ega4.c04.prop.11-9-7}
\expandafter\def\csname sga@target@IV.11.9.7.1\endcsname{ega4.c04.lem.11-9-7-1}
\expandafter\def\csname sga@target@IV.11.9.8\endcsname{ega4.c04.subsec.11-9-8}
\expandafter\def\csname sga@target@IV.11.9.9\endcsname{ega4.c04.cor.11-9-9}
\expandafter\def\csname sga@target@IV.12.0.1\endcsname{ega4.c04.statement.12-0-1}
\expandafter\def\csname sga@target@IV.12.0.2\endcsname{ega4.c04.statement.12-0-2}
\expandafter\def\csname sga@target@IV.12.1.1\endcsname{ega4.c04.thm.12-1-1}
\expandafter\def\csname sga@target@IV.12.1.1.1\endcsname{ega4.c04.statement.12-1-1-1}
\expandafter\def\csname sga@target@IV.12.1.1.2\endcsname{ega4.c04.statement.12-1-1-2}
\expandafter\def\csname sga@target@IV.12.1.1.3\endcsname{ega4.c04.statement.12-1-1-3}
\expandafter\def\csname sga@target@IV.12.1.1.4\endcsname{ega4.c04.statement.12-1-1-4}
\expandafter\def\csname sga@target@IV.12.1.1.5\endcsname{ega4.c04.prop.12-1-1-5}
\expandafter\def\csname sga@target@IV.12.1.2\endcsname{ega4.c04.item.12-1-2}
\expandafter\def\csname sga@target@IV.12.1.3\endcsname{ega4.c04.cor.12-1-3}
\expandafter\def\csname sga@target@IV.12.1.4\endcsname{ega4.c04.rem.12-1-4}
\expandafter\def\csname sga@target@IV.12.1.5\endcsname{ega4.c04.cor.12-1-5}
\expandafter\def\csname sga@target@IV.12.1.6\endcsname{ega4.c04.thm.12-1-6}
\expandafter\def\csname sga@target@IV.12.1.7\endcsname{ega4.c04.cor.12-1-7}
\expandafter\def\csname sga@target@IV.12.1.8\endcsname{ega4.c04.item.12-1-8}
\expandafter\def\csname sga@target@IV.12.2.1\endcsname{ega4.c04.thm.12-2-1}
\expandafter\def\csname sga@target@IV.12.2.1.1\endcsname{ega4.c04.statement.12-2-1-1}
\expandafter\def\csname sga@target@IV.12.2.1.2\endcsname{ega4.c04.statement.12-2-1-2}
\expandafter\def\csname sga@target@IV.12.2.1.3\endcsname{ega4.c04.statement.12-2-1-3}
\expandafter\def\csname sga@target@IV.12.2.1.4\endcsname{ega4.c04.lem.12-2-1-4}
\expandafter\def\csname sga@target@IV.12.2.2\endcsname{ega4.c04.cor.12-2-2}
\expandafter\def\csname sga@target@IV.12.2.3\endcsname{ega4.c04.item.12-2-3}
\expandafter\def\csname sga@target@IV.12.2.4\endcsname{ega4.c04.thm.12-2-4}
\expandafter\def\csname sga@target@IV.12.3.1\endcsname{ega4.c04.lem.12-3-1}
\expandafter\def\csname sga@target@IV.12.3.2\endcsname{ega4.c04.prop.12-3-2}
\expandafter\def\csname sga@target@IV.12.3.2.1\endcsname{ega4.c04.eq.12-3-2-1}
\expandafter\def\csname sga@target@IV.12.3.3\endcsname{ega4.c04.prop.12-3-3}
\expandafter\def\csname sga@target@IV.12.3.3.1\endcsname{ega4.c04.lem.12-3-3-1}
\expandafter\def\csname sga@target@IV.12.3.3.2\endcsname{ega4.c04.eq.12-3-3-2}
\expandafter\def\csname sga@target@IV.12.3.3.3\endcsname{ega4.c04.lem.12-3-3-3}
\expandafter\def\csname sga@target@IV.12.3.3.4\endcsname{ega4.c04.eq.12-3-3-4}
\expandafter\def\csname sga@target@IV.12.3.3.5\endcsname{ega4.c04.lem.12-3-3-5}
\expandafter\def\csname sga@target@IV.12.3.4\endcsname{ega4.c04.cor.12-3-4}
\expandafter\def\csname sga@target@IV.13.1.1\endcsname{ega4.c04.lem.13-1-1}
\expandafter\def\csname sga@target@IV.13.1.2\endcsname{ega4.c04.lem.13-1-2}
\expandafter\def\csname sga@target@IV.13.1.3\endcsname{ega4.c04.thm.13-1-3}
\expandafter\def\csname sga@target@IV.13.1.4\endcsname{ega4.c04.cor.13-1-4}
\expandafter\def\csname sga@target@IV.13.1.5\endcsname{ega4.c04.cor.13-1-5}
\expandafter\def\csname sga@target@IV.13.1.6\endcsname{ega4.c04.cor.13-1-6}
\expandafter\def\csname sga@target@IV.13.1.7\endcsname{ega4.c04.prop.13-1-7}
\expandafter\def\csname sga@target@IV.13.2.1\endcsname{ega4.c04.statement.13-2-1}
\expandafter\def\csname sga@target@IV.13.2.10\endcsname{ega4.c04.prop.13-2-10}
\expandafter\def\csname sga@target@IV.13.2.10.1\endcsname{ega4.c04.eq.13-2-10-1}
\expandafter\def\csname sga@target@IV.13.2.10.2\endcsname{ega4.c04.eq.13-2-10-2}
\expandafter\def\csname sga@target@IV.13.2.11\endcsname{ega4.c04.cor.13-2-11}
\expandafter\def\csname sga@target@IV.13.2.12\endcsname{ega4.c04.item.13-2-12}
\expandafter\def\csname sga@target@IV.13.2.2\endcsname{ega4.c04.def.13-2-2}
\expandafter\def\csname sga@target@IV.13.2.3\endcsname{ega4.c04.prop.13-2-3}
\expandafter\def\csname sga@target@IV.13.2.3.1\endcsname{ega4.c04.prop.13-2-3-1}
\expandafter\def\csname sga@target@IV.13.2.4\endcsname{ega4.c04.statement.13-2-4}
\expandafter\def\csname sga@target@IV.13.2.4.1\endcsname{ega4.c04.statement.13-2-4-1}
\expandafter\def\csname sga@target@IV.13.2.4.2\endcsname{ega4.c04.eq.13-2-4-2}
\expandafter\def\csname sga@target@IV.13.2.4.3\endcsname{ega4.c04.eq.13-2-4-3}
\expandafter\def\csname sga@target@IV.13.2.4.4\endcsname{ega4.c04.eq.13-2-4-4}
\expandafter\def\csname sga@target@IV.13.2.4.5\endcsname{ega4.c04.eq.13-2-4-5}
\expandafter\def\csname sga@target@IV.13.2.5\endcsname{ega4.c04.lem.13-2-5}
\expandafter\def\csname sga@target@IV.13.2.6\endcsname{ega4.c04.statement.13-2-6}
\expandafter\def\csname sga@target@IV.13.2.6.1\endcsname{ega4.c04.statement.13-2-6-1}
\expandafter\def\csname sga@target@IV.13.2.7\endcsname{ega4.c04.prop.13-2-7}
\expandafter\def\csname sga@target@IV.13.2.7.1\endcsname{ega4.c04.prop.13-2-7-1}
\expandafter\def\csname sga@target@IV.13.2.8\endcsname{ega4.c04.cor.13-2-8}
\expandafter\def\csname sga@target@IV.13.2.9\endcsname{ega4.c04.cor.13-2-9}
\expandafter\def\csname sga@target@IV.13.2.9.1\endcsname{ega4.c04.eq.13-2-9-1}
\expandafter\def\csname sga@target@IV.13.2.9.2\endcsname{ega4.c04.eq.13-2-9-2}
\expandafter\def\csname sga@target@IV.13.3.1\endcsname{ega4.c04.prop.13-3-1}
\expandafter\def\csname sga@target@IV.13.3.1.1\endcsname{ega4.c04.lem.13-3-1-1}
\expandafter\def\csname sga@target@IV.13.3.2\endcsname{ega4.c04.def.13-3-2}
\expandafter\def\csname sga@target@IV.13.3.3\endcsname{ega4.c04.prop.13-3-3}
\expandafter\def\csname sga@target@IV.13.3.4\endcsname{ega4.c04.cor.13-3-4}
\expandafter\def\csname sga@target@IV.13.3.4.1\endcsname{ega4.c04.cor.13-3-4-1}
\expandafter\def\csname sga@target@IV.13.3.5\endcsname{ega4.c04.cor.13-3-5}
\expandafter\def\csname sga@target@IV.13.3.6\endcsname{ega4.c04.prop.13-3-6}
\expandafter\def\csname sga@target@IV.13.3.6.1\endcsname{ega4.c04.prop.13-3-6-1}
\expandafter\def\csname sga@target@IV.13.3.6.2\endcsname{ega4.c04.eq.13-3-6-2}
\expandafter\def\csname sga@target@IV.13.3.7\endcsname{ega4.c04.prop.13-3-7}
\expandafter\def\csname sga@target@IV.13.3.7.1\endcsname{ega4.c04.prop.13-3-7-1}
\expandafter\def\csname sga@target@IV.13.3.8\endcsname{ega4.c04.prop.13-3-8}
\expandafter\def\csname sga@target@IV.13.3.9\endcsname{ega4.c04.statement.13-3-9}
\expandafter\def\csname sga@target@IV.14.1.1\endcsname{ega4.c04.statement.14-1-1}
\expandafter\def\csname sga@target@IV.14.1.2\endcsname{ega4.c04.prop.14-1-2}
\expandafter\def\csname sga@target@IV.14.1.3\endcsname{ega4.c04.item.14-1-3}
\expandafter\def\csname sga@target@IV.14.1.4\endcsname{ega4.c04.prop.14-1-4}
\expandafter\def\csname sga@target@IV.14.1.5\endcsname{ega4.c04.rem.14-1-5}
\expandafter\def\csname sga@target@IV.14.1.6\endcsname{ega4.c04.statement.14-1-6}
\expandafter\def\csname sga@target@IV.14.2.1\endcsname{ega4.c04.thm.14-2-1}
\expandafter\def\csname sga@target@IV.14.2.1.1\endcsname{ega4.c04.eq.14-2-1-1}
\expandafter\def\csname sga@target@IV.14.2.2\endcsname{ega4.c04.cor.14-2-2}
\expandafter\def\csname sga@target@IV.14.2.3\endcsname{ega4.c04.prop.14-2-3}
\expandafter\def\csname sga@target@IV.14.2.4\endcsname{ega4.c04.cor.14-2-4}
\expandafter\def\csname sga@target@IV.14.2.5\endcsname{ega4.c04.cor.14-2-5}
\expandafter\def\csname sga@target@IV.14.2.6\endcsname{ega4.c04.cor.14-2-6}
\expandafter\def\csname sga@target@IV.14.2.7\endcsname{ega4.c04.rem.14-2-7}
\expandafter\def\csname sga@target@IV.14.3\endcsname{ega4.c04.subsec.14-3.key-iv-14-3}
\expandafter\def\csname sga@target@IV.14.3.1\endcsname{ega4.c04.statement.14-3-1}
\expandafter\def\csname sga@target@IV.14.3.10\endcsname{ega4.c04.lem.14-3-10}
\expandafter\def\csname sga@target@IV.14.3.11\endcsname{ega4.c04.prop.14-3-11}
\expandafter\def\csname sga@target@IV.14.3.12\endcsname{ega4.c04.rem.14-3-12}
\expandafter\def\csname sga@target@IV.14.3.13\endcsname{ega4.c04.prop.14-3-13}
\expandafter\def\csname sga@target@IV.14.3.14\endcsname{ega4.c04.rem.14-3-14}
\expandafter\def\csname sga@target@IV.14.3.2\endcsname{ega4.c04.prop.14-3-2}
\expandafter\def\csname sga@target@IV.14.3.3\endcsname{ega4.c04.def.14-3-3}
\expandafter\def\csname sga@target@IV.14.3.3.1\endcsname{ega4.c04.subsec.14-3-3-1}
\expandafter\def\csname sga@target@IV.14.3.4\endcsname{ega4.c04.prop.14-3-4}
\expandafter\def\csname sga@target@IV.14.3.5\endcsname{ega4.c04.prop.14-3-5}
\expandafter\def\csname sga@target@IV.14.3.6\endcsname{ega4.c04.prop.14-3-6}
\expandafter\def\csname sga@target@IV.14.3.7\endcsname{ega4.c04.cor.14-3-7}
\expandafter\def\csname sga@target@IV.14.3.8\endcsname{ega4.c04.prop.14-3-8}
\expandafter\def\csname sga@target@IV.14.3.9\endcsname{ega4.c04.item.14-3-9}
\expandafter\def\csname sga@target@IV.14.4.1\endcsname{ega4.c04.thm.14-4-1}
\expandafter\def\csname sga@target@IV.14.4.1.1\endcsname{ega4.c04.lem.14-4-1-1}
\expandafter\def\csname sga@target@IV.14.4.1.2\endcsname{ega4.c04.lem.14-4-1-2}
\expandafter\def\csname sga@target@IV.14.4.1.3\endcsname{ega4.c04.lem.14-4-1-3}
\expandafter\def\csname sga@target@IV.14.4.10\endcsname{ega4.c04.item.14-4-10}
\expandafter\def\csname sga@target@IV.14.4.2\endcsname{ega4.c04.cor.14-4-2}
\expandafter\def\csname sga@target@IV.14.4.3\endcsname{ega4.c04.cor.14-4-3}
\expandafter\def\csname sga@target@IV.14.4.4\endcsname{ega4.c04.cor.14-4-4}
\expandafter\def\csname sga@target@IV.14.4.5\endcsname{ega4.c04.rem.14-4-5}
\expandafter\def\csname sga@target@IV.14.4.6\endcsname{ega4.c04.cor.14-4-6}
\expandafter\def\csname sga@target@IV.14.4.7\endcsname{ega4.c04.prop.14-4-7}
\expandafter\def\csname sga@target@IV.14.4.7.1\endcsname{ega4.c04.eq.14-4-7-1}
\expandafter\def\csname sga@target@IV.14.4.7.2\endcsname{ega4.c04.eq.14-4-7-2}
\expandafter\def\csname sga@target@IV.14.4.8\endcsname{ega4.c04.prop.14-4-8}
\expandafter\def\csname sga@target@IV.14.4.8.1\endcsname{ega4.c04.lem.14-4-8-1}
\expandafter\def\csname sga@target@IV.14.4.9\endcsname{ega4.c04.cor.14-4-9}
\expandafter\def\csname sga@target@IV.14.5\endcsname{ega4.c04.subsec.14-5.key-iv-14-5}
\expandafter\def\csname sga@target@IV.14.5.1\endcsname{ega4.c04.lem.14-5-1}
\expandafter\def\csname sga@target@IV.14.5.1.1\endcsname{ega4.c04.eq.14-5-1-1}
\expandafter\def\csname sga@target@IV.14.5.1.2\endcsname{ega4.c04.eq.14-5-1-2}
\expandafter\def\csname sga@target@IV.14.5.1.3\endcsname{ega4.c04.eq.14-5-1-3}
\expandafter\def\csname sga@target@IV.14.5.1.4\endcsname{ega4.c04.eq.14-5-1-4}
\expandafter\def\csname sga@target@IV.14.5.10\endcsname{ega4.c04.prop.14-5-10}
\expandafter\def\csname sga@target@IV.14.5.11\endcsname{ega4.c04.cor.14-5-11}
\expandafter\def\csname sga@target@IV.14.5.11.1\endcsname{ega4.c04.lem.14-5-11-1}
\expandafter\def\csname sga@target@IV.14.5.12\endcsname{ega4.c04.cor.14-5-12}
\expandafter\def\csname sga@target@IV.14.5.2\endcsname{ega4.c04.cor.14-5-2}
\expandafter\def\csname sga@target@IV.14.5.2.1\endcsname{ega4.c04.eq.14-5-2-1}
\expandafter\def\csname sga@target@IV.14.5.3\endcsname{ega4.c04.prop.14-5-3}
\expandafter\def\csname sga@target@IV.14.5.4\endcsname{ega4.c04.cor.14-5-4}
\expandafter\def\csname sga@target@IV.14.5.5\endcsname{ega4.c04.rem.14-5-5}
\expandafter\def\csname sga@target@IV.14.5.6\endcsname{ega4.c04.prop.14-5-6}
\expandafter\def\csname sga@target@IV.14.5.7\endcsname{ega4.c04.rem.14-5-7}
\expandafter\def\csname sga@target@IV.14.5.8\endcsname{ega4.c04.prop.14-5-8}
\expandafter\def\csname sga@target@IV.14.5.9\endcsname{ega4.c04.prop.14-5-9}
\expandafter\def\csname sga@target@IV.15.1.1\endcsname{ega4.c04.prop.15-1-1}
\expandafter\def\csname sga@target@IV.15.1.1-proof\endcsname{ega4.c04.statement.15-1-1-proof}
\expandafter\def\csname sga@target@IV.15.1.1.1\endcsname{ega4.c04.eq.15-1-1-1}
\expandafter\def\csname sga@target@IV.15.1.1.2\endcsname{ega4.c04.eq.15-1-1-2}
\expandafter\def\csname sga@target@IV.15.1.1.3\endcsname{ega4.c04.eq.15-1-1-3}
\expandafter\def\csname sga@target@IV.15.1.1.4\endcsname{ega4.c04.eq.15-1-1-4}
\expandafter\def\csname sga@target@IV.15.1.1.5\endcsname{ega4.c04.eq.15-1-1-5}
\expandafter\def\csname sga@target@IV.15.1.1.6\endcsname{ega4.c04.lem.15-1-1-6}
\expandafter\def\csname sga@target@IV.15.1.2\endcsname{ega4.c04.cor.15-1-2}
\expandafter\def\csname sga@target@IV.15.1.3\endcsname{ega4.c04.cor.15-1-3}
\expandafter\def\csname sga@target@IV.15.1.4\endcsname{ega4.c04.statement.15-1-4}
\expandafter\def\csname sga@target@IV.15.2.1\endcsname{ega4.c04.statement.15-2-1}
\expandafter\def\csname sga@target@IV.15.2.2\endcsname{ega4.c04.thm.15-2-2}
\expandafter\def\csname sga@target@IV.15.2.2-proof\endcsname{ega4.c04.item.15-2-2-proof}
\expandafter\def\csname sga@target@IV.15.2.2.1\endcsname{ega4.c04.lem.15-2-2-1}
\expandafter\def\csname sga@target@IV.15.2.3\endcsname{ega4.c04.cor.15-2-3}
\expandafter\def\csname sga@target@IV.15.2.4\endcsname{ega4.c04.item.15-2-4}
\expandafter\def\csname sga@target@IV.15.3.1\endcsname{ega4.c04.prop.15-3-1}
\expandafter\def\csname sga@target@IV.15.3.1-proof\endcsname{ega4.c04.item.15-3-1-proof}
\expandafter\def\csname sga@target@IV.15.3.2\endcsname{ega4.c04.prop.15-3-2}
\expandafter\def\csname sga@target@IV.15.3.2-proof\endcsname{ega4.c04.item.15-3-2-proof}
\expandafter\def\csname sga@target@IV.15.3.3\endcsname{ega4.c04.cor.15-3-3}
\expandafter\def\csname sga@target@IV.15.3.3-proof\endcsname{ega4.c04.item.15-3-3-proof}
\expandafter\def\csname sga@target@IV.15.3.4\endcsname{ega4.c04.item.15-3-4}
\expandafter\def\csname sga@target@IV.15.4.1\endcsname{ega4.c04.prop.15-4-1}
\expandafter\def\csname sga@target@IV.15.4.2\endcsname{ega4.c04.prop.15-4-2}
\expandafter\def\csname sga@target@IV.15.4.3\endcsname{ega4.c04.prop.15-4-3}
\expandafter\def\csname sga@target@IV.15.4.3-proof\endcsname{ega4.c04.item.15-4-3-proof}
\expandafter\def\csname sga@target@IV.15.5.1\endcsname{ega4.c04.prop.15-5-1}
\expandafter\def\csname sga@target@IV.15.5.10\endcsname{ega4.c04.item.15-5-10}
\expandafter\def\csname sga@target@IV.15.5.2\endcsname{ega4.c04.cor.15-5-2}
\expandafter\def\csname sga@target@IV.15.5.2.1\endcsname{ega4.c04.eq.15-5-2-1}
\expandafter\def\csname sga@target@IV.15.5.2.2\endcsname{ega4.c04.lem.15-5-2-2}
\expandafter\def\csname sga@target@IV.15.5.2.3\endcsname{ega4.c04.eq.15-5-2-3}
\expandafter\def\csname sga@target@IV.15.5.2.4\endcsname{ega4.c04.eq.15-5-2-4}
\expandafter\def\csname sga@target@IV.15.5.3\endcsname{ega4.c04.prop.15-5-3}
\expandafter\def\csname sga@target@IV.15.5.4\endcsname{ega4.c04.cor.15-5-4}
\expandafter\def\csname sga@target@IV.15.5.5\endcsname{ega4.c04.statement.15-5-5}
\expandafter\def\csname sga@target@IV.15.5.6\endcsname{ega4.c04.lem.15-5-6}
\expandafter\def\csname sga@target@IV.15.5.6.1\endcsname{ega4.c04.lem.15-5-6-1}
\expandafter\def\csname sga@target@IV.15.5.7\endcsname{ega4.c04.prop.15-5-7}
\expandafter\def\csname sga@target@IV.15.5.7.1\endcsname{ega4.c04.eq.15-5-7-1}
\expandafter\def\csname sga@target@IV.15.5.8\endcsname{ega4.c04.statement.15-5-8}
\expandafter\def\csname sga@target@IV.15.5.9\endcsname{ega4.c04.prop.15-5-9}
\expandafter\def\csname sga@target@IV.15.6.1\endcsname{ega4.c04.prop.15-6-1}
\expandafter\def\csname sga@target@IV.15.6.1.1\endcsname{ega4.c04.eq.15-6-1-1}
\expandafter\def\csname sga@target@IV.15.6.10\endcsname{ega4.c04.item.15-6-10}
\expandafter\def\csname sga@target@IV.15.6.2\endcsname{ega4.c04.cor.15-6-2}
\expandafter\def\csname sga@target@IV.15.6.3\endcsname{ega4.c04.prop.15-6-3}
\expandafter\def\csname sga@target@IV.15.6.4\endcsname{ega4.c04.prop.15-6-4}
\expandafter\def\csname sga@target@IV.15.6.5\endcsname{ega4.c04.cor.15-6-5}
\expandafter\def\csname sga@target@IV.15.6.6\endcsname{ega4.c04.prop.15-6-6}
\expandafter\def\csname sga@target@IV.15.6.7\endcsname{ega4.c04.cor.15-6-7}
\expandafter\def\csname sga@target@IV.15.6.8\endcsname{ega4.c04.prop.15-6-8}
\expandafter\def\csname sga@target@IV.15.6.9\endcsname{ega4.c04.cor.15-6-9}
\expandafter\def\csname sga@target@IV.15.7.1\endcsname{ega4.c04.def.15-7-1}
\expandafter\def\csname sga@target@IV.15.7.10\endcsname{ega4.c04.cor.15-7-10}
\expandafter\def\csname sga@target@IV.15.7.11\endcsname{ega4.c04.cor.15-7-11}
\expandafter\def\csname sga@target@IV.15.7.11.1\endcsname{ega4.c04.lem.15-7-11-1}
\expandafter\def\csname sga@target@IV.15.7.2\endcsname{ega4.c04.prop.15-7-2}
\expandafter\def\csname sga@target@IV.15.7.3\endcsname{ega4.c04.rem.15-7-3}
\expandafter\def\csname sga@target@IV.15.7.4\endcsname{ega4.c04.cor.15-7-4}
\expandafter\def\csname sga@target@IV.15.7.5\endcsname{ega4.c04.cor.15-7-5}
\expandafter\def\csname sga@target@IV.15.7.6\endcsname{ega4.c04.cor.15-7-6}
\expandafter\def\csname sga@target@IV.15.7.7\endcsname{ega4.c04.item.15-7-7}
\expandafter\def\csname sga@target@IV.15.7.8\endcsname{ega4.c04.statement.15-7-8}
\expandafter\def\csname sga@target@IV.15.7.9\endcsname{ega4.c04.prop.15-7-9}
\expandafter\def\csname sga@target@IV.16.1\endcsname{ega4.c04.subsec.16-1}
\expandafter\def\csname sga@target@IV.16.1.1\endcsname{ega4.c04.statement.16-1-1}
\expandafter\def\csname sga@target@IV.16.1.10\endcsname{ega4.c04.cor.16-1-10}
\expandafter\def\csname sga@target@IV.16.1.11\endcsname{ega4.c04.rem.16-1-11}
\expandafter\def\csname sga@target@IV.16.1.2\endcsname{ega4.c04.def.16-1-2}
\expandafter\def\csname sga@target@IV.16.1.2.1\endcsname{ega4.c04.def.16-1-2-1}
\expandafter\def\csname sga@target@IV.16.1.2.2\endcsname{ega4.c04.eq.16-1-2-2}
\expandafter\def\csname sga@target@IV.16.1.3\endcsname{ega4.c04.statement.16-1-3}
\expandafter\def\csname sga@target@IV.16.1.4\endcsname{ega4.c04.statement.16-1-4}
\expandafter\def\csname sga@target@IV.16.1.5\endcsname{ega4.c04.prop.16-1-5}
\expandafter\def\csname sga@target@IV.16.1.6\endcsname{ega4.c04.cor.16-1-6}
\expandafter\def\csname sga@target@IV.16.1.7\endcsname{ega4.c04.cor.16-1-7}
\expandafter\def\csname sga@target@IV.16.1.8\endcsname{ega4.c04.prop.16-1-8}
\expandafter\def\csname sga@target@IV.16.1.9\endcsname{ega4.c04.prop.16-1-9}
\expandafter\def\csname sga@target@IV.16.10\endcsname{ega4.c04.subsec.16-10}
\expandafter\def\csname sga@target@IV.16.10.1\endcsname{ega4.c04.def.16-10-1}
\expandafter\def\csname sga@target@IV.16.10.2\endcsname{ega4.c04.prop.16-10-2}
\expandafter\def\csname sga@target@IV.16.10.3\endcsname{ega4.c04.prop.16-10-3}
\expandafter\def\csname sga@target@IV.16.10.4\endcsname{ega4.c04.prop.16-10-4}
\expandafter\def\csname sga@target@IV.16.10.5\endcsname{ega4.c04.prop.16-10-5}
\expandafter\def\csname sga@target@IV.16.10.6\endcsname{ega4.c04.rem.16-10-6}
\expandafter\def\csname sga@target@IV.16.11\endcsname{ega4.c04.subsec.16-11}
\expandafter\def\csname sga@target@IV.16.11.1\endcsname{ega4.c04.statement.16-11-1}
\expandafter\def\csname sga@target@IV.16.11.1.1\endcsname{ega4.c04.statement.16-11-1-1}
\expandafter\def\csname sga@target@IV.16.11.1.2\endcsname{ega4.c04.statement.16-11-1-2}
\expandafter\def\csname sga@target@IV.16.11.1.3\endcsname{ega4.c04.eq.16-11-1-3}
\expandafter\def\csname sga@target@IV.16.11.1.4\endcsname{ega4.c04.eq.16-11-1-4}
\expandafter\def\csname sga@target@IV.16.11.1.5\endcsname{ega4.c04.eq.16-11-1-5}
\expandafter\def\csname sga@target@IV.16.11.1.6\endcsname{ega4.c04.eq.16-11-1-6}
\expandafter\def\csname sga@target@IV.16.11.1.7\endcsname{ega4.c04.eq.16-11-1-7}
\expandafter\def\csname sga@target@IV.16.11.2\endcsname{ega4.c04.thm.16-11-2}
\expandafter\def\csname sga@target@IV.16.11.2.1\endcsname{ega4.c04.eq.16-11-2-1}
\expandafter\def\csname sga@target@IV.16.11.2.2\endcsname{ega4.c04.eq.16-11-2-2}
\expandafter\def\csname sga@target@IV.16.11.2.3\endcsname{ega4.c04.eq.16-11-2-3}
\expandafter\def\csname sga@target@IV.16.11.3\endcsname{ega4.c04.statement.16-11-3}
\expandafter\def\csname sga@target@IV.16.12\endcsname{ega4.c04.subsec.16-12}
\expandafter\def\csname sga@target@IV.16.12.1\endcsname{ega4.c04.statement.16-12-1}
\expandafter\def\csname sga@target@IV.16.12.2\endcsname{ega4.c04.thm.16-12-2}
\expandafter\def\csname sga@target@IV.16.12.3\endcsname{ega4.c04.statement.16-12-3}
\expandafter\def\csname sga@target@IV.16.2\endcsname{ega4.c04.subsec.16-2}
\expandafter\def\csname sga@target@IV.16.2.1\endcsname{ega4.c04.statement.16-2-1}
\expandafter\def\csname sga@target@IV.16.2.1.1\endcsname{ega4.c04.statement.16-2-1-1}
\expandafter\def\csname sga@target@IV.16.2.1.2\endcsname{ega4.c04.eq.16-2-1-2}
\expandafter\def\csname sga@target@IV.16.2.1.3\endcsname{ega4.c04.eq.16-2-1-3}
\expandafter\def\csname sga@target@IV.16.2.1.4\endcsname{ega4.c04.eq.16-2-1-4}
\expandafter\def\csname sga@target@IV.16.2.2\endcsname{ega4.c04.prop.16-2-2}
\expandafter\def\csname sga@target@IV.16.2.3\endcsname{ega4.c04.cor.16-2-3}
\expandafter\def\csname sga@target@IV.16.2.3.1\endcsname{ega4.c04.eq.16-2-3-1}
\expandafter\def\csname sga@target@IV.16.2.3.2\endcsname{ega4.c04.eq.16-2-3-2}
\expandafter\def\csname sga@target@IV.16.2.4\endcsname{ega4.c04.cor.16-2-4}
\expandafter\def\csname sga@target@IV.16.2.5\endcsname{ega4.c04.statement.16-2-5}
\expandafter\def\csname sga@target@IV.16.2.6\endcsname{ega4.c04.statement.16-2-6}
\expandafter\def\csname sga@target@IV.16.2.6.1\endcsname{ega4.c04.statement.16-2-6-1}
\expandafter\def\csname sga@target@IV.16.2.6.2\endcsname{ega4.c04.eq.16-2-6-2}
\expandafter\def\csname sga@target@IV.16.2.7\endcsname{ega4.c04.prop.16-2-7}
\expandafter\def\csname sga@target@IV.16.2.7.1\endcsname{ega4.c04.prop.16-2-7-1}
\expandafter\def\csname sga@target@IV.16.3\endcsname{ega4.c04.subsec.16-3}
\expandafter\def\csname sga@target@IV.16.3.1\endcsname{ega4.c04.def.16-3-1}
\expandafter\def\csname sga@target@IV.16.3.1.1\endcsname{ega4.c04.eq.16-3-1-1}
\expandafter\def\csname sga@target@IV.16.3.2\endcsname{ega4.c04.statement.16-3-2}
\expandafter\def\csname sga@target@IV.16.3.3\endcsname{ega4.c04.statement.16-3-3}
\expandafter\def\csname sga@target@IV.16.3.3.1\endcsname{ega4.c04.statement.16-3-3-1}
\expandafter\def\csname sga@target@IV.16.3.4\endcsname{ega4.c04.prop.16-3-4}
\expandafter\def\csname sga@target@IV.16.3.5\endcsname{ega4.c04.statement.16-3-5}
\expandafter\def\csname sga@target@IV.16.3.6\endcsname{ega4.c04.def.16-3-6}
\expandafter\def\csname sga@target@IV.16.3.6.1\endcsname{ega4.c04.eq.16-3-6-1}
\expandafter\def\csname sga@target@IV.16.3.7\endcsname{ega4.c04.statement.16-3-7}
\expandafter\def\csname sga@target@IV.16.3.7.1\endcsname{ega4.c04.eq.16-3-7-1}
\expandafter\def\csname sga@target@IV.16.3.8\endcsname{ega4.c04.prop.16-3-8}
\expandafter\def\csname sga@target@IV.16.3.9\endcsname{ega4.c04.prop.16-3-9}
\expandafter\def\csname sga@target@IV.16.4\endcsname{ega4.c04.subsec.16-4}
\expandafter\def\csname sga@target@IV.16.4.1\endcsname{ega4.c04.statement.16-4-1}
\expandafter\def\csname sga@target@IV.16.4.1.1\endcsname{ega4.c04.statement.16-4-1-1}
\expandafter\def\csname sga@target@IV.16.4.1.2\endcsname{ega4.c04.statement.16-4-1-2}
\expandafter\def\csname sga@target@IV.16.4.1.3\endcsname{ega4.c04.eq.16-4-1-3}
\expandafter\def\csname sga@target@IV.16.4.1.4\endcsname{ega4.c04.eq.16-4-1-4}
\expandafter\def\csname sga@target@IV.16.4.1.5\endcsname{ega4.c04.eq.16-4-1-5}
\expandafter\def\csname sga@target@IV.16.4.10\endcsname{ega4.c04.cor.16-4-10}
\expandafter\def\csname sga@target@IV.16.4.11\endcsname{ega4.c04.prop.16-4-11}
\expandafter\def\csname sga@target@IV.16.4.11.1\endcsname{ega4.c04.prop.16-4-11-1}
\expandafter\def\csname sga@target@IV.16.4.11.2\endcsname{ega4.c04.prop.16-4-11-2}
\expandafter\def\csname sga@target@IV.16.4.11.3\endcsname{ega4.c04.eq.16-4-11-3}
\expandafter\def\csname sga@target@IV.16.4.11.4\endcsname{ega4.c04.eq.16-4-11-4}
\expandafter\def\csname sga@target@IV.16.4.12\endcsname{ega4.c04.cor.16-4-12}
\expandafter\def\csname sga@target@IV.16.4.13\endcsname{ega4.c04.cor.16-4-13}
\expandafter\def\csname sga@target@IV.16.4.14\endcsname{ega4.c04.prop.16-4-14}
\expandafter\def\csname sga@target@IV.16.4.14.1\endcsname{ega4.c04.prop.16-4-14-1}
\expandafter\def\csname sga@target@IV.16.4.15\endcsname{ega4.c04.cor.16-4-15}
\expandafter\def\csname sga@target@IV.16.4.15.1\endcsname{ega4.c04.cor.16-4-15-1}
\expandafter\def\csname sga@target@IV.16.4.15.2\endcsname{ega4.c04.eq.16-4-15-2}
\expandafter\def\csname sga@target@IV.16.4.16\endcsname{ega4.c04.cor.16-4-16}
\expandafter\def\csname sga@target@IV.16.4.16.1\endcsname{ega4.c04.cor.16-4-16-1}
\expandafter\def\csname sga@target@IV.16.4.16.2\endcsname{ega4.c04.eq.16-4-16-2}
\expandafter\def\csname sga@target@IV.16.4.17\endcsname{ega4.c04.cor.16-4-17}
\expandafter\def\csname sga@target@IV.16.4.18\endcsname{ega4.c04.statement.16-4-18}
\expandafter\def\csname sga@target@IV.16.4.18.1\endcsname{ega4.c04.statement.16-4-18-1}
\expandafter\def\csname sga@target@IV.16.4.18.2\endcsname{ega4.c04.statement.16-4-18-2}
\expandafter\def\csname sga@target@IV.16.4.18.3\endcsname{ega4.c04.eq.16-4-18-3}
\expandafter\def\csname sga@target@IV.16.4.19\endcsname{ega4.c04.cor.16-4-19}
\expandafter\def\csname sga@target@IV.16.4.19.1\endcsname{ega4.c04.cor.16-4-19-1}
\expandafter\def\csname sga@target@IV.16.4.2\endcsname{ega4.c04.statement.16-4-2}
\expandafter\def\csname sga@target@IV.16.4.20\endcsname{ega4.c04.prop.16-4-20}
\expandafter\def\csname sga@target@IV.16.4.21\endcsname{ega4.c04.cor.16-4-21}
\expandafter\def\csname sga@target@IV.16.4.21.1\endcsname{ega4.c04.cor.16-4-21-1}
\expandafter\def\csname sga@target@IV.16.4.22\endcsname{ega4.c04.cor.16-4-22}
\expandafter\def\csname sga@target@IV.16.4.23\endcsname{ega4.c04.prop.16-4-23}
\expandafter\def\csname sga@target@IV.16.4.23.1\endcsname{ega4.c04.prop.16-4-23-1}
\expandafter\def\csname sga@target@IV.16.4.24\endcsname{ega4.c04.statement.16-4-24}
\expandafter\def\csname sga@target@IV.16.4.3\endcsname{ega4.c04.statement.16-4-3}
\expandafter\def\csname sga@target@IV.16.4.3.1\endcsname{ega4.c04.statement.16-4-3-1}
\expandafter\def\csname sga@target@IV.16.4.3.2\endcsname{ega4.c04.statement.16-4-3-2}
\expandafter\def\csname sga@target@IV.16.4.3.3\endcsname{ega4.c04.eq.16-4-3-3}
\expandafter\def\csname sga@target@IV.16.4.3.4\endcsname{ega4.c04.eq.16-4-3-4}
\expandafter\def\csname sga@target@IV.16.4.3.5\endcsname{ega4.c04.eq.16-4-3-5}
\expandafter\def\csname sga@target@IV.16.4.3.6\endcsname{ega4.c04.eq.16-4-3-6}
\expandafter\def\csname sga@target@IV.16.4.3.7\endcsname{ega4.c04.eq.16-4-3-7}
\expandafter\def\csname sga@target@IV.16.4.4\endcsname{ega4.c04.statement.16-4-4}
\expandafter\def\csname sga@target@IV.16.4.5\endcsname{ega4.c04.prop.16-4-5}
\expandafter\def\csname sga@target@IV.16.4.6\endcsname{ega4.c04.cor.16-4-6}
\expandafter\def\csname sga@target@IV.16.4.7\endcsname{ega4.c04.statement.16-4-7}
\expandafter\def\csname sga@target@IV.16.4.7.1\endcsname{ega4.c04.statement.16-4-7-1}
\expandafter\def\csname sga@target@IV.16.4.7.2\endcsname{ega4.c04.eq.16-4-7-2}
\expandafter\def\csname sga@target@IV.16.4.7.3\endcsname{ega4.c04.eq.16-4-7-3}
\expandafter\def\csname sga@target@IV.16.4.8\endcsname{ega4.c04.prop.16-4-8}
\expandafter\def\csname sga@target@IV.16.4.8.1\endcsname{ega4.c04.prop.16-4-8-1}
\expandafter\def\csname sga@target@IV.16.4.8.2\endcsname{ega4.c04.eq.16-4-8-2}
\expandafter\def\csname sga@target@IV.16.4.8.3\endcsname{ega4.c04.lem.16-4-8-3}
\expandafter\def\csname sga@target@IV.16.4.8.4\endcsname{ega4.c04.eq.16-4-8-4}
\expandafter\def\csname sga@target@IV.16.4.9\endcsname{ega4.c04.cor.16-4-9}
\expandafter\def\csname sga@target@IV.16.4.9.1\endcsname{ega4.c04.cor.16-4-9-1}
\expandafter\def\csname sga@target@IV.16.5\endcsname{ega4.c04.subsec.16-5}
\expandafter\def\csname sga@target@IV.16.5.1\endcsname{ega4.c04.statement.16-5-1}
\expandafter\def\csname sga@target@IV.16.5.1.1\endcsname{ega4.c04.eq.16-5-1-1}
\expandafter\def\csname sga@target@IV.16.5.1.2\endcsname{ega4.c04.eq.16-5-1-2}
\expandafter\def\csname sga@target@IV.16.5.10\endcsname{ega4.c04.statement.16-5-10}
\expandafter\def\csname sga@target@IV.16.5.10.1\endcsname{ega4.c04.statement.16-5-10-1}
\expandafter\def\csname sga@target@IV.16.5.10.2\endcsname{ega4.c04.statement.16-5-10-2}
\expandafter\def\csname sga@target@IV.16.5.11\endcsname{ega4.c04.prop.16-5-11}
\expandafter\def\csname sga@target@IV.16.5.12\endcsname{ega4.c04.statement.16-5-12}
\expandafter\def\csname sga@target@IV.16.5.12.1\endcsname{ega4.c04.statement.16-5-12-1}
\expandafter\def\csname sga@target@IV.16.5.12.2\endcsname{ega4.c04.eq.16-5-12-2}
\expandafter\def\csname sga@target@IV.16.5.13\endcsname{ega4.c04.statement.16-5-13}
\expandafter\def\csname sga@target@IV.16.5.13.1\endcsname{ega4.c04.statement.16-5-13-1}
\expandafter\def\csname sga@target@IV.16.5.13.2\endcsname{ega4.c04.eq.16-5-13-2}
\expandafter\def\csname sga@target@IV.16.5.13.3\endcsname{ega4.c04.eq.16-5-13-3}
\expandafter\def\csname sga@target@IV.16.5.13.4\endcsname{ega4.c04.eq.16-5-13-4}
\expandafter\def\csname sga@target@IV.16.5.13.5\endcsname{ega4.c04.eq.16-5-13-5}
\expandafter\def\csname sga@target@IV.16.5.14\endcsname{ega4.c04.statement.16-5-14}
\expandafter\def\csname sga@target@IV.16.5.14.1\endcsname{ega4.c04.statement.16-5-14-1}
\expandafter\def\csname sga@target@IV.16.5.14.2\endcsname{ega4.c04.eq.16-5-14-2}
\expandafter\def\csname sga@target@IV.16.5.14.3\endcsname{ega4.c04.eq.16-5-14-3}
\expandafter\def\csname sga@target@IV.16.5.15\endcsname{ega4.c04.statement.16-5-15}
\expandafter\def\csname sga@target@IV.16.5.16\endcsname{ega4.c04.prop.16-5-16}
\expandafter\def\csname sga@target@IV.16.5.17\endcsname{ega4.c04.prop.16-5-17}
\expandafter\def\csname sga@target@IV.16.5.18\endcsname{ega4.c04.cor.16-5-18}
\expandafter\def\csname sga@target@IV.16.5.2\endcsname{ega4.c04.prop.16-5-2}
\expandafter\def\csname sga@target@IV.16.5.3\endcsname{ega4.c04.prop.16-5-3}
\expandafter\def\csname sga@target@IV.16.5.3.1\endcsname{ega4.c04.eq.16-5-3-1}
\expandafter\def\csname sga@target@IV.16.5.4\endcsname{ega4.c04.statement.16-5-4}
\expandafter\def\csname sga@target@IV.16.5.5\endcsname{ega4.c04.cor.16-5-5}
\expandafter\def\csname sga@target@IV.16.5.5.1\endcsname{ega4.c04.cor.16-5-5-1}
\expandafter\def\csname sga@target@IV.16.5.6\endcsname{ega4.c04.cor.16-5-6}
\expandafter\def\csname sga@target@IV.16.5.7\endcsname{ega4.c04.statement.16-5-7}
\expandafter\def\csname sga@target@IV.16.5.7.1\endcsname{ega4.c04.statement.16-5-7-1}
\expandafter\def\csname sga@target@IV.16.5.8\endcsname{ega4.c04.statement.16-5-8}
\expandafter\def\csname sga@target@IV.16.5.8.1\endcsname{ega4.c04.eq.16-5-8-1}
\expandafter\def\csname sga@target@IV.16.5.8.2\endcsname{ega4.c04.eq.16-5-8-2}
\expandafter\def\csname sga@target@IV.16.5.9\endcsname{ega4.c04.statement.16-5-9}
\expandafter\def\csname sga@target@IV.16.5.9.1\endcsname{ega4.c04.statement.16-5-9-1}
\expandafter\def\csname sga@target@IV.16.6\endcsname{ega4.c04.subsec.16-6}
\expandafter\def\csname sga@target@IV.16.6.1\endcsname{ega4.c04.statement.16-6-1}
\expandafter\def\csname sga@target@IV.16.6.1.1\endcsname{ega4.c04.statement.16-6-1-1}
\expandafter\def\csname sga@target@IV.16.6.1.2\endcsname{ega4.c04.statement.16-6-1-2}
\expandafter\def\csname sga@target@IV.16.6.2\endcsname{ega4.c04.thm.16-6-2}
\expandafter\def\csname sga@target@IV.16.6.2.1\endcsname{ega4.c04.eq.16-6-2-1}
\expandafter\def\csname sga@target@IV.16.6.2.2\endcsname{ega4.c04.eq.16-6-2-2}
\expandafter\def\csname sga@target@IV.16.6.2.3\endcsname{ega4.c04.eq.16-6-2-3}
\expandafter\def\csname sga@target@IV.16.6.2.4\endcsname{ega4.c04.eq.16-6-2-4}
\expandafter\def\csname sga@target@IV.16.6.2.5\endcsname{ega4.c04.eq.16-6-2-5}
\expandafter\def\csname sga@target@IV.16.6.3\endcsname{ega4.c04.def.16-6-3}
\expandafter\def\csname sga@target@IV.16.6.4\endcsname{ega4.c04.prop.16-6-4}
\expandafter\def\csname sga@target@IV.16.6.4.1\endcsname{ega4.c04.prop.16-6-4-1}
\expandafter\def\csname sga@target@IV.16.6.5\endcsname{ega4.c04.statement.16-6-5}
\expandafter\def\csname sga@target@IV.16.6.5.1\endcsname{ega4.c04.statement.16-6-5-1}
\expandafter\def\csname sga@target@IV.16.7\endcsname{ega4.c04.subsec.16-7}
\expandafter\def\csname sga@target@IV.16.7.1\endcsname{ega4.c04.statement.16-7-1}
\expandafter\def\csname sga@target@IV.16.7.1.1\endcsname{ega4.c04.eq.16-7-1-1}
\expandafter\def\csname sga@target@IV.16.7.1.2\endcsname{ega4.c04.eq.16-7-1-2}
\expandafter\def\csname sga@target@IV.16.7.10\endcsname{ega4.c04.rem.16-7-10}
\expandafter\def\csname sga@target@IV.16.7.2\endcsname{ega4.c04.statement.16-7-2}
\expandafter\def\csname sga@target@IV.16.7.2.1\endcsname{ega4.c04.statement.16-7-2-1}
\expandafter\def\csname sga@target@IV.16.7.3\endcsname{ega4.c04.prop.16-7-3}
\expandafter\def\csname sga@target@IV.16.7.4\endcsname{ega4.c04.statement.16-7-4}
\expandafter\def\csname sga@target@IV.16.7.4.1\endcsname{ega4.c04.statement.16-7-4-1}
\expandafter\def\csname sga@target@IV.16.7.5\endcsname{ega4.c04.statement.16-7-5}
\expandafter\def\csname sga@target@IV.16.7.5.1\endcsname{ega4.c04.statement.16-7-5-1}
\expandafter\def\csname sga@target@IV.16.7.5.2\endcsname{ega4.c04.statement.16-7-5-2}
\expandafter\def\csname sga@target@IV.16.7.5.3\endcsname{ega4.c04.statement.16-7-5-3}
\expandafter\def\csname sga@target@IV.16.7.5.4\endcsname{ega4.c04.statement.16-7-5-4}
\expandafter\def\csname sga@target@IV.16.7.6\endcsname{ega4.c04.prop.16-7-6}
\expandafter\def\csname sga@target@IV.16.7.7\endcsname{ega4.c04.statement.16-7-7}
\expandafter\def\csname sga@target@IV.16.7.7.1\endcsname{ega4.c04.eq.16-7-7-1}
\expandafter\def\csname sga@target@IV.16.7.7.2\endcsname{ega4.c04.eq.16-7-7-2}
\expandafter\def\csname sga@target@IV.16.7.8\endcsname{ega4.c04.statement.16-7-8}
\expandafter\def\csname sga@target@IV.16.7.8.1\endcsname{ega4.c04.statement.16-7-8-1}
\expandafter\def\csname sga@target@IV.16.7.8.2\endcsname{ega4.c04.eq.16-7-8-2}
\expandafter\def\csname sga@target@IV.16.7.9\endcsname{ega4.c04.statement.16-7-9}
\expandafter\def\csname sga@target@IV.16.7.9.1\endcsname{ega4.c04.statement.16-7-9-1}
\expandafter\def\csname sga@target@IV.16.8\endcsname{ega4.c04.subsec.16-8}
\expandafter\def\csname sga@target@IV.16.8.1\endcsname{ega4.c04.def.16-8-1}
\expandafter\def\csname sga@target@IV.16.8.10\endcsname{ega4.c04.cor.16-8-10}
\expandafter\def\csname sga@target@IV.16.8.11\endcsname{ega4.c04.statement.16-8-11}
\expandafter\def\csname sga@target@IV.16.8.2\endcsname{ega4.c04.statement.16-8-2}
\expandafter\def\csname sga@target@IV.16.8.3\endcsname{ega4.c04.statement.16-8-3}
\expandafter\def\csname sga@target@IV.16.8.3.1\endcsname{ega4.c04.eq.16-8-3-1}
\expandafter\def\csname sga@target@IV.16.8.4\endcsname{ega4.c04.prop.16-8-4}
\expandafter\def\csname sga@target@IV.16.8.4.1\endcsname{ega4.c04.prop.16-8-4-1}
\expandafter\def\csname sga@target@IV.16.8.5\endcsname{ega4.c04.statement.16-8-5}
\expandafter\def\csname sga@target@IV.16.8.5.1\endcsname{ega4.c04.eq.16-8-5-1}
\expandafter\def\csname sga@target@IV.16.8.5.2\endcsname{ega4.c04.eq.16-8-5-2}
\expandafter\def\csname sga@target@IV.16.8.6\endcsname{ega4.c04.prop.16-8-6}
\expandafter\def\csname sga@target@IV.16.8.7\endcsname{ega4.c04.statement.16-8-7}
\expandafter\def\csname sga@target@IV.16.8.7.1\endcsname{ega4.c04.eq.16-8-7-1}
\expandafter\def\csname sga@target@IV.16.8.8\endcsname{ega4.c04.prop.16-8-8}
\expandafter\def\csname sga@target@IV.16.8.8.1\endcsname{ega4.c04.eq.16-8-8-1}
\expandafter\def\csname sga@target@IV.16.8.8.2\endcsname{ega4.c04.eq.16-8-8-2}
\expandafter\def\csname sga@target@IV.16.8.9\endcsname{ega4.c04.prop.16-8-9}
\expandafter\def\csname sga@target@IV.16.8.9.1\endcsname{ega4.c04.lem.16-8-9-1}
\expandafter\def\csname sga@target@IV.16.8.9.2\endcsname{ega4.c04.lem.16-8-9-2}
\expandafter\def\csname sga@target@IV.16.8.9.3\endcsname{ega4.c04.lem.16-8-9-3}
\expandafter\def\csname sga@target@IV.16.8.9.4\endcsname{ega4.c04.eq.16-8-9-4}
\expandafter\def\csname sga@target@IV.16.9\endcsname{ega4.c04.subsec.16-9}
\expandafter\def\csname sga@target@IV.16.9.1\endcsname{ega4.c04.def.16-9-1}
\expandafter\def\csname sga@target@IV.16.9.10\endcsname{ega4.c04.prop.16-9-10}
\expandafter\def\csname sga@target@IV.16.9.11\endcsname{ega4.c04.cor.16-9-11}
\expandafter\def\csname sga@target@IV.16.9.12\endcsname{ega4.c04.statement.16-9-12}
\expandafter\def\csname sga@target@IV.16.9.13\endcsname{ega4.c04.prop.16-9-13}
\expandafter\def\csname sga@target@IV.16.9.13.1\endcsname{ega4.c04.prop.16-9-13-1}
\expandafter\def\csname sga@target@IV.16.9.13.2\endcsname{ega4.c04.lem.16-9-13-2}
\expandafter\def\csname sga@target@IV.16.9.13.3\endcsname{ega4.c04.lem.16-9-13-3}
\expandafter\def\csname sga@target@IV.16.9.13.4\endcsname{ega4.c04.eq.16-9-13-4}
\expandafter\def\csname sga@target@IV.16.9.13.5\endcsname{ega4.c04.cor.16-9-13-5}
\expandafter\def\csname sga@target@IV.16.9.13.6\endcsname{ega4.c04.cor.16-9-13-6}
\expandafter\def\csname sga@target@IV.16.9.13.7\endcsname{ega4.c04.eq.16-9-13-7}
\expandafter\def\csname sga@target@IV.16.9.2\endcsname{ega4.c04.def.16-9-2}
\expandafter\def\csname sga@target@IV.16.9.3\endcsname{ega4.c04.prop.16-9-3}
\expandafter\def\csname sga@target@IV.16.9.4\endcsname{ega4.c04.cor.16-9-4}
\expandafter\def\csname sga@target@IV.16.9.4.1\endcsname{ega4.c04.eq.16-9-4-1}
\expandafter\def\csname sga@target@IV.16.9.5\endcsname{ega4.c04.cor.16-9-5}
\expandafter\def\csname sga@target@IV.16.9.6\endcsname{ega4.c04.statement.16-9-6}
\expandafter\def\csname sga@target@IV.16.9.7\endcsname{ega4.c04.statement.16-9-7}
\expandafter\def\csname sga@target@IV.16.9.8\endcsname{ega4.c04.prop.16-9-8}
\expandafter\def\csname sga@target@IV.16.9.9\endcsname{ega4.c04.cor.16-9-9}
\expandafter\def\csname sga@target@IV.17.1.1\endcsname{ega4.c04.def.17-1-1}
\expandafter\def\csname sga@target@IV.17.1.1.1\endcsname{ega4.c04.def.17-1-1-1}
\expandafter\def\csname sga@target@IV.17.1.2\endcsname{ega4.c04.rem.17-1-2}
\expandafter\def\csname sga@target@IV.17.1.2.1\endcsname{ega4.c04.eq.17-1-2-1}
\expandafter\def\csname sga@target@IV.17.1.3\endcsname{ega4.c04.prop.17-1-3}
\expandafter\def\csname sga@target@IV.17.1.4\endcsname{ega4.c04.prop.17-1-4}
\expandafter\def\csname sga@target@IV.17.1.5\endcsname{ega4.c04.cor.17-1-5}
\expandafter\def\csname sga@target@IV.17.1.6\endcsname{ega4.c04.prop.17-1-6}
\expandafter\def\csname sga@target@IV.17.10.1\endcsname{ega4.c04.def.17-10-1}
\expandafter\def\csname sga@target@IV.17.10.2\endcsname{ega4.c04.prop.17-10-2}
\expandafter\def\csname sga@target@IV.17.10.3\endcsname{ega4.c04.cor.17-10-3}
\expandafter\def\csname sga@target@IV.17.10.3.1\endcsname{ega4.c04.cor.17-10-3-1}
\expandafter\def\csname sga@target@IV.17.10.4\endcsname{ega4.c04.cor.17-10-4}
\expandafter\def\csname sga@target@IV.17.10.4.1\endcsname{ega4.c04.eq.17-10-4-1}
\expandafter\def\csname sga@target@IV.17.11.1\endcsname{ega4.c04.thm.17-11-1}
\expandafter\def\csname sga@target@IV.17.11.1.1\endcsname{ega4.c04.eq.17-11-1-1}
\expandafter\def\csname sga@target@IV.17.11.1.2\endcsname{ega4.c04.eq.17-11-1-2}
\expandafter\def\csname sga@target@IV.17.11.2\endcsname{ega4.c04.cor.17-11-2}
\expandafter\def\csname sga@target@IV.17.11.3\endcsname{ega4.c04.cor.17-11-3}
\expandafter\def\csname sga@target@IV.17.11.4\endcsname{ega4.c04.cor.17-11-4}
\expandafter\def\csname sga@target@IV.17.11.5\endcsname{ega4.c04.prop.17-11-5}
\expandafter\def\csname sga@target@IV.17.12.1\endcsname{ega4.c04.thm.17-12-1}
\expandafter\def\csname sga@target@IV.17.12.2\endcsname{ega4.c04.cor.17-12-2}
\expandafter\def\csname sga@target@IV.17.12.3\endcsname{ega4.c04.cor.17-12-3}
\expandafter\def\csname sga@target@IV.17.12.4\endcsname{ega4.c04.prop.17-12-4}
\expandafter\def\csname sga@target@IV.17.12.5\endcsname{ega4.c04.prop.17-12-5}
\expandafter\def\csname sga@target@IV.17.12.6\endcsname{ega4.c04.cor.17-12-6}
\expandafter\def\csname sga@target@IV.17.12.7\endcsname{ega4.c04.ex.17-12-7}
\expandafter\def\csname sga@target@IV.17.12.8\endcsname{ega4.c04.rem.17-12-8}
\expandafter\def\csname sga@target@IV.17.13\endcsname{ega4.c04.subsec.17-13.key-iv-17-13}
\expandafter\def\csname sga@target@IV.17.13.1\endcsname{ega4.c04.statement.17-13-1}
\expandafter\def\csname sga@target@IV.17.13.1.1\endcsname{ega4.c04.statement.17-13-1-1}
\expandafter\def\csname sga@target@IV.17.13.1.2\endcsname{ega4.c04.statement.17-13-1-2}
\expandafter\def\csname sga@target@IV.17.13.1.3\endcsname{ega4.c04.eq.17-13-1-3}
\expandafter\def\csname sga@target@IV.17.13.10\endcsname{ega4.c04.cor.17-13-10}
\expandafter\def\csname sga@target@IV.17.13.11\endcsname{ega4.c04.statement.17-13-11}
\expandafter\def\csname sga@target@IV.17.13.11.1\endcsname{ega4.c04.statement.17-13-11-1}
\expandafter\def\csname sga@target@IV.17.13.12\endcsname{ega4.c04.cor.17-13-12}
\expandafter\def\csname sga@target@IV.17.13.12.1\endcsname{ega4.c04.cor.17-13-12-1}
\expandafter\def\csname sga@target@IV.17.13.12.2\endcsname{ega4.c04.eq.17-13-12-2}
\expandafter\def\csname sga@target@IV.17.13.2\endcsname{ega4.c04.prop.17-13-2}
\expandafter\def\csname sga@target@IV.17.13.2.1\endcsname{ega4.c04.eq.17-13-2-1}
\expandafter\def\csname sga@target@IV.17.13.2.2\endcsname{ega4.c04.eq.17-13-2-2}
\expandafter\def\csname sga@target@IV.17.13.3\endcsname{ega4.c04.def.17-13-3}
\expandafter\def\csname sga@target@IV.17.13.4\endcsname{ega4.c04.statement.17-13-4}
\expandafter\def\csname sga@target@IV.17.13.5\endcsname{ega4.c04.statement.17-13-5}
\expandafter\def\csname sga@target@IV.17.13.5.1\endcsname{ega4.c04.statement.17-13-5-1}
\expandafter\def\csname sga@target@IV.17.13.5.2\endcsname{ega4.c04.eq.17-13-5-2}
\expandafter\def\csname sga@target@IV.17.13.5.3\endcsname{ega4.c04.eq.17-13-5-3}
\expandafter\def\csname sga@target@IV.17.13.5.4\endcsname{ega4.c04.eq.17-13-5-4}
\expandafter\def\csname sga@target@IV.17.13.6\endcsname{ega4.c04.cor.17-13-6}
\expandafter\def\csname sga@target@IV.17.13.6.1\endcsname{ega4.c04.eq.17-13-6-1}
\expandafter\def\csname sga@target@IV.17.13.6.2\endcsname{ega4.c04.eq.17-13-6-2}
\expandafter\def\csname sga@target@IV.17.13.7\endcsname{ega4.c04.statement.17-13-7}
\expandafter\def\csname sga@target@IV.17.13.7.1\endcsname{ega4.c04.eq.17-13-7-1}
\expandafter\def\csname sga@target@IV.17.13.8\endcsname{ega4.c04.prop.17-13-8}
\expandafter\def\csname sga@target@IV.17.13.8.1\endcsname{ega4.c04.eq.17-13-8-1}
\expandafter\def\csname sga@target@IV.17.13.9\endcsname{ega4.c04.statement.17-13-9}
\expandafter\def\csname sga@target@IV.17.13.9.1\endcsname{ega4.c04.eq.17-13-9-1}
\expandafter\def\csname sga@target@IV.17.13.9.2\endcsname{ega4.c04.eq.17-13-9-2}
\expandafter\def\csname sga@target@IV.17.14.1\endcsname{ega4.c04.prop.17-14-1}
\expandafter\def\csname sga@target@IV.17.14.2\endcsname{ega4.c04.prop.17-14-2}
\expandafter\def\csname sga@target@IV.17.14.3\endcsname{ega4.c04.rem.17-14-3}
\expandafter\def\csname sga@target@IV.17.15.1\endcsname{ega4.c04.prop.17-15-1}
\expandafter\def\csname sga@target@IV.17.15.10\endcsname{ega4.c04.statement.17-15-10}
\expandafter\def\csname sga@target@IV.17.15.11\endcsname{ega4.c04.prop.17-15-11}
\expandafter\def\csname sga@target@IV.17.15.12\endcsname{ega4.c04.prop.17-15-12}
\expandafter\def\csname sga@target@IV.17.15.13\endcsname{ega4.c04.cor.17-15-13}
\expandafter\def\csname sga@target@IV.17.15.14\endcsname{ega4.c04.prop.17-15-14}
\expandafter\def\csname sga@target@IV.17.15.14.1\endcsname{ega4.c04.lem.17-15-14-1}
\expandafter\def\csname sga@target@IV.17.15.14.2\endcsname{ega4.c04.lem.17-15-14-2}
\expandafter\def\csname sga@target@IV.17.15.15\endcsname{ega4.c04.prop.17-15-15}
\expandafter\def\csname sga@target@IV.17.15.2\endcsname{ega4.c04.cor.17-15-2}
\expandafter\def\csname sga@target@IV.17.15.3\endcsname{ega4.c04.prop.17-15-3}
\expandafter\def\csname sga@target@IV.17.15.4\endcsname{ega4.c04.cor.17-15-4}
\expandafter\def\csname sga@target@IV.17.15.5\endcsname{ega4.c04.prop.17-15-5}
\expandafter\def\csname sga@target@IV.17.15.6\endcsname{ega4.c04.cor.17-15-6}
\expandafter\def\csname sga@target@IV.17.15.7\endcsname{ega4.c04.cor.17-15-7}
\expandafter\def\csname sga@target@IV.17.15.8\endcsname{ega4.c04.prop.17-15-8}
\expandafter\def\csname sga@target@IV.17.15.8.1\endcsname{ega4.c04.eq.17-15-8-1}
\expandafter\def\csname sga@target@IV.17.15.8.2\endcsname{ega4.c04.eq.17-15-8-2}
\expandafter\def\csname sga@target@IV.17.15.9\endcsname{ega4.c04.cor.17-15-9}
\expandafter\def\csname sga@target@IV.17.16\endcsname{ega4.c04.subsec.17-16.key-iv-17-16}
\expandafter\def\csname sga@target@IV.17.16.1\endcsname{ega4.c04.prop.17-16-1}
\expandafter\def\csname sga@target@IV.17.16.2\endcsname{ega4.c04.cor.17-16-2}
\expandafter\def\csname sga@target@IV.17.16.3\endcsname{ega4.c04.cor.17-16-3}
\expandafter\def\csname sga@target@IV.17.16.4\endcsname{ega4.c04.prop.17-16-4}
\expandafter\def\csname sga@target@IV.17.16.5\endcsname{ega4.c04.cor.17-16-5}
\expandafter\def\csname sga@target@IV.17.16.6\endcsname{ega4.c04.prop.17-16-6}
\expandafter\def\csname sga@target@IV.17.2.1\endcsname{ega4.c04.prop.17-2-1}
\expandafter\def\csname sga@target@IV.17.2.2\endcsname{ega4.c04.cor.17-2-2}
\expandafter\def\csname sga@target@IV.17.2.3\endcsname{ega4.c04.prop.17-2-3}
\expandafter\def\csname sga@target@IV.17.2.3.1\endcsname{ega4.c04.eq.17-2-3-1}
\expandafter\def\csname sga@target@IV.17.2.4\endcsname{ega4.c04.cor.17-2-4}
\expandafter\def\csname sga@target@IV.17.2.5\endcsname{ega4.c04.prop.17-2-5}
\expandafter\def\csname sga@target@IV.17.2.5.1\endcsname{ega4.c04.prop.17-2-5-1}
\expandafter\def\csname sga@target@IV.17.2.6\endcsname{ega4.c04.prop.17-2-6}
\expandafter\def\csname sga@target@IV.17.3.1\endcsname{ega4.c04.def.17-3-1}
\expandafter\def\csname sga@target@IV.17.3.2\endcsname{ega4.c04.subsec.17-3-2}
\expandafter\def\csname sga@target@IV.17.3.3\endcsname{ega4.c04.prop.17-3-3}
\expandafter\def\csname sga@target@IV.17.3.4\endcsname{ega4.c04.prop.17-3-4}
\expandafter\def\csname sga@target@IV.17.3.5\endcsname{ega4.c04.cor.17-3-5}
\expandafter\def\csname sga@target@IV.17.3.6\endcsname{ega4.c04.prop.17-3-6}
\expandafter\def\csname sga@target@IV.17.3.7\endcsname{ega4.c04.def.17-3-7}
\expandafter\def\csname sga@target@IV.17.3.8\endcsname{ega4.c04.prop.17-3-8}
\expandafter\def\csname sga@target@IV.17.3.9\endcsname{ega4.c04.cor.17-3-9}
\expandafter\def\csname sga@target@IV.17.4.1\endcsname{ega4.c04.thm.17-4-1}
\expandafter\def\csname sga@target@IV.17.4.1.1\endcsname{ega4.c04.eq.17-4-1-1}
\expandafter\def\csname sga@target@IV.17.4.1.2\endcsname{ega4.c04.rem.17-4-1-2}
\expandafter\def\csname sga@target@IV.17.4.10\endcsname{ega4.c04.rem.17-4-10}
\expandafter\def\csname sga@target@IV.17.4.2\endcsname{ega4.c04.cor.17-4-2}
\expandafter\def\csname sga@target@IV.17.4.3\endcsname{ega4.c04.cor.17-4-3}
\expandafter\def\csname sga@target@IV.17.4.4\endcsname{ega4.c04.prop.17-4-4}
\expandafter\def\csname sga@target@IV.17.4.5\endcsname{ega4.c04.statement.17-4-5}
\expandafter\def\csname sga@target@IV.17.4.6\endcsname{ega4.c04.prop.17-4-6}
\expandafter\def\csname sga@target@IV.17.4.7\endcsname{ega4.c04.cor.17-4-7}
\expandafter\def\csname sga@target@IV.17.4.8\endcsname{ega4.c04.cor.17-4-8}
\expandafter\def\csname sga@target@IV.17.4.9\endcsname{ega4.c04.prop.17-4-9}
\expandafter\def\csname sga@target@IV.17.5.1\endcsname{ega4.c04.thm.17-5-1}
\expandafter\def\csname sga@target@IV.17.5.1.1\endcsname{ega4.c04.eq.17-5-1-1}
\expandafter\def\csname sga@target@IV.17.5.1.2\endcsname{ega4.c04.eq.17-5-1-2}
\expandafter\def\csname sga@target@IV.17.5.2\endcsname{ega4.c04.cor.17-5-2}
\expandafter\def\csname sga@target@IV.17.5.3\endcsname{ega4.c04.prop.17-5-3}
\expandafter\def\csname sga@target@IV.17.5.4\endcsname{ega4.c04.rem.17-5-4}
\expandafter\def\csname sga@target@IV.17.5.5\endcsname{ega4.c04.prop.17-5-5}
\expandafter\def\csname sga@target@IV.17.5.6\endcsname{ega4.c04.cor.17-5-6}
\expandafter\def\csname sga@target@IV.17.5.7\endcsname{ega4.c04.prop.17-5-7}
\expandafter\def\csname sga@target@IV.17.5.8\endcsname{ega4.c04.prop.17-5-8}
\expandafter\def\csname sga@target@IV.17.6.1\endcsname{ega4.c04.thm.17-6-1}
\expandafter\def\csname sga@target@IV.17.6.2\endcsname{ega4.c04.cor.17-6-2}
\expandafter\def\csname sga@target@IV.17.6.3\endcsname{ega4.c04.prop.17-6-3}
\expandafter\def\csname sga@target@IV.17.6.4\endcsname{ega4.c04.prop.17-6-4}
\expandafter\def\csname sga@target@IV.17.7\endcsname{ega4.c04.subsec.17-7.key-iv-17-7}
\expandafter\def\csname sga@target@IV.17.7.1\endcsname{ega4.c04.prop.17-7-1}
\expandafter\def\csname sga@target@IV.17.7.10\endcsname{ega4.c04.prop.17-7-10}
\expandafter\def\csname sga@target@IV.17.7.11\endcsname{ega4.c04.prop.17-7-11}
\expandafter\def\csname sga@target@IV.17.7.11.1\endcsname{ega4.c04.lem.17-7-11-1}
\expandafter\def\csname sga@target@IV.17.7.2\endcsname{ega4.c04.cor.17-7-2}
\expandafter\def\csname sga@target@IV.17.7.3\endcsname{ega4.c04.cor.17-7-3}
\expandafter\def\csname sga@target@IV.17.7.4\endcsname{ega4.c04.prop.17-7-4}
\expandafter\def\csname sga@target@IV.17.7.5\endcsname{ega4.c04.lem.17-7-5}
\expandafter\def\csname sga@target@IV.17.7.6\endcsname{ega4.c04.rem.17-7-6}
\expandafter\def\csname sga@target@IV.17.7.7\endcsname{ega4.c04.prop.17-7-7}
\expandafter\def\csname sga@target@IV.17.7.8\endcsname{ega4.c04.prop.17-7-8}
\expandafter\def\csname sga@target@IV.17.7.9\endcsname{ega4.c04.cor.17-7-9}
\expandafter\def\csname sga@target@IV.17.8\endcsname{ega4.c04.subsec.17-8.key-iv-17-8}
\expandafter\def\csname sga@target@IV.17.8.1\endcsname{ega4.c04.prop.17-8-1}
\expandafter\def\csname sga@target@IV.17.8.2\endcsname{ega4.c04.prop.17-8-2}
\expandafter\def\csname sga@target@IV.17.8.3\endcsname{ega4.c04.rem.17-8-3}
\expandafter\def\csname sga@target@IV.17.9\endcsname{ega4.c04.subsec.17-9.key-iv-17-9}
\expandafter\def\csname sga@target@IV.17.9.1\endcsname{ega4.c04.thm.17-9-1}
\expandafter\def\csname sga@target@IV.17.9.2\endcsname{ega4.c04.cor.17-9-2}
\expandafter\def\csname sga@target@IV.17.9.3\endcsname{ega4.c04.cor.17-9-3}
\expandafter\def\csname sga@target@IV.17.9.4\endcsname{ega4.c04.cor.17-9-4}
\expandafter\def\csname sga@target@IV.17.9.5\endcsname{ega4.c04.cor.17-9-5}
\expandafter\def\csname sga@target@IV.17.9.6\endcsname{ega4.c04.prop.17-9-6}
\expandafter\def\csname sga@target@IV.18\endcsname{ega4.c04.sec.18.key-iv-18}
\expandafter\def\csname sga@target@IV.18.1.1\endcsname{ega4.c04.prop.18-1-1}
\expandafter\def\csname sga@target@IV.18.1.1.1\endcsname{ega4.c04.eq.18-1-1-1}
\expandafter\def\csname sga@target@IV.18.1.2\endcsname{ega4.c04.thm.18-1-2}
\expandafter\def\csname sga@target@IV.18.1.3\endcsname{ega4.c04.cor.18-1-3}
\expandafter\def\csname sga@target@IV.18.10.1\endcsname{ega4.c04.thm.18-10-1}
\expandafter\def\csname sga@target@IV.18.10.10\endcsname{ega4.c04.statement.18-10-10}
\expandafter\def\csname sga@target@IV.18.10.11\endcsname{ega4.c04.rem.18-10-11}
\expandafter\def\csname sga@target@IV.18.10.12\endcsname{ega4.c04.cor.18-10-12}
\expandafter\def\csname sga@target@IV.18.10.13\endcsname{ega4.c04.prop.18-10-13}
\expandafter\def\csname sga@target@IV.18.10.14\endcsname{ega4.c04.cor.18-10-14}
\expandafter\def\csname sga@target@IV.18.10.15\endcsname{ega4.c04.statement.18-10-15}
\expandafter\def\csname sga@target@IV.18.10.16\endcsname{ega4.c04.thm.18-10-16}
\expandafter\def\csname sga@target@IV.18.10.17\endcsname{ega4.c04.rem.18-10-17}
\expandafter\def\csname sga@target@IV.18.10.17.1\endcsname{ega4.c04.statement.18-10-17-1}
\expandafter\def\csname sga@target@IV.18.10.17.2\endcsname{ega4.c04.statement.18-10-17-2}
\expandafter\def\csname sga@target@IV.18.10.18\endcsname{ega4.c04.lem.18-10-18}
\expandafter\def\csname sga@target@IV.18.10.19\endcsname{ega4.c04.prop.18-10-19}
\expandafter\def\csname sga@target@IV.18.10.2\endcsname{ega4.c04.rem.18-10-2}
\expandafter\def\csname sga@target@IV.18.10.20\endcsname{ega4.c04.rem.18-10-20}
\expandafter\def\csname sga@target@IV.18.10.21\endcsname{ega4.c04.prop.18-10-21}
\expandafter\def\csname sga@target@IV.18.10.3\endcsname{ega4.c04.cor.18-10-3}
\expandafter\def\csname sga@target@IV.18.10.3.1\endcsname{ega4.c04.rem.18-10-3-1}
\expandafter\def\csname sga@target@IV.18.10.4\endcsname{ega4.c04.cor.18-10-4}
\expandafter\def\csname sga@target@IV.18.10.5\endcsname{ega4.c04.cor.18-10-5}
\expandafter\def\csname sga@target@IV.18.10.6\endcsname{ega4.c04.lem.18-10-6}
\expandafter\def\csname sga@target@IV.18.10.7\endcsname{ega4.c04.prop.18-10-7}
\expandafter\def\csname sga@target@IV.18.10.8\endcsname{ega4.c04.prop.18-10-8}
\expandafter\def\csname sga@target@IV.18.10.9\endcsname{ega4.c04.cor.18-10-9}
\expandafter\def\csname sga@target@IV.18.11.1\endcsname{ega4.c04.lem.18-11-1}
\expandafter\def\csname sga@target@IV.18.11.10\endcsname{ega4.c04.thm.18-11-10}
\expandafter\def\csname sga@target@IV.18.11.11\endcsname{ega4.c04.rem.18-11-11}
\expandafter\def\csname sga@target@IV.18.11.12\endcsname{ega4.c04.cor.18-11-12}
\expandafter\def\csname sga@target@IV.18.11.12.1\endcsname{ega4.c04.lem.18-11-12-1}
\expandafter\def\csname sga@target@IV.18.11.2\endcsname{ega4.c04.lem.18-11-2}
\expandafter\def\csname sga@target@IV.18.11.2.1\endcsname{ega4.c04.eq.18-11-2-1}
\expandafter\def\csname sga@target@IV.18.11.2.2\endcsname{ega4.c04.eq.18-11-2-2}
\expandafter\def\csname sga@target@IV.18.11.3\endcsname{ega4.c04.lem.18-11-3}
\expandafter\def\csname sga@target@IV.18.11.3.1\endcsname{ega4.c04.eq.18-11-3-1}
\expandafter\def\csname sga@target@IV.18.11.3.2\endcsname{ega4.c04.eq.18-11-3-2}
\expandafter\def\csname sga@target@IV.18.11.3.3\endcsname{ega4.c04.lem.18-11-3-3}
\expandafter\def\csname sga@target@IV.18.11.3.4\endcsname{ega4.c04.lem.18-11-3-4}
\expandafter\def\csname sga@target@IV.18.11.3.5\endcsname{ega4.c04.lem.18-11-3-5}
\expandafter\def\csname sga@target@IV.18.11.3.6\endcsname{ega4.c04.eq.18-11-3-6}
\expandafter\def\csname sga@target@IV.18.11.4\endcsname{ega4.c04.lem.18-11-4}
\expandafter\def\csname sga@target@IV.18.11.5\endcsname{ega4.c04.prop.18-11-5}
\expandafter\def\csname sga@target@IV.18.11.5.1\endcsname{ega4.c04.lem.18-11-5-1}
\expandafter\def\csname sga@target@IV.18.11.5.2\endcsname{ega4.c04.lem.18-11-5-2}
\expandafter\def\csname sga@target@IV.18.11.6\endcsname{ega4.c04.rem.18-11-6}
\expandafter\def\csname sga@target@IV.18.11.7\endcsname{ega4.c04.cor.18-11-7}
\expandafter\def\csname sga@target@IV.18.11.8\endcsname{ega4.c04.cor.18-11-8}
\expandafter\def\csname sga@target@IV.18.11.9\endcsname{ega4.c04.cor.18-11-9}
\expandafter\def\csname sga@target@IV.18.12.1\endcsname{ega4.c04.thm.18-12-1}
\expandafter\def\csname sga@target@IV.18.12.10\endcsname{ega4.c04.cor.18-12-10}
\expandafter\def\csname sga@target@IV.18.12.11\endcsname{ega4.c04.cor.18-12-11}
\expandafter\def\csname sga@target@IV.18.12.12\endcsname{ega4.c04.prop.18-12-12}
\expandafter\def\csname sga@target@IV.18.12.13\endcsname{ega4.c04.cor.18-12-13}
\expandafter\def\csname sga@target@IV.18.12.14\endcsname{ega4.c04.rem.18-12-14}
\expandafter\def\csname sga@target@IV.18.12.15\endcsname{ega4.c04.prop.18-12-15}
\expandafter\def\csname sga@target@IV.18.12.16\endcsname{ega4.c04.prop.18-12-16}
\expandafter\def\csname sga@target@IV.18.12.17\endcsname{ega4.c04.cor.18-12-17}
\expandafter\def\csname sga@target@IV.18.12.2\endcsname{ega4.c04.rem.18-12-2}
\expandafter\def\csname sga@target@IV.18.12.3\endcsname{ega4.c04.cor.18-12-3}
\expandafter\def\csname sga@target@IV.18.12.4\endcsname{ega4.c04.cor.18-12-4}
\expandafter\def\csname sga@target@IV.18.12.5\endcsname{ega4.c04.rem.18-12-5}
\expandafter\def\csname sga@target@IV.18.12.6\endcsname{ega4.c04.cor.18-12-6}
\expandafter\def\csname sga@target@IV.18.12.7\endcsname{ega4.c04.prop.18-12-7}
\expandafter\def\csname sga@target@IV.18.12.7.1\endcsname{ega4.c04.lem.18-12-7-1}
\expandafter\def\csname sga@target@IV.18.12.8\endcsname{ega4.c04.prop.18-12-8}
\expandafter\def\csname sga@target@IV.18.12.9\endcsname{ega4.c04.rem.18-12-9}
\expandafter\def\csname sga@target@IV.18.2.1\endcsname{ega4.c04.statement.18-2-1}
\expandafter\def\csname sga@target@IV.18.2.1.1\endcsname{ega4.c04.statement.18-2-1-1}
\expandafter\def\csname sga@target@IV.18.2.1.2\endcsname{ega4.c04.eq.18-2-1-2}
\expandafter\def\csname sga@target@IV.18.2.10\endcsname{ega4.c04.statement.18-2-10}
\expandafter\def\csname sga@target@IV.18.2.2\endcsname{ega4.c04.statement.18-2-2}
\expandafter\def\csname sga@target@IV.18.2.2.1\endcsname{ega4.c04.eq.18-2-2-1}
\expandafter\def\csname sga@target@IV.18.2.2.2\endcsname{ega4.c04.eq.18-2-2-2}
\expandafter\def\csname sga@target@IV.18.2.2.3\endcsname{ega4.c04.eq.18-2-2-3}
\expandafter\def\csname sga@target@IV.18.2.2.4\endcsname{ega4.c04.eq.18-2-2-4}
\expandafter\def\csname sga@target@IV.18.2.3\endcsname{ega4.c04.prop.18-2-3}
\expandafter\def\csname sga@target@IV.18.2.4\endcsname{ega4.c04.cor.18-2-4}
\expandafter\def\csname sga@target@IV.18.2.5\endcsname{ega4.c04.cor.18-2-5}
\expandafter\def\csname sga@target@IV.18.2.6\endcsname{ega4.c04.cor.18-2-6}
\expandafter\def\csname sga@target@IV.18.2.7\endcsname{ega4.c04.item.18-2-7}
\expandafter\def\csname sga@target@IV.18.2.7.1\endcsname{ega4.c04.eq.18-2-7-1}
\expandafter\def\csname sga@target@IV.18.2.7.2\endcsname{ega4.c04.eq.18-2-7-2}
\expandafter\def\csname sga@target@IV.18.2.8\endcsname{ega4.c04.prop.18-2-8}
\expandafter\def\csname sga@target@IV.18.2.9\endcsname{ega4.c04.cor.18-2-9}
\expandafter\def\csname sga@target@IV.18.3.1\endcsname{ega4.c04.prop.18-3-1}
\expandafter\def\csname sga@target@IV.18.3.2\endcsname{ega4.c04.prop.18-3-2}
\expandafter\def\csname sga@target@IV.18.3.2.1\endcsname{ega4.c04.lem.18-3-2-1}
\expandafter\def\csname sga@target@IV.18.3.2.2\endcsname{ega4.c04.eq.18-3-2-2}
\expandafter\def\csname sga@target@IV.18.3.3\endcsname{ega4.c04.statement.18-3-3}
\expandafter\def\csname sga@target@IV.18.3.4\endcsname{ega4.c04.thm.18-3-4}
\expandafter\def\csname sga@target@IV.18.4.1\endcsname{ega4.c04.lem.18-4-1}
\expandafter\def\csname sga@target@IV.18.4.10\endcsname{ega4.c04.prop.18-4-10}
\expandafter\def\csname sga@target@IV.18.4.11\endcsname{ega4.c04.statement.18-4-11}
\expandafter\def\csname sga@target@IV.18.4.12\endcsname{ega4.c04.cor.18-4-12}
\expandafter\def\csname sga@target@IV.18.4.13\endcsname{ega4.c04.cor.18-4-13}
\expandafter\def\csname sga@target@IV.18.4.14\endcsname{ega4.c04.cor.18-4-14}
\expandafter\def\csname sga@target@IV.18.4.2\endcsname{ega4.c04.prop.18-4-2}
\expandafter\def\csname sga@target@IV.18.4.3\endcsname{ega4.c04.cor.18-4-3}
\expandafter\def\csname sga@target@IV.18.4.4\endcsname{ega4.c04.lem.18-4-4}
\expandafter\def\csname sga@target@IV.18.4.5\endcsname{ega4.c04.prop.18-4-5}
\expandafter\def\csname sga@target@IV.18.4.6\endcsname{ega4.c04.thm.18-4-6}
\expandafter\def\csname sga@target@IV.18.4.6.1\endcsname{ega4.c04.lem.18-4-6-1}
\expandafter\def\csname sga@target@IV.18.4.7\endcsname{ega4.c04.cor.18-4-7}
\expandafter\def\csname sga@target@IV.18.4.8\endcsname{ega4.c04.cor.18-4-8}
\expandafter\def\csname sga@target@IV.18.4.9\endcsname{ega4.c04.prop.18-4-9}
\expandafter\def\csname sga@target@IV.18.5.1\endcsname{ega4.c04.statement.18-5-1}
\expandafter\def\csname sga@target@IV.18.5.1.1\endcsname{ega4.c04.eq.18-5-1-1}
\expandafter\def\csname sga@target@IV.18.5.10\endcsname{ega4.c04.prop.18-5-10}
\expandafter\def\csname sga@target@IV.18.5.11\endcsname{ega4.c04.thm.18-5-11}
\expandafter\def\csname sga@target@IV.18.5.12\endcsname{ega4.c04.cor.18-5-12}
\expandafter\def\csname sga@target@IV.18.5.13\endcsname{ega4.c04.rem.18-5-13}
\expandafter\def\csname sga@target@IV.18.5.13.1\endcsname{ega4.c04.lem.18-5-13-1}
\expandafter\def\csname sga@target@IV.18.5.14\endcsname{ega4.c04.prop.18-5-14}
\expandafter\def\csname sga@target@IV.18.5.15\endcsname{ega4.c04.prop.18-5-15}
\expandafter\def\csname sga@target@IV.18.5.16\endcsname{ega4.c04.item.18-5-16}
\expandafter\def\csname sga@target@IV.18.5.17\endcsname{ega4.c04.thm.18-5-17}
\expandafter\def\csname sga@target@IV.18.5.18\endcsname{ega4.c04.rem.18-5-18}
\expandafter\def\csname sga@target@IV.18.5.19\endcsname{ega4.c04.prop.18-5-19}
\expandafter\def\csname sga@target@IV.18.5.2\endcsname{ega4.c04.statement.18-5-2}
\expandafter\def\csname sga@target@IV.18.5.3\endcsname{ega4.c04.lem.18-5-3}
\expandafter\def\csname sga@target@IV.18.5.3.1\endcsname{ega4.c04.lem.18-5-3-1}
\expandafter\def\csname sga@target@IV.18.5.4\endcsname{ega4.c04.prop.18-5-4}
\expandafter\def\csname sga@target@IV.18.5.4.1\endcsname{ega4.c04.eq.18-5-4-1}
\expandafter\def\csname sga@target@IV.18.5.4.2\endcsname{ega4.c04.eq.18-5-4-2}
\expandafter\def\csname sga@target@IV.18.5.4.3\endcsname{ega4.c04.lem.18-5-4-3}
\expandafter\def\csname sga@target@IV.18.5.4.4\endcsname{ega4.c04.rem.18-5-4-4}
\expandafter\def\csname sga@target@IV.18.5.5\endcsname{ega4.c04.def.18-5-5}
\expandafter\def\csname sga@target@IV.18.5.6\endcsname{ega4.c04.prop.18-5-6}
\expandafter\def\csname sga@target@IV.18.5.7\endcsname{ega4.c04.rem.18-5-7}
\expandafter\def\csname sga@target@IV.18.5.8\endcsname{ega4.c04.def.18-5-8}
\expandafter\def\csname sga@target@IV.18.5.9\endcsname{ega4.c04.prop.18-5-9}
\expandafter\def\csname sga@target@IV.18.6.1\endcsname{ega4.c04.statement.18-6-1}
\expandafter\def\csname sga@target@IV.18.6.10\endcsname{ega4.c04.cor.18-6-10}
\expandafter\def\csname sga@target@IV.18.6.11\endcsname{ega4.c04.cor.18-6-11}
\expandafter\def\csname sga@target@IV.18.6.12\endcsname{ega4.c04.prop.18-6-12}
\expandafter\def\csname sga@target@IV.18.6.13\endcsname{ega4.c04.cor.18-6-13}
\expandafter\def\csname sga@target@IV.18.6.14\endcsname{ega4.c04.prop.18-6-14}
\expandafter\def\csname sga@target@IV.18.6.15\endcsname{ega4.c04.cor.18-6-15}
\expandafter\def\csname sga@target@IV.18.6.2\endcsname{ega4.c04.lem.18-6-2}
\expandafter\def\csname sga@target@IV.18.6.3\endcsname{ega4.c04.lem.18-6-3}
\expandafter\def\csname sga@target@IV.18.6.4\endcsname{ega4.c04.statement.18-6-4}
\expandafter\def\csname sga@target@IV.18.6.5\endcsname{ega4.c04.def.18-6-5}
\expandafter\def\csname sga@target@IV.18.6.6\endcsname{ega4.c04.thm.18-6-6}
\expandafter\def\csname sga@target@IV.18.6.7\endcsname{ega4.c04.statement.18-6-7}
\expandafter\def\csname sga@target@IV.18.6.8\endcsname{ega4.c04.prop.18-6-8}
\expandafter\def\csname sga@target@IV.18.6.9\endcsname{ega4.c04.thm.18-6-9}
\expandafter\def\csname sga@target@IV.18.6.9.1\endcsname{ega4.c04.lem.18-6-9-1}
\expandafter\def\csname sga@target@IV.18.6.9.2\endcsname{ega4.c04.lem.18-6-9-2}
\expandafter\def\csname sga@target@IV.18.7.1\endcsname{ega4.c04.statement.18-7-1}
\expandafter\def\csname sga@target@IV.18.7.2\endcsname{ega4.c04.prop.18-7-2}
\expandafter\def\csname sga@target@IV.18.7.3\endcsname{ega4.c04.cor.18-7-3}
\expandafter\def\csname sga@target@IV.18.7.4\endcsname{ega4.c04.cor.18-7-4}
\expandafter\def\csname sga@target@IV.18.7.5\endcsname{ega4.c04.prop.18-7-5}
\expandafter\def\csname sga@target@IV.18.7.5.1\endcsname{ega4.c04.lem.18-7-5-1}
\expandafter\def\csname sga@target@IV.18.7.6\endcsname{ega4.c04.cor.18-7-6}
\expandafter\def\csname sga@target@IV.18.7.7\endcsname{ega4.c04.rem.18-7-7}
\expandafter\def\csname sga@target@IV.18.8.1\endcsname{ega4.c04.prop.18-8-1}
\expandafter\def\csname sga@target@IV.18.8.10\endcsname{ega4.c04.prop.18-8-10}
\expandafter\def\csname sga@target@IV.18.8.11\endcsname{ega4.c04.rem.18-8-11}
\expandafter\def\csname sga@target@IV.18.8.12\endcsname{ega4.c04.prop.18-8-12}
\expandafter\def\csname sga@target@IV.18.8.13\endcsname{ega4.c04.cor.18-8-13}
\expandafter\def\csname sga@target@IV.18.8.14\endcsname{ega4.c04.cor.18-8-14}
\expandafter\def\csname sga@target@IV.18.8.15\endcsname{ega4.c04.prop.18-8-15}
\expandafter\def\csname sga@target@IV.18.8.16\endcsname{ega4.c04.cor.18-8-16}
\expandafter\def\csname sga@target@IV.18.8.17\endcsname{ega4.c04.prop.18-8-17}
\expandafter\def\csname sga@target@IV.18.8.18\endcsname{ega4.c04.prop.18-8-18}
\expandafter\def\csname sga@target@IV.18.8.2\endcsname{ega4.c04.def.18-8-2}
\expandafter\def\csname sga@target@IV.18.8.3\endcsname{ega4.c04.rem.18-8-3}
\expandafter\def\csname sga@target@IV.18.8.4\endcsname{ega4.c04.lem.18-8-4}
\expandafter\def\csname sga@target@IV.18.8.5\endcsname{ega4.c04.lem.18-8-5}
\expandafter\def\csname sga@target@IV.18.8.6\endcsname{ega4.c04.statement.18-8-6}
\expandafter\def\csname sga@target@IV.18.8.6.1\endcsname{ega4.c04.eq.18-8-6-1}
\expandafter\def\csname sga@target@IV.18.8.6.2\endcsname{ega4.c04.eq.18-8-6-2}
\expandafter\def\csname sga@target@IV.18.8.7\endcsname{ega4.c04.def.18-8-7}
\expandafter\def\csname sga@target@IV.18.8.8\endcsname{ega4.c04.prop.18-8-8}
\expandafter\def\csname sga@target@IV.18.8.8.1\endcsname{ega4.c04.statement.18-8-8-1}
\expandafter\def\csname sga@target@IV.18.8.9\endcsname{ega4.c04.statement.18-8-9}
\expandafter\def\csname sga@target@IV.18.9.1\endcsname{ega4.c04.thm.18-9-1}
\expandafter\def\csname sga@target@IV.18.9.10\endcsname{ega4.c04.rem.18-9-10}
\expandafter\def\csname sga@target@IV.18.9.11\endcsname{ega4.c04.cor.18-9-11}
\expandafter\def\csname sga@target@IV.18.9.2\endcsname{ega4.c04.cor.18-9-2}
\expandafter\def\csname sga@target@IV.18.9.3\endcsname{ega4.c04.cor.18-9-3}
\expandafter\def\csname sga@target@IV.18.9.4\endcsname{ega4.c04.cor.18-9-4}
\expandafter\def\csname sga@target@IV.18.9.5\endcsname{ega4.c04.cor.18-9-5}
\expandafter\def\csname sga@target@IV.18.9.5.1\endcsname{ega4.c04.eq.18-9-5-1}
\expandafter\def\csname sga@target@IV.18.9.6\endcsname{ega4.c04.statement.18-9-6}
\expandafter\def\csname sga@target@IV.18.9.7\endcsname{ega4.c04.thm.18-9-7}
\expandafter\def\csname sga@target@IV.18.9.7.1\endcsname{ega4.c04.rem.18-9-7-1}
\expandafter\def\csname sga@target@IV.18.9.7.2\endcsname{ega4.c04.statement.18-9-7-2}
\expandafter\def\csname sga@target@IV.18.9.7.3\endcsname{ega4.c04.lem.18-9-7-3}
\expandafter\def\csname sga@target@IV.18.9.7.4\endcsname{ega4.c04.eq.18-9-7-4}
\expandafter\def\csname sga@target@IV.18.9.7.5\endcsname{ega4.c04.lem.18-9-7-5}
\expandafter\def\csname sga@target@IV.18.9.7.6\endcsname{ega4.c04.rem.18-9-7-6}
\expandafter\def\csname sga@target@IV.18.9.7.7\endcsname{ega4.c04.lem.18-9-7-7}
\expandafter\def\csname sga@target@IV.18.9.7.8\endcsname{ega4.c04.rem.18-9-7-8}
\expandafter\def\csname sga@target@IV.18.9.8\endcsname{ega4.c04.cor.18-9-8}
\expandafter\def\csname sga@target@IV.18.9.9\endcsname{ega4.c04.cor.18-9-9}
\expandafter\def\csname sga@target@IV.19.1.1\endcsname{ega4.c04.prop.19-1-1}
\expandafter\def\csname sga@target@IV.19.1.2\endcsname{ega4.c04.cor.19-1-2}
\expandafter\def\csname sga@target@IV.19.1.3\endcsname{ega4.c04.def.19-1-3}
\expandafter\def\csname sga@target@IV.19.1.4\endcsname{ega4.c04.prop.19-1-4}
\expandafter\def\csname sga@target@IV.19.1.4.1\endcsname{ega4.c04.eq.19-1-4-1}
\expandafter\def\csname sga@target@IV.19.1.5\endcsname{ega4.c04.prop.19-1-5}
\expandafter\def\csname sga@target@IV.19.1.5.1\endcsname{ega4.c04.item.19-1-5-1}
\expandafter\def\csname sga@target@IV.19.1.6\endcsname{ega4.c04.cor.19-1-6}
\expandafter\def\csname sga@target@IV.19.1.6.1\endcsname{ega4.c04.lem.19-1-6-1}
\expandafter\def\csname sga@target@IV.19.1.7\endcsname{ega4.c04.item.19-1-7}
\expandafter\def\csname sga@target@IV.19.1.8\endcsname{ega4.c04.statement.19-1-8}
\expandafter\def\csname sga@target@IV.19.2.1\endcsname{ega4.c04.def.19-2-1}
\expandafter\def\csname sga@target@IV.19.2.2\endcsname{ega4.c04.def.19-2-2}
\expandafter\def\csname sga@target@IV.19.2.3\endcsname{ega4.c04.prop.19-2-3}
\expandafter\def\csname sga@target@IV.19.2.4\endcsname{ega4.c04.prop.19-2-4}
\expandafter\def\csname sga@target@IV.19.2.5\endcsname{ega4.c04.item.19-2-5}
\expandafter\def\csname sga@target@IV.19.2.6\endcsname{ega4.c04.cor.19-2-6}
\expandafter\def\csname sga@target@IV.19.2.7\endcsname{ega4.c04.prop.19-2-7}
\expandafter\def\csname sga@target@IV.19.2.8\endcsname{ega4.c04.prop.19-2-8}
\expandafter\def\csname sga@target@IV.19.2.9\endcsname{ega4.c04.prop.19-2-9}
\expandafter\def\csname sga@target@IV.19.3.1\endcsname{ega4.c04.def.19-3-1}
\expandafter\def\csname sga@target@IV.19.3.10\endcsname{ega4.c04.cor.19-3-10}
\expandafter\def\csname sga@target@IV.19.3.2\endcsname{ega4.c04.prop.19-3-2}
\expandafter\def\csname sga@target@IV.19.3.2.1\endcsname{ega4.c04.lem.19-3-2-1}
\expandafter\def\csname sga@target@IV.19.3.3\endcsname{ega4.c04.cor.19-3-3}
\expandafter\def\csname sga@target@IV.19.3.4\endcsname{ega4.c04.cor.19-3-4}
\expandafter\def\csname sga@target@IV.19.3.5\endcsname{ega4.c04.cor.19-3-5}
\expandafter\def\csname sga@target@IV.19.3.6\endcsname{ega4.c04.def.19-3-6}
\expandafter\def\csname sga@target@IV.19.3.7\endcsname{ega4.c04.prop.19-3-7}
\expandafter\def\csname sga@target@IV.19.3.8\endcsname{ega4.c04.cor.19-3-8}
\expandafter\def\csname sga@target@IV.19.3.9\endcsname{ega4.c04.prop.19-3-9}
\expandafter\def\csname sga@target@IV.19.4.1\endcsname{ega4.c04.statement.19-4-1}
\expandafter\def\csname sga@target@IV.19.4.1.1\endcsname{ega4.c04.subsec.19-4-1-1}
\expandafter\def\csname sga@target@IV.19.4.10\endcsname{ega4.c04.rem.19-4-10}
\expandafter\def\csname sga@target@IV.19.4.11\endcsname{ega4.c04.prop.19-4-11}
\expandafter\def\csname sga@target@IV.19.4.11.1\endcsname{ega4.c04.eq.19-4-11-1}
\expandafter\def\csname sga@target@IV.19.4.2\endcsname{ega4.c04.lem.19-4-2}
\expandafter\def\csname sga@target@IV.19.4.3\endcsname{ega4.c04.cor.19-4-3}
\expandafter\def\csname sga@target@IV.19.4.4\endcsname{ega4.c04.prop.19-4-4}
\expandafter\def\csname sga@target@IV.19.4.5\endcsname{ega4.c04.cor.19-4-5}
\expandafter\def\csname sga@target@IV.19.4.6\endcsname{ega4.c04.prop.19-4-6}
\expandafter\def\csname sga@target@IV.19.4.7\endcsname{ega4.c04.cor.19-4-7}
\expandafter\def\csname sga@target@IV.19.4.8\endcsname{ega4.c04.prop.19-4-8}
\expandafter\def\csname sga@target@IV.19.4.9\endcsname{ega4.c04.cor.19-4-9}
\expandafter\def\csname sga@target@IV.19.5.1\endcsname{ega4.c04.thm.19-5-1}
\expandafter\def\csname sga@target@IV.19.5.1.1\endcsname{ega4.c04.eq.19-5-1-1}
\expandafter\def\csname sga@target@IV.19.5.2\endcsname{ega4.c04.cor.19-5-2}
\expandafter\def\csname sga@target@IV.19.5.3\endcsname{ega4.c04.prop.19-5-3}
\expandafter\def\csname sga@target@IV.19.5.4\endcsname{ega4.c04.statement.19-5-4}
\expandafter\def\csname sga@target@IV.19.5.5\endcsname{ega4.c04.thm.19-5-5}
\expandafter\def\csname sga@target@IV.19.5.5.1\endcsname{ega4.c04.statement.19-5-5-1}
\expandafter\def\csname sga@target@IV.19.5.5.2\endcsname{ega4.c04.statement.19-5-5-2}
\expandafter\def\csname sga@target@IV.19.5.5.3\endcsname{ega4.c04.statement.19-5-5-3}
\expandafter\def\csname sga@target@IV.19.5.5.4\endcsname{ega4.c04.statement.19-5-5-4}
\expandafter\def\csname sga@target@IV.19.5.5.5\endcsname{ega4.c04.statement.19-5-5-5}
\expandafter\def\csname sga@target@IV.19.5.6\endcsname{ega4.c04.statement.19-5-6}
\expandafter\def\csname sga@target@IV.19.6.1\endcsname{ega4.c04.statement.19-6-1}
\expandafter\def\csname sga@target@IV.19.6.1.1\endcsname{ega4.c04.eq.19-6-1-1}
\expandafter\def\csname sga@target@IV.19.6.2\endcsname{ega4.c04.statement.19-6-2}
\expandafter\def\csname sga@target@IV.19.6.2.1\endcsname{ega4.c04.statement.19-6-2-1}
\expandafter\def\csname sga@target@IV.19.6.3\endcsname{ega4.c04.prop.19-6-3}
\expandafter\def\csname sga@target@IV.19.6.4\endcsname{ega4.c04.cor.19-6-4}
\expandafter\def\csname sga@target@IV.19.7\endcsname{ega4.c04.subsec.19-7.key-iv-19-7}
\expandafter\def\csname sga@target@IV.19.7.1\endcsname{ega4.c04.thm.19-7-1}
\expandafter\def\csname sga@target@IV.19.7.1.1\endcsname{ega4.c04.statement.19-7-1-1}
\expandafter\def\csname sga@target@IV.19.7.1.10\endcsname{ega4.c04.statement.19-7-1-10}
\expandafter\def\csname sga@target@IV.19.7.1.11\endcsname{ega4.c04.statement.19-7-1-11}
\expandafter\def\csname sga@target@IV.19.7.1.12\endcsname{ega4.c04.statement.19-7-1-12}
\expandafter\def\csname sga@target@IV.19.7.1.13\endcsname{ega4.c04.statement.19-7-1-13}
\expandafter\def\csname sga@target@IV.19.7.1.14\endcsname{ega4.c04.statement.19-7-1-14}
\expandafter\def\csname sga@target@IV.19.7.1.15\endcsname{ega4.c04.statement.19-7-1-15}
\expandafter\def\csname sga@target@IV.19.7.1.16\endcsname{ega4.c04.statement.19-7-1-16}
\expandafter\def\csname sga@target@IV.19.7.1.17\endcsname{ega4.c04.statement.19-7-1-17}
\expandafter\def\csname sga@target@IV.19.7.1.18\endcsname{ega4.c04.statement.19-7-1-18}
\expandafter\def\csname sga@target@IV.19.7.1.19\endcsname{ega4.c04.statement.19-7-1-19}
\expandafter\def\csname sga@target@IV.19.7.1.2\endcsname{ega4.c04.statement.19-7-1-2}
\expandafter\def\csname sga@target@IV.19.7.1.20\endcsname{ega4.c04.statement.19-7-1-20}
\expandafter\def\csname sga@target@IV.19.7.1.21\endcsname{ega4.c04.statement.19-7-1-21}
\expandafter\def\csname sga@target@IV.19.7.1.22\endcsname{ega4.c04.statement.19-7-1-22}
\expandafter\def\csname sga@target@IV.19.7.1.23\endcsname{ega4.c04.statement.19-7-1-23}
\expandafter\def\csname sga@target@IV.19.7.1.3\endcsname{ega4.c04.statement.19-7-1-3}
\expandafter\def\csname sga@target@IV.19.7.1.5\endcsname{ega4.c04.statement.19-7-1-5}
\expandafter\def\csname sga@target@IV.19.7.1.6\endcsname{ega4.c04.statement.19-7-1-6}
\expandafter\def\csname sga@target@IV.19.7.1.7\endcsname{ega4.c04.lem.19-7-1-7}
\expandafter\def\csname sga@target@IV.19.7.1.8\endcsname{ega4.c04.statement.19-7-1-8}
\expandafter\def\csname sga@target@IV.19.7.1.9\endcsname{ega4.c04.statement.19-7-1-9}
\expandafter\def\csname sga@target@IV.19.7.2\endcsname{ega4.c04.statement.19-7-2}
\expandafter\def\csname sga@target@IV.19.7.3\endcsname{ega4.c04.cor.19-7-3}
\expandafter\def\csname sga@target@IV.19.7.3.1\endcsname{ega4.c04.eq.19-7-3-1}
\expandafter\def\csname sga@target@IV.19.7.4\endcsname{ega4.c04.cor.19-7-4}
\expandafter\def\csname sga@target@IV.19.7.4.1\endcsname{ega4.c04.eq.19-7-4-1}
\expandafter\def\csname sga@target@IV.19.7.5\endcsname{ega4.c04.cor.19-7-5}
\expandafter\def\csname sga@target@IV.19.8.1\endcsname{ega4.c04.prop.19-8-1}
\expandafter\def\csname sga@target@IV.19.8.2\endcsname{ega4.c04.prop.19-8-2}
\expandafter\def\csname sga@target@IV.19.8.3\endcsname{ega4.c04.rem.19-8-3}
\expandafter\def\csname sga@target@IV.19.9.1\endcsname{ega4.c04.statement.19-9-1}
\expandafter\def\csname sga@target@IV.19.9.1.1\endcsname{ega4.c04.eq.19-9-1-1}
\expandafter\def\csname sga@target@IV.19.9.1.2\endcsname{ega4.c04.eq.19-9-1-2}
\expandafter\def\csname sga@target@IV.19.9.1.3\endcsname{ega4.c04.eq.19-9-1-3}
\expandafter\def\csname sga@target@IV.19.9.2\endcsname{ega4.c04.lem.19-9-2}
\expandafter\def\csname sga@target@IV.19.9.3\endcsname{ega4.c04.prop.19-9-3}
\expandafter\def\csname sga@target@IV.19.9.4\endcsname{ega4.c04.cor.19-9-4}
\expandafter\def\csname sga@target@IV.19.9.5\endcsname{ega4.c04.prop.19-9-5}
\expandafter\def\csname sga@target@IV.19.9.6\endcsname{ega4.c04.prop.19-9-6}
\expandafter\def\csname sga@target@IV.19.9.7\endcsname{ega4.c04.cor.19-9-7}
\expandafter\def\csname sga@target@IV.19.9.8\endcsname{ega4.c04.prop.19-9-8}
\expandafter\def\csname sga@target@IV.19.9.9\endcsname{ega4.c04.rem.19-9-9}
\expandafter\def\csname sga@target@IV.2.1.1\endcsname{ega4.c04.statement.2-1-1}
\expandafter\def\csname sga@target@IV.2.1.10\endcsname{ega4.c04.cor.2-1-10}
\expandafter\def\csname sga@target@IV.2.1.11\endcsname{ega4.c04.prop.2-1-11}
\expandafter\def\csname sga@target@IV.2.1.12\endcsname{ega4.c04.prop.2-1-12}
\expandafter\def\csname sga@target@IV.2.1.13\endcsname{ega4.c04.prop.2-1-13}
\expandafter\def\csname sga@target@IV.2.1.14\endcsname{ega4.c04.prop.2-1-14}
\expandafter\def\csname sga@target@IV.2.1.2\endcsname{ega4.c04.prop.2-1-2}
\expandafter\def\csname sga@target@IV.2.1.3\endcsname{ega4.c04.prop.2-1-3}
\expandafter\def\csname sga@target@IV.2.1.4\endcsname{ega4.c04.prop.2-1-4}
\expandafter\def\csname sga@target@IV.2.1.5\endcsname{ega4.c04.prop.2-1-5}
\expandafter\def\csname sga@target@IV.2.1.6\endcsname{ega4.c04.cor.2-1-6}
\expandafter\def\csname sga@target@IV.2.1.7\endcsname{ega4.c04.cor.2-1-7}
\expandafter\def\csname sga@target@IV.2.1.8\endcsname{ega4.c04.prop.2-1-8}
\expandafter\def\csname sga@target@IV.2.1.9\endcsname{ega4.c04.cor.2-1-9}
\expandafter\def\csname sga@target@IV.2.2.1\endcsname{ega4.c04.prop.2-2-1}
\expandafter\def\csname sga@target@IV.2.2.10\endcsname{ega4.c04.prop.2-2-10}
\expandafter\def\csname sga@target@IV.2.2.10.1\endcsname{ega4.c04.eq.2-2-10-1}
\expandafter\def\csname sga@target@IV.2.2.11\endcsname{ega4.c04.cor.2-2-11}
\expandafter\def\csname sga@target@IV.2.2.12\endcsname{ega4.c04.cor.2-2-12}
\expandafter\def\csname sga@target@IV.2.2.13\endcsname{ega4.c04.cor.2-2-13}
\expandafter\def\csname sga@target@IV.2.2.14\endcsname{ega4.c04.prop.2-2-14}
\expandafter\def\csname sga@target@IV.2.2.15\endcsname{ega4.c04.prop.2-2-15}
\expandafter\def\csname sga@target@IV.2.2.16\endcsname{ega4.c04.cor.2-2-16}
\expandafter\def\csname sga@target@IV.2.2.17\endcsname{ega4.c04.prop.2-2-17}
\expandafter\def\csname sga@target@IV.2.2.2\endcsname{ega4.c04.prop.2-2-2}
\expandafter\def\csname sga@target@IV.2.2.2.1\endcsname{ega4.c04.eq.2-2-2-1}
\expandafter\def\csname sga@target@IV.2.2.3\endcsname{ega4.c04.cor.2-2-3}
\expandafter\def\csname sga@target@IV.2.2.4\endcsname{ega4.c04.def.2-2-4}
\expandafter\def\csname sga@target@IV.2.2.5\endcsname{ega4.c04.cor.2-2-5}
\expandafter\def\csname sga@target@IV.2.2.6\endcsname{ega4.c04.cor.2-2-6}
\expandafter\def\csname sga@target@IV.2.2.7\endcsname{ega4.c04.prop.2-2-7}
\expandafter\def\csname sga@target@IV.2.2.8\endcsname{ega4.c04.cor.2-2-8}
\expandafter\def\csname sga@target@IV.2.2.8.1\endcsname{ega4.c04.eq.2-2-8-1}
\expandafter\def\csname sga@target@IV.2.2.9\endcsname{ega4.c04.rem.2-2-9}
\expandafter\def\csname sga@target@IV.2.3.1\endcsname{ega4.c04.lem.2-3-1}
\expandafter\def\csname sga@target@IV.2.3.1.1\endcsname{ega4.c04.eq.2-3-1-1}
\expandafter\def\csname sga@target@IV.2.3.10\endcsname{ega4.c04.thm.2-3-10}
\expandafter\def\csname sga@target@IV.2.3.11\endcsname{ega4.c04.cor.2-3-11}
\expandafter\def\csname sga@target@IV.2.3.12\endcsname{ega4.c04.cor.2-3-12}
\expandafter\def\csname sga@target@IV.2.3.13\endcsname{ega4.c04.cor.2-3-13}
\expandafter\def\csname sga@target@IV.2.3.14\endcsname{ega4.c04.cor.2-3-14}
\expandafter\def\csname sga@target@IV.2.3.15\endcsname{ega4.c04.rem.2-3-15}
\expandafter\def\csname sga@target@IV.2.3.2\endcsname{ega4.c04.prop.2-3-2}
\expandafter\def\csname sga@target@IV.2.3.3\endcsname{ega4.c04.statement.2-3-3}
\expandafter\def\csname sga@target@IV.2.3.4\endcsname{ega4.c04.prop.2-3-4}
\expandafter\def\csname sga@target@IV.2.3.5\endcsname{ega4.c04.cor.2-3-5}
\expandafter\def\csname sga@target@IV.2.3.6\endcsname{ega4.c04.prop.2-3-6}
\expandafter\def\csname sga@target@IV.2.3.7\endcsname{ega4.c04.prop.2-3-7}
\expandafter\def\csname sga@target@IV.2.3.8\endcsname{ega4.c04.cor.2-3-8}
\expandafter\def\csname sga@target@IV.2.3.9\endcsname{ega4.c04.rem.2-3-9}
\expandafter\def\csname sga@target@IV.2.4.1\endcsname{ega4.c04.statement.2-4-1}
\expandafter\def\csname sga@target@IV.2.4.10\endcsname{ega4.c04.cor.2-4-10}
\expandafter\def\csname sga@target@IV.2.4.11\endcsname{ega4.c04.rem.2-4-11}
\expandafter\def\csname sga@target@IV.2.4.2\endcsname{ega4.c04.def.2-4-2}
\expandafter\def\csname sga@target@IV.2.4.3\endcsname{ega4.c04.prop.2-4-3}
\expandafter\def\csname sga@target@IV.2.4.4\endcsname{ega4.c04.prop.2-4-4}
\expandafter\def\csname sga@target@IV.2.4.5\endcsname{ega4.c04.prop.2-4-5}
\expandafter\def\csname sga@target@IV.2.4.6\endcsname{ega4.c04.thm.2-4-6}
\expandafter\def\csname sga@target@IV.2.4.7\endcsname{ega4.c04.cor.2-4-7}
\expandafter\def\csname sga@target@IV.2.4.8\endcsname{ega4.c04.rem.2-4-8}
\expandafter\def\csname sga@target@IV.2.4.9\endcsname{ega4.c04.prop.2-4-9}
\expandafter\def\csname sga@target@IV.2.5.1\endcsname{ega4.c04.prop.2-5-1}
\expandafter\def\csname sga@target@IV.2.5.2\endcsname{ega4.c04.prop.2-5-2}
\expandafter\def\csname sga@target@IV.2.5.3\endcsname{ega4.c04.lem.2-5-3}
\expandafter\def\csname sga@target@IV.2.5.4\endcsname{ega4.c04.rem.2-5-4}
\expandafter\def\csname sga@target@IV.2.5.4.1\endcsname{ega4.c04.lem.2-5-4-1}
\expandafter\def\csname sga@target@IV.2.5.5\endcsname{ega4.c04.prop.2-5-5}
\expandafter\def\csname sga@target@IV.2.5.5.1\endcsname{ega4.c04.statement.2-5-5-1}
\expandafter\def\csname sga@target@IV.2.5.5.2\endcsname{ega4.c04.lem.2-5-5-2}
\expandafter\def\csname sga@target@IV.2.5.5.3\endcsname{ega4.c04.statement.2-5-5-3}
\expandafter\def\csname sga@target@IV.2.5.5.4\endcsname{ega4.c04.rem.2-5-5-4}
\expandafter\def\csname sga@target@IV.2.5.6\endcsname{ega4.c04.lem.2-5-6}
\expandafter\def\csname sga@target@IV.2.5.7\endcsname{ega4.c04.cor.2-5-7}
\expandafter\def\csname sga@target@IV.2.5.8\endcsname{ega4.c04.prop.2-5-8}
\expandafter\def\csname sga@target@IV.2.5.8.1\endcsname{ega4.c04.statement.2-5-8-1}
\expandafter\def\csname sga@target@IV.2.5.8.2\endcsname{ega4.c04.statement.2-5-8-2}
\expandafter\def\csname sga@target@IV.2.5.8.3\endcsname{ega4.c04.statement.2-5-8-3}
\expandafter\def\csname sga@target@IV.2.5.8.4\endcsname{ega4.c04.statement.2-5-8-4}
\expandafter\def\csname sga@target@IV.2.5.8.5\endcsname{ega4.c04.statement.2-5-8-5}
\expandafter\def\csname sga@target@IV.2.5.8.6\endcsname{ega4.c04.statement.2-5-8-6}
\expandafter\def\csname sga@target@IV.2.5.8.7\endcsname{ega4.c04.statement.2-5-8-7}
\expandafter\def\csname sga@target@IV.2.5.8.8\endcsname{ega4.c04.statement.2-5-8-8}
\expandafter\def\csname sga@target@IV.2.5.8.9\endcsname{ega4.c04.statement.2-5-8-9}
\expandafter\def\csname sga@target@IV.2.6.1\endcsname{ega4.c04.prop.2-6-1}
\expandafter\def\csname sga@target@IV.2.6.2\endcsname{ega4.c04.prop.2-6-2}
\expandafter\def\csname sga@target@IV.2.6.3\endcsname{ega4.c04.rem.2-6-3}
\expandafter\def\csname sga@target@IV.2.6.4\endcsname{ega4.c04.cor.2-6-4}
\expandafter\def\csname sga@target@IV.2.7.1\endcsname{ega4.c04.prop.2-7-1}
\expandafter\def\csname sga@target@IV.2.7.1.1\endcsname{ega4.c04.lem.2-7-1-1}
\expandafter\def\csname sga@target@IV.2.7.1.viii\endcsname{ega4.c04.item.2-7-1-viii}
\expandafter\def\csname sga@target@IV.2.7.2\endcsname{ega4.c04.cor.2-7-2}
\expandafter\def\csname sga@target@IV.2.7.3\endcsname{ega4.c04.sec.2-7-3}
\expandafter\def\csname sga@target@IV.2.8.1\endcsname{ega4.c04.prop.2-8-1}
\expandafter\def\csname sga@target@IV.2.8.1.1\endcsname{ega4.c04.lem.2-8-1-1}
\expandafter\def\csname sga@target@IV.2.8.2\endcsname{ega4.c04.cor.2-8-2}
\expandafter\def\csname sga@target@IV.2.8.3\endcsname{ega4.c04.statement.2-8-3}
\expandafter\def\csname sga@target@IV.2.8.4\endcsname{ega4.c04.prop.2-8-4}
\expandafter\def\csname sga@target@IV.2.8.4.1\endcsname{ega4.c04.statement.2-8-4-1}
\expandafter\def\csname sga@target@IV.2.8.5\endcsname{ega4.c04.prop.2-8-5}
\expandafter\def\csname sga@target@IV.2.8.6\endcsname{ega4.c04.cor.2-8-6}
\expandafter\def\csname sga@target@IV.2.8.6.1\endcsname{ega4.c04.statement.2-8-6-1}
\expandafter\def\csname sga@target@IV.20\endcsname{ega4.c04.sec.20}
\expandafter\def\csname sga@target@IV.20.0\endcsname{ega4.c04.subsec.20-0}
\expandafter\def\csname sga@target@IV.20.1\endcsname{ega4.c04.subsec.20-1}
\expandafter\def\csname sga@target@IV.20.1.1\endcsname{ega4.c04.statement.20-1-1}
\expandafter\def\csname sga@target@IV.20.1.1.1\endcsname{ega4.c04.eq.20-1-1-1}
\expandafter\def\csname sga@target@IV.20.1.10\endcsname{ega4.c04.statement.20-1-10}
\expandafter\def\csname sga@target@IV.20.1.11\endcsname{ega4.c04.statement.20-1-11}
\expandafter\def\csname sga@target@IV.20.1.11.1\endcsname{ega4.c04.eq.20-1-11-1}
\expandafter\def\csname sga@target@IV.20.1.11.2\endcsname{ega4.c04.eq.20-1-11-2}
\expandafter\def\csname sga@target@IV.20.1.12\endcsname{ega4.c04.prop.20-1-12}
\expandafter\def\csname sga@target@IV.20.1.13\endcsname{ega4.c04.rem.20-1-13}
\expandafter\def\csname sga@target@IV.20.1.2\endcsname{ega4.c04.statement.20-1-2}
\expandafter\def\csname sga@target@IV.20.1.2.1\endcsname{ega4.c04.statement.20-1-2-1}
\expandafter\def\csname sga@target@IV.20.1.2.2\endcsname{ega4.c04.statement.20-1-2-2}
\expandafter\def\csname sga@target@IV.20.1.3\endcsname{ega4.c04.statement.20-1-3}
\expandafter\def\csname sga@target@IV.20.1.3.1\endcsname{ega4.c04.statement.20-1-3-1}
\expandafter\def\csname sga@target@IV.20.1.4\endcsname{ega4.c04.statement.20-1-4}
\expandafter\def\csname sga@target@IV.20.1.4.1\endcsname{ega4.c04.statement.20-1-4-1}
\expandafter\def\csname sga@target@IV.20.1.5\endcsname{ega4.c04.statement.20-1-5}
\expandafter\def\csname sga@target@IV.20.1.5.1\endcsname{ega4.c04.statement.20-1-5-1}
\expandafter\def\csname sga@target@IV.20.1.6\endcsname{ega4.c04.prop.20-1-6}
\expandafter\def\csname sga@target@IV.20.1.7\endcsname{ega4.c04.statement.20-1-7}
\expandafter\def\csname sga@target@IV.20.1.8\endcsname{ega4.c04.statement.20-1-8}
\expandafter\def\csname sga@target@IV.20.1.9\endcsname{ega4.c04.statement.20-1-9}
\expandafter\def\csname sga@target@IV.20.2\endcsname{ega4.c04.subsec.20-2}
\expandafter\def\csname sga@target@IV.20.2.1\endcsname{ega4.c04.statement.20-2-1}
\expandafter\def\csname sga@target@IV.20.2.10\endcsname{ega4.c04.statement.20-2-10}
\expandafter\def\csname sga@target@IV.20.2.10.1\endcsname{ega4.c04.eq.20-2-10-1}
\expandafter\def\csname sga@target@IV.20.2.10.2\endcsname{ega4.c04.eq.20-2-10-2}
\expandafter\def\csname sga@target@IV.20.2.11\endcsname{ega4.c04.prop.20-2-11}
\expandafter\def\csname sga@target@IV.20.2.11.1\endcsname{ega4.c04.eq.20-2-11-1}
\expandafter\def\csname sga@target@IV.20.2.12\endcsname{ega4.c04.cor.20-2-12}
\expandafter\def\csname sga@target@IV.20.2.13\endcsname{ega4.c04.statement.20-2-13}
\expandafter\def\csname sga@target@IV.20.2.14\endcsname{ega4.c04.statement.20-2-14}
\expandafter\def\csname sga@target@IV.20.2.15\endcsname{ega4.c04.prop.20-2-15}
\expandafter\def\csname sga@target@IV.20.2.16\endcsname{ega4.c04.rem.20-2-16}
\expandafter\def\csname sga@target@IV.20.2.2\endcsname{ega4.c04.statement.20-2-2}
\expandafter\def\csname sga@target@IV.20.2.3\endcsname{ega4.c04.statement.20-2-3}
\expandafter\def\csname sga@target@IV.20.2.3.1\endcsname{ega4.c04.statement.20-2-3-1}
\expandafter\def\csname sga@target@IV.20.2.4\endcsname{ega4.c04.prop.20-2-4}
\expandafter\def\csname sga@target@IV.20.2.5\endcsname{ega4.c04.cor.20-2-5}
\expandafter\def\csname sga@target@IV.20.2.6\endcsname{ega4.c04.cor.20-2-6}
\expandafter\def\csname sga@target@IV.20.2.7\endcsname{ega4.c04.rem.20-2-7}
\expandafter\def\csname sga@target@IV.20.2.8\endcsname{ega4.c04.statement.20-2-8}
\expandafter\def\csname sga@target@IV.20.2.9\endcsname{ega4.c04.lem.20-2-9}
\expandafter\def\csname sga@target@IV.20.3\endcsname{ega4.c04.subsec.20-3}
\expandafter\def\csname sga@target@IV.20.3.1\endcsname{ega4.c04.statement.20-3-1}
\expandafter\def\csname sga@target@IV.20.3.10\endcsname{ega4.c04.lem.20-3-10}
\expandafter\def\csname sga@target@IV.20.3.11\endcsname{ega4.c04.prop.20-3-11}
\expandafter\def\csname sga@target@IV.20.3.11.1\endcsname{ega4.c04.eq.20-3-11-1}
\expandafter\def\csname sga@target@IV.20.3.12\endcsname{ega4.c04.cor.20-3-12}
\expandafter\def\csname sga@target@IV.20.3.13\endcsname{ega4.c04.rem.20-3-13}
\expandafter\def\csname sga@target@IV.20.3.2\endcsname{ega4.c04.statement.20-3-2}
\expandafter\def\csname sga@target@IV.20.3.2.1\endcsname{ega4.c04.eq.20-3-2-1}
\expandafter\def\csname sga@target@IV.20.3.3\endcsname{ega4.c04.statement.20-3-3}
\expandafter\def\csname sga@target@IV.20.3.4\endcsname{ega4.c04.statement.20-3-4}
\expandafter\def\csname sga@target@IV.20.3.5\endcsname{ega4.c04.prop.20-3-5}
\expandafter\def\csname sga@target@IV.20.3.6\endcsname{ega4.c04.statement.20-3-6}
\expandafter\def\csname sga@target@IV.20.3.6.1\endcsname{ega4.c04.eq.20-3-6-1}
\expandafter\def\csname sga@target@IV.20.3.7\endcsname{ega4.c04.prop.20-3-7}
\expandafter\def\csname sga@target@IV.20.3.8\endcsname{ega4.c04.prop.20-3-8}
\expandafter\def\csname sga@target@IV.20.3.8.1\endcsname{ega4.c04.eq.20-3-8-1}
\expandafter\def\csname sga@target@IV.20.3.8.2\endcsname{ega4.c04.lem.20-3-8-2}
\expandafter\def\csname sga@target@IV.20.3.8.3\endcsname{ega4.c04.lem.20-3-8-3}
\expandafter\def\csname sga@target@IV.20.3.8.4\endcsname{ega4.c04.lem.20-3-8-4}
\expandafter\def\csname sga@target@IV.20.3.8.5\endcsname{ega4.c04.sec.20-3-8-5}
\expandafter\def\csname sga@target@IV.20.3.8.6\endcsname{ega4.c04.sec.20-3-8-6}
\expandafter\def\csname sga@target@IV.20.3.8.7\endcsname{ega4.c04.sec.20-3-8-7}
\expandafter\def\csname sga@target@IV.20.3.8.8\endcsname{ega4.c04.lem.20-3-8-8}
\expandafter\def\csname sga@target@IV.20.3.9\endcsname{ega4.c04.cor.20-3-9}
\expandafter\def\csname sga@target@IV.20.4\endcsname{ega4.c04.subsec.20-4}
\expandafter\def\csname sga@target@IV.20.4.1\endcsname{ega4.c04.statement.20-4-1}
\expandafter\def\csname sga@target@IV.20.4.10\endcsname{ega4.c04.cor.20-4-10}
\expandafter\def\csname sga@target@IV.20.4.11\endcsname{ega4.c04.prop.20-4-11}
\expandafter\def\csname sga@target@IV.20.4.12\endcsname{ega4.c04.cor.20-4-12}
\expandafter\def\csname sga@target@IV.20.4.2\endcsname{ega4.c04.statement.20-4-2}
\expandafter\def\csname sga@target@IV.20.4.3\endcsname{ega4.c04.statement.20-4-3}
\expandafter\def\csname sga@target@IV.20.4.4\endcsname{ega4.c04.prop.20-4-4}
\expandafter\def\csname sga@target@IV.20.4.5\endcsname{ega4.c04.prop.20-4-5}
\expandafter\def\csname sga@target@IV.20.4.6\endcsname{ega4.c04.prop.20-4-6}
\expandafter\def\csname sga@target@IV.20.4.7\endcsname{ega4.c04.cor.20-4-7}
\expandafter\def\csname sga@target@IV.20.4.8\endcsname{ega4.c04.rem.20-4-8}
\expandafter\def\csname sga@target@IV.20.4.9\endcsname{ega4.c04.cor.20-4-9}
\expandafter\def\csname sga@target@IV.20.5\endcsname{ega4.c04.subsec.20-5}
\expandafter\def\csname sga@target@IV.20.5.1\endcsname{ega4.c04.statement.20-5-1}
\expandafter\def\csname sga@target@IV.20.5.2\endcsname{ega4.c04.statement.20-5-2}
\expandafter\def\csname sga@target@IV.20.5.2.1\endcsname{ega4.c04.statement.20-5-2-1}
\expandafter\def\csname sga@target@IV.20.5.3\endcsname{ega4.c04.prop.20-5-3}
\expandafter\def\csname sga@target@IV.20.5.4\endcsname{ega4.c04.statement.20-5-4}
\expandafter\def\csname sga@target@IV.20.5.5\endcsname{ega4.c04.statement.20-5-5}
\expandafter\def\csname sga@target@IV.20.5.6\endcsname{ega4.c04.statement.20-5-6}
\expandafter\def\csname sga@target@IV.20.5.7\endcsname{ega4.c04.statement.20-5-7}
\expandafter\def\csname sga@target@IV.20.5.8\endcsname{ega4.c04.statement.20-5-8}
\expandafter\def\csname sga@target@IV.20.6\endcsname{ega4.c04.subsec.20-6}
\expandafter\def\csname sga@target@IV.20.6.1\endcsname{ega4.c04.statement.20-6-1}
\expandafter\def\csname sga@target@IV.20.6.1.1\endcsname{ega4.c04.eq.20-6-1-1}
\expandafter\def\csname sga@target@IV.20.6.2\endcsname{ega4.c04.statement.20-6-2}
\expandafter\def\csname sga@target@IV.20.6.2.1\endcsname{ega4.c04.statement.20-6-2-1}
\expandafter\def\csname sga@target@IV.20.6.2.2\endcsname{ega4.c04.eq.20-6-2-2}
\expandafter\def\csname sga@target@IV.20.6.3\endcsname{ega4.c04.prop.20-6-3}
\expandafter\def\csname sga@target@IV.20.6.4\endcsname{ega4.c04.prop.20-6-4}
\expandafter\def\csname sga@target@IV.20.6.5\endcsname{ega4.c04.statement.20-6-5}
\expandafter\def\csname sga@target@IV.20.6.5.1\endcsname{ega4.c04.eq.20-6-5-1}
\expandafter\def\csname sga@target@IV.20.6.6\endcsname{ega4.c04.prop.20-6-6}
\expandafter\def\csname sga@target@IV.20.6.7\endcsname{ega4.c04.statement.20-6-7}
\expandafter\def\csname sga@target@IV.20.6.8\endcsname{ega4.c04.prop.20-6-8}
\expandafter\def\csname sga@target@IV.20.6.8.1\endcsname{ega4.c04.eq.20-6-8-1}
\expandafter\def\csname sga@target@IV.20.6.9\endcsname{ega4.c04.statement.20-6-9}
\expandafter\def\csname sga@target@IV.20.6.9.1\endcsname{ega4.c04.statement.20-6-9-1}
\expandafter\def\csname sga@target@IV.21.1\endcsname{ega4.c04.subsec.21-1.key-iv-21-1}
\expandafter\def\csname sga@target@IV.21.1.1\endcsname{ega4.c04.statement.21-1-1}
\expandafter\def\csname sga@target@IV.21.1.10\endcsname{ega4.c04.prop.21-1-10}
\expandafter\def\csname sga@target@IV.21.1.2\endcsname{ega4.c04.statement.21-1-2}
\expandafter\def\csname sga@target@IV.21.1.3\endcsname{ega4.c04.statement.21-1-3}
\expandafter\def\csname sga@target@IV.21.1.3.1\endcsname{ega4.c04.statement.21-1-3-1}
\expandafter\def\csname sga@target@IV.21.1.3.2\endcsname{ega4.c04.statement.21-1-3-2}
\expandafter\def\csname sga@target@IV.21.1.3.3\endcsname{ega4.c04.statement.21-1-3-3}
\expandafter\def\csname sga@target@IV.21.1.3.4\endcsname{ega4.c04.statement.21-1-3-4}
\expandafter\def\csname sga@target@IV.21.1.4\endcsname{ega4.c04.statement.21-1-4}
\expandafter\def\csname sga@target@IV.21.1.4.1\endcsname{ega4.c04.eq.21-1-4-1}
\expandafter\def\csname sga@target@IV.21.1.4.2\endcsname{ega4.c04.eq.21-1-4-2}
\expandafter\def\csname sga@target@IV.21.1.4.3\endcsname{ega4.c04.eq.21-1-4-3}
\expandafter\def\csname sga@target@IV.21.1.5\endcsname{ega4.c04.statement.21-1-5}
\expandafter\def\csname sga@target@IV.21.1.5.1\endcsname{ega4.c04.statement.21-1-5-1}
\expandafter\def\csname sga@target@IV.21.1.5.2\endcsname{ega4.c04.eq.21-1-5-2}
\expandafter\def\csname sga@target@IV.21.1.6\endcsname{ega4.c04.def.21-1-6}
\expandafter\def\csname sga@target@IV.21.1.6.1\endcsname{ega4.c04.eq.21-1-6-1}
\expandafter\def\csname sga@target@IV.21.1.6.2\endcsname{ega4.c04.eq.21-1-6-2}
\expandafter\def\csname sga@target@IV.21.1.6.3\endcsname{ega4.c04.eq.21-1-6-3}
\expandafter\def\csname sga@target@IV.21.1.6.4\endcsname{ega4.c04.eq.21-1-6-4}
\expandafter\def\csname sga@target@IV.21.1.7\endcsname{ega4.c04.statement.21-1-7}
\expandafter\def\csname sga@target@IV.21.1.7.1\endcsname{ega4.c04.statement.21-1-7-1}
\expandafter\def\csname sga@target@IV.21.1.8\endcsname{ega4.c04.prop.21-1-8}
\expandafter\def\csname sga@target@IV.21.1.9\endcsname{ega4.c04.cor.21-1-9}
\expandafter\def\csname sga@target@IV.21.10.1\endcsname{ega4.c04.statement.21-10-1}
\expandafter\def\csname sga@target@IV.21.10.1.1\endcsname{ega4.c04.eq.21-10-1-1}
\expandafter\def\csname sga@target@IV.21.10.1.2\endcsname{ega4.c04.eq.21-10-1-2}
\expandafter\def\csname sga@target@IV.21.10.10\endcsname{ega4.c04.statement.21-10-10}
\expandafter\def\csname sga@target@IV.21.10.11\endcsname{ega4.c04.statement.21-10-11}
\expandafter\def\csname sga@target@IV.21.10.11.1\endcsname{ega4.c04.eq.21-10-11-1}
\expandafter\def\csname sga@target@IV.21.10.12\endcsname{ega4.c04.rem.21-10-12}
\expandafter\def\csname sga@target@IV.21.10.13\endcsname{ega4.c04.lem.21-10-13}
\expandafter\def\csname sga@target@IV.21.10.13.1\endcsname{ega4.c04.prop.21-10-13-1}
\expandafter\def\csname sga@target@IV.21.10.13bis\endcsname{ega4.c04.prop.21-10-13bis}
\expandafter\def\csname sga@target@IV.21.10.14\endcsname{ega4.c04.statement.21-10-14}
\expandafter\def\csname sga@target@IV.21.10.14.1\endcsname{ega4.c04.eq.21-10-14-1}
\expandafter\def\csname sga@target@IV.21.10.14.2\endcsname{ega4.c04.eq.21-10-14-2}
\expandafter\def\csname sga@target@IV.21.10.15\endcsname{ega4.c04.prop.21-10-15}
\expandafter\def\csname sga@target@IV.21.10.16\endcsname{ega4.c04.prop.21-10-16}
\expandafter\def\csname sga@target@IV.21.10.16.1\endcsname{ega4.c04.prop.21-10-16-1}
\expandafter\def\csname sga@target@IV.21.10.16.2\endcsname{ega4.c04.eq.21-10-16-2}
\expandafter\def\csname sga@target@IV.21.10.17\endcsname{ega4.c04.prop.21-10-17}
\expandafter\def\csname sga@target@IV.21.10.17.1\endcsname{ega4.c04.prop.21-10-17-1}
\expandafter\def\csname sga@target@IV.21.10.17.10\endcsname{ega4.c04.eq.21-10-17-10}
\expandafter\def\csname sga@target@IV.21.10.17.2\endcsname{ega4.c04.eq.21-10-17-2}
\expandafter\def\csname sga@target@IV.21.10.17.3\endcsname{ega4.c04.lem.21-10-17-3}
\expandafter\def\csname sga@target@IV.21.10.17.4\endcsname{ega4.c04.lem.21-10-17-4}
\expandafter\def\csname sga@target@IV.21.10.17.5\endcsname{ega4.c04.eq.21-10-17-5}
\expandafter\def\csname sga@target@IV.21.10.17.6\endcsname{ega4.c04.eq.21-10-17-6}
\expandafter\def\csname sga@target@IV.21.10.17.7\endcsname{ega4.c04.lem.21-10-17-7}
\expandafter\def\csname sga@target@IV.21.10.17.8\endcsname{ega4.c04.lem.21-10-17-8}
\expandafter\def\csname sga@target@IV.21.10.17.9\endcsname{ega4.c04.eq.21-10-17-9}
\expandafter\def\csname sga@target@IV.21.10.18\endcsname{ega4.c04.prop.21-10-18}
\expandafter\def\csname sga@target@IV.21.10.18.1\endcsname{ega4.c04.prop.21-10-18-1}
\expandafter\def\csname sga@target@IV.21.10.19\endcsname{ega4.c04.cor.21-10-19}
\expandafter\def\csname sga@target@IV.21.10.19.1\endcsname{ega4.c04.cor.21-10-19-1}
\expandafter\def\csname sga@target@IV.21.10.2\endcsname{ega4.c04.statement.21-10-2}
\expandafter\def\csname sga@target@IV.21.10.3\endcsname{ega4.c04.statement.21-10-3}
\expandafter\def\csname sga@target@IV.21.10.3.1\endcsname{ega4.c04.statement.21-10-3-1}
\expandafter\def\csname sga@target@IV.21.10.4\endcsname{ega4.c04.prop.21-10-4}
\expandafter\def\csname sga@target@IV.21.10.4.1\endcsname{ega4.c04.prop.21-10-4-1}
\expandafter\def\csname sga@target@IV.21.10.5\endcsname{ega4.c04.statement.21-10-5}
\expandafter\def\csname sga@target@IV.21.10.6\endcsname{ega4.c04.prop.21-10-6}
\expandafter\def\csname sga@target@IV.21.10.7\endcsname{ega4.c04.rem.21-10-7}
\expandafter\def\csname sga@target@IV.21.10.8\endcsname{ega4.c04.prop.21-10-8}
\expandafter\def\csname sga@target@IV.21.10.9\endcsname{ega4.c04.statement.21-10-9}
\expandafter\def\csname sga@target@IV.21.11\endcsname{ega4.c04.subsec.21-11.key-iv-21-11}
\expandafter\def\csname sga@target@IV.21.11.1\endcsname{ega4.c04.thm.21-11-1}
\expandafter\def\csname sga@target@IV.21.11.1.1\endcsname{ega4.c04.eq.21-11-1-1}
\expandafter\def\csname sga@target@IV.21.11.2\endcsname{ega4.c04.lem.21-11-2}
\expandafter\def\csname sga@target@IV.21.12\endcsname{ega4.c04.subsec.21-12.key-iv-21-12}
\expandafter\def\csname sga@target@IV.21.12.1\endcsname{ega4.c04.statement.21-12-1}
\expandafter\def\csname sga@target@IV.21.12.1.1\endcsname{ega4.c04.eq.21-12-1-1}
\expandafter\def\csname sga@target@IV.21.12.10\endcsname{ega4.c04.prop.21-12-10}
\expandafter\def\csname sga@target@IV.21.12.12\endcsname{ega4.c04.thm.21-12-12}
\expandafter\def\csname sga@target@IV.21.12.13\endcsname{ega4.c04.cor.21-12-13}
\expandafter\def\csname sga@target@IV.21.12.14\endcsname{ega4.c04.item.21-12-14}
\expandafter\def\csname sga@target@IV.21.12.15\endcsname{ega4.c04.prop.21-12-15}
\expandafter\def\csname sga@target@IV.21.12.16\endcsname{ega4.c04.cor.21-12-16}
\expandafter\def\csname sga@target@IV.21.12.17\endcsname{ega4.c04.item.21-12-17}
\expandafter\def\csname sga@target@IV.21.12.2\endcsname{ega4.c04.prop.21-12-2}
\expandafter\def\csname sga@target@IV.21.12.2.1\endcsname{ega4.c04.prop.21-12-2-1}
\expandafter\def\csname sga@target@IV.21.12.3\endcsname{ega4.c04.cor.21-12-3}
\expandafter\def\csname sga@target@IV.21.12.4\endcsname{ega4.c04.eq.21-12-4}
\expandafter\def\csname sga@target@IV.21.12.5\endcsname{ega4.c04.statement.21-12-5}
\expandafter\def\csname sga@target@IV.21.12.5.1\endcsname{ega4.c04.eq.21-12-5-1}
\expandafter\def\csname sga@target@IV.21.12.5.2\endcsname{ega4.c04.eq.21-12-5-2}
\expandafter\def\csname sga@target@IV.21.12.5.3\endcsname{ega4.c04.eq.21-12-5-3}
\expandafter\def\csname sga@target@IV.21.12.6\endcsname{ega4.c04.thm.21-12-6}
\expandafter\def\csname sga@target@IV.21.12.7\endcsname{ega4.c04.cor.21-12-7}
\expandafter\def\csname sga@target@IV.21.12.8\endcsname{ega4.c04.statement.21-12-8}
\expandafter\def\csname sga@target@IV.21.12.9\endcsname{ega4.c04.statement.21-12-9}
\expandafter\def\csname sga@target@IV.21.13\endcsname{ega4.c04.subsec.21-13.key-iv-21-13}
\expandafter\def\csname sga@target@IV.21.13.1\endcsname{ega4.c04.def.21-13-1}
\expandafter\def\csname sga@target@IV.21.13.10\endcsname{ega4.c04.prop.21-13-10}
\expandafter\def\csname sga@target@IV.21.13.11\endcsname{ega4.c04.cor.21-13-11}
\expandafter\def\csname sga@target@IV.21.13.12\endcsname{ega4.c04.prop.21-13-12}
\expandafter\def\csname sga@target@IV.21.13.12.1\endcsname{ega4.c04.statement.21-13-12-1}
\expandafter\def\csname sga@target@IV.21.13.13\endcsname{ega4.c04.rem.21-13-13}
\expandafter\def\csname sga@target@IV.21.13.14\endcsname{ega4.c04.prop.21-13-14}
\expandafter\def\csname sga@target@IV.21.13.15\endcsname{ega4.c04.cor.21-13-15}
\expandafter\def\csname sga@target@IV.21.13.16\endcsname{ega4.c04.rem.21-13-16}
\expandafter\def\csname sga@target@IV.21.13.2\endcsname{ega4.c04.lem.21-13-2}
\expandafter\def\csname sga@target@IV.21.13.2.1\endcsname{ega4.c04.eq.21-13-2-1}
\expandafter\def\csname sga@target@IV.21.13.3\endcsname{ega4.c04.cor.21-13-3}
\expandafter\def\csname sga@target@IV.21.13.4\endcsname{ega4.c04.lem.21-13-4}
\expandafter\def\csname sga@target@IV.21.13.5\endcsname{ega4.c04.prop.21-13-5}
\expandafter\def\csname sga@target@IV.21.13.6\endcsname{ega4.c04.cor.21-13-6}
\expandafter\def\csname sga@target@IV.21.13.7\endcsname{ega4.c04.def.21-13-7}
\expandafter\def\csname sga@target@IV.21.13.8\endcsname{ega4.c04.prop.21-13-8}
\expandafter\def\csname sga@target@IV.21.13.9\endcsname{ega4.c04.statement.21-13-9}
\expandafter\def\csname sga@target@IV.21.13.9.1\endcsname{ega4.c04.eq.21-13-9-1}
\expandafter\def\csname sga@target@IV.21.14\endcsname{ega4.c04.subsec.21-14.key-iv-21-14}
\expandafter\def\csname sga@target@IV.21.14.1\endcsname{ega4.c04.thm.21-14-1}
\expandafter\def\csname sga@target@IV.21.14.2\endcsname{ega4.c04.cor.21-14-2}
\expandafter\def\csname sga@target@IV.21.14.3\endcsname{ega4.c04.prop.21-14-3}
\expandafter\def\csname sga@target@IV.21.14.4\endcsname{ega4.c04.statement.21-14-4}
\expandafter\def\csname sga@target@IV.21.14.4bis\endcsname{ega4.c04.statement.21-14-4bis}
\expandafter\def\csname sga@target@IV.21.15\endcsname{ega4.c04.subsec.21-15.key-iv-21-15}
\expandafter\def\csname sga@target@IV.21.15.1\endcsname{ega4.c04.statement.21-15-1}
\expandafter\def\csname sga@target@IV.21.15.2\endcsname{ega4.c04.def.21-15-2}
\expandafter\def\csname sga@target@IV.21.15.3\endcsname{ega4.c04.statement.21-15-3}
\expandafter\def\csname sga@target@IV.21.15.3.1\endcsname{ega4.c04.eq.21-15-3-1}
\expandafter\def\csname sga@target@IV.21.15.3.2\endcsname{ega4.c04.eq.21-15-3-2}
\expandafter\def\csname sga@target@IV.21.15.3bis\endcsname{ega4.c04.statement.21-15-3bis}
\expandafter\def\csname sga@target@IV.21.15.4\endcsname{ega4.c04.prop.21-15-4}
\expandafter\def\csname sga@target@IV.21.15.5\endcsname{ega4.c04.prop.21-15-5}
\expandafter\def\csname sga@target@IV.21.15.6\endcsname{ega4.c04.prop.21-15-6}
\expandafter\def\csname sga@target@IV.21.15.7\endcsname{ega4.c04.statement.21-15-7}
\expandafter\def\csname sga@target@IV.21.15.8\endcsname{ega4.c04.statement.21-15-8}
\expandafter\def\csname sga@target@IV.21.15.9\endcsname{ega4.c04.statement.21-15-9}
\expandafter\def\csname sga@target@IV.21.15.9.1\endcsname{ega4.c04.eq.21-15-9-1}
\expandafter\def\csname sga@target@IV.21.2.1\endcsname{ega4.c04.statement.21-2-1}
\expandafter\def\csname sga@target@IV.21.2.10\endcsname{ega4.c04.statement.21-2-10}
\expandafter\def\csname sga@target@IV.21.2.10.1\endcsname{ega4.c04.eq.21-2-10-1}
\expandafter\def\csname sga@target@IV.21.2.11\endcsname{ega4.c04.prop.21-2-11}
\expandafter\def\csname sga@target@IV.21.2.11.1\endcsname{ega4.c04.prop.21-2-11-1}
\expandafter\def\csname sga@target@IV.21.2.12\endcsname{ega4.c04.statement.21-2-12}
\expandafter\def\csname sga@target@IV.21.2.2\endcsname{ega4.c04.prop.21-2-2}
\expandafter\def\csname sga@target@IV.21.2.3\endcsname{ega4.c04.cor.21-2-3}
\expandafter\def\csname sga@target@IV.21.2.4\endcsname{ega4.c04.statement.21-2-4}
\expandafter\def\csname sga@target@IV.21.2.5\endcsname{ega4.c04.statement.21-2-5}
\expandafter\def\csname sga@target@IV.21.2.5.1\endcsname{ega4.c04.statement.21-2-5-1}
\expandafter\def\csname sga@target@IV.21.2.5.2\endcsname{ega4.c04.eq.21-2-5-2}
\expandafter\def\csname sga@target@IV.21.2.5.3\endcsname{ega4.c04.eq.21-2-5-3}
\expandafter\def\csname sga@target@IV.21.2.5.4\endcsname{ega4.c04.eq.21-2-5-4}
\expandafter\def\csname sga@target@IV.21.2.6\endcsname{ega4.c04.prop.21-2-6}
\expandafter\def\csname sga@target@IV.21.2.6.1\endcsname{ega4.c04.eq.21-2-6-1}
\expandafter\def\csname sga@target@IV.21.2.7\endcsname{ega4.c04.statement.21-2-7}
\expandafter\def\csname sga@target@IV.21.2.7.1\endcsname{ega4.c04.statement.21-2-7-1}
\expandafter\def\csname sga@target@IV.21.2.7.2\endcsname{ega4.c04.eq.21-2-7-2}
\expandafter\def\csname sga@target@IV.21.2.8\endcsname{ega4.c04.statement.21-2-8}
\expandafter\def\csname sga@target@IV.21.2.8.1\endcsname{ega4.c04.statement.21-2-8-1}
\expandafter\def\csname sga@target@IV.21.2.8.2\endcsname{ega4.c04.statement.21-2-8-2}
\expandafter\def\csname sga@target@IV.21.2.8.3\endcsname{ega4.c04.statement.21-2-8-3}
\expandafter\def\csname sga@target@IV.21.2.9\endcsname{ega4.c04.statement.21-2-9}
\expandafter\def\csname sga@target@IV.21.2.9.1\endcsname{ega4.c04.statement.21-2-9-1}
\expandafter\def\csname sga@target@IV.21.2.9.2\endcsname{ega4.c04.eq.21-2-9-2}
\expandafter\def\csname sga@target@IV.21.2.9.3\endcsname{ega4.c04.eq.21-2-9-3}
\expandafter\def\csname sga@target@IV.21.2.9.4\endcsname{ega4.c04.eq.21-2-9-4}
\expandafter\def\csname sga@target@IV.21.3.1\endcsname{ega4.c04.statement.21-3-1}
\expandafter\def\csname sga@target@IV.21.3.2\endcsname{ega4.c04.statement.21-3-2}
\expandafter\def\csname sga@target@IV.21.3.2.1\endcsname{ega4.c04.statement.21-3-2-1}
\expandafter\def\csname sga@target@IV.21.3.2.2\endcsname{ega4.c04.eq.21-3-2-2}
\expandafter\def\csname sga@target@IV.21.3.2.3\endcsname{ega4.c04.eq.21-3-2-3}
\expandafter\def\csname sga@target@IV.21.3.2.4\endcsname{ega4.c04.eq.21-3-2-4}
\expandafter\def\csname sga@target@IV.21.3.3\endcsname{ega4.c04.prop.21-3-3}
\expandafter\def\csname sga@target@IV.21.3.3.1\endcsname{ega4.c04.eq.21-3-3-1}
\expandafter\def\csname sga@target@IV.21.3.4\endcsname{ega4.c04.prop.21-3-4}
\expandafter\def\csname sga@target@IV.21.3.5\endcsname{ega4.c04.cor.21-3-5}
\expandafter\def\csname sga@target@IV.21.3.6\endcsname{ega4.c04.rem.21-3-6}
\expandafter\def\csname sga@target@IV.21.3.6.1\endcsname{ega4.c04.eq.21-3-6-1}
\expandafter\def\csname sga@target@IV.21.4\endcsname{ega4.c04.subsec.21-4.key-iv-21-4}
\expandafter\def\csname sga@target@IV.21.4.1\endcsname{ega4.c04.statement.21-4-1}
\expandafter\def\csname sga@target@IV.21.4.10\endcsname{ega4.c04.cor.21-4-10}
\expandafter\def\csname sga@target@IV.21.4.11\endcsname{ega4.c04.cor.21-4-11}
\expandafter\def\csname sga@target@IV.21.4.12\endcsname{ega4.c04.rem.21-4-12}
\expandafter\def\csname sga@target@IV.21.4.13\endcsname{ega4.c04.cor.21-4-13}
\expandafter\def\csname sga@target@IV.21.4.2\endcsname{ega4.c04.def.21-4-2}
\expandafter\def\csname sga@target@IV.21.4.2.1\endcsname{ega4.c04.eq.21-4-2-1}
\expandafter\def\csname sga@target@IV.21.4.3\endcsname{ega4.c04.statement.21-4-3}
\expandafter\def\csname sga@target@IV.21.4.4\endcsname{ega4.c04.statement.21-4-4}
\expandafter\def\csname sga@target@IV.21.4.5\endcsname{ega4.c04.prop.21-4-5}
\expandafter\def\csname sga@target@IV.21.4.6\endcsname{ega4.c04.cor.21-4-6}
\expandafter\def\csname sga@target@IV.21.4.6.1\endcsname{ega4.c04.eq.21-4-6-1}
\expandafter\def\csname sga@target@IV.21.4.7\endcsname{ega4.c04.statement.21-4-7}
\expandafter\def\csname sga@target@IV.21.4.8\endcsname{ega4.c04.prop.21-4-8}
\expandafter\def\csname sga@target@IV.21.4.9\endcsname{ega4.c04.prop.21-4-9}
\expandafter\def\csname sga@target@IV.21.5\endcsname{ega4.c04.subsec.21-5.key-iv-21-5}
\expandafter\def\csname sga@target@IV.21.5.1\endcsname{ega4.c04.statement.21-5-1}
\expandafter\def\csname sga@target@IV.21.5.2\endcsname{ega4.c04.lem.21-5-2}
\expandafter\def\csname sga@target@IV.21.5.3\endcsname{ega4.c04.statement.21-5-3}
\expandafter\def\csname sga@target@IV.21.5.3.1\endcsname{ega4.c04.eq.21-5-3-1}
\expandafter\def\csname sga@target@IV.21.5.3.2\endcsname{ega4.c04.eq.21-5-3-2}
\expandafter\def\csname sga@target@IV.21.5.4\endcsname{ega4.c04.statement.21-5-4}
\expandafter\def\csname sga@target@IV.21.5.5\endcsname{ega4.c04.def.21-5-5}
\expandafter\def\csname sga@target@IV.21.5.5.1\endcsname{ega4.c04.eq.21-5-5-1}
\expandafter\def\csname sga@target@IV.21.5.5.2\endcsname{ega4.c04.eq.21-5-5-2}
\expandafter\def\csname sga@target@IV.21.5.5.3\endcsname{ega4.c04.eq.21-5-5-3}
\expandafter\def\csname sga@target@IV.21.5.6\endcsname{ega4.c04.prop.21-5-6}
\expandafter\def\csname sga@target@IV.21.5.6.1\endcsname{ega4.c04.prop.21-5-6-1}
\expandafter\def\csname sga@target@IV.21.5.7\endcsname{ega4.c04.prop.21-5-7}
\expandafter\def\csname sga@target@IV.21.5.7.1\endcsname{ega4.c04.prop.21-5-7-1}
\expandafter\def\csname sga@target@IV.21.5.7.2\endcsname{ega4.c04.lem.21-5-7-2}
\expandafter\def\csname sga@target@IV.21.5.7.3\endcsname{ega4.c04.lem.21-5-7-3}
\expandafter\def\csname sga@target@IV.21.5.7.4\endcsname{ega4.c04.eq.21-5-7-4}
\expandafter\def\csname sga@target@IV.21.5.8\endcsname{ega4.c04.prop.21-5-8}
\expandafter\def\csname sga@target@IV.21.5.8.1\endcsname{ega4.c04.eq.21-5-8-1}
\expandafter\def\csname sga@target@IV.21.6\endcsname{ega4.c04.subsec.21-6.key-iv-21-6}
\expandafter\def\csname sga@target@IV.21.6.1\endcsname{ega4.c04.statement.21-6-1}
\expandafter\def\csname sga@target@IV.21.6.10\endcsname{ega4.c04.cor.21-6-10}
\expandafter\def\csname sga@target@IV.21.6.10.1\endcsname{ega4.c04.eq.21-6-10-1}
\expandafter\def\csname sga@target@IV.21.6.11\endcsname{ega4.c04.cor.21-6-11}
\expandafter\def\csname sga@target@IV.21.6.12\endcsname{ega4.c04.prop.21-6-12}
\expandafter\def\csname sga@target@IV.21.6.12.1\endcsname{ega4.c04.eq.21-6-12-1}
\expandafter\def\csname sga@target@IV.21.6.12.2\endcsname{ega4.c04.eq.21-6-12-2}
\expandafter\def\csname sga@target@IV.21.6.13\endcsname{ega4.c04.cor.21-6-13}
\expandafter\def\csname sga@target@IV.21.6.14\endcsname{ega4.c04.cor.21-6-14}
\expandafter\def\csname sga@target@IV.21.6.2\endcsname{ega4.c04.statement.21-6-2}
\expandafter\def\csname sga@target@IV.21.6.3\endcsname{ega4.c04.statement.21-6-3}
\expandafter\def\csname sga@target@IV.21.6.4\endcsname{ega4.c04.statement.21-6-4}
\expandafter\def\csname sga@target@IV.21.6.4.1\endcsname{ega4.c04.statement.21-6-4-1}
\expandafter\def\csname sga@target@IV.21.6.4.2\endcsname{ega4.c04.eq.21-6-4-2}
\expandafter\def\csname sga@target@IV.21.6.5\endcsname{ega4.c04.statement.21-6-5}
\expandafter\def\csname sga@target@IV.21.6.5.1\endcsname{ega4.c04.eq.21-6-5-1}
\expandafter\def\csname sga@target@IV.21.6.6\endcsname{ega4.c04.prop.21-6-6}
\expandafter\def\csname sga@target@IV.21.6.6.2\endcsname{ega4.c04.statement.21-6-6-2}
\expandafter\def\csname sga@target@IV.21.6.7\endcsname{ega4.c04.statement.21-6-7}
\expandafter\def\csname sga@target@IV.21.6.7.1\endcsname{ega4.c04.eq.21-6-7-1}
\expandafter\def\csname sga@target@IV.21.6.7.2\endcsname{ega4.c04.eq.21-6-7-2}
\expandafter\def\csname sga@target@IV.21.6.7.3\endcsname{ega4.c04.eq.21-6-7-3}
\expandafter\def\csname sga@target@IV.21.6.7.4\endcsname{ega4.c04.eq.21-6-7-4}
\expandafter\def\csname sga@target@IV.21.6.7.5\endcsname{ega4.c04.eq.21-6-7-5}
\expandafter\def\csname sga@target@IV.21.6.8\endcsname{ega4.c04.statement.21-6-8}
\expandafter\def\csname sga@target@IV.21.6.9\endcsname{ega4.c04.thm.21-6-9}
\expandafter\def\csname sga@target@IV.21.6.9.1\endcsname{ega4.c04.eq.21-6-9-1}
\expandafter\def\csname sga@target@IV.21.6.9.2\endcsname{ega4.c04.statement.21-6-9-2}
\expandafter\def\csname sga@target@IV.21.7\endcsname{ega4.c04.subsec.21-7.key-iv-21-7}
\expandafter\def\csname sga@target@IV.21.7.1\endcsname{ega4.c04.statement.21-7-1}
\expandafter\def\csname sga@target@IV.21.7.2\endcsname{ega4.c04.prop.21-7-2}
\expandafter\def\csname sga@target@IV.21.7.2.1\endcsname{ega4.c04.eq.21-7-2-1}
\expandafter\def\csname sga@target@IV.21.7.3\endcsname{ega4.c04.cor.21-7-3}
\expandafter\def\csname sga@target@IV.21.7.3.1\endcsname{ega4.c04.statement.21-7-3-1}
\expandafter\def\csname sga@target@IV.21.7.4\endcsname{ega4.c04.prop.21-7-4}
\expandafter\def\csname sga@target@IV.21.7.5\endcsname{ega4.c04.rem.21-7-5}
\expandafter\def\csname sga@target@IV.21.7.6\endcsname{ega4.c04.cor.21-7-6}
\expandafter\def\csname sga@target@IV.21.7.7\endcsname{ega4.c04.rem.21-7-7}
\expandafter\def\csname sga@target@IV.21.8.1\endcsname{ega4.c04.lem.21-8-1}
\expandafter\def\csname sga@target@IV.21.8.2\endcsname{ega4.c04.cor.21-8-2}
\expandafter\def\csname sga@target@IV.21.8.3\endcsname{ega4.c04.prop.21-8-3}
\expandafter\def\csname sga@target@IV.21.8.3.1\endcsname{ega4.c04.eq.21-8-3-1}
\expandafter\def\csname sga@target@IV.21.8.4\endcsname{ega4.c04.cor.21-8-4}
\expandafter\def\csname sga@target@IV.21.8.5\endcsname{ega4.c04.prop.21-8-5}
\expandafter\def\csname sga@target@IV.21.8.5.1\endcsname{ega4.c04.eq.21-8-5-1}
\expandafter\def\csname sga@target@IV.21.8.5.2\endcsname{ega4.c04.eq.21-8-5-2}
\expandafter\def\csname sga@target@IV.21.8.5.3\endcsname{ega4.c04.eq.21-8-5-3}
\expandafter\def\csname sga@target@IV.21.8.6\endcsname{ega4.c04.rem.21-8-6}
\expandafter\def\csname sga@target@IV.21.9\endcsname{ega4.c04.subsec.21-9.key-iv-21-9}
\expandafter\def\csname sga@target@IV.21.9.1\endcsname{ega4.c04.statement.21-9-1}
\expandafter\def\csname sga@target@IV.21.9.1.1\endcsname{ega4.c04.eq.21-9-1-1}
\expandafter\def\csname sga@target@IV.21.9.1.2\endcsname{ega4.c04.eq.21-9-1-2}
\expandafter\def\csname sga@target@IV.21.9.10\endcsname{ega4.c04.statement.21-9-10}
\expandafter\def\csname sga@target@IV.21.9.11\endcsname{ega4.c04.prop.21-9-11}
\expandafter\def\csname sga@target@IV.21.9.12\endcsname{ega4.c04.cor.21-9-12}
\expandafter\def\csname sga@target@IV.21.9.13\endcsname{ega4.c04.rem.21-9-13}
\expandafter\def\csname sga@target@IV.21.9.2\endcsname{ega4.c04.lem.21-9-2}
\expandafter\def\csname sga@target@IV.21.9.3\endcsname{ega4.c04.rem.21-9-3}
\expandafter\def\csname sga@target@IV.21.9.4\endcsname{ega4.c04.prop.21-9-4}
\expandafter\def\csname sga@target@IV.21.9.4.1\endcsname{ega4.c04.prop.21-9-4-1}
\expandafter\def\csname sga@target@IV.21.9.5\endcsname{ega4.c04.cor.21-9-5}
\expandafter\def\csname sga@target@IV.21.9.6\endcsname{ega4.c04.cor.21-9-6}
\expandafter\def\csname sga@target@IV.21.9.7\endcsname{ega4.c04.cor.21-9-7}
\expandafter\def\csname sga@target@IV.21.9.8\endcsname{ega4.c04.statement.21-9-8}
\expandafter\def\csname sga@target@IV.21.9.9\endcsname{ega4.c04.prop.21-9-9}
\expandafter\def\csname sga@target@IV.3.1.1\endcsname{ega4.c04.def.3-1-1}
\expandafter\def\csname sga@target@IV.3.1.1.1\endcsname{ega4.c04.eq.3-1-1-1}
\expandafter\def\csname sga@target@IV.3.1.10\endcsname{ega4.c04.prop.3-1-10}
\expandafter\def\csname sga@target@IV.3.1.10.1\endcsname{ega4.c04.lem.3-1-10-1}
\expandafter\def\csname sga@target@IV.3.1.10.2\endcsname{ega4.c04.rem.3-1-10-2}
\expandafter\def\csname sga@target@IV.3.1.10.3\endcsname{ega4.c04.eq.3-1-10-3}
\expandafter\def\csname sga@target@IV.3.1.10.4\endcsname{ega4.c04.eq.3-1-10-4}
\expandafter\def\csname sga@target@IV.3.1.10.5\endcsname{ega4.c04.rem.3-1-10-5}
\expandafter\def\csname sga@target@IV.3.1.11\endcsname{ega4.c04.cor.3-1-11}
\expandafter\def\csname sga@target@IV.3.1.12\endcsname{ega4.c04.rem.3-1-12}
\expandafter\def\csname sga@target@IV.3.1.13\endcsname{ega4.c04.prop.3-1-13}
\expandafter\def\csname sga@target@IV.3.1.2\endcsname{ega4.c04.prop.3-1-2}
\expandafter\def\csname sga@target@IV.3.1.3\endcsname{ega4.c04.prop.3-1-3}
\expandafter\def\csname sga@target@IV.3.1.4\endcsname{ega4.c04.cor.3-1-4}
\expandafter\def\csname sga@target@IV.3.1.5\endcsname{ega4.c04.cor.3-1-5}
\expandafter\def\csname sga@target@IV.3.1.6\endcsname{ega4.c04.prop.3-1-6}
\expandafter\def\csname sga@target@IV.3.1.7\endcsname{ega4.c04.prop.3-1-7}
\expandafter\def\csname sga@target@IV.3.1.8\endcsname{ega4.c04.prop.3-1-8}
\expandafter\def\csname sga@target@IV.3.1.9\endcsname{ega4.c04.cor.3-1-9}
\expandafter\def\csname sga@target@IV.3.2.1\endcsname{ega4.c04.prop.3-2-1}
\expandafter\def\csname sga@target@IV.3.2.2\endcsname{ega4.c04.def.3-2-2}
\expandafter\def\csname sga@target@IV.3.2.3\endcsname{ega4.c04.prop.3-2-3}
\expandafter\def\csname sga@target@IV.3.2.4\endcsname{ega4.c04.def.3-2-4}
\expandafter\def\csname sga@target@IV.3.2.5\endcsname{ega4.c04.def.3-2-5}
\expandafter\def\csname sga@target@IV.3.2.6\endcsname{ega4.c04.prop.3-2-6}
\expandafter\def\csname sga@target@IV.3.2.7\endcsname{ega4.c04.cor.3-2-7}
\expandafter\def\csname sga@target@IV.3.2.8\endcsname{ega4.c04.cor.3-2-8}
\expandafter\def\csname sga@target@IV.3.2.9\endcsname{ega4.c04.rem.3-2-9}
\expandafter\def\csname sga@target@IV.3.3.1\endcsname{ega4.c04.prop.3-3-1}
\expandafter\def\csname sga@target@IV.3.3.1.1\endcsname{ega4.c04.prop.3-3-1-1}
\expandafter\def\csname sga@target@IV.3.3.2\endcsname{ega4.c04.cor.3-3-2}
\expandafter\def\csname sga@target@IV.3.3.3\endcsname{ega4.c04.cor.3-3-3}
\expandafter\def\csname sga@target@IV.3.3.4\endcsname{ega4.c04.prop.3-3-4}
\expandafter\def\csname sga@target@IV.3.3.5\endcsname{ega4.c04.cor.3-3-5}
\expandafter\def\csname sga@target@IV.3.3.6\endcsname{ega4.c04.prop.3-3-6}
\expandafter\def\csname sga@target@IV.3.3.6.1\endcsname{ega4.c04.eq.3-3-6-1}
\expandafter\def\csname sga@target@IV.3.3.6.2\endcsname{ega4.c04.eq.3-3-6-2}
\expandafter\def\csname sga@target@IV.3.3.6.3\endcsname{ega4.c04.eq.3-3-6-3}
\expandafter\def\csname sga@target@IV.3.3.6.4\endcsname{ega4.c04.eq.3-3-6-4}
\expandafter\def\csname sga@target@IV.3.3.7\endcsname{ega4.c04.cor.3-3-7}
\expandafter\def\csname sga@target@IV.3.3.7.1\endcsname{ega4.c04.cor.3-3-7-1}
\expandafter\def\csname sga@target@IV.3.3.8\endcsname{ega4.c04.sec.3-3-8}
\expandafter\def\csname sga@target@IV.3.3.9\endcsname{ega4.c04.cor.3-3-9}
\expandafter\def\csname sga@target@IV.3.3.9.1\endcsname{ega4.c04.cor.3-3-9-1}
\expandafter\def\csname sga@target@IV.3.4\endcsname{ega4.c04.subsec.3-4.key-iv-3-4}
\expandafter\def\csname sga@target@IV.3.4.1\endcsname{ega4.c04.prop.3-4-1}
\expandafter\def\csname sga@target@IV.3.4.1.1\endcsname{ega4.c04.eq.3-4-1-1}
\expandafter\def\csname sga@target@IV.3.4.1.2\endcsname{ega4.c04.lem.3-4-1-2}
\expandafter\def\csname sga@target@IV.3.4.1.3\endcsname{ega4.c04.lem.3-4-1-3}
\expandafter\def\csname sga@target@IV.3.4.1.4\endcsname{ega4.c04.lem.3-4-1-4}
\expandafter\def\csname sga@target@IV.3.4.1.5\endcsname{ega4.c04.eq.3-4-1-5}
\expandafter\def\csname sga@target@IV.3.4.2\endcsname{ega4.c04.cor.3-4-2}
\expandafter\def\csname sga@target@IV.3.4.3\endcsname{ega4.c04.prop.3-4-3}
\expandafter\def\csname sga@target@IV.3.4.4\endcsname{ega4.c04.prop.3-4-4}
\expandafter\def\csname sga@target@IV.3.4.5\endcsname{ega4.c04.prop.3-4-5}
\expandafter\def\csname sga@target@IV.3.4.6\endcsname{ega4.c04.prop.3-4-6}
\expandafter\def\csname sga@target@IV.3.4.6.1\endcsname{ega4.c04.lem.3-4-6-1}
\expandafter\def\csname sga@target@IV.3.4.7\endcsname{ega4.c04.cor.3-4-7}
\expandafter\def\csname sga@target@IV.3.4.8\endcsname{ega4.c04.rem.3-4-8}
\expandafter\def\csname sga@target@IV.3.4.9\endcsname{ega4.c04.prop.3-4-9}
\expandafter\def\csname sga@target@IV.4.1.1\endcsname{ega4.c04.def.4-1-1}
\expandafter\def\csname sga@target@IV.4.1.1.1\endcsname{ega4.c04.eq.4-1-1-1}
\expandafter\def\csname sga@target@IV.4.1.1.2\endcsname{ega4.c04.eq.4-1-1-2}
\expandafter\def\csname sga@target@IV.4.1.1.3\endcsname{ega4.c04.eq.4-1-1-3}
\expandafter\def\csname sga@target@IV.4.1.2\endcsname{ega4.c04.thm.4-1-2}
\expandafter\def\csname sga@target@IV.4.1.2.1\endcsname{ega4.c04.cor.4-1-2-1}
\expandafter\def\csname sga@target@IV.4.1.3\endcsname{ega4.c04.rem.4-1-3}
\expandafter\def\csname sga@target@IV.4.1.4\endcsname{ega4.c04.cor.4-1-4}
\expandafter\def\csname sga@target@IV.4.1.5\endcsname{ega4.c04.cor.4-1-5}
\expandafter\def\csname sga@target@IV.4.2.1\endcsname{ega4.c04.prop.4-2-1}
\expandafter\def\csname sga@target@IV.4.2.1.1\endcsname{ega4.c04.eq.4-2-1-1}
\expandafter\def\csname sga@target@IV.4.2.1.2\endcsname{ega4.c04.eq.4-2-1-2}
\expandafter\def\csname sga@target@IV.4.2.1.3\endcsname{ega4.c04.eq.4-2-1-3}
\expandafter\def\csname sga@target@IV.4.2.1.4\endcsname{ega4.c04.rem.4-2-1-4}
\expandafter\def\csname sga@target@IV.4.2.1.5\endcsname{ega4.c04.eq.4-2-1-5}
\expandafter\def\csname sga@target@IV.4.2.10\endcsname{ega4.c04.cor.4-2-10}
\expandafter\def\csname sga@target@IV.4.2.2\endcsname{ega4.c04.cor.4-2-2}
\expandafter\def\csname sga@target@IV.4.2.3\endcsname{ega4.c04.cor.4-2-3}
\expandafter\def\csname sga@target@IV.4.2.4\endcsname{ega4.c04.prop.4-2-4}
\expandafter\def\csname sga@target@IV.4.2.4.i\endcsname{ega4.c04.sec.4-2-4-i}
\expandafter\def\csname sga@target@IV.4.2.4.ii\endcsname{ega4.c04.sec.4-2-4-ii}
\expandafter\def\csname sga@target@IV.4.2.4.iii\endcsname{ega4.c04.sec.4-2-4-iii}
\expandafter\def\csname sga@target@IV.4.2.5\endcsname{ega4.c04.prop.4-2-5}
\expandafter\def\csname sga@target@IV.4.2.6\endcsname{ega4.c04.cor.4-2-6}
\expandafter\def\csname sga@target@IV.4.2.7\endcsname{ega4.c04.prop.4-2-7}
\expandafter\def\csname sga@target@IV.4.2.7.i\endcsname{ega4.c04.sec.4-2-7-i}
\expandafter\def\csname sga@target@IV.4.2.7.ii\endcsname{ega4.c04.sec.4-2-7-ii}
\expandafter\def\csname sga@target@IV.4.2.7.iii\endcsname{ega4.c04.sec.4-2-7-iii}
\expandafter\def\csname sga@target@IV.4.2.7.iv\endcsname{ega4.c04.sec.4-2-7-iv}
\expandafter\def\csname sga@target@IV.4.2.8\endcsname{ega4.c04.cor.4-2-8}
\expandafter\def\csname sga@target@IV.4.2.9\endcsname{ega4.c04.prop.4-2-9}
\expandafter\def\csname sga@target@IV.4.3.1\endcsname{ega4.c04.statement.4-3-1}
\expandafter\def\csname sga@target@IV.4.3.2\endcsname{ega4.c04.prop.4-3-2}
\expandafter\def\csname sga@target@IV.4.3.3\endcsname{ega4.c04.cor.4-3-3}
\expandafter\def\csname sga@target@IV.4.3.4\endcsname{ega4.c04.cor.4-3-4}
\expandafter\def\csname sga@target@IV.4.3.5\endcsname{ega4.c04.prop.4-3-5}
\expandafter\def\csname sga@target@IV.4.3.6\endcsname{ega4.c04.cor.4-3-6}
\expandafter\def\csname sga@target@IV.4.3.7\endcsname{ega4.c04.cor.4-3-7}
\expandafter\def\csname sga@target@IV.4.3.8\endcsname{ega4.c04.cor.4-3-8}
\expandafter\def\csname sga@target@IV.4.3.9\endcsname{ega4.c04.prop.4-3-9}
\expandafter\def\csname sga@target@IV.4.4.1\endcsname{ega4.c04.statement.4-4-1}
\expandafter\def\csname sga@target@IV.4.4.2\endcsname{ega4.c04.lem.4-4-2}
\expandafter\def\csname sga@target@IV.4.4.2.i\endcsname{ega4.c04.sec.4-4-2-i}
\expandafter\def\csname sga@target@IV.4.4.2.ii\endcsname{ega4.c04.sec.4-4-2-ii}
\expandafter\def\csname sga@target@IV.4.4.3\endcsname{ega4.c04.rem.4-4-3}
\expandafter\def\csname sga@target@IV.4.4.4\endcsname{ega4.c04.thm.4-4-4}
\expandafter\def\csname sga@target@IV.4.4.5\endcsname{ega4.c04.cor.4-4-5}
\expandafter\def\csname sga@target@IV.4.4.6\endcsname{ega4.c04.cor.4-4-6}
\expandafter\def\csname sga@target@IV.4.5\endcsname{ega4.c04.subsec.4-5.key-iv-4-5}
\expandafter\def\csname sga@target@IV.4.5.1\endcsname{ega4.c04.prop.4-5-1}
\expandafter\def\csname sga@target@IV.4.5.10\endcsname{ega4.c04.cor.4-5-10}
\expandafter\def\csname sga@target@IV.4.5.11\endcsname{ega4.c04.cor.4-5-11}
\expandafter\def\csname sga@target@IV.4.5.12\endcsname{ega4.c04.sec.4-5-12}
\expandafter\def\csname sga@target@IV.4.5.12.i\endcsname{ega4.c04.sec.4-5-12-i}
\expandafter\def\csname sga@target@IV.4.5.12.ii\endcsname{ega4.c04.sec.4-5-12-ii}
\expandafter\def\csname sga@target@IV.4.5.13\endcsname{ega4.c04.prop.4-5-13}
\expandafter\def\csname sga@target@IV.4.5.13.1\endcsname{ega4.c04.cor.4-5-13-1}
\expandafter\def\csname sga@target@IV.4.5.13.1.i\endcsname{ega4.c04.sec.4-5-13-1-i}
\expandafter\def\csname sga@target@IV.4.5.13.1.ii\endcsname{ega4.c04.sec.4-5-13-1-ii}
\expandafter\def\csname sga@target@IV.4.5.14\endcsname{ega4.c04.cor.4-5-14}
\expandafter\def\csname sga@target@IV.4.5.15\endcsname{ega4.c04.prop.4-5-15}
\expandafter\def\csname sga@target@IV.4.5.15.i\endcsname{ega4.c04.sec.4-5-15-i}
\expandafter\def\csname sga@target@IV.4.5.15.ii\endcsname{ega4.c04.sec.4-5-15-ii}
\expandafter\def\csname sga@target@IV.4.5.16\endcsname{ega4.c04.cor.4-5-16}
\expandafter\def\csname sga@target@IV.4.5.16.i\endcsname{ega4.c04.sec.4-5-16-i}
\expandafter\def\csname sga@target@IV.4.5.16.ii\endcsname{ega4.c04.sec.4-5-16-ii}
\expandafter\def\csname sga@target@IV.4.5.17\endcsname{ega4.c04.cor.4-5-17}
\expandafter\def\csname sga@target@IV.4.5.18\endcsname{ega4.c04.rem.4-5-18}
\expandafter\def\csname sga@target@IV.4.5.19\endcsname{ega4.c04.prop.4-5-19}
\expandafter\def\csname sga@target@IV.4.5.19.1\endcsname{ega4.c04.lem.4-5-19-1}
\expandafter\def\csname sga@target@IV.4.5.19.2\endcsname{ega4.c04.eq.4-5-19-2}
\expandafter\def\csname sga@target@IV.4.5.19.3\endcsname{ega4.c04.cor.4-5-19-3}
\expandafter\def\csname sga@target@IV.4.5.2\endcsname{ega4.c04.def.4-5-2}
\expandafter\def\csname sga@target@IV.4.5.2.a\endcsname{ega4.c04.sec.4-5-2-a}
\expandafter\def\csname sga@target@IV.4.5.2.b\endcsname{ega4.c04.sec.4-5-2-b}
\expandafter\def\csname sga@target@IV.4.5.2.c\endcsname{ega4.c04.sec.4-5-2-c}
\expandafter\def\csname sga@target@IV.4.5.20\endcsname{ega4.c04.rem.4-5-20}
\expandafter\def\csname sga@target@IV.4.5.21\endcsname{ega4.c04.prop.4-5-21}
\expandafter\def\csname sga@target@IV.4.5.3\endcsname{ega4.c04.statement.4-5-3}
\expandafter\def\csname sga@target@IV.4.5.4\endcsname{ega4.c04.prop.4-5-4}
\expandafter\def\csname sga@target@IV.4.5.5\endcsname{ega4.c04.def.4-5-5}
\expandafter\def\csname sga@target@IV.4.5.6\endcsname{ega4.c04.prop.4-5-6}
\expandafter\def\csname sga@target@IV.4.5.6.i\endcsname{ega4.c04.sec.4-5-6-i}
\expandafter\def\csname sga@target@IV.4.5.6.ii\endcsname{ega4.c04.sec.4-5-6-ii}
\expandafter\def\csname sga@target@IV.4.5.7\endcsname{ega4.c04.prop.4-5-7}
\expandafter\def\csname sga@target@IV.4.5.8\endcsname{ega4.c04.cor.4-5-8}
\expandafter\def\csname sga@target@IV.4.5.8.i\endcsname{ega4.c04.sec.4-5-8-i}
\expandafter\def\csname sga@target@IV.4.5.8.ii\endcsname{ega4.c04.sec.4-5-8-ii}
\expandafter\def\csname sga@target@IV.4.5.9\endcsname{ega4.c04.prop.4-5-9}
\expandafter\def\csname sga@target@IV.4.5.9.a\endcsname{ega4.c04.sec.4-5-9-a}
\expandafter\def\csname sga@target@IV.4.5.9.b\endcsname{ega4.c04.sec.4-5-9-b}
\expandafter\def\csname sga@target@IV.4.5.9.c\endcsname{ega4.c04.sec.4-5-9-c}
\expandafter\def\csname sga@target@IV.4.6\endcsname{ega4.c04.subsec.4-6.key-iv-4-6}
\expandafter\def\csname sga@target@IV.4.6.1\endcsname{ega4.c04.prop.4-6-1}
\expandafter\def\csname sga@target@IV.4.6.1.a\endcsname{ega4.c04.sec.4-6-1-a}
\expandafter\def\csname sga@target@IV.4.6.1.b\endcsname{ega4.c04.sec.4-6-1-b}
\expandafter\def\csname sga@target@IV.4.6.1.c\endcsname{ega4.c04.sec.4-6-1-c}
\expandafter\def\csname sga@target@IV.4.6.1.d\endcsname{ega4.c04.sec.4-6-1-d}
\expandafter\def\csname sga@target@IV.4.6.1.e\endcsname{ega4.c04.sec.4-6-1-e}
\expandafter\def\csname sga@target@IV.4.6.10\endcsname{ega4.c04.prop.4-6-10}
\expandafter\def\csname sga@target@IV.4.6.11\endcsname{ega4.c04.cor.4-6-11}
\expandafter\def\csname sga@target@IV.4.6.12\endcsname{ega4.c04.prop.4-6-12}
\expandafter\def\csname sga@target@IV.4.6.12.a\endcsname{ega4.c04.sec.4-6-12-a}
\expandafter\def\csname sga@target@IV.4.6.12.b\endcsname{ega4.c04.sec.4-6-12-b}
\expandafter\def\csname sga@target@IV.4.6.12.c\endcsname{ega4.c04.sec.4-6-12-c}
\expandafter\def\csname sga@target@IV.4.6.12.d\endcsname{ega4.c04.sec.4-6-12-d}
\expandafter\def\csname sga@target@IV.4.6.12.e\endcsname{ega4.c04.sec.4-6-12-e}
\expandafter\def\csname sga@target@IV.4.6.13\endcsname{ega4.c04.cor.4-6-13}
\expandafter\def\csname sga@target@IV.4.6.13.f\endcsname{ega4.c04.sec.4-6-13-f}
\expandafter\def\csname sga@target@IV.4.6.14\endcsname{ega4.c04.statement.4-6-14}
\expandafter\def\csname sga@target@IV.4.6.15\endcsname{ega4.c04.prop.4-6-15}
\expandafter\def\csname sga@target@IV.4.6.15.i\endcsname{ega4.c04.sec.4-6-15-i}
\expandafter\def\csname sga@target@IV.4.6.15.ii\endcsname{ega4.c04.sec.4-6-15-ii}
\expandafter\def\csname sga@target@IV.4.6.16\endcsname{ega4.c04.prop.4-6-16}
\expandafter\def\csname sga@target@IV.4.6.16.a\endcsname{ega4.c04.sec.4-6-16-a}
\expandafter\def\csname sga@target@IV.4.6.16.b\endcsname{ega4.c04.sec.4-6-16-b}
\expandafter\def\csname sga@target@IV.4.6.16.c\endcsname{ega4.c04.sec.4-6-16-c}
\expandafter\def\csname sga@target@IV.4.6.16.d\endcsname{ega4.c04.sec.4-6-16-d}
\expandafter\def\csname sga@target@IV.4.6.17\endcsname{ega4.c04.def.4-6-17}
\expandafter\def\csname sga@target@IV.4.6.18\endcsname{ega4.c04.prop.4-6-18}
\expandafter\def\csname sga@target@IV.4.6.18.a\endcsname{ega4.c04.sec.4-6-18-a}
\expandafter\def\csname sga@target@IV.4.6.18.b\endcsname{ega4.c04.sec.4-6-18-b}
\expandafter\def\csname sga@target@IV.4.6.18.c\endcsname{ega4.c04.sec.4-6-18-c}
\expandafter\def\csname sga@target@IV.4.6.18.d\endcsname{ega4.c04.sec.4-6-18-d}
\expandafter\def\csname sga@target@IV.4.6.19\endcsname{ega4.c04.def.4-6-19}
\expandafter\def\csname sga@target@IV.4.6.2\endcsname{ega4.c04.def.4-6-2}
\expandafter\def\csname sga@target@IV.4.6.20\endcsname{ega4.c04.prop.4-6-20}
\expandafter\def\csname sga@target@IV.4.6.21\endcsname{ega4.c04.prop.4-6-21}
\expandafter\def\csname sga@target@IV.4.6.21.i\endcsname{ega4.c04.sec.4-6-21-i}
\expandafter\def\csname sga@target@IV.4.6.21.ii\endcsname{ega4.c04.sec.4-6-21-ii}
\expandafter\def\csname sga@target@IV.4.6.22\endcsname{ega4.c04.def.4-6-22}
\expandafter\def\csname sga@target@IV.4.6.3\endcsname{ega4.c04.cor.4-6-3}
\expandafter\def\csname sga@target@IV.4.6.4\endcsname{ega4.c04.cor.4-6-4}
\expandafter\def\csname sga@target@IV.4.6.5\endcsname{ega4.c04.prop.4-6-5}
\expandafter\def\csname sga@target@IV.4.6.5.i\endcsname{ega4.c04.sec.4-6-5-i}
\expandafter\def\csname sga@target@IV.4.6.5.ii\endcsname{ega4.c04.sec.4-6-5-ii}
\expandafter\def\csname sga@target@IV.4.6.6\endcsname{ega4.c04.prop.4-6-6}
\expandafter\def\csname sga@target@IV.4.6.7\endcsname{ega4.c04.cor.4-6-7}
\expandafter\def\csname sga@target@IV.4.6.8\endcsname{ega4.c04.cor.4-6-8}
\expandafter\def\csname sga@target@IV.4.6.9\endcsname{ega4.c04.def.4-6-9}
\expandafter\def\csname sga@target@IV.4.7.1\endcsname{ega4.c04.lem.4-7-1}
\expandafter\def\csname sga@target@IV.4.7.1.1\endcsname{ega4.c04.eq.4-7-1-1}
\expandafter\def\csname sga@target@IV.4.7.10\endcsname{ega4.c04.prop.4-7-10}
\expandafter\def\csname sga@target@IV.4.7.11\endcsname{ega4.c04.prop.4-7-11}
\expandafter\def\csname sga@target@IV.4.7.12\endcsname{ega4.c04.def.4-7-12}
\expandafter\def\csname sga@target@IV.4.7.13\endcsname{ega4.c04.prop.4-7-13}
\expandafter\def\csname sga@target@IV.4.7.2\endcsname{ega4.c04.cor.4-7-2}
\expandafter\def\csname sga@target@IV.4.7.2.1\endcsname{ega4.c04.eq.4-7-2-1}
\expandafter\def\csname sga@target@IV.4.7.3\endcsname{ega4.c04.prop.4-7-3}
\expandafter\def\csname sga@target@IV.4.7.3.1\endcsname{ega4.c04.eq.4-7-3-1}
\expandafter\def\csname sga@target@IV.4.7.3.i\endcsname{ega4.c04.sec.4-7-3-i}
\expandafter\def\csname sga@target@IV.4.7.3.ii\endcsname{ega4.c04.sec.4-7-3-ii}
\expandafter\def\csname sga@target@IV.4.7.3.iii\endcsname{ega4.c04.sec.4-7-3-iii}
\expandafter\def\csname sga@target@IV.4.7.4\endcsname{ega4.c04.def.4-7-4}
\expandafter\def\csname sga@target@IV.4.7.5\endcsname{ega4.c04.def.4-7-5}
\expandafter\def\csname sga@target@IV.4.7.6\endcsname{ega4.c04.prop.4-7-6}
\expandafter\def\csname sga@target@IV.4.7.6.1\endcsname{ega4.c04.eq.4-7-6-1}
\expandafter\def\csname sga@target@IV.4.7.7\endcsname{ega4.c04.statement.4-7-7}
\expandafter\def\csname sga@target@IV.4.7.8\endcsname{ega4.c04.prop.4-7-8}
\expandafter\def\csname sga@target@IV.4.7.9\endcsname{ega4.c04.cor.4-7-9}
\expandafter\def\csname sga@target@IV.4.8.1\endcsname{ega4.c04.statement.4-8-1}
\expandafter\def\csname sga@target@IV.4.8.10\endcsname{ega4.c04.cor.4-8-10}
\expandafter\def\csname sga@target@IV.4.8.11\endcsname{ega4.c04.cor.4-8-11}
\expandafter\def\csname sga@target@IV.4.8.12\endcsname{ega4.c04.cor.4-8-12}
\expandafter\def\csname sga@target@IV.4.8.13\endcsname{ega4.c04.prop.4-8-13}
\expandafter\def\csname sga@target@IV.4.8.14\endcsname{ega4.c04.prop.4-8-14}
\expandafter\def\csname sga@target@IV.4.8.15\endcsname{ega4.c04.sec.4-8-15}
\expandafter\def\csname sga@target@IV.4.8.15.i\endcsname{ega4.c04.sec.4-8-15-i}
\expandafter\def\csname sga@target@IV.4.8.15.ii\endcsname{ega4.c04.sec.4-8-15-ii}
\expandafter\def\csname sga@target@IV.4.8.2\endcsname{ega4.c04.statement.4-8-2}
\expandafter\def\csname sga@target@IV.4.8.2.1\endcsname{ega4.c04.eq.4-8-2-1}
\expandafter\def\csname sga@target@IV.4.8.2.2\endcsname{ega4.c04.eq.4-8-2-2}
\expandafter\def\csname sga@target@IV.4.8.2.3\endcsname{ega4.c04.eq.4-8-2-3}
\expandafter\def\csname sga@target@IV.4.8.2.4\endcsname{ega4.c04.eq.4-8-2-4}
\expandafter\def\csname sga@target@IV.4.8.2.5\endcsname{ega4.c04.eq.4-8-2-5}
\expandafter\def\csname sga@target@IV.4.8.2.6\endcsname{ega4.c04.eq.4-8-2-6}
\expandafter\def\csname sga@target@IV.4.8.3\endcsname{ega4.c04.prop.4-8-3}
\expandafter\def\csname sga@target@IV.4.8.4\endcsname{ega4.c04.def.4-8-4}
\expandafter\def\csname sga@target@IV.4.8.5\endcsname{ega4.c04.prop.4-8-5}
\expandafter\def\csname sga@target@IV.4.8.6\endcsname{ega4.c04.lem.4-8-6}
\expandafter\def\csname sga@target@IV.4.8.6.a\endcsname{ega4.c04.sec.4-8-6-a}
\expandafter\def\csname sga@target@IV.4.8.6.b\endcsname{ega4.c04.sec.4-8-6-b}
\expandafter\def\csname sga@target@IV.4.8.6.c\endcsname{ega4.c04.sec.4-8-6-c}
\expandafter\def\csname sga@target@IV.4.8.7\endcsname{ega4.c04.cor.4-8-7}
\expandafter\def\csname sga@target@IV.4.8.8\endcsname{ega4.c04.lem.4-8-8}
\expandafter\def\csname sga@target@IV.4.8.8.i\endcsname{ega4.c04.sec.4-8-8-i}
\expandafter\def\csname sga@target@IV.4.8.8.ii\endcsname{ega4.c04.sec.4-8-8-ii}
\expandafter\def\csname sga@target@IV.4.8.8.iii\endcsname{ega4.c04.sec.4-8-8-iii}
\expandafter\def\csname sga@target@IV.4.8.9\endcsname{ega4.c04.prop.4-8-9}
\expandafter\def\csname sga@target@IV.4.9.1\endcsname{ega4.c04.prop.4-9-1}
\expandafter\def\csname sga@target@IV.4.9.1.a\endcsname{ega4.c04.sec.4-9-1-a}
\expandafter\def\csname sga@target@IV.4.9.1.b\endcsname{ega4.c04.sec.4-9-1-b}
\expandafter\def\csname sga@target@IV.4.9.1.bprime\endcsname{ega4.c04.sec.4-9-1-bprime}
\expandafter\def\csname sga@target@IV.4.9.1.c\endcsname{ega4.c04.sec.4-9-1-c}
\expandafter\def\csname sga@target@IV.4.9.2\endcsname{ega4.c04.cor.4-9-2}
\expandafter\def\csname sga@target@IV.4.9.3\endcsname{ega4.c04.rem.4-9-3}
\expandafter\def\csname sga@target@IV.4.9.4\endcsname{ega4.c04.prop.4-9-4}
\expandafter\def\csname sga@target@IV.4.9.5\endcsname{ega4.c04.cor.4-9-5}
\expandafter\def\csname sga@target@IV.4.9.5.a\endcsname{ega4.c04.sec.4-9-5-a}
\expandafter\def\csname sga@target@IV.4.9.5.b\endcsname{ega4.c04.sec.4-9-5-b}
\expandafter\def\csname sga@target@IV.4.9.6\endcsname{ega4.c04.prop.4-9-6}
\expandafter\def\csname sga@target@IV.4.9.6.a\endcsname{ega4.c04.sec.4-9-6-a}
\expandafter\def\csname sga@target@IV.4.9.6.aprime\endcsname{ega4.c04.sec.4-9-6-aprime}
\expandafter\def\csname sga@target@IV.4.9.6.b\endcsname{ega4.c04.sec.4-9-6-b}
\expandafter\def\csname sga@target@IV.4.9.6.bprime\endcsname{ega4.c04.sec.4-9-6-bprime}
\expandafter\def\csname sga@target@IV.4.9.7\endcsname{ega4.c04.cor.4-9-7}
\expandafter\def\csname sga@target@IV.5.1.1\endcsname{ega4.c04.statement.5-1-1}
\expandafter\def\csname sga@target@IV.5.1.1.1\endcsname{ega4.c04.eq.5-1-1-1}
\expandafter\def\csname sga@target@IV.5.1.1.2\endcsname{ega4.c04.eq.5-1-1-2}
\expandafter\def\csname sga@target@IV.5.1.1.3\endcsname{ega4.c04.eq.5-1-1-3}
\expandafter\def\csname sga@target@IV.5.1.1.4\endcsname{ega4.c04.sec.5-1-1-4}
\expandafter\def\csname sga@target@IV.5.1.1.5\endcsname{ega4.c04.sec.5-1-1-5}
\expandafter\def\csname sga@target@IV.5.1.1.6\endcsname{ega4.c04.sec.5-1-1-6}
\expandafter\def\csname sga@target@IV.5.1.10\endcsname{ega4.c04.prop.5-1-10}
\expandafter\def\csname sga@target@IV.5.1.10.1\endcsname{ega4.c04.lem.5-1-10-1}
\expandafter\def\csname sga@target@IV.5.1.10.i\endcsname{ega4.c04.sec.5-1-10-i}
\expandafter\def\csname sga@target@IV.5.1.10.ii\endcsname{ega4.c04.sec.5-1-10-ii}
\expandafter\def\csname sga@target@IV.5.1.11\endcsname{ega4.c04.cor.5-1-11}
\expandafter\def\csname sga@target@IV.5.1.12\endcsname{ega4.c04.statement.5-1-12}
\expandafter\def\csname sga@target@IV.5.1.12.1\endcsname{ega4.c04.eq.5-1-12-1}
\expandafter\def\csname sga@target@IV.5.1.12.2\endcsname{ega4.c04.eq.5-1-12-2}
\expandafter\def\csname sga@target@IV.5.1.12.3\endcsname{ega4.c04.eq.5-1-12-3}
\expandafter\def\csname sga@target@IV.5.1.13\endcsname{ega4.c04.prop.5-1-13}
\expandafter\def\csname sga@target@IV.5.1.13.1\endcsname{ega4.c04.eq.5-1-13-1}
\expandafter\def\csname sga@target@IV.5.1.2\endcsname{ega4.c04.prop.5-1-2}
\expandafter\def\csname sga@target@IV.5.1.2.1\endcsname{ega4.c04.eq.5-1-2-1}
\expandafter\def\csname sga@target@IV.5.1.3\endcsname{ega4.c04.cor.5-1-3}
\expandafter\def\csname sga@target@IV.5.1.3.1\endcsname{ega4.c04.eq.5-1-3-1}
\expandafter\def\csname sga@target@IV.5.1.3.2\endcsname{ega4.c04.eq.5-1-3-2}
\expandafter\def\csname sga@target@IV.5.1.4\endcsname{ega4.c04.prop.5-1-4}
\expandafter\def\csname sga@target@IV.5.1.4.1\endcsname{ega4.c04.eq.5-1-4-1}
\expandafter\def\csname sga@target@IV.5.1.4.2\endcsname{ega4.c04.eq.5-1-4-2}
\expandafter\def\csname sga@target@IV.5.1.5\endcsname{ega4.c04.cor.5-1-5}
\expandafter\def\csname sga@target@IV.5.1.6\endcsname{ega4.c04.statement.5-1-6}
\expandafter\def\csname sga@target@IV.5.1.7\endcsname{ega4.c04.prop.5-1-7}
\expandafter\def\csname sga@target@IV.5.1.8\endcsname{ega4.c04.cor.5-1-8}
\expandafter\def\csname sga@target@IV.5.1.9\endcsname{ega4.c04.prop.5-1-9}
\expandafter\def\csname sga@target@IV.5.1.9.1\endcsname{ega4.c04.statement.5-1-9-1}
\expandafter\def\csname sga@target@IV.5.10.1\endcsname{ega4.c04.statement.5-10-1}
\expandafter\def\csname sga@target@IV.5.10.1.1\endcsname{ega4.c04.eq.5-10-1-1}
\expandafter\def\csname sga@target@IV.5.10.10\endcsname{ega4.c04.prop.5-10-10}
\expandafter\def\csname sga@target@IV.5.10.11\endcsname{ega4.c04.statement.5-10-11}
\expandafter\def\csname sga@target@IV.5.10.11.1\endcsname{ega4.c04.eq.5-10-11-1}
\expandafter\def\csname sga@target@IV.5.10.12\endcsname{ega4.c04.prop.5-10-12}
\expandafter\def\csname sga@target@IV.5.10.13\endcsname{ega4.c04.statement.5-10-13}
\expandafter\def\csname sga@target@IV.5.10.14\endcsname{ega4.c04.prop.5-10-14}
\expandafter\def\csname sga@target@IV.5.10.15\endcsname{ega4.c04.cor.5-10-15}
\expandafter\def\csname sga@target@IV.5.10.16\endcsname{ega4.c04.prop.5-10-16}
\expandafter\def\csname sga@target@IV.5.10.17\endcsname{ega4.c04.prop.5-10-17}
\expandafter\def\csname sga@target@IV.5.10.2\endcsname{ega4.c04.prop.5-10-2}
\expandafter\def\csname sga@target@IV.5.10.3\endcsname{ega4.c04.cor.5-10-3}
\expandafter\def\csname sga@target@IV.5.10.4\endcsname{ega4.c04.cor.5-10-4}
\expandafter\def\csname sga@target@IV.5.10.5\endcsname{ega4.c04.thm.5-10-5}
\expandafter\def\csname sga@target@IV.5.10.6\endcsname{ega4.c04.cor.5-10-6}
\expandafter\def\csname sga@target@IV.5.10.7\endcsname{ega4.c04.thm.5-10-7}
\expandafter\def\csname sga@target@IV.5.10.8\endcsname{ega4.c04.cor.5-10-8}
\expandafter\def\csname sga@target@IV.5.10.8.1\endcsname{ega4.c04.lem.5-10-8-1}
\expandafter\def\csname sga@target@IV.5.10.9\endcsname{ega4.c04.cor.5-10-9}
\expandafter\def\csname sga@target@IV.5.11.1\endcsname{ega4.c04.prop.5-11-1}
\expandafter\def\csname sga@target@IV.5.11.1.1\endcsname{ega4.c04.lem.5-11-1-1}
\expandafter\def\csname sga@target@IV.5.11.2\endcsname{ega4.c04.cor.5-11-2}
\expandafter\def\csname sga@target@IV.5.11.3\endcsname{ega4.c04.rem.5-11-3}
\expandafter\def\csname sga@target@IV.5.11.4\endcsname{ega4.c04.cor.5-11-4}
\expandafter\def\csname sga@target@IV.5.11.5\endcsname{ega4.c04.cor.5-11-5}
\expandafter\def\csname sga@target@IV.5.11.6\endcsname{ega4.c04.cor.5-11-6}
\expandafter\def\csname sga@target@IV.5.11.7\endcsname{ega4.c04.item.5-11-7}
\expandafter\def\csname sga@target@IV.5.12\endcsname{ega4.c04.subsec.5-12.key-iv-5-12}
\expandafter\def\csname sga@target@IV.5.12.1\endcsname{ega4.c04.prop.5-12-1}
\expandafter\def\csname sga@target@IV.5.12.1.1\endcsname{ega4.c04.eq.5-12-1-1}
\expandafter\def\csname sga@target@IV.5.12.10\endcsname{ega4.c04.cor.5-12-10}
\expandafter\def\csname sga@target@IV.5.12.2\endcsname{ega4.c04.prop.5-12-2}
\expandafter\def\csname sga@target@IV.5.12.3\endcsname{ega4.c04.rem.5-12-3}
\expandafter\def\csname sga@target@IV.5.12.4\endcsname{ega4.c04.cor.5-12-4}
\expandafter\def\csname sga@target@IV.5.12.5\endcsname{ega4.c04.prop.5-12-5}
\expandafter\def\csname sga@target@IV.5.12.6\endcsname{ega4.c04.rem.5-12-6}
\expandafter\def\csname sga@target@IV.5.12.7\endcsname{ega4.c04.cor.5-12-7}
\expandafter\def\csname sga@target@IV.5.12.8\endcsname{ega4.c04.prop.5-12-8}
\expandafter\def\csname sga@target@IV.5.12.9\endcsname{ega4.c04.rem.5-12-9}
\expandafter\def\csname sga@target@IV.5.13.1\endcsname{ega4.c04.prop.5-13-1}
\expandafter\def\csname sga@target@IV.5.13.2\endcsname{ega4.c04.cor.5-13-2}
\expandafter\def\csname sga@target@IV.5.13.3\endcsname{ega4.c04.prop.5-13-3}
\expandafter\def\csname sga@target@IV.5.13.4\endcsname{ega4.c04.cor.5-13-4}
\expandafter\def\csname sga@target@IV.5.13.5\endcsname{ega4.c04.statement.5-13-5}
\expandafter\def\csname sga@target@IV.5.13.6\endcsname{ega4.c04.cor.5-13-6}
\expandafter\def\csname sga@target@IV.5.13.7\endcsname{ega4.c04.prop.5-13-7}
\expandafter\def\csname sga@target@IV.5.13.7.1\endcsname{ega4.c04.lem.5-13-7-1}
\expandafter\def\csname sga@target@IV.5.13.8\endcsname{ega4.c04.cor.5-13-8}
\expandafter\def\csname sga@target@IV.5.2.1\endcsname{ega4.c04.prop.5-2-1}
\expandafter\def\csname sga@target@IV.5.2.1.1\endcsname{ega4.c04.eq.5-2-1-1}
\expandafter\def\csname sga@target@IV.5.2.2\endcsname{ega4.c04.cor.5-2-2}
\expandafter\def\csname sga@target@IV.5.2.3\endcsname{ega4.c04.cor.5-2-3}
\expandafter\def\csname sga@target@IV.5.2.3.1\endcsname{ega4.c04.eq.5-2-3-1}
\expandafter\def\csname sga@target@IV.5.2.4\endcsname{ega4.c04.cor.5-2-4}
\expandafter\def\csname sga@target@IV.5.2.5\endcsname{ega4.c04.sec.5-2-5}
\expandafter\def\csname sga@target@IV.5.2.5.i\endcsname{ega4.c04.sec.5-2-5-i}
\expandafter\def\csname sga@target@IV.5.2.5.ii\endcsname{ega4.c04.sec.5-2-5-ii}
\expandafter\def\csname sga@target@IV.5.3.1\endcsname{ega4.c04.prop.5-3-1}
\expandafter\def\csname sga@target@IV.5.3.1.1\endcsname{ega4.c04.eq.5-3-1-1}
\expandafter\def\csname sga@target@IV.5.3.1.2\endcsname{ega4.c04.eq.5-3-1-2}
\expandafter\def\csname sga@target@IV.5.3.1.3\endcsname{ega4.c04.lem.5-3-1-3}
\expandafter\def\csname sga@target@IV.5.4.1\endcsname{ega4.c04.prop.5-4-1}
\expandafter\def\csname sga@target@IV.5.4.1.i\endcsname{ega4.c04.sec.5-4-1-i}
\expandafter\def\csname sga@target@IV.5.4.1.ii\endcsname{ega4.c04.sec.5-4-1-ii}
\expandafter\def\csname sga@target@IV.5.4.2\endcsname{ega4.c04.cor.5-4-2}
\expandafter\def\csname sga@target@IV.5.4.3\endcsname{ega4.c04.sec.5-4-3}
\expandafter\def\csname sga@target@IV.5.4.3.1\endcsname{ega4.c04.eq.5-4-3-1}
\expandafter\def\csname sga@target@IV.5.4.3.i\endcsname{ega4.c04.sec.5-4-3-i}
\expandafter\def\csname sga@target@IV.5.4.3.ii\endcsname{ega4.c04.sec.5-4-3-ii}
\expandafter\def\csname sga@target@IV.5.5\endcsname{ega4.c04.sec.5-5}
\expandafter\def\csname sga@target@IV.5.5.1\endcsname{ega4.c04.statement.5-5-1}
\expandafter\def\csname sga@target@IV.5.5.1.1\endcsname{ega4.c04.eq.5-5-1-1}
\expandafter\def\csname sga@target@IV.5.5.2\endcsname{ega4.c04.prop.5-5-2}
\expandafter\def\csname sga@target@IV.5.5.2.1\endcsname{ega4.c04.eq.5-5-2-1}
\expandafter\def\csname sga@target@IV.5.5.2.2\endcsname{ega4.c04.eq.5-5-2-2}
\expandafter\def\csname sga@target@IV.5.5.3\endcsname{ega4.c04.prop.5-5-3}
\expandafter\def\csname sga@target@IV.5.5.3.1\endcsname{ega4.c04.eq.5-5-3-1}
\expandafter\def\csname sga@target@IV.5.5.4\endcsname{ega4.c04.cor.5-5-4}
\expandafter\def\csname sga@target@IV.5.5.4.1\endcsname{ega4.c04.eq.5-5-4-1}
\expandafter\def\csname sga@target@IV.5.5.5\endcsname{ega4.c04.cor.5-5-5}
\expandafter\def\csname sga@target@IV.5.5.6\endcsname{ega4.c04.cor.5-5-6}
\expandafter\def\csname sga@target@IV.5.5.6.1\endcsname{ega4.c04.eq.5-5-6-1}
\expandafter\def\csname sga@target@IV.5.5.7\endcsname{ega4.c04.lem.5-5-7}
\expandafter\def\csname sga@target@IV.5.5.7.1\endcsname{ega4.c04.eq.5-5-7-1}
\expandafter\def\csname sga@target@IV.5.5.7.2\endcsname{ega4.c04.statement.5-5-7-2}
\expandafter\def\csname sga@target@IV.5.5.7.3\endcsname{ega4.c04.eq.5-5-7-3}
\expandafter\def\csname sga@target@IV.5.5.7.i\endcsname{ega4.c04.sec.5-5-7-i}
\expandafter\def\csname sga@target@IV.5.5.7.ii\endcsname{ega4.c04.sec.5-5-7-ii}
\expandafter\def\csname sga@target@IV.5.5.8\endcsname{ega4.c04.thm.5-5-8}
\expandafter\def\csname sga@target@IV.5.5.8.1\endcsname{ega4.c04.eq.5-5-8-1}
\expandafter\def\csname sga@target@IV.5.5.8.2\endcsname{ega4.c04.eq.5-5-8-2}
\expandafter\def\csname sga@target@IV.5.5.8.3\endcsname{ega4.c04.eq.5-5-8-3}
\expandafter\def\csname sga@target@IV.5.6\endcsname{ega4.c04.sec.5-6}
\expandafter\def\csname sga@target@IV.5.6.1\endcsname{ega4.c04.prop.5-6-1}
\expandafter\def\csname sga@target@IV.5.6.1.1\endcsname{ega4.c04.eq.5-6-1-1}
\expandafter\def\csname sga@target@IV.5.6.1.2\endcsname{ega4.c04.eq.5-6-1-2}
\expandafter\def\csname sga@target@IV.5.6.1.3\endcsname{ega4.c04.statement.5-6-1-3}
\expandafter\def\csname sga@target@IV.5.6.1.4\endcsname{ega4.c04.statement.5-6-1-4}
\expandafter\def\csname sga@target@IV.5.6.1.5\endcsname{ega4.c04.statement.5-6-1-5}
\expandafter\def\csname sga@target@IV.5.6.1.a\endcsname{ega4.c04.sec.5-6-1-a}
\expandafter\def\csname sga@target@IV.5.6.1.b\endcsname{ega4.c04.sec.5-6-1-b}
\expandafter\def\csname sga@target@IV.5.6.1.c\endcsname{ega4.c04.sec.5-6-1-c}
\expandafter\def\csname sga@target@IV.5.6.10\endcsname{ega4.c04.prop.5-6-10}
\expandafter\def\csname sga@target@IV.5.6.11\endcsname{ega4.c04.ex.5-6-11}
\expandafter\def\csname sga@target@IV.5.6.11.1\endcsname{ega4.c04.lem.5-6-11-1}
\expandafter\def\csname sga@target@IV.5.6.11.1.i\endcsname{ega4.c04.sec.5-6-11-1-i}
\expandafter\def\csname sga@target@IV.5.6.11.1.ii\endcsname{ega4.c04.sec.5-6-11-1-ii}
\expandafter\def\csname sga@target@IV.5.6.2\endcsname{ega4.c04.def.5-6-2}
\expandafter\def\csname sga@target@IV.5.6.3\endcsname{ega4.c04.sec.5-6-3}
\expandafter\def\csname sga@target@IV.5.6.3.i\endcsname{ega4.c04.sec.5-6-3-i}
\expandafter\def\csname sga@target@IV.5.6.3.ii\endcsname{ega4.c04.sec.5-6-3-ii}
\expandafter\def\csname sga@target@IV.5.6.3.iii\endcsname{ega4.c04.sec.5-6-3-iii}
\expandafter\def\csname sga@target@IV.5.6.3.iv\endcsname{ega4.c04.sec.5-6-3-iv}
\expandafter\def\csname sga@target@IV.5.6.4\endcsname{ega4.c04.prop.5-6-4}
\expandafter\def\csname sga@target@IV.5.6.5\endcsname{ega4.c04.prop.5-6-5}
\expandafter\def\csname sga@target@IV.5.6.5.1\endcsname{ega4.c04.eq.5-6-5-1}
\expandafter\def\csname sga@target@IV.5.6.5.2\endcsname{ega4.c04.eq.5-6-5-2}
\expandafter\def\csname sga@target@IV.5.6.5.3\endcsname{ega4.c04.eq.5-6-5-3}
\expandafter\def\csname sga@target@IV.5.6.6\endcsname{ega4.c04.cor.5-6-6}
\expandafter\def\csname sga@target@IV.5.6.6.1\endcsname{ega4.c04.eq.5-6-6-1}
\expandafter\def\csname sga@target@IV.5.6.6.2\endcsname{ega4.c04.eq.5-6-6-2}
\expandafter\def\csname sga@target@IV.5.6.7\endcsname{ega4.c04.cor.5-6-7}
\expandafter\def\csname sga@target@IV.5.6.7.1\endcsname{ega4.c04.eq.5-6-7-1}
\expandafter\def\csname sga@target@IV.5.6.8\endcsname{ega4.c04.cor.5-6-8}
\expandafter\def\csname sga@target@IV.5.6.8.1\endcsname{ega4.c04.eq.5-6-8-1}
\expandafter\def\csname sga@target@IV.5.6.8.2\endcsname{ega4.c04.eq.5-6-8-2}
\expandafter\def\csname sga@target@IV.5.6.9\endcsname{ega4.c04.sec.5-6-9}
\expandafter\def\csname sga@target@IV.5.6.9.i\endcsname{ega4.c04.sec.5-6-9-i}
\expandafter\def\csname sga@target@IV.5.6.9.ii\endcsname{ega4.c04.sec.5-6-9-ii}
\expandafter\def\csname sga@target@IV.5.7.1\endcsname{ega4.c04.def.5-7-1}
\expandafter\def\csname sga@target@IV.5.7.1.1\endcsname{ega4.c04.eq.5-7-1-1}
\expandafter\def\csname sga@target@IV.5.7.10\endcsname{ega4.c04.cor.5-7-10}
\expandafter\def\csname sga@target@IV.5.7.11\endcsname{ega4.c04.cor.5-7-11}
\expandafter\def\csname sga@target@IV.5.7.11.i\endcsname{ega4.c04.sec.5-7-11-i}
\expandafter\def\csname sga@target@IV.5.7.11.ii\endcsname{ega4.c04.sec.5-7-11-ii}
\expandafter\def\csname sga@target@IV.5.7.12\endcsname{ega4.c04.statement.5-7-12}
\expandafter\def\csname sga@target@IV.5.7.13\endcsname{ega4.c04.cor.5-7-13}
\expandafter\def\csname sga@target@IV.5.7.13.1\endcsname{ega4.c04.eq.5-7-13-1}
\expandafter\def\csname sga@target@IV.5.7.2\endcsname{ega4.c04.def.5-7-2}
\expandafter\def\csname sga@target@IV.5.7.2.1\endcsname{ega4.c04.eq.5-7-2-1}
\expandafter\def\csname sga@target@IV.5.7.2.2\endcsname{ega4.c04.eq.5-7-2-2}
\expandafter\def\csname sga@target@IV.5.7.3\endcsname{ega4.c04.sec.5-7-3}
\expandafter\def\csname sga@target@IV.5.7.3.i\endcsname{ega4.c04.sec.5-7-3-i}
\expandafter\def\csname sga@target@IV.5.7.3.ii\endcsname{ega4.c04.sec.5-7-3-ii}
\expandafter\def\csname sga@target@IV.5.7.3.iii\endcsname{ega4.c04.sec.5-7-3-iii}
\expandafter\def\csname sga@target@IV.5.7.3.iv\endcsname{ega4.c04.sec.5-7-3-iv}
\expandafter\def\csname sga@target@IV.5.7.3.v\endcsname{ega4.c04.sec.5-7-3-v}
\expandafter\def\csname sga@target@IV.5.7.3.vi\endcsname{ega4.c04.sec.5-7-3-vi}
\expandafter\def\csname sga@target@IV.5.7.4\endcsname{ega4.c04.prop.5-7-4}
\expandafter\def\csname sga@target@IV.5.7.4.1\endcsname{ega4.c04.eq.5-7-4-1}
\expandafter\def\csname sga@target@IV.5.7.4.2\endcsname{ega4.c04.eq.5-7-4-2}
\expandafter\def\csname sga@target@IV.5.7.4.i\endcsname{ega4.c04.sec.5-7-4-i}
\expandafter\def\csname sga@target@IV.5.7.4.ii\endcsname{ega4.c04.sec.5-7-4-ii}
\expandafter\def\csname sga@target@IV.5.7.5\endcsname{ega4.c04.prop.5-7-5}
\expandafter\def\csname sga@target@IV.5.7.6\endcsname{ega4.c04.cor.5-7-6}
\expandafter\def\csname sga@target@IV.5.7.7\endcsname{ega4.c04.cor.5-7-7}
\expandafter\def\csname sga@target@IV.5.7.8\endcsname{ega4.c04.sec.5-7-8}
\expandafter\def\csname sga@target@IV.5.7.9\endcsname{ega4.c04.prop.5-7-9}
\expandafter\def\csname sga@target@IV.5.7.9.1\endcsname{ega4.c04.eq.5-7-9-1}
\expandafter\def\csname sga@target@IV.5.7.9.2\endcsname{ega4.c04.eq.5-7-9-2}
\expandafter\def\csname sga@target@IV.5.8.1\endcsname{ega4.c04.statement.5-8-1}
\expandafter\def\csname sga@target@IV.5.8.2\endcsname{ega4.c04.def.5-8-2}
\expandafter\def\csname sga@target@IV.5.8.3\endcsname{ega4.c04.prop.5-8-3}
\expandafter\def\csname sga@target@IV.5.8.4\endcsname{ega4.c04.statement.5-8-4}
\expandafter\def\csname sga@target@IV.5.8.5\endcsname{ega4.c04.prop.5-8-5}
\expandafter\def\csname sga@target@IV.5.8.6\endcsname{ega4.c04.thm.5-8-6}
\expandafter\def\csname sga@target@IV.5.8.6.1\endcsname{ega4.c04.eq.5-8-6-1}
\expandafter\def\csname sga@target@IV.5.8.6.i\endcsname{ega4.c04.item.5-8-6-i}
\expandafter\def\csname sga@target@IV.5.8.6.ii\endcsname{ega4.c04.item.5-8-6-ii}
\expandafter\def\csname sga@target@IV.5.9.1\endcsname{ega4.c04.statement.5-9-1}
\expandafter\def\csname sga@target@IV.5.9.1.1\endcsname{ega4.c04.eq.5-9-1-1}
\expandafter\def\csname sga@target@IV.5.9.1.2\endcsname{ega4.c04.eq.5-9-1-2}
\expandafter\def\csname sga@target@IV.5.9.1.3\endcsname{ega4.c04.eq.5-9-1-3}
\expandafter\def\csname sga@target@IV.5.9.10\endcsname{ega4.c04.cor.5-9-10}
\expandafter\def\csname sga@target@IV.5.9.10.i\endcsname{ega4.c04.item.5-9-10-i}
\expandafter\def\csname sga@target@IV.5.9.10.ii\endcsname{ega4.c04.item.5-9-10-ii}
\expandafter\def\csname sga@target@IV.5.9.11\endcsname{ega4.c04.prop.5-9-11}
\expandafter\def\csname sga@target@IV.5.9.12\endcsname{ega4.c04.item.5-9-12}
\expandafter\def\csname sga@target@IV.5.9.12.i\endcsname{ega4.c04.item.5-9-12-i}
\expandafter\def\csname sga@target@IV.5.9.12.ii\endcsname{ega4.c04.item.5-9-12-ii}
\expandafter\def\csname sga@target@IV.5.9.12.iii\endcsname{ega4.c04.item.5-9-12-iii}
\expandafter\def\csname sga@target@IV.5.9.2\endcsname{ega4.c04.prop.5-9-2}
\expandafter\def\csname sga@target@IV.5.9.2.i\endcsname{ega4.c04.item.5-9-2-i}
\expandafter\def\csname sga@target@IV.5.9.2.ii\endcsname{ega4.c04.item.5-9-2-ii}
\expandafter\def\csname sga@target@IV.5.9.3\endcsname{ega4.c04.rem.5-9-3}
\expandafter\def\csname sga@target@IV.5.9.3.1\endcsname{ega4.c04.eq.5-9-3-1}
\expandafter\def\csname sga@target@IV.5.9.4\endcsname{ega4.c04.prop.5-9-4}
\expandafter\def\csname sga@target@IV.5.9.4.1\endcsname{ega4.c04.eq.5-9-4-1}
\expandafter\def\csname sga@target@IV.5.9.5\endcsname{ega4.c04.cor.5-9-5}
\expandafter\def\csname sga@target@IV.5.9.6\endcsname{ega4.c04.cor.5-9-6}
\expandafter\def\csname sga@target@IV.5.9.6.1\endcsname{ega4.c04.eq.5-9-6-1}
\expandafter\def\csname sga@target@IV.5.9.7\endcsname{ega4.c04.statement.5-9-7}
\expandafter\def\csname sga@target@IV.5.9.7.1\endcsname{ega4.c04.eq.5-9-7-1}
\expandafter\def\csname sga@target@IV.5.9.8\endcsname{ega4.c04.prop.5-9-8}
\expandafter\def\csname sga@target@IV.5.9.8.1\endcsname{ega4.c04.eq.5-9-8-1}
\expandafter\def\csname sga@target@IV.5.9.8.a\endcsname{ega4.c04.item.5-9-8-a}
\expandafter\def\csname sga@target@IV.5.9.8.a-prime\endcsname{ega4.c04.item.5-9-8-a-prime}
\expandafter\def\csname sga@target@IV.5.9.8.b\endcsname{ega4.c04.item.5-9-8-b}
\expandafter\def\csname sga@target@IV.5.9.8.b-prime\endcsname{ega4.c04.item.5-9-8-b-prime}
\expandafter\def\csname sga@target@IV.5.9.9\endcsname{ega4.c04.def.5-9-9}
\expandafter\def\csname sga@target@IV.6.1.1\endcsname{ega4.c04.prop.6-1-1}
\expandafter\def\csname sga@target@IV.6.1.1.1\endcsname{ega4.c04.eq.6-1-1-1}
\expandafter\def\csname sga@target@IV.6.1.2\endcsname{ega4.c04.cor.6-1-2}
\expandafter\def\csname sga@target@IV.6.1.2.1\endcsname{ega4.c04.eq.6-1-2-1}
\expandafter\def\csname sga@target@IV.6.1.3\endcsname{ega4.c04.cor.6-1-3}
\expandafter\def\csname sga@target@IV.6.1.4\endcsname{ega4.c04.cor.6-1-4}
\expandafter\def\csname sga@target@IV.6.1.4.1\endcsname{ega4.c04.eq.6-1-4-1}
\expandafter\def\csname sga@target@IV.6.1.5\endcsname{ega4.c04.prop.6-1-5}
\expandafter\def\csname sga@target@IV.6.10.1\endcsname{ega4.c04.statement.6-10-1}
\expandafter\def\csname sga@target@IV.6.10.2\endcsname{ega4.c04.prop.6-10-2}
\expandafter\def\csname sga@target@IV.6.10.3\endcsname{ega4.c04.prop.6-10-3}
\expandafter\def\csname sga@target@IV.6.10.3.1\endcsname{ega4.c04.eq.6-10-3-1}
\expandafter\def\csname sga@target@IV.6.10.3.2\endcsname{ega4.c04.eq.6-10-3-2}
\expandafter\def\csname sga@target@IV.6.10.3.3\endcsname{ega4.c04.eq.6-10-3-3}
\expandafter\def\csname sga@target@IV.6.10.4\endcsname{ega4.c04.cor.6-10-4}
\expandafter\def\csname sga@target@IV.6.10.5\endcsname{ega4.c04.prop.6-10-5}
\expandafter\def\csname sga@target@IV.6.10.5.1\endcsname{ega4.c04.eq.6-10-5-1}
\expandafter\def\csname sga@target@IV.6.10.5.2\endcsname{ega4.c04.eq.6-10-5-2}
\expandafter\def\csname sga@target@IV.6.10.5.3\endcsname{ega4.c04.eq.6-10-5-3}
\expandafter\def\csname sga@target@IV.6.10.5.4\endcsname{ega4.c04.eq.6-10-5-4}
\expandafter\def\csname sga@target@IV.6.10.6\endcsname{ega4.c04.prop.6-10-6}
\expandafter\def\csname sga@target@IV.6.10.6.1\endcsname{ega4.c04.eq.6-10-6-1}
\expandafter\def\csname sga@target@IV.6.10.6.2\endcsname{ega4.c04.eq.6-10-6-2}
\expandafter\def\csname sga@target@IV.6.10.6.3\endcsname{ega4.c04.eq.6-10-6-3}
\expandafter\def\csname sga@target@IV.6.11\endcsname{ega4.c04.subsec.6-11.key-iv-6-11}
\expandafter\def\csname sga@target@IV.6.11.1\endcsname{ega4.c04.lem.6-11-1}
\expandafter\def\csname sga@target@IV.6.11.2\endcsname{ega4.c04.prop.6-11-2}
\expandafter\def\csname sga@target@IV.6.11.2.1\endcsname{ega4.c04.eq.6-11-2-1}
\expandafter\def\csname sga@target@IV.6.11.2.2\endcsname{ega4.c04.eq.6-11-2-2}
\expandafter\def\csname sga@target@IV.6.11.3\endcsname{ega4.c04.cor.6-11-3}
\expandafter\def\csname sga@target@IV.6.11.4\endcsname{ega4.c04.statement.6-11-4}
\expandafter\def\csname sga@target@IV.6.11.5\endcsname{ega4.c04.prop.6-11-5}
\expandafter\def\csname sga@target@IV.6.11.5.1\endcsname{ega4.c04.eq.6-11-5-1}
\expandafter\def\csname sga@target@IV.6.11.6\endcsname{ega4.c04.prop.6-11-6}
\expandafter\def\csname sga@target@IV.6.11.6.1\endcsname{ega4.c04.lem.6-11-6-1}
\expandafter\def\csname sga@target@IV.6.11.7\endcsname{ega4.c04.cor.6-11-7}
\expandafter\def\csname sga@target@IV.6.11.8\endcsname{ega4.c04.prop.6-11-8}
\expandafter\def\csname sga@target@IV.6.11.8.1\endcsname{ega4.c04.eq.6-11-8-1}
\expandafter\def\csname sga@target@IV.6.11.9\endcsname{ega4.c04.statement.6-11-9}
\expandafter\def\csname sga@target@IV.6.12\endcsname{ega4.c04.subsec.6-12.key-iv-6-12}
\expandafter\def\csname sga@target@IV.6.12.1\endcsname{ega4.c04.statement.6-12-1}
\expandafter\def\csname sga@target@IV.6.12.10\endcsname{ega4.c04.prop.6-12-10}
\expandafter\def\csname sga@target@IV.6.12.2\endcsname{ega4.c04.prop.6-12-2}
\expandafter\def\csname sga@target@IV.6.12.3\endcsname{ega4.c04.cor.6-12-3}
\expandafter\def\csname sga@target@IV.6.12.4\endcsname{ega4.c04.thm.6-12-4}
\expandafter\def\csname sga@target@IV.6.12.4.1\endcsname{ega4.c04.lem.6-12-4-1}
\expandafter\def\csname sga@target@IV.6.12.5\endcsname{ega4.c04.cor.6-12-5}
\expandafter\def\csname sga@target@IV.6.12.6\endcsname{ega4.c04.cor.6-12-6}
\expandafter\def\csname sga@target@IV.6.12.7\endcsname{ega4.c04.thm.6-12-7}
\expandafter\def\csname sga@target@IV.6.12.8\endcsname{ega4.c04.cor.6-12-8}
\expandafter\def\csname sga@target@IV.6.12.9\endcsname{ega4.c04.prop.6-12-9}
\expandafter\def\csname sga@target@IV.6.13\endcsname{ega4.c04.subsec.6-13.key-iv-6-13}
\expandafter\def\csname sga@target@IV.6.13.1\endcsname{ega4.c04.statement.6-13-1}
\expandafter\def\csname sga@target@IV.6.13.2\endcsname{ega4.c04.prop.6-13-2}
\expandafter\def\csname sga@target@IV.6.13.3\endcsname{ega4.c04.cor.6-13-3}
\expandafter\def\csname sga@target@IV.6.13.4\endcsname{ega4.c04.prop.6-13-4}
\expandafter\def\csname sga@target@IV.6.13.5\endcsname{ega4.c04.cor.6-13-5}
\expandafter\def\csname sga@target@IV.6.13.6\endcsname{ega4.c04.prop.6-13-6}
\expandafter\def\csname sga@target@IV.6.13.7\endcsname{ega4.c04.prop.6-13-7}
\expandafter\def\csname sga@target@IV.6.13.7.1\endcsname{ega4.c04.lem.6-13-7-1}
\expandafter\def\csname sga@target@IV.6.14.1\endcsname{ega4.c04.prop.6-14-1}
\expandafter\def\csname sga@target@IV.6.14.1.1\endcsname{ega4.c04.lem.6-14-1-1}
\expandafter\def\csname sga@target@IV.6.14.1.2\endcsname{ega4.c04.lem.6-14-1-2}
\expandafter\def\csname sga@target@IV.6.14.1.2.env\endcsname{ega4.c04.statement.6-14-1-2-env}
\expandafter\def\csname sga@target@IV.6.14.1.3\endcsname{ega4.c04.eq.6-14-1-3}
\expandafter\def\csname sga@target@IV.6.14.1.4\endcsname{ega4.c04.lem.6-14-1-4}
\expandafter\def\csname sga@target@IV.6.14.2\endcsname{ega4.c04.cor.6-14-2}
\expandafter\def\csname sga@target@IV.6.14.3\endcsname{ega4.c04.cor.6-14-3}
\expandafter\def\csname sga@target@IV.6.14.4\endcsname{ega4.c04.prop.6-14-4}
\expandafter\def\csname sga@target@IV.6.14.4.1\endcsname{ega4.c04.lem.6-14-4-1}
\expandafter\def\csname sga@target@IV.6.14.5\endcsname{ega4.c04.cor.6-14-5}
\expandafter\def\csname sga@target@IV.6.15.1\endcsname{ega4.c04.statement.6-15-1}
\expandafter\def\csname sga@target@IV.6.15.10\endcsname{ega4.c04.prop.6-15-10}
\expandafter\def\csname sga@target@IV.6.15.11\endcsname{ega4.c04.item.6-15-11}
\expandafter\def\csname sga@target@IV.6.15.2\endcsname{ega4.c04.statement.6-15-2}
\expandafter\def\csname sga@target@IV.6.15.3\endcsname{ega4.c04.statement.6-15-3}
\expandafter\def\csname sga@target@IV.6.15.3.1\endcsname{ega4.c04.lem.6-15-3-1}
\expandafter\def\csname sga@target@IV.6.15.4\endcsname{ega4.c04.statement.6-15-4}
\expandafter\def\csname sga@target@IV.6.15.4.1\endcsname{ega4.c04.lem.6-15-4-1}
\expandafter\def\csname sga@target@IV.6.15.5\endcsname{ega4.c04.prop.6-15-5}
\expandafter\def\csname sga@target@IV.6.15.6\endcsname{ega4.c04.prop.6-15-6}
\expandafter\def\csname sga@target@IV.6.15.7\endcsname{ega4.c04.prop.6-15-7}
\expandafter\def\csname sga@target@IV.6.15.8\endcsname{ega4.c04.lem.6-15-8}
\expandafter\def\csname sga@target@IV.6.15.9\endcsname{ega4.c04.cor.6-15-9}
\expandafter\def\csname sga@target@IV.6.2.1\endcsname{ega4.c04.prop.6-2-1}
\expandafter\def\csname sga@target@IV.6.2.1.1\endcsname{ega4.c04.eq.6-2-1-1}
\expandafter\def\csname sga@target@IV.6.2.1.2\endcsname{ega4.c04.eq.6-2-1-2}
\expandafter\def\csname sga@target@IV.6.2.2\endcsname{ega4.c04.cor.6-2-2}
\expandafter\def\csname sga@target@IV.6.2.3\endcsname{ega4.c04.prop.6-2-3}
\expandafter\def\csname sga@target@IV.6.2.3.1\endcsname{ega4.c04.eq.6-2-3-1}
\expandafter\def\csname sga@target@IV.6.2.4\endcsname{ega4.c04.prop.6-2-4}
\expandafter\def\csname sga@target@IV.6.2.4.1\endcsname{ega4.c04.eq.6-2-4-1}
\expandafter\def\csname sga@target@IV.6.3.1\endcsname{ega4.c04.prop.6-3-1}
\expandafter\def\csname sga@target@IV.6.3.1.1\endcsname{ega4.c04.eq.6-3-1-1}
\expandafter\def\csname sga@target@IV.6.3.2\endcsname{ega4.c04.cor.6-3-2}
\expandafter\def\csname sga@target@IV.6.3.2.1\endcsname{ega4.c04.eq.6-3-2-1}
\expandafter\def\csname sga@target@IV.6.3.3\endcsname{ega4.c04.cor.6-3-3}
\expandafter\def\csname sga@target@IV.6.3.4\endcsname{ega4.c04.cor.6-3-4}
\expandafter\def\csname sga@target@IV.6.3.4.1\endcsname{ega4.c04.eq.6-3-4-1}
\expandafter\def\csname sga@target@IV.6.3.4.2\endcsname{ega4.c04.eq.6-3-4-2}
\expandafter\def\csname sga@target@IV.6.3.5\endcsname{ega4.c04.cor.6-3-5}
\expandafter\def\csname sga@target@IV.6.3.6\endcsname{ega4.c04.cor.6-3-6}
\expandafter\def\csname sga@target@IV.6.3.7\endcsname{ega4.c04.cor.6-3-7}
\expandafter\def\csname sga@target@IV.6.3.8\endcsname{ega4.c04.prop.6-3-8}
\expandafter\def\csname sga@target@IV.6.4.1\endcsname{ega4.c04.prop.6-4-1}
\expandafter\def\csname sga@target@IV.6.4.2\endcsname{ega4.c04.cor.6-4-2}
\expandafter\def\csname sga@target@IV.6.4.3\endcsname{ega4.c04.rem.6-4-3}
\expandafter\def\csname sga@target@IV.6.5.1\endcsname{ega4.c04.prop.6-5-1}
\expandafter\def\csname sga@target@IV.6.5.2\endcsname{ega4.c04.cor.6-5-2}
\expandafter\def\csname sga@target@IV.6.5.3\endcsname{ega4.c04.prop.6-5-3}
\expandafter\def\csname sga@target@IV.6.5.4\endcsname{ega4.c04.cor.6-5-4}
\expandafter\def\csname sga@target@IV.6.5.5\endcsname{ega4.c04.item.6-5-5}
\expandafter\def\csname sga@target@IV.6.6.1\endcsname{ega4.c04.prop.6-6-1}
\expandafter\def\csname sga@target@IV.6.6.2\endcsname{ega4.c04.rem.6-6-2}
\expandafter\def\csname sga@target@IV.6.7.1\endcsname{ega4.c04.prop.6-7-1}
\expandafter\def\csname sga@target@IV.6.7.1.1\endcsname{ega4.c04.lem.6-7-1-1}
\expandafter\def\csname sga@target@IV.6.7.2\endcsname{ega4.c04.cor.6-7-2}
\expandafter\def\csname sga@target@IV.6.7.3\endcsname{ega4.c04.cor.6-7-3}
\expandafter\def\csname sga@target@IV.6.7.4\endcsname{ega4.c04.prop.6-7-4}
\expandafter\def\csname sga@target@IV.6.7.4.1\endcsname{ega4.c04.lem.6-7-4-1}
\expandafter\def\csname sga@target@IV.6.7.5\endcsname{ega4.c04.item.6-7-5}
\expandafter\def\csname sga@target@IV.6.7.6\endcsname{ega4.c04.def.6-7-6}
\expandafter\def\csname sga@target@IV.6.7.7\endcsname{ega4.c04.prop.6-7-7}
\expandafter\def\csname sga@target@IV.6.7.8\endcsname{ega4.c04.cor.6-7-8}
\expandafter\def\csname sga@target@IV.6.8.1\endcsname{ega4.c04.def.6-8-1}
\expandafter\def\csname sga@target@IV.6.8.2\endcsname{ega4.c04.prop.6-8-2}
\expandafter\def\csname sga@target@IV.6.8.3\endcsname{ega4.c04.prop.6-8-3}
\expandafter\def\csname sga@target@IV.6.8.4\endcsname{ega4.c04.item.6-8-4}
\expandafter\def\csname sga@target@IV.6.8.5\endcsname{ega4.c04.prop.6-8-5}
\expandafter\def\csname sga@target@IV.6.8.6\endcsname{ega4.c04.thm.6-8-6}
\expandafter\def\csname sga@target@IV.6.8.7\endcsname{ega4.c04.cor.6-8-7}
\expandafter\def\csname sga@target@IV.6.8.8\endcsname{ega4.c04.rem.6-8-8}
\expandafter\def\csname sga@target@IV.6.9.1\endcsname{ega4.c04.thm.6-9-1}
\expandafter\def\csname sga@target@IV.6.9.2\endcsname{ega4.c04.lem.6-9-2}
\expandafter\def\csname sga@target@IV.6.9.3\endcsname{ega4.c04.cor.6-9-3}
\expandafter\def\csname sga@target@IV.7.1.1\endcsname{ega4.c04.def.7-1-1}
\expandafter\def\csname sga@target@IV.7.1.10\endcsname{ega4.c04.prop.7-1-10}
\expandafter\def\csname sga@target@IV.7.1.11\endcsname{ega4.c04.prop.7-1-11}
\expandafter\def\csname sga@target@IV.7.1.12\endcsname{ega4.c04.item.7-1-12}
\expandafter\def\csname sga@target@IV.7.1.13\endcsname{ega4.c04.cor.7-1-13}
\expandafter\def\csname sga@target@IV.7.1.2\endcsname{ega4.c04.prop.7-1-2}
\expandafter\def\csname sga@target@IV.7.1.3\endcsname{ega4.c04.prop.7-1-3}
\expandafter\def\csname sga@target@IV.7.1.3.1\endcsname{ega4.c04.eq.7-1-3-1}
\expandafter\def\csname sga@target@IV.7.1.3.2\endcsname{ega4.c04.eq.7-1-3-2}
\expandafter\def\csname sga@target@IV.7.1.3.3\endcsname{ega4.c04.eq.7-1-3-3}
\expandafter\def\csname sga@target@IV.7.1.3.4\endcsname{ega4.c04.eq.7-1-3-4}
\expandafter\def\csname sga@target@IV.7.1.3.5\endcsname{ega4.c04.eq.7-1-3-5}
\expandafter\def\csname sga@target@IV.7.1.3.6\endcsname{ega4.c04.eq.7-1-3-6}
\expandafter\def\csname sga@target@IV.7.1.3.7\endcsname{ega4.c04.eq.7-1-3-7}
\expandafter\def\csname sga@target@IV.7.1.4\endcsname{ega4.c04.cor.7-1-4}
\expandafter\def\csname sga@target@IV.7.1.5\endcsname{ega4.c04.cor.7-1-5}
\expandafter\def\csname sga@target@IV.7.1.6\endcsname{ega4.c04.rem.7-1-6}
\expandafter\def\csname sga@target@IV.7.1.7\endcsname{ega4.c04.lem.7-1-7}
\expandafter\def\csname sga@target@IV.7.1.8\endcsname{ega4.c04.thm.7-1-8}
\expandafter\def\csname sga@target@IV.7.1.9\endcsname{ega4.c04.def.7-1-9}
\expandafter\def\csname sga@target@IV.7.2.1\endcsname{ega4.c04.statement.7-2-1}
\expandafter\def\csname sga@target@IV.7.2.1.1\endcsname{ega4.c04.statement.7-2-1-1}
\expandafter\def\csname sga@target@IV.7.2.10\endcsname{ega4.c04.statement.7-2-10}
\expandafter\def\csname sga@target@IV.7.2.2\endcsname{ega4.c04.prop.7-2-2}
\expandafter\def\csname sga@target@IV.7.2.3\endcsname{ega4.c04.prop.7-2-3}
\expandafter\def\csname sga@target@IV.7.2.4\endcsname{ega4.c04.prop.7-2-4}
\expandafter\def\csname sga@target@IV.7.2.5\endcsname{ega4.c04.thm.7-2-5}
\expandafter\def\csname sga@target@IV.7.2.5.1\endcsname{ega4.c04.eq.7-2-5-1}
\expandafter\def\csname sga@target@IV.7.2.6\endcsname{ega4.c04.def.7-2-6}
\expandafter\def\csname sga@target@IV.7.2.7\endcsname{ega4.c04.cor.7-2-7}
\expandafter\def\csname sga@target@IV.7.2.8\endcsname{ega4.c04.cor.7-2-8}
\expandafter\def\csname sga@target@IV.7.2.9\endcsname{ega4.c04.cor.7-2-9}
\expandafter\def\csname sga@target@IV.7.3.1\endcsname{ega4.c04.statement.7-3-1}
\expandafter\def\csname sga@target@IV.7.3.10\endcsname{ega4.c04.statement.7-3-10}
\expandafter\def\csname sga@target@IV.7.3.11\endcsname{ega4.c04.rem.7-3-11}
\expandafter\def\csname sga@target@IV.7.3.12\endcsname{ega4.c04.prop.7-3-12}
\expandafter\def\csname sga@target@IV.7.3.12.1\endcsname{ega4.c04.eq.7-3-12-1}
\expandafter\def\csname sga@target@IV.7.3.13\endcsname{ega4.c04.statement.7-3-13}
\expandafter\def\csname sga@target@IV.7.3.14\endcsname{ega4.c04.prop.7-3-14}
\expandafter\def\csname sga@target@IV.7.3.15\endcsname{ega4.c04.prop.7-3-15}
\expandafter\def\csname sga@target@IV.7.3.16\endcsname{ega4.c04.prop.7-3-16}
\expandafter\def\csname sga@target@IV.7.3.16.1\endcsname{ega4.c04.eq.7-3-16-1}
\expandafter\def\csname sga@target@IV.7.3.16.2\endcsname{ega4.c04.lem.7-3-16-2}
\expandafter\def\csname sga@target@IV.7.3.17\endcsname{ega4.c04.prop.7-3-17}
\expandafter\def\csname sga@target@IV.7.3.18\endcsname{ega4.c04.prop.7-3-18}
\expandafter\def\csname sga@target@IV.7.3.19\endcsname{ega4.c04.statement.7-3-19}
\expandafter\def\csname sga@target@IV.7.3.2\endcsname{ega4.c04.lem.7-3-2}
\expandafter\def\csname sga@target@IV.7.3.3\endcsname{ega4.c04.cor.7-3-3}
\expandafter\def\csname sga@target@IV.7.3.4\endcsname{ega4.c04.statement.7-3-4}
\expandafter\def\csname sga@target@IV.7.3.5\endcsname{ega4.c04.item.7-3-5}
\expandafter\def\csname sga@target@IV.7.3.6\endcsname{ega4.c04.statement.7-3-6}
\expandafter\def\csname sga@target@IV.7.3.7\endcsname{ega4.c04.lem.7-3-7}
\expandafter\def\csname sga@target@IV.7.3.8\endcsname{ega4.c04.statement.7-3-8}
\expandafter\def\csname sga@target@IV.7.3.9\endcsname{ega4.c04.item.7-3-9}
\expandafter\def\csname sga@target@IV.7.4.1\endcsname{ega4.c04.statement.7-4-1}
\expandafter\def\csname sga@target@IV.7.4.2\endcsname{ega4.c04.lem.7-4-2}
\expandafter\def\csname sga@target@IV.7.4.3\endcsname{ega4.c04.cor.7-4-3}
\expandafter\def\csname sga@target@IV.7.4.4\endcsname{ega4.c04.thm.7-4-4}
\expandafter\def\csname sga@target@IV.7.4.4.1\endcsname{ega4.c04.lem.7-4-4-1}
\expandafter\def\csname sga@target@IV.7.4.4.2\endcsname{ega4.c04.lem.7-4-4-2}
\expandafter\def\csname sga@target@IV.7.4.5\endcsname{ega4.c04.cor.7-4-5}
\expandafter\def\csname sga@target@IV.7.4.6\endcsname{ega4.c04.prop.7-4-6}
\expandafter\def\csname sga@target@IV.7.4.7\endcsname{ega4.c04.cor.7-4-7}
\expandafter\def\csname sga@target@IV.7.4.8\endcsname{ega4.c04.rem.7-4-8}
\expandafter\def\csname sga@target@IV.7.5.0\endcsname{ega4.c04.statement.7-5-0}
\expandafter\def\csname sga@target@IV.7.5.1\endcsname{ega4.c04.thm.7-5-1}
\expandafter\def\csname sga@target@IV.7.5.1.1\endcsname{ega4.c04.lem.7-5-1-1}
\expandafter\def\csname sga@target@IV.7.5.2\endcsname{ega4.c04.cor.7-5-2}
\expandafter\def\csname sga@target@IV.7.5.3\endcsname{ega4.c04.item.7-5-3}
\expandafter\def\csname sga@target@IV.7.5.4\endcsname{ega4.c04.item.7-5-4}
\expandafter\def\csname sga@target@IV.7.5.5\endcsname{ega4.c04.lem.7-5-5}
\expandafter\def\csname sga@target@IV.7.5.6\endcsname{ega4.c04.prop.7-5-6}
\expandafter\def\csname sga@target@IV.7.5.7\endcsname{ega4.c04.cor.7-5-7}
\expandafter\def\csname sga@target@IV.7.5.7.1\endcsname{ega4.c04.lem.7-5-7-1}
\expandafter\def\csname sga@target@IV.7.5.8\endcsname{ega4.c04.rem.7-5-8}
\expandafter\def\csname sga@target@IV.7.6.1\endcsname{ega4.c04.prop.7-6-1}
\expandafter\def\csname sga@target@IV.7.6.2\endcsname{ega4.c04.cor.7-6-2}
\expandafter\def\csname sga@target@IV.7.6.3\endcsname{ega4.c04.cor.7-6-3}
\expandafter\def\csname sga@target@IV.7.6.4\endcsname{ega4.c04.thm.7-6-4}
\expandafter\def\csname sga@target@IV.7.6.4.1\endcsname{ega4.c04.lem.7-6-4-1}
\expandafter\def\csname sga@target@IV.7.6.5\endcsname{ega4.c04.cor.7-6-5}
\expandafter\def\csname sga@target@IV.7.6.6\endcsname{ega4.c04.cor.7-6-6}
\expandafter\def\csname sga@target@IV.7.6.7\endcsname{ega4.c04.item.7-6-7}
\expandafter\def\csname sga@target@IV.7.7.1\endcsname{ega4.c04.statement.7-7-1}
\expandafter\def\csname sga@target@IV.7.7.2\endcsname{ega4.c04.thm.7-7-2}
\expandafter\def\csname sga@target@IV.7.7.3\endcsname{ega4.c04.cor.7-7-3}
\expandafter\def\csname sga@target@IV.7.7.4\endcsname{ega4.c04.cor.7-7-4}
\expandafter\def\csname sga@target@IV.7.8.1\endcsname{ega4.c04.statement.7-8-1}
\expandafter\def\csname sga@target@IV.7.8.2\endcsname{ega4.c04.def.7-8-2}
\expandafter\def\csname sga@target@IV.7.8.3\endcsname{ega4.c04.statement.7-8-3}
\expandafter\def\csname sga@target@IV.7.8.3.1\endcsname{ega4.c04.statement.7-8-3-1}
\expandafter\def\csname sga@target@IV.7.8.4\endcsname{ega4.c04.item.7-8-4}
\expandafter\def\csname sga@target@IV.7.8.4.1\endcsname{ega4.c04.statement.7-8-4-1}
\expandafter\def\csname sga@target@IV.7.8.5\endcsname{ega4.c04.statement.7-8-5}
\expandafter\def\csname sga@target@IV.7.8.6\endcsname{ega4.c04.prop.7-8-6}
\expandafter\def\csname sga@target@IV.7.9.1\endcsname{ega4.c04.statement.7-9-1}
\expandafter\def\csname sga@target@IV.7.9.10\endcsname{ega4.c04.item.7-9-10}
\expandafter\def\csname sga@target@IV.7.9.2\endcsname{ega4.c04.statement.7-9-2}
\expandafter\def\csname sga@target@IV.7.9.3\endcsname{ega4.c04.prop.7-9-3}
\expandafter\def\csname sga@target@IV.7.9.3.1\endcsname{ega4.c04.lem.7-9-3-1}
\expandafter\def\csname sga@target@IV.7.9.4\endcsname{ega4.c04.cor.7-9-4}
\expandafter\def\csname sga@target@IV.7.9.5\endcsname{ega4.c04.prop.7-9-5}
\expandafter\def\csname sga@target@IV.7.9.6\endcsname{ega4.c04.rem.7-9-6}
\expandafter\def\csname sga@target@IV.7.9.7\endcsname{ega4.c04.statement.7-9-7}
\expandafter\def\csname sga@target@IV.7.9.8\endcsname{ega4.c04.prop.7-9-8}
\expandafter\def\csname sga@target@IV.7.9.8.1\endcsname{ega4.c04.statement.7-9-8-1}
\expandafter\def\csname sga@target@IV.7.9.9\endcsname{ega4.c04.cor.7-9-9}
\expandafter\def\csname sga@target@IV.8.1.1\endcsname{ega4.c04.statement.8-1-1}
\expandafter\def\csname sga@target@IV.8.1.2\endcsname{ega4.c04.statement.8-1-2}
\expandafter\def\csname sga@target@IV.8.1.3\endcsname{ega4.c04.statement.8-1-3}
\expandafter\def\csname sga@target@IV.8.10.1\endcsname{ega4.c04.prop.8-10-1}
\expandafter\def\csname sga@target@IV.8.10.2\endcsname{ega4.c04.cor.8-10-2}
\expandafter\def\csname sga@target@IV.8.10.3\endcsname{ega4.c04.prop.8-10-3}
\expandafter\def\csname sga@target@IV.8.10.3.1\endcsname{ega4.c04.rem.8-10-3-1}
\expandafter\def\csname sga@target@IV.8.10.4\endcsname{ega4.c04.prop.8-10-4}
\expandafter\def\csname sga@target@IV.8.10.5\endcsname{ega4.c04.thm.8-10-5}
\expandafter\def\csname sga@target@IV.8.10.5.1\endcsname{ega4.c04.prop.8-10-5-1}
\expandafter\def\csname sga@target@IV.8.10.5.2\endcsname{ega4.c04.lem.8-10-5-2}
\expandafter\def\csname sga@target@IV.8.11\endcsname{ega4.c04.subsec.8-11.key-iv-8-11}
\expandafter\def\csname sga@target@IV.8.11.1\endcsname{ega4.c04.thm.8-11-1}
\expandafter\def\csname sga@target@IV.8.11.2\endcsname{ega4.c04.thm.8-11-2}
\expandafter\def\csname sga@target@IV.8.11.3\endcsname{ega4.c04.rem.8-11-3}
\expandafter\def\csname sga@target@IV.8.11.4\endcsname{ega4.c04.statement.8-11-4}
\expandafter\def\csname sga@target@IV.8.11.5\endcsname{ega4.c04.prop.8-11-5}
\expandafter\def\csname sga@target@IV.8.11.5.1\endcsname{ega4.c04.rem.8-11-5-1}
\expandafter\def\csname sga@target@IV.8.11.6\endcsname{ega4.c04.prop.8-11-6}
\expandafter\def\csname sga@target@IV.8.12\endcsname{ega4.c04.subsec.8-12.key-iv-8-12}
\expandafter\def\csname sga@target@IV.8.12.1\endcsname{ega4.c04.lem.8-12-1}
\expandafter\def\csname sga@target@IV.8.12.10\endcsname{ega4.c04.cor.8-12-10}
\expandafter\def\csname sga@target@IV.8.12.10.1\endcsname{ega4.c04.lem.8-12-10-1}
\expandafter\def\csname sga@target@IV.8.12.11\endcsname{ega4.c04.cor.8-12-11}
\expandafter\def\csname sga@target@IV.8.12.12\endcsname{ega4.c04.cor.8-12-12}
\expandafter\def\csname sga@target@IV.8.12.2\endcsname{ega4.c04.lem.8-12-2}
\expandafter\def\csname sga@target@IV.8.12.3\endcsname{ega4.c04.prop.8-12-3}
\expandafter\def\csname sga@target@IV.8.12.4\endcsname{ega4.c04.cor.8-12-4}
\expandafter\def\csname sga@target@IV.8.12.5\endcsname{ega4.c04.statement.8-12-5}
\expandafter\def\csname sga@target@IV.8.12.6\endcsname{ega4.c04.thm.8-12-6}
\expandafter\def\csname sga@target@IV.8.12.6.1\endcsname{ega4.c04.eq.8-12-6-1}
\expandafter\def\csname sga@target@IV.8.12.7\endcsname{ega4.c04.rem.8-12-7}
\expandafter\def\csname sga@target@IV.8.12.8\endcsname{ega4.c04.cor.8-12-8}
\expandafter\def\csname sga@target@IV.8.12.9\endcsname{ega4.c04.cor.8-12-9}
\expandafter\def\csname sga@target@IV.8.13.1\endcsname{ega4.c04.prop.8-13-1}
\expandafter\def\csname sga@target@IV.8.13.1.1\endcsname{ega4.c04.eq.8-13-1-1}
\expandafter\def\csname sga@target@IV.8.13.2\endcsname{ega4.c04.cor.8-13-2}
\expandafter\def\csname sga@target@IV.8.13.2.1\endcsname{ega4.c04.eq.8-13-2-1}
\expandafter\def\csname sga@target@IV.8.13.3\endcsname{ega4.c04.statement.8-13-3}
\expandafter\def\csname sga@target@IV.8.13.4\endcsname{ega4.c04.statement.8-13-4}
\expandafter\def\csname sga@target@IV.8.13.4.1\endcsname{ega4.c04.eq.8-13-4-1}
\expandafter\def\csname sga@target@IV.8.13.4.2\endcsname{ega4.c04.eq.8-13-4-2}
\expandafter\def\csname sga@target@IV.8.13.5\endcsname{ega4.c04.prop.8-13-5}
\expandafter\def\csname sga@target@IV.8.13.6\endcsname{ega4.c04.rem.8-13-6}
\expandafter\def\csname sga@target@IV.8.14.1\endcsname{ega4.c04.statement.8-14-1}
\expandafter\def\csname sga@target@IV.8.14.2\endcsname{ega4.c04.prop.8-14-2}
\expandafter\def\csname sga@target@IV.8.14.2.1\endcsname{ega4.c04.eq.8-14-2-1}
\expandafter\def\csname sga@target@IV.8.14.2.2\endcsname{ega4.c04.cor.8-14-2-2}
\expandafter\def\csname sga@target@IV.8.14.2.3\endcsname{ega4.c04.eq.8-14-2-3}
\expandafter\def\csname sga@target@IV.8.14.2.4\endcsname{ega4.c04.eq.8-14-2-4}
\expandafter\def\csname sga@target@IV.8.14.3\endcsname{ega4.c04.rem.8-14-3}
\expandafter\def\csname sga@target@IV.8.2.1\endcsname{ega4.c04.statement.8-2-1}
\expandafter\def\csname sga@target@IV.8.2.1.1\endcsname{ega4.c04.eq.8-2-1-1}
\expandafter\def\csname sga@target@IV.8.2.10\endcsname{ega4.c04.cor.8-2-10}
\expandafter\def\csname sga@target@IV.8.2.11\endcsname{ega4.c04.cor.8-2-11}
\expandafter\def\csname sga@target@IV.8.2.12\endcsname{ega4.c04.cor.8-2-12}
\expandafter\def\csname sga@target@IV.8.2.12.1\endcsname{ega4.c04.eq.8-2-12-1}
\expandafter\def\csname sga@target@IV.8.2.13\endcsname{ega4.c04.prop.8-2-13}
\expandafter\def\csname sga@target@IV.8.2.14\endcsname{ega4.c04.rem.8-2-14}
\expandafter\def\csname sga@target@IV.8.2.2\endcsname{ega4.c04.statement.8-2-2}
\expandafter\def\csname sga@target@IV.8.2.3\endcsname{ega4.c04.prop.8-2-3}
\expandafter\def\csname sga@target@IV.8.2.3.1\endcsname{ega4.c04.eq.8-2-3-1}
\expandafter\def\csname sga@target@IV.8.2.4\endcsname{ega4.c04.lem.8-2-4}
\expandafter\def\csname sga@target@IV.8.2.5\endcsname{ega4.c04.prop.8-2-5}
\expandafter\def\csname sga@target@IV.8.2.6\endcsname{ega4.c04.lem.8-2-6}
\expandafter\def\csname sga@target@IV.8.2.7\endcsname{ega4.c04.rem.8-2-7}
\expandafter\def\csname sga@target@IV.8.2.8\endcsname{ega4.c04.statement.8-2-8}
\expandafter\def\csname sga@target@IV.8.2.9\endcsname{ega4.c04.prop.8-2-9}
\expandafter\def\csname sga@target@IV.8.3.1\endcsname{ega4.c04.statement.8-3-1}
\expandafter\def\csname sga@target@IV.8.3.10\endcsname{ega4.c04.statement.8-3-10}
\expandafter\def\csname sga@target@IV.8.3.10.1\endcsname{ega4.c04.eq.8-3-10-1}
\expandafter\def\csname sga@target@IV.8.3.11\endcsname{ega4.c04.thm.8-3-11}
\expandafter\def\csname sga@target@IV.8.3.12\endcsname{ega4.c04.cor.8-3-12}
\expandafter\def\csname sga@target@IV.8.3.13\endcsname{ega4.c04.prop.8-3-13}
\expandafter\def\csname sga@target@IV.8.3.2\endcsname{ega4.c04.thm.8-3-2}
\expandafter\def\csname sga@target@IV.8.3.2.1\endcsname{ega4.c04.eq.8-3-2-1}
\expandafter\def\csname sga@target@IV.8.3.3\endcsname{ega4.c04.cor.8-3-3}
\expandafter\def\csname sga@target@IV.8.3.4\endcsname{ega4.c04.cor.8-3-4}
\expandafter\def\csname sga@target@IV.8.3.5\endcsname{ega4.c04.cor.8-3-5}
\expandafter\def\csname sga@target@IV.8.3.6\endcsname{ega4.c04.cor.8-3-6}
\expandafter\def\csname sga@target@IV.8.3.7\endcsname{ega4.c04.cor.8-3-7}
\expandafter\def\csname sga@target@IV.8.3.7.1\endcsname{ega4.c04.eq.8-3-7-1}
\expandafter\def\csname sga@target@IV.8.3.8\endcsname{ega4.c04.prop.8-3-8}
\expandafter\def\csname sga@target@IV.8.3.9\endcsname{ega4.c04.statement.8-3-9}
\expandafter\def\csname sga@target@IV.8.3.9.1\endcsname{ega4.c04.eq.8-3-9-1}
\expandafter\def\csname sga@target@IV.8.3.9.2\endcsname{ega4.c04.eq.8-3-9-2}
\expandafter\def\csname sga@target@IV.8.3.9.3\endcsname{ega4.c04.eq.8-3-9-3}
\expandafter\def\csname sga@target@IV.8.3.9.4\endcsname{ega4.c04.eq.8-3-9-4}
\expandafter\def\csname sga@target@IV.8.3.9.5\endcsname{ega4.c04.eq.8-3-9-5}
\expandafter\def\csname sga@target@IV.8.4.1\endcsname{ega4.c04.prop.8-4-1}
\expandafter\def\csname sga@target@IV.8.4.2\endcsname{ega4.c04.prop.8-4-2}
\expandafter\def\csname sga@target@IV.8.4.2.1\endcsname{ega4.c04.lem.8-4-2-1}
\expandafter\def\csname sga@target@IV.8.4.3\endcsname{ega4.c04.cor.8-4-3}
\expandafter\def\csname sga@target@IV.8.4.4\endcsname{ega4.c04.prop.8-4-4}
\expandafter\def\csname sga@target@IV.8.4.5\endcsname{ega4.c04.cor.8-4-5}
\expandafter\def\csname sga@target@IV.8.4.6\endcsname{ega4.c04.rem.8-4-6}
\expandafter\def\csname sga@target@IV.8.5.1\endcsname{ega4.c04.statement.8-5-1}
\expandafter\def\csname sga@target@IV.8.5.1.1\endcsname{ega4.c04.eq.8-5-1-1}
\expandafter\def\csname sga@target@IV.8.5.1.2\endcsname{ega4.c04.eq.8-5-1-2}
\expandafter\def\csname sga@target@IV.8.5.1.3\endcsname{ega4.c04.eq.8-5-1-3}
\expandafter\def\csname sga@target@IV.8.5.1.4\endcsname{ega4.c04.eq.8-5-1-4}
\expandafter\def\csname sga@target@IV.8.5.10\endcsname{ega4.c04.statement.8-5-10}
\expandafter\def\csname sga@target@IV.8.5.10.1\endcsname{ega4.c04.eq.8-5-10-1}
\expandafter\def\csname sga@target@IV.8.5.10.2\endcsname{ega4.c04.eq.8-5-10-2}
\expandafter\def\csname sga@target@IV.8.5.11\endcsname{ega4.c04.prop.8-5-11}
\expandafter\def\csname sga@target@IV.8.5.2\endcsname{ega4.c04.thm.8-5-2}
\expandafter\def\csname sga@target@IV.8.5.2.1\endcsname{ega4.c04.lem.8-5-2-1}
\expandafter\def\csname sga@target@IV.8.5.2.2\endcsname{ega4.c04.eq.8-5-2-2}
\expandafter\def\csname sga@target@IV.8.5.2.3\endcsname{ega4.c04.eq.8-5-2-3}
\expandafter\def\csname sga@target@IV.8.5.2.4\endcsname{ega4.c04.cor.8-5-2-4}
\expandafter\def\csname sga@target@IV.8.5.2.5\endcsname{ega4.c04.cor.8-5-2-5}
\expandafter\def\csname sga@target@IV.8.5.3\endcsname{ega4.c04.statement.8-5-3}
\expandafter\def\csname sga@target@IV.8.5.4\endcsname{ega4.c04.cor.8-5-4}
\expandafter\def\csname sga@target@IV.8.5.5\endcsname{ega4.c04.prop.8-5-5}
\expandafter\def\csname sga@target@IV.8.5.6\endcsname{ega4.c04.prop.8-5-6}
\expandafter\def\csname sga@target@IV.8.5.7\endcsname{ega4.c04.cor.8-5-7}
\expandafter\def\csname sga@target@IV.8.5.8\endcsname{ega4.c04.cor.8-5-8}
\expandafter\def\csname sga@target@IV.8.5.9\endcsname{ega4.c04.lem.8-5-9}
\expandafter\def\csname sga@target@IV.8.6.1\endcsname{ega4.c04.statement.8-6-1}
\expandafter\def\csname sga@target@IV.8.6.2\endcsname{ega4.c04.statement.8-6-2}
\expandafter\def\csname sga@target@IV.8.6.2.1\endcsname{ega4.c04.eq.8-6-2-1}
\expandafter\def\csname sga@target@IV.8.6.2.2\endcsname{ega4.c04.eq.8-6-2-2}
\expandafter\def\csname sga@target@IV.8.6.2.3\endcsname{ega4.c04.eq.8-6-2-3}
\expandafter\def\csname sga@target@IV.8.6.3\endcsname{ega4.c04.prop.8-6-3}
\expandafter\def\csname sga@target@IV.8.6.3.1\endcsname{ega4.c04.cor.8-6-3-1}
\expandafter\def\csname sga@target@IV.8.6.4\endcsname{ega4.c04.cor.8-6-4}
\expandafter\def\csname sga@target@IV.8.7.1\endcsname{ega4.c04.prop.8-7-1}
\expandafter\def\csname sga@target@IV.8.7.2\endcsname{ega4.c04.prop.8-7-2}
\expandafter\def\csname sga@target@IV.8.7.3\endcsname{ega4.c04.cor.8-7-3}
\expandafter\def\csname sga@target@IV.8.7.4\endcsname{ega4.c04.rem.8-7-4}
\expandafter\def\csname sga@target@IV.8.8.1\endcsname{ega4.c04.statement.8-8-1}
\expandafter\def\csname sga@target@IV.8.8.1.1\endcsname{ega4.c04.eq.8-8-1-1}
\expandafter\def\csname sga@target@IV.8.8.2\endcsname{ega4.c04.thm.8-8-2}
\expandafter\def\csname sga@target@IV.8.8.2.1\endcsname{ega4.c04.lem.8-8-2-1}
\expandafter\def\csname sga@target@IV.8.8.2.2\endcsname{ega4.c04.lem.8-8-2-2}
\expandafter\def\csname sga@target@IV.8.8.2.3\endcsname{ega4.c04.lem.8-8-2-3}
\expandafter\def\csname sga@target@IV.8.8.2.4\endcsname{ega4.c04.cor.8-8-2-4}
\expandafter\def\csname sga@target@IV.8.8.2.5\endcsname{ega4.c04.cor.8-8-2-5}
\expandafter\def\csname sga@target@IV.8.8.3\endcsname{ega4.c04.statement.8-8-3}
\expandafter\def\csname sga@target@IV.8.9.1\endcsname{ega4.c04.prop.8-9-1}
\expandafter\def\csname sga@target@IV.8.9.2\endcsname{ega4.c04.statement.8-9-2}
\expandafter\def\csname sga@target@IV.8.9.3\endcsname{ega4.c04.prop.8-9-3}
\expandafter\def\csname sga@target@IV.8.9.4\endcsname{ega4.c04.prop.8-9-4}
\expandafter\def\csname sga@target@IV.8.9.4.1\endcsname{ega4.c04.lem.8-9-4-1}
\expandafter\def\csname sga@target@IV.8.9.5\endcsname{ega4.c04.cor.8-9-5}
\expandafter\def\csname sga@target@IV.9.1.1\endcsname{ega4.c04.prop.9-1-1}
\expandafter\def\csname sga@target@IV.9.1.2\endcsname{ega4.c04.cor.9-1-2}
\expandafter\def\csname sga@target@IV.9.1.3\endcsname{ega4.c04.rem.9-1-3}
\expandafter\def\csname sga@target@IV.9.1.4\endcsname{ega4.c04.prop.9-1-4}
\expandafter\def\csname sga@target@IV.9.1.5\endcsname{ega4.c04.cor.9-1-5}
\expandafter\def\csname sga@target@IV.9.1.6\endcsname{ega4.c04.prop.9-1-6}
\expandafter\def\csname sga@target@IV.9.2\endcsname{ega4.c04.subsec.9-2.key-iv-9-2}
\expandafter\def\csname sga@target@IV.9.2.1\endcsname{ega4.c04.def.9-2-1}
\expandafter\def\csname sga@target@IV.9.2.2\endcsname{ega4.c04.item.9-2-2}
\expandafter\def\csname sga@target@IV.9.2.3\endcsname{ega4.c04.prop.9-2-3}
\expandafter\def\csname sga@target@IV.9.2.4\endcsname{ega4.c04.cor.9-2-4}
\expandafter\def\csname sga@target@IV.9.2.5\endcsname{ega4.c04.rem.9-2-5}
\expandafter\def\csname sga@target@IV.9.2.6\endcsname{ega4.c04.prop.9-2-6}
\expandafter\def\csname sga@target@IV.9.2.6.1\endcsname{ega4.c04.prop.9-2-6-1}
\expandafter\def\csname sga@target@IV.9.2.6.2\endcsname{ega4.c04.cor.9-2-6-2}
\expandafter\def\csname sga@target@IV.9.3.1\endcsname{ega4.c04.prop.9-3-1}
\expandafter\def\csname sga@target@IV.9.3.2\endcsname{ega4.c04.prop.9-3-2}
\expandafter\def\csname sga@target@IV.9.3.3\endcsname{ega4.c04.prop.9-3-3}
\expandafter\def\csname sga@target@IV.9.3.4\endcsname{ega4.c04.rem.9-3-4}
\expandafter\def\csname sga@target@IV.9.3.5\endcsname{ega4.c04.prop.9-3-5}
\expandafter\def\csname sga@target@IV.9.3.6\endcsname{ega4.c04.statement.9-3-6}
\expandafter\def\csname sga@target@IV.9.4.1\endcsname{ega4.c04.statement.9-4-1}
\expandafter\def\csname sga@target@IV.9.4.2\endcsname{ega4.c04.prop.9-4-2}
\expandafter\def\csname sga@target@IV.9.4.3\endcsname{ega4.c04.cor.9-4-3}
\expandafter\def\csname sga@target@IV.9.4.3.1\endcsname{ega4.c04.eq.9-4-3-1}
\expandafter\def\csname sga@target@IV.9.4.4\endcsname{ega4.c04.prop.9-4-4}
\expandafter\def\csname sga@target@IV.9.4.5\endcsname{ega4.c04.cor.9-4-5}
\expandafter\def\csname sga@target@IV.9.4.6\endcsname{ega4.c04.cor.9-4-6}
\expandafter\def\csname sga@target@IV.9.4.7\endcsname{ega4.c04.prop.9-4-7}
\expandafter\def\csname sga@target@IV.9.4.7.1\endcsname{ega4.c04.lem.9-4-7-1}
\expandafter\def\csname sga@target@IV.9.4.7.2\endcsname{ega4.c04.rem.9-4-7-2}
\expandafter\def\csname sga@target@IV.9.4.8\endcsname{ega4.c04.prop.9-4-8}
\expandafter\def\csname sga@target@IV.9.4.9\endcsname{ega4.c04.rem.9-4-9}
\expandafter\def\csname sga@target@IV.9.5.1\endcsname{ega4.c04.prop.9-5-1}
\expandafter\def\csname sga@target@IV.9.5.2\endcsname{ega4.c04.cor.9-5-2}
\expandafter\def\csname sga@target@IV.9.5.3\endcsname{ega4.c04.prop.9-5-3}
\expandafter\def\csname sga@target@IV.9.5.3.case2\endcsname{ega4.c04.item.9-5-3-case2}
\expandafter\def\csname sga@target@IV.9.5.4\endcsname{ega4.c04.cor.9-5-4}
\expandafter\def\csname sga@target@IV.9.5.5\endcsname{ega4.c04.prop.9-5-5}
\expandafter\def\csname sga@target@IV.9.5.6\endcsname{ega4.c04.cor.9-5-6}
\expandafter\def\csname sga@target@IV.9.5.7\endcsname{ega4.c04.rem.9-5-7}
\expandafter\def\csname sga@target@IV.9.6.1\endcsname{ega4.c04.prop.9-6-1}
\expandafter\def\csname sga@target@IV.9.6.2\endcsname{ega4.c04.prop.9-6-2}
\expandafter\def\csname sga@target@IV.9.6.3\endcsname{ega4.c04.prop.9-6-3}
\expandafter\def\csname sga@target@IV.9.6.3.1\endcsname{ega4.c04.eq.9-6-3-1}
\expandafter\def\csname sga@target@IV.9.6.4\endcsname{ega4.c04.cor.9-6-4}
\expandafter\def\csname sga@target@IV.9.6.5\endcsname{ega4.c04.cor.9-6-5}
\expandafter\def\csname sga@target@IV.9.6.6\endcsname{ega4.c04.sec.9-6-6}
\expandafter\def\csname sga@target@IV.9.6.7\endcsname{ega4.c04.sec.9-6-7}
\expandafter\def\csname sga@target@IV.9.7\endcsname{ega4.c04.subsec.9-7.key-iv-9-7}
\expandafter\def\csname sga@target@IV.9.7.1\endcsname{ega4.c04.lem.9-7-1}
\expandafter\def\csname sga@target@IV.9.7.10\endcsname{ega4.c04.cor.9-7-10}
\expandafter\def\csname sga@target@IV.9.7.11\endcsname{ega4.c04.rem.9-7-11}
\expandafter\def\csname sga@target@IV.9.7.12\endcsname{ega4.c04.prop.9-7-12}
\expandafter\def\csname sga@target@IV.9.7.2\endcsname{ega4.c04.lem.9-7-2}
\expandafter\def\csname sga@target@IV.9.7.3\endcsname{ega4.c04.subsec.9-7-3}
\expandafter\def\csname sga@target@IV.9.7.4\endcsname{ega4.c04.lem.9-7-4}
\expandafter\def\csname sga@target@IV.9.7.5\endcsname{ega4.c04.lem.9-7-5}
\expandafter\def\csname sga@target@IV.9.7.5.1\endcsname{ega4.c04.lem.9-7-5-1}
\expandafter\def\csname sga@target@IV.9.7.6\endcsname{ega4.c04.prop.9-7-6}
\expandafter\def\csname sga@target@IV.9.7.7\endcsname{ega4.c04.thm.9-7-7}
\expandafter\def\csname sga@target@IV.9.7.8\endcsname{ega4.c04.prop.9-7-8}
\expandafter\def\csname sga@target@IV.9.7.9\endcsname{ega4.c04.cor.9-7-9}
\expandafter\def\csname sga@target@IV.9.8\endcsname{ega4.c04.subsec.9-8.key-iv-9-8}
\expandafter\def\csname sga@target@IV.9.8.1\endcsname{ega4.c04.prop.9-8-1}
\expandafter\def\csname sga@target@IV.9.8.1-proof\endcsname{ega4.c04.subsec.9-8-1-proof}
\expandafter\def\csname sga@target@IV.9.8.2\endcsname{ega4.c04.subsec.9-8-2}
\expandafter\def\csname sga@target@IV.9.8.3\endcsname{ega4.c04.prop.9-8-3}
\expandafter\def\csname sga@target@IV.9.8.3-proof\endcsname{ega4.c04.subsec.9-8-3-proof}
\expandafter\def\csname sga@target@IV.9.8.4\endcsname{ega4.c04.cor.9-8-4}
\expandafter\def\csname sga@target@IV.9.8.4-proof\endcsname{ega4.c04.item.9-8-4-proof}
\expandafter\def\csname sga@target@IV.9.8.5\endcsname{ega4.c04.prop.9-8-5}
\expandafter\def\csname sga@target@IV.9.8.5-proof\endcsname{ega4.c04.item.9-8-5-proof}
\expandafter\def\csname sga@target@IV.9.8.6\endcsname{ega4.c04.prop.9-8-6}
\expandafter\def\csname sga@target@IV.9.8.6-proof\endcsname{ega4.c04.item.9-8-6-proof}
\expandafter\def\csname sga@target@IV.9.8.7\endcsname{ega4.c04.cor.9-8-7}
\expandafter\def\csname sga@target@IV.9.8.7-proof\endcsname{ega4.c04.item.9-8-7-proof}
\expandafter\def\csname sga@target@IV.9.8.8\endcsname{ega4.c04.prop.9-8-8}
\expandafter\def\csname sga@target@IV.9.8.8-proof\endcsname{ega4.c04.item.9-8-8-proof}
\expandafter\def\csname sga@target@IV.9.8.9\endcsname{ega4.c04.rem.9-8-9}
\expandafter\def\csname sga@target@IV.9.8.9.1\endcsname{ega4.c04.item.9-8-9-1}
\expandafter\def\csname sga@target@IV.9.9\endcsname{ega4.c04.subsec.9-9.key-iv-9-9}
\expandafter\def\csname sga@target@IV.9.9.1\endcsname{ega4.c04.prop.9-9-1}
\expandafter\def\csname sga@target@IV.9.9.1-proof\endcsname{ega4.c04.item.9-9-1-proof}
\expandafter\def\csname sga@target@IV.9.9.1.1\endcsname{ega4.c04.eq.9-9-1-1}
\expandafter\def\csname sga@target@IV.9.9.2\endcsname{ega4.c04.prop.9-9-2}
\expandafter\def\csname sga@target@IV.9.9.2-proof\endcsname{ega4.c04.item.9-9-2-proof}
\expandafter\def\csname sga@target@IV.9.9.2.1\endcsname{ega4.c04.eq.9-9-2-1}
\expandafter\def\csname sga@target@IV.9.9.2.2\endcsname{ega4.c04.eq.9-9-2-2}
\expandafter\def\csname sga@target@IV.9.9.3\endcsname{ega4.c04.cor.9-9-3}
\expandafter\def\csname sga@target@IV.9.9.3-proof\endcsname{ega4.c04.statement.9-9-3-proof}
\expandafter\def\csname sga@target@IV.9.9.4\endcsname{ega4.c04.prop.9-9-4}
\expandafter\def\csname sga@target@IV.9.9.4-proof\endcsname{ega4.c04.item.9-9-4-proof}
\expandafter\def\csname sga@target@IV.9.9.5\endcsname{ega4.c04.cor.9-9-5}
\expandafter\def\csname sga@target@IV.9.9.5-proof\endcsname{ega4.c04.item.9-9-5-proof}
\expandafter\def\csname sga@target@IV.9.9.6\endcsname{ega4.c04.prop.9-9-6}
\expandafter\def\csname sga@target@IV.9.9.6-proof\endcsname{ega4.c04.item.9-9-6-proof}
\expandafter\def\csname sga@target@section:IV.1\endcsname{ega4.c04.sec.1}
\expandafter\def\csname sga@target@section:IV.10\endcsname{ega4.c04.sec.10}
\expandafter\def\csname sga@target@section:IV.11\endcsname{ega4.c04.sec.11}
\expandafter\def\csname sga@target@section:IV.12\endcsname{ega4.c04.sec.12}
\expandafter\def\csname sga@target@section:IV.13\endcsname{ega4.c04.sec.13}
\expandafter\def\csname sga@target@section:IV.14\endcsname{ega4.c04.sec.14}
\expandafter\def\csname sga@target@section:IV.15\endcsname{ega4.c04.sec.15}
\expandafter\def\csname sga@target@section:IV.16\endcsname{ega4.c04.sec.16}
\expandafter\def\csname sga@target@section:IV.17\endcsname{ega4.c04.sec.17}
\expandafter\def\csname sga@target@section:IV.18\endcsname{ega4.c04.sec.18}
\expandafter\def\csname sga@target@section:IV.19\endcsname{ega4.c04.sec.19}
\expandafter\def\csname sga@target@section:IV.2\endcsname{ega4.c04.sec.2}
\expandafter\def\csname sga@target@section:IV.21\endcsname{ega4.c04.sec.21}
\expandafter\def\csname sga@target@section:IV.3\endcsname{ega4.c04.sec.3}
\expandafter\def\csname sga@target@section:IV.4\endcsname{ega4.c04.sec.4}
\expandafter\def\csname sga@target@section:IV.5\endcsname{ega4.c04.sec.5}
\expandafter\def\csname sga@target@section:IV.6\endcsname{ega4.c04.sec.6}
\expandafter\def\csname sga@target@section:IV.7\endcsname{ega4.c04.sec.7}
\expandafter\def\csname sga@target@section:IV.8\endcsname{ega4.c04.sec.8}
\expandafter\def\csname sga@target@section:IV.9\endcsname{ega4.c04.sec.9}
\expandafter\def\csname sga@target@subsection:IV.1.1\endcsname{ega4.c04.subsec.1-1}
\expandafter\def\csname sga@target@subsection:IV.1.10\endcsname{ega4.c04.subsec.1-10}
\expandafter\def\csname sga@target@subsection:IV.1.2\endcsname{ega4.c04.subsec.1-2}
\expandafter\def\csname sga@target@subsection:IV.1.3\endcsname{ega4.c04.subsec.1-3}
\expandafter\def\csname sga@target@subsection:IV.1.4\endcsname{ega4.c04.subsec.1-4}
\expandafter\def\csname sga@target@subsection:IV.1.5\endcsname{ega4.c04.subsec.1-5}
\expandafter\def\csname sga@target@subsection:IV.1.6\endcsname{ega4.c04.subsec.1-6}
\expandafter\def\csname sga@target@subsection:IV.1.7\endcsname{ega4.c04.subsec.1-7}
\expandafter\def\csname sga@target@subsection:IV.1.8\endcsname{ega4.c04.subsec.1-8}
\expandafter\def\csname sga@target@subsection:IV.1.9\endcsname{ega4.c04.subsec.1-9}
\expandafter\def\csname sga@target@subsection:IV.10.1\endcsname{ega4.c04.subsec.10-1}
\expandafter\def\csname sga@target@subsection:IV.10.10\endcsname{ega4.c04.subsec.10-10}
\expandafter\def\csname sga@target@subsection:IV.10.2\endcsname{ega4.c04.subsec.10-2}
\expandafter\def\csname sga@target@subsection:IV.10.3\endcsname{ega4.c04.subsec.10-3}
\expandafter\def\csname sga@target@subsection:IV.10.4\endcsname{ega4.c04.subsec.10-4}
\expandafter\def\csname sga@target@subsection:IV.10.5\endcsname{ega4.c04.subsec.10-5}
\expandafter\def\csname sga@target@subsection:IV.10.6\endcsname{ega4.c04.subsec.10-6}
\expandafter\def\csname sga@target@subsection:IV.10.7\endcsname{ega4.c04.subsec.10-7}
\expandafter\def\csname sga@target@subsection:IV.10.8\endcsname{ega4.c04.subsec.10-8}
\expandafter\def\csname sga@target@subsection:IV.10.9\endcsname{ega4.c04.subsec.10-9}
\expandafter\def\csname sga@target@subsection:IV.11.1\endcsname{ega4.c04.subsec.11-1}
\expandafter\def\csname sga@target@subsection:IV.11.10\endcsname{ega4.c04.subsec.11-10}
\expandafter\def\csname sga@target@subsection:IV.11.2\endcsname{ega4.c04.subsec.11-2}
\expandafter\def\csname sga@target@subsection:IV.11.3\endcsname{ega4.c04.subsec.11-3}
\expandafter\def\csname sga@target@subsection:IV.11.4\endcsname{ega4.c04.subsec.11-4}
\expandafter\def\csname sga@target@subsection:IV.11.5\endcsname{ega4.c04.subsec.11-5}
\expandafter\def\csname sga@target@subsection:IV.11.6\endcsname{ega4.c04.subsec.11-6}
\expandafter\def\csname sga@target@subsection:IV.11.7\endcsname{ega4.c04.subsec.11-7}
\expandafter\def\csname sga@target@subsection:IV.11.8\endcsname{ega4.c04.subsec.11-8}
\expandafter\def\csname sga@target@subsection:IV.11.9\endcsname{ega4.c04.subsec.11-9}
\expandafter\def\csname sga@target@subsection:IV.12.0\endcsname{ega4.c04.subsec.12-0}
\expandafter\def\csname sga@target@subsection:IV.12.1\endcsname{ega4.c04.subsec.12-1}
\expandafter\def\csname sga@target@subsection:IV.12.2\endcsname{ega4.c04.subsec.12-2}
\expandafter\def\csname sga@target@subsection:IV.12.3\endcsname{ega4.c04.subsec.12-3}
\expandafter\def\csname sga@target@subsection:IV.13.1\endcsname{ega4.c04.subsec.13-1}
\expandafter\def\csname sga@target@subsection:IV.13.2\endcsname{ega4.c04.subsec.13-2}
\expandafter\def\csname sga@target@subsection:IV.13.3\endcsname{ega4.c04.subsec.13-3}
\expandafter\def\csname sga@target@subsection:IV.14.1\endcsname{ega4.c04.subsec.14-1}
\expandafter\def\csname sga@target@subsection:IV.14.2\endcsname{ega4.c04.subsec.14-2}
\expandafter\def\csname sga@target@subsection:IV.14.3\endcsname{ega4.c04.subsec.14-3}
\expandafter\def\csname sga@target@subsection:IV.14.4\endcsname{ega4.c04.subsec.14-4}
\expandafter\def\csname sga@target@subsection:IV.14.5\endcsname{ega4.c04.subsec.14-5}
\expandafter\def\csname sga@target@subsection:IV.15.1\endcsname{ega4.c04.subsec.15-1}
\expandafter\def\csname sga@target@subsection:IV.15.2\endcsname{ega4.c04.subsec.15-2}
\expandafter\def\csname sga@target@subsection:IV.15.3\endcsname{ega4.c04.subsec.15-3}
\expandafter\def\csname sga@target@subsection:IV.15.4\endcsname{ega4.c04.subsec.15-4}
\expandafter\def\csname sga@target@subsection:IV.15.5\endcsname{ega4.c04.subsec.15-5}
\expandafter\def\csname sga@target@subsection:IV.15.6\endcsname{ega4.c04.subsec.15-6}
\expandafter\def\csname sga@target@subsection:IV.15.7\endcsname{ega4.c04.subsec.15-7}
\expandafter\def\csname sga@target@subsection:IV.17.1\endcsname{ega4.c04.subsec.17-1}
\expandafter\def\csname sga@target@subsection:IV.17.10\endcsname{ega4.c04.subsec.17-10}
\expandafter\def\csname sga@target@subsection:IV.17.11\endcsname{ega4.c04.subsec.17-11}
\expandafter\def\csname sga@target@subsection:IV.17.12\endcsname{ega4.c04.subsec.17-12}
\expandafter\def\csname sga@target@subsection:IV.17.13\endcsname{ega4.c04.subsec.17-13}
\expandafter\def\csname sga@target@subsection:IV.17.14\endcsname{ega4.c04.subsec.17-14}
\expandafter\def\csname sga@target@subsection:IV.17.15\endcsname{ega4.c04.subsec.17-15}
\expandafter\def\csname sga@target@subsection:IV.17.16\endcsname{ega4.c04.subsec.17-16}
\expandafter\def\csname sga@target@subsection:IV.17.2\endcsname{ega4.c04.subsec.17-2}
\expandafter\def\csname sga@target@subsection:IV.17.3\endcsname{ega4.c04.subsec.17-3}
\expandafter\def\csname sga@target@subsection:IV.17.4\endcsname{ega4.c04.subsec.17-4}
\expandafter\def\csname sga@target@subsection:IV.17.5\endcsname{ega4.c04.subsec.17-5}
\expandafter\def\csname sga@target@subsection:IV.17.6\endcsname{ega4.c04.subsec.17-6}
\expandafter\def\csname sga@target@subsection:IV.17.7\endcsname{ega4.c04.subsec.17-7}
\expandafter\def\csname sga@target@subsection:IV.17.8\endcsname{ega4.c04.subsec.17-8}
\expandafter\def\csname sga@target@subsection:IV.17.9\endcsname{ega4.c04.subsec.17-9}
\expandafter\def\csname sga@target@subsection:IV.18.1\endcsname{ega4.c04.subsec.18-1}
\expandafter\def\csname sga@target@subsection:IV.18.10\endcsname{ega4.c04.subsec.18-10}
\expandafter\def\csname sga@target@subsection:IV.18.11\endcsname{ega4.c04.subsec.18-11}
\expandafter\def\csname sga@target@subsection:IV.18.12\endcsname{ega4.c04.subsec.18-12}
\expandafter\def\csname sga@target@subsection:IV.18.2\endcsname{ega4.c04.subsec.18-2}
\expandafter\def\csname sga@target@subsection:IV.18.3\endcsname{ega4.c04.subsec.18-3}
\expandafter\def\csname sga@target@subsection:IV.18.4\endcsname{ega4.c04.subsec.18-4}
\expandafter\def\csname sga@target@subsection:IV.18.5\endcsname{ega4.c04.subsec.18-5}
\expandafter\def\csname sga@target@subsection:IV.18.6\endcsname{ega4.c04.subsec.18-6}
\expandafter\def\csname sga@target@subsection:IV.18.7\endcsname{ega4.c04.subsec.18-7}
\expandafter\def\csname sga@target@subsection:IV.18.8\endcsname{ega4.c04.subsec.18-8}
\expandafter\def\csname sga@target@subsection:IV.18.9\endcsname{ega4.c04.subsec.18-9}
\expandafter\def\csname sga@target@subsection:IV.19.1\endcsname{ega4.c04.subsec.19-1}
\expandafter\def\csname sga@target@subsection:IV.19.2\endcsname{ega4.c04.subsec.19-2}
\expandafter\def\csname sga@target@subsection:IV.19.3\endcsname{ega4.c04.subsec.19-3}
\expandafter\def\csname sga@target@subsection:IV.19.4\endcsname{ega4.c04.subsec.19-4}
\expandafter\def\csname sga@target@subsection:IV.19.5\endcsname{ega4.c04.subsec.19-5}
\expandafter\def\csname sga@target@subsection:IV.19.6\endcsname{ega4.c04.subsec.19-6}
\expandafter\def\csname sga@target@subsection:IV.19.7\endcsname{ega4.c04.subsec.19-7}
\expandafter\def\csname sga@target@subsection:IV.19.8\endcsname{ega4.c04.subsec.19-8}
\expandafter\def\csname sga@target@subsection:IV.19.9\endcsname{ega4.c04.subsec.19-9}
\expandafter\def\csname sga@target@subsection:IV.2.1\endcsname{ega4.c04.subsec.2-1}
\expandafter\def\csname sga@target@subsection:IV.2.2\endcsname{ega4.c04.subsec.2-2}
\expandafter\def\csname sga@target@subsection:IV.2.3\endcsname{ega4.c04.subsec.2-3}
\expandafter\def\csname sga@target@subsection:IV.2.4\endcsname{ega4.c04.subsec.2-4}
\expandafter\def\csname sga@target@subsection:IV.2.5\endcsname{ega4.c04.subsec.2-5}
\expandafter\def\csname sga@target@subsection:IV.2.6\endcsname{ega4.c04.subsec.2-6}
\expandafter\def\csname sga@target@subsection:IV.2.7\endcsname{ega4.c04.subsec.2-7}
\expandafter\def\csname sga@target@subsection:IV.2.8\endcsname{ega4.c04.subsec.2-8}
\expandafter\def\csname sga@target@subsection:IV.21.1\endcsname{ega4.c04.subsec.21-1}
\expandafter\def\csname sga@target@subsection:IV.21.10\endcsname{ega4.c04.subsec.21-10}
\expandafter\def\csname sga@target@subsection:IV.21.11\endcsname{ega4.c04.subsec.21-11}
\expandafter\def\csname sga@target@subsection:IV.21.12\endcsname{ega4.c04.subsec.21-12}
\expandafter\def\csname sga@target@subsection:IV.21.13\endcsname{ega4.c04.subsec.21-13}
\expandafter\def\csname sga@target@subsection:IV.21.14\endcsname{ega4.c04.subsec.21-14}
\expandafter\def\csname sga@target@subsection:IV.21.15\endcsname{ega4.c04.subsec.21-15}
\expandafter\def\csname sga@target@subsection:IV.21.2\endcsname{ega4.c04.subsec.21-2}
\expandafter\def\csname sga@target@subsection:IV.21.3\endcsname{ega4.c04.subsec.21-3}
\expandafter\def\csname sga@target@subsection:IV.21.4\endcsname{ega4.c04.subsec.21-4}
\expandafter\def\csname sga@target@subsection:IV.21.5\endcsname{ega4.c04.subsec.21-5}
\expandafter\def\csname sga@target@subsection:IV.21.6\endcsname{ega4.c04.subsec.21-6}
\expandafter\def\csname sga@target@subsection:IV.21.7\endcsname{ega4.c04.subsec.21-7}
\expandafter\def\csname sga@target@subsection:IV.21.8\endcsname{ega4.c04.subsec.21-8}
\expandafter\def\csname sga@target@subsection:IV.21.9\endcsname{ega4.c04.subsec.21-9}
\expandafter\def\csname sga@target@subsection:IV.3.1\endcsname{ega4.c04.subsec.3-1}
\expandafter\def\csname sga@target@subsection:IV.3.2\endcsname{ega4.c04.subsec.3-2}
\expandafter\def\csname sga@target@subsection:IV.3.3\endcsname{ega4.c04.subsec.3-3}
\expandafter\def\csname sga@target@subsection:IV.3.4\endcsname{ega4.c04.subsec.3-4}
\expandafter\def\csname sga@target@subsection:IV.4.1\endcsname{ega4.c04.subsec.4-1}
\expandafter\def\csname sga@target@subsection:IV.4.2\endcsname{ega4.c04.subsec.4-2}
\expandafter\def\csname sga@target@subsection:IV.4.3\endcsname{ega4.c04.subsec.4-3}
\expandafter\def\csname sga@target@subsection:IV.4.4\endcsname{ega4.c04.subsec.4-4}
\expandafter\def\csname sga@target@subsection:IV.4.5\endcsname{ega4.c04.subsec.4-5}
\expandafter\def\csname sga@target@subsection:IV.4.6\endcsname{ega4.c04.subsec.4-6}
\expandafter\def\csname sga@target@subsection:IV.4.7\endcsname{ega4.c04.subsec.4-7}
\expandafter\def\csname sga@target@subsection:IV.4.8\endcsname{ega4.c04.subsec.4-8}
\expandafter\def\csname sga@target@subsection:IV.4.9\endcsname{ega4.c04.subsec.4-9}
\expandafter\def\csname sga@target@subsection:IV.5.1\endcsname{ega4.c04.subsec.5-1}
\expandafter\def\csname sga@target@subsection:IV.5.10\endcsname{ega4.c04.subsec.5-10}
\expandafter\def\csname sga@target@subsection:IV.5.11\endcsname{ega4.c04.subsec.5-11}
\expandafter\def\csname sga@target@subsection:IV.5.12\endcsname{ega4.c04.subsec.5-12}
\expandafter\def\csname sga@target@subsection:IV.5.13\endcsname{ega4.c04.subsec.5-13}
\expandafter\def\csname sga@target@subsection:IV.5.2\endcsname{ega4.c04.subsec.5-2}
\expandafter\def\csname sga@target@subsection:IV.5.3\endcsname{ega4.c04.subsec.5-3}
\expandafter\def\csname sga@target@subsection:IV.5.4\endcsname{ega4.c04.subsec.5-4}
\expandafter\def\csname sga@target@subsection:IV.5.5\endcsname{ega4.c04.subsec.5-5}
\expandafter\def\csname sga@target@subsection:IV.5.6\endcsname{ega4.c04.subsec.5-6}
\expandafter\def\csname sga@target@subsection:IV.5.7\endcsname{ega4.c04.subsec.5-7}
\expandafter\def\csname sga@target@subsection:IV.5.8\endcsname{ega4.c04.subsec.5-8}
\expandafter\def\csname sga@target@subsection:IV.5.9\endcsname{ega4.c04.subsec.5-9}
\expandafter\def\csname sga@target@subsection:IV.6.1\endcsname{ega4.c04.subsec.6-1}
\expandafter\def\csname sga@target@subsection:IV.6.10\endcsname{ega4.c04.subsec.6-10}
\expandafter\def\csname sga@target@subsection:IV.6.11\endcsname{ega4.c04.subsec.6-11}
\expandafter\def\csname sga@target@subsection:IV.6.12\endcsname{ega4.c04.subsec.6-12}
\expandafter\def\csname sga@target@subsection:IV.6.13\endcsname{ega4.c04.subsec.6-13}
\expandafter\def\csname sga@target@subsection:IV.6.14\endcsname{ega4.c04.subsec.6-14}
\expandafter\def\csname sga@target@subsection:IV.6.15\endcsname{ega4.c04.subsec.6-15}
\expandafter\def\csname sga@target@subsection:IV.6.2\endcsname{ega4.c04.subsec.6-2}
\expandafter\def\csname sga@target@subsection:IV.6.3\endcsname{ega4.c04.subsec.6-3}
\expandafter\def\csname sga@target@subsection:IV.6.4\endcsname{ega4.c04.subsec.6-4}
\expandafter\def\csname sga@target@subsection:IV.6.5\endcsname{ega4.c04.subsec.6-5}
\expandafter\def\csname sga@target@subsection:IV.6.6\endcsname{ega4.c04.subsec.6-6}
\expandafter\def\csname sga@target@subsection:IV.6.7\endcsname{ega4.c04.subsec.6-7}
\expandafter\def\csname sga@target@subsection:IV.6.8\endcsname{ega4.c04.subsec.6-8}
\expandafter\def\csname sga@target@subsection:IV.6.9\endcsname{ega4.c04.subsec.6-9}
\expandafter\def\csname sga@target@subsection:IV.7.1\endcsname{ega4.c04.subsec.7-1}
\expandafter\def\csname sga@target@subsection:IV.7.2\endcsname{ega4.c04.subsec.7-2}
\expandafter\def\csname sga@target@subsection:IV.7.3\endcsname{ega4.c04.subsec.7-3}
\expandafter\def\csname sga@target@subsection:IV.7.4\endcsname{ega4.c04.subsec.7-4}
\expandafter\def\csname sga@target@subsection:IV.7.5\endcsname{ega4.c04.subsec.7-5}
\expandafter\def\csname sga@target@subsection:IV.7.6\endcsname{ega4.c04.subsec.7-6}
\expandafter\def\csname sga@target@subsection:IV.7.7\endcsname{ega4.c04.subsec.7-7}
\expandafter\def\csname sga@target@subsection:IV.7.8\endcsname{ega4.c04.subsec.7-8}
\expandafter\def\csname sga@target@subsection:IV.7.9\endcsname{ega4.c04.subsec.7-9}
\expandafter\def\csname sga@target@subsection:IV.8.1\endcsname{ega4.c04.subsec.8-1}
\expandafter\def\csname sga@target@subsection:IV.8.10\endcsname{ega4.c04.subsec.8-10}
\expandafter\def\csname sga@target@subsection:IV.8.11\endcsname{ega4.c04.subsec.8-11}
\expandafter\def\csname sga@target@subsection:IV.8.12\endcsname{ega4.c04.subsec.8-12}
\expandafter\def\csname sga@target@subsection:IV.8.13\endcsname{ega4.c04.subsec.8-13}
\expandafter\def\csname sga@target@subsection:IV.8.14\endcsname{ega4.c04.subsec.8-14}
\expandafter\def\csname sga@target@subsection:IV.8.2\endcsname{ega4.c04.subsec.8-2}
\expandafter\def\csname sga@target@subsection:IV.8.3\endcsname{ega4.c04.subsec.8-3}
\expandafter\def\csname sga@target@subsection:IV.8.4\endcsname{ega4.c04.subsec.8-4}
\expandafter\def\csname sga@target@subsection:IV.8.5\endcsname{ega4.c04.subsec.8-5}
\expandafter\def\csname sga@target@subsection:IV.8.6\endcsname{ega4.c04.subsec.8-6}
\expandafter\def\csname sga@target@subsection:IV.8.7\endcsname{ega4.c04.subsec.8-7}
\expandafter\def\csname sga@target@subsection:IV.8.8\endcsname{ega4.c04.subsec.8-8}
\expandafter\def\csname sga@target@subsection:IV.8.9\endcsname{ega4.c04.subsec.8-9}
\expandafter\def\csname sga@target@subsection:IV.9.1\endcsname{ega4.c04.subsec.9-1}
\expandafter\def\csname sga@target@subsection:IV.9.2\endcsname{ega4.c04.subsec.9-2}
\expandafter\def\csname sga@target@subsection:IV.9.3\endcsname{ega4.c04.subsec.9-3}
\expandafter\def\csname sga@target@subsection:IV.9.4\endcsname{ega4.c04.subsec.9-4}
\expandafter\def\csname sga@target@subsection:IV.9.5\endcsname{ega4.c04.subsec.9-5}
\expandafter\def\csname sga@target@subsection:IV.9.6\endcsname{ega4.c04.subsec.9-6}
\expandafter\def\csname sga@target@subsection:IV.9.7\endcsname{ega4.c04.subsec.9-7}
\expandafter\def\csname sga@target@subsection:IV.9.8\endcsname{ega4.c04.subsec.9-8}
\expandafter\def\csname sga@target@subsection:IV.9.9\endcsname{ega4.c04.subsec.9-9}
\expandafter\def\csname sga@bib@I-4\endcsname{ega4.reader.bib.i-4}
\expandafter\def\csname sga@bib@IV-32\endcsname{ega4.reader.bib.iv-32}
\expandafter\def\csname sga@bib@IV-34\endcsname{ega4.reader.bib.iv-34}
\expandafter\def\csname sga@bib@IV-37\endcsname{ega4.reader.bib.iv-37}
\expandafter\def\csname sga@bib@IV-41\endcsname{ega4.reader.bib.iv-41}
\expandafter\def\csname sga@bib@IV-42\endcsname{ega4.reader.bib.iv-42}
\expandafter\def\csname sga@bib@IV-43\endcsname{ega4.reader.bib.iv-43}
\expandafter\def\csname sga@bib@IV-44\endcsname{ega4.reader.bib.iv-44}
\expandafter\def\csname sga@bib@SGA1\endcsname{ega4.reader.bib.sga1}
\let\sgaOriginalLabel\label
\renewcommand{\label}[1]{%
  \sgaOriginalLabel{#1}%
  \ifcsname sga@target@#1\endcsname
    \expandafter\sgaStableTarget\expandafter{\csname sga@target@#1\endcsname}%
  \fi
}
\let\sgaOriginalBibitem\bibitem
\newcommand{\sgaBibTarget}[1]{%
  \ifcsname sga@bib@#1\endcsname
    \expandafter\sgaStableTarget\expandafter{\csname sga@bib@#1\endcsname}%
  \fi
}
\newcommand{\sgaBibitemPlain}[1]{%
  \sgaOriginalBibitem{#1}\sgaBibTarget{#1}}
\newcommand{\sgaBibitemOptional}[2][]{%
  \sgaOriginalBibitem[#1]{#2}\sgaBibTarget{#2}}
\renewcommand{\bibitem}{%
  \@ifnextchar[\sgaBibitemOptional\sgaBibitemPlain}
\makeatother
% END NATIVE INPUT EGA4_reference_v2_registry
% BEGIN NATIVE INPUT EGA4_reference_v2_links
% Generated stable-link dispatch for the EGA IV reader.
% Stable destinations are emitted by EGA4_reference_v2_registry.tex.
\makeatletter
\expandafter\def\csname sga@target@IV.1.1.1\endcsname{ega4.c04.def.1-1-1}
\expandafter\def\csname sga@target@IV.1.1.2\endcsname{ega4.c04.prop.1-1-2}
\expandafter\def\csname sga@target@IV.1.1.3\endcsname{ega4.c04.prop.1-1-3}
\expandafter\def\csname sga@target@IV.1.1.4\endcsname{ega4.c04.cor.1-1-4}
\expandafter\def\csname sga@target@IV.1.1.5\endcsname{ega4.c04.prop.1-1-5}
\expandafter\def\csname sga@target@IV.1.1.6\endcsname{ega4.c04.prop.1-1-6}
\expandafter\def\csname sga@target@IV.1.10.1\endcsname{ega4.c04.thm.1-10-1}
\expandafter\def\csname sga@target@IV.1.10.2\endcsname{ega4.c04.cor.1-10-2}
\expandafter\def\csname sga@target@IV.1.10.3\endcsname{ega4.c04.prop.1-10-3}
\expandafter\def\csname sga@target@IV.1.10.4\endcsname{ega4.c04.cor.1-10-4}
\expandafter\def\csname sga@target@IV.1.2.1\endcsname{ega4.c04.def.1-2-1}
\expandafter\def\csname sga@target@IV.1.2.2\endcsname{ega4.c04.prop.1-2-2}
\expandafter\def\csname sga@target@IV.1.2.3\endcsname{ega4.c04.cor.1-2-3}
\expandafter\def\csname sga@target@IV.1.2.4\endcsname{ega4.c04.prop.1-2-4}
\expandafter\def\csname sga@target@IV.1.2.4.1\endcsname{ega4.c04.eq.1-2-4-1}
\expandafter\def\csname sga@target@IV.1.2.5\endcsname{ega4.c04.prop.1-2-5}
\expandafter\def\csname sga@target@IV.1.2.6\endcsname{ega4.c04.prop.1-2-6}
\expandafter\def\csname sga@target@IV.1.2.7\endcsname{ega4.c04.prop.1-2-7}
\expandafter\def\csname sga@target@IV.1.2.8\endcsname{ega4.c04.cor.1-2-8}
\expandafter\def\csname sga@target@IV.1.2.9\endcsname{ega4.c04.prop.1-2-9}
\expandafter\def\csname sga@target@IV.1.3.1\endcsname{ega4.c04.statement.1-3-1}
\expandafter\def\csname sga@target@IV.1.3.10\endcsname{ega4.c04.statement.1-3-10}
\expandafter\def\csname sga@target@IV.1.3.11\endcsname{ega4.c04.prop.1-3-11}
\expandafter\def\csname sga@target@IV.1.3.2\endcsname{ega4.c04.def.1-3-2}
\expandafter\def\csname sga@target@IV.1.3.3\endcsname{ega4.c04.prop.1-3-3}
\expandafter\def\csname sga@target@IV.1.3.4\endcsname{ega4.c04.prop.1-3-4}
\expandafter\def\csname sga@target@IV.1.3.5\endcsname{ega4.c04.cor.1-3-5}
\expandafter\def\csname sga@target@IV.1.3.6\endcsname{ega4.c04.prop.1-3-6}
\expandafter\def\csname sga@target@IV.1.3.7\endcsname{ega4.c04.prop.1-3-7}
\expandafter\def\csname sga@target@IV.1.3.8\endcsname{ega4.c04.statement.1-3-8}
\expandafter\def\csname sga@target@IV.1.3.9\endcsname{ega4.c04.prop.1-3-9}
\expandafter\def\csname sga@target@IV.1.4.1\endcsname{ega4.c04.statement.1-4-1}
\expandafter\def\csname sga@target@IV.1.4.2\endcsname{ega4.c04.def.1-4-2}
\expandafter\def\csname sga@target@IV.1.4.3\endcsname{ega4.c04.prop.1-4-3}
\expandafter\def\csname sga@target@IV.1.4.3.1\endcsname{ega4.c04.cor.1-4-3-1}
\expandafter\def\csname sga@target@IV.1.4.4\endcsname{ega4.c04.prop.1-4-4}
\expandafter\def\csname sga@target@IV.1.4.4.1\endcsname{ega4.c04.lem.1-4-4-1}
\expandafter\def\csname sga@target@IV.1.4.5\endcsname{ega4.c04.cor.1-4-5}
\expandafter\def\csname sga@target@IV.1.4.6\endcsname{ega4.c04.prop.1-4-6}
\expandafter\def\csname sga@target@IV.1.4.7\endcsname{ega4.c04.prop.1-4-7}
\expandafter\def\csname sga@target@IV.1.4.7.1\endcsname{ega4.c04.lem.1-4-7-1}
\expandafter\def\csname sga@target@IV.1.5.1\endcsname{ega4.c04.prop.1-5-1}
\expandafter\def\csname sga@target@IV.1.5.2\endcsname{ega4.c04.def.1-5-2}
\expandafter\def\csname sga@target@IV.1.5.3\endcsname{ega4.c04.prop.1-5-3}
\expandafter\def\csname sga@target@IV.1.5.4\endcsname{ega4.c04.prop.1-5-4}
\expandafter\def\csname sga@target@IV.1.5.5\endcsname{ega4.c04.cor.1-5-5}
\expandafter\def\csname sga@target@IV.1.5.6\endcsname{ega4.c04.cor.1-5-6}
\expandafter\def\csname sga@target@IV.1.5.7\endcsname{ega4.c04.prop.1-5-7}
\expandafter\def\csname sga@target@IV.1.6.1\endcsname{ega4.c04.def.1-6-1}
\expandafter\def\csname sga@target@IV.1.6.2\endcsname{ega4.c04.prop.1-6-2}
\expandafter\def\csname sga@target@IV.1.6.3\endcsname{ega4.c04.cor.1-6-3}
\expandafter\def\csname sga@target@IV.1.6.4\endcsname{ega4.c04.rem.1-6-4}
\expandafter\def\csname sga@target@IV.1.6.5\endcsname{ega4.c04.prop.1-6-5}
\expandafter\def\csname sga@target@IV.1.7.1\endcsname{ega4.c04.statement.1-7-1}
\expandafter\def\csname sga@target@IV.1.7.10\endcsname{ega4.c04.statement.1-7-10}
\expandafter\def\csname sga@target@IV.1.7.11\endcsname{ega4.c04.statement.1-7-11}
\expandafter\def\csname sga@target@IV.1.7.12\endcsname{ega4.c04.statement.1-7-12}
\expandafter\def\csname sga@target@IV.1.7.13\endcsname{ega4.c04.statement.1-7-13}
\expandafter\def\csname sga@target@IV.1.7.14\endcsname{ega4.c04.statement.1-7-14}
\expandafter\def\csname sga@target@IV.1.7.15\endcsname{ega4.c04.statement.1-7-15}
\expandafter\def\csname sga@target@IV.1.7.16\endcsname{ega4.c04.statement.1-7-16}
\expandafter\def\csname sga@target@IV.1.7.17\endcsname{ega4.c04.statement.1-7-17}
\expandafter\def\csname sga@target@IV.1.7.18\endcsname{ega4.c04.statement.1-7-18}
\expandafter\def\csname sga@target@IV.1.7.19\endcsname{ega4.c04.statement.1-7-19}
\expandafter\def\csname sga@target@IV.1.7.2\endcsname{ega4.c04.statement.1-7-2}
\expandafter\def\csname sga@target@IV.1.7.20\endcsname{ega4.c04.statement.1-7-20}
\expandafter\def\csname sga@target@IV.1.7.21\endcsname{ega4.c04.statement.1-7-21}
\expandafter\def\csname sga@target@IV.1.7.22\endcsname{ega4.c04.statement.1-7-22}
\expandafter\def\csname sga@target@IV.1.7.3\endcsname{ega4.c04.statement.1-7-3}
\expandafter\def\csname sga@target@IV.1.7.4\endcsname{ega4.c04.statement.1-7-4}
\expandafter\def\csname sga@target@IV.1.7.5\endcsname{ega4.c04.statement.1-7-5}
\expandafter\def\csname sga@target@IV.1.7.6\endcsname{ega4.c04.statement.1-7-6}
\expandafter\def\csname sga@target@IV.1.7.7\endcsname{ega4.c04.statement.1-7-7}
\expandafter\def\csname sga@target@IV.1.7.8\endcsname{ega4.c04.statement.1-7-8}
\expandafter\def\csname sga@target@IV.1.7.9\endcsname{ega4.c04.statement.1-7-9}
\expandafter\def\csname sga@target@IV.1.8.1\endcsname{ega4.c04.statement.1-8-1}
\expandafter\def\csname sga@target@IV.1.8.2\endcsname{ega4.c04.prop.1-8-2}
\expandafter\def\csname sga@target@IV.1.8.3\endcsname{ega4.c04.lem.1-8-3}
\expandafter\def\csname sga@target@IV.1.8.4\endcsname{ega4.c04.thm.1-8-4}
\expandafter\def\csname sga@target@IV.1.8.4.1\endcsname{ega4.c04.lem.1-8-4-1}
\expandafter\def\csname sga@target@IV.1.8.4.2\endcsname{ega4.c04.lem.1-8-4-2}
\expandafter\def\csname sga@target@IV.1.8.5\endcsname{ega4.c04.cor.1-8-5}
\expandafter\def\csname sga@target@IV.1.8.5.1\endcsname{ega4.c04.lem.1-8-5-1}
\expandafter\def\csname sga@target@IV.1.8.6\endcsname{ega4.c04.cor.1-8-6}
\expandafter\def\csname sga@target@IV.1.8.7\endcsname{ega4.c04.prop.1-8-7}
\expandafter\def\csname sga@target@IV.1.8.7.1\endcsname{ega4.c04.lem.1-8-7-1}
\expandafter\def\csname sga@target@IV.1.9.1\endcsname{ega4.c04.lem.1-9-1}
\expandafter\def\csname sga@target@IV.1.9.10\endcsname{ega4.c04.prop.1-9-10}
\expandafter\def\csname sga@target@IV.1.9.11\endcsname{ega4.c04.prop.1-9-11}
\expandafter\def\csname sga@target@IV.1.9.12\endcsname{ega4.c04.cor.1-9-12}
\expandafter\def\csname sga@target@IV.1.9.13\endcsname{ega4.c04.statement.1-9-13}
\expandafter\def\csname sga@target@IV.1.9.14\endcsname{ega4.c04.prop.1-9-14}
\expandafter\def\csname sga@target@IV.1.9.15\endcsname{ega4.c04.prop.1-9-15}
\expandafter\def\csname sga@target@IV.1.9.16\endcsname{ega4.c04.statement.1-9-16}
\expandafter\def\csname sga@target@IV.1.9.2\endcsname{ega4.c04.prop.1-9-2}
\expandafter\def\csname sga@target@IV.1.9.2.1\endcsname{ega4.c04.statement.1-9-2-1}
\expandafter\def\csname sga@target@IV.1.9.2.i\endcsname{ega4.c04.statement.1-9-2-i}
\expandafter\def\csname sga@target@IV.1.9.2.ii\endcsname{ega4.c04.statement.1-9-2-ii}
\expandafter\def\csname sga@target@IV.1.9.3\endcsname{ega4.c04.prop.1-9-3}
\expandafter\def\csname sga@target@IV.1.9.3.1\endcsname{ega4.c04.lem.1-9-3-1}
\expandafter\def\csname sga@target@IV.1.9.3.2\endcsname{ega4.c04.lem.1-9-3-2}
\expandafter\def\csname sga@target@IV.1.9.4\endcsname{ega4.c04.def.1-9-4}
\expandafter\def\csname sga@target@IV.1.9.5\endcsname{ega4.c04.prop.1-9-5}
\expandafter\def\csname sga@target@IV.1.9.6\endcsname{ega4.c04.statement.1-9-6}
\expandafter\def\csname sga@target@IV.1.9.7\endcsname{ega4.c04.thm.1-9-7}
\expandafter\def\csname sga@target@IV.1.9.8\endcsname{ega4.c04.cor.1-9-8}
\expandafter\def\csname sga@target@IV.1.9.9\endcsname{ega4.c04.cor.1-9-9}
\expandafter\def\csname sga@target@IV.10.1.1\endcsname{ega4.c04.statement.10-1-1}
\expandafter\def\csname sga@target@IV.10.1.2\endcsname{ega4.c04.prop.10-1-2}
\expandafter\def\csname sga@target@IV.10.1.3\endcsname{ega4.c04.def.10-1-3}
\expandafter\def\csname sga@target@IV.10.1.4\endcsname{ega4.c04.cor.10-1-4}
\expandafter\def\csname sga@target@IV.10.10.1\endcsname{ega4.c04.statement.10-10-1}
\expandafter\def\csname sga@target@IV.10.10.2\endcsname{ega4.c04.statement.10-10-2}
\expandafter\def\csname sga@target@IV.10.10.3\endcsname{ega4.c04.statement.10-10-3}
\expandafter\def\csname sga@target@IV.10.10.4\endcsname{ega4.c04.statement.10-10-4}
\expandafter\def\csname sga@target@IV.10.10.5\endcsname{ega4.c04.statement.10-10-5}
\expandafter\def\csname sga@target@IV.10.2.1\endcsname{ega4.c04.prop.10-2-1}
\expandafter\def\csname sga@target@IV.10.2.2\endcsname{ega4.c04.def.10-2-2}
\expandafter\def\csname sga@target@IV.10.2.3\endcsname{ega4.c04.cor.10-2-3}
\expandafter\def\csname sga@target@IV.10.2.4\endcsname{ega4.c04.cor.10-2-4}
\expandafter\def\csname sga@target@IV.10.2.5\endcsname{ega4.c04.cor.10-2-5}
\expandafter\def\csname sga@target@IV.10.2.6\endcsname{ega4.c04.cor.10-2-6}
\expandafter\def\csname sga@target@IV.10.2.7\endcsname{ega4.c04.prop.10-2-7}
\expandafter\def\csname sga@target@IV.10.2.8\endcsname{ega4.c04.item.10-2-8}
\expandafter\def\csname sga@target@IV.10.2.8.1\endcsname{ega4.c04.eq.10-2-8-1}
\expandafter\def\csname sga@target@IV.10.3.1\endcsname{ega4.c04.def.10-3-1}
\expandafter\def\csname sga@target@IV.10.3.2\endcsname{ega4.c04.prop.10-3-2}
\expandafter\def\csname sga@target@IV.10.3.3\endcsname{ega4.c04.prop.10-3-3}
\expandafter\def\csname sga@target@IV.10.4.1\endcsname{ega4.c04.def.10-4-1}
\expandafter\def\csname sga@target@IV.10.4.10\endcsname{ega4.c04.statement.10-4-10}
\expandafter\def\csname sga@target@IV.10.4.11\endcsname{ega4.c04.statement.10-4-11}
\expandafter\def\csname sga@target@IV.10.4.11.1\endcsname{ega4.c04.lem.10-4-11-1}
\expandafter\def\csname sga@target@IV.10.4.12\endcsname{ega4.c04.cor.10-4-12}
\expandafter\def\csname sga@target@IV.10.4.12.1\endcsname{ega4.c04.sec.10-4-12-1}
\expandafter\def\csname sga@target@IV.10.4.12.1.errata\endcsname{ega4.c04.eq.10-4-12-1-errata}
\expandafter\def\csname sga@target@IV.10.4.13\endcsname{ega4.c04.cor.10-4-13}
\expandafter\def\csname sga@target@IV.10.4.14\endcsname{ega4.c04.rem.10-4-14}
\expandafter\def\csname sga@target@IV.10.4.2\endcsname{ega4.c04.prop.10-4-2}
\expandafter\def\csname sga@target@IV.10.4.3\endcsname{ega4.c04.statement.10-4-3}
\expandafter\def\csname sga@target@IV.10.4.4\endcsname{ega4.c04.prop.10-4-4}
\expandafter\def\csname sga@target@IV.10.4.5\endcsname{ega4.c04.prop.10-4-5}
\expandafter\def\csname sga@target@IV.10.4.5.1\endcsname{ega4.c04.lem.10-4-5-1}
\expandafter\def\csname sga@target@IV.10.4.6\endcsname{ega4.c04.cor.10-4-6}
\expandafter\def\csname sga@target@IV.10.4.7\endcsname{ega4.c04.cor.10-4-7}
\expandafter\def\csname sga@target@IV.10.4.8\endcsname{ega4.c04.cor.10-4-8}
\expandafter\def\csname sga@target@IV.10.4.9\endcsname{ega4.c04.statement.10-4-9}
\expandafter\def\csname sga@target@IV.10.5.1\endcsname{ega4.c04.prop.10-5-1}
\expandafter\def\csname sga@target@IV.10.5.10\endcsname{ega4.c04.prop.10-5-10}
\expandafter\def\csname sga@target@IV.10.5.2\endcsname{ega4.c04.cor.10-5-2}
\expandafter\def\csname sga@target@IV.10.5.3\endcsname{ega4.c04.cor.10-5-3}
\expandafter\def\csname sga@target@IV.10.5.4\endcsname{ega4.c04.rem.10-5-4}
\expandafter\def\csname sga@target@IV.10.5.5\endcsname{ega4.c04.prop.10-5-5}
\expandafter\def\csname sga@target@IV.10.5.6\endcsname{ega4.c04.rem.10-5-6}
\expandafter\def\csname sga@target@IV.10.5.7\endcsname{ega4.c04.prop.10-5-7}
\expandafter\def\csname sga@target@IV.10.5.8\endcsname{ega4.c04.cor.10-5-8}
\expandafter\def\csname sga@target@IV.10.5.9\endcsname{ega4.c04.cor.10-5-9}
\expandafter\def\csname sga@target@IV.10.6.1\endcsname{ega4.c04.prop.10-6-1}
\expandafter\def\csname sga@target@IV.10.6.1.1\endcsname{ega4.c04.eq.10-6-1-1}
\expandafter\def\csname sga@target@IV.10.6.1.2\endcsname{ega4.c04.eq.10-6-1-2}
\expandafter\def\csname sga@target@IV.10.6.1.3\endcsname{ega4.c04.eq.10-6-1-3}
\expandafter\def\csname sga@target@IV.10.6.1.4\endcsname{ega4.c04.eq.10-6-1-4}
\expandafter\def\csname sga@target@IV.10.6.2\endcsname{ega4.c04.cor.10-6-2}
\expandafter\def\csname sga@target@IV.10.6.3\endcsname{ega4.c04.prop.10-6-3}
\expandafter\def\csname sga@target@IV.10.6.3.1\endcsname{ega4.c04.eq.10-6-3-1}
\expandafter\def\csname sga@target@IV.10.6.3.2\endcsname{ega4.c04.eq.10-6-3-2}
\expandafter\def\csname sga@target@IV.10.6.3.3\endcsname{ega4.c04.eq.10-6-3-3}
\expandafter\def\csname sga@target@IV.10.6.4\endcsname{ega4.c04.cor.10-6-4}
\expandafter\def\csname sga@target@IV.10.6.4.1\endcsname{ega4.c04.eq.10-6-4-1}
\expandafter\def\csname sga@target@IV.10.7.1\endcsname{ega4.c04.statement.10-7-1}
\expandafter\def\csname sga@target@IV.10.7.2\endcsname{ega4.c04.statement.10-7-2}
\expandafter\def\csname sga@target@IV.10.7.3\endcsname{ega4.c04.statement.10-7-3}
\expandafter\def\csname sga@target@IV.10.8.1\endcsname{ega4.c04.def.10-8-1}
\expandafter\def\csname sga@target@IV.10.8.1.1\endcsname{ega4.c04.eq.10-8-1-1}
\expandafter\def\csname sga@target@IV.10.8.1.2\endcsname{ega4.c04.eq.10-8-1-2}
\expandafter\def\csname sga@target@IV.10.8.2\endcsname{ega4.c04.statement.10-8-2}
\expandafter\def\csname sga@target@IV.10.8.3\endcsname{ega4.c04.lem.10-8-3}
\expandafter\def\csname sga@target@IV.10.8.3.1\endcsname{ega4.c04.eq.10-8-3-1}
\expandafter\def\csname sga@target@IV.10.8.4\endcsname{ega4.c04.cor.10-8-4}
\expandafter\def\csname sga@target@IV.10.8.5\endcsname{ega4.c04.prop.10-8-5}
\expandafter\def\csname sga@target@IV.10.8.6\endcsname{ega4.c04.cor.10-8-6}
\expandafter\def\csname sga@target@IV.10.8.7\endcsname{ega4.c04.prop.10-8-7}
\expandafter\def\csname sga@target@IV.10.8.7.1\endcsname{ega4.c04.eq.10-8-7-1}
\expandafter\def\csname sga@target@IV.10.8.8\endcsname{ega4.c04.cor.10-8-8}
\expandafter\def\csname sga@target@IV.10.9.1\endcsname{ega4.c04.statement.10-9-1}
\expandafter\def\csname sga@target@IV.10.9.2\endcsname{ega4.c04.statement.10-9-2}
\expandafter\def\csname sga@target@IV.10.9.3\endcsname{ega4.c04.statement.10-9-3}
\expandafter\def\csname sga@target@IV.10.9.4\endcsname{ega4.c04.prop.10-9-4}
\expandafter\def\csname sga@target@IV.10.9.5\endcsname{ega4.c04.statement.10-9-5}
\expandafter\def\csname sga@target@IV.10.9.6\endcsname{ega4.c04.prop.10-9-6}
\expandafter\def\csname sga@target@IV.11.1.1\endcsname{ega4.c04.thm.11-1-1}
\expandafter\def\csname sga@target@IV.11.1.1.1\endcsname{ega4.c04.sec.11-1-1-1}
\expandafter\def\csname sga@target@IV.11.1.1.2\endcsname{ega4.c04.lem.11-1-1-2}
\expandafter\def\csname sga@target@IV.11.1.2\endcsname{ega4.c04.cor.11-1-2}
\expandafter\def\csname sga@target@IV.11.1.3\endcsname{ega4.c04.cor.11-1-3}
\expandafter\def\csname sga@target@IV.11.1.4\endcsname{ega4.c04.cor.11-1-4}
\expandafter\def\csname sga@target@IV.11.1.5\endcsname{ega4.c04.cor.11-1-5}
\expandafter\def\csname sga@target@IV.11.1.6\endcsname{ega4.c04.cor.11-1-6}
\expandafter\def\csname sga@target@IV.11.10.1\endcsname{ega4.c04.prop.11-10-1}
\expandafter\def\csname sga@target@IV.11.10.1.1\endcsname{ega4.c04.eq.11-10-1-1}
\expandafter\def\csname sga@target@IV.11.10.1.2\endcsname{ega4.c04.eq.11-10-1-2}
\expandafter\def\csname sga@target@IV.11.10.10\endcsname{ega4.c04.prop.11-10-10}
\expandafter\def\csname sga@target@IV.11.10.2\endcsname{ega4.c04.def.11-10-2}
\expandafter\def\csname sga@target@IV.11.10.3\endcsname{ega4.c04.item.11-10-3}
\expandafter\def\csname sga@target@IV.11.10.4\endcsname{ega4.c04.prop.11-10-4}
\expandafter\def\csname sga@target@IV.11.10.5\endcsname{ega4.c04.thm.11-10-5}
\expandafter\def\csname sga@target@IV.11.10.6\endcsname{ega4.c04.prop.11-10-6}
\expandafter\def\csname sga@target@IV.11.10.7\endcsname{ega4.c04.cor.11-10-7}
\expandafter\def\csname sga@target@IV.11.10.8\endcsname{ega4.c04.def.11-10-8}
\expandafter\def\csname sga@target@IV.11.10.9\endcsname{ega4.c04.thm.11-10-9}
\expandafter\def\csname sga@target@IV.11.2.1\endcsname{ega4.c04.statement.11-2-1}
\expandafter\def\csname sga@target@IV.11.2.1.1\endcsname{ega4.c04.eq.11-2-1-1}
\expandafter\def\csname sga@target@IV.11.2.1.2\endcsname{ega4.c04.eq.11-2-1-2}
\expandafter\def\csname sga@target@IV.11.2.1.3\endcsname{ega4.c04.eq.11-2-1-3}
\expandafter\def\csname sga@target@IV.11.2.1.4\endcsname{ega4.c04.eq.11-2-1-4}
\expandafter\def\csname sga@target@IV.11.2.1.5\endcsname{ega4.c04.eq.11-2-1-5}
\expandafter\def\csname sga@target@IV.11.2.10\endcsname{ega4.c04.rem.11-2-10}
\expandafter\def\csname sga@target@IV.11.2.2\endcsname{ega4.c04.statement.11-2-2}
\expandafter\def\csname sga@target@IV.11.2.2.1\endcsname{ega4.c04.eq.11-2-2-1}
\expandafter\def\csname sga@target@IV.11.2.2.2\endcsname{ega4.c04.lem.11-2-2-2}
\expandafter\def\csname sga@target@IV.11.2.3\endcsname{ega4.c04.statement.11-2-3}
\expandafter\def\csname sga@target@IV.11.2.3.1\endcsname{ega4.c04.eq.11-2-3-1}
\expandafter\def\csname sga@target@IV.11.2.3.2\endcsname{ega4.c04.eq.11-2-3-2}
\expandafter\def\csname sga@target@IV.11.2.3.3\endcsname{ega4.c04.eq.11-2-3-3}
\expandafter\def\csname sga@target@IV.11.2.4\endcsname{ega4.c04.lem.11-2-4}
\expandafter\def\csname sga@target@IV.11.2.5\endcsname{ega4.c04.lem.11-2-5}
\expandafter\def\csname sga@target@IV.11.2.6\endcsname{ega4.c04.thm.11-2-6}
\expandafter\def\csname sga@target@IV.11.2.6.1\endcsname{ega4.c04.cor.11-2-6-1}
\expandafter\def\csname sga@target@IV.11.2.7\endcsname{ega4.c04.cor.11-2-7}
\expandafter\def\csname sga@target@IV.11.2.8\endcsname{ega4.c04.prop.11-2-8}
\expandafter\def\csname sga@target@IV.11.2.9\endcsname{ega4.c04.thm.11-2-9}
\expandafter\def\csname sga@target@IV.11.2.9.1\endcsname{ega4.c04.statement.11-2-9-1}
\expandafter\def\csname sga@target@IV.11.2.9.2\endcsname{ega4.c04.lem.11-2-9-2}
\expandafter\def\csname sga@target@IV.11.2.9.3\endcsname{ega4.c04.lem.11-2-9-3}
\expandafter\def\csname sga@target@IV.11.2.9.4\endcsname{ega4.c04.cor.11-2-9-4}
\expandafter\def\csname sga@target@IV.11.2.9.5\endcsname{ega4.c04.lem.11-2-9-5}
\expandafter\def\csname sga@target@IV.11.2.9.6\endcsname{ega4.c04.lem.11-2-9-6}
\expandafter\def\csname sga@target@IV.11.2.9.7\endcsname{ega4.c04.eq.11-2-9-7}
\expandafter\def\csname sga@target@IV.11.2.9.8\endcsname{ega4.c04.eq.11-2-9-8}
\expandafter\def\csname sga@target@IV.11.2.9.9\endcsname{ega4.c04.eq.11-2-9-9}
\expandafter\def\csname sga@target@IV.11.3\endcsname{ega4.c04.subsec.11-3.key-iv-11-3}
\expandafter\def\csname sga@target@IV.11.3.1\endcsname{ega4.c04.thm.11-3-1}
\expandafter\def\csname sga@target@IV.11.3.10\endcsname{ega4.c04.thm.11-3-10}
\expandafter\def\csname sga@target@IV.11.3.10.1\endcsname{ega4.c04.lem.11-3-10-1}
\expandafter\def\csname sga@target@IV.11.3.10.2\endcsname{ega4.c04.lem.11-3-10-2}
\expandafter\def\csname sga@target@IV.11.3.10.3\endcsname{ega4.c04.eq.11-3-10-3}
\expandafter\def\csname sga@target@IV.11.3.10.4\endcsname{ega4.c04.sec.11-3-10-4}
\expandafter\def\csname sga@target@IV.11.3.11\endcsname{ega4.c04.cor.11-3-11}
\expandafter\def\csname sga@target@IV.11.3.12\endcsname{ega4.c04.rem.11-3-12}
\expandafter\def\csname sga@target@IV.11.3.13\endcsname{ega4.c04.prop.11-3-13}
\expandafter\def\csname sga@target@IV.11.3.14\endcsname{ega4.c04.cor.11-3-14}
\expandafter\def\csname sga@target@IV.11.3.15\endcsname{ega4.c04.prop.11-3-15}
\expandafter\def\csname sga@target@IV.11.3.16\endcsname{ega4.c04.prop.11-3-16}
\expandafter\def\csname sga@target@IV.11.3.17\endcsname{ega4.c04.cor.11-3-17}
\expandafter\def\csname sga@target@IV.11.3.2\endcsname{ega4.c04.cor.11-3-2}
\expandafter\def\csname sga@target@IV.11.3.3\endcsname{ega4.c04.cor.11-3-3}
\expandafter\def\csname sga@target@IV.11.3.4\endcsname{ega4.c04.prop.11-3-4}
\expandafter\def\csname sga@target@IV.11.3.5\endcsname{ega4.c04.cor.11-3-5}
\expandafter\def\csname sga@target@IV.11.3.6\endcsname{ega4.c04.prop.11-3-6}
\expandafter\def\csname sga@target@IV.11.3.7\endcsname{ega4.c04.prop.11-3-7}
\expandafter\def\csname sga@target@IV.11.3.8\endcsname{ega4.c04.thm.11-3-8}
\expandafter\def\csname sga@target@IV.11.3.9\endcsname{ega4.c04.prop.11-3-9}
\expandafter\def\csname sga@target@IV.11.3.9.1\endcsname{ega4.c04.lem.11-3-9-1}
\expandafter\def\csname sga@target@IV.11.4.1\endcsname{ega4.c04.thm.11-4-1}
\expandafter\def\csname sga@target@IV.11.4.1.1\endcsname{ega4.c04.lem.11-4-1-1}
\expandafter\def\csname sga@target@IV.11.4.1.2\endcsname{ega4.c04.lem.11-4-1-2}
\expandafter\def\csname sga@target@IV.11.4.10\endcsname{ega4.c04.cor.11-4-10}
\expandafter\def\csname sga@target@IV.11.4.11\endcsname{ega4.c04.prop.11-4-11}
\expandafter\def\csname sga@target@IV.11.4.12\endcsname{ega4.c04.prop.11-4-12}
\expandafter\def\csname sga@target@IV.11.4.13\endcsname{ega4.c04.cor.11-4-13}
\expandafter\def\csname sga@target@IV.11.4.2\endcsname{ega4.c04.rem.11-4-2}
\expandafter\def\csname sga@target@IV.11.4.3\endcsname{ega4.c04.cor.11-4-3}
\expandafter\def\csname sga@target@IV.11.4.4\endcsname{ega4.c04.cor.11-4-4}
\expandafter\def\csname sga@target@IV.11.4.5\endcsname{ega4.c04.prop.11-4-5}
\expandafter\def\csname sga@target@IV.11.4.6\endcsname{ega4.c04.cor.11-4-6}
\expandafter\def\csname sga@target@IV.11.4.7\endcsname{ega4.c04.prop.11-4-7}
\expandafter\def\csname sga@target@IV.11.4.8\endcsname{ega4.c04.rem.11-4-8}
\expandafter\def\csname sga@target@IV.11.4.9\endcsname{ega4.c04.cor.11-4-9}
\expandafter\def\csname sga@target@IV.11.5.1\endcsname{ega4.c04.thm.11-5-1}
\expandafter\def\csname sga@target@IV.11.5.2\endcsname{ega4.c04.cor.11-5-2}
\expandafter\def\csname sga@target@IV.11.5.3\endcsname{ega4.c04.prop.11-5-3}
\expandafter\def\csname sga@target@IV.11.5.4\endcsname{ega4.c04.cor.11-5-4}
\expandafter\def\csname sga@target@IV.11.5.5\endcsname{ega4.c04.prop.11-5-5}
\expandafter\def\csname sga@target@IV.11.5.5.1\endcsname{ega4.c04.lem.11-5-5-1}
\expandafter\def\csname sga@target@IV.11.5.5.2\endcsname{ega4.c04.lem.11-5-5-2}
\expandafter\def\csname sga@target@IV.11.6.1\endcsname{ega4.c04.thm.11-6-1}
\expandafter\def\csname sga@target@IV.11.6.2\endcsname{ega4.c04.cor.11-6-2}
\expandafter\def\csname sga@target@IV.11.7.1\endcsname{ega4.c04.statement.11-7-1}
\expandafter\def\csname sga@target@IV.11.7.2\endcsname{ega4.c04.statement.11-7-2}
\expandafter\def\csname sga@target@IV.11.7.3\endcsname{ega4.c04.statement.11-7-3}
\expandafter\def\csname sga@target@IV.11.7.4\endcsname{ega4.c04.statement.11-7-4}
\expandafter\def\csname sga@target@IV.11.7.5\endcsname{ega4.c04.statement.11-7-5}
\expandafter\def\csname sga@target@IV.11.8.1\endcsname{ega4.c04.thm.11-8-1}
\expandafter\def\csname sga@target@IV.11.8.2\endcsname{ega4.c04.subsec.11-8-2}
\expandafter\def\csname sga@target@IV.11.9.1\endcsname{ega4.c04.statement.11-9-1}
\expandafter\def\csname sga@target@IV.11.9.1.1\endcsname{ega4.c04.statement.11-9-1-1}
\expandafter\def\csname sga@target@IV.11.9.10\endcsname{ega4.c04.thm.11-9-10}
\expandafter\def\csname sga@target@IV.11.9.11\endcsname{ega4.c04.rem.11-9-11}
\expandafter\def\csname sga@target@IV.11.9.12\endcsname{ega4.c04.lem.11-9-12}
\expandafter\def\csname sga@target@IV.11.9.13\endcsname{ega4.c04.prop.11-9-13}
\expandafter\def\csname sga@target@IV.11.9.14\endcsname{ega4.c04.statement.11-9-14}
\expandafter\def\csname sga@target@IV.11.9.15\endcsname{ega4.c04.prop.11-9-15}
\expandafter\def\csname sga@target@IV.11.9.16\endcsname{ega4.c04.thm.11-9-16}
\expandafter\def\csname sga@target@IV.11.9.17\endcsname{ega4.c04.prop.11-9-17}
\expandafter\def\csname sga@target@IV.11.9.17.1\endcsname{ega4.c04.lem.11-9-17-1}
\expandafter\def\csname sga@target@IV.11.9.18\endcsname{ega4.c04.item.11-9-18}
\expandafter\def\csname sga@target@IV.11.9.2\endcsname{ega4.c04.statement.11-9-2}
\expandafter\def\csname sga@target@IV.11.9.3\endcsname{ega4.c04.statement.11-9-3}
\expandafter\def\csname sga@target@IV.11.9.4\endcsname{ega4.c04.statement.11-9-4}
\expandafter\def\csname sga@target@IV.11.9.5\endcsname{ega4.c04.statement.11-9-5}
\expandafter\def\csname sga@target@IV.11.9.6\endcsname{ega4.c04.statement.11-9-6}
\expandafter\def\csname sga@target@IV.11.9.7\endcsname{ega4.c04.prop.11-9-7}
\expandafter\def\csname sga@target@IV.11.9.7.1\endcsname{ega4.c04.lem.11-9-7-1}
\expandafter\def\csname sga@target@IV.11.9.8\endcsname{ega4.c04.subsec.11-9-8}
\expandafter\def\csname sga@target@IV.11.9.9\endcsname{ega4.c04.cor.11-9-9}
\expandafter\def\csname sga@target@IV.12.0.1\endcsname{ega4.c04.statement.12-0-1}
\expandafter\def\csname sga@target@IV.12.0.2\endcsname{ega4.c04.statement.12-0-2}
\expandafter\def\csname sga@target@IV.12.1.1\endcsname{ega4.c04.thm.12-1-1}
\expandafter\def\csname sga@target@IV.12.1.1.1\endcsname{ega4.c04.statement.12-1-1-1}
\expandafter\def\csname sga@target@IV.12.1.1.2\endcsname{ega4.c04.statement.12-1-1-2}
\expandafter\def\csname sga@target@IV.12.1.1.3\endcsname{ega4.c04.statement.12-1-1-3}
\expandafter\def\csname sga@target@IV.12.1.1.4\endcsname{ega4.c04.statement.12-1-1-4}
\expandafter\def\csname sga@target@IV.12.1.1.5\endcsname{ega4.c04.prop.12-1-1-5}
\expandafter\def\csname sga@target@IV.12.1.2\endcsname{ega4.c04.item.12-1-2}
\expandafter\def\csname sga@target@IV.12.1.3\endcsname{ega4.c04.cor.12-1-3}
\expandafter\def\csname sga@target@IV.12.1.4\endcsname{ega4.c04.rem.12-1-4}
\expandafter\def\csname sga@target@IV.12.1.5\endcsname{ega4.c04.cor.12-1-5}
\expandafter\def\csname sga@target@IV.12.1.6\endcsname{ega4.c04.thm.12-1-6}
\expandafter\def\csname sga@target@IV.12.1.7\endcsname{ega4.c04.cor.12-1-7}
\expandafter\def\csname sga@target@IV.12.1.8\endcsname{ega4.c04.item.12-1-8}
\expandafter\def\csname sga@target@IV.12.2.1\endcsname{ega4.c04.thm.12-2-1}
\expandafter\def\csname sga@target@IV.12.2.1.1\endcsname{ega4.c04.statement.12-2-1-1}
\expandafter\def\csname sga@target@IV.12.2.1.2\endcsname{ega4.c04.statement.12-2-1-2}
\expandafter\def\csname sga@target@IV.12.2.1.3\endcsname{ega4.c04.statement.12-2-1-3}
\expandafter\def\csname sga@target@IV.12.2.1.4\endcsname{ega4.c04.lem.12-2-1-4}
\expandafter\def\csname sga@target@IV.12.2.2\endcsname{ega4.c04.cor.12-2-2}
\expandafter\def\csname sga@target@IV.12.2.3\endcsname{ega4.c04.item.12-2-3}
\expandafter\def\csname sga@target@IV.12.2.4\endcsname{ega4.c04.thm.12-2-4}
\expandafter\def\csname sga@target@IV.12.3.1\endcsname{ega4.c04.lem.12-3-1}
\expandafter\def\csname sga@target@IV.12.3.2\endcsname{ega4.c04.prop.12-3-2}
\expandafter\def\csname sga@target@IV.12.3.2.1\endcsname{ega4.c04.eq.12-3-2-1}
\expandafter\def\csname sga@target@IV.12.3.3\endcsname{ega4.c04.prop.12-3-3}
\expandafter\def\csname sga@target@IV.12.3.3.1\endcsname{ega4.c04.lem.12-3-3-1}
\expandafter\def\csname sga@target@IV.12.3.3.2\endcsname{ega4.c04.eq.12-3-3-2}
\expandafter\def\csname sga@target@IV.12.3.3.3\endcsname{ega4.c04.lem.12-3-3-3}
\expandafter\def\csname sga@target@IV.12.3.3.4\endcsname{ega4.c04.eq.12-3-3-4}
\expandafter\def\csname sga@target@IV.12.3.3.5\endcsname{ega4.c04.lem.12-3-3-5}
\expandafter\def\csname sga@target@IV.12.3.4\endcsname{ega4.c04.cor.12-3-4}
\expandafter\def\csname sga@target@IV.13.1.1\endcsname{ega4.c04.lem.13-1-1}
\expandafter\def\csname sga@target@IV.13.1.2\endcsname{ega4.c04.lem.13-1-2}
\expandafter\def\csname sga@target@IV.13.1.3\endcsname{ega4.c04.thm.13-1-3}
\expandafter\def\csname sga@target@IV.13.1.4\endcsname{ega4.c04.cor.13-1-4}
\expandafter\def\csname sga@target@IV.13.1.5\endcsname{ega4.c04.cor.13-1-5}
\expandafter\def\csname sga@target@IV.13.1.6\endcsname{ega4.c04.cor.13-1-6}
\expandafter\def\csname sga@target@IV.13.1.7\endcsname{ega4.c04.prop.13-1-7}
\expandafter\def\csname sga@target@IV.13.2.1\endcsname{ega4.c04.statement.13-2-1}
\expandafter\def\csname sga@target@IV.13.2.10\endcsname{ega4.c04.prop.13-2-10}
\expandafter\def\csname sga@target@IV.13.2.10.1\endcsname{ega4.c04.eq.13-2-10-1}
\expandafter\def\csname sga@target@IV.13.2.10.2\endcsname{ega4.c04.eq.13-2-10-2}
\expandafter\def\csname sga@target@IV.13.2.11\endcsname{ega4.c04.cor.13-2-11}
\expandafter\def\csname sga@target@IV.13.2.12\endcsname{ega4.c04.item.13-2-12}
\expandafter\def\csname sga@target@IV.13.2.2\endcsname{ega4.c04.def.13-2-2}
\expandafter\def\csname sga@target@IV.13.2.3\endcsname{ega4.c04.prop.13-2-3}
\expandafter\def\csname sga@target@IV.13.2.3.1\endcsname{ega4.c04.prop.13-2-3-1}
\expandafter\def\csname sga@target@IV.13.2.4\endcsname{ega4.c04.statement.13-2-4}
\expandafter\def\csname sga@target@IV.13.2.4.1\endcsname{ega4.c04.statement.13-2-4-1}
\expandafter\def\csname sga@target@IV.13.2.4.2\endcsname{ega4.c04.eq.13-2-4-2}
\expandafter\def\csname sga@target@IV.13.2.4.3\endcsname{ega4.c04.eq.13-2-4-3}
\expandafter\def\csname sga@target@IV.13.2.4.4\endcsname{ega4.c04.eq.13-2-4-4}
\expandafter\def\csname sga@target@IV.13.2.4.5\endcsname{ega4.c04.eq.13-2-4-5}
\expandafter\def\csname sga@target@IV.13.2.5\endcsname{ega4.c04.lem.13-2-5}
\expandafter\def\csname sga@target@IV.13.2.6\endcsname{ega4.c04.statement.13-2-6}
\expandafter\def\csname sga@target@IV.13.2.6.1\endcsname{ega4.c04.statement.13-2-6-1}
\expandafter\def\csname sga@target@IV.13.2.7\endcsname{ega4.c04.prop.13-2-7}
\expandafter\def\csname sga@target@IV.13.2.7.1\endcsname{ega4.c04.prop.13-2-7-1}
\expandafter\def\csname sga@target@IV.13.2.8\endcsname{ega4.c04.cor.13-2-8}
\expandafter\def\csname sga@target@IV.13.2.9\endcsname{ega4.c04.cor.13-2-9}
\expandafter\def\csname sga@target@IV.13.2.9.1\endcsname{ega4.c04.eq.13-2-9-1}
\expandafter\def\csname sga@target@IV.13.2.9.2\endcsname{ega4.c04.eq.13-2-9-2}
\expandafter\def\csname sga@target@IV.13.3.1\endcsname{ega4.c04.prop.13-3-1}
\expandafter\def\csname sga@target@IV.13.3.1.1\endcsname{ega4.c04.lem.13-3-1-1}
\expandafter\def\csname sga@target@IV.13.3.2\endcsname{ega4.c04.def.13-3-2}
\expandafter\def\csname sga@target@IV.13.3.3\endcsname{ega4.c04.prop.13-3-3}
\expandafter\def\csname sga@target@IV.13.3.4\endcsname{ega4.c04.cor.13-3-4}
\expandafter\def\csname sga@target@IV.13.3.4.1\endcsname{ega4.c04.cor.13-3-4-1}
\expandafter\def\csname sga@target@IV.13.3.5\endcsname{ega4.c04.cor.13-3-5}
\expandafter\def\csname sga@target@IV.13.3.6\endcsname{ega4.c04.prop.13-3-6}
\expandafter\def\csname sga@target@IV.13.3.6.1\endcsname{ega4.c04.prop.13-3-6-1}
\expandafter\def\csname sga@target@IV.13.3.6.2\endcsname{ega4.c04.eq.13-3-6-2}
\expandafter\def\csname sga@target@IV.13.3.7\endcsname{ega4.c04.prop.13-3-7}
\expandafter\def\csname sga@target@IV.13.3.7.1\endcsname{ega4.c04.prop.13-3-7-1}
\expandafter\def\csname sga@target@IV.13.3.8\endcsname{ega4.c04.prop.13-3-8}
\expandafter\def\csname sga@target@IV.13.3.9\endcsname{ega4.c04.statement.13-3-9}
\expandafter\def\csname sga@target@IV.14.1.1\endcsname{ega4.c04.statement.14-1-1}
\expandafter\def\csname sga@target@IV.14.1.2\endcsname{ega4.c04.prop.14-1-2}
\expandafter\def\csname sga@target@IV.14.1.3\endcsname{ega4.c04.item.14-1-3}
\expandafter\def\csname sga@target@IV.14.1.4\endcsname{ega4.c04.prop.14-1-4}
\expandafter\def\csname sga@target@IV.14.1.5\endcsname{ega4.c04.rem.14-1-5}
\expandafter\def\csname sga@target@IV.14.1.6\endcsname{ega4.c04.statement.14-1-6}
\expandafter\def\csname sga@target@IV.14.2.1\endcsname{ega4.c04.thm.14-2-1}
\expandafter\def\csname sga@target@IV.14.2.1.1\endcsname{ega4.c04.eq.14-2-1-1}
\expandafter\def\csname sga@target@IV.14.2.2\endcsname{ega4.c04.cor.14-2-2}
\expandafter\def\csname sga@target@IV.14.2.3\endcsname{ega4.c04.prop.14-2-3}
\expandafter\def\csname sga@target@IV.14.2.4\endcsname{ega4.c04.cor.14-2-4}
\expandafter\def\csname sga@target@IV.14.2.5\endcsname{ega4.c04.cor.14-2-5}
\expandafter\def\csname sga@target@IV.14.2.6\endcsname{ega4.c04.cor.14-2-6}
\expandafter\def\csname sga@target@IV.14.2.7\endcsname{ega4.c04.rem.14-2-7}
\expandafter\def\csname sga@target@IV.14.3\endcsname{ega4.c04.subsec.14-3.key-iv-14-3}
\expandafter\def\csname sga@target@IV.14.3.1\endcsname{ega4.c04.statement.14-3-1}
\expandafter\def\csname sga@target@IV.14.3.10\endcsname{ega4.c04.lem.14-3-10}
\expandafter\def\csname sga@target@IV.14.3.11\endcsname{ega4.c04.prop.14-3-11}
\expandafter\def\csname sga@target@IV.14.3.12\endcsname{ega4.c04.rem.14-3-12}
\expandafter\def\csname sga@target@IV.14.3.13\endcsname{ega4.c04.prop.14-3-13}
\expandafter\def\csname sga@target@IV.14.3.14\endcsname{ega4.c04.rem.14-3-14}
\expandafter\def\csname sga@target@IV.14.3.2\endcsname{ega4.c04.prop.14-3-2}
\expandafter\def\csname sga@target@IV.14.3.3\endcsname{ega4.c04.def.14-3-3}
\expandafter\def\csname sga@target@IV.14.3.3.1\endcsname{ega4.c04.subsec.14-3-3-1}
\expandafter\def\csname sga@target@IV.14.3.4\endcsname{ega4.c04.prop.14-3-4}
\expandafter\def\csname sga@target@IV.14.3.5\endcsname{ega4.c04.prop.14-3-5}
\expandafter\def\csname sga@target@IV.14.3.6\endcsname{ega4.c04.prop.14-3-6}
\expandafter\def\csname sga@target@IV.14.3.7\endcsname{ega4.c04.cor.14-3-7}
\expandafter\def\csname sga@target@IV.14.3.8\endcsname{ega4.c04.prop.14-3-8}
\expandafter\def\csname sga@target@IV.14.3.9\endcsname{ega4.c04.item.14-3-9}
\expandafter\def\csname sga@target@IV.14.4.1\endcsname{ega4.c04.thm.14-4-1}
\expandafter\def\csname sga@target@IV.14.4.1.1\endcsname{ega4.c04.lem.14-4-1-1}
\expandafter\def\csname sga@target@IV.14.4.1.2\endcsname{ega4.c04.lem.14-4-1-2}
\expandafter\def\csname sga@target@IV.14.4.1.3\endcsname{ega4.c04.lem.14-4-1-3}
\expandafter\def\csname sga@target@IV.14.4.10\endcsname{ega4.c04.item.14-4-10}
\expandafter\def\csname sga@target@IV.14.4.2\endcsname{ega4.c04.cor.14-4-2}
\expandafter\def\csname sga@target@IV.14.4.3\endcsname{ega4.c04.cor.14-4-3}
\expandafter\def\csname sga@target@IV.14.4.4\endcsname{ega4.c04.cor.14-4-4}
\expandafter\def\csname sga@target@IV.14.4.5\endcsname{ega4.c04.rem.14-4-5}
\expandafter\def\csname sga@target@IV.14.4.6\endcsname{ega4.c04.cor.14-4-6}
\expandafter\def\csname sga@target@IV.14.4.7\endcsname{ega4.c04.prop.14-4-7}
\expandafter\def\csname sga@target@IV.14.4.7.1\endcsname{ega4.c04.eq.14-4-7-1}
\expandafter\def\csname sga@target@IV.14.4.7.2\endcsname{ega4.c04.eq.14-4-7-2}
\expandafter\def\csname sga@target@IV.14.4.8\endcsname{ega4.c04.prop.14-4-8}
\expandafter\def\csname sga@target@IV.14.4.8.1\endcsname{ega4.c04.lem.14-4-8-1}
\expandafter\def\csname sga@target@IV.14.4.9\endcsname{ega4.c04.cor.14-4-9}
\expandafter\def\csname sga@target@IV.14.5\endcsname{ega4.c04.subsec.14-5.key-iv-14-5}
\expandafter\def\csname sga@target@IV.14.5.1\endcsname{ega4.c04.lem.14-5-1}
\expandafter\def\csname sga@target@IV.14.5.1.1\endcsname{ega4.c04.eq.14-5-1-1}
\expandafter\def\csname sga@target@IV.14.5.1.2\endcsname{ega4.c04.eq.14-5-1-2}
\expandafter\def\csname sga@target@IV.14.5.1.3\endcsname{ega4.c04.eq.14-5-1-3}
\expandafter\def\csname sga@target@IV.14.5.1.4\endcsname{ega4.c04.eq.14-5-1-4}
\expandafter\def\csname sga@target@IV.14.5.10\endcsname{ega4.c04.prop.14-5-10}
\expandafter\def\csname sga@target@IV.14.5.11\endcsname{ega4.c04.cor.14-5-11}
\expandafter\def\csname sga@target@IV.14.5.11.1\endcsname{ega4.c04.lem.14-5-11-1}
\expandafter\def\csname sga@target@IV.14.5.12\endcsname{ega4.c04.cor.14-5-12}
\expandafter\def\csname sga@target@IV.14.5.2\endcsname{ega4.c04.cor.14-5-2}
\expandafter\def\csname sga@target@IV.14.5.2.1\endcsname{ega4.c04.eq.14-5-2-1}
\expandafter\def\csname sga@target@IV.14.5.3\endcsname{ega4.c04.prop.14-5-3}
\expandafter\def\csname sga@target@IV.14.5.4\endcsname{ega4.c04.cor.14-5-4}
\expandafter\def\csname sga@target@IV.14.5.5\endcsname{ega4.c04.rem.14-5-5}
\expandafter\def\csname sga@target@IV.14.5.6\endcsname{ega4.c04.prop.14-5-6}
\expandafter\def\csname sga@target@IV.14.5.7\endcsname{ega4.c04.rem.14-5-7}
\expandafter\def\csname sga@target@IV.14.5.8\endcsname{ega4.c04.prop.14-5-8}
\expandafter\def\csname sga@target@IV.14.5.9\endcsname{ega4.c04.prop.14-5-9}
\expandafter\def\csname sga@target@IV.15.1.1\endcsname{ega4.c04.prop.15-1-1}
\expandafter\def\csname sga@target@IV.15.1.1-proof\endcsname{ega4.c04.statement.15-1-1-proof}
\expandafter\def\csname sga@target@IV.15.1.1.1\endcsname{ega4.c04.eq.15-1-1-1}
\expandafter\def\csname sga@target@IV.15.1.1.2\endcsname{ega4.c04.eq.15-1-1-2}
\expandafter\def\csname sga@target@IV.15.1.1.3\endcsname{ega4.c04.eq.15-1-1-3}
\expandafter\def\csname sga@target@IV.15.1.1.4\endcsname{ega4.c04.eq.15-1-1-4}
\expandafter\def\csname sga@target@IV.15.1.1.5\endcsname{ega4.c04.eq.15-1-1-5}
\expandafter\def\csname sga@target@IV.15.1.1.6\endcsname{ega4.c04.lem.15-1-1-6}
\expandafter\def\csname sga@target@IV.15.1.2\endcsname{ega4.c04.cor.15-1-2}
\expandafter\def\csname sga@target@IV.15.1.3\endcsname{ega4.c04.cor.15-1-3}
\expandafter\def\csname sga@target@IV.15.1.4\endcsname{ega4.c04.statement.15-1-4}
\expandafter\def\csname sga@target@IV.15.2.1\endcsname{ega4.c04.statement.15-2-1}
\expandafter\def\csname sga@target@IV.15.2.2\endcsname{ega4.c04.thm.15-2-2}
\expandafter\def\csname sga@target@IV.15.2.2-proof\endcsname{ega4.c04.item.15-2-2-proof}
\expandafter\def\csname sga@target@IV.15.2.2.1\endcsname{ega4.c04.lem.15-2-2-1}
\expandafter\def\csname sga@target@IV.15.2.3\endcsname{ega4.c04.cor.15-2-3}
\expandafter\def\csname sga@target@IV.15.2.4\endcsname{ega4.c04.item.15-2-4}
\expandafter\def\csname sga@target@IV.15.3.1\endcsname{ega4.c04.prop.15-3-1}
\expandafter\def\csname sga@target@IV.15.3.1-proof\endcsname{ega4.c04.item.15-3-1-proof}
\expandafter\def\csname sga@target@IV.15.3.2\endcsname{ega4.c04.prop.15-3-2}
\expandafter\def\csname sga@target@IV.15.3.2-proof\endcsname{ega4.c04.item.15-3-2-proof}
\expandafter\def\csname sga@target@IV.15.3.3\endcsname{ega4.c04.cor.15-3-3}
\expandafter\def\csname sga@target@IV.15.3.3-proof\endcsname{ega4.c04.item.15-3-3-proof}
\expandafter\def\csname sga@target@IV.15.3.4\endcsname{ega4.c04.item.15-3-4}
\expandafter\def\csname sga@target@IV.15.4.1\endcsname{ega4.c04.prop.15-4-1}
\expandafter\def\csname sga@target@IV.15.4.2\endcsname{ega4.c04.prop.15-4-2}
\expandafter\def\csname sga@target@IV.15.4.3\endcsname{ega4.c04.prop.15-4-3}
\expandafter\def\csname sga@target@IV.15.4.3-proof\endcsname{ega4.c04.item.15-4-3-proof}
\expandafter\def\csname sga@target@IV.15.5.1\endcsname{ega4.c04.prop.15-5-1}
\expandafter\def\csname sga@target@IV.15.5.10\endcsname{ega4.c04.item.15-5-10}
\expandafter\def\csname sga@target@IV.15.5.2\endcsname{ega4.c04.cor.15-5-2}
\expandafter\def\csname sga@target@IV.15.5.2.1\endcsname{ega4.c04.eq.15-5-2-1}
\expandafter\def\csname sga@target@IV.15.5.2.2\endcsname{ega4.c04.lem.15-5-2-2}
\expandafter\def\csname sga@target@IV.15.5.2.3\endcsname{ega4.c04.eq.15-5-2-3}
\expandafter\def\csname sga@target@IV.15.5.2.4\endcsname{ega4.c04.eq.15-5-2-4}
\expandafter\def\csname sga@target@IV.15.5.3\endcsname{ega4.c04.prop.15-5-3}
\expandafter\def\csname sga@target@IV.15.5.4\endcsname{ega4.c04.cor.15-5-4}
\expandafter\def\csname sga@target@IV.15.5.5\endcsname{ega4.c04.statement.15-5-5}
\expandafter\def\csname sga@target@IV.15.5.6\endcsname{ega4.c04.lem.15-5-6}
\expandafter\def\csname sga@target@IV.15.5.6.1\endcsname{ega4.c04.lem.15-5-6-1}
\expandafter\def\csname sga@target@IV.15.5.7\endcsname{ega4.c04.prop.15-5-7}
\expandafter\def\csname sga@target@IV.15.5.7.1\endcsname{ega4.c04.eq.15-5-7-1}
\expandafter\def\csname sga@target@IV.15.5.8\endcsname{ega4.c04.statement.15-5-8}
\expandafter\def\csname sga@target@IV.15.5.9\endcsname{ega4.c04.prop.15-5-9}
\expandafter\def\csname sga@target@IV.15.6.1\endcsname{ega4.c04.prop.15-6-1}
\expandafter\def\csname sga@target@IV.15.6.1.1\endcsname{ega4.c04.eq.15-6-1-1}
\expandafter\def\csname sga@target@IV.15.6.10\endcsname{ega4.c04.item.15-6-10}
\expandafter\def\csname sga@target@IV.15.6.2\endcsname{ega4.c04.cor.15-6-2}
\expandafter\def\csname sga@target@IV.15.6.3\endcsname{ega4.c04.prop.15-6-3}
\expandafter\def\csname sga@target@IV.15.6.4\endcsname{ega4.c04.prop.15-6-4}
\expandafter\def\csname sga@target@IV.15.6.5\endcsname{ega4.c04.cor.15-6-5}
\expandafter\def\csname sga@target@IV.15.6.6\endcsname{ega4.c04.prop.15-6-6}
\expandafter\def\csname sga@target@IV.15.6.7\endcsname{ega4.c04.cor.15-6-7}
\expandafter\def\csname sga@target@IV.15.6.8\endcsname{ega4.c04.prop.15-6-8}
\expandafter\def\csname sga@target@IV.15.6.9\endcsname{ega4.c04.cor.15-6-9}
\expandafter\def\csname sga@target@IV.15.7.1\endcsname{ega4.c04.def.15-7-1}
\expandafter\def\csname sga@target@IV.15.7.10\endcsname{ega4.c04.cor.15-7-10}
\expandafter\def\csname sga@target@IV.15.7.11\endcsname{ega4.c04.cor.15-7-11}
\expandafter\def\csname sga@target@IV.15.7.11.1\endcsname{ega4.c04.lem.15-7-11-1}
\expandafter\def\csname sga@target@IV.15.7.2\endcsname{ega4.c04.prop.15-7-2}
\expandafter\def\csname sga@target@IV.15.7.3\endcsname{ega4.c04.rem.15-7-3}
\expandafter\def\csname sga@target@IV.15.7.4\endcsname{ega4.c04.cor.15-7-4}
\expandafter\def\csname sga@target@IV.15.7.5\endcsname{ega4.c04.cor.15-7-5}
\expandafter\def\csname sga@target@IV.15.7.6\endcsname{ega4.c04.cor.15-7-6}
\expandafter\def\csname sga@target@IV.15.7.7\endcsname{ega4.c04.item.15-7-7}
\expandafter\def\csname sga@target@IV.15.7.8\endcsname{ega4.c04.statement.15-7-8}
\expandafter\def\csname sga@target@IV.15.7.9\endcsname{ega4.c04.prop.15-7-9}
\expandafter\def\csname sga@target@IV.16.1\endcsname{ega4.c04.subsec.16-1}
\expandafter\def\csname sga@target@IV.16.1.1\endcsname{ega4.c04.statement.16-1-1}
\expandafter\def\csname sga@target@IV.16.1.10\endcsname{ega4.c04.cor.16-1-10}
\expandafter\def\csname sga@target@IV.16.1.11\endcsname{ega4.c04.rem.16-1-11}
\expandafter\def\csname sga@target@IV.16.1.2\endcsname{ega4.c04.def.16-1-2}
\expandafter\def\csname sga@target@IV.16.1.2.1\endcsname{ega4.c04.def.16-1-2-1}
\expandafter\def\csname sga@target@IV.16.1.2.2\endcsname{ega4.c04.eq.16-1-2-2}
\expandafter\def\csname sga@target@IV.16.1.3\endcsname{ega4.c04.statement.16-1-3}
\expandafter\def\csname sga@target@IV.16.1.4\endcsname{ega4.c04.statement.16-1-4}
\expandafter\def\csname sga@target@IV.16.1.5\endcsname{ega4.c04.prop.16-1-5}
\expandafter\def\csname sga@target@IV.16.1.6\endcsname{ega4.c04.cor.16-1-6}
\expandafter\def\csname sga@target@IV.16.1.7\endcsname{ega4.c04.cor.16-1-7}
\expandafter\def\csname sga@target@IV.16.1.8\endcsname{ega4.c04.prop.16-1-8}
\expandafter\def\csname sga@target@IV.16.1.9\endcsname{ega4.c04.prop.16-1-9}
\expandafter\def\csname sga@target@IV.16.10\endcsname{ega4.c04.subsec.16-10}
\expandafter\def\csname sga@target@IV.16.10.1\endcsname{ega4.c04.def.16-10-1}
\expandafter\def\csname sga@target@IV.16.10.2\endcsname{ega4.c04.prop.16-10-2}
\expandafter\def\csname sga@target@IV.16.10.3\endcsname{ega4.c04.prop.16-10-3}
\expandafter\def\csname sga@target@IV.16.10.4\endcsname{ega4.c04.prop.16-10-4}
\expandafter\def\csname sga@target@IV.16.10.5\endcsname{ega4.c04.prop.16-10-5}
\expandafter\def\csname sga@target@IV.16.10.6\endcsname{ega4.c04.rem.16-10-6}
\expandafter\def\csname sga@target@IV.16.11\endcsname{ega4.c04.subsec.16-11}
\expandafter\def\csname sga@target@IV.16.11.1\endcsname{ega4.c04.statement.16-11-1}
\expandafter\def\csname sga@target@IV.16.11.1.1\endcsname{ega4.c04.statement.16-11-1-1}
\expandafter\def\csname sga@target@IV.16.11.1.2\endcsname{ega4.c04.statement.16-11-1-2}
\expandafter\def\csname sga@target@IV.16.11.1.3\endcsname{ega4.c04.eq.16-11-1-3}
\expandafter\def\csname sga@target@IV.16.11.1.4\endcsname{ega4.c04.eq.16-11-1-4}
\expandafter\def\csname sga@target@IV.16.11.1.5\endcsname{ega4.c04.eq.16-11-1-5}
\expandafter\def\csname sga@target@IV.16.11.1.6\endcsname{ega4.c04.eq.16-11-1-6}
\expandafter\def\csname sga@target@IV.16.11.1.7\endcsname{ega4.c04.eq.16-11-1-7}
\expandafter\def\csname sga@target@IV.16.11.2\endcsname{ega4.c04.thm.16-11-2}
\expandafter\def\csname sga@target@IV.16.11.2.1\endcsname{ega4.c04.eq.16-11-2-1}
\expandafter\def\csname sga@target@IV.16.11.2.2\endcsname{ega4.c04.eq.16-11-2-2}
\expandafter\def\csname sga@target@IV.16.11.2.3\endcsname{ega4.c04.eq.16-11-2-3}
\expandafter\def\csname sga@target@IV.16.11.3\endcsname{ega4.c04.statement.16-11-3}
\expandafter\def\csname sga@target@IV.16.12\endcsname{ega4.c04.subsec.16-12}
\expandafter\def\csname sga@target@IV.16.12.1\endcsname{ega4.c04.statement.16-12-1}
\expandafter\def\csname sga@target@IV.16.12.2\endcsname{ega4.c04.thm.16-12-2}
\expandafter\def\csname sga@target@IV.16.12.3\endcsname{ega4.c04.statement.16-12-3}
\expandafter\def\csname sga@target@IV.16.2\endcsname{ega4.c04.subsec.16-2}
\expandafter\def\csname sga@target@IV.16.2.1\endcsname{ega4.c04.statement.16-2-1}
\expandafter\def\csname sga@target@IV.16.2.1.1\endcsname{ega4.c04.statement.16-2-1-1}
\expandafter\def\csname sga@target@IV.16.2.1.2\endcsname{ega4.c04.eq.16-2-1-2}
\expandafter\def\csname sga@target@IV.16.2.1.3\endcsname{ega4.c04.eq.16-2-1-3}
\expandafter\def\csname sga@target@IV.16.2.1.4\endcsname{ega4.c04.eq.16-2-1-4}
\expandafter\def\csname sga@target@IV.16.2.2\endcsname{ega4.c04.prop.16-2-2}
\expandafter\def\csname sga@target@IV.16.2.3\endcsname{ega4.c04.cor.16-2-3}
\expandafter\def\csname sga@target@IV.16.2.3.1\endcsname{ega4.c04.eq.16-2-3-1}
\expandafter\def\csname sga@target@IV.16.2.3.2\endcsname{ega4.c04.eq.16-2-3-2}
\expandafter\def\csname sga@target@IV.16.2.4\endcsname{ega4.c04.cor.16-2-4}
\expandafter\def\csname sga@target@IV.16.2.5\endcsname{ega4.c04.statement.16-2-5}
\expandafter\def\csname sga@target@IV.16.2.6\endcsname{ega4.c04.statement.16-2-6}
\expandafter\def\csname sga@target@IV.16.2.6.1\endcsname{ega4.c04.statement.16-2-6-1}
\expandafter\def\csname sga@target@IV.16.2.6.2\endcsname{ega4.c04.eq.16-2-6-2}
\expandafter\def\csname sga@target@IV.16.2.7\endcsname{ega4.c04.prop.16-2-7}
\expandafter\def\csname sga@target@IV.16.2.7.1\endcsname{ega4.c04.prop.16-2-7-1}
\expandafter\def\csname sga@target@IV.16.3\endcsname{ega4.c04.subsec.16-3}
\expandafter\def\csname sga@target@IV.16.3.1\endcsname{ega4.c04.def.16-3-1}
\expandafter\def\csname sga@target@IV.16.3.1.1\endcsname{ega4.c04.eq.16-3-1-1}
\expandafter\def\csname sga@target@IV.16.3.2\endcsname{ega4.c04.statement.16-3-2}
\expandafter\def\csname sga@target@IV.16.3.3\endcsname{ega4.c04.statement.16-3-3}
\expandafter\def\csname sga@target@IV.16.3.3.1\endcsname{ega4.c04.statement.16-3-3-1}
\expandafter\def\csname sga@target@IV.16.3.4\endcsname{ega4.c04.prop.16-3-4}
\expandafter\def\csname sga@target@IV.16.3.5\endcsname{ega4.c04.statement.16-3-5}
\expandafter\def\csname sga@target@IV.16.3.6\endcsname{ega4.c04.def.16-3-6}
\expandafter\def\csname sga@target@IV.16.3.6.1\endcsname{ega4.c04.eq.16-3-6-1}
\expandafter\def\csname sga@target@IV.16.3.7\endcsname{ega4.c04.statement.16-3-7}
\expandafter\def\csname sga@target@IV.16.3.7.1\endcsname{ega4.c04.eq.16-3-7-1}
\expandafter\def\csname sga@target@IV.16.3.8\endcsname{ega4.c04.prop.16-3-8}
\expandafter\def\csname sga@target@IV.16.3.9\endcsname{ega4.c04.prop.16-3-9}
\expandafter\def\csname sga@target@IV.16.4\endcsname{ega4.c04.subsec.16-4}
\expandafter\def\csname sga@target@IV.16.4.1\endcsname{ega4.c04.statement.16-4-1}
\expandafter\def\csname sga@target@IV.16.4.1.1\endcsname{ega4.c04.statement.16-4-1-1}
\expandafter\def\csname sga@target@IV.16.4.1.2\endcsname{ega4.c04.statement.16-4-1-2}
\expandafter\def\csname sga@target@IV.16.4.1.3\endcsname{ega4.c04.eq.16-4-1-3}
\expandafter\def\csname sga@target@IV.16.4.1.4\endcsname{ega4.c04.eq.16-4-1-4}
\expandafter\def\csname sga@target@IV.16.4.1.5\endcsname{ega4.c04.eq.16-4-1-5}
\expandafter\def\csname sga@target@IV.16.4.10\endcsname{ega4.c04.cor.16-4-10}
\expandafter\def\csname sga@target@IV.16.4.11\endcsname{ega4.c04.prop.16-4-11}
\expandafter\def\csname sga@target@IV.16.4.11.1\endcsname{ega4.c04.prop.16-4-11-1}
\expandafter\def\csname sga@target@IV.16.4.11.2\endcsname{ega4.c04.prop.16-4-11-2}
\expandafter\def\csname sga@target@IV.16.4.11.3\endcsname{ega4.c04.eq.16-4-11-3}
\expandafter\def\csname sga@target@IV.16.4.11.4\endcsname{ega4.c04.eq.16-4-11-4}
\expandafter\def\csname sga@target@IV.16.4.12\endcsname{ega4.c04.cor.16-4-12}
\expandafter\def\csname sga@target@IV.16.4.13\endcsname{ega4.c04.cor.16-4-13}
\expandafter\def\csname sga@target@IV.16.4.14\endcsname{ega4.c04.prop.16-4-14}
\expandafter\def\csname sga@target@IV.16.4.14.1\endcsname{ega4.c04.prop.16-4-14-1}
\expandafter\def\csname sga@target@IV.16.4.15\endcsname{ega4.c04.cor.16-4-15}
\expandafter\def\csname sga@target@IV.16.4.15.1\endcsname{ega4.c04.cor.16-4-15-1}
\expandafter\def\csname sga@target@IV.16.4.15.2\endcsname{ega4.c04.eq.16-4-15-2}
\expandafter\def\csname sga@target@IV.16.4.16\endcsname{ega4.c04.cor.16-4-16}
\expandafter\def\csname sga@target@IV.16.4.16.1\endcsname{ega4.c04.cor.16-4-16-1}
\expandafter\def\csname sga@target@IV.16.4.16.2\endcsname{ega4.c04.eq.16-4-16-2}
\expandafter\def\csname sga@target@IV.16.4.17\endcsname{ega4.c04.cor.16-4-17}
\expandafter\def\csname sga@target@IV.16.4.18\endcsname{ega4.c04.statement.16-4-18}
\expandafter\def\csname sga@target@IV.16.4.18.1\endcsname{ega4.c04.statement.16-4-18-1}
\expandafter\def\csname sga@target@IV.16.4.18.2\endcsname{ega4.c04.statement.16-4-18-2}
\expandafter\def\csname sga@target@IV.16.4.18.3\endcsname{ega4.c04.eq.16-4-18-3}
\expandafter\def\csname sga@target@IV.16.4.19\endcsname{ega4.c04.cor.16-4-19}
\expandafter\def\csname sga@target@IV.16.4.19.1\endcsname{ega4.c04.cor.16-4-19-1}
\expandafter\def\csname sga@target@IV.16.4.2\endcsname{ega4.c04.statement.16-4-2}
\expandafter\def\csname sga@target@IV.16.4.20\endcsname{ega4.c04.prop.16-4-20}
\expandafter\def\csname sga@target@IV.16.4.21\endcsname{ega4.c04.cor.16-4-21}
\expandafter\def\csname sga@target@IV.16.4.21.1\endcsname{ega4.c04.cor.16-4-21-1}
\expandafter\def\csname sga@target@IV.16.4.22\endcsname{ega4.c04.cor.16-4-22}
\expandafter\def\csname sga@target@IV.16.4.23\endcsname{ega4.c04.prop.16-4-23}
\expandafter\def\csname sga@target@IV.16.4.23.1\endcsname{ega4.c04.prop.16-4-23-1}
\expandafter\def\csname sga@target@IV.16.4.24\endcsname{ega4.c04.statement.16-4-24}
\expandafter\def\csname sga@target@IV.16.4.3\endcsname{ega4.c04.statement.16-4-3}
\expandafter\def\csname sga@target@IV.16.4.3.1\endcsname{ega4.c04.statement.16-4-3-1}
\expandafter\def\csname sga@target@IV.16.4.3.2\endcsname{ega4.c04.statement.16-4-3-2}
\expandafter\def\csname sga@target@IV.16.4.3.3\endcsname{ega4.c04.eq.16-4-3-3}
\expandafter\def\csname sga@target@IV.16.4.3.4\endcsname{ega4.c04.eq.16-4-3-4}
\expandafter\def\csname sga@target@IV.16.4.3.5\endcsname{ega4.c04.eq.16-4-3-5}
\expandafter\def\csname sga@target@IV.16.4.3.6\endcsname{ega4.c04.eq.16-4-3-6}
\expandafter\def\csname sga@target@IV.16.4.3.7\endcsname{ega4.c04.eq.16-4-3-7}
\expandafter\def\csname sga@target@IV.16.4.4\endcsname{ega4.c04.statement.16-4-4}
\expandafter\def\csname sga@target@IV.16.4.5\endcsname{ega4.c04.prop.16-4-5}
\expandafter\def\csname sga@target@IV.16.4.6\endcsname{ega4.c04.cor.16-4-6}
\expandafter\def\csname sga@target@IV.16.4.7\endcsname{ega4.c04.statement.16-4-7}
\expandafter\def\csname sga@target@IV.16.4.7.1\endcsname{ega4.c04.statement.16-4-7-1}
\expandafter\def\csname sga@target@IV.16.4.7.2\endcsname{ega4.c04.eq.16-4-7-2}
\expandafter\def\csname sga@target@IV.16.4.7.3\endcsname{ega4.c04.eq.16-4-7-3}
\expandafter\def\csname sga@target@IV.16.4.8\endcsname{ega4.c04.prop.16-4-8}
\expandafter\def\csname sga@target@IV.16.4.8.1\endcsname{ega4.c04.prop.16-4-8-1}
\expandafter\def\csname sga@target@IV.16.4.8.2\endcsname{ega4.c04.eq.16-4-8-2}
\expandafter\def\csname sga@target@IV.16.4.8.3\endcsname{ega4.c04.lem.16-4-8-3}
\expandafter\def\csname sga@target@IV.16.4.8.4\endcsname{ega4.c04.eq.16-4-8-4}
\expandafter\def\csname sga@target@IV.16.4.9\endcsname{ega4.c04.cor.16-4-9}
\expandafter\def\csname sga@target@IV.16.4.9.1\endcsname{ega4.c04.cor.16-4-9-1}
\expandafter\def\csname sga@target@IV.16.5\endcsname{ega4.c04.subsec.16-5}
\expandafter\def\csname sga@target@IV.16.5.1\endcsname{ega4.c04.statement.16-5-1}
\expandafter\def\csname sga@target@IV.16.5.1.1\endcsname{ega4.c04.eq.16-5-1-1}
\expandafter\def\csname sga@target@IV.16.5.1.2\endcsname{ega4.c04.eq.16-5-1-2}
\expandafter\def\csname sga@target@IV.16.5.10\endcsname{ega4.c04.statement.16-5-10}
\expandafter\def\csname sga@target@IV.16.5.10.1\endcsname{ega4.c04.statement.16-5-10-1}
\expandafter\def\csname sga@target@IV.16.5.10.2\endcsname{ega4.c04.statement.16-5-10-2}
\expandafter\def\csname sga@target@IV.16.5.11\endcsname{ega4.c04.prop.16-5-11}
\expandafter\def\csname sga@target@IV.16.5.12\endcsname{ega4.c04.statement.16-5-12}
\expandafter\def\csname sga@target@IV.16.5.12.1\endcsname{ega4.c04.statement.16-5-12-1}
\expandafter\def\csname sga@target@IV.16.5.12.2\endcsname{ega4.c04.eq.16-5-12-2}
\expandafter\def\csname sga@target@IV.16.5.13\endcsname{ega4.c04.statement.16-5-13}
\expandafter\def\csname sga@target@IV.16.5.13.1\endcsname{ega4.c04.statement.16-5-13-1}
\expandafter\def\csname sga@target@IV.16.5.13.2\endcsname{ega4.c04.eq.16-5-13-2}
\expandafter\def\csname sga@target@IV.16.5.13.3\endcsname{ega4.c04.eq.16-5-13-3}
\expandafter\def\csname sga@target@IV.16.5.13.4\endcsname{ega4.c04.eq.16-5-13-4}
\expandafter\def\csname sga@target@IV.16.5.13.5\endcsname{ega4.c04.eq.16-5-13-5}
\expandafter\def\csname sga@target@IV.16.5.14\endcsname{ega4.c04.statement.16-5-14}
\expandafter\def\csname sga@target@IV.16.5.14.1\endcsname{ega4.c04.statement.16-5-14-1}
\expandafter\def\csname sga@target@IV.16.5.14.2\endcsname{ega4.c04.eq.16-5-14-2}
\expandafter\def\csname sga@target@IV.16.5.14.3\endcsname{ega4.c04.eq.16-5-14-3}
\expandafter\def\csname sga@target@IV.16.5.15\endcsname{ega4.c04.statement.16-5-15}
\expandafter\def\csname sga@target@IV.16.5.16\endcsname{ega4.c04.prop.16-5-16}
\expandafter\def\csname sga@target@IV.16.5.17\endcsname{ega4.c04.prop.16-5-17}
\expandafter\def\csname sga@target@IV.16.5.18\endcsname{ega4.c04.cor.16-5-18}
\expandafter\def\csname sga@target@IV.16.5.2\endcsname{ega4.c04.prop.16-5-2}
\expandafter\def\csname sga@target@IV.16.5.3\endcsname{ega4.c04.prop.16-5-3}
\expandafter\def\csname sga@target@IV.16.5.3.1\endcsname{ega4.c04.eq.16-5-3-1}
\expandafter\def\csname sga@target@IV.16.5.4\endcsname{ega4.c04.statement.16-5-4}
\expandafter\def\csname sga@target@IV.16.5.5\endcsname{ega4.c04.cor.16-5-5}
\expandafter\def\csname sga@target@IV.16.5.5.1\endcsname{ega4.c04.cor.16-5-5-1}
\expandafter\def\csname sga@target@IV.16.5.6\endcsname{ega4.c04.cor.16-5-6}
\expandafter\def\csname sga@target@IV.16.5.7\endcsname{ega4.c04.statement.16-5-7}
\expandafter\def\csname sga@target@IV.16.5.7.1\endcsname{ega4.c04.statement.16-5-7-1}
\expandafter\def\csname sga@target@IV.16.5.8\endcsname{ega4.c04.statement.16-5-8}
\expandafter\def\csname sga@target@IV.16.5.8.1\endcsname{ega4.c04.eq.16-5-8-1}
\expandafter\def\csname sga@target@IV.16.5.8.2\endcsname{ega4.c04.eq.16-5-8-2}
\expandafter\def\csname sga@target@IV.16.5.9\endcsname{ega4.c04.statement.16-5-9}
\expandafter\def\csname sga@target@IV.16.5.9.1\endcsname{ega4.c04.statement.16-5-9-1}
\expandafter\def\csname sga@target@IV.16.6\endcsname{ega4.c04.subsec.16-6}
\expandafter\def\csname sga@target@IV.16.6.1\endcsname{ega4.c04.statement.16-6-1}
\expandafter\def\csname sga@target@IV.16.6.1.1\endcsname{ega4.c04.statement.16-6-1-1}
\expandafter\def\csname sga@target@IV.16.6.1.2\endcsname{ega4.c04.statement.16-6-1-2}
\expandafter\def\csname sga@target@IV.16.6.2\endcsname{ega4.c04.thm.16-6-2}
\expandafter\def\csname sga@target@IV.16.6.2.1\endcsname{ega4.c04.eq.16-6-2-1}
\expandafter\def\csname sga@target@IV.16.6.2.2\endcsname{ega4.c04.eq.16-6-2-2}
\expandafter\def\csname sga@target@IV.16.6.2.3\endcsname{ega4.c04.eq.16-6-2-3}
\expandafter\def\csname sga@target@IV.16.6.2.4\endcsname{ega4.c04.eq.16-6-2-4}
\expandafter\def\csname sga@target@IV.16.6.2.5\endcsname{ega4.c04.eq.16-6-2-5}
\expandafter\def\csname sga@target@IV.16.6.3\endcsname{ega4.c04.def.16-6-3}
\expandafter\def\csname sga@target@IV.16.6.4\endcsname{ega4.c04.prop.16-6-4}
\expandafter\def\csname sga@target@IV.16.6.4.1\endcsname{ega4.c04.prop.16-6-4-1}
\expandafter\def\csname sga@target@IV.16.6.5\endcsname{ega4.c04.statement.16-6-5}
\expandafter\def\csname sga@target@IV.16.6.5.1\endcsname{ega4.c04.statement.16-6-5-1}
\expandafter\def\csname sga@target@IV.16.7\endcsname{ega4.c04.subsec.16-7}
\expandafter\def\csname sga@target@IV.16.7.1\endcsname{ega4.c04.statement.16-7-1}
\expandafter\def\csname sga@target@IV.16.7.1.1\endcsname{ega4.c04.eq.16-7-1-1}
\expandafter\def\csname sga@target@IV.16.7.1.2\endcsname{ega4.c04.eq.16-7-1-2}
\expandafter\def\csname sga@target@IV.16.7.10\endcsname{ega4.c04.rem.16-7-10}
\expandafter\def\csname sga@target@IV.16.7.2\endcsname{ega4.c04.statement.16-7-2}
\expandafter\def\csname sga@target@IV.16.7.2.1\endcsname{ega4.c04.statement.16-7-2-1}
\expandafter\def\csname sga@target@IV.16.7.3\endcsname{ega4.c04.prop.16-7-3}
\expandafter\def\csname sga@target@IV.16.7.4\endcsname{ega4.c04.statement.16-7-4}
\expandafter\def\csname sga@target@IV.16.7.4.1\endcsname{ega4.c04.statement.16-7-4-1}
\expandafter\def\csname sga@target@IV.16.7.5\endcsname{ega4.c04.statement.16-7-5}
\expandafter\def\csname sga@target@IV.16.7.5.1\endcsname{ega4.c04.statement.16-7-5-1}
\expandafter\def\csname sga@target@IV.16.7.5.2\endcsname{ega4.c04.statement.16-7-5-2}
\expandafter\def\csname sga@target@IV.16.7.5.3\endcsname{ega4.c04.statement.16-7-5-3}
\expandafter\def\csname sga@target@IV.16.7.5.4\endcsname{ega4.c04.statement.16-7-5-4}
\expandafter\def\csname sga@target@IV.16.7.6\endcsname{ega4.c04.prop.16-7-6}
\expandafter\def\csname sga@target@IV.16.7.7\endcsname{ega4.c04.statement.16-7-7}
\expandafter\def\csname sga@target@IV.16.7.7.1\endcsname{ega4.c04.eq.16-7-7-1}
\expandafter\def\csname sga@target@IV.16.7.7.2\endcsname{ega4.c04.eq.16-7-7-2}
\expandafter\def\csname sga@target@IV.16.7.8\endcsname{ega4.c04.statement.16-7-8}
\expandafter\def\csname sga@target@IV.16.7.8.1\endcsname{ega4.c04.statement.16-7-8-1}
\expandafter\def\csname sga@target@IV.16.7.8.2\endcsname{ega4.c04.eq.16-7-8-2}
\expandafter\def\csname sga@target@IV.16.7.9\endcsname{ega4.c04.statement.16-7-9}
\expandafter\def\csname sga@target@IV.16.7.9.1\endcsname{ega4.c04.statement.16-7-9-1}
\expandafter\def\csname sga@target@IV.16.8\endcsname{ega4.c04.subsec.16-8}
\expandafter\def\csname sga@target@IV.16.8.1\endcsname{ega4.c04.def.16-8-1}
\expandafter\def\csname sga@target@IV.16.8.10\endcsname{ega4.c04.cor.16-8-10}
\expandafter\def\csname sga@target@IV.16.8.11\endcsname{ega4.c04.statement.16-8-11}
\expandafter\def\csname sga@target@IV.16.8.2\endcsname{ega4.c04.statement.16-8-2}
\expandafter\def\csname sga@target@IV.16.8.3\endcsname{ega4.c04.statement.16-8-3}
\expandafter\def\csname sga@target@IV.16.8.3.1\endcsname{ega4.c04.eq.16-8-3-1}
\expandafter\def\csname sga@target@IV.16.8.4\endcsname{ega4.c04.prop.16-8-4}
\expandafter\def\csname sga@target@IV.16.8.4.1\endcsname{ega4.c04.prop.16-8-4-1}
\expandafter\def\csname sga@target@IV.16.8.5\endcsname{ega4.c04.statement.16-8-5}
\expandafter\def\csname sga@target@IV.16.8.5.1\endcsname{ega4.c04.eq.16-8-5-1}
\expandafter\def\csname sga@target@IV.16.8.5.2\endcsname{ega4.c04.eq.16-8-5-2}
\expandafter\def\csname sga@target@IV.16.8.6\endcsname{ega4.c04.prop.16-8-6}
\expandafter\def\csname sga@target@IV.16.8.7\endcsname{ega4.c04.statement.16-8-7}
\expandafter\def\csname sga@target@IV.16.8.7.1\endcsname{ega4.c04.eq.16-8-7-1}
\expandafter\def\csname sga@target@IV.16.8.8\endcsname{ega4.c04.prop.16-8-8}
\expandafter\def\csname sga@target@IV.16.8.8.1\endcsname{ega4.c04.eq.16-8-8-1}
\expandafter\def\csname sga@target@IV.16.8.8.2\endcsname{ega4.c04.eq.16-8-8-2}
\expandafter\def\csname sga@target@IV.16.8.9\endcsname{ega4.c04.prop.16-8-9}
\expandafter\def\csname sga@target@IV.16.8.9.1\endcsname{ega4.c04.lem.16-8-9-1}
\expandafter\def\csname sga@target@IV.16.8.9.2\endcsname{ega4.c04.lem.16-8-9-2}
\expandafter\def\csname sga@target@IV.16.8.9.3\endcsname{ega4.c04.lem.16-8-9-3}
\expandafter\def\csname sga@target@IV.16.8.9.4\endcsname{ega4.c04.eq.16-8-9-4}
\expandafter\def\csname sga@target@IV.16.9\endcsname{ega4.c04.subsec.16-9}
\expandafter\def\csname sga@target@IV.16.9.1\endcsname{ega4.c04.def.16-9-1}
\expandafter\def\csname sga@target@IV.16.9.10\endcsname{ega4.c04.prop.16-9-10}
\expandafter\def\csname sga@target@IV.16.9.11\endcsname{ega4.c04.cor.16-9-11}
\expandafter\def\csname sga@target@IV.16.9.12\endcsname{ega4.c04.statement.16-9-12}
\expandafter\def\csname sga@target@IV.16.9.13\endcsname{ega4.c04.prop.16-9-13}
\expandafter\def\csname sga@target@IV.16.9.13.1\endcsname{ega4.c04.prop.16-9-13-1}
\expandafter\def\csname sga@target@IV.16.9.13.2\endcsname{ega4.c04.lem.16-9-13-2}
\expandafter\def\csname sga@target@IV.16.9.13.3\endcsname{ega4.c04.lem.16-9-13-3}
\expandafter\def\csname sga@target@IV.16.9.13.4\endcsname{ega4.c04.eq.16-9-13-4}
\expandafter\def\csname sga@target@IV.16.9.13.5\endcsname{ega4.c04.cor.16-9-13-5}
\expandafter\def\csname sga@target@IV.16.9.13.6\endcsname{ega4.c04.cor.16-9-13-6}
\expandafter\def\csname sga@target@IV.16.9.13.7\endcsname{ega4.c04.eq.16-9-13-7}
\expandafter\def\csname sga@target@IV.16.9.2\endcsname{ega4.c04.def.16-9-2}
\expandafter\def\csname sga@target@IV.16.9.3\endcsname{ega4.c04.prop.16-9-3}
\expandafter\def\csname sga@target@IV.16.9.4\endcsname{ega4.c04.cor.16-9-4}
\expandafter\def\csname sga@target@IV.16.9.4.1\endcsname{ega4.c04.eq.16-9-4-1}
\expandafter\def\csname sga@target@IV.16.9.5\endcsname{ega4.c04.cor.16-9-5}
\expandafter\def\csname sga@target@IV.16.9.6\endcsname{ega4.c04.statement.16-9-6}
\expandafter\def\csname sga@target@IV.16.9.7\endcsname{ega4.c04.statement.16-9-7}
\expandafter\def\csname sga@target@IV.16.9.8\endcsname{ega4.c04.prop.16-9-8}
\expandafter\def\csname sga@target@IV.16.9.9\endcsname{ega4.c04.cor.16-9-9}
\expandafter\def\csname sga@target@IV.17.1.1\endcsname{ega4.c04.def.17-1-1}
\expandafter\def\csname sga@target@IV.17.1.1.1\endcsname{ega4.c04.def.17-1-1-1}
\expandafter\def\csname sga@target@IV.17.1.2\endcsname{ega4.c04.rem.17-1-2}
\expandafter\def\csname sga@target@IV.17.1.2.1\endcsname{ega4.c04.eq.17-1-2-1}
\expandafter\def\csname sga@target@IV.17.1.3\endcsname{ega4.c04.prop.17-1-3}
\expandafter\def\csname sga@target@IV.17.1.4\endcsname{ega4.c04.prop.17-1-4}
\expandafter\def\csname sga@target@IV.17.1.5\endcsname{ega4.c04.cor.17-1-5}
\expandafter\def\csname sga@target@IV.17.1.6\endcsname{ega4.c04.prop.17-1-6}
\expandafter\def\csname sga@target@IV.17.10.1\endcsname{ega4.c04.def.17-10-1}
\expandafter\def\csname sga@target@IV.17.10.2\endcsname{ega4.c04.prop.17-10-2}
\expandafter\def\csname sga@target@IV.17.10.3\endcsname{ega4.c04.cor.17-10-3}
\expandafter\def\csname sga@target@IV.17.10.3.1\endcsname{ega4.c04.cor.17-10-3-1}
\expandafter\def\csname sga@target@IV.17.10.4\endcsname{ega4.c04.cor.17-10-4}
\expandafter\def\csname sga@target@IV.17.10.4.1\endcsname{ega4.c04.eq.17-10-4-1}
\expandafter\def\csname sga@target@IV.17.11.1\endcsname{ega4.c04.thm.17-11-1}
\expandafter\def\csname sga@target@IV.17.11.1.1\endcsname{ega4.c04.eq.17-11-1-1}
\expandafter\def\csname sga@target@IV.17.11.1.2\endcsname{ega4.c04.eq.17-11-1-2}
\expandafter\def\csname sga@target@IV.17.11.2\endcsname{ega4.c04.cor.17-11-2}
\expandafter\def\csname sga@target@IV.17.11.3\endcsname{ega4.c04.cor.17-11-3}
\expandafter\def\csname sga@target@IV.17.11.4\endcsname{ega4.c04.cor.17-11-4}
\expandafter\def\csname sga@target@IV.17.11.5\endcsname{ega4.c04.prop.17-11-5}
\expandafter\def\csname sga@target@IV.17.12.1\endcsname{ega4.c04.thm.17-12-1}
\expandafter\def\csname sga@target@IV.17.12.2\endcsname{ega4.c04.cor.17-12-2}
\expandafter\def\csname sga@target@IV.17.12.3\endcsname{ega4.c04.cor.17-12-3}
\expandafter\def\csname sga@target@IV.17.12.4\endcsname{ega4.c04.prop.17-12-4}
\expandafter\def\csname sga@target@IV.17.12.5\endcsname{ega4.c04.prop.17-12-5}
\expandafter\def\csname sga@target@IV.17.12.6\endcsname{ega4.c04.cor.17-12-6}
\expandafter\def\csname sga@target@IV.17.12.7\endcsname{ega4.c04.ex.17-12-7}
\expandafter\def\csname sga@target@IV.17.12.8\endcsname{ega4.c04.rem.17-12-8}
\expandafter\def\csname sga@target@IV.17.13\endcsname{ega4.c04.subsec.17-13.key-iv-17-13}
\expandafter\def\csname sga@target@IV.17.13.1\endcsname{ega4.c04.statement.17-13-1}
\expandafter\def\csname sga@target@IV.17.13.1.1\endcsname{ega4.c04.statement.17-13-1-1}
\expandafter\def\csname sga@target@IV.17.13.1.2\endcsname{ega4.c04.statement.17-13-1-2}
\expandafter\def\csname sga@target@IV.17.13.1.3\endcsname{ega4.c04.eq.17-13-1-3}
\expandafter\def\csname sga@target@IV.17.13.10\endcsname{ega4.c04.cor.17-13-10}
\expandafter\def\csname sga@target@IV.17.13.11\endcsname{ega4.c04.statement.17-13-11}
\expandafter\def\csname sga@target@IV.17.13.11.1\endcsname{ega4.c04.statement.17-13-11-1}
\expandafter\def\csname sga@target@IV.17.13.12\endcsname{ega4.c04.cor.17-13-12}
\expandafter\def\csname sga@target@IV.17.13.12.1\endcsname{ega4.c04.cor.17-13-12-1}
\expandafter\def\csname sga@target@IV.17.13.12.2\endcsname{ega4.c04.eq.17-13-12-2}
\expandafter\def\csname sga@target@IV.17.13.2\endcsname{ega4.c04.prop.17-13-2}
\expandafter\def\csname sga@target@IV.17.13.2.1\endcsname{ega4.c04.eq.17-13-2-1}
\expandafter\def\csname sga@target@IV.17.13.2.2\endcsname{ega4.c04.eq.17-13-2-2}
\expandafter\def\csname sga@target@IV.17.13.3\endcsname{ega4.c04.def.17-13-3}
\expandafter\def\csname sga@target@IV.17.13.4\endcsname{ega4.c04.statement.17-13-4}
\expandafter\def\csname sga@target@IV.17.13.5\endcsname{ega4.c04.statement.17-13-5}
\expandafter\def\csname sga@target@IV.17.13.5.1\endcsname{ega4.c04.statement.17-13-5-1}
\expandafter\def\csname sga@target@IV.17.13.5.2\endcsname{ega4.c04.eq.17-13-5-2}
\expandafter\def\csname sga@target@IV.17.13.5.3\endcsname{ega4.c04.eq.17-13-5-3}
\expandafter\def\csname sga@target@IV.17.13.5.4\endcsname{ega4.c04.eq.17-13-5-4}
\expandafter\def\csname sga@target@IV.17.13.6\endcsname{ega4.c04.cor.17-13-6}
\expandafter\def\csname sga@target@IV.17.13.6.1\endcsname{ega4.c04.eq.17-13-6-1}
\expandafter\def\csname sga@target@IV.17.13.6.2\endcsname{ega4.c04.eq.17-13-6-2}
\expandafter\def\csname sga@target@IV.17.13.7\endcsname{ega4.c04.statement.17-13-7}
\expandafter\def\csname sga@target@IV.17.13.7.1\endcsname{ega4.c04.eq.17-13-7-1}
\expandafter\def\csname sga@target@IV.17.13.8\endcsname{ega4.c04.prop.17-13-8}
\expandafter\def\csname sga@target@IV.17.13.8.1\endcsname{ega4.c04.eq.17-13-8-1}
\expandafter\def\csname sga@target@IV.17.13.9\endcsname{ega4.c04.statement.17-13-9}
\expandafter\def\csname sga@target@IV.17.13.9.1\endcsname{ega4.c04.eq.17-13-9-1}
\expandafter\def\csname sga@target@IV.17.13.9.2\endcsname{ega4.c04.eq.17-13-9-2}
\expandafter\def\csname sga@target@IV.17.14.1\endcsname{ega4.c04.prop.17-14-1}
\expandafter\def\csname sga@target@IV.17.14.2\endcsname{ega4.c04.prop.17-14-2}
\expandafter\def\csname sga@target@IV.17.14.3\endcsname{ega4.c04.rem.17-14-3}
\expandafter\def\csname sga@target@IV.17.15.1\endcsname{ega4.c04.prop.17-15-1}
\expandafter\def\csname sga@target@IV.17.15.10\endcsname{ega4.c04.statement.17-15-10}
\expandafter\def\csname sga@target@IV.17.15.11\endcsname{ega4.c04.prop.17-15-11}
\expandafter\def\csname sga@target@IV.17.15.12\endcsname{ega4.c04.prop.17-15-12}
\expandafter\def\csname sga@target@IV.17.15.13\endcsname{ega4.c04.cor.17-15-13}
\expandafter\def\csname sga@target@IV.17.15.14\endcsname{ega4.c04.prop.17-15-14}
\expandafter\def\csname sga@target@IV.17.15.14.1\endcsname{ega4.c04.lem.17-15-14-1}
\expandafter\def\csname sga@target@IV.17.15.14.2\endcsname{ega4.c04.lem.17-15-14-2}
\expandafter\def\csname sga@target@IV.17.15.15\endcsname{ega4.c04.prop.17-15-15}
\expandafter\def\csname sga@target@IV.17.15.2\endcsname{ega4.c04.cor.17-15-2}
\expandafter\def\csname sga@target@IV.17.15.3\endcsname{ega4.c04.prop.17-15-3}
\expandafter\def\csname sga@target@IV.17.15.4\endcsname{ega4.c04.cor.17-15-4}
\expandafter\def\csname sga@target@IV.17.15.5\endcsname{ega4.c04.prop.17-15-5}
\expandafter\def\csname sga@target@IV.17.15.6\endcsname{ega4.c04.cor.17-15-6}
\expandafter\def\csname sga@target@IV.17.15.7\endcsname{ega4.c04.cor.17-15-7}
\expandafter\def\csname sga@target@IV.17.15.8\endcsname{ega4.c04.prop.17-15-8}
\expandafter\def\csname sga@target@IV.17.15.8.1\endcsname{ega4.c04.eq.17-15-8-1}
\expandafter\def\csname sga@target@IV.17.15.8.2\endcsname{ega4.c04.eq.17-15-8-2}
\expandafter\def\csname sga@target@IV.17.15.9\endcsname{ega4.c04.cor.17-15-9}
\expandafter\def\csname sga@target@IV.17.16\endcsname{ega4.c04.subsec.17-16.key-iv-17-16}
\expandafter\def\csname sga@target@IV.17.16.1\endcsname{ega4.c04.prop.17-16-1}
\expandafter\def\csname sga@target@IV.17.16.2\endcsname{ega4.c04.cor.17-16-2}
\expandafter\def\csname sga@target@IV.17.16.3\endcsname{ega4.c04.cor.17-16-3}
\expandafter\def\csname sga@target@IV.17.16.4\endcsname{ega4.c04.prop.17-16-4}
\expandafter\def\csname sga@target@IV.17.16.5\endcsname{ega4.c04.cor.17-16-5}
\expandafter\def\csname sga@target@IV.17.16.6\endcsname{ega4.c04.prop.17-16-6}
\expandafter\def\csname sga@target@IV.17.2.1\endcsname{ega4.c04.prop.17-2-1}
\expandafter\def\csname sga@target@IV.17.2.2\endcsname{ega4.c04.cor.17-2-2}
\expandafter\def\csname sga@target@IV.17.2.3\endcsname{ega4.c04.prop.17-2-3}
\expandafter\def\csname sga@target@IV.17.2.3.1\endcsname{ega4.c04.eq.17-2-3-1}
\expandafter\def\csname sga@target@IV.17.2.4\endcsname{ega4.c04.cor.17-2-4}
\expandafter\def\csname sga@target@IV.17.2.5\endcsname{ega4.c04.prop.17-2-5}
\expandafter\def\csname sga@target@IV.17.2.5.1\endcsname{ega4.c04.prop.17-2-5-1}
\expandafter\def\csname sga@target@IV.17.2.6\endcsname{ega4.c04.prop.17-2-6}
\expandafter\def\csname sga@target@IV.17.3.1\endcsname{ega4.c04.def.17-3-1}
\expandafter\def\csname sga@target@IV.17.3.2\endcsname{ega4.c04.subsec.17-3-2}
\expandafter\def\csname sga@target@IV.17.3.3\endcsname{ega4.c04.prop.17-3-3}
\expandafter\def\csname sga@target@IV.17.3.4\endcsname{ega4.c04.prop.17-3-4}
\expandafter\def\csname sga@target@IV.17.3.5\endcsname{ega4.c04.cor.17-3-5}
\expandafter\def\csname sga@target@IV.17.3.6\endcsname{ega4.c04.prop.17-3-6}
\expandafter\def\csname sga@target@IV.17.3.7\endcsname{ega4.c04.def.17-3-7}
\expandafter\def\csname sga@target@IV.17.3.8\endcsname{ega4.c04.prop.17-3-8}
\expandafter\def\csname sga@target@IV.17.3.9\endcsname{ega4.c04.cor.17-3-9}
\expandafter\def\csname sga@target@IV.17.4.1\endcsname{ega4.c04.thm.17-4-1}
\expandafter\def\csname sga@target@IV.17.4.1.1\endcsname{ega4.c04.eq.17-4-1-1}
\expandafter\def\csname sga@target@IV.17.4.1.2\endcsname{ega4.c04.rem.17-4-1-2}
\expandafter\def\csname sga@target@IV.17.4.10\endcsname{ega4.c04.rem.17-4-10}
\expandafter\def\csname sga@target@IV.17.4.2\endcsname{ega4.c04.cor.17-4-2}
\expandafter\def\csname sga@target@IV.17.4.3\endcsname{ega4.c04.cor.17-4-3}
\expandafter\def\csname sga@target@IV.17.4.4\endcsname{ega4.c04.prop.17-4-4}
\expandafter\def\csname sga@target@IV.17.4.5\endcsname{ega4.c04.statement.17-4-5}
\expandafter\def\csname sga@target@IV.17.4.6\endcsname{ega4.c04.prop.17-4-6}
\expandafter\def\csname sga@target@IV.17.4.7\endcsname{ega4.c04.cor.17-4-7}
\expandafter\def\csname sga@target@IV.17.4.8\endcsname{ega4.c04.cor.17-4-8}
\expandafter\def\csname sga@target@IV.17.4.9\endcsname{ega4.c04.prop.17-4-9}
\expandafter\def\csname sga@target@IV.17.5.1\endcsname{ega4.c04.thm.17-5-1}
\expandafter\def\csname sga@target@IV.17.5.1.1\endcsname{ega4.c04.eq.17-5-1-1}
\expandafter\def\csname sga@target@IV.17.5.1.2\endcsname{ega4.c04.eq.17-5-1-2}
\expandafter\def\csname sga@target@IV.17.5.2\endcsname{ega4.c04.cor.17-5-2}
\expandafter\def\csname sga@target@IV.17.5.3\endcsname{ega4.c04.prop.17-5-3}
\expandafter\def\csname sga@target@IV.17.5.4\endcsname{ega4.c04.rem.17-5-4}
\expandafter\def\csname sga@target@IV.17.5.5\endcsname{ega4.c04.prop.17-5-5}
\expandafter\def\csname sga@target@IV.17.5.6\endcsname{ega4.c04.cor.17-5-6}
\expandafter\def\csname sga@target@IV.17.5.7\endcsname{ega4.c04.prop.17-5-7}
\expandafter\def\csname sga@target@IV.17.5.8\endcsname{ega4.c04.prop.17-5-8}
\expandafter\def\csname sga@target@IV.17.6.1\endcsname{ega4.c04.thm.17-6-1}
\expandafter\def\csname sga@target@IV.17.6.2\endcsname{ega4.c04.cor.17-6-2}
\expandafter\def\csname sga@target@IV.17.6.3\endcsname{ega4.c04.prop.17-6-3}
\expandafter\def\csname sga@target@IV.17.6.4\endcsname{ega4.c04.prop.17-6-4}
\expandafter\def\csname sga@target@IV.17.7\endcsname{ega4.c04.subsec.17-7.key-iv-17-7}
\expandafter\def\csname sga@target@IV.17.7.1\endcsname{ega4.c04.prop.17-7-1}
\expandafter\def\csname sga@target@IV.17.7.10\endcsname{ega4.c04.prop.17-7-10}
\expandafter\def\csname sga@target@IV.17.7.11\endcsname{ega4.c04.prop.17-7-11}
\expandafter\def\csname sga@target@IV.17.7.11.1\endcsname{ega4.c04.lem.17-7-11-1}
\expandafter\def\csname sga@target@IV.17.7.2\endcsname{ega4.c04.cor.17-7-2}
\expandafter\def\csname sga@target@IV.17.7.3\endcsname{ega4.c04.cor.17-7-3}
\expandafter\def\csname sga@target@IV.17.7.4\endcsname{ega4.c04.prop.17-7-4}
\expandafter\def\csname sga@target@IV.17.7.5\endcsname{ega4.c04.lem.17-7-5}
\expandafter\def\csname sga@target@IV.17.7.6\endcsname{ega4.c04.rem.17-7-6}
\expandafter\def\csname sga@target@IV.17.7.7\endcsname{ega4.c04.prop.17-7-7}
\expandafter\def\csname sga@target@IV.17.7.8\endcsname{ega4.c04.prop.17-7-8}
\expandafter\def\csname sga@target@IV.17.7.9\endcsname{ega4.c04.cor.17-7-9}
\expandafter\def\csname sga@target@IV.17.8\endcsname{ega4.c04.subsec.17-8.key-iv-17-8}
\expandafter\def\csname sga@target@IV.17.8.1\endcsname{ega4.c04.prop.17-8-1}
\expandafter\def\csname sga@target@IV.17.8.2\endcsname{ega4.c04.prop.17-8-2}
\expandafter\def\csname sga@target@IV.17.8.3\endcsname{ega4.c04.rem.17-8-3}
\expandafter\def\csname sga@target@IV.17.9\endcsname{ega4.c04.subsec.17-9.key-iv-17-9}
\expandafter\def\csname sga@target@IV.17.9.1\endcsname{ega4.c04.thm.17-9-1}
\expandafter\def\csname sga@target@IV.17.9.2\endcsname{ega4.c04.cor.17-9-2}
\expandafter\def\csname sga@target@IV.17.9.3\endcsname{ega4.c04.cor.17-9-3}
\expandafter\def\csname sga@target@IV.17.9.4\endcsname{ega4.c04.cor.17-9-4}
\expandafter\def\csname sga@target@IV.17.9.5\endcsname{ega4.c04.cor.17-9-5}
\expandafter\def\csname sga@target@IV.17.9.6\endcsname{ega4.c04.prop.17-9-6}
\expandafter\def\csname sga@target@IV.18\endcsname{ega4.c04.sec.18.key-iv-18}
\expandafter\def\csname sga@target@IV.18.1.1\endcsname{ega4.c04.prop.18-1-1}
\expandafter\def\csname sga@target@IV.18.1.1.1\endcsname{ega4.c04.eq.18-1-1-1}
\expandafter\def\csname sga@target@IV.18.1.2\endcsname{ega4.c04.thm.18-1-2}
\expandafter\def\csname sga@target@IV.18.1.3\endcsname{ega4.c04.cor.18-1-3}
\expandafter\def\csname sga@target@IV.18.10.1\endcsname{ega4.c04.thm.18-10-1}
\expandafter\def\csname sga@target@IV.18.10.10\endcsname{ega4.c04.statement.18-10-10}
\expandafter\def\csname sga@target@IV.18.10.11\endcsname{ega4.c04.rem.18-10-11}
\expandafter\def\csname sga@target@IV.18.10.12\endcsname{ega4.c04.cor.18-10-12}
\expandafter\def\csname sga@target@IV.18.10.13\endcsname{ega4.c04.prop.18-10-13}
\expandafter\def\csname sga@target@IV.18.10.14\endcsname{ega4.c04.cor.18-10-14}
\expandafter\def\csname sga@target@IV.18.10.15\endcsname{ega4.c04.statement.18-10-15}
\expandafter\def\csname sga@target@IV.18.10.16\endcsname{ega4.c04.thm.18-10-16}
\expandafter\def\csname sga@target@IV.18.10.17\endcsname{ega4.c04.rem.18-10-17}
\expandafter\def\csname sga@target@IV.18.10.17.1\endcsname{ega4.c04.statement.18-10-17-1}
\expandafter\def\csname sga@target@IV.18.10.17.2\endcsname{ega4.c04.statement.18-10-17-2}
\expandafter\def\csname sga@target@IV.18.10.18\endcsname{ega4.c04.lem.18-10-18}
\expandafter\def\csname sga@target@IV.18.10.19\endcsname{ega4.c04.prop.18-10-19}
\expandafter\def\csname sga@target@IV.18.10.2\endcsname{ega4.c04.rem.18-10-2}
\expandafter\def\csname sga@target@IV.18.10.20\endcsname{ega4.c04.rem.18-10-20}
\expandafter\def\csname sga@target@IV.18.10.21\endcsname{ega4.c04.prop.18-10-21}
\expandafter\def\csname sga@target@IV.18.10.3\endcsname{ega4.c04.cor.18-10-3}
\expandafter\def\csname sga@target@IV.18.10.3.1\endcsname{ega4.c04.rem.18-10-3-1}
\expandafter\def\csname sga@target@IV.18.10.4\endcsname{ega4.c04.cor.18-10-4}
\expandafter\def\csname sga@target@IV.18.10.5\endcsname{ega4.c04.cor.18-10-5}
\expandafter\def\csname sga@target@IV.18.10.6\endcsname{ega4.c04.lem.18-10-6}
\expandafter\def\csname sga@target@IV.18.10.7\endcsname{ega4.c04.prop.18-10-7}
\expandafter\def\csname sga@target@IV.18.10.8\endcsname{ega4.c04.prop.18-10-8}
\expandafter\def\csname sga@target@IV.18.10.9\endcsname{ega4.c04.cor.18-10-9}
\expandafter\def\csname sga@target@IV.18.11.1\endcsname{ega4.c04.lem.18-11-1}
\expandafter\def\csname sga@target@IV.18.11.10\endcsname{ega4.c04.thm.18-11-10}
\expandafter\def\csname sga@target@IV.18.11.11\endcsname{ega4.c04.rem.18-11-11}
\expandafter\def\csname sga@target@IV.18.11.12\endcsname{ega4.c04.cor.18-11-12}
\expandafter\def\csname sga@target@IV.18.11.12.1\endcsname{ega4.c04.lem.18-11-12-1}
\expandafter\def\csname sga@target@IV.18.11.2\endcsname{ega4.c04.lem.18-11-2}
\expandafter\def\csname sga@target@IV.18.11.2.1\endcsname{ega4.c04.eq.18-11-2-1}
\expandafter\def\csname sga@target@IV.18.11.2.2\endcsname{ega4.c04.eq.18-11-2-2}
\expandafter\def\csname sga@target@IV.18.11.3\endcsname{ega4.c04.lem.18-11-3}
\expandafter\def\csname sga@target@IV.18.11.3.1\endcsname{ega4.c04.eq.18-11-3-1}
\expandafter\def\csname sga@target@IV.18.11.3.2\endcsname{ega4.c04.eq.18-11-3-2}
\expandafter\def\csname sga@target@IV.18.11.3.3\endcsname{ega4.c04.lem.18-11-3-3}
\expandafter\def\csname sga@target@IV.18.11.3.4\endcsname{ega4.c04.lem.18-11-3-4}
\expandafter\def\csname sga@target@IV.18.11.3.5\endcsname{ega4.c04.lem.18-11-3-5}
\expandafter\def\csname sga@target@IV.18.11.3.6\endcsname{ega4.c04.eq.18-11-3-6}
\expandafter\def\csname sga@target@IV.18.11.4\endcsname{ega4.c04.lem.18-11-4}
\expandafter\def\csname sga@target@IV.18.11.5\endcsname{ega4.c04.prop.18-11-5}
\expandafter\def\csname sga@target@IV.18.11.5.1\endcsname{ega4.c04.lem.18-11-5-1}
\expandafter\def\csname sga@target@IV.18.11.5.2\endcsname{ega4.c04.lem.18-11-5-2}
\expandafter\def\csname sga@target@IV.18.11.6\endcsname{ega4.c04.rem.18-11-6}
\expandafter\def\csname sga@target@IV.18.11.7\endcsname{ega4.c04.cor.18-11-7}
\expandafter\def\csname sga@target@IV.18.11.8\endcsname{ega4.c04.cor.18-11-8}
\expandafter\def\csname sga@target@IV.18.11.9\endcsname{ega4.c04.cor.18-11-9}
\expandafter\def\csname sga@target@IV.18.12.1\endcsname{ega4.c04.thm.18-12-1}
\expandafter\def\csname sga@target@IV.18.12.10\endcsname{ega4.c04.cor.18-12-10}
\expandafter\def\csname sga@target@IV.18.12.11\endcsname{ega4.c04.cor.18-12-11}
\expandafter\def\csname sga@target@IV.18.12.12\endcsname{ega4.c04.prop.18-12-12}
\expandafter\def\csname sga@target@IV.18.12.13\endcsname{ega4.c04.cor.18-12-13}
\expandafter\def\csname sga@target@IV.18.12.14\endcsname{ega4.c04.rem.18-12-14}
\expandafter\def\csname sga@target@IV.18.12.15\endcsname{ega4.c04.prop.18-12-15}
\expandafter\def\csname sga@target@IV.18.12.16\endcsname{ega4.c04.prop.18-12-16}
\expandafter\def\csname sga@target@IV.18.12.17\endcsname{ega4.c04.cor.18-12-17}
\expandafter\def\csname sga@target@IV.18.12.2\endcsname{ega4.c04.rem.18-12-2}
\expandafter\def\csname sga@target@IV.18.12.3\endcsname{ega4.c04.cor.18-12-3}
\expandafter\def\csname sga@target@IV.18.12.4\endcsname{ega4.c04.cor.18-12-4}
\expandafter\def\csname sga@target@IV.18.12.5\endcsname{ega4.c04.rem.18-12-5}
\expandafter\def\csname sga@target@IV.18.12.6\endcsname{ega4.c04.cor.18-12-6}
\expandafter\def\csname sga@target@IV.18.12.7\endcsname{ega4.c04.prop.18-12-7}
\expandafter\def\csname sga@target@IV.18.12.7.1\endcsname{ega4.c04.lem.18-12-7-1}
\expandafter\def\csname sga@target@IV.18.12.8\endcsname{ega4.c04.prop.18-12-8}
\expandafter\def\csname sga@target@IV.18.12.9\endcsname{ega4.c04.rem.18-12-9}
\expandafter\def\csname sga@target@IV.18.2.1\endcsname{ega4.c04.statement.18-2-1}
\expandafter\def\csname sga@target@IV.18.2.1.1\endcsname{ega4.c04.statement.18-2-1-1}
\expandafter\def\csname sga@target@IV.18.2.1.2\endcsname{ega4.c04.eq.18-2-1-2}
\expandafter\def\csname sga@target@IV.18.2.10\endcsname{ega4.c04.statement.18-2-10}
\expandafter\def\csname sga@target@IV.18.2.2\endcsname{ega4.c04.statement.18-2-2}
\expandafter\def\csname sga@target@IV.18.2.2.1\endcsname{ega4.c04.eq.18-2-2-1}
\expandafter\def\csname sga@target@IV.18.2.2.2\endcsname{ega4.c04.eq.18-2-2-2}
\expandafter\def\csname sga@target@IV.18.2.2.3\endcsname{ega4.c04.eq.18-2-2-3}
\expandafter\def\csname sga@target@IV.18.2.2.4\endcsname{ega4.c04.eq.18-2-2-4}
\expandafter\def\csname sga@target@IV.18.2.3\endcsname{ega4.c04.prop.18-2-3}
\expandafter\def\csname sga@target@IV.18.2.4\endcsname{ega4.c04.cor.18-2-4}
\expandafter\def\csname sga@target@IV.18.2.5\endcsname{ega4.c04.cor.18-2-5}
\expandafter\def\csname sga@target@IV.18.2.6\endcsname{ega4.c04.cor.18-2-6}
\expandafter\def\csname sga@target@IV.18.2.7\endcsname{ega4.c04.item.18-2-7}
\expandafter\def\csname sga@target@IV.18.2.7.1\endcsname{ega4.c04.eq.18-2-7-1}
\expandafter\def\csname sga@target@IV.18.2.7.2\endcsname{ega4.c04.eq.18-2-7-2}
\expandafter\def\csname sga@target@IV.18.2.8\endcsname{ega4.c04.prop.18-2-8}
\expandafter\def\csname sga@target@IV.18.2.9\endcsname{ega4.c04.cor.18-2-9}
\expandafter\def\csname sga@target@IV.18.3.1\endcsname{ega4.c04.prop.18-3-1}
\expandafter\def\csname sga@target@IV.18.3.2\endcsname{ega4.c04.prop.18-3-2}
\expandafter\def\csname sga@target@IV.18.3.2.1\endcsname{ega4.c04.lem.18-3-2-1}
\expandafter\def\csname sga@target@IV.18.3.2.2\endcsname{ega4.c04.eq.18-3-2-2}
\expandafter\def\csname sga@target@IV.18.3.3\endcsname{ega4.c04.statement.18-3-3}
\expandafter\def\csname sga@target@IV.18.3.4\endcsname{ega4.c04.thm.18-3-4}
\expandafter\def\csname sga@target@IV.18.4.1\endcsname{ega4.c04.lem.18-4-1}
\expandafter\def\csname sga@target@IV.18.4.10\endcsname{ega4.c04.prop.18-4-10}
\expandafter\def\csname sga@target@IV.18.4.11\endcsname{ega4.c04.statement.18-4-11}
\expandafter\def\csname sga@target@IV.18.4.12\endcsname{ega4.c04.cor.18-4-12}
\expandafter\def\csname sga@target@IV.18.4.13\endcsname{ega4.c04.cor.18-4-13}
\expandafter\def\csname sga@target@IV.18.4.14\endcsname{ega4.c04.cor.18-4-14}
\expandafter\def\csname sga@target@IV.18.4.2\endcsname{ega4.c04.prop.18-4-2}
\expandafter\def\csname sga@target@IV.18.4.3\endcsname{ega4.c04.cor.18-4-3}
\expandafter\def\csname sga@target@IV.18.4.4\endcsname{ega4.c04.lem.18-4-4}
\expandafter\def\csname sga@target@IV.18.4.5\endcsname{ega4.c04.prop.18-4-5}
\expandafter\def\csname sga@target@IV.18.4.6\endcsname{ega4.c04.thm.18-4-6}
\expandafter\def\csname sga@target@IV.18.4.6.1\endcsname{ega4.c04.lem.18-4-6-1}
\expandafter\def\csname sga@target@IV.18.4.7\endcsname{ega4.c04.cor.18-4-7}
\expandafter\def\csname sga@target@IV.18.4.8\endcsname{ega4.c04.cor.18-4-8}
\expandafter\def\csname sga@target@IV.18.4.9\endcsname{ega4.c04.prop.18-4-9}
\expandafter\def\csname sga@target@IV.18.5.1\endcsname{ega4.c04.statement.18-5-1}
\expandafter\def\csname sga@target@IV.18.5.1.1\endcsname{ega4.c04.eq.18-5-1-1}
\expandafter\def\csname sga@target@IV.18.5.10\endcsname{ega4.c04.prop.18-5-10}
\expandafter\def\csname sga@target@IV.18.5.11\endcsname{ega4.c04.thm.18-5-11}
\expandafter\def\csname sga@target@IV.18.5.12\endcsname{ega4.c04.cor.18-5-12}
\expandafter\def\csname sga@target@IV.18.5.13\endcsname{ega4.c04.rem.18-5-13}
\expandafter\def\csname sga@target@IV.18.5.13.1\endcsname{ega4.c04.lem.18-5-13-1}
\expandafter\def\csname sga@target@IV.18.5.14\endcsname{ega4.c04.prop.18-5-14}
\expandafter\def\csname sga@target@IV.18.5.15\endcsname{ega4.c04.prop.18-5-15}
\expandafter\def\csname sga@target@IV.18.5.16\endcsname{ega4.c04.item.18-5-16}
\expandafter\def\csname sga@target@IV.18.5.17\endcsname{ega4.c04.thm.18-5-17}
\expandafter\def\csname sga@target@IV.18.5.18\endcsname{ega4.c04.rem.18-5-18}
\expandafter\def\csname sga@target@IV.18.5.19\endcsname{ega4.c04.prop.18-5-19}
\expandafter\def\csname sga@target@IV.18.5.2\endcsname{ega4.c04.statement.18-5-2}
\expandafter\def\csname sga@target@IV.18.5.3\endcsname{ega4.c04.lem.18-5-3}
\expandafter\def\csname sga@target@IV.18.5.3.1\endcsname{ega4.c04.lem.18-5-3-1}
\expandafter\def\csname sga@target@IV.18.5.4\endcsname{ega4.c04.prop.18-5-4}
\expandafter\def\csname sga@target@IV.18.5.4.1\endcsname{ega4.c04.eq.18-5-4-1}
\expandafter\def\csname sga@target@IV.18.5.4.2\endcsname{ega4.c04.eq.18-5-4-2}
\expandafter\def\csname sga@target@IV.18.5.4.3\endcsname{ega4.c04.lem.18-5-4-3}
\expandafter\def\csname sga@target@IV.18.5.4.4\endcsname{ega4.c04.rem.18-5-4-4}
\expandafter\def\csname sga@target@IV.18.5.5\endcsname{ega4.c04.def.18-5-5}
\expandafter\def\csname sga@target@IV.18.5.6\endcsname{ega4.c04.prop.18-5-6}
\expandafter\def\csname sga@target@IV.18.5.7\endcsname{ega4.c04.rem.18-5-7}
\expandafter\def\csname sga@target@IV.18.5.8\endcsname{ega4.c04.def.18-5-8}
\expandafter\def\csname sga@target@IV.18.5.9\endcsname{ega4.c04.prop.18-5-9}
\expandafter\def\csname sga@target@IV.18.6.1\endcsname{ega4.c04.statement.18-6-1}
\expandafter\def\csname sga@target@IV.18.6.10\endcsname{ega4.c04.cor.18-6-10}
\expandafter\def\csname sga@target@IV.18.6.11\endcsname{ega4.c04.cor.18-6-11}
\expandafter\def\csname sga@target@IV.18.6.12\endcsname{ega4.c04.prop.18-6-12}
\expandafter\def\csname sga@target@IV.18.6.13\endcsname{ega4.c04.cor.18-6-13}
\expandafter\def\csname sga@target@IV.18.6.14\endcsname{ega4.c04.prop.18-6-14}
\expandafter\def\csname sga@target@IV.18.6.15\endcsname{ega4.c04.cor.18-6-15}
\expandafter\def\csname sga@target@IV.18.6.2\endcsname{ega4.c04.lem.18-6-2}
\expandafter\def\csname sga@target@IV.18.6.3\endcsname{ega4.c04.lem.18-6-3}
\expandafter\def\csname sga@target@IV.18.6.4\endcsname{ega4.c04.statement.18-6-4}
\expandafter\def\csname sga@target@IV.18.6.5\endcsname{ega4.c04.def.18-6-5}
\expandafter\def\csname sga@target@IV.18.6.6\endcsname{ega4.c04.thm.18-6-6}
\expandafter\def\csname sga@target@IV.18.6.7\endcsname{ega4.c04.statement.18-6-7}
\expandafter\def\csname sga@target@IV.18.6.8\endcsname{ega4.c04.prop.18-6-8}
\expandafter\def\csname sga@target@IV.18.6.9\endcsname{ega4.c04.thm.18-6-9}
\expandafter\def\csname sga@target@IV.18.6.9.1\endcsname{ega4.c04.lem.18-6-9-1}
\expandafter\def\csname sga@target@IV.18.6.9.2\endcsname{ega4.c04.lem.18-6-9-2}
\expandafter\def\csname sga@target@IV.18.7.1\endcsname{ega4.c04.statement.18-7-1}
\expandafter\def\csname sga@target@IV.18.7.2\endcsname{ega4.c04.prop.18-7-2}
\expandafter\def\csname sga@target@IV.18.7.3\endcsname{ega4.c04.cor.18-7-3}
\expandafter\def\csname sga@target@IV.18.7.4\endcsname{ega4.c04.cor.18-7-4}
\expandafter\def\csname sga@target@IV.18.7.5\endcsname{ega4.c04.prop.18-7-5}
\expandafter\def\csname sga@target@IV.18.7.5.1\endcsname{ega4.c04.lem.18-7-5-1}
\expandafter\def\csname sga@target@IV.18.7.6\endcsname{ega4.c04.cor.18-7-6}
\expandafter\def\csname sga@target@IV.18.7.7\endcsname{ega4.c04.rem.18-7-7}
\expandafter\def\csname sga@target@IV.18.8.1\endcsname{ega4.c04.prop.18-8-1}
\expandafter\def\csname sga@target@IV.18.8.10\endcsname{ega4.c04.prop.18-8-10}
\expandafter\def\csname sga@target@IV.18.8.11\endcsname{ega4.c04.rem.18-8-11}
\expandafter\def\csname sga@target@IV.18.8.12\endcsname{ega4.c04.prop.18-8-12}
\expandafter\def\csname sga@target@IV.18.8.13\endcsname{ega4.c04.cor.18-8-13}
\expandafter\def\csname sga@target@IV.18.8.14\endcsname{ega4.c04.cor.18-8-14}
\expandafter\def\csname sga@target@IV.18.8.15\endcsname{ega4.c04.prop.18-8-15}
\expandafter\def\csname sga@target@IV.18.8.16\endcsname{ega4.c04.cor.18-8-16}
\expandafter\def\csname sga@target@IV.18.8.17\endcsname{ega4.c04.prop.18-8-17}
\expandafter\def\csname sga@target@IV.18.8.18\endcsname{ega4.c04.prop.18-8-18}
\expandafter\def\csname sga@target@IV.18.8.2\endcsname{ega4.c04.def.18-8-2}
\expandafter\def\csname sga@target@IV.18.8.3\endcsname{ega4.c04.rem.18-8-3}
\expandafter\def\csname sga@target@IV.18.8.4\endcsname{ega4.c04.lem.18-8-4}
\expandafter\def\csname sga@target@IV.18.8.5\endcsname{ega4.c04.lem.18-8-5}
\expandafter\def\csname sga@target@IV.18.8.6\endcsname{ega4.c04.statement.18-8-6}
\expandafter\def\csname sga@target@IV.18.8.6.1\endcsname{ega4.c04.eq.18-8-6-1}
\expandafter\def\csname sga@target@IV.18.8.6.2\endcsname{ega4.c04.eq.18-8-6-2}
\expandafter\def\csname sga@target@IV.18.8.7\endcsname{ega4.c04.def.18-8-7}
\expandafter\def\csname sga@target@IV.18.8.8\endcsname{ega4.c04.prop.18-8-8}
\expandafter\def\csname sga@target@IV.18.8.8.1\endcsname{ega4.c04.statement.18-8-8-1}
\expandafter\def\csname sga@target@IV.18.8.9\endcsname{ega4.c04.statement.18-8-9}
\expandafter\def\csname sga@target@IV.18.9.1\endcsname{ega4.c04.thm.18-9-1}
\expandafter\def\csname sga@target@IV.18.9.10\endcsname{ega4.c04.rem.18-9-10}
\expandafter\def\csname sga@target@IV.18.9.11\endcsname{ega4.c04.cor.18-9-11}
\expandafter\def\csname sga@target@IV.18.9.2\endcsname{ega4.c04.cor.18-9-2}
\expandafter\def\csname sga@target@IV.18.9.3\endcsname{ega4.c04.cor.18-9-3}
\expandafter\def\csname sga@target@IV.18.9.4\endcsname{ega4.c04.cor.18-9-4}
\expandafter\def\csname sga@target@IV.18.9.5\endcsname{ega4.c04.cor.18-9-5}
\expandafter\def\csname sga@target@IV.18.9.5.1\endcsname{ega4.c04.eq.18-9-5-1}
\expandafter\def\csname sga@target@IV.18.9.6\endcsname{ega4.c04.statement.18-9-6}
\expandafter\def\csname sga@target@IV.18.9.7\endcsname{ega4.c04.thm.18-9-7}
\expandafter\def\csname sga@target@IV.18.9.7.1\endcsname{ega4.c04.rem.18-9-7-1}
\expandafter\def\csname sga@target@IV.18.9.7.2\endcsname{ega4.c04.statement.18-9-7-2}
\expandafter\def\csname sga@target@IV.18.9.7.3\endcsname{ega4.c04.lem.18-9-7-3}
\expandafter\def\csname sga@target@IV.18.9.7.4\endcsname{ega4.c04.eq.18-9-7-4}
\expandafter\def\csname sga@target@IV.18.9.7.5\endcsname{ega4.c04.lem.18-9-7-5}
\expandafter\def\csname sga@target@IV.18.9.7.6\endcsname{ega4.c04.rem.18-9-7-6}
\expandafter\def\csname sga@target@IV.18.9.7.7\endcsname{ega4.c04.lem.18-9-7-7}
\expandafter\def\csname sga@target@IV.18.9.7.8\endcsname{ega4.c04.rem.18-9-7-8}
\expandafter\def\csname sga@target@IV.18.9.8\endcsname{ega4.c04.cor.18-9-8}
\expandafter\def\csname sga@target@IV.18.9.9\endcsname{ega4.c04.cor.18-9-9}
\expandafter\def\csname sga@target@IV.19.1.1\endcsname{ega4.c04.prop.19-1-1}
\expandafter\def\csname sga@target@IV.19.1.2\endcsname{ega4.c04.cor.19-1-2}
\expandafter\def\csname sga@target@IV.19.1.3\endcsname{ega4.c04.def.19-1-3}
\expandafter\def\csname sga@target@IV.19.1.4\endcsname{ega4.c04.prop.19-1-4}
\expandafter\def\csname sga@target@IV.19.1.4.1\endcsname{ega4.c04.eq.19-1-4-1}
\expandafter\def\csname sga@target@IV.19.1.5\endcsname{ega4.c04.prop.19-1-5}
\expandafter\def\csname sga@target@IV.19.1.5.1\endcsname{ega4.c04.item.19-1-5-1}
\expandafter\def\csname sga@target@IV.19.1.6\endcsname{ega4.c04.cor.19-1-6}
\expandafter\def\csname sga@target@IV.19.1.6.1\endcsname{ega4.c04.lem.19-1-6-1}
\expandafter\def\csname sga@target@IV.19.1.7\endcsname{ega4.c04.item.19-1-7}
\expandafter\def\csname sga@target@IV.19.1.8\endcsname{ega4.c04.statement.19-1-8}
\expandafter\def\csname sga@target@IV.19.2.1\endcsname{ega4.c04.def.19-2-1}
\expandafter\def\csname sga@target@IV.19.2.2\endcsname{ega4.c04.def.19-2-2}
\expandafter\def\csname sga@target@IV.19.2.3\endcsname{ega4.c04.prop.19-2-3}
\expandafter\def\csname sga@target@IV.19.2.4\endcsname{ega4.c04.prop.19-2-4}
\expandafter\def\csname sga@target@IV.19.2.5\endcsname{ega4.c04.item.19-2-5}
\expandafter\def\csname sga@target@IV.19.2.6\endcsname{ega4.c04.cor.19-2-6}
\expandafter\def\csname sga@target@IV.19.2.7\endcsname{ega4.c04.prop.19-2-7}
\expandafter\def\csname sga@target@IV.19.2.8\endcsname{ega4.c04.prop.19-2-8}
\expandafter\def\csname sga@target@IV.19.2.9\endcsname{ega4.c04.prop.19-2-9}
\expandafter\def\csname sga@target@IV.19.3.1\endcsname{ega4.c04.def.19-3-1}
\expandafter\def\csname sga@target@IV.19.3.10\endcsname{ega4.c04.cor.19-3-10}
\expandafter\def\csname sga@target@IV.19.3.2\endcsname{ega4.c04.prop.19-3-2}
\expandafter\def\csname sga@target@IV.19.3.2.1\endcsname{ega4.c04.lem.19-3-2-1}
\expandafter\def\csname sga@target@IV.19.3.3\endcsname{ega4.c04.cor.19-3-3}
\expandafter\def\csname sga@target@IV.19.3.4\endcsname{ega4.c04.cor.19-3-4}
\expandafter\def\csname sga@target@IV.19.3.5\endcsname{ega4.c04.cor.19-3-5}
\expandafter\def\csname sga@target@IV.19.3.6\endcsname{ega4.c04.def.19-3-6}
\expandafter\def\csname sga@target@IV.19.3.7\endcsname{ega4.c04.prop.19-3-7}
\expandafter\def\csname sga@target@IV.19.3.8\endcsname{ega4.c04.cor.19-3-8}
\expandafter\def\csname sga@target@IV.19.3.9\endcsname{ega4.c04.prop.19-3-9}
\expandafter\def\csname sga@target@IV.19.4.1\endcsname{ega4.c04.statement.19-4-1}
\expandafter\def\csname sga@target@IV.19.4.1.1\endcsname{ega4.c04.subsec.19-4-1-1}
\expandafter\def\csname sga@target@IV.19.4.10\endcsname{ega4.c04.rem.19-4-10}
\expandafter\def\csname sga@target@IV.19.4.11\endcsname{ega4.c04.prop.19-4-11}
\expandafter\def\csname sga@target@IV.19.4.11.1\endcsname{ega4.c04.eq.19-4-11-1}
\expandafter\def\csname sga@target@IV.19.4.2\endcsname{ega4.c04.lem.19-4-2}
\expandafter\def\csname sga@target@IV.19.4.3\endcsname{ega4.c04.cor.19-4-3}
\expandafter\def\csname sga@target@IV.19.4.4\endcsname{ega4.c04.prop.19-4-4}
\expandafter\def\csname sga@target@IV.19.4.5\endcsname{ega4.c04.cor.19-4-5}
\expandafter\def\csname sga@target@IV.19.4.6\endcsname{ega4.c04.prop.19-4-6}
\expandafter\def\csname sga@target@IV.19.4.7\endcsname{ega4.c04.cor.19-4-7}
\expandafter\def\csname sga@target@IV.19.4.8\endcsname{ega4.c04.prop.19-4-8}
\expandafter\def\csname sga@target@IV.19.4.9\endcsname{ega4.c04.cor.19-4-9}
\expandafter\def\csname sga@target@IV.19.5.1\endcsname{ega4.c04.thm.19-5-1}
\expandafter\def\csname sga@target@IV.19.5.1.1\endcsname{ega4.c04.eq.19-5-1-1}
\expandafter\def\csname sga@target@IV.19.5.2\endcsname{ega4.c04.cor.19-5-2}
\expandafter\def\csname sga@target@IV.19.5.3\endcsname{ega4.c04.prop.19-5-3}
\expandafter\def\csname sga@target@IV.19.5.4\endcsname{ega4.c04.statement.19-5-4}
\expandafter\def\csname sga@target@IV.19.5.5\endcsname{ega4.c04.thm.19-5-5}
\expandafter\def\csname sga@target@IV.19.5.5.1\endcsname{ega4.c04.statement.19-5-5-1}
\expandafter\def\csname sga@target@IV.19.5.5.2\endcsname{ega4.c04.statement.19-5-5-2}
\expandafter\def\csname sga@target@IV.19.5.5.3\endcsname{ega4.c04.statement.19-5-5-3}
\expandafter\def\csname sga@target@IV.19.5.5.4\endcsname{ega4.c04.statement.19-5-5-4}
\expandafter\def\csname sga@target@IV.19.5.5.5\endcsname{ega4.c04.statement.19-5-5-5}
\expandafter\def\csname sga@target@IV.19.5.6\endcsname{ega4.c04.statement.19-5-6}
\expandafter\def\csname sga@target@IV.19.6.1\endcsname{ega4.c04.statement.19-6-1}
\expandafter\def\csname sga@target@IV.19.6.1.1\endcsname{ega4.c04.eq.19-6-1-1}
\expandafter\def\csname sga@target@IV.19.6.2\endcsname{ega4.c04.statement.19-6-2}
\expandafter\def\csname sga@target@IV.19.6.2.1\endcsname{ega4.c04.statement.19-6-2-1}
\expandafter\def\csname sga@target@IV.19.6.3\endcsname{ega4.c04.prop.19-6-3}
\expandafter\def\csname sga@target@IV.19.6.4\endcsname{ega4.c04.cor.19-6-4}
\expandafter\def\csname sga@target@IV.19.7\endcsname{ega4.c04.subsec.19-7.key-iv-19-7}
\expandafter\def\csname sga@target@IV.19.7.1\endcsname{ega4.c04.thm.19-7-1}
\expandafter\def\csname sga@target@IV.19.7.1.1\endcsname{ega4.c04.statement.19-7-1-1}
\expandafter\def\csname sga@target@IV.19.7.1.10\endcsname{ega4.c04.statement.19-7-1-10}
\expandafter\def\csname sga@target@IV.19.7.1.11\endcsname{ega4.c04.statement.19-7-1-11}
\expandafter\def\csname sga@target@IV.19.7.1.12\endcsname{ega4.c04.statement.19-7-1-12}
\expandafter\def\csname sga@target@IV.19.7.1.13\endcsname{ega4.c04.statement.19-7-1-13}
\expandafter\def\csname sga@target@IV.19.7.1.14\endcsname{ega4.c04.statement.19-7-1-14}
\expandafter\def\csname sga@target@IV.19.7.1.15\endcsname{ega4.c04.statement.19-7-1-15}
\expandafter\def\csname sga@target@IV.19.7.1.16\endcsname{ega4.c04.statement.19-7-1-16}
\expandafter\def\csname sga@target@IV.19.7.1.17\endcsname{ega4.c04.statement.19-7-1-17}
\expandafter\def\csname sga@target@IV.19.7.1.18\endcsname{ega4.c04.statement.19-7-1-18}
\expandafter\def\csname sga@target@IV.19.7.1.19\endcsname{ega4.c04.statement.19-7-1-19}
\expandafter\def\csname sga@target@IV.19.7.1.2\endcsname{ega4.c04.statement.19-7-1-2}
\expandafter\def\csname sga@target@IV.19.7.1.20\endcsname{ega4.c04.statement.19-7-1-20}
\expandafter\def\csname sga@target@IV.19.7.1.21\endcsname{ega4.c04.statement.19-7-1-21}
\expandafter\def\csname sga@target@IV.19.7.1.22\endcsname{ega4.c04.statement.19-7-1-22}
\expandafter\def\csname sga@target@IV.19.7.1.23\endcsname{ega4.c04.statement.19-7-1-23}
\expandafter\def\csname sga@target@IV.19.7.1.3\endcsname{ega4.c04.statement.19-7-1-3}
\expandafter\def\csname sga@target@IV.19.7.1.5\endcsname{ega4.c04.statement.19-7-1-5}
\expandafter\def\csname sga@target@IV.19.7.1.6\endcsname{ega4.c04.statement.19-7-1-6}
\expandafter\def\csname sga@target@IV.19.7.1.7\endcsname{ega4.c04.lem.19-7-1-7}
\expandafter\def\csname sga@target@IV.19.7.1.8\endcsname{ega4.c04.statement.19-7-1-8}
\expandafter\def\csname sga@target@IV.19.7.1.9\endcsname{ega4.c04.statement.19-7-1-9}
\expandafter\def\csname sga@target@IV.19.7.2\endcsname{ega4.c04.statement.19-7-2}
\expandafter\def\csname sga@target@IV.19.7.3\endcsname{ega4.c04.cor.19-7-3}
\expandafter\def\csname sga@target@IV.19.7.3.1\endcsname{ega4.c04.eq.19-7-3-1}
\expandafter\def\csname sga@target@IV.19.7.4\endcsname{ega4.c04.cor.19-7-4}
\expandafter\def\csname sga@target@IV.19.7.4.1\endcsname{ega4.c04.eq.19-7-4-1}
\expandafter\def\csname sga@target@IV.19.7.5\endcsname{ega4.c04.cor.19-7-5}
\expandafter\def\csname sga@target@IV.19.8.1\endcsname{ega4.c04.prop.19-8-1}
\expandafter\def\csname sga@target@IV.19.8.2\endcsname{ega4.c04.prop.19-8-2}
\expandafter\def\csname sga@target@IV.19.8.3\endcsname{ega4.c04.rem.19-8-3}
\expandafter\def\csname sga@target@IV.19.9.1\endcsname{ega4.c04.statement.19-9-1}
\expandafter\def\csname sga@target@IV.19.9.1.1\endcsname{ega4.c04.eq.19-9-1-1}
\expandafter\def\csname sga@target@IV.19.9.1.2\endcsname{ega4.c04.eq.19-9-1-2}
\expandafter\def\csname sga@target@IV.19.9.1.3\endcsname{ega4.c04.eq.19-9-1-3}
\expandafter\def\csname sga@target@IV.19.9.2\endcsname{ega4.c04.lem.19-9-2}
\expandafter\def\csname sga@target@IV.19.9.3\endcsname{ega4.c04.prop.19-9-3}
\expandafter\def\csname sga@target@IV.19.9.4\endcsname{ega4.c04.cor.19-9-4}
\expandafter\def\csname sga@target@IV.19.9.5\endcsname{ega4.c04.prop.19-9-5}
\expandafter\def\csname sga@target@IV.19.9.6\endcsname{ega4.c04.prop.19-9-6}
\expandafter\def\csname sga@target@IV.19.9.7\endcsname{ega4.c04.cor.19-9-7}
\expandafter\def\csname sga@target@IV.19.9.8\endcsname{ega4.c04.prop.19-9-8}
\expandafter\def\csname sga@target@IV.19.9.9\endcsname{ega4.c04.rem.19-9-9}
\expandafter\def\csname sga@target@IV.2.1.1\endcsname{ega4.c04.statement.2-1-1}
\expandafter\def\csname sga@target@IV.2.1.10\endcsname{ega4.c04.cor.2-1-10}
\expandafter\def\csname sga@target@IV.2.1.11\endcsname{ega4.c04.prop.2-1-11}
\expandafter\def\csname sga@target@IV.2.1.12\endcsname{ega4.c04.prop.2-1-12}
\expandafter\def\csname sga@target@IV.2.1.13\endcsname{ega4.c04.prop.2-1-13}
\expandafter\def\csname sga@target@IV.2.1.14\endcsname{ega4.c04.prop.2-1-14}
\expandafter\def\csname sga@target@IV.2.1.2\endcsname{ega4.c04.prop.2-1-2}
\expandafter\def\csname sga@target@IV.2.1.3\endcsname{ega4.c04.prop.2-1-3}
\expandafter\def\csname sga@target@IV.2.1.4\endcsname{ega4.c04.prop.2-1-4}
\expandafter\def\csname sga@target@IV.2.1.5\endcsname{ega4.c04.prop.2-1-5}
\expandafter\def\csname sga@target@IV.2.1.6\endcsname{ega4.c04.cor.2-1-6}
\expandafter\def\csname sga@target@IV.2.1.7\endcsname{ega4.c04.cor.2-1-7}
\expandafter\def\csname sga@target@IV.2.1.8\endcsname{ega4.c04.prop.2-1-8}
\expandafter\def\csname sga@target@IV.2.1.9\endcsname{ega4.c04.cor.2-1-9}
\expandafter\def\csname sga@target@IV.2.2.1\endcsname{ega4.c04.prop.2-2-1}
\expandafter\def\csname sga@target@IV.2.2.10\endcsname{ega4.c04.prop.2-2-10}
\expandafter\def\csname sga@target@IV.2.2.10.1\endcsname{ega4.c04.eq.2-2-10-1}
\expandafter\def\csname sga@target@IV.2.2.11\endcsname{ega4.c04.cor.2-2-11}
\expandafter\def\csname sga@target@IV.2.2.12\endcsname{ega4.c04.cor.2-2-12}
\expandafter\def\csname sga@target@IV.2.2.13\endcsname{ega4.c04.cor.2-2-13}
\expandafter\def\csname sga@target@IV.2.2.14\endcsname{ega4.c04.prop.2-2-14}
\expandafter\def\csname sga@target@IV.2.2.15\endcsname{ega4.c04.prop.2-2-15}
\expandafter\def\csname sga@target@IV.2.2.16\endcsname{ega4.c04.cor.2-2-16}
\expandafter\def\csname sga@target@IV.2.2.17\endcsname{ega4.c04.prop.2-2-17}
\expandafter\def\csname sga@target@IV.2.2.2\endcsname{ega4.c04.prop.2-2-2}
\expandafter\def\csname sga@target@IV.2.2.2.1\endcsname{ega4.c04.eq.2-2-2-1}
\expandafter\def\csname sga@target@IV.2.2.3\endcsname{ega4.c04.cor.2-2-3}
\expandafter\def\csname sga@target@IV.2.2.4\endcsname{ega4.c04.def.2-2-4}
\expandafter\def\csname sga@target@IV.2.2.5\endcsname{ega4.c04.cor.2-2-5}
\expandafter\def\csname sga@target@IV.2.2.6\endcsname{ega4.c04.cor.2-2-6}
\expandafter\def\csname sga@target@IV.2.2.7\endcsname{ega4.c04.prop.2-2-7}
\expandafter\def\csname sga@target@IV.2.2.8\endcsname{ega4.c04.cor.2-2-8}
\expandafter\def\csname sga@target@IV.2.2.8.1\endcsname{ega4.c04.eq.2-2-8-1}
\expandafter\def\csname sga@target@IV.2.2.9\endcsname{ega4.c04.rem.2-2-9}
\expandafter\def\csname sga@target@IV.2.3.1\endcsname{ega4.c04.lem.2-3-1}
\expandafter\def\csname sga@target@IV.2.3.1.1\endcsname{ega4.c04.eq.2-3-1-1}
\expandafter\def\csname sga@target@IV.2.3.10\endcsname{ega4.c04.thm.2-3-10}
\expandafter\def\csname sga@target@IV.2.3.11\endcsname{ega4.c04.cor.2-3-11}
\expandafter\def\csname sga@target@IV.2.3.12\endcsname{ega4.c04.cor.2-3-12}
\expandafter\def\csname sga@target@IV.2.3.13\endcsname{ega4.c04.cor.2-3-13}
\expandafter\def\csname sga@target@IV.2.3.14\endcsname{ega4.c04.cor.2-3-14}
\expandafter\def\csname sga@target@IV.2.3.15\endcsname{ega4.c04.rem.2-3-15}
\expandafter\def\csname sga@target@IV.2.3.2\endcsname{ega4.c04.prop.2-3-2}
\expandafter\def\csname sga@target@IV.2.3.3\endcsname{ega4.c04.statement.2-3-3}
\expandafter\def\csname sga@target@IV.2.3.4\endcsname{ega4.c04.prop.2-3-4}
\expandafter\def\csname sga@target@IV.2.3.5\endcsname{ega4.c04.cor.2-3-5}
\expandafter\def\csname sga@target@IV.2.3.6\endcsname{ega4.c04.prop.2-3-6}
\expandafter\def\csname sga@target@IV.2.3.7\endcsname{ega4.c04.prop.2-3-7}
\expandafter\def\csname sga@target@IV.2.3.8\endcsname{ega4.c04.cor.2-3-8}
\expandafter\def\csname sga@target@IV.2.3.9\endcsname{ega4.c04.rem.2-3-9}
\expandafter\def\csname sga@target@IV.2.4.1\endcsname{ega4.c04.statement.2-4-1}
\expandafter\def\csname sga@target@IV.2.4.10\endcsname{ega4.c04.cor.2-4-10}
\expandafter\def\csname sga@target@IV.2.4.11\endcsname{ega4.c04.rem.2-4-11}
\expandafter\def\csname sga@target@IV.2.4.2\endcsname{ega4.c04.def.2-4-2}
\expandafter\def\csname sga@target@IV.2.4.3\endcsname{ega4.c04.prop.2-4-3}
\expandafter\def\csname sga@target@IV.2.4.4\endcsname{ega4.c04.prop.2-4-4}
\expandafter\def\csname sga@target@IV.2.4.5\endcsname{ega4.c04.prop.2-4-5}
\expandafter\def\csname sga@target@IV.2.4.6\endcsname{ega4.c04.thm.2-4-6}
\expandafter\def\csname sga@target@IV.2.4.7\endcsname{ega4.c04.cor.2-4-7}
\expandafter\def\csname sga@target@IV.2.4.8\endcsname{ega4.c04.rem.2-4-8}
\expandafter\def\csname sga@target@IV.2.4.9\endcsname{ega4.c04.prop.2-4-9}
\expandafter\def\csname sga@target@IV.2.5.1\endcsname{ega4.c04.prop.2-5-1}
\expandafter\def\csname sga@target@IV.2.5.2\endcsname{ega4.c04.prop.2-5-2}
\expandafter\def\csname sga@target@IV.2.5.3\endcsname{ega4.c04.lem.2-5-3}
\expandafter\def\csname sga@target@IV.2.5.4\endcsname{ega4.c04.rem.2-5-4}
\expandafter\def\csname sga@target@IV.2.5.4.1\endcsname{ega4.c04.lem.2-5-4-1}
\expandafter\def\csname sga@target@IV.2.5.5\endcsname{ega4.c04.prop.2-5-5}
\expandafter\def\csname sga@target@IV.2.5.5.1\endcsname{ega4.c04.statement.2-5-5-1}
\expandafter\def\csname sga@target@IV.2.5.5.2\endcsname{ega4.c04.lem.2-5-5-2}
\expandafter\def\csname sga@target@IV.2.5.5.3\endcsname{ega4.c04.statement.2-5-5-3}
\expandafter\def\csname sga@target@IV.2.5.5.4\endcsname{ega4.c04.rem.2-5-5-4}
\expandafter\def\csname sga@target@IV.2.5.6\endcsname{ega4.c04.lem.2-5-6}
\expandafter\def\csname sga@target@IV.2.5.7\endcsname{ega4.c04.cor.2-5-7}
\expandafter\def\csname sga@target@IV.2.5.8\endcsname{ega4.c04.prop.2-5-8}
\expandafter\def\csname sga@target@IV.2.5.8.1\endcsname{ega4.c04.statement.2-5-8-1}
\expandafter\def\csname sga@target@IV.2.5.8.2\endcsname{ega4.c04.statement.2-5-8-2}
\expandafter\def\csname sga@target@IV.2.5.8.3\endcsname{ega4.c04.statement.2-5-8-3}
\expandafter\def\csname sga@target@IV.2.5.8.4\endcsname{ega4.c04.statement.2-5-8-4}
\expandafter\def\csname sga@target@IV.2.5.8.5\endcsname{ega4.c04.statement.2-5-8-5}
\expandafter\def\csname sga@target@IV.2.5.8.6\endcsname{ega4.c04.statement.2-5-8-6}
\expandafter\def\csname sga@target@IV.2.5.8.7\endcsname{ega4.c04.statement.2-5-8-7}
\expandafter\def\csname sga@target@IV.2.5.8.8\endcsname{ega4.c04.statement.2-5-8-8}
\expandafter\def\csname sga@target@IV.2.5.8.9\endcsname{ega4.c04.statement.2-5-8-9}
\expandafter\def\csname sga@target@IV.2.6.1\endcsname{ega4.c04.prop.2-6-1}
\expandafter\def\csname sga@target@IV.2.6.2\endcsname{ega4.c04.prop.2-6-2}
\expandafter\def\csname sga@target@IV.2.6.3\endcsname{ega4.c04.rem.2-6-3}
\expandafter\def\csname sga@target@IV.2.6.4\endcsname{ega4.c04.cor.2-6-4}
\expandafter\def\csname sga@target@IV.2.7.1\endcsname{ega4.c04.prop.2-7-1}
\expandafter\def\csname sga@target@IV.2.7.1.1\endcsname{ega4.c04.lem.2-7-1-1}
\expandafter\def\csname sga@target@IV.2.7.1.viii\endcsname{ega4.c04.item.2-7-1-viii}
\expandafter\def\csname sga@target@IV.2.7.2\endcsname{ega4.c04.cor.2-7-2}
\expandafter\def\csname sga@target@IV.2.7.3\endcsname{ega4.c04.sec.2-7-3}
\expandafter\def\csname sga@target@IV.2.8.1\endcsname{ega4.c04.prop.2-8-1}
\expandafter\def\csname sga@target@IV.2.8.1.1\endcsname{ega4.c04.lem.2-8-1-1}
\expandafter\def\csname sga@target@IV.2.8.2\endcsname{ega4.c04.cor.2-8-2}
\expandafter\def\csname sga@target@IV.2.8.3\endcsname{ega4.c04.statement.2-8-3}
\expandafter\def\csname sga@target@IV.2.8.4\endcsname{ega4.c04.prop.2-8-4}
\expandafter\def\csname sga@target@IV.2.8.4.1\endcsname{ega4.c04.statement.2-8-4-1}
\expandafter\def\csname sga@target@IV.2.8.5\endcsname{ega4.c04.prop.2-8-5}
\expandafter\def\csname sga@target@IV.2.8.6\endcsname{ega4.c04.cor.2-8-6}
\expandafter\def\csname sga@target@IV.2.8.6.1\endcsname{ega4.c04.statement.2-8-6-1}
\expandafter\def\csname sga@target@IV.20\endcsname{ega4.c04.sec.20}
\expandafter\def\csname sga@target@IV.20.0\endcsname{ega4.c04.subsec.20-0}
\expandafter\def\csname sga@target@IV.20.1\endcsname{ega4.c04.subsec.20-1}
\expandafter\def\csname sga@target@IV.20.1.1\endcsname{ega4.c04.statement.20-1-1}
\expandafter\def\csname sga@target@IV.20.1.1.1\endcsname{ega4.c04.eq.20-1-1-1}
\expandafter\def\csname sga@target@IV.20.1.10\endcsname{ega4.c04.statement.20-1-10}
\expandafter\def\csname sga@target@IV.20.1.11\endcsname{ega4.c04.statement.20-1-11}
\expandafter\def\csname sga@target@IV.20.1.11.1\endcsname{ega4.c04.eq.20-1-11-1}
\expandafter\def\csname sga@target@IV.20.1.11.2\endcsname{ega4.c04.eq.20-1-11-2}
\expandafter\def\csname sga@target@IV.20.1.12\endcsname{ega4.c04.prop.20-1-12}
\expandafter\def\csname sga@target@IV.20.1.13\endcsname{ega4.c04.rem.20-1-13}
\expandafter\def\csname sga@target@IV.20.1.2\endcsname{ega4.c04.statement.20-1-2}
\expandafter\def\csname sga@target@IV.20.1.2.1\endcsname{ega4.c04.statement.20-1-2-1}
\expandafter\def\csname sga@target@IV.20.1.2.2\endcsname{ega4.c04.statement.20-1-2-2}
\expandafter\def\csname sga@target@IV.20.1.3\endcsname{ega4.c04.statement.20-1-3}
\expandafter\def\csname sga@target@IV.20.1.3.1\endcsname{ega4.c04.statement.20-1-3-1}
\expandafter\def\csname sga@target@IV.20.1.4\endcsname{ega4.c04.statement.20-1-4}
\expandafter\def\csname sga@target@IV.20.1.4.1\endcsname{ega4.c04.statement.20-1-4-1}
\expandafter\def\csname sga@target@IV.20.1.5\endcsname{ega4.c04.statement.20-1-5}
\expandafter\def\csname sga@target@IV.20.1.5.1\endcsname{ega4.c04.statement.20-1-5-1}
\expandafter\def\csname sga@target@IV.20.1.6\endcsname{ega4.c04.prop.20-1-6}
\expandafter\def\csname sga@target@IV.20.1.7\endcsname{ega4.c04.statement.20-1-7}
\expandafter\def\csname sga@target@IV.20.1.8\endcsname{ega4.c04.statement.20-1-8}
\expandafter\def\csname sga@target@IV.20.1.9\endcsname{ega4.c04.statement.20-1-9}
\expandafter\def\csname sga@target@IV.20.2\endcsname{ega4.c04.subsec.20-2}
\expandafter\def\csname sga@target@IV.20.2.1\endcsname{ega4.c04.statement.20-2-1}
\expandafter\def\csname sga@target@IV.20.2.10\endcsname{ega4.c04.statement.20-2-10}
\expandafter\def\csname sga@target@IV.20.2.10.1\endcsname{ega4.c04.eq.20-2-10-1}
\expandafter\def\csname sga@target@IV.20.2.10.2\endcsname{ega4.c04.eq.20-2-10-2}
\expandafter\def\csname sga@target@IV.20.2.11\endcsname{ega4.c04.prop.20-2-11}
\expandafter\def\csname sga@target@IV.20.2.11.1\endcsname{ega4.c04.eq.20-2-11-1}
\expandafter\def\csname sga@target@IV.20.2.12\endcsname{ega4.c04.cor.20-2-12}
\expandafter\def\csname sga@target@IV.20.2.13\endcsname{ega4.c04.statement.20-2-13}
\expandafter\def\csname sga@target@IV.20.2.14\endcsname{ega4.c04.statement.20-2-14}
\expandafter\def\csname sga@target@IV.20.2.15\endcsname{ega4.c04.prop.20-2-15}
\expandafter\def\csname sga@target@IV.20.2.16\endcsname{ega4.c04.rem.20-2-16}
\expandafter\def\csname sga@target@IV.20.2.2\endcsname{ega4.c04.statement.20-2-2}
\expandafter\def\csname sga@target@IV.20.2.3\endcsname{ega4.c04.statement.20-2-3}
\expandafter\def\csname sga@target@IV.20.2.3.1\endcsname{ega4.c04.statement.20-2-3-1}
\expandafter\def\csname sga@target@IV.20.2.4\endcsname{ega4.c04.prop.20-2-4}
\expandafter\def\csname sga@target@IV.20.2.5\endcsname{ega4.c04.cor.20-2-5}
\expandafter\def\csname sga@target@IV.20.2.6\endcsname{ega4.c04.cor.20-2-6}
\expandafter\def\csname sga@target@IV.20.2.7\endcsname{ega4.c04.rem.20-2-7}
\expandafter\def\csname sga@target@IV.20.2.8\endcsname{ega4.c04.statement.20-2-8}
\expandafter\def\csname sga@target@IV.20.2.9\endcsname{ega4.c04.lem.20-2-9}
\expandafter\def\csname sga@target@IV.20.3\endcsname{ega4.c04.subsec.20-3}
\expandafter\def\csname sga@target@IV.20.3.1\endcsname{ega4.c04.statement.20-3-1}
\expandafter\def\csname sga@target@IV.20.3.10\endcsname{ega4.c04.lem.20-3-10}
\expandafter\def\csname sga@target@IV.20.3.11\endcsname{ega4.c04.prop.20-3-11}
\expandafter\def\csname sga@target@IV.20.3.11.1\endcsname{ega4.c04.eq.20-3-11-1}
\expandafter\def\csname sga@target@IV.20.3.12\endcsname{ega4.c04.cor.20-3-12}
\expandafter\def\csname sga@target@IV.20.3.13\endcsname{ega4.c04.rem.20-3-13}
\expandafter\def\csname sga@target@IV.20.3.2\endcsname{ega4.c04.statement.20-3-2}
\expandafter\def\csname sga@target@IV.20.3.2.1\endcsname{ega4.c04.eq.20-3-2-1}
\expandafter\def\csname sga@target@IV.20.3.3\endcsname{ega4.c04.statement.20-3-3}
\expandafter\def\csname sga@target@IV.20.3.4\endcsname{ega4.c04.statement.20-3-4}
\expandafter\def\csname sga@target@IV.20.3.5\endcsname{ega4.c04.prop.20-3-5}
\expandafter\def\csname sga@target@IV.20.3.6\endcsname{ega4.c04.statement.20-3-6}
\expandafter\def\csname sga@target@IV.20.3.6.1\endcsname{ega4.c04.eq.20-3-6-1}
\expandafter\def\csname sga@target@IV.20.3.7\endcsname{ega4.c04.prop.20-3-7}
\expandafter\def\csname sga@target@IV.20.3.8\endcsname{ega4.c04.prop.20-3-8}
\expandafter\def\csname sga@target@IV.20.3.8.1\endcsname{ega4.c04.eq.20-3-8-1}
\expandafter\def\csname sga@target@IV.20.3.8.2\endcsname{ega4.c04.lem.20-3-8-2}
\expandafter\def\csname sga@target@IV.20.3.8.3\endcsname{ega4.c04.lem.20-3-8-3}
\expandafter\def\csname sga@target@IV.20.3.8.4\endcsname{ega4.c04.lem.20-3-8-4}
\expandafter\def\csname sga@target@IV.20.3.8.5\endcsname{ega4.c04.sec.20-3-8-5}
\expandafter\def\csname sga@target@IV.20.3.8.6\endcsname{ega4.c04.sec.20-3-8-6}
\expandafter\def\csname sga@target@IV.20.3.8.7\endcsname{ega4.c04.sec.20-3-8-7}
\expandafter\def\csname sga@target@IV.20.3.8.8\endcsname{ega4.c04.lem.20-3-8-8}
\expandafter\def\csname sga@target@IV.20.3.9\endcsname{ega4.c04.cor.20-3-9}
\expandafter\def\csname sga@target@IV.20.4\endcsname{ega4.c04.subsec.20-4}
\expandafter\def\csname sga@target@IV.20.4.1\endcsname{ega4.c04.statement.20-4-1}
\expandafter\def\csname sga@target@IV.20.4.10\endcsname{ega4.c04.cor.20-4-10}
\expandafter\def\csname sga@target@IV.20.4.11\endcsname{ega4.c04.prop.20-4-11}
\expandafter\def\csname sga@target@IV.20.4.12\endcsname{ega4.c04.cor.20-4-12}
\expandafter\def\csname sga@target@IV.20.4.2\endcsname{ega4.c04.statement.20-4-2}
\expandafter\def\csname sga@target@IV.20.4.3\endcsname{ega4.c04.statement.20-4-3}
\expandafter\def\csname sga@target@IV.20.4.4\endcsname{ega4.c04.prop.20-4-4}
\expandafter\def\csname sga@target@IV.20.4.5\endcsname{ega4.c04.prop.20-4-5}
\expandafter\def\csname sga@target@IV.20.4.6\endcsname{ega4.c04.prop.20-4-6}
\expandafter\def\csname sga@target@IV.20.4.7\endcsname{ega4.c04.cor.20-4-7}
\expandafter\def\csname sga@target@IV.20.4.8\endcsname{ega4.c04.rem.20-4-8}
\expandafter\def\csname sga@target@IV.20.4.9\endcsname{ega4.c04.cor.20-4-9}
\expandafter\def\csname sga@target@IV.20.5\endcsname{ega4.c04.subsec.20-5}
\expandafter\def\csname sga@target@IV.20.5.1\endcsname{ega4.c04.statement.20-5-1}
\expandafter\def\csname sga@target@IV.20.5.2\endcsname{ega4.c04.statement.20-5-2}
\expandafter\def\csname sga@target@IV.20.5.2.1\endcsname{ega4.c04.statement.20-5-2-1}
\expandafter\def\csname sga@target@IV.20.5.3\endcsname{ega4.c04.prop.20-5-3}
\expandafter\def\csname sga@target@IV.20.5.4\endcsname{ega4.c04.statement.20-5-4}
\expandafter\def\csname sga@target@IV.20.5.5\endcsname{ega4.c04.statement.20-5-5}
\expandafter\def\csname sga@target@IV.20.5.6\endcsname{ega4.c04.statement.20-5-6}
\expandafter\def\csname sga@target@IV.20.5.7\endcsname{ega4.c04.statement.20-5-7}
\expandafter\def\csname sga@target@IV.20.5.8\endcsname{ega4.c04.statement.20-5-8}
\expandafter\def\csname sga@target@IV.20.6\endcsname{ega4.c04.subsec.20-6}
\expandafter\def\csname sga@target@IV.20.6.1\endcsname{ega4.c04.statement.20-6-1}
\expandafter\def\csname sga@target@IV.20.6.1.1\endcsname{ega4.c04.eq.20-6-1-1}
\expandafter\def\csname sga@target@IV.20.6.2\endcsname{ega4.c04.statement.20-6-2}
\expandafter\def\csname sga@target@IV.20.6.2.1\endcsname{ega4.c04.statement.20-6-2-1}
\expandafter\def\csname sga@target@IV.20.6.2.2\endcsname{ega4.c04.eq.20-6-2-2}
\expandafter\def\csname sga@target@IV.20.6.3\endcsname{ega4.c04.prop.20-6-3}
\expandafter\def\csname sga@target@IV.20.6.4\endcsname{ega4.c04.prop.20-6-4}
\expandafter\def\csname sga@target@IV.20.6.5\endcsname{ega4.c04.statement.20-6-5}
\expandafter\def\csname sga@target@IV.20.6.5.1\endcsname{ega4.c04.eq.20-6-5-1}
\expandafter\def\csname sga@target@IV.20.6.6\endcsname{ega4.c04.prop.20-6-6}
\expandafter\def\csname sga@target@IV.20.6.7\endcsname{ega4.c04.statement.20-6-7}
\expandafter\def\csname sga@target@IV.20.6.8\endcsname{ega4.c04.prop.20-6-8}
\expandafter\def\csname sga@target@IV.20.6.8.1\endcsname{ega4.c04.eq.20-6-8-1}
\expandafter\def\csname sga@target@IV.20.6.9\endcsname{ega4.c04.statement.20-6-9}
\expandafter\def\csname sga@target@IV.20.6.9.1\endcsname{ega4.c04.statement.20-6-9-1}
\expandafter\def\csname sga@target@IV.21.1\endcsname{ega4.c04.subsec.21-1.key-iv-21-1}
\expandafter\def\csname sga@target@IV.21.1.1\endcsname{ega4.c04.statement.21-1-1}
\expandafter\def\csname sga@target@IV.21.1.10\endcsname{ega4.c04.prop.21-1-10}
\expandafter\def\csname sga@target@IV.21.1.2\endcsname{ega4.c04.statement.21-1-2}
\expandafter\def\csname sga@target@IV.21.1.3\endcsname{ega4.c04.statement.21-1-3}
\expandafter\def\csname sga@target@IV.21.1.3.1\endcsname{ega4.c04.statement.21-1-3-1}
\expandafter\def\csname sga@target@IV.21.1.3.2\endcsname{ega4.c04.statement.21-1-3-2}
\expandafter\def\csname sga@target@IV.21.1.3.3\endcsname{ega4.c04.statement.21-1-3-3}
\expandafter\def\csname sga@target@IV.21.1.3.4\endcsname{ega4.c04.statement.21-1-3-4}
\expandafter\def\csname sga@target@IV.21.1.4\endcsname{ega4.c04.statement.21-1-4}
\expandafter\def\csname sga@target@IV.21.1.4.1\endcsname{ega4.c04.eq.21-1-4-1}
\expandafter\def\csname sga@target@IV.21.1.4.2\endcsname{ega4.c04.eq.21-1-4-2}
\expandafter\def\csname sga@target@IV.21.1.4.3\endcsname{ega4.c04.eq.21-1-4-3}
\expandafter\def\csname sga@target@IV.21.1.5\endcsname{ega4.c04.statement.21-1-5}
\expandafter\def\csname sga@target@IV.21.1.5.1\endcsname{ega4.c04.statement.21-1-5-1}
\expandafter\def\csname sga@target@IV.21.1.5.2\endcsname{ega4.c04.eq.21-1-5-2}
\expandafter\def\csname sga@target@IV.21.1.6\endcsname{ega4.c04.def.21-1-6}
\expandafter\def\csname sga@target@IV.21.1.6.1\endcsname{ega4.c04.eq.21-1-6-1}
\expandafter\def\csname sga@target@IV.21.1.6.2\endcsname{ega4.c04.eq.21-1-6-2}
\expandafter\def\csname sga@target@IV.21.1.6.3\endcsname{ega4.c04.eq.21-1-6-3}
\expandafter\def\csname sga@target@IV.21.1.6.4\endcsname{ega4.c04.eq.21-1-6-4}
\expandafter\def\csname sga@target@IV.21.1.7\endcsname{ega4.c04.statement.21-1-7}
\expandafter\def\csname sga@target@IV.21.1.7.1\endcsname{ega4.c04.statement.21-1-7-1}
\expandafter\def\csname sga@target@IV.21.1.8\endcsname{ega4.c04.prop.21-1-8}
\expandafter\def\csname sga@target@IV.21.1.9\endcsname{ega4.c04.cor.21-1-9}
\expandafter\def\csname sga@target@IV.21.10.1\endcsname{ega4.c04.statement.21-10-1}
\expandafter\def\csname sga@target@IV.21.10.1.1\endcsname{ega4.c04.eq.21-10-1-1}
\expandafter\def\csname sga@target@IV.21.10.1.2\endcsname{ega4.c04.eq.21-10-1-2}
\expandafter\def\csname sga@target@IV.21.10.10\endcsname{ega4.c04.statement.21-10-10}
\expandafter\def\csname sga@target@IV.21.10.11\endcsname{ega4.c04.statement.21-10-11}
\expandafter\def\csname sga@target@IV.21.10.11.1\endcsname{ega4.c04.eq.21-10-11-1}
\expandafter\def\csname sga@target@IV.21.10.12\endcsname{ega4.c04.rem.21-10-12}
\expandafter\def\csname sga@target@IV.21.10.13\endcsname{ega4.c04.lem.21-10-13}
\expandafter\def\csname sga@target@IV.21.10.13.1\endcsname{ega4.c04.prop.21-10-13-1}
\expandafter\def\csname sga@target@IV.21.10.13bis\endcsname{ega4.c04.prop.21-10-13bis}
\expandafter\def\csname sga@target@IV.21.10.14\endcsname{ega4.c04.statement.21-10-14}
\expandafter\def\csname sga@target@IV.21.10.14.1\endcsname{ega4.c04.eq.21-10-14-1}
\expandafter\def\csname sga@target@IV.21.10.14.2\endcsname{ega4.c04.eq.21-10-14-2}
\expandafter\def\csname sga@target@IV.21.10.15\endcsname{ega4.c04.prop.21-10-15}
\expandafter\def\csname sga@target@IV.21.10.16\endcsname{ega4.c04.prop.21-10-16}
\expandafter\def\csname sga@target@IV.21.10.16.1\endcsname{ega4.c04.prop.21-10-16-1}
\expandafter\def\csname sga@target@IV.21.10.16.2\endcsname{ega4.c04.eq.21-10-16-2}
\expandafter\def\csname sga@target@IV.21.10.17\endcsname{ega4.c04.prop.21-10-17}
\expandafter\def\csname sga@target@IV.21.10.17.1\endcsname{ega4.c04.prop.21-10-17-1}
\expandafter\def\csname sga@target@IV.21.10.17.10\endcsname{ega4.c04.eq.21-10-17-10}
\expandafter\def\csname sga@target@IV.21.10.17.2\endcsname{ega4.c04.eq.21-10-17-2}
\expandafter\def\csname sga@target@IV.21.10.17.3\endcsname{ega4.c04.lem.21-10-17-3}
\expandafter\def\csname sga@target@IV.21.10.17.4\endcsname{ega4.c04.lem.21-10-17-4}
\expandafter\def\csname sga@target@IV.21.10.17.5\endcsname{ega4.c04.eq.21-10-17-5}
\expandafter\def\csname sga@target@IV.21.10.17.6\endcsname{ega4.c04.eq.21-10-17-6}
\expandafter\def\csname sga@target@IV.21.10.17.7\endcsname{ega4.c04.lem.21-10-17-7}
\expandafter\def\csname sga@target@IV.21.10.17.8\endcsname{ega4.c04.lem.21-10-17-8}
\expandafter\def\csname sga@target@IV.21.10.17.9\endcsname{ega4.c04.eq.21-10-17-9}
\expandafter\def\csname sga@target@IV.21.10.18\endcsname{ega4.c04.prop.21-10-18}
\expandafter\def\csname sga@target@IV.21.10.18.1\endcsname{ega4.c04.prop.21-10-18-1}
\expandafter\def\csname sga@target@IV.21.10.19\endcsname{ega4.c04.cor.21-10-19}
\expandafter\def\csname sga@target@IV.21.10.19.1\endcsname{ega4.c04.cor.21-10-19-1}
\expandafter\def\csname sga@target@IV.21.10.2\endcsname{ega4.c04.statement.21-10-2}
\expandafter\def\csname sga@target@IV.21.10.3\endcsname{ega4.c04.statement.21-10-3}
\expandafter\def\csname sga@target@IV.21.10.3.1\endcsname{ega4.c04.statement.21-10-3-1}
\expandafter\def\csname sga@target@IV.21.10.4\endcsname{ega4.c04.prop.21-10-4}
\expandafter\def\csname sga@target@IV.21.10.4.1\endcsname{ega4.c04.prop.21-10-4-1}
\expandafter\def\csname sga@target@IV.21.10.5\endcsname{ega4.c04.statement.21-10-5}
\expandafter\def\csname sga@target@IV.21.10.6\endcsname{ega4.c04.prop.21-10-6}
\expandafter\def\csname sga@target@IV.21.10.7\endcsname{ega4.c04.rem.21-10-7}
\expandafter\def\csname sga@target@IV.21.10.8\endcsname{ega4.c04.prop.21-10-8}
\expandafter\def\csname sga@target@IV.21.10.9\endcsname{ega4.c04.statement.21-10-9}
\expandafter\def\csname sga@target@IV.21.11\endcsname{ega4.c04.subsec.21-11.key-iv-21-11}
\expandafter\def\csname sga@target@IV.21.11.1\endcsname{ega4.c04.thm.21-11-1}
\expandafter\def\csname sga@target@IV.21.11.1.1\endcsname{ega4.c04.eq.21-11-1-1}
\expandafter\def\csname sga@target@IV.21.11.2\endcsname{ega4.c04.lem.21-11-2}
\expandafter\def\csname sga@target@IV.21.12\endcsname{ega4.c04.subsec.21-12.key-iv-21-12}
\expandafter\def\csname sga@target@IV.21.12.1\endcsname{ega4.c04.statement.21-12-1}
\expandafter\def\csname sga@target@IV.21.12.1.1\endcsname{ega4.c04.eq.21-12-1-1}
\expandafter\def\csname sga@target@IV.21.12.10\endcsname{ega4.c04.prop.21-12-10}
\expandafter\def\csname sga@target@IV.21.12.12\endcsname{ega4.c04.thm.21-12-12}
\expandafter\def\csname sga@target@IV.21.12.13\endcsname{ega4.c04.cor.21-12-13}
\expandafter\def\csname sga@target@IV.21.12.14\endcsname{ega4.c04.item.21-12-14}
\expandafter\def\csname sga@target@IV.21.12.15\endcsname{ega4.c04.prop.21-12-15}
\expandafter\def\csname sga@target@IV.21.12.16\endcsname{ega4.c04.cor.21-12-16}
\expandafter\def\csname sga@target@IV.21.12.17\endcsname{ega4.c04.item.21-12-17}
\expandafter\def\csname sga@target@IV.21.12.2\endcsname{ega4.c04.prop.21-12-2}
\expandafter\def\csname sga@target@IV.21.12.2.1\endcsname{ega4.c04.prop.21-12-2-1}
\expandafter\def\csname sga@target@IV.21.12.3\endcsname{ega4.c04.cor.21-12-3}
\expandafter\def\csname sga@target@IV.21.12.4\endcsname{ega4.c04.eq.21-12-4}
\expandafter\def\csname sga@target@IV.21.12.5\endcsname{ega4.c04.statement.21-12-5}
\expandafter\def\csname sga@target@IV.21.12.5.1\endcsname{ega4.c04.eq.21-12-5-1}
\expandafter\def\csname sga@target@IV.21.12.5.2\endcsname{ega4.c04.eq.21-12-5-2}
\expandafter\def\csname sga@target@IV.21.12.5.3\endcsname{ega4.c04.eq.21-12-5-3}
\expandafter\def\csname sga@target@IV.21.12.6\endcsname{ega4.c04.thm.21-12-6}
\expandafter\def\csname sga@target@IV.21.12.7\endcsname{ega4.c04.cor.21-12-7}
\expandafter\def\csname sga@target@IV.21.12.8\endcsname{ega4.c04.statement.21-12-8}
\expandafter\def\csname sga@target@IV.21.12.9\endcsname{ega4.c04.statement.21-12-9}
\expandafter\def\csname sga@target@IV.21.13\endcsname{ega4.c04.subsec.21-13.key-iv-21-13}
\expandafter\def\csname sga@target@IV.21.13.1\endcsname{ega4.c04.def.21-13-1}
\expandafter\def\csname sga@target@IV.21.13.10\endcsname{ega4.c04.prop.21-13-10}
\expandafter\def\csname sga@target@IV.21.13.11\endcsname{ega4.c04.cor.21-13-11}
\expandafter\def\csname sga@target@IV.21.13.12\endcsname{ega4.c04.prop.21-13-12}
\expandafter\def\csname sga@target@IV.21.13.12.1\endcsname{ega4.c04.statement.21-13-12-1}
\expandafter\def\csname sga@target@IV.21.13.13\endcsname{ega4.c04.rem.21-13-13}
\expandafter\def\csname sga@target@IV.21.13.14\endcsname{ega4.c04.prop.21-13-14}
\expandafter\def\csname sga@target@IV.21.13.15\endcsname{ega4.c04.cor.21-13-15}
\expandafter\def\csname sga@target@IV.21.13.16\endcsname{ega4.c04.rem.21-13-16}
\expandafter\def\csname sga@target@IV.21.13.2\endcsname{ega4.c04.lem.21-13-2}
\expandafter\def\csname sga@target@IV.21.13.2.1\endcsname{ega4.c04.eq.21-13-2-1}
\expandafter\def\csname sga@target@IV.21.13.3\endcsname{ega4.c04.cor.21-13-3}
\expandafter\def\csname sga@target@IV.21.13.4\endcsname{ega4.c04.lem.21-13-4}
\expandafter\def\csname sga@target@IV.21.13.5\endcsname{ega4.c04.prop.21-13-5}
\expandafter\def\csname sga@target@IV.21.13.6\endcsname{ega4.c04.cor.21-13-6}
\expandafter\def\csname sga@target@IV.21.13.7\endcsname{ega4.c04.def.21-13-7}
\expandafter\def\csname sga@target@IV.21.13.8\endcsname{ega4.c04.prop.21-13-8}
\expandafter\def\csname sga@target@IV.21.13.9\endcsname{ega4.c04.statement.21-13-9}
\expandafter\def\csname sga@target@IV.21.13.9.1\endcsname{ega4.c04.eq.21-13-9-1}
\expandafter\def\csname sga@target@IV.21.14\endcsname{ega4.c04.subsec.21-14.key-iv-21-14}
\expandafter\def\csname sga@target@IV.21.14.1\endcsname{ega4.c04.thm.21-14-1}
\expandafter\def\csname sga@target@IV.21.14.2\endcsname{ega4.c04.cor.21-14-2}
\expandafter\def\csname sga@target@IV.21.14.3\endcsname{ega4.c04.prop.21-14-3}
\expandafter\def\csname sga@target@IV.21.14.4\endcsname{ega4.c04.statement.21-14-4}
\expandafter\def\csname sga@target@IV.21.14.4bis\endcsname{ega4.c04.statement.21-14-4bis}
\expandafter\def\csname sga@target@IV.21.15\endcsname{ega4.c04.subsec.21-15.key-iv-21-15}
\expandafter\def\csname sga@target@IV.21.15.1\endcsname{ega4.c04.statement.21-15-1}
\expandafter\def\csname sga@target@IV.21.15.2\endcsname{ega4.c04.def.21-15-2}
\expandafter\def\csname sga@target@IV.21.15.3\endcsname{ega4.c04.statement.21-15-3}
\expandafter\def\csname sga@target@IV.21.15.3.1\endcsname{ega4.c04.eq.21-15-3-1}
\expandafter\def\csname sga@target@IV.21.15.3.2\endcsname{ega4.c04.eq.21-15-3-2}
\expandafter\def\csname sga@target@IV.21.15.3bis\endcsname{ega4.c04.statement.21-15-3bis}
\expandafter\def\csname sga@target@IV.21.15.4\endcsname{ega4.c04.prop.21-15-4}
\expandafter\def\csname sga@target@IV.21.15.5\endcsname{ega4.c04.prop.21-15-5}
\expandafter\def\csname sga@target@IV.21.15.6\endcsname{ega4.c04.prop.21-15-6}
\expandafter\def\csname sga@target@IV.21.15.7\endcsname{ega4.c04.statement.21-15-7}
\expandafter\def\csname sga@target@IV.21.15.8\endcsname{ega4.c04.statement.21-15-8}
\expandafter\def\csname sga@target@IV.21.15.9\endcsname{ega4.c04.statement.21-15-9}
\expandafter\def\csname sga@target@IV.21.15.9.1\endcsname{ega4.c04.eq.21-15-9-1}
\expandafter\def\csname sga@target@IV.21.2.1\endcsname{ega4.c04.statement.21-2-1}
\expandafter\def\csname sga@target@IV.21.2.10\endcsname{ega4.c04.statement.21-2-10}
\expandafter\def\csname sga@target@IV.21.2.10.1\endcsname{ega4.c04.eq.21-2-10-1}
\expandafter\def\csname sga@target@IV.21.2.11\endcsname{ega4.c04.prop.21-2-11}
\expandafter\def\csname sga@target@IV.21.2.11.1\endcsname{ega4.c04.prop.21-2-11-1}
\expandafter\def\csname sga@target@IV.21.2.12\endcsname{ega4.c04.statement.21-2-12}
\expandafter\def\csname sga@target@IV.21.2.2\endcsname{ega4.c04.prop.21-2-2}
\expandafter\def\csname sga@target@IV.21.2.3\endcsname{ega4.c04.cor.21-2-3}
\expandafter\def\csname sga@target@IV.21.2.4\endcsname{ega4.c04.statement.21-2-4}
\expandafter\def\csname sga@target@IV.21.2.5\endcsname{ega4.c04.statement.21-2-5}
\expandafter\def\csname sga@target@IV.21.2.5.1\endcsname{ega4.c04.statement.21-2-5-1}
\expandafter\def\csname sga@target@IV.21.2.5.2\endcsname{ega4.c04.eq.21-2-5-2}
\expandafter\def\csname sga@target@IV.21.2.5.3\endcsname{ega4.c04.eq.21-2-5-3}
\expandafter\def\csname sga@target@IV.21.2.5.4\endcsname{ega4.c04.eq.21-2-5-4}
\expandafter\def\csname sga@target@IV.21.2.6\endcsname{ega4.c04.prop.21-2-6}
\expandafter\def\csname sga@target@IV.21.2.6.1\endcsname{ega4.c04.eq.21-2-6-1}
\expandafter\def\csname sga@target@IV.21.2.7\endcsname{ega4.c04.statement.21-2-7}
\expandafter\def\csname sga@target@IV.21.2.7.1\endcsname{ega4.c04.statement.21-2-7-1}
\expandafter\def\csname sga@target@IV.21.2.7.2\endcsname{ega4.c04.eq.21-2-7-2}
\expandafter\def\csname sga@target@IV.21.2.8\endcsname{ega4.c04.statement.21-2-8}
\expandafter\def\csname sga@target@IV.21.2.8.1\endcsname{ega4.c04.statement.21-2-8-1}
\expandafter\def\csname sga@target@IV.21.2.8.2\endcsname{ega4.c04.statement.21-2-8-2}
\expandafter\def\csname sga@target@IV.21.2.8.3\endcsname{ega4.c04.statement.21-2-8-3}
\expandafter\def\csname sga@target@IV.21.2.9\endcsname{ega4.c04.statement.21-2-9}
\expandafter\def\csname sga@target@IV.21.2.9.1\endcsname{ega4.c04.statement.21-2-9-1}
\expandafter\def\csname sga@target@IV.21.2.9.2\endcsname{ega4.c04.eq.21-2-9-2}
\expandafter\def\csname sga@target@IV.21.2.9.3\endcsname{ega4.c04.eq.21-2-9-3}
\expandafter\def\csname sga@target@IV.21.2.9.4\endcsname{ega4.c04.eq.21-2-9-4}
\expandafter\def\csname sga@target@IV.21.3.1\endcsname{ega4.c04.statement.21-3-1}
\expandafter\def\csname sga@target@IV.21.3.2\endcsname{ega4.c04.statement.21-3-2}
\expandafter\def\csname sga@target@IV.21.3.2.1\endcsname{ega4.c04.statement.21-3-2-1}
\expandafter\def\csname sga@target@IV.21.3.2.2\endcsname{ega4.c04.eq.21-3-2-2}
\expandafter\def\csname sga@target@IV.21.3.2.3\endcsname{ega4.c04.eq.21-3-2-3}
\expandafter\def\csname sga@target@IV.21.3.2.4\endcsname{ega4.c04.eq.21-3-2-4}
\expandafter\def\csname sga@target@IV.21.3.3\endcsname{ega4.c04.prop.21-3-3}
\expandafter\def\csname sga@target@IV.21.3.3.1\endcsname{ega4.c04.eq.21-3-3-1}
\expandafter\def\csname sga@target@IV.21.3.4\endcsname{ega4.c04.prop.21-3-4}
\expandafter\def\csname sga@target@IV.21.3.5\endcsname{ega4.c04.cor.21-3-5}
\expandafter\def\csname sga@target@IV.21.3.6\endcsname{ega4.c04.rem.21-3-6}
\expandafter\def\csname sga@target@IV.21.3.6.1\endcsname{ega4.c04.eq.21-3-6-1}
\expandafter\def\csname sga@target@IV.21.4\endcsname{ega4.c04.subsec.21-4.key-iv-21-4}
\expandafter\def\csname sga@target@IV.21.4.1\endcsname{ega4.c04.statement.21-4-1}
\expandafter\def\csname sga@target@IV.21.4.10\endcsname{ega4.c04.cor.21-4-10}
\expandafter\def\csname sga@target@IV.21.4.11\endcsname{ega4.c04.cor.21-4-11}
\expandafter\def\csname sga@target@IV.21.4.12\endcsname{ega4.c04.rem.21-4-12}
\expandafter\def\csname sga@target@IV.21.4.13\endcsname{ega4.c04.cor.21-4-13}
\expandafter\def\csname sga@target@IV.21.4.2\endcsname{ega4.c04.def.21-4-2}
\expandafter\def\csname sga@target@IV.21.4.2.1\endcsname{ega4.c04.eq.21-4-2-1}
\expandafter\def\csname sga@target@IV.21.4.3\endcsname{ega4.c04.statement.21-4-3}
\expandafter\def\csname sga@target@IV.21.4.4\endcsname{ega4.c04.statement.21-4-4}
\expandafter\def\csname sga@target@IV.21.4.5\endcsname{ega4.c04.prop.21-4-5}
\expandafter\def\csname sga@target@IV.21.4.6\endcsname{ega4.c04.cor.21-4-6}
\expandafter\def\csname sga@target@IV.21.4.6.1\endcsname{ega4.c04.eq.21-4-6-1}
\expandafter\def\csname sga@target@IV.21.4.7\endcsname{ega4.c04.statement.21-4-7}
\expandafter\def\csname sga@target@IV.21.4.8\endcsname{ega4.c04.prop.21-4-8}
\expandafter\def\csname sga@target@IV.21.4.9\endcsname{ega4.c04.prop.21-4-9}
\expandafter\def\csname sga@target@IV.21.5\endcsname{ega4.c04.subsec.21-5.key-iv-21-5}
\expandafter\def\csname sga@target@IV.21.5.1\endcsname{ega4.c04.statement.21-5-1}
\expandafter\def\csname sga@target@IV.21.5.2\endcsname{ega4.c04.lem.21-5-2}
\expandafter\def\csname sga@target@IV.21.5.3\endcsname{ega4.c04.statement.21-5-3}
\expandafter\def\csname sga@target@IV.21.5.3.1\endcsname{ega4.c04.eq.21-5-3-1}
\expandafter\def\csname sga@target@IV.21.5.3.2\endcsname{ega4.c04.eq.21-5-3-2}
\expandafter\def\csname sga@target@IV.21.5.4\endcsname{ega4.c04.statement.21-5-4}
\expandafter\def\csname sga@target@IV.21.5.5\endcsname{ega4.c04.def.21-5-5}
\expandafter\def\csname sga@target@IV.21.5.5.1\endcsname{ega4.c04.eq.21-5-5-1}
\expandafter\def\csname sga@target@IV.21.5.5.2\endcsname{ega4.c04.eq.21-5-5-2}
\expandafter\def\csname sga@target@IV.21.5.5.3\endcsname{ega4.c04.eq.21-5-5-3}
\expandafter\def\csname sga@target@IV.21.5.6\endcsname{ega4.c04.prop.21-5-6}
\expandafter\def\csname sga@target@IV.21.5.6.1\endcsname{ega4.c04.prop.21-5-6-1}
\expandafter\def\csname sga@target@IV.21.5.7\endcsname{ega4.c04.prop.21-5-7}
\expandafter\def\csname sga@target@IV.21.5.7.1\endcsname{ega4.c04.prop.21-5-7-1}
\expandafter\def\csname sga@target@IV.21.5.7.2\endcsname{ega4.c04.lem.21-5-7-2}
\expandafter\def\csname sga@target@IV.21.5.7.3\endcsname{ega4.c04.lem.21-5-7-3}
\expandafter\def\csname sga@target@IV.21.5.7.4\endcsname{ega4.c04.eq.21-5-7-4}
\expandafter\def\csname sga@target@IV.21.5.8\endcsname{ega4.c04.prop.21-5-8}
\expandafter\def\csname sga@target@IV.21.5.8.1\endcsname{ega4.c04.eq.21-5-8-1}
\expandafter\def\csname sga@target@IV.21.6\endcsname{ega4.c04.subsec.21-6.key-iv-21-6}
\expandafter\def\csname sga@target@IV.21.6.1\endcsname{ega4.c04.statement.21-6-1}
\expandafter\def\csname sga@target@IV.21.6.10\endcsname{ega4.c04.cor.21-6-10}
\expandafter\def\csname sga@target@IV.21.6.10.1\endcsname{ega4.c04.eq.21-6-10-1}
\expandafter\def\csname sga@target@IV.21.6.11\endcsname{ega4.c04.cor.21-6-11}
\expandafter\def\csname sga@target@IV.21.6.12\endcsname{ega4.c04.prop.21-6-12}
\expandafter\def\csname sga@target@IV.21.6.12.1\endcsname{ega4.c04.eq.21-6-12-1}
\expandafter\def\csname sga@target@IV.21.6.12.2\endcsname{ega4.c04.eq.21-6-12-2}
\expandafter\def\csname sga@target@IV.21.6.13\endcsname{ega4.c04.cor.21-6-13}
\expandafter\def\csname sga@target@IV.21.6.14\endcsname{ega4.c04.cor.21-6-14}
\expandafter\def\csname sga@target@IV.21.6.2\endcsname{ega4.c04.statement.21-6-2}
\expandafter\def\csname sga@target@IV.21.6.3\endcsname{ega4.c04.statement.21-6-3}
\expandafter\def\csname sga@target@IV.21.6.4\endcsname{ega4.c04.statement.21-6-4}
\expandafter\def\csname sga@target@IV.21.6.4.1\endcsname{ega4.c04.statement.21-6-4-1}
\expandafter\def\csname sga@target@IV.21.6.4.2\endcsname{ega4.c04.eq.21-6-4-2}
\expandafter\def\csname sga@target@IV.21.6.5\endcsname{ega4.c04.statement.21-6-5}
\expandafter\def\csname sga@target@IV.21.6.5.1\endcsname{ega4.c04.eq.21-6-5-1}
\expandafter\def\csname sga@target@IV.21.6.6\endcsname{ega4.c04.prop.21-6-6}
\expandafter\def\csname sga@target@IV.21.6.6.2\endcsname{ega4.c04.statement.21-6-6-2}
\expandafter\def\csname sga@target@IV.21.6.7\endcsname{ega4.c04.statement.21-6-7}
\expandafter\def\csname sga@target@IV.21.6.7.1\endcsname{ega4.c04.eq.21-6-7-1}
\expandafter\def\csname sga@target@IV.21.6.7.2\endcsname{ega4.c04.eq.21-6-7-2}
\expandafter\def\csname sga@target@IV.21.6.7.3\endcsname{ega4.c04.eq.21-6-7-3}
\expandafter\def\csname sga@target@IV.21.6.7.4\endcsname{ega4.c04.eq.21-6-7-4}
\expandafter\def\csname sga@target@IV.21.6.7.5\endcsname{ega4.c04.eq.21-6-7-5}
\expandafter\def\csname sga@target@IV.21.6.8\endcsname{ega4.c04.statement.21-6-8}
\expandafter\def\csname sga@target@IV.21.6.9\endcsname{ega4.c04.thm.21-6-9}
\expandafter\def\csname sga@target@IV.21.6.9.1\endcsname{ega4.c04.eq.21-6-9-1}
\expandafter\def\csname sga@target@IV.21.6.9.2\endcsname{ega4.c04.statement.21-6-9-2}
\expandafter\def\csname sga@target@IV.21.7\endcsname{ega4.c04.subsec.21-7.key-iv-21-7}
\expandafter\def\csname sga@target@IV.21.7.1\endcsname{ega4.c04.statement.21-7-1}
\expandafter\def\csname sga@target@IV.21.7.2\endcsname{ega4.c04.prop.21-7-2}
\expandafter\def\csname sga@target@IV.21.7.2.1\endcsname{ega4.c04.eq.21-7-2-1}
\expandafter\def\csname sga@target@IV.21.7.3\endcsname{ega4.c04.cor.21-7-3}
\expandafter\def\csname sga@target@IV.21.7.3.1\endcsname{ega4.c04.statement.21-7-3-1}
\expandafter\def\csname sga@target@IV.21.7.4\endcsname{ega4.c04.prop.21-7-4}
\expandafter\def\csname sga@target@IV.21.7.5\endcsname{ega4.c04.rem.21-7-5}
\expandafter\def\csname sga@target@IV.21.7.6\endcsname{ega4.c04.cor.21-7-6}
\expandafter\def\csname sga@target@IV.21.7.7\endcsname{ega4.c04.rem.21-7-7}
\expandafter\def\csname sga@target@IV.21.8.1\endcsname{ega4.c04.lem.21-8-1}
\expandafter\def\csname sga@target@IV.21.8.2\endcsname{ega4.c04.cor.21-8-2}
\expandafter\def\csname sga@target@IV.21.8.3\endcsname{ega4.c04.prop.21-8-3}
\expandafter\def\csname sga@target@IV.21.8.3.1\endcsname{ega4.c04.eq.21-8-3-1}
\expandafter\def\csname sga@target@IV.21.8.4\endcsname{ega4.c04.cor.21-8-4}
\expandafter\def\csname sga@target@IV.21.8.5\endcsname{ega4.c04.prop.21-8-5}
\expandafter\def\csname sga@target@IV.21.8.5.1\endcsname{ega4.c04.eq.21-8-5-1}
\expandafter\def\csname sga@target@IV.21.8.5.2\endcsname{ega4.c04.eq.21-8-5-2}
\expandafter\def\csname sga@target@IV.21.8.5.3\endcsname{ega4.c04.eq.21-8-5-3}
\expandafter\def\csname sga@target@IV.21.8.6\endcsname{ega4.c04.rem.21-8-6}
\expandafter\def\csname sga@target@IV.21.9\endcsname{ega4.c04.subsec.21-9.key-iv-21-9}
\expandafter\def\csname sga@target@IV.21.9.1\endcsname{ega4.c04.statement.21-9-1}
\expandafter\def\csname sga@target@IV.21.9.1.1\endcsname{ega4.c04.eq.21-9-1-1}
\expandafter\def\csname sga@target@IV.21.9.1.2\endcsname{ega4.c04.eq.21-9-1-2}
\expandafter\def\csname sga@target@IV.21.9.10\endcsname{ega4.c04.statement.21-9-10}
\expandafter\def\csname sga@target@IV.21.9.11\endcsname{ega4.c04.prop.21-9-11}
\expandafter\def\csname sga@target@IV.21.9.12\endcsname{ega4.c04.cor.21-9-12}
\expandafter\def\csname sga@target@IV.21.9.13\endcsname{ega4.c04.rem.21-9-13}
\expandafter\def\csname sga@target@IV.21.9.2\endcsname{ega4.c04.lem.21-9-2}
\expandafter\def\csname sga@target@IV.21.9.3\endcsname{ega4.c04.rem.21-9-3}
\expandafter\def\csname sga@target@IV.21.9.4\endcsname{ega4.c04.prop.21-9-4}
\expandafter\def\csname sga@target@IV.21.9.4.1\endcsname{ega4.c04.prop.21-9-4-1}
\expandafter\def\csname sga@target@IV.21.9.5\endcsname{ega4.c04.cor.21-9-5}
\expandafter\def\csname sga@target@IV.21.9.6\endcsname{ega4.c04.cor.21-9-6}
\expandafter\def\csname sga@target@IV.21.9.7\endcsname{ega4.c04.cor.21-9-7}
\expandafter\def\csname sga@target@IV.21.9.8\endcsname{ega4.c04.statement.21-9-8}
\expandafter\def\csname sga@target@IV.21.9.9\endcsname{ega4.c04.prop.21-9-9}
\expandafter\def\csname sga@target@IV.3.1.1\endcsname{ega4.c04.def.3-1-1}
\expandafter\def\csname sga@target@IV.3.1.1.1\endcsname{ega4.c04.eq.3-1-1-1}
\expandafter\def\csname sga@target@IV.3.1.10\endcsname{ega4.c04.prop.3-1-10}
\expandafter\def\csname sga@target@IV.3.1.10.1\endcsname{ega4.c04.lem.3-1-10-1}
\expandafter\def\csname sga@target@IV.3.1.10.2\endcsname{ega4.c04.rem.3-1-10-2}
\expandafter\def\csname sga@target@IV.3.1.10.3\endcsname{ega4.c04.eq.3-1-10-3}
\expandafter\def\csname sga@target@IV.3.1.10.4\endcsname{ega4.c04.eq.3-1-10-4}
\expandafter\def\csname sga@target@IV.3.1.10.5\endcsname{ega4.c04.rem.3-1-10-5}
\expandafter\def\csname sga@target@IV.3.1.11\endcsname{ega4.c04.cor.3-1-11}
\expandafter\def\csname sga@target@IV.3.1.12\endcsname{ega4.c04.rem.3-1-12}
\expandafter\def\csname sga@target@IV.3.1.13\endcsname{ega4.c04.prop.3-1-13}
\expandafter\def\csname sga@target@IV.3.1.2\endcsname{ega4.c04.prop.3-1-2}
\expandafter\def\csname sga@target@IV.3.1.3\endcsname{ega4.c04.prop.3-1-3}
\expandafter\def\csname sga@target@IV.3.1.4\endcsname{ega4.c04.cor.3-1-4}
\expandafter\def\csname sga@target@IV.3.1.5\endcsname{ega4.c04.cor.3-1-5}
\expandafter\def\csname sga@target@IV.3.1.6\endcsname{ega4.c04.prop.3-1-6}
\expandafter\def\csname sga@target@IV.3.1.7\endcsname{ega4.c04.prop.3-1-7}
\expandafter\def\csname sga@target@IV.3.1.8\endcsname{ega4.c04.prop.3-1-8}
\expandafter\def\csname sga@target@IV.3.1.9\endcsname{ega4.c04.cor.3-1-9}
\expandafter\def\csname sga@target@IV.3.2.1\endcsname{ega4.c04.prop.3-2-1}
\expandafter\def\csname sga@target@IV.3.2.2\endcsname{ega4.c04.def.3-2-2}
\expandafter\def\csname sga@target@IV.3.2.3\endcsname{ega4.c04.prop.3-2-3}
\expandafter\def\csname sga@target@IV.3.2.4\endcsname{ega4.c04.def.3-2-4}
\expandafter\def\csname sga@target@IV.3.2.5\endcsname{ega4.c04.def.3-2-5}
\expandafter\def\csname sga@target@IV.3.2.6\endcsname{ega4.c04.prop.3-2-6}
\expandafter\def\csname sga@target@IV.3.2.7\endcsname{ega4.c04.cor.3-2-7}
\expandafter\def\csname sga@target@IV.3.2.8\endcsname{ega4.c04.cor.3-2-8}
\expandafter\def\csname sga@target@IV.3.2.9\endcsname{ega4.c04.rem.3-2-9}
\expandafter\def\csname sga@target@IV.3.3.1\endcsname{ega4.c04.prop.3-3-1}
\expandafter\def\csname sga@target@IV.3.3.1.1\endcsname{ega4.c04.prop.3-3-1-1}
\expandafter\def\csname sga@target@IV.3.3.2\endcsname{ega4.c04.cor.3-3-2}
\expandafter\def\csname sga@target@IV.3.3.3\endcsname{ega4.c04.cor.3-3-3}
\expandafter\def\csname sga@target@IV.3.3.4\endcsname{ega4.c04.prop.3-3-4}
\expandafter\def\csname sga@target@IV.3.3.5\endcsname{ega4.c04.cor.3-3-5}
\expandafter\def\csname sga@target@IV.3.3.6\endcsname{ega4.c04.prop.3-3-6}
\expandafter\def\csname sga@target@IV.3.3.6.1\endcsname{ega4.c04.eq.3-3-6-1}
\expandafter\def\csname sga@target@IV.3.3.6.2\endcsname{ega4.c04.eq.3-3-6-2}
\expandafter\def\csname sga@target@IV.3.3.6.3\endcsname{ega4.c04.eq.3-3-6-3}
\expandafter\def\csname sga@target@IV.3.3.6.4\endcsname{ega4.c04.eq.3-3-6-4}
\expandafter\def\csname sga@target@IV.3.3.7\endcsname{ega4.c04.cor.3-3-7}
\expandafter\def\csname sga@target@IV.3.3.7.1\endcsname{ega4.c04.cor.3-3-7-1}
\expandafter\def\csname sga@target@IV.3.3.8\endcsname{ega4.c04.sec.3-3-8}
\expandafter\def\csname sga@target@IV.3.3.9\endcsname{ega4.c04.cor.3-3-9}
\expandafter\def\csname sga@target@IV.3.3.9.1\endcsname{ega4.c04.cor.3-3-9-1}
\expandafter\def\csname sga@target@IV.3.4\endcsname{ega4.c04.subsec.3-4.key-iv-3-4}
\expandafter\def\csname sga@target@IV.3.4.1\endcsname{ega4.c04.prop.3-4-1}
\expandafter\def\csname sga@target@IV.3.4.1.1\endcsname{ega4.c04.eq.3-4-1-1}
\expandafter\def\csname sga@target@IV.3.4.1.2\endcsname{ega4.c04.lem.3-4-1-2}
\expandafter\def\csname sga@target@IV.3.4.1.3\endcsname{ega4.c04.lem.3-4-1-3}
\expandafter\def\csname sga@target@IV.3.4.1.4\endcsname{ega4.c04.lem.3-4-1-4}
\expandafter\def\csname sga@target@IV.3.4.1.5\endcsname{ega4.c04.eq.3-4-1-5}
\expandafter\def\csname sga@target@IV.3.4.2\endcsname{ega4.c04.cor.3-4-2}
\expandafter\def\csname sga@target@IV.3.4.3\endcsname{ega4.c04.prop.3-4-3}
\expandafter\def\csname sga@target@IV.3.4.4\endcsname{ega4.c04.prop.3-4-4}
\expandafter\def\csname sga@target@IV.3.4.5\endcsname{ega4.c04.prop.3-4-5}
\expandafter\def\csname sga@target@IV.3.4.6\endcsname{ega4.c04.prop.3-4-6}
\expandafter\def\csname sga@target@IV.3.4.6.1\endcsname{ega4.c04.lem.3-4-6-1}
\expandafter\def\csname sga@target@IV.3.4.7\endcsname{ega4.c04.cor.3-4-7}
\expandafter\def\csname sga@target@IV.3.4.8\endcsname{ega4.c04.rem.3-4-8}
\expandafter\def\csname sga@target@IV.3.4.9\endcsname{ega4.c04.prop.3-4-9}
\expandafter\def\csname sga@target@IV.4.1.1\endcsname{ega4.c04.def.4-1-1}
\expandafter\def\csname sga@target@IV.4.1.1.1\endcsname{ega4.c04.eq.4-1-1-1}
\expandafter\def\csname sga@target@IV.4.1.1.2\endcsname{ega4.c04.eq.4-1-1-2}
\expandafter\def\csname sga@target@IV.4.1.1.3\endcsname{ega4.c04.eq.4-1-1-3}
\expandafter\def\csname sga@target@IV.4.1.2\endcsname{ega4.c04.thm.4-1-2}
\expandafter\def\csname sga@target@IV.4.1.2.1\endcsname{ega4.c04.cor.4-1-2-1}
\expandafter\def\csname sga@target@IV.4.1.3\endcsname{ega4.c04.rem.4-1-3}
\expandafter\def\csname sga@target@IV.4.1.4\endcsname{ega4.c04.cor.4-1-4}
\expandafter\def\csname sga@target@IV.4.1.5\endcsname{ega4.c04.cor.4-1-5}
\expandafter\def\csname sga@target@IV.4.2.1\endcsname{ega4.c04.prop.4-2-1}
\expandafter\def\csname sga@target@IV.4.2.1.1\endcsname{ega4.c04.eq.4-2-1-1}
\expandafter\def\csname sga@target@IV.4.2.1.2\endcsname{ega4.c04.eq.4-2-1-2}
\expandafter\def\csname sga@target@IV.4.2.1.3\endcsname{ega4.c04.eq.4-2-1-3}
\expandafter\def\csname sga@target@IV.4.2.1.4\endcsname{ega4.c04.rem.4-2-1-4}
\expandafter\def\csname sga@target@IV.4.2.1.5\endcsname{ega4.c04.eq.4-2-1-5}
\expandafter\def\csname sga@target@IV.4.2.10\endcsname{ega4.c04.cor.4-2-10}
\expandafter\def\csname sga@target@IV.4.2.2\endcsname{ega4.c04.cor.4-2-2}
\expandafter\def\csname sga@target@IV.4.2.3\endcsname{ega4.c04.cor.4-2-3}
\expandafter\def\csname sga@target@IV.4.2.4\endcsname{ega4.c04.prop.4-2-4}
\expandafter\def\csname sga@target@IV.4.2.4.i\endcsname{ega4.c04.sec.4-2-4-i}
\expandafter\def\csname sga@target@IV.4.2.4.ii\endcsname{ega4.c04.sec.4-2-4-ii}
\expandafter\def\csname sga@target@IV.4.2.4.iii\endcsname{ega4.c04.sec.4-2-4-iii}
\expandafter\def\csname sga@target@IV.4.2.5\endcsname{ega4.c04.prop.4-2-5}
\expandafter\def\csname sga@target@IV.4.2.6\endcsname{ega4.c04.cor.4-2-6}
\expandafter\def\csname sga@target@IV.4.2.7\endcsname{ega4.c04.prop.4-2-7}
\expandafter\def\csname sga@target@IV.4.2.7.i\endcsname{ega4.c04.sec.4-2-7-i}
\expandafter\def\csname sga@target@IV.4.2.7.ii\endcsname{ega4.c04.sec.4-2-7-ii}
\expandafter\def\csname sga@target@IV.4.2.7.iii\endcsname{ega4.c04.sec.4-2-7-iii}
\expandafter\def\csname sga@target@IV.4.2.7.iv\endcsname{ega4.c04.sec.4-2-7-iv}
\expandafter\def\csname sga@target@IV.4.2.8\endcsname{ega4.c04.cor.4-2-8}
\expandafter\def\csname sga@target@IV.4.2.9\endcsname{ega4.c04.prop.4-2-9}
\expandafter\def\csname sga@target@IV.4.3.1\endcsname{ega4.c04.statement.4-3-1}
\expandafter\def\csname sga@target@IV.4.3.2\endcsname{ega4.c04.prop.4-3-2}
\expandafter\def\csname sga@target@IV.4.3.3\endcsname{ega4.c04.cor.4-3-3}
\expandafter\def\csname sga@target@IV.4.3.4\endcsname{ega4.c04.cor.4-3-4}
\expandafter\def\csname sga@target@IV.4.3.5\endcsname{ega4.c04.prop.4-3-5}
\expandafter\def\csname sga@target@IV.4.3.6\endcsname{ega4.c04.cor.4-3-6}
\expandafter\def\csname sga@target@IV.4.3.7\endcsname{ega4.c04.cor.4-3-7}
\expandafter\def\csname sga@target@IV.4.3.8\endcsname{ega4.c04.cor.4-3-8}
\expandafter\def\csname sga@target@IV.4.3.9\endcsname{ega4.c04.prop.4-3-9}
\expandafter\def\csname sga@target@IV.4.4.1\endcsname{ega4.c04.statement.4-4-1}
\expandafter\def\csname sga@target@IV.4.4.2\endcsname{ega4.c04.lem.4-4-2}
\expandafter\def\csname sga@target@IV.4.4.2.i\endcsname{ega4.c04.sec.4-4-2-i}
\expandafter\def\csname sga@target@IV.4.4.2.ii\endcsname{ega4.c04.sec.4-4-2-ii}
\expandafter\def\csname sga@target@IV.4.4.3\endcsname{ega4.c04.rem.4-4-3}
\expandafter\def\csname sga@target@IV.4.4.4\endcsname{ega4.c04.thm.4-4-4}
\expandafter\def\csname sga@target@IV.4.4.5\endcsname{ega4.c04.cor.4-4-5}
\expandafter\def\csname sga@target@IV.4.4.6\endcsname{ega4.c04.cor.4-4-6}
\expandafter\def\csname sga@target@IV.4.5\endcsname{ega4.c04.subsec.4-5.key-iv-4-5}
\expandafter\def\csname sga@target@IV.4.5.1\endcsname{ega4.c04.prop.4-5-1}
\expandafter\def\csname sga@target@IV.4.5.10\endcsname{ega4.c04.cor.4-5-10}
\expandafter\def\csname sga@target@IV.4.5.11\endcsname{ega4.c04.cor.4-5-11}
\expandafter\def\csname sga@target@IV.4.5.12\endcsname{ega4.c04.sec.4-5-12}
\expandafter\def\csname sga@target@IV.4.5.12.i\endcsname{ega4.c04.sec.4-5-12-i}
\expandafter\def\csname sga@target@IV.4.5.12.ii\endcsname{ega4.c04.sec.4-5-12-ii}
\expandafter\def\csname sga@target@IV.4.5.13\endcsname{ega4.c04.prop.4-5-13}
\expandafter\def\csname sga@target@IV.4.5.13.1\endcsname{ega4.c04.cor.4-5-13-1}
\expandafter\def\csname sga@target@IV.4.5.13.1.i\endcsname{ega4.c04.sec.4-5-13-1-i}
\expandafter\def\csname sga@target@IV.4.5.13.1.ii\endcsname{ega4.c04.sec.4-5-13-1-ii}
\expandafter\def\csname sga@target@IV.4.5.14\endcsname{ega4.c04.cor.4-5-14}
\expandafter\def\csname sga@target@IV.4.5.15\endcsname{ega4.c04.prop.4-5-15}
\expandafter\def\csname sga@target@IV.4.5.15.i\endcsname{ega4.c04.sec.4-5-15-i}
\expandafter\def\csname sga@target@IV.4.5.15.ii\endcsname{ega4.c04.sec.4-5-15-ii}
\expandafter\def\csname sga@target@IV.4.5.16\endcsname{ega4.c04.cor.4-5-16}
\expandafter\def\csname sga@target@IV.4.5.16.i\endcsname{ega4.c04.sec.4-5-16-i}
\expandafter\def\csname sga@target@IV.4.5.16.ii\endcsname{ega4.c04.sec.4-5-16-ii}
\expandafter\def\csname sga@target@IV.4.5.17\endcsname{ega4.c04.cor.4-5-17}
\expandafter\def\csname sga@target@IV.4.5.18\endcsname{ega4.c04.rem.4-5-18}
\expandafter\def\csname sga@target@IV.4.5.19\endcsname{ega4.c04.prop.4-5-19}
\expandafter\def\csname sga@target@IV.4.5.19.1\endcsname{ega4.c04.lem.4-5-19-1}
\expandafter\def\csname sga@target@IV.4.5.19.2\endcsname{ega4.c04.eq.4-5-19-2}
\expandafter\def\csname sga@target@IV.4.5.19.3\endcsname{ega4.c04.cor.4-5-19-3}
\expandafter\def\csname sga@target@IV.4.5.2\endcsname{ega4.c04.def.4-5-2}
\expandafter\def\csname sga@target@IV.4.5.2.a\endcsname{ega4.c04.sec.4-5-2-a}
\expandafter\def\csname sga@target@IV.4.5.2.b\endcsname{ega4.c04.sec.4-5-2-b}
\expandafter\def\csname sga@target@IV.4.5.2.c\endcsname{ega4.c04.sec.4-5-2-c}
\expandafter\def\csname sga@target@IV.4.5.20\endcsname{ega4.c04.rem.4-5-20}
\expandafter\def\csname sga@target@IV.4.5.21\endcsname{ega4.c04.prop.4-5-21}
\expandafter\def\csname sga@target@IV.4.5.3\endcsname{ega4.c04.statement.4-5-3}
\expandafter\def\csname sga@target@IV.4.5.4\endcsname{ega4.c04.prop.4-5-4}
\expandafter\def\csname sga@target@IV.4.5.5\endcsname{ega4.c04.def.4-5-5}
\expandafter\def\csname sga@target@IV.4.5.6\endcsname{ega4.c04.prop.4-5-6}
\expandafter\def\csname sga@target@IV.4.5.6.i\endcsname{ega4.c04.sec.4-5-6-i}
\expandafter\def\csname sga@target@IV.4.5.6.ii\endcsname{ega4.c04.sec.4-5-6-ii}
\expandafter\def\csname sga@target@IV.4.5.7\endcsname{ega4.c04.prop.4-5-7}
\expandafter\def\csname sga@target@IV.4.5.8\endcsname{ega4.c04.cor.4-5-8}
\expandafter\def\csname sga@target@IV.4.5.8.i\endcsname{ega4.c04.sec.4-5-8-i}
\expandafter\def\csname sga@target@IV.4.5.8.ii\endcsname{ega4.c04.sec.4-5-8-ii}
\expandafter\def\csname sga@target@IV.4.5.9\endcsname{ega4.c04.prop.4-5-9}
\expandafter\def\csname sga@target@IV.4.5.9.a\endcsname{ega4.c04.sec.4-5-9-a}
\expandafter\def\csname sga@target@IV.4.5.9.b\endcsname{ega4.c04.sec.4-5-9-b}
\expandafter\def\csname sga@target@IV.4.5.9.c\endcsname{ega4.c04.sec.4-5-9-c}
\expandafter\def\csname sga@target@IV.4.6\endcsname{ega4.c04.subsec.4-6.key-iv-4-6}
\expandafter\def\csname sga@target@IV.4.6.1\endcsname{ega4.c04.prop.4-6-1}
\expandafter\def\csname sga@target@IV.4.6.1.a\endcsname{ega4.c04.sec.4-6-1-a}
\expandafter\def\csname sga@target@IV.4.6.1.b\endcsname{ega4.c04.sec.4-6-1-b}
\expandafter\def\csname sga@target@IV.4.6.1.c\endcsname{ega4.c04.sec.4-6-1-c}
\expandafter\def\csname sga@target@IV.4.6.1.d\endcsname{ega4.c04.sec.4-6-1-d}
\expandafter\def\csname sga@target@IV.4.6.1.e\endcsname{ega4.c04.sec.4-6-1-e}
\expandafter\def\csname sga@target@IV.4.6.10\endcsname{ega4.c04.prop.4-6-10}
\expandafter\def\csname sga@target@IV.4.6.11\endcsname{ega4.c04.cor.4-6-11}
\expandafter\def\csname sga@target@IV.4.6.12\endcsname{ega4.c04.prop.4-6-12}
\expandafter\def\csname sga@target@IV.4.6.12.a\endcsname{ega4.c04.sec.4-6-12-a}
\expandafter\def\csname sga@target@IV.4.6.12.b\endcsname{ega4.c04.sec.4-6-12-b}
\expandafter\def\csname sga@target@IV.4.6.12.c\endcsname{ega4.c04.sec.4-6-12-c}
\expandafter\def\csname sga@target@IV.4.6.12.d\endcsname{ega4.c04.sec.4-6-12-d}
\expandafter\def\csname sga@target@IV.4.6.12.e\endcsname{ega4.c04.sec.4-6-12-e}
\expandafter\def\csname sga@target@IV.4.6.13\endcsname{ega4.c04.cor.4-6-13}
\expandafter\def\csname sga@target@IV.4.6.13.f\endcsname{ega4.c04.sec.4-6-13-f}
\expandafter\def\csname sga@target@IV.4.6.14\endcsname{ega4.c04.statement.4-6-14}
\expandafter\def\csname sga@target@IV.4.6.15\endcsname{ega4.c04.prop.4-6-15}
\expandafter\def\csname sga@target@IV.4.6.15.i\endcsname{ega4.c04.sec.4-6-15-i}
\expandafter\def\csname sga@target@IV.4.6.15.ii\endcsname{ega4.c04.sec.4-6-15-ii}
\expandafter\def\csname sga@target@IV.4.6.16\endcsname{ega4.c04.prop.4-6-16}
\expandafter\def\csname sga@target@IV.4.6.16.a\endcsname{ega4.c04.sec.4-6-16-a}
\expandafter\def\csname sga@target@IV.4.6.16.b\endcsname{ega4.c04.sec.4-6-16-b}
\expandafter\def\csname sga@target@IV.4.6.16.c\endcsname{ega4.c04.sec.4-6-16-c}
\expandafter\def\csname sga@target@IV.4.6.16.d\endcsname{ega4.c04.sec.4-6-16-d}
\expandafter\def\csname sga@target@IV.4.6.17\endcsname{ega4.c04.def.4-6-17}
\expandafter\def\csname sga@target@IV.4.6.18\endcsname{ega4.c04.prop.4-6-18}
\expandafter\def\csname sga@target@IV.4.6.18.a\endcsname{ega4.c04.sec.4-6-18-a}
\expandafter\def\csname sga@target@IV.4.6.18.b\endcsname{ega4.c04.sec.4-6-18-b}
\expandafter\def\csname sga@target@IV.4.6.18.c\endcsname{ega4.c04.sec.4-6-18-c}
\expandafter\def\csname sga@target@IV.4.6.18.d\endcsname{ega4.c04.sec.4-6-18-d}
\expandafter\def\csname sga@target@IV.4.6.19\endcsname{ega4.c04.def.4-6-19}
\expandafter\def\csname sga@target@IV.4.6.2\endcsname{ega4.c04.def.4-6-2}
\expandafter\def\csname sga@target@IV.4.6.20\endcsname{ega4.c04.prop.4-6-20}
\expandafter\def\csname sga@target@IV.4.6.21\endcsname{ega4.c04.prop.4-6-21}
\expandafter\def\csname sga@target@IV.4.6.21.i\endcsname{ega4.c04.sec.4-6-21-i}
\expandafter\def\csname sga@target@IV.4.6.21.ii\endcsname{ega4.c04.sec.4-6-21-ii}
\expandafter\def\csname sga@target@IV.4.6.22\endcsname{ega4.c04.def.4-6-22}
\expandafter\def\csname sga@target@IV.4.6.3\endcsname{ega4.c04.cor.4-6-3}
\expandafter\def\csname sga@target@IV.4.6.4\endcsname{ega4.c04.cor.4-6-4}
\expandafter\def\csname sga@target@IV.4.6.5\endcsname{ega4.c04.prop.4-6-5}
\expandafter\def\csname sga@target@IV.4.6.5.i\endcsname{ega4.c04.sec.4-6-5-i}
\expandafter\def\csname sga@target@IV.4.6.5.ii\endcsname{ega4.c04.sec.4-6-5-ii}
\expandafter\def\csname sga@target@IV.4.6.6\endcsname{ega4.c04.prop.4-6-6}
\expandafter\def\csname sga@target@IV.4.6.7\endcsname{ega4.c04.cor.4-6-7}
\expandafter\def\csname sga@target@IV.4.6.8\endcsname{ega4.c04.cor.4-6-8}
\expandafter\def\csname sga@target@IV.4.6.9\endcsname{ega4.c04.def.4-6-9}
\expandafter\def\csname sga@target@IV.4.7.1\endcsname{ega4.c04.lem.4-7-1}
\expandafter\def\csname sga@target@IV.4.7.1.1\endcsname{ega4.c04.eq.4-7-1-1}
\expandafter\def\csname sga@target@IV.4.7.10\endcsname{ega4.c04.prop.4-7-10}
\expandafter\def\csname sga@target@IV.4.7.11\endcsname{ega4.c04.prop.4-7-11}
\expandafter\def\csname sga@target@IV.4.7.12\endcsname{ega4.c04.def.4-7-12}
\expandafter\def\csname sga@target@IV.4.7.13\endcsname{ega4.c04.prop.4-7-13}
\expandafter\def\csname sga@target@IV.4.7.2\endcsname{ega4.c04.cor.4-7-2}
\expandafter\def\csname sga@target@IV.4.7.2.1\endcsname{ega4.c04.eq.4-7-2-1}
\expandafter\def\csname sga@target@IV.4.7.3\endcsname{ega4.c04.prop.4-7-3}
\expandafter\def\csname sga@target@IV.4.7.3.1\endcsname{ega4.c04.eq.4-7-3-1}
\expandafter\def\csname sga@target@IV.4.7.3.i\endcsname{ega4.c04.sec.4-7-3-i}
\expandafter\def\csname sga@target@IV.4.7.3.ii\endcsname{ega4.c04.sec.4-7-3-ii}
\expandafter\def\csname sga@target@IV.4.7.3.iii\endcsname{ega4.c04.sec.4-7-3-iii}
\expandafter\def\csname sga@target@IV.4.7.4\endcsname{ega4.c04.def.4-7-4}
\expandafter\def\csname sga@target@IV.4.7.5\endcsname{ega4.c04.def.4-7-5}
\expandafter\def\csname sga@target@IV.4.7.6\endcsname{ega4.c04.prop.4-7-6}
\expandafter\def\csname sga@target@IV.4.7.6.1\endcsname{ega4.c04.eq.4-7-6-1}
\expandafter\def\csname sga@target@IV.4.7.7\endcsname{ega4.c04.statement.4-7-7}
\expandafter\def\csname sga@target@IV.4.7.8\endcsname{ega4.c04.prop.4-7-8}
\expandafter\def\csname sga@target@IV.4.7.9\endcsname{ega4.c04.cor.4-7-9}
\expandafter\def\csname sga@target@IV.4.8.1\endcsname{ega4.c04.statement.4-8-1}
\expandafter\def\csname sga@target@IV.4.8.10\endcsname{ega4.c04.cor.4-8-10}
\expandafter\def\csname sga@target@IV.4.8.11\endcsname{ega4.c04.cor.4-8-11}
\expandafter\def\csname sga@target@IV.4.8.12\endcsname{ega4.c04.cor.4-8-12}
\expandafter\def\csname sga@target@IV.4.8.13\endcsname{ega4.c04.prop.4-8-13}
\expandafter\def\csname sga@target@IV.4.8.14\endcsname{ega4.c04.prop.4-8-14}
\expandafter\def\csname sga@target@IV.4.8.15\endcsname{ega4.c04.sec.4-8-15}
\expandafter\def\csname sga@target@IV.4.8.15.i\endcsname{ega4.c04.sec.4-8-15-i}
\expandafter\def\csname sga@target@IV.4.8.15.ii\endcsname{ega4.c04.sec.4-8-15-ii}
\expandafter\def\csname sga@target@IV.4.8.2\endcsname{ega4.c04.statement.4-8-2}
\expandafter\def\csname sga@target@IV.4.8.2.1\endcsname{ega4.c04.eq.4-8-2-1}
\expandafter\def\csname sga@target@IV.4.8.2.2\endcsname{ega4.c04.eq.4-8-2-2}
\expandafter\def\csname sga@target@IV.4.8.2.3\endcsname{ega4.c04.eq.4-8-2-3}
\expandafter\def\csname sga@target@IV.4.8.2.4\endcsname{ega4.c04.eq.4-8-2-4}
\expandafter\def\csname sga@target@IV.4.8.2.5\endcsname{ega4.c04.eq.4-8-2-5}
\expandafter\def\csname sga@target@IV.4.8.2.6\endcsname{ega4.c04.eq.4-8-2-6}
\expandafter\def\csname sga@target@IV.4.8.3\endcsname{ega4.c04.prop.4-8-3}
\expandafter\def\csname sga@target@IV.4.8.4\endcsname{ega4.c04.def.4-8-4}
\expandafter\def\csname sga@target@IV.4.8.5\endcsname{ega4.c04.prop.4-8-5}
\expandafter\def\csname sga@target@IV.4.8.6\endcsname{ega4.c04.lem.4-8-6}
\expandafter\def\csname sga@target@IV.4.8.6.a\endcsname{ega4.c04.sec.4-8-6-a}
\expandafter\def\csname sga@target@IV.4.8.6.b\endcsname{ega4.c04.sec.4-8-6-b}
\expandafter\def\csname sga@target@IV.4.8.6.c\endcsname{ega4.c04.sec.4-8-6-c}
\expandafter\def\csname sga@target@IV.4.8.7\endcsname{ega4.c04.cor.4-8-7}
\expandafter\def\csname sga@target@IV.4.8.8\endcsname{ega4.c04.lem.4-8-8}
\expandafter\def\csname sga@target@IV.4.8.8.i\endcsname{ega4.c04.sec.4-8-8-i}
\expandafter\def\csname sga@target@IV.4.8.8.ii\endcsname{ega4.c04.sec.4-8-8-ii}
\expandafter\def\csname sga@target@IV.4.8.8.iii\endcsname{ega4.c04.sec.4-8-8-iii}
\expandafter\def\csname sga@target@IV.4.8.9\endcsname{ega4.c04.prop.4-8-9}
\expandafter\def\csname sga@target@IV.4.9.1\endcsname{ega4.c04.prop.4-9-1}
\expandafter\def\csname sga@target@IV.4.9.1.a\endcsname{ega4.c04.sec.4-9-1-a}
\expandafter\def\csname sga@target@IV.4.9.1.b\endcsname{ega4.c04.sec.4-9-1-b}
\expandafter\def\csname sga@target@IV.4.9.1.bprime\endcsname{ega4.c04.sec.4-9-1-bprime}
\expandafter\def\csname sga@target@IV.4.9.1.c\endcsname{ega4.c04.sec.4-9-1-c}
\expandafter\def\csname sga@target@IV.4.9.2\endcsname{ega4.c04.cor.4-9-2}
\expandafter\def\csname sga@target@IV.4.9.3\endcsname{ega4.c04.rem.4-9-3}
\expandafter\def\csname sga@target@IV.4.9.4\endcsname{ega4.c04.prop.4-9-4}
\expandafter\def\csname sga@target@IV.4.9.5\endcsname{ega4.c04.cor.4-9-5}
\expandafter\def\csname sga@target@IV.4.9.5.a\endcsname{ega4.c04.sec.4-9-5-a}
\expandafter\def\csname sga@target@IV.4.9.5.b\endcsname{ega4.c04.sec.4-9-5-b}
\expandafter\def\csname sga@target@IV.4.9.6\endcsname{ega4.c04.prop.4-9-6}
\expandafter\def\csname sga@target@IV.4.9.6.a\endcsname{ega4.c04.sec.4-9-6-a}
\expandafter\def\csname sga@target@IV.4.9.6.aprime\endcsname{ega4.c04.sec.4-9-6-aprime}
\expandafter\def\csname sga@target@IV.4.9.6.b\endcsname{ega4.c04.sec.4-9-6-b}
\expandafter\def\csname sga@target@IV.4.9.6.bprime\endcsname{ega4.c04.sec.4-9-6-bprime}
\expandafter\def\csname sga@target@IV.4.9.7\endcsname{ega4.c04.cor.4-9-7}
\expandafter\def\csname sga@target@IV.5.1.1\endcsname{ega4.c04.statement.5-1-1}
\expandafter\def\csname sga@target@IV.5.1.1.1\endcsname{ega4.c04.eq.5-1-1-1}
\expandafter\def\csname sga@target@IV.5.1.1.2\endcsname{ega4.c04.eq.5-1-1-2}
\expandafter\def\csname sga@target@IV.5.1.1.3\endcsname{ega4.c04.eq.5-1-1-3}
\expandafter\def\csname sga@target@IV.5.1.1.4\endcsname{ega4.c04.sec.5-1-1-4}
\expandafter\def\csname sga@target@IV.5.1.1.5\endcsname{ega4.c04.sec.5-1-1-5}
\expandafter\def\csname sga@target@IV.5.1.1.6\endcsname{ega4.c04.sec.5-1-1-6}
\expandafter\def\csname sga@target@IV.5.1.10\endcsname{ega4.c04.prop.5-1-10}
\expandafter\def\csname sga@target@IV.5.1.10.1\endcsname{ega4.c04.lem.5-1-10-1}
\expandafter\def\csname sga@target@IV.5.1.10.i\endcsname{ega4.c04.sec.5-1-10-i}
\expandafter\def\csname sga@target@IV.5.1.10.ii\endcsname{ega4.c04.sec.5-1-10-ii}
\expandafter\def\csname sga@target@IV.5.1.11\endcsname{ega4.c04.cor.5-1-11}
\expandafter\def\csname sga@target@IV.5.1.12\endcsname{ega4.c04.statement.5-1-12}
\expandafter\def\csname sga@target@IV.5.1.12.1\endcsname{ega4.c04.eq.5-1-12-1}
\expandafter\def\csname sga@target@IV.5.1.12.2\endcsname{ega4.c04.eq.5-1-12-2}
\expandafter\def\csname sga@target@IV.5.1.12.3\endcsname{ega4.c04.eq.5-1-12-3}
\expandafter\def\csname sga@target@IV.5.1.13\endcsname{ega4.c04.prop.5-1-13}
\expandafter\def\csname sga@target@IV.5.1.13.1\endcsname{ega4.c04.eq.5-1-13-1}
\expandafter\def\csname sga@target@IV.5.1.2\endcsname{ega4.c04.prop.5-1-2}
\expandafter\def\csname sga@target@IV.5.1.2.1\endcsname{ega4.c04.eq.5-1-2-1}
\expandafter\def\csname sga@target@IV.5.1.3\endcsname{ega4.c04.cor.5-1-3}
\expandafter\def\csname sga@target@IV.5.1.3.1\endcsname{ega4.c04.eq.5-1-3-1}
\expandafter\def\csname sga@target@IV.5.1.3.2\endcsname{ega4.c04.eq.5-1-3-2}
\expandafter\def\csname sga@target@IV.5.1.4\endcsname{ega4.c04.prop.5-1-4}
\expandafter\def\csname sga@target@IV.5.1.4.1\endcsname{ega4.c04.eq.5-1-4-1}
\expandafter\def\csname sga@target@IV.5.1.4.2\endcsname{ega4.c04.eq.5-1-4-2}
\expandafter\def\csname sga@target@IV.5.1.5\endcsname{ega4.c04.cor.5-1-5}
\expandafter\def\csname sga@target@IV.5.1.6\endcsname{ega4.c04.statement.5-1-6}
\expandafter\def\csname sga@target@IV.5.1.7\endcsname{ega4.c04.prop.5-1-7}
\expandafter\def\csname sga@target@IV.5.1.8\endcsname{ega4.c04.cor.5-1-8}
\expandafter\def\csname sga@target@IV.5.1.9\endcsname{ega4.c04.prop.5-1-9}
\expandafter\def\csname sga@target@IV.5.1.9.1\endcsname{ega4.c04.statement.5-1-9-1}
\expandafter\def\csname sga@target@IV.5.10.1\endcsname{ega4.c04.statement.5-10-1}
\expandafter\def\csname sga@target@IV.5.10.1.1\endcsname{ega4.c04.eq.5-10-1-1}
\expandafter\def\csname sga@target@IV.5.10.10\endcsname{ega4.c04.prop.5-10-10}
\expandafter\def\csname sga@target@IV.5.10.11\endcsname{ega4.c04.statement.5-10-11}
\expandafter\def\csname sga@target@IV.5.10.11.1\endcsname{ega4.c04.eq.5-10-11-1}
\expandafter\def\csname sga@target@IV.5.10.12\endcsname{ega4.c04.prop.5-10-12}
\expandafter\def\csname sga@target@IV.5.10.13\endcsname{ega4.c04.statement.5-10-13}
\expandafter\def\csname sga@target@IV.5.10.14\endcsname{ega4.c04.prop.5-10-14}
\expandafter\def\csname sga@target@IV.5.10.15\endcsname{ega4.c04.cor.5-10-15}
\expandafter\def\csname sga@target@IV.5.10.16\endcsname{ega4.c04.prop.5-10-16}
\expandafter\def\csname sga@target@IV.5.10.17\endcsname{ega4.c04.prop.5-10-17}
\expandafter\def\csname sga@target@IV.5.10.2\endcsname{ega4.c04.prop.5-10-2}
\expandafter\def\csname sga@target@IV.5.10.3\endcsname{ega4.c04.cor.5-10-3}
\expandafter\def\csname sga@target@IV.5.10.4\endcsname{ega4.c04.cor.5-10-4}
\expandafter\def\csname sga@target@IV.5.10.5\endcsname{ega4.c04.thm.5-10-5}
\expandafter\def\csname sga@target@IV.5.10.6\endcsname{ega4.c04.cor.5-10-6}
\expandafter\def\csname sga@target@IV.5.10.7\endcsname{ega4.c04.thm.5-10-7}
\expandafter\def\csname sga@target@IV.5.10.8\endcsname{ega4.c04.cor.5-10-8}
\expandafter\def\csname sga@target@IV.5.10.8.1\endcsname{ega4.c04.lem.5-10-8-1}
\expandafter\def\csname sga@target@IV.5.10.9\endcsname{ega4.c04.cor.5-10-9}
\expandafter\def\csname sga@target@IV.5.11.1\endcsname{ega4.c04.prop.5-11-1}
\expandafter\def\csname sga@target@IV.5.11.1.1\endcsname{ega4.c04.lem.5-11-1-1}
\expandafter\def\csname sga@target@IV.5.11.2\endcsname{ega4.c04.cor.5-11-2}
\expandafter\def\csname sga@target@IV.5.11.3\endcsname{ega4.c04.rem.5-11-3}
\expandafter\def\csname sga@target@IV.5.11.4\endcsname{ega4.c04.cor.5-11-4}
\expandafter\def\csname sga@target@IV.5.11.5\endcsname{ega4.c04.cor.5-11-5}
\expandafter\def\csname sga@target@IV.5.11.6\endcsname{ega4.c04.cor.5-11-6}
\expandafter\def\csname sga@target@IV.5.11.7\endcsname{ega4.c04.item.5-11-7}
\expandafter\def\csname sga@target@IV.5.12\endcsname{ega4.c04.subsec.5-12.key-iv-5-12}
\expandafter\def\csname sga@target@IV.5.12.1\endcsname{ega4.c04.prop.5-12-1}
\expandafter\def\csname sga@target@IV.5.12.1.1\endcsname{ega4.c04.eq.5-12-1-1}
\expandafter\def\csname sga@target@IV.5.12.10\endcsname{ega4.c04.cor.5-12-10}
\expandafter\def\csname sga@target@IV.5.12.2\endcsname{ega4.c04.prop.5-12-2}
\expandafter\def\csname sga@target@IV.5.12.3\endcsname{ega4.c04.rem.5-12-3}
\expandafter\def\csname sga@target@IV.5.12.4\endcsname{ega4.c04.cor.5-12-4}
\expandafter\def\csname sga@target@IV.5.12.5\endcsname{ega4.c04.prop.5-12-5}
\expandafter\def\csname sga@target@IV.5.12.6\endcsname{ega4.c04.rem.5-12-6}
\expandafter\def\csname sga@target@IV.5.12.7\endcsname{ega4.c04.cor.5-12-7}
\expandafter\def\csname sga@target@IV.5.12.8\endcsname{ega4.c04.prop.5-12-8}
\expandafter\def\csname sga@target@IV.5.12.9\endcsname{ega4.c04.rem.5-12-9}
\expandafter\def\csname sga@target@IV.5.13.1\endcsname{ega4.c04.prop.5-13-1}
\expandafter\def\csname sga@target@IV.5.13.2\endcsname{ega4.c04.cor.5-13-2}
\expandafter\def\csname sga@target@IV.5.13.3\endcsname{ega4.c04.prop.5-13-3}
\expandafter\def\csname sga@target@IV.5.13.4\endcsname{ega4.c04.cor.5-13-4}
\expandafter\def\csname sga@target@IV.5.13.5\endcsname{ega4.c04.statement.5-13-5}
\expandafter\def\csname sga@target@IV.5.13.6\endcsname{ega4.c04.cor.5-13-6}
\expandafter\def\csname sga@target@IV.5.13.7\endcsname{ega4.c04.prop.5-13-7}
\expandafter\def\csname sga@target@IV.5.13.7.1\endcsname{ega4.c04.lem.5-13-7-1}
\expandafter\def\csname sga@target@IV.5.13.8\endcsname{ega4.c04.cor.5-13-8}
\expandafter\def\csname sga@target@IV.5.2.1\endcsname{ega4.c04.prop.5-2-1}
\expandafter\def\csname sga@target@IV.5.2.1.1\endcsname{ega4.c04.eq.5-2-1-1}
\expandafter\def\csname sga@target@IV.5.2.2\endcsname{ega4.c04.cor.5-2-2}
\expandafter\def\csname sga@target@IV.5.2.3\endcsname{ega4.c04.cor.5-2-3}
\expandafter\def\csname sga@target@IV.5.2.3.1\endcsname{ega4.c04.eq.5-2-3-1}
\expandafter\def\csname sga@target@IV.5.2.4\endcsname{ega4.c04.cor.5-2-4}
\expandafter\def\csname sga@target@IV.5.2.5\endcsname{ega4.c04.sec.5-2-5}
\expandafter\def\csname sga@target@IV.5.2.5.i\endcsname{ega4.c04.sec.5-2-5-i}
\expandafter\def\csname sga@target@IV.5.2.5.ii\endcsname{ega4.c04.sec.5-2-5-ii}
\expandafter\def\csname sga@target@IV.5.3.1\endcsname{ega4.c04.prop.5-3-1}
\expandafter\def\csname sga@target@IV.5.3.1.1\endcsname{ega4.c04.eq.5-3-1-1}
\expandafter\def\csname sga@target@IV.5.3.1.2\endcsname{ega4.c04.eq.5-3-1-2}
\expandafter\def\csname sga@target@IV.5.3.1.3\endcsname{ega4.c04.lem.5-3-1-3}
\expandafter\def\csname sga@target@IV.5.4.1\endcsname{ega4.c04.prop.5-4-1}
\expandafter\def\csname sga@target@IV.5.4.1.i\endcsname{ega4.c04.sec.5-4-1-i}
\expandafter\def\csname sga@target@IV.5.4.1.ii\endcsname{ega4.c04.sec.5-4-1-ii}
\expandafter\def\csname sga@target@IV.5.4.2\endcsname{ega4.c04.cor.5-4-2}
\expandafter\def\csname sga@target@IV.5.4.3\endcsname{ega4.c04.sec.5-4-3}
\expandafter\def\csname sga@target@IV.5.4.3.1\endcsname{ega4.c04.eq.5-4-3-1}
\expandafter\def\csname sga@target@IV.5.4.3.i\endcsname{ega4.c04.sec.5-4-3-i}
\expandafter\def\csname sga@target@IV.5.4.3.ii\endcsname{ega4.c04.sec.5-4-3-ii}
\expandafter\def\csname sga@target@IV.5.5\endcsname{ega4.c04.sec.5-5}
\expandafter\def\csname sga@target@IV.5.5.1\endcsname{ega4.c04.statement.5-5-1}
\expandafter\def\csname sga@target@IV.5.5.1.1\endcsname{ega4.c04.eq.5-5-1-1}
\expandafter\def\csname sga@target@IV.5.5.2\endcsname{ega4.c04.prop.5-5-2}
\expandafter\def\csname sga@target@IV.5.5.2.1\endcsname{ega4.c04.eq.5-5-2-1}
\expandafter\def\csname sga@target@IV.5.5.2.2\endcsname{ega4.c04.eq.5-5-2-2}
\expandafter\def\csname sga@target@IV.5.5.3\endcsname{ega4.c04.prop.5-5-3}
\expandafter\def\csname sga@target@IV.5.5.3.1\endcsname{ega4.c04.eq.5-5-3-1}
\expandafter\def\csname sga@target@IV.5.5.4\endcsname{ega4.c04.cor.5-5-4}
\expandafter\def\csname sga@target@IV.5.5.4.1\endcsname{ega4.c04.eq.5-5-4-1}
\expandafter\def\csname sga@target@IV.5.5.5\endcsname{ega4.c04.cor.5-5-5}
\expandafter\def\csname sga@target@IV.5.5.6\endcsname{ega4.c04.cor.5-5-6}
\expandafter\def\csname sga@target@IV.5.5.6.1\endcsname{ega4.c04.eq.5-5-6-1}
\expandafter\def\csname sga@target@IV.5.5.7\endcsname{ega4.c04.lem.5-5-7}
\expandafter\def\csname sga@target@IV.5.5.7.1\endcsname{ega4.c04.eq.5-5-7-1}
\expandafter\def\csname sga@target@IV.5.5.7.2\endcsname{ega4.c04.statement.5-5-7-2}
\expandafter\def\csname sga@target@IV.5.5.7.3\endcsname{ega4.c04.eq.5-5-7-3}
\expandafter\def\csname sga@target@IV.5.5.7.i\endcsname{ega4.c04.sec.5-5-7-i}
\expandafter\def\csname sga@target@IV.5.5.7.ii\endcsname{ega4.c04.sec.5-5-7-ii}
\expandafter\def\csname sga@target@IV.5.5.8\endcsname{ega4.c04.thm.5-5-8}
\expandafter\def\csname sga@target@IV.5.5.8.1\endcsname{ega4.c04.eq.5-5-8-1}
\expandafter\def\csname sga@target@IV.5.5.8.2\endcsname{ega4.c04.eq.5-5-8-2}
\expandafter\def\csname sga@target@IV.5.5.8.3\endcsname{ega4.c04.eq.5-5-8-3}
\expandafter\def\csname sga@target@IV.5.6\endcsname{ega4.c04.sec.5-6}
\expandafter\def\csname sga@target@IV.5.6.1\endcsname{ega4.c04.prop.5-6-1}
\expandafter\def\csname sga@target@IV.5.6.1.1\endcsname{ega4.c04.eq.5-6-1-1}
\expandafter\def\csname sga@target@IV.5.6.1.2\endcsname{ega4.c04.eq.5-6-1-2}
\expandafter\def\csname sga@target@IV.5.6.1.3\endcsname{ega4.c04.statement.5-6-1-3}
\expandafter\def\csname sga@target@IV.5.6.1.4\endcsname{ega4.c04.statement.5-6-1-4}
\expandafter\def\csname sga@target@IV.5.6.1.5\endcsname{ega4.c04.statement.5-6-1-5}
\expandafter\def\csname sga@target@IV.5.6.1.a\endcsname{ega4.c04.sec.5-6-1-a}
\expandafter\def\csname sga@target@IV.5.6.1.b\endcsname{ega4.c04.sec.5-6-1-b}
\expandafter\def\csname sga@target@IV.5.6.1.c\endcsname{ega4.c04.sec.5-6-1-c}
\expandafter\def\csname sga@target@IV.5.6.10\endcsname{ega4.c04.prop.5-6-10}
\expandafter\def\csname sga@target@IV.5.6.11\endcsname{ega4.c04.ex.5-6-11}
\expandafter\def\csname sga@target@IV.5.6.11.1\endcsname{ega4.c04.lem.5-6-11-1}
\expandafter\def\csname sga@target@IV.5.6.11.1.i\endcsname{ega4.c04.sec.5-6-11-1-i}
\expandafter\def\csname sga@target@IV.5.6.11.1.ii\endcsname{ega4.c04.sec.5-6-11-1-ii}
\expandafter\def\csname sga@target@IV.5.6.2\endcsname{ega4.c04.def.5-6-2}
\expandafter\def\csname sga@target@IV.5.6.3\endcsname{ega4.c04.sec.5-6-3}
\expandafter\def\csname sga@target@IV.5.6.3.i\endcsname{ega4.c04.sec.5-6-3-i}
\expandafter\def\csname sga@target@IV.5.6.3.ii\endcsname{ega4.c04.sec.5-6-3-ii}
\expandafter\def\csname sga@target@IV.5.6.3.iii\endcsname{ega4.c04.sec.5-6-3-iii}
\expandafter\def\csname sga@target@IV.5.6.3.iv\endcsname{ega4.c04.sec.5-6-3-iv}
\expandafter\def\csname sga@target@IV.5.6.4\endcsname{ega4.c04.prop.5-6-4}
\expandafter\def\csname sga@target@IV.5.6.5\endcsname{ega4.c04.prop.5-6-5}
\expandafter\def\csname sga@target@IV.5.6.5.1\endcsname{ega4.c04.eq.5-6-5-1}
\expandafter\def\csname sga@target@IV.5.6.5.2\endcsname{ega4.c04.eq.5-6-5-2}
\expandafter\def\csname sga@target@IV.5.6.5.3\endcsname{ega4.c04.eq.5-6-5-3}
\expandafter\def\csname sga@target@IV.5.6.6\endcsname{ega4.c04.cor.5-6-6}
\expandafter\def\csname sga@target@IV.5.6.6.1\endcsname{ega4.c04.eq.5-6-6-1}
\expandafter\def\csname sga@target@IV.5.6.6.2\endcsname{ega4.c04.eq.5-6-6-2}
\expandafter\def\csname sga@target@IV.5.6.7\endcsname{ega4.c04.cor.5-6-7}
\expandafter\def\csname sga@target@IV.5.6.7.1\endcsname{ega4.c04.eq.5-6-7-1}
\expandafter\def\csname sga@target@IV.5.6.8\endcsname{ega4.c04.cor.5-6-8}
\expandafter\def\csname sga@target@IV.5.6.8.1\endcsname{ega4.c04.eq.5-6-8-1}
\expandafter\def\csname sga@target@IV.5.6.8.2\endcsname{ega4.c04.eq.5-6-8-2}
\expandafter\def\csname sga@target@IV.5.6.9\endcsname{ega4.c04.sec.5-6-9}
\expandafter\def\csname sga@target@IV.5.6.9.i\endcsname{ega4.c04.sec.5-6-9-i}
\expandafter\def\csname sga@target@IV.5.6.9.ii\endcsname{ega4.c04.sec.5-6-9-ii}
\expandafter\def\csname sga@target@IV.5.7.1\endcsname{ega4.c04.def.5-7-1}
\expandafter\def\csname sga@target@IV.5.7.1.1\endcsname{ega4.c04.eq.5-7-1-1}
\expandafter\def\csname sga@target@IV.5.7.10\endcsname{ega4.c04.cor.5-7-10}
\expandafter\def\csname sga@target@IV.5.7.11\endcsname{ega4.c04.cor.5-7-11}
\expandafter\def\csname sga@target@IV.5.7.11.i\endcsname{ega4.c04.sec.5-7-11-i}
\expandafter\def\csname sga@target@IV.5.7.11.ii\endcsname{ega4.c04.sec.5-7-11-ii}
\expandafter\def\csname sga@target@IV.5.7.12\endcsname{ega4.c04.statement.5-7-12}
\expandafter\def\csname sga@target@IV.5.7.13\endcsname{ega4.c04.cor.5-7-13}
\expandafter\def\csname sga@target@IV.5.7.13.1\endcsname{ega4.c04.eq.5-7-13-1}
\expandafter\def\csname sga@target@IV.5.7.2\endcsname{ega4.c04.def.5-7-2}
\expandafter\def\csname sga@target@IV.5.7.2.1\endcsname{ega4.c04.eq.5-7-2-1}
\expandafter\def\csname sga@target@IV.5.7.2.2\endcsname{ega4.c04.eq.5-7-2-2}
\expandafter\def\csname sga@target@IV.5.7.3\endcsname{ega4.c04.sec.5-7-3}
\expandafter\def\csname sga@target@IV.5.7.3.i\endcsname{ega4.c04.sec.5-7-3-i}
\expandafter\def\csname sga@target@IV.5.7.3.ii\endcsname{ega4.c04.sec.5-7-3-ii}
\expandafter\def\csname sga@target@IV.5.7.3.iii\endcsname{ega4.c04.sec.5-7-3-iii}
\expandafter\def\csname sga@target@IV.5.7.3.iv\endcsname{ega4.c04.sec.5-7-3-iv}
\expandafter\def\csname sga@target@IV.5.7.3.v\endcsname{ega4.c04.sec.5-7-3-v}
\expandafter\def\csname sga@target@IV.5.7.3.vi\endcsname{ega4.c04.sec.5-7-3-vi}
\expandafter\def\csname sga@target@IV.5.7.4\endcsname{ega4.c04.prop.5-7-4}
\expandafter\def\csname sga@target@IV.5.7.4.1\endcsname{ega4.c04.eq.5-7-4-1}
\expandafter\def\csname sga@target@IV.5.7.4.2\endcsname{ega4.c04.eq.5-7-4-2}
\expandafter\def\csname sga@target@IV.5.7.4.i\endcsname{ega4.c04.sec.5-7-4-i}
\expandafter\def\csname sga@target@IV.5.7.4.ii\endcsname{ega4.c04.sec.5-7-4-ii}
\expandafter\def\csname sga@target@IV.5.7.5\endcsname{ega4.c04.prop.5-7-5}
\expandafter\def\csname sga@target@IV.5.7.6\endcsname{ega4.c04.cor.5-7-6}
\expandafter\def\csname sga@target@IV.5.7.7\endcsname{ega4.c04.cor.5-7-7}
\expandafter\def\csname sga@target@IV.5.7.8\endcsname{ega4.c04.sec.5-7-8}
\expandafter\def\csname sga@target@IV.5.7.9\endcsname{ega4.c04.prop.5-7-9}
\expandafter\def\csname sga@target@IV.5.7.9.1\endcsname{ega4.c04.eq.5-7-9-1}
\expandafter\def\csname sga@target@IV.5.7.9.2\endcsname{ega4.c04.eq.5-7-9-2}
\expandafter\def\csname sga@target@IV.5.8.1\endcsname{ega4.c04.statement.5-8-1}
\expandafter\def\csname sga@target@IV.5.8.2\endcsname{ega4.c04.def.5-8-2}
\expandafter\def\csname sga@target@IV.5.8.3\endcsname{ega4.c04.prop.5-8-3}
\expandafter\def\csname sga@target@IV.5.8.4\endcsname{ega4.c04.statement.5-8-4}
\expandafter\def\csname sga@target@IV.5.8.5\endcsname{ega4.c04.prop.5-8-5}
\expandafter\def\csname sga@target@IV.5.8.6\endcsname{ega4.c04.thm.5-8-6}
\expandafter\def\csname sga@target@IV.5.8.6.1\endcsname{ega4.c04.eq.5-8-6-1}
\expandafter\def\csname sga@target@IV.5.8.6.i\endcsname{ega4.c04.item.5-8-6-i}
\expandafter\def\csname sga@target@IV.5.8.6.ii\endcsname{ega4.c04.item.5-8-6-ii}
\expandafter\def\csname sga@target@IV.5.9.1\endcsname{ega4.c04.statement.5-9-1}
\expandafter\def\csname sga@target@IV.5.9.1.1\endcsname{ega4.c04.eq.5-9-1-1}
\expandafter\def\csname sga@target@IV.5.9.1.2\endcsname{ega4.c04.eq.5-9-1-2}
\expandafter\def\csname sga@target@IV.5.9.1.3\endcsname{ega4.c04.eq.5-9-1-3}
\expandafter\def\csname sga@target@IV.5.9.10\endcsname{ega4.c04.cor.5-9-10}
\expandafter\def\csname sga@target@IV.5.9.10.i\endcsname{ega4.c04.item.5-9-10-i}
\expandafter\def\csname sga@target@IV.5.9.10.ii\endcsname{ega4.c04.item.5-9-10-ii}
\expandafter\def\csname sga@target@IV.5.9.11\endcsname{ega4.c04.prop.5-9-11}
\expandafter\def\csname sga@target@IV.5.9.12\endcsname{ega4.c04.item.5-9-12}
\expandafter\def\csname sga@target@IV.5.9.12.i\endcsname{ega4.c04.item.5-9-12-i}
\expandafter\def\csname sga@target@IV.5.9.12.ii\endcsname{ega4.c04.item.5-9-12-ii}
\expandafter\def\csname sga@target@IV.5.9.12.iii\endcsname{ega4.c04.item.5-9-12-iii}
\expandafter\def\csname sga@target@IV.5.9.2\endcsname{ega4.c04.prop.5-9-2}
\expandafter\def\csname sga@target@IV.5.9.2.i\endcsname{ega4.c04.item.5-9-2-i}
\expandafter\def\csname sga@target@IV.5.9.2.ii\endcsname{ega4.c04.item.5-9-2-ii}
\expandafter\def\csname sga@target@IV.5.9.3\endcsname{ega4.c04.rem.5-9-3}
\expandafter\def\csname sga@target@IV.5.9.3.1\endcsname{ega4.c04.eq.5-9-3-1}
\expandafter\def\csname sga@target@IV.5.9.4\endcsname{ega4.c04.prop.5-9-4}
\expandafter\def\csname sga@target@IV.5.9.4.1\endcsname{ega4.c04.eq.5-9-4-1}
\expandafter\def\csname sga@target@IV.5.9.5\endcsname{ega4.c04.cor.5-9-5}
\expandafter\def\csname sga@target@IV.5.9.6\endcsname{ega4.c04.cor.5-9-6}
\expandafter\def\csname sga@target@IV.5.9.6.1\endcsname{ega4.c04.eq.5-9-6-1}
\expandafter\def\csname sga@target@IV.5.9.7\endcsname{ega4.c04.statement.5-9-7}
\expandafter\def\csname sga@target@IV.5.9.7.1\endcsname{ega4.c04.eq.5-9-7-1}
\expandafter\def\csname sga@target@IV.5.9.8\endcsname{ega4.c04.prop.5-9-8}
\expandafter\def\csname sga@target@IV.5.9.8.1\endcsname{ega4.c04.eq.5-9-8-1}
\expandafter\def\csname sga@target@IV.5.9.8.a\endcsname{ega4.c04.item.5-9-8-a}
\expandafter\def\csname sga@target@IV.5.9.8.a-prime\endcsname{ega4.c04.item.5-9-8-a-prime}
\expandafter\def\csname sga@target@IV.5.9.8.b\endcsname{ega4.c04.item.5-9-8-b}
\expandafter\def\csname sga@target@IV.5.9.8.b-prime\endcsname{ega4.c04.item.5-9-8-b-prime}
\expandafter\def\csname sga@target@IV.5.9.9\endcsname{ega4.c04.def.5-9-9}
\expandafter\def\csname sga@target@IV.6.1.1\endcsname{ega4.c04.prop.6-1-1}
\expandafter\def\csname sga@target@IV.6.1.1.1\endcsname{ega4.c04.eq.6-1-1-1}
\expandafter\def\csname sga@target@IV.6.1.2\endcsname{ega4.c04.cor.6-1-2}
\expandafter\def\csname sga@target@IV.6.1.2.1\endcsname{ega4.c04.eq.6-1-2-1}
\expandafter\def\csname sga@target@IV.6.1.3\endcsname{ega4.c04.cor.6-1-3}
\expandafter\def\csname sga@target@IV.6.1.4\endcsname{ega4.c04.cor.6-1-4}
\expandafter\def\csname sga@target@IV.6.1.4.1\endcsname{ega4.c04.eq.6-1-4-1}
\expandafter\def\csname sga@target@IV.6.1.5\endcsname{ega4.c04.prop.6-1-5}
\expandafter\def\csname sga@target@IV.6.10.1\endcsname{ega4.c04.statement.6-10-1}
\expandafter\def\csname sga@target@IV.6.10.2\endcsname{ega4.c04.prop.6-10-2}
\expandafter\def\csname sga@target@IV.6.10.3\endcsname{ega4.c04.prop.6-10-3}
\expandafter\def\csname sga@target@IV.6.10.3.1\endcsname{ega4.c04.eq.6-10-3-1}
\expandafter\def\csname sga@target@IV.6.10.3.2\endcsname{ega4.c04.eq.6-10-3-2}
\expandafter\def\csname sga@target@IV.6.10.3.3\endcsname{ega4.c04.eq.6-10-3-3}
\expandafter\def\csname sga@target@IV.6.10.4\endcsname{ega4.c04.cor.6-10-4}
\expandafter\def\csname sga@target@IV.6.10.5\endcsname{ega4.c04.prop.6-10-5}
\expandafter\def\csname sga@target@IV.6.10.5.1\endcsname{ega4.c04.eq.6-10-5-1}
\expandafter\def\csname sga@target@IV.6.10.5.2\endcsname{ega4.c04.eq.6-10-5-2}
\expandafter\def\csname sga@target@IV.6.10.5.3\endcsname{ega4.c04.eq.6-10-5-3}
\expandafter\def\csname sga@target@IV.6.10.5.4\endcsname{ega4.c04.eq.6-10-5-4}
\expandafter\def\csname sga@target@IV.6.10.6\endcsname{ega4.c04.prop.6-10-6}
\expandafter\def\csname sga@target@IV.6.10.6.1\endcsname{ega4.c04.eq.6-10-6-1}
\expandafter\def\csname sga@target@IV.6.10.6.2\endcsname{ega4.c04.eq.6-10-6-2}
\expandafter\def\csname sga@target@IV.6.10.6.3\endcsname{ega4.c04.eq.6-10-6-3}
\expandafter\def\csname sga@target@IV.6.11\endcsname{ega4.c04.subsec.6-11.key-iv-6-11}
\expandafter\def\csname sga@target@IV.6.11.1\endcsname{ega4.c04.lem.6-11-1}
\expandafter\def\csname sga@target@IV.6.11.2\endcsname{ega4.c04.prop.6-11-2}
\expandafter\def\csname sga@target@IV.6.11.2.1\endcsname{ega4.c04.eq.6-11-2-1}
\expandafter\def\csname sga@target@IV.6.11.2.2\endcsname{ega4.c04.eq.6-11-2-2}
\expandafter\def\csname sga@target@IV.6.11.3\endcsname{ega4.c04.cor.6-11-3}
\expandafter\def\csname sga@target@IV.6.11.4\endcsname{ega4.c04.statement.6-11-4}
\expandafter\def\csname sga@target@IV.6.11.5\endcsname{ega4.c04.prop.6-11-5}
\expandafter\def\csname sga@target@IV.6.11.5.1\endcsname{ega4.c04.eq.6-11-5-1}
\expandafter\def\csname sga@target@IV.6.11.6\endcsname{ega4.c04.prop.6-11-6}
\expandafter\def\csname sga@target@IV.6.11.6.1\endcsname{ega4.c04.lem.6-11-6-1}
\expandafter\def\csname sga@target@IV.6.11.7\endcsname{ega4.c04.cor.6-11-7}
\expandafter\def\csname sga@target@IV.6.11.8\endcsname{ega4.c04.prop.6-11-8}
\expandafter\def\csname sga@target@IV.6.11.8.1\endcsname{ega4.c04.eq.6-11-8-1}
\expandafter\def\csname sga@target@IV.6.11.9\endcsname{ega4.c04.statement.6-11-9}
\expandafter\def\csname sga@target@IV.6.12\endcsname{ega4.c04.subsec.6-12.key-iv-6-12}
\expandafter\def\csname sga@target@IV.6.12.1\endcsname{ega4.c04.statement.6-12-1}
\expandafter\def\csname sga@target@IV.6.12.10\endcsname{ega4.c04.prop.6-12-10}
\expandafter\def\csname sga@target@IV.6.12.2\endcsname{ega4.c04.prop.6-12-2}
\expandafter\def\csname sga@target@IV.6.12.3\endcsname{ega4.c04.cor.6-12-3}
\expandafter\def\csname sga@target@IV.6.12.4\endcsname{ega4.c04.thm.6-12-4}
\expandafter\def\csname sga@target@IV.6.12.4.1\endcsname{ega4.c04.lem.6-12-4-1}
\expandafter\def\csname sga@target@IV.6.12.5\endcsname{ega4.c04.cor.6-12-5}
\expandafter\def\csname sga@target@IV.6.12.6\endcsname{ega4.c04.cor.6-12-6}
\expandafter\def\csname sga@target@IV.6.12.7\endcsname{ega4.c04.thm.6-12-7}
\expandafter\def\csname sga@target@IV.6.12.8\endcsname{ega4.c04.cor.6-12-8}
\expandafter\def\csname sga@target@IV.6.12.9\endcsname{ega4.c04.prop.6-12-9}
\expandafter\def\csname sga@target@IV.6.13\endcsname{ega4.c04.subsec.6-13.key-iv-6-13}
\expandafter\def\csname sga@target@IV.6.13.1\endcsname{ega4.c04.statement.6-13-1}
\expandafter\def\csname sga@target@IV.6.13.2\endcsname{ega4.c04.prop.6-13-2}
\expandafter\def\csname sga@target@IV.6.13.3\endcsname{ega4.c04.cor.6-13-3}
\expandafter\def\csname sga@target@IV.6.13.4\endcsname{ega4.c04.prop.6-13-4}
\expandafter\def\csname sga@target@IV.6.13.5\endcsname{ega4.c04.cor.6-13-5}
\expandafter\def\csname sga@target@IV.6.13.6\endcsname{ega4.c04.prop.6-13-6}
\expandafter\def\csname sga@target@IV.6.13.7\endcsname{ega4.c04.prop.6-13-7}
\expandafter\def\csname sga@target@IV.6.13.7.1\endcsname{ega4.c04.lem.6-13-7-1}
\expandafter\def\csname sga@target@IV.6.14.1\endcsname{ega4.c04.prop.6-14-1}
\expandafter\def\csname sga@target@IV.6.14.1.1\endcsname{ega4.c04.lem.6-14-1-1}
\expandafter\def\csname sga@target@IV.6.14.1.2\endcsname{ega4.c04.lem.6-14-1-2}
\expandafter\def\csname sga@target@IV.6.14.1.2.env\endcsname{ega4.c04.statement.6-14-1-2-env}
\expandafter\def\csname sga@target@IV.6.14.1.3\endcsname{ega4.c04.eq.6-14-1-3}
\expandafter\def\csname sga@target@IV.6.14.1.4\endcsname{ega4.c04.lem.6-14-1-4}
\expandafter\def\csname sga@target@IV.6.14.2\endcsname{ega4.c04.cor.6-14-2}
\expandafter\def\csname sga@target@IV.6.14.3\endcsname{ega4.c04.cor.6-14-3}
\expandafter\def\csname sga@target@IV.6.14.4\endcsname{ega4.c04.prop.6-14-4}
\expandafter\def\csname sga@target@IV.6.14.4.1\endcsname{ega4.c04.lem.6-14-4-1}
\expandafter\def\csname sga@target@IV.6.14.5\endcsname{ega4.c04.cor.6-14-5}
\expandafter\def\csname sga@target@IV.6.15.1\endcsname{ega4.c04.statement.6-15-1}
\expandafter\def\csname sga@target@IV.6.15.10\endcsname{ega4.c04.prop.6-15-10}
\expandafter\def\csname sga@target@IV.6.15.11\endcsname{ega4.c04.item.6-15-11}
\expandafter\def\csname sga@target@IV.6.15.2\endcsname{ega4.c04.statement.6-15-2}
\expandafter\def\csname sga@target@IV.6.15.3\endcsname{ega4.c04.statement.6-15-3}
\expandafter\def\csname sga@target@IV.6.15.3.1\endcsname{ega4.c04.lem.6-15-3-1}
\expandafter\def\csname sga@target@IV.6.15.4\endcsname{ega4.c04.statement.6-15-4}
\expandafter\def\csname sga@target@IV.6.15.4.1\endcsname{ega4.c04.lem.6-15-4-1}
\expandafter\def\csname sga@target@IV.6.15.5\endcsname{ega4.c04.prop.6-15-5}
\expandafter\def\csname sga@target@IV.6.15.6\endcsname{ega4.c04.prop.6-15-6}
\expandafter\def\csname sga@target@IV.6.15.7\endcsname{ega4.c04.prop.6-15-7}
\expandafter\def\csname sga@target@IV.6.15.8\endcsname{ega4.c04.lem.6-15-8}
\expandafter\def\csname sga@target@IV.6.15.9\endcsname{ega4.c04.cor.6-15-9}
\expandafter\def\csname sga@target@IV.6.2.1\endcsname{ega4.c04.prop.6-2-1}
\expandafter\def\csname sga@target@IV.6.2.1.1\endcsname{ega4.c04.eq.6-2-1-1}
\expandafter\def\csname sga@target@IV.6.2.1.2\endcsname{ega4.c04.eq.6-2-1-2}
\expandafter\def\csname sga@target@IV.6.2.2\endcsname{ega4.c04.cor.6-2-2}
\expandafter\def\csname sga@target@IV.6.2.3\endcsname{ega4.c04.prop.6-2-3}
\expandafter\def\csname sga@target@IV.6.2.3.1\endcsname{ega4.c04.eq.6-2-3-1}
\expandafter\def\csname sga@target@IV.6.2.4\endcsname{ega4.c04.prop.6-2-4}
\expandafter\def\csname sga@target@IV.6.2.4.1\endcsname{ega4.c04.eq.6-2-4-1}
\expandafter\def\csname sga@target@IV.6.3.1\endcsname{ega4.c04.prop.6-3-1}
\expandafter\def\csname sga@target@IV.6.3.1.1\endcsname{ega4.c04.eq.6-3-1-1}
\expandafter\def\csname sga@target@IV.6.3.2\endcsname{ega4.c04.cor.6-3-2}
\expandafter\def\csname sga@target@IV.6.3.2.1\endcsname{ega4.c04.eq.6-3-2-1}
\expandafter\def\csname sga@target@IV.6.3.3\endcsname{ega4.c04.cor.6-3-3}
\expandafter\def\csname sga@target@IV.6.3.4\endcsname{ega4.c04.cor.6-3-4}
\expandafter\def\csname sga@target@IV.6.3.4.1\endcsname{ega4.c04.eq.6-3-4-1}
\expandafter\def\csname sga@target@IV.6.3.4.2\endcsname{ega4.c04.eq.6-3-4-2}
\expandafter\def\csname sga@target@IV.6.3.5\endcsname{ega4.c04.cor.6-3-5}
\expandafter\def\csname sga@target@IV.6.3.6\endcsname{ega4.c04.cor.6-3-6}
\expandafter\def\csname sga@target@IV.6.3.7\endcsname{ega4.c04.cor.6-3-7}
\expandafter\def\csname sga@target@IV.6.3.8\endcsname{ega4.c04.prop.6-3-8}
\expandafter\def\csname sga@target@IV.6.4.1\endcsname{ega4.c04.prop.6-4-1}
\expandafter\def\csname sga@target@IV.6.4.2\endcsname{ega4.c04.cor.6-4-2}
\expandafter\def\csname sga@target@IV.6.4.3\endcsname{ega4.c04.rem.6-4-3}
\expandafter\def\csname sga@target@IV.6.5.1\endcsname{ega4.c04.prop.6-5-1}
\expandafter\def\csname sga@target@IV.6.5.2\endcsname{ega4.c04.cor.6-5-2}
\expandafter\def\csname sga@target@IV.6.5.3\endcsname{ega4.c04.prop.6-5-3}
\expandafter\def\csname sga@target@IV.6.5.4\endcsname{ega4.c04.cor.6-5-4}
\expandafter\def\csname sga@target@IV.6.5.5\endcsname{ega4.c04.item.6-5-5}
\expandafter\def\csname sga@target@IV.6.6.1\endcsname{ega4.c04.prop.6-6-1}
\expandafter\def\csname sga@target@IV.6.6.2\endcsname{ega4.c04.rem.6-6-2}
\expandafter\def\csname sga@target@IV.6.7.1\endcsname{ega4.c04.prop.6-7-1}
\expandafter\def\csname sga@target@IV.6.7.1.1\endcsname{ega4.c04.lem.6-7-1-1}
\expandafter\def\csname sga@target@IV.6.7.2\endcsname{ega4.c04.cor.6-7-2}
\expandafter\def\csname sga@target@IV.6.7.3\endcsname{ega4.c04.cor.6-7-3}
\expandafter\def\csname sga@target@IV.6.7.4\endcsname{ega4.c04.prop.6-7-4}
\expandafter\def\csname sga@target@IV.6.7.4.1\endcsname{ega4.c04.lem.6-7-4-1}
\expandafter\def\csname sga@target@IV.6.7.5\endcsname{ega4.c04.item.6-7-5}
\expandafter\def\csname sga@target@IV.6.7.6\endcsname{ega4.c04.def.6-7-6}
\expandafter\def\csname sga@target@IV.6.7.7\endcsname{ega4.c04.prop.6-7-7}
\expandafter\def\csname sga@target@IV.6.7.8\endcsname{ega4.c04.cor.6-7-8}
\expandafter\def\csname sga@target@IV.6.8.1\endcsname{ega4.c04.def.6-8-1}
\expandafter\def\csname sga@target@IV.6.8.2\endcsname{ega4.c04.prop.6-8-2}
\expandafter\def\csname sga@target@IV.6.8.3\endcsname{ega4.c04.prop.6-8-3}
\expandafter\def\csname sga@target@IV.6.8.4\endcsname{ega4.c04.item.6-8-4}
\expandafter\def\csname sga@target@IV.6.8.5\endcsname{ega4.c04.prop.6-8-5}
\expandafter\def\csname sga@target@IV.6.8.6\endcsname{ega4.c04.thm.6-8-6}
\expandafter\def\csname sga@target@IV.6.8.7\endcsname{ega4.c04.cor.6-8-7}
\expandafter\def\csname sga@target@IV.6.8.8\endcsname{ega4.c04.rem.6-8-8}
\expandafter\def\csname sga@target@IV.6.9.1\endcsname{ega4.c04.thm.6-9-1}
\expandafter\def\csname sga@target@IV.6.9.2\endcsname{ega4.c04.lem.6-9-2}
\expandafter\def\csname sga@target@IV.6.9.3\endcsname{ega4.c04.cor.6-9-3}
\expandafter\def\csname sga@target@IV.7.1.1\endcsname{ega4.c04.def.7-1-1}
\expandafter\def\csname sga@target@IV.7.1.10\endcsname{ega4.c04.prop.7-1-10}
\expandafter\def\csname sga@target@IV.7.1.11\endcsname{ega4.c04.prop.7-1-11}
\expandafter\def\csname sga@target@IV.7.1.12\endcsname{ega4.c04.item.7-1-12}
\expandafter\def\csname sga@target@IV.7.1.13\endcsname{ega4.c04.cor.7-1-13}
\expandafter\def\csname sga@target@IV.7.1.2\endcsname{ega4.c04.prop.7-1-2}
\expandafter\def\csname sga@target@IV.7.1.3\endcsname{ega4.c04.prop.7-1-3}
\expandafter\def\csname sga@target@IV.7.1.3.1\endcsname{ega4.c04.eq.7-1-3-1}
\expandafter\def\csname sga@target@IV.7.1.3.2\endcsname{ega4.c04.eq.7-1-3-2}
\expandafter\def\csname sga@target@IV.7.1.3.3\endcsname{ega4.c04.eq.7-1-3-3}
\expandafter\def\csname sga@target@IV.7.1.3.4\endcsname{ega4.c04.eq.7-1-3-4}
\expandafter\def\csname sga@target@IV.7.1.3.5\endcsname{ega4.c04.eq.7-1-3-5}
\expandafter\def\csname sga@target@IV.7.1.3.6\endcsname{ega4.c04.eq.7-1-3-6}
\expandafter\def\csname sga@target@IV.7.1.3.7\endcsname{ega4.c04.eq.7-1-3-7}
\expandafter\def\csname sga@target@IV.7.1.4\endcsname{ega4.c04.cor.7-1-4}
\expandafter\def\csname sga@target@IV.7.1.5\endcsname{ega4.c04.cor.7-1-5}
\expandafter\def\csname sga@target@IV.7.1.6\endcsname{ega4.c04.rem.7-1-6}
\expandafter\def\csname sga@target@IV.7.1.7\endcsname{ega4.c04.lem.7-1-7}
\expandafter\def\csname sga@target@IV.7.1.8\endcsname{ega4.c04.thm.7-1-8}
\expandafter\def\csname sga@target@IV.7.1.9\endcsname{ega4.c04.def.7-1-9}
\expandafter\def\csname sga@target@IV.7.2.1\endcsname{ega4.c04.statement.7-2-1}
\expandafter\def\csname sga@target@IV.7.2.1.1\endcsname{ega4.c04.statement.7-2-1-1}
\expandafter\def\csname sga@target@IV.7.2.10\endcsname{ega4.c04.statement.7-2-10}
\expandafter\def\csname sga@target@IV.7.2.2\endcsname{ega4.c04.prop.7-2-2}
\expandafter\def\csname sga@target@IV.7.2.3\endcsname{ega4.c04.prop.7-2-3}
\expandafter\def\csname sga@target@IV.7.2.4\endcsname{ega4.c04.prop.7-2-4}
\expandafter\def\csname sga@target@IV.7.2.5\endcsname{ega4.c04.thm.7-2-5}
\expandafter\def\csname sga@target@IV.7.2.5.1\endcsname{ega4.c04.eq.7-2-5-1}
\expandafter\def\csname sga@target@IV.7.2.6\endcsname{ega4.c04.def.7-2-6}
\expandafter\def\csname sga@target@IV.7.2.7\endcsname{ega4.c04.cor.7-2-7}
\expandafter\def\csname sga@target@IV.7.2.8\endcsname{ega4.c04.cor.7-2-8}
\expandafter\def\csname sga@target@IV.7.2.9\endcsname{ega4.c04.cor.7-2-9}
\expandafter\def\csname sga@target@IV.7.3.1\endcsname{ega4.c04.statement.7-3-1}
\expandafter\def\csname sga@target@IV.7.3.10\endcsname{ega4.c04.statement.7-3-10}
\expandafter\def\csname sga@target@IV.7.3.11\endcsname{ega4.c04.rem.7-3-11}
\expandafter\def\csname sga@target@IV.7.3.12\endcsname{ega4.c04.prop.7-3-12}
\expandafter\def\csname sga@target@IV.7.3.12.1\endcsname{ega4.c04.eq.7-3-12-1}
\expandafter\def\csname sga@target@IV.7.3.13\endcsname{ega4.c04.statement.7-3-13}
\expandafter\def\csname sga@target@IV.7.3.14\endcsname{ega4.c04.prop.7-3-14}
\expandafter\def\csname sga@target@IV.7.3.15\endcsname{ega4.c04.prop.7-3-15}
\expandafter\def\csname sga@target@IV.7.3.16\endcsname{ega4.c04.prop.7-3-16}
\expandafter\def\csname sga@target@IV.7.3.16.1\endcsname{ega4.c04.eq.7-3-16-1}
\expandafter\def\csname sga@target@IV.7.3.16.2\endcsname{ega4.c04.lem.7-3-16-2}
\expandafter\def\csname sga@target@IV.7.3.17\endcsname{ega4.c04.prop.7-3-17}
\expandafter\def\csname sga@target@IV.7.3.18\endcsname{ega4.c04.prop.7-3-18}
\expandafter\def\csname sga@target@IV.7.3.19\endcsname{ega4.c04.statement.7-3-19}
\expandafter\def\csname sga@target@IV.7.3.2\endcsname{ega4.c04.lem.7-3-2}
\expandafter\def\csname sga@target@IV.7.3.3\endcsname{ega4.c04.cor.7-3-3}
\expandafter\def\csname sga@target@IV.7.3.4\endcsname{ega4.c04.statement.7-3-4}
\expandafter\def\csname sga@target@IV.7.3.5\endcsname{ega4.c04.item.7-3-5}
\expandafter\def\csname sga@target@IV.7.3.6\endcsname{ega4.c04.statement.7-3-6}
\expandafter\def\csname sga@target@IV.7.3.7\endcsname{ega4.c04.lem.7-3-7}
\expandafter\def\csname sga@target@IV.7.3.8\endcsname{ega4.c04.statement.7-3-8}
\expandafter\def\csname sga@target@IV.7.3.9\endcsname{ega4.c04.item.7-3-9}
\expandafter\def\csname sga@target@IV.7.4.1\endcsname{ega4.c04.statement.7-4-1}
\expandafter\def\csname sga@target@IV.7.4.2\endcsname{ega4.c04.lem.7-4-2}
\expandafter\def\csname sga@target@IV.7.4.3\endcsname{ega4.c04.cor.7-4-3}
\expandafter\def\csname sga@target@IV.7.4.4\endcsname{ega4.c04.thm.7-4-4}
\expandafter\def\csname sga@target@IV.7.4.4.1\endcsname{ega4.c04.lem.7-4-4-1}
\expandafter\def\csname sga@target@IV.7.4.4.2\endcsname{ega4.c04.lem.7-4-4-2}
\expandafter\def\csname sga@target@IV.7.4.5\endcsname{ega4.c04.cor.7-4-5}
\expandafter\def\csname sga@target@IV.7.4.6\endcsname{ega4.c04.prop.7-4-6}
\expandafter\def\csname sga@target@IV.7.4.7\endcsname{ega4.c04.cor.7-4-7}
\expandafter\def\csname sga@target@IV.7.4.8\endcsname{ega4.c04.rem.7-4-8}
\expandafter\def\csname sga@target@IV.7.5.0\endcsname{ega4.c04.statement.7-5-0}
\expandafter\def\csname sga@target@IV.7.5.1\endcsname{ega4.c04.thm.7-5-1}
\expandafter\def\csname sga@target@IV.7.5.1.1\endcsname{ega4.c04.lem.7-5-1-1}
\expandafter\def\csname sga@target@IV.7.5.2\endcsname{ega4.c04.cor.7-5-2}
\expandafter\def\csname sga@target@IV.7.5.3\endcsname{ega4.c04.item.7-5-3}
\expandafter\def\csname sga@target@IV.7.5.4\endcsname{ega4.c04.item.7-5-4}
\expandafter\def\csname sga@target@IV.7.5.5\endcsname{ega4.c04.lem.7-5-5}
\expandafter\def\csname sga@target@IV.7.5.6\endcsname{ega4.c04.prop.7-5-6}
\expandafter\def\csname sga@target@IV.7.5.7\endcsname{ega4.c04.cor.7-5-7}
\expandafter\def\csname sga@target@IV.7.5.7.1\endcsname{ega4.c04.lem.7-5-7-1}
\expandafter\def\csname sga@target@IV.7.5.8\endcsname{ega4.c04.rem.7-5-8}
\expandafter\def\csname sga@target@IV.7.6.1\endcsname{ega4.c04.prop.7-6-1}
\expandafter\def\csname sga@target@IV.7.6.2\endcsname{ega4.c04.cor.7-6-2}
\expandafter\def\csname sga@target@IV.7.6.3\endcsname{ega4.c04.cor.7-6-3}
\expandafter\def\csname sga@target@IV.7.6.4\endcsname{ega4.c04.thm.7-6-4}
\expandafter\def\csname sga@target@IV.7.6.4.1\endcsname{ega4.c04.lem.7-6-4-1}
\expandafter\def\csname sga@target@IV.7.6.5\endcsname{ega4.c04.cor.7-6-5}
\expandafter\def\csname sga@target@IV.7.6.6\endcsname{ega4.c04.cor.7-6-6}
\expandafter\def\csname sga@target@IV.7.6.7\endcsname{ega4.c04.item.7-6-7}
\expandafter\def\csname sga@target@IV.7.7.1\endcsname{ega4.c04.statement.7-7-1}
\expandafter\def\csname sga@target@IV.7.7.2\endcsname{ega4.c04.thm.7-7-2}
\expandafter\def\csname sga@target@IV.7.7.3\endcsname{ega4.c04.cor.7-7-3}
\expandafter\def\csname sga@target@IV.7.7.4\endcsname{ega4.c04.cor.7-7-4}
\expandafter\def\csname sga@target@IV.7.8.1\endcsname{ega4.c04.statement.7-8-1}
\expandafter\def\csname sga@target@IV.7.8.2\endcsname{ega4.c04.def.7-8-2}
\expandafter\def\csname sga@target@IV.7.8.3\endcsname{ega4.c04.statement.7-8-3}
\expandafter\def\csname sga@target@IV.7.8.3.1\endcsname{ega4.c04.statement.7-8-3-1}
\expandafter\def\csname sga@target@IV.7.8.4\endcsname{ega4.c04.item.7-8-4}
\expandafter\def\csname sga@target@IV.7.8.4.1\endcsname{ega4.c04.statement.7-8-4-1}
\expandafter\def\csname sga@target@IV.7.8.5\endcsname{ega4.c04.statement.7-8-5}
\expandafter\def\csname sga@target@IV.7.8.6\endcsname{ega4.c04.prop.7-8-6}
\expandafter\def\csname sga@target@IV.7.9.1\endcsname{ega4.c04.statement.7-9-1}
\expandafter\def\csname sga@target@IV.7.9.10\endcsname{ega4.c04.item.7-9-10}
\expandafter\def\csname sga@target@IV.7.9.2\endcsname{ega4.c04.statement.7-9-2}
\expandafter\def\csname sga@target@IV.7.9.3\endcsname{ega4.c04.prop.7-9-3}
\expandafter\def\csname sga@target@IV.7.9.3.1\endcsname{ega4.c04.lem.7-9-3-1}
\expandafter\def\csname sga@target@IV.7.9.4\endcsname{ega4.c04.cor.7-9-4}
\expandafter\def\csname sga@target@IV.7.9.5\endcsname{ega4.c04.prop.7-9-5}
\expandafter\def\csname sga@target@IV.7.9.6\endcsname{ega4.c04.rem.7-9-6}
\expandafter\def\csname sga@target@IV.7.9.7\endcsname{ega4.c04.statement.7-9-7}
\expandafter\def\csname sga@target@IV.7.9.8\endcsname{ega4.c04.prop.7-9-8}
\expandafter\def\csname sga@target@IV.7.9.8.1\endcsname{ega4.c04.statement.7-9-8-1}
\expandafter\def\csname sga@target@IV.7.9.9\endcsname{ega4.c04.cor.7-9-9}
\expandafter\def\csname sga@target@IV.8.1.1\endcsname{ega4.c04.statement.8-1-1}
\expandafter\def\csname sga@target@IV.8.1.2\endcsname{ega4.c04.statement.8-1-2}
\expandafter\def\csname sga@target@IV.8.1.3\endcsname{ega4.c04.statement.8-1-3}
\expandafter\def\csname sga@target@IV.8.10.1\endcsname{ega4.c04.prop.8-10-1}
\expandafter\def\csname sga@target@IV.8.10.2\endcsname{ega4.c04.cor.8-10-2}
\expandafter\def\csname sga@target@IV.8.10.3\endcsname{ega4.c04.prop.8-10-3}
\expandafter\def\csname sga@target@IV.8.10.3.1\endcsname{ega4.c04.rem.8-10-3-1}
\expandafter\def\csname sga@target@IV.8.10.4\endcsname{ega4.c04.prop.8-10-4}
\expandafter\def\csname sga@target@IV.8.10.5\endcsname{ega4.c04.thm.8-10-5}
\expandafter\def\csname sga@target@IV.8.10.5.1\endcsname{ega4.c04.prop.8-10-5-1}
\expandafter\def\csname sga@target@IV.8.10.5.2\endcsname{ega4.c04.lem.8-10-5-2}
\expandafter\def\csname sga@target@IV.8.11\endcsname{ega4.c04.subsec.8-11.key-iv-8-11}
\expandafter\def\csname sga@target@IV.8.11.1\endcsname{ega4.c04.thm.8-11-1}
\expandafter\def\csname sga@target@IV.8.11.2\endcsname{ega4.c04.thm.8-11-2}
\expandafter\def\csname sga@target@IV.8.11.3\endcsname{ega4.c04.rem.8-11-3}
\expandafter\def\csname sga@target@IV.8.11.4\endcsname{ega4.c04.statement.8-11-4}
\expandafter\def\csname sga@target@IV.8.11.5\endcsname{ega4.c04.prop.8-11-5}
\expandafter\def\csname sga@target@IV.8.11.5.1\endcsname{ega4.c04.rem.8-11-5-1}
\expandafter\def\csname sga@target@IV.8.11.6\endcsname{ega4.c04.prop.8-11-6}
\expandafter\def\csname sga@target@IV.8.12\endcsname{ega4.c04.subsec.8-12.key-iv-8-12}
\expandafter\def\csname sga@target@IV.8.12.1\endcsname{ega4.c04.lem.8-12-1}
\expandafter\def\csname sga@target@IV.8.12.10\endcsname{ega4.c04.cor.8-12-10}
\expandafter\def\csname sga@target@IV.8.12.10.1\endcsname{ega4.c04.lem.8-12-10-1}
\expandafter\def\csname sga@target@IV.8.12.11\endcsname{ega4.c04.cor.8-12-11}
\expandafter\def\csname sga@target@IV.8.12.12\endcsname{ega4.c04.cor.8-12-12}
\expandafter\def\csname sga@target@IV.8.12.2\endcsname{ega4.c04.lem.8-12-2}
\expandafter\def\csname sga@target@IV.8.12.3\endcsname{ega4.c04.prop.8-12-3}
\expandafter\def\csname sga@target@IV.8.12.4\endcsname{ega4.c04.cor.8-12-4}
\expandafter\def\csname sga@target@IV.8.12.5\endcsname{ega4.c04.statement.8-12-5}
\expandafter\def\csname sga@target@IV.8.12.6\endcsname{ega4.c04.thm.8-12-6}
\expandafter\def\csname sga@target@IV.8.12.6.1\endcsname{ega4.c04.eq.8-12-6-1}
\expandafter\def\csname sga@target@IV.8.12.7\endcsname{ega4.c04.rem.8-12-7}
\expandafter\def\csname sga@target@IV.8.12.8\endcsname{ega4.c04.cor.8-12-8}
\expandafter\def\csname sga@target@IV.8.12.9\endcsname{ega4.c04.cor.8-12-9}
\expandafter\def\csname sga@target@IV.8.13.1\endcsname{ega4.c04.prop.8-13-1}
\expandafter\def\csname sga@target@IV.8.13.1.1\endcsname{ega4.c04.eq.8-13-1-1}
\expandafter\def\csname sga@target@IV.8.13.2\endcsname{ega4.c04.cor.8-13-2}
\expandafter\def\csname sga@target@IV.8.13.2.1\endcsname{ega4.c04.eq.8-13-2-1}
\expandafter\def\csname sga@target@IV.8.13.3\endcsname{ega4.c04.statement.8-13-3}
\expandafter\def\csname sga@target@IV.8.13.4\endcsname{ega4.c04.statement.8-13-4}
\expandafter\def\csname sga@target@IV.8.13.4.1\endcsname{ega4.c04.eq.8-13-4-1}
\expandafter\def\csname sga@target@IV.8.13.4.2\endcsname{ega4.c04.eq.8-13-4-2}
\expandafter\def\csname sga@target@IV.8.13.5\endcsname{ega4.c04.prop.8-13-5}
\expandafter\def\csname sga@target@IV.8.13.6\endcsname{ega4.c04.rem.8-13-6}
\expandafter\def\csname sga@target@IV.8.14.1\endcsname{ega4.c04.statement.8-14-1}
\expandafter\def\csname sga@target@IV.8.14.2\endcsname{ega4.c04.prop.8-14-2}
\expandafter\def\csname sga@target@IV.8.14.2.1\endcsname{ega4.c04.eq.8-14-2-1}
\expandafter\def\csname sga@target@IV.8.14.2.2\endcsname{ega4.c04.cor.8-14-2-2}
\expandafter\def\csname sga@target@IV.8.14.2.3\endcsname{ega4.c04.eq.8-14-2-3}
\expandafter\def\csname sga@target@IV.8.14.2.4\endcsname{ega4.c04.eq.8-14-2-4}
\expandafter\def\csname sga@target@IV.8.14.3\endcsname{ega4.c04.rem.8-14-3}
\expandafter\def\csname sga@target@IV.8.2.1\endcsname{ega4.c04.statement.8-2-1}
\expandafter\def\csname sga@target@IV.8.2.1.1\endcsname{ega4.c04.eq.8-2-1-1}
\expandafter\def\csname sga@target@IV.8.2.10\endcsname{ega4.c04.cor.8-2-10}
\expandafter\def\csname sga@target@IV.8.2.11\endcsname{ega4.c04.cor.8-2-11}
\expandafter\def\csname sga@target@IV.8.2.12\endcsname{ega4.c04.cor.8-2-12}
\expandafter\def\csname sga@target@IV.8.2.12.1\endcsname{ega4.c04.eq.8-2-12-1}
\expandafter\def\csname sga@target@IV.8.2.13\endcsname{ega4.c04.prop.8-2-13}
\expandafter\def\csname sga@target@IV.8.2.14\endcsname{ega4.c04.rem.8-2-14}
\expandafter\def\csname sga@target@IV.8.2.2\endcsname{ega4.c04.statement.8-2-2}
\expandafter\def\csname sga@target@IV.8.2.3\endcsname{ega4.c04.prop.8-2-3}
\expandafter\def\csname sga@target@IV.8.2.3.1\endcsname{ega4.c04.eq.8-2-3-1}
\expandafter\def\csname sga@target@IV.8.2.4\endcsname{ega4.c04.lem.8-2-4}
\expandafter\def\csname sga@target@IV.8.2.5\endcsname{ega4.c04.prop.8-2-5}
\expandafter\def\csname sga@target@IV.8.2.6\endcsname{ega4.c04.lem.8-2-6}
\expandafter\def\csname sga@target@IV.8.2.7\endcsname{ega4.c04.rem.8-2-7}
\expandafter\def\csname sga@target@IV.8.2.8\endcsname{ega4.c04.statement.8-2-8}
\expandafter\def\csname sga@target@IV.8.2.9\endcsname{ega4.c04.prop.8-2-9}
\expandafter\def\csname sga@target@IV.8.3.1\endcsname{ega4.c04.statement.8-3-1}
\expandafter\def\csname sga@target@IV.8.3.10\endcsname{ega4.c04.statement.8-3-10}
\expandafter\def\csname sga@target@IV.8.3.10.1\endcsname{ega4.c04.eq.8-3-10-1}
\expandafter\def\csname sga@target@IV.8.3.11\endcsname{ega4.c04.thm.8-3-11}
\expandafter\def\csname sga@target@IV.8.3.12\endcsname{ega4.c04.cor.8-3-12}
\expandafter\def\csname sga@target@IV.8.3.13\endcsname{ega4.c04.prop.8-3-13}
\expandafter\def\csname sga@target@IV.8.3.2\endcsname{ega4.c04.thm.8-3-2}
\expandafter\def\csname sga@target@IV.8.3.2.1\endcsname{ega4.c04.eq.8-3-2-1}
\expandafter\def\csname sga@target@IV.8.3.3\endcsname{ega4.c04.cor.8-3-3}
\expandafter\def\csname sga@target@IV.8.3.4\endcsname{ega4.c04.cor.8-3-4}
\expandafter\def\csname sga@target@IV.8.3.5\endcsname{ega4.c04.cor.8-3-5}
\expandafter\def\csname sga@target@IV.8.3.6\endcsname{ega4.c04.cor.8-3-6}
\expandafter\def\csname sga@target@IV.8.3.7\endcsname{ega4.c04.cor.8-3-7}
\expandafter\def\csname sga@target@IV.8.3.7.1\endcsname{ega4.c04.eq.8-3-7-1}
\expandafter\def\csname sga@target@IV.8.3.8\endcsname{ega4.c04.prop.8-3-8}
\expandafter\def\csname sga@target@IV.8.3.9\endcsname{ega4.c04.statement.8-3-9}
\expandafter\def\csname sga@target@IV.8.3.9.1\endcsname{ega4.c04.eq.8-3-9-1}
\expandafter\def\csname sga@target@IV.8.3.9.2\endcsname{ega4.c04.eq.8-3-9-2}
\expandafter\def\csname sga@target@IV.8.3.9.3\endcsname{ega4.c04.eq.8-3-9-3}
\expandafter\def\csname sga@target@IV.8.3.9.4\endcsname{ega4.c04.eq.8-3-9-4}
\expandafter\def\csname sga@target@IV.8.3.9.5\endcsname{ega4.c04.eq.8-3-9-5}
\expandafter\def\csname sga@target@IV.8.4.1\endcsname{ega4.c04.prop.8-4-1}
\expandafter\def\csname sga@target@IV.8.4.2\endcsname{ega4.c04.prop.8-4-2}
\expandafter\def\csname sga@target@IV.8.4.2.1\endcsname{ega4.c04.lem.8-4-2-1}
\expandafter\def\csname sga@target@IV.8.4.3\endcsname{ega4.c04.cor.8-4-3}
\expandafter\def\csname sga@target@IV.8.4.4\endcsname{ega4.c04.prop.8-4-4}
\expandafter\def\csname sga@target@IV.8.4.5\endcsname{ega4.c04.cor.8-4-5}
\expandafter\def\csname sga@target@IV.8.4.6\endcsname{ega4.c04.rem.8-4-6}
\expandafter\def\csname sga@target@IV.8.5.1\endcsname{ega4.c04.statement.8-5-1}
\expandafter\def\csname sga@target@IV.8.5.1.1\endcsname{ega4.c04.eq.8-5-1-1}
\expandafter\def\csname sga@target@IV.8.5.1.2\endcsname{ega4.c04.eq.8-5-1-2}
\expandafter\def\csname sga@target@IV.8.5.1.3\endcsname{ega4.c04.eq.8-5-1-3}
\expandafter\def\csname sga@target@IV.8.5.1.4\endcsname{ega4.c04.eq.8-5-1-4}
\expandafter\def\csname sga@target@IV.8.5.10\endcsname{ega4.c04.statement.8-5-10}
\expandafter\def\csname sga@target@IV.8.5.10.1\endcsname{ega4.c04.eq.8-5-10-1}
\expandafter\def\csname sga@target@IV.8.5.10.2\endcsname{ega4.c04.eq.8-5-10-2}
\expandafter\def\csname sga@target@IV.8.5.11\endcsname{ega4.c04.prop.8-5-11}
\expandafter\def\csname sga@target@IV.8.5.2\endcsname{ega4.c04.thm.8-5-2}
\expandafter\def\csname sga@target@IV.8.5.2.1\endcsname{ega4.c04.lem.8-5-2-1}
\expandafter\def\csname sga@target@IV.8.5.2.2\endcsname{ega4.c04.eq.8-5-2-2}
\expandafter\def\csname sga@target@IV.8.5.2.3\endcsname{ega4.c04.eq.8-5-2-3}
\expandafter\def\csname sga@target@IV.8.5.2.4\endcsname{ega4.c04.cor.8-5-2-4}
\expandafter\def\csname sga@target@IV.8.5.2.5\endcsname{ega4.c04.cor.8-5-2-5}
\expandafter\def\csname sga@target@IV.8.5.3\endcsname{ega4.c04.statement.8-5-3}
\expandafter\def\csname sga@target@IV.8.5.4\endcsname{ega4.c04.cor.8-5-4}
\expandafter\def\csname sga@target@IV.8.5.5\endcsname{ega4.c04.prop.8-5-5}
\expandafter\def\csname sga@target@IV.8.5.6\endcsname{ega4.c04.prop.8-5-6}
\expandafter\def\csname sga@target@IV.8.5.7\endcsname{ega4.c04.cor.8-5-7}
\expandafter\def\csname sga@target@IV.8.5.8\endcsname{ega4.c04.cor.8-5-8}
\expandafter\def\csname sga@target@IV.8.5.9\endcsname{ega4.c04.lem.8-5-9}
\expandafter\def\csname sga@target@IV.8.6.1\endcsname{ega4.c04.statement.8-6-1}
\expandafter\def\csname sga@target@IV.8.6.2\endcsname{ega4.c04.statement.8-6-2}
\expandafter\def\csname sga@target@IV.8.6.2.1\endcsname{ega4.c04.eq.8-6-2-1}
\expandafter\def\csname sga@target@IV.8.6.2.2\endcsname{ega4.c04.eq.8-6-2-2}
\expandafter\def\csname sga@target@IV.8.6.2.3\endcsname{ega4.c04.eq.8-6-2-3}
\expandafter\def\csname sga@target@IV.8.6.3\endcsname{ega4.c04.prop.8-6-3}
\expandafter\def\csname sga@target@IV.8.6.3.1\endcsname{ega4.c04.cor.8-6-3-1}
\expandafter\def\csname sga@target@IV.8.6.4\endcsname{ega4.c04.cor.8-6-4}
\expandafter\def\csname sga@target@IV.8.7.1\endcsname{ega4.c04.prop.8-7-1}
\expandafter\def\csname sga@target@IV.8.7.2\endcsname{ega4.c04.prop.8-7-2}
\expandafter\def\csname sga@target@IV.8.7.3\endcsname{ega4.c04.cor.8-7-3}
\expandafter\def\csname sga@target@IV.8.7.4\endcsname{ega4.c04.rem.8-7-4}
\expandafter\def\csname sga@target@IV.8.8.1\endcsname{ega4.c04.statement.8-8-1}
\expandafter\def\csname sga@target@IV.8.8.1.1\endcsname{ega4.c04.eq.8-8-1-1}
\expandafter\def\csname sga@target@IV.8.8.2\endcsname{ega4.c04.thm.8-8-2}
\expandafter\def\csname sga@target@IV.8.8.2.1\endcsname{ega4.c04.lem.8-8-2-1}
\expandafter\def\csname sga@target@IV.8.8.2.2\endcsname{ega4.c04.lem.8-8-2-2}
\expandafter\def\csname sga@target@IV.8.8.2.3\endcsname{ega4.c04.lem.8-8-2-3}
\expandafter\def\csname sga@target@IV.8.8.2.4\endcsname{ega4.c04.cor.8-8-2-4}
\expandafter\def\csname sga@target@IV.8.8.2.5\endcsname{ega4.c04.cor.8-8-2-5}
\expandafter\def\csname sga@target@IV.8.8.3\endcsname{ega4.c04.statement.8-8-3}
\expandafter\def\csname sga@target@IV.8.9.1\endcsname{ega4.c04.prop.8-9-1}
\expandafter\def\csname sga@target@IV.8.9.2\endcsname{ega4.c04.statement.8-9-2}
\expandafter\def\csname sga@target@IV.8.9.3\endcsname{ega4.c04.prop.8-9-3}
\expandafter\def\csname sga@target@IV.8.9.4\endcsname{ega4.c04.prop.8-9-4}
\expandafter\def\csname sga@target@IV.8.9.4.1\endcsname{ega4.c04.lem.8-9-4-1}
\expandafter\def\csname sga@target@IV.8.9.5\endcsname{ega4.c04.cor.8-9-5}
\expandafter\def\csname sga@target@IV.9.1.1\endcsname{ega4.c04.prop.9-1-1}
\expandafter\def\csname sga@target@IV.9.1.2\endcsname{ega4.c04.cor.9-1-2}
\expandafter\def\csname sga@target@IV.9.1.3\endcsname{ega4.c04.rem.9-1-3}
\expandafter\def\csname sga@target@IV.9.1.4\endcsname{ega4.c04.prop.9-1-4}
\expandafter\def\csname sga@target@IV.9.1.5\endcsname{ega4.c04.cor.9-1-5}
\expandafter\def\csname sga@target@IV.9.1.6\endcsname{ega4.c04.prop.9-1-6}
\expandafter\def\csname sga@target@IV.9.2\endcsname{ega4.c04.subsec.9-2.key-iv-9-2}
\expandafter\def\csname sga@target@IV.9.2.1\endcsname{ega4.c04.def.9-2-1}
\expandafter\def\csname sga@target@IV.9.2.2\endcsname{ega4.c04.item.9-2-2}
\expandafter\def\csname sga@target@IV.9.2.3\endcsname{ega4.c04.prop.9-2-3}
\expandafter\def\csname sga@target@IV.9.2.4\endcsname{ega4.c04.cor.9-2-4}
\expandafter\def\csname sga@target@IV.9.2.5\endcsname{ega4.c04.rem.9-2-5}
\expandafter\def\csname sga@target@IV.9.2.6\endcsname{ega4.c04.prop.9-2-6}
\expandafter\def\csname sga@target@IV.9.2.6.1\endcsname{ega4.c04.prop.9-2-6-1}
\expandafter\def\csname sga@target@IV.9.2.6.2\endcsname{ega4.c04.cor.9-2-6-2}
\expandafter\def\csname sga@target@IV.9.3.1\endcsname{ega4.c04.prop.9-3-1}
\expandafter\def\csname sga@target@IV.9.3.2\endcsname{ega4.c04.prop.9-3-2}
\expandafter\def\csname sga@target@IV.9.3.3\endcsname{ega4.c04.prop.9-3-3}
\expandafter\def\csname sga@target@IV.9.3.4\endcsname{ega4.c04.rem.9-3-4}
\expandafter\def\csname sga@target@IV.9.3.5\endcsname{ega4.c04.prop.9-3-5}
\expandafter\def\csname sga@target@IV.9.3.6\endcsname{ega4.c04.statement.9-3-6}
\expandafter\def\csname sga@target@IV.9.4.1\endcsname{ega4.c04.statement.9-4-1}
\expandafter\def\csname sga@target@IV.9.4.2\endcsname{ega4.c04.prop.9-4-2}
\expandafter\def\csname sga@target@IV.9.4.3\endcsname{ega4.c04.cor.9-4-3}
\expandafter\def\csname sga@target@IV.9.4.3.1\endcsname{ega4.c04.eq.9-4-3-1}
\expandafter\def\csname sga@target@IV.9.4.4\endcsname{ega4.c04.prop.9-4-4}
\expandafter\def\csname sga@target@IV.9.4.5\endcsname{ega4.c04.cor.9-4-5}
\expandafter\def\csname sga@target@IV.9.4.6\endcsname{ega4.c04.cor.9-4-6}
\expandafter\def\csname sga@target@IV.9.4.7\endcsname{ega4.c04.prop.9-4-7}
\expandafter\def\csname sga@target@IV.9.4.7.1\endcsname{ega4.c04.lem.9-4-7-1}
\expandafter\def\csname sga@target@IV.9.4.7.2\endcsname{ega4.c04.rem.9-4-7-2}
\expandafter\def\csname sga@target@IV.9.4.8\endcsname{ega4.c04.prop.9-4-8}
\expandafter\def\csname sga@target@IV.9.4.9\endcsname{ega4.c04.rem.9-4-9}
\expandafter\def\csname sga@target@IV.9.5.1\endcsname{ega4.c04.prop.9-5-1}
\expandafter\def\csname sga@target@IV.9.5.2\endcsname{ega4.c04.cor.9-5-2}
\expandafter\def\csname sga@target@IV.9.5.3\endcsname{ega4.c04.prop.9-5-3}
\expandafter\def\csname sga@target@IV.9.5.3.case2\endcsname{ega4.c04.item.9-5-3-case2}
\expandafter\def\csname sga@target@IV.9.5.4\endcsname{ega4.c04.cor.9-5-4}
\expandafter\def\csname sga@target@IV.9.5.5\endcsname{ega4.c04.prop.9-5-5}
\expandafter\def\csname sga@target@IV.9.5.6\endcsname{ega4.c04.cor.9-5-6}
\expandafter\def\csname sga@target@IV.9.5.7\endcsname{ega4.c04.rem.9-5-7}
\expandafter\def\csname sga@target@IV.9.6.1\endcsname{ega4.c04.prop.9-6-1}
\expandafter\def\csname sga@target@IV.9.6.2\endcsname{ega4.c04.prop.9-6-2}
\expandafter\def\csname sga@target@IV.9.6.3\endcsname{ega4.c04.prop.9-6-3}
\expandafter\def\csname sga@target@IV.9.6.3.1\endcsname{ega4.c04.eq.9-6-3-1}
\expandafter\def\csname sga@target@IV.9.6.4\endcsname{ega4.c04.cor.9-6-4}
\expandafter\def\csname sga@target@IV.9.6.5\endcsname{ega4.c04.cor.9-6-5}
\expandafter\def\csname sga@target@IV.9.6.6\endcsname{ega4.c04.sec.9-6-6}
\expandafter\def\csname sga@target@IV.9.6.7\endcsname{ega4.c04.sec.9-6-7}
\expandafter\def\csname sga@target@IV.9.7\endcsname{ega4.c04.subsec.9-7.key-iv-9-7}
\expandafter\def\csname sga@target@IV.9.7.1\endcsname{ega4.c04.lem.9-7-1}
\expandafter\def\csname sga@target@IV.9.7.10\endcsname{ega4.c04.cor.9-7-10}
\expandafter\def\csname sga@target@IV.9.7.11\endcsname{ega4.c04.rem.9-7-11}
\expandafter\def\csname sga@target@IV.9.7.12\endcsname{ega4.c04.prop.9-7-12}
\expandafter\def\csname sga@target@IV.9.7.2\endcsname{ega4.c04.lem.9-7-2}
\expandafter\def\csname sga@target@IV.9.7.3\endcsname{ega4.c04.subsec.9-7-3}
\expandafter\def\csname sga@target@IV.9.7.4\endcsname{ega4.c04.lem.9-7-4}
\expandafter\def\csname sga@target@IV.9.7.5\endcsname{ega4.c04.lem.9-7-5}
\expandafter\def\csname sga@target@IV.9.7.5.1\endcsname{ega4.c04.lem.9-7-5-1}
\expandafter\def\csname sga@target@IV.9.7.6\endcsname{ega4.c04.prop.9-7-6}
\expandafter\def\csname sga@target@IV.9.7.7\endcsname{ega4.c04.thm.9-7-7}
\expandafter\def\csname sga@target@IV.9.7.8\endcsname{ega4.c04.prop.9-7-8}
\expandafter\def\csname sga@target@IV.9.7.9\endcsname{ega4.c04.cor.9-7-9}
\expandafter\def\csname sga@target@IV.9.8\endcsname{ega4.c04.subsec.9-8.key-iv-9-8}
\expandafter\def\csname sga@target@IV.9.8.1\endcsname{ega4.c04.prop.9-8-1}
\expandafter\def\csname sga@target@IV.9.8.1-proof\endcsname{ega4.c04.subsec.9-8-1-proof}
\expandafter\def\csname sga@target@IV.9.8.2\endcsname{ega4.c04.subsec.9-8-2}
\expandafter\def\csname sga@target@IV.9.8.3\endcsname{ega4.c04.prop.9-8-3}
\expandafter\def\csname sga@target@IV.9.8.3-proof\endcsname{ega4.c04.subsec.9-8-3-proof}
\expandafter\def\csname sga@target@IV.9.8.4\endcsname{ega4.c04.cor.9-8-4}
\expandafter\def\csname sga@target@IV.9.8.4-proof\endcsname{ega4.c04.item.9-8-4-proof}
\expandafter\def\csname sga@target@IV.9.8.5\endcsname{ega4.c04.prop.9-8-5}
\expandafter\def\csname sga@target@IV.9.8.5-proof\endcsname{ega4.c04.item.9-8-5-proof}
\expandafter\def\csname sga@target@IV.9.8.6\endcsname{ega4.c04.prop.9-8-6}
\expandafter\def\csname sga@target@IV.9.8.6-proof\endcsname{ega4.c04.item.9-8-6-proof}
\expandafter\def\csname sga@target@IV.9.8.7\endcsname{ega4.c04.cor.9-8-7}
\expandafter\def\csname sga@target@IV.9.8.7-proof\endcsname{ega4.c04.item.9-8-7-proof}
\expandafter\def\csname sga@target@IV.9.8.8\endcsname{ega4.c04.prop.9-8-8}
\expandafter\def\csname sga@target@IV.9.8.8-proof\endcsname{ega4.c04.item.9-8-8-proof}
\expandafter\def\csname sga@target@IV.9.8.9\endcsname{ega4.c04.rem.9-8-9}
\expandafter\def\csname sga@target@IV.9.8.9.1\endcsname{ega4.c04.item.9-8-9-1}
\expandafter\def\csname sga@target@IV.9.9\endcsname{ega4.c04.subsec.9-9.key-iv-9-9}
\expandafter\def\csname sga@target@IV.9.9.1\endcsname{ega4.c04.prop.9-9-1}
\expandafter\def\csname sga@target@IV.9.9.1-proof\endcsname{ega4.c04.item.9-9-1-proof}
\expandafter\def\csname sga@target@IV.9.9.1.1\endcsname{ega4.c04.eq.9-9-1-1}
\expandafter\def\csname sga@target@IV.9.9.2\endcsname{ega4.c04.prop.9-9-2}
\expandafter\def\csname sga@target@IV.9.9.2-proof\endcsname{ega4.c04.item.9-9-2-proof}
\expandafter\def\csname sga@target@IV.9.9.2.1\endcsname{ega4.c04.eq.9-9-2-1}
\expandafter\def\csname sga@target@IV.9.9.2.2\endcsname{ega4.c04.eq.9-9-2-2}
\expandafter\def\csname sga@target@IV.9.9.3\endcsname{ega4.c04.cor.9-9-3}
\expandafter\def\csname sga@target@IV.9.9.3-proof\endcsname{ega4.c04.statement.9-9-3-proof}
\expandafter\def\csname sga@target@IV.9.9.4\endcsname{ega4.c04.prop.9-9-4}
\expandafter\def\csname sga@target@IV.9.9.4-proof\endcsname{ega4.c04.item.9-9-4-proof}
\expandafter\def\csname sga@target@IV.9.9.5\endcsname{ega4.c04.cor.9-9-5}
\expandafter\def\csname sga@target@IV.9.9.5-proof\endcsname{ega4.c04.item.9-9-5-proof}
\expandafter\def\csname sga@target@IV.9.9.6\endcsname{ega4.c04.prop.9-9-6}
\expandafter\def\csname sga@target@IV.9.9.6-proof\endcsname{ega4.c04.item.9-9-6-proof}
\expandafter\def\csname sga@target@section:IV.1\endcsname{ega4.c04.sec.1}
\expandafter\def\csname sga@target@section:IV.10\endcsname{ega4.c04.sec.10}
\expandafter\def\csname sga@target@section:IV.11\endcsname{ega4.c04.sec.11}
\expandafter\def\csname sga@target@section:IV.12\endcsname{ega4.c04.sec.12}
\expandafter\def\csname sga@target@section:IV.13\endcsname{ega4.c04.sec.13}
\expandafter\def\csname sga@target@section:IV.14\endcsname{ega4.c04.sec.14}
\expandafter\def\csname sga@target@section:IV.15\endcsname{ega4.c04.sec.15}
\expandafter\def\csname sga@target@section:IV.16\endcsname{ega4.c04.sec.16}
\expandafter\def\csname sga@target@section:IV.17\endcsname{ega4.c04.sec.17}
\expandafter\def\csname sga@target@section:IV.18\endcsname{ega4.c04.sec.18}
\expandafter\def\csname sga@target@section:IV.19\endcsname{ega4.c04.sec.19}
\expandafter\def\csname sga@target@section:IV.2\endcsname{ega4.c04.sec.2}
\expandafter\def\csname sga@target@section:IV.21\endcsname{ega4.c04.sec.21}
\expandafter\def\csname sga@target@section:IV.3\endcsname{ega4.c04.sec.3}
\expandafter\def\csname sga@target@section:IV.4\endcsname{ega4.c04.sec.4}
\expandafter\def\csname sga@target@section:IV.5\endcsname{ega4.c04.sec.5}
\expandafter\def\csname sga@target@section:IV.6\endcsname{ega4.c04.sec.6}
\expandafter\def\csname sga@target@section:IV.7\endcsname{ega4.c04.sec.7}
\expandafter\def\csname sga@target@section:IV.8\endcsname{ega4.c04.sec.8}
\expandafter\def\csname sga@target@section:IV.9\endcsname{ega4.c04.sec.9}
\expandafter\def\csname sga@target@subsection:IV.1.1\endcsname{ega4.c04.subsec.1-1}
\expandafter\def\csname sga@target@subsection:IV.1.10\endcsname{ega4.c04.subsec.1-10}
\expandafter\def\csname sga@target@subsection:IV.1.2\endcsname{ega4.c04.subsec.1-2}
\expandafter\def\csname sga@target@subsection:IV.1.3\endcsname{ega4.c04.subsec.1-3}
\expandafter\def\csname sga@target@subsection:IV.1.4\endcsname{ega4.c04.subsec.1-4}
\expandafter\def\csname sga@target@subsection:IV.1.5\endcsname{ega4.c04.subsec.1-5}
\expandafter\def\csname sga@target@subsection:IV.1.6\endcsname{ega4.c04.subsec.1-6}
\expandafter\def\csname sga@target@subsection:IV.1.7\endcsname{ega4.c04.subsec.1-7}
\expandafter\def\csname sga@target@subsection:IV.1.8\endcsname{ega4.c04.subsec.1-8}
\expandafter\def\csname sga@target@subsection:IV.1.9\endcsname{ega4.c04.subsec.1-9}
\expandafter\def\csname sga@target@subsection:IV.10.1\endcsname{ega4.c04.subsec.10-1}
\expandafter\def\csname sga@target@subsection:IV.10.10\endcsname{ega4.c04.subsec.10-10}
\expandafter\def\csname sga@target@subsection:IV.10.2\endcsname{ega4.c04.subsec.10-2}
\expandafter\def\csname sga@target@subsection:IV.10.3\endcsname{ega4.c04.subsec.10-3}
\expandafter\def\csname sga@target@subsection:IV.10.4\endcsname{ega4.c04.subsec.10-4}
\expandafter\def\csname sga@target@subsection:IV.10.5\endcsname{ega4.c04.subsec.10-5}
\expandafter\def\csname sga@target@subsection:IV.10.6\endcsname{ega4.c04.subsec.10-6}
\expandafter\def\csname sga@target@subsection:IV.10.7\endcsname{ega4.c04.subsec.10-7}
\expandafter\def\csname sga@target@subsection:IV.10.8\endcsname{ega4.c04.subsec.10-8}
\expandafter\def\csname sga@target@subsection:IV.10.9\endcsname{ega4.c04.subsec.10-9}
\expandafter\def\csname sga@target@subsection:IV.11.1\endcsname{ega4.c04.subsec.11-1}
\expandafter\def\csname sga@target@subsection:IV.11.10\endcsname{ega4.c04.subsec.11-10}
\expandafter\def\csname sga@target@subsection:IV.11.2\endcsname{ega4.c04.subsec.11-2}
\expandafter\def\csname sga@target@subsection:IV.11.3\endcsname{ega4.c04.subsec.11-3}
\expandafter\def\csname sga@target@subsection:IV.11.4\endcsname{ega4.c04.subsec.11-4}
\expandafter\def\csname sga@target@subsection:IV.11.5\endcsname{ega4.c04.subsec.11-5}
\expandafter\def\csname sga@target@subsection:IV.11.6\endcsname{ega4.c04.subsec.11-6}
\expandafter\def\csname sga@target@subsection:IV.11.7\endcsname{ega4.c04.subsec.11-7}
\expandafter\def\csname sga@target@subsection:IV.11.8\endcsname{ega4.c04.subsec.11-8}
\expandafter\def\csname sga@target@subsection:IV.11.9\endcsname{ega4.c04.subsec.11-9}
\expandafter\def\csname sga@target@subsection:IV.12.0\endcsname{ega4.c04.subsec.12-0}
\expandafter\def\csname sga@target@subsection:IV.12.1\endcsname{ega4.c04.subsec.12-1}
\expandafter\def\csname sga@target@subsection:IV.12.2\endcsname{ega4.c04.subsec.12-2}
\expandafter\def\csname sga@target@subsection:IV.12.3\endcsname{ega4.c04.subsec.12-3}
\expandafter\def\csname sga@target@subsection:IV.13.1\endcsname{ega4.c04.subsec.13-1}
\expandafter\def\csname sga@target@subsection:IV.13.2\endcsname{ega4.c04.subsec.13-2}
\expandafter\def\csname sga@target@subsection:IV.13.3\endcsname{ega4.c04.subsec.13-3}
\expandafter\def\csname sga@target@subsection:IV.14.1\endcsname{ega4.c04.subsec.14-1}
\expandafter\def\csname sga@target@subsection:IV.14.2\endcsname{ega4.c04.subsec.14-2}
\expandafter\def\csname sga@target@subsection:IV.14.3\endcsname{ega4.c04.subsec.14-3}
\expandafter\def\csname sga@target@subsection:IV.14.4\endcsname{ega4.c04.subsec.14-4}
\expandafter\def\csname sga@target@subsection:IV.14.5\endcsname{ega4.c04.subsec.14-5}
\expandafter\def\csname sga@target@subsection:IV.15.1\endcsname{ega4.c04.subsec.15-1}
\expandafter\def\csname sga@target@subsection:IV.15.2\endcsname{ega4.c04.subsec.15-2}
\expandafter\def\csname sga@target@subsection:IV.15.3\endcsname{ega4.c04.subsec.15-3}
\expandafter\def\csname sga@target@subsection:IV.15.4\endcsname{ega4.c04.subsec.15-4}
\expandafter\def\csname sga@target@subsection:IV.15.5\endcsname{ega4.c04.subsec.15-5}
\expandafter\def\csname sga@target@subsection:IV.15.6\endcsname{ega4.c04.subsec.15-6}
\expandafter\def\csname sga@target@subsection:IV.15.7\endcsname{ega4.c04.subsec.15-7}
\expandafter\def\csname sga@target@subsection:IV.17.1\endcsname{ega4.c04.subsec.17-1}
\expandafter\def\csname sga@target@subsection:IV.17.10\endcsname{ega4.c04.subsec.17-10}
\expandafter\def\csname sga@target@subsection:IV.17.11\endcsname{ega4.c04.subsec.17-11}
\expandafter\def\csname sga@target@subsection:IV.17.12\endcsname{ega4.c04.subsec.17-12}
\expandafter\def\csname sga@target@subsection:IV.17.13\endcsname{ega4.c04.subsec.17-13}
\expandafter\def\csname sga@target@subsection:IV.17.14\endcsname{ega4.c04.subsec.17-14}
\expandafter\def\csname sga@target@subsection:IV.17.15\endcsname{ega4.c04.subsec.17-15}
\expandafter\def\csname sga@target@subsection:IV.17.16\endcsname{ega4.c04.subsec.17-16}
\expandafter\def\csname sga@target@subsection:IV.17.2\endcsname{ega4.c04.subsec.17-2}
\expandafter\def\csname sga@target@subsection:IV.17.3\endcsname{ega4.c04.subsec.17-3}
\expandafter\def\csname sga@target@subsection:IV.17.4\endcsname{ega4.c04.subsec.17-4}
\expandafter\def\csname sga@target@subsection:IV.17.5\endcsname{ega4.c04.subsec.17-5}
\expandafter\def\csname sga@target@subsection:IV.17.6\endcsname{ega4.c04.subsec.17-6}
\expandafter\def\csname sga@target@subsection:IV.17.7\endcsname{ega4.c04.subsec.17-7}
\expandafter\def\csname sga@target@subsection:IV.17.8\endcsname{ega4.c04.subsec.17-8}
\expandafter\def\csname sga@target@subsection:IV.17.9\endcsname{ega4.c04.subsec.17-9}
\expandafter\def\csname sga@target@subsection:IV.18.1\endcsname{ega4.c04.subsec.18-1}
\expandafter\def\csname sga@target@subsection:IV.18.10\endcsname{ega4.c04.subsec.18-10}
\expandafter\def\csname sga@target@subsection:IV.18.11\endcsname{ega4.c04.subsec.18-11}
\expandafter\def\csname sga@target@subsection:IV.18.12\endcsname{ega4.c04.subsec.18-12}
\expandafter\def\csname sga@target@subsection:IV.18.2\endcsname{ega4.c04.subsec.18-2}
\expandafter\def\csname sga@target@subsection:IV.18.3\endcsname{ega4.c04.subsec.18-3}
\expandafter\def\csname sga@target@subsection:IV.18.4\endcsname{ega4.c04.subsec.18-4}
\expandafter\def\csname sga@target@subsection:IV.18.5\endcsname{ega4.c04.subsec.18-5}
\expandafter\def\csname sga@target@subsection:IV.18.6\endcsname{ega4.c04.subsec.18-6}
\expandafter\def\csname sga@target@subsection:IV.18.7\endcsname{ega4.c04.subsec.18-7}
\expandafter\def\csname sga@target@subsection:IV.18.8\endcsname{ega4.c04.subsec.18-8}
\expandafter\def\csname sga@target@subsection:IV.18.9\endcsname{ega4.c04.subsec.18-9}
\expandafter\def\csname sga@target@subsection:IV.19.1\endcsname{ega4.c04.subsec.19-1}
\expandafter\def\csname sga@target@subsection:IV.19.2\endcsname{ega4.c04.subsec.19-2}
\expandafter\def\csname sga@target@subsection:IV.19.3\endcsname{ega4.c04.subsec.19-3}
\expandafter\def\csname sga@target@subsection:IV.19.4\endcsname{ega4.c04.subsec.19-4}
\expandafter\def\csname sga@target@subsection:IV.19.5\endcsname{ega4.c04.subsec.19-5}
\expandafter\def\csname sga@target@subsection:IV.19.6\endcsname{ega4.c04.subsec.19-6}
\expandafter\def\csname sga@target@subsection:IV.19.7\endcsname{ega4.c04.subsec.19-7}
\expandafter\def\csname sga@target@subsection:IV.19.8\endcsname{ega4.c04.subsec.19-8}
\expandafter\def\csname sga@target@subsection:IV.19.9\endcsname{ega4.c04.subsec.19-9}
\expandafter\def\csname sga@target@subsection:IV.2.1\endcsname{ega4.c04.subsec.2-1}
\expandafter\def\csname sga@target@subsection:IV.2.2\endcsname{ega4.c04.subsec.2-2}
\expandafter\def\csname sga@target@subsection:IV.2.3\endcsname{ega4.c04.subsec.2-3}
\expandafter\def\csname sga@target@subsection:IV.2.4\endcsname{ega4.c04.subsec.2-4}
\expandafter\def\csname sga@target@subsection:IV.2.5\endcsname{ega4.c04.subsec.2-5}
\expandafter\def\csname sga@target@subsection:IV.2.6\endcsname{ega4.c04.subsec.2-6}
\expandafter\def\csname sga@target@subsection:IV.2.7\endcsname{ega4.c04.subsec.2-7}
\expandafter\def\csname sga@target@subsection:IV.2.8\endcsname{ega4.c04.subsec.2-8}
\expandafter\def\csname sga@target@subsection:IV.21.1\endcsname{ega4.c04.subsec.21-1}
\expandafter\def\csname sga@target@subsection:IV.21.10\endcsname{ega4.c04.subsec.21-10}
\expandafter\def\csname sga@target@subsection:IV.21.11\endcsname{ega4.c04.subsec.21-11}
\expandafter\def\csname sga@target@subsection:IV.21.12\endcsname{ega4.c04.subsec.21-12}
\expandafter\def\csname sga@target@subsection:IV.21.13\endcsname{ega4.c04.subsec.21-13}
\expandafter\def\csname sga@target@subsection:IV.21.14\endcsname{ega4.c04.subsec.21-14}
\expandafter\def\csname sga@target@subsection:IV.21.15\endcsname{ega4.c04.subsec.21-15}
\expandafter\def\csname sga@target@subsection:IV.21.2\endcsname{ega4.c04.subsec.21-2}
\expandafter\def\csname sga@target@subsection:IV.21.3\endcsname{ega4.c04.subsec.21-3}
\expandafter\def\csname sga@target@subsection:IV.21.4\endcsname{ega4.c04.subsec.21-4}
\expandafter\def\csname sga@target@subsection:IV.21.5\endcsname{ega4.c04.subsec.21-5}
\expandafter\def\csname sga@target@subsection:IV.21.6\endcsname{ega4.c04.subsec.21-6}
\expandafter\def\csname sga@target@subsection:IV.21.7\endcsname{ega4.c04.subsec.21-7}
\expandafter\def\csname sga@target@subsection:IV.21.8\endcsname{ega4.c04.subsec.21-8}
\expandafter\def\csname sga@target@subsection:IV.21.9\endcsname{ega4.c04.subsec.21-9}
\expandafter\def\csname sga@target@subsection:IV.3.1\endcsname{ega4.c04.subsec.3-1}
\expandafter\def\csname sga@target@subsection:IV.3.2\endcsname{ega4.c04.subsec.3-2}
\expandafter\def\csname sga@target@subsection:IV.3.3\endcsname{ega4.c04.subsec.3-3}
\expandafter\def\csname sga@target@subsection:IV.3.4\endcsname{ega4.c04.subsec.3-4}
\expandafter\def\csname sga@target@subsection:IV.4.1\endcsname{ega4.c04.subsec.4-1}
\expandafter\def\csname sga@target@subsection:IV.4.2\endcsname{ega4.c04.subsec.4-2}
\expandafter\def\csname sga@target@subsection:IV.4.3\endcsname{ega4.c04.subsec.4-3}
\expandafter\def\csname sga@target@subsection:IV.4.4\endcsname{ega4.c04.subsec.4-4}
\expandafter\def\csname sga@target@subsection:IV.4.5\endcsname{ega4.c04.subsec.4-5}
\expandafter\def\csname sga@target@subsection:IV.4.6\endcsname{ega4.c04.subsec.4-6}
\expandafter\def\csname sga@target@subsection:IV.4.7\endcsname{ega4.c04.subsec.4-7}
\expandafter\def\csname sga@target@subsection:IV.4.8\endcsname{ega4.c04.subsec.4-8}
\expandafter\def\csname sga@target@subsection:IV.4.9\endcsname{ega4.c04.subsec.4-9}
\expandafter\def\csname sga@target@subsection:IV.5.1\endcsname{ega4.c04.subsec.5-1}
\expandafter\def\csname sga@target@subsection:IV.5.10\endcsname{ega4.c04.subsec.5-10}
\expandafter\def\csname sga@target@subsection:IV.5.11\endcsname{ega4.c04.subsec.5-11}
\expandafter\def\csname sga@target@subsection:IV.5.12\endcsname{ega4.c04.subsec.5-12}
\expandafter\def\csname sga@target@subsection:IV.5.13\endcsname{ega4.c04.subsec.5-13}
\expandafter\def\csname sga@target@subsection:IV.5.2\endcsname{ega4.c04.subsec.5-2}
\expandafter\def\csname sga@target@subsection:IV.5.3\endcsname{ega4.c04.subsec.5-3}
\expandafter\def\csname sga@target@subsection:IV.5.4\endcsname{ega4.c04.subsec.5-4}
\expandafter\def\csname sga@target@subsection:IV.5.5\endcsname{ega4.c04.subsec.5-5}
\expandafter\def\csname sga@target@subsection:IV.5.6\endcsname{ega4.c04.subsec.5-6}
\expandafter\def\csname sga@target@subsection:IV.5.7\endcsname{ega4.c04.subsec.5-7}
\expandafter\def\csname sga@target@subsection:IV.5.8\endcsname{ega4.c04.subsec.5-8}
\expandafter\def\csname sga@target@subsection:IV.5.9\endcsname{ega4.c04.subsec.5-9}
\expandafter\def\csname sga@target@subsection:IV.6.1\endcsname{ega4.c04.subsec.6-1}
\expandafter\def\csname sga@target@subsection:IV.6.10\endcsname{ega4.c04.subsec.6-10}
\expandafter\def\csname sga@target@subsection:IV.6.11\endcsname{ega4.c04.subsec.6-11}
\expandafter\def\csname sga@target@subsection:IV.6.12\endcsname{ega4.c04.subsec.6-12}
\expandafter\def\csname sga@target@subsection:IV.6.13\endcsname{ega4.c04.subsec.6-13}
\expandafter\def\csname sga@target@subsection:IV.6.14\endcsname{ega4.c04.subsec.6-14}
\expandafter\def\csname sga@target@subsection:IV.6.15\endcsname{ega4.c04.subsec.6-15}
\expandafter\def\csname sga@target@subsection:IV.6.2\endcsname{ega4.c04.subsec.6-2}
\expandafter\def\csname sga@target@subsection:IV.6.3\endcsname{ega4.c04.subsec.6-3}
\expandafter\def\csname sga@target@subsection:IV.6.4\endcsname{ega4.c04.subsec.6-4}
\expandafter\def\csname sga@target@subsection:IV.6.5\endcsname{ega4.c04.subsec.6-5}
\expandafter\def\csname sga@target@subsection:IV.6.6\endcsname{ega4.c04.subsec.6-6}
\expandafter\def\csname sga@target@subsection:IV.6.7\endcsname{ega4.c04.subsec.6-7}
\expandafter\def\csname sga@target@subsection:IV.6.8\endcsname{ega4.c04.subsec.6-8}
\expandafter\def\csname sga@target@subsection:IV.6.9\endcsname{ega4.c04.subsec.6-9}
\expandafter\def\csname sga@target@subsection:IV.7.1\endcsname{ega4.c04.subsec.7-1}
\expandafter\def\csname sga@target@subsection:IV.7.2\endcsname{ega4.c04.subsec.7-2}
\expandafter\def\csname sga@target@subsection:IV.7.3\endcsname{ega4.c04.subsec.7-3}
\expandafter\def\csname sga@target@subsection:IV.7.4\endcsname{ega4.c04.subsec.7-4}
\expandafter\def\csname sga@target@subsection:IV.7.5\endcsname{ega4.c04.subsec.7-5}
\expandafter\def\csname sga@target@subsection:IV.7.6\endcsname{ega4.c04.subsec.7-6}
\expandafter\def\csname sga@target@subsection:IV.7.7\endcsname{ega4.c04.subsec.7-7}
\expandafter\def\csname sga@target@subsection:IV.7.8\endcsname{ega4.c04.subsec.7-8}
\expandafter\def\csname sga@target@subsection:IV.7.9\endcsname{ega4.c04.subsec.7-9}
\expandafter\def\csname sga@target@subsection:IV.8.1\endcsname{ega4.c04.subsec.8-1}
\expandafter\def\csname sga@target@subsection:IV.8.10\endcsname{ega4.c04.subsec.8-10}
\expandafter\def\csname sga@target@subsection:IV.8.11\endcsname{ega4.c04.subsec.8-11}
\expandafter\def\csname sga@target@subsection:IV.8.12\endcsname{ega4.c04.subsec.8-12}
\expandafter\def\csname sga@target@subsection:IV.8.13\endcsname{ega4.c04.subsec.8-13}
\expandafter\def\csname sga@target@subsection:IV.8.14\endcsname{ega4.c04.subsec.8-14}
\expandafter\def\csname sga@target@subsection:IV.8.2\endcsname{ega4.c04.subsec.8-2}
\expandafter\def\csname sga@target@subsection:IV.8.3\endcsname{ega4.c04.subsec.8-3}
\expandafter\def\csname sga@target@subsection:IV.8.4\endcsname{ega4.c04.subsec.8-4}
\expandafter\def\csname sga@target@subsection:IV.8.5\endcsname{ega4.c04.subsec.8-5}
\expandafter\def\csname sga@target@subsection:IV.8.6\endcsname{ega4.c04.subsec.8-6}
\expandafter\def\csname sga@target@subsection:IV.8.7\endcsname{ega4.c04.subsec.8-7}
\expandafter\def\csname sga@target@subsection:IV.8.8\endcsname{ega4.c04.subsec.8-8}
\expandafter\def\csname sga@target@subsection:IV.8.9\endcsname{ega4.c04.subsec.8-9}
\expandafter\def\csname sga@target@subsection:IV.9.1\endcsname{ega4.c04.subsec.9-1}
\expandafter\def\csname sga@target@subsection:IV.9.2\endcsname{ega4.c04.subsec.9-2}
\expandafter\def\csname sga@target@subsection:IV.9.3\endcsname{ega4.c04.subsec.9-3}
\expandafter\def\csname sga@target@subsection:IV.9.4\endcsname{ega4.c04.subsec.9-4}
\expandafter\def\csname sga@target@subsection:IV.9.5\endcsname{ega4.c04.subsec.9-5}
\expandafter\def\csname sga@target@subsection:IV.9.6\endcsname{ega4.c04.subsec.9-6}
\expandafter\def\csname sga@target@subsection:IV.9.7\endcsname{ega4.c04.subsec.9-7}
\expandafter\def\csname sga@target@subsection:IV.9.8\endcsname{ega4.c04.subsec.9-8}
\expandafter\def\csname sga@target@subsection:IV.9.9\endcsname{ega4.c04.subsec.9-9}
\let\sgaOriginalRef\ref
\let\sgaOriginalPageref\pageref
\renewcommand{\ref}[1]{%
  \ifcsname sga@target@#1\endcsname
    \hyperlink{\csname sga@target@#1\endcsname}{\sgaOriginalRef*{#1}}%
  \else
    \sgaOriginalRef{#1}%
  \fi
}
\renewcommand{\pageref}[1]{%
  \ifcsname sga@target@#1\endcsname
    \hyperlink{\csname sga@target@#1\endcsname}{\sgaOriginalPageref*{#1}}%
  \else
    \sgaOriginalPageref{#1}%
  \fi
}
\renewcommand{\eqref}[1]{%
  \ifcsname sga@target@#1\endcsname
    \hyperlink{\csname sga@target@#1\endcsname}{\textup{\tagform@{\sgaOriginalRef*{#1}}}}%
  \else
    \textup{\tagform@{\sgaOriginalRef{#1}}}%
  \fi
}
% Stable semantic routing for the reader's prose reference helpers.
\newcommand{\sgaStableOrOriginalLink}[2]{%
  \ifcsname sga@target@#1\endcsname
    \hyperlink{\csname sga@target@#1\endcsname}{#2}%
  \else
    \@ifundefined{r@#1}{#2}{\hyperref[#1]{#2}}%
  \fi
}
\RenewDocumentCommand{\sref}{o m o}{%
  \IfNoValueTF{#1}%
    {\IfNoValueTF{#3}%
      {\sgaStableOrOriginalLink{#2}{\normalfont{(\StrBehind{#2}{.})}}}%
      {\sgaStableOrOriginalLink{#2}{\normalfont{(\StrBehind{#2}{.},~{#3})}}}}%
    {\IfNoValueTF{#3}%
      {\sgaStableOrOriginalLink{#2}{\normalfont{(\textbf{#1},~\StrBehind{#2}{.})}}}%
      {\sgaStableOrOriginalLink{#2}{\normalfont{(\textbf{#1},~\StrBehind{#2}{.},~{#3})}}}}%
}
\renewcommand{\eref}[1]{%
  \sgaStableOrOriginalLink{#1}{\normalfont{(\StrBehind{#1}{.})}}}
\makeatother
% END NATIVE INPUT EGA4_reference_v2_links
% Six Chapter 0 targets cited in inherited IV-4 source are outside this
% volume-only reader.  Supply printable, non-linking label stubs locally so
% the references remain visible without claiming that their targets are here.
\makeatletter
\newcommand{\EGAIVFourExternalLabel}[1]{%
  \expandafter\gdef\csname r@#1\endcsname{{#1}{0}{}{#1}{}}}
\EGAIVFourExternalLabel{0.4.3.3}
\EGAIVFourExternalLabel{0.5.4.3}
\EGAIVFourExternalLabel{0.5.4.5}
\EGAIVFourExternalLabel{0.5.4.7}
\EGAIVFourExternalLabel{0.6.1.2}
\EGAIVFourExternalLabel{0.16.4.6}
\makeatother
\setlength{\emergencystretch}{3em}

\hypersetup{
  pdftitle={Elements of Algebraic Geometry IV, Part Four -- English Edition},
  pdfauthor={Alexandre Grothendieck and Jean Dieudonne; English translation},
  pdfsubject={Standalone English reader of EGA IV, Part Four}
}

\begin{document}

\title{Elements of Algebraic Geometry IV\\Part Four}
\author{Alexandre Grothendieck and Jean Dieudonn\'e}
\date{}
\maketitle

% BEGIN SCHOLARLY TRANSLATION ATTRIBUTION

\begingroup
\small
\begin{quote}
\textbf{Translation sources and contributors.}
The original work is by Alexandre Grothendieck and Jean Dieudonn\'e.
The EGA English edition series draws on the human community translation
\texttt{github.com/ryankeleti/ega}, source version \texttt{3d7afe28}.
Its named contributors are Mateo Carmona, Tim Holzschuh, Tim Hosgood,
Ryan Keleti, Thiago Solovera \& Nery, and Wayne Yang.
The inherited English text includes section 16, the opening and section 17.1, and the opening and section 20.1, with subsequent editorial changes.
The additional translation is AI-translated; subsequent corrections, editing,
typesetting and assembly are AI-produced. The publisher did not translate
the text. Recorded AI work includes OpenAI Codex,
GPT-6 Astra, Ultra.
The accompanying \texttt{ATTRIBUTION.md} gives source links, contribution
scope and rights information. Attribution does not assert endorsement
or a licence grant from the upstream contributors.

\medskip
\textbf{Rights and availability.}
Rights in the original works and in any protected contributions to these
derivative translations and transcriptions remain with their respective
rights holders. This project does not dedicate those works or contributions
to the public domain or grant a licence over them. Existing third-party
notices and any permissions already granted remain applicable.

These editions are available free of charge for non-commercial, purely
educational purposes, on the presumption that the respective rights holders
would not object to these translations and transcriptions being freely
available alongside the French originals on NUMDAM
(\texttt{www.numdam.org/journals/PMIHES/}). This presumption is not a statement
that permission has been obtained; free access to the originals does not
itself grant permission to translate, adapt or republish them.

Rights holders who object to the availability of their material may open
an issue at \texttt{github.com/KokunoYumeto/ega-en/issues}. The publisher
will take down the affected material at the rights holder's request.
\end{quote}
\endgroup
% END SCHOLARLY TRANSLATION ATTRIBUTION


\begin{center}
\large Standalone English reader
\end{center}

\tableofcontents

\setcounter{section}{15}
% BEGIN NATIVE INPUT source_aligned/ega4-16
\section{Differential invariants. Differentially smooth morphisms}
\label{section:IV.16}

\oldpage[IV-4]{5}
In this section we present, in global form, some notions of differential calculus that are particularly useful in algebraic geometry.
We pass over many developments classical in differential geometry (connections, infinitesimal transformations associated with a vector field, jets, etc.), although these notions admit particularly natural formulations in the setting of schemes.
We likewise pass over here the phenomena special to characteristic $p>0$ (some of which are studied, in the affine setting, in \sref[0]{0.21}).
For further material on the differential formalism for preschemes, the reader may consult Expos\'es~II and VII of \cite{IV-42}, as well as later chapters of this treatise.

\subsection{Normal invariants of an immersion}
\label{IV.16.1}

\begin{env}[16.1.1]
\label{IV.16.1.1}
Let $(X, \sh{O}_X)$ and $(Y, \sh{O}_Y)$ be two ringed spaces, and let $f = (\psi, \theta): Y \to X$ be a morphism of ringed spaces \sref[0]{0.4.1.1} such that the homomorphism
\[
  \theta^\#: \psi^*(\sh{O}_X) \to \sh{O}_Y
\]
is surjective, so that $\sh{O}_Y$ is identified with a sheaf of quotient rings $\psi^*(\sh{O}_X)/\sh{I}_f$. 
We can then endow $\psi^*(\sh{O}_X)$ with the $\sh{I}_f$-preadic filtration.
\end{env}

\begin{definition}[16.1.2]
\label{IV.16.1.2}
The $\sh{O}_Y$-augmented sheaf of rings $\psi^*(\sh{O}_X)/\sh{I}_f^{n+1}$ is called the $n$-th \emph{normal invariant} of $f$;
the ringed space $(Y, \psi^*(\sh{O}_X)/\sh{I}_f^{n+1})$ is called the $n$-th \emph{infinitesimal neighbourhood} of $Y$ for the morphism $f$ and is denoted by $Y^{(n)}_f$ or simply $Y^{(n)}$.
The sheaf of graded rings associated to the sheaf of filtered rings $\psi^*(\sh{O}_X)$
\[
  \label{IV.16.1.2.1}
  \shGr_\bullet(f) = \bigoplus_{n \geq 0}(\sh{I}_f^{n}/\sh{I}_f^{n+1} )
  \tag{16.1.2.1}
\]
is called the sheaf of graded rings \emph{associated with} $f$. The sheaf $\shGr_1(f) = \sh{I}_f/\sh{I}_f^{2}$ is called the \emph{conormal sheaf} of $f$ (and is also denoted by $\sh{N}_{Y/X}$ when no confusion can result).
\end{definition}

It is clear that the $\sh{O}_{Y^{(n)}} = \psi^*(\sh{O}_X)/\sh{I}_f^{n+1}$ (also denoted $\sh{O}_{Y_f^{(n)}}$) form a
\oldpage[IV-4]{6}
projective system of sheaves of rings on $Y$; for $n \leq m$, the transition homomorphism $\vphi_{nm}:\sh{O}_{Y^{(m)}} \to \sh{O}_{Y^{(n)}}$ identifies $\sh{O}_{Y^{(n)}}$ with the quotient of $\sh{O}_{Y^{(m)}}$ by the power $(\sh{I}_f/\sh{I}_f^{m+1})^{n+1}$ of the \emph{augmentation ideal} of $\sh{O}_{Y^{(m)}}$, the kernel of $\vphi_{0m}: \sh{O}_{Y^{(m)}} \to \sh{O}_{Y}$.
The $Y^{(n)}$ therefore form an inductive system of ringed spaces, all having underlying space $Y$, and we have canonical morphisms of ringed spaces $h_n: Y^{(n)} \to X$ equal to $(\psi, \theta_n)$, where $\theta^\#_n$ is the canonical morphism $\psi^*(\sh{O}_X) \to \psi^*(\sh{O}_X)/\sh{I}_f^{n+1}$.
It is clear that the sheaf $\shGr_\bullet(f)$ is a sheaf of graded algebras over the sheaf of rings $\sh{O}_Y = \shGr_0(f)$, and that the $\shGr_k(f)$ are $\sh{O}_Y$-modules.

As with every sheaf of filtered rings, we have a \emph{canonical surjective homomorphism} of graded $\sh{O}_Y$-algebras
\[
  \label{IV.16.1.2.2}
  \bb{S}_{\sh{O}_Y}^\bullet(\shGr_1(f)) \to \shGr_\bullet(f)
  \tag{16.1.2.2}
\]
which coincides in degrees $0$ and $1$ with the identities.

\begin{examples}[16.1.3]
\label{IV.16.1.3}
\medskip\noindent
\begin{enumerate}
  \item[(i)] Suppose that $X$ is a locally ringed space, $Y$ is reduced to a single point $y$ (endowed with a ring $\sh{O}_y$) and that, if $x = \psi(y)$, $\theta^\#:\sh{O}_x \to \sh{O}_y$ is a \emph{surjective} homomorphism of rings having as kernel the maximal ideal $\mathfrak{m}_x$ of $\sh{O}_x$.
  So the $\sh{O}_{Y^{(n)}}$ are identified with the rings $\sh{O}_x/\mathfrak{m}_x^{n+1}$ and $\shGr_\bullet(f)$ with the graded ring associated with the local ring $\sh{O}_x$ endowed with the $\mathfrak{m}_x$-preadic filtration.
  \item[(ii)] Suppose that $Y$ is a closed subset of an open subspace $U$ of $X$ and that $\sh{O}_Y$ is induced on $Y$ by a quotient sheaf $\sh{O}_U/\sh{I}$, where $\sh{I}$ is an ideal of $\sh{O}_U$ such that $\sh{I}_x = \sh{O}_x$ for every $x \not\in Y$;
  if $X$ is a locally ringed space, suppose further that $\sh{I}_x \neq \sh{O}_x$ for $x \in Y$, so that $(Y, \sh{O}_Y)$ is again a locally ringed space.
  
  Let $\psi_0: Y \to U$ be the canonical injection and denote by $\theta_0: \sh{O}_U \to (\psi_0)_*(\sh{O}_Y)$ the homomorphism such that $\theta_0^\#$ is the canonical homomorphism $\psi^*_0(\sh{O}_U) = \sh{O}_U|Y \to (\sh{O}_U/\sh{I})|Y$, so that $j_0=(\psi_0, \theta_0):Y \to U$ is a morphism of ringed spaces (and of locally ringed spaces if $X$ is a locally ringed space);
  if $i:U \to X$ is the canonical injection (morphism of ringed spaces), $j = i\circ j_0$ is the morphism $(\psi, \theta)$ of $Y$ to $X$ where $\psi: Y \to X$ is the canonical injection and $\theta:\sh{O}_X \to \psi_*(\sh{O}_Y)$ is the homomorphism such that $\theta^\# = \theta_0^\#$.
  Since $\theta^\#$ is surjective we can apply the previous definitions;
  $\sh{O}_{Y^{(n)}}$ is equal to $\psi^*_0(\sh{O}_U/\sh{I}^{n+1})$, and we have $(\psi_0)_*(\sh{O}_{Y^{(n)}} ) = \sh{O}_U/\sh{I}^{n+1}$, and $\shGr_n(j) = \shGr_n(j_0) = \psi^*_0(\sh{I}^n/\sh{I}^{n+1}) = j^*_0(\sh{I}^n/\sh{I}^{n+1})$.
\end{enumerate}
\end{examples}

\begin{env}[16.1.4]
\label{IV.16.1.4}
The example \sref{IV.16.1.3}[(ii)] shows that in general the $\sh{O}_{Y^{(n)}}$ are \emph{not canonically endowed with a structure of an $\sh{O}_Y$-module}, or \emph{a fortiori} with a structure of an $\sh{O}_Y$-algebra.
The data of such a structure is equivalent to the data of a homomorphism of sheaves of rings $\lambda_n:\sh{O}_Y \to \sh{O}_{Y^{(n)}}$, right inverse to the augmentation morphism $\vphi_{0n}$;
it is also equivalent to the data of a morphism of ringed spaces $(1_Y, \lambda_n): Y^{(n)} \to Y$ left inverse to the canonical morphism $(1_Y, \vphi_{0n}): Y \to Y^{(n)}$.
\end{env}

\begin{proposition}[16.1.5]
\label{IV.16.1.5}
Let $f = (\psi, \theta): Y \to X$ be an immersion of preschemes.
Then:
\begin{enumerate}
  \item[{\rm(i)}] $\shGr_\bullet(f)$ is a quasi-coherent graded $\sh{O}_Y$-algebra.
\oldpage[IV-4]{7}
  \item[{\rm(ii)}] The $Y^{(n)}$ are preschemes, canonically isomorphic to subpreschemes of $X$.
  \item[{\rm(iii)}] Every homomorphism of sheaves of rings $\lambda_n: \sh{O}_Y \to \sh{O}_{Y^{(n)}}$, right inverse to the augmentation homomorphism $\vphi_{0n}$, makes $\sh{O}_{Y^{(n)}}$ and the $\sh{O}_{Y^{(k)}}$ for $k\leq n$ quasi-coherent $\sh{O}_Y$-algebras;
  the $\sh{O}_Y$-module structures induced from the above structures on the $\shGr_k(f)$ for $k \leq n$ coincide with the ones defined in \sref{IV.16.1.2}.
\end{enumerate}
\end{proposition}

\begin{proof}
(i) Since the question is local on $X$ and $Y$, we can reduce to the case where $Y$ is a closed subprescheme of $X$ defined by a quasi-coherent ideal $\sh{I}$ of $\sh{O}_X$;
since $\sh{O}_Y$ is the restriction to $Y$ of $\sh{O}_X/\sh{I}$, assertion (i) is evident, and $Y^{(n)}$ is the closed subprescheme of $X$ defined by the quasi-coherent ideal $\sh{I}^{n+1}$ of $\sh{O}_X$.
Finally, to prove (iii), note that the data of $\lambda_n$ makes the ideal $\sh{I}/\sh{I}^{n+1}$ of the augmentation $\vphi_{0n}$ and its quotients $\sh{I}/\sh{I}^{k+1}$ ($1\leq k \leq n$) $\sh{O}_Y$-modules; it suffices to prove by induction on $k$ that the $\sh{I}/\sh{I}^{k+1}$ are quasi-coherent $\sh{O}_Y$-modules and that the quotient $\sh{O}_Y$-module structure thereby induced on $\sh{I}^k/\sh{I}^{k+1}$ is the same as that defined in \sref{IV.16.1.2}.
The second assertion is immediate, $\sh{I}^k/\sh{I}^{k+1}$ being killed by $\sh{I}/\sh{I}^{n+1}$;
the first follows by induction on $k$, since it is trivial for $k=1$ and $\sh{I}/\sh{I}^{k+1}$ is an extension of $\sh{I}/\sh{I}^{k}$ by $\sh{I}^k/\sh{I}^{k+1}$ \sref{III.1.4.17}.
\end{proof}

\begin{corollary}[16.1.6]
\label{IV.16.1.6}
Under the general hypotheses of \sref{IV.16.1.5}, if the immersion $f$ is locally of finite presentation then the $\shGr_n(f)$ are quasi-coherent $\sh{O}_Y$-modules of finite type.
\end{corollary}

\begin{proof}
Indeed, with the notation from the proof of \sref{IV.16.1.5}, $\sh{I}$ is an ideal of finite type of $\sh{O}_X$ \sref{IV.1.4.7}, therefore the $\sh{I}^n/\sh{I}^{n+1}$ are $\sh{O}_Y$-modules of finite type, hence the conclusion.
\end{proof}

\begin{corollary}[16.1.7]
\label{IV.16.1.7}
Under the general hypotheses of \sref{IV.16.1.5}, let $g:X \to Y$ be a morphism of preschemes, left inverse to $f$.
Then, for every $n$, the composite morphism $(1, \lambda_n): Y^{(n)}\xrightarrow{h_n} X \xrightarrow{g} Y$ defines a homomorphism of sheaves of rings $\lambda_n: \sh{O}_Y \to \sh{O}_{Y^{(n)}}$ right inverse to the augmentation $\vphi_{0n}$, making $\sh{O}_{Y^{(n)}}$ a quasi-coherent $\sh{O}_Y$-algebra;
via these homomorphisms, the transition homomorphisms $\vphi_{nm}:\sh{O}_{Y^{(m)}} \to \sh{O}_{Y^{(n)}}$ ($n\leq m$) are homomorphisms of $\sh{O}_Y$-algebras.
Also, if $g$ is locally of finite type, then the $\sh{O}_{Y^{(n)}}$ are quasi-coherent $\sh{O}_Y$-modules of finite type.
\end{corollary}

\begin{proof}
The first assertion follows immediately from the definitions and \sref{IV.16.1.5}.
On the other hand, if $g$ is locally of finite type, then $f$ is locally of finite presentation \sref{IV.1.4.3}[(v)];
the $\shGr_n(f)$ being then quasi-coherent $\sh{O}_Y$-modules of finite type by \sref{IV.16.1.6}, the same goes for the $\sh{O}_Y$-modules $\sh{I}/\sh{I}^{n+1}$, being extensions of a finite number of the $\shGr_k(f)$ \sref[III]{III.1.4.17}.
\end{proof}

\begin{proposition}[16.1.8]
\label{IV.16.1.8}
Let $X$ be a locally Noetherian prescheme and let $j:Y \to X$ be an immersion.
Then the $Y^{(n)}$ are locally Noetherian preschemes, the $\shGr_n(j)$ are coherent $\sh{O}_Y$-modules, and $\shGr_\bullet(j)$ is a coherent sheaf of rings on the space $Y$.
\end{proposition}

\begin{proof}
Everything is local on $X$ and $Y$, so we reduce to the case where $X$ is affine and $j$ is a closed immersion. All the assertions are then evident except the last; this follows from the fact that, if $A$ is a Noetherian ring and $\mathfrak{I}$ is an ideal of $A$, then $\gr_\mathfrak{I}^\bullet(A)$ is a Noetherian ring, taking into account the exactness of the functor $\psi^*$ and \sref[0]{0.5.3.7}.
\end{proof}

\begin{proposition}[16.1.9]
\label{IV.16.1.9}
\oldpage[IV-4]{8}
Let $X$ be a prescheme, let $j: Y \to X$ be an immersion locally of finite presentation, and let $y$ be a point of $Y$. The following conditions are equivalent:
\begin{enumerate}
  \item[(a)] There exists an open neighbourhood $U$ of $y$ in $Y$ such that $j|U$ is a homeomorphism of $U$ onto an open subset of $X$.
  \item[(b)] There is an integer $n>0$ such that the canonical homomorphism
  \[
    % The authority prints \vphi_{n,n-1}; the displayed map and the convention in 16.1.2 require \vphi_{n-1,n}.
    (\vphi_{n-1,n})_y: \sh{O}_{Y^{(n)},y} \to \sh{O}_{Y^{(n-1)},y}
  \]
  is bijective.
  \item[(c)] There is an integer $n>0$ such that $(\shGr_n(j))_y = 0$.
  
  In addition, if the integer $n$ satisfies \emph{(b)} or \emph{(c)}, then there is a neighbourhood $V$ of $y$ in $Y$ such that $\shGr_m(j)|V = 0$ for $m \geq n$ and that $\vphi_{nm}|V: \sh{O}_{Y^{(m)}}|V \to \sh{O}_{Y^{(n)}}|V$ is bijective for $m \geq n$. 
\end{enumerate}
\end{proposition}

\begin{proof}
Since the question is local on $Y$, we can restrict ourselves to the case where $j$ is a closed immersion, with $Y$ defined by a quasi-coherent ideal \emph{of finite type} $\sh{I}$ of $\sh{O}_X$.
The equivalence of (b) and (c), for a given $n$, is immediate;
also, since $\sh{I}^n/\sh{I}^{n+1}$ is an $\sh{O}_X$-module of finite type, there is an open neighbourhood $U$ of $y$ in $X$ such that $\sh{I}^n|U = \sh{I}^{n+1}|U$ \sref[0]{0.5.2.2}, and hence $\sh{I}^n|U = \sh{I}^m|U$ for $m \geq n$, proving the last assertions.
To prove that (a) implies (b), we can restrict ourselves to the case where the underlying space of $Y$ is equal to the underlying space of $X$ and where $\sh{I}$ is generated by finitely many sections over $X$:
since $\sh{I}$ is contained in the nilradical $\sh{N}$ of $\sh{O}_X$ \sref[I]{I.5.1.2}, it is nilpotent, which proves (b).
Finally, to prove that (b) implies (a), we can likewise restrict ourselves to the case where $\sh{I}^n = \sh{I}^{n+1}$;
therefore, for every $y \in Y$, since $\sh{I}_y \subset \mathfrak{m}_y$, maximal ideal of $\sh{O}_{X,y}$, we must have $\sh{I}^n_y = 0$ because of Nakayama's lemma, since $\sh{I}_y$ is an ideal of finite type.
The set of $x \in X$ such that $\sh{I}^n_x = 0$ is therefore an open subset $U$ of $X$ containing $Y$ \sref[0]{0.5.2.2};
since on the other hand $\sh{I}_x \neq 0$ for $x \notin Y$, we must have $U = Y$.
\end{proof}

\begin{corollary}[16.1.10]
\label{IV.16.1.10}
For the restriction of the immersion $j$ to a neighbourhood of $y$ in $Y$ to be an open immersion (in other words, for $j$ to be a \emph{local isomorphism} at the point $y$), it is necessary and sufficient that $(\shGr_1(j))_y = (\sh{N}_{Y/X})_y = 0$.
\end{corollary}

\begin{proof}
The condition is clearly necessary, and the previous reasoning applied to $n=1$ proves that it is sufficient.
\end{proof}

\begin{remark}[16.1.11]
\label{IV.16.1.11}
\medskip\noindent
\begin{enumerate}
  \item[(i)] Under the conditions of the definition \sref{IV.16.1.1}, the projective limit of the projective system $(\sh{O}_{Y^{(n)}}, \vphi_{nm})$ of sheaves of rings on $Y$ is called the \emph{normal invariant of infinite order} of $f$, and is sometimes denoted by $\sh{O}_{Y^{(\infty)}}$.
  When $X$ is a locally Noetherian prescheme and $j:Y \to X$ is a closed immersion, so that $Y$ is a closed subprescheme of $X$ defined by a coherent ideal $\sh{I}$, the sheaf $\sh{O}_{Y^{(\infty)}}$ is precisely the \emph{formal completion} of $\sh{O}_X$ along $Y$ \sref[I]{I.10.8.4}, and $Y^{(\infty)} = (Y, \sh{O}_{Y^{(\infty)}})$ is the formal prescheme obtained by \emph{completing} $X$ along $Y$ \sref[I]{I.10.8.5}.
  In all cases, one may say that $Y^{(\infty)}$ is the \emph{formal neighbourhood} of $Y$ in $X$ (for the morphism $f$).
  In the particular case we have just considered, it is the formal prescheme that is the inductive limit of the infinitesimal neighbourhoods of order $n$.
  \item[(ii)] Note that for a morphism of preschemes $f=(\psi, \theta): Y \to X$, it can happen that the homomorphism $\theta^\#:\psi^*(\sh{O}_X) \to \sh{O}_Y$ is surjective without $f$ being a local 
\oldpage[IV-4]{9}
  immersion and without $f$ being injective.
  We obtain an example by taking $Y$ to be a sum of preschemes $Y_\lambda$, all isomorphic to $\Spec(\sh{O}_x)$, where $x \in X$, and taking $f$ to be the morphism equal to the canonical morphism on each $Y_\lambda$.
\end{enumerate}
\end{remark}

\subsection{Functorial properties of the normal invariants of an immersion}
\label{IV.16.2}

\begin{env}[16.2.1]
\label{IV.16.2.1}
Let $f = (\psi, \theta): Y \to X$ and $f' = (\psi', \theta'): Y' \to X'$ be two morphisms of ringed spaces such that $\theta^\#$ and $\theta'^\#$ are surjective;
consider a commutative diagram of morphisms of ringed spaces
\[
  \label{IV.16.2.1.1}
  \xymatrix{
    Y \ar[r]^f & X \\
    Y'\ar[r]_{f'} \ar[u]^u & X'\ar[u]_v  
  }
  \tag{16.2.1.1}
\]

Let $u = (\rho, \lambda), v = (\sigma, \mu)$. 
We have $\rho^*(\psi^*(\sh{O}_X)) = \psi'^*(\sigma^*(\sh{O}_X))$, and hence a commutative diagram of homomorphisms of sheaves of rings on $Y'$:
\[
  \xymatrix{
    \rho^*(\psi^*(\sh{O}_X)) = \psi'^*(\sigma^*(\sh{O}_X)) \ar[r]^-{\psi'^*(\mu^\#)}\ar[d]_{\rho^*(\theta^\#)} & \psi'^*(\sh{O}_{X'}) \ar[d]^{\theta'^\#} \\
    \rho^*(\sh{O}_Y)\ar[r]_-{\lambda^\#}  & \sh{O}_{Y'}  
  }
\]
from which, if $\sh{I}$ and $\sh{I'}$ are the kernels of $\theta^\#$ and $\theta'^\#$, we conclude that $\psi'^*(\mu^\#)(\rho^*(\sh{I})) \subset \sh{I'}$, by the exactness of the functor $\rho^*$.
It follows immediately that, for every integer $n$, $\psi'^*(\mu^\#)(\rho^*(\sh{I}^n)) \subset \sh{I'}^n$, which shows that $\psi'^*(\mu^\#)$ defines, on passing to the quotients, a homomorphism of sheaves of rings
\[
  \label{IV.16.2.1.2}
  \nu_n: \rho^*(\psi^*(\sh{O}_X)/\sh{I}^{n+1}) \to \psi'^*(\sh{O}_{X'})/\sh{I'}^{n+1}
  \tag{16.2.1.2}
\]
and therefore a morphism of ringed spaces $w_n = (\rho, \nu_n): Y'^{(n)} \to Y^{(n)}$ (which, for $n = 0$, is none other than $u$).
It follows immediately from this definition that the diagrams
\[
  \xymatrix@R=1pc{
    Y^{(n)} \ar[r]^-{h_{mn}} & Y^{(m)} \ar[r]^-{h_m} & X \\
    & & & (n \leq m) \\
    Y'^{(n)} \ar[r]_-{h'_{mn}} \ar[uu]^-{w_n} & Y'^{(m)} \ar[r]_-{h'_m} \ar[uu]^-{w_m} & X' \ar[uu]_-v \\
  }
\]
(where the horizontal arrows are the canonical morphisms \sref{IV.16.1.2}) are commutative.

By passage to the quotients via the morphisms \sref{IV.16.2.1.2}, and taking into
\oldpage[IV-4]{10}
account the exactness of the functor $\rho^*$, we obtain a di-homomorphism of graded algebras (relative to the homomorphism $\lambda^\#: \rho^*(\sh{O}_Y) \to \sh{O}_{Y'}$)
\[
  \label{IV.16.2.1.3}
  \gr(u): \rho^*(\shGr_\bullet(f)) \to \shGr_\bullet(f')
  \tag{16.2.1.3}
\]
(or, if one prefers, a $\rho$-morphism \sref[0]{0.3.5.1} $\shGr_\bullet(f) \to \shGr_\bullet(f')$), and, in particular, a di-homomorphism of conormal sheaves
\[
  \gr_1(u): \rho^*(\shGr_1(f)) \to \shGr_1(f').
\]

It is also immediate that these homomorphisms give rise to a commutative diagram
\[
  \label{IV.16.2.1.4}
  \xymatrix{
    \rho^*(\bb{S}_{\sh{O}_Y}^\bullet(\shGr_1(f))) \ar[r] \ar[d]_-{\bb{S}(\gr_1(u))} & \rho^*(\shGr_\bullet(f)) \ar[d]^-{\gr(u)}\\
    \bb{S}_{\sh{O}_{Y'}}^\bullet(\shGr_1(f')) \ar[r] & \shGr_\bullet(f')
  }
  \tag{16.2.1.4}
\]
where the horizontal arrows are the canonical homomorphisms \sref{IV.16.1.2.2}.

Finally, if we have a commutative diagram of morphisms of ringed spaces
\[
  \xymatrix{
    Y \ar[r]^{f} & X \\
    Y' \ar[r]_{f'} \ar[u]^u & X' \ar[u]_v\\
    Y'' \ar[r]_{f''} \ar[u]^{u'} & X'' \ar[u]_{v'}\\
  }
\]
where $f'' = (\psi'', \theta'')$ is such that $\theta''^\#$ is surjective, and if $w_n'$ and $w_n''$ are defined, respectively, from $u',v'$ and from $u''=u\circ u',v''=v\circ v'$, then $w_n''=w_n\circ w_n'$, as follows immediately from the definitions and from \sref[0]{0.3.5.5};
we have also $\gr(u'') = \gr(u') \circ \rho'^*(\gr(u))$ if $u' = (\rho', \lambda')$.
Therefore we can say that $Y^{(n)}$ and $\shGr_\bullet(f)$ \emph{depend functorially} on $f$. 
\end{env}

\begin{proposition}[16.2.2]
\label{IV.16.2.2}
With the notation and hypotheses of \sref{IV.16.2.1}, suppose also that $f$, $f'$, $u$, and $v$ are morphisms of preschemes. We have:
\begin{enumerate}
  \item[{\rm(i)}] The morphisms $w_n:Y'^{(n)} \to Y^{(n)}$ are morphisms of preschemes.
  \item[{\rm(ii)}] If $Y'=Y\times_X X'$, with $u$ and $f'$ the canonical projections, and if $f$ is an immersion or $v$ is flat, then $Y'^{(n)}=Y^{(n)}\times_X X'$.
  \item[{\rm(iii)}] If $Y' = Y \times_X X'$ and if $v$ is flat (resp. if $f$ is an immersion), the homomorphism 
  \[
    \Gr(u)=\gr(u)\otimes 1 : \shGr_\bullet(f)\otimes_{\sh{O}_Y}\sh{O}_{Y'} \to \shGr_\bullet(f')
  \]
  is bijective (resp. surjective).
\end{enumerate}
\end{proposition}

\begin{proof}
\medskip\noindent
\begin{enumerate}
  \item[(i)] The hypotheses immediately imply that, for every $y' \in Y'$, $\rho_{y'}^*(\theta_{\psi'(y')}^\#)$ is a \emph{local} homomorphism \sref[I]{I.1.6.2}, so $w_n$ is a morphism of preschemes \sref[I]{I.2.2.1}.
  \oldpage[IV-4]{11}
  \item[(ii) and (iii)] If $f$ is an immersion, we may restrict to the case where $f$ is a closed immersion, with $Y$ defined by a quasi-coherent ideal $\sh{I}$ of $\sh{O}_X$ and $Y^{(n)}$ by the ideal $\sh{I}^{n+1}$;
  the assertions then follow from \sref[I]{I.4.4.5}.

  Second, suppose that $v$ is flat;
  we can restrict ourselves to the case where $X = \Spec(A)$, $Y = \Spec(B)$, $X' = \Spec(A')$ are affines, $A'$ being a flat $A$-module;
  so $Y' = \Spec(B')$ where $B' = B \otimes_A A'$;
  moreover, if $\mathfrak{I}$ is the kernel of the homomorphism $A\to B$, the kernel $\mathfrak{I}'$ of $A'\to B'$ is identified with $\mathfrak{I}\otimes_A A'$ by flatness, and
  \[
    \mathfrak{I}'^n/\mathfrak{I}'^{n+1}=(\mathfrak{I}^n/\mathfrak{I}^{n+1})\otimes_A A'.
  \]
  It follows at once from \sref[0]{0.4.3.3} that the $\sh{O}_{Y'}$-module $\sh{I}'^n/\sh{I}'^{n+1}$ is equal to
  \begin{align*}
    \psi'^*(\sigma^*((\mathfrak{I}^n/\mathfrak{I}^{n+1})\supertilde) \otimes_{\sigma^*(\sh{O}_X)} \sh{O}_{X'}) =& \\
    \psi'^*(\sigma^*((\mathfrak{I}^n/\mathfrak{I}^{n+1})\supertilde)) \otimes_{\psi'^*(\sigma^*(\sh{O}_X))} &\psi'^*(\sh{O}_{X'}) = \rho^*(\sh{I}^n/\sh{I}^{n+1})\otimes_{\rho^*(\psi^*(\sh{O}_X))} \psi'^*(\sh{O}_{X'})
  \end{align*}
  and in particular for $n = 0$, we have
  \[
    \sh{O}_{Y'} = \rho^*(\sh{O}_Y) \otimes_{\rho^*(\psi^*(\sh{O}_X))} \psi'^*(\sh{O}_{X'})
  \]
  whence we obtain a canonical isomorphism from $\sh{I}'^n/\sh{I}'^{n+1}$ onto
  \[
    \rho^*(\sh{I}^n/\sh{I}^{n+1})\otimes_{\rho^*(\sh{O}_Y)} \sh{O}_{Y'} = (\sh{I}^n/\sh{I}^{n+1}) \otimes_{\sh{O}_Y} \sh{O}_{Y'}
  \]
  which proves (iii).
  Set now $C_n=\Gamma(Y,\sh{O}_{Y^{(n)}})$ and $C'_n=\Gamma(Y',\sh{O}_{Y'^{(n)}})$.
  Since $Y^{(n)}$ and $Y'^{(n)}$ are affine schemes \sref{IV.16.1.5}, the kernel $\mathfrak{K}_n$ (resp. $\mathfrak{K}'_n$) of the homomorphism $C_n\to C_{n-1}$ (resp. $C'_n\to C'_{n-1}$) is $\Gamma(Y,\sh{I}^n/\sh{I}^{n+1})$ (resp. $\Gamma(Y',\sh{I}'^n/\sh{I}'^{n+1})$);
  therefore we can deduce from the above results that $\mathfrak{K}'_n = \mathfrak{K}_n \otimes_A A'$.
  Now, we have a commutative diagram
  \[
    \xymatrix{
      0 \ar[r] & \mathfrak{K}_n \ar[d]^-r\ar[r] \otimes_A A' & C_n \otimes_A A' \ar[d]^-{s_n}\ar[r] & C_{n-1} \otimes_A A' \ar[d]^-{s_{n-1}}\ar[r] & 0 \\
      0 \ar[r] & \mathfrak{K}'_n \ar[r] & C'_n \ar[r] & C'_{n-1} \ar[r] & 0
    }
  \]
  where the left vertical arrow is bijective and both rows are exact (since $A'$ is a flat $A$-module).
  It follows by induction that $s_n$ is bijective for every $n$: this is true by hypothesis for $n=0$, and the general case follows from the five lemma.
  This proves the second assertion of (ii).
\end{enumerate}
\end{proof}

\begin{corollary}[16.2.3]
\label{IV.16.2.3}
Let $g:X\to Y$ and $u:Y'\to Y$ be two morphisms of preschemes, let $X'=X\times_Y Y'$, and let $g':X'\to Y'$ and $v:X'\to X$ be the canonical projections. Let $f:Y\to X$ be a $Y$-section of $X$ (and hence an immersion), and let $f'=f_{(Y')}:Y'\to X'$ be the $Y'$-section of $X'$ deduced from $f$ by the base change $u$.
We have:
\begin{enumerate}
  \item[{\rm(i)}] The morphism $w_n:{Y'_{f'}}^{(n)}\to Y_f^{(n)}$ corresponding to $f,f',u,v$ \sref{IV.16.2.1} and the canonical morphism $h'_n:{Y'_{f'}}^{(n)}\to X'$ jointly identify ${Y'_{f'}}^{(n)}$ with the product $Y_f^{(n)}\times_X X'$.
  \item[{\rm(ii)}] If we endow $\sh{O}_{Y_f^{(n)}}$ (resp. $\sh{O}_{{Y'_{f'}}^{(n)}}$) with the structure of an $\sh{O}_Y$-algebra defined by $g$ (resp. with the structure of an $\sh{O}_{Y'}$-algebra defined by $g'$ ) \sref{IV.16.1.5}[(iii)],
  % The original citation is IV.16.1.6, but he clearly meant 16.1.5 item (iii)
  then the homomorphism of $\sh{O}_{Y'}$-algebras
  \[
    \label{IV.16.2.3.1}
    \rho^*(\sh{O}_{Y_f^{(n)}})\otimes_{\sh{O}_Y} \sh{O}_{Y'} \to \sh{O}_{{Y'_{f'}}^{(n)}}
    \tag{16.2.3.1}
  \]
\oldpage[IV-4]{12}
  induced by the homomorphism $\nu_n$ \sref{IV.16.2.1.2} is bijective.
  Also, the homomorphism of $\sh{O}_{Y'}$-modules
  \[
    \label{IV.16.2.3.2}
    \Gr_1(u): \shGr_1(f)\otimes_{\sh{O}_Y} \sh{O}_{Y'} \to \shGr_1(f')
    \tag{16.2.3.2}
  \]
  is bijective.
 \end{enumerate} 
\end{corollary}

\begin{proof}
\medskip\noindent
\begin{enumerate}
  \item[(i)] First note that $f':Y'\to X'$ and $u:Y'\to Y$ jointly identify $Y'$ with the product $Y\times_X X'$ (for the structure morphisms $f:Y\to X$ and $v:X'\to X$) \sref{IV.14.5.11.1}.\footnote{The printed source cites (14.5.12.1); the relevant product lemma is {\hyperlink{ega4.c04.lem.14-5-11-1}{(14.5.11.1)}}.}
  The conclusion of (i) now follows from \sref{IV.16.2.2}[(ii)], the morphism $g$ being an immersion.
  \item[(ii)] The commutative diagram
  \[
  \xymatrix{
    Y_f^{(n)} \ar[d]^{h_n}  & {Y'_{f'}}^{(n)} \ar[d]^-{h'_n} \ar[l]^{w_n}\\  
    X         \ar[d]^{g}    & X'              \ar[d]^{g'} \ar[l]^v \\  
    Y                       & Y' \ar[l]^u \\  
  }
  \]
  identifies ${Y'_{f'}}^{(n)}$ with the product $Y_f^{(n)}\times_X X'$ and identifies $X'$ with the product $X\times_Y Y'$; hence, by \sref[I]{I.3.3.9}, the morphisms $g'\circ h'_n$ and $w_n$ identify ${Y'_{f'}}^{(n)}$ with the product $Y_f^{(n)}\times_Y Y'$.
  Since $Y_f^{(n)}$ (resp. ${Y'_{f'}}^{(n)}$) is the affine prescheme over $Y$ (resp. over $Y'$) associated with the $\sh{O}_Y$-algebra $\sh{O}_{Y_f^{(n)}}$ (resp. the $\sh{O}_{Y'}$-algebra $\sh{O}_{{Y'_{f'}}^{(n)}}$), the bijectivity of the canonical homomorphism \sref{IV.16.2.3.1} follows from \sref[II]{II.1.5.2}.
  Finally, the canonical homomorphism \sref{IV.16.2.3.1} is compatible with the augmentations $\sh{O}_{Y_f^{(n)}} \to \sh{O}_Y$ and $\sh{O}_{{Y'_{f'}}^{(n)}} \to \sh{O}_{Y'}$;
  since $\sh{O}_{Y_f^{(n)}}$ is the direct sum (as an $\sh{O}_Y$-module) of $\sh{O}_Y$ and the augmentation ideal $\sh{I}/\sh{I}^{n+1}$, it follows that the canonical homomorphism \sref{IV.16.2.3.1}, restricted to $(\sh{I}/\sh{I}^{n+1})\otimes_{\sh{O}_Y}\sh{O}_{Y'}$, maps this module bijectively onto $\sh{I}'/\sh{I}'^{n+1}$. For $n=1$, this shows that $\Gr_1(u)$ is bijective.
\end{enumerate}
\end{proof}

We note that, under the hypotheses of \sref{IV.16.2.3}, the homomorphisms $\Gr_n(u)$ are \emph{surjective} in view of the above, but are not bijective in general for $n \geq 2$.
However:

\begin{corollary}[16.2.4]
\label{IV.16.2.4}
Under the hypotheses of \sref{IV.16.2.3}, suppose that $u: Y' \to Y$ is a flat morphism (resp. that the $\shGr_n(f)$ are flat $\sh{O}_Y$-modules for $n \leq m$).
Then the homomorphism 
\[
  \Gr_n(u):\shGr_n(f) \otimes_{\sh{O}_Y}\sh{O}_{Y'} \to \shGr_n(f')
\]
is bijective for all $n$ (resp. for $n \leq m$).
\end{corollary}

\begin{proof}
If $u$ is flat, then $v:X'\to X$ is flat by base change, and in this case we already know that $\Gr(u)$ is bijective \sref{IV.16.2.2}[(iii)].
If the $\shGr_n(f)$ are flat for $n\leq m$, then we first see by induction on $n$ that the same holds for $\sh{I}/\sh{I}^{n+1}$ for $n\leq m$, because of the exact sequences
  \[
    \xymatrix{
      0 \ar[r] & \sh{I}^n/\sh{I}^{n+1} \ar[r] & \sh{I}/\sh{I}^{n+1} \ar[r] & \sh{I}/\sh{I}^{n} \ar[r] & 0
    }
  \]
\oldpage[IV-4]{13}
  \sref[0]{0.6.1.2};
  in addition, we have the commutative diagram
  \[
    \xymatrix{
      0 \ar[r] & (\sh{I}^n/\sh{I}^{n+1}) \otimes_{\sh{O}_Y}\sh{O}_{Y'} \ar[d]\ar[r] &( \sh{I}/\sh{I}^{n+1}) \otimes_{\sh{O}_Y}\sh{O}_{Y'} \ar[d]\ar[r] & (\sh{I}/\sh{I}^{n}) \otimes_{\sh{O}_Y}\sh{O}_{Y'} \ar[d]\ar[r] & 0 \\
      0 \ar[r] & \sh{I}'^n/\sh{I}'^{n+1} \ar[r] & \sh{I}'/\sh{I}'^{n+1} \ar[r] & \sh{I}'/\sh{I}'^{n} \ar[r] & 0
    }
  \]
  in which the rows are exact (the first by flatness \sref[0]{0.6.1.2}) and the last two vertical arrows are bijective by \sref{IV.16.2.2}[(ii)];
  hence the conclusion.
\end{proof}

\begin{remarks}[16.2.5]
\label{IV.16.2.5}
\medskip\noindent
\begin{enumerate}
  \item[(i)] The reasoning of \sref{IV.16.2.2}[(i)] still applies when the maps in \sref{IV.16.2.1.1} are morphisms of \emph{locally ringed spaces} ($\mathbf{Err_{II}}$, \sref[I]{I.1.8.2}).
  \item[(ii)] In \sref{IV.16.2.2}[(ii)], the conclusion is no longer necessarily valid if we only suppose that $v$ and $f$ are morphisms of preschemes ($f$ satisfying the condition of \sref{IV.16.1.1}).
  For example (with the notation of the proof of \sref{IV.16.2.2}[(ii)]), it can happen that $\mathfrak{I} = 0$ but the kernel $\mathfrak{I}'$ of $A' \to B' = B \otimes_A A'$ is not zero and that $B' \neq 0$, in which case we have $Y^{(n)} = Y$ for all $n$, but ${Y'}^{(n)} \neq Y'$.
  We have an example of this by taking $A = \bb{Z}$, $B = \bb{Q}$, $A' = \prod_{h = 1}^\infty (\bb{Z}/m^h\bb{Z})$ where $m>1$.
\end{enumerate}
\end{remarks}

\begin{env}[16.2.6]
\label{IV.16.2.6}
Consider the special case of diagram \sref{IV.16.2.1.1} in which $X'=X$, $v$ is the identity, $X$ is a prescheme, $Y$ is a subprescheme of $X$, $Y'$ is a subprescheme of $Y$, and $f,u,f'=f\circ u$ are the canonical injections;
by tensoring with $\sh{O}_{Y'}$ over $\rho^*(\sh{O}_Y)$, the di-homomorphism \sref{IV.16.2.1.3} gives a homomorphism of graded $\sh{O}_{Y'}$-algebras
\[
  \label{IV.16.2.6.1}
  u^*(\shGr_\bullet(f)) \to \shGr_\bullet(f').
  \tag{16.2.6.1}
\]
On the other hand, we identify $\sh{O}_Y$ with $\psi^*(\sh{O}_X)/\sh{I}_f$ and $\sh{O}_{Y'}$ with $\rho^*(\sh{O}_Y)/\sh{I}_u$;
since $\rho^*$ is exact, we have $\rho^*(\sh{O}_Y)=\rho^*(\psi^*(\sh{O}_X))/\rho^*(\sh{I}_f)=\psi'^*(\sh{O}_X)/\rho^*(\sh{I}_f)$, and since $\sh{O}_{Y'}$ is also identified with $\psi'^*(\sh{O}_X)/\sh{I}_{f'}$, we see that $\sh{I}_u=\sh{I}_{f'}/\rho^*(\sh{I}_f)$.
Thus, for every integer $n$, there is a canonical homomorphism $\sh{I}_{f'}^n/\sh{I}_{f'}^{n+1}\to \sh{I}_{u}^n/\sh{I}_{u}^{n+1}$, and hence a canonical homomorphism of graded $\sh{O}_{Y'}$-algebras
\[
  \label{IV.16.2.6.2}
  \shGr_\bullet(f') \to \shGr_\bullet(u).
  \tag{16.2.6.2}
\]
\end{env}

\begin{proposition}[16.2.7]
\label{IV.16.2.7}
Let $X$ be a prescheme, let $Y$ be a subprescheme of $X$, let $Y'$ be a subprescheme of $Y$, and let $j:Y'\to Y$ be the canonical injection.
We then have an exact sequence of conormal sheaves ($\sh{O}_{Y'}$-modules)
\[
  \label{IV.16.2.7.1}
  \xymatrix{
    j^*(\sh{N}_{Y/X}) \ar[r] & \sh{N}_{Y'/X} \ar[r] & \sh{N}_{Y'/Y} \ar[r] & 0
  }
  \tag{16.2.7.1}
\]
where the arrows are the degree $1$ components of the canonical homomorphisms \sref{IV.16.2.6.1} and \sref{IV.16.2.6.2}.
\end{proposition}

\begin{proof}
Since the question is local, we may restrict to the case where $X=\Spec(A)$, $Y=\Spec(A/\mathfrak{I})$, and $Y'=\Spec(A/\mathfrak{K})$, where $\mathfrak{I}$ and $\mathfrak{K}$ are ideals of $A$ such that $\mathfrak{I}\subset\mathfrak{K}$;
everything reduces to seeing 
\oldpage[IV-4]{14}
that the sequence of canonical morphisms $\mathfrak{I}/\mathfrak{K}\mathfrak{I} \to \mathfrak{K}/\mathfrak{K}^2 \to (\mathfrak{K}/\mathfrak{I})/(\mathfrak{K}/\mathfrak{I})^2 \to 0$ is exact, which is immediate given that the image of $\mathfrak{I}/\mathfrak{K}\mathfrak{I}$ in $\mathfrak{K}/\mathfrak{K}^2$ is $(\mathfrak{I} + \mathfrak{K}^2)/\mathfrak{K}^2$ and that $(\mathfrak{K}/\mathfrak{I})/(\mathfrak{K}/\mathfrak{I})^2$ is identified with $\mathfrak{K}/(\mathfrak{I} + \mathfrak{K}^2)$.
\end{proof}

It is easy to give examples where the sequence \sref{IV.16.2.7.1} extended on the left by $0$ is not exact;
with the above notation, it suffices to take $A = k[T]$, $\mathfrak{I} = AT^2$, $\mathfrak{K} = AT$, because then $(\mathfrak{I} + \mathfrak{K}^2)/\mathfrak{K}^2 = 0$ and $\mathfrak{I}/\mathfrak{K}\mathfrak{I} \neq 0$.
See however \sref{IV.16.9.13} and \sref{IV.19.1.5} for some cases where the extended sequence is indeed exact.

\subsection{Fundamental differential invariants of morphisms of preschemes}
\label{IV.16.3}

\begin{definition}[16.3.1]
\label{IV.16.3.1}
Let $f:X\to S$ be a morphism of preschemes, and let $\Delta_f:X\to X\times_S X$ be the corresponding diagonal morphism, which is an immersion (\sref[I]{I.5.3.9} and $\mathbf{Err_{III}}$, 10).
We denote by $\sh{P}_f^n$ or $\sh{P}_{X/S}^n$, and call the \emph{sheaf of principal parts of order $n$ of the $S$-prescheme $X$}, the $\sh{O}_X$-augmented sheaf of rings that is the $n$-th normal invariant of $\Delta_f$ \sref{IV.16.1.2}.
We will also write $\sh{P}_f^\infty = \sh{P}_{X/S}^\infty = \varprojlim_n \sh{P}_{X/S}^n$, $\shGr_n(\sh{P}_f) = \shGr_n(\sh{P}_{X/S}) = \shGr_n(\Delta_f)$ \sref{IV.16.1.2};
the $\sh{O}_X$-module $\shGr_1(\Delta_f)$, the augmentation ideal of $\sh{P}_{X/S}^1$, is also denoted by $\Omega_f^1$ or $\Omega_{X/S}^1$, and is called the $\sh{O}_X$-module of \emph{$1$-differentials} of $f$, or of $X$ with respect to $S$, or of the $S$-prescheme $X$.
\end{definition}

It follows from this definition that $\sh{P}_{X/S}^0$ is canonically identified with $\sh{O}_X$ \sref{IV.16.1.2}.

We have \sref{IV.16.1.2.2} a canonical surjective homomorphism of graded $\sh{O}_X$-algebras
\[
  \label{IV.16.3.1.1}
  \bb{S}_{\sh{O}_X}^\bullet(\Omega_{X/S}^1) \to \shGr_\bullet(\sh{P}_{X/S}).
  \tag{16.3.1.1}
\]
It also follows from Definition~\sref{IV.16.3.1} that for every open $U$ of $X$ we have $\sh{P}_{f|U}^n=\sh{P}_f^n|U$, $\sh{P}_{f|U}^\infty=\sh{P}_f^\infty|U$, $\shGr_n(\sh{P}_{f|U})=\shGr_n(\sh{P}_f)|U$, and $\Omega_{f|U}^1=\Omega_f^1|U$ (in other words, the notions introduced are \emph{local} on $X$).

\begin{env}[16.3.2]
\label{IV.16.3.2}
Let $p_1,p_2$ be the two canonical projections of the product $X\times_S X$;
since $\Delta_f$ is an $X$-section of $X\times_S X$ for both $p_1$ and $p_2$, \emph{each} of these morphisms defines, for every $n$, a homomorphism of sheaves of rings $\sh{O}_X\to\sh{P}_{X/S}^n$ that is right inverse to the augmentation $\sh{P}_{X/S}^n\to\sh{O}_X$ \sref{IV.16.1.7};
equivalently, this defines on $\sh{P}_{X/S}^n$ \emph{two structures of quasi-coherent augmented $\sh{O}_X$-algebra};
the corresponding $\sh{O}_X$-module structures on $\shGr_n(\sh{P}_{X/S})$ are the same.
We also have, by passing to the limit, two $\sh{O}_X$-algebra structures on $\sh{P}_{X/S}^\infty$.
\end{env}

\begin{env}[16.3.3]
\label{IV.16.3.3}
The morphism $s = (p_2, p_1)_S: X \times_S X \to X \times_S X$ is an \emph{involutive automorphism} of $X \times_S X$, called the \emph{canonical symmetry}, such that
\[
  \label{IV.16.3.3.1}
  p_1 \circ s = p_2, \qquad p_2 \circ s = p_1, \qquad s \circ \Delta_f = \Delta_f.
  \tag{16.3.3.1}
\]

If we set $s=(\rho,\lambda)$, $p_i=(\pi_i,\mu_i)$ ($i=1,2$), and $\Delta_f=(\delta,\nu)$, then $\lambda^\#$ is an isomorphism from $\rho^*(\pi_1^*(\sh{O}_X))$ onto $\pi_2^*(\sh{O}_X)$, and $\delta^*(\lambda^\#)$ leaves invariant $\delta^*(\sh{O}_{X\times_S X})$ and the kernel $\sh{I}$ of the homomorphism $\nu^\#: \delta^*(\sh{O}_{X\times_S X})\to\sh{O}_X$.
Therefore:
\end{env}

\begin{proposition}[16.3.4]
\label{IV.16.3.4}
The homomorphism $\sigma = \delta^*(\lambda^\#)$ induced from $s$ (and also called the \emph{canonical symmetry}) is an involutive automorphism of the projective system $(\sh{P}_{X/S}^n)$ of $\sh{O}_X$-augmented 
\oldpage[IV-4]{15}
sheaves of rings, and as a result also of the projective limit $\sh{P}_{X/S}^\infty$.
This automorphism interchanges the two $\sh{O}_X$-algebra structures on $\sh{P}_{X/S}^n$ and on $\sh{P}_{X/S}^\infty$.
\end{proposition}

\begin{env}[16.3.5]
\label{IV.16.3.5}
In what follows, the two $\sh{O}_X$-algebra structures defined on the $\sh{P}_{X/S}^n$ and on $\sh{P}_{X/S}^\infty$ will play very different roles:
\emph{we will now agree, unless said otherwise, that when $\sh{P}_{X/S}^n$ or $\sh{P}_{X/S}^\infty$ is considered as an $\sh{O}_X$-algebra, it is the algebra structure induced by $p_1$}.
\end{env}

For every open $U$ of $X$ and every section $t \in \Gamma(U, \sh{O}_X)$, we will simply denote by $t.1$ or even $t$ the image of $t$ under the structure morphism $\Gamma(U, \sh{O}_X) \to \Gamma(U, \sh{P}_{X/S}^n)$ (resp. $\Gamma(U, \sh{O}_X) \to \Gamma(U, \sh{P}_{X/S}^\infty)$) (that is to say, the homomorphism corresponding to $p_1$).

\begin{definition}[16.3.6]
\label{IV.16.3.6}
We denote by $d_f^n$, or $d_{X/S}^n$ (resp. $d_f^\infty$, or $d_{X/S}^\infty$), or simply $d^n$ (resp. $d^\infty$), the homomorphism of sheaves of rings $\sh{O}_X \to \sh{P}_{f}^n = \sh{P}_{X/S}^n$ (resp. $\sh{O}_X \to \sh{P}_{f}^\infty = \sh{P}_{X/S}^\infty$) induced by $p_2$.
For every open $U$ of $X$, and every $t \in \Gamma(U, \sh{O}_X)$, $d^n t$ (resp. $d^\infty t$) is called the \emph{principal part of order $n$} (resp. \emph{principal part of infinite order}) of $t$.
We set $dt = d^1 t - t$, and we say that $dt$ is the differential of $t$ (an element of $\Gamma(U, \Omega_{X/S}^1)$, also denoted $d_{X/S}(t)$).
\end{definition}  

It follows immediately
\footnote{
  [Translator's note.] Locally, this is \sref[0]{0.20.1.1}.
}
from this definition that we have
\[
  \label{IV.16.3.6.1}
  d(t_1 t_2) = t_1 dt_2 + t_2 dt_1
  \tag{16.3.6.1} 
\]
for every $t_1,t_2$ in $\Gamma(U,\sh{O}_X)$; in other words, $d$ is a \emph{derivation} of the ring $\Gamma(U,\sh{O}_X)$ with values in the $\Gamma(U,\sh{O}_X)$-module $\Gamma(U,\Omega_{X/S}^1)$.

In all the notation introduced in \sref{IV.16.3.1} and \sref{IV.16.3.6}, we will sometimes replace $S$ by $A$ when $S=\Spec(A)$.

\begin{env}[16.3.7]
\label{IV.16.3.7}
Suppose in particular that $S = \Spec(A)$ and $X = \Spec(B)$ are affine schemes, $B$ then being an $A$-algebra.
Then $\Delta_f$ corresponds to the canonical surjective homomorphism $\pi: B \otimes_A B \to B$ such that $\pi(b\otimes b') = bb'$, with kernel $\mathfrak{I} = \mathfrak{I}_{B/A}$ \sref[0]{0.20.4.1};
$\sh{P}_{f}^n$ is the structure sheaf of the prescheme $\Spec(P_{B/A}^n)$, where
\[
  P_{B/A}^n = (B \otimes_A B)/\mathfrak{I}^{n+1};
\]
$\shGr_\bullet(\sh{P}_f)$ is the quasi-coherent $\sh{O}_X$-module corresponding to the graded $B$-module
\[
  \gr_\mathfrak{I}^\bullet(B \otimes_A B) = \bigoplus_{n \geq 0} (\mathfrak{I}^n/\mathfrak{I}^{n+1});
\]
in particular $\Omega_f^1 = \Omega_{X/S}^1$ is the quasi-coherent $\sh{O}_X$-module corresponding to the $B$-module of $1$-differentials of $B$ over $A$, $\Omega_{B/A}^1$ \sref[0]{0.20.4.3}.
The projection morphisms $p_1:X\times_S X\to X$ and $p_2:X\times_S X\to X$ correspond to the two homomorphisms of rings $j_1:B\to B\otimes_A B$ and $j_2:B\to B\otimes_A B$ such that $j_1(b)=b\otimes 1$ and $j_2(b)=1\otimes b$; thus (by the convention of \sref{IV.16.3.5}), $P_{B/A}^n$ is always considered as a $B$-algebra via the composite homomorphism $B\xrightarrow{j_1}B\otimes_A B\to P_{B/A}^n$;
the ring homomorphism $B \xrightarrow{j_2} B \otimes_A B \to P_{B/A}^n$ is denoted by $d_{B/A}^n$ and corresponds to $d_{X/S}^n$ acting on $\Gamma(X, \sh{O}_X)$;
for every $t\in B$, $dt$ is equal to $d_{B/A}t$, defined in \sref[0]{0.20.4.6} (cf. $\mathbf{Err_{IV}}$, 11).

If $\pi_n:B\otimes_A B\to P_{B/A}^n$ is the canonical homomorphism, then the preceding definitions give
\[
  \label{IV.16.3.7.1}
  \pi_n(b\otimes b') = b \cdot \pi_n(1\otimes b') = b \cdot d_{B/A}^n(b') \quad \text{for } b \in B, b' \in B.
  \tag{16.3.7.1}
\]
\end{env}

\begin{proposition}[16.3.8]
\label{IV.16.3.8}
\oldpage[IV-4]{16}
The image of the canonical homomorphism $d_{X/S}^n: \sh{O}_X \to \sh{P}_{X/S}^n$ generates the $\sh{O}_X$-module $\sh{P}_{X/S}^n$.
\end{proposition}

\begin{proof}
We immediately reduce to the case where $X=\Spec(B)$ and $S=\Spec(A)$ are affine; the proposition then follows from \sref{IV.16.3.7.1}, since $\pi_n$ is surjective.
We note that in general $d_{X/S}^n$ \emph{is not surjective} (even for $n = 1$).
\end{proof}

\begin{proposition}[16.3.9]
\label{IV.16.3.9}
Suppose that $f:X \to S$ is a morphism locally of finite type.
Then the $\sh{P}_{f}^n$ and the $\shGr_n(\sh{P}_{f})$ are quasi-coherent $\sh{O}_X$-modules of finite type.
\end{proposition}

\begin{proof}
This follows from \sref{IV.16.1.6} and from the fact that $\Delta_f$ is locally of finite presentation \sref[I]{I.4.3.1}.
\end{proof}

\subsection{Functorial properties of differential invariants}
\label{IV.16.4}

\begin{env}[16.4.1]
\label{IV.16.4.1}
Consider a commutative diagram of morphisms of preschemes
\[
  \label{IV.16.4.1.1}
  \xymatrix{
    X \ar[d]_-{f} & X' \ar[l]_-u \ar[d]^-{f'}\\
    S & S' \ar[l]^-w
  }
  \tag{16.4.1.1}
\]
We deduce a commutative diagram
\[
  \xymatrix{
    X \ar[d]_-{\Delta_f} & X' \ar[l]_-u \ar[d]^-{\Delta_{f'}}\\
    X \times_S X & X' \times_{S'} X' \ar[l]^-v
  }
\]
where $v$ is the composite morphism \sref[I]{I.5.3.5} and \sref[I]{I.5.3.15}:
\[
  \label{IV.16.4.1.2}
  X' \times_{S'} X' \xrightarrow{(p'_1, p'_2)_S} X' \times_{S} X' \xrightarrow{u \times_S u} X \times_S X.
  \tag{16.4.1.2}
\]

Thus $u$ and $v$ give, as explained in \sref{IV.16.2.1}, homomorphisms of augmented sheaves of rings
\[
  \label{IV.16.4.1.3}
  \nu_n: \rho^*(\sh{P}_{X/S}^n) \to \sh{P}_{X'/S'}^n
  \tag{16.4.1.3}
\]
(where we put $u = (\rho,\lambda)$);
these homomorphisms form a projective system and therefore give, in the limit, a homomorphism of augmented sheaves of rings
\[
  \label{IV.16.4.1.4}
  \nu_\infty: \rho^*(\sh{P}_{X/S}^\infty) \to \sh{P}_{X'/S'}^\infty;
  \tag{16.4.1.4}
\]
on the other hand, by passing to the quotients, the homomorphisms $\nu_n$ give rise to a di-homomorphism of graded algebras (relative to $\lambda^\#$):
\[
  \label{IV.16.4.1.5}
  \gr(u): \rho^*(\shGr_\bullet(\sh{P}_{X/S})) \to \shGr_\bullet(\sh{P}_{X'/S'}).
  \tag{16.4.1.5}
\]
\end{env}

\begin{env}[16.4.2]
\label{IV.16.4.2}
\oldpage[IV-4]{17}
If we have a commutative diagram
\[
  \xymatrix{
    X \ar[d]_-{f} & X' \ar[l]_-u \ar[d]^-{f'} & X'' \ar[l]_-{u'} \ar[d]^-{f''}\\
    S & S' \ar[l]^-w & S'' \ar[l]^-{w'} 
  }
\]
we deduce a commutative diagram
\[
  \xymatrix{
    X \ar[d]_-{\Delta_f} & X' \ar[l]_-u \ar[d]^-{\Delta_{f'}} & X'' \ar[l]_-{u'} \ar[d]^-{\Delta_{f''}}\\
    X \times_S X & X' \times_{S'} X' \ar[l]^-v & X'' \times_{S''} X'' \ar[l]^-{v'}
  }
\]
where $v'$ is defined from $u'$, $w'$, $f'$, $f''$ as $v$ is from $u$, $w$, $f$, $f'$.
We verify immediately that if $u''=u\circ u'$ and $w''=w\circ w'$, then the composite morphism $v\circ v'$ is equal to the morphism $v''$ deduced from $u'',w'',f,f''$ as $v$ is deduced from $u,w,f,f'$.
If we set $u'=(\rho',\lambda')$ and $u''=(\rho'',\lambda'')$, it follows from \sref{IV.16.2.1} that the homomorphism $\nu''_n:{\rho''}^*(\sh{P}_{X/S}^n)\to\sh{P}_{X''/S''}^n$ is equal to the composite
\[
  {\rho'}^*(\rho^*(\sh{P}_{X/S}^n)) \xrightarrow{{\rho'}^*(\nu_n^\#)} {\rho'}^*(\sh{P}_{X'/S'}^n) \xrightarrow{\nu'_n} \sh{P}_{X''/S''}^n,
\]
and there are analogous transitivity properties for the homomorphisms \sref{IV.16.4.1.4} and \sref{IV.16.4.1.5}, allowing us to say that $\sh{P}_{X/S}^n$, $\sh{P}_{X/S}^\infty$, and $\shGr_\bullet(\sh{P}_{X/S})$ \emph{depend functorially on $f$}.
\end{env}

\begin{env}[16.4.3]
\label{IV.16.4.3}
We verify immediately (for example, by reducing to the affine case using \sref{IV.16.3.7}) that, with the notation of \sref{IV.16.4.1}, the diagram
\[
  \label{IV.16.4.3.1}
  \xymatrix{
    \rho^*(\sh{O}_X) \ar[r]^-{\lambda^\#} \ar[d] & \sh{O}_{X'} \ar[d] \\
    \rho^*(\sh{P}_{X/S}^n) \ar[r]_-{\nu_n} & \sh{P}_{X'/S'}^n  
  }
  \tag{16.4.3.1}
\]
is commutative, where the vertical arrows define the algebra structures chosen in \sref{IV.16.3.5} (that is, those coming from the first projections);
the same goes for the diagram
\[
  \label{IV.16.4.3.2}
  \xymatrix{
    \rho^*(\sh{O}_X) \ar[r]^-{\lambda^\#} \ar[d]_-{\rho^*(d_{X/S}^n)} & \sh{O}_{X'} \ar[d]^-{d_{X'/S'}^n} \\
    \rho^*(\sh{P}_{X/S}^n) \ar[r]_-{\nu_n} & \sh{P}_{X'/S'}^n  
  }
  \tag{16.4.3.2}
\]
\oldpage[IV-4]{18}
the vertical arrows here define the algebra structures coming from the second projections;
besides, if $\sigma$ and $\sigma'$ are the canonical symmetries corresponding to $f$ and $f'$ \sref{IV.16.3.4}, we have
\[
  \nu_n \circ \rho^*(\sigma) = \sigma' \circ \nu_n 
\] 
which carries one of the preceding diagrams to the other.
We therefore deduce from \sref{IV.16.4.3.1} a canonical homomorphism of \emph{augmented $\sh{O}_{X'}$-algebras}
\[
  \label{IV.16.4.3.3}
  P^n(u): u^*(\sh{P}_{X/S}^n) = \sh{P}_{X/S}^n \otimes_{\sh{O}_X} \sh{O}_{X'} \to \sh{P}_{X'/S'}^n 
  \tag{16.4.3.3}
\]
and it follows from \sref{IV.16.4.3.2} that the diagram 
\[
  \label{IV.16.4.3.4}
  \xymatrix{
    \sh{O}_{X'} \ar[r]^-{\operatorname{id}} \ar[d]_-{u^*(d_{X/S}^n)}  & \sh{O}_{X'} \ar[d]^-{d_{X'/S'}^n} \\
    u^*(\sh{P}_{X/S}^n) \ar[r]_{P^n(u)} & \sh{P}_{X'/S'}^n
  }
  \tag{16.4.3.4}
\]
is commutative.
We thus obtain a homomorphism of \emph{graded $\sh{O}_{X'}$-algebras}
\[
  \label{IV.16.4.3.5}
  \Gr_\bullet(u):u^*(\shGr_\bullet(\sh{P}_{X/S})) \to \shGr_\bullet(\sh{P}_{X'/S'})
  \tag{16.4.3.5}
\]
and in particular a homomorphism of $\sh{O}_{X'}$-modules 
\[
  \label{IV.16.4.3.6}
  \Gr_1(u):\Omega_{X/S}^1 \otimes_{\sh{O}_X} \sh{O}_{X'} \to \Omega_{X'/S'}^1
  \tag{16.4.3.6}
\]
giving rise to a commutative diagram
\[
  \label{IV.16.4.3.7}
  \xymatrix{
    \sh{O}_{X'} \ar[r]^-{\operatorname{id}} \ar[d]_-{d_{X/S} \otimes 1}  & \sh{O}_{X'} \ar[d]^-{d_{X'/S'}} \\
    \Omega_{X/S}^1 \otimes_{\sh{O}_X} \sh{O}_{X'} \ar[r] & \Omega_{X'/S'}^1
  }
  \tag{16.4.3.7}
\]
\end{env}

\begin{env}[16.4.4]
\label{IV.16.4.4}
When $S = \Spec(A)$, $S' = \Spec(A')$, $X = \Spec(B)$, $X' = \Spec(B')$ are affine, so that we have a commutative diagram of ring homomorphisms
\[
  \xymatrix{
    B \ar[r] & B' \\
    A \ar[u] \ar[r] & A' \ar[u]
  }
\]
the image of $\mathfrak{I}_{B/A}$ in $B' \otimes_{A'} B'$ is contained in $\mathfrak{I}_{B'/A'}$, and the homomorphism $\nu_n$ corresponds to the homomorphism of rings $P_{B/A}^n \to P_{B'/A'}^n$ induced from the homomorphism $B \otimes_A B \to B' \otimes_{A'} B'$ by passing to quotients.
The homomorphism \sref{IV.16.4.3.6} corresponds to the homomorphism defined in \sref[0]{0.20.5.4.1}, and the commutative diagram \sref{IV.16.4.3.7} to the diagram \sref[0]{0.20.5.4.2}.
\end{env}

\begin{proposition}[16.4.5]
\label{IV.16.4.5}
\oldpage[IV-4]{19}
Suppose that $X' = X \times_S S'$, $f'$ and $u$ the canonical projections.
Then the canonical homomorphisms $P^n(u)$ \sref{IV.16.4.3.3} and $\Gr_1(u)$ \sref{IV.16.4.3.6} are bijective.
\end{proposition}

\begin{proof}
We have $X' \times_{S'} X' = (X \times_S X) \times_S S'$, and it suffices to apply \sref{IV.16.2.3}[(ii)], replacing $g$ by the first projection $p_1:X\times_S X \to X$ and $f$ by the diagonal $\Delta_f$.
\end{proof}

We note that under the hypotheses of \sref{IV.16.4.5} the homomorphism $\Gr_\bullet(u)$ \sref{IV.16.4.3.5} is \emph{surjective}, but not bijective in general. 
However \sref{IV.16.2.4}:

\begin{corollary}[16.4.6]
\label{IV.16.4.6}
Under the hypotheses of \sref{IV.16.4.5}, suppose in addition that $w: S' \to S$ is flat (resp. that $\shGr_n(\sh{P}_{X/S})$ are flat $\sh{O}_X$-modules for $n \leq m$);
then the homomorphism
\[
  \Gr_n(u):u^*(\shGr_n(\sh{P}_{X/S})) \to \shGr_n(\sh{P}_{X'/S'})
\]
is bijective for each $n$ (resp. for $n \leq m$).
\end{corollary}

\begin{proof}
Indeed, if $w$ is flat, then so is $v: X' \times_{S'} X' \to X \times_S X$, so the conclusion follows from \sref{IV.16.2.4}.
\end{proof}

\begin{env}[16.4.7]
\label{IV.16.4.7}
Let $S$ be a prescheme, let $\sh{E}$ be a quasi-coherent $\sh{O}_S$-module, and set $X = \bb{V}(\sh{E})$ \sref[II]{II.1.7.8}, the vector bundle associated with $\sh{E}$, equal to $\Spec(\bb{S}_{\sh{O}_S}(\sh{E}))$.
Let $f:X \to S$ be the structure morphism.
For every open $U$ of $S$ and every section $t \in \Gamma(U, \sh{E})$, $t$ is identified with a section of $\bb{S}_{\sh{O}_S}(\sh{E})$ over $U$;
let $t'$ be its image in $\Gamma(f^{-1}(U), \sh{O}_X) = \Gamma(U, f_*(\sh{O}_X)) = \Gamma(U, \bb{S}_{\sh{O}_S}(\sh{E}))$, and set
\[
  \label{IV.16.4.7.1}
  \delta(t) = d_{X/S}^n(t') - t' \in \Gamma(f^{-1}(U), \sh{P}_{X/S}^n);
  \tag{16.4.7.1}
\]
it is clear that $\delta$ is a di-homomorphism of modules (corresponding to the homomorphism of rings $\Gamma(U, \sh{O}_S) \to \Gamma(f^{-1}(U), \sh{O}_X)$) from $\Gamma(U, \sh{E})$ to $\Gamma(f^{-1}(U), \sh{P}_{X/S}^n)$, whose image moreover belongs to the augmentation ideal of $\Gamma(f^{-1}(U), \sh{P}_{X/S}^n)$.
We deduce (by varying $U$) a canonical homomorphism of $\sh{O}_X$-algebras
\[
  \label{IV.16.4.7.2}
  f^*(\bb{S}_{\sh{O}_S}(\sh{E})) \to \sh{P}_{X/S}^n
  \tag{16.4.7.2}
\]
and in view of the preceding remark, if $\sh{K}$ is the kernel ideal of the augmentation $\bb{S}_{\sh{O}_S}(\sh{E}) \to \sh{O}_S$, the image of $\sh{K}^{n+1}$ under \sref{IV.16.4.7.2} is zero; hence, after factoring by $\sh{K}^{n+1}$, we finally obtain a canonical homomorphism
\[
  \label{IV.16.4.7.3}
  \delta_n: f^*(\bb{S}_{\sh{O}_S}(\sh{E})/\sh{K}^{n+1}) \to \sh{P}_{X/S}^n.
  \tag{16.4.7.3}
\]
\end{env} 

\begin{proposition}[16.4.8]
\label{IV.16.4.8}
Under the conditions of \sref{IV.16.4.7}, the homomorphisms $\delta_n$ are bijective and form a projective system of isomorphisms;
we deduce an isomorphism of graded $\sh{O}_X$-algebras 
\[
  \label{IV.16.4.8.1}
  f^*(\bb{S}_{\sh{O}_S}^\bullet(\sh{E})) \to \shGr_\bullet(\sh{P}_{X/S}).
  \tag{16.4.8.1}
\]
\end{proposition}

\begin{proof}
The fact that homomorphisms \sref{IV.16.4.7.3} form a projective system follows immediately from their definition.
To prove they are isomorphisms, it suffices to 
\oldpage[IV-4]{20}
prove that \sref{IV.16.4.8.1} is an isomorphism, since the filtrations on both sides of \sref{IV.16.4.7.3} are finite (Bourbaki, \emph{Alg. comm.}, chap.~III, \textsection2, no 8, Corollary 3 to Theorem 1).
To do this, consider the split exact sequence of $\sh{O}_S$-modules
\[
  \label{IV.16.4.8.2}
  \xymatrix{
    0 \ar[r] & \sh{E} \ar[r]^-u & \sh{E} \oplus \sh{E} \ar[r]^-v & \sh{E} \ar[r] & 0
  }
  \tag{16.4.8.2}
\] 
where, for every pair of sections $s$, $t$ of $\sh{E}$ over an open $U$ of $S$, we take $u(s) = (-s, s)$ and $v(s,t) = s + t$.
We have
\[
  X \times_S X = \Spec(\bb{S}_{\sh{O}_S}(\sh{E}) \otimes_{\sh{O}_S} \bb{S}_{\sh{O}_S}(\sh{E})) = \Spec(\bb{S}_{\sh{O}_S}(\sh{E} \oplus \sh{E}))
\]
(\sref[II]{II.1.4.6} and \sref[II]{II.1.7.11}), and the diagonal morphism $X \to X \times_S X$ corresponds \sref[II]{II.1.2.7} to the homomorphism of $\sh{O}_S$-algebras $\bb{S}(v): \bb{S}_{\sh{O}_S}(\sh{E} \oplus \sh{E}) \to \bb{S}_{\sh{O}_S}(\sh{E})$ \sref[II]{II.1.7.4}; thus, if $\sh{I}$ is the kernel of this homomorphism, then
\[
  \sh{P}_{X/S}^n = f^*(\bb{S}_{\sh{O}_S}(\sh{E} \oplus \sh{E})/\sh{I}^{n+1}).
\]
The proposition now will be a consequence of the following lemma:
\end{proof}

\begin{lemma}[16.4.8.3]
\label{IV.16.4.8.3}
Let $Y$ be a ringed space, $0 \to \sh{F}'  \xrightarrow{u} \sh{F} \xrightarrow{v} \sh{F}'' \to 0$ an exact sequence of $\sh{O}_Y$-modules such that each point $y \in Y$ has an open neighbourhood $V$ such that the sequence $0 \to \sh{F}'|V \to \sh{F}|V \to \sh{F}''|V \to 0$ is split.
Let $\sh{I}$ be the kernel ideal of $\bb{S}(v)$:
\[
  \bb{S}_{\sh{O}_Y}(\sh{F}) \to \bb{S}_{\sh{O}_Y}(\sh{F}''),
\]
and let $\gr_{\sh{I}}^\bullet(\bb{S}_{\sh{O}_Y}(\sh{F}))$ be the graded $\sh{O}_Y$-algebra associated to the $\sh{O}_Y$-algebra $\bb{S}_{\sh{O}_Y}(\sh{F})$ \emph{endowed with the $\sh{I}$-preadic filtration}.
Then the homomorphism of graded $\sh{O}_Y$-algebras
\[
  \label{IV.16.4.8.4}
  \bb{S}_{\sh{O}_Y}^\bullet(\sh{F}') \otimes_{\sh{O}_Y} \bb{S}_{\sh{O}_Y}^\bullet(\sh{F}'') \to \gr_{\sh{I}}^\bullet(\bb{S}_{\sh{O}_Y}(\sh{F}))
  \tag{16.4.8.4}
\]
(where the left-hand side is the graded tensor product of symmetric $\sh{O}_Y$-algebras endowed with their canonical gradings \sref[II]{II.1.7.4} and \sref[II]{II.2.1.2}), induced by the canonical injection
\[
  \sh{F}' \to \sh{I}/\sh{I}^2 = \gr_{\sh{I}}^1(\bb{S}_{\sh{O}_Y}(\sh{F})),
\]
is bijective.
\end{lemma}

\begin{proof}
The injection $\sh{F}' \to \sh{I}$ indeed canonically gives a homomorphism of graded $\sh{O}_Y$-algebras $\bb{S}_{\sh{O}_Y}^\bullet(\sh{F}') \to \gr_\sh{I}^\bullet(\bb{S}_{\sh{O}_Y}(\sh{F}))$, and since the right-hand side is by definition a graded $\bb{S}_{\sh{O}_Y}^\bullet(\sh{F}'')$-algebra, we obtain the canonical homomorphism \sref{IV.16.4.8.4} by tensoring the preceding one with $\bb{S}_{\sh{O}_Y}^\bullet(\sh{F}'')$.
To prove the lemma we can, being a local problem, restrict to the case where $\sh{F} = \sh{F}' \oplus \sh{F}''$, $u$ and $v$ the canonical homomorphisms.
Then the graded algebra $\bb{S}_{\sh{O}_Y}^\bullet(\sh{F})$ is canonically identified with the graded tensor product $\bb{S}_{\sh{O}_Y}^\bullet(\sh{F}') \otimes_{\sh{O}_Y} \bb{S}_{\sh{O}_Y}^\bullet(\sh{F}'')$ \sref[II]{II.1.7.4}, and it is immediate that $\sh{I}$ is therefore the ideal $\sh{I}' \otimes_{\sh{O}_Y} \bb{S}_{\sh{O}_Y}^\bullet(\sh{F}'')$, where $\sh{I}'$ is the augmentation ideal of $\bb{S}_{\sh{O}_Y}^\bullet(\sh{F}')$, that is to say the (direct) sum of the $\bb{S}_{\sh{O}_Y}^m(\sh{F}')$ for $m \geq 1$.
We conclude that $\sh{I}^n = \sh{I}'^n \otimes_{\sh{O}_Y} \bb{S}_{\sh{O}_Y}^\bullet(\sh{F}'')$, where this time $\sh{I}'^n$ is the direct sum of the $\bb{S}_{\sh{O}_Y}^m(\sh{F}')$ for $m \geq n$;
we therefore have $\sh{I}^n/\sh{I}^{n+1} = \bb{S}_{\sh{O}_Y}^n(\sh{F}') \otimes_{\sh{O}_Y} \bb{S}_{\sh{O}_Y}^\bullet(\sh{F}'')$, which proves that \sref{IV.16.4.8.4} is bijective.
\end{proof}

\oldpage[IV-4]{21}
Having proved the lemma, it remains to see that the homomorphism \sref{IV.16.4.8.1} is indeed the image under $f^*$ of the homomorphism \sref{IV.16.4.8.4} corresponding to the exact sequence \sref{IV.16.4.8.2};
this follows readily from the definitions of $u$ \sref{IV.16.4.8.2} and $\delta$ \sref{IV.16.4.7.1}, together with the definitions of the $\sh{O}_X$-algebra structure on $\sh{P}_{X/S}^n$ and of $d_{X/S}^n$ \sref{IV.16.3.5} and \sref{IV.16.3.6}.

In particular:

\begin{corollary}[16.4.9]
\label{IV.16.4.9}
Under the conditions of \sref{IV.16.4.7}, we have a canonical isomorphism
\[
\label{IV.16.4.9.1}
  \gr_1(\delta): f^*(\sh{E}) \isoto \Omega_{X/S}^1.
  \tag{16.4.9.1}
\]
\end{corollary}

\begin{corollary}[16.4.10]
\label{IV.16.4.10}
If $S = \Spec(A)$, $\sh{E} = \sh{O}_S^m$, so that
\[
  X = \Spec(A[T_1, \dots, T_m]),
\]
then $\sh{P}_{X/S}^n$ is canonically identified with the $\sh{O}_X$-algebra corresponding to the quotient $A[T_1, \dots, T_m]$-algebra\\
$A[T_1, \dots, T_m, U_1, \dots, U_m]/\mathfrak{K}^{n+1}$, where the $U_i$ ($1 \leq i \leq m$) are $m$ new indeterminates and $\mathfrak{K}$ is the ideal generated by $U_1, \dots, U_m$.
\end{corollary}

We thus recover in particular the structure of $\Omega_{X/S}^1$ in this case \sref[0]{0.20.5.13}.

Moreover, $d_{X/S}^n$ then sends a polynomial $F(T_1, \dots, T_m)$ to the class modulo $\mathfrak{K}^{n+1}$ of $F(T_1 + U_1, \dots, T_m + U_m)$, as follows from the definition \sref{IV.16.4.7.1}.

\begin{proposition}[16.4.11]
\label{IV.16.4.11}  
Let $f:X \to S$ be a morphism, let $g: S \to X$ be an $S$-section of $X$, and let $S_g^{(n)}$ be the $n$-th infinitesimal neighbourhood of $S$ with respect to the immersion $g$ \sref{IV.16.1.2}.
Then there exists a unique isomorphism of $\sh{O}_S$-algebras
\[
  \label{IV.16.4.11.1}
  \varpi_n: g^*(\sh{P}_{X/S}^n) \to \sh{O}_{S_g^{(n)}}
  \tag{16.4.11.1}
\]  
(via the $\sh{O}_S$-algebra structure on $\sh{O}_{S_g^{(n)}}$ defined by $f$ \sref{IV.16.1.7}), making the diagram
\[
  \label{IV.16.4.11.2}
  \xymatrix{
    \sh{O}_S = g^*(\sh{O}_X) \ar[rr]^-{\lambda_n} \ar[dr]_-{g^*(d_{X/S}^n)} && \sh{O}_{S_g^{(n)}} \\
    & g^*(\sh{P}_{X/S}^n) \ar[ur]_-{\varpi_n}
  }
  \tag{16.4.11.2}
\]
commutative (where $\lambda_n$ is the structure morphism).
\end{proposition} 

\begin{proof}
In light of \sref[I]{I.5.3.7}, where we replace $X$, $Y$, $S$ by $S$, $X$, $S$ respectively and $f$ by $g$, the diagrams
\[
  \label{IV.16.4.11.3}
  \xymatrix@=3pc{
    S \ar[r]^-g \ar[d]_-g &  X \ar[d]^-{\Delta_f} & S \ar[r]^-g  \ar[d]_-g & X \ar[d]^-{\Delta_f}  \\
    X \ar[r]_{(g\circ f, 1_X)_S} & X\times_S X & X \ar[r]_{(1_X, g\circ f)_S} & X \times_S X
  }
  \tag{16.4.11.3}
\]
\oldpage[IV-4]{22}
identify $S$ with the product of the $(X \times_S X)$-preschemes $X$ and $X$ via the morphisms $\Delta_f$ and $(g\circ f, 1_X)_S$ (resp. $(1_X, g\circ f)_S$).
On the other hand, the diagrams
\[
  \label{IV.16.4.11.4}
  \xymatrix@=3pc{
    X \ar[r]^{(g\circ f, 1_X)_S} \ar[d]_-f &  X \ar[d]^-{p_1} &  X \ar[r]^{(1_X, g\circ f)_S} \ar[d]_-f &  X \ar[d]^-{p_2} \\
    S \ar[r]_{g} & X & S \ar[r]_{g} & X
  }
  \tag{16.4.11.4}
\]
identify $X$ with the product of the $X$-preschemes $S$ and $X \times_S X$ via the morphisms $g$ and $p_1$ (resp. $p_2$) (a special case of the associativity formula \sref[I]{I.3.3.9.1}).
We can say that $\Delta_f$, considered as an $X$-section of $X \times_S X$ (relative to $p_1$ or $p_2$) plays the role of a \emph{universal section} for the $S$-sections of $X$:
each such section $g$ is indeed obtained from it by the \emph{base change} $(g \circ f, 1_X)_S: X \to X \times_S X$.
The definition of the homomorphism $\varpi_n$ and the fact that it is bijective therefore follow from these remarks and from \sref{IV.16.2.3}[(ii)] applied to the first diagram \sref{IV.16.4.11.4}.
The commutativity of diagram \sref{IV.16.4.11.2} likewise follows from \sref{IV.16.2.3}[(ii)], this time applied to the second diagram \sref{IV.16.4.11.4}.
To explain $\varpi_n$, we can restrict ourselves to the case where $g$ is a closed immersion:
Indeed, for every $s \in S$, there is an open neighbourhood $W$ of $s$ in $S$ such that $g(W)$ is closed in an open set $U$ of $X$, and it is clear that $g|W$ is a $W$-section of the morphism $U \cap f^{-1}(W) \to W$, the restriction of $f$, and that $g(W)$ is \emph{a fortiori} closed in $U \cap f^{-1}(W)$.
We can then suppose that $S$ is a closed subprescheme of $X$ defined by a quasi-coherent ideal $\sh{K}$.
Then the preceding definitions show that if $W$ is an open of $S$, $t$ is a section of $\sh{O}_X$ over $f^{-1}(W)$, $\varpi_n(d^n t|W)$ is equal to the canonical image of $t$ in $\Gamma(W, (\sh{O}_X/\sh{K}^{n+1})|W)$. 
The uniqueness of $\varpi_n$ then follows since the image of $\sh{O}_X$ under $d_{X/S}^n$ generates the $\sh{O}_X$-module $\sh{P}_{X/S}^n$ \sref{IV.16.3.8}.
\end{proof}

\begin{corollary}[16.4.12]
\label{IV.16.4.12}
Let $k$ be a field, $X$ a $k$-prescheme, $x$ a point of $X$ \emph{rational over $k$}.
Then $(\sh{P}_{X/k}^n)_x \otimes_{\sh{O}_x} \kres(x)$ is canonically isomorphic (as an augmented $\kres(x)$-algebra) to $\sh{O}_x/\mathfrak{m}_x^{n+1}$.
\end{corollary}

\begin{proof}
It suffices to consider the unique $k$-section $g$ of $X$ such that $g(\Spec(k)) = \{x\}$.
\end{proof}

\begin{corollary}[16.4.13]
\label{IV.16.4.13}
Let $f: X \to S$ be a morphism, $s$ a point of $S$, $X_s = X \times_S \Spec(\kres(s))$ the fibre of $f$ in $s$.
If $x \in X_s$ is \emph{rational over $\kres(s)$}, $(\sh{P}_{X/S}^n)_x \otimes_{\sh{O}_s} \kres(s)$ is canonically isomorphic to $\sh{O}_{X_s, x}/{\mathfrak{m}'_x}^{n+1}$, where $\mathfrak{m}'_x$ is the maximal ideal of $\sh{O}_{X_s, x}$;
more precisely, this isomorphism sends $(d^n t)_x \otimes 1$ \emph{(where $t$ is a section of $\sh{O}_X$ over an open neighbourhood of $x$ in $X$)} to the class of $t_x \otimes 1$ modulo ${\mathfrak{m}'_x}^{n+1}$.
\end{corollary}

\begin{proof}
This follows from \sref{IV.16.4.5} and \sref{IV.16.4.12}.
\end{proof}

The preceding corollaries justify the terminology ``sheaf of principal parts of order $n$''.

\begin{proposition}[16.4.14]
\label{IV.16.4.14}
Let $\rho:A \to B$ be a homomorphism of rings, and let $S$ be a multiplicative subset of $B$.
Then the canonical homomorphisms 
\[
  \label{IV.16.4.14.1}
  S^{-1}P_{B/A}^n \to P_{S^{-1}B/A}^n
  \tag{16.4.14.1}
\] 
\oldpage[IV-4]{23}
deduced from the canonical homomorphisms $P_{B/A}^n \to P_{S^{-1}B/A}^n$ \sref{IV.16.4.4}, form a projective system and are bijective.
\end{proposition}

\begin{proof}
It suffices to remark that $S^{-1}((B \otimes_A B)/\mathfrak{I}^{n+1}) = S^{-1}(B \otimes_A B)/(S^{-1}\mathfrak{I})^{n+1}$ by flatness, and that $S^{-1}(B \otimes_A B) = (S^{-1}B)\otimes_A (S^{-1}B) $ \sref[I]{I.1.3.4}.
\end{proof}

\begin{corollary}[16.4.15]
\label{IV.16.4.15}
The notation being that of \sref{IV.16.4.14}, let $R$ be a multiplicative subset of $A$ such that $\rho(R) \subset S$.
Then we have canonical isomorphisms 
\[
  \label{IV.16.4.15.1}
  S^{-1}P_{B/A}^n \isoto P_{S^{-1}B/R^{-1}A}^n
  \tag{16.4.15.1}
\]
forming a projective system.
\end{corollary}

\begin{proof}
It evidently suffices to define canonical isomorphisms 
\[
  \label{IV.16.4.15.2}
  P_{S^{-1}B/A}^n \isoto P_{S^{-1}B/R^{-1}A}^n
  \tag{16.4.15.2}
\]
that is to say that we reduce to the case where $\rho(R)$ consists of \emph{invertible} elements of $B$.
But then the isomorphism \sref{IV.16.4.15.2} is simply induced by the canonical isomorphism $B \otimes_A B \to B \otimes_{R^{-1}A} B$ by passing to quotients \sref[I]{I.1.5.3}.
\end{proof}

\begin{corollary}[16.4.16]
\label{IV.16.4.16}
Let $f:X \to S$ be a morphism of preschemes, $x$ a point of $X$, $s = f(x)$.
Then we have canonical isomorphisms
\[
  \label{IV.16.4.16.1}
  (\sh{P}_{X/S}^n)_x \isoto P_{\sh{O}_x/\sh{O}_s}^n
  \tag{16.4.16.1}
\]
forming a projective system.
\end{corollary}

We deduce corresponding isomorphisms for the associated graded rings, and in particular a canonical isomorphism
\[
  \label{IV.16.4.16.2}
  (\Omega_{X/S}^1)_x \isoto \Omega_{\sh{O}_x/\sh{O}_s}^1.
  \tag{16.4.16.2}
\]

\begin{corollary}[16.4.17]
\label{IV.16.4.17}
Let $k$ be a field, $K$ the field of rational functions $k(T_1, \dots, T_r)$.
Then, for every integer $n$, the homomorphism from $K[U_1, \dots, U_r]$ ($U_i$ indeterminates) to $P_{K/k}^n$ which sends $U_i$ to $d^nT_i - T_i.1$ is surjective and induces an isomorphism from the quotient $K[U_1, \dots, U_r]/\mathfrak{m}^{n+1}$ (where $\mathfrak{m}$ is the ideal generated by the $U_i$) to $P_{K/k}^n$.
\end{corollary}

\begin{proof}
This follows from \sref{IV.16.4.8}, \sref{IV.16.4.10}  and \sref{IV.16.4.14}, where we take $A = k$, $B = k[T_1, \dots, T_r]$ and $S = B \setmin \{0\}$.
\end{proof}

We thus recover the fact that the $dT_i$ form a basis of the $K$-vector space $\Omega_{K/k}^1$ \sref[0]{0.20.5.10}.

\begin{env}[16.4.18]
\label{IV.16.4.18}
Let $f:X \to Y$, $g:Y \to Z$ be two morphisms of preschemes, and consider the canonical homomorphisms of augmented $\sh{O}_X$-algebras \sref{IV.16.4.3.3}
\[
  \label{IV.16.4.18.1}
  g_{X/Y/Z}:\sh{P}_{X/Z}^n \to \sh{P}_{X/Y}^n
  \tag{16.4.18.1}
\]
\[
  \label{IV.16.4.18.2}
  f_{X/Y/Z}:f^*(\sh{P}_{Y/Z}^n) \to \sh{P}_{X/Z}^n.
  \tag{16.4.18.2}
\]
Then $g_{X/Y/Z}$ is surjective, and its kernel is the sheaf of ideals generated by the image under $f_{X/Y/Z}$ of the augmentation ideal of $f^*(\sh{P}_{Y/Z}^n)$.
\end{env}

\begin{proof}
\oldpage[IV-4]{24}
First note that $g_{X/Y/Z}$ corresponds to the case in \sref{IV.16.4.3.3} where $X' = X$, $S' = Y$ and $S = Z$, $u = 1_X$, and $f_{X/Y/Z}$ to the case where we replace $X', X, S, S'$ by $X, Y, Z, Z$ respectively and $u$, $f$ by $f$, $g$ respectively.

We have a commutative diagram \sref[I]{I.5.3.5}
\[
  \label{IV.16.4.18.3}
  \xymatrix{
    X \ar[r]^-{\Delta_f} \ar[dr]_-f & X \times_Y X \ar[d]^-p \ar[r]^-j & X \times_Z X \ar[d]^-{f \times_Z f} \\
    &   Y \ar[r]_-{\Delta_g}  & Y \times_Z Y
  }
  \tag{16.4.18.3}
\] 
where $j = (1_X, 1_X)_Z$ is an immersion, $j \circ \Delta_f = \Delta_{g \circ f}$, and $p$ is the structure morphism.
Since we can restrict ourselves to the case where $X$, $Y$ and $Z$ are affine, we can suppose that the immersions $\Delta_f$, $\Delta_g$ and $j$ are closed, so that $\sh{O}_X$ and $\sh{O}_{X \times_Y X}$ are identified respectively with $\sh{O}_{X \times_Z X} / \sh{I}$ and $\sh{O}_{X \times_Z X}/\sh{L}$, where $\sh{I} \supset \sh{L}$ are two quasi-coherent ideals corresponding respectively to the immersions $\Delta_{g \circ f} $ and $j$.
The $\sh{O}_X$-algebra $\sh{P}_{X/Z}^n$ is identified with $\sh{O}_{X \times_Z X} / \sh{I}^{n+1}$, and $\sh{P}_{X/Y}^n$ is identified with $\sh{O}_{X \times_Y X}/(\sh{I}/\sh{L})^{n+1}$, which is to say with $\sh{O}_{X \times_Z X}/(\sh{I}^{n+1} + \sh{L})$, and therefore with the quotient of $\sh{P}_{X/Z}^n$ by $(\sh{I}^{n+1} + \sh{L})/\sh{I}^{n+1}$.
But we know (\emph{loc. cit}) that if $p$ and $j$ make $X \times_Y X$ the product of the $(Y \times_Z Y)$-preschemes $Y$ and $X \times_Z X$, so if $\sh{O}_Y$ is identified to $\sh{O}_{Y \times_Z Y}/\sh{K}$, where $\sh{K}$ is the ideal corresponding to $\Delta_g$, $\sh{L}$ is equal to $(f\times_Z f)^*(\sh{K}).\sh{O}_{X \times_Z X}$ \sref[I]{I.4.4.5}.
Since $(\sh{I}^{n+1} + \sh{L})/\sh{I}^{n+1}$ is the ideal of $\sh{P}_{X/Z}^n$ generated by the image of $\sh{L}$, we deduce the proposition.
\end{proof}

\begin{corollary}[16.4.19]
\label{IV.16.4.19}
With the notation of \sref{IV.16.4.18}, we have an exact sequence of quasi-coherent $\sh{O}_X$-modules 
\[
  \label{IV.16.4.19.1}
  \xymatrix{
    f^*(\Omega_{Y/Z}^1) \ar[r]^-{f_{X/Y/Z}} & \Omega_{X/Z}^1 \ar[r]^-{g_{X/Y/Z}} & \Omega_{X/Y}^1 \ar[r] & 0.
  }
  \tag{16.4.19.1}
\]
\end{corollary}

When $X$, $Y$, $Z$ are affine, we recover the exact sequence \sref[0]{0.20.5.7.1}.

\begin{proposition}[16.4.20]
\label{IV.16.4.20}
Let $f:Y \to Z$ be a morphism, $j: X \to Y$ a closed immersion, $\sh{K}$ the quasi-coherent sheaf of ideals of $\sh{O}_Y$ corresponding to $j$.
It follows that $\sh{P}_{X/Y}^n = \sh{O}_X = \sh{O}_Y/\sh{K}$, the canonical homomorphism $j_{X/Y/Z}: j^*(\sh{P}_{Y/Z}^n) \to \sh{P}_{X/Z}^n$ is surjective, and its kernel is the ideal of $j^*(\sh{P}_{Y/Z}^n)$ generated by $j^*(\sh{O}_Y \cdot d_{Y/Z}^n(\sh{K}))$ (it should be noted that $d_{Y/Z}^n(\sh{K})$ is a subsheaf of abelian groups of $\sh{P}_{Y/Z}^n$, but not an $\sh{O}_Y$-module in general).
\end{proposition}

\begin{proof}
We know \sref[I]{I.5.3.8} that the diagonal $\Delta_j:X \to X \times_Y X$ is an isomorphism, from which the first assertion follows.
If $\varpi_1$ and $\varpi_2$ are the two canonical homomorphisms of algebras $\sh{O}_Y \to \sh{P}_{Y/Z}^n$ corresponding respectively to the two canonical projections $p_1$, $p_2$ of $Y \times_Z Y \to Y$, recall that by definition (\sref{IV.16.3.5} and \sref{IV.16.3.6}) $\varpi_1$ is the structure homomorphism of the $\sh{O}_Y$-algebra $\sh{P}_{Y/Z}^n$ and $\varpi_2 = d_{Y/Z}^n$.
The $\sh{O}_X$-algebra $j^*(\sh{P}_{Y/Z}^n)$ is therefore identified with $\sh{P}_{Y/Z}^n/\varpi_1(\sh{K})\sh{P}_{Y/Z}^n$ and its quotient by the ideal generated by $j^*(d_{Y/Z}^n(\sh{K}))$ to $\sh{P}_{Y/Z}^n/(\varpi_1(\sh{K}) + \varpi_2(\sh{K}))\sh{P}_{Y/Z}^n$.
Now note that we have a commutative diagram
\oldpage[IV-4]{25}
\[
  \xymatrix{
  Y \ar[d]_-{\Delta_f} & X \ar[d]^-{\Delta_{f \circ j}} \ar[l]_-{j}  \\
  Y \times_Z Y & X \times_Z X \ar[l]^-{j \times_Z j} 
  }
\]
identifying $X$ with the product of the $(Y \times_Z Y)$-preschemes $Y$ and $X \times_Z X$ \sref[I]{I.5.3.7}.
Since $j \times_Z j$ is an immersion, we therefore deduce from this remark and from \sref{IV.16.2.2} that if $\Delta_{Y/Z}^n$ and $\Delta_{X/Z}^n$ denote the infinitesimal neighbourhoods of order $n$ of $Y$ and $X$ by the canonical immersions $\Delta_f$ and $\Delta_{f\circ j}$ respectively, then we have a diagram
\[
  \xymatrix{
  \Delta_{Y/Z}^n \ar[d] & \Delta_{X/Z}^n \ar[d] \ar[l]  \\
  Y \times_Z Y & X \times_Z X \ar[l]^-{j \times_Z j} 
  }
\]
making $\Delta_{X/Z}^n$ the product of the $(Y \times_Z Y)$-preschemes $\Delta_{Y/Z}^n$ and $X \times_Z X$.
We can also say that $\sh{P}_{X/Z}^n$ is identified with the sheaf of rings $\sh{P}_{Y/Z}^n \otimes_{\sh{O}_{Y \times_Z Y}} \sh{O}_{X \times_Z X}$.
But we see immediately (for example, by restricting to the affine case) that $\sh{O}_{X \times_Z X} = \sh{O}_{Y \times_Z Y}/(p_1^*(\sh{K}) + p_2^*(\sh{K}))\sh{O}_{Y \times_Z Y}$.
Therefore $\sh{P}_{X/Z}^n$ is identified with the quotient of $\sh{P}_{Y/Z}^n$ by the ideal generated by the image in $\sh{P}_{Y/Z}^n$ of $p_1^*(\sh{K}) + p_2^*(\sh{K})$.
But by definition this ideal is generated by $\varpi_1(\sh{K}) + \varpi_2(\sh{K})$. 
\end{proof}

\begin{corollary}[16.4.21]
\label{IV.16.4.21}
Let $f:Y \to Z$ be a morphism, $j: X \to Y$ an immersion.
We have an exact sequence of quasi-coherent $\sh{O}_X$-modules
\[
  \label{IV.16.4.21.1}
  \xymatrix{
    \sh{N}_{X/Y} \ar[r] & j^*(\Omega_{Y/Z}^1) \ar[r] & \Omega_{X/Z}^1 \ar[r] & 0.
  }
  \tag{16.4.21.1}
\]
\end{corollary}

When $X$, $Y$, $Z$ are affine, we recover the exact sequence \sref[0]{0.20.5.12.1}.

\begin{corollary}[16.4.22]
\label{IV.16.4.22}
If $f: X \to S$ is a morphism locally of finite presentation, $\sh{P}_{X/S}^n$ and $\Omega_{X/S}^1$ are quasi-coherent $\sh{O}_X$-modules of finite presentation.
\end{corollary}

\begin{proof}
We immediately reduce to the case where $S = \Spec(A)$ is affine, $X = \Spec(B)$, where $B = A[T_1, \dots, T_r]/\mathfrak{K}$, $\mathfrak{K}$ being an ideal of finite type of $C = A[T_1, \dots, T_r]$.
We then apply \sref{IV.16.4.20} with $Z = S$, $Y = \Spec(C)$ and $\sh{K} = \widetilde{\mathfrak{K}}$.
Then $j^*(\sh{P}_{Y/Z}^n)$ is a free $\sh{O}_X$-module of finite rank \sref{IV.16.4.10} and the hypothesis on $\mathfrak{K}$ implies that $j^*(\sh{O}_Y.d_{Y/Z}^n(\sh{K}))$ generates a quasi-coherent $\sh{O}_X$-module of finite type;
hence the conclusion.
\end{proof}

\begin{proposition}[16.4.23]
\label{IV.16.4.23}
Let $X$, $Y$ be two $S$-preschemes, $Z = X \times_S Y$ their product, $p:X \times_S Y \to X$ and $q:X \times_S Y \to Y$ the canonical projections.
Then the canonical homomorphism
\[
  \label{IV.16.4.23.1}
  p_{Z/X/S} \oplus q_{Z/Y/S}: p^*(\Omega_{X/S}^1) \oplus  q^*(\Omega_{Y/S}^1) \to \Omega_{(X\times_S Y)/S}^1
  \tag{16.4.23.1}
\]
is bijective.
\end{proposition}   

\begin{proof}
\oldpage[IV-4]{26}
The commutative diagram
\[
  \xymatrix{
    Y \ar[d]_-g & X \times_S Y \ar[l]_-q \ar[d]_-h & X \times_S Y \ar[l]_-{\operatorname{id}} \ar[d]^-p\\
    S & S \ar[l]^-{\operatorname{id}} & X \ar[l]^-f
  }
\]
gives us a factorization of the canonical \emph{isomorphism} $P^n(p)$ \sref{IV.16.4.5}
\[
  p^*(\sh{P}_{X/S}^n) \to \sh{P}_{Z/S}^n \to \sh{P}_{Z/Y}^n
\]
and similarly, switching $X$ with $Y$, we have a factorization of the \emph{isomorphism} $P^n(q)$
\[
  q^*(\sh{P}_{Y/S}^n) \to \sh{P}_{Z/S}^n \to \sh{P}_{Z/X}^n.
\]
This proves that the canonical homomorphism \sref{IV.16.4.18.1}
\[
  p_{Z/X/S}:p^*(\sh{P}_{X/S}^n) \to \sh{P}_{Z/S}^n \quad \text{(resp. $q_{Z/Y/S}:q^*(\sh{P}_{Y/S}^n) \to \sh{P}_{Z/S}^n$)}
\] 
is \emph{injective}, and that the kernel of the canonical surjective homomorphism  \sref{IV.16.4.18.2}
\[
  \sh{P}_{Z/S}^n \to \sh{P}_{Z/Y}^n \quad \text{(resp. $\sh{P}_{Z/S}^n \to \sh{P}_{Z/X}^n$)}
\] 
is complementary to the image of $p_{Z/X/S}$ (resp. $q_{Z/Y/S}$).
On the other hand, by \sref{IV.16.4.18}, this kernel is generated by the image under $q_{Z/Y/S}$ (resp. $p_{Z/X/S}$) of the augmentation ideal of $q^*(\sh{P}_{Y/S}^n)$ (resp. $p^*(\sh{P}_{X/S}^n)$).
We conclude the proposition by considering the case $n = 1$.
\end{proof}

We immediately generalize \sref{IV.16.4.23} to the case of a product of any finite number of $S$-preschemes.

\begin{remarks}[16.4.24]
\label{IV.16.4.24}
\medskip\noindent
\begin{enumerate}
  \item[(i)] We will see \sref{IV.17.2.3} that when the morphism $f: X \to Y$ in \sref{IV.16.4.18} is \emph{smooth}, the homomorphism $f_{X/Y/Z}$ in \sref{IV.16.4.19.1} is locally \emph{left invertible} and in particular injective.
  Similarly, when the morphism $f \circ j: X \to Z$ of \sref{IV.16.4.20} is \emph{smooth}, the homomorphism on the left in \sref{IV.16.4.21.1} is locally \emph{left invertible} and \textit{a fortiori} injective \sref{IV.17.2.5}.
  In Chapter~V, we will also give a variant, in the case of modules over a prescheme, of the ``imperfection modules'' studied in \sref[0]{0.20.6}, and the exact sequences where they occur.
  \item[(ii)] Let $X$ be a topological space, $\sh{A}$ a sheaf of rings over $X$ and $\sh{B}$ an $\sh{A}$-algebra over $X$.
  Then it is clear that 
  \[
    U \mapsto P_{\Gamma(U, \sh{B})/\Gamma(U, \sh{A})}^n \quad \text{($U$ open in $X$)}
  \]
  is a presheaf of augmented $\Gamma(U, \sh{B})$-algebras, and therefore the associated sheaf $\sh{P}_{\sh{B}/\sh{A}}^n$ is an augmented $\sh{B}$-algebra.
  In the particular case where $X$ is a prescheme, $f = (\psi, \theta): X \to S$ a morphism of preschemes, it follows easily from \sref{IV.16.4.16} and from the exactness of the functor $\varinjlim$ that $\sh{P}_{X/S}^n$ is canonically isomorphic to $\sh{P}_{\sh{O}_X/\psi^*(\sh{O}_S)}^n$.
  It follows that the formalism developed in the present paragraph could be considered as a
\oldpage[IV-4]{27}
  particular case of a differential formalism for ringed spaces endowed with a sheaf of algebras over the structure sheaf.
  However, we did not start with this point of view, which is less intuitive and less convenient for applications.
  It also seems that, for various kinds of ``varieties'', the ``global'' constructions of the $\sh{P}^n$ analogous to those we have used here are also better suited for applications.
\end{enumerate}
\end{remarks}

\subsection{Relative tangent sheaves and bundles; derivations.}
\label{IV.16.5}

\begin{env}[16.5.1]
\label{IV.16.5.1}
Let $f = (\psi, \theta) : X \to S$ be a morphism of ringed spaces.
For every $\sh{O}_X$-module $\sh{F}$, an \emph{$S$-derivation} (or \emph{$(X/S)$-derivation}, or \emph{$f$-derivation}) \emph{from $\sh{O}_X$ to $\sh{F}$} is any homomorphism of \emph{sheaves of additive groups} $D : \sh{O}_X \to \sh{F}$ satisfying the following conditions:
\begin{enumerate}
  \item[(a)] for every open $V$ of $X$, and every pair of sections $(t_1, t_2)$ of $\sh{O}_X$ over $V$, we have
  \[
  \label{IV.16.5.1.1}
    D(t_1 t_2) = t_1 D(t_2) + D(t_1)t_2;
    \tag{16.5.1.1}
  \]
  \item[(b)] for every open $V$ of $X$, every section $t$ of $\sh{O}_X$ over $V$, and every section $s$ of $\sh{O}_S$ over an open $U$ of $S$ such that $V \subset f^{-1}(U)$, we have
  \[
  \label{IV.16.5.1.2}
    D((s|V)t) = (s|V)D(t).
    \tag{16.5.1.2}
  \]
\end{enumerate}
It is clear that this amounts to saying that, for all $x \in X$, the homomorphism of additive groups $D_x : \sh{O}_x \to \sh{F}_x$ is an \emph{$\sh{O}_{f(x)}$-derivation}.

Another interpretation consists of considering the $\sh{O}_X$-algebra $\sh{D}_{\sh{O}_X}(\sh{F})$ as equal to $\sh{O}_X \oplus \sh{F}$, the algebra structure being defined by the condition that for every open $V$ of $X$, the product of two sections of $\sh{O}_X$ (resp. of a section of $\sh{O}_X$ and a section of $\sh{F}$) over $V$ is defined by the ring structure of $\Gamma(V, \sh{O}_X)$ (resp. the $\Gamma(V, \sh{O}_X)$-module structure on $\Gamma(V, \sh{F})$), and the product of two sections of $\sh{F}$ over $V$ is chosen to be zero;
then $\sh{F}$ is an ideal of $\sh{D}_{\sh{O}_X}(\sh{F})$, the kernel of the canonical augmentation $\sh{D}_{\sh{O}_X}(\sh{F}) \to \sh{O}_X$, and to say that $D$ is an $S$-derivation from $\sh{O}_X$ to $\sh{F}$ means that $\iota_{\sh{O}_X} + D$ is an \emph{$\sh{O}_S$-homomorphism of algebras} from $\sh{O}_X$ to $\sh{D}_{\sh{O}_X}(\sh{F})$, which, composed with the augmentation, gives $\iota_{\sh{O}_X}$.

The $S$-derivations from $\sh{O}_X$ to $\sh{F}$ clearly form a \emph{$\Gamma(X, \sh{O}_X)$-module} $\Der_S(\sh{O}_X, \sh{F})$.

When $\sh{F} = \sh{O}_X$, an $S$-derivation of $\sh{O}_X$ to itself is simply called an \emph{$S$-derivation of $\sh{O}_X$}.
\end{env}

\begin{proposition}[16.5.2]
\label{IV.16.5.2}
Let $A$ be a ring, $B$ an $A$-algebra, $L$ a $B$-module;
let $S = \Spec(A)$, $X=\Spec(B)$, $\sh{F} = \widetilde{L}$.
Then the map $D \mapsto \Gamma(D)$ which sends every $S$-derivation $D$ from $\sh{O}_X$ to $\sh{F}$ to the map $\Gamma(D) : t \mapsto D(t)$ from $B$ to $L$, is an isomorphism of $B$-modules from $\Der_S(\sh{O}_X, \sh{F})$ to $\Der_A(B, L)$ (cf.~\sref[0]{0.20.1.2}).
\end{proposition}

\begin{proof}
This follows immediately from the given interpretation of $S$-derivations in terms
\oldpage[IV-4]{28}
of homomorphisms of algebras, analogous to the interpretation given in \sref[0]{0.20.1.6}, and from the canonical correspondence between homomorphisms of $\sh{O}_X$-algebras and homomorphisms of $B$-algebras (\sref[I]{I.1.3.13} and \sref[I]{I.1.3.8}).
\end{proof}

\begin{proposition}[16.5.3]
\label{IV.16.5.3}
Let $f = (\psi, \theta) : X \to S$ be a morphism of preschemes.
\begin{enumerate}
  \item[{\rm(i)}] The differential $d_{X/S} : \sh{O}_X \to \Omega_{X/S}^1$ \sref{IV.16.3.6} is an $S$-derivation.
  \item[{\rm(ii)}] For every $\sh{O}_X$-module $\sh{F}$, the map $u \mapsto u \circ d_{X/S}$ is an isomorphism of $\Gamma(X, \sh{O}_X)$-modules
  \[
  \label{IV.16.5.3.1}
    \Hom_{\sh{O}_X}(\Omega_{X/S}^1, \sh{F}) \isoto \Der_S(\sh{O}_X, \sh{F}).
    \tag{16.5.3.1}
  \]
\end{enumerate}
\end{proposition}

\begin{proof}
Assertion (i) has already been noted in \sref{IV.16.3.6}.
On the other hand, it follows immediately (in light of \sref[0]{0.20.4.8}) that $u \mapsto u\circ d_{X/S}$ is injective, by considering the restrictions of both sides to a stalk $\sh{O}_x$ and using \sref{IV.16.4.16.2}.
To see that the homomorphism \sref{IV.16.5.3.1} is surjective, consider an $S$-derivation $D : \sh{O}_X \to \sh{F}$;
for every affine open $V = \Spec(B)$ of $X$, such that $f(V)$ is contained in an affine open $U = \Spec(A)$ of $S$, $D_V : B \to \Gamma(V, \sh{F}) $ is an $A$-derivation, and therefore there exists a \emph{unique} $B$-homomorphism $u_V : \Omega_{B/A}^1 \to \Gamma(V, \sh{F})$ such that $D_V = u_V \circ d_{B/A}$ \sref[0]{0.20.4.8};
in addition, the uniqueness of $u_V$ shows immediately that for an affine open $W \subset V$ we have $u_W = u_V|W$, and therefore the $u_V$ define a homomorphism of $\sh{O}_X$-modules $u : \Omega_{X/S}^1 \to \sh{F}$ answering the question.
\end{proof}

\begin{env}[16.5.4]
\label{IV.16.5.4}
With the notation of \sref{IV.16.5.1}, for every open $U$ of $X$, $\Der_S(\sh{O}_U, \sh{F}|U)$ is a $\Gamma(U, \sh{O}_X)$-module and it is clear that the map $U \mapsto \Der_S(\sh{O}_U, \sh{F}|U)$ is a presheaf;
in fact, it is even a \emph{sheaf} (and therefore an $\sh{O}_X$-module), in light of the pointwise characterization of $S$-derivations, seen in \sref{IV.16.5.1}.
This $\sh{O}_X$-module is denoted by $\shDer_S(\sh{O}_X, \sh{F})$ and is called the \emph{sheaf of $S$-derivations from $\sh{O}_X$ to $\sh{F}$}, and what we have just seen is also expressed by the following corollary:
\end{env}

\begin{corollary}[16.5.5]
\label{IV.16.5.5}
For every $\sh{O}_X$-module $\sh{F}$, the homomorphism of $\sh{O}_X$-modules induced by $u \mapsto u \circ d_{X/S}$
\[
\label{IV.16.5.5.1}
  \shHom_{\sh{O}_X}(\Omega_{X/S}^1, \sh{F}) \to \shDer_S(\sh{O}_X, \sh{F})
  \tag{16.5.5.1}
\]
is bijective.
\end{corollary}

\begin{corollary}[16.5.6]
\label{IV.16.5.6}
\begin{enumerate}
  \item[{\rm(i)}] If the morphism $f : X \to S$ is locally of finite presentation and if $\sh{F}$ is a quasi-coherent $\sh{O}_X$-module, then $\shDer_S(\sh{O}_X, \sh{F})$ is a quasi-coherent $\sh{O}_X$-module.
  \item[{\rm(ii)}] If in addition $S$ is locally Noetherian and if $\sh{F}$ is coherent, then $\shDer_S(\sh{O}_X, \sh{F})$ is a coherent $\sh{O}_X$-module.
\end{enumerate}
\end{corollary}

\begin{proof}
The assertion (i) follows from the isomorphism \sref{IV.16.5.5.1}, from \sref{IV.16.4.22}, and \sref[I]{I.1.3.12};
the assertion (ii) follows from \sref[0]{0.5.3.5}.
\end{proof}

\begin{env}[16.5.7]
\label{IV.16.5.7}
We set
\[
\label{IV.16.5.7.1}
  \mathfrak{G}_{X/S} = \shHom_{\sh{O}_X}(\Omega_{X/S}^1, \sh{O}_X) = \shDer_S(\sh{O}_X, \sh{O}_X),
  \tag{16.5.7.1}
\]
and say that it is the \emph{sheaf of $S$-derivations of $\sh{O}_X$}, or even the \emph{tangent sheaf of $X$ relative to $S$}:
it is therefore the \emph{dual} of the $\sh{O}_X$-module $\Omega_{X/S}^1$.
If $f$ is locally of finite presentation,
\oldpage[IV-4]{29}
$\mathfrak{G}_{X/S}$ is a quasi-coherent $\sh{O}_X$-module;
if in addition $S$ is locally Noetherian, then $\mathfrak{G}_{X/S}$ is coherent \sref{IV.16.5.6}.
\end{env}

\begin{env}[16.5.8]
\label{IV.16.5.8}
Suppose in particular that $\Omega_{X/S}^1$ is a \emph{locally free} $\sh{O}_X$-module (of finite rank) (which will be the case when $f$ is \emph{smooth} \sref{IV.17.2.3});
then $\mathfrak{G}_{X/S}$ is a locally free $\sh{O}_X$-module of the same rank as $\Omega_{X/S}^1$ at each point.
More specifically, suppose that $\Omega_{X/S}^1$ is of rank $n$ at a point $x$;
then there are $n$ sections $s_i$ ($1 \leq i \leq n$) of $\sh{O}_X$ over an affine neighbourhood $U$ of $x$ such that the canonical images of the $ds_i$ in $\Omega_{X/S}^1 \otimes_{\sh{O}_X} \kres(x)$ form a basis of this $\kres(x)$-vector space;
by virtue of Nakayama's lemma, the germs $(ds_i)_x = d(s_i)_x$ of the $ds_i$ at the point $x$ form a basis of the $\sh{O}_x$-module $(\Omega_{X/S}^1)_x$, and therefore, by restricting $U$, we can suppose that the $ds_i$ form a \emph{basis} of the $\Gamma(U, \sh{O}_X)$-module $\Gamma(U, \Omega_{X/S}^1)$.
So the $\Gamma(U, \sh{O}_X)$-module $\Gamma(U, \mathfrak{G}_{X/S})$ is dual to the above;
we denote by $(D_i)_{1 \leq i \leq n}$ or $\left(\frac{\partial}{\partial s_i}\right)_{1 \leq i \leq n}$ the dual basis of $(ds_i)_{1 \leq i \leq n}$, so that, by \sref{IV.16.5.3}, we have
\[
\label{IV.16.5.8.1}
  D_i s_j = \langle D_i, ds_j \rangle = \left\langle \frac{\partial}{\partial s_i}, ds_j \right\rangle = \delta_{ij} \quad \text{(Kronecker's symbol)}.
  \tag{16.5.8.1}
\]
Every section of $\mathfrak{G}_{X/S}$ over $U$ is therefore written uniquely as
\[
  D = \sum_{i = 1}^n a_i D_i = \sum_{i = 1}^n a_i \frac{\partial}{\partial s_i},
\] 
where the $a_i$ ($1 \leq i \leq n$) are sections of $\sh{O}_X$ over $U$.
For every section $g \in \Gamma(U, \sh{O}_X)$, if we put $dg = \sum_{i = 1}^n c_i ds_i$, then we have $c_i = \langle D_i, dg \rangle = D_i g$ by virtue of \sref{IV.16.5.8.1}, in other words,
\[
\label{IV.16.5.8.2}
  dg = \sum_{i = 1}^n (D_i g) ds_i = \sum_{i = 1}^n \frac{\partial g}{\partial s_i} ds_i.
  \tag{16.5.8.2}
\]
\end{env}

\begin{env}[16.5.9]
\label{IV.16.5.9}
Let $D_1$, $D_2$ be two $S$-derivations of $\sh{O}_X$. 
For every open $U$ of $X$, if $D_1^U$, $D_2^U$ are the corresponding derivations of the ring $\Gamma(U, \sh{O}_X)$, the \emph{bracket}
\[
  [D_1^U, D_2^U] = D_1^U \circ D_2^U - D_2^U \circ D_1^U
\]
is also a derivation of this ring, and therefore the $\psi^*(\sh{O}_S)$-endomorphism of $\sh{O}_X$
\[
\label{IV.16.5.9.1}
  [D_1, D_2] = D_1 \circ D_2 - D_2 \circ D_1
  \tag{16.5.9.1}
\]
is also an $S$-derivation;
as we immediately check that this bracket satisfies the Jacobi identity, we have thus defined on $\Der_S(\sh{O}_X, \sh{O}_X)$ a \emph{$\Gamma(S, \sh{O}_S)$-Lie algebra} structure.
Since the definition of this structure commutes with the restriction to an open of $X$, we thus see that $\mathfrak{G}_{X/S}$ is canonically equipped with a \emph{ $\psi^*(\sh{O}_S)$-Lie algebra} structure.
Note that the mapping $(D_1, D_2) \mapsto [D_1, D_2]$ is \emph{not $\Gamma(X, \sh{O}_X)$-bilinear}.
\end{env} 

\begin{env}[16.5.10]
\label{IV.16.5.10}
For every base change $g : S' \to S$, if we set $X' = X \times_S S'$, then we have seen in \sref{IV.16.4.5} that there is a canonical isomorphism
\[
\label{IV.16.5.10.1}
  \Omega_{X/S}^1 \otimes_S S' \isoto \Omega_{X'/S'}^1,
  \tag{16.5.10.1}
\]
\oldpage[IV-4]{30}  
from which we deduce, by \sref{IV.16.5.10.1}, a canonical homomorphism (Bourbaki, \emph{Alg.}, chap.~II, 3\textsuperscript{rd}ed., \textsection5, n\textsuperscript{o}3)
\[
\label{IV.16.5.10.2}
  \mathfrak{G}_{X/S} \otimes_{\sh{O}_S} \sh{O}_{S'} \to \mathfrak{G}_{X'/S'},
  \tag{16.5.10.2}
\]
which is neither injective nor surjective in general.
However:
\end{env}

\begin{proposition}[16.5.11]
\label{IV.16.5.11}
\begin{enumerate}
  \item[{\rm(i)}] If $g : S' \to S$ is a flat morphism and if $f$ is locally of finite type (resp. locally of finite presentation), then the homomorphism \sref{IV.16.5.10.2} is injective (resp. bijective).
  \item[{\rm(ii)}] If $\Omega_{X/S}^1$ is a locally free $\sh{O}_X$-module of finite type, then the homomorphism \sref{IV.16.5.10.2} is bijective.
\end{enumerate}
\end{proposition}

\begin{proof}
The assertion (ii) follows from Bourbaki, \emph{Alg.}, chap.~II, 3\textsuperscript{rd}~ed., \textsection5, n\textsuperscript{o}3, prop.~7.
The assertion (i) follows similarly from Bourbaki, \emph{Alg. Comm.}, chap.~I, \textsection2, n\textsuperscript{o}10, prop.~11 and from the fact that if $f$ is locally of finite type (resp. locally of finite presentation), then $\Omega_{X/S}^1$ is an $\sh{O}_X$-module of finite type (resp. of finite presentation) (\sref{IV.16.3.9} and \sref{IV.16.4.22}).
\end{proof}

\begin{env}[16.5.12]
\label{IV.16.5.12}
Since $\Omega_{X/S}^1$ is a quasi-coherent $\sh{O}_X$-module, we can consider the vector bundle over $X$ defined by $\Omega_{X/S}^1$ \sref[II]{II.1.7.8}
\[
  \label{IV.16.5.12.1}
  T_{X/S} = \bb{V}(\Omega_{X/S}^1)
  \tag{16.5.12.1}
\]
which is called the \emph{tangent bundle of $X$ relative to $S$}.
We have therefore a canonical bijection \sref[II]{II.1.7.9}
\[
  \Gamma(T_{X/S}/S) \isoto \Hom_{\sh{O}_X}(\Omega_{X/S}^1, \sh{O}_X) = \Gamma(X, \mathfrak{G}_{X/S})
\]
by definition of $\mathfrak{G}_{X/S}$, and we can replace $X$ by an open set $U$ of $X$ in this isomorphism;
so we can say that the \emph{tangent sheaf} of $X$ relative to $S$ is isomorphic to the \emph{sheaf of germs of $S$-sections} of the tangent bundle of $X$ relative to $S$.
If $f:X \to Y$ is an $S$-morphism, we saw \sref{IV.16.4.19} that we have a canonical homomorphism $f_{X/Y/S}:f^*(\Omega_{Y/S}^1) \to \Omega_{X/S}^1$, which, having in mind that
\[
  \bb{V}(f^*(\Omega_{Y/S}^1)) = \bb{V}(\Omega_{Y/S}^1) \times_Y X \quad \text{\sref[II]{II.1.7.11} },
\]
gives us an $X$-morphism $T_{X/S}(f): T_{X/S} \to T_{Y/S} \times_Y X$.
If $g:Y \to Z$ is a second $S$-morphism, we have $T_{X/S}(g \circ f) = (T_{Y/S}(g) \times 1_X ) \circ T_{X/S}(f)$ \sref[0]{0.20.5.4.1}.

It follows from \sref{IV.16.5.10.1} and from \sref[II]{II.1.7.11} that for every base change $g:S' \to S$ we have a canonical isomorphism
\[
  \label{IV.16.5.12.2}
  T_{X'/S'} \isoto T_{X/S} \times_S S' = T_{X/S} \times_X X'.
  \tag{16.5.12.2}
\]
\end{env}

\begin{env}[16.5.13]
\label{IV.16.5.13}
For every point $x \in X$, we define the \emph{tangent space of $X$ at the point $x$} (relative to $S$) to be the set of points in the fibre $T_{X/S} \times_X \Spec(\kres(x))$ that are \emph{rational over $\kres(x)$}, that is, the set
\[
\label{IV.16.5.13.1}
  T_{X/S}(x) = \Hom_{\kres(x)}(\Omega_{X/S}^1 \otimes_{\sh{O}_x} \kres(x), \kres(x)),
  \tag{16.5.13.1}
\]
which is the \emph{dual} of the $\kres(x)$-vector space $\Omega_{\sh{O}_x/\sh{O}_s}^1/\mathfrak{m}_x \cdot \Omega_{\sh{O}_x/\sh{O}_s}^1$.
When $\Omega_{X/S}^1$ is an $\sh{O}_X$-module of \emph{finite type}, then $T_{X/S}(x)$ is a vector space of finite rank over $\kres(x)$, and for every base change
\oldpage[IV-4]{31}
$g : S' \to S$, and every point $x' \in X' = X \times_S S'$ over $x$, we have a canonical isomorphism
\[
  \label{IV.16.5.13.2}
  T_{X'/S'}(x') \isoto T_{X/S}(x) \otimes_{\kres(x)} \kres(x').
  \tag{16.5.13.2}
\]
If $x$ is \emph{rational over $\kres(s)$}, where $s = f(x)$ (so that $\kres(s) \to \kres(x)$ is an isomorphism), it follows from \sref{IV.16.4.13} that we have a canonical isomorphism
\[
\label{IV.16.5.13.3}
  T_{X/S}(x) = T_{X_s/\kres(s)}(x) = \Hom_{\kres(s)}(\mathfrak{m}'_x/{\mathfrak{m}'_x}^{2}, \kres(s)),
  \tag{16.5.13.3}
\]
where $\mathfrak{m}'_x$ is the maximal ideal of $\sh{O}_{X_s,x} = \sh{O}_{X, x}/\mathfrak{m}_s \sh{O}_{X,x}$.
In the case where $S$ is the spectrum of a field $k$, we recover the definition of the Zariski tangent space of a point $x \in X$ \emph{rational over $k$}, as the dual of $\mathfrak{m}_x / \mathfrak{m}_x^2$.

Let $Y$ be a second $S$-prescheme and let $g : Y \to X$ be an $S$-morphism;
then we have a canonical homomorphism of $\sh{O}_Y$-modules \sref{IV.16.4.19}
\[
\label{IV.16.5.13.4}
  g_{Y/X/S} : g^*(\Omega_{X/S}^1) \to \Omega_{Y/S}^1.
  \tag{16.5.13.4}
\]
Now note that if $y \in Y$ and $x = g(y)$, then we have
\[
  g^*(\Omega_{X/S}^1) \otimes_{\sh{O}_Y} \kres(y) = (\Omega_{X/S}^1 \otimes_{\sh{O}_X} \kres(x)) \otimes_{\kres(x)} \kres(y) 
\]
and consequently, if $\Omega_{X/S}^1$ is an $\sh{O}_X$-module \emph{of finite type}, then we can identify
\[
  \Hom_{\kres(y)}(g^*(\Omega_{X/S}^1) \otimes_{\sh{O}_Y} \kres(y), \kres(y))
\]
with $T_{X/S}(x) \otimes_{\kres(x)} \kres(y)$.
We therefore deduce from the homomorphism \sref{IV.16.5.13.4} a homomorphism of \emph{$\kres(y)$-vector spaces}
\[
\label{IV.16.5.13.5}
  T_y(g) : T_{Y/S}(y) \to T_{X/S}(x) \otimes_{\kres(x)} \kres(y)
  \tag{16.5.13.5}
\]
called the \emph{linear map tangent to $g$ at the point $y$}.
When $y$ is \emph{rational over $\kres(s)$}, we can identify $\kres(s)$, $\kres(y)$, and $\kres(x)$, and $T_y(g)$ is then a homomorphism of $\kres(s)$-vector spaces $T_{Y/S}(y) \to T_{X/S}(x)$;
also note that in this case, $g^*(\Omega_{X/S}^1) \otimes_{\sh{O}_Y} \kres(y)$ is identified with $\Omega_{X/S}^1 \otimes_{\sh{O}_X} \kres(x)$, and the above homomorphism is therefore defined without any finiteness conditions on $\Omega_{X/S}^1$ and it is none other than the homomorphism $T_{Y/S}(g)$ \sref{IV.16.5.12} restricted to the fibre of $T_{Y/S}$ at the point $y$.
\end{env}

\begin{env}[16.5.14]
\label{IV.16.5.14}
The interpretation of derivations of an $A$-algebra $B$ to a $B$-module $L$, given in \sref[0]{0.20.1.1}, translates to the language of preschemes in the following way.

Consider two morphisms of preschemes $f : X \to S$, $g : Y \to S$, and a closed subprescheme $Y_0$ of $Y$ defined by a \emph{square-zero} ideal $\sh{J}$ of $\sh{O}_Y$ (so that $Y$ and $Y_0$ have the \emph{same underlying topological space}).
Suppose we are given an $S$-morphism $u_0 : Y_0 \to X$, so that we have a commutative diagram
\[
\label{IV.16.5.14.1}
  \xymatrix{
    X \ar[d]_-f & Y_0 \ar[l]_-{u_0}\ar[d]^j \\
    S & Y \ar@{-->}[ul]_-u \ar[l]^-g
  }
  \tag{16.5.14.1}
\]
\oldpage[IV-4]{32}
and we seek \emph{$S$-morphisms} $u : Y \to X$ such that $u_0 = u \circ j$ (in other words, we ask whether the preceding diagram can be completed by the dotted arrow $u$ while remaining \emph{commutative}).

For that, consider an affine open $U = \Spec(C)$ of $Y$;
its inverse image $j^{-1}(U)$ is the affine open $U_0 = \Spec(C/\mathfrak{L})$, where $\mathfrak{L} = \Gamma(U, \sh{J})$, a \emph{square-zero} ideal in $C$;
suppose that $U$ is small enough so that $u_0(U_0)$ is contained in an affine open $V = \Spec(B)$ of $X$ and that $g(U) = f(u_0(U_0))$ is contained in an affine open $W = \Spec(A)$ of $S$, so that $B$ and $C$ are $A$-algebras and $u_0|U_0$ corresponds to an $A$-homomorphism $\psi$ from $B$ to $C/\mathfrak{L}$;
Let $P(U_0)$ be the set of restrictions $u|U$ of the sought morphisms; it corresponds canonically to the $A$-algebra homomorphisms $\vphi : B \to C$ such that the \emph{composite $B \xrightarrow{\vphi} C \to C/\mathfrak{L}$ is equal to $\psi$}.
So we know \sref[0]{0.20.1.1} that the set of such homomorphisms is either empty or of the form $\vphi_1 + \Der_A(B, \mathfrak{L})$;
when $P(U_0)$ is not empty, the additive group $\Der_A(B, \mathfrak{L})$ acts by addition on $P(U_0)$, which is therefore an \emph{affine space} for the additive group $\Der_A(B, \mathfrak{L})$ (or even a \emph{principal homogeneous space} (or \emph{torsor}) \emph{under} $\Der_A(B, \mathfrak{L})$).

Now notice that, since $\mathfrak{L}$ is equipped with a $B$-module structure via $\psi$, we have an \emph{isomorphism} $v \mapsto v \circ d_{B/A}$ of $\Hom_B(\Omega_{B/A}^1, \mathfrak{L})$ onto $\Der_A(B, \mathfrak{L})$ \sref[0]{0.20.4.8}.
Besides, as $\mathfrak{L}$ is square-zero, therefore a $(C/\mathfrak{L})$-module, every $B$-homomorphism $v : \Omega_{B/A}^1 \to \mathfrak{L}$ can be considered as a $(C/\mathfrak{L})$-homomorphism $\Omega_{B/A}^1 \otimes_B (C/\mathfrak{L}) \to \mathfrak{L}$.
As $\sh{J}$ is square-zero, it can be considered as a quasi-coherent $\sh{O}_{Y_0}$-module;
let's introduce the $\sh{O}_{Y_0}$-module
\[
\label{IV.16.5.14.2}
  \sh{G} = \shHom(u_0^*(\Omega_{X/S}^1), \sh{J});
  \tag{16.5.14.2}
\]
it follows from the fact that $\Omega_{B/A}^1 = \Gamma(V, \Omega_{X/S}^1)$ \sref{IV.16.3.7} that we can write $\Der_A(B, \mathfrak{L}) = \Gamma(U_0, \sh{G})$.

As $P(U_0)$ is defined as a set of $S$-morphisms $U \to X$, it is clear that $U_0 \mapsto P(U_0)$ is a \emph{sheaf of sets} $\sh{P}$ on $Y_0$.
We can use this fact to prove that the map $h : \Gamma(U_0, \sh{G}) \times P(U_0) \to P(U_0)$ defining the torsor structure on $P(U_0)$ is independent of choice of $V$ and $W$ and also that, if $U' \subset U$ is a second affine open of $Y$, $U'_0$ its inverse image in $Y_0$, then the diagram
\[
\label{IV.16.5.14.3}
  \xymatrix{
    \Gamma(U_0, \sh{G}) \times P(U_0) \ar[d] \ar[r]^-h & P(U_0) \ar[d] \\
    \Gamma(U'_0, \sh{G}) \times P(U'_0) \ar[r]_-{h'} & P(U'_0)
  }
  \tag{16.5.14.3}
\]
is commutative (the vertical arrows being the restrictions).
In light of the above remark, we reduce to proving the commutativity of the above diagram when $h$ is defined as such from affine opens $V$, $W$ and $h'$ from affine
\oldpage[IV-4]{33}
opens $V' \subset V$ and $W' \subset W$.
But because of the preceding description of $h$, this follows from the commutativity of the diagram \sref[0]{0.20.5.4.2}.

The mapping $\Gamma(U_0, \sh{G}) \times P(U_0) \to P(U_0)$ therefore defines a homomorphism of \emph{sheaves of sets} $m:\sh{G} \times \sh{P} \to \sh{P}$ such that, for all open sets $U_0$ for which $\Gamma(U_0, \sh{P}) \neq \emp$, $m_{U_0}:\Gamma(U_0, \sh{G}) \times \Gamma(U_0, \sh{P}) \to \Gamma(U_0, \sh{P})$ is an external law defining in $\Gamma(U_0, \sh{P})$ a torsor structure for the group $\Gamma(U_0, \sh{G})$.
\end{env}

\begin{env}[16.5.15]
\label{IV.16.5.15}
In general, when we are given a sheaf of sets $\sh{P}$ over a topological space $Z$, a sheaf of groups $\sh{G}$ (not necessarily commutative), and a homomorphism of sheaves of sets $m : \sh{G} \times \sh{P} \to \sh{P}$ such that, for every open $U \subset Z$ such that $\Gamma(U, \sh{P}) \neq \emp$, $m_U : \Gamma(U, \sh{G}) \times \Gamma(U, \sh{P}) \to \Gamma(U, \sh{P})$ makes $\Gamma(U, \sh{P})$ a \emph{torsor} under the group $\Gamma(U,\sh{G})$, then we say that $\sh{P}$ is a \emph{pseudo-torsor} (or \emph{formally principal homogeneous sheaf}) under the sheaf of groups $\sh{G}$.
We say that $\sh{P}$ is a \emph{torsor} (or \emph{principal homogeneous sheaf}) under $\sh{G}$
\footnote{
[Trans.] This is nowadays more commonly called a $\sh{G}$-torsor rather than a torsor \emph{under} $\sh{G}$.
}
if in addition $\Gamma(U, \sh{P}) \neq \emp$ for every open $U \neq \emp$ in a suitable basis for the topology of $Z$.

For the general theory of torsors, we refer to \cite{IV-42};
we will limit ourselves to recalling the canonical correspondence between isomorphism classes of torsors (for a \emph{given} $\sh{G}$) and elements of the cohomology set $\HH^1(Z, \sh{G})$.
Consider a torsor $\sh{P}$ under $\sh{G}$ and an open cover $(U_\lambda)$ of $Z$ such that $\Gamma(U_\lambda, \sh{P}) \neq \emp$ for every $\lambda$;
denote by $p_\lambda$ an element of $\Gamma(U_\lambda, \sh{P})$.
For every pair of indices $\lambda$, $\mu$ such that $U_\lambda \cap U_\mu \neq \emp$, there then exists a unique element $\gamma_{\lambda\mu}$ of $\Gamma(U_\lambda \cap U_\mu, \sh{G})$ such that $\gamma_{\lambda\mu} \cdot (p_\mu|U_\lambda \cap U_\mu) = p_\lambda|U_\lambda \cap U_\mu$;
in addition, if $\lambda$, $\mu$, $\nu$ are three indices such that $U_\lambda \cap U_\mu \cap U_\nu \neq \emp$, then the restrictions $\gamma'_{\lambda\mu}$, $\gamma'_{\mu\nu}$, $\gamma'_{\lambda\nu}$ of $\gamma_{\lambda\mu}$, $\gamma_{\mu\nu}$, $\gamma_{\lambda\nu}$ to $U_\lambda \cap U_\mu \cap U_\nu$ satisfy the condition $\gamma'_{\lambda\nu} = \gamma'_{\lambda\mu} \gamma'_{\mu\nu}$;
in other words, $(\lambda, \mu) \mapsto \gamma_{\lambda\mu}$ is a \emph{$1$-cocycle} of the cover $(U_\lambda)$ with values in $\sh{G}$.
If, for every $\lambda$, $p'_\lambda$ is a second element of $\Gamma(U_\lambda, \sh{P})$, then there exists a unique element $\beta_\lambda \in \Gamma(U_\lambda, \sh{G})$ such that $p'_\lambda = \beta_\lambda \cdot p_\lambda$, and the $1$-cocycle $(\gamma'_{\lambda\mu})$ corresponding to the family $(p'_\lambda)$ is given by $\gamma'_{\lambda\mu} = \beta_\lambda \gamma_{\lambda\mu} \beta_\mu^{-1}$, that is, it is \emph{cohomologous} to $\gamma_{\lambda\mu}$.
Conversely, the data of a $1$-cocycle $(\gamma_{\lambda\mu})$ defines, for every pair $(\lambda, \mu)$, an automorphism $\theta_{\lambda\mu}$ of the sheaf \emph{of sets} $\sh{G}|U_\lambda \cap U_\mu$, namely the right translation by $\gamma_{\lambda\mu}$, and the fact that it is a cocycle shows that we can \emph{glue} the sheaves of sets $\sh{G}|U_\lambda$ via the automorphisms $\theta_{\lambda\mu}$ \sref[0]{0.3.3.1};
we thus obtain a torsor under $\sh{G}$, denoted $\sh{P}$, and if we take for $p_\lambda$ the unit section over $U_\lambda$, then the corresponding $1$-cocycle is none other than the given $1$-cocycle $(\gamma_{\lambda\mu})$;
in addition, if we replace $(\gamma_{\lambda\mu})$ by a $1$-cocycle $\gamma'_{\lambda\mu} = \beta_\lambda \gamma_{\lambda\mu} \beta_\mu^{-1}$ cohomologous to it, then we check immediately that the torsor obtained is isomorphic to $\sh{P}$.

In particular, if $(\gamma_{\lambda\mu})$ is a \emph{$1$-coboundary}, in other words of the form $\gamma_{\lambda\mu} = \beta_\lambda \beta_\mu^{-1}$, then the torsor $\sh{P}$ obtained is \emph{isomorphic to $\sh{G}$} (considered as a torsor under itself by left translations);
we say in this case that $\sh{P}$ is \emph{trivial}, and the converse is evident.

In particular, it follows from \sref[III]{III.1.3.1} that we have:
\end{env}

\begin{proposition}[16.5.16]
\label{IV.16.5.16}
Let $Z$ be an affine scheme, $\sh{G}$ a quasi-coherent $\sh{O}_Z$-module;
then every torsor under $\sh{G}$ is trivial.
\end{proposition}

\oldpage[IV-4]{34}
Returning to the problem considered in \sref{IV.16.5.14}, we thus obtain:

\begin{proposition}[16.5.17]
\label{IV.16.5.17}
Let $X$, $Y$ be two $S$-preschemes, $Y_0$ a closed subprescheme of $Y$ defined by a quasi-coherent ideal $\sh{J}$ of $\sh{O}_Y$ such that $\sh{J}^2 = 0$, $j : Y_0 \to Y$ the canonical injection.
Let $u_0:Y_0 \to X$ be an $S$-morphism, and $\sh{P}$ the sheaf of sets on $Y$ such that, for every open $U$ of $Y$, $\Gamma(U, \sh{P})$ is the set of $S$-morphisms $u : U \to X$ such that $u_0|U_0 = u \circ (j|U_0)$, where $U_0 = j^{-1}(U)$.
Then $\sh{P}$ has the structure of a pseudo-torsor under the $\sh{O}_{Y_0}$-module $\sh{G} = \shHom_{\sh{O}_{Y_0}}(u_0^*(\Omega_{X/S}^1), \sh{J})$.
\end{proposition}

In particular:

\begin{corollary}[16.5.18]
\label{IV.16.5.18}
With the notation of \sref{IV.16.5.17}, suppose that $Y$ is affine and $\Omega_{X/S}^1$ is of finite presentation;
if there is an open cover $(U_\alpha)$ of $Y$, and, for every index $\alpha$, an $S$-morphism $v_\alpha : U_\alpha \to X$ such that, if $U_\alpha^0 = j^{-1}(U_\alpha)$, we have $v_\alpha \circ (j|U_\alpha^0) = u_0|U_\alpha^0$, then there is an $S$-morphism $u : Y \to X$ such that $u \circ j = u_0$.
\end{corollary}

\begin{proof}
Indeed, $\sh{G}$ is a quasi-coherent $\sh{O}_{Y_0}$-module \sref[I]{I.1.3.12};
by \sref{IV.16.5.16} and the fact that $Y_0$ is then affine, the sheaf $\sh{P}$, which is by hypothesis a torsor under $\sh{G}$, and not only a pseudo-torsor, is \emph{trivial};
if $w$ is an isomorphism from $\sh{G}$ to $\sh{P}$ (as torsors under $\sh{G}$), the image under $w$ of the zero section of $\sh{G}$ is the $S$-morphism sought.
\end{proof}

\subsection{Sheaf of $p$-differentials and exterior differentials.}
\label{IV.16.6}

\begin{env}[16.6.1]
\label{IV.16.6.1}
Let $f:X \to S$ be a morphism of preschemes.
We define the \emph{sheaf of $p$-differentials of $X$ relative to $S$} ($p$ an integer) to be the \emph{$p$th exterior power} \sref[0]{0.4.1.5} of the $\sh{O}_X$-module $\Omega_{X/S}^1$, denoted by
\[
  \label{IV.16.6.1.1}
  \Omega_{X/S}^p = \bigwedge^p(\Omega_{X/S}^1).
  \tag{16.6.1.1}
\] 
So we have $\Omega_{X/S}^0 = \sh{O}_X$, and $\Omega_{X/S}^p = 0$ for $p < 0$;
the $\Omega_{X/S}^p$ are the homogeneous components of the exterior algebra of $\Omega_{X/S}^1$
\[  
  \label{IV.16.6.1.2}
  \Omega_{X/S}^\bullet = \bigwedge(\Omega_{X/S}^1) = \bigoplus_{p \in \bb{Z}}\bigwedge^p(\Omega_{X/S}^1),  
  \tag{16.6.1.2}
\] 
which is therefore a quasi-coherent graded anti-commutative $\sh{O}_X$-algebra whose elements of degree $1$ are square-zero.
For every affine $U$ of $X$, we have $\Gamma(U, \Omega_{X/S}^\bullet) = \bigwedge (\Gamma(U, \Omega_{X/S}^1))$, where $\Gamma(U, \Omega_{X/S}^1)$ is considered as a $\Gamma(U, \sh{O}_X)$-module.

When $S = \Spec(A)$ and $X = \Spec(B)$ are affine, so that $B$ is an $A$-algebra, we have \sref[0]{0.4.1.5} $\Omega_{X/S}^p = (\Omega_{B/A}^p)^\sim$, where $\Omega_{B/A}^p = \bigwedge^p\Omega^1_{B/A}$. 
\end{env}

\begin{theorem}[16.6.2]
\label{IV.16.6.2}
There is one and only one endomorphism $d$ of the sheaf of additive groups $\Omega_{X/S}^\bullet$ with the following properties:
\begin{enumerate}
  \item[\rm{(i)}] $d \circ d = 0$.
  \item[\rm{(ii)}] For every open set $U$ of $X$ and every section $f \in \Gamma(U, \sh{O}_X)$ we have $df = d_{X/S}f$. 
\oldpage[IV-4]{35}
  \item[\rm{(iii)}] For every open set $U$ of $X$, every pair of integers $p$, $q$ and every pair of sections $\omega'_p \in \Gamma(U, \Omega_{X/S}^p)$, $\omega''_q \in \Gamma(U, \Omega_{X/S}^q)$, we have
  \[
    \label{IV.16.6.2.1}
    d(\omega'_p \wedge \omega''_q) = (d\omega'_p) \wedge \omega''_q + (-1)^p \omega'_p \wedge d\omega''_q.
    \tag{16.6.2.1}
  \]
\end{enumerate}
Also, $d$ is an endomorphism of graded $\psi^*(\sh{O}_S)$-modules of degree $+1$.
\end{theorem}

\begin{proof}
Suppose that we have proved the existence of an endomorphism $d$.
For every affine open $U$ of $X$, every section of $\Omega_{X/S}^p$ over $U$ is (by \rm{(ii)}) a linear combination of finitely many elements of the form $g(df_1 \wedge df_2 \wedge \dots \wedge df_p)$, where $g$ and the $f_i$ are sections of $\sh{O}_X$ over $U$ \sref[0]{0.20.4.7}.
The conditions \rm{(i)} and \rm{(iii)} then show, by induction on $p$, that we necessarily have
\[
  \label{IV.16.6.2.2}
  d(g(df_1 \wedge df_2 \wedge \dots \wedge df_p)) = dg \wedge df_1 \wedge df_2 \wedge \dots \wedge df_p.
  \tag{16.6.2.2}
\]
This therefore proves the \emph{uniqueness} of $d$ and the last claim of the theorem.
By virtue of this uniqueness property, to show the existence of $d$, we can restrict ourselves to the case where $S = \Spec(A)$ and $X = \Spec(B)$ are affines.
Now (Bourbaki, \emph{Alg.}, chap.~III, 3\textsuperscript{rd} ed., \textsection10), to define an $A$-antiderivation $D$ of degree $+1$ on an exterior algebra $\bigwedge(M)$ (where $M$ is a $B$-module and $B$ an $A$-algebra), with values in a \emph{graded anticommutative} $A$-algebra $C = \bigoplus_{n=0}^\infty C_n$ whose elements of degree $1$ are square-zero, it suffices to choose an $A$-derivation $D_0$ from $B$ to $C_1$ and an $A$-homomorphism $D_1$ from $M$ to $C_2$ satisfying $D_1(bm)=D_0(b)\wedge m+bD_1(m)$ for $b\in B$ and $m\in M$, with their structural images in $C$ understood;
there then exists a unique $A$-antiderivation $D$ from $\bigwedge(M)$ to $C$ agreeing with $D_0$ on $B$ and with $D_1$ on $M$. 

In the present case, $D_0$ is necessarily equal to $d_{B/A}$ by \rm{(ii)};
we reduce to seeing, having \sref{IV.16.6.2.2} in mind, that there is an $A$-homomorphism $u$ of $\Omega_{B/A}^1$ in $\Omega_{B/A}^2$ such that 
\[
  \label{IV.16.6.2.3}
  u(g.df) = dg \wedge df
  \tag{16.6.2.3}
\]
for all $f$, $g$ in $B$;
it suffices to show that there is an $A$-homomorphism $v:B \otimes_A \Omega_{B/A}^1 \to \Omega_{B/A}^2$ such that
\[
  \label{IV.16.6.2.4}
  v(g.\omega) = dg \wedge \omega
  \tag{16.6.2.4}
\]
for $g \in B$ and $\omega \in \Omega_{B/A}^1$.
Finally, since $\Omega_{B/A}^1 = \mathfrak{I}/\mathfrak{I}^2$ (where $\mathfrak{I} = \mathfrak{I}_{B/A}$ is the kernel of the canonical homomorphism $B \otimes_A B \to B$) and $\Omega_{B/A}^1$ is generated by elements of the form $g.df$, it is enough to define an $A$-homomorphism $w: B \otimes_A (B \otimes_A B) \to \Omega_{B/A}^2$ such that
\[
  \label{IV.16.6.2.5}
  w(g' \otimes g \otimes f) = dg' \wedge (g.df)
  \tag{16.6.2.5}
\]
and such that $w$ is zero on the image of $B \otimes_A \mathfrak{I}^2$.
Since the right-hand side of \sref{IV.16.6.2.5} is $A$-trilinear in $g'$, $g$, and $f$, the existence of $w$ satisfying \sref{IV.16.6.2.5} is immediate.
Since, on the other hand, $\mathfrak{I}$ is generated by elements of the form $1 \otimes x - x \otimes 1$ ($x \in B$), we reduce to checking that when $z = (1 \otimes x - x \otimes 1)(1 \otimes y - y \otimes 1)$ we have $w(g'\otimes z) = 0$. 
Now, since $z = 1 \otimes (xy) + (xy) \otimes 1 - x \otimes y - y \otimes x$, formula \sref{IV.16.6.2.5} shows that it is 
\oldpage[IV-4]{36}
enough to see that we have $d(xy) - x.dy - y.dx = 0$, which is to say that $d$ is a derivation.

It remains to be shown that $d$ satisfies the condition \rm{(i)}.
Now, the square of an antiderivation is a derivation (Bourbaki, \emph{loc. cit.}), and since $\Omega_{B/A}^\bullet$ is generated by $\Omega_{B/A}^1$ as a $B$-algebra, it is enough to verify that $d(dz) = 0$ when $z \in B$ and when $z \in \Omega_{B/A}^1$;
in the first case, this follows from the formula \sref{IV.16.6.2.3} when $g = 1$;
for the second, we can restrict ourselves to the case where $z = g.df$ with $f$, $g$ in $B$, and then we have, because of \sref{IV.16.6.2.1} and \sref{IV.16.6.2.3}, 
\[
  d(d(g.df)) = d(dg \wedge df) = (d(dg)) \wedge (df) - (dg) \wedge (d(df)) = 0.
\] 
\end{proof}

\begin{definition}[16.6.3]
\label{IV.16.6.3}
The antiderivation $d$ defined in \sref{IV.16.6.2} (also denoted by $d_{X/S}$) is called the \emph{exterior differential} on $X$ (relative to $S$). 
\end{definition}

\begin{proposition}[16.6.4]
\label{IV.16.6.4}
For every base change $g:S' \to S$, if we put $X' = X \times_S S'$, the canonical morphism
\[
  \label{IV.16.6.4.1}
  \Omega_{X/S}^\bullet \otimes_S S' \to \Omega_{X'/S'}^\bullet
  \tag{16.6.4.1}
\]
deduced from the isomorphism \sref{IV.16.5.10.1} is bijective.
Also, if $s$ is a section of $\Omega_{X/S}^\bullet$ over an open set $U$ of $X$, $s \otimes 1$ its inverse image, section of $\Omega_{X'/S'}^\bullet$ over the inverse image $U'$ of $U$ in $X'$, we have $d_{X'/S'}(s \otimes 1) = d_{X/S}(s) \otimes 1$.
\end{proposition}

\begin{proof}
The first claim is immediate, the formation of the exterior algebra of a module commutes with extending the scalar ring.
To prove the second, we can, because of \sref{IV.16.6.2.2}, restrict ourselves to the case where $s \in \Gamma(U, \sh{O}_X)$, and in this case the claim has already been proven \sref{IV.16.4.3.7}.
\end{proof}

\begin{env}[16.6.5]
\label{IV.16.6.5}
Suppose that $\Omega_{X/S}^1$ is a locally free $\sh{O}_X$-module of rank $n$ at a point $x$, so that we have $n$ sections $s_i \in \Gamma(U, \sh{O}_X)$ such that the $ds_i$ form a basis for the $\Gamma(U, \sh{O}_X)$-module $\Gamma(U, \Omega_{X/S}^1)$ \sref{IV.16.5.8}.
Then, for every integer $p \geq 1$, the $p$-differentials $ds_{i_1} \wedge ds_{i_2} \wedge \dots \wedge ds_{i_p}$ (for $i_1 < i_2 < \dots < i_p$ in $[1,n]$) form a basis of $\binom{n}{p}$ elements of $\Gamma(U, \Omega_{X/S}^p)$ over $\Gamma(U, \sh{O}_X)$.
Also the formula \sref{IV.16.6.2.2} shows that for every section $g \in \Gamma(U, \sh{O}_X)$, we have
\[
  \label{IV.16.6.5.1}
  d(g.ds_{i_1} \wedge ds_{i_2} \wedge \dots \wedge ds_{i_p}) = \sum_k (-1)^r \frac{\partial g}{\partial s_k} ds_{i_1} \wedge \dots \wedge ds_{i_r} \wedge ds_k \wedge ds_{i_{r+1}} \wedge \dots \wedge ds_{i_p}
  \tag{16.6.5.1}
\]
where, on the right-hand side, $k$ ranges over the $n-p$ indices distinct from the $i_h$, and $r$ is the number of indices $h$ such that $i_h<k$ ($0\leq r\leq p$); the initial wedge through $ds_{i_r}$ is empty when $r=0$, and the final wedge from $ds_{i_{r+1}}$ is empty when $r=p$.

We note that the relation $d(dg) = 0$ for every section $g \in \Gamma(U, \sh{O}_X)$ expresses itself in the form
\[
  D_i(D_j g) = D_j(D_i g) \quad \text{for $i \neq j$};
\]
in other words, the derivations $D_i$ defined in \sref{IV.16.5.8} commute with each other.
\end{env}

\subsection{The ~$\mathcal{P}_{X/S}^n(\mathcal{F})$.}
\label{IV.16.7}

\begin{env}[16.7.1]
\label{IV.16.7.1}
Let $f:X \to S$ be a morphism of preschemes, $\sh{F}$ an $\sh{O}_X$-module.
We denote by $X_{\Delta_f}^{(n)}$ the \emph{$n$th infinitesimal neighbourhood} of $X$ for the diagonal morphism
\oldpage[IV-4]{37}
$\Delta_f: X \to X \times_S X$, by $h_n:X_{\Delta_f}^{(n)} \to X \times_S X$ the canonical morphism \sref{IV.16.1.2}, and consider the two composite morphisms
\[
  p_1^{(n)}:X_{\Delta_f}^{(n)} \xrightarrow{h_n} X \times_S X \xrightarrow{p_1} X, \quad 
  p_2^{(n)}:X_{\Delta_f}^{(n)} \xrightarrow{h_n} X \times_S X \xrightarrow{p_2} X
\]
so that, by definition, $p_1^{(n)}$ corresponds to the homomorphism of sheaves of rings $\sh{O}_X \to \sh{P}_{X/S}^n$ which we have chosen to define the $\sh{O}_X$-algebra structure on $\sh{P}_{X/S}^n$ \sref{IV.16.3.5}, and $p_2^{(n)}$ to the homomorphism of sheaves of rings $d^n_{X/S}:\sh{O}_X \to \sh{P}_{X/S}^n$ \sref{IV.16.3.6}.
Since $X_{\Delta_f}^{(n)}$ and $X$ have the same underlying space, we can write
\[
  \label{IV.16.7.1.1}
  \sh{P}_{X/S}^n = (p_1^{(n)})_* ((p_2^{(n)})^* (\sh{O}_X)).
  \tag{16.7.1.1}
\]
More generally, we define
\[
  \label{IV.16.7.1.2}
  \sh{P}_{X/S}^n(\sh{F}) = (p_1^{(n)})_* ((p_2^{(n)})^* (\sh{F})).
  \tag{16.7.1.2}
\]
so that $\sh{P}_{X/S}^n = \sh{P}_{X/S}^n(\sh{O}_X)$;
by definition, $\sh{P}_{X/S}^n(\sh{F})$ is an $\sh{O}_X$-module.
\end{env}

\begin{env}[16.7.2]
\label{IV.16.7.2}
Returning to the definition of the inverse image of modules on ringed spaces \sref[0]{0.4.3.1}, and bearing in mind that $X_{\Delta_f}^{(n)}$ and $X$ have the same underlying space, we see that definition \sref{IV.16.7.1.2} can be written in the form
\[
  \label{IV.16.7.2.1}
  \sh{P}_{X/S}^n(\sh{F}) = \sh{P}_{X/S}^n \otimes_{\sh{O}_X} \sh{F},
  \tag{16.7.2.1}
\]
where it must be kept in mind that, in interpreting the symbol $\otimes$, $\sh{P}_{X/S}^n$ is endowed with the $\sh{O}_X$-module structure defined by \emph{the homomorphism of sheaves of rings $d_{X/S}^n:\sh{O}_X \to \sh{P}_{X/S}^n$}.
It follows immediately from such formula (or directly from \sref{IV.16.7.1.2}) that $\sh{P}_{X/S}^n(\sh{F})$ is canonically equipped with a $\sh{P}_{X/S}^n$-module structure. 
\end{env}

\begin{proposition}[16.7.3]
\label{IV.16.7.3}
\begin{enumerate}
  \item[\rm{(i)}] The functor $\sh{F} \mapsto \sh{P}_{X/S}^n(\sh{F})$ from the category of $\sh{O}_X$-modules to the category of $\sh{P}_{X/S}^n$-modules is right exact, and commutes with arbitrary inductive limits;
  it is exact when $\sh{P}_{X/S}^n$ is flat.
  \item[\rm{(ii)}] If $\sh{F}$ is a quasi-coherent $\sh{O}_X$-module (resp. of finite type, resp. of finite presentation), then $\sh{P}_{X/S}^n(\sh{F})$ is quasi-coherent (resp. of finite type, resp. of finite presentation).
\end{enumerate}
\end{proposition}

\begin{proof}
The claims from \rm{(i)} follow from \sref{IV.16.7.2.1} and the consideration of the symmetry of $\sh{P}_{X/S}^n$ \sref{IV.16.3.4}.
The claims from \rm{(ii)} follow from the right exactness of the functor $\sh{F} \mapsto \sh{P}_{X/S}^n(\sh{F})$.
\end{proof}

\begin{env}[16.7.4]
\label{IV.16.7.4}
The two $\sh{O}_X$-module structures on $\sh{P}_{X/S}^n$ define two commuting $\sh{O}_X$-module structures on $\sh{P}_{X/S}^n(\sh{F})$, and hence an $\sh{O}_X$-bimodule structure.
It is convenient to denote the structure coming from the structure homomorphism $\sh{O}_X \to \sh{P}_{X/S}^n$ (chosen in \sref{IV.16.3.5}) on the \emph{left} and the one coming from the homomorphism $d_{X/S}^n: \sh{O}_X \to \sh{P}_{X/S}^n$ on the \emph{right}.
In other words, for every open $U$ of $X$, and every triple $a \in \Gamma(U, \sh{O}_X)$, $b \in \Gamma(U, \sh{P}_{X/S}^n)$, $t \in \Gamma(U, \sh{F})$, we have by definition
\[
  \label{IV.16.7.4.1}
  a(b \otimes t) = (ab) \otimes t, \quad (b \otimes t)a = (b.d^n a) \otimes t = b \otimes (at) = (d^na).(b \otimes t).
  \tag{16.7.4.1}
\]
The $\sh{O}_X$-module structure coming from the definition \sref{IV.16.7.1.2} is therefore, under these conventions, the \emph{left} $\sh{O}_X$-module structure.
\oldpage[IV-4]{38}
If $\sh{F}$ is a quasi-coherent $\sh{O}_X$-module, then the same is true for $\sh{P}_{X/S}^n(\sh{F})$ for any one of its $\sh{O}_X$-module structures.
If also $\sh{F}$ is of finite type (resp. of finite presentation) and $f:X \to S$ is locally of finite type (resp. locally of finite presentation), $\sh{P}_{X/S}^n(\sh{F})$ is (for any one of its $\sh{O}_X$-module structures) of finite type (resp. of finite presentation), which is a consequence of \sref{IV.16.3.9} and \sref{IV.16.4.22}.
\end{env}

\begin{env}[16.7.5]
\label{IV.16.7.5}
The definition \sref{IV.16.7.2.1} entails the existence of a homomorphism of sheaves of commutative groups
\[
  \label{IV.16.7.5.1}
  d_{X/S, \sh{F}}^n: \sh{F} \to \sh{P}_{X/S}^n(\sh{F}) \quad (\text{also denoted } d_{X/S}^n)
  \tag{16.7.5.1}
\]
such that, in the notations of \sref{IV.16.7.4}, we have
\[
  \label{IV.16.7.5.2}
  d_{X/S, \sh{F}}^n(t) = 1 \otimes t
  \tag{16.7.5.2}
\]
and consequently, because of \sref{IV.16.7.4.1}
\[
  \label{IV.16.7.5.3}
  d_{X/S, \sh{F}}^n(at) = (1 \otimes t) a = (d_{X/S, \sh{F}}^n(t)).a
  \tag{16.7.5.3}
\]
\[
  \label{IV.16.7.5.4}
  d_{X/S, \sh{F}}^n(at) = (d_{X/S}^n(a)).(1 \otimes t) = (d_{X/S}^n(a))(d_{X/S, \sh{F}}^n(t)).
  \tag{16.7.5.4}
\]
Therefore it is $\sh{O}_X$-linear for the \emph{right} $\sh{O}_X$-module structure on $\sh{P}_{X/S}^n(\sh{F})$, and \emph{semilinear} (relative to the automorphism $\sigma$ \sref{IV.16.3.4}) for the \emph{left} $\sh{O}_X$-module structure.
\end{env}

\begin{proposition}[16.7.6]
\label{IV.16.7.6}
The left $\sh{O}_X$-module $\sh{P}_{X/S}^n(\sh{F})$ is generated by the image of $\sh{F}$ under the homomorphism $d_{X/S, \sh{F}}^n$.
\end{proposition}

\begin{proof}
This is an immediate consequence of \sref{IV.16.7.5.3} and of the particular case $\sh{F} = \sh{O}_X$ \sref{IV.16.3.8}.
\end{proof}

\begin{env}[16.7.7]
\label{IV.16.7.7}
The canonical homomorphisms of sheaves of rings
\[
  \vphi_{nm}:\sh{P}_{X/S}^m \to \sh{P}_{X/S}^n 
\]
for $n \leq m$ \sref{IV.16.1.2} define, because of \sref{IV.16.7.2.1}, canonical homomorphisms 
\[
  \sh{P}_{X/S}^m(\sh{F}) \to \sh{P}_{X/S}^n(\sh{F}) \quad (n \leq m)
\]
which are homomorphisms of $\sh{O}_X$-bimodules in light of \sref{IV.16.1.6} and \sref{IV.16.7.4.1};
also we have commutative diagrams
\[
  \xymatrix{
    \sh{P}_{X/S}^m(\sh{F}) \ar[rr] && \sh{P}_{X/S}^n(\sh{F}) \\
      & \sh{F} \ar[ul]^-{d_{X/S,\sh{F}}^m} \ar[ur]_-{d_{X/S,\sh{F}}^n}
  }
\]
We have therefore a projective system of $\sh{O}_X$-bimodules ($\sh{P}_{X/S}^n(\sh{F})$), and we define
\[
  \label{IV.16.7.7.1}
  \sh{P}_{X/S}^\infty(\sh{F}) = \varprojlim \sh{P}_{X/S}^n(\sh{F}).
  \tag{16.7.7.1}
\]
Also, this shows that the homomorphisms \sref{IV.16.7.5.1} form a projective system of homomorphisms, and therefore define a canonical homomorphism
\[
  \label{IV.16.7.7.2}
  d_{X/S,\sh{F}}^\infty: \sh{F} \to \sh{P}_{X/S}^\infty(\sh{F}).
  \tag{16.7.7.2}
\]
\end{env}

\oldpage[IV-4]{39}

\begin{env}[16.7.8]
\label{IV.16.7.8}
Let $\sh{F}$, $\sh{G}$ be two $\sh{O}_X$-modules;
it follows immediately from the definition \sref{IV.16.7.2.1} that we have a canonical isomorphism of $\sh{P}_{X/S}^n$-modules
\[
  \label{IV.16.7.8.1}
  \sh{P}_{X/S}^n(\sh{F} \otimes_{\sh{O}_X} \sh{G}) \isoto \sh{P}_{X/S}^n(\sh{F}) \otimes_{\sh{P}_{X/S}^n} \sh{P}_{X/S}^n(\sh{G})
  \tag{16.7.8.1}
\]
(Bourbaki, \emph{Alg.}, chap.~II, 3\textsuperscript{rd} ed., \textsection~5, no.~1, prop.~3).

We conclude in particular (or see directly from definition \sref{IV.16.7.2.1}) that if $\sh{F}$ has an $\sh{O}_X$-algebra structure (not necessarily associative), $\sh{P}_{X/S}^n(\sh{F})$ has a canonical $\sh{P}_{X/S}^n$-algebra structure;
the latter is associative (resp. commutative, resp. unital, resp. a Lie algebra) if $\sh{F}$ is so.
Moreover, the canonical homomorphisms $\sh{P}_{X/S}^m(\sh{F}) \to \sh{P}_{X/S}^n(\sh{F})$ for $n \leq m$ \sref{IV.16.7.7} are then algebra di-homomorphisms;
similarly, \sref{IV.16.7.5.1} is then an $\sh{O}_X$-algebra homomorphism when $\sh{P}_{X/S}^n(\sh{F})$ is equipped with the $\sh{O}_X$-algebra structure arising from its \emph{right} $\sh{O}_X$-module structure.

With the same notation, we also have a canonical homomorphism of $\sh{P}_{X/S}^n$-modules
\[
  \label{IV.16.7.8.2}
  \sh{P}_{X/S}^n(\shHom_{\sh{O}_X}(\sh{F}, \sh{G})) \to \shHom_{\sh{P}_{X/S}^n}(\sh{P}_{X/S}^n(\sh{F}), \sh{P}_{X/S}^n(\sh{G}))
  \tag{16.7.8.2}
\]
(Bourbaki, \emph{Alg.}, chap.~II, 3\textsuperscript{rd} ed., \textsection~5, no.~3), which is bijective when $\sh{P}_{X/S}^n$ is an $\sh{O}_X$-module \emph{locally free of finite type} (\emph{loc. cit.}, prop.~7).
\end{env}

\begin{env}[16.7.9]
\label{IV.16.7.9}
Suppose we are in the situation described in \sref{IV.16.4.1};
then from the canonical homomorphism $P^n(u)$ \sref{IV.16.4.3.3} we deduce immediately a canonical homomorphism of $\sh{O}_{X'}$-bimodules 
\[
  \label{IV.16.7.9.1}
  u^*(\sh{P}_{X/S}^n(\sh{F})) \to \sh{P}_{X'/S'}^n(u^*(\sh{F})).
  \tag{16.7.9.1}
\]
We leave it to the reader to extend to this homomorphism the properties established in \sref{IV.16.4} for the case $\sh{F} = \sh{O}_X$.
\end{env}

\begin{remark}[16.7.10]
\label{IV.16.7.10}
The definition of $\sh{P}_{X/S}^n(\sh{F})$ in the form \sref{IV.16.7.1.2} still makes sense when $\sh{F}$ is a sheaf of sets (the inverse image of a sheaf of sets by $p_2^{(n)}$ being defined in \sref[0]{0.3.7.1});
a variant of this definition allows one to define the ``jet scheme'' (relative to $S$) of an arbitrary $X$-prescheme.
\end{remark}


\subsection{Differential operators.}
\footnote{
  For a more general formalism, see Expos\'e VII of \cite{IV-42} (due to P. Gabriel).
}
\label{IV.16.8}

\begin{definition}[16.8.1]
\label{IV.16.8.1}
Let $f = (\psi, \theta): X \to S$ be a morphism of preschemes, $\sh{F}$, $\sh{G}$ two $\sh{O}_X$-modules, $n$ an integer $\geq 0$.
We say that a morphism $D: \sh{F} \to \sh{G}$ of sheaves of additive groups is a differential operator of order $\leq n$ (relative to $S$) if there is a homomorphism of $\sh{O}_X$-modules $u: \sh{P}_{X/S}^n(\sh{F}) \to \sh{G}$ (where $\sh{P}_{X/S}^n(\sh{F})$ is equipped with its left $\sh{O}_X$-module structure \sref{IV.16.7.4}) such that $D = u \circ d_{X/S, \sh{F}}^n$.
\end{definition}

It is clear, because of the existence of canonical morphisms
\[
  \sh{P}_{X/S}^m(\sh{F}) \to \sh{P}_{X/S}^n(\sh{F})
\] 
\oldpage[IV-4]{40}
for $n \leq m$ \sref{IV.16.7.7}, that a differential operator of order $\leq n$ is a differential operator of order $\leq m$ for $n \leq m$.
If $D:\sh{F} \to \sh{G}$ is a differential operator of order $\leq n$, then, for every open set $U$ of $X$, $D|U: \sh{F}|U \to \sh{G}|U$ is a differential operator of order $\leq n$.

We say that a homomorphism $D: \sh{F} \to \sh{G}$ of the sheaves of additive groups underlying $\sh{F}$ and $\sh{G}$ is a \emph{differential operator} (relative to $S$) if, for every $x \in X$, there is an open neighbourhood $U$ of $x$ and an integer $n \geq 0$ such that $D|U: \sh{F}|U \to \sh{G}|U$ is a differential operator of order $\leq n$.
The \emph{order} of a differential operator $D$ is the infimum of the integers $n$ such that $D$ is a differential operator of order $\leq n$ (and therefore $+\infty$ if there is no such integer);
this order is always finite if $X$ is \emph{quasi-compact}.
The differential operators of order $0$ are exactly the homomorphisms of $\sh{O}_X$-modules $\sh{F} \to \sh{G}$;
the operators of order $<0$ are \emph{zero} by convention.
For $n \geq 0$, a differential operator of order $\leq n$ is not in general a homomorphism of $\sh{O}_X$-modules, but is always a homomorphism of $\psi^*(\sh{O}_S)$-modules.

When $\sh{F} = \sh{O}_X$, a differential operator of order $\leq 1$ from $\sh{O}_X$ to $\sh{G}$ can be written uniquely in the form $v + D$, where $v: \sh{O}_X \to \sh{G}$ is an $\sh{O}_X$-homomorphism and $D$ is an \emph{$S$-derivation} \sref{IV.16.5.1} from $\sh{O}_X$ to $\sh{G}$; this follows from the structure of $P_{B/A}^1$ \sref[0]{0.20.4.8}.

\begin{env}[16.8.2]
\label{IV.16.8.2}
To describe in a more precise manner a differential operator of order $\leq n$, $D: \sh{F} \to \sh{G}$, it suffices, for every open set $U$ of $X$ whose image in $S$ is contained in an affine open set $V$, to characterize the homomorphism $D = D_U: \Gamma(U, \sh{F}) \to \Gamma(U, \sh{G})$.
If we put $\Gamma(V, \sh{O}_S) = A$, $\Gamma(U, \sh{O}_X) = B$, so that $B$ is an $A$-algebra, we have $\Gamma(U, \sh{P}_{X/S}^n) = (B \otimes_A B)/\mathfrak{I}^{n+1}$, where, for brevity, $\mathfrak{I} = \mathfrak{I}_{B/A}$. 
Also put $M = \Gamma(U, \sh{F})$, $N = \Gamma(U, \sh{G})$;
then the definition of $D$ means that for every pair $(U,V)$ satisfying the above, the $A$-homomorphism $D:M \to N$ factors through
\[
  M \to ((B \otimes_A B)/\mathfrak{I}^{n+1}) \otimes_B M \xrightarrow{v} N
\]
where the first arrow is the canonical homomorphism $t \mapsto 1 \otimes t$, and $v$ is a $B$-homomorphism, the $B$-module structure on its source coming from the first factor $B$ (whereas, in forming the tensor product over $B$, the $B$-module structure on $(B \otimes_A B)/\mathfrak{I}^{n+1}$ comes from the second factor $B$).
Note also that the $B$-module $((B \otimes_A B)/\mathfrak{I}^{n+1}) \otimes_B M$ is isomorphic to $(B \otimes_A M)/\mathfrak{I}^{n+1}(B \otimes_A M)$, where $B \otimes_A M$ is regarded as a $(B \otimes_A B)$-module and its $B$-module structure comes from the homomorphism $b \mapsto b \otimes 1$ from $B$ to $B \otimes_A B$.  
Let $D'$ then be the $B$-homomorphism from $B \otimes_A M$ to $N$ such that $D'(b \otimes t) = bD(t)$; the factorization condition on $D$ is equivalently that $D'$ vanish on the $B$-submodule $\mathfrak{I}^{n+1}(B \otimes_A M)$.
\end{env}

\begin{env}[16.8.3]
\label{IV.16.8.3}
It is clear that the set of differential operators of order $\leq n$ from $\sh{F}$ to $\sh{G}$ forms an additive group, denoted by $\Diff_{X/S}^n(\sh{F}, \sh{G})$;
when $\sh{F} = \sh{G} = \sh{O}_X$, we also write $\Diff_{X/S}^n$ instead of $\Diff_{X/S}^n(\sh{O}_X, \sh{O}_X)$.

We saw in \sref{IV.16.8.1} that for two open sets $U \supset V$ of $X$, there is a canonical restriction homomorphism
\[
  \Diff_{U/S}^n(\sh{F}|U, \sh{G}|U) \to \Diff_{V/S}^n(\sh{F}|V, \sh{G}|V)
\]
\oldpage[IV-4]{41}
from which we deduce that $U \mapsto \Diff_{U/S}^n(\sh{F}|U, \sh{G}|U)$ is a presheaf of additive groups;
in fact, it is actually a \emph{sheaf}, since for an open set $U$ varying in $X$, the homomorphisms $u \mapsto u \circ d_{U/S, \sh{F}|U}^n$ are \emph{isomorphisms} of additive groups
\[
  \label{IV.16.8.3.1}
  \Hom_{\sh{O}_U}(\sh{P}_{U/S}^n(\sh{F}|U), \sh{G}|U) \isoto \Diff_{U/S}^n(\sh{F}|U, \sh{G}|U),
  \tag{16.8.3.1}
\]
because of the fact that the image of $\sh{F}$ by $d_{X/S, \sh{F}}^n$ generates $\sh{P}_{X/S}^n(\sh{F})$ \sref{IV.16.7.6}.
We denote this sheaf by $\shDiff_{X/S}^n(\sh{F}, \sh{G})$, and we have:
\end{env}

\begin{proposition}[16.8.4]
\label{IV.16.8.4}
The isomorphisms \sref{IV.16.8.3.1} define an isomorphism of sheaves of additive groups
\[
  \label{IV.16.8.4.1}
  \shHom_{\sh{O}_X}(\sh{P}_{X/S}^n(\sh{F}), \sh{G}) \isoto \shDiff_{X/S}^n(\sh{F}, \sh{G}).
  \tag{16.8.4.1}
\]
\end{proposition}

When $\sh{F} = \sh{G} = \sh{O}_X$, we also write $\shDiff_{X/S}^n$ instead of $\shDiff_{X/S}^n(\sh{O}_X, \sh{O}_X)$;
it follows from \sref{IV.16.8.4} that $\shDiff_{X/S}^n$ is canonically identified with the \emph{dual} of the $\sh{O}_X$-module $\sh{P}_{X/S}^n$;
so we also write $\langle t, D \rangle$ instead of $u(t)$ if $t$ is a section of $\sh{P}_{X/S}^n$ over an open set and if $u$ is a homomorphism from $\sh{P}_{X/S}^n$ to $\sh{O}_X$ corresponding to $D$.

\begin{env}[16.8.5]
\label{IV.16.8.5}
Since $\sh{P}_{X/S}^n(\sh{F})$ has an $\sh{O}_X$-bimodule structure \sref{IV.16.7.4}, we canonically obtain an $\sh{O}_X$-bimodule structure on $\shHom_{\sh{O}_X}(\sh{P}_{X/S}^n(\sh{F}), \sh{G})$, and hence on $\shDiff_{X/S}^n(\sh{F}, \sh{G})$ by \sref{IV.16.8.4.1}.
More precisely, to the left $\sh{O}_X$-module structure on $\sh{P}_{X/S}^n(\sh{F})$ corresponds, because of \sref{IV.16.8.1}, the left $\sh{O}_X$-module structure on $\shDiff_{X/S}^n(\sh{F}, \sh{G})$ explained as follows:
for every open set $U$ of $X$, every section $a \in \Gamma(U, \sh{O}_X)$ and every differential operator $D: \sh{F}|U \to \sh{G}|U$, $aD$ is the differential operator which, for every section $t \in \Gamma(U, \sh{F})$, makes correspond the section
\[
  \label{IV.16.8.5.1}
  (aD)(t) = a(D(t))
  \tag{16.8.5.1}
\]
of $\Gamma(U, \sh{G})$.
Similarly, the right $\sh{O}_X$-module structure on $\shDiff_{X/S}^n(\sh{F}, \sh{G})$ is made explicit as follows:
under the same notations as above, $Da$ is the operator which, to every $t \in \Gamma(U, \sh{F})$, makes correspond the section
\[
  \label{IV.16.8.5.2}
  (Da)(t) = D(at).
  \tag{16.8.5.2}
\]
\end{env}

\begin{proposition}[16.8.6]
\label{IV.16.8.6}
If $f:X \to S$ is a morphism locally of finite presentation, $\sh{F}$ a quasi-coherent $\sh{O}_X$-module of finite presentation and $\sh{G}$ a quasi-coherent $\sh{O}_X$-module, then $\shDiff_{X/S}^n(\sh{F}, \sh{G})$ is a quasi-coherent $\sh{O}_X$-module for either of the structures defined in \sref{IV.16.8.5}.
\end{proposition}

\begin{proof}
The proposition follows from the fact that, under these hypotheses, $\sh{P}_{X/S}^n(\sh{F})$ is a quasi-coherent $\sh{O}_X$-module of finite presentation \sref{IV.16.7.4}, together with \sref[I]{I.1.3.12}.
\end{proof}

\begin{env}[16.8.7]
\label{IV.16.8.7}
The set of differential operators (of unspecified order \sref{IV.16.8.1}) is denoted by $\Diff_{X/S}(\sh{F}, \sh{G})$;
we also see as in \sref{IV.16.8.3} that $U \mapsto \Diff_{U/S}(\sh{F}|U, \sh{G}|U)$ is a sheaf of additive groups, which we will denote by $\shDiff_{X/S}(\sh{F}, \sh{G})$.
It is immediate that $\shDiff_{X/S}(\sh{F}, \sh{G})$ is the union of the increasing filtered family of its subsheaves $\shDiff_{X/S}^n(\sh{F}, \sh{G})$;
if $X$ is quasi-compact, $\Diff_{X/S}(\sh{F}, \sh{G})$ 
\oldpage[IV-4]{42}
is similarly the union of its subgroups $\Diff_{X/S}^n(\sh{F}, \sh{G})$ \sref{IV.16.8.1}.
The $\sh{O}_X$-bimodule structures on the $\shDiff_{X/S}^n(\sh{F}, \sh{G})$ therefore induce an $\sh{O}_X$-bimodule structure on $\shDiff_{X/S}(\sh{F}, \sh{G})$, again described by \sref{IV.16.8.5.1} and \sref{IV.16.8.5.2}.

Note that, for $n \leq m$, we have a commutative diagram
\[
  \label{IV.16.8.7.1}
  \xymatrix{
    \shHom_{\sh{O}_X}(\sh{P}_{X/S}^n(\sh{F}), \sh{G}) \ar[r]^-{\sim} \ar[d] & \shDiff_{X/S}^n(\sh{F}, \sh{G}) \ar[d] \\
    \shHom_{\sh{O}_X}(\sh{P}_{X/S}^m(\sh{F}), \sh{G}) \ar[r]^-{\sim} & \shDiff_{X/S}^m(\sh{F}, \sh{G})
  }
  \tag{16.8.7.1}
\] 
where the horizontal arrows are the isomorphisms \sref{IV.16.8.4.1} and the left vertical arrow comes from the canonical homomorphism $\sh{P}_{X/S}^m(\sh{F}) \to \sh{P}_{X/S}^n(\sh{F})$ \sref{IV.16.7.7}.
For every open set $U$ of $X$, we then equip $\Gamma(U, \sh{P}_{X/S}^\infty(\sh{F})) = \varprojlim \Gamma(U, \sh{P}_{X/S}^n(\sh{F}))$ with the projective-limit topology of the discrete topologies on the groups $\Gamma(U, \sh{P}_{X/S}^n(\sh{F}))$.  This makes $\Gamma(U, \sh{P}_{X/S}^\infty(\sh{F}))$ a topological $\Gamma(U, \sh{O}_X)$-bimodule, so that $\sh{P}_{X/S}^\infty(\sh{F})$ becomes a sheaf with values in the category of topological commutative groups \sref[0]{0.3.2.6}.
Then, by \cite[II, 1.11]{I-4}, the inductive limit of the system of sheaves of commutative groups $\bigl(\shHom_{\sh{O}_X}(\sh{P}_{X/S}^n(\sh{F}), \sh{G})\bigr)$ is precisely the sheaf of germs of \emph{continuous} homomorphisms from $\sh{P}_{X/S}^\infty(\sh{F})$ to $\sh{G}$ (the latter equipped with the discrete topology):
indeed, the continuous homomorphisms from $\Gamma(U, \sh{P}_{X/S}^\infty(\sh{F}))$ to the discrete group $\Gamma(U, \sh{G})$ correspond bijectively to the inductive systems of group homomorphisms $\Gamma(U, \sh{P}_{X/S}^n(\sh{F})) \to \Gamma(U, \sh{G})$.
We can furthermore express \sref{IV.16.8.4} by saying there is a canonical isomorphism
\[
  \shHomcont_{\sh{O}_X} (\sh{P}_{X/S}^\infty(\sh{F}), \sh{G}) \isoto \shDiff_{X/S}(\sh{F}, \sh{G})
\]
where the left-hand side denotes the sheaf of germs of continuous homomorphisms from $\sh{P}_{X/S}^\infty(\sh{F})$ to $\sh{G}$.
\end{env}

\begin{proposition}[16.8.8]
\label{IV.16.8.8}
Let $\sh{F}$, $\sh{G}$ be two $\sh{O}_X$-modules, $D: \sh{F} \to \sh{G}$ a homomorphism of $\psi^*(\sh{O}_S)$-modules, $n$ an integer $\geq 0$.
The following conditions are equivalent:
\begin{enumerate}
  \item[(a)] $D$ is a differential operator of order $\leq n$.
  \item[(b)] For all sections $a$ of $\sh{O}_X$ over an open set $U$, the homomorphism $D_a: \sh{F}|U \to \sh{G}|U$ such that, for every section $t$ of $\sh{F}$ over an open set $V \subset U$, we have
  \[  
    \label{IV.16.8.8.1}
    D_a(t) = D(at) - a D(t)
    \tag{16.8.8.1}
  \]
  is a differential operator of order $\leq n-1$.
  \item[(c)] For every open set $U$ of $X$, every family $(a_i)_{1 \leq i \leq n+1}$ of $n+1$ sections of $\sh{O}_X$ over $U$ and every section $t$ of $\sh{F}$ over $U$, we have the identity
  \[
    \label{IV.16.8.8.2}
    \sum_{H \subset I_{n+1}} (-1)^{\operatorname{Card}(H)} (\prod_{i \in H} a_i) D ((\prod_{i \notin H} a_i ) t ) = 0
    \tag{16.8.8.2}
  \]
  (where $I_{n+1}$ is the interval $1 \leq i \leq n+1$ of $\mathbb{N}$).
\end{enumerate}
\oldpage[IV-4]{43}
\end{proposition}

\begin{proof}
Let us first prove the equivalence of (a) and (c).
By definition, to prove that $D$ is a differential operator of order $\leq n$, it suffices to prove this for the restriction $D|U: \sh{F}|U \to \sh{G}|U$ to every affine open subset $U$ of $X$; on the other hand, property (c) holds for every open subset $U$ of $X$ if it holds for every affine open subset.
We may therefore restrict ourselves to the case where $S = \Spec(A)$ and $X = \Spec(B)$ are affine.
By \sref{IV.16.8.2} (whose notation we retain), condition (a) then means that the $A$-homomorphism $D': B \otimes_A M \to N$ defined by $D'(b \otimes t) = b D(t)$ vanishes on $\mathfrak{I}^{n+1}(B \otimes_A M)$; by \sref[0]{0.20.4.4}, this is equivalent to saying that $D'$ vanishes on every element of the form
\[
  (\prod_{i=1}^{n+1} (a_i \otimes 1 - 1 \otimes a_i) ) . (1 \otimes t)
\]
where $a_i \in B$ and $t \in M$.
Now this element can be written, up to the common sign $(-1)^{n+1}$, as $\sum_{H \subset I_{n+1}} (-1)^{\operatorname{Card}(H)} (\prod_{i \in H} a_i) \otimes ((\prod_{i \notin H} a_i ) t )$, and its image under $D'$ is, up to that common sign, the left-hand side of \sref{IV.16.8.8.2}; this proves the equivalence of (a) and (c).

Let us now prove the equivalence of (b) and (c).
We argue by induction on $n$, the assertion being trivial for $n=0$.
Writing $a_{n+1}$ in place of $a$ in condition (b), we see from the induction hypothesis that condition (b) means that, for every family $(a_i)_{1 \leq i \leq n}$ of $n$ sections of $\sh{O}_X$ over $U$ and every section $t$ of $\sh{F}$ over $U$, one has
\[
  \sum_{H' \subset I_n} (-1)^{\operatorname{Card}(H')}  (\prod_{i \in H'} a_i) D_{a_{n+1}} ((\prod_{i \notin H'} a_i ) t ) = 0.
\]
But if in this relation we replace $D_{a_{n+1}}$ by its definition \sref{IV.16.8.8.1}, we see at once that, up to sign, we obtain the left-hand side of \sref{IV.16.8.8.2}; hence the conclusion.
\end{proof}

\begin{proposition}[16.8.9]
\label{IV.16.8.9}
If $D: \sh{F} \to \sh{G}$ is a differential operator of order $\leq n$, and $D': \sh{G} \to \sh{H}$ a differential operator of order $\leq n'$, then $D' \circ D: \sh{F} \to \sh{H}$ is a differential operator of order $\leq n + n'$.
\end{proposition}

\begin{proof}
By hypothesis, we can write $D = u \circ d_{X/S, \sh{F}}^n$ and $D' = v \circ d_{X/S, \sh{G}}^{n'}$, where $u: \sh{P}_{X/S}^n \otimes_{\sh{O}_X} \sh{F} \to \sh{G}$ and $v: \sh{P}_{X/S}^{n'} \otimes_{\sh{O}_X} \sh{G} \to \sh{H}$ are $\sh{O}_X$-homomorphisms.
It is therefore enough to show that the composite homomorphism of sheaves of additive groups
\[
  \xymatrix{
    \sh{F} \ar[r]^-{d_{X/S, \sh{F}}^n} & \sh{P}_{X/S}^n \otimes_{\sh{O}_X} \sh{F} \ar[r]^-u & \sh{G} \ar[r]^-{d_{X/S, \sh{G}}^{n'}} & \sh{P}_{X/S}^{n'} \otimes_{\sh{O}_X} \sh{G}
  }
\]
factors as
\[
  \xymatrix{
    \sh{F} \ar[r]^-{d_{X/S, \sh{F}}^{n + n'}} & \sh{P}_{X/S}^{n + n'} \otimes_{\sh{O}_X} \sh{F} \ar[r]^-w &  \sh{P}_{X/S}^{n'} \otimes_{\sh{O}_X} \sh{G}
  }
\]
where $w$ is an $\sh{O}_X$-homomorphism.
It suffices to prove the following lemma.
\begin{lemma}[16.8.9.1]
\label{IV.16.8.9.1}
There is one and only one $\sh{O}_X$-homomorphism
\[
  \label{IV.16.8.9.2}
  \delta: \sh{P}_{X/S}^{n + n'} \to \sh{P}_{X/S}^{n'}(\sh{P}_{X/S}^n) = \sh{P}_{X/S}^{n'} \otimes_{\sh{O}_X} \sh{P}_{X/S}^n 
  \tag{16.8.9.2}
\]
\oldpage[IV-4]{44}
making the following diagram commute
\[
  \label{IV.16.8.9.3}
  \xymatrix{
    \sh{O}_X \ar[r]^-{d_{X/S}^{n + n'}} \ar[d]_-{d_{X/S}^{n}} & \sh{P}_{X/S}^{n + n'} \ar[d]^-\delta \\
    \sh{P}_{X/S}^n \ar[r]_-{d_{X/S, \sh{P}_{X/S}^n}^{n'}} & \sh{P}_{X/S}^{n'} (\sh{P}_{X/S}^n)
  }
  \tag{16.8.9.3}
\]
\end{lemma}

We will then have, indeed, a commutative diagram deduced from \sref{IV.16.8.9.3} by tensorization with $\sh{F}$
\[
  \xymatrix{
    \sh{F} \ar[r]^-{d_{X/S, \sh{F}}^{n + n'}} \ar[d]_-{d_{X/S, \sh{F}}^{n}} & \sh{P}_{X/S}^{n + n'}(\sh{F}) \ar[d]^-{\delta \otimes 1} \\
    \sh{P}_{X/S}^n(\sh{F}) \ar[r]_-{d_{X/S, \sh{P}_{X/S}^n(\sh{F})}^{n'}} & \sh{P}_{X/S}^{n'}(\sh{P}_{X/S}^n(\sh{F}))
  }
\]
and on the other hand, we verify immediately from the definition \sref{IV.16.7.5} that the diagram
\[
  \xymatrix{
    \sh{P}_{X/S}^n(\sh{F}) \ar[r]^-{u} \ar[d]_-{d_{X/S, \sh{P}_{X/S}^n(\sh{F})}^{n'}} & \sh{G} \ar[d]^-{d_{X/S, \sh{G}}^{n'}} \\
    \sh{P}_{X/S}^{n'}(\sh{P}_{X/S}^n(\sh{F})) \ar[r]_-{1 \otimes u} & \sh{P}_{X/S}^{n'}(\sh{G})
  }
\]
is commutative.
We finish the proof by taking $w$ to be the composite $\sh{O}_X$-homomorphism
\[
  \xymatrix{
    \sh{P}_{X/S}^{n+n'}(\sh{F}) \ar[r]^-{\delta \otimes 1} & \sh{P}_{X/S}^{n'}(\sh{P}_{X/S}^n(\sh{F})) \ar[r]^-{1 \otimes u} & \sh{P}_{X/S}^{n'}(\sh{G}).
  }
\]
\end{proof}

\begin{proof}[Proof of {\hyperlink{ega4.c04.lem.16-8-9-1}{Lemma 16.8.9.1}}]
It remains to prove Lemma \sref{IV.16.8.9.1}.
In view of \sref{IV.16.7.6}, which proves the uniqueness of $\delta$, we are reduced to the case where $S = \Spec(A)$ and $X = \Spec(B)$ are affine.
Putting $\mathfrak{I} = \mathfrak{I}_{B/A}$, it suffices to define a canonical homomorphism of $B$-modules
\[
  \vphi: (B \otimes_A B)/\mathfrak{I}^{n+n'+1} \to ((B \otimes_A B)/\mathfrak{I}^{n' + 1}) \otimes_B ((B \otimes_A B)/\mathfrak{I}^{n + 1})
\]
the $B$-module structures on both sides being induced by the first factor $B$.
Recall that in the tensor product on the right, $(B \otimes_A B)/\mathfrak{I}^{n' + 1}$ must be regarded
\oldpage[IV-4]{45}
as a right $B$-module by its second $B$ factor, and $(B \otimes_A B)/\mathfrak{I}^{n + 1}$ as a left $B$-module by its first $B$ factor \sref{IV.16.7.2}.
This is in turn equivalent to defining a homomorphism of $B$-modules
\[
  \vphi_0: B \otimes_A B \to ((B \otimes_A B)/\mathfrak{I}^{n' + 1}) \otimes_B ((B \otimes_A B)/\mathfrak{I}^{n + 1})
\]
and proving that it vanishes on $\mathfrak{I}^{n+n'+1}$.
Such a homomorphism is defined at once by the condition
\[
  \vphi_0(b \otimes b') = \pi_{n'}(b \otimes 1) \otimes \pi_n(1 \otimes b') \quad \text{for $b$, $b'$ in $B$}
\]
with the notation of \sref{IV.16.3.7}.
Moreover, it is immediate that $\vphi_0$ is a homomorphism of \emph{rings}.
Now we may write
\[
  \vphi_0(b \otimes 1 - 1 \otimes b) = \pi_{n'}(b \otimes 1 - 1 \otimes b) \otimes \pi_n(1 \otimes 1) + \pi_{n'}(1 \otimes b) \otimes \pi_n(1 \otimes 1) - \pi_{n'}(1 \otimes 1) \otimes \pi_n(1 \otimes b)
\]
and we have 
\[
  \begin{aligned}
  \pi_{n'}(1 \otimes b) \otimes \pi_n(1 \otimes 1)
  &=\pi_{n'}(1 \otimes 1)b \otimes \pi_n(1 \otimes 1)\\
  &=\pi_{n'}(1 \otimes 1) \otimes b\pi_n(1 \otimes 1)\\
  &=\pi_{n'}(1 \otimes 1) \otimes \pi_n(b \otimes 1)
  \end{aligned}
\]
whence, finally,
\[
  \label{IV.16.8.9.4}
  \vphi_0(b \otimes 1 - 1 \otimes b) = \pi_{n'}(b \otimes 1 - 1 \otimes b) \otimes \pi_n(1 \otimes 1) + \pi_{n'}(1 \otimes 1) \otimes \pi_n(b \otimes 1 - 1 \otimes b). 
  \tag{16.8.9.4}
\]
A product of $n + n' + 1$ factors of the form \sref{IV.16.8.9.4} is therefore necessarily zero, since this is so for a product of $n+1$ factors of the form $\pi_n(b \otimes 1 - 1 \otimes b)$ and for a product of $n' + 1$ factors of the form $\pi_{n'}(b \otimes 1 - 1 \otimes b)$.
The conclusion therefore follows from \sref[0]{0.20.4.4}.
\end{proof}

\begin{corollary}[16.8.10]
\label{IV.16.8.10}
The sheaf $\shDiff_{X/S}(\sh{O}_X, \sh{O}_X)$ (also denoted $\shDiff_{X/S}$) is canonically endowed with a structure of sheaf of rings, and the $\shDiff_{X/S}^n$ form an increasing filtration compatible with this structure.
\end{corollary}

In particular, $\shDiff_{X/S}^0$ is a sheaf of subrings of $\shDiff_{X/S}$, which is canonically identified with $\sh{O}_X$ \sref{IV.16.8.1}.
The formulas \sref{IV.16.8.5.1} and \sref{IV.16.8.5.2} show that the structure of $\sh{O}_X$-bimodule of $\shDiff_{X/S}$ comes from the multiplication on the left and on the right by sections of $\sh{O}_X$ considered as a sheaf of subrings of $\shDiff_{X/S}$.

\begin{remarks}[16.8.11]
\label{IV.16.8.11}
\begin{enumerate}
  \item[\rm{(i)}] Suppose that $\sh{F} = \oplus_{\lambda \in L} \sh{F}_\lambda$;
  then it is clear from \sref{IV.16.7.2.1} that $\sh{P}_{X/S}^n(\sh{F}) = \oplus_{\lambda \in L} \sh{P}_{X/S}^n(\sh{F}_\lambda)$.
  Since the functor $\sh{F} \mapsto \Gamma(U, \sh{F})$ commutes with the formation of arbitrary direct sums, $d_{X/S, \sh{F}}^n$ is the homomorphism whose restriction to each $\sh{F}_\lambda$ is $d_{X/S, \sh{F}_\lambda}^n: \sh{F}_\lambda \to \sh{P}_{X/S}^n(\sh{F}_\lambda)$;
  it follows at once that
  \[
    \Diff_{X/S}^n (\sh{F}, \sh{G}) = \prod_{\lambda \in L} \Diff_{X/S}^n (\sh{F}_\lambda, \sh{G}),
  \]
  and therefore also \sref[0]{0.3.2.6}
  \[
    \shDiff_{X/S}^n (\sh{F}, \sh{G}) = \prod_{\lambda \in L} \shDiff_{X/S}^n (\sh{F}_\lambda, \sh{G}).
  \]
  Moreover, if $\sh{G} = \prod_{\mu \in M} \sh{G}_\mu$ \sref[0]{0.3.2.6}, we have
  \[
    \Hom_{\sh{O}_X}(\sh{P}_{X/S}^n(\sh{F}), \sh{G}) = \prod_{\mu \in M} \Hom_{\sh{O}_X}(\sh{P}_{X/S}^n(\sh{F}), \sh{G}_\mu),
  \]
  \oldpage[IV-4]{46}
  Every homomorphism $u$ from $\sh{P}_{X/S}^n(\sh{F})$ to $\sh{G}$ corresponds bijectively to the family of its composites $u_\mu: \sh{P}_{X/S}^n(\sh{F}) \to \sh{G} \to \sh{G}_\mu$.
  Thus
  \[
    \Diff_{X/S}^n(\sh{F}, \sh{G}) = \prod_{\mu \in M} \Diff_{X/S}^n(\sh{F}, \sh{G}_\mu),
  \]
  and consequently also
  \[
    \shDiff_{X/S}^n(\sh{F}, \sh{G}) = \prod_{\mu \in M} \shDiff_{X/S}^n(\sh{F}, \sh{G}_\mu).
  \]
  \item[\rm{(ii)}]
  So far, we have encountered differential operators $\sh{F} \to \sh{G}$ almost exclusively when $\sh{F}$ and $\sh{G}$ are locally free of finite rank; in that case, by \rm{(i)}, their structure is reduced locally to that of the sheaf $\shDiff_{X/S}$.
  The latter will be studied below, in a particular case, in \sref{IV.16.11}.
\end{enumerate}
\end{remarks}

\subsection{Regular and quasi-regular immersions}
\label{IV.16.9}

\begin{definition}[16.9.1]
\label{IV.16.9.1}
Let $X$ be a ringed space.
We say that an ideal $\sh{I}$ of $\sh{O}_X$ is regular (resp. quasi-regular) if, for every point $x \in \Supp(\sh{O}_X/\sh{I})$, there is an open neighbourhood $U$ of $x$ in $X$ and a regular sequence \sref[0]{0.15.2.2} (resp. quasi-regular sequence \sref[0]{0.15.2.2}) of elements of $\Gamma(U, \sh{O}_X)$ which generates $\sh{I}|U$.
\end{definition}

A regular (resp. quasi-regular) sequence of sections of $\sh{O}_X$ over $U$ which generates $\sh{I}|U$ is called a \emph{regular system} (resp. \emph{quasi-regular system}) of generators of $\sh{I}|U$.

\begin{definition}[16.9.2]
\label{IV.16.9.2}
Let $j: Y \to X$ be an immersion of preschemes, and let $U$ be an open subset of $X$ such that $j(Y) \subset U$ and $j$ is a closed immersion of $Y$ into $U$.
We say that $j$ is regular (resp. quasi-regular) if the closed subprescheme $j(Y)$ of $U$ associated to $j$ is defined by a regular (resp. quasi-regular) ideal of $\sh{O}_U$ (condition independent of the choice of $U$).
\end{definition}

We say that a subprescheme $Y$ of a prescheme $X$ is \emph{regularly immersed} (resp. \emph{quasi-regularly immersed}) if the canonical injection $j:Y \to X$ is a regular immersion (resp. quasi-regular immersion). 
If $Y$ is a closed subprescheme and $\sh{I}$ is the ideal of $\sh{O}_X$ that defines $Y$, it is the same as asking $\sh{I}$ to be \emph{regular} (resp. \emph{quasi-regular}).

For example, if $A$ is an \emph{integral} ring and $f$ is a nonzero element of $A$, the closed subprescheme $V(f)$ of $\Spec(A)$ (isomorphic to $\Spec(A/(f))$) is \emph{regularly immersed} in $\Spec(A)$.

Every regular ideal is quasi-regular \sref[0]{0.15.2.2};
every regular immersion is quasi-regular (cf. \sref{IV.16.9.11} for a converse).  

\begin{proposition}[16.9.3]
\label{IV.16.9.3}
Let $X$ be a ringed space, $\sh{I}$ an ideal of $\sh{O}_X$, $(f_i)_{1 \leq i \leq m}$ a finite sequence of sections of $\sh{O}_X$ over $X$ generating $\sh{I}$.
For $(f_i)$ to be a quasi-regular sequence \sref[0]{0.15.2.2}, it is necessary and sufficient that the following conditions be satisfied:
\begin{enumerate}
  \item[\rm{(i)}] The canonical images of $f_i$ in $\sh{I}/\sh{I}^2$ form a basis of this $\sh{O}_X/\sh{I}$-module.
  \item[\rm{(ii)}] The canonical surjective homomorphism \sref{IV.16.1.2.2}
  \[
    \bb{S}_{\sh{O}_X/\sh{I}}^\bullet (\sh{I}/\sh{I}^2) \to \shGr_{\sh{I}}^\bullet(\sh{O}_X) 
  \]
  is bijective.
\end{enumerate}

Moreover, if this is so, every sequence $(f'_i)_{1 \leq i \leq n}$ of $n$ sections of $\sh{I}$ over $X$ which generates $\sh{I}$ is quasi-regular.
\end{proposition}

\oldpage[IV-4]{47}

\begin{proof}
The two conditions in the statement merely express the definition given in \sref[0]{0.15.2.2}, taking into account the definition of the canonical homomorphisms \sref[0]{0.15.2.1.1}.
The last claim follows from the fact that, if a module $M$ over a commutative ring $A$ admits a basis of $n$ elements, every system of $n$ generators of $M$ is a basis (Bourbaki, \emph{Alg. comm.}, chap.~II, \textsection3, cor. 5 of th. 1).
\end{proof}

\begin{corollary}[16.9.4]
\label{IV.16.9.4}
Let $X$ be a locally ringed space, $\sh{I}$ an ideal of $\sh{O}_X$.
For $\sh{I}$ to be quasi-regular, it is necessary and sufficient that the following conditions hold:
\begin{enumerate}
  \item[\rm{(i)}] $\sh{I}$ is of finite type.
  \item[\rm{(ii)}] $\sh{I}/\sh{I}^2$ is a locally free $\sh{O}_X/\sh{I}$-module.
  \item[\rm{(iii)}] The canonical homomorphism
  \[
    \label{IV.16.9.4.1}
    \bb{S}_{\sh{O}_X/\sh{I}}^\bullet (\sh{I}/\sh{I}^2) \to \shGr_{\sh{I}}^\bullet(\sh{O}_X) 
    \tag{16.9.4.1}
  \]
  is bijective.
\end{enumerate}
\end{corollary}

\begin{proof}
The necessity of the conditions follows immediately from \sref{IV.16.9.3}.
To see that they are sufficient, it is enough, by \sref{IV.16.9.3}, to prove the following: if, at a point $x \in \Supp(\sh{O}_X/\sh{I})$, there are an open neighbourhood $U$ of $x$ in $X$ and $n$ sections $f_i$ $(1 \leq i \leq n)$ of $\sh{I}$ over $U$ whose canonical images in $\sh{I}/\sh{I}^2$ form a basis of $(\sh{I}/\sh{I}^2)|U$ over $(\sh{O}_X/\sh{I})|U$, then there is an open neighbourhood $V \subset U$ of $x$ such that the $f_i|V$ generate $\sh{I}|V$.
Now, by hypothesis, $\sh{I}_x \neq \sh{O}_x$, so $\sh{I}_x$ is contained in the maximal ideal of $\sh{O}_x$.
Since $\sh{I}_x$ is an $\sh{O}_x$-module of finite type and the classes of $(f_i)_x$ in $\sh{I}_x/\sh{I}_x^2$ generate this $(\sh{O}_x/\sh{I}_x)$-module, Nakayama's lemma shows that the $(f_i)_x$ generate $\sh{I}_x$.
Since $\sh{I}$ is of finite type, we conclude by \sref[0]{0.5.2.2}.
\end{proof} 

\begin{corollary}[16.9.5]
\label{IV.16.9.5}
Let $X$ be a locally ringed space, $\sh{I}$ a quasi-regular ideal of $\sh{O}_X$, $(f_i)_{1 \leq i \leq n}$ a sequence of sections of $\sh{I}$ over $X$, $x$ a point of $\Supp(\sh{O}_X/\sh{I})$. 
The following conditions are equivalent:
\begin{enumerate}
  \item[(a)] There is an open neighbourhood $U$ of $x$ in $X$ such that the $f_i|U$ form a quasi-regular sequence of elements of $\Gamma(U, \sh{O}_X)$ generating $\sh{I}|U$.
  \item[(b)] The $(f_i)_x$ form a system of generators of $\sh{I}_x$ whose size is as small as possible.
  \item[(b')] The $(f_i)_x$ form a minimal set of generators of $\sh{I}_x$.
  \item[(c)] If $\bar f_i$ is the canonical image of $f_i$ in $\Gamma(X, \sh{I}/\sh{I}^2)$, the $(\bar f_i)_x$ form a basis for the $(\sh{O}_x/\sh{I}_x)$-module $\sh{I}_x/\sh{I}_x^2$.
\end{enumerate}
\end{corollary}

\begin{proof}
By hypothesis, $\sh{O}_x$ is a local ring, $\sh{I}_x$ an ideal of finite type of $\sh{O}_x$ contained in the maximal ideal of $\sh{O}_x$;
the equivalence of (b), (b') and (c) follows from Nakayama's lemma (Bourbaki, \emph{Alg. comm.}, chap.~II, \textsection3, n\textsuperscript{o}2, prop. 5).
It is clear that (a) implies (c) because of \sref{IV.16.9.3};
on the other hand, from \sref[0]{0.5.2.2} it follows that if condition (c) is satisfied (and therefore so is (b)), there is a neighbourhood $U$ of $x$ in $X$ such that $(\sh{I}/\sh{I}^2)|U$ has constant rank equal to $n$, and that the $f_i|U$ generate $\sh{I}|U$;
  it suffices now to apply the last assertion of \sref{IV.16.9.3} to $U$.
\end{proof}

\begin{remarks}[16.9.6]
\label{IV.16.9.6}
\begin{enumerate}
  \item[\rm{(i)}] Under the general hypothesis of \sref{IV.16.9.5}, for the sequence $(f_i)$ to generate $\sh{I}$, it is not enough that the $(\bar f_i)_y$ form a basis of the $(\sh{O}_y/\sh{I}_y)$-module $(\sh{I}_y/\sh{I}_y^2)$ for all $y \in X$.
  An example is obtained by
  \oldpage[IV-4]{48}
  taking $X = \Spec(A)$, where $A$ is a Dedekind ring, and $\sh{I} = \widetilde{\mathfrak{I}}$, where $\mathfrak{I}$ is a \emph{non-principal} prime ideal of $A$.
  Then $\sh{I}_y/\sh{I}_y^2 = 0$ at every point $y$ other than the point $x \in X$ corresponding to $\mathfrak{I}$, while $\sh{I}_x/\sh{I}_x^2$ has rank $1$ over the field $\sh{O}_x/\sh{I}_x$;
  moreover, $\sh{I}$ is plainly a regular ideal.

  \item[\rm{(ii)}] In \sref{IV.16.9.5}, one cannot replace ``quasi-regular'' by ``regular'', even when $X$ is a prescheme (cf. \sref{IV.16.9.12}).
  Indeed, let $B$ denote the ring of germs at $0$ in $\mathbb{R}$ of infinitely differentiable functions.
  It has a maximal ideal $\mathfrak{m}$ generated by the germ $t$ of the identity map of $\mathbb{R}$ at $0$, and the intersection $\mathfrak{n}$ of the $\mathfrak{m}^k$ for $k>0$ is nonzero.
  Now let $A$ be the quotient ring $B[T]/\mathfrak{n}TB[T]$, and let $f_1$, $f_2$ be the canonical images in $A$ of the elements $t$ and $T$ of $B[T]$.
  The sequence $(f_1, f_2)$ is \emph{regular} in $A$: indeed, $f_1$ is not a zero divisor in $A$, because the relation $tP[T] \in \mathfrak{n}TB[T]$, for a polynomial $P \in B[T]$, implies that the products of $t$ with the coefficients of $P$ belong to $\mathfrak{n}$; it follows at once that the coefficients themselves belong to $\mathfrak{n}$, and hence $P[T] \in \mathfrak{n}TB[T]$.
  Since $B/tB$ is isomorphic to $\mathbb{R}$, $A/f_1A$ is isomorphic to the polynomial ring $\mathbb{R}[T]$ and is therefore integral; the image of $f_2$ in $A/f_1A$, being equal to $T$, is consequently not a zero divisor, proving the assertion.
  However, $f_2$ is a zero divisor in $A$, since for every nonzero element $x \in \mathfrak{n}$, the image of $x$ in $A$ is nonzero whereas the image of $xT$ is zero.
  We conclude that the sequence $(f_2, f_1)$ is \emph{not regular} in $A$;
  on the other hand, the ideal $\mathfrak{I} = f_1A + f_2A$ is distinct from $A$, so the conditions (b), (b') and (c) of \sref{IV.16.9.5} do not imply the condition (a) when we replace ``quasi-regular'' by ``regular''.
\end{enumerate}
\end{remarks}

\begin{env}[16.9.7]
\label{IV.16.9.7}
If $X = \Spec(A)$ is an affine scheme, we say that an ideal $\mathfrak{I}$ of $A$ is \emph{regular} (resp. \emph{quasi-regular}) if the ideal $\sh{I} = \widetilde{\mathfrak{I}}$ of $\sh{O}_X$ is regular (resp. quasi-regular).
We note that this notion is \emph{local} and does not imply the existence of a \emph{system of generators} of $\mathfrak{I}$ forming a regular (resp. quasi-regular) sequence in $A$, as Example \sref{IV.16.9.6}[(i)] shows;
it does, however, imply this when $A$ is \emph{local} \sref{IV.16.9.5}.

The proposition \sref{IV.16.9.4} can be translated in terms of quasi-regular immersions in the following manner:
\end{env}

\begin{proposition}[16.9.8]
\label{IV.16.9.8}
Let $j:Y \to X$ be a morphism of preschemes;
for $j$ to be a quasi-regular immersion, it is necessary and sufficient that $j$ satisfies the following conditions:
\begin{enumerate}
  \item[(i)] $j$ is an immersion locally of finite presentation.
  \item[(ii)] The conormal sheaf $\shGr^1(j) = \sh{N}_{Y/X}$ \sref{IV.16.1.2} is a locally free $\sh{O}_Y$-module.
  \item[(iii)] The canonical homomorphism
  \[
    \bb{S}_{\sh{O}_Y}^\bullet(\shGr^1(j)) \to \shGr^\bullet(j)
  \]
  \sref{IV.16.1.2.2} is bijective.
\end{enumerate}
\end{proposition}

\begin{proof}
Since the question is local on $Y$, we may restrict ourselves to the case where $j$ is the canonical injection of a closed subprescheme $Y$ of $X$; the translation of \sref{IV.16.9.4} into \sref{IV.16.9.8} then follows from the description of $\shGr^1(j)$ and $\shGr^\bullet(j)$ in terms of the ideal $\sh{I}$ of $\sh{O}_X$ defining $Y$ \sref{IV.16.1.3}[(ii)].
\end{proof}

\begin{corollary}[16.9.9]
\label{IV.16.9.9}
Let $Y$ be a prescheme, $X$ a $Y$-prescheme, and $j:Y \to X$ a $Y$-section of $X$, so that the $n$-th normal invariant $\sh{A}^{(n)}$ of $j$ \sref{IV.16.1.2} is an augmented $\sh{O}_Y$-algebra \sref{IV.16.1.7};
take $\sh{A}^{(\infty)} = \varprojlim \sh{A}^{(n)}$.
For $j$ to be a quasi-regular immersion, it is necessary and sufficient that $j$ be locally of finite presentation and that every $y \in Y$ admit an affine open neighbourhood $U$ with coordinate ring $C$ such that $\sh{A}^{(\infty)}|U$ is isomorphic, as an augmented topological $\sh{O}_U$-algebra, to $\sh{O}_U[[T_1, \dots, T_n]]$.
\end{corollary}

\begin{proof}
We may reduce to the case where $j$ is a closed immersion by restricting to a sufficiently small neighbourhood of $y$ (see the argument of \sref{IV.16.4.11}); then $\sh{O}_Y$ is identified with the quotient algebra $\sh{O}_X/\sh{I}$, and the canonical surjective homomorphism $\sh{O}_X \to \sh{O}_Y$ admits a
\oldpage[IV-4]{49}
right inverse \sref{IV.16.1.7}.
We may therefore suppose that $X = \Spec(B)$ and $Y = \Spec(A)$ are affine, with $B$ an augmented $A$-algebra whose augmentation ideal $\mathfrak{I}$ is of finite type.
Since $\sh{A}^{(n)}$ is then identified with $(B/\mathfrak{I}^{n+1})^{\sim}$, the corollary follows from the equivalence of (b) and (c) in \sref[0]{0.19.5.4}, since $B/\mathfrak{I} = A$.
\end{proof}

We note that, in the affine case under consideration, the fact that $j$ is a quasi-regular immersion is also equivalent, by \sref[0]{0.19.5.4}, to saying that $B$ is a \emph{formally smooth} $A$-algebra for the $\mathfrak{I}$-preadic topology.

We also note that the condition that $j$ is locally of finite presentation is always satisfied when $X \to Y$ is locally of finite type \sref{IV.1.4.3}[(v)].

\begin{proposition}[16.9.10]
\label{IV.16.9.10}
Let $X$ be a locally Noetherian prescheme, $Y$ a subprescheme of $X$, $j:Y \to X$ the canonical injection, $y$ a point of $Y$.
\begin{enumerate}
  \item[\rm{(i)}] For there to exist an open neighbourhood $U$ of $y$ in $X$ such that the restriction $Y \cap U \to U$ of $j$ be a regular immersion, it is necessary and sufficient that the kernel $\sh{I}_y$ of the surjective homomorphism $\sh{O}_{X,y} \to \sh{O}_{Y,y}$ is generated by a regular sequence of elements of $\sh{O}_{X,y}$.

  \item[\rm{(ii)}] For the immersion $j$ to be regular, it is necessary and sufficient that it is quasi-regular.
\end{enumerate}
\end{proposition}

\begin{proof}
  \begin{enumerate}
    \item[\rm{(i)}] We can reduce to the case where $Y$ is a closed subprescheme of $X$ defined by a \emph{coherent} ideal $\sh{I}$ of $\sh{O}_X$.
    The condition is clearly necessary.
    Conversely, if $\sh{I}_y$ is generated by a regular sequence $(s_i)_y$, where the $s_i$ are sections of $\sh{I}$ over an open neighbourhood $U$ of $y$ in $X$, we may suppose that the $s_i$ generate $\sh{I}|U$ \sref[0]{0.5.2.2} and form a regular sequence \sref[0]{0.15.2.4}, whence the assertion.
    \item[\rm{(ii)}] The fact that a quasi-regular immersion is regular follows from (i) and the identification of quasi-regular and regular sequences in $\sh{O}_{X,y}$ formed by elements of the maximal ideal \sref[0]{0.15.1.11}.
  \end{enumerate}
\end{proof}

When, without any Noetherian hypothesis on $X$, the kernel $\sh{I}_y$ of $\sh{O}_{X,y} \to \sh{O}_{Y,y}$ is generated by a regular sequence of elements of $\sh{O}_{X,y}$, we say that the immersion $j$ is \emph{regular at the point $y$}.

\begin{corollary}[16.9.11]
\label{IV.16.9.11}
Let $X$ be a locally Noetherian prescheme;
then every quasi-regular ideal of $\sh{O}_X$ is regular.
\end{corollary}

\begin{remarks}[16.9.12]
\label{IV.16.9.12}
\begin{enumerate}
  \item[\rm{(i)}] We note that a regular immersion is not in general a flat morphism, and hence \textit{a fortiori} is not a regular morphism in the sense of \sref{IV.6.8.1}.
  \item[\rm{(ii)}] Let $A$ be a local Noetherian ring;
  it follows immediately from \sref{IV.16.9.4} and from \sref[0]{0.17.1.1} that for $A$ to be \emph{regular}, it is necessary and sufficient that its maximal ideal $\mathfrak{m}$ is \emph{quasi-regular} (or regular, which amounts to the same thing given that $A$ is Noetherian).
  For an affine Noetherian scheme $X$ to be \emph{regular}, it is necessary and sufficient that for every closed point $x \in X$, the canonical injection $\Spec(\kres(x)) \to X$ be a \emph{regular} immersion.
\end{enumerate}
\end{remarks}

\begin{proposition}[16.9.13]
\label{IV.16.9.13}
Let $X$ be a locally Noetherian prescheme, $Y$ a subprescheme of $X$, $Y'$ a subprescheme of $Y$, such that the canonical injection $j:Y' \to Y$ is regular. 
Then the sequence of $\sh{O}_{Y'}$-modules
\[
  \label{IV.16.9.13.1}
  \xymatrix{
    0 \ar[r] & j^*(\sh{N}_{Y/X}) \ar[r] & \sh{N}_{Y'/X} \ar[r] & \sh{N}_{Y'/Y} \ar[r] & 0
  }
  \tag{16.9.13.1}
\]
\oldpage[IV-4]{50}
is exact;
furthermore, for every $x \in X$, there is an open neighbourhood $U$ of $x$ such that the restrictions to $U$ of the homomorphisms of \sref{IV.16.9.13.1} form a split exact sequence.
\end{proposition}

Let us first prove the following lemma:

\begin{lemma}[16.9.13.2]
\label{IV.16.9.13.2}
Let $A$ be a ring, $\mathfrak{I}$ an ideal of $A$, $A' = A/\mathfrak{I}$, $(f_i)_{1 \leq i \leq r}$ a sequence of elements of $A$ which is $A'$-regular, $\mathfrak{K} = \sum_i f_i A$, $\mathfrak{L} = \mathfrak{I} + \mathfrak{K}$, $\mathfrak{K}' = \sum_i f_i A'$, so that $C = A/\mathfrak{L}$ is isomorphic to $A'/\mathfrak{K}'$.
For every integer $n > 0$, and every integer $N \geq n$, we have the relation
\[
  \label{IV.16.9.13.3}
  \mathfrak{I} \cap \mathfrak{K}^n = \mathfrak{I}\mathfrak{K}^n + (\mathfrak{I} \cap \mathfrak{K}^N).
  \tag{16.9.13.3}
\]
\end{lemma}

\begin{proof}
It is plainly enough to prove that every element of the left-hand side belongs to the right-hand side, and an induction reduces us to the case $N = n + 1$.
An element of the left-hand side of \sref{IV.16.9.13.3}, being in $\mathfrak{K}^n$, can be written as $P(f_1, \dots, f_r)$, where $P \in A[T_1, \dots, T_r]$ is homogeneous of degree $n$.
If $f'_i$ is the canonical image of $f_i$ in $A'$, the hypothesis $P(f_1, \dots, f_r) \in \mathfrak{I}$ means that $P(f'_1, \dots, f'_r) = 0$.
But $P(f'_1, \dots, f'_r) \in \mathfrak{K}'^n$, so the canonical image of $P(f'_1, \dots, f'_r)$ in $\mathfrak{K}'^n/\mathfrak{K}'^{n+1}$ is zero.
Now the hypothesis that $(f_i)$ is $A'$-regular implies that the canonical homomorphism $\bb{S}_C^n(\mathfrak{K}'/\mathfrak{K}'^2) \to \mathfrak{K}'^n/\mathfrak{K}'^{n+1}$ is bijective \sref[0]{0.15.1.9};
we conclude that the coefficients of $P$ are in $\mathfrak{L} = \mathfrak{I} + \mathfrak{K}$.
It follows at once that $P(f_1, \dots, f_r) \in \mathfrak{I}\mathfrak{K}^n + \mathfrak{K}^{n+1}$, and since $P(f_1, \dots, f_r) \in \mathfrak{I}$, we finally have $P(f_1, \dots, f_r) \in \mathfrak{I}\mathfrak{K}^n + (\mathfrak{I} \cap \mathfrak{K}^{n+1})$, which proves the lemma.
\end{proof}

Taking the quotients of both sides of \sref{IV.16.9.13.3} by $\mathfrak{I}\mathfrak{K}^{n}$, we see that the relations \sref{IV.16.9.13.3}, for $N \geq n$, imply
\[
  \label{IV.16.9.13.4}
  (\mathfrak{I} \cap \mathfrak{K}^n)/\mathfrak{I}\mathfrak{K}^n \subset \bigcap_{N \geq n} \mathfrak{K}^N \cdot (A/(\mathfrak{I}\mathfrak{K}^n)).
  \tag{16.9.13.4}
\]

We deduce the following corollary.

\begin{corollary}[16.9.13.5]
\label{IV.16.9.13.5}
Suppose that the hypotheses of \sref{IV.16.9.13.2} are satisfied and also that the ring $A$ is Noetherian and $\mathfrak{K}$ is contained in the radical of $A$.
Then for every integer $n > 0$,
\[
  \label{IV.16.9.13.6}
  \mathfrak{I} \cap \mathfrak{K}^n = \mathfrak{I}\mathfrak{K}^n.
  \tag{16.9.13.6}
\]
\end{corollary}

\begin{proof}
Indeed, the second member of \sref{IV.16.9.13.4} is then zero, given that $A/\mathfrak{I}\mathfrak{K}^n$ is an $A$-module of finite type (Bourbaki, \emph{Alg. comm.}, chap.~III, \textsection3, n\textsuperscript{o}3, prop. 6).
\end{proof}

Taking in particular $n = 2$ in \sref{IV.16.9.13.6}, and observing that $\mathfrak{L}^2 = \mathfrak{I}^2 + \mathfrak{I}\mathfrak{K} + \mathfrak{K}^2 = \mathfrak{I}\mathfrak{L} + \mathfrak{K}^2$, we have $\mathfrak{I}\mathfrak{L} \subset \mathfrak{I} \cap \mathfrak{L}^2$ and therefore
\[
  \mathfrak{I} \cap \mathfrak{L}^2 = \mathfrak{I}\mathfrak{L} + (\mathfrak{I} \cap \mathfrak{K}^2) = \mathfrak{I}\mathfrak{L} + \mathfrak{I}\mathfrak{K}^2 = \mathfrak{I}\mathfrak{L},
\]
in other words, 
\[
  \label{IV.16.9.13.7}
  \mathfrak{I} \cap \mathfrak{L}^2 = \mathfrak{I}\mathfrak{L},
  \tag{16.9.13.7}
\]
which may also be expressed by saying that the canonical homomorphism
\[
  \mathfrak{I}/\mathfrak{I}\mathfrak{L} \to (\mathfrak{I} + \mathfrak{L}^2)/\mathfrak{L}^2
\]
is bijective.

\begin{proof}[Proof {\hyperlink{ega4.c04.prop.16-9-13}{(16.9.13)}}]
Having demonstrated the lemmas, let us prove the first claim of \sref{IV.16.9.13}:
It is clearly
\oldpage[IV-4]{51}
enough to prove that, at a point $x \in Y'$, the sequence of stalks of the sheaves appearing in \sref{IV.16.9.13.1} is exact.
Put $A = \sh{O}_{X,x}$. We may write $\sh{O}_{Y,x} = A' = A/\mathfrak{I}$, where $\mathfrak{I}$ is an ideal contained in the maximal ideal of $A$, and then $\sh{O}_{Y',x} = A'/\mathfrak{K}'$, where $\mathfrak{K}'$ is generated by an \emph{$A'$-regular} sequence of elements of $A'$ which are themselves images of the elements of an $A'$-regular sequence of elements of $A$ belonging to the maximal ideal of $A$.
If $\mathfrak{K}$ is the ideal generated by the latter elements and $\mathfrak{L} = \mathfrak{I} + \mathfrak{K}$, then $\sh{O}_{Y',x} = A/\mathfrak{L}$; since we are in the situation of \sref{IV.16.9.13.5}, the canonical homomorphism $\mathfrak{I}/\mathfrak{I}\mathfrak{L} \to (\mathfrak{I} + \mathfrak{L}^2)/\mathfrak{L}^2$ is bijective.
But this shows that the sequence
\[
  \xymatrix{
    0 \ar[r] & \mathfrak{I}/\mathfrak{I}\mathfrak{L} \ar[r] & \mathfrak{L}/\mathfrak{L}^2 \ar[r] &  (\mathfrak{L}/\mathfrak{I})/(\mathfrak{L}/\mathfrak{I})^2 \ar[r] & 0
  }
\]
is exact (see the proof of \sref{IV.16.2.7}), and the modules appearing in this sequence are precisely the stalks at $x$ of the sheaves in \sref{IV.16.9.13.1}.
The second claim follows from the fact that $\sh{N}_{Y'/Y}$ is a locally free $\sh{O}_{Y'}$-module \sref{IV.16.9.8} and Bourbaki, \emph{Alg.}, chap.~II, 3\textsuperscript{rd} ed., \textsection1, n\textsuperscript{o}11, prop.\ 21.
\end{proof}

\subsection{Differentially smooth morphisms.}
\label{IV.16.10}

\begin{definition}[16.10.1]
\label{IV.16.10.1}
We say that a morphism of preschemes $f: X \to S$ is differentially smooth (or that $X$ is differentially smooth over $S$) if it satisfies the following conditions:
\begin{enumerate}
  \item[\rm{(i)}] $\Omega_{X/S}^1$ is a locally projective $\sh{O}_X$-module, which is to say that every point $x \in X$ admits an affine open neighbourhood $U$ such that $\Gamma(U, \Omega_{X/S}^1)$ is a projective $\Gamma(U, \sh{O}_X)$-module (not necessarily of finite type).
  \item[\rm{(ii)}] The canonical morphism \sref{IV.16.3.1.1}
  \[
    \bb{S}_{\sh{O}_X}^\bullet(\Omega_{X/S}^1) \to \shGr_\bullet(\sh{P}_{X/S})
  \]
  is bijective.
\end{enumerate}

In particular, if $\Omega_{X/S}^1$ is locally free of finite rank, the $\sh{P}_{X/S}^n$ are locally free of finite rank $\sh{O}_X$-modules (being extensions of such modules).
\end{definition}

We say that $f$ is \emph{differentially smooth at a point $x \in X$} if there is an open neighbourhood $U$ of $x$ in $X$ such that $f|U$ is differentially smooth.

We will see later \sref{IV.17.12.4} that a smooth morphism is differentially smooth, which justifies the terminology;
but the converse is not true;
indeed, a \emph{monomorphism} $f:X \to S$ is differentially smooth, since $\Omega^1_{X/S} = 0$ because of \sref[I]{I.5.3.8}, and consequently the surjective homomorphism \sref{IV.16.3.1.1} is clearly bijective;
however a monomorphism is not even necessarily flat, neither \textit{a fortiori} smooth. 
Let us limit ourselves to proving the following proposition:

\begin{proposition}[16.10.2]
\label{IV.16.10.2}
Let $A$ be a ring, $B$ a formally smooth $A$-algebra for the discrete topology \sref[0]{0.19.3.1}.
Then $\Spec(B)$ is differentially smooth over $\Spec(A)$.
\end{proposition}

\begin{proof}
Indeed, $B \otimes_A B$ is then (with the discrete topology) a formally smooth $B$-algebra (by one or the other canonical homomorphisms $b \mapsto b \otimes 1$, $b \mapsto 1 \otimes b$ of $B$
\oldpage[IV-4]{52}
to $B \otimes_A B$
) \sref[0]{0.19.3.5}[(iii)];
therefore $B \otimes_A B$ is a formally smooth $A$-algebra with the discrete topology \sref[0]{0.19.3.5}[(ii)].
Taking $\mathfrak{I} = \mathfrak{I}_{B/A}$, it follows that $B \otimes_A B$ is a formally smooth $A$-algebra for the $\mathfrak{I}$-preadic topology \sref[0]{0.19.3.8};
since by hypothesis $B = (B \otimes_A B)/\mathfrak{I}$ is a formally smooth $A$-algebra for the discrete topology, the proposition follows from the equivalence of (a) and (b) in \sref[0]{0.19.5.4}.
\end{proof}

\begin{proposition}[16.10.3]
\label{IV.16.10.3}
For a morphism $f:X \to S$ to be differentially smooth, it is necessary and sufficient that for every $x \in X$, there is an affine open neighbourhood $U$ of $x$, with coordinate ring $A$, such that $\Gamma(U, \sh{P}_{X/S}^\infty)$ is an augmented topological $A$-algebra isomorphic to the completion $\widehat{B}$, where $B = \bb{S}_A(V)$, $V$ being a projective $A$-module and $B$ being endowed with the $B^+$-preadic topology (where $B^+$ is the augmentation ideal).
If $\Omega_{X/S}^1$ is locally free of finite rank, we can replace $\widehat{B}$ with the ring of formal series $A[[T_1, \dots, T_n]]$.
\end{proposition}

\begin{proof}
The notion of a differentially smooth morphism is clearly local on $X$, so we reduce to the case where $S = \Spec(B)$, $X = \Spec(C)$.
Consider $C \otimes_B C$ as a $C$-algebra (by the first factor);
put $\mathfrak{I} = \mathfrak{I}_{C/B}$ and endow $C \otimes_B C$ with the $\mathfrak{I}$-preadic topology;
we can apply to the topological $C$-algebra $C \otimes_B C$ and to the ideal $\mathfrak{I}$ of $C \otimes_B C$ the equivalence of (b) and (c) of \sref[0]{0.19.5.4}, given that $(C \otimes_B C)/\mathfrak{I} = C$ is clearly a formally smooth $C$-algebra for the discrete topologies.
The topology of $\Gamma(U, \sh{P}_{X/S}^\infty)$ is clearly the projective limit topology of this ring \sref{IV.16.1.11}.
\end{proof}

We note that the integer $n$ of the proposition \sref{IV.16.10.3} is the \emph{rank} of $\Omega_{X/S}^1$ in the point $x$.
We shall see \sref{IV.17.13.5} that when $f$ is differentially smooth and locally of finite type, $n$ is equal to the dimension of the fibre $f^{-1}(f(x))$. 

\begin{proposition}[16.10.4]
\label{IV.16.10.4}
Let $f:X \to S$, $g:S' \to S$ be two morphisms, and take $X' = X \times_S S'$, $f' = f_{(S')}:X' \to S'$.
\begin{enumerate}
  \item[\rm{(i)}] If $f$ is differentially smooth, the same is true for $f'$. 
  \item[\rm{(ii)}] Conversely, if $g$ is faithfully flat and quasi-compact, and if $f'$ is differentially smooth and $\Omega_{X'/S'}^1$ is a finite type $\sh{O}_{X'}$-module, $f$ is differentially smooth and $\Omega_{X/S}^1$ is a finite type $\sh{O}_X$-module.
\end{enumerate}
\end{proposition}

\begin{proof}
Indeed, if $f$ is differentially smooth, the $\shGr_n(\sh{P}_{X/S})$ are \emph{flat} $\sh{O}_X$-modules;
therefore, by \sref{IV.16.4.6}, the homomorphisms $\shGr_n(\sh{P}_{X/S}) \otimes_{\sh{O}_X} \sh{O}_{X'} \to \shGr_n(\sh{P}_{X'/S'})$ are bijective for every $n$. In view of the commutativity of diagram \sref{IV.16.2.1.3}, it follows from Definition \sref{IV.16.10.1} that $f'$ is differentially smooth.
On the other hand, if $g$ is faithfully flat and quasi-compact, it also follows from \sref{IV.16.4.6} that $\shGr_n(\sh{P}_{X/S}) \otimes_{\sh{O}_X} \sh{O}_{X'} \to \shGr_n(\sh{P}_{X'/S'})$ is bijective for every $n$.
Suppose also that $f'$ is differentially smooth and $\Omega_{X'/S'}^1$ is of finite rank.
Since the canonical projection $X' \to X$ is a faithfully flat and quasi-compact morphism, it follows first from \sref{IV.2.5.2} that $\Omega_{X/S}^1$ is an $\sh{O}_X$-module locally free of finite rank, then from \sref{IV.2.2.7} that the canonical homomorphism \sref{IV.16.3.1.1} is bijective, and therefore $f$ is differentially smooth.
\end{proof}

\begin{proposition}[16.10.5]
\label{IV.16.10.5}
For a morphism locally of finite type $f:X \to S$ to be differentially smooth, it is necessary and sufficient that the diagonal immersion $\Delta_f:X \to X \times_S X$ be quasi-regular.
\end{proposition}
  
\oldpage[IV-4]{53}

\begin{proof}
Being a local problem, we can reduce to the case where $S$ and $X$ are affines, and therefore the diagonal subprescheme of $X \times_S X$ is closed.
The hypothesis that $f$ is locally of finite type implies that $\Delta_f$ is locally of finite presentation \sref[I]{I.4.3.1}, therefore the diagonal prescheme of $X \times_S X$ is defined by an ideal $\sh{I}$ of finite type, and $\Omega_{X/S}^1 = \sh{I}/\sh{I}^2$ is an $\sh{O}_X$-module of finite type.
The proposition is now immediate from the comparison of the conditions of \sref{IV.16.10.1} and \sref{IV.16.9.4}. 
\end{proof}

\begin{remark}[16.10.6]
\label{IV.16.10.6}
Let $f:X \to S$ be a morphism such that the $\sh{O}_X$-module $\Omega_{X/S}^1$ is locally free of finite rank.
It follows from \sref[0]{0.20.4.7} that every $x \in X$ has an open neighbourhood $U$ such that there is a finite family $(z_\lambda)_{\lambda \in L}$ of sections of $\sh{O}_X$ over $U$ for which $(dz_\lambda)_{\lambda \in L}$ forms a \emph{basis} of the $\Gamma(U, \sh{O}_X)$-module $\Gamma(U, \Omega_{X/S}^1)$.
\end{remark}

\subsection{Differential operators on a differentially smooth $S$-prescheme}
\label{IV.16.11}

\begin{env}[16.11.1]
\label{IV.16.11.1}
Let $f:X \to S$ be a morphism, $U$ an open set of $X$, $(z_\lambda)_{\lambda \in L}$ a family of sections of $\sh{O}_X$ over $U$ such that the $dz_\lambda$ form a system of generators of $\Omega_{X/S}^1|U = \Omega_{U/S}^1$.
Let $m$ be an integer or the symbol $\infty$, and take, for every $\lambda$
\[
  \label{IV.16.11.1.1}
  \zeta_\lambda = \delta z_\lambda = d^m z_\lambda - z_\lambda \in \Gamma(U, \sh{P}_{X/S}^m).
  \tag{16.11.1.1}
\]

We will also use the usual notations from analysis;
for every $\bb{p} = (p_\lambda) \in \mathbb{N}^{(L)}$ (so that $p_\lambda = 0$ except for a finite number of indices), we put
\[
  \label{IV.16.11.1.2}
  |\bb{p}| = \sum_\lambda p_\lambda, \quad \bb{p}! = \prod_\lambda (p_\lambda!)
  \tag{16.11.1.2}
\]
\[
  \label{IV.16.11.1.3}
  \binom{\bb p}{\bb q} = \bb{p}!/(\bb{q}!(\bb{p} - \bb{q})!), \quad \text{for $\bb{p}$, $\bb{q}$ in $\mathbb{N}^{(L)}$, $\bb{q} \leq \bb{p}$}
  \tag{16.11.1.3}
\]
and we adopt the convention that $\binom{\bb p}{\bb q} = 0$ when $\bb{q} \not\leq \bb{p}$,
\[
  \label{IV.16.11.1.4}
  \bb{z}^\bb{p} = \prod_\lambda (z_\lambda)^{p_\lambda}, \quad {\pmb\zeta}^{\bb p} = \prod_\lambda (\zeta_\lambda)^{p_\lambda}.
  \tag{16.11.1.4}
\]

We therefore have, with this notation
\[
  \label{IV.16.11.1.5}
  d^m(\bb{z}^{\bb p}) = (d^m(\bb{z}))^{\bb p} = (\pmb{\zeta} + \bb{z})^{\bb p} = \sum_{\bb{q} \leq \bb{p}} \binom{\bb p}{\bb q} \bb{z}^{\bb{p} - \bb{q}} \pmb{\zeta}^{\bb q}
  \tag{16.11.1.5}
\]
\[
  \label{IV.16.11.1.6}
  \pmb{\zeta}^{\bb p} = (d^m \bb{z} - \bb{z})^{\bb p} =  \sum_{\bb{q} \leq \bb{p}} (-1)^{|\bb{p} - \bb{q}|} \binom{\bb p}{\bb q} \bb{z}^{\bb{p} - \bb{q}} d^m(\bb{z}^{\bb q}).
  \tag{16.11.1.6}
\]
  
Since the $dz_\lambda$ generate $\Omega_{X/S}^1$, are the images of the $\delta z_\lambda$, and the canonical morphism \sref{IV.16.3.1.1} is surjective, we conclude that, for finite $m$, the $\delta z_\lambda$ generate the $\sh{O}_U$-algebra $\sh{P}_{U/S}^m$ (Bourbaki, \emph{Alg. comm.}, chap.~III, \textsection2, n\textsuperscript{o}8, cor.\ 2 to th.\ 1).
Therefore the $\pmb{\zeta}^{\bb p}$ (for $|\bb{p}| \leq m$) generate the $\sh{O}_U$-module $\sh{P}_{U/S}^m$.
A differential operator $D \in \Diff_{U/S}^m$ is consequently entirely determined by the values $\langle \pmb{\zeta}^{\bb p}, D \rangle$ for $|\bb{p}| \leq m$, or, equivalently by \sref{IV.16.11.1.5} and \sref{IV.16.11.1.6}, by the values
\oldpage[IV-4]{54}
of $\langle d^m(\bb{z}^{\bb p}), D \rangle = D(\bb{z}^{\bb p})$ for $|\bb{p}| \leq m$;
more precisely, it follows from \sref{IV.16.11.1.5} that
\[
  \label{IV.16.11.1.7}
  D(\bb{z}^{\bb p}) = \langle d^m({\bb z}^{\bb p}), D \rangle = \sum_{\bb{q} \leq \bb{p}} \binom{\bb p}{\bb q} \langle \pmb{\zeta}^{\bb q} , D \rangle \bb{z}^{\bb{p} - \bb{q}}.
  \tag{16.11.1.7}
\]
\end{env}

\begin{theorem}[16.11.2]
\label{IV.16.11.2}
Let $f:X \to S$ be a morphism, $U$ an open set of $X$, $(z_\lambda)_{\lambda \in L}$ a family of sections of $\sh{O}_X$ over $U$ such that the family $(dz_\lambda)_{\lambda \in L}$ generates $\Omega_{X/S}^1|U = \Omega_{U/S}^1$.
The following conditions are equivalent:
\begin{enumerate}
  \item[(a)] $f|U$ is differentially smooth and $(dz_\lambda)$ is a basis of the $\sh{O}_U$-module $\Omega_{U/S}^1$.
  \item[(b)] There is a family $(D_{\bb p})_{\bb{p} \in \mathbb{N}^{(L)}}$ of differential operators of $\sh{O}_U$ to itself verifying the conditions
  \[
    \label{IV.16.11.2.1}
    D_{\bb p} (\bb{z}^{\bb q}) = \binom{\bb q}{\bb p} \bb{z}^{\bb{q} - \bb{p}}, \quad (\bb{p}, \bb{q} \text{ in } \mathbb{N}^{(L)}).
    \tag{16.11.2.1}
  \]
\end{enumerate}

Also, when these conditions are satisfied, the family $(D_\bb{p})$ is uniquely determined by the conditions \sref{IV.16.11.2.1} and satisfies the relations
\[
  \label{IV.16.11.2.2}
  D_{\bb q} \circ D_{\bb p} = D_{\bb p} \circ D_{\bb q} = \frac{(\bb{p} + \bb{q})!}{\bb{p}!\bb{q}!}D_{\bb{p} + \bb{q}} \quad (\bb{p}, \bb{q} \text{ in }  \mathbb{N}^{(L)}).
  \tag{16.11.2.2}
\]

Finally, if $L$ is finite, for every integer $m$, the $D_{\bb p}$ such that $|\bb{p}| \leq m$ form a basis of the $\sh{O}_U$-module $\shDiff_{U/S}^m$, in other words, every differential operator of order $\leq m$ on $U$ can be written uniquely as
\[
  D = \sum_{|\bb{p}| \leq m} a_{\bb p} D_{\bb p}
\] 
where the $a_{\bb p}$ are sections of $\sh{O}_X$ over $U$.
\end{theorem}

\begin{proof}
Note first that because of \sref{IV.16.11.1.6} and \sref{IV.16.11.1.5}, we verify immediately that the conditions \sref{IV.16.11.2.1} are equivalent to 
\[
  \label{IV.16.11.2.3}
  \langle \pmb{\zeta}^{\bb p}, D_{\bb q} \rangle = \delta_{\bb{p} \bb{q}} \quad (\text{Kronecker's symbol}).
  \tag{16.11.2.3}
\]

The existence of the family $(D_\bb{p})$ verifying these conditions implies first (by taking $|\bb{p}| = 1$) that the $dz_\lambda$ are linearly independent, and therefore form a basis of the $\sh{O}_U$-module $\Omega_{U/S}^1$.
Then, for every integer $m \geq 1$, we deduce similarly from \sref{IV.16.11.2.3} that the $\pmb{\zeta}^{\bb p}$ such that $|\bb{p}| \leq m$ are linearly independent;
it follows that the canonical homomorphism \sref{IV.16.3.1.1} is injective, and therefore bijective, which proves that (b) implies (a).
The converse follows immediately from Definition \sref{IV.16.10.1}: the fact that the $\pmb{\zeta}^{\bb p}$ form a basis of $\sh{P}_{U/S}^m$ for $|\bb{p}| \leq m$ implies the existence and uniqueness of a family of homomorphisms $u_{\bb{q},m}: \sh{P}_{U/S}^m \to \sh{O}_U$ ($|\bb{q}| \leq m$) such that $\langle \pmb{\zeta}^{\bb p},u_{\bb{q},m}\rangle = \delta_{\bb p\bb q}$ for $|\bb{p}| \leq m$, $|\bb{q}| \leq m$.
For a given value of $\bb{q}$, the differential operators corresponding to $u_{\bb{q}, m}$ for $m \geq |\bb{q}|$ are identified with the same operator $D_\bb{q}$.
This proves that (a) implies (b), and also that the family $(D_\bb{q})$ is uniquely determined and that, if $L$ is finite, the $D_\bb{p}$ with $|\bb{p}| \leq m$ form a basis of the $\sh{O}_U$-module $\shDiff_{U/S}^m$, the dual of $\sh{P}_{U/S}^m$.
Finally, the relations \sref{IV.16.11.2.2} follow immediately by evaluating the three operators in question on the $\bb{z}^\bb{r}$ and using the fact that the $\pmb{\zeta}^\bb{r}$ for $|\bb{r}| \leq m$ generate $\sh{P}_{U/S}^m$.
\end{proof}

\oldpage[IV-4]{55}

\begin{remarks}[16.11.3]
\label{IV.16.11.3}
\begin{enumerate}
  \item[\rm{(i)}] The fact that, by \sref{IV.16.11.2.2}, the $D_\bb{p}$ commute pairwise naturally does not imply that the $\sh{O}_U$-algebra $\shDiff_{U/S}$ is commutative: the $D_\bb{p}$ do not commute with multiplication by sections of $\sh{O}_U$ unless $\bb{p} = 0$.
  \item[\rm{(ii)}] The indices $\bb{p}$ such that $|\bb{p}| = 1$ are the $\pmb{\epsilon}_\lambda =  (\epsilon_{\lambda\mu})_{\mu \in L}$ where $\epsilon_{\lambda\mu} = 0$ if $\mu \neq \lambda$ and $\epsilon_{\lambda\lambda} = 1$;
  when $L$ is finite, the operators $D_{\pmb{\epsilon}_\lambda}$ are exactly the $S$-derivations $D_\lambda$ introduced in \sref{IV.16.5.7}.
  We note that in general (contrary to what happens in classical analysis), it is not true that a differential operator of any order can be written as a linear combination of powers of $D_i$ (cf. \sref{IV.16.12}). 
  \item[\rm{(iii)}] For every integer $r \geq 1$, we can define the notion of \emph{differentially smooth up to order $r$} by replacing in \sref{IV.16.10.1} the condition (ii) by the condition that the homomorphisms
  \[
    \bb{S}_{\sh{O}_X}^m(\Omega_{X/S}^1) \to \shGr_m(\sh{P}_{X/S})
  \]
  are bijective \emph{for every $m \leq r$}.
  The argument of \sref{IV.16.11.2} proves also that if, in the condition (a), we replace ``differentially smooth'' by ``differentially smooth up to order $r$'', this condition is equivalent to (b) by restricting ourselves to $\bb{p} \in \mathbb{N}^{(L)}$, $\bb{q} \in \mathbb{N}^{(L)}$  such that $|\bb{p}| \leq r$, $|\bb{q}| \leq r$.
\end{enumerate}
\end{remarks}

\subsection{Case of characteristic zero: Jacobian criterion for differentially smooth morphisms.}
\label{IV.16.12}

\begin{env}[16.12.1]
\label{IV.16.12.1}
We say that a prescheme $X$ is \emph{of characteristic $p$} ($p$ equal to zero or a prime number) if, for every affine open set $U$ of $X$, the ring $\Gamma(U, \sh{O}_X)$ is of characteristic $p$ \sref[0]{0.21.1.1}.
It follows from \sref[0]{0.21.1.3} that for $X$ to be of characteristic $0$, it is necessary and sufficient that for every \emph{closed} point $x$ of $X$, the residue field $\kres(x)$ is of characteristic $0$, or even that $X$ can be given a structure of $\mathbb{Q}$-prescheme (necessarily unique).
\end{env}

\begin{theorem}[16.12.2]
\label{IV.16.12.2}
Let $X$ be a prescheme of characteristic $0$, $f:X \to S$ a morphism.
If $\Omega_{X/S}^1$ is a locally free $\sh{O}_X$-module (not necessarily of finite type), $f$ is differentially smooth.
\end{theorem}

\begin{proof}
The problem being local on $X$, we can suppose there is a family $(z_\lambda)$ of sections of $\sh{O}_X$ over $X$ such that the $(dz_\lambda)$ is a basis for the $\sh{O}_X$-module $\Omega_{X/S}^1$.
Applying the criterion \sref{IV.16.11.2}, it is enough for the operators 
\[
  D_\bb{p} = (\bb{p}!)^{-1} \prod_\lambda D_\lambda^{p_\lambda}
\]
(where the $D_\lambda$ are the coordinate forms corresponding to the basis $(dz_\lambda)$) to verify the relations \sref{IV.16.11.2.1}, which is a consequence of the fact that the $D_\lambda$ are derivations.
\end{proof}

\begin{env}[16.12.3]
\label{IV.16.12.3}
The theorem above is not true if we discard the hypothesis that $X$ is of characteristic $0$.
For example, if $S = \Spec(k)$, where $k$ is a field of characteristic $p > 0$, $X = \Spec(K)$ where $K = k(\alpha)$ where $\alpha \notin k$, $\alpha^p \in k$, we verify immediately 
\oldpage[IV-4]{56}
that $\Omega_{X/S}^1$ has rank $1$, and that the morphism $X \to S$ is differentially smooth up to order $p - 1$ \sref{IV.16.11.3}[(iii)], but not up to order $p$.
However, the proof of \sref{IV.16.12.2} proves that if $\Omega_{X/S}^1$ is locally free, and if $n! 1_{\sh{O}_X}$ is invertible in $\Gamma(X, \sh{O}_X)$, then $X$ is differentially smooth over $S$ up to order $n$.
\end{env}
% END NATIVE INPUT source_aligned/ega4-16
% BEGIN NATIVE INPUT source_aligned/ega4-17
\section{Smooth morphisms, unramified (or net) morphisms, and \'etale morphisms.}
\label{section:IV.17}

In this section, we return to the notions studied in \sref[0]{0.19}, expressed in the geometric language of schemes and from the global point of view, for preschemes locally of finite presentation over a given base prescheme.

Most of the results (except those of \sref{IV.17.7}, \sref{IV.17.8}, \sref{IV.17.9}, \sref{IV.17.13}, and \sref{IV.17.16}) reduce, moreover, to variants of properties already encountered in \sref[0]{0.19}.

For more specific results on \'etale morphisms, the reader should consult \sref{IV.18}.

\subsection{Formally smooth morphisms, formally unramified morphisms, formally \'etale morphisms.}
\label{subsection:IV.17.1}

\begin{definition}[17.1.1]
\label{IV.17.1.1}
Let $f:X\to Y$ be a morphism of preschemes.
We say that $f$ is \emph{formally smooth} (resp. \emph{formally unramified}, resp. \emph{formally \'etale}) if, for all affine schemes $Y'$, all closed subschemes $Y_0'$ of $Y'$ defined by a nilpotent ideal $\sh{J}$ of $\sh{O}_{Y'}$, and every morphism $Y'\to Y$, the map 
\[
\label{IV.17.1.1.1}
  \Hom_Y(Y',X)\to\Hom_Y(Y_0',X)
  \tag{17.1.1.1}
\]
induced by the canonical map $Y_0'\to Y'$, is \emph{surjective} (resp. \emph{injective}, resp. \emph{bijective}).

One also says that $X$ is \emph{formally smooth} (resp. \emph{formally unramified}, resp. \emph{formally \'etale}) over $Y$.

It is clear that for $f$ to be formally \'etale, it is necessary and sufficient for $f$ to be formally smooth and formally unramified.
\end{definition}

\begin{remark}[17.1.2]
\label{IV.17.1.2}
\medskip\noindent
\begin{enumerate}
  \item[(i)] Suppose that $Y=\Spec(A)$ and $X=\Spec(B)$ are affine, so that $f$ comes from a homomorphism of rings $\vphi:A\to B$. 
    According to \sref[0]{0.19.3.1} and \sref[0]{0.19.10.1}, saying that $f$ is formally smooth (resp. formally unramified, resp. formally \'etale) means that, via $\vphi$, $B$ is a \emph{formally smooth} (resp. \emph{formally unramified}, resp. \emph{formally \'etale}) $A$-algebra, for the \emph{discrete} topologies on $A$ and $B$.
  \item[(ii)] To verify that $f$ is formally smooth (resp. formally unramified, resp. formally \'etale), we can, in Definition~\sref{IV.17.1.1}, restrict to the case where $\sh{J}^2=0$.
  To see this, if $f$ satisfies the corresponding condition of Definition~\sref{IV.17.1.1} in the particular case $\sh{J}^2=0$, and if we have $\sh{J}^n=0$, then we consider the closed subscheme $Y_j'$ of $Y'$ defined by the sheaf of ideals $\sh{J}^{j+1}$ for $0\leq j\leq n-1$, so that $Y_j'$ is a closed subscheme of $Y_{j+1}'$ defined by a square-zero sheaf of ideals;
the hypotheses imply that each of the maps
\[
  \Hom_Y(Y_{j+1}',X)\to\Hom_Y(Y_j',X)\quad(0\leq j\leq n-1) 
\]
\oldpage[IV]{57}
is surjective (resp. injective, resp. bijective);
by composition, we conclude that the same holds for \sref{IV.17.1.1.1}.
\item[(iii)] Note that the properties of the morphism $f$ defined in \sref{IV.17.1.1} are properties of the \emph{representable functor} \sref[\textbf{0}\textsubscript{III}]{0.8.1.8}
\[
  Y'\mapsto\Hom_Y(Y',X) 
\]
from the category of $Y$-preschemes to the category of sets;
they keep a meaning for \emph{any} contravariant functor with the same domain and codomain, representable or not.
\item[(iv)] Assume that the morphism $f$ is formally unramified (resp. formally \'etale);
consider an \emph{arbitrary} $Y$-prescheme $Z$ and a closed subprescheme $Z_0$ of $Z$ defined by a \emph{locally nilpotent} sheaf of ideals $\sh{J}$ of $\sh{O}_Z$. 
Then the map
\[
\label{IV.17.1.2.1}  
  \Hom_Y(Z,X)\to\Hom_Y(Z_0,X)
  \tag{17.1.2.1}
\]
induced by the canonical injection $Z_0\to Z$, is still injective (resp. bijective).
To see this, let $(U_\alpha)$ be an affine open covering of $Z$ such that the sheaves of ideals $\sh{J}|U_\alpha$ are nilpotent, and for each $\alpha$, let $U_\alpha^0$ be the inverse image of $U_\alpha$ in $Z_0$, which is the closed subprescheme of $U_\alpha$ defined by $\sh{J}|U_\alpha$.
Let $f_0:Z_0\to X$ be a $Y$-morphism;
by hypothesis, for each $\alpha$, there is at most one (resp. one and only one) $Y$-morphism $f_\alpha:U_\alpha\to X$ whose restriction to $Z_0$ coincides with $f_0|U_\alpha$.
We immediately conclude that if $f_\alpha$ and $f_\beta$ are defined, then, for each affine open $V\subset U_\alpha\cap U_\beta$, we have $f_\alpha|V=f_\beta|V$, as the restrictions of these morphisms to the inverse image $V_0$ of $V$ in $Z_0$ coincide.
There is therefore at most one (resp. one and only one) $Y$-morphism $f:Z\to X$ whose restriction to $Z_0$ coincides with $f_0$.
\end{enumerate}
\end{remark}

\begin{proposition}[17.1.3]
\label{IV.17.1.3}
\medskip\noindent
\begin{enumerate}
  \item[{\rm(i)}] A monomorphism of preschemes is formally unramified;
    an open immersion is formally \'etale.
  \item[{\rm(ii)}] The composition of two formally smooth (resp. formally unramified, resp. formally \'etale) morphisms is formally smooth (resp. formally unramified, resp. formally \'etale).
  \item[{\rm(iii)}] If $f:X\to Y$ is a formally smooth (resp. formally unramified, resp. formally \'etale) $S$-morphism, then so is $f_{(S')}:X_{(S')}\to Y_{(S')}$ for any base extension $S'\to S$.
  \item[{\rm(iv)}] If $f:X\to X'$ and $g:Y\to Y'$ are two formally smooth (resp. formally unramified, resp. formally \'etale) $S$-morphisms, then so is $f\times_S g:X\times_S Y\to X'\times_S Y'$.
  \item[{\rm(v)}] Let $f:X\to Y$ and $g:Y\to Z$ be two morphisms;
  if $g\circ f$ is formally unramified, then so is $f$.
  \item[{\rm(vi)}] If $f:X\to Y$ is a formally unramified morphism, then so is $f_\red:X_\red\to Y_\red$.
\end{enumerate}
\end{proposition}

\begin{proof}
According to \sref[I]{I.5.5.12}, it suffices to prove (i), (ii), and (iii). 
The assertions in (i) are both trivial.
To prove (ii), consider two morphisms $f:X\to Y$, $g:Y\to Z$, an affine scheme $Z'$, a closed subscheme $Z_0'$ of $Z'$ defined by a nilpotent ideal, and a morphism $Z'\to Z$.
Suppose that $f$ and $g$ are formally smooth, and consider a $Z$-morphism
\oldpage[IV]{58}
$u_0:Z_0'\to X$;
the hypothesis on $g$ implies that there exists a $Z$-morphism $v:Z'\to Y$ such that $f\circ u_0=v\circ j$ (where $j:Z_0'\to Z'$ is the canonical injection);
the hypothesis on $f$ then implies that there exists a morphism $u:Z'\to X$ such that $f\circ u=v$ and $u\circ j=u_0$, therefore $(g\circ f)\circ u$ is equal to the given morphism $Z'\to Z$ and $u\circ j=u_0$, which proves that $g\circ f$ is formally smooth;
we argue the same way when we suppose that $f$ and $g$ are formally unramified.

Finally, to prove (iii), let $X'=X_{S'}$, $Y'=Y_{S'}$, $f'=f_{S'}$;
consider an affine scheme $Y''$, a closed subscheme $Y_0''$ defined by a nilpotent sheaf of ideals, and a morphism $g:Y''\to Y'$ making $Y''$ a $Y'$-prescheme;
we then know by \sref[I]{I.3.3.8} that $\Hom_{Y'}(Y'',X')$ is canonically identified with $\Hom_Y(Y'',X)$, and $\Hom_{Y'}(Y_0'',X')$ with $\Hom_Y(Y_0'',X)$, and the conclusion follows immediately from Definition~\sref{IV.17.1.1}.
\end{proof}

We note that a \emph{closed immersion} is not necessarily formally smooth.
\begin{proposition}[17.1.4]
\label{IV.17.1.4}
Let $f:X\to Y$ and $g:Y\to Z$ be two morphisms, and suppose that $g$ is formally unramified.
Then, if $g\circ f$ is formally smooth (resp. formally \'etale), so is $f$.
\end{proposition}

\begin{proof}
Let $Y'$ be an affine scheme, $Y_0'$ a closed subscheme of $Y'$ defined by a nilpotent sheaf of ideals, $h:Y'\to Y$ a morphism, $j:Y_0'\to Y'$ the canonical injection, and $u_0:Y_0'\to X$ a $Y$-morphism, so that $f\circ u_0=h\circ j$.
Suppose that $g\circ f$ is formally smooth;
then there exists a morphism $u:Y'\to X$ such that $u\circ j=u_0$ and $(g\circ f)\circ u=g\circ h$. 
But these two relations imply that $f\circ u$ and $h$ are $Z$-morphisms from $Y'$ to $Y$ such that $(f\circ u)\circ j=h\circ j$;
by virtue of the hypothesis that $g$ is formally unramified, we get that $f\circ u=h$, in other words that $u$ is a $Y$-morphism;
thus $f$ is formally smooth.
Taking into account \sref{IV.17.1.3}[(v)], this proves the proposition.
\end{proof}

\begin{corollary}[17.1.5]
\label{IV.17.1.5}
Suppose that $g$ is formally \'etale;
then, for $g\circ f$ to be formally smooth (resp. formally unramified, resp. formally \'etale), it is necessary and sufficient that $f$ is.
\end{corollary}

\begin{proof}
This follows from \sref{IV.17.1.4} and \sref{IV.17.1.3}[(ii) and (v)].
\end{proof}

\begin{proposition}[17.1.6]
\label{IV.17.1.6}
Let $f:X\to Y$ be a morphism of preschemes.
\begin{enumerate}
  \item[{\rm(i)}] Let $(U_\alpha)$ be an open covering of $X$ and, for each $\alpha$, let $i_\alpha:U_\alpha\to X$ be the canonical injection.
    For $f$ to be formally smooth (resp. formally unramified, resp. formally \'etale), it is necessary and sufficient that each $f\circ i_\alpha$ is.
  \item[{\rm(ii)}] Let $(V_\lambda)$ be an open covering of $Y$.
    For $f$ to be formally smooth (resp. formally unramified, resp. formally \'etale), it is necessary and sufficient that each of the restrictions $f^{-1}(V_\lambda)\to V_\lambda$ of $f$ is.
\end{enumerate}
\end{proposition}

\begin{proof}
First note that (ii) is a consequence of (i): if $j_\lambda:V_\lambda\to Y$ and $i_\lambda:f^{-1}(V_\lambda)\to X$ are the canonical injections, then the restriction $f_\lambda:f^{-1}(V_\lambda)\to V_\lambda$ of $f$ is such that $j_\lambda\circ f_\lambda=f\circ i_\lambda$;
if $f$ is formally smooth (resp. formally unramified), then so is $f\circ i_\lambda$ since $i_\lambda$ is formally \'etale \sref{IV.17.1.3};
but since $j_\lambda$ is formally \'etale, this means that $f_\lambda$ is formally smooth (resp. formally unramified), by virtue of \sref{IV.17.1.5}.
Conversely, if all the $f_\lambda$ are formally smooth (resp. formally unramified), the same applies to $j_\lambda\circ f_\lambda$ \sref{IV.17.1.3}, so also to $f$ in virtue of (i).

\oldpage[IV]{59}
If we take into account that the $i_\alpha$ are formally \'etale, everything comes down to proving that if the $f\circ i_\alpha$ are formally smooth (resp. formally unramified), then the same applies to $f$.

Therefore let $Y'$ be an affine scheme, $Y_0'$ a closed subscheme of $Y'$ defined by a nilpotent ideal $\sh{J}$, which we may assume to satisfy $\sh{J}^2=0$ \sref{IV.17.1.2}[(ii)], and finally let $g:Y'\to Y$ be a morphism. 
Suppose we are given a $Y$-morphism $u_0:Y_0'\to X$;
denote by $W_\alpha$ (resp. $W_\alpha^0$) the prescheme induced by $Y'$ (resp. $Y_0'$) on the open subset $u_0^{-1}(U_\alpha)$ (we recall that $Y'$ and $Y_0'$ share the \emph{same underlying topological space}). 
Let us first suppose that the $f\circ i_\alpha$ are \emph{formally unramified}, and show that, if $u'$ and $u''$ are two $Y$-morphisms from $Y'$ to $X$ whose restrictions to $Y_0'$ coincide, then we have $u'=u''$. 
Indeed, taking into account \sref{IV.17.1.2}[(iv)], the hypothesis that the $f\circ i_\alpha$ are formally unramified implies that for all $\alpha$, we have $u'|W_\alpha=u''|W_\alpha$, since the restrictions of both $Y$-morphisms to $W_\alpha^0$ coincide. 
Hence the conclusion follows.

Now suppose that the $f\circ i_\alpha$ are \emph{formally smooth} and prove the existence of a $Y$-morphism $u:Y'\to X$ whose restriction to $Y_0'$ is $u_0$.
Now, since $Y'$ is an \emph{affine scheme}, we can apply \sref{IV.16.5.17}, the hypotheses of which are satisfied, and the conclusion of which precisely proves the existence of $u$.
\end{proof}

We can therefore say that the notions introduced in \sref{IV.17.1.1} are \emph{local} on $X$ and $Y$, which always allows, in virtue of \sref{IV.17.1.2}[(i)], to be reduced to the study of formally smooth (resp. formally unramified, resp. formally \'etale) \emph{algebras}.

\subsection{General properties of differentials}
\label{subsection:IV.17.2}

\begin{proposition}[17.2.1]
\label{IV.17.2.1}
For a morphism $f:X\to Y$ to be formally unramified, it is necessary and sufficient that $\Omega_f^1=0$ (what we still write $\Omega_{X/Y}^1=0$ \sref{IV.16.3.1}).  
\end{proposition}

\begin{proof}
Taking into account \sref{IV.17.1.6}, we reduce to the case where $Y=\Spec(A)$ and $X=\Spec(B)$ are affine, and the conclusion then follows from \sref[0]{0.20.7.4} and the interpretation of $\Omega_{X/Y}^1$ in this case \sref{IV.16.3.7}.
\end{proof}

\begin{corollary}[17.2.2]
\label{IV.17.2.2}
Let $f:X\to Y$ and $g:Y\to Z$ be two morphisms. 
For $f$ to be formally unramified, it is necessary and sufficient that the canonical morphism \sref{IV.16.4.19}
\[
  f^*(\Omega^1_{Y/Z})\to \Omega^1_{X/Z}
\] is surjective.
\end{corollary}

\begin{proof}
This is an immediate consequence of \sref{IV.17.2.1} and the exact sequence \sref{IV.16.4.19.1}.
\end{proof}

\begin{proposition}[17.2.3]
\label{IV.17.2.3}
Let $f:X\to Y$ be a formally smooth morphism.
\begin{enumerate}
  \item[{\rm(i)}] The $\sh{O}_X$-module $\Omega_{X/Y}^1$ is locally projective \sref{IV.16.10.1}.
    If $f$ is locally of finite type, then $\Omega_{X/Y}^1$ is locally free and of finite type.
  \item[{\rm(ii)}] For all morphisms $g:Y\to Z$, the sequence \sref{IV.16.4.19} of $\sh{O}_X$-modules
\[
\label{IV.17.2.3.1}
  0\to f^*(\Omega_{Y/Z}^1)\to\Omega_{X/Z}^1\to\Omega_{X/Y}^1\to 0
  \tag{17.2.3.1}
\]
is exact; moreover, for each $x\in X$, there exists an open neighbourhood $U$ of $x$ such that the restrictions to $U$ of the homomorphisms in \sref{IV.17.2.3.1} form a \emph{split} exact sequence.
\end{enumerate}
\end{proposition}

\oldpage[IV]{60}
\begin{proof}
\medskip\noindent
\begin{enumerate}
  \item[(i)] We know \sref{IV.16.3.9} that if $f$ is locally of finite type, then $\Omega_f^1$ is an $\sh{O}_X$-module of finite type. 
    To prove that, in all cases, it is locally projective, we can reduce, by virtue of \sref{IV.17.1.6}, to the case where $Y=\Spec(A)$ and $X=\Spec(B)$ are affine, and the result follows from the hypothesis on $f$ and from \sref[0]{0.20.4.9} and \sref[0]{0.19.2.1}.
  \item[(ii)] Again, we can restrict to the case where $X$, $Y$, and $Z$ are affine \sref{IV.17.1.6}, and the conclusion in this case follows from the interpretation of the sheaves of modules in the sequence \sref{IV.17.2.3.1} and from \sref[0]{0.20.5.7}.
\end{enumerate}
\end{proof}

\begin{corollary}[17.2.4]
\label{IV.17.2.4}
If $f:X\to Y$ is formally \'etale, then, for all morphisms $g:Y\to Z$, the canonical homomorphism of $\sh{O}_X$-modules
\[
  f^*(\Omega_{Y/Z}^1)\to\Omega_{X/Z}^1
\]
is bijective.
\end{corollary}

\begin{proof}
This follows from the exactness of the sequence \sref{IV.17.2.3.1} and from the fact that we then have $\Omega_{X/Y}^1=0$ \sref{IV.17.2.1}.
\end{proof}

\begin{proposition}[17.2.5]
\label{IV.17.2.5}
Let $f:X\to Y$ be a morphism, $X'$ a subprescheme of $X$ such that the composite morphism $X'\xrightarrow{j}X\xrightarrow{f}Y$ (where $j$ is the canonical injection) is formally smooth.
Then the sequence of $\sh{O}_{X'}$-modules \sref{IV.16.4.21}
\[
\label{IV.17.2.5.1}
  0\to\sh{N}_{X'/X}\to\Omega_{X/Y}^1\otimes_{\sh{O}_X}\sh{O}_{X'}\to\Omega_{X'/Y}^1\to 0
\tag{17.2.5.1}
\]
is exact; moreover, for each $x\in X$, there exists an open neighbourhood $U$ of $x$ such that the restrictions to $U$ of the homomorphisms in \sref{IV.17.2.5.1} form a \emph{split} exact sequence.
\end{proposition}

\begin{proof}
By virtue of \sref{IV.17.1.6}, we reduce to the case where $Y=\Spec(A)$ and $X=\Spec(B)$ are affine, and $X'=\Spec(B/\mathfrak{J})$, where $\mathfrak{J}$ is an ideal of $B$. 
The conormal sheaf $\sh{N}_{X'/X}$ then corresponds to the $B$-module $\mathfrak{J}/\mathfrak{J}^2$ \sref{IV.16.1.3}, and the conclusion follows from \sref[0]{0.20.5.14}.
\end{proof}

\begin{proposition}[17.2.6]
\label{IV.17.2.6}
Let $X$ and $Y$ be two preschemes, $f:X\to Y$ a morphism locally of finite type. 
The following conditions are equivalent:
\begin{enumerate}
  \item[{\rm(a)}] $f$ is a monomorphism.
  \item[{\rm(b)}] $f$ is radicial and formally unramified.
  \item[{\rm(c)}] For each $y\in Y$, the fibre $f^{-1}(y)$ is empty or $\kres(y)$-isomorphic to $\Spec(\kres(y))$ (in other words, it is reduced to a single point $z$ such that $\kres(y)\to\sh{O}_z/\mathfrak{m}_y\sh{O}_z$ is an isomorphism).
\end{enumerate}
\end{proposition}

\begin{proof}
The fact that (a) implies (c) follows from \sref{IV.8.11.5.1}. 
It is clear that (c) implies that $f$ is radicial;
let us prove that it also follows from (c) that $\Omega_{X/Y}^1=0$, which will prove that (c) implies (b) \sref{IV.17.2.1}.
Note that the $\sh{O}_X$-module $\Omega_{X/Y}^1$ is quasi-coherent of finite type \sref{IV.16.3.9}.
It follows from \sref[I]{I.9.1.13.1} that, for $(\Omega_{X/Y}^1)_x=0$, it is necessary and sufficient that if we set $Y_1=\Spec(\kres(y))$, $X_1=f^{-1}(y)=X\times_Y Y_1$, then we have $(\Omega_{X_1/Y_1}^1)_x=0$;
but as the morphism $f_1:X_1\to Y_1$ induced by $f$ is formally unramified by virtue of the hypothesis (c) \sref{IV.17.1.3}, the conclusion follows from \sref{IV.17.2.1}.
Finally, let us prove that (b) implies (a);
for this, consider the diagonal morphism $g=\Delta_f:X\to X\times_Y X$;
since $f$ is radicial, $g$ is surjective \sref[I]{I.8.7.1};
on the other hand, $\Omega_{X/Y}^1$ is by definition the conormal sheaf $\shGr_1(g)$ of the immersion $g$ \sref{IV.16.3.1}, and to say that $f$ is formally unramified therefore means that
\oldpage[IV]{61}
$\shGr_1(g)=0$ \sref{IV.17.2.1}. 
In addition, $g$ is locally of finite presentation \sref[I]{I.4.3.1};
therefore the hypothesis $\shGr_1(g)=0$ implies that $g$ is an open immersion \sref{IV.16.1.10};
being surjective, this immersion is an isomorphism, hence $f$ is a monomorphism \sref[I]{I.5.3.8}.
\end{proof}

\subsection{Smooth morphisms, unramified morphisms, \'etale morphisms}
\label{subsection:IV.17.3}

\begin{definition}[17.3.1]
\label{IV.17.3.1}
We say that a morphism $f:X\to Y$ is \emph{smooth} (resp. \emph{unramified}, or \emph{net}
\footnote{The words ``net'' and ``formally net'' seem preferable to the established terminology ``unramified'' (resp. ``formally unramified'') and will be used almost exclusively beginning with Chapter~V.
In this chapter, we have retained the old terminology so as not to conflict with \sref[0]{0.19.10}.}
resp. \emph{\'etale})
if it is locally of finite presentation and formally smooth (resp. formally unramified, resp. formally \'etale).
\end{definition}

We then also say that $X$ is \emph{smooth} (resp. \emph{unramified}, resp. \emph{\'etale}) \emph{over $Y$}.

We will see later \sref{IV.17.5.2} that this definition of a smooth morphism coincides with the definition already given in \sref{IV.6.8.1};
until then, we will exclusively use definition \sref{IV.17.3.1}.

It is clear that saying that $f$ is \'etale means that it is \emph{both} smooth and unramified.

\begin{remarks}[17.3.2]
\label{IV.17.3.2}
\medskip\noindent
\begin{enumerate}
  \item[(i)] Note that definition \sref{IV.17.3.1} can be phrased using only the functor
    \[
      Y'\mapsto\Hom_Y(Y',X)
    \]
    considered in \sref{IV.17.1.2}[(iii)] because to say that $f$ is locally of finite presentation is equivalent to saying that the preceding functor \emph{commutes with projective limits of affine schemes} \sref{IV.8.14.2}.
  \item[(ii)] Let $A$ be a ring and $B$ an $A$-algebra. 
   We say that $B$ is a \emph{smooth} (resp. \emph{unramified}, resp. \emph{\'etale}) $A$-algebra if the corresponding morphism $\Spec(B)\to\Spec(A)$ is smooth (resp. unramified, resp. \'etale).
   It is equivalent to saying that $B$ is an $A$-algebra \emph{of finite presentation} \sref[I]{I.4.4.6} that is furthermore formally smooth (resp. formally unramified, resp. formally \'etale) for the discrete topologies.
  \item[(iii)] It follows from \sref{IV.17.1.6} and the definition of a morphism locally of finite presentation \sref[I]{I.4.4.2} that the notion of a smooth (resp. unramified, resp. \'etale) morphism is \emph{local on $X$ and on $Y$}.
\end{enumerate}
\end{remarks}

\begin{proposition}[17.3.3]
\label{IV.17.3.3}
\medskip\noindent
\begin{enumerate}
  \item[{\rm(i)}] An open immersion is \'etale.
    For an immersion to be unramified, it is necessary and sufficient for it to be locally of finite presentation.
  \item[{\rm(ii)}] The composition of two smooth (resp. unramified, resp. \'etale) morphisms is smooth (resp. unramified, resp. \'etale).
  \item[{\rm(iii)}] If $f:X\to Y$ is a smooth (resp. unramified, resp. \'etale) $S$-morphism, then so is $f_{(S')}:X_{(S')}\to Y_{(S')}$ for any base extension $S'\to S$. 
  \item[{\rm(iv)}] If $f:X\to X'$ and $g:Y\to Y'$ are smooth (resp. unramified, resp. \'etale) $S$-morphisms, then so is $f\times_S g:X\times_S Y\to X'\times_S Y'$.
\oldpage[IV]{62}
  \item[{\rm(v)}] Let $f:X\to Y$ and $g:Y\to Z$ be two morphisms;
    if $g$ is locally of finite type and if $g\circ f$ is unramified, then $f$ is unramified.
\end{enumerate}
\end{proposition}

\begin{proof}
This follows from \sref[I]{I.4.3} and \sref{IV.17.1.3}.
\end{proof}

\begin{proposition}[17.3.4]
\label{IV.17.3.4}
Let $f:X\to Y$ and $g:Y\to Z$ be two morphisms, and suppose that $g$ is unramified.
Then, if $g\circ f$ is smooth (resp. unramified, resp. \'etale), so is $f$.
\end{proposition}

\begin{proof}
As $g$ and $g\circ f$ are locally of finite presentation, so is $f$ \sref[I]{I.4.4.3}[(v)];
the conclusion thus follows from \sref{IV.17.1.4} and \sref{IV.17.1.3}[(v)].
\end{proof}

\begin{corollary}[17.3.5]
\label{IV.17.3.5}
Suppose that $g$ is \'etale;
then, for $f$ to be smooth (resp. unramified, resp. \'etale) it is necessary and sufficient that $g\circ f$ is.
\end{corollary}

\begin{proof}
This follows from \sref{IV.17.3.4} and \sref{IV.17.3.3}[(ii)].
\end{proof}

\begin{proposition}[17.3.6]
\label{IV.17.3.6}
Let $g:Y\to S$ and $h:X\to S$ be two morphisms locally of finite presentation.
For an $S$-morphism $f:X\to Y$ to be unramified, it is necessary and sufficient that the canonical homomorphism \sref{IV.16.4.19}
\[
  f^*(\Omega_{Y/S}^1)\to\Omega_{X/S}^1
\]
is surjective.
\end{proposition}

\begin{proof}
As $f$ is locally of finite presentation \sref[I]{I.4.4.3}[(v)], the proposition follows from \sref{IV.17.2.2}.
\end{proof}

\begin{definition}[17.3.7]
\label{IV.17.3.7}
Let $f:X\to Y$ be a morphism.
We say that $f$ is \emph{smooth} (resp. \emph{unramified}, resp. \emph{\'etale}) at a point $x\in X$, if there exists an open neighbourhood $U$ of $x$ in $X$ such that the restriction $f|U$ is a smooth (resp. unramified, resp. \'etale) morphism from $U$ to $Y$.
\end{definition}

We then also say that $X$ is \emph{smooth} (resp. \emph{unramified}, resp. \emph{\'etale}) \emph{over $Y$ at the point $x$}.

Taking Remark \sref{IV.17.3.2}[(iii)] into account, saying that $f$ is smooth (resp. unramified, resp. \'etale) is equivalent to saying that it is smooth (resp. unramified, resp. \'etale) at every point of $X$.

It is clear that the set of points of $X$ at which the morphism $f:X\to Y$ is smooth (resp. unramified, resp. \'etale) is \emph{open} in $X$.

\begin{proposition}[17.3.8]
\label{IV.17.3.8}
For all preschemes $Y$ and all locally free $\sh{O}_Y$-modules $\sh{E}$ of finite type, the vector bundle prescheme $\bb{V}(\sh{E})$ \sref[II]{II.1.7.8} associated to $\sh{E}$ is a smooth $Y$-prescheme.
\end{proposition}

\begin{proof}
By \sref{IV.17.3.2}[(iii)], we may restrict to the case where $Y=\Spec(A)$ is affine and $\bb{V}(\sh{E})=\Spec(A[T_1,\dots,T_r])$;
since $A[T_1,\dots,T_r]$ is a formally smooth $A$-algebra for the discrete topologies \sref[0]{0.19.3.2} and is of finite presentation, this proves the proposition \sref{IV.17.3.2}[(ii)].
\end{proof}

\begin{corollary}[17.3.9]
\label{IV.17.3.9}
Under the hypotheses of \sref{IV.17.3.8}, the projective prescheme $\bb{P}(\sh{E})$ \sref[II]{II.4.1.1} is a smooth $Y$-prescheme.
\end{corollary}

\begin{proof}
We can still restrict to the case where $Y=\Spec(A)$ is affine and $\bb{P}(\sh{E})=\bb{P}_Y^r$.
We then know \sref[II]{II.2.3.14} that we have a finite open cover of $\bb{P}^r_{A}$ by the $D_+(T_i)$ ($0\leq i\leq r$), respectively equal to the spectra of the rings $S_{(f)}$, where $S$ is replaced by $A[T_0,T_1,\dots,T_r]$ and $f$ by $T_i$;
but it follows immediately from the definition of $S_{(f)}$ \sref[II]{II.2.2.1}, that this ring, in this case, is isomorphic to $A[T_0,\dots,T_{i-1},T_{i+1},\dots,T_r]$;
hence the corollary follows by \sref{IV.17.3.8}.
\end{proof}

\subsection{Characterizations of unramified morphisms.}
\label{subsection:IV.17.4}

\oldpage[IV]{63}

\begin{theorem}[17.4.1]
\label{IV.17.4.1}
Let $f:X\to Y$ be a morphism locally of finite presentation, $x$ a point of $X$.
The following properties are equivalent:
\begin{enumerate}[label=(\alph*)]
  \item $f$ is unramified at the point $x$.
  \item The diagonal morphism $\Delta_f:X\to X\times_Y X$ is a local isomorphism at the point $x$.
  \item[b')] If one sets $\Delta_f=(\psi,\theta)$, $Z=X\times_Y X$ and $z=\psi(x)$, the homomorphism $\theta_z^\#:\sh{O}_{Z,z}\to\sh{O}_{X,x}$ is bijective.
  \item[b'')] For every morphism $g:Y'\to Y$, every point $y'\in Y'$ above $y=f(x)$, and every $Y'$-section $s'$ of $X'=X\times_Y Y'$ such that $x'=s'(y')$ lies above $x$, the section $s'$ is a local isomorphism at $y'$.
  \item $(\Omega^1_{X/Y})_x=0$.
  \item The $k(y)$-prescheme $f^{-1}(y)$ is unramified over $k(y)$ at the point $x$.
  \item[d')] The point $x$ is isolated in $X_y=f^{-1}(y)$ (equivalently (\textbf{Err}$_\textbf{III}$,~20), the morphism $f$ is quasi-finite at $x$), and the ring $\sh{O}_{X_y,x}$ is a field that is a separable extension of $k(y)$.
  \item[d'')] The ring $\sh{O}_{X_y,x}=\sh{O}_{X,x}/\mathfrak{m}_y\sh{O}_{X,x}$ is a field that is a finite separable extension of $k(y)$.
  \item The ring $\sh{O}_{X,x}$ is a formally unramified $\sh{O}_{Y,y}$-algebra for the discrete topologies.
\end{enumerate}
\end{theorem}

\begin{proof}
Since $f$ is locally of finite type, the $\sh{O}_X$-module $\Omega^1_{X/Y}$ is of finite type \sref{IV.16.3.9}; hence $(\Omega^1_{X/Y})_x=0$ is equivalent to the existence of an open neighbourhood $U$ of $x$ such that $\Omega^1_{X/Y}|U=0$.
In view of \sref{IV.17.2.1}, this proves the equivalence of a) and c).
On the other hand, if $A=\sh{O}_{Y,y}$ and $B=\sh{O}_{X,x}$, then $(\Omega^1_{X/Y})_x=\Omega^1_{B/A}$ \sref{IV.16.4.15}, and the equivalence of c) and e) follows from \sref[0]{0.20.7.4}.

Since d') involves only properties of the morphism $X_y\to\Spec(k(y))$, the equivalence of a) and d') will imply, \emph{ipso facto}, the equivalence of d) and d').
Moreover, d') and d'') are equivalent, since, $X_y$ being a $k(y)$-prescheme locally of finite type, saying that $\sh{O}_{X_y,x}$ is a \emph{finite} $k(y)$-algebra is equivalent to saying that $x$ is an isolated point of $X_y$ \sref[I]{I.6.4.4}.

We now prove the equivalence of $b)$ and $b')$.
We may restrict to the case where $Y=\Spec(R)$ and $X=\Spec(S)$ are affine and $f$ is of finite presentation. Then $Z=\Spec(S\otimes_R S)$, and $\Delta_f$ corresponds to the canonical surjective homomorphism $S\otimes_R S\to S$, whose kernel $\mathfrak{J}$ is an ideal of finite type \sref[0]{0.20.4.4}. If $\sh{J}=\widetilde{\mathfrak{J}}$, the $\sh{O}_Z$-module $\psi_*(\sh{O}_X)=\sh{O}_Z/\sh{J}$ is therefore of finite presentation; the hypothesis that $\theta_z^\#:\sh{O}_{Z,z}\to(\psi_*(\sh{O}_X))_z=\sh{O}_{X,x}$ is bijective then implies that, after replacing $X$ if necessary by a suitable open neighbourhood of $x$, the homomorphism $\theta:\sh{O}_Z\to\psi_*(\sh{O}_X)$ is itself bijective \sref[0]{0.5.2.7}.
This shows that $b')$ implies $b)$; the converse is obvious.

On the other hand, the equivalence of b) and b'') follows from \sref[I]{I.5.3.7}, without any finiteness hypothesis on $f$: giving a $Y'$-section $s':Y'\to X'$ is equivalent to giving a
\oldpage[IV]{64}
$Y$-morphism $h=g'\circ s':Y'\to X$ (where $g':X'\to X$ is the canonical projection), so that $s'=(1_{Y'},h)_X$, and then the diagram
\[
\tag{17.4.1.1}
\label{IV.17.4.1.1}
\xymatrix{
Y' \ar[r]^-{s'} \ar[d]_h & X'=Y'\times_Y X \ar[d]^{h\times 1_X}\\
X \ar[r]_-{\Delta_f} & X\times_Y X
}
\]
identifies $Y'$ with the $(X\times_Y X)$-prescheme product of $X$ and $X'$.
Hence (\sref[I]{I.4.3.2}) if $\Delta_f$ is an isomorphism local at the point $x$, $s'$ is an isomorphism local at the point $y'$ (since $x=h(y')$), which proves that $b)$ implies $b'')$.
The converse is obtained by applying $b'')$ to the case $Y'=X$, $y'=x$, $g=f$, and $s'=\Delta_f$.

To finish the proof of {\hyperlink{ega4.c04.thm.17-4-1}{(17.4.1)}}, it suffices to prove the implications $d'')\Rightarrow c)\Rightarrow b)\Rightarrow d'')$.

$d'')\Rightarrow c)$: Since $\Omega^1_{X/Y}$ is an $\sh{O}_X$-module of finite type, Nakayama's lemma shows that condition c) is equivalent to $(\Omega^1_{X/Y})_x/\mathfrak{m}_y(\Omega^1_{X/Y})_x=0$, that is, by \sref{IV.16.4.5}, to $(\Omega^1_{X_y/\Spec(k(y))})_x=0$.
We are thus reduced to the case where $Y$ is the spectrum of a field $k$ and $X$ is an algebraic $k$-prescheme.
The hypothesis that $\sh{O}_{X,x}$ is a field $k'$, finite extension of $k$, implies first that $x$ is a \emph{closed} point in $X$ (\sref[I]{I.6.4.2}), then that $x$ is a \emph{maximal} point, hence is an isolated point of $X$.
Replacing $X$ by the open $\{x\}$ of $X$, one may suppose $X=\Spec(k')$; but then the hypothesis that $k'$ is a finite separable extension of $k$ implies $\Omega^1_{k'/k}=0$ ($\mathbf{0}$,~20.6.20), which proves $c)$.

$c)\Rightarrow b)$: We saw above that one then has $\Omega^1_{X/Y}|U=0$ for a neighbourhood open $U$ of $x$ in $X$; the assertion $b)$ follows then from the definition of $\Omega^1_{X/Y}$ {\hyperlink{ega4.c04.def.16-3-1}{(16.3.1)}} and {\hyperlink{ega4.c04.prop.16-1-9}{(16.1.9)}}.

$b)\Rightarrow d'')$: Replacing $X$ by an open neighbourhood of $x$, we may suppose that $\Delta_f$ is an open immersion. If $f_y:X_y\to\Spec(k(y))$ denotes the morphism obtained from $f$ by base change, then $\Delta_{f_y}$ is also an open immersion \sref[I]{I.5.3.4}. Since condition d'') concerns only $X_y$, we may reduce to the case where $Y$ is the spectrum of a field $k$ and $X$ is the spectrum of a $k$-algebra $A$ of finite type. What must be proved is that $A$ is a \emph{finite separable} $k$-algebra, that is, a finite product of finite separable extensions of $k$.
If $K$ is an algebraically closed extension of $k$, saying this is equivalent to saying that $A\otimes_k K$ is a finite separable $K$-algebra \sref{IV.4.6.1}; we may therefore suppose that $k$ is algebraically closed. We first show that $A$ is finite over $k$. It is enough to show that every closed point $x$ of $X$ is isolated: the set of such points will then be open and discrete in $X$, and hence finite because it is quasi-compact ($X$ being Noetherian), which proves the assertion by \sref[I]{I.6.4.4}. But then $k(x)=k$, since $k$ is algebraically closed \sref[I]{I.6.4.3}; consequently there is a $Y$-section $s$ of $X$ such that $s(Y)=\{x\}$. By \sref{IV.17.4.1.1}, $\{x\}$ is the inverse image of the diagonal $\Delta_X(X)$ under a morphism $X\to X\times_Y X$, and therefore $\{x\}$ is open in $X$ by hypothesis b). Thus we have shown that $A$ is a
\oldpage[IV]{65}
\emph{finite} $k$-algebra.
Since a finite $k$-algebra is a finite product of finite local $k$-algebras, to express that $\Delta_f$ is an open immersion we may reduce to the case where $A$ is a finite \emph{local} $k$-algebra. The residue field of $A$, being a finite extension of $k$, is necessarily equal to $k$; hence \sref[I]{I.3.4.9} the underlying space of $X\times_k X$ consists of a single point, and $\Delta_f$ is necessarily an \emph{isomorphism}.
But since $A$ is a $k$-algebra, the canonical homomorphism $A\otimes_k A\to A$ can only be bijective if $A=k$.
\end{proof}

\begin{remark}[17.4.1.2]
\label{IV.17.4.1.2}
Suppose only that $f$ is \emph{locally of finite type}.
Then $\Omega^1_{X/Y}$ is still an $\sh{O}_X$-module of finite type \sref{IV.16.3.9}, and $\Delta_f$ is a morphism locally of finite presentation \sref{IV.1.4.3.1}.
The entire proof of \sref{IV.17.4.1} remains valid, provided a) is replaced by: \emph{the restriction of $f$ to a suitable neighbourhood of $x$ is a formally unramified morphism}.
One sees moreover that in this case the restriction of $f$ to a suitable neighbourhood of $x$ is a morphism \emph{locally quasi-finite}.
\end{remark}

\begin{corollary}[17.4.2]
\label{IV.17.4.2}
Let $f:X\to Y$ be a morphism locally of finite presentation.
The following properties are equivalent:
\begin{enumerate}[label=(\alph*)]
  \item $f$ is unramified.
  \item The diagonal morphism $\Delta_f:X\to X\times_Y X$ is an open immersion.
  \item[b')] For every morphism $Y'\to Y$, every $Y'$-section of $X'=X\times_Y Y'$ is an open immersion.
  \item $\Omega^1_{X/Y}=0$.
  \item For every $y\in Y$, the $k(y)$-prescheme $f^{-1}(y)$ is unramified over $k(y)$.
  \item[d')] For every $y\in Y$, the $k(y)$-prescheme $f^{-1}(y)$ is isomorphic to a prescheme of the form $\coprod_{\lambda\in L}\Spec(K_\lambda)$, where each $K_\lambda$ is a finite separable extension of $k(y)$.
  \item For every $x\in X$, the ring $\sh{O}_{X,x}$ is a formally unramified $\sh{O}_{Y,f(x)}$-algebra for the discrete topologies.
\end{enumerate}
\end{corollary}

\begin{corollary}[17.4.3]
\label{IV.17.4.3}
If $f:X\to Y$ is unramified, then $f$ is locally quasi-finite (\textbf{Err}$_\textbf{III}$,~20).
\end{corollary}

\begin{proposition}[17.4.4]
\label{IV.17.4.4}
Let $Y$ be a locally Noetherian prescheme, $f:X\to Y$ a morphism locally of finite type, $x$ a point of $X$, $y=f(x)$.
Set $A=\sh{O}_{Y,y}$, $B=\sh{O}_{X,x}$, which are local Noetherian rings, and let $k$ be the residue field of $A$.
Then the conditions equivalent $a)$ to $e)$ of theorem {\hyperlink{ega4.c04.thm.17-4-1}{(17.4.1)}} are also equivalent to each of the following:
\begin{enumerate}
  \item[f)] $\hat{B}\otimes_{\hat{A}}k$ is a field, finite separable extension of $k$ (which implies that $\hat{B}$ is an $\hat{A}$-algebra finite).
  \item[f')] $\hat{B}$ is a formally unramified $\hat{A}$-algebra for the adic topologies.
\end{enumerate}
If moreover $k(x)=k(y)$, or if $k$ is separably closed, these conditions are also equivalent to:
\begin{enumerate}
  \item[f'')] The homomorphism $\hat{A}\to\hat{B}$ is surjective.
\end{enumerate}
\end{proposition}

\begin{proof}
First note, by the same argument as in \sref[0]{0.19.3.6}, that saying that $B$ is a formally unramified $A$-algebra for the preadic topologies is equivalent to saying that $\widehat{B}$ is a formally unramified $\widehat{A}$-algebra for the adic topologies.
On the other hand, the hypothesis that $f$ is locally of finite type implies that
\oldpage[IV]{66}
$\Omega^1_{B/A}$ is a $B$-module of finite type \sref{IV.16.3.9}, and hence is separated for the $\mathfrak{n}$-preadic topology (where $\mathfrak{n}$ is the maximal ideal of $B$) \sref[0]{0.7.3.5}; consequently, $\Omega^1_{B/A}=0$ is equivalent to $\widehat{\Omega^1_{B/A}}=0$. By \sref[0]{0.20.7.4}, it is therefore equivalent to say that $B$ is a formally unramified $A$-algebra for the discrete topologies or for the preadic topologies. This proves the equivalence of e) and f').
If $\mathfrak{m}$ is the maximal ideal of $A$, then
\[
k=A/\mathfrak{m}=\widehat{A}/\mathfrak{m}\widehat{A},\qquad
\widehat{B}\otimes_{\widehat{A}}k=\widehat{B}/\mathfrak{m}\widehat{B}
=\widehat{B}\otimes_B(B/\mathfrak{m}B).
\]
It follows from \sref[0]{0.7.3.5} that $\widehat{B}/\mathfrak{m}\widehat{B}$ is the completion of $B/\mathfrak{m}B=B\otimes_A k$ for the $\mathfrak{n}$-preadic topology; this proves the equivalence of d'') and f).
Finally, when $k(x)=k(y)$, or when $k$ is separably closed, condition f) implies that the homomorphism $\widehat{A}/\mathfrak{m}\widehat{A}\to\widehat{B}/\mathfrak{m}\widehat{B}$ is bijective. Condition f) also implies that $\widehat{B}$ is a quasi-finite $\widehat{A}$-algebra \sref[0]{0.7.4.4}, hence a finite one, since $\widehat{A}$ is complete and $\widehat{B}$ is separated for the $\mathfrak{m}$-preadic topology, $\mathfrak{m}\widehat{B}$ being an ideal of definition of $\widehat{B}$ \sref[0]{0.7.4.1}. Nakayama's lemma therefore shows that $\widehat{A}\to\widehat{B}$ is surjective. Thus f) implies f''), and the converse is immediate.
\end{proof}

\begin{env}[17.4.5]
\label{IV.17.4.5}
Given an $S$-prescheme $Y$ and two $S$-morphisms $f:X\to Y$, $g:X\to Y$, one canonically deduces an $S$-morphism $(f,g)_S:X\to Y\times_S Y$.
We call the \emph{prescheme of coincidences} of $f$ and $g$ the inverse image under $(f,g)_S$ of the diagonal $\Delta_{Y|S}$; it is therefore a subprescheme of $X$, and is a \emph{closed} subprescheme when $Y$ is an $S$-scheme \sref[I]{I.5.4.1}.
\end{env}

\begin{proposition}[17.4.6]
\label{IV.17.4.6}
Let $h:Y\to S$ be an unramified morphism, and let $f:X\to Y$, $g:X\to Y$ be two $S$-morphisms.
Then the prescheme $C$ of coincidences of $f$ and $g$ is a sub-prescheme induced on an open of $X$; if moreover $Y$ is an $S$-scheme (\sref[I]{I.5.4.1}), $C$ is a closed sub-prescheme of $X$.
\end{proposition}

\begin{proof}
Indeed, since $\Delta_h:Y\to Y\times_S Y$ is an open immersion {\hyperlink{ega4.c04.cor.17-4-2}{(17.4.2)}}, the inverse image by $(f,g)_S$ of $\Delta_{Y|S}$ is a sub-prescheme induced on an open of $X$ (\sref[I]{I.4.4.1}).
The last assertion follows from {\hyperlink{ega4.c04.statement.17-4-5}{(17.4.5)}}.
\end{proof}

\begin{corollary}[17.4.7]
\label{IV.17.4.7}
Under the hypotheses of {\hyperlink{ega4.c04.prop.17-4-6}{(17.4.6)}}, let $x$ be a point of $X$ such that the two composed morphisms $\Spec(k(x))\to X\xrightarrow{f} Y$ and $\Spec(k(x))\to X\xrightarrow{g} Y$ are equal.
Then there exists an open neighbourhood $U$ of $x$ such that $f|U=g|U$.
If moreover $Y$ is an $S$-scheme, there exists a closed open neighbourhood $X'$ of $x$ in $X$ such that $f|X'=g|X'$.
If finally $X$ is connected, then $f=g$.
\end{corollary}

This follows from {\hyperlink{ega4.c04.prop.17-4-6}{(17.4.6)}} and (\sref[I]{I.5.3.17}).

\begin{corollary}[17.4.8]
\label{IV.17.4.8}
Under the hypotheses of {\hyperlink{ega4.c04.prop.17-4-6}{(17.4.6)}}, suppose that the structural morphism $\varphi=h\circ f=h\circ g$ of $X$ in $S$ is closed.
Let $s$ be a point of $S$; let $X_s$ be the $k(s)$-prescheme $\varphi^{-1}(s)$ and suppose that the two morphisms composed $X_s\to X\xrightarrow{f} Y$ and $X_s\to X\xrightarrow{g} Y$ are equal.
Then there exists a neighbourhood open $V$ of $s$ in $S$ such that $f|\varphi^{-1}(V)=g|\varphi^{-1}(V)$.
If moreover $Y$ is an $S$-scheme and $\varphi$ is open, one can take $V$ open and closed.
If finally $S$ is connected, $f=g$.
\end{corollary}

\begin{proof}
It follows from {\hyperlink{ega4.c04.cor.17-4-7}{(17.4.7)}} that the prescheme $C$ of coincidences of $f$ and $g$ is induced on an open subset of $X$ and contains $X_s$.
Since $\varphi$ is closed, there exists an open neighbourhood $V$ of $s$ such that $\varphi^{-1}(V)\subset C$.
If moreover $Y$ is an $S$-scheme, then $C$ is closed; hence, when $\varphi$ is also open, $\varphi(X-C)$ is
\oldpage[IV]{67}
both open and closed in $S$, and its complement $V$ in $S$ is therefore an open-and-closed neighbourhood of $s$ such that $\varphi^{-1}(V)\subset C$.
If finally $S$ is connected, then $V=S$, and consequently $f=g$.
\end{proof}

\begin{proposition}[17.4.9]
\label{IV.17.4.9}
Let $Y$ be a connected prescheme, $f:X\to Y$ an unramified and separated morphism.
Then every $Y$-section $g$ of $X$ is an isomorphism of $Y$ onto a connected component open of $X$, and the application $g\mapsto g(Y)$ is a bijection of $\Gamma(X/Y)$ onto the set of connected components $Z$ of $X$ (necessarily \emph{open} in $X$) such that the restriction of $f$ to $Z$ is an isomorphism.
In particular, if $g'$ and $g''$ are two $Y$-sections of $X$ such that $g'(y)=g''(y)$ for some $y\in Y$, then $g'=g''$.
\end{proposition}

\begin{proof}
It follows from ({\hyperlink{ega4.c04.thm.17-4-1}{17.4.1}},~$b'')$) that a $Y$-section $s$ of $X$ is also an open immersion, and since $X$ is a $Y$-scheme, $s$ is also a closed immersion (\sref[I]{I.5.4.7}); it follows that $s$ is an isomorphism of $Y$ onto a sub-prescheme of $X$ that is both open and closed, hence a connected component open of $X$, and since $s(Y)$ is connected, it is necessarily a connected component.
The rest of the proposition is immediate.
\end{proof}

\begin{remark}[17.4.10]
\label{IV.17.4.10}
In view of Remark \sref{IV.17.4.1.2}, in statements \sref{IV.17.4.6} through \sref{IV.17.4.9} the word ``unramified'' may everywhere be replaced by ``formally unramified and locally of finite type''.
\end{remark}

\subsection{Characterizations of smooth morphisms.}
\label{subsection:IV.17.5}

\begin{theorem}[17.5.1]
\label{IV.17.5.1}
Let $f:X\to Y$ be a morphism locally of finite presentation, $x$ a point of $X$, $y=f(x)$.
The following conditions are equivalent:
\begin{enumerate}[label=(\alph*)]
  \item $f$ is smooth at the point $x$.
  \item $f$ is flat at the point $x$ and the $k(y)$-prescheme $f^{-1}(y)$ is smooth over $k(y)$ at the point $x$.
  \item[b')] $f$ is regular at the point $x$ {\hyperlink{ega4.c04.def.6-8-1}{(6.8.1)}}.
  \item $\sh{O}_{X,x}$ is an $\sh{O}_{Y,y}$-algebra formally smooth for the discrete topologies.
\end{enumerate}
\end{theorem}

\begin{proof}
One may reduce to the case where $Y=\Spec(A)$, $X=\Spec(C)$, where $C=B/\mathfrak{J}$, $B=A[T_1,\ldots,T_n]$ being a polynomial algebra and $\mathfrak{J}$ an ideal \emph{of finite type}.
The equivalence of $a)$ and $c)$ follows from ($\mathbf{0}$,~22.6.4).
On the other hand, applying this result to the morphism locally of finite type $f^{-1}(y)\to\Spec(k(y))$, one sees that the equivalence of $b)$ and $b')$ follows from the equivalence of $a)$ and $b)$ in {\hyperlink{ega4.c04.thm.6-8-6}{(6.8.6)}}.
It remains to prove the equivalence of $a)$ and $b)$.

Let us first show that $a)$ implies $b)$.
Denote by $\mathfrak{p}$ the prime ideal $\mathfrak{i}_x$ of $C$, and by $\mathfrak{r}$ the prime ideal $\mathfrak{i}_y$ of $A$; then $\mathfrak{p}=\mathfrak{q}/\mathfrak{J}$, where $\mathfrak{q}$ is a prime ideal of $B$, and $\mathfrak{r}$ is the inverse image of $\mathfrak{q}$ in $A$.
Hypothesis $a)$ first implies, by {\hyperlink{ega4.c04.prop.17-3-3}{(17.3.3)}}, that $f^{-1}(y)$ is smooth over $k(y)$ at $x$, and it remains to prove that $C_\mathfrak{p}=\sh{O}_{X,x}$ is a \emph{flat} $A_\mathfrak{r}$-module.
Since $C_\mathfrak{p}$ is a formally smooth $A_\mathfrak{r}$-algebra and $B$ is a formally smooth $A$-algebra (for the discrete topologies), the Jacobian criterion ($\mathbf{0}$,~22.6.4), together with ($\mathbf{0}$,~19.1.12), implies that there exist in $\mathfrak{J}$ a system of $r$ polynomials $u_i$ ($1\leqslant i\leqslant r$) and $r$ indices $j_i$ ($1\leqslant i\leqslant r$) such that the images of the $u_i$ in $\mathfrak{J}_\mathfrak{q}/\mathfrak{J}_\mathfrak{q}^2$ generate this $B_\mathfrak{q}$-module and
\[
\label{IV.17.5.1.1}
\tag{17.5.1.1}
\det\left(\frac{\partial u_i}{\partial T_{j_k}}\right)\notin\mathfrak{q}.
\]

\oldpage[IV]{68}

Now let $g:\Spec(B)\to\Spec(A)$ be the structural morphism.
The fibre $g^{-1}(y)$ is the spectrum of the regular ring $k(y)[T_1,\ldots,T_n]$ ($\mathbf{0}$,~17.3.7); hence the Noetherian local ring $B_\mathfrak{q}/\mathfrak{r}B_\mathfrak{q}$ at a point of this fibre is regular.
Condition {\hyperlink{ega4.c04.eq.17-5-1-1}{(17.5.1.1)}} implies that the canonical images $v_i$ of the $u_i$ in the maximal ideal $\mathfrak{m}$ of $B_\mathfrak{q}/\mathfrak{r}B_\mathfrak{q}$ are \emph{linearly independent} modulo $\mathfrak{m}^2$.
Indeed, otherwise there would exist polynomials $w_i\in B$ ($1\leqslant i\leqslant r$), not all belonging to $\mathfrak{q}$, such that $\sum_{i=1}^r w_i u_i\in\mathfrak{q}^2$.
Differentiating with respect to the $T_{j_k}$ would then give $\sum_{i=1}^r w_i(\partial u_i/\partial T_{j_k})\in\mathfrak{q}$ for $1\leqslant k\leqslant r$, contradicting {\hyperlink{ega4.c04.eq.17-5-1-1}{(17.5.1.1)}}, since $\mathfrak{q}$ is prime.
It follows from ($\mathbf{0}$,~17.1.7) that $(v_i)$ is a \emph{regular sequence} in $B_\mathfrak{q}/\mathfrak{r}B_\mathfrak{q}$.
But since $g$ is locally of finite presentation and $B$ is a flat $A$-module, {\hyperlink{ega4.c04.thm.11-3-8}{(11.3.8)}} shows that the canonical images $u_i'$ of the $u_i$ in $B_\mathfrak{q}$ also form a regular sequence and that $B_\mathfrak{q}/(\sum_i u_i'B_\mathfrak{q})$ is a flat $A_\mathfrak{r}$-module.
Since the images of the $u_i'$ in $\mathfrak{J}_\mathfrak{q}/\mathfrak{J}_\mathfrak{q}^2$ generate this $B_\mathfrak{q}$-module, Nakayama's lemma gives $\sum_i u_i'B_\mathfrak{q}=\mathfrak{J}_\mathfrak{q}$; hence $C_\mathfrak{p}=B_\mathfrak{q}/\mathfrak{J}_\mathfrak{q}$ is indeed a flat $A_\mathfrak{r}$-module.

Finally, let us prove that $b)$ implies $a)$.
With the same notation, the hypothesis that $B_\mathfrak{q}/\mathfrak{J}_\mathfrak{q}=C_\mathfrak{p}$ is a flat $A_\mathfrak{r}$-module implies that the canonical homomorphism $\mathfrak{J}_\mathfrak{q}/\mathfrak{r}\mathfrak{J}_\mathfrak{q}\to B_\mathfrak{q}/\mathfrak{r}B_\mathfrak{q}$ is injective ($\mathbf{0}_\mathbf{I}$,~6.1.2), so that $\mathfrak{J}_\mathfrak{q}/\mathfrak{r}\mathfrak{J}_\mathfrak{q}$ is identified with an ideal of $B_\mathfrak{q}/\mathfrak{r}B_\mathfrak{q}$.
Since $B_\mathfrak{q}/\mathfrak{r}B_\mathfrak{q}$ is a formally smooth $A_\mathfrak{r}$-algebra for the discrete topologies, the Jacobian criterion ($\mathbf{0}$,~22.6.4), applied to
\[
C_\mathfrak{p}=(B_\mathfrak{q}/\mathfrak{r}B_\mathfrak{q})/(\mathfrak{J}_\mathfrak{q}/\mathfrak{r}\mathfrak{J}_\mathfrak{q}),
\]
together with ($\mathbf{0}$,~19.1.12), proves, by hypothesis $b)$, the existence of $r$ polynomials $v_i\in k(y)[T_1,\ldots,T_n]$ and $r$ indices $j_i$ ($1\leqslant i\leqslant r$) such that their images in
\[
(\mathfrak{J}_\mathfrak{q}/\mathfrak{r}\mathfrak{J}_\mathfrak{q})/(\mathfrak{J}_\mathfrak{q}/\mathfrak{r}\mathfrak{J}_\mathfrak{q})^2
\]
generate this $(B_\mathfrak{q}/\mathfrak{r}B_\mathfrak{q})$-module and
\[
\label{IV.17.5.1.2}
\tag{17.5.1.2}
\det\left(\frac{\partial v_i}{\partial T_{j_k}}\right)\notin\mathfrak{q}B_\mathfrak{q}/\mathfrak{r}B_\mathfrak{q}.
\]
For each $i$, choose an element $u_i\in\mathfrak{J}$ whose canonical image is $v_i$.
It follows from {\hyperlink{ega4.c04.eq.17-5-1-2}{(17.5.1.2)}} that the $u_i$ satisfy {\hyperlink{ega4.c04.eq.17-5-1-1}{(17.5.1.1)}}; on the other hand, Nakayama's lemma shows that the images $u_i'$ of the $u_i$ in $\mathfrak{J}_\mathfrak{q}$ generate this $B_\mathfrak{q}$-module.
The Jacobian criterion ($\mathbf{0}$,~22.6.4), together with ($\mathbf{0}$,~19.1.12), then proves that $C_\mathfrak{p}=B_\mathfrak{q}/\mathfrak{J}_\mathfrak{q}$ is a formally smooth $A_\mathfrak{r}$-algebra for the discrete topologies.
Q.E.D.
\end{proof}

\begin{corollary}[17.5.2]
\label{IV.17.5.2}
Let $f:X\to Y$ be a morphism locally of finite presentation.
For $f$ to be smooth (in the sense of {\hyperlink{ega4.c04.def.17-3-1}{(17.3.1)}}), it is necessary and sufficient that $f$ be regular {\hyperlink{ega4.c04.def.6-8-1}{(6.8.1)}}, i.e.\ that $f$ be flat and that for every $y\in Y$, $f^{-1}(y)$ be a $k(y)$-prescheme geometrically regular {\hyperlink{ega4.c04.def.6-7-6}{(6.7.6)}}.
\end{corollary}

One has thus established the equivalence of the two definitions of ``smooth morphism'' given in {\hyperlink{ega4.c04.def.6-8-1}{(6.8.1)}} and {\hyperlink{ega4.c04.def.17-3-1}{(17.3.1)}}.

\begin{proposition}[17.5.3]
\label{IV.17.5.3}
Let $Y$ be a locally Noetherian prescheme, $f:X\to Y$ a morphism locally of finite type, $x$ a point of $X$, $y=f(x)$.
Set $A=\sh{O}_{Y,y}$, $B=\sh{O}_{X,x}$, which are local Noetherian rings.
Then the conditions equivalent $a)$ to $c)$ of {\hyperlink{ega4.c04.thm.17-5-1}{(17.5.1)}} are also equivalent to each of the following:

\oldpage[IV]{69}

\begin{enumerate}
  \item[d)] $B$ is an $A$-algebra formally smooth for the preadic topologies.
  \item[d')] $\hat{B}$ is an $\hat{A}$-algebra formally smooth for the adic topologies.
\end{enumerate}
If moreover $k(x)=k(y)$, these conditions are also equivalent to:
\begin{enumerate}
  \item[d'')] $\hat{B}$ is an $\hat{A}$-algebra isomorphic to a formal power series algebra $\hat{A}[[T_1,\ldots,T_n]]$.
\end{enumerate}
\end{proposition}

\begin{proof}
The equivalence of condition $c)$ of {\hyperlink{ega4.c04.thm.17-5-1}{(17.5.1)}} and condition $d)$ follows from the equivalence of $a)$ and $d)$ in the Jacobian criterion ($\mathbf{0}$,~22.6.4), and the equivalence of $d)$ and $d')$ follows from ($\mathbf{0}$,~19.3.6).
Moreover, $d'')$ implies $d')$ without any hypothesis on the residue fields ($\mathbf{0}$,~19.3.4).
Finally, if $\mathfrak{m}$ denotes the maximal ideal of $\hat{A}$, hypothesis $d')$ implies that $\hat{B}/\mathfrak{m}\hat{B}$ is a complete Noetherian local $k(y)$-algebra, formally smooth for its adic topology ($\mathbf{0}$,~19.3.5); the hypothesis $k(y)=k(x)$ therefore implies that $\hat{B}/\mathfrak{m}\hat{B}$ is $k(y)$-isomorphic to a formal power series algebra $k(y)[[T_1,\ldots,T_n]]$ ($\mathbf{0}$,~19.6.4).
On the other hand, $\hat{A}[[T_1,\ldots,T_n]]$ is a flat $\hat{A}$-module and a complete Noetherian local $\hat{A}$-algebra; it follows from ($\mathbf{0}$,~19.7.1.5) that this algebra is isomorphic to $\hat{B}$.
Thus $d')$ implies $d'')$ under the additional hypothesis $k(x)=k(y)$.
\end{proof}

\begin{remark}[17.5.4]
\label{IV.17.5.4}
Suppose that $Y$ is a locally Noetherian prescheme, and $f:X\to Y$ a morphism locally of finite type.
The criterion ({\hyperlink{ega4.c04.prop.17-5-3}{17.5.3}},~$d)$), together with ($\mathbf{0}$,~22.1.4) shows that for proving $f$ is smooth, one can apply the definition {\hyperlink{ega4.c04.def.17-1-1}{(17.1.1)}}, \emph{restricting} to the case where the affine scheme $Y'$ is the spectrum of a local Artinian ring.
\end{remark}

\begin{proposition}[17.5.5]
\label{IV.17.5.5}
Let $Y$ be a locally Noetherian prescheme, $f:X\to Y$ a morphism locally of finite type, $x$ a point of $X$, $y=f(x)$.
Suppose that $Y$ is reduced at the point $y$.
Then, for $f$ to be smooth at the point $x$, it is necessary and sufficient that $f$ be universally open in a neighbourhood of $x$ in $f^{-1}(y)$ and that $f^{-1}(y)$ be a $k(y)$-prescheme geometrically regular at the point $x$.
\end{proposition}

\begin{proof}
By {\hyperlink{ega4.c04.thm.17-5-1}{(17.5.1)}}, everything reduces to seeing that, if $f^{-1}(y)$ is a $k(y)$-prescheme geometrically regular at the point $x$, it is equivalent to say that $f$ is flat or universally open in a neighbourhood of $x$ in $f^{-1}(y)$.
But if $f$ is flat at the point $x$, it is in a neighbourhood of $x$ in $X$ {\hyperlink{ega4.c04.thm.11-1-1}{(11.1.1)}} and hence universally open in this neighbourhood {\hyperlink{ega4.c04.thm.2-4-6}{(2.4.6)}}.
Reciprocally, the hypothesis that $f$ is universally open in a neighbourhood of $x$ in $f^{-1}(y)$ and that $f^{-1}(y)$ is a $k(y)$-prescheme geometrically regular at the point $x$ implies that $f$ is flat at the point $x$, since $\sh{O}_{Y,y}$ is reduced {\hyperlink{ega4.c04.thm.15-2-2}{(15.2.2)}}.
\end{proof}

\begin{corollary}[17.5.6]
\label{IV.17.5.6}
Let $Y$ be a locally Noetherian prescheme, $f:X\to Y$ a morphism locally of finite type, $x$ a point of $X$, $y=f(x)$.
Suppose that $Y$ is reduced and geometrically unibranch {\hyperlink{ega4.c04.statement.6-15-1}{(6.15.1)}} at the point $y$.
Then, for $f$ to be smooth at the point $x$, it is necessary and sufficient that $f$ be equidimensional at the point $x$ and that $f^{-1}(y)$ be a $k(y)$-prescheme geometrically regular at the point $x$.
\end{corollary}

Observing that the set of points at which $f$ is equidimensional is open {\hyperlink{ega4.c04.def.13-3-2}{(13.3.2)}}, one sees that the corollary follows from \sref{IV.17.5.5} and the Chevalley criterion {\hyperlink{ega4.c04.cor.14-4-4}{(14.4.4)}}.

The fact that $f$ is smooth at a point implies in particular that $f$ satisfies
\oldpage[IV]{70}
at that point all the properties defined in {\hyperlink{ega4.c04.def.6-8-1}{(6.8.1)}}. We therefore have the following properties, which we recall for ease of reference:

\begin{proposition}[17.5.7]
\label{IV.17.5.7}
Let $f:X\to Y$ be a morphism locally of finite presentation, smooth at a point $x\in X$; set $y=f(x)$.
Then, for $\sh{O}_{X,x}$ to be reduced (resp.\ integrally closed, resp.\ geometrically unibranch), it is necessary and sufficient that $\sh{O}_{Y,y}$ be so.
\end{proposition}

This has been proved in {\hyperlink{ega4.c04.prop.11-3-13}{(11.3.13)}} and {\hyperlink{ega4.c04.cor.11-3-14}{(11.3.14)}} completed by \textbf{Err}$_\textbf{IV}$,~30.

\begin{proposition}[17.5.8]
\label{IV.17.5.8}
Let $Y$ be a locally Noetherian prescheme, $f:X\to Y$ a morphism locally of finite type, smooth at a point $x\in X$; set $y=f(x)$.
Then:
\begin{enumerate}[label=(\roman*)]
  \item $\dim(\sh{O}_{X,x})=\dim(\sh{O}_{Y,y})+\dim(\sh{O}_{X,x}\otimes_{\sh{O}_{Y,y}}k(y))$.
  \item $\operatorname{coprof}(\sh{O}_{X,x})=\operatorname{coprof}(\sh{O}_{Y,y})$.
  \item For $\sh{O}_{X,x}$ to have the property $(\mathrm{S}_n)$ (5.7.2) (resp.\ $(\mathrm{R}_n)$ (5.8.2)), it is necessary and sufficient that $\sh{O}_{Y,y}$ have it.
  In particular, for $\sh{O}_{X,x}$ to be regular, it is necessary and sufficient that $\sh{O}_{Y,y}$ be so.
\end{enumerate}
\end{proposition}

These are special cases of {\hyperlink{ega4.c04.cor.6-1-2}{(6.1.2)}}, {\hyperlink{ega4.c04.cor.6-3-2}{(6.3.2)}}, {\hyperlink{ega4.c04.prop.6-4-1}{(6.4.1)}} and {\hyperlink{ega4.c04.prop.6-5-3}{(6.5.3)}}.


\subsection{Characterizations of \'etale morphisms.}
\label{subsection:IV.17.6}

\begin{theorem}[17.6.1]
\label{IV.17.6.1}
Let $f:X\to Y$ be a morphism locally of finite presentation, $x$ a point of $X$, $y=f(x)$.
The following conditions are equivalent:
\begin{enumerate}[label=(\alph*)]
  \item[a)] $f$ is \'etale at the point $x$.
  \item[a')] $f$ is smooth at the point $x$ and non ramified at the point $x$.
  \item[b)] $f$ is smooth at the point $x$ and quasi-finite at the point $x$ (\textbf{Err}$_\textbf{III}$,~20).
  \item[c)] $f$ is flat at the point $x$ and non ramified at the point $x$.
  \item[c')] $f$ is flat at the point $x$ and the ring $\sh{O}_{X,x}/\mathfrak{m}_y\sh{O}_{X,x}$ is a field, extension finite separable of $k(y)$.
  \item[d)] $\sh{O}_{X,x}$ is an $\sh{O}_{Y,y}$-algebra formally \'etale for the discrete topologies.
\end{enumerate}
\end{theorem}

\begin{proof}
The equivalence of $a)$ and $a')$ follows immediately from the definitions; that of $a)$ and $d)$ follows from the equivalence of $a)$ and $e)$ in {\hyperlink{ega4.c04.thm.17-4-1}{(17.4.1)}} and the equivalence of $a)$ and $d)$ in {\hyperlink{ega4.c04.thm.17-5-1}{(17.5.1)}}.
The equivalence of $c)$ and $c')$ follows from the equivalence of $a)$ and $d')$ in {\hyperlink{ega4.c04.thm.17-4-1}{(17.4.1)}}.
The fact that $a')$ implies $c')$ follows from {\hyperlink{ega4.c04.thm.17-5-1}{(17.5.1)}}; conversely, if $c')$ holds, then $f$ is regular (hence smooth by {\hyperlink{ega4.c04.thm.17-5-1}{(17.5.1)}}) at $x$, because if $K$ is a finite separable extension of a field $k$, then, for every extension $k'$ of $k$, $\Spec(K\otimes_k k')$ is regular, being a finite sum of spectra of fields.
The fact that $a')$ implies $b)$ follows from ({\hyperlink{ega4.c04.thm.17-4-1}{17.4.1}},~$d')$) and ({\hyperlink{ega4.c04.thm.17-5-1}{17.5.1}},~$b)$).
It remains to show that $b)$ implies $c)$; since one already knows from {\hyperlink{ega4.c04.thm.17-5-1}{(17.5.1)}} that $f$ is flat at $x$, it is enough to prove that $f^{-1}(y)$ is a $k(y)$-prescheme unramified over $k(y)$.
In other words, one is reduced to proving that $b)$ implies $c)$ when $Y=\Spec(k)$ is the spectrum of a field $k$ and $X=\Spec(A)$, where $A$ is a finite local $k$-algebra ($\mathbf{0}_\mathbf{I}$,~7.4.1).
By hypothesis $b)$, $A$ is a formally smooth $k$-algebra for the discrete topologies, which here coincide with the preadic topologies; hence ($\mathbf{0}$,~19.6.5) $A$ is a regular local ring.

\oldpage[IV]{71}

It is therefore a field since it is Artinian.
It then follows from ($\mathbf{0}$,~19.6.5.1) that $A$ is a finite separable extension of $k$, which completes the proof {\hyperlink{ega4.c04.thm.17-4-1}{(17.4.1)}}.
\end{proof}

\begin{corollary}[17.6.2]
\label{IV.17.6.2}
Let $f:X\to Y$ be a morphism locally of finite presentation.
The following conditions are equivalent:
\begin{enumerate}[label=(\alph*)]
  \item[a)] $f$ is \'etale.
  \item[a')] $f$ is smooth and non ramified.
  \item[b)] $f$ is smooth and locally quasi-finite (\textbf{Err}$_\textbf{III}$,~20).
  \item[c)] $f$ is flat and non ramified.
  \item[c')] $f$ is flat, and every fibre $f^{-1}(y)$ is a $k(y)$-scheme sum of spectra of fields, extensions finite and separable of $k(y)$.
  \item[c'')] $f$ is flat, and for every $y\in Y$ and every algebraically closed extension $k'$ of $k(y)$, the ``geometric fibre'' $f^{-1}(y)\otimes_{k(y)}k'$ is a sum of spectra of fields isomorphic to $k'$.
\end{enumerate}
\end{corollary}

\begin{proof}
The only point left to prove is the equivalence of $c')$ and $c'')$.
It is clear from base change {\hyperlink{ega4.c04.prop.17-3-3}{(17.3.3)}} that $c')$ implies $c'')$.
Conversely, since the projection morphism
\[
f^{-1}(y)\otimes_{k(y)}k'\longrightarrow f^{-1}(y)
\]
is open {\hyperlink{ega4.c04.cor.2-4-10}{(2.4.10)}}, hypothesis $c'')$ implies that the space $f^{-1}(y)$ is discrete.
Thus, for every point $x\in X_y=f^{-1}(y)$, the local ring $\sh{O}_{X_y,x}=A$ is a finite $k(y)$-algebra, hence an Artinian local ring.
Moreover, it follows from $c'')$ that $\Spec(A\otimes_{k(y)}k')$ is a sum of spectra of fields isomorphic to $k'$, which is possible only if $A$ is a field that is a finite separable extension of $k(y)$ {\hyperlink{ega4.c04.prop.4-6-1}{(4.6.1)}}.
\end{proof}

\begin{proposition}[17.6.3]
\label{IV.17.6.3}
Let $Y$ be a locally Noetherian prescheme, $f:X\to Y$ a morphism locally of finite type, $x$ a point of $X$, $y=f(x)$.
Set $A=\sh{O}_{Y,y}$, $B=\sh{O}_{X,x}$, which are local Noetherian rings, and let $k$ be the residue field of $A$.
Then the conditions equivalent $a)$ to $d)$ of {\hyperlink{ega4.c04.thm.17-6-1}{(17.6.1)}} are also equivalent to each of the following:
\begin{enumerate}
  \item[e)] $\hat{B}$ is an $\hat{A}$-algebra formally \'etale for the adic topologies.
  \item[e')] $\hat{B}$ is an $\hat{A}$-module free and $\hat{B}\otimes_{\hat{A}}k$ is a field, extension finite and separable of $k$ (which implies that $\hat{B}$ is an $\hat{A}$-algebra finite).
\end{enumerate}
If moreover $k(x)=k(y)$, or if $k$ is separably closed, these conditions are also equivalent to:
\begin{enumerate}
  \item[e'')] The canonical homomorphism $\hat{A}\to\hat{B}$ is bijective.
\end{enumerate}
\end{proposition}

\begin{proof}
The equivalence of $e)$ with each of the conditions of {\hyperlink{ega4.c04.thm.17-6-1}{(17.6.1)}} follows at once from ({\hyperlink{ega4.c04.prop.17-4-4}{17.4.4}},~$f')$) and ({\hyperlink{ega4.c04.prop.17-5-3}{17.5.3}},~$d')$).
That $e)$ implies $e')$ follows from ({\hyperlink{ega4.c04.prop.17-4-4}{17.4.4}},~$f)$) and ($\mathbf{0}$,~19.7.1), since $\hat{B}$ is then a finite $\hat{A}$-algebra {\hyperlink{ega4.c04.prop.17-4-4}{(17.4.4)}}, and saying that it is a flat $\hat{A}$-module is equivalent to saying that it is a free $\hat{A}$-module ($\mathbf{0}_\mathbf{III}$,~10.1.3).
Conversely, the implication $e')\Rightarrow e)$ follows from {\hyperlink{ega4.c04.prop.17-4-4}{(17.4.4)}} and ($\mathbf{0}$,~19.7.1).
Finally, $e')$ implies that the homomorphism $\hat{A}\to\hat{B}$ is injective; if $k(x)=k(y)$, or if $k$ is separably closed, this homomorphism is surjective by {\hyperlink{ega4.c04.prop.17-4-4}{(17.4.4)}}.
The converse is immediate.
\end{proof}

\begin{proposition}[17.6.4]
\label{IV.17.6.4}
Under the hypotheses of {\hyperlink{ega4.c04.prop.17-6-3}{(17.6.3)}}, if $f$ is \'etale at the point $x$, then $\dim(\sh{O}_{X,x})=\dim(\sh{O}_{Y,y})$.
\end{proposition}

This is a particular case of ({\hyperlink{ega4.c04.prop.17-5-8}{17.5.8}},~(i)) since $x$ is isolated in its fibre $f^{-1}(y)$.


\subsection{Properties of descent, passage to the limit, and constructibility.}
\label{subsection:IV.17.7}
\label{IV.17.7}

\oldpage[IV]{72}

\begin{proposition}[17.7.1]
\label{IV.17.7.1}
Let $f:X\to Y$ be a morphism locally of finite presentation, $g:Y'\to Y$ a morphism, $X'=X\times_Y Y'$, $f'=f_{(Y')}:X'\to Y'$ and $g':X'\to X$ the canonical projections.
Let $x'$ be a point of $X'$ and set $x=g'(x')$, $y'=f'(x')$.
\begin{enumerate}[label=(\roman*)]
  \item If $f'$ is non ramified at the point $x'$, then $f$ is non ramified at the point $x$.
  \item Suppose moreover that $g$ is flat at the point $y'$.
  Then, if $f'$ is smooth (resp.\ \'etale) at the point $x'$, $f$ is smooth (resp.\ \'etale) at the point $x$.
\end{enumerate}
\end{proposition}

\begin{proof}
Set $y=f(x)=g(y')$, so that $f'^{-1}(y')=f^{-1}(y)\otimes_{k(y)}k(y')$; note that $f'$ is locally of finite presentation.

(i) Since the property for a morphism (locally of finite presentation) of being non ramified at a point only involves the fibre of the morphism at that point ({\hyperlink{ega4.c04.thm.17-4-1}{17.4.1}},~$d$)), one may reduce to the case where $Y$ and $Y'$ are spectra of fields.
But then it amounts to the same thing to say that $f$ (resp.\ $f'$) is non ramified at $x$ (resp.\ $x'$) or that it is \'etale at that point (17.6.1,~$c$)); thus (i) is a consequence of (ii).

(ii) Since $f'$ is flat at $x'$ {\hyperlink{ega4.c04.thm.17-5-1}{(17.5.1)}}, the hypothesis that $g$ is flat at $y'$ implies that $f$ is flat at $x$: indeed, the projection $g':X'\to X$ is flat at $x'$, and $f\circ g'=g\circ f'$ is flat at $x'$, whence the conclusion ({\hyperlink{ega4.c04.cor.2-2-11}{2.2.11}},~(iv)).
Whether $f$ is smooth (resp.\ \'etale) at $x$ then only involves the fibre $f^{-1}(y)$ (17.5.1 and 17.6.1), and one is again reduced to the case where $Y=\Spec(k)$ and $Y'=\Spec(k')$ are spectra of fields.
To say that $f'$ is smooth at $x'$ then means {\hyperlink{ega4.c04.thm.17-5-1}{(17.5.1)}} that $X'$ is geometrically regular over $k'$ at $x'$, and this implies {\hyperlink{ega4.c04.cor.6-7-8}{(6.7.8)}} that $X$ is geometrically regular over $k$ at $x$; hence $f$ is smooth at $x$.
Suppose moreover that $f'$ is \'etale at $x'$, so that $x'$ is isolated in $f'^{-1}(y')$ {\hyperlink{ega4.c04.thm.17-6-1}{(17.6.1)}}.
Since the projection $f'^{-1}(y')\to f^{-1}(y)$ is an open morphism {\hyperlink{ega4.c04.cor.2-4-10}{(2.4.10)}}, $x$ is isolated in $f^{-1}(y)$; as $f$ is already known to be smooth at $x$, it is \'etale at that point {\hyperlink{ega4.c04.thm.17-6-1}{(17.6.1)}}.
\end{proof}

\begin{corollary}[17.7.2]
\label{IV.17.7.2}
\begin{enumerate}[label=(\roman*)]
  \item With the notation of (17.7.1), let $U$ (resp.\ $U'$) be the set of points of $X$ (resp.\ $X'$) where $f$ (resp.\ $f'$) is non ramified; then $U'=g'^{-1}(U)$.
  \item Suppose moreover that $g$ is flat, and let $V$ (resp.\ $V'$) be the set of points of $X$ (resp.\ $X'$) where $f$ (resp.\ $f'$) is smooth; then $V'=g'^{-1}(V)$.
\end{enumerate}
\end{corollary}

\begin{corollary}[17.7.3]
\label{IV.17.7.3}
\begin{enumerate}[label=(\roman*)]
  \item Suppose $g$ surjective; then, for $f$ to be non ramified, it is necessary and sufficient that $f'$ be so.
  \item Let $f:X\to Y$ be an $S$-morphism and $g:S'\to S$ a faithfully flat morphism.
  Suppose that $f$ is locally of finite presentation, or that $g$ is quasi-compact.
  Then, for $f$ to be smooth (resp.\ non ramified, resp.\ \'etale), it is necessary and sufficient that $f'$ be so.
\end{enumerate}
\end{corollary}

In (ii), the case where $f$ is locally of finite presentation follows from {\hyperlink{ega4.c04.prop.17-7-1}{(17.7.1)}}.
If $g$ is quasi-compact and $f'$ is smooth (resp.\ non ramified, resp.\ \'etale), then $f'$ is locally of finite presentation by (2.7.1,~(iv)), and one is reduced to the first case.

\begin{proposition}[17.7.4]
\label{IV.17.7.4}
Let $f:X\to Y$ be a morphism, $g:Y'\to Y$ a morphism flat and locally of finite presentation, $X'=X\times_Y Y'$, $f'=f_{(Y')}:X'\to Y'$ and $g':X'\to X$ the

\oldpage[IV]{73}

canonical projections. Let $V$ (resp.\ $V'$) be the set of $x\in X$ (resp.\ $x'\in X'$) where $f$ (resp.\ $f'$) has one of the following properties:
\begin{enumerate}[label=(\roman*)]
  \item locally of finite type;
  \item locally of finite presentation;
  \item flat;
  \item non ramified;
  \item smooth;
  \item \'etale.
\end{enumerate}
Then $V'=g'^{-1}(V)$ (in other words, for $f'$ to have the property considered at a point $x'$, it is necessary and sufficient that $f$ have this property at $x$).
\end{proposition}

\begin{proof}
Property (iii) is immediate and does not even require the hypothesis that $g$ be locally of finite presentation ({\hyperlink{ega4.c04.cor.2-2-11}{2.2.11}},~(iv)), since the projection $g':X'\to X$ is flat.
By {\hyperlink{ega4.c04.cor.17-7-2}{(17.7.2)}}, the assertions concerning properties (iv), (v), and (vi) are consequences of the assertion concerning property (ii).
It is therefore enough to consider cases (i) and (ii).
It is clear that $V'\supset g'^{-1}(V)$; it remains to prove the reverse inclusion.
The sets $V$ and $V'$ are open, and $W=g'(V')$ is also open in $X$ by {\hyperlink{ega4.c04.thm.2-4-6}{(2.4.6)}}.
It is thus enough to prove that the morphism $f|W:W\to Y$ is locally of finite type (resp.\ locally of finite presentation).
By hypothesis, the composite
\[
V'\xrightarrow{g''}W\xrightarrow{f|W}Y,
\]
where $g''$ is the restriction of $g'$, is locally of finite type (resp.\ locally of finite presentation), because it is equal to the composite $V'\xrightarrow{f'}Y'\xrightarrow{g}Y$; moreover, $g''$ is surjective, hence faithfully flat.
The assertion is therefore reduced to the following lemma, which strengthens {\hyperlink{ega4.c04.prop.11-3-16}{(11.3.16)}}.
\end{proof}

\begin{lemma}[17.7.5]
\label{IV.17.7.5}
Let $f:X\to Y$ be a faithfully flat morphism and locally of finite presentation, $g:Y\to Z$ a morphism such that $g\circ f:X\to Z$ has one of the following properties: being
\begin{enumerate}[label=(\roman*)]
  \item locally of finite type;
  \item locally of finite presentation;
  \item of finite type.
\end{enumerate}
Then $g$ has the same property.
If moreover $f$ is quasi-compact or $g$ quasi-separated, the same conclusion holds for the property:
\begin{enumerate}
  \item[(iv)] of finite presentation.
\end{enumerate}
\end{lemma}

\begin{proof}
In cases (i) and (ii), one must show that, for every $y\in Y$, there are an affine open neighbourhood $V$ of $y$ in $Y$ and an affine open neighbourhood $W$ of $z=g(y)$ in $Z$, containing $g(V)$, such that the morphism $V\to W$ induced by $g$ is of finite type (resp.\ of finite presentation).
By hypothesis, there exist a point $x\in X$ such that $f(x)=y$, an affine open neighbourhood $U$ of $x$ in $X$, an affine open neighbourhood $V'$ of $y$ in $Y$ containing $f(U)$, and an affine open neighbourhood $W$ of $z$ in $Z$ containing $g(V')$, such that the morphism $f_1:U\to V'$ induced by $f$ is flat and of finite presentation, while the composite of $f_1$ with the restriction $g_1:V'\to W$ of $g$ is of finite type (resp.\ of finite presentation).
Then $f(U)$ is open in $Y$ {\hyperlink{ega4.c04.thm.2-4-6}{(2.4.6)}}; if $V\subset f(U)$ is an affine neighbourhood of $y$, the morphism
\[
f_2:f_1^{-1}(V)\longrightarrow V,
\]

\oldpage[IV]{74}

obtained by restricting $f_1$, is still of finite presentation and is moreover faithfully flat.
Furthermore, if $g_2=g_1|V$, then $g_2\circ f_2$ is of finite type (resp.\ of finite presentation), while the open set $f_1^{-1}(V)$ is quasi-compact (resp.\ quasi-compact and quasi-separated).
One is thus reduced to the hypotheses of (11.3.16), which proves that $g_2$ is of finite type (resp.\ of finite presentation).

In case (iii), the question is local on $Z$, so one may suppose $Z$ affine; then $X$ is quasi-compact, and so is $Y=f(X)$.
Case (iii) is therefore a consequence of case (i).

In case (iv), under the additional hypotheses on $f$ or $g$, one may again suppose $Z$ affine, hence $X$ and $Y$ quasi-compact.
Moreover, $g$ is already known to be locally of finite presentation and quasi-compact (I.1.3); thus everything reduces to showing that $g$ is quasi-separated, and it is enough to prove this when $f$ is quasi-compact.
But since $g\circ f$ is quasi-separated, the same is true of $f$ (I.2.2,~(v)); as $f$ is quasi-compact and locally of finite presentation, it is of finite presentation (I.6.1).
It then suffices to repeat the argument of the first paragraph of the proof of {\hyperlink{ega4.c04.prop.11-3-16}{(11.3.16)}}.
\end{proof}

\begin{remark}[17.7.6]
\label{IV.17.7.6}
Regarding the assertion relative to property (iv), one cannot suppress the hypothesis that $f$ is quasi-compact; without this, it would imply that every morphism $g:Y\to Z$ quasi-compact and locally of finite presentation would be of finite presentation, a conclusion known to be erroneous \sref[I]{I.6.4}.
Indeed, one can only restrict to the case where $Z$ is affine, hence $Y$ quasi-compact; \mbox{there} is then a finite covering of $Y$ by affine opens $U_i$ such that the restrictions $g|U_i$ are of finite presentation; it would then suffice to take for $X$ the prescheme \emph{sum} of the $U_i$, for $f:X\to Y$ the canonical morphism, which is evidently faithfully flat and locally of finite presentation \sref[I]{I.6.5}; $g\circ f$ would be of finite presentation by virtue of the choice of the $U_i$ and of \sref[I]{I.6.5}, whence our assertion.
\end{remark}

\begin{proposition}[17.7.7]
\label{IV.17.7.7}
Let $f:Y\to S$ and $h:X\to S$ be two morphisms locally of finite presentation, $g:X\to Y$ an $S$-morphism, $x$ a point of $X$, $y=g(x)$.
Suppose that $g$ is flat at the point $x$.
Then, if $h$ is smooth (resp.\ non ramified, resp.\ \'etale) at the point $x$, $f$ is smooth (resp.\ non ramified, resp.\ \'etale) at the point $y$.
\end{proposition}

\begin{proof}
Set $s=h(x)=f(y)$.
To say that $h$ is non ramified at $x$ (resp.\ that $f$ is non ramified at $y$) amounts to saying that $h^{-1}(s)$ is \'etale over $k(s)$ at $x$ (resp.\ that $f^{-1}(s)$ is \'etale over $k(s)$ at $y$) (17.4.1 and 17.6.1); since the morphism $g_s:h^{-1}(s)\to f^{-1}(s)$ induced by $g$ is flat at $x$, one may restrict to proving the proposition when $h$ is smooth or \'etale at $x$.
Moreover, since $h$ is then flat at $x$ {\hyperlink{ega4.c04.thm.17-5-1}{(17.5.1)}}, $f$ is flat at $y$ by ({\hyperlink{ega4.c04.cor.2-2-11}{2.2.11}},~(iv)).
It therefore remains to say that $f^{-1}(s)$ is smooth (resp.\ \'etale) over $k(s)$ at $y$, and one is thus reduced to the case where $S=\Spec(k)$ is the spectrum of a field.

\emph{(i) Case of smooth morphisms.}
Since $g$ is locally of finite presentation (I.4.3,~(v)), there is an open neighbourhood $U$ of $x$ in $X$ on which $g$ is flat {\hyperlink{ega4.c04.thm.11-3-1}{(11.3.1)}} and $h$ is smooth.
Moreover, $g(U)$ is an open neighbourhood of $y$ in $Y$ {\hyperlink{ega4.c04.thm.2-4-6}{(2.4.6)}}; replacing $X$ by $U$ and $Y$ by $g(U)$, one may therefore suppose that $g$ is faithfully flat and that $h$ is smooth, and one is reduced to proving that $f$ is smooth.
If $k'$ is an algebraically

\oldpage[IV]{75}

closed extension of $k$, then $X\otimes_k k'$ is smooth over $k'$; by ({\hyperlink{ega4.c04.cor.17-7-3}{17.7.3}},~(ii)), it is enough to prove that $Y\otimes_k k'$ is smooth over $k'$ (since $\Spec(k')\to\Spec(k)$ is faithfully flat).
Thus one may suppose that $k$ is algebraically closed.
The set of $k$-rational points of $Y$ is then very dense in $Y$ {\hyperlink{ega4.c04.cor.10-4-8}{(10.4.8)}}, and the locus where $Y$ is smooth over $k$ is open; it is therefore enough to prove that $Y$ is smooth over $k$ at every $k$-rational point $y$.
At such a point, saying that $Y$ is smooth over $k$ is equivalent to saying that $Y$ is regular at $y$ ({\hyperlink{ega4.c04.thm.17-5-1}{17.5.1}} and {\hyperlink{ega4.c04.cor.6-7-8}{6.7.8}}).
Now $y=g(x)$ for some $x\in X$, and by hypothesis $X$ is regular at $x$ {\hyperlink{ega4.c04.thm.17-5-1}{(17.5.1)}}; since $X$ and $Y$ are locally Noetherian and $g$ is flat, $Y$ is indeed regular at $y$ ({\hyperlink{ega4.c04.prop.6-5-1}{6.5.1}},~(i)).

\emph{(ii) Case of \'etale morphisms.}
By (i), one already knows that $f$ is smooth at $y$; by {\hyperlink{ega4.c04.thm.17-6-1}{(17.6.1)}}, it is therefore enough to prove that $f$ is quasi-finite at $y$, or equivalently that $\sh{O}_{Y,y}$ is a finite $k$-algebra.
Since $\sh{O}_{X,x}$ is a faithfully flat $\sh{O}_{Y,y}$-module ($\mathbf{0}_\mathbf{I}$,~6.6.2), $\sh{O}_{Y,y}$ is identified with a $k$-submodule of $\sh{O}_{X,x}$; and since $\sh{O}_{X,x}$ is, by hypothesis, a finite $k$-algebra, the same holds for $\sh{O}_{Y,y}$.
\end{proof}

\begin{proposition}[17.7.8]
\label{IV.17.7.8}
With the notation of {\hyperlink{ega4.c04.statement.8-8-1}{(8.8.1)}}, suppose that $X_\alpha$ and $Y_\alpha$ are locally of finite presentation over $S_\alpha$.
Let $f_\alpha:X_\alpha\to Y_\alpha$ be an $S_\alpha$-morphism and $f:X\to Y$ the corresponding $S$-morphism.
\begin{enumerate}[label=(\roman*)]
  \item Let $x$ be a point of $X$, $x_\lambda$ its canonical projection in $X_\lambda$.
  For $f$ to be smooth (resp.\ non ramified, resp.\ \'etale) at $x$, it is necessary and sufficient that there exist $\lambda\geq\alpha$ such that $f_\lambda$ is smooth (resp.\ non ramified, resp.\ \'etale) at $x_\lambda$.
  \item Suppose moreover that $X_\alpha$ is quasi-compact.
  For $f$ to be smooth (resp.\ non ramified, resp.\ \'etale), it is necessary and sufficient that there exist $\lambda\geq\alpha$ such that $f_\lambda$ is smooth (resp.\ non ramified, resp.\ \'etale).
\end{enumerate}
\end{proposition}

\begin{proof}
(i) If $y=f(x)$ and $y_\lambda=f_\lambda(x_\lambda)$ is the canonical projection of $y$ in $Y_\lambda$, then
\[
f^{-1}(y)=f_\lambda^{-1}(y_\lambda)\otimes_{k(y_\lambda)}k(y).
\]
The assertion concerning non-ramified morphisms therefore follows from ({\hyperlink{ega4.c04.prop.17-7-1}{17.7.1}},~(i)), and it is enough to consider smooth morphisms.
Since $f$ and $f_\lambda$ are locally of finite presentation, {\hyperlink{ega4.c04.thm.17-5-1}{(17.5.1)}} says that it amounts to the same thing to require that $f^{-1}(y)$ be geometrically regular at $x$ or that $f_\lambda^{-1}(y_\lambda)$ be geometrically regular at $x_\lambda$ {\hyperlink{ega4.c04.cor.6-7-8}{(6.7.8)}}.
The proposition then follows from {\hyperlink{ega4.c04.thm.17-5-1}{(17.5.1)}} and {\hyperlink{ega4.c04.thm.11-2-6}{(11.2.6)}}.

(ii) For every $\lambda$, let $U_\lambda$ be the open set of points $x_\lambda\in X_\lambda$ at which $f_\lambda$ is smooth (resp.\ non ramified, resp.\ \'etale), and let $V_\lambda$ be its inverse image in $X$.
By (i), for every $x\in X$ there exists a $\lambda$ such that $f_\lambda$ has the required property at $x_\lambda$; hence $X$ is the union of the $V_\lambda$.
Moreover, by {\hyperlink{ega4.c04.prop.17-3-3}{(17.3.3)}}, $V_\lambda\subset V_\mu$ for $\lambda\leq\mu$; since $X$ is quasi-compact, there is therefore an index $\mu$ such that $X=V_\mu$.
Since the $X_\lambda$ are quasi-compact, it follows from (8.3.4) that there exists an index $\nu\geq\mu$ such that the inverse image of $U_\mu$ in $X_\nu$ is all of $X_\nu$; this means that $f_\nu$ is smooth (resp.\ non ramified, resp.\ \'etale), by (17.3.3).
\end{proof}

\begin{corollary}[17.7.9]
\label{IV.17.7.9}
Let $S=\operatorname{Spec}(A)$ be an affine scheme, $f:X\to S$ a morphism.
The following conditions are equivalent:
\begin{enumerate}
  \item[a)] $f$ is a morphism of finite presentation and smooth (resp.\ non ramified, resp.\ \'etale).
  
  \oldpage[IV]{76}
  
  \item[b)] There exists an affine Noetherian scheme $S_0=\operatorname{Spec}(A_0)$, a morphism of finite type $f_0:X_0\to S_0$ and a morphism $S\to S_0$ such that the $S$-prescheme $X_0\otimes_{S_0} S$ is $S$-isomorphic to $X$ and $f_0$ is smooth (resp.\ non ramified, resp.\ \'etale).
  \item[c)] The conditions of $b$) are satisfied, and moreover $A_0$ is a sub-$\bb{Z}$-algebra of finite type of $A$, and the morphism $S\to S_0$ corresponds to the canonical injection $A_0\to A$.
\end{enumerate}
\end{corollary}

\begin{proof}
The proof from \sref{IV.17.7.8} is the same as that of \sref{IV.11.2.7} from \sref{IV.11.2.6}.
\end{proof}

\begin{proposition}[17.7.10]
\label{IV.17.7.10}
Let $f:Y\to S$, $h:X\to S$ be two morphisms locally of finite presentation, $g:X\to Y$ an S-morphism, $x$ a point of $X$, $y=g(x)$.
Suppose that $h$ is flat at the point $x$ and that $f$ is non ramified at the point $y$.
Then $f$ is \'etale at the point $y$ and $g$ is flat at the point $x$.
\end{proposition}

\begin{proof}
(One will see further on \sref{IV.18.4.9} that one can in fact dispense with the hypothesis that $h$ is locally of finite presentation.)

The question being local on $S$, $X$, and $Y$, one may suppose that $S$, $X$, and $Y$ are affine, that $f$, $g$, and $h$ are of finite presentation, that $h$ is flat, and that $f$ is non ramified.
In view of {\hyperlink{ega4.c04.cor.11-2-7}{(11.2.7)}} and {\hyperlink{ega4.c04.cor.17-7-9}{(17.7.9)}}, one may moreover suppose that $S$, $X$, and $Y$ are Noetherian; finally, one may restrict to the case $S=\Spec(\sh{O}_{S,s})$, where $s=f(y)=h(x)$.
Since $A=\sh{O}_{S,s}$ is a Noetherian local ring, there exist a complete Noetherian local ring $B$ with algebraically closed residue field and a local homomorphism $A\to B$ making $B$ a faithfully flat $A$-module ($\mathbf{0}_\mathbf{III}$,~10.3.1).
Replacing $X$ and $Y$ by $X\times_S\Spec(B)$ and $Y\times_S\Spec(B)$, one concludes from {\hyperlink{ega4.c04.prop.2-5-1}{(2.5.1)}} and {\hyperlink{ega4.c04.prop.17-7-1}{(17.7.1)}} that one may restrict to the case where $A$ is complete and has algebraically closed residue field.
The hypothesis that $f$ is non ramified at $y$ then implies {\hyperlink{ega4.c04.prop.17-4-4}{(17.4.4)}} that $\sh{O}_{Y,y}$ is a finite $\sh{O}_{S,s}$-algebra and a complete Noetherian local ring ($\mathbf{0}_\mathbf{I}$,~7.4.2); since the residue field of $\sh{O}_{S,s}$ is algebraically closed, the homomorphism $\sh{O}_{S,s}\to\sh{O}_{Y,y}$ is surjective {\hyperlink{ega4.c04.prop.17-4-4}{(17.4.4)}}.
On the other hand, the hypothesis that $h$ is flat at $x$ implies that the composite homomorphism $\sh{O}_{S,s}\to\sh{O}_{Y,y}\to\sh{O}_{X,x}$ is injective ($\mathbf{0}_\mathbf{I}$,~6.5.1); consequently $\sh{O}_{S,s}\to\sh{O}_{Y,y}$ is bijective, which proves that $f$ is \'etale at $y$ {\hyperlink{ega4.c04.prop.17-6-3}{(17.6.3)}}.
Furthermore, $\sh{O}_{X,x}$ is a flat $\sh{O}_{Y,y}$-module, so $g$ is flat at $x$.
\end{proof}

\begin{proposition}[17.7.11]
\label{IV.17.7.11}
Let $S$ be a prescheme and $X$, $Y$ two $S$-preschemes locally of finite presentation over $S$; let $f:X\to Y$ be an $S$-morphism.
For every $s\in S$, denote $X_s$, $Y_s$, $f_s$ the preschemes and morphism deduced from $X$, $Y$, $f$ by the change of base $\operatorname{Spec}(k(s))\to S$.
Then:
\begin{enumerate}[label=(\roman*)]
  \item The set of $x\in X$ such that, if $s$ is the image of $x$ in $S$, $f_s:X_s\to Y_s$ is smooth (resp.\ non ramified, resp.\ \'etale, resp.\ differentially smooth) at $x$, is locally constructible.
  \item Suppose moreover that $f$ is of finite presentation.
  The set of $y\in Y$ such that, if $s$ is the image of $y$ in $S$, $f_s$ is smooth (resp.\ non ramified, resp.\ \'etale, resp.\ differentially smooth) at all the points of $f_s^{-1}(y)$, is locally constructible.
  \item Suppose $X$ and $Y$ of finite presentation on $S$.
  The set of $s\in S$ such that $f_s$ is smooth (resp.\ non ramified, resp.\ \'etale, resp.\ differentially smooth) is locally constructible.
\end{enumerate}
\end{proposition}

\begin{proof}
Let $E$ be the set of $x\in X$ satisfying the property considered in (i).
The set $F$ of points $y\in Y$ satisfying the corresponding property in (ii) is none

\oldpage[IV]{77}

other than $Y-f(X-E)$; hence (ii) follows from (i) and Chevalley's theorem when $f$ is of finite presentation {\hyperlink{ega4.c04.thm.1-8-4}{(1.8.4)}}.
Likewise, if $h:Y\to S$ is the structural morphism, the set of $s\in S$ satisfying the corresponding property in (iii) is $S-h(Y-F)$; thus (iii) again follows from (i) and Chevalley's theorem when $f$ and $h$ are of finite presentation.
It therefore remains to prove the assertions in (i).

Let us first prove (i) for the property of being smooth.
The question is local on $X$, so one may restrict to the case where $S=\Spec(A)$, $X=\Spec(B)$, and $Y=\Spec(C)$ are affine, with $B$ and $C$ algebras of finite presentation.
Reasoning as at the beginning of {\hyperlink{ega4.c04.prop.9-9-1}{(9.9.1)}} and using ({\hyperlink{ega4.c04.cor.17-7-2}{17.7.2}},~(ii)), one is reduced to the case where $A$ is Noetherian.
By ($\mathbf{0}_\mathbf{III}$,~9.2.3), it is enough to show that, if $x\in E$ (resp.\ $x\notin E$), there is a neighbourhood $V$ of $x$ in $\overline{\{x\}}$ contained in $E$ (resp.\ in $X-E$).

Denote by $g:X\to S$ and $h:Y\to S$ the structural morphisms.
One may first replace $S$ by the reduced sub-prescheme $S'$ having $\overline{\{g(x)\}}$ as underlying space, $X$ by $X'=g^{-1}(S')$, and $Y$ by $Y'=h^{-1}(S')$; the fibres of $X$ and $X'$ (resp.\ of $Y$ and $Y'$) at the points of $S'$ are the same.
Thus one may restrict to the case where $S$ is integral and where $\eta=g(x)=h(y)$, with $y=f(x)$, is its generic point.

$1\textsuperscript{o}$ Suppose first that $x\in E$.
The local rings $\sh{O}_{X_\eta,x}$ and $\sh{O}_{Y_\eta,y}$ are respectively equal to $\sh{O}_{X,x}$ and $\sh{O}_{Y,y}$.
Since smoothness of a morphism of finite presentation at a point depends only on the local ring at that point and the local ring at its image {\hyperlink{ega4.c04.thm.17-5-1}{(17.5.1)}}, the hypothesis $x\in E$ means that $f$ is smooth at $x$.
It then has this property at every point of an open neighbourhood of $x$ in $X$, and ({\hyperlink{ega4.c04.prop.17-3-3}{17.3.3}},~(iii)) yields the conclusion.

$2\textsuperscript{o}$ Suppose next that $x\in X-E$ and that $f_\eta$ is not flat at $x$.
The conclusion follows from the following lemma, which sharpens {\hyperlink{ega4.c04.prop.11-2-8}{(11.2.8)}}:

\begin{lemma}[17.7.11.1]
\label{IV.17.7.11.1}
Let $g:X\to S$ and $h:Y\to S$ be two morphisms locally of finite presentation, $f:X\to Y$ an $S$-morphism, and $\sh{F}$ a quasi-coherent $\sh{O}_X$-module of finite presentation.
Then the set $E$ of $x\in X$ such that $\sh{F}_{g(x)}$ is $f_{g(x)}$-flat at $x$ is locally constructible.
\end{lemma}

\begin{proof}[Proof of {\hyperlink{ega4.c04.lem.17-7-11-1}{Lemma 17.7.11.1}}]
Reasoning again as at the beginning of {\hyperlink{ega4.c04.prop.9-9-1}{(9.9.1)}} and using {\hyperlink{ega4.c04.prop.2-5-1}{(2.5.1)}}, one is reduced to the case where $S$, $X$, and $Y$ are Noetherian.
Then, as above, one reduces to the case where $S$ is integral, where $\eta=g(x)=h(f(x))$ is its generic point, and where it is enough to prove that, if $x\in E$ (resp.\ $x\notin E$), there is a neighbourhood $V$ of $x$ in $\overline{\{x\}}$ contained in $E$ (resp.\ in $X-E$).
The case $x\in E$ is an immediate consequence of {\hyperlink{ega4.c04.thm.11-1-1}{(11.1.1)}}.
To treat the case $x\notin E$, one reasons as in {\hyperlink{ega4.c04.lem.9-4-7-1}{(9.4.7.1)}}: after possibly replacing $Y$ and $X$ by neighbourhoods of $f(x)$ and $x$, respectively, there exist two coherent $\sh{O}_Y$-modules $\sh{G}$ and $\sh{H}$ and an $\sh{O}_Y$-homomorphism $u:\sh{G}\to\sh{H}$ such that, for every $s\in S$, the homomorphism $u_s:\sh{G}_s\to\sh{H}_s$ is injective, whereas
\[
1\otimes u_\eta:\sh{F}\otimes_{\sh{O}_{Y_\eta}}\sh{G}_\eta
\longrightarrow
\sh{F}\otimes_{\sh{O}_{Y_\eta}}\sh{H}_\eta
\]
is not injective at $x$; equivalently, $x\in\Supp(\Ker(1\otimes u_\eta))$.
If $T=\overline{\{x\}}$, then $\Supp(\Ker(1\otimes u_\eta))\supset T_\eta$.
One may suppose that, for every $s\in S$, base change gives
\[
\Ker(1\otimes u_s)=(\Ker(1\otimes u))_s
\]
{\hyperlink{ega4.c04.prop.9-4-2}{(9.4.2)}}, and hence

\oldpage[IV]{78}

\[
\Supp(\Ker(1\otimes u_s))=(\Supp(\Ker(1\otimes u)))_s
\]
(I.9.1.13.1).
It follows from {\hyperlink{ega4.c04.cor.9-5-2}{(9.5.2)}} that, for $s$ in a neighbourhood of $\eta$, one has $\Supp(\Ker(1\otimes u_s))\supset T_s$, which proves the lemma.
\end{proof}

$3\textsuperscript{o}$ Suppose now that $x\in X-E$, that $f_\eta$ is flat at $x$, but that $f_\eta$ is not smooth at $x$.
To say that $f_\eta$ is flat at $x$ is equivalent to saying that $f$ itself is flat at $x$; replacing $X$ by a neighbourhood of $x$, one may therefore suppose that $f$ is flat {\hyperlink{ega4.c04.thm.11-1-1}{(11.1.1)}}.
It follows that $f_s$ is flat for every $s\in S$; and since, for every $y\in Y$, one has $f^{-1}(y)=f_{h(y)}^{-1}(y)$, it amounts to the same thing to say that $f_{g(x')}$ is smooth at $x'$ or that $f$ is smooth at $x'$.
But the locus where $f$ is smooth is open in $X$ {\hyperlink{ega4.c04.cor.12-1-7}{(12.1.7)}}; hence the locus where $f$ is not smooth is closed, and since it contains $x$ by hypothesis, it also contains $\overline{\{x\}}$.
This completes the proof for the first property considered in (i).

Let us next prove (i) for the property of being \'etale.
For $f_s$ to be \'etale at $x$ means that it is both smooth and quasi-finite at $x$ {\hyperlink{ega4.c04.thm.17-6-1}{(17.6.1)}}.
Now it amounts to the same thing to say that $f_{g(x)}$ is quasi-finite at $x$ or that $f$ itself is quasi-finite at that point; it follows from {\hyperlink{ega4.c04.cor.13-1-4}{(13.1.4)}} that the set of points $x$ at which $f_{g(x)}$ is quasi-finite is open in $X$, hence a fortiori locally constructible.
The conclusion therefore follows from the local constructibility of the set of points $x$ at which $f_{g(x)}$ is smooth.

Let us pass to the proof of (i) for the property of being differentially smooth.
Let $p:X\times_Y X\to X$ be the second canonical projection; for every $s\in S$, the second projection $X_s\times_{Y_s}X_s\to X_s$ is just $p_s$.
It follows from {\hyperlink{ega4.c04.prop.17-12-5}{(17.12.5.1)}}\footnote{The printed French text cites (17.12.5.1), but \S17.12 has no such item. Proposition (17.12.5), specifically the equivalence of its conditions (a) and (c), is the result used here. The reader may verify that the result of {\hyperlink{ega4.c04.prop.17-7-11}{(17.7.11)}} is not used in the remainder of \S17, so there is no vicious circle.} that $f_{g(x)}$ is differentially smooth at $x$ if and only if $p_{g(x)}$ is smooth at $\Delta_f(x)$.
Since $p$ is locally of finite presentation, the set of points $z\in X\times_Y X$ at which $p_{g(p(z))}$ is smooth is locally constructible; its intersection with the locally closed subset $\Delta_f(X)$ is therefore locally constructible in $\Delta_f(X)$ {\hyperlink{ega4.c04.prop.1-8-2}{(1.8.2)}}.
As the restriction of $p$ to $\Delta_f(X)$ is an isomorphism onto $X$, the set of $x\in X$ at which $f_{g(x)}$ is differentially smooth is locally constructible.

Finally, consider the property of being non ramified.
The diagonal morphism $\Delta_f:X\to X\times_Y X$ is an immersion locally of finite presentation (I.4.3.1), and for every $s\in S$ the diagonal $\Delta_{f_s}:X_s\to X_s\times_{Y_s}X_s$ is just $(\Delta_f)_s$.
To say that $f_{g(x)}$ is non ramified at $x$ amounts to saying that $(\Delta_f)_{g(x)}$ is a local isomorphism at $x$ {\hyperlink{ega4.c04.thm.17-4-1}{(17.4.1)}}; since this is an immersion locally of finite presentation, it amounts to the same thing {\hyperlink{ega4.c04.thm.17-9-1}{(17.9.1)}}\footnote{The reader may verify that the result of {\hyperlink{ega4.c04.prop.17-7-11}{(17.7.11)}} is not used in the remainder of \S17, so there is no vicious circle.} to say that $(\Delta_f)_{g(x)}$ is \'etale at $x$.
It is therefore enough to apply what has just been proved for the property of being \'etale.
\end{proof}


\subsection{Criteria for smoothness and non-ramification by fibres.}
\label{subsection:IV.17.8}
\label{IV.17.8}

\oldpage[IV]{79}

\begin{proposition}[17.8.1]
\label{IV.17.8.1}
Let $g:Y\to S$, $h:X\to S$ be two morphisms locally of finite presentation.
For an S-morphism $f:X\to Y$ to be non ramified, it is necessary and sufficient that for every $s\in S$, the morphism $f_s:h^{-1}(s)\to g^{-1}(s)$ deduced from $f$ by the change of base $\operatorname{Spec}(k(s))\to S$, be non ramified.
\end{proposition}

\begin{proof}
This follows immediately from the fact that, for a morphism locally of finite presentation, being non ramified is a property of the \emph{fibres} of this morphism (\sref{IV.17.4.1},~$d$)), and from the fact that for every $y\in Y$, one has $f^{-1}(y)=f_s^{-1}(y)$ where $s=g(y)$.
\end{proof}

\begin{proposition}[17.8.2]
\label{IV.17.8.2}
Let $g:Y\to S$, $h:X\to S$ be two morphisms locally of finite presentation; suppose moreover that $h$ is flat.
For an $S$-morphism $f:X\to Y$ to be smooth (resp.\ \'etale), it is necessary and sufficient that, for every $s\in S$, the morphism $f_s:h^{-1}(s)\to g^{-1}(s)$ induced from $f$ by the base change $\Spec(k(s))\to S$ be smooth (resp.\ \'etale).
When this is the case, the morphism $g$ is flat at all points of $f(X)$.
\end{proposition}

\begin{proof}
Indeed, one knows from {\hyperlink{ega4.c04.thm.11-3-10}{(11.3.10)}} that $f$ is flat if and only if $f_s$ is flat for every $s\in S$, and that in this case $g$ is flat at the points of $f(X)$.
But for a flat morphism locally of finite presentation, smoothness is a property of its fibres ({\hyperlink{ega4.c04.thm.17-5-1}{17.5.1}},~$b$)); and for every $y\in Y$, one has $f^{-1}(y)=f_s^{-1}(y)$ when $s=g(y)$.
\end{proof}

\begin{remark}[17.8.3]
\label{IV.17.8.3}
The preceding proofs show (taking (11.3.10) into account) that, under the same hypotheses on $g$ and $h$, for $f$ to be non ramified (resp.\ smooth, resp.\ \'etale) \emph{at a point} $x\in X$, it is enough that, on setting $s=h(x)$, $f_s$ be non ramified (resp.\ smooth, resp.\ \'etale) at $x$.
\end{remark}


\subsection{\'Etale morphisms and open immersions.}
\label{subsection:IV.17.9}
\label{IV.17.9}

\begin{theorem}[17.9.1]
\label{IV.17.9.1}
Let $f:X\to Y$ be a morphism.
The following properties are equivalent:
\begin{enumerate}
  \item[a)] $f$ is an open immersion.
  \item[b)] $f$ is a flat monomorphism locally of finite presentation.
  \item[c)] $f$ is \'etale and radiciel.
\end{enumerate}
\end{theorem}

\begin{proof}
It results from \sref[I]{I.4.3},~(i)) that $a)$ implies $b)$.
The condition $b)$ implies that for every $y\in Y$, the fibre $f^{-1}(y)$ is either empty or isomorphic to $\operatorname{Spec}(k(y))$ \sref{IV.8.11.5.1}, hence $b)$ implies $c)$, by virtue of \sref{IV.17.6.2},~$c$)).
It remains to show that $c)$ implies $a)$.
The question being local on $X$ and on $Y$ (since $f$ is injective), one can reduce to the case where $Y$ is affine and $f$ of finite presentation.
As $f$ is flat, this is an \emph{open} morphism \sref{IV.2.4.6}, hence, replacing $Y$ by $f(X)$, one can suppose that $f$ is \emph{surjective}.
For every morphism $Y'\to Y$, $f'=f_{(Y')}:X_{(Y')}\to Y'$ is again \'etale, radiciel, surjective, and of finite presentation; it is therefore open, hence a homeomorphism.
In other words, $f$ is a \emph{universal homeomorphism}; and since it is of finite type and separated (I.8.7.1,~(ii)), it is proper.
The hypothesis that $f$ is radiciel and of finite type implies that $f$ is quasi-finite; hence \sref{IV.8.11.1} $f$ is a finite morphism.
To prove that $f$ is an isomorphism, one can reduce to the case

\oldpage[IV]{80}

where $Y=\operatorname{Spec}(A)$, $A$ being a local ring.
As $f$ is of finite presentation, one has $X=\Spec(B)$, where $B$ is a flat $A$-module of finite presentation (I.4.7), hence free (Bourbaki, \emph{Alg.~comm.}, chap.~II, \S5, no.~2, Cor.~2 to Th.~1).
Moreover, if $\mathfrak{m}$ is the maximal ideal of $A$ and $k$ its residue field, $B/\mathfrak{m}B$ is by hypothesis a field, which is both radiciel extension and \'etale extension of $k$ \sref{IV.17.6.1}; hence $B/\mathfrak{m}B$ is isomorphic to $k$, and as $B$ is a free A-module, $B$ is isomorphic to $A$.
C.Q.F.D.
\end{proof}

\begin{corollary}[17.9.2]
\label{IV.17.9.2}
Let $X$ be a connected prescheme; if $f:X\to Y$ is a closed \'etale immersion, then $f$ is an isomorphism of $X$ onto an open connected component of $Y$.
\end{corollary}

\begin{proof}
Indeed $f$ is \'etale and radiciel, hence an isomorphism from $X$ onto a sub-prescheme of $Y$; but by hypothesis $f(X)$ is closed in $Y$, hence since open and closed, and since $f(X)$ is connected, it is a connected component of $Y$.
\end{proof}

\begin{corollary}[17.9.3]
\label{IV.17.9.3}
Let $f:X\to Y$ be a morphism \'etale (resp.\ \'etale and separated).
Then every Y-section $g:Y\to X$ of $X$ is an open immersion (resp.\ open and closed); moreover the application $g\mapsto g(Y)$ is a bijection of the set $\Gamma(X/Y)$ of Y-sections of $X$ to the set of open (resp.\ open and closed) parts $Z$ of $X$ such that the restriction of $f$ to $Z$ is a surjective morphism and radiciel of $Z$ onto $Y$.
\end{corollary}

\begin{proof}
Indeed, the fact that $g$ is an open immersion follows already from the fact that $f$ is non ramified \sref{IV.17.4.1},~$b''$)), and the restriction of $f$ to the open $g(Y)$ of $X$ is an isomorphism.
Conversely, if $Z$ is an open subset of $X$ such that $f|Z$ is surjective and radiciel, then, since $f$ is \'etale, $f|Z$ is an isomorphism by {\hyperlink{ega4.c04.thm.17-9-1}{(17.9.1)}}.
If moreover $f$ is separated, then $g$ is a closed immersion (I.5.4.6), which completes the proof.
\end{proof}

\begin{corollary}[17.9.4]
\label{IV.17.9.4}
Let $Y$ be a connected prescheme and $f:X\to Y$ an \'etale and separated morphism.
Then every Y-section $g$ of $X$ is an isomorphism of $Y$ onto an open connected component of $X$, and the application $g\mapsto g(Y)$ is a bijection of $\Gamma(X/Y)$ to the set of open connected components of $X$ such that the restriction of $f$ to $Z$ is a surjective morphism and radiciel of $Z$ onto $Y$.
\end{corollary}

\begin{corollary}[17.9.5]
\label{IV.17.9.5}
Let $g:Y\to S$, $h:X\to S$ be two morphisms locally of finite presentation; suppose moreover that $h$ is flat.
For an S-morphism $f:X\to Y$ to be an open immersion (resp.\ isomorphism), it is necessary and sufficient that for every $s\in S$, the morphism $f_s:h^{-1}(s)\to g^{-1}(s)$ deduced from $f$ by the change of base $\operatorname{Spec}(k(s))\to S$, be an open immersion (resp.\ isomorphism).
\end{corollary}

\begin{proof}
Indeed, if $f_s$ is an open immersion for every $s\in S$, it results from \sref{IV.17.8.3} that $f$ is \'etale, and as for every $y\in Y$, $f^{-1}(y)=f_s^{-1}(y)$ with $s=g(y)$, $f$ is radiciel; hence $f$ is an open immersion by virtue of \sref{IV.17.9.1}.
If moreover $f_s$ is surjective for every $s\in S$, $f$ is surjective, hence an isomorphism.
\end{proof}

The following proposition specifies \sref{IV.10.4.11}:
\begin{proposition}[17.9.6]
\label{IV.17.9.6}
Let $S$ be a prescheme and $X$ an $S$-prescheme of finite presentation.
Every $S$-endomorphism of $X$ which is a monomorphism is an automorphism of $X$.
\end{proposition}

\begin{proof}
Let $f:X\to S$ be the structural morphism and $g$ the $S$-endomorphism under consideration.
The question being local on $S$, one can suppose $S$ affine.
Using \sref{IV.8.9.1} and \sref{IV.8.10.5},~(i \emph{bis})), one is reduced to the case where $S=\operatorname{Spec}(A)$, where $A$ is a $\bb{Z}$-algebra of finite type, and hence $X$ is a $\bb{Z}$-prescheme of finite type.
It follows already from \sref{IV.10.4.11} that $g$ is a bijective morphism, since a monomorphism is radiciel (\sref{IV.8.11.5.1}); it will thus suffice to show

\oldpage[IV]{81}

that $g$ is an open immersion, and since $g$ is radiciel, it will suffice, by virtue of \sref{IV.17.9.1}, to prove that $g$ is \'etale.
Moreover, as the set of points where $g$ is \'etale is open and $X$ is a Jacobson prescheme (\sref{IV.10.4.7}), it suffices to show that $g$ is \'etale at every \emph{closed} point $z$ of $X$ (\sref{IV.10.3.1}).
Set $z'=g(z)$; it follows from \sref{IV.10.4.11},~(i) that $z'$ is also a closed point of $X$, and since $g$ is a monomorphism, the map $k(z')\to k(z)$ deduced from $g$ is an isomorphism.
For $g$ to be \'etale at the point $z$, it is thus necessary and sufficient that the canonical homomorphism $\hat{\sh{O}}_{X,z'}\to\hat{\sh{O}}_{X,z}$ be bijective (\sref{IV.17.6.3},~$e''$)).
We will show for this purpose that for every integer $n$, the canonical homomorphism $\sh{O}_{X,z'}/\mathfrak{m}_{z'}^{n+1}\to\sh{O}_{X,z}/\mathfrak{m}_z^{n+1}$ is bijective, from which the conclusion will immediately follow.
One can suppose that $z$ belongs to the finite set $T_{p,d}$ of closed points $t$ of $X$ such that $k(t)$ is an extension of $\bb{F}_p$ whose degree divides $d$; in which case the same holds for $z'$.
Denote by $\sh{J}$ the coherent Ideal of $\sh{O}_X$ such that $\sh{J}_x=\sh{O}_x$ for $x\notin T_{p,d}$, and $\sh{J}_x=\mathfrak{m}_x^{n+1}$ for $x\in T_{p,d}$.
Let $T_{p,d,n}$ be the closed sub-prescheme of $X$ defined by $\sh{J}$ and $j:T_{p,d,n}\to X$ the canonical injection; the composite $g\circ j:T_{p,d,n}\to X$ maps $T_{p,d}$ into itself, and for every $x\in T_{p,d}$ the homomorphism $\sh{O}_{X,g(x)}\to\sh{O}_{X,x}/\mathfrak{m}_x^{n+1}$ factors as
\[
  \sh{O}_{X,g(x)} \to \sh{O}_{X,g(x)}/\mathfrak{m}_{g(x)}^{n+1} \to \sh{O}_{X,x}/\mathfrak{m}_x^{n+1}.
\]
By (\sref[I]{I.4.9}), there therefore exists a unique endomorphism $g'$ of $T_{p,d,n}$ such that $j\circ g'=g\circ j$.
As $j$ is a monomorphism, and $g$ is also one by hypothesis, one deduces that $g'$ is also a monomorphism, and the proposition will be proven if one shows that $g'$ is an \emph{automorphism} of $T_{p,d,n}$.
But, for every $x\in T_{p,d,n}$, we have seen that $k(x)$ is a finite field, and since $\sh{O}_{X,x}$ is Noetherian, each quotient $\mathfrak{m}_x^h/\mathfrak{m}_x^{h+1}$ is a finite-dimensional $k(x)$-vector space.
It follows immediately that the local ring $\sh{O}_{X,x}/\mathfrak{m}_x^{n+1}$ of $T_{p,d,n}$ at $x$ has finitely many elements and, a fortiori, is a finite-type $\bb{Z}$-module.
As $T_{p,d,n}$ is the sum of the preschemes $\operatorname{Spec}(\sh{O}_{X,x}/\mathfrak{m}_x^{n+1})$ in finite number, it is a finite $\bb{Z}$-prescheme; the endomorphism $g'$ is therefore also finite (\textbf{II},~6.1.5,~(v)), hence proper.
But a proper monomorphism is a closed immersion (\sref{IV.8.11.5}); if $T_{p,d,n}=\operatorname{Spec}(B)$, $g'$ corresponds to a surjective endomorphism $\varphi:B\to B$ of the ring $B$.
But the set $B$ is finite, hence $\varphi$ is necessarily bijective.
\end{proof}


\subsection{Relative dimension of a smooth prescheme over another.}
\label{subsection:IV.17.10}

\begin{definition}[17.10.1]
\label{IV.17.10.1}
Let $f:X\to Y$ be a morphism locally of finite type.
The \emph{relative dimension} of $f$ at a point $x\in X$ (or the \emph{dimension of $X$ relative to $Y$ at the point $x$}) is the positive integer $\dim_x f=\dim_x(f^{-1}(f(x)))$.
\end{definition}

To say that $f$ is \emph{quasi-finite} at the point $x$ ($\mathbf{Err_{III}}$, 20) is equivalent to saying that $\dim_x f=0$.
One has seen \sref{IV.13.1.3} that the function $x\mapsto\dim_x f$ is upper semi-continuous.
One notes that, even when the morphism $f$ has the property $(S_1)$ (that is, $f$ is flat \sref{IV.6.8.1} and its fibres have no embedded associated prime cycle), the function $x\mapsto\dim_x f$ is not necessarily continuous, as shown by the example where $Y=\operatorname{Spec}(k)$, where $k$ is a field, and $X=\operatorname{Spec}(k[\mathrm{U},\mathrm{V},\mathrm{W}]/\mathfrak{p}\mathfrak{q})$, where $\mathfrak{p}=(\mathrm{W})$ and $\mathfrak{q}=(\mathrm{U})+(\mathrm{V}-\mathrm{W})$, the prime ideals of $k[\mathrm{U},\mathrm{V},\mathrm{W}]$ ($X$ being thus the reunion, in the affine space of $3$ dimensions, of a plane and of a line non parallel to this plane).

To say that a morphism locally of finite presentation $f:X\to Y$ is \'etale at a point $x\in X$ means that $f$ is smooth at $x$ and that $\dim_x f=0$ {\hyperlink{ega4.c04.thm.17-6-1}{(17.6.1)}}.

\begin{proposition}[17.10.2]
\label{IV.17.10.2}
Let $f:X\to Y$ be a smooth morphism.
For every $x\in X$, the locally free $\sh{O}_X$-module $\Omega^1_f$ {\hyperlink{ega4.c04.prop.17-2-3}{(17.2.3)}} has rank at $x$ equal to $\dim_x f$ (which implies that $x\mapsto\dim_x f$ is a continuous function on $X$).
\end{proposition}

\begin{proof}
Indeed, if $y=f(x)$, then $X_y=f^{-1}(y)$ is smooth over $k(y)$; if $f_y:X_y\to\Spec(k(y))$ is the structural morphism, one has $\Omega^1_{f_y}=\Omega^1_f\otimes_{\sh{O}_X}\sh{O}_{X_y}$ {\hyperlink{ega4.c04.prop.16-4-5}{(16.4.5)}}.
Thus one may reduce to the case where $Y=\Spec(k)$ is the spectrum of a field.
Moreover, by {\hyperlink{ega4.c04.prop.16-4-5}{(16.4.5)}} and {\hyperlink{ega4.c04.prop.17-7-1}{(17.7.1)}}, one may replace $k$ by an algebraically closed extension; in other words,

\oldpage[IV]{82}

one may suppose that $k$ is algebraically closed.
The set of $k$-rational points of $X$ is then dense in $X$ {\hyperlink{ega4.c04.cor.10-4-8}{(10.4.8)}}, so it is enough to consider the case where $x$ is rational over $k$.
Now {\hyperlink{ega4.c04.cor.16-4-12}{(16.4.12)}} identifies $(\Omega^1_{X/k})_x\otimes_{\sh{O}_{X,x}}k(x)$ with $\mathfrak{m}_x/\mathfrak{m}_x^2$; since $\sh{O}_{X,x}$ is a regular local ring {\hyperlink{ega4.c04.thm.17-5-1}{(17.5.1)}}, one has $\operatorname{rg}_{k(x)}(\mathfrak{m}_x/\mathfrak{m}_x^2)=\dim(\sh{O}_{X,x})$ ($\mathbf{0}$,~17.1.1).
Hence $\operatorname{rg}_{\sh{O}_{X,x}}(\Omega^1_{X/k})_x=\dim(\sh{O}_{X,x})$.
Since $k(x)=k$, one also has $\dim_x f=\dim(\sh{O}_{X,x})$ {\hyperlink{ega4.c04.cor.5-2-3}{(5.2.3)}}.
Q.E.D.
\end{proof}

We will prove a reciprocal of this result further on \sref{IV.17.15.5}.

\begin{corollary}[17.10.3]
\label{IV.17.10.3}
Let $f:X\to Y$, $g:Y\to Z$ be two smooth morphisms.
Then for every $x\in X$, one has
\[
\label{IV.17.10.3.1}
\tag{17.10.3.1}
\dim_x(g\circ f) = \dim_x f + \dim_{f(x)} g.
\]
\end{corollary}

\begin{proof}
Indeed, $g\circ f$ is smooth {\hyperlink{ega4.c04.prop.17-3-3}{(17.3.3)}}, so the three $\sh{O}_X$-modules $\Omega^1_{X/Y}$, $\Omega^1_{X/Z}$, and $f^*(\Omega^1_{Y/Z})$ are locally free ({\hyperlink{ega4.c04.prop.17-2-3}{17.2.3}} and $\mathbf{0}_\mathbf{I}$,~5.4.5); moreover, the rank at $x$ of $f^*(\Omega^1_{Y/Z})$ equals the rank at $y=f(x)$ of $\Omega^1_{Y/Z}$.
Equality {\hyperlink{ega4.c04.cor.17-10-3-1}{(17.10.3.1)}} is therefore a consequence of {\hyperlink{ega4.c04.prop.17-10-2}{(17.10.2)}} and the exactness of the sequence {\hyperlink{ega4.c04.eq.17-2-3-1}{(17.2.3.1)}}.
\end{proof}

\begin{corollary}[17.10.4]
\label{IV.17.10.4}
Let $f:X\to Y$ be a smooth morphism and $X'$ a sub-prescheme of $X$ such that the composite $X'\xrightarrow{j}X\xrightarrow{f}Y$ (where $j$ is the canonical injection) is smooth.
Then the conormal sheaf $\sh{N}_{X'/X}$ is a locally free $\sh{O}_{X'}$-module, and for every $x\in X'$ one has
\begin{equation}
\label{IV.17.10.4.1}
\tag{17.10.4.1}
  \dim_x f = \dim_x(f\circ j) + \operatorname{rg}_{\sh{O}_{X',x}}(\sh{N}_{X'/X})_x.
\end{equation}
\end{corollary}

\begin{proof}
Indeed, $\Omega^1_{X/Y}\otimes_{\sh{O}_X}\sh{O}_{X'}$ and $\Omega^1_{X'/Y}$ are locally free, and the exact sequence {\hyperlink{ega4.c04.prop.17-2-5-1}{(17.2.5.1)}} splits locally in a suitable neighbourhood of every point of $X'$.
Thus $\sh{N}_{X'/X}$ is locally free, and relation {\hyperlink{ega4.c04.eq.17-10-4-1}{(17.10.4.1)}} follows at once from the exactness of {\hyperlink{ega4.c04.prop.17-2-5-1}{(17.2.5.1)}}.
\end{proof}


\subsection{Smooth morphisms of smooth preschemes.}
\label{subsection:IV.17.11}

\begin{theorem}[17.11.1]
\label{IV.17.11.1}
Let $f:Y\to S$ and $h:X\to S$ be two morphisms locally of finite presentation, $g:X\to Y$ an S-morphism, $x$ a point of $X$; set $y=g(x)$, $s=f(y)=h(x)$.
The following conditions are equivalent:
\begin{enumerate}
  \item[a)] $f$ is smooth at $y$ and $g$ is smooth at $x$.
  \item[b)] $g$ and $h$ are smooth at $x$.
  \item[c)] $h$ is smooth at $x$, and the canonical homomorphism \sref{IV.16.4.18}
\end{enumerate}
\begin{equation}
\label{IV.17.11.1.1}
\tag{17.11.1.1}
  (g^*(\Omega^1_{Y/S}))_x \to (\Omega^1_{X/S})_x
\end{equation}
is left invertible (in other words is an isomorphism onto a direct factor of $(\Omega^1_{X/S})_x$).
\begin{enumerate}
  \item[c')] $h$ is smooth at $x$, and the canonical homomorphism
\end{enumerate}
\begin{equation}
\label{IV.17.11.1.2}
\tag{17.11.1.2}
  ((\Omega^1_{Y/S})_y\otimes_{\sh{O}_{Y,y}} k(y))\otimes_{k(y)} k(x) \to (\Omega^1_{X/S})_x\otimes_{\sh{O}_{X,x}} k(x)
\end{equation}
is injective.

Suppose moreover that the homomorphism $k(y)\to k(x)$ is bijective.
Then the preceding conditions are also equivalent to the following:

\oldpage[IV]{83}

\begin{enumerate}
  \item[d)] $h$ is smooth at $x$, and the canonical map
  $\mathrm{T}_{X/S}(x)\to\mathrm{T}_{Y/S}(y)$ \sref{IV.16.5.12}
  from the tangent vector space to $X$ at $x$ to the tangent vector space to $Y$ at $y$ is surjective.
\end{enumerate}
\end{theorem}

\begin{proof}
That $a)$ implies $b)$ follows at once from \sref{IV.17.3.3},~(iii), and $b)$ implies $c)$ by \sref{IV.17.2.3},~(ii).
To see that $c)$ and $c')$ are equivalent, note that under either condition $\Omega^1_{X/S}$ is locally free of finite type, by \sref{IV.17.2.3},~(i); it is therefore enough to apply ($\mathbf{0}$, \sref{0.19.1.12}) to the local ring $\sh{O}_{X,x}$ and to the homomorphism \sref{IV.17.11.1.1} of finite-type $\sh{O}_{X,x}$-modules.
When $k(x)=k(y)$, the linear tangent map $\mathrm{T}_{X/S}(x)\to\mathrm{T}_{Y/S}(y)$ is, by \sref{IV.16.5.12}, the transpose of \sref{IV.17.11.1.2}; this proves the equivalence of $c')$ and $d)$ in that case.

It remains to prove that $c)$ implies $a)$.
We may restrict ourselves to the case in which $S=\operatorname{Spec}(A)$, $Y=\operatorname{Spec}(B)$, and $X=\operatorname{Spec}(C)$ are affine.
The hypothesis implies, by \sref{IV.17.2.3},~(i), that
$(\Omega^1_{X/S})_x=(\Omega^1_{C/A})_{\mathfrak m}$ is a free $C_{\mathfrak m}$-module.
By \sref{IV.16.10.6}, there are elements $t_i$ $(1\leq i\leq r)$ of
$B_{\mathfrak p}=\sh{O}_{Y,y}$ such that their differentials generate
$(\Omega^1_{B/A})_{\mathfrak p}$ and their images in
$(\Omega^1_{C/A})_{\mathfrak m}$ are part of a basis of that free $C_{\mathfrak m}$-module.
Since $\Omega^1_{Y/S}$ and $\Omega^1_{X/S}$ are of finite presentation \sref{IV.16.4.22}, after replacing $X$ and $Y$ by suitable affine open neighbourhoods of $x$ and $y$, we may suppose that $\Omega^1_{C/A}$ is free and that the $t_i$ are the images of elements $s_i\in B$ whose differentials $d_{B/A}(s_i)$ generate $\Omega^1_{B/A}$ and whose images in $\Omega^1_{C/A}$ are part of a basis of that $C$-module (Bourbaki, \emph{Alg. comm.}, chap.~II, \S~5, n\textsuperscript{o}~1, prop.~2).

Let $\varphi:B'=A[\mathrm{T}_1,\ldots,\mathrm{T}_r]\to B$ be the $A$-homomorphism such that $\varphi(\mathrm{T}_i)=s_i$.
The corresponding dihomomorphism
$\Omega^1_{B'/A}\to\Omega^1_{B/A}$ ($\mathbf{0}$, \sref{0.20.5.2}) sends the basis elements $d_{B'/A}(\mathrm{T}_i)$ ($\mathbf{0}$, \sref{0.20.4.13}) to the $d_{B/A}(s_i)$, and is therefore surjective.
If $Y'=\operatorname{Spec}(B')$ and $u:Y\to Y'$ is the S-morphism corresponding to $\varphi$, it follows from \sref{IV.17.2.2} that $u$ is unramified.
If $u\circ g:X\to Y'$ is smooth at $x$, then \sref{IV.17.7.10} shows that $u$ is \'etale at $y$, and \sref{IV.17.3.5} then shows that $g$ is smooth at $x$; moreover, since the structural morphism $f':Y'\to S$ is smooth \sref{IV.17.3.8}, $f=f'\circ u$ is smooth at $y$.
The canonical images of the $d_{B'/A}(\mathrm{T}_i)$ in $\Omega^1_{C/A}$ are the images of the $d_{B/A}(s_i)$ and thus form part of a basis.
Consequently, in proving that $c)$ implies $a)$, we may restrict ourselves to the case $Y'=Y$, and hence suppose that $f$ is smooth.

By \sref{IV.17.8.2}, we may then reduce to the case in which $S=\operatorname{Spec}(k)$ is the spectrum of a field; by \sref{IV.17.7.1},~(ii), we may moreover suppose that $k$ is algebraically closed.
After replacing $X$ and $Y$ by suitable open neighbourhoods of $x$ and $y$, we may suppose that $f$ and $h$ are both smooth and that the canonical homomorphism
\[
g^*(\Omega^1_{Y/S}) \to \Omega^1_{X/S}
\]
is left invertible ($\mathbf{0}$, \sref{0.19.1.12}).
The set of $k$-rational points of $X$ is then very dense in $X$ \sref{IV.10.4.8}; thus, to prove that $g$ is smooth, it is enough to prove smoothness at every $k$-rational point $x\in X$, or equivalently that $\sh{O}_{X,x}$ is flat over $\sh{O}_{Y,y}$ and that $\sh{O}_{X,x}/\mathfrak m_y\sh{O}_{X,x}$ is regular \sref{IV.17.5.1}.

Now $y=g(x)$ is

\oldpage[IV]{84}

\emph{a fortiori} $k$-rational.
Since $\sh{O}_{Y,y}$ is then a formally smooth $k$-algebra for the discrete topologies and $\sh{O}_{Y,y}/\mathfrak m_y=k$, it follows from ($\mathbf{0}$, \sref{0.20.5.14}) that the canonical homomorphism
\[
  \mathfrak m_y/\mathfrak m_y^2
  \longrightarrow
  (\Omega^1_{Y/S})_y\otimes_{\sh{O}_{Y,y}}k(y)
\]
is bijective; likewise the canonical homomorphism
\[
  \mathfrak m_x/\mathfrak m_x^2
  \longrightarrow
  (\Omega^1_{X/S})_x\otimes_{\sh{O}_{X,x}}k(x)
\]
is bijective.
Condition $c')$, which is equivalent to $c)$, therefore says here that the canonical homomorphism
\[
  (\mathfrak m_y/\mathfrak m_y^2)\otimes_{k(y)}k(x)
  \longrightarrow \mathfrak m_x/\mathfrak m_x^2
\]
is injective.
Since $\sh{O}_{X,x}$ is regular, the conclusion follows from ($\mathbf{0}$, \sref{0.17.3.3}).
\end{proof}

\begin{corollary}[17.11.2]
\label{IV.17.11.2}
Under the general hypotheses of \sref{IV.17.11.1}, the following conditions are equivalent:
\begin{enumerate}
  \item[a)] $f$ is smooth at $y$ and $g$ is \'etale at $x$.
  \item[b)] $h$ is smooth at $x$ and $g$ is \'etale at $x$.
  \item[c)] $h$ is smooth at $x$, and the canonical homomorphism
\[
(g^*(\Omega^1_{Y/S}))_x \to (\Omega^1_{X/S})_x
\]
is bijective.
  \item[c')] $h$ is smooth at $x$, and the canonical homomorphism
\[
((\Omega^1_{Y/S})_y\otimes_{\sh{O}_{Y,y}} k(y))\otimes_{k(y)} k(x) \to (\Omega^1_{X/S})_x\otimes_{\sh{O}_{X,x}} k(x)
\]
is bijective.
\end{enumerate}
Suppose moreover that the homomorphism $k(y)\to k(x)$ is bijective.
Then the preceding conditions are also equivalent to the following:
\begin{enumerate}
  \item[d)] $h$ is smooth at $x$, and the canonical map $\mathrm{T}_{X/S}(x)\to\mathrm{T}_{Y/S}(y)$ \sref{IV.16.5.12} is bijective.
\end{enumerate}
\end{corollary}

\begin{proof}
Each of conditions $a)$, $b)$, and $c)$ of \sref{IV.17.11.2} is equivalent to the conjunction of the corresponding condition of \sref{IV.17.11.1} and the assertion that $g$ is unramified at $x$, taking \sref{IV.17.2.2} into account for condition $c)$; hence $a)$, $b)$, and $c)$ are equivalent.
The equivalence of $c)$ and $c')$ follows from the equivalence of the corresponding conditions of \sref{IV.17.11.1} and Nakayama's lemma; when $k(x)=k(y)$, the equivalence of $c')$ and $d)$ follows immediately by transposition \sref{IV.16.5.12}.
\end{proof}

\begin{corollary}[17.11.3]
\label{IV.17.11.3}
Let $h:X\to S$ be a smooth morphism, let $s_i$ $(1\leq i\leq r)$ be sections of $\sh{O}_X$ over $X$ (and thus also sections of $h_*(\sh{O}_X)$ over $S$), and let
$g:X\to S[\mathrm{T}_1,\ldots,\mathrm{T}_r]=\bb{V}^r_S$
be the S-morphism corresponding to the homomorphism of $\sh{O}_S$-modules
$\sh{O}^r_S\to h_*(\sh{O}_X)$ defined by these sections (\textbf{II}, \sref{IV.1.2.7}).
For $g$ to be smooth (resp.\ \'etale) at a point $x\in X$, it is necessary and sufficient that the $(d_{X/S}(s_i))_x$ form part of a basis (resp.\ form a basis) of the $\sh{O}_{X,x}$-module $(\Omega^1_{X/S})_x$.
\end{corollary}

\begin{proof}
It suffices to apply \sref{IV.17.11.1} (resp.\ \sref{IV.17.11.2}) taking for $f$ the structural morphism $S[\mathrm{T}_1,\ldots,\mathrm{T}_r]\to S$.
\end{proof}

\begin{corollary}[17.11.4]
\label{IV.17.11.4}
For a morphism $h:X\to S$ to be smooth at a point $x\in X$, it is necessary and sufficient that there exist an open neighbourhood $U$ of $x$, an integer $r$, and an \'etale S-morphism
$U\to S[\mathrm{T}_1,\ldots,\mathrm{T}_r]$.
\end{corollary}

\oldpage[IV]{85}

\begin{proof}
The condition is plainly sufficient, since the structural morphism $S[\mathrm{T}_1,\ldots,\mathrm{T}_r]\to S$ is smooth \sref{IV.17.3.8}.
Conversely, since $h$ is smooth at $x$, there is an open neighbourhood $U$ of $x$ such that $\Omega^1_{X/S}|U$ is locally free \sref{IV.17.2.3}; using \sref{IV.16.10.6} and restricting $U$ if necessary, we obtain sections $s_i$ of $\sh{O}_X$ over $U$ such that the $d_{X/S}(s_i)$ form a basis of $\Omega^1_{X/S}|U$.
The conclusion follows from \sref{IV.17.11.3}.
\end{proof}

\begin{proposition}[17.11.5]
\label{IV.17.11.5}
Let $f:Y\to S$, $h:X\to S$ be two smooth morphisms.
For an S-morphism $g:X\to Y$ to be an open immersion, it is necessary and sufficient that $g$ be a monomorphism of preschemes and that for every $x\in X$, if one sets $y=g(x)$, one has $\dim_y(f)=\dim_x(h)$ (for a generalisation of this proposition, see \sref{IV.18.10.5}).
\end{proposition}

\begin{proof}
The condition is evidently necessary; let us prove that it is sufficient.
By \sref{IV.17.9.5}, we are immediately reduced to the case in which $S=\operatorname{Spec}(k)$ is the spectrum of a field, and by \sref{IV.2.7.1},~(x), we may suppose that $k$ is algebraically closed.
In view of \sref{IV.17.9.1}, it is then enough to prove that $g$ is \'etale; since the locus where $g$ is \'etale is open, it is enough to prove this at the closed points of $X$, or equivalently at its $k$-rational points \sref{IV.10.4.8}.

Let $x$ be such a point and set $y=g(x)$; then $y$ too is $k$-rational.
Put $A=\sh{O}_{Y,y}$, a regular local ring with residue field $k$, let $d=\dim(A)$, and let $(t_i)_{1\leq i\leq d}$ be a regular system of parameters of $A$.
Put $B=\sh{O}_{X,x}$ and $C=B/\mathfrak m_yB$.
Since $g$ is a monomorphism, so is the morphism
$\operatorname{Spec}(C)\to\operatorname{Spec}(k(y))=\operatorname{Spec}(k)$ obtained from $g$ by base change ($\mathbf{I}$, \sref{I.3.3.12}).
This means that the corresponding homomorphism $u:k\to C$ is surjective, and hence bijective: it has a left inverse $v:C\to k$, while $u\circ v$ and the identity of $C$, after composition with $u$, give the same homomorphism $u:k\to C$; the monomorphism property therefore gives $u\circ v=1_C$.

By hypothesis $B$ is a regular ring of dimension $d$.
The images of the $t_i$ in $B$ consequently form a regular system of parameters of $B$ ($\mathbf{0}$, \sref{0.17.1.7}); condition \sref{IV.17.6.3},~$e''$), is therefore satisfied by $g$ at $x$ ($\mathbf{0}$, \sref{0.17.1.1}, and Bourbaki, \emph{Alg. comm.}, chap.~III, \S~2, n\textsuperscript{o}~8, cor.~3 to th.~1), which completes the proof.
\end{proof}


\subsection{Smooth subpreschemes of a smooth prescheme. Smooth morphisms and differentially smooth morphisms.}
\label{subsection:IV.17.12}

\begin{theorem}[17.12.1]
\label{IV.17.12.1}
Let $f:X\to S$, $h:Y\to S$ be two morphisms locally of finite presentation, $j:Y\to X$ an immersion, $y$ a point of $Y$, $x=j(y)$.
The following conditions are equivalent:
\begin{enumerate}
  \item[a)] $h$ is smooth at $y$ and $f$ is smooth at $x$.
  \item[b)] $f$ is smooth at $x$, and the canonical homomorphism \sref{IV.16.4.21}
\[
(\sh{N}_{Y/X})_y \to (j^*(\Omega^1_{X/S}))_y
\]
is left invertible.
  \item[c)] $h$ is smooth at $y$ and there exists an open neighbourhood $U$ of $y$ in $Y$ such that $j|U:U\to X$ is a regular immersion \sref{IV.16.9.2}.
\oldpage[IV]{86}

  \item[c')] $h$ is smooth at $y$ and there exists an open neighbourhood $U$ of $y$ in $Y$ such that $j|U:U\to X$ is a quasi-regular immersion; in other words \sref{IV.16.9.8}, there is a neighbourhood $U$ of $y$ in $Y$ on which $\sh{N}_{Y/X}$ is locally free and the canonical homomorphism
  \[
    \mathbf{S}_{\sh{O}_Y}(\sh{N}_{Y/X})
    \longrightarrow \operatorname{gr}^{\bullet}(j)
  \]
  is bijective.
\end{enumerate}
\end{theorem}

\begin{proof}
To prove the equivalence of $a)$ and $b)$, we may restrict ourselves to the case in which $S=\operatorname{Spec}(A)$ and $X=\operatorname{Spec}(B)$ are affine, with
$Y=\operatorname{Spec}(B/\mathfrak J)$, where $\mathfrak J$ is a finite-type ideal of $B$, and where $B$ is a smooth $A$-algebra \sref{IV.17.3.2},~(ii).
It is then enough to apply the Jacobian criterion ($\mathbf{0}$, \sref{0.22.6.1}) together with ($\mathbf{0}$, \sref{0.19.1.14}), bearing in mind that $\Omega^1_{X/S}$ is locally free in a neighbourhood of $x$ and that $\sh{N}_{Y/X}$ is of finite type \sref{IV.16.1.6}.

We next prove that $a)$ implies $c')$.
We may again suppose that $S$, $X$, and $Y$ are affine, and that $B$ and $B/\mathfrak J$ are smooth $A$-algebras.
It follows from \sref{IV.17.7.9} and \sref{IV.8.10.5},~(iv), that there exist a Noetherian subring $A'$ of $A$, a finite-type $A'$-algebra $B'$, and an ideal $\mathfrak J'$ of $B'$ such that
\[
  B=B'\otimes_{A'}A,\qquad \mathfrak J=\mathfrak J'B,
\]
and such that $B'$ and $B'/\mathfrak J'$ are smooth $A'$-algebras.
By ($\mathbf{0}$, \sref{0.19.3.8}), $B'$ is formally smooth when $A'$ is given the discrete topology and $B'$ the $\mathfrak J'$-preadic topology.
Consequently ($\mathbf{0}$, \sref{0.19.5.4}), $\mathfrak J'/\mathfrak J'^2$ is a projective $B'/\mathfrak J'$-module and the canonical homomorphism
\[
  \mathbf{S}_{B'/\mathfrak J'}(\mathfrak J'/\mathfrak J'^2)
  \longrightarrow \operatorname{gr}^{\bullet}_{\mathfrak J'}(B')
\]
is bijective.
Thus the immersion $\operatorname{Spec}(B'/\mathfrak J')\to\operatorname{Spec}(B')$ is quasi-regular \sref{IV.16.9.8}, and hence regular because $B'$ is Noetherian \sref{IV.16.9.10}.
Moreover, $B'/\mathfrak J'$ is flat over $A'$ \sref{IV.17.5.1} and $B'$ is an $A'$-algebra of finite presentation; applying \sref{IV.11.3.8} with $\widetilde{\mathfrak J'}$ in place of $\sh F$ and then changing base shows that $j$ is a regular immersion, and a fortiori a quasi-regular immersion.
The fact that $B$ is an $A$-algebra of finite presentation and is flat over $A$ also shows, by \sref{IV.11.3.8}, that conditions $c')$ and $c)$ are equivalent and are stable under base change.

It remains to prove that $c)$ implies $a)$.
Since condition $b)$ of \sref{IV.11.3.8} implies condition $c)$ there, condition $c)$ above already shows that $f$ is flat at $x$.
By \sref{IV.17.5.1}, it is therefore enough to prove that the fibre $f^{-1}(s)$, where $s=h(y)=f(x)$, is smooth over $k(s)$ at $x$.
Since condition $c)$ is stable under base change, we reduce to the case in which $S=\operatorname{Spec}(k)$ is the spectrum of a field; in view of \sref{IV.17.7.1},~(ii), we may suppose that $k$ is algebraically closed.
It is enough to prove that $\sh{O}_{X,x}$ is regular \sref{IV.17.5.1}.
By hypothesis,
\[
  \sh{O}_{Y,y}=\sh{O}_{X,x}/\mathfrak J_x,
\]
where $\mathfrak J_x$ is generated by an $\sh{O}_{X,x}$-regular sequence $(t_i)$, and $\sh{O}_{Y,y}$ is regular.
The $t_i$ form part of a system of parameters of $\sh{O}_{X,x}$ ($\mathbf{0}$, \sref{0.16.4.1}); the conclusion follows from ($\mathbf{0}$, \sref{0.17.1.7}).
\end{proof}

\begin{corollary}[17.12.2]
\label{IV.17.12.2}
With the notation of \sref{IV.17.12.1}, suppose that $f$ is smooth at $x$, and let $Y$ be a closed subprescheme of $X$, defined by an ideal $\sh{J}$ of $\sh{O}_X$, which is smooth over $S$ at $x$.
Let $(g_i)_{1\leq i\leq r}$ be sections of $\sh{J}$ over $X$.
The following conditions are equivalent:
\begin{enumerate}
  \item[a)] The $(g_i)_x$ form a system of generators of the $\sh{O}_{X,x}$-module $\sh{J}_x$ having the smallest possible number of elements.
  \item[b)] If $g'_i$ denotes the canonical image of $g_i$ in $\Gamma(X,\sh{J}/\sh{J}^2)$, the $(g'_i)_x$ form a basis of the $(\sh{O}_{X,x}/\sh{J}_x)$-module $\sh{J}_x/\sh{J}_x^2$.
\oldpage[IV]{87}

  \item[c)] The images of the $d_{X/S}g_i$ in
  $(\Omega^1_{X/S})_x\otimes_{\sh{O}_{X,x}}k(x)$ are linearly independent over $k(x)$, and the $(g_i)_x$ generate $\sh{J}_x$.
  \item[d)] There exist an open neighbourhood $U$ of $x$ in $X$ and sections $g_j$ $(r+1\leq j\leq n)$ of $\sh{O}_X$ over $U$ such that the S-morphism
  \[
    u:U\longrightarrow S[\mathrm{T}_1,\ldots,\mathrm{T}_n]
  \]
  corresponding to the homomorphism $\sh{O}^n_S\to f_*(\sh{O}_U)$ defined by the sections $g_i|U$ $(1\leq i\leq r)$ and $g_j$ $(r+1\leq j\leq n)$ (\textbf{II}, \sref{IV.1.2.7}) is \'etale, and the inverse image under this morphism of the subprescheme
  $Y'=S[\mathrm{T}_{r+1},\ldots,\mathrm{T}_n]$
  is the subprescheme induced by $j(Y)$ on the open set $U\cap j(Y)$.
\end{enumerate}
\end{corollary}

\begin{proof}
Since $(g'_i)_x$ is the canonical image of $(g_i)_x$, the equivalence of $a)$ and $b)$ follows from Nakayama's lemma: $\sh{J}_x$ is of finite type and $\sh{J}_x/\sh{J}_x^2$ is a free $(\sh{O}_{X,x}/\sh{J}_x)$-module \sref{IV.17.10.4} (Bourbaki, \emph{Alg. comm.}, chap.~II, \S~3, n\textsuperscript{o}~2, prop.~5).
By \sref{IV.17.12.1},~$b)$, the module $\sh{J}_x/\sh{J}_x^2$ is canonically identified with a direct factor of the free $(\sh{O}_{X,x}/\sh{J}_x)$-module
\[
  (\Omega^1_{X/S})_x\otimes_{\sh{O}_{X,x}}(\sh{O}_{X,x}/\sh{J}_x),
\]
and the equivalence of $b)$ and $c)$ follows from Bourbaki, \emph{loc. cit.}

If $a)$ holds, $(g'_i)_x$ is identified with $(d_{X/S}g_i)_x\otimes1$ for $1\leq i\leq r$.
Since $(\Omega^1_{X/S})_x$ is free of rank $n$, Bourbaki, \emph{loc. cit.}, shows that, after replacing $X$ by an open neighbourhood of $x$, there exist $n-r$ sections $g_i$ $(r+1\leq i\leq n)$ of $\sh{O}_X$ over $X$ such that the $(d_{X/S}g_i)_x$ for $1\leq i\leq n$ form a basis of $(\Omega^1_{X/S})_x$.
The corresponding morphism $u$ is then \'etale at $x$ by \sref{IV.17.11.3}; after a further restriction, we may suppose it is \'etale.
For the same reason, the restriction $u^{-1}(Y')\to Y'$ is \'etale at $x$, and we may suppose both morphisms \'etale near $x$.
Moreover, $Y$ (identified for simplicity with the closed subprescheme $j(Y)$ of $X$) is plainly a closed subprescheme of $u^{-1}(Y')$; by the choice of the $g_i$ for $i>r$, together with \sref{IV.17.2.5} and \sref{IV.17.11.2}, the restriction of $u$ to $Y$ may also be supposed \'etale.
It follows that for every $y'\in Y'$, the immersion
$Y\cap u^{-1}(y')\to u^{-1}(y')$ is open \sref{IV.17.9.6}; hence $Y\to u^{-1}(Y')$ is an open immersion.
This proves that $a)$ implies $d)$.
The converse $d)\Rightarrow a)$ follows at once from \sref{IV.17.11.3}.
\end{proof}

\begin{corollary}[17.12.3]
\label{IV.17.12.3}
Let $f:X\to S$ be a morphism locally of finite presentation, $u:S\to X$ an S-section of $X$.
For $f$ to be smooth at a point $x=u(s)$, where $s\in S$, it is necessary and sufficient that $u$ be a quasi-regular immersion in a neighbourhood of $s=f(x)$.
\end{corollary}

\begin{proof}
The conclusion follows from \sref{IV.17.12.1}, since $f\circ u=1_S$ is smooth.
\end{proof}

\begin{proposition}[17.12.4]
\label{IV.17.12.4}
Every smooth morphism $f:X\to S$ is differentially smooth (in other words \sref{IV.16.10.5} $\Delta_f:X\to X\times_S X$ is a quasi-regular immersion).
In particular, the $\sh{P}^n_f=\sh{P}^n_{X/S}$ and the $\operatorname{gr}_n(\sh{P}_{X/S})$ are locally free $\sh{O}_X$-modules of finite type.
\end{proposition}

\begin{proof}
It is enough to note that the structural morphism $p_1:X\times_S X\to X$ is smooth \sref{IV.17.3.3},~(iii), and that $\Delta_f$ is an X-section of $p_1$; the conclusion follows from \sref{IV.17.12.3}.
\end{proof}

\begin{proposition}[17.12.5]
\label{IV.17.12.5}
Let $f:X\to Y$ be a morphism locally of finite presentation.
The following conditions are equivalent:
\begin{enumerate}
  \item[a)] $f$ is differentially smooth.
  \item[b)] For every morphism $g:Y'\to Y$, if $X'=X\times_Y Y'$ and $f'=f_{(Y')}:X'\to Y'$, then for every Y'-section $s'$ of $X'$, $f'$ is smooth at all points of $s'(Y')$.
  \item[c)] The second projection $p_2:X\times_Y X\to X$ is smooth at all points of the diagonal $\Delta_f(X)$.
\end{enumerate}
\end{proposition}

\oldpage[IV]{88}

\begin{proof}
Condition $c)$ is a special case of $b)$: take $Y'=X$ and $g=f$ in $b)$; then $f'$ is the second projection $p_2$, and $\Delta_f$ is an X-section of $X\times_Y X$.
Next, $c)$ implies $a)$, since $p_2$ is locally of finite presentation and, if $p_2$ is smooth at the points of $\Delta_f(X)$, then $\Delta_f$ is a quasi-regular immersion by \sref{IV.17.12.3}; hence $f$ is differentially smooth.

Finally, suppose $a)$ and prove $b)$.
If $\Delta_f$ is quasi-regular, then $p_2$ is smooth at every point of $\Delta_f(X)$ by \sref{IV.17.12.3}.
Let $g':X'\to X$ be the canonical projection and let
$v=(g',g')_Y:X'\to X\times_Y X$.
Then the diagram
\[
\xymatrix{
  X\times_Y X \ar[d]_{p_2} & X' \ar[l]_-v \ar[d]^{f'}\\
  X & Y' \ar[l]^{h=g'\circ s'}
}
\]
is commutative and identifies $X'$ with the product
$(X\times_Y X)\times_XY'$ ($\mathbf{I}$, \sref{I.3.3.9}).
Taking into account that diagram \sref{IV.17.4.1.1} identifies $Y'$ with the product of the $(X\times_Y X)$-preschemes $X$ and $X'$, we conclude from the smoothness of $p_2$ at every point of $\Delta_f(X)$ that $f'$ is smooth at every point of $s'(Y')$ \sref{IV.17.3.3},~(iii).
\end{proof}

\begin{corollary}[17.12.6]
\label{IV.17.12.6}
With the notations being those of \sref{IV.8.8.1}, suppose $X_\alpha$ quasi-compact, and $f_\alpha:X_\alpha\to S_\alpha$ locally of finite presentation.
For $f:X\to S$ to be differentially smooth, it is necessary and sufficient that there exist $\lambda\geq\alpha$ such that $f_\lambda:X_\lambda\to S_\lambda$ is differentially smooth (in which case $f_\mu:X_\mu\to S_\mu$ is differentially smooth for $\mu\geq\lambda$).
\end{corollary}

\begin{proof}
The sufficiency of the condition and the last assertion follow from \sref{IV.16.10.4}.
To prove necessity, note that $\Delta_f:X\to X\times_SX$ and $p_2:X\times_SX\to X$ are obtained by base change from
$\Delta_{f_\lambda}:X_\lambda\to X_\lambda\times_{S_\lambda}X_\lambda$ and
$p_{2,\lambda}:X_\lambda\times_{S_\lambda}X_\lambda\to X_\lambda$.
By \sref{IV.17.12.5}, for every $z\in\Delta_f(X)$ there is an affine open neighbourhood $U(z)$ of $z$ in $X\times_SX$ on which $p_2$ is smooth.
By \sref{IV.8.2.11} and \sref{IV.17.7.8}, there exist an index $\lambda(z)$ and a neighbourhood $U_{\lambda(z)}$ of the projection of $z$ in
$X_{\lambda(z)}\times_{S_{\lambda(z)}}X_{\lambda(z)}$ such that $U(z)$ is the inverse image of $U_{\lambda(z)}$ and $p_{2,\lambda(z)}$ is smooth on $U_{\lambda(z)}$.
Since $X$ is quasi-compact, $\Delta_f(X)$ is covered by finitely many such neighbourhoods $U(z_i)$.
By \sref{IV.8.3.4}, there is an index $\lambda$ greater than all the $\lambda(z_i)$ such that the inverse images of the $U_{\lambda(z_i)}$ in $X_\lambda\times_{S_\lambda}X_\lambda$ cover $\Delta_{f_\lambda}(X_\lambda)$.
It follows from \sref{IV.17.12.5} that this $\lambda$ has the required property.
\end{proof}

\begin{example}[17.12.7]
\label{IV.17.12.7}
Let $S$ be a prescheme, $G$ an S-\emph{prescheme in groups}, locally of finite presentation on $S$.
For $G$ to be differentially smooth over $S$, it is necessary and sufficient that $G$ be smooth over $S$ at all points of the ``identity section'' $e$ of $G$.
The condition is indeed necessary by \sref{IV.17.12.5}.
Conversely, suppose it holds.
For every morphism $S'\to S$, $G'=G\times_SS'$ is an $S'$-prescheme in groups locally of finite presentation over $S'$; by \sref{IV.17.12.5}, it is therefore enough to prove that for every S-section $s$ of $G$, the structural morphism is smooth at the points of $s(S)$.
If $m:G\times_SG\to G$ is the multiplication, then
$m\circ(s\times1_G):S\times_SG\to G\times_SG\to G$

\oldpage[IV]{89}

is an S-isomorphism of $G$, the ``left translation'' $\tau_s$, carrying the points of the identity section of $G$ to the points of $s(S)$.
Thus $G$ is smooth over $S$ at the points of $s(S)$.
\end{example}

\begin{remark}[17.12.8]
\label{IV.17.12.8}
An S-prescheme in groups $G$, of finite type (and even finite) over a locally Noetherian prescheme $S$, can be differentially smooth over $S$ without being smooth (or even \emph{flat}) over $S$.
For example, take $S$ to be the spectrum of the dual-number algebra
$D=k[\mathrm{T}]/\mathrm{T}^2k[\mathrm{T}]$ over a field $k$, and take $G$ to be the spectrum of the direct-product $D$-algebra $E=D\oplus k$.
Define an S-group-prescheme structure on $G$ by a ``diagonal map'' $\delta:E\to E\otimes_DE$ as follows.
If $e'$ and $e''$ are the idempotents given by the canonical images of $1$ of $D$ in the two factors $D$ and $k$ of $E$, then
\[
  \delta(e')=e'\otimes e'+e''\otimes e'',\qquad
  \delta(e'')=e'\otimes e''+e''\otimes e'.
\]
The identity section of the group scheme $G$ corresponds to the homomorphism $E\to D$ that is the identity on $D$ and zero on $k$.
It is an isomorphism of $S$ onto a connected component of $G$; hence $G$ is differentially smooth over $S$ \sref{IV.17.12.6}, whereas it is clear that $G$ is not flat over $S$.
\end{remark}


\subsection{Transversal morphisms.}
\label{subsection:IV.17.13}
\label{IV.17.13}

\begin{env}[17.13.1]
\label{IV.17.13.1}
Let $S$ be a prescheme, let $X$, $Y$, and $X'$ be three S-preschemes, let $i:Y\to X$ be an S-immersion, and let $f:X'\to X$ be an S-morphism.
Set $Y'=Y\times_X X'$, and let $g:Y'\to Y$, $j:Y'\to X'$ be the canonical projections, so that one has the commutative diagram
\[
\label{IV.17.13.1.1}
\tag{17.13.1.1}
\xymatrix{
  Y \ar[d]_i & Y' \ar[l]_g \ar[d]^j\\
  X & X' \ar[l]^f
}
\]
and $j$ is an S-immersion.
There is then a commutative diagram of quasi-coherent $\sh{O}_{Y'}$-modules
\[
\label{IV.17.13.1.2}
\tag{17.13.1.2}
\begin{tikzcd}
  g^*(\sh{N}_{Y/X}) \ar{r} \ar{d}{\operatorname{gr}_1(g)} & f^*(\Omega^1_{X/S})\otimes_{\sh{O}_{X'}} \sh{O}_{Y'} \ar{r} \ar{d} & g^*(\Omega^1_{Y/S}) \ar{r} \ar{d} & 0 \\
  \sh{N}_{Y'/X'} \ar{r} & \Omega^1_{X'/S}\otimes_{\sh{O}_{X'}} \sh{O}_{Y'} \ar{r} & \Omega^1_{Y'/S} \ar{r} & 0
\end{tikzcd}
\]
Here the lower row is the exact sequence \sref{IV.16.4.21} applied to $j$; the upper row is obtained from the same exact sequence for $i$ by applying the right-exact functor $g^*$, which therefore preserves its exactness.
The homomorphism $\operatorname{gr}_1(g)$ is defined in \sref{IV.16.2.1}, and the commutativity of the two squares follows from ($\mathbf{0}$, \sref{0.20.5.7.3}) and ($\mathbf{0}$, \sref{0.20.5.11.3}).

\oldpage[IV]{90}

Moreover, we saw in \sref{IV.16.2.2},~(iii), that $\operatorname{gr}_1(g)$ is surjective here; hence diagram \sref{IV.17.13.1.2} gives the exact sequence
\[
\label{IV.17.13.1.3}
\tag{17.13.1.3}
  g^*(\sh{N}_{Y/X}) \xrightarrow{\alpha} \Omega^1_{X'/S}\otimes_{\sh{O}_{X'}} \sh{O}_{Y'} \to \Omega^1_{Y'/S} \to 0.
\]
\end{env}

\begin{proposition}[17.13.2]
\label{IV.17.13.2}
With the notation of \sref{IV.17.13.1}, let $x'$ be a point of $Y'$ and let $x=g(x')$ be its image in $Y$; suppose that $X$ and $Y$ are smooth over $S$ at $x$, and that $X'$ is smooth over $S$ at $x'$.
Let $m$ and $m-c$ be the relative dimensions of $X$ and $Y$ over $S$ at $x$ \sref{IV.17.10.1}, and let $n$ be that of $X'$ over $S$ at $x'$.
The following conditions are equivalent:
\begin{enumerate}
  \item[a)] $Y'$ is smooth over $S$ at $x'$ and has relative dimension $n-c$.
  \item[b)] The homomorphism $\alpha\otimes 1:g^*(\sh{N}_{Y/X})\otimes_{\sh{O}_{Y'}} k(x') \to \Omega^1_{X'/S}\otimes_{\sh{O}_{X'}} k(x')$ deduced from the homomorphism $\alpha$ of \sref{IV.17.13.1.3} is injective.
\end{enumerate}
Let $s$ be the canonical image of $x$ in $S$, and suppose that $x'$ is rational over $k(s)$.
Then the conditions $a)$ and $b)$ are also equivalent to the following:
\begin{enumerate}
  \item[b')] The transposed homomorphism of $\alpha\otimes 1$
\[
\mathrm{T}_{X'/S}(x') \to (\sh{N}_{Y/X}\otimes_{\sh{O}_Y} k(x'))^\vee = \mathrm{T}_{X/S}(x)/\mathrm{T}_{Y/S}(x)
\]
(cf.\ \sref{IV.16.5.12}) is surjective.
\end{enumerate}
Moreover, when conditions $a)$ and $b)$ hold at $x'$, they hold throughout a neighbourhood of $x'$ in $Y'$; after replacing $X'$ by a neighbourhood of $x'$, the homomorphism
\[
  \operatorname{gr}_1(g):g^*(\sh{N}_{Y/X})\longrightarrow\sh{N}_{Y'/X'}
\]
is bijective, and the sequence
\begin{equation}
\label{IV.17.13.2.1}
\tag{17.13.2.1}
  0\longrightarrow g^*(\sh{N}_{Y/X})
  \xrightarrow{\alpha}
  \Omega^1_{X'/S}\otimes_{\sh{O}_{X'}}\sh{O}_{Y'}
  \longrightarrow\Omega^1_{Y'/S}\longrightarrow0
\end{equation}
obtained by adjoining a zero to \sref{IV.17.13.1.3} is exact.
\end{proposition}

\begin{proof}
If the equivalent conditions $a)$ and $b)$ hold at $x'$, they also hold in a neighbourhood of $x'$ in $Y'$: this follows from the openness of the smooth locus \sref{IV.17.3.7} and from \sref{IV.17.10.2}.

By \sref{IV.17.12.1}, applied to $X'$ and $Y'$, saying that $Y'$ is smooth over $S$ at $x'$ is equivalent to saying that the homomorphism
\[
  \delta\otimes1:
  \sh{N}_{Y'/X'}\otimes_{\sh{O}_{Y'}}k(x')
  \longrightarrow
  \Omega^1_{X'/S}\otimes_{\sh{O}_{X'}}k(x')
\]
is injective, taking account of ($\mathbf{0}$, \sref{0.19.1.12}) and of the fact that $\Omega^1_{X'/S}$ is locally free at $x'$ \sref{IV.17.2.3}.
When this holds, $\Omega^1_{Y'/S}$ is locally free in a neighbourhood of $x'$ and the sequence
\begin{equation}
\label{IV.17.13.2.2}
\tag{17.13.2.2}
  0\longrightarrow\sh{N}_{Y'/X'}
  \longrightarrow\Omega^1_{X'/S}\otimes_{\sh{O}_{X'}}\sh{O}_{Y'}
  \longrightarrow\Omega^1_{Y'/S}\longrightarrow0
\end{equation}
is exact \sref{IV.17.2.5}.
Thus, by \sref{IV.17.10.2}, saying in addition that $Y'$ has relative dimension $n-c$ over $S$ at $x'$ means that $\sh{N}_{Y'/X'}$, which is locally free near $x'$, has rank $c$ at $x'$.
By the same reasoning and the hypotheses on $X$ and $Y$, $\sh{N}_{Y/X}$ is locally free of rank $c$ at $x$; hence $g^*(\sh{N}_{Y/X})$ is locally free of rank $c$ at $x'$.
Since $\operatorname{gr}_1(g):g^*(\sh{N}_{Y/X})\to\sh{N}_{Y'/X'}$ is surjective, the preceding conditions are equivalent to its being

\oldpage[IV]{91}

bijective at $x'$ (Bourbaki, \emph{Alg. comm.}, chap.~II, \S~3, n\textsuperscript{o}~2, cor. to prop.~6), and therefore also in a neighbourhood of $x'$ ($\mathbf{0}_{\mathrm I}$, \sref{0.5.2.7}).
This plainly implies $b)$ and the final assertion of the statement, by exactness of \sref{IV.17.13.2.2}.

Conversely, $\alpha\otimes1$ factors as
\[
  g^*(\sh{N}_{Y/X})\otimes k(x')
  \xrightarrow{\operatorname{gr}_1(g)\otimes1}
  \sh{N}_{Y'/X'}\otimes k(x')
  \xrightarrow{\delta\otimes1}
  \Omega^1_{X'/S}\otimes k(x').
\]
Since $\operatorname{gr}_1(g)$ is surjective, injectivity of $\alpha\otimes1$ implies both that $\delta\otimes1$ is injective and that $\operatorname{gr}_1(g)\otimes1$ is bijective.
It follows from \sref{IV.17.12.1} that $Y'$ is smooth over $S$ at $x'$, and $\operatorname{gr}_1(g)$ is bijective in a neighbourhood of $x'$ (Bourbaki, \emph{loc. cit.}).
The exact sequence \sref{IV.17.13.2.2} then shows that $\sh{N}_{Y'/X'}$, isomorphic to $g^*(\sh{N}_{Y/X})$, has rank $c$ at $x'$, while $\Omega^1_{Y'/S}$ has rank $n-c$ there.
By \sref{IV.17.10.2}, this completes the proof of the equivalence of $a)$ and $b)$.

It remains to prove the equivalence of $b)$ and $b')$ when $x'$ is rational over $k(s)$; this also implies that $x$ is rational over $k(s)$.
Then $g^*(\sh{N}_{Y/X})\otimes_{\sh{O}_{Y'}}k(x')$ identifies with
$\sh{N}_{Y/X}\otimes_{\sh{O}_Y}k(x')$.
Since the sequence
\[
  0\longrightarrow\sh{N}_{Y/X}
  \longrightarrow\Omega^1_{X/S}\otimes_{\sh{O}_X}\sh{O}_Y
  \longrightarrow\Omega^1_{Y/S}\longrightarrow0
\]
is exact at $x$ and consists there of locally free $\sh{O}_Y$-modules, the dual of
$\sh{N}_{Y/X}\otimes_{\sh{O}_Y}k(x')$ identifies with
$\mathrm{T}_{X/S}(x)/\mathrm{T}_{Y/S}(x)$ \sref{IV.16.5.12}.
Thus injectivity in $b)$ is equivalent to surjectivity of the transposed homomorphism in $b')$.
\end{proof}

\begin{definition}[17.13.3]
\label{IV.17.13.3}
With the notation of \sref{IV.17.13.1}, we say that the morphism $f$ is \emph{transversal} to $Y$ at $x'$, relative to $S$, if $X$ and $Y$ are smooth over $S$ at $x$, $X'$ is smooth over $S$ at $x'$, and the equivalent conditions $a)$ and $b)$ of \sref{IV.17.13.2} hold.
When there is no confusion to fear, we suppress the mention of the prescheme $S$ and simply say that $f$ is transversal to $Y$ at the point $x'$.
\end{definition}

\begin{remarks}[17.13.4]
\label{IV.17.13.4}
(i)
Suppose that $X$, $Y$ and $X'$ are flat and locally of finite presentation on $S$.
For every $s\in S$, denote $X_s$, $Y_s$, $X'_s$, $Y'_s$ the fibres of $X$, $Y$, $X'$, $Y'$ at the point $s$, $f_s:X'_s\to X_s$ the morphism deduced from $f$ by change of base.
It follows from \sref{IV.17.5.1} that for $f$ to be transversal to $Y$ at $x'$, it is necessary and sufficient that, if $s$ is the image of $x$ in $S$, $f_s$ be transversal to $Y_s$ at $x'$, relative to $k(s)$.\footnote{[Trans.] The original reference (17.5.9) is an erratum in the original; there is no item~17.5.9. The correct reference is {\hyperlink{ega4.c04.thm.17-5-1}{(17.5.1)}}.}

(ii)
We saw in \sref{IV.17.13.2} that the set of points $x'\in Y'$ at which $f$ is transversal to $Y$ is open in $Y'$.
If $Y'$ is moreover \emph{proper} over $S$, it follows that the set of $s\in S$ such that $f$ is transversal to $Y$ at all points of $Y'_s$, relative to $S$, is open in $S$.

When $X$, $Y$, and $X'$ are flat and locally of finite presentation over $S$, it follows from (i) that the set of $x'\in Y'$ such that $f_s$ is transversal to $Y_s$ at $x'$ (where $s$ is the image of $x'$ in $S$), relative to $k(s)$, is open in $Y'$.
If, in addition, $Y'$ is proper over $S$, the set of $s\in S$ such that $f_s$ is transversal to $Y_s$ at all points of $Y'_s$ (relative to $k(s)$) is open in $S$.

(iii)
The property of being smooth over $S$ at a point, as well as the notion of relative dimension over $S$ at a point, is stable under every base change $S'\to S$ (\sref{IV.17.3.3} and \sref{IV.4.2.7}); the same is therefore true of the property that a morphism $f:X'\to X$ be transversal to a subprescheme of $X$ at a point, relative to $S$.

\oldpage[IV]{92}

(iv)
Condition $b)$ of \sref{IV.17.13.2} also says that the homomorphism
$\alpha:g^*(\sh{N}_{Y/X})\to\Omega^1_{X'/S}\otimes_{\sh{O}_{X'}}\sh{O}_{Y'}$
is \emph{universally injective} relative to $Y'$ \sref{IV.11.9.18}.
\end{remarks}

\begin{env}[17.13.5]
\label{IV.17.13.5}
Consider a prescheme $S$, three S-preschemes $X$, $Y$, and $Z$, and two S-morphisms $f:Y\to X$, $g:Z\to X$; put $T=Y\times_XZ$.
By (\textbf{I}, 5.3.5), there is a commutative diagram
\[
\label{IV.17.13.5.1}
\tag{17.13.5.1}
\xymatrix{
  X \ar[d]^{\Delta} & T=Y\times_X Z \ar[l] \ar[d]^{j} \\
  X\times_S X & Y\times_S Z \ar[l]^{u}
}
\]
which makes $T$ the product of the $(X\times_SX)$-preschemes $X$ and $Y\times_SZ$, where $u=f\times_Sg$.
Since $\Delta$ is an S-immersion (\textbf{I}, 5.3.9), we are in the situation of diagram \sref{IV.17.13.1.1}.
The object corresponding there to $\Omega^1_{X'/S}$ is, by \sref{IV.16.4.23},
\[
  (\Omega^1_{Y/S}\otimes_{\sh{O}_Y}\sh{O}_{Y\times_SZ})
  \oplus
  (\Omega^1_{Z/S}\otimes_{\sh{O}_Z}\sh{O}_{Y\times_SZ}).
\]
On the other hand, the object corresponding to the $\sh{O}_Y$-module $\sh{N}_{Y/X}$ in \sref{IV.17.13.1} is here, by definition, $\Omega^1_{X/S}$ \sref{IV.16.3.1}.
Thus \sref{IV.17.13.1.3} becomes the exact sequence
\[
\label{IV.17.13.5.2}
\tag{17.13.5.2}
\Omega^1_{X/S}\otimes_{\sh{O}_X}\sh{O}_T \xrightarrow{\rho}
(\Omega^1_{Y/S}\otimes_{\sh{O}_Y}\sh{O}_T)\oplus
(\Omega^1_{Z/S}\otimes_{\sh{O}_Z}\sh{O}_T)
\xrightarrow{\sigma} \Omega^1_{T/S} \to 0
\]
where it remains to describe $\rho$ and $\sigma$ explicitly.
First, taking account of \sref{IV.16.4.23} and ($\mathbf{0}$, 20.5.2), if $p:T\to Y$ and $q:T\to Z$ are the canonical projections, then, with the notation of \sref{IV.16.4.19},
\[
\label{IV.17.13.5.3}
\tag{17.13.5.3}
\sigma = p_{T/Y/S} + q_{T/Z/S}.
\]
To evaluate $\rho$, use the commutativity of the left square in \sref{IV.17.13.1.2}.
In the present case, this first requires the canonical homomorphism
$\rho': \Omega^1_{X/S} \to \Omega^1_{(X\times_S X)/S}\otimes_{\sh{O}_{X\times_S X}}\sh{O}_X$
defined in \sref{IV.16.4.21} applied to the immersion $\Delta$.
We may reduce to the affine case $S=\Spec(A)$, $X=\Spec(B)$.
Then $\rho'$ corresponds to the homomorphism $\delta$ of ($\mathbf{0}$, 20.5.11.2), with $B$ there replaced by $B\otimes_AB$ and $C$ by
$B=(B\otimes_AB)/\mathfrak J_{B/A}$.
It sends the class of
$x\otimes1-1\otimes x\pmod{\mathfrak J_{B/A}^2}$, for $x\in B$, to the image of
\[
  (x\otimes1-1\otimes x)\otimes(1\otimes1)
  -(1\otimes1)\otimes(x\otimes1-1\otimes x)
\]
in $\Omega^1_{(B\otimes_AB)/A}\otimes_{B\otimes_AB}B$.
The preceding element may be written
\[
\begin{split}
 &((x\otimes1)\otimes(1\otimes1)-(1\otimes1)\otimes(x\otimes1))\\
 &\quad-((1\otimes x)\otimes(1\otimes1)-(1\otimes1)\otimes(1\otimes x)).
\end{split}
\]
Thus $\rho'$ is the difference of the two homomorphisms $\pi^{(1)}$ and $\pi^{(2)}$ from $\Omega^1_{X/S}$ to
$\Omega^1_{(X\times_SX)/S}\otimes_{\sh{O}_{X\times_SX}}\sh{O}_X$
corresponding, by \sref{IV.16.4.3.3}, to the first and second projections $X\times_SX\to X$.
To obtain $\rho$, consider the homomorphism
$\rho'':\Omega^1_{(X\times_SX)/S}\otimes_{\sh{O}_{X\times_SX}}\sh{O}_{Y\times_SZ}\to\Omega^1_{(Y\times_SZ)/S}$
corresponding by \sref{IV.16.4.3.3} to the morphism $u$ of \sref{IV.17.13.5.1}; after tensoring with $\sh{O}_T$, form the composite
$(\rho''\otimes1)\circ(\rho'\otimes1)$.
It follows from the preceding computation that, with the notation of \sref{IV.16.4.18},
\[
\label{IV.17.13.5.4}
\tag{17.13.5.4}
\rho = (f_{Y/X/S}\otimes 1_{\sh{O}_T})\oplus(-g_{Z/X/S}\otimes 1_{\sh{O}_T}).
\]

\oldpage[IV]{93}
Applying \sref{IV.17.13.2} to diagram \sref{IV.17.13.5.1} (taking account of \sref{IV.17.3.3},~(iv), which shows that $X\times_SX$ is smooth over $S$ at $(x,x)$ when $X$ is smooth over $S$ at $x$) gives the following corollary.
\end{env}

\begin{corollary}[17.13.6]
\label{IV.17.13.6}
Let $S$ be a prescheme, $X$, $Y$, $Z$ three $S$-preschemes, $f:Y\to X$, $g:Z\to X$ two $S$-morphisms; set $T=Y\times_X Z$, and let $p:T\to Y$, $q:T\to Z$ be the canonical projections.
Let $t$ be a point of $T$, $y=p(t)$, $z=q(t)$, $x=f(y)=g(z)$.
Suppose that $X$ is smooth over $S$ at the point $x$, of relative dimension $m$, $Y$ (resp.\ $Z$) smooth over $S$ at the point $y$ (resp.\ $z$), of relative dimension $m+a$ (resp.\ $m+b$), $a$ and $b$ being positive or negative integers.
Then the following conditions are equivalent:
\begin{enumerate}[label=\alph*)]
  \item $T$ is smooth over $S$ at the point $t$, of relative dimension $m+a+b$.
  \item The homomorphism
  \[
    \rho\otimes 1: \Omega^1_{X/S}\otimes_{\sh{O}_X}k(t) \to (\Omega^1_{Y/S}\otimes_{\sh{O}_Y}k(t))\oplus(\Omega^1_{Z/S}\otimes_{\sh{O}_Z}k(t))
  \]
  where $\rho$ is given by \sref{IV.17.13.5.4}, is injective.
  \item The morphism $u:Y\times_S Z\to X\times_S X$ is transversal to the diagonal of $X\times_S X$ at the point $t$.
\end{enumerate}
If $t$ is rational over the residue field of its image $s$ in $S$, these conditions are also equivalent to the following:
\begin{enumerate}
  \item[b${}'$)] The homomorphism
  \[
  \label{IV.17.13.6.1}
  \tag{17.13.6.1}
    \mathrm{T}_y(f)-\mathrm{T}_z(g): \mathrm{T}_{Y/S}(y)\oplus\mathrm{T}_{Z/S}(z) \to \mathrm{T}_{X/S}(x)
  \]
  (cf.\ \sref{IV.16.5.13.5}) is surjective.\footnote{The printed text cites (16.5.12.5); the tangent map used here is defined in {\hyperlink{ega4.c04.eq.16-5-13-5}{(16.5.13.5)}}.}
\end{enumerate}
Moreover, when conditions \emph{a)} and \emph{b)} are verified at $t$, they are so in a neighbourhood of $t$ in $T$, and restricting $T$ to such a neighbourhood, the sequence
\[
\label{IV.17.13.6.2}
\tag{17.13.6.2}
  0 \to \Omega^1_{X/S}\otimes_{\sh{O}_X}\sh{O}_T \xrightarrow{\rho} (\Omega^1_{Y/S}\otimes_{\sh{O}_Y}\sh{O}_T)\oplus(\Omega^1_{Z/S}\otimes_{\sh{O}_Z}\sh{O}_T) \xrightarrow{\sigma} \Omega^1_{T/S} \to 0
\]
(where $\sigma$ is given by \sref{IV.17.13.5.3}) is exact.

The only point in \sref{IV.17.13.6} still requiring proof is $b')$; it follows from the fact that, when $t$ is rational over $k(s)$, the homomorphism $\rho\otimes1$ is the \emph{transpose} of $\mathrm{T}_y(f)-\mathrm{T}_z(g)$, by \sref{IV.17.13.5.4}.

When the equivalent conditions of \sref{IV.17.13.6} hold, we say that $f$ and $g$ form a \emph{pair of transversal morphisms at $t$, relative to $S$}; here again, the mention of $S$ is often omitted.
\end{corollary}

\begin{env}[17.13.7]
\label{IV.17.13.7}
Consider in particular the case where $Y$ and $Z$ are \emph{sub-preschemes} of $X$, $f$ and $g$ being the canonical injections; $T$ is then the sub-prescheme ``intersection''

\oldpage[IV]{94}

$\inf(Y,Z)$ of $X$ (\textbf{I}, 4.4.3), and one has $t=y=z=x$.
Instead of saying that the couple $(f,g)$ is transversal at the point $x$, one says then that $Y$ and $Z$ \emph{cut transversally at the point $x$} (relatively to $S$).
Let $X_s$, $Y_s$, $Z_s$, and $T_s$ denote the fibres of $X$, $Y$, $Z$, and $T$ at $s$, the image of $x$ in $S$.
Taking account of \sref{IV.5.2.3}, \sref{IV.5.1.9}, and ($\mathbf{0}$, 16.5.12), we see that, when $X$, $Y$, and $Z$ are smooth over $S$ at $x$, they cut transversally at $x$ if and only if $T$ is smooth over $S$ at $x$ and
\[
\label{IV.17.13.7.1}
\tag{17.13.7.1}
\codim_x(T_s,X_s) = \codim_x(Y_s,X_s) + \codim_x(Z_s,X_s).
\]
\end{env}

\begin{proposition}[17.13.8]
\label{IV.17.13.8}
Let $S$ be a prescheme, $X$ an $S$-prescheme locally of finite presentation over $S$, $Y$, $Z$ two sub-preschemes of $X$, $T=\inf(Y,Z)$ the sub-prescheme ``intersection'' of $Y$ and $Z$, $x$ a point of $T$.
The following conditions are equivalent:
\begin{enumerate}
  \item[a)] The canonical injection $i:Y\to X$ is a transversal morphism to $Z$ at the point $x$, relatively to $S$ \sref{IV.17.13.3}.
  \item[a${}'$)] The canonical injection $j:Z\to X$ is a transversal morphism to $Y$ at the point $x$, relatively to $S$.
  \item[a${}''$)] $Y$ and $Z$ cut transversally at the point $x$, relatively to $S$.
  \item[b)] $X$, $Y$, and $Z$ are smooth over $S$ at $x$, and the homomorphism
  \[
    \rho\otimes 1: \Omega^1_{X/S}\otimes_{\sh{O}_X}k(x) \to (\Omega^1_{Y/S}\otimes_{\sh{O}_Y}k(x))\oplus(\Omega^1_{Z/S}\otimes_{\sh{O}_Z}k(x))
  \]
  is injective.
\end{enumerate}
When $x$ is rational over the residue field $k(s)$ ($s$ image of $x$ in $S$), these conditions are also equivalent to the following:
\begin{enumerate}
  \item[b${}'$)] The homomorphism
  \[
    \mathrm{T}_x(i)-\mathrm{T}_x(j): \mathrm{T}_{Y/S}(x)\oplus\mathrm{T}_{Z/S}(x) \to \mathrm{T}_{X/S}(x)
  \]
  is surjective.
\end{enumerate}
Moreover, when conditions \emph{a)} and \emph{b)} are verified at $x$, they are so in a neighbourhood of $x$ in $T$, and restricting $X$ to a neighbourhood of $x$, the sequence \sref{IV.17.13.6.2} is exact, and one has a canonical isomorphism
\[
\label{IV.17.13.8.1}
\tag{17.13.8.1}
\sh{N}_{T/X} \xrightarrow{\sim} (\sh{N}_{Y/X}\otimes_{\sh{O}_Y}\sh{O}_T)\oplus(\sh{N}_{Z/X}\otimes_{\sh{O}_Z}\sh{O}_T).
\]
\end{proposition}

\begin{proof}
Conditions $a)$, $a')$, and $a'')$ all imply that $X$, $Y$, and $Z$ are smooth over $S$ at $x$.
Let $m$, $m-a$, and $m-b$ be their respective relative dimensions over $S$ at $x$.
Applying \sref{IV.17.13.2} with $Z$ in place of $X'$ and $T=Y\times_XZ$ in place of $Y'$, conditions $a)$ and $a')$ are both equivalent to saying that $T$ is smooth over $S$ at $x$ and has relative dimension $m-a-b$ there.
By \sref{IV.17.13.7}, this is precisely condition $a'')$; hence $a)$, $a')$, and $a'')$ are equivalent.
The equivalence of $a'')$ and $b)$ (and of $b')$ when $x$ is rational over $k(s)$), the persistence of these conditions in a neighbourhood of $x$, and the exactness of \sref{IV.17.13.6.2} there were proved in \sref{IV.17.13.6}.
It remains to define the canonical isomorphism \sref{IV.17.13.8.1}.

\oldpage[IV]{95}

Denote by $\alpha$ and $-\beta$ the two homomorphisms in the middle term of \sref{IV.17.13.5.4}, and by $\gamma$ and $\delta$ the two homomorphisms in the middle term of \sref{IV.17.13.5.2}.
Exactness of \sref{IV.17.13.6.2}, that is, of
\[
0\to\Omega^1_{X/S}\otimes_{\sh{O}_X}\sh{O}_T
\xrightarrow{\alpha\oplus(-\beta)}
(\Omega^1_{Y/S}\otimes_{\sh{O}_Y}\sh{O}_T)\oplus
(\Omega^1_{Z/S}\otimes_{\sh{O}_Z}\sh{O}_T)
\xrightarrow{\gamma+\delta}\Omega^1_{T/S}\to0,
\]
means that, in the category of $\sh{O}_T$-modules,
$\Omega^1_{X/S}\otimes_{\sh{O}_X}\sh{O}_T$ is canonically the fibre product of
$\Omega^1_{Y/S}\otimes_{\sh{O}_Y}\sh{O}_T$ and
$\Omega^1_{Z/S}\otimes_{\sh{O}_Z}\sh{O}_T$ over $\Omega^1_{T/S}$, for $\gamma$ and $\delta$.
The same reasoning as in ($\mathbf{0}$, 18.1.2 and 18.1.3), replacing rings and ideals by modules and submodules, gives the commutative diagram
{\scriptsize
\[
\begin{tikzcd}[column sep=small,row sep=large]
0 \ar[dr] & & 0 \ar[d] & 0 \ar[d] & \\
 & (\sh{N}_{Y/X}\otimes\sh{O}_T)\oplus(\sh{N}_{Z/X}\otimes\sh{O}_T)
   \ar[r] \ar[d] \ar[dr,"\iota" description]
 & \sh{N}_{Z/X}\otimes\sh{O}_T \ar[r,"\sim"] \ar[d]
 & \sh{N}_{T/Y} \ar[d] & \\
0 \ar[r] & \sh{N}_{Y/X}\otimes\sh{O}_T \ar[r] \ar[d]
 & \Omega^1_{X/S}\otimes\sh{O}_T \ar[r,"\alpha"] \ar[d,"\beta"] \ar[dr,"\varepsilon" description]
 & \Omega^1_{Y/S}\otimes\sh{O}_T \ar[r] \ar[d,"\gamma"] & 0 \\
0 \ar[r] & \sh{N}_{T/Z} \ar[r]
 & \Omega^1_{Z/S}\otimes\sh{O}_T \ar[r,"\delta"] \ar[d]
 & \Omega^1_{T/S} \ar[r] \ar[d] \ar[dr] & 0 \\
 & & 0 & 0 & 0
\end{tikzcd}
\]
}
whose third and fourth rows and columns are exact by the smoothness hypotheses.
The identity $\varepsilon=\gamma\circ\alpha=\delta\circ\beta$, and the fact that this is the homomorphism corresponding to the canonical injection $T\to X$, follow from ($\mathbf{0}$, 20.5.2.7).
Since the diagonal of the preceding diagram is exact, $\ker(\varepsilon)$ identifies canonically with $\sh{N}_{T/X}$, taking account of the fact that $T$ is smooth over $S$ \sref{IV.17.2.5}.
Thus the canonical isomorphism \sref{IV.17.13.8.1} is the inverse of $\iota$.
\end{proof}

\begin{env}[17.13.9]
\label{IV.17.13.9}
The results of \sref{IV.17.13.5} and \sref{IV.17.13.6} may be generalized to any finite family of S-morphisms
$f_i:Y_i\to X$ $(i\in I)$.
Let $X^I$ be the product of the family $(X_i)_{i\in I}$ of S-preschemes all identical to $X$ (\textbf{I}, 3.3.5), and let $p_i:X^I\to X$ be the canonical projections.
The diagonal morphism $\Delta:X\to X^I$ is defined, as in (\textbf{I}, 5.3.1), by the identities $p_i\circ\Delta=1_X$ for every $i\in I$; it is an X-section of $X^I$ for $p_1$, and hence an immersion (\textbf{I}, 5.3.11).

\oldpage[IV]{96}

Let $Y$ be the product of the S-preschemes $Y_i$ for the composites $Y_i\xrightarrow{f_i}X\to S$, and let $T$ be the product of the X-preschemes $Y_i$ for the morphisms $f_i$.
As in (\textbf{I}, 5.3.5), there is a commutative diagram
\[
\label{IV.17.13.9.1}
\tag{17.13.9.1}
\xymatrix{
  X \ar[d]^{\Delta} & T \ar[l]_-v \ar[d]^{j} \\
  X^I & Y \ar[l]^{u}
}
\]
(where $u$ is the product over $S$ of the $f_i$), which makes $T$ the product of the $X^I$-preschemes $X$ and $Y$.
By induction on the number of elements of $I$, \sref{IV.16.4.23} shows that $\Omega^1_{X^I/S}$ identifies canonically with the direct sum of the $p_i^*(\Omega^1_{X/S})$; similarly, $\Omega^1_{Y/S}$ identifies with the direct sum of the $q_i^*(\Omega^1_{Y_i/S})$, where $q_i:Y\to Y_i$ are the canonical projections.

Suppose now that $X$ is smooth over $S$ at a point $x$, and hence in a neighbourhood of $x$.
Then $X^I$ is smooth at the corresponding point \sref{IV.17.3.3},~(iv), and, after restricting $X$ to a neighbourhood of $x$, \sref{IV.17.2.5} gives an exact sequence of locally free $\sh{O}_X$-modules
\[
  0\to\operatorname{gr}_1(\Delta)\to(\Omega^1_{X/S})^I\to\Omega^1_{X/S}\to0.
\]
Put $\sh{J}=\operatorname{gr}_1(\Delta)$.
The preceding sequence gives an exact sequence of locally free $\sh{O}_T$-modules
\[
  0 \to \sh{J}\otimes_{\sh{O}_X}\sh{O}_T \to (\Omega^1_{X/S})^I\otimes_{\sh{O}_X}\sh{O}_T \xrightarrow{\sigma} \Omega^1_{X/S}\otimes_{\sh{O}_X}\sh{O}_T \to 0
\]
which is the first row of diagram \sref{IV.17.13.1.2}, and where $\sigma$ is the canonical homomorphism sending each family of sections $(t'_i)_{i\in I}$ of $\Omega^1_{X/S}\otimes_{\sh{O}_X}\sh{O}_T$ over an open subset of $T$ to its \emph{sum}.
On the other hand, the homomorphism corresponding here to the second vertical arrow of diagram \sref{IV.17.13.1.2} is
\[
\label{IV.17.13.9.2}
\tag{17.13.9.2}
\tau: (\Omega^1_{X/S})^I\otimes_{\sh{O}_X}\sh{O}_T \to \bigoplus_{i\in I}(\Omega^1_{Y_i/S}\otimes_{\sh{O}_{Y_i}}\sh{O}_T)
\]
which sends a family $(t'_i\otimes1)_{i\in I}$, where the $t'_i$ are sections of $\Omega^1_{X/S}$, to the family of the
$((f_i)_{Y_i/X/S}(t'_i))\otimes1$, with the notation of \sref{IV.16.4.18}.
The homomorphism $\rho$ corresponding to $\alpha$ of \sref{IV.17.13.1.3} is therefore the restriction of $\tau$ to $\sh{J}\otimes_{\sh{O}_X}\sh{O}_T$.
With these hypotheses and notation, Proposition \sref{IV.17.13.2} applied to \sref{IV.17.13.9.1} gives the following corollary.
\end{env}

\begin{corollary}[17.13.10]
\label{IV.17.13.10}
Under the general hypotheses of \sref{IV.17.13.9}, let $t$ be a point of $T$, $y_i$ ($i\in I$) its projections in the $Y_i$, $x$ the point of $X$ equal to each of the $f_i(y_i)$.
Suppose that $X$ is smooth over $S$ at the point $x$ and of relative dimension $d$, and that each of the $Y_i$ is smooth over $S$ at the point $y_i$; let $d+c_i$ ($c_i$ a positive or negative integer) be the relative dimension of $Y_i$ at the point $y_i$.
Then the following conditions are equivalent:

\oldpage[IV]{97}

\begin{enumerate}[label=\alph*)]
  \item $T$ is smooth over $S$ at the point $t$, of relative dimension $d+\sum_{i\in I} c_i$.
  \item The homomorphism
  \[
    \rho\otimes 1: \sh{J}\otimes_{\sh{O}_X}k(t) \to \bigoplus_{i\in I}(\Omega^1_{Y_i/S}\otimes_{\sh{O}_{Y_i}}k(t))
  \]
  is injective.
\end{enumerate}
When $t$ is rational over the residue field $k(s)$ ($s$ the image of $t$ in $S$), these conditions are also equivalent to the following:
\begin{enumerate}
  \item[b${}'$)] The homomorphism transpose of $\rho\otimes 1$
  \[
    \bigoplus_{i\in I}\mathrm{T}_{Y_i/S}(y_i) \to (\mathrm{T}_{X/S}(x))^I/\delta(\mathrm{T}_{X/S}(x))
  \]
  (where $\delta$ is the diagonal application) is surjective.
\end{enumerate}
It suffices to note that $X^I$ is smooth over $S$ at the point $x$ and of relative dimension $d\cdot\operatorname{Card}(I)$, and $Y$ smooth over $S$ at the point $t$ and of relative dimension $d\cdot\operatorname{Card}(I)+\sum_{i\in I} c_i$.

When the equivalent conditions of \sref{IV.17.13.10} are verified, one says further that the $f_i$ form a \emph{family of transversal morphisms at the point $t$, relatively to} $S$.
One sees further that the set of points $t\in T$ where this holds is \emph{open} in $T$.
The remark (\sref{IV.17.13.4}, (ii)) shows moreover that if $X$ and the $Y_i$ are flat and locally of finite presentation over $S$, and if moreover $T$ is \emph{proper} over $S$, the set of $s\in S$ such that the $f_i$ form a family of transversal morphisms at \emph{all} the points of $T_s$, relatively to $S$, is \emph{open} in $S$.
\end{corollary}

\begin{env}[17.13.11]
\label{IV.17.13.11}
Consider in particular the case, generalising \sref{IV.17.13.7}, where the $Y_i$ ($i\in I$) are \emph{sub-preschemes} of $X$, the $f_i$ being the canonical injections, so that $T=\inf_{i\in I}(Y_i)$ is also the sub-prescheme ``intersection'' of the $Y_i$, and $t=y_i=x$ for every $i$; instead of saying that the $f_i$ form a family of transversal morphisms at the point $x$, one says further that the $Y_i$ \emph{cut transversally at the point} $x$ (relatively to $S$).
The condition $a)$ of \sref{IV.17.13.10} is expressed further by the relation that generalises \sref{IV.17.13.7.1}:
\[
\label{IV.17.13.11.1}
\tag{17.13.11.1}
\codim_x(T_s,X_s) = \sum_{i\in I}\codim_x((Y_i)_s,X_s).
\]

Moreover, one has the following property which extends \sref{IV.17.13.8.1}, and of which we give another demonstration when $I$ has $2$ elements:
\end{env}

\begin{corollary}[17.13.12]
\label{IV.17.13.12}
When the $Y_i$ cut transversally at the point $x$, one has a canonical isomorphism
\[
\label{IV.17.13.12.1}
\tag{17.13.12.1}
\sh{N}_{T/X} \xrightarrow{\sim} \bigoplus_{i\in I}(\sh{N}_{Y_i/X}\otimes_{\sh{O}_{Y_i}}\sh{O}_T).
\]
\end{corollary}

\begin{proof}
One can reduce to the case where the $Y_i$ are closed sub-preschemes of $X$, defined by quasi-coherent Ideals $\sh{J}_i$, so that $T$ is defined by the Ideal $\sh{J}=\sum_i\sh{J}_i$.
By definition of the conormal sheaf of an immersion \sref{IV.16.1.3}, the canonical homomorphism $\bigoplus_i\sh{J}_i\to\sh{J}$ gives, by passage to quotients, a \emph{surjective} homomorphism
\[
\label{IV.17.13.12.2}
\tag{17.13.12.2}
\bigoplus_{i\in I}(\sh{N}_{Y_i/X}\otimes_{\sh{O}_{Y_i}}\sh{O}_T) \to \sh{N}_{T/X}.
\]
\oldpage[IV]{98}
Here the $\sh{O}_T$-modules on both sides of \sref{IV.17.13.12.2} are locally free and have the same rank $\sum_i c_i$, where $c_i$ is the rank of $\sh{N}_{Y_i/X}$, by \sref{IV.17.2.5} and condition $a)$ of \sref{IV.17.13.10}.
Bourbaki, \emph{Alg. comm.}, chap.~II, \S~3, n\textsuperscript{o}~2, cor. to prop.~6, therefore shows that \sref{IV.17.13.12.2} is bijective; \sref{IV.17.13.12.1} is its inverse.
\end{proof}

\subsection{Local and infinitesimal characterisations of smooth, unramified, and \'etale morphisms.}
\label{subsection:IV.17.14}
\begin{proposition}[17.14.1]
\label{IV.17.14.1}
Let $f:X\to Y$ be a morphism locally of finite presentation, $x$ a point of $X$, $y=f(x)$.
For $f$ to be smooth (resp.\ unramified, resp.\ \'etale) at $x$, it is necessary and sufficient that the following condition hold:

For every local scheme $Y'=\operatorname{Spec}(A')$ with closed point $y'$, every morphism $h:Y'\to Y$ such that $h(y')=y$, every closed subscheme $Y'_0=\operatorname{Spec}(A'/\mathfrak J')$ of $Y'$, where $(\mathfrak J')^2=0$, and every Y-morphism $g_0:Y'_0\to X$ such that $g_0(y')=x$, there exists \emph{at least one} (resp.\ \emph{at most one}, resp.\ \emph{one and only one}) Y-morphism $g:Y'\to X$ whose restriction to $Y'_0$ is $g_0$.
\end{proposition}

\begin{proof}
In view of Definitions \sref{IV.17.1.1} and \sref{IV.17.3.1}, it remains only to show that the stated condition is sufficient in the smooth and unramified cases.

\emph{(i) The smooth case.}
We may restrict ourselves to the affine case $Y=\operatorname{Spec}(A)$ and $X=\operatorname{Spec}(C)$, where
$C=B/\mathfrak R$, $B=A[\mathrm{T}_1,\ldots,\mathrm{T}_n]$, and $\mathfrak R$ is a finite-type ideal of $B$.
To prove that $f$ is smooth at $x$, it is enough to prove that $\sh{O}_{X,x}$ is formally smooth over $\sh{O}_{Y,y}$ for the discrete topologies \sref{IV.17.5.1}.
Now $\sh{O}_{X,x}=B_{\mathfrak q}/\mathfrak RB_{\mathfrak q}$, where $\mathfrak q$ is a prime ideal of $B$; if $\mathfrak p$ is its inverse image in $A$, then $B_{\mathfrak q}$ is formally smooth over $A_{\mathfrak p}$ for the discrete topologies ($\mathbf{0}$, 19.3.2 and 19.3.5).
By the Jacobian criterion ($\mathbf{0}$, 22.6.1 and 20.5.12), it is therefore enough to show that
$B_{\mathfrak q}/\mathfrak R^2B_{\mathfrak q}$ is an $A_{\mathfrak p}$-trivial extension of
$B_{\mathfrak q}/\mathfrak RB_{\mathfrak q}$ by
$\mathfrak RB_{\mathfrak q}/\mathfrak R^2B_{\mathfrak q}$.
But this is precisely the hypothesis applied to
\[
 A'=B_{\mathfrak q}/\mathfrak R^2B_{\mathfrak q},\qquad
 \mathfrak J'=\mathfrak RB_{\mathfrak q}/\mathfrak R^2B_{\mathfrak q},
\]
and to the morphism $g_0$ corresponding to the identity automorphism of
$A'/\mathfrak J'=\sh{O}_{X,x}$.

\emph{(ii) The unramified case.}
Put $A=\sh{O}_{Y,y}$ and $B=\sh{O}_{X,x}$.
By \sref{IV.17.4.1},~$c')$, it is enough to prove that $\Omega^1_{B/A}=0$.
With the notation of ($\mathbf{0}$, 20.4.1), the ring
$C=\sh{P}^1_{B/A}=(B\otimes_AB)/\mathfrak J^2$ is local, since
$C/(\mathfrak J/\mathfrak J^2)$ is isomorphic to $B$.
Applying the hypothesis to $A'=C$ and $\mathfrak J'=\mathfrak J/\mathfrak J^2$ shows, by definition, that the map
$d_{B/A}:B\to\Omega^1_{B/A}$ is zero ($\mathbf{0}$, 20.4.6); hence $\Omega^1_{B/A}=0$ by ($\mathbf{0}$, 20.4.7).
\end{proof}

\begin{proposition}[17.14.2]
\label{IV.17.14.2}
Let $Y$ be a prescheme locally Noetherian, $f:X\to Y$ a morphism locally of finite type, $x$ a point of $X$, $y=f(x)$.
For $f$ to be smooth (resp.\ unramified, resp.\ \'etale) at $x$, it is necessary and sufficient that the condition of \sref{IV.17.14.1} hold when one additionally requires that $A'$ be Artinian, that $\mathfrak m'\mathfrak J'=0$, where $\mathfrak m'$ is the maximal ideal of $A'$, and that the residue field $A'/\mathfrak m'$ be equal to $k(x)$.
\end{proposition}

\oldpage[IV]{99}

\begin{proof}
It remains again only to prove sufficiency in the smooth and unramified cases.

\emph{(i) The smooth case.}
Put $A=\sh{O}_{Y,y}$ and $B=\sh{O}_{X,x}$.  In view of \sref{IV.17.5.3}, the point is to prove that $B$ is a formally smooth $A$-algebra for the preadic topologies.  By ($\mathbf{0}$, 22.1.4), this follows from the hypothesis if, in it, the condition $\mathfrak m'\mathfrak J'=0$ is replaced by $\mathfrak J'\subset\mathfrak m'$.  Since $\mathfrak m'$ is nilpotent, however, the condition in the statement is already sufficient: consider the rings $A'/\mathfrak m'^j\mathfrak J'$ and argue by induction on $j$, as in ($\mathbf{0}$, 19.4.3).

\emph{(ii) The unramified case.}
Resume the notation of the proof of \sref{IV.17.14.1},~(ii).  To prove that the ideal $\mathfrak J/\mathfrak J^2$ of $C$ is zero, note first that $C$ is a Noetherian local ring: it is the local ring of an affine open subscheme of $X\times_YX$, which is locally of finite type over $X$, and therefore also over $Y$.  If $\mathfrak r$ is the maximal ideal of $C$, it is consequently enough to verify that
\[
 \mathfrak J/\mathfrak J^2\subset\mathfrak r^n
 \quad\text{for every }n
 \qquad(\mathbf{0}_{\mathrm I}, 7.3.5),
\]
or, equivalently, that $\mathfrak r^n=\mathfrak r^n+(\mathfrak J/\mathfrak J^2)$.  With the notation of ($\mathbf{0}$, 20.4.6), this says that for every $n$ the two composite maps
\[
 B\xrightarrow{p_1}C\longrightarrow C/\mathfrak r^n,
 \qquad
 B\xrightarrow{p_2}C\longrightarrow C/\mathfrak r^n
\]
are identical.  But $C/\mathfrak r^n$ is Artinian, and this identity follows by induction on $n$ from the hypothesis applied to
\[
 A'=C/\mathfrak r^n,
 \qquad
 \mathfrak J'=\bigl(\mathfrak r^n+(\mathfrak J/\mathfrak J^2)\bigr)/\mathfrak r^n.
\]
\end{proof}

\begin{remark}[17.14.3]
\label{IV.17.14.3}
Under the hypotheses of \sref{IV.17.14.2}, if in addition the residue field $k(x)$ is a \emph{finite} extension of $k(y)$, then the $A'$-modules $\mathfrak m'^j/\mathfrak m'^{j+1}$ are finite-dimensional $k(y)$-vector spaces; hence, \emph{a fortiori}, $A'$ is a finite $\sh{O}_{Y,y}$-algebra.
\end{remark}


\subsection{The case of preschemes over a base field.}
\label{subsection:IV.17.15}

We first recall \sref{IV.6.7.7}, \sref{IV.6.7.8}, and \sref{IV.6.8.1}:

\begin{proposition}[17.15.1]
\label{IV.17.15.1}
Let $k$ be a field, $X$ a prescheme locally of finite type over $k$, and $x$ a point of $X$.
For $X$ to be smooth at $x$, it is necessary and sufficient that, for every radicial extension $k'$ of $k$ (or only for every finite radicial extension $k'$ of $k$), the local ring $\sh{O}_{X,x}\otimes_k k'$ be regular.
If $k(x)$ is a separable extension of $k$ (in particular if $k$ is perfect), this amounts to saying that $\sh{O}_{X,x}$ is regular.
\end{proposition}

\begin{corollary}[17.15.2]
\label{IV.17.15.2}
Under the conditions of \sref{IV.17.15.1}, for $X$ to be smooth over $k$, it is necessary and sufficient that $X$ be geometrically regular over $k$ \sref{IV.6.7.6}; this implies that $X$ is reduced.
If $k$ is perfect, $X$ is smooth over $k$ if and only if $X$ is regular.
\end{corollary}

\begin{proposition}[17.15.3]
\label{IV.17.15.3}
Let $k$ be a field, $X$ a prescheme locally of finite type over $k$, $x$ a point of $X$, and $n=\dim_x(X)$.
Let $(g_i)_{1\leq i\leq n}$ be a family of $n$ sections of $\sh{O}_X$ over $X$, and let $f:X\to Y=\operatorname{Spec}(k[\mathrm{T}_1,\ldots,\mathrm{T}_n])$ be the morphism corresponding to the $k$-homomorphism $k[\mathrm{T}_1,\ldots,\mathrm{T}_n]\to\Gamma(X,\sh{O}_X)$ sending $\mathrm{T}_i$ to $g_i$ \sref[I]{I.2.2.4}.
The following conditions are equivalent (and imply that $X$ is smooth over $k$ at $x$, and therefore regular at $x$):
\begin{enumerate}
  \item[a)] $f$ is \'etale at the point $x$.
  \item[b)] The images $d_{X/k}g_i$ form a basis of the $\sh{O}_{X,x}$-module $(\Omega^1_{X/k})_x$.
  \item[c)] The images $d_{X/k}g_i$ generate the $\sh{O}_{X,x}$-module $(\Omega^1_{X/k})_x$.
\end{enumerate}
If $f$ is \'etale at $x$, then $X$ is smooth over $k$ at $x$, since $Y=\operatorname{Spec}(k[\mathrm{T}_1,\ldots,\mathrm{T}_n])$ is smooth over $k$ \sref{IV.17.3.8}; the implication $a)\Rightarrow b)$ is therefore a special case of \sref{IV.17.11.3}.
As $b)$ trivially implies $c)$, it remains to see that $c)$ implies $a)$.  First note
\oldpage[IV]{100}
\end{proposition}

\begin{proof}
that hypothesis $c)$ makes the homomorphism $(f^*\Omega^1_{Y/k})_x\to(\Omega^1_{X/k})_x$ surjective.  After replacing $X$ by an open neighbourhood of $x$, we may therefore suppose that $f^*\Omega^1_{Y/k}\to\Omega^1_{X/k}$ is surjective (Bourbaki, \emph{Alg.~comm.}, chap.~II, \S~5, n\textsuperscript{o}~1, prop.~2); hence $f$ is unramified \sref{IV.17.2.2}.
Let us first show that one can reduce to the case where $x$ is rational over $k$.
Indeed, put $k'=k(x)$, $X'=X\otimes_k k'$, and $Y'=Y\otimes_k k'=\operatorname{Spec}(k'[\mathrm{T}_1,\ldots,\mathrm{T}_n])$.  There exists a point $x'\in X'$ above $x$ such that $k(x')=k'$.
To prove that $f$ is \'etale at $x$, it suffices to show that $f'=f_{(k')}:X'\to Y'$ is \'etale at $x'$ \sref{IV.17.7.1},~(ii); moreover, $f'$ is unramified \sref{IV.17.3.3},~(iii), and $\dim_{x'}(X')=n$ \sref{IV.4.2.7}.
In the same way, replacing $k'$ by an algebraically closed extension of $k'$, one can suppose that $k$ is algebraically closed.
Now put $y=f(x)$, $A=\sh{O}_{Y,y}$, and $B=\sh{O}_{X,x}$.  Since the residue field of $B$ is $k$, the same is true of the residue field of $A$; thus $x$ (resp.\ $y$) is a closed point of $X$ (resp.\ $Y$) \sref[I]{I.6.4.2}, and consequently $\dim A=\dim B=n$.
Since $f$ is unramified, the homomorphism $\widehat A\to\widehat B$ is surjective \sref{IV.17.4.4},~$f'')$.  But $\dim\widehat A=\dim\widehat B=n$, and $\widehat A$, being regular, is an integral domain; hence $\widehat A\to\widehat B$ is also injective ($\mathbf{0}$, 16.3.10).  It is therefore bijective, which completes the proof that $f$ is \'etale at $x$ \sref{IV.17.6.3},~$e'')$.
\end{proof}

\begin{corollary}[17.15.4]
\label{IV.17.15.4}
Under the general hypotheses of \sref{IV.17.15.3}, suppose moreover that $k(x)$ is a finite and separable extension of $k$ and that the germs $(g_i)_x$ belong to $\mathfrak{m}_x$.
Then the conditions $a)$, $b)$, $c)$ of \sref{IV.17.15.3} are also equivalent to each of the following:
\begin{enumerate}
  \item[d)] The $n$ germs $(g_i)_x$ generate the maximal ideal $\mathfrak m_x$ of $\sh{O}_{X,x}$.
  \item[d')] The ring $\sh{O}_{X,x}$ is regular and the $(g_i)_x$ form a regular system of parameters for this ring ($\mathbf{0}$, 17.1.6).
\end{enumerate}
\end{corollary}

\begin{proof}
Clearly $d')$ implies $d)$.
Since $k(x)$ is a finite separable extension of $k$, one has $\Omega^1_{k(x)/k}=0$ ($\mathbf{0}$, 20.6.20), and $k(x)=\sh{O}_{X,x}/\mathfrak m_x$ is a formally smooth $k$-algebra for the discrete topologies ($\mathbf{0}$, 19.6.1).  The exact sequence ($\mathbf{0}$, 20.5.14.1), applied with $A=k$, $B=\sh{O}_{X,x}$, and $\mathfrak R=\mathfrak m_x$, therefore gives the canonical isomorphism
\[
\delta:\mathfrak m_x/\mathfrak m_x^2\xrightarrow{\sim} (\Omega^1_{X/k})_x\otimes_{\sh{O}_{X,x}} k(x).
\]
It follows first, by Nakayama's lemma, that $c)$ and $d)$ are equivalent.
If $f$ is \'etale at $x$, the ring $\sh{O}_{X,x}$ is regular of dimension $n$, since $x$ is closed \sref[I]{I.6.4.2}; and $n$ elements which generate its maximal ideal necessarily form a regular system of parameters ($\mathbf{0}$, 17.1.6).  Thus $a)$ implies $d')$.
\end{proof}

\begin{proposition}[17.15.5]
\label{IV.17.15.5}
Let $k$ be a field, $X$ a prescheme locally of finite type over $k$, $x$ a point of $X$, and $n=\dim_x(X)$.
The following conditions are equivalent:
\begin{enumerate}
  \item[a)] $X$ is smooth over $k$ at $x$.
  \item[b)] $X$ is differentially smooth over $k$ at $x$.
  \item[c)] $(\Omega^1_{X/k})_x$ is a free $\sh{O}_{X,x}$-module of rank $n$.
  \item[d)] $(\Omega^1_{X/k})_x$ admits a system of $n$ generators as an $\sh{O}_{X,x}$-module.
\oldpage[IV]{101}
  \item[e)] There exist a perfect extension $k'$ of $k$ and an open neighbourhood $U$ of $x$ in $X$ such that the prescheme $U\otimes_k k'$ is regular.
\end{enumerate}
\end{proposition}

\begin{proof}
If $X$ is smooth over $k$ at $x$, there is an open neighbourhood $U$ of $x$ which is smooth over $k$; consequently $U\otimes_k k'$ is regular for every extension $k'$ of $k$ \sref{IV.17.15.2}.  This proves $a)\Rightarrow e)$.  Conversely, $e)$ implies $a)$, because $U\otimes_k k'$ is then smooth over the perfect field $k'$ \sref{IV.17.15.2}, and hence $U$ is smooth over $k$ \sref{IV.17.7.1},~(ii).

We have already proved $a)\Rightarrow c)$ \sref{IV.17.10.2}; clearly $c)\Rightarrow d)$, and $d)\Rightarrow a)$ follows from ($\mathbf{0}$, 20.4.7) and \sref{IV.17.15.3},~$c)$.

Finally, $a)\Rightarrow b)$ was proved in \sref{IV.17.12.4}.  Conversely, to prove $b)\Rightarrow a)$, we may reduce to the case where $x$ is rational over $k$: as in the proof of \sref{IV.17.15.3}, take a point $x'$ of $X'=X\otimes_k k'$, above $x$, with $k(x')=k'=k(x)$, using again \sref{IV.17.7.1},~(ii), and the fact that differential smoothness at $x$ implies differential smoothness of $X'$ at $x'$ \sref{IV.16.10.4}.  When $x$ is rational over $k$, the smoothness of the structural morphism $X\to\operatorname{Spec}(k)$ at $x$ follows from $b)$ and \sref{IV.17.12.5}, applied to the $k$-section $u:\operatorname{Spec}(k)\to X$ such that $u(\operatorname{Spec}(k))=\{x\}$.
\end{proof}

\begin{corollary}[17.15.6]
\label{IV.17.15.6}
Let $X$ be a prescheme of finite type over a field $k$.
For $X$ to be smooth over $k$, it is necessary and sufficient that the $\sh{O}_X$-module $\Omega^1_{X/k}$ be locally free and that the local rings of $X$ at its maximal points be fields which are separable extensions of $k$ (the latter condition is automatically satisfied if $k$ is perfect and $X$ is \emph{reduced}).
\end{corollary}

\begin{proof}
The conditions are necessary.  Indeed, if $X$ is smooth over $k$, then $\Omega^1_{X/k}$ is locally free by \sref{IV.17.10.2}.  Moreover, $X$ is regular and hence, \emph{a fortiori}, reduced; at every maximal point $x$ of $X$, therefore, $\sh{O}_{X,x}$ is a field.  This field must be a formally smooth $k$-algebra for the discrete topologies \sref{IV.17.5.1}, and is consequently a separable extension of $k$ ($\mathbf{0}$, 19.6.1).

The conditions are sufficient.
We may reduce to the case where $X$ is connected; then $\Omega^1_{X/k}$ is locally free of constant rank $n$.  At every maximal point $x$ of $X$ one has
$\operatorname{rank}_{k(x)}\Omega^1_{k(x)/k}=n$, since $\sh{O}_{X,x}=k(x)$.  By hypothesis $k(x)$ is separable over $k$, so $\Upsilon_{k(x)/k}=0$ ($\mathbf{0}$, 20.6.3), and Cartier's equality gives $\deg.\operatorname{tr}_k k(x)=n$ ($\mathbf{0}$, 21.7.1).  All irreducible components of $X$ therefore have the same dimension $n$ \sref{IV.5.2.1}, and the implication $c)\Rightarrow a)$ of \sref{IV.17.15.5} shows that $X$ is smooth over $k$ at every point.
\end{proof}

\begin{corollary}[17.15.7]
\label{IV.17.15.7}
If $k$ has characteristic $0$, then $X$ is smooth over $k$ at $x$ if and only if $(\Omega^1_{X/k})_x$ is a free $\sh{O}_{X,x}$-module.
\end{corollary}

\begin{proof}
This results from \sref{IV.16.12.2} and the equivalence of $b)$ and $a)$ in \sref{IV.17.15.5}.
\end{proof}

\begin{proposition}[17.15.8]
\label{IV.17.15.8}
Let $k$ be a field, $X$ a prescheme locally of finite type over $k$, and $x$ a point of $X$; put $n=\dim_x(X)$ and $r=\dim(\sh{O}_{X,x})$, so that $n-r=\deg.\operatorname{tr}_k k(x)$ \sref{IV.5.2.3}.
Let $(g_i)_{1\leq i\leq n}$ be a family of $n$ sections of $\sh{O}_X$ over $X$, with $(g_i)_x\in\mathfrak m_x$ for $1\leq i\leq r$; let $f:X\to Y=\operatorname{Spec}(k[\mathrm{T}_1,\ldots,\mathrm{T}_n])$ be the morphism corresponding to the $k$-homomorphism which sends $\mathrm{T}_i$ to $g_i$.
The following conditions are equivalent (and imply that $X$ is smooth over $k$ at $x$ and that $k(x)$ is a separable extension of $k$):
\begin{enumerate}
  \item[a)] $f$ is \'etale at the point $x$.
\oldpage[IV]{102}
  \item[b)] The germs $(g_i)_x$ for $1\leq i\leq r$ generate $\mathfrak m_x$ (and hence form a regular system of parameters for $\sh{O}_{X,x}$), and the images in $\Omega^1_{k(x)/k}$ of the elements $(d_{X/k}g_i)_x$ for $r+1\leq i\leq n$ generate $\Omega^1_{k(x)/k}$.
\end{enumerate}
\end{proposition}

\begin{proof}
There is an exact sequence of $k(x)$-modules ($\mathbf{0}$, 20.5.12.1)
\[
\label{IV.17.15.8.1}
\tag{17.15.8.1}
 \mathfrak m_x/\mathfrak m_x^2
 \longrightarrow (\Omega^1_{X/k})_x\otimes_{\sh{O}_{X,x}}k(x)
 \longrightarrow \Omega^1_{k(x)/k}\longrightarrow0.
\]
Condition $b)$ and Nakayama's lemma therefore imply that the $(d_{X/k}g_i)_x$ generate $(\Omega^1_{X/k})_x$; hence $b)\Rightarrow a)$ by \sref{IV.17.15.3}.

Conversely, suppose $a)$ holds.  By \sref{IV.17.15.3}, the $(d_{X/k}g_i)_x$, $1\leq i\leq n$, form a basis of $(\Omega^1_{X/k})_x$.  If $t_i$ is the image of $(g_i)_x$ in $k(x)$, it follows that the $dt_i$ generate $\Omega^1_{k(x)/k}$.  Since $t_i=0$ for $1\leq i\leq r$, the $dt_i$ for $r+1\leq i\leq n$ already generate $\Omega^1_{k(x)/k}$.  As $\deg.\operatorname{tr}_k k(x)=n-r$, Cartier's equality gives $\Upsilon_{k(x)/k}=0$; thus $k(x)$ is separable over $k$ ($\mathbf{0}$, 20.6.3), and the sequence
\[
\label{IV.17.15.8.2}
\tag{17.15.8.2}
 0\longrightarrow\mathfrak m_x/\mathfrak m_x^2
 \longrightarrow(\Omega^1_{X/k})_x\otimes_{\sh{O}_{X,x}}k(x)
 \longrightarrow\Omega^1_{k(x)/k}\longrightarrow0
\]
is exact (($\mathbf{0}$, 20.5.14) and ($\mathbf{0}$, 19.6.1)).  Moreover the $dt_i$, $r+1\leq i\leq n$, form a basis of $\Omega^1_{k(x)/k}$, so none of the corresponding $t_i$ is zero.  It follows that the images of the $(g_i)_x$, $1\leq i\leq r$, generate $\mathfrak m_x/\mathfrak m_x^2$; by Nakayama's lemma the $(g_i)_x$, $1\leq i\leq r$, generate $\mathfrak m_x$.  This proves $a)\Rightarrow b)$.
\end{proof}

\begin{corollary}[17.15.9]
\label{IV.17.15.9}
Let $X$ be a prescheme locally of finite type over a field $k$, and $x$ a point of $X$; put $n=\dim_x(X)$ and $r=\dim(\sh{O}_{X,x})$.
The following conditions are equivalent:
\begin{enumerate}
  \item[a)] The ring $\sh{O}_{X,x}$ is regular and $k(x)$ is a separable extension of $k$.
  \item[b)] $X$ is smooth over $k$ at $x$, and the canonical homomorphism
\[
\mathfrak m_x/\mathfrak m_x^2 \longrightarrow (\Omega^1_{X/k})_x\otimes_{\sh{O}_{X,x}} k(x)
\]
is injective.
  \item[c)] There exist $n$ sections $g_i$ of $\sh{O}_X$ over an open neighbourhood $U$ of $x$, with $(g_i)_x\in\mathfrak m_x$ for $1\leq i\leq r$, such that the morphism $U\to\operatorname{Spec}(k[\mathrm{T}_1,\ldots,\mathrm{T}_n])$ corresponding to the $g_i$ (cf.\ \sref{IV.17.15.8}) is \'etale at $x$.
  \item[d)] There exist $n$ sections $g_i$ $(1\leq i\leq n)$ of $\sh{O}_X$ over an open neighbourhood $U$ of $x$, such that the $(g_i)_x$, $1\leq i\leq r$, generate $\mathfrak m_x$, and the images in $\Omega^1_{k(x)/k}$ of the $(d_{U/k}g_i)_x$, $r+1\leq i\leq n$, generate $\Omega^1_{k(x)/k}$.
\end{enumerate}
\end{corollary}

\begin{proof}
The equivalence of $c)$ and $d)$ follows from \sref{IV.17.15.8}.  Its proof also shows that $c)$ implies that $X$ is smooth over $k$ at $x$ and that the sequence \sref{IV.17.15.8.2} is exact; hence $c)\Rightarrow b)$.

Condition $b)$ implies that $\sh{O}_{X,x}$ is regular, so $\mathfrak m_x/\mathfrak m_x^2$ has rank $r$; moreover $(\Omega^1_{X/k})_x$ is then free of rank $n$ \sref{IV.17.10.2}.  Since the sequence \sref{IV.17.15.8.1} is exact by hypothesis, $\Omega^1_{k(x)/k}$ has rank $n-r=\deg.\operatorname{tr}_k k(x)$.  Cartier's equality therefore gives $\Upsilon_{k(x)/k}=0$, and $k(x)$ is separable over $k$ ($\mathbf{0}$, 20.6.3).  Thus $b)\Rightarrow a)$.

Finally, if $a)$
\oldpage[IV]{103}
holds, Cartier's equality again shows that $\Omega^1_{k(x)/k}$ has rank $n-r$.  Since $\sh{O}_{X,x}$ is regular, the existence of sections $g_i$ satisfying $d)$ is immediate.
\end{proof}

\begin{remarks}[17.15.10]
\label{IV.17.15.10}
\begin{enumerate}[label=(\roman*)]
  \item For a prescheme $X$ of finite type over $k$ to be smooth over $k$, it is not enough that $\Omega^1_{X/k}$ be locally free.  This is shown by $X=\operatorname{Spec}(K)$, where $K$ is a finite inseparable extension of $k$.
  \item When $k$ is not perfect, $X$ may be smooth over $k$ even though $k(x)$ is not separable over $k$.  For example, take $X=\operatorname{Spec}(k[\mathrm{T}])$ and let $x$ correspond to the principal prime ideal $(\mathrm{T}^p-\lambda)$, where $p>0$ is the characteristic of $k$ and $\lambda\in k-k^p$.
  \item Nevertheless, if $X$ is smooth over $k$, the set of closed points $x$ of $X$ for which $k(x)$ is a separable (and finite) extension of $k$ is \emph{dense} in $X$.  Indeed, let $f:X\to\operatorname{Spec}(k)$ be the structural morphism.  For every $x_0\in X$, there is an open neighbourhood $U$ of $x_0$ and a factorisation
  \[
   f|U:U\xrightarrow{g}\operatorname{Spec}(k[\mathrm{T}_1,\ldots,\mathrm{T}_n])\xrightarrow{h}\operatorname{Spec}(k),
  \]
  in which $g$ is \'etale \sref{IV.17.11.4}.  We may restrict to the case where $k$ is not perfect, and hence is infinite.  The open set $g(U)$ then contains a point $y$ rational over $k$; if $x\in U$ lies above $y$, then $x$ is closed in $X$ and $k(x)$ is separable over $k(y)=k$ \sref{IV.17.6.2}.
\end{enumerate}
\end{remarks}

\begin{proposition}[17.15.11]
\label{IV.17.15.11}
Let $X$ be a prescheme of finite type over a field $k$.
The following conditions are equivalent:
\begin{enumerate}
  \item[a)] $X$ is \'etale over $k$.
  \item[b)] $X$ is unramified over $k$.
  \item[c)] $X$ is isomorphic to $\operatorname{Spec}(A)$, where $A$ is a finite and separable $k$-algebra.
\end{enumerate}
\end{proposition}

\begin{proof}
This follows from \sref{IV.17.6.2} and \sref{IV.17.4.2}, noting that condition $b)$ here implies that $X$ is finite over $k$.
\end{proof}

\begin{proposition}[17.15.12]
\label{IV.17.15.12}
Let $k$ be a field, $X$ a prescheme locally of finite type over $k$.
There exists an everywhere dense open subset $U$ of $X$ which is smooth over $k$ if and only if, at every maximal point $x$ of $X$, the prescheme $X$ is reduced and $k(x)$ is a separable extension of $k$.
There exists an everywhere dense open subset $U$ of $X$ such that $U_{\mathrm{red}}$ is smooth over $k$ if and only if $k(x)$ is a separable extension of $k$ for every maximal point $x$ of $X$.
\end{proposition}

\begin{proof}
The second assertion follows immediately from the first.  To say that there is an everywhere dense open subset $U$ smooth over $k$ means that $X$ is smooth over $k$ at each of its maximal points $x$.  Necessarily $\sh{O}_{X,x}$ is regular \sref{IV.17.15.1}, and \emph{a fortiori} reduced; it is therefore the field $k(x)$.  Moreover, \sref{IV.17.15.1} requires $k(x)\otimes_k k'$ to be regular for every radicial extension $k'$ of $k$, so $k(x)$ must be separable over $k$ \sref{IV.4.3.5}.  The converse follows at once from \sref{IV.17.15.1}.
\end{proof}

\begin{corollary}[17.15.13]
\label{IV.17.15.13}
Let $X$ be an algebraic prescheme over a field $k$.
Then there exist an everywhere dense open subset $U$ of $X$ and a finite radicial extension $k'$ of $k$ such that $(U_{(k')})_{\mathrm{red}}$ is smooth over $k'$.
\end{corollary}

\begin{proof}
There is such an extension $k'$ for which $(X_{(k')})_{\mathrm{red}}$ is geometrically reduced over $k'$ \sref{IV.4.6.6}; equivalently \sref{IV.4.6.1}, $k(x')$ is separable over $k'$ at every maximal point $x'$ of $X_{(k')}$.  It follows from \sref{IV.17.15.12} that there is an everywhere dense open subset $U'$ of $X_{(k')}$ such that $U'_{\mathrm{red}}$ is smooth over $k'$.  The morphism
\oldpage[IV]{104}
$X_{(k')}\to X$ is a homeomorphism \sref{IV.2.4.5}, so $U'=U_{(k')}$ for an everywhere dense open subset $U$ of $X$.
\end{proof}

\begin{proposition}[17.15.14]
\label{IV.17.15.14}
Let $X$ be an algebraic prescheme over a field $k$, of dimension $\leq 1$.
Then there exists a finite radicial extension $k'$ of $k$ such that the normalization ($\mathbf{II}$, 6.3.8) of $(X_{(k')})_{\mathrm{red}}$ is smooth over $k'$.
\end{proposition}

We first prove the following two lemmas.

\begin{lemma}[17.15.14.1]
\label{IV.17.15.14.1}
Let $X$ be a reduced prescheme whose set of irreducible components is locally finite, and let $f:Y\to X$ be a morphism.  Then $Y$ is $X$-isomorphic to the normalization $X'$ of $X$ ($\mathbf{II}$, 6.3.8) if and only if the following conditions hold:
\begin{enumerate}[label=(\roman*)]
 \item $Y$ is normal;
 \item $f$ is integral and birational.
\end{enumerate}
If $X$ contains a dense normal open subset (as is always the case when $X$ is excellent \sref{IV.7.8.3}, in particular when $X$ is locally of finite type over a field), condition (ii) may be replaced by:
\begin{enumerate}[label=(\roman* bis)]
 \setcounter{enumi}{1}
 \item $f$ is integral, and there exists a dense open subset $U$ of $X$ such that $f^{-1}(U)$ is dense in $Y$ and the restriction $f^{-1}(U)\to U$ is an isomorphism.
\end{enumerate}
\end{lemma}

\begin{proof}
The question is local on $X$, so we may suppose that $X$ has only finitely many irreducible components.  Let $X_i$ $(1\leq i\leq m)$ be the reduced subpreschemes of $X$ having these components as their underlying spaces.  The normalization $X'$ is the sum of the normalizations $X_i'$ of the $X_i$ ($\mathbf{II}$, 6.3.6).  Each structural morphism $f_i:X_i'\to X_i$ is integral and birational, so $f$ is integral and birational.  If $X$ has a dense normal open subset $V$, condition (ii) plainly implies (ii bis), with $U=V$.

Conversely, suppose (i) and (ii) hold.  The prescheme $Y$ has only finitely many irreducible components $Y_i$ $(1\leq i\leq m)$, with $Y_i$ dominating $X_i$ for each $i$ \sref{IV.6.15.4}.  Since $Y$ is normal it is the sum of the $Y_i$, and the morphism $g_i:Y_i\to X_i$ through which the restriction $Y_i\to X$ of $f$ factors is birational \sref[I]{I.5.2.2}.  As $X_i$ and $Y_i$ are integral, $g_i$ is integral, and $Y_i$ is normal, for every affine open $V$ of $X_i$ the ring $\Gamma(g_i^{-1}(V),\sh{O}_{Y_i})$ is the integral closure of $\Gamma(V,\sh{O}_{X_i})$.  Thus $Y$ identifies canonically with $X'$ ($\mathbf{II}$, 6.3.4).

Finally, suppose (ii bis).  We may take $U$ to be a disjoint union of irreducible open subsets $U_i$ $(1\leq i\leq m)$; then $f^{-1}(U)$ is the disjoint union of the $V_i=f^{-1}(U_i)$.  Since the $V_i$ are irreducible and $f^{-1}(U)$ is dense in $Y$, they are the irreducible components of $Y$, and $f$ is therefore birational.  This proves the lemma.
\end{proof}

\begin{lemma}[17.15.14.2]
\label{IV.17.15.14.2}
Let $k$ be a field and $X$ an algebraic prescheme over $k$.  There exists a finite radicial extension $k'$ of $k$ such that $(X_{(k')})_{\mathrm{red}}$ is geometrically reduced over $k'$ and its normalization is geometrically normal over $k'$.
\end{lemma}

\begin{proof}
In view of \sref{IV.4.6.6} and \sref[I]{I.5.1.8}, we may reduce to the case where $X$ is already geometrically reduced over $k$.  Let $p$ be the characteristic exponent of $k$, and let $k_1=k^{p^{-\infty}}$ be the smallest perfect extension of $k$; thus $k_1$ is the filtered direct limit of the finite radicial extensions of $k$.  Put $X_1=X_{(k_1)}$, which is reduced by hypothesis, and let $Y_1$ be its normalization.  If $f_1:Y_1\to X_1$ is the structural morphism, then $f_1$ is finite \sref{IV.7.8.3},~(vi), and surjective.  There is a dense open subset $U_1$ of $X_1$ such that $V_1=f_1^{-1}(U_1)$ is dense in $Y_1$ and
\oldpage[IV]{105}
$V_1\to U_1$ is an isomorphism \sref{IV.17.15.14.1}.

Applying \sref{IV.8.9.1} and \sref{IV.8.10.5},~(vi) and (x), we first obtain a finite radicial extension $k''$ of $k$ and a finite surjective morphism $Y''\to X''=X_{(k'')}$ such that
\[
 Y_1=Y''\times_{X''}X_1=Y''\otimes_{k''}k_1.
\]
Since $Y_1$ is normal and $k_1$ is perfect, \sref{IV.6.7.7} shows that $Y''$ is geometrically normal over $k''$.  The projections $X_1\to X''$ and $Y_1\to Y''$ are integral, surjective, and radicial, hence homeomorphisms \sref{IV.2.4.5}.  If $U''$ and $V''$ are the respective images of $U_1$ and $V_1$, they are dense open subsets and $V''=f''^{-1}(U'')$, where $f'':Y''\to X''$ is the structural morphism.  Since $U_1=U''\otimes_{k''}k_1$, \sref{IV.8.10.5},~(i), gives a finite radicial extension $k'$ of $k''$ such that, after putting
\[
 X'=X\otimes_k k'=X''\otimes_{k''}k',\qquad
 Y'=Y''\otimes_{k''}k',
\]
and writing $f':Y'\to X'$ for the resulting morphism, the images $U'$ and $V'$ of $U''$ and $V''$ satisfy that $V'\to U'$ is an isomorphism.  Since $Y'$ is normal and $f'$ is integral and birational, \sref{IV.17.15.14.1} identifies $Y'$ with the normalization of $X'$.  This proves the lemma, because $Y'$ is geometrically normal.
\end{proof}

\begin{proof}[Proof of Proposition \textup{{\hyperlink{ega4.c04.prop.17-15-14}{17.15.14}}}]
Apply \sref{IV.17.15.14.2} to $k$ and $X$.  Since $\dim X\leq1$, one also has $\dim X_{(k')}\leq1$ \sref{IV.4.1.4}, and the normalization $Y'$ of $X_{(k')}$ has dimension at most $1$ (\sref{IV.5.4.2} and $\mathbf{II}$, 7.4.6).  By the definitions and ($\mathbf{II}$, 7.4.5), geometric normality of $Y'$ is then equivalent to geometric regularity over $k'$, and hence to smoothness over $k'$ \sref{IV.17.5.1}.
\end{proof}

\begin{proposition}[17.15.15]
\label{IV.17.15.15}
Let $f:X\to Y$ be a morphism locally of finite presentation.
For $f$ to be smooth at a point $x\in X$, it is necessary and sufficient that $f$ be flat at $x$ and that $\Omega^1_f$ be a locally free $\sh{O}_X$-module in a neighbourhood of $x$, of rank at $x$ equal to $\dim_x f$.
\end{proposition}

\begin{proof}
Necessity follows from \sref{IV.17.5.1} and \sref{IV.17.10.2}.  Conversely, suppose the conditions hold and put $y=f(x)$.  By \sref{IV.17.5.1}, it is enough to show that $f^{-1}(y)$ is smooth over $k(y)$ at $x$; this follows from the definition of $\dim_x f$ \sref{IV.17.10.1}, from \sref{IV.16.4.5}, and from \sref{IV.17.15.5}.
\end{proof}


\subsection{Quasi-sections of flat or smooth morphisms.}
\label{subsection:IV.17.16}
\label{IV.17.16}

The statements of this number complement those of \sref{IV.14.5}, by means of hypotheses of flatness.

\begin{proposition}[17.16.1]
\label{IV.17.16.1}
Let $f:X\to S$ be a flat morphism locally of finite presentation.  Let $s\in S$, and let $x$ be a closed point of $X_s$ such that $\sh{O}_{X_s,x}$ is a Cohen--Macaulay ring.  Then there exist an open neighbourhood $U$ of $s$ in $S$ and a subprescheme $X'\subset f^{-1}(U)$ of $X$, containing $x$, such that the restriction $X'\to U$ of $f$ is flat, quasi-finite, and of finite presentation.
\end{proposition}

\begin{proof}
The question is local on $X$ and $S$, so we may suppose that $f$ is of finite presentation.  Let $(t_i)_{1\leq i\leq m}$ be a system of parameters of the local ring $\sh{O}_{X_s,x}$.  There exist an affine open neighbourhood $V=\operatorname{Spec}(A)$ of $x$ in $X$ and $m$ sections $g_i$ of $\sh{O}_X$ over $V$ such that

\oldpage[IV]{106}
the images of the germs $(g_i)_x\in\sh{O}_{X,x}$ in
\[
 \sh{O}_{X_s,x}=\sh{O}_{X,x}/\mathfrak m_s\sh{O}_{X,x}
\]
are the $t_i$.  Let $X'$ be the closed subprescheme of $V$ defined by the ideal of $A$ generated by the $g_i$.  The sequence $(t_i)$ is regular by hypothesis ($\mathbf{0}$, 16.5.7); it therefore follows from \sref{IV.11.3.8}, after replacing $V$ if necessary by a smaller open neighbourhood, that the restriction $X'\to S$ of $f$ is flat and of finite presentation.  Since the $t_i$ form a system of parameters of $\sh{O}_{X_s,x}$, the ring $\sh{O}_{X'_s,x}$ is Artinian.  As $x$ is closed in $X'_s$, it is consequently isolated in $X'_s$.  Replacing $V$ once more by a smaller neighbourhood of $x$, \sref{IV.13.1.4} shows that $X'\to S$ is quasi-finite.
\end{proof}

\begin{corollary}[17.16.2]
\label{IV.17.16.2}
Let $f:X\to S$ be faithfully flat and locally of finite presentation.  Then there exists a faithfully flat, locally finitely presented, and locally quasi-finite morphism $g:S'\to S$ such that there is an $S$-morphism $S'\to X$ (equivalently, an $S'$-section of $X'=X\times_SS'$; $\mathbf{I}$, 3.3.14).
If $S$ is quasi-compact (resp.\ quasi-compact and quasi-separated), $S'$ may be chosen affine (resp.\ $S'$ may be chosen affine and $g$ of finite presentation and quasi-finite).
\end{corollary}

\begin{proof}
For every $s\in S$, the fibre $X_s$ is nonempty and is a prescheme locally of finite type over $k(s)$.  The set $U_s$ of points $x$ of $X_s$ for which $\sh{O}_{X_s,x}$ is Cohen--Macaulay is open in $X_s$ \sref{IV.6.11.3} and nonempty, since it contains the maximal points of $X_s$ ($\mathbf{0}$, 16.5.1).  It therefore contains a point $x_s$ closed in $X_s$ \sref{IV.10.4.7}.  Let $X'(s)$ be an affine subprescheme of $X$, containing $x_s$, and satisfying the conclusion of \sref{IV.17.16.1}.

To obtain a prescheme $S'$ satisfying the assertion, it is enough to take the sum of the $X'(s)$ as $s$ runs through $S$.  Indeed, since $X'(s)\to S$ is flat and locally of finite presentation, its image $U(s)$ is open in $S$ \sref{IV.2.4.6}.  If $S$ is quasi-compact, it can therefore be covered by finitely many of the $U(s_i)$, and the sum of the corresponding $X'(s_i)$ is still suitable and is affine.  If, in addition, $S$ is quasi-separated, the open immersions $U(s_i)\to S$ may be supposed quasi-compact \sref[I]{I.2.7}, hence of finite presentation \sref[I]{I.6.2}.  The morphisms $X'(s_i)\to S$ are then of finite presentation, and so is $S'\to S$.
\end{proof}

\begin{corollary}[17.16.3]
\label{IV.17.16.3}
(i) Let $f:X\to S$ be smooth.  Let $s\in S$, and let $x$ be a closed point of $X_s$ such that the residue field $k(x)$ is separable over $k(s)$.  In the conclusion of \sref{IV.17.16.1}, the restriction $X'\to U$ of $f$ may then be chosen \'etale.

(ii)
Let $f:X\to S$ be smooth and surjective.  In the conclusion of \sref{IV.17.16.2}, one may moreover suppose that $g:S'\to S$ is \'etale.
\end{corollary}

\begin{proof}
Assertion (ii) follows from (i) just as \sref{IV.17.16.2} follows from \sref{IV.17.16.1}: for every $s\in S$, the nonempty smooth $k(s)$-prescheme $X_s$ has a closed point $x$ for which $k(x)$ is separable over $k(s)$ \sref{IV.17.15.10},~(iii).  It remains to prove (i).  The ring $\sh{O}_{X_s,x}$ is regular.  Repeat the construction of \sref{IV.17.16.1}, taking for $(t_i)$ a regular system of parameters of $\sh{O}_{X_s,x}$.  The ring $\sh{O}_{X'_s,x}$ is then a field isomorphic to $k(x)$, and hence separable over $k(s)$ by hypothesis.  The conclusion follows from \sref{IV.17.6.1},~$c')$.
\end{proof}

\begin{proposition}[17.16.4]
\label{IV.17.16.4}
Let $S$ be a quasi-compact and quasi-separated prescheme, and let $f:X\to S$ be a surjective morphism locally of finite presentation.  Then there exists a finite family $(S_i)_{i\in I}$
\oldpage[IV]{107}
of affine subpreschemes of $S$, of finite presentation over $S$, pairwise disjoint, and with union $S$, having the following property: for every $i\in I$, there exist two finite, finitely presented, and surjective morphisms
\[
 S_i''\xrightarrow{g_i}S_i'\xrightarrow{h_i}S_i,
\]
where $h_i$ is \'etale and $g_i$ is flat and radicial, and an $S$-morphism $S_i''\to X$ (in other words, an $S_i''$-section of $X_i''=X\times_SS_i''$).
\end{proposition}

\begin{proof}
There is a finite covering $(U_i)_{1\leq i\leq n}$ of $S$ by affine open subsets.  For every $i$, the intersection
\[
 W_i=U_i\cap\Bigl(\bigcup_{j<i}U_j\Bigr)
\]
is quasi-compact \sref[I]{I.2.7}; hence there exists a closed subprescheme $V_i$ of $U_i$ whose underlying space is $U_i-W_i$ and which is defined by an ideal of finite type in the ring of $U_i$.  Thus $V_i$ is an affine subprescheme of $S$ of finite presentation over $S$ \sref[I]{I.6.2}.  It is clear that it is enough to prove the proposition after replacing $S$ by each $V_i$ and $X$ by $X\times_SV_i$.  In other words, we may already suppose that $S$ is affine.

On the other hand, $X$ is the union of affine open subsets $W_\alpha$, and the $f(W_\alpha)$ are constructible in $S$ \sref[I]{I.8.4} and cover $S$.  Consequently \sref[I]{I.9.9}, there is a finite subfamily $(W_{\alpha_j})_{1\leq j\leq m}$ whose images already cover $S$.  Let $X'$ be the sum prescheme of the subpreschemes induced on the open subsets $W_{\alpha_j}$ of $X$.  It is immediate that it is enough to prove the proposition after replacing $X$ by $X'$, since an $S$-morphism $S_i''\to X'$ gives, by composition, an $S$-morphism $S_i''\to X'\to X$.  We may therefore suppose that $X$ is affine and that $f$ is of finite presentation.

We may then write \sref{IV.8.9.1}
\[
 X=X_0\times_{S_0}S,
\]
where $S_0$ is Noetherian and $X_0\to S_0$ is a morphism of finite type which may be supposed surjective \sref{IV.8.10.5},~(vi).  If the proposition is proved for this morphism, it follows at once for $X$ by the base change $S\to S_0$.  We may thus suppose in addition that $S$ is Noetherian and that $f$ is of finite type.

Let $s_h$ $(1\leq h\leq r)$ be the maximal points of $S$.  We show that it is enough to prove the assertion after replacing $S$ by a sufficiently small open neighbourhood $U_h$ of each $s_h$, and $X$ by the inverse image of $U_h$ in $X$.  Indeed, suppose the proposition established in this case, and argue by Noetherian induction ($\mathbf{0}_{\mathbf I}$, 2.2.2), assuming the assertion established for every closed subprescheme of $S$ whose underlying space is not $S$.  We may suppose that the $U_h$ are pairwise disjoint.  If $T$ is a closed subprescheme whose underlying space is
\[
 S-\bigcup_hU_h,
\]
the induction hypothesis implies that the assertion is true for $T$; since it is also true for each $U_h$, it is plainly true for $S$.

Since we may clearly suppose $S$ reduced (by replacing $S$ by $S_{\mathrm{red}}$ and $X$ by $X\times_SS_{\mathrm{red}}$), each $\sh{O}_{S,s_h}$ is a field $k(s_h)$.  Suppose that the proposition has been proved when $S$ is the spectrum of a field and $f$ is of finite type.  The existence of the $U_h$ then follows from the method of \sref{IV.8.1.2},~a), together with \sref{IV.8.8.2},~(i) and (ii), \sref{IV.8.10.5},~(vi), (vii), and (x), \sref{IV.11.2.6},~(ii), and \sref{IV.17.7.8},~(ii).

We may therefore suppose that $S=\operatorname{Spec}(k)$, where $k$ is a field, and that $X$ is a nonempty prescheme of finite type over $k$.  Since $X$ is Noetherian, it has a closed point $x$ ($\mathbf{0}_{\mathbf I}$, 2.1.3), and $k(x)$ is then a finite extension of $k$ \sref[I]{I.6.4.2}.  There is consequently a finite separable extension $k'$ of $k$ such that $k''=k(x)$ is a finite radicial extension of $k'$.  The assertion follows by taking $I$ to have one element, $S_i'=\operatorname{Spec}(k')$, and $S_i''=\operatorname{Spec}(k'')$.
\end{proof}

\oldpage[IV]{108}

\begin{corollary}[17.16.5]
\label{IV.17.16.5}
Let $f:X\to S$ be a surjective morphism locally of finite presentation.
Then there exists a surjective morphism $g:S'\to S$, locally of finite presentation and locally quasi-finite, such that there is an $S$-morphism $S'\to X$ (in other words, an $S'$-section of $X'=X\times_S S'$).
If $S$ is quasi-compact (resp.\ quasi-compact and quasi-separated), $S'$ may be supposed affine (resp.\ $S'$ may be supposed affine and $g$ of finite presentation and quasi-finite).
\end{corollary}

\begin{proof}
It is enough to prove the corollary when $S$ is affine.  Indeed, $S$ is the union of a family $(S_\alpha)$ of affine open subsets.  If, for every $\alpha$, $S'_\alpha$ is affine and a morphism $g_\alpha:S'_\alpha\to S_\alpha$ satisfies the assertion and is of finite presentation and quasi-finite, take for $S'$ the sum prescheme of the $S'_\alpha$, with $g:S'\to S$ agreeing with $g_\alpha$ on each $S'_\alpha$.  This morphism has the required properties.  Moreover, if $S$ is quasi-compact, the family $(S_\alpha)$ may be supposed finite, so $S'$ is affine; if $S$ is also quasi-separated, the immersions $S_\alpha\to S$ are of finite presentation \sref[I]{I.6.2}, and consequently so is $g$.

Now suppose that $S$ is affine.  Form the finite morphisms of finite presentation $S_i''\to S_i$ supplied by \sref{IV.17.16.4}.  Since the immersions $S_i\to S$ are of finite presentation (the $S_i$ being affine), the composite morphisms
\[
 g_i'':S_i''\longrightarrow S_i\longrightarrow S
\]
are of finite presentation and quasi-finite.  The assertion follows by taking $S'$ to be the sum of the $S_i''$ and $g$ to agree with $g_i''$ on each $S_i''$.
\end{proof}

\begin{proposition}[17.16.6]
\label{IV.17.16.6}
Let $S$ be a quasi-compact and quasi-separated prescheme, and let $f:X\to S$ be a morphism of finite presentation such that, for every $s\in S$, the fibre $X_s=f^{-1}(s)$ is a prescheme proper over $k(s)$.
Then there exists a finite family $(S_i)_{i\in I}$ of affine subpreschemes of $S$, of finite presentation over $S$, pairwise disjoint and with union $S$, such that, for every $i$, the morphism $X_i=X\times_S S_i\to S_i$ is proper and flat.
If, moreover, $X_s$ is finite over $k(s)$ for every $s\in S$, the $S_i$ may be chosen so that each morphism $X_i\to S_i$ factors as
\[
 X_i \xrightarrow{u_i} X'_i \xrightarrow{h_i} S_i,
\]
where $u_i$ and $h_i$ are finite and locally free \sref{IV.18.2.7}, $h_i$ is \'etale, and $u_i$ is radicial and surjective.
\end{proposition}

\begin{proof}
We follow a course analogous to that of \sref{IV.17.16.4}.  First reduce to the case in which $S$ is affine and $f$ is flat, using the generic-flatness theorem \sref{IV.8.9.5}; note that the subpreschemes $S_i'$ defined in the proof of \sref{IV.8.9.5} are affine.  Since $f$ is of finite presentation, we may then write
\[
 X=X_0\times_{S_0}S,
\]
where $S_0$ is Noetherian and $X_0\to S_0$ is a morphism of finite type and flat \sref{IV.11.2.7}.  Moreover, every fibre $(X_0)_{s_0}$ is still proper over $k(s_0)$, by \sref{IV.2.7.1},~(vii), and we are therefore reduced to the case in which $S$ is Noetherian.  By Noetherian induction and application of the method of \sref{IV.8.1.2},~a) (using also \sref{IV.8.10.5},~(xii) on this occasion), the proof of the first assertion is finally reduced to the case in which $S$ is the spectrum of a field; it then follows trivially from the hypothesis, with $S_i=S$.

The case in which $X_s$ is supposed finite over $k(s)$ for every $s\in S$ is treated in the same way, now using \sref{IV.2.7.1},~(xv), \sref{IV.8.10.5},~(x), (iv), (vi), and (vii), \sref{IV.17.7.8}, and \sref{IV.2.1.12}.  We are reduced to proving the last assertion when $S=\operatorname{Spec}(k)$ is the spectrum of a field.  Then $X=\operatorname{Spec}(A)$, where $A$ is a finite $k$-algebra \sref[I]{I.6.4.4}; thus $A$ is a direct product of finite $k$-algebras $A_j$ which are local rings $(1\leq j\leq m)$.  Since $A_j$ is an Artinian ring and a
\oldpage[IV]{109}
$k$-algebra, it contains a subfield $K'_j$ canonically isomorphic to the residue field of $A_j$ ($\mathbf{0}$, 19.6.2).  Take $X_i=X$ and take $X'_i$ to be the sum of the $\operatorname{Spec}(K'_j)$.  This answers the question, since for every $j$ there are homomorphisms
\[
 K'_j\longrightarrow A_j\longrightarrow k(A_j)
\]
whose composite is an isomorphism.
\end{proof}
% END NATIVE INPUT source_aligned/ega4-17
% BEGIN NATIVE INPUT source_aligned/ega4-18
\section{Supplement on \'etale morphisms. Henselian local rings and strictly local rings}
\label{section:IV.18}
\label{IV.18}

In the present section we study various properties peculiar to \'etale morphisms.  Moreover, the notion of an \'etale morphism makes it possible to develop, in a very natural way, Nagata's theory of Henselian rings, as well as that of strictly local rings.  These rings play an important role in many recent developments, appearing whenever one needs a process of ``localisation'' finer than that furnished by the Zariski topology (cf.\ for example~[43]), pending the appearance of the chapter of our Treatise devoted to the study of the ``\'etale topology.''

\subsection{A remarkable equivalence of categories.}
\label{subsection:IV.18.1}

\begin{proposition}[18.1.1]
\label{IV.18.1.1}
Let $S$ be a prescheme, $S_0$ a closed subprescheme of $S$, $X_0$ an $S_0$-prescheme smooth (resp.\ \'etale) over $S_0$, and $x_0$ a point of $X_0$.  Then there exist an open neighbourhood $U_0$ of $x_0$ in $X_0$, an $S$-prescheme $U$ smooth (resp.\ \'etale) over $S$, and an $S_0$-isomorphism
\[
 U\times_SS_0\xrightarrow{\sim}U_0.
\]
\end{proposition}

\begin{proof}
Observe that if $X_0$ is \'etale over $S_0$ at $x_0$, it is \emph{a fortiori} unramified over $S_0$ at that point.  If an $S$-prescheme $U$ smooth over $S$ has been constructed such that $U\times_SS_0$ is isomorphic to $U_0$, then the fibres of the morphisms $U_0\to S_0$ and $U\to S$ containing $x_0$ are isomorphic.  It follows that $U$ is unramified over $S$ at $x_0$ \sref{IV.17.4.1},~d), hence also in a neighbourhood of $x_0$; replacing $U$ by this neighbourhood, we conclude that $U$ is \'etale over $S$.  It is therefore enough to prove the proposition when $X_0$ is merely supposed smooth over $S_0$.

The question being local on $S$ and on $X_0$, we may suppose that $S=\operatorname{Spec}(A)$ and $X_0=\operatorname{Spec}(C_0)$ are affine, so that $S_0=\operatorname{Spec}(A_0)$, where $A_0$ is a quotient ring of $A$, and
\[
 C_0=B_0/\mathfrak{J}_0,\qquad B_0=A_0[T_1,\ldots,T_n],
\]
where $\mathfrak{J}_0$ is an ideal of finite type in $B_0$; finally, $C_0$ is a formally smooth $A_0$-algebra for the discrete topologies.  Let $\mathfrak p_0$ be the ideal $\mathfrak j_{x_0}$ in $C_0$; then $\mathfrak p_0=\mathfrak q_0/\mathfrak{J}_0$, where $\mathfrak q_0$ is a prime ideal of $B_0$.  The Jacobian criterion ($\mathbf{0}$, 22.6.4), together with ($\mathbf{0}$, 19.1.12), shows that there exist in $\mathfrak{J}_0$ a family of $r$ polynomials $u_i$ $(1\leq i\leq r)$ and indices $j_h$ $(1\leq h\leq r)$ such that the images of the $u_i$ in
\[
 (\mathfrak{J}_0)_{\mathfrak q_0}/(\mathfrak{J}_0^2)_{\mathfrak q_0}
\]
generate this $(B_0)_{\mathfrak q_0}$-module and
\begin{equation}
 \det\left(\frac{\partial u_i}{\partial T_{j_h}}\right)\notin\mathfrak q_0.
 \tag{18.1.1.1}\label{IV.18.1.1.1}
\end{equation}

Since $(B_0)_{\mathfrak q_0}$ is a local ring, Nakayama's lemma shows that we may suppose that the images of the $u_i$ in $(\mathfrak{J}_0)_{\mathfrak q_0}$ generate this $(B_0)_{\mathfrak q_0}$-module and then, replacing $X_0$ if necessary by an affine open neighbourhood of $x_0$, that the $u_i$ generate $\mathfrak{J}_0$ ($\mathbf{0}_{\mathbf I}$, 5.2.2).  Put $B=A[T_1,\ldots,T_n]$; then $B_0$ is a quotient ring
\oldpage[IV]{110}
of $B$, $\mathfrak q_0$ is the image of a prime ideal $\mathfrak q$ of $B$, and $\mathfrak q$ is the inverse image of $\mathfrak q_0$.  For every $i$, let $v_i\in B$ be an element whose image in $B_0$ is $u_i$, and let $\mathfrak{J}$ be the ideal of $B$ generated by the $v_i$, so that $\mathfrak{J}_0$ is the image of $\mathfrak{J}$ in $B_0$.

The proposition will be proved by taking for $U$ an open neighbourhood of the point of $\operatorname{Spec}(B/\mathfrak{J})$ corresponding to the prime ideal $\mathfrak p=\mathfrak q/\mathfrak{J}$, provided that we show that $(B/\mathfrak{J})_{\mathfrak q/\mathfrak{J}}$ is a formally smooth $A$-algebra for the discrete topologies.  This follows from the Jacobian criterion: the images of the $v_i$ in $\mathfrak{J}_{\mathfrak q}/\mathfrak{J}_{\mathfrak q}^2$ generate this $B_{\mathfrak q}$-module, and \eqref{IV.18.1.1.1} gives
\[
 \det\left(\frac{\partial v_i}{\partial T_{j_h}}\right)\notin\mathfrak q.
\]
\end{proof}

\begin{theorem}[18.1.2]
\label{IV.18.1.2}
Let $S$ be a prescheme, $S_0$ a closed sub-prescheme of $S$ whose underlying space is identical to that of $S$.
Then the functor
\[
X\mapsto X\times_S S_0
\]
from the category of $S$-preschemes \'etale over $S$, to the category of $S_0$-preschemes \'etale over $S_0$, is an equivalence of categories.
\end{theorem}

\begin{proof}
Let us first show that this functor is \emph{fully faithful}.
Let $X$, $Y$ be two $S$-preschemes \'etale over $S$, and set $X_0=X\times_S S_0$, $Y_0=Y\times_S S_0$.
If $Z=X\times_S Y$, the set $\operatorname{Hom}_S(X,Y)$ is in canonical bijective correspondence with the set of $X$-sections $\Gamma(Z/X)$, and likewise $\operatorname{Hom}_{S_0}(X_0,Y_0)$ is in canonical bijective correspondence with $\Gamma(Z_0/X_0)$, where $Z_0=Z\times_S S_0=X_0\times_{S_0}Y_0$.
Now $Z$ is \'etale over $X$, $Z_0$ is \'etale over $X_0$, and $X_0$ (resp.\ $Z_0$) is a closed subprescheme of $X$ (resp.\ $Z$) having the same underlying space.  The open subsets of $Z$ on which the restriction of $Z\to X$ is surjective and radicial are therefore the \emph{same} as the open subsets of $Z_0$ having the corresponding properties, and our assertion follows from \sref{IV.17.9.3}.

To complete the proof, it suffices to show that for every $S_0$-prescheme $X_0$ \'etale over $S_0$, there exists an $S$-prescheme $X$ \'etale over $S$ and an $S_0$-isomorphism $X_0\xrightarrow{\sim}X\times_S S_0$.
By virtue of Proposition~\sref{IV.18.1.1}, there exists an open covering $(U_\alpha)$ of $X_0$ and for each $\alpha$, an $S$-prescheme $V_\alpha$ which is \'etale over $S$, and finally an $S_0$-isomorphism $\theta_\alpha:U_\alpha\xrightarrow{\sim}V_\alpha\times_S S_0$.
Furthermore, by virtue of the first part of the proof, the identity of $U_\alpha\cap U_\beta$ lifts uniquely, via $\theta_\alpha$ and $\theta_\beta$, to an $S$-isomorphism $\varphi_{\alpha\beta}$ between the corresponding open subpreschemes of $V_\alpha$ and $V_\beta$, and it is immediate, for the same reason, that these isomorphisms satisfy the glueing condition (\sref[0]{0.4.1.7}).
Hence there exists an $S$-prescheme $X$ such that the $V_\alpha$ identify canonically with sub-preschemes induced on the open subsets of $X$, the $\theta_\alpha$ identify with $S_0$-isomorphisms which coincide on the intersections $U_\alpha\cap U_\beta$, and define an $S_0$-isomorphism $X_0\xrightarrow{\sim}X\times_S S_0$.
It is clear that $X$ is \'etale over $S$ \sref{IV.17.3.2}, which completes the proof.
\end{proof}

\begin{corollary}[18.1.3]
\label{IV.18.1.3}
Let $S$ be a prescheme, $X$ an $S$-prescheme \'etale over $S$, $S'$ an $S$-prescheme, $S_0'$ a closed sub-prescheme of $S'$ having the same underlying space.
Then the canonical map $X(S')_S\to X(S_0')_S$ \sref[I]{I.3.4.3} is bijective.
\end{corollary}

\begin{proof}
Indeed, set $X'=X\times_S S'$, $X_0'=X\times_S S_0'$, so that $X(S')_S=\Gamma(X'/S')$ and $X(S_0')_S=\Gamma(X_0'/S_0')$;
the corollary follows from the fact that the functor defined in \sref{IV.18.1.2} (with $S$ and $S_0$ replaced by $S'$ and $S_0'$ respectively) is fully faithful (or directly from \sref{IV.17.9.3}).
\end{proof}

\subsection{\'Etale coverings.}
\label{subsection:IV.18.2}

\oldpage[IV]{111}

\begin{env}[18.2.1]
\label{IV.18.2.1}
Given a ring $A$ and a commutative $A$-algebra $B$ which is \emph{finite} and is a \emph{free} $A$-module, recall (Bourbaki, \emph{Alg.}, ch.~VIII, \textsection~12, no.~2) that one defines on $B$ an \emph{$A$-linear form} $\operatorname{Tr}_{B/A}$, the ``trace form''; from it one obtains a \emph{symmetric $A$-bilinear form} (also called the ``trace form'')
\[
\label{IV.18.2.1.1}
(x,y)\mapsto\operatorname{Tr}_{B/A}(xy)
\tag{18.2.1.1}
\]
whose datum is equivalent to that of the associated $A$-linear map $\operatorname{astr}_{B/A}:B\to\check{B}$ from the $A$-module $B$ into its \emph{dual} $\check{B}$; this map is equal to its transpose.  When $A$ is a \emph{field}, saying that this bilinear form is \emph{non-degenerate} is equivalent to saying that $B$ is a \emph{separable} $A$-algebra (Bourbaki, \emph{Alg.}, ch.~IX, \textsection~2, prop.~5).

Let $f:A\to A'$ be a ring homomorphism;
set $B'=B\otimes_A A'$, and let $g:B\to B'$ be the canonical homomorphism;
the image under $g$ of a basis of the $A$-module $B$ is then a basis of the $A'$-module $B'$, and it follows from the definitions that, for every $x\in B$,
\[
\label{IV.18.2.1.2}
\operatorname{Tr}_{B'/A'}(g(x))=f(\operatorname{Tr}_{B/A}(x)).
\tag{18.2.1.2}
\]
\end{env}

\begin{env}[18.2.2]
\label{IV.18.2.2}
Consider now a \emph{ringed space} $(X,\sh{O}_X)$ and let $\sh{B}$ be an $\sh{O}_X$-algebra which, as an $\sh{O}_X$-module, is \emph{locally free of finite rank}.  For every open subset $U\subset X$ such that $\sh{B}|_U$ is, as an $\sh{O}_U$-module, isomorphic to $\sh{O}_U^n$ (for some $n$ depending on $U$), $\Gamma(U,\sh{B})$ is a $\Gamma(U,\sh{O}_X)$-algebra which, as a $\Gamma(U,\sh{O}_X)$-module, is free of finite rank.  It therefore defines a $\Gamma(U,\sh{O}_X)$-linear form $\operatorname{Tr}_{\Gamma(U,\sh{B})/\Gamma(U,\sh{O}_X)}$, which we also denote by $\operatorname{Tr}_{\sh{B}|\sh{O}_X,U}$, and hence an associated linear map
\[
\operatorname{astr}_{\sh{B}|\sh{O}_X,U}:\Gamma(U,\sh{B})\to\Gamma(U,\sh{B})^\vee=\Gamma(U,\check{\sh{B}}).
\]

Furthermore, it follows from \sref{IV.18.2.1.2} that these linear maps are compatible with restriction from $U$ to a smaller open subset and therefore define, on the one hand, an \emph{$\sh{O}_X$-module homomorphism}, again called the \emph{trace homomorphism},
\[
\label{IV.18.2.2.1}
\operatorname{Tr}_{\sh{B}|\sh{O}_X}:\sh{B}\to\sh{O}_X
\tag{18.2.2.1}
\]
and, on the other hand, an \emph{$\sh{O}_X$-module homomorphism}
\[
\label{IV.18.2.2.2}
\operatorname{astr}_{\sh{B}|\sh{O}_X}:\sh{B}\to\check{\sh{B}}
\tag{18.2.2.2}
\]
called the \emph{homomorphism associated to the trace}, and equal to its transpose.
It also follows from \sref{IV.18.2.1.2} that for every $x\in X$, we have
\[
\label{IV.18.2.2.3}
(\operatorname{Tr}_{\sh{B}|\sh{O}_X})_x=\operatorname{Tr}_{\sh{B}_x|\sh{O}_{X,x}}
\tag{18.2.2.3}
\]
\[
\label{IV.18.2.2.4}
(\operatorname{astr}_{\sh{B}|\sh{O}_X})_x=\operatorname{astr}_{\sh{B}_x|\sh{O}_{X,x}}.
\tag{18.2.2.4}
\]

Finally, under the conditions of \sref{IV.18.2.1}, if we set $X=\operatorname{Spec}(A)$ and $\sh{B}=\widetilde{B}$, the form $\operatorname{Tr}_{\sh{B}|\sh{O}_X}$ (resp.\ the homomorphism $\operatorname{astr}_{\sh{B}|\sh{O}_X}$) corresponds to the form $\operatorname{Tr}_{B/A}$ (resp.\ the $A$-module homomorphism $\operatorname{astr}_{B/A}$), as also follows from \sref{IV.18.2.1.2}.
\end{env}

\begin{proposition}[18.2.3]
\label{IV.18.2.3}
Let $f:X\to Y$ be a finite morphism of preschemes and set $\sh{B}=f_*(\sh{O}_X)$.
The following conditions are equivalent:
\begin{enumerate}
  \item[a)] $f$ is \'etale.
  \item[a$'$)] $f$ is a flat morphism of finite presentation, and for every $x\in X$, if $y=f(x)$, the ring $\sh{O}_{X,x}/\mathfrak{m}_y\sh{O}_{X,x}$ is a field which is a finite separable extension of $\kappa(y)$.
  \item[b)] $\sh{B}$ is a locally free $\sh{O}_Y$-module and, for every $y\in Y$, $\sh{B}_y\otimes_{\sh{O}_{Y,y}}\kappa(y)$ is a finite separable $\kappa(y)$-algebra (hence a direct product of finitely many fields which are finite separable extensions of $\kappa(y)$).
  \item[c)] $\sh{B}$ is a locally free $\sh{O}_Y$-module and the homomorphism $\operatorname{astr}_{\sh{B}|\sh{O}_Y}:\sh{B}\to\check{\sh{B}}$ \sref{IV.18.2.2} is bijective.
\end{enumerate}
\end{proposition}

\begin{proof}
\oldpage[IV]{112}
Since $f$ is quasi-compact, the equivalence of $a)$ and $a')$ has already been shown \sref{IV.17.6.2}.
To prove the rest of the proposition, we may restrict to the case where $Y=\operatorname{Spec}(A)$ and $X=\operatorname{Spec}(B)$ are affine, $B$ being a finite $A$-algebra and $\sh{B}=\widetilde{B}$.
To say that $f$ is a morphism of finite presentation is equivalent to saying that $B$ is an $A$-module of finite presentation \sref{IV.1.4.7}.
If, in addition, $f$ is flat, then $B$ is a flat $A$-module; it is known (Bourbaki, \emph{Alg.\ comm.}, ch.~II, \textsection~5, no.~2, cor.~2 to th.~1) that $B$ is then a \emph{projective} $A$-module.  Thus $\sh{B}$ is a locally free $\sh{O}_Y$-module (\emph{loc.\ cit.}, no.~2, th.~1), and the converse is immediate.
On the other hand, $f^{-1}(y)$ is nothing but the spectrum of the $\kappa(y)$-algebra $\sh{B}(y)=\sh{B}_y/\mathfrak{m}_y\sh{B}_y$, which completes the proof of the equivalence of $a')$ and $b)$.
To see that $b)$ is equivalent to $c)$, note that the second assertion of $b)$ is equivalent to the bijectivity of
\[
 \operatorname{astr}_{\sh{B}(y)|\kappa(y)}:\sh{B}(y)\longrightarrow\check{\sh{B}(y)}.
\]
Since
\[
 \sh{B}(y)=\sh{B}_y/\mathfrak{m}_y\sh{B}_y,
 \qquad
 \check{\sh{B}(y)}=\check{\sh{B}}_y/\mathfrak{m}_y\check{\sh{B}}_y,
\]
and $\sh{B}_y$ and $\check{\sh{B}}_y$ are free $\sh{O}_{Y,y}$-modules, it follows from \sref{IV.18.2.2.4} and Bourbaki, \emph{Alg.\ comm.}, ch.~II, \textsection~3, no.~3, cor.\ to prop.~6, that the homomorphism $\operatorname{astr}_{\sh{B}_y|\sh{O}_{Y,y}}:\sh{B}_y\to\check{\sh{B}}_y$ is itself bijective;
the converse being evident, this completes the proof.

When an $\sh{O}_X$-Algebra $\sh{B}$ satisfies the equivalent conditions $b)$ and $c)$ of \sref{IV.18.2.3}, we say that $\sh{B}$ is a \emph{finite \'etale $\sh{O}_X$-Algebra}.
When $X=\operatorname{Spec}(A)$ is affine and $\sh{B}=\widetilde{B}$, this amounts, by virtue of \sref{IV.18.2.3}, to saying that $B$ is a finite \'etale $A$-algebra, or that $B$ is a finite $A$-algebra that is \'etale (in the sense of \sref{IV.17.3.2}).
\end{proof}

\begin{corollary}[18.2.4]
\label{IV.18.2.4}
Let $f:X\to Y$ be a finite morphism of finite presentation, and set $\sh{B}=f_*(\sh{O}_X)$.
Let $y$ be a point of $Y$.
The following conditions are equivalent:
\begin{enumerate}[label=\emph{\alph*)}]
  \item There exists an open neighbourhood $U$ of $y$ in $Y$ such that the restriction $f^{-1}(U)\to U$ of $f$ is an \'etale morphism.
  \item $\sh{B}_y$ is a free $\sh{O}_{Y,y}$-module of finite type and $\sh{B}_y\otimes_{\sh{O}_{Y,y}}\kappa(y)$ is a separable $\kappa(y)$-algebra.
\end{enumerate}
\end{corollary}

\begin{proof}
It is clear that $a)$ implies $b)$ by virtue of \sref{IV.18.2.3}.
On the other hand, $\sh{B}$ is an $\sh{O}_Y$-module of finite presentation \sref{IV.1.6.3} and \sref{IV.1.4.7}; hence, if $\sh{B}_y$ is a free $\sh{O}_{Y,y}$-module, there exists an open neighbourhood $U$ of $y$ in $Y$ such that $\sh{B}|_U$ is a locally free $\sh{O}_U$-module ($\mathbf{0}_{\mathbf I}$, 5.2.7).  Moreover, since by hypothesis the homomorphism $\operatorname{astr}_{\sh{B}_y|\sh{O}_{Y,y}}:\sh{B}_y\to\check{\sh{B}}_y$ is bijective, the same result shows that $U$ may be chosen so that

\oldpage[IV]{113}

the homomorphism $\operatorname{astr}_{\sh{B}|_U/\sh{O}_U}$ is bijective.
The fact that $b)$ implies $a)$ then follows from \sref{IV.18.2.3}.
\end{proof}

\begin{corollary}[18.2.5]
\label{IV.18.2.5}
Let $Y$ be a quasi-compact or locally Noetherian prescheme, and let $f:X\to Y$ be a finite morphism of finite presentation;
set $\sh{B}=f_*(\sh{O}_X)$.
Suppose that for every \emph{closed} point $y$ of $Y$, $\sh{B}_y$ is a free $\sh{O}_{Y,y}$-module and $\sh{B}_y\otimes_{\sh{O}_{Y,y}}\kappa(y)$ is a separable $\kappa(y)$-algebra.
Then $f$ is \'etale.
\end{corollary}

\begin{proof}
Indeed, it follows from \sref{IV.18.2.4} that every closed point $y$ of $Y$ has an open neighbourhood $U$ such that the restriction $f^{-1}(U)\to U$ of $f$ is \'etale;
the conclusion follows from the fact that in the two cases considered, every non-empty closed subset of $Y$ contains a closed point (\sref{IV.5.1.11} and \sref[0]{0.2.1.3}).
\end{proof}

\begin{corollary}[18.2.6]
\label{IV.18.2.6}
If $f:X\to Y$ is a surjective finite \'etale morphism, and if we set $\sh{B}=f_*(\sh{O}_X)$, then the $\sh{O}_Y$-homomorphism $\operatorname{Tr}_{\sh{B}|\sh{O}_Y}:\sh{B}\to\sh{O}_Y$ \sref{IV.18.2.2} (also denoted $\operatorname{Tr}_f$) is surjective.
\end{corollary}

\begin{proof}
The question being local, we may, by virtue of \sref{IV.18.2.3}, suppose that $Y=\operatorname{Spec}(A)$, $X=\operatorname{Spec}(B)$, with $B$ a free $A$-module.
The trace pairing, perfect by \sref{IV.18.2.3}, identifies $B$ with $\check{B}$ and sends $1$ to $\operatorname{Tr}_{B/A}$; since $f$ is surjective, $1$ is unimodular, so the trace is surjective.
\end{proof}

\begin{remarks}[18.2.7]
\label{IV.18.2.7}
\begin{enumerate}
  \item[(i)] When $f:X\to Y$ is a finite morphism such that $f_*(\sh{O}_X)$ is a locally free $\sh{O}_Y$-Module (resp.\ locally free of rank $n$), we say further that $f$ is a \emph{finite locally free} morphism (resp.\ \emph{locally free of rank $n$}).
This condition, by virtue of \sref{IV.18.2.3}, is satisfied if $f$ is a finite \'etale morphism, but does not imply by itself that $f$ is \'etale, as shown by the example where $X=\operatorname{Spec}(K)$ and $Y=\operatorname{Spec}(k)$ are spectra of fields, $K$ being a finite non-separable extension of $k$.
When $f:X\to Y$ is a finite \'etale morphism, we say further that $X$ is an \emph{\'etale covering} of $Y$.
One notes that in this case, $f$ is universally open and universally closed, and in particular $f(X)$ is a subset that is both open and closed of $Y$.

We say that an \'etale covering $X$ of $Y$ is \emph{trivial} if $X$ is a \emph{sum} of a finite number of preschemes \emph{isomorphic} to $Y$.
We say that an \'etale covering $X$ of $Y$ is \emph{locally trivial} if the morphism $f:X\to Y$ is such that for every point $y\in Y$ there is an open neighbourhood $U$ for which the covering $f^{-1}(U)$ of $U$ is trivial.

  \item[(ii)] Let $f:X\to Y$ be a finite morphism, locally free of rank $n$;
set $\sh{B}=f_*(\sh{O}_X)$ and let $u=\operatorname{astr}_{\sh{B}|\sh{O}_Y}:\sh{B}\to\check{\sh{B}}$;
we deduce from this an $n$-th exterior power homomorphism
$\bigwedge^n u:\bigwedge^n\sh{B}\to\bigwedge^n\check{\sh{B}}=(\bigwedge^n\sh{B})^\vee$ of \emph{invertible} $\sh{O}_Y$-Modules, and by \sref[0]{0.5.4.2} an element
\[
\label{IV.18.2.7.1}
d_{X/Y}\in\Gamma(Y,(\bigwedge^n\check{\sh{B}})\otimes_{\sh{O}_Y}(\bigwedge^n\check{\sh{B}}))
\tag{18.2.7.1}
\]
which we call the \emph{discriminant} of $X$ over $Y$.
Furthermore, since $(\bigwedge^n\check{\sh{B}})\otimes_{\sh{O}_Y}(\bigwedge^n\check{\sh{B}})$ is the dual of $(\bigwedge^n\sh{B})\otimes_{\sh{O}_Y}(\bigwedge^n\sh{B})$, $d_{X/Y}$ can also be identified with a homomorphism
\[
\label{IV.18.2.7.2}
(\bigwedge^n\sh{B})\otimes_{\sh{O}_Y}(\bigwedge^n\sh{B})\to\sh{O}_Y
\tag{18.2.7.2}
\]
and we denote by $\mathscr{D}_{X/Y}$ the quasi-coherent ideal of $\sh{O}_Y$ of finite type, \emph{image} of the homomorphism \sref{IV.18.2.7.2}, also called the \emph{discriminant ideal} of $X$ over $Y$.
\end{enumerate}
\end{remarks}

\oldpage[IV]{114}
For the homomorphism $u$ to be bijective, it is necessary and sufficient that $\bigwedge^n u$ be bijective, or again that the section $d_{X/Y}$ have an \emph{invertible} germ at every point $y\in Y$, which is also written $d_{X/Y}(y)\neq 0$ for every $y$ (\sref[0]{0.5.5.2}).
It amounts to the same to say that the discriminant ideal $\mathscr{D}_{X/Y}$ is \emph{equal} to $\sh{O}_Y$.

The terminology ``covering'' introduced in \sref{IV.18.2.7}, (i), is justified by the following proposition:

\begin{proposition}[18.2.8]
\label{IV.18.2.8}
Let $f:X\to Y$ be a morphism \'etale, separated and of finite type, and for every $y\in Y$, let $n(y)$ be the geometric number of points of $f^{-1}(y)$.
Then the function $y\mapsto n(y)$ is lower semi-continuous on $Y$.
For it to be continuous at a point $y$ (and hence constant in a neighbourhood of $y$), it is necessary and sufficient that there exists an open neighbourhood $U$ of $y$ such that the restriction $f^{-1}(U)\to U$ of $f$ is a \emph{finite} (\'etale) morphism.
\end{proposition}

\begin{proof}
Since $f$ is quasi-finite \sref{IV.17.6.1} and locally of finite presentation, it amounts to saying that $f$ is finite or that $f$ is proper \sref{IV.8.11.1}:
moreover, each fibre $f^{-1}(y)$ is geometrically reduced over $\kappa(y)$.
The conclusions therefore follow from \sref{IV.15.5.9}, (i) and (ii) and from the fact that $f$ is flat.
\end{proof}

\begin{corollary}[18.2.9]
\label{IV.18.2.9}
Let $Y$ be a connected prescheme, $f:X\to Y$ a morphism \'etale, separated and of finite type.
For $f$ to be finite (in other words, for $X$ to be an \'etale covering of $Y$), it is necessary and sufficient that all the fibres of $f$ have the same geometric number of points.
\end{corollary}

\begin{remarks}[18.2.10]
\label{IV.18.2.10}
\begin{enumerate}
  \item[(i)] The example of the ``affine line with a doubled point'' (\sref[I]{I.5.5.11}) shows that a morphism \'etale of finite type of Noetherian preschemes may not be separated;
for this example, the first assertion of \sref{IV.18.2.8} is no longer exact.

  \item[(ii)] For a separated morphism of finite type and \'etale $f:X\to Y$ to make $X$ a \emph{locally trivial} covering, it is necessary and sufficient that for every $x\in X$ there exist an open neighbourhood $U$ of $y=f(x)$ and a $U$-section $g$ of $f^{-1}(U)$ such that $g(y)=x$.
Indeed, the condition is obviously necessary;
the fact that it is also sufficient follows from the fact that every fibre $f^{-1}(y)$ is finite (\sref{IV.17.6.1}) and \sref[I]{I.6.2.2}, from the characterisation of sections of a $Y$-scheme \'etale \sref{IV.17.9.3} and from Proposition~\sref{IV.18.2.8}.
\end{enumerate}
\end{remarks}


\subsection{Finite algebras and \'etale algebras.}
\label{subsection:IV.18.3}

\begin{proposition}[18.3.1]
\label{IV.18.3.1}
Let $A$ be a ring, $B$ an $A$-algebra of finite presentation.
\begin{enumerate}
  \item[(i)] For $B$ to be an unramified $A$-algebra, it is necessary and sufficient that $B$ be a \emph{projective} $(B\otimes_A B)$-module.
  \item[(ii)] Suppose moreover that $B$ is a finite $A$-algebra.  For $B$ to be an \'etale $A$-algebra, it is necessary and sufficient that $B$ be a projective $A$-module and a projective $(B\otimes_A B)$-module.
\end{enumerate}
\end{proposition}

\begin{proof}
\oldpage[IV]{115}
The $(B\otimes_A B)$-module structure on $B$ is that coming from the canonical ring homomorphism $B\otimes_A B\to B$, which is surjective and has kernel $\mathfrak{J}_{B/A}$ (\sref[0]{0.20.4.1}).

(i) To say that the morphism $\operatorname{Spec}(B)\to\operatorname{Spec}(A)$ is locally of finite presentation is equivalent to saying that $B$ is an $A$-algebra of finite presentation \sref{IV.1.4.6}.  To say that $B$ is an unramified $A$-algebra means \sref{IV.17.4.2} that
\[
 \operatorname{Spec}\bigl((B\otimes_A B)/\mathfrak{J}_{B/A}\bigr)
\]
is the subprescheme induced on an open and closed subset of $\operatorname{Spec}(B\otimes_A B)$.  This holds if and only if $\mathfrak{J}_{B/A}$ is a \emph{direct-factor ideal} of $B\otimes_A B$ (Bourbaki, \emph{Alg.\ comm.}, ch.~II, \textsection~4, no.~3, prop.~15);
but this amounts to saying that the $(B\otimes_A B)$-module quotient $(B\otimes_A B)/\mathfrak{J}_{B/A}$ is projective (Bourbaki, \emph{Alg.}, ch.~II, $3^{\text{e}}$ \'ed., \textsection~2, no.~2, prop.~4).

(ii) Recall that a flat $A$-module of finite presentation is projective, and conversely (Bourbaki, \emph{Alg.\ comm.}, ch.~II, \textsection~5, no.~2, cor.~2 du th.~1), so the assertion (ii) follows from that of (i) and from \sref{IV.17.6.2}.
\end{proof}

\begin{proposition}[18.3.2]
\label{IV.18.3.2}
Let $A$ be a ring, $\mathfrak{J}$ an ideal of $A$ such that, for the $\mathfrak{J}$-preadic topology, $A$ is separated and complete;
set $A_0=A/\mathfrak{J}$.
Then the functor
\[
B\mapsto B\otimes_A A_0
\]
is an equivalence of the category of finite \'etale $A$-algebras, to the category of finite \'etale $A_0$-algebras.
\end{proposition}

\begin{proof}
We will first prove the following lemma:

\begin{lemma}[18.3.2.1]
\label{IV.18.3.2.1}
Let $A$ be a ring, $\mathfrak{J}$ an ideal of $A$ such that, for the $\mathfrak{J}$-preadic topology, $A$ is separated and complete.
\begin{enumerate}
  \item[(i)] Every projective $A$-module of finite type $M$ is separated and complete for the $\mathfrak{J}$-preadic topology, and hence is the projective limit of the projective $(A/\mathfrak{J}^{n+1})$-modules $M/\mathfrak{J}^{n+1}M$.
  \item[(ii)] Conversely, set $A_n=A/\mathfrak{J}^{n+1}$, and let $(M_n)$ be a projective system of $A_n$-modules such that, for every $n$, the homomorphism $M_{n+1}\otimes_{A_{n+1}}A_n\to M_n$ induced by the transition homomorphism $M_{n+1}\to M_n$ is bijective.
Suppose moreover that the $M_n$ are projective and $M_0$ of finite type.
Then $M=\varprojlim M_n$ is a projective $A$-module of finite type such that the canonical homomorphism $M\otimes_A A_0\to M_0$ is bijective.
\end{enumerate}
\end{lemma}

\begin{proof}
(i) There is a free $A$-module of finite type $L$ such that $M$ is isomorphic to a direct factor of $L$;
as $L$ is separated for the $\mathfrak{J}$-preadic topology, so is every sub-module $N$ of $L$ since $\mathfrak{J}^{n+1}N\subset\mathfrak{J}^{n+1}L$;
in particular $M$ is separated for this topology, and since the surjective homomorphism $f:L\to M$ is continuous for the $\mathfrak{J}$-preadic topology, its kernel $N$ is closed for the topology induced by that of $L$;
since $L$ is complete and $f$ a strict morphism, we conclude that $M$ is complete (Bourbaki, \emph{Top.\ g\'en.}, ch.~IX, $2^{\text{e}}$ \'ed., \textsection~3, no.~1, prop.~4).

(ii) It follows from Nakayama's lemma that, if $M_0$ is generated by a finite family $(x_{i0})$ of $r$ elements and if, for every $n\geq 1$, $x_{in}$ is an element of $M_n$ whose image in $M_{n-1}$ is $x_{i,n-1}$, then $(x_{in})$ $(1\leq i\leq r)$ is a system of generators of $M_n$ (Bourbaki, \emph{Alg.\ comm.}, ch.~II, \textsection~3, cor.~2 to prop.~4).  For every $n$, put $L_n=A_n^r$.  If $(e_{in})_{1\leq i\leq r}$ is the canonical basis of $L_n$, let $u_n:L_n\to M_n$ be the $A_n$-linear map such that $u_n(e_{in})=x_{in}$ for every $i$.
By hypothesis, we have a split exact sequence
\[
0\to N_n\xrightarrow{v_n} L_n\xrightarrow{u_n} M_n\to 0
\]
\oldpage[IV]{116}
and since $L_n=L_{n+1}/\mathfrak{J}^{n+1}L_{n+1}$ and $M_n=M_{n+1}/\mathfrak{J}^{n+1}M_{n+1}$, the vertical arrows in the commutative diagram
\[
\xymatrix{
0\ar[r] & N_{n+1}\ar[r]^{v_{n+1}}\ar[d] & L_{n+1}\ar[r]^{u_{n+1}}\ar[d] & M_{n+1}\ar[r]\ar[d] & 0\\
0\ar[r] & N_n\ar[r]_{v_n} & L_n\ar[r]_{u_n} & M_n\ar[r] & 0
}
\]
are all three surjective.
Now we have $M=\varprojlim M_n$ and $L=A^r=\varprojlim L_n$;
setting $N=\varprojlim N_n$, $u=\varprojlim u_n$, $v=\varprojlim v_n$, by (EGA $0_{\mathrm{III}}$, 13.2.2) the exact sequence
\[
\label{IV.18.3.2.2}
0\to N\xrightarrow{v} L\xrightarrow{u} M\to 0
\tag{18.3.2.2}
\]
is exact.
Moreover, since for every $n$, $v_n$ is left-invertible and $M_n$ is a projective $A_n$-module, by virtue of (\sref[0]{0.19.1.8}) the exact sequence \sref{IV.18.3.2.2} is split, which proves the lemma.
\end{proof}

This being established, let us first show that the functor of \sref{IV.18.3.2} is \emph{fully faithful}.
Set, as in the lemma, $A_n=A/\mathfrak{J}^{n+1}$;
let $B$, $C$ be two finite \'etale $A$-algebras, and set, for every $n$, $B_n=B\otimes_A A_n$, $C_n=C\otimes_A A_n$;
by virtue of \sref{IV.18.3.1} and \sref{IV.18.3.2.1}, $B$ and $C$ are separated and complete for the $\mathfrak{J}$-preadic topology, and we have $B=\varprojlim B_n$, $C=\varprojlim C_n$;
moreover, every $A$-algebra homomorphism $u:B\to C$ is continuous for the $\mathfrak{J}$-preadic topologies, and hence gives a projective system of $A_n$-algebra homomorphisms $u_n=u\otimes 1:B_n\to C_n$, of which it is the projective limit;
the converse being evident, we have a canonical bijection
\[
\operatorname{Hom}_{A\text{-}\mathrm{alg.}}(B,C)\xrightarrow{\sim}\varprojlim\operatorname{Hom}_{A_n\text{-}\mathrm{alg.}}(B_n,C_n).
\]
But since $B$ and $C$ are finite \'etale $A$-algebras, it follows also from \sref{IV.18.1.2} that the canonical map
\[
\operatorname{Hom}_{A_{n+1}\text{-}\mathrm{alg.}}(B_{n+1},C_{n+1})\to\operatorname{Hom}_{A_n\text{-}\mathrm{alg.}}(B_n,C_n)
\]
is \emph{bijective} for $n\geq 0$, which proves that the canonical map $\operatorname{Hom}_{A\text{-}\mathrm{alg.}}(B,C)\to\operatorname{Hom}_{A_0\text{-}\mathrm{alg.}}(B_0,C_0)$ is bijective.

To complete the proof of \sref{IV.18.3.2}, it suffices to show that for every finite \'etale $A_0$-algebra $B_0$, there exists a finite \'etale $A$-algebra $B$ and an $A_0$-isomorphism $B_0\xrightarrow{\sim}B\otimes_A A_0$.
Now, by \sref{IV.18.1.2}, there is a projective system $(B_n)$ such that $B_n$ is an \'etale $A_n$-algebra and that the homomorphisms $B_{n+1}\otimes_{A_{n+1}}A_n\to B_n$ are bijective.
It follows from \sref{IV.18.3.1} and \sref{IV.18.3.2.1} that the $A$-algebra $B=\varprojlim B_n$ is a projective $A$-module of finite type and that $B_0$ is isomorphic to $B\otimes_A A_0$.
To prove that $B$ is a finite \'etale $A$-algebra, it suffices, by \sref{IV.18.2.5}, to show that for every maximal ideal $\mathfrak{m}$ of $A$, $B_\mathfrak{m}/\mathfrak{m}B_\mathfrak{m}$ is a separable $(A/\mathfrak{m})$-algebra.
Now, since $\mathfrak{J}$ is contained in the radical of $A$ (\sref[0]{0.7.1.10}),
\oldpage[IV]{117}
we have $\mathfrak{J}\subset\mathfrak{m}$, so $\mathfrak{m}_0=\mathfrak{m}/\mathfrak{J}$, we have $A_0/\mathfrak{m}_0=A/\mathfrak{m}$ and $B_\mathfrak{m}/\mathfrak{m}B_\mathfrak{m}=(B_0)_{\mathfrak{m}_0}/\mathfrak{m}_0(B_0)_{\mathfrak{m}_0}$;
the conclusion therefore follows from the fact that $B_0$ is a finite \'etale $A_0$-algebra \sref{IV.18.2.5}.
\end{proof}

\begin{env}[18.3.3]
\label{IV.18.3.3}
\emph{Example.} ---
Proposition \sref{IV.18.3.2} applies in particular when $A$ is a \emph{local} ring, separated and complete, $\mathfrak{J}$ being the maximal ideal of $A$, so that $A_0$ is a field and the finite \'etale $A_0$-algebras are precisely the finite separable $A_0$-algebras, hence the direct products of finite separable field extensions of $A_0$.
In particular, if the field $A_0$ is \emph{separably closed}, these extensions are all identical to $A_0$, and hence every \'etale covering of $\operatorname{Spec}(A)$ is \emph{trivial} \sref{IV.18.2.7} by virtue of \sref{IV.18.3.2}.
\end{env}

\begin{theorem}[18.3.4]
\label{IV.18.3.4}
Let $A$ be a Noetherian ring, $\mathfrak{J}$ an ideal of $A$ such that $A$ is separated and complete for the $\mathfrak{J}$-preadic topology, $A_0=A/\mathfrak{J}$.
Let $S=\operatorname{Spec}(A)$, $S_0=\operatorname{Spec}(A_0)$.
Let $X$ be a proper $S$-scheme, and set $X_0=X\times_S S_0$.
Then the functor
\[
Z\mapsto Z\times_X X_0
\]
from the category of finite \'etale $X$-schemes over $X$, to the category of finite \'etale $X_0$-schemes over $X_0$, is an equivalence of categories.
\end{theorem}

\begin{proof}
Let us first show that this functor is \emph{fully faithful}.
Let $Z'$ and $Z''$ be two finite \'etale $X$-schemes over $X$.
Set $S_n=\operatorname{Spec}(A/\mathfrak{J}^{n+1})$, $X_n=X\times_S S_n$, $Z'_n=Z'\times_S S_n$, $Z''_n=Z''\times_S S_n$.
It follows from (\sref[III]{III.5.4.1}) that we have a canonical bijection $\operatorname{Hom}_X(Z',Z'')\xrightarrow{\sim}\varprojlim\operatorname{Hom}_{X_n}(Z'_n,Z''_n)$.
Now, by virtue of \sref{IV.18.1.2}, the canonical map $\operatorname{Hom}_{X_{n+1}}(Z'_{n+1},Z''_{n+1})\to\operatorname{Hom}_{X_n}(Z'_n,Z''_n)$ is bijective, which proves our assertion.

It remains to show that if $\sh{B}_0$ is a finite \'etale $\sh{O}_{X_0}$-algebra, there exists a finite \'etale $\sh{O}_X$-algebra $\sh{B}$ and an isomorphism $\sh{B}_0\xrightarrow{\sim}\sh{B}\otimes_{\sh{O}_X}\sh{O}_{X_0}$.
It follows from \sref{IV.18.1.2} that there is a projective system $(\sh{B}_n)$, where $\sh{B}_n$ is a finite \'etale $\sh{O}_{X_n}$-algebra, and the second comparison theorem (\sref[III]{III.5.1.4}) proves the existence of a coherent $\sh{O}_X$-module $\sh{B}$ and of a projective system of isomorphisms $\sh{B}_n\xrightarrow{\sim}\sh{B}\otimes_{\sh{O}_X}\sh{O}_{X_n}$.
The datum of an algebra structure on an $\sh{O}_X$-module $\sh{F}$ is equivalent to that of a unit homomorphism $\sh{O}_X\to\sh{F}$ and a multiplication homomorphism $\sh{F}\otimes_{\sh{O}_X}\sh{F}\to\sh{F}$ making certain diagrams involving only tensor powers of $\sh{F}$ commute;
by \sref[III]{III.5.1.3}, \sref[I]{I.10.11.4} and \sref[I]{I.10.11.7}, $\sh{B}$ is therefore endowed in a natural way with a structure of $\sh{O}_X$-algebra for which the isomorphism $\sh{B}_0\xrightarrow{\sim}\sh{B}\otimes_{\sh{O}_X}\sh{O}_{X_0}$ is an isomorphism of algebras.
Moreover, $\sh{B}$ is a \emph{locally free} $\sh{O}_X$-module;
this also follows from \sref[III]{III.5.1.3}, \sref[I]{I.10.11.4}, \sref[I]{I.10.11.7} and from the fact that, in the category of coherent $\sh{O}_X$-modules, the locally free $\sh{O}_X$-modules $\sh{F}$ can be characterized as those for which the functors $\sh{G}\mapsto\shHom_{\sh{O}_X}(\sh{F},\sh{G})$ are \emph{exact}.
Finally, to see that $\sh{B}$ is a finite \'etale $\sh{O}_X$-Algebra, it suffices \sref{IV.18.2.5} to show that for every \emph{closed} point $x\in X$, $\sh{B}_x\otimes_{\sh{O}_{X,x}}\kappa(x)$ is a separable $\kappa(x)$-algebra.
But since the structural morphism $f:X\to S$ is \emph{proper}, $f(x)$ is a closed point of $S$, so it belongs to $S_0$, since $\mathfrak{J}$ is contained in the radical of $A$ (\sref[0]{0.7.1.10});
the conclusion therefore follows from the fact that $X_0=f^{-1}(S_0)$ and that $\sh{B}_0$ is a finite \'etale $\sh{O}_{X_0}$-algebra.
\end{proof}


\subsection{Local structure of unramified morphisms and of \'etale morphisms.}
\label{subsection:IV.18.4}

\oldpage[IV]{118}

\begin{lemma}[18.4.1]
\label{IV.18.4.1}
Let $A$ be a ring, $B$ a finite monogenic $A$-algebra, $u$ a generator of the $A$-algebra $B$, $F\in A[\mathrm{T}]$ a polynomial such that $F(u)=0$, $F'$ the derivative polynomial;
set $u'=F'(u)$.
Then the annihilator ideal of $\Omega^1_{B/A}$ in $B$ contains $u'B$;
it is equal to $u'B$ if the ideal $\mathfrak{J}$ of polynomials $G\in A[\mathrm{T}]$ such that $G(u)=0$ is generated by $F$, in other words if the canonical surjection $\varphi:A[\mathrm{T}]/F A[\mathrm{T}]\to B$, which sends the class of $\mathrm{T}$ to $u$, is bijective.
\end{lemma}

\begin{proof}
Set $C=A[\mathrm{T}]$, so that $B=C/\mathfrak{J}$.
We have the exact sequence (\sref[0]{0.20.5.12.1})
\[
\mathfrak{J}/\mathfrak{J}^2\to\Omega^1_{C/A}\otimes_C B\to\Omega^1_{B/A}\to 0
\]
and $\Omega^1_{B/A}$ therefore identifies with the quotient $B/\mathfrak{J}'$, where $\mathfrak{J}'$ is the ideal generated by the elements $G'(u)$, where $G$ ranges over a system of generators of the ideal $\mathfrak{J}$ (\sref[0]{0.20.5.13});
whence the lemma immediately follows.
\end{proof}

\begin{proposition}[18.4.2]
\label{IV.18.4.2}
Let the notation be those of \sref{IV.18.4.1}, let $\mathfrak{q}$ be a prime ideal of $B$.
Then:
\begin{enumerate}
  \item[(i)] If $\mathfrak{q}$ does not contain $u'$, $B_\mathfrak{q}$ is a formally unramified $A_\mathfrak{p}$-algebra ($\mathfrak{p}$ being the inverse image of $\mathfrak{q}$ in $A$);
in other terms, $\operatorname{Spec}(B_{u'})$ is formally unramified over $\operatorname{Spec}(A)$.
  \item[(ii)] Suppose moreover that $F$ is monic and generates $\mathfrak{J}$.
Then, for $\operatorname{Spec}(B)$ to be \'etale over $\operatorname{Spec}(A)$ at the point $\mathfrak{q}$, it is necessary and sufficient that $u'\notin\mathfrak{q}$.
\end{enumerate}
\end{proposition}

\begin{proof}
The hypothesis that $u'\notin\mathfrak{q}$ implies that $\Omega^1_{B_\mathfrak{q}/A_\mathfrak{p}}=0$ (\sref[0]{0.20.5.9}), hence (i) follows from \sref{IV.17.2.1}.
Moreover, under the hypotheses of (ii), $B$ is a free $A$-module by virtue of the Euclidean division;
as the annihilator $\mathfrak{J}'$ of $\Omega^1_{B/A}$ is then equal to $u'B$ by virtue of \sref{IV.18.4.1} and $\Omega^1_{B/A}$ is a $B$-module of finite presentation \sref{IV.16.4.22}, the annihilator of $\Omega^1_{B_\mathfrak{q}/A_\mathfrak{p}}$ is equal to $u'B_\mathfrak{q}$ (Bourbaki, \emph{Alg.\ comm.}, ch.~II, \textsection~2, no.~4, formule (g)), and (ii) follows from (i) and from the implication $c)\Rightarrow a)$ in \sref{IV.17.6.1}.
\end{proof}

\begin{corollary}[18.4.3]
\label{IV.18.4.3}
With the notation of \sref{IV.18.4.2}, suppose that $F$ is monic and generates $\mathfrak{J}$.
Then, for $B$ to be an \'etale $A$-algebra, it is necessary and sufficient that $u'$ be invertible in $B$ (or, what amounts to the same, that the ideal of $A[\mathrm{T}]$ generated by $F$ and $F'$ be equal to $A[\mathrm{T}]$).
\end{corollary}

\begin{proof}
In view of \sref{IV.18.4.2},~(ii), saying that $\operatorname{Spec}(B)$ is \'etale over $\operatorname{Spec}(A)$ means that $u'$ does not belong to any prime ideal of $B$, i.e.\ that it is invertible in $B$.
\end{proof}

We say that a monic polynomial $F\in A[\mathrm{T}]$ such that the ideal of $A[\mathrm{T}]$ generated by $F$ and $F'$ is equal to $A[\mathrm{T}]$ is \emph{separable};
it is immediate that this definition coincides with the usual definition (Bourbaki, \emph{Alg.}, ch.~V, \textsection~7, no.~6) when $A$ is a field.

\begin{lemma}[18.4.4]
\label{IV.18.4.4}
Let $A$ be a local ring, $B$ a finite monogenic $A$-algebra, $u$ a generator of the $A$-algebra $B$.
Let $\mathfrak{n}_i$ ($1\leq i\leq r$) be the maximal ideals of $B$ such that $\operatorname{Spec}(B)$ is formally unramified over $\operatorname{Spec}(A)$ at the points $\mathfrak{n}_i$.
Then there exists a monic polynomial $F\in A[\mathrm{T}]$ such that $F(u)=0$ and $F'(u)\notin\mathfrak{n}_i$ for every $i$ ($1\leq i\leq r$).
Moreover, if $k$ is the
\oldpage[IV]{119}
residue field of $A$ and $f\in k[\mathrm{T}]$ is the minimal polynomial of the image of $u$ in $B\otimes_A k$, one can choose $F$ so that its canonical image in $k[\mathrm{T}]$ is $f$.
\end{lemma}

\begin{proof}
The maximal ideal $\mathfrak{m}$ of $A$ is the inverse image of each $\mathfrak{n}_i$ (\sref[II]{II.6.1.10}).
Set $L=B\otimes_A k$, a finite $k$-algebra, and let $\xi$ be the image of $u$ in $L$; if $n$ is the rank of $L$ over $k$, the minimal polynomial $f\in k[\mathrm{T}]$ of $\xi$ over $k$ has degree $n$, and $L$ is isomorphic to $k[\mathrm{T}]/f k[\mathrm{T}]$.
If $\mathfrak{n}'_i=\mathfrak{n}_i/\mathfrak{m}B$, the hypothesis implies that $\operatorname{Spec}(L)$ is \'etale over $k$ at the points $\mathfrak{n}'_i$ ($1\leq i\leq r$), by \sref{IV.17.6.1},~c); hence $f'(\xi)\notin\mathfrak{n}'_i$ by \sref{IV.18.4.2},~(ii).
Now $\mathfrak{m}B$ is contained in the radical of $B$; since $1,\xi,\ldots,\xi^{n-1}$ form a basis of $L$ over $k$, Nakayama's lemma shows that $1,u,\ldots,u^{n-1}$ generate the $A$-module $B$.
Consequently there exists a monic polynomial $F\in A[\mathrm{T}]$ of degree $n$ such that $F(u)=0$; since $\xi$ is a root of the canonical image of $F$ in $k[\mathrm{T}]$, this image is necessarily $f$.
The image of $F'(u)$ in $L$ is then $f'(\xi)$, and therefore $F'(u)\notin\mathfrak{n}_i$ for every $i$.
\end{proof}

\begin{proposition}[18.4.5]
\label{IV.18.4.5}
Let $A$ be a local ring, $k$ its residue field, and $B$ a finite $A$-algebra (respectively, a finite $A$-algebra of finite presentation).
Suppose moreover that the residue field $k$ is \emph{infinite}, or else that $B$ is a local ring.
Let $n$ be the rank of $L=B\otimes_A k$ over $k$.
For $B$ to be a formally unramified $A$-algebra (resp.\ \'etale), it is necessary and sufficient that there exist a monic separable polynomial $F\in A[\mathrm{T}]$ such that $B$ is isomorphic to a quotient of $A[\mathrm{T}]/F.A[\mathrm{T}]$ (resp.\ isomorphic to $A[\mathrm{T}]/F.A[\mathrm{T}]$).
Moreover, one can suppose that $F$ has degree $n$ (respectively, $F$ necessarily has degree $n$).
\end{proposition}

\begin{proof}
The conditions are sufficient by virtue of \sref{IV.18.4.2}, without assuming either that $k$ is infinite or that $B$ is local.
To see that the conditions are necessary, note that if $B$ is a formally unramified $A$-algebra, $L$ is a finite separable $k$-algebra, hence a direct product of finitely many finite separable field extensions $k_j$ of $k$ ($1\leq j\leq r$), each generated by an element $\bar{\xi}_j$ with minimal polynomial $f_j$ (Bourbaki, \emph{Alg.}, ch.~V, \textsection~11, no.~4, prop.~4).
Let us show that by virtue of the hypotheses made on $k$ or $B$, there exists an element $\bar{\xi}$ of $L$ generating the $k$-algebra $L$.
This is immediate if $B$ is local, since then $r=1$.
Otherwise, $k$ being assumed infinite, one can suppose that the irreducible polynomials $f_j\in k[\mathrm{T}]$ are all \emph{distinct} (by replacing each $\bar{\xi}_j$ by $\bar{\xi}_j+a_j$, for a suitable element $a_j\in k$);
if we set $f=f_1 f_2\ldots f_r$, it is clear that $L$ is isomorphic to $k[\mathrm{T}]/f k[\mathrm{T}]$ in both cases, hence generated by an element $\bar{\xi}$ of minimal polynomial $f\in k[\mathrm{T}]$ of degree $n$.
Choose $u\in B$ whose image in $L$ is $\bar{\xi}$.
Since $\mathfrak{m}B$ is contained in the radical of $B$ and $1,\bar{\xi},\ldots,\bar{\xi}^{n-1}$ form a basis of $L$ over $k$, Nakayama's lemma shows that $1,u,\ldots,u^{n-1}$ generate the $A$-module $B$; hence there exists a monic polynomial $F\in A[\mathrm{T}]$ of degree $n$ such that $F(u)=0$;
moreover, since $\bar{\xi}$ is a root of the canonical image of $F$ in $k[\mathrm{T}]$, this image is necessarily equal to $f$.
Thus $u$ generates the $A$-algebra $B$, which is consequently a quotient of $A[\mathrm{T}]/F A[\mathrm{T}]$.
Moreover, $B$ is semilocal and, at each of its maximal ideals $\mathfrak{n}_i$, \sref{IV.18.4.4} gives $F'(u)\notin\mathfrak{n}_i$; hence $F'(u)$ is invertible in $B$, and $F$ is separable.
Finally, if $B$ is an \'etale $A$-algebra, then $B$, being a flat $A$-module of finite presentation (\sref{IV.1.4.7}), is free, and $1,u,\ldots,u^{n-1}$ form a basis of the $A$-module $B$ (Bourbaki, \emph{Alg.\ comm.}, ch.~II, \textsection~3, no.~2, prop.~5); in other words,
\oldpage[IV]{120}
the $A$-algebra $B$ is isomorphic to $A[\mathrm{T}]/F A[\mathrm{T}]$; and any other monic polynomial $G\in A[\mathrm{T}]$ such that $B\simeq A[\mathrm{T}]/G A[\mathrm{T}]$ necessarily has degree $n$.
\end{proof}

\begin{theorem}[18.4.6]
\label{IV.18.4.6}
(Chevalley). ---
\begin{enumerate}
  \item[(i)] Let $f:X\to Y$ be a morphism locally of finite type, $x$ a point of $X$, $y=f(x)$, and set $A=\sh{O}_{Y,y}$.
For $\sh{O}_{X,x}$ to be a formally \'etale (respectively, formally unramified) $A$-algebra, it is necessary and sufficient that there exist a monic polynomial $F\in A[\mathrm{T}]$ and a maximal ideal $\mathfrak{n}$ of $B=A[\mathrm{T}]/F A[\mathrm{T}]$ (respectively, of a quotient algebra $B$ of $A[\mathrm{T}]/F A[\mathrm{T}]$) such that $\sh{O}_{X,x}$ is $A$-isomorphic to $B_\mathfrak{n}$ and, if $u$ is the image of $\mathrm{T}$ in $B$, one has $F'(u)\notin\mathfrak{n}$.
  \item[(ii)] Suppose in addition that $f$ is locally of finite presentation.
Then, for $f$ to be \'etale at the point $x$, it is necessary and sufficient that we can moreover take $B=A[\mathrm{T}]/F.A[\mathrm{T}]$.
\end{enumerate}
\end{theorem}

\begin{proof}
The conditions are sufficient by virtue of \sref{IV.18.4.2}.
To see that they are necessary, we may plainly reduce to the case where $X$ and $Y$ are affine and, in view of \sref{IV.17.4.1.2}, to the case where $f$ is formally unramified and quasi-finite.
It then follows from \sref{IV.8.12.8} that there exist a finite $A$-algebra $C$ and a maximal ideal $\mathfrak{r}$ of $C$ (necessarily lying above the maximal ideal $\mathfrak{m}$ of $A$) such that $\sh{O}_{X,x}$ is $A$-isomorphic to $C_\mathfrak{r}$.
Moreover, by \sref{IV.17.4.1.2}, the residue field $C/\mathfrak{r}=C_\mathfrak{r}/\mathfrak{r}C_\mathfrak{r}$ is a finite separable extension of $k=A/\mathfrak{m}$, and hence is of the form $k[v]$, where $v$ is separable over $k$.
Let $\mathfrak{r}_i$ ($1\leq i\leq h$) be the maximal ideals of the semilocal ring $C$ other than $\mathfrak{r}$.
There exists an element $u\in C$ belonging to every $\mathfrak{r}_i$ and whose image in $C/\mathfrak{r}$ is $v$ (Bourbaki, \emph{Alg.\ comm.}, ch.~II, \textsection~1, no.~2, prop.~5).
We shall show that the $A$-subalgebra $B=A[u]$ of $C$ and the ideal $\mathfrak{n}=\mathfrak{r}\cap B$ (necessarily maximal, since $C$ is finite over $B$) answer the question.

For the case where $\sh{O}_{X,x}$ is a formally unramified $A$-algebra, it suffices to prove that $B_\mathfrak{n}$ is isomorphic to $C_\mathfrak{r}$; indeed, $B_\mathfrak{n}$ will then be formally unramified over $A$, and the existence of a polynomial $F\in A[\mathrm{T}]$ having the asserted properties will follow from \sref{IV.18.4.4}.
Since $f$ is formally unramified, $C/\mathfrak{r}$ is isomorphic to the $A$-algebra $C_\mathfrak{r}/\mathfrak{m}C_\mathfrak{r}$ (\sref{IV.17.4.1.2}).
We are thus reduced to proving the following lemma.

\begin{lemma}[18.4.6.1]
\label{IV.18.4.6.1}
Let $A$ be a local ring with maximal ideal $\mathfrak{m}$, let $C$ be a finite $A$-algebra, and let $\mathfrak{r}$ be a maximal ideal of $C$.
Let $u\in C$ belong to every maximal ideal of $C$ distinct from $\mathfrak{r}$, but not to $\mathfrak{r}$, and suppose that $C_\mathfrak{r}/\mathfrak{m}C_\mathfrak{r}$ is a monogenic algebra over $k=A/\mathfrak{m}$, generated by the image of $u$.
Set $B=A[u]$ and $\mathfrak{n}=\mathfrak{r}\cap B$.
Then the canonical homomorphism $B_\mathfrak{n}\to C_\mathfrak{r}$ is an isomorphism.
\end{lemma}

\begin{proof}
Set $R=B\setminus\mathfrak{n}$ and $S=C\setminus\mathfrak{r}$, so that $B_\mathfrak{n}=R^{-1}B$ and $C_\mathfrak{r}=S^{-1}C$; the canonical homomorphism $B_\mathfrak{n}\to C_\mathfrak{r}$ factors as
\[
R^{-1}B\xrightarrow{g}R^{-1}C\xrightarrow{h}S^{-1}C.
\]
It suffices to prove that both homomorphisms are bijective.
First, $h:R^{-1}C\to S^{-1}C=C_\mathfrak{r}$ is bijective.
Indeed, it suffices to see that the images in $R^{-1}C$ of the elements of $S$ are invertible, or equivalently that every maximal ideal $\mathfrak{p}$ of $R^{-1}C$ has inverse image in $C$ disjoint from $S$, and hence equal to $\mathfrak{r}$.
\oldpage[IV]{121}
Since $R^{-1}C$ is finite over the local ring $R^{-1}B=B_\mathfrak{n}$, the inverse image of $\mathfrak{p}$ in $R^{-1}B$ is its unique maximal ideal $\mathfrak{n}B_\mathfrak{n}$; if $\mathfrak{q}$ is the inverse image of $\mathfrak{p}$ in $C$, then $\mathfrak{q}\cap B=\mathfrak{n}$ and, since $\mathfrak{n}$ is maximal in $B$ and $C$ is finite over $B$, $\mathfrak{q}$ is one of the maximal ideals of $C$.
Furthermore $u\notin\mathfrak{q}$, since $u\notin\mathfrak{r}$ and $u\in B$; the hypothesis therefore forces $\mathfrak{q}=\mathfrak{r}$.

Since $B\subset C$, the homomorphism $g:R^{-1}B\to R^{-1}C$ is injective (EGA $0_{\mathrm{I}}$, 1.3.2).
To prove that it is surjective, observe that $R^{-1}C$ is an $R^{-1}B$-module of finite type and that $\mathfrak{m}R^{-1}B$ is contained in the maximal ideal of the local ring $R^{-1}B$.
By Nakayama's lemma, it suffices to prove that
\[
R^{-1}B/\mathfrak{m}R^{-1}B\longrightarrow R^{-1}C/\mathfrak{m}R^{-1}C
\]
is surjective.
But, by the first part of the proof, $R^{-1}C/\mathfrak{m}R^{-1}C$ identifies with $C_\mathfrak{r}/\mathfrak{m}C_\mathfrak{r}$; by hypothesis this $k$-algebra is generated by the image of $u$, and hence it is the image of $R^{-1}B/\mathfrak{m}R^{-1}B$.
\end{proof}

Consider next the case where $f$ is \'etale at $x$.
Replacing $X$ by a neighbourhood of $x$, we may suppose that $X$ is a neighbourhood of $\mathfrak{n}$ in $\operatorname{Spec}(B)$ (\sref[I]{I.7.2}).
Set $B'=A[\mathrm{T}]/F A[\mathrm{T}]$ and let $\mathfrak{n}'$ be the inverse image of $\mathfrak{n}$ in $B'$.
Since the image of $F'(\mathrm{T})$ in $B'$ does not belong to $\mathfrak{n}'$, the morphism $\operatorname{Spec}(B')\to\operatorname{Spec}(A)$ is \'etale at $\mathfrak{n}'$ by \sref{IV.18.4.2},~(ii).
As $\operatorname{Spec}(B)\to\operatorname{Spec}(A)$ is \'etale at $\mathfrak{n}$, \sref{IV.17.3.4} shows that $\operatorname{Spec}(B)\to\operatorname{Spec}(B')$ is \'etale at $\mathfrak{n}$; but this morphism is an immersion, so it can be \'etale at a point only if it is a local isomorphism there (\sref{IV.17.9.1}).
Thus $B_\mathfrak{n}$ and $B'_{\mathfrak{n}'}$ are isomorphic.

Finally, suppose that $\sh{O}_{X,x}$ is a formally \'etale $A$-algebra.
With the preceding notation, \sref{IV.17.1.5} shows that $B_\mathfrak{n}$ is a formally \'etale $B'_{\mathfrak{n}'}$-algebra; since the homomorphism $B'_{\mathfrak{n}'}\to B_\mathfrak{n}$ is surjective, this is possible only if it is \emph{bijective} (\sref[0]{0.19.10.3},~(i)).
This completes the proof of \sref{IV.18.4.6}.
\end{proof}

\begin{corollary}[18.4.7]
\label{IV.18.4.7}
Let $f:X\to Y$ be a morphism locally of finite type, $x$ a point of $X$.
For $f$ to be formally unramified at the point $x$, it is necessary and sufficient that there exist an open neighbourhood $U$ of $x$ such that $f|U$ factorises as $U\xrightarrow{j}X'\xrightarrow{h}Y$, where $h$ is an \'etale morphism and $j$ is a closed immersion.
\end{corollary}

\begin{proof}
We can obviously reduce to the case where $Y=\operatorname{Spec}(R)$ is affine and $f$ is of finite type.
If $A=\sh{O}_{Y,y}$, with $y=f(x)$, the condition for $f$ to be formally unramified at the point $x$ is equivalent, by virtue of \sref{IV.17.4.1.2}, to saying that $\sh{O}_{X,x}$ is a formally unramified $A$-algebra.
If it is so, we can apply \sref{IV.18.4.6}, (i);
replacing $Y$ as needed by an affine neighbourhood of $y$, we may suppose (with the notation of \sref{IV.18.4.6}) that the polynomial $F$ is the image in $A[\mathrm{T}]$ of a monic polynomial $G\in R[\mathrm{T}]$.
We then set $X'=\operatorname{Spec}(R[\mathrm{T}]/G.R[\mathrm{T}])$;
let $x'$ be the image of the point $\mathfrak{n}$ of $\operatorname{Spec}(B)$ by the morphism corresponding to the composite homomorphism $R[\mathrm{T}]/G.R[\mathrm{T}]\to A[\mathrm{T}]/F.A[\mathrm{T}]\to B$.
It follows from \sref{IV.18.4.2} that the canonical morphism $h:X'\to Y$ is \'etale at the point $x'$, hence, restricting $X'$ and $Y$ to open neighbourhoods of $x'$ and $y$ if necessary, we may suppose that $h$ is \'etale.
On the other hand, by virtue of \sref[I]{I.6.5.1},~(ii) and \sref[I]{I.7.2}, the local homomorphism $\varphi:\sh{O}_{X',x'}\to\sh{O}_{X,x}$ corresponds to a morphism
\oldpage[IV]{122}
$j:U\to X'$ from an open neighbourhood $U$ of $x$ in $X$;
in applying \sref[I]{I.6.5.1}, (i), suppose that $h\circ j=f|U$, hence $j$ is of finite type (\sref[I]{I.6.3.4}, (v));
finally, the homomorphism $\varphi$ is surjective by \sref{IV.18.4.6},~(i), hence it follows from \sref[I]{I.6.5.4}, (i), that we can, restricting further $U$ and $X'$, suppose that $j$ is a closed immersion.
This proves the necessity of the condition stated;
the sufficiency is immediate \sref{IV.17.1.3}, (i) and (ii).
\end{proof}

\begin{corollary}[18.4.8]
\label{IV.18.4.8}
Let $f:X\to Y$ be a morphism locally of finite presentation.
For $f$ to be unramified (respectively, \'etale), it is necessary and sufficient that there exist a family of flat morphisms $g_\alpha:Y'_\alpha\to Y$ and, for each $\alpha$, an open subset $U_\alpha$ of $X'_\alpha=X\times_Y Y'_\alpha$, such that, if $g'_\alpha:X'_\alpha\to X$ and $f'_\alpha:X'_\alpha\to Y'_\alpha$ are the canonical projections, the $g'_\alpha(U_\alpha)$ cover $X$, and each composite $U_\alpha\to X'_\alpha\xrightarrow{f'_\alpha}Y'_\alpha$ is a closed immersion (respectively, an open immersion).
One can moreover take the $g_\alpha$ to be \'etale.
\end{corollary}

\begin{proof}
The necessity of the condition for \'etale morphisms is immediate: take a single $Y'_\alpha=X$, let $g_\alpha=f$, and take for $U_\alpha\subset X\times_Y X$ the diagonal.
When $f$ is unramified, necessity follows from \sref{IV.18.4.7}: choose an open covering $(V_\alpha)$ of $X$ such that, for each $\alpha$, $f|V_\alpha$ factors as $V_\alpha\xrightarrow{j_\alpha}Y'_\alpha\xrightarrow{g_\alpha}Y$, where $j_\alpha$ is a closed immersion and $g_\alpha$ is \'etale.
Then $j_\alpha:V_\alpha\to Y'_\alpha$ factors as $V_\alpha\xrightarrow{s_\alpha}X'_\alpha\xrightarrow{f'_\alpha}Y'_\alpha$, where $s_\alpha$ is a $V_\alpha$-section of $X'_\alpha$; since $g'_\alpha:X'_\alpha\to X$ is \'etale, $s_\alpha$ is an open immersion (\sref{IV.17.4.1}), and it suffices to take $U_\alpha=s_\alpha(V_\alpha)$.

The sufficiency of the conditions follows from \sref{IV.17.7.1};
one concludes at each point $x$ of $g'_\alpha(U_\alpha)$, hence in all of $X$ since the $g'_\alpha(U_\alpha)$ cover $X$.
\end{proof}

\begin{proposition}[18.4.9]
\label{IV.18.4.9}
Let $S$ be a prescheme, $f:X\to S$ a morphism, $h:Y\to S$ a morphism locally of finite presentation, $g:X\to Y$ an $S$-morphism, $x$ a point of $X$, and $y=g(x)$.
The following conditions are equivalent:
\begin{enumerate}[label=\emph{\alph*)}]
  \item $h$ is \'etale at the point $y$ and $g$ is flat at the point $x$.
  \item $h$ is unramified at the point $y$ and $f$ is flat at the point $x$.
\end{enumerate}
\end{proposition}

\begin{proof}
Since $f=h\circ g$, condition $a)$ implies that $f$ is flat at $x$ (\sref{IV.2.1.6}); moreover, an \'etale morphism is unramified, so $a)$ implies $b)$.

To prove that $b)$ implies $a)$, we may first suppose that $h$ is unramified, after replacing $Y$ by an open neighbourhood of $y$.
Then, after replacing $S$ by an open neighbourhood of $s=h(y)=f(x)$, we may suppose that there exist an \'etale morphism $u:S'\to S$, a point $y'$ of $Y'=Y_{(S')}$ above $y$, and an open neighbourhood $V$ of $y'$ in $Y'$ such that, if $h'=h_{(S')}:Y'\to S'$, the restriction $h'|V$ is a closed immersion (\sref{IV.18.4.8}).
If one proves that $h'$ is \'etale at the point $y'$, it will follow that $h$ is \'etale at the point $y$ \sref{IV.17.7.1}, (ii);
moreover, $f'=f_{(S')}:X_{(S')}\to S'$ is flat at every point $x'$ above $x$.
As the projections $v:Y'\to Y$, $w:X'\to X$ are \'etale morphisms (hence flat), if one proves that $g'=g_{(S')}:X_{(S')}\to Y_{(S')}$ is flat at the point $x'$, it will follow from \sref{IV.2.2.11}, (iv) that $g$ is flat at the point $x$.
We may therefore reduce to the case where $h$ is a \emph{closed immersion}
\oldpage[IV]{123}
of finite presentation, with $f$ flat at $x$.
Let $\sh{J}$ be the quasi-coherent ideal of finite type in $\sh{O}_S$ defining the closed subprescheme $Y$ of $S$.
The hypothesis that $f$ is flat at $x$ implies that the homomorphism $\sh{O}_{S,s}\to\sh{O}_{X,x}$ is injective (EGA $0_{\mathrm{I}}$, 6.5.1);
\emph{a fortiori} the homomorphism $\sh{O}_{S,s}\to\sh{O}_{Y,y}$ is injective, which means that $\sh{J}_s=0$, since it is the kernel of the preceding homomorphism.
Since $\sh{J}$ is of finite type, there exists an open neighbourhood $U$ of $s$ in $S$ such that $\sh{J}|U=0$ (\sref[0]{0.5.2.2}).
We can hence suppose that $h$ is an open immersion, and it is clear that $g$ is flat at the point $x$.
\end{proof}

\begin{proposition}[18.4.10]
\label{IV.18.4.10}
Let us denote by $\mathbf{P}(f,x)$ a property satisfying the following conditions:
\begin{enumerate}
  \item[$1^\circ$] For every morphism $f:X\to Y$ and every local isomorphism $h:Y\to Z$, $\mathbf{P}(f,x)$ is equivalent to $\mathbf{P}(h\circ f,x)$ for $x\in X$.
  \item[$2^\circ$] For every morphism $f:X\to Y$, every \'etale morphism $g:Y'\to Y$, every point $x\in X$, if we set $X'=X\times_Y Y'$, $f'=f_{(Y')}:X'\to Y'$ and if $x'\in X'$ is above $x$, the properties $\mathbf{P}(f,x)$ and $\mathbf{P}(f',x')$ are equivalent (``invariance by \'etale base change'').
\end{enumerate}
Let then $S$ be a prescheme, $f:X\to S$ and $h:Y\to S$ two morphisms, $g:X\to Y$ an $S$-morphism, $x$ a point of $X$, $y=g(x)$, and suppose $h$ \'etale at the point $y$.
Then the properties $\mathbf{P}(f,x)$ and $\mathbf{P}(g,x)$ are equivalent.
\end{proposition}

\begin{proof}
With the same notation as in the proof of \sref{IV.18.4.9}, and after replacing $S$ (respectively, $Y$) by a neighbourhood of $s=h(y)$ (respectively, of $y$), \sref{IV.18.4.8} gives an \'etale morphism $u:S'\to S$ such that $h'$ is an open immersion.
If $x'\in X'$ lies above $x$, the hypothesis makes $\mathbf{P}(f',x')$ (respectively, $\mathbf{P}(g',x')$) equivalent to $\mathbf{P}(f,x)$ (respectively, $\mathbf{P}(g,x)$).
We may therefore reduce to the case where $h$ is an open immersion; condition $1^\circ$ then makes $\mathbf{P}(g,x)$ equivalent to $\mathbf{P}(h\circ g,x)=\mathbf{P}(f,x)$.
\end{proof}

\begin{env}[18.4.11]
\label{IV.18.4.11}
\emph{Examples.} ---
One can take for the property $\mathbf{P}(f,x)$ any one of the following, taking into account \sref{IV.17.7.4},~(ii):
\begin{enumerate}
  \item[(i)] $f$ is flat at the point $x$ \sref{IV.2.2.11}, (iv);
  \item[(ii)] $f$ is locally of finite presentation and has codepth $\leq n$ at the point $x$ (\sref{IV.6.8.1} and \sref{IV.6.7.8});
  \item[(iii)] $f$ is locally of finite presentation and Cohen--Macaulay at the point $x$ (\sref{IV.6.8.1} and \sref{IV.6.7.8});
  \item[(iv)] $f$ is locally of finite presentation and possesses the property $(S_n)$ at the point $x$ (\sref{IV.6.8.1} and \sref{IV.6.7.8});
  \item[(v)] $f$ is locally of finite presentation and possesses the property $(R_n)$ at the point $x$ (\sref{IV.6.8.1} and \sref{IV.6.7.8});
  \item[(vi)] $f$ is locally of finite presentation and normal at the point $x$ (\sref{IV.6.8.1} and \sref{IV.6.7.8});
  \item[(vii)] $f$ is locally of finite presentation and reduced at the point $x$ (\sref{IV.6.8.1} and \sref{IV.6.7.8});
  \item[(viii)] $f$ is unramified at the point $x$ \sref{IV.17.7.4};
  \item[(ix)] $f$ is smooth at the point $x$ \sref{IV.17.7.4};
  \item[(x)] $f$ is \'etale at the point $x$ \sref{IV.17.7.4}.
\end{enumerate}
\end{env}

\oldpage[IV]{124}
\begin{corollary}[18.4.12]
\label{IV.18.4.12}
\begin{enumerate}
  \item[(i)] Let $f:X\to Y$ be a morphism locally of finite type, $x$ a point of $X$.
If $f$ is flat and formally unramified at the point $x$, then, in every factorisation $f|U:U\xrightarrow{j}X'\xrightarrow{h}Y$, where $U$ is an open neighbourhood of $x$, $j$ a closed immersion and $h$ an \'etale morphism, the homomorphism $\sh{O}_{X',j(x)}\to\sh{O}_{X,x}$ corresponding to $j$ is bijective (in particular $\sh{O}_{X,x}$ is an essentially \'etale $\sh{O}_{Y,f(x)}$-algebra \sref{IV.18.6.1}).
  \item[(ii)] For a morphism $f:X\to Y$ to be \'etale, it is necessary and sufficient that it be locally of finite type, formally unramified and flat.
\end{enumerate}
\end{corollary}

\begin{proof}
(i) Since the homomorphism $\sh{O}_{X',j(x)}\to\sh{O}_{X,x}$ is surjective, it suffices to prove that it is injective, and for that it suffices to prove that $\sh{O}_{X,x}$ is an $\sh{O}_{X',j(x)}$-module \emph{faithfully flat}, or merely \emph{flat} (\sref[0]{0.6.5.1} and \sref{IV.6.6.2});
in other words, it amounts to showing that $j$ is a flat morphism at the point $x$;
but since $h\circ j=f$ is flat at the point $x$ by hypothesis and $h$ is \'etale, this follows from \sref{IV.18.4.10} and \sref{IV.18.4.11}, (i).

(ii) It remains only to prove the sufficiency of the stated conditions.
For every $x\in X$, there is an open neighbourhood $U$ of $x$ on which $f|U$ admits a factorization having the properties in (i).
Since $f$ is flat and formally unramified at \emph{every} point of $U$, part (i) applies not only at $x$ but at every $z\in U$.
Thus, if $\sh{J}$ is the quasi-coherent ideal of $\sh{O}_{X'}$ defining the closed subprescheme $U$ associated with $j$, then $\sh{J}_{j(z)}=0$ for every $z\in U$; consequently $j$ is an open immersion, and $f$ is \'etale at every point of $U$, hence at every point of $X$.
\end{proof}

\begin{corollary}[18.4.13]
\label{IV.18.4.13}
Let $f:X\to Y$ be a morphism locally of finite type, $x$ a point of $X$, $y=f(x)$.
Suppose that $y$ admits an open neighbourhood which is a reduced prescheme and has only a finite number of irreducible components.
Then, for $f$ to be \'etale at the point $x$, it is necessary and sufficient that $f$ be flat and formally unramified at the point $x$.
\end{corollary}

\begin{proof}
There is only the sufficiency of the condition to prove.
The question being local on $X$ and $Y$, one can suppose \sref{IV.18.4.7} that $f$ factorises as $X\xrightarrow{j}X'\xrightarrow{h}Y$ where $h$ is \'etale and $j$ a closed immersion.
Then $X'$ is reduced \sref{IV.17.5.7} and, replacing $X'$ by an open neighbourhood of $j(x)$, one can suppose that $X'$ has only a finite number of irreducible components.
Then the maximal points of $X'$ are above the maximal points of $Y$ \sref{IV.2.3.4} and since they are finite, their number is finite.
It all amounts to showing, with the notations of the proof of \sref{IV.18.4.12}, that $\sh{J}_{x'}=0$ for all points $x'$ in a neighbourhood of $j(x)$ in $X'$, knowing that $\sh{J}_{j(x)}=0$.
Now, replacing $X'$ by an affine neighbourhood of $j(x)$, we can suppose that all the irreducible components of $X'$ contain $j(x)$;
if $X'=\operatorname{Spec}(A')$, and if $\mathfrak{p}'$ is the prime ideal of $A'$ corresponding to the point $j(x)$, the morphism $\operatorname{Spec}(A'_{\mathfrak{p}'})\to\operatorname{Spec}(A')$ is dominant, hence the corresponding homomorphism $A'\to A'_{\mathfrak{p}'}$ is \emph{injective} since $A'$ is reduced \sref[I]{I.1.2.7}.
If $\sh{J}=\widetilde{\mathfrak{J}}$, where $\mathfrak{J}$ is an ideal of $A'$, then $\mathfrak{J}$ identifies with a submodule of $A'_{\mathfrak{p}'}$; since $\mathfrak{J}_{\mathfrak{p}'}=0$ by hypothesis, it follows that $\mathfrak{J}=0$.
\end{proof}

\begin{corollary}[18.4.14]
\label{IV.18.4.14}
Let $A$ be a local ring with maximal ideal $\mathfrak{m}$, residue field $k$, $B$ a finite $A$-algebra.
\oldpage[IV]{125}
\begin{enumerate}
  \item[(i)] For $B$ to be a formally unramified $A$-algebra, it is necessary and sufficient that $B\otimes_A k$ be a $k$-algebra \'etale.
  \item[(ii)] For $B$ to be a finite \'etale $A$-algebra, it is necessary and sufficient that $B\otimes_A k$ be an \'etale $k$-algebra and that $B$ be a flat $A$-module (which is equivalent to saying (\sref[0]{0.6.3.3}) that for every maximal ideal $\mathfrak{n}$ of $B$ (necessarily above $\mathfrak{m}$), $B_\mathfrak{n}$ is a flat $A$-module).
\end{enumerate}
\end{corollary}

\begin{proof}
Only sufficiency needs proof.
It is clear that (ii) follows from (i) and \sref{IV.18.4.12},~(ii), taking into account \sref{IV.17.1.2},~(i).
To prove (i), let us note that if $B\otimes_A k$ is an \'etale $k$-algebra, we have $\Omega^1_{B\otimes_A k/k}=0$ \sref{IV.17.2.1}.
Since $\Omega^1_{B\otimes_A k/k}=\Omega^1_{B/A}\otimes_B(B\otimes_A k)$ (\sref[0]{0.20.5.5}) and since $B$ is an $A$-algebra of finite type, $\Omega^1_{B/A}$ is a $B$-module of finite type (\sref[0]{0.20.4.7}).
But since $B$ is a finite $A$-algebra, $\mathfrak{m}B$ is contained in the radical of $B$ (Bourbaki, \emph{Alg.\ comm.}, ch.~V, \textsection~2, no.~1, prop.~1), hence the lemma of Nakayama proves that $\Omega^1_{B/A}=0$, and hence $B$ is a formally unramified $A$-algebra \sref{IV.17.2.1}.
\end{proof}


\subsection{Henselian local rings.}
\label{subsection:IV.18.5}
\footnote{The notion of a Henselian local ring is due to Azumaya, and that of Henselization to Nagata, to whom the principal results of this theory are also due.}

\begin{env}[18.5.1]
\label{IV.18.5.1}
Let $X$ be a prescheme and $\sh{E}$ a locally free $\sh{O}_X$-module of finite rank.
Its dual $\check{\sh{E}}=\shHom_{\sh{O}_X}(\sh{E},\sh{O}_X)$ is therefore a locally free $\sh{O}_X$-module whose rank at every point equals that of $\sh{E}$ there, and the canonical homomorphism $\sh{E}\to\shHom_{\sh{O}_X}(\check{\sh{E}},\sh{O}_X)=(\check{\sh{E}})^\vee$ is \emph{bijective}.
For every morphism $X'\to X$, set $\sh{E}_{(X')}=\sh{E}\otimes_{\sh{O}_X}\sh{O}_{X'}$, which is a locally free $\sh{O}_{X'}$-module, and consider the \emph{set} $\Gamma(X',\sh{E}_{(X')})$ of its sections over $X'$.
We shall see that this defines a \emph{contravariant representable functor}
\[
\label{IV.18.5.1.1}
{}^t\! V_{\sh{E}}:X'\mapsto\Gamma(X',\sh{E}_{(X')})
\tag{18.5.1.1}
\]
from the category of $X$-preschemes to the category of \emph{sets} (EGA $0_{\mathrm{III}}$, 8.1.8).

First of all, this is indeed a well-defined functor, for if $f:X''\to X'$ is an $X$-morphism, we have $\sh{E}_{(X'')}=f^*(\sh{E}_{(X')})$, whence (\sref[0]{0.4.3.2}) a map $\Gamma(X',\sh{E}_{(X')})\to\Gamma(X'',\sh{E}_{(X'')})$ which completes the definition of the functor ${}^t\! V$.
Let us then show that the $X$-prescheme $\mathbf{V}(\check{\sh{E}})$ (\sref[II]{II.1.7.8}) represents the functor ${}^t\! V$.
It is immediate that $(\sh{E}_{(X')})^\vee=(\check{\sh{E}})_{(X')}$, hence $\mathbf{V}(\check{\sh{E}})\times_X X'=\mathbf{V}((\check{\sh{E}})_{(X')})$.
Taking \sref[I]{I.3.3.14} into account, it suffices to define a bijection
\[
\Gamma(\mathbf{V}(\check{\sh{E}})/X)\xrightarrow{\sim}\Gamma(X,\sh{E})
\]
and to verify that the corresponding bijection $\Gamma(\mathbf{V}((\check{\sh{E}})_{(X')})/X')\xrightarrow{\sim}\Gamma(X',\sh{E}_{(X')})$ is functorial in $X'$.
There is a canonical bijection from $\Gamma(\mathbf{V}(\check{\sh{E}})/X)$ to $\operatorname{Hom}_{\sh{O}_X}(\check{\sh{E}},\sh{O}_X)$ (\sref[II]{II.1.7.8}), and transposition $u\mapsto{}^t u$ gives a canonical bijection from $\operatorname{Hom}_{\sh{O}_X}(\check{\sh{E}},\sh{O}_X)$ to $\operatorname{Hom}_{\sh{O}_X}(\sh{O}_X,\sh{E})=\Gamma(X,\sh{E})$, because $(\check{\sh{E}})^\vee$ is identified with $\sh{E}$.
The functoriality is immediate.

\oldpage[IV]{126}
Observe also that, in accordance with the general theory (EGA $0_{\mathrm{III}}$, 8.1.6), the identity automorphism of $\mathbf{V}(\check{\sh{E}})$ corresponds canonically to a section $c$ of $\sh{E}_{(\mathbf{V}(\check{\sh{E}}))}$ over $\mathbf{V}(\check{\sh{E}})$.
Equivalently (\sref[II]{II.1.4.1}), it corresponds to a homomorphism of $\mathbf{S}_{\sh{O}_X}(\check{\sh{E}})$-modules
\[
u:\mathbf{S}_{\sh{O}_X}(\check{\sh{E}})\longrightarrow
\mathbf{S}_{\sh{O}_X}(\check{\sh{E}})\otimes_{\sh{O}_X}\sh{E}.
\]
For every affine open $W\subset X$, if $\Gamma(W,\sh{O}_X)=A$ and $\Gamma(W,\mathbf{S}_{\sh{O}_X}(\check{\sh{E}}))$ is identified with $C=A[\mathrm{T}_1,\ldots,\mathrm{T}_n]$ so that the $\mathrm{T}_i$ form a basis of $\Gamma(W,\check{\sh{E}})$, and if $(e_i)$ denotes the dual basis of $\Gamma(W,\sh{E})$, then $u$ corresponds to the $C$-module homomorphism satisfying
\[
u(1)=\sum_{i=1}^n \mathrm{T}_i\otimes e_i.
\]
If $X'=\operatorname{Spec}(A)$ is affine and $\sh{E}_{(X')}\simeq\sh{O}_{X'}^n$, then $\Gamma(X',\sh{E}_{(X')})$ is a free $A$-module of rank $n$; figuratively, $\mathbf{V}(\check{\sh{E}})$ represents ``the set of points of the twisted affine space over $X$ defined by $\sh{E}$''.
\end{env}

\begin{env}[18.5.2]
\label{IV.18.5.2}
Recall (\sref[II]{II.1.7.8}) that by definition $\mathbf{V}(\check{\sh{E}})=\operatorname{Spec}(\mathbf{S}_{\sh{O}_X}(\check{\sh{E}}))$.
Let $\sh{J}$ be a quasi-coherent ideal of $\mathbf{S}_{\sh{O}_X}(\check{\sh{E}})$, so that $\operatorname{Spec}(\mathbf{S}_{\sh{O}_X}(\check{\sh{E}})/\sh{J})$ is a closed subprescheme of $\mathbf{V}(\check{\sh{E}})$; we shall interpret it as \emph{representing a functor} from the category of $X$-preschemes to that of sets.
A section $u\in\Gamma(X,\sh{E})$ identifies canonically with an $\sh{O}_X$-homomorphism $u:\sh{O}_X\to\sh{E}$, to which correspond by transposition an $\sh{O}_X$-homomorphism ${}^t u:\check{\sh{E}}\to\sh{O}_X$ and hence an $\sh{O}_X$-algebra homomorphism $v:\mathbf{S}_{\sh{O}_X}(\check{\sh{E}})\to\sh{O}_X$.
Let $\mathrm{Al}(X,\sh{E},\sh{J})$ denote the set of $u\in\Gamma(X,\sh{E})$ such that $\sh{J}$ is \emph{contained in the kernel} of $v$;
it follows immediately that
\[
X'\longmapsto\mathrm{Al}(X',\sh{E}_{(X')},\sh{J}\otimes_{\sh{O}_X}\sh{O}_{X'})
\]
is represented by $\operatorname{Spec}(\mathbf{S}_{\sh{O}_X}(\check{\sh{E}})/\sh{J})$.
If $X'=\operatorname{Spec}(A)$ is affine, $\sh{E}_{(X')}\simeq\sh{O}_{X'}^n$, and $\sh{J}\otimes_{\sh{O}_X}\sh{O}_{X'}=\widetilde{\mathfrak{J}}$ for an ideal $\mathfrak{J}$ of $A[\mathrm{T}_1,\ldots,\mathrm{T}_n]$, then this set identifies with the subset of $A^n$ consisting of the points $(t_1,\ldots,t_n)$ such that $P(t_1,\ldots,t_n)=0$ for every $P\in\mathfrak{J}$.
Thus $\operatorname{Spec}(\mathbf{S}_{\sh{O}_X}(\check{\sh{E}})/\sh{J})$ represents ``the algebraic subset of the twisted affine space defined by $\sh{E}$ consisting of the points annihilating $\sh{J}$''.
We also denote this $X$-prescheme by $\mathbf{Al}(\sh{E},\sh{J})$.
If $\sh{J}$ is an ideal of finite type in $\mathbf{S}_{\sh{O}_X}(\check{\sh{E}})$, then $\mathbf{Al}(\sh{E},\sh{J})$ is an $X$-prescheme of finite presentation, since $\mathbf{S}_{\sh{O}_X}(\check{\sh{E}})$ is an $\sh{O}_X$-algebra of finite presentation.
\end{env}

\begin{lemma}[18.5.3]
\label{IV.18.5.3}
Let $S$ be a prescheme, $f:X\to S$ a finite and locally free morphism (\sref{IV.18.2.7}).
Consider the contravariant functor from the category of $S$-preschemes to the category of sets
\[
\label{IV.18.5.3.1}
S'\mapsto\operatorname{Of}(X\times_S S')
\tag{18.5.3.1}
\]
where $\operatorname{Of}(X\times_S S')$ denotes the set of subsets that are both open and closed of the underlying space of $X\times_S S'$.
Then this functor is representable by an $S$-prescheme $\mathbf{Of}(X)$, which is affine, \'etale and of finite presentation over $S$.
\end{lemma}

\begin{proof}
\oldpage[IV]{127}
By hypothesis $X=\operatorname{Spec}(\sh{B})$, where $\sh{B}=f_*(\sh{O}_X)$ is a finite and locally free $\sh{O}_S$-algebra.
Set $X'=X\times_S S'$, $f'=f_{(S')}$, $\sh{B}'=\sh{B}\otimes_{\sh{O}_S}\sh{O}_{S'}=f'_*(\sh{O}_{X'})$, so that $X'=\operatorname{Spec}(\sh{B}')$;
there is then a canonical bijection, functorial in $S'$, from the set $\operatorname{Of}(X')$ to the set $\operatorname{Id}(\sh{B}')$ of \emph{idempotents} of the ring $\Gamma(S',\sh{B}')=\Gamma(X',\sh{O}_{X'})$.
Indeed, by virtue of the equivalence of the category of affine $S'$-schemes and the opposite category of quasi-coherent $\sh{O}_{S'}$-algebras (\sref[II]{II.1.2.7} and \sref[II]{II.1.3.1}), there is a canonical bijective correspondence between decompositions of $X'$ as a \emph{disjoint union} $X'_1\sqcup X'_2$ of two induced open subpreschemes and decompositions of $\sh{B}'$ as a \emph{direct sum} of two ideals $\sh{B}'_1$, $\sh{B}'_2$;
these latter in turn form a set in canonical bijective correspondence with $\operatorname{Id}(\sh{B}')$.
It suffices therefore to prove the lemma for the functor $S'\mapsto\operatorname{Id}(\sh{B}')$.

For this, we shall show that there exists an ideal $\sh{J}$ of finite type in $\mathbf{S}_{\sh{O}_S}(\check{\sh{B}})$ such that $\operatorname{Id}(\sh{B}')$ is of the form $\mathrm{Al}(S',\sh{B}',\sh{J}\otimes_{\sh{O}_S}\sh{O}_{S'})$ (\sref{IV.18.5.2}).
Since $\sh{B}$ is an $\sh{O}_S$-algebra, its inverse image $\sh{B}_{\mathbf{V}(\check{\sh{B}})}$ is an $\sh{O}_{\mathbf{V}(\check{\sh{B}})}$-algebra, and we may form in this algebra the square $c^2$ of the canonical section $c$ (\sref{IV.18.5.1});
it corresponds canonically to a homomorphism $u^{(2)}:\mathbf{S}_{\sh{O}_S}(\check{\sh{B}})\to\mathbf{S}_{\sh{O}_S}(\check{\sh{B}})\otimes_{\sh{O}_S}\sh{B}$, and one verifies immediately that for every affine open subset $W$ of $S$, with the notations of \sref{IV.18.5.1}, $u^{(2)}$ corresponds to the $C$-module homomorphism such that $u^{(2)}(1)=\sum_k(\sum_{i,j}c_{ijk}\mathrm{T}_i\mathrm{T}_j)\otimes e_k$, where $(c_{ijk})$ is the table of multiplication of the algebra $\Gamma(W,\sh{B})$.
Let $\sh{J}$ be the ideal of $\mathbf{S}_{\sh{O}_S}(\check{\sh{B}})$ locally generated by the coefficients of $(u-u^{(2)})(1)$ with respect to a basis of $\sh{B}$; it answers the question.
Indeed, $\Gamma(W,\sh{J})$ is generated by the polynomials $\mathrm{P}_k(\mathrm{T}_1,\ldots,\mathrm{T}_n)=\mathrm{T}_k-\sum_{i,j}c_{ijk}\mathrm{T}_i\mathrm{T}_j$, and the idempotents of $\Gamma(W,\sh{B})$ are precisely the elements $\sum_i t_i e_i$ such that $(t_1,\ldots,t_n)$ annihilates every $\mathrm{P}_k$ ($1\leq k\leq n$).
We deduce from \sref{IV.18.5.2} that the $S$-prescheme affine $\mathbf{Of}(X)=\mathbf{Al}(\sh{B},\sh{J})$ does represent the functor \sref{IV.18.5.3.1}.
It remains to show that $\mathbf{Of}(X)$ is \'etale over $S$, or equivalently, that it is \emph{formally \'etale} over $S$.
But if $S'$ is an $S$-prescheme, $S_0'$ a closed sub-prescheme of $S'$ defined by a locally nilpotent Ideal (and having the same underlying space as $S'$), it is clear that $X'=X\times_S S'$ and $X_0'=X\times_S S_0'$ have the same underlying space, hence the canonical map $\operatorname{Of}(X')\to\operatorname{Of}(X_0')$ is \emph{bijective}, which completes the proof \sref{IV.17.1.1}.
\end{proof}

\begin{proposition}[18.5.4]
\label{IV.18.5.4}
Let $S$ be a prescheme, $S_0$ a closed sub-prescheme of $S$;
consider the following properties:
\begin{enumerate}[label=\emph{\alph*)}]
  \item For every finite morphism $g:S'\to S$, the canonical map
\[
\label{IV.18.5.4.1}
\operatorname{Of}(S')\to\operatorname{Of}(S'\times_S S_0)
\tag{18.5.4.1}
  \]
is bijective {\upshape(}cf.\ \sref{IV.18.5.3}{\upshape)}.
  \item[a$'$)] For every finite and locally free morphism $g:S'\to S$, the map \sref{IV.18.5.4.1} is bijective.
\oldpage[IV]{128}
  \item For every \'etale and separated morphism $g:S'\to S$, the canonical map
\[
\label{IV.18.5.4.2}
\Gamma(S'/S)\to\Gamma(S'\times_S S_0/S_0)
\tag{18.5.4.2}
\]
is bijective.
\end{enumerate}
\emph{Condition $b)$ implies a$'$);
if moreover $S$ is quasi-compact and quasi-separated, condition $a)$ implies $b)$.}
\end{proposition}

\begin{proof}
Let us first prove that $b)$ implies $a')$.
Suppose therefore $b)$ satisfied, and let $g:S'\to S$ be a finite and locally free morphism;
set $S_0'=S'\times_S S_0$, so that $g_0=g_{(S_0)}:S_0'\to S_0$ is finite and locally free.
Then it follows from \sref{IV.18.5.3} that $\mathrm{P}=\mathbf{Of}(S')$ is an $S$-prescheme \emph{\'etale and separated};
moreover, by the definition of the functor $\operatorname{Of}$, we see immediately that if we set $\mathrm{P}_0=\mathbf{Of}(S_0')$ (for the category of $S_0$-preschemes), we have $\mathrm{P}_0=\mathrm{P}\times_S S_0$.
This being so, by definition we have the commutative diagram
\[
\xymatrix{
\Gamma(\mathrm{P}/S)\ar[r]\ar[d]^{\sim} & \Gamma(\mathrm{P}_0/S_0)\ar[d]^{\sim}\\
\operatorname{Of}(S')\ar[r] & \operatorname{Of}(S_0')
}
\]
where the vertical arrows are the bijections.
Since hypothesis $b)$, applied to the morphism $\mathrm{P}\to S$, implies that the first row is a bijection, the same holds for the second, which establishes our assertion.

Before proving that $a)$ implies $b)$ when $S$ is quasi-compact and quasi-separated, we will establish the following lemma.

\begin{lemma}[18.5.4.3]
\label{IV.18.5.4.3}
If $S$ and $S_0$ satisfy condition $a)$ of \sref{IV.18.5.4}, then, for every finite morphism $g:S'\to S$, $S'$ is the unique neighbourhood of $S_0'=g^{-1}(S_0)=S'\times_S S_0$ in $S'$.
\end{lemma}

\begin{proof}
Indeed, it amounts to saying that if $T'$ is a closed subset of $S'$ such that $T'\cap S_0'=\varnothing$, then $T'=\varnothing$.
Now, if we also denote by $T'$ the closed sub-prescheme of $S'$ having $T'$ as underlying space, the composite morphism $h:T'\to S'\xrightarrow{g}S$ is finite and $h^{-1}(S_0)$ is empty;
condition $a)$, applied to the morphism $h$, implies that $T'$ is necessarily empty.
\end{proof}

This lemma being established, let us first prove that, under hypothesis $a)$, the map \sref{IV.18.5.4.2} is \emph{injective}.
Indeed, if $u'$, $u''$ are two $S$-sections of $S'$, the fact that the morphism $S'\to S$ is unramified implies that the prescheme of coincidences of $u'$ and $u''$ is induced on an open subset $U$ of $S$ \sref{IV.17.4.6}.
If the restrictions to $S_0$ of $u'$ and $u''$ are equal, then the fact that $u'$ and $u''$ are open immersions \sref{IV.17.4.1} implies that $U$ contains $S_0$, and hence is equal to $S$ by virtue of Lemma \sref{IV.18.5.4.3}, applied to the case $S'=S$.

\oldpage[IV]{129}
It remains to show that, under hypothesis $a)$, the map \sref{IV.18.5.4.2} is \emph{surjective} (when $S$ is quasi-compact and quasi-separated).
Let therefore $u_0:S_0\to S_0'$ be an $S_0$-section of $S_0'$;
$u_0(S_0)$ being quasi-compact in $S'$, can be covered by a finite number of open affine subsets $V_i$ such that the restriction $g|V_i$ is a morphism of finite type;
if $V$ is the union of the $V_i$, we conclude that $g|V:V\to S$ is a morphism \emph{of finite presentation}, being separated and locally of finite presentation \sref{IV.1.6.1}.
Replacing $S'$ by $V$, one can hence suppose that $g$ is of finite presentation.
Since $S$ is quasi-compact and quasi-separated, and $g$ is quasi-finite and separated \sref{IV.17.6.1}, it follows from the ``Main theorem'' \sref{IV.8.12.6} that $g$ factorises as $S'\xrightarrow{j}S''\xrightarrow{f}S$, where $j$ is an open immersion and $f$ a \emph{finite} morphism.
Set $S_0''=S''\times_S S_0$, $j_0=j_{(S_0)}:S_0'\to S_0''$, which is an open immersion, and $f_0=f_{(S_0)}:S_0''\to S_0$, which is a finite morphism.
Then $u_0$ is also an $S_0$-section of $S_0''$.
Since $g_0:S_0'\to S_0$ is \'etale, $u_0$ is an open immersion of $S_0$ into $S_0'$ \sref{IV.17.4.1}, so $u_0(S_0)$ is open in $S_0'$, and \emph{a fortiori} in $S_0''$;
on the other hand, since $f_0$ is finite and hence separated, $u_0$ is a closed immersion of $S_0$ into $S_0''$ \sref[I]{I.5.4.6}.
Thus $X_0=u_0(S_0)$ is \emph{both open and closed} in $S_0''$.
In virtue of condition $a)$, there exists an open and closed subset $X$ of $S''$ such that $X\cap S_0''=X_0$.
Let us show that the morphism $f:S''\to S$ is \emph{\'etale at the points of $X$}:
indeed, the set of points of $X$ where $f|X$ is \'etale is open and contains $X_0\subset S_0'$.
But the lemma \sref{IV.18.5.4.3} applied to the finite morphism $f|X$ proves that $U=X$.
On the other hand, $S'\cap X$ is open in $X$ and contains $X_0$, hence by the same reasoning proves that $S'\cap X=X$, i.e.\ $X\subset S'$.
It remains to show that, for every $s\in S$, the geometric number $n(s)$ of points of $X\cap f^{-1}(s)$ is equal to $1$;
it will then follow that $f|X$ is radicial and surjective, and, since $f|X$ is \'etale, we shall have proved \sref{IV.17.9.1} that $f|X$ is an isomorphism from the open subset $X\subset S'$ onto $S$, whose inverse will be the desired $S$-section extending $u_0$.
But since $f|X$ is \'etale and finite, $s\mapsto n(s)$ is continuous on $S$ \sref{IV.18.2.8}, and since $X\cap S_0''=X_0$, we have $n(s)=1$ on $S_0$;
the set of points $s\in S$ such that $n(s)=1$ is therefore open in $S$ and contains $S_0$, so it is equal to $S$ by \sref{IV.18.5.4.3}.
\end{proof}

\begin{remark}[18.5.4.4]
\label{IV.18.5.4.4}
One can show that the statement \sref{IV.18.5.4} remains valid when, in condition $b)$, one supposes merely the morphism $g$ \'etale (but not necessarily separated) \cite[exp.~XII]{SGA1}.
\end{remark}

\begin{definition}[18.5.5]
\label{IV.18.5.5}
We say that a prescheme $S$ and a closed sub-prescheme $S_0$ of $S$ form a \emph{Henselian couple} if they satisfy condition $a)$ of \sref{IV.18.5.4}.
\end{definition}

Taking into account \sref[I]{I.5.1.8}, it amounts to the same to say that $(S,S_0)$ is a Henselian couple or that $(S_{\mathrm{red}},(S_0)_{\mathrm{red}})$ is one.

\begin{proposition}[18.5.6]
\label{IV.18.5.6}
\begin{enumerate}
  \item[(i)] If $(S,S_0)$ is a Henselian couple, then, for every finite morphism $f:S'\to S$, if $S_0'$ is the closed sub-prescheme $f^{-1}(S_0)$ of $S'$, the couple $(S',S_0')$ is Henselian.
  \item[(ii)] Let $S=\coprod_\alpha S^{(\alpha)}$ be a sum of preschemes, $S_0$ a closed sub-prescheme of $S$, the sum of sub-preschemes $S_0^{(\alpha)}$ of the $S^{(\alpha)}$.
For the couple $(S,S_0)$ to be Henselian, it is necessary and sufficient that each of the couples $(S^{(\alpha)},S_0^{(\alpha)})$ be Henselian.
\end{enumerate}
\end{proposition}

\begin{proof}
Assertion (i) is an immediate consequence of the definition, since for every finite morphism $g:S''\to S'$, $f\circ g:S''\to S$ is a finite morphism.
Similarly, under the conditions of (ii), for a morphism $g:S'\to S$ to be finite, it is necessary and sufficient that each \oldpage[IV]{130}of its restrictions $g^{(\alpha)}:S'^{(\alpha)}=g^{-1}(S^{(\alpha)})\to S^{(\alpha)}$ be finite.
If we set $S_0'=g^{-1}(S_0)$ and $S_0'^{(\alpha)}=(g^{(\alpha)})^{-1}(S_0^{(\alpha)})$, the open-and-closed subsets $U$ (resp.\ $U_0$) of $S'$ (resp.\ $S_0'$) correspond bijectively to the families $(U^{(\alpha)})$ (resp.\ $(U_0^{(\alpha)})$), where $U^{(\alpha)}$ (resp.\ $U_0^{(\alpha)}$) is an open-and-closed subset of $S'^{(\alpha)}$ (resp.\ $S_0'^{(\alpha)}$); assertion (ii) follows.
\end{proof}

\begin{remark}[18.5.7]
\label{IV.18.5.7}
Let $S=\operatorname{Spec}(A)$ be an affine scheme, $S_0$ a closed sub-prescheme defined by an ideal $\mathfrak{J}$ of $A$.
Then, if the couple $(S,S_0)$ is Henselian, the ideal $\mathfrak{J}$ is necessarily contained in the \emph{radical} of $A$.
Indeed, if $\mathfrak{m}$ is a maximal ideal of $A$, $\mathfrak{m}$ must belong to $V(\mathfrak{J})=S_0$, by virtue of \sref{IV.18.5.4.3}, whence the conclusion.
In particular, suppose that $S_0$ is \emph{reduced to a point}, then $\mathfrak{J}$ must be the radical of $A$, so $A$ must be a \emph{local} ring, and $S_0$ the unique closed point of $\operatorname{Spec}(A)$.
\end{remark}

\begin{definition}[18.5.8]
\label{IV.18.5.8}
We say that a ring $A$ is \emph{Henselian} if it is semi-local and if, denoting by $\mathfrak{r}$ the radical of $A$, the couple $(\operatorname{Spec}(A),\operatorname{Spec}(A/\mathfrak{r}))$ is Henselian.
We call a \emph{Henselian local scheme} a scheme isomorphic to the spectrum of a Henselian local ring.
\end{definition}

\begin{proposition}[18.5.9]
\label{IV.18.5.9}
\begin{enumerate}
  \item[(i)] For a semi-local ring $A$ to be Henselian, it is necessary and sufficient that it be a direct product of Henselian local rings.
  \item[(ii)] For a local ring $A$ to be Henselian, it is necessary and sufficient that every finite $A$-algebra $B$ be isomorphic to a product of local rings.
\end{enumerate}
\end{proposition}

\begin{proof}
(i) Indeed, the definition, applied to $S=\operatorname{Spec}(A)$ and $S_0=\operatorname{Spec}(A/\mathfrak{r})$, shows that $S_0$ is a finite discrete set, closed in $S$, and $S$ is the union of a finite number of open and closed subsets $S_i$ ($1\leq i\leq n$), pairwise disjoint, each containing exactly one of the maximal ideals $\mathfrak{m}_i$ of $A$;
the conclusion follows from \sref{IV.18.5.6}, (ii) and the remark \sref{IV.18.5.7}.

(ii) Every finite morphism $S'\to\operatorname{Spec}(A)$ is of the form $\operatorname{Spec}(B)\to\operatorname{Spec}(A)$, where $B$ is a finite $A$-algebra.
If $k$ is the residue field of $A$, $\operatorname{Spec}(B\otimes_A k)$ is a finite Artinian spectrum, hence finite and discrete.
To say that the couple $(\operatorname{Spec}(A),\operatorname{Spec}(k))$ is Henselian means that $B$ is a direct product of rings $A_i$ such that $\operatorname{Spec}(A_i\otimes_A k)$ is reduced to a point, i.e.\ that each $A_i$ (which is a finite $A$-algebra) must have only a single maximal ideal (Bourbaki, \emph{Alg.\ comm.}, ch.~V, \textsection~2, no.~1, prop.~1).
\end{proof}

The study of Henselian rings is thus essentially reduced to that of Henselian \emph{local} rings.

\begin{proposition}[18.5.10]
\label{IV.18.5.10}
If $A$ is a Henselian ring, every finite $A$-algebra $B$ is a Henselian ring (hence a direct product of Henselian local rings \sref{IV.18.5.9}).
\end{proposition}

\begin{proof}
Indeed, if $\mathfrak{r}'$ is the radical of $B$, the inverse image of $\mathfrak{r}'$ in $A$ is the radical $\mathfrak{r}$ of $A$, and every prime ideal of $B$ lying over a maximal ideal of $A$ is a maximal ideal of $B$;
hence the set $V(\mathfrak{r}')$ in $\operatorname{Spec}(B)$ is the inverse image of the set $V(\mathfrak{r})$ in $\operatorname{Spec}(A)$.
The proposition is then a consequence of \sref{IV.18.5.6}, (i).
\end{proof}

\begin{theorem}[18.5.11]
\label{IV.18.5.11}
Let $A$ be a local ring, $\mathfrak{m}$ its maximal ideal.
The following conditions are equivalent:
\begin{enumerate}[label=\emph{\alph*)}]
  \item $A$ is Henselian, in other words, every finite $A$-algebra $B$ is isomorphic to a product of local rings.
\oldpage[IV]{131}
  \item[a$'$)] Condition $a)$ is satisfied for all $A$-algebras $B$ of the form $A[\mathrm{T}]/F.A[\mathrm{T}]$, where $F\in A[\mathrm{T}]$ is a monic polynomial.
  \item Let $S=\operatorname{Spec}(A)$, $S_0=\operatorname{Spec}(k)$.
For every \'etale morphism $g:S'\to S$, if we set $S_0'=S'\otimes_A k$, every $S_0$-section $u_0$ of $S_0'$ is the restriction of an $S$-section $u$ of $S'$.
  \item For every morphism locally of finite type $f:X\to S$, separated, and every point $x\in X$ such that $f(x)$ is equal to the closed point $s$ of $S$ and that $f$ is quasi-finite at the point $x$ (\sref[III]{III.Err.20}), $X$ is the sum of two preschemes $X'$, $X''$ such that $X'=\operatorname{Spec}(\sh{O}_{X,x})$ and that $f|X':X'\to S$ is a finite morphism.
  \item[c$'$)] For every morphism locally of finite type $f:X\to S$, and every point $x\in X$ such that $f$ be quasi-finite at the point $x$ and $f(x)$ equal to the closed point $s$ of $S$, $\sh{O}_{X,x}$ is a finite algebra over $\sh{O}_{S,s}=A$.
  \item[c$''$)] For every morphism locally of finite presentation $f:X\to S$, and every point $x\in X$ such that $f$ be quasi-finite at the point $x$ and $f(x)$ equal to the closed point $s$ of $S$, $\sh{O}_{X,x}$ is a finite algebra and of finite presentation over $\sh{O}_{S,s}=A$.
\end{enumerate}
\end{theorem}

\begin{proof}
Let us first note that $c')$ (resp.\ $c''$) is equivalent to the same condition where moreover $f$ is supposed \emph{separated}, the question being local on $X$.
Similarly, condition $b)$ is equivalent to the same condition where moreover $g$ is supposed \emph{separated}:
indeed, it suffices to apply the latter to the restriction of $g$ to an open neighbourhood affine of the point $u_0(S_0)$ in $S'$.
Let us henceforth restrict to the case where, in $b)$, $c')$ and $c''$, the morphisms are separated.

The fact that $a)$ implies $b)$ and that $b)$ implies $a')$ is a particular case of \sref{IV.18.5.4}.
Let us moreover show that $a')$ implies $a)$.
It is a matter of proving that if $e_0$ is an idempotent of $C=B\otimes_A k$, there is an idempotent $e\in B$ whose image is $e_0$ in $C$.
If $b$ is an element of $B$ whose image in $C$ is $e_0$, the sub-$A$-algebra $A[b]=B'$ of $B$ is finite, and the canonical image $C'$ of $B'\otimes_A k$ in $C$ contains $e_0$.
Now, $B'\otimes_A k$ is a finite $k$-algebra, hence a direct product of finite local $k$-algebras, and consequently $e_0$ is the image in $C'$ of an idempotent $e_0'$ of $B'\otimes_A k$.
We are thus reduced to the case where $B$ is monogenic, and hence by $a')$, a quotient of an algebra of the form $A[\mathrm{T}]/F.A[\mathrm{T}]$;
now, by $a')$, $A[\mathrm{T}]/F.A[\mathrm{T}]$ is a direct product of local rings, hence so is each of its quotient algebras, which proves the existence of the idempotent $e$.

It is immediate that $c)$ implies $a)$, as one sees by induction on the number of maximal ideals of the semi-local ring $B$.
To see that $a)$ implies $c)$, we may, by \sref{IV.13.1.4}, restrict to the case where $f$ is affine and quasi-finite.
Then, by the ``Main theorem'' \sref{IV.8.12.8}, $f$ can be written as the composite
\[
X\xrightarrow{j}Y\xrightarrow{g}S,
\]
where $g$ is finite and $j$ is an open immersion.
Since $Y=\operatorname{Spec}(B)$, where $B$ is a finite $A$-algebra, condition $a)$ implies that $B$ is a direct product of local rings, which are evidently finite $A$-algebras, and $\sh{O}_{X,x}$ identifies with one of these local rings because $f(x)=s$;
moreover, every open subset of $Y$ containing $x$ necessarily contains $\operatorname{Spec}(\sh{O}_{X,x})$.

It is trivial that $c)$ implies $c')$, by virtue of the remark at the beginning.
Let us prove that $c')$ implies $c'')$.
Suppose $c')$ satisfied and let us prove that, under the hypotheses \oldpage[IV]{132}of $c'')$, the subset $Z=\operatorname{Spec}(\sh{O}_{X,x})$ identifies with an \emph{open and closed} subset of $X$, which will establish $c'')$ \sref[I]{I.2.4.2}.
First, the composite morphism $Z\xrightarrow{j}X\xrightarrow{f}S$, where $j$ is the canonical morphism \sref[I]{I.2.4.1}, is finite and of finite presentation by hypothesis;
and since $f$ is separated and locally of finite presentation, $j$ is also finite and of finite presentation (\sref[II]{II.6.1.5} and \sref[I]{I.4.3}), hence $j$ is a closed morphism \sref[II]{II.6.1.10}, which proves that $Z$ is closed in $X$.
It then follows from \sref[I]{I.2.4.2} and \sref[I]{I.4.2.2} that $j$ is a closed immersion.
If $\sh{J}$ is the ideal of $\sh{O}_X$ defining $Z$, then by hypothesis $\sh{J}_x=0$, and hence $\sh{J}_z=0$ at every point $z$ of some open neighbourhood $V$ of $x$ in $X$, since $\sh{J}$ is a quasi-coherent $\sh{O}_X$-module of finite type (\sref{IV.1.4.7} and \sref[0]{0.5.2.2}).
Such a neighbourhood contains $Z$, so $Z$ is open in $X$, since $\sh{J}|V=0$.

Finally, $c'')$ implies $a')$:
indeed, if $B=A[\mathrm{T}]/F.A[\mathrm{T}]$, the morphism $X=\operatorname{Spec}(B)\to\operatorname{Spec}(A)=Y$ is finite and of finite presentation, and the preceding argument shows that $B$ is a direct product of finite $A$-algebras $\sh{O}_{X,x_i}$, where the $x_i$ are the points of the fibre of the closed point of $Y$.
\end{proof}

\begin{corollary}[18.5.12]
\label{IV.18.5.12}
Let $A$ be a semi-local ring, $\mathfrak{r}$ its radical;
set $S=\operatorname{Spec}(A)$, $S_0=\operatorname{Spec}(A/\mathfrak{r})$.
For $A$ to be Henselian, it is necessary and sufficient that, for every finite morphism $f:X\to S$ and every \'etale and separated morphism $g:Y\to S$, if we set $X_0=X\times_S S_0$ and $Y_0=Y\times_S S_0$, the canonical map
\[
\operatorname{Hom}_S(X,Y)\to\operatorname{Hom}_{S_0}(X_0,Y_0)
\]
is bijective.
\end{corollary}

\begin{proof}
The sufficiency of the condition follows from the equivalence of $a)$ and $b)$ in \sref{IV.18.5.11}, by applying this condition to the case where $f=1_S$.
To see that the condition is necessary, note that if $A$ is Henselian, the couple $(X,X_0)$ is Henselian by \sref{IV.18.5.6}, (i);
moreover $\operatorname{Hom}_S(X,Y)=\Gamma(X\times_S Y/X)$ and $\operatorname{Hom}_{S_0}(X_0,Y_0)=\Gamma(X_0\times_{S_0}Y_0/X_0)$;
since $X\times_S Y$ is \'etale and separated over $X$, the conclusion follows from the fact that $a)$ implies $b)$ in \sref{IV.18.5.4}.
\end{proof}

\begin{remark}[18.5.13]
\label{IV.18.5.13}
The equivalent conditions $a)$, $a')$, $b)$ of Theorem~\sref{IV.18.5.11} are also equivalent to the following (``Hensel's lemma''):
\begin{enumerate}
  \item[a$''$)] For every monic polynomial $F\in A[\mathrm{T}]$, with canonical image $F_0\in k[\mathrm{T}]$, and every decomposition $F_0=G_0 H_0$ of $F_0$ into a product of two monic polynomials that are coprime $G_0$, $H_0$ in $k[\mathrm{T}]$, there exists a unique couple $(G,H)$ of monic polynomials of $A[\mathrm{T}]$ possessing the following properties:
$G_0$ and $H_0$ are the canonical images respectively of $G$ and $H$, we have $F=GH$, and the ideal of $A[\mathrm{T}]$ generated by $G$ and $H$ is equal to $A[\mathrm{T}]$.
\end{enumerate}
\end{remark}

\begin{lemma}[18.5.13.1]
\label{IV.18.5.13.1}
Let $A$ be a local ring with residue field $k$, $F\in A[\mathrm{T}]$ a monic polynomial, $B$ the $A$-algebra $A[\mathrm{T}]/F.A[\mathrm{T}]$.
There exists a canonical bijective correspondence between the decompositions of $B$ into the direct product of two quotient $A$-algebras $B'$, $B''$ and the decompositions $F=GH$ of $F$ into a product of two monic polynomials $G$, $H$ of $A[\mathrm{T}]$, such that the ideal generated by $G$ and $H$ is equal to $A[\mathrm{T}]$;
the quotient algebras $B'$, $B''$ corresponding to such a couple of polynomials $G$, $H$ are respectively $A[\mathrm{T}]/H.A[\mathrm{T}]$ and $A[\mathrm{T}]/G.A[\mathrm{T}]$.
\end{lemma}

\begin{proof}
\oldpage[IV]{133}
If $F=GH$ and $G$ and $H$ generate the ideal $A[\mathrm{T}]$, there are two polynomials $P$, $Q$ of $A[\mathrm{T}]$ such that $1=PG+QH$.
We deduce that the intersection of the principal ideals $\mathfrak{a}=G.A[\mathrm{T}]$ and $\mathfrak{b}=H.A[\mathrm{T}]$ is equal to $\mathfrak{c}=F.A[\mathrm{T}]$:
indeed, if $R\in G.A[\mathrm{T}]\cap H.A[\mathrm{T}]$, we can write $R=PRG+QRH$;
since $RG$ (resp.\ $RH$) is a multiple of $F$ since $R$ is a multiple of $H$ (resp.\ $G$), therefore $R\in F.A[\mathrm{T}]$.
Since $A[\mathrm{T}]=\mathfrak{a}+\mathfrak{b}$, $A[\mathrm{T}]/\mathfrak{c}$ is the direct sum of ideals $\mathfrak{a}/(\mathfrak{a}\cap\mathfrak{b})$ and $\mathfrak{b}/(\mathfrak{a}\cap\mathfrak{b})$, canonically isomorphic respectively to $A[\mathrm{T}]/\mathfrak{b}$ and $A[\mathrm{T}]/\mathfrak{a}$.

Conversely, suppose given a decomposition of $B$ into a direct product of two $A$-algebras $B'$, $B''$, which identify canonically with two ideals $e'B$, $e''B$ of $B$, corresponding to a decomposition $1=e'+e''$ into orthogonal idempotents $e'$, $e''$ of $B$.
Set further $B_0=B\otimes_A k=k[\mathrm{T}]/F_0.k[\mathrm{T}]$, where $F_0$ is the canonical image of $F$ in $k[\mathrm{T}]$, of the same degree $n$ as $F$;
if $e_0'$, $e_0''$ are the canonical images of $e'$, $e''$ in $B_0$, these are two orthogonal idempotents such that $1=e_0'+e_0''$, and $B_0$ is the direct product of $B_0'=e_0'B_0$ and $B_0''=e_0''B_0$.
Let $t$ and $t_0$ be the canonical images of $\mathrm{T}$ in $B$ and $B_0$;
$B_0'$ (resp.\ $B_0''$) is a $k$-algebra generated by $t_0'=e_0't_0$ (resp.\ $t_0''=e_0''t_0$), and it admits a basis of the form $\{e_0',t_0',t_0'^2,\ldots,t_0'^{s-1}\}$ (resp.\ $\{e_0'',t_0'',t_0''^2,\ldots,t_0''^{r-1}\}$) with $r+s=n$.
On the other hand, $B$ being a free $A$-module (with basis $\{1,t,\ldots,t^{n-1}\}$), $B'$ and $B''$ are projective $A$-modules, hence free since $A$ is a local ring (Bourbaki, \emph{Alg.\ comm.}, ch.~II, \textsection~5, no.~3, cor.\ of prop.~5);
it follows from the above and from Bourbaki, \emph{Alg.\ comm.}, ch.~II, \textsection~3, no.~3, prop.~5, that if we set $t'=e't$, $t''=e''t$, $\{e',t',\ldots,t'^{s-1}\}$ (resp.\ $\{e'',t'',\ldots,t''^{r-1}\}$) is a basis of the $A$-module $B'$ (resp.\ $B''$).
There is hence a monic polynomial $H$ (resp.\ $G$) of degree $s$ (resp.\ $r$) in $A[\mathrm{T}]$ such that $e'H(t')=0$ and $e''G(t'')=0$;
as $t^h=t'^h+t''^h$ for every integer $h\geq 1$, with $t'^h=e't'^h$ and $t''^h=e''t''^h$, we also have $G(t)=e'G(t')$ and $H(t)=e''H(t'')$, whence $G(t)H(t)=0$;
we conclude that the polynomial $G(T)H(T)$ is divisible by $F(T)$;
but since these two monic polynomials have the same degree, we have $GH=F$.
Moreover, $B'$ (resp.\ $B''$) is isomorphic to $A[\mathrm{T}]/H.A[\mathrm{T}]$ (resp.\ $A[\mathrm{T}]/G.A[\mathrm{T}]$).
Finally, there are two polynomials $R$, $S$ in $A[\mathrm{T}]$ such that $e'=R(t)$ and $e''=S(t)$;
as $R(t)=e'R(t')+e''R(t'')$, we necessarily have $e''R(t'')=0$, and similarly $e'S(t')=0$, so that, by definition of $G$ and $H$, $R=QH$ and $S=PG$, where $P$, $Q$ belong to $A[\mathrm{T}]$;
the relation $1=R(t)+S(t)$ in $B$ gives $PG+QH=1+LF$ for some polynomial $L\in A[\mathrm{T}]$, and since $F=GH$, this proves that the ideal generated by $G$ and $H$ is equal to $A[\mathrm{T}]$, and completes the proof of the lemma.
\end{proof}

This lemma being established, it suffices to apply it, to the ring $A$ on the one hand, to the field $k$ on the other, to see immediately that conditions $a')$ and $a'')$ are equivalent.

\begin{proposition}[18.5.14]
\label{IV.18.5.14}
Every semi-local ring, separated and complete for the $\mathfrak{r}$-preadic topology (where $\mathfrak{r}$ is the radical of $A$) is Henselian.
\end{proposition}

\begin{proof}
Indeed, $A$ is a direct product of separated and complete local rings (Bourbaki, \emph{Alg.\ comm.}, ch.~III, \textsection~2, no.~13, cor.\ of prop.~19), so we are reduced to the case where $A$ is a local ring.
Let us verify criterion $a')$ of \sref{IV.18.5.11}.
Since $B$ is a free $A$-module of finite type, $B$ is separated and complete for the $\mathfrak{r}$-adic topology, \oldpage[IV]{134}which is also the $\mathfrak{s}$-adic topology, where $\mathfrak{s}$ is the radical of the semi-local ring $B$, because $B/\mathfrak{r}B$ is an Artinian ring with radical $\mathfrak{s}/\mathfrak{r}B$.
It is then known (Bourbaki, \emph{Alg.\ comm.}, ch.~III, \textsection~2, no.~13, cor.\ of prop.~19) that $B$ is a direct product of local rings.
\end{proof}

\begin{proposition}[18.5.15]
\label{IV.18.5.15}
Let $A$ be a Henselian ring, $\mathfrak{r}$ its radical.
Then the functor $B\mapsto B/\mathfrak{r}B$ is an equivalence of the category of finite \'etale $A$-algebras with the category of finite \'etale $(A/\mathfrak{r})$-algebras.
\end{proposition}

\begin{proof}
The fact that the functor is fully faithful is a particular case of \sref{IV.18.5.12}.
To show that this functor is an equivalence of categories, we may reduce to the case where $A$ is local;
it suffices then to apply \sref{IV.18.1.1} (for \'etale morphisms), with $S=\operatorname{Spec}(A)$ and $S_0=\operatorname{Spec}(A/\mathfrak{r})$, reduced to a single point.
\end{proof}

\begin{remarks}[18.5.16]
\label{IV.18.5.16}
\begin{enumerate}
  \item[(i)] We do not know if Proposition~\sref{IV.18.5.15} generalises to a Henselian couple $(S,S_0)$, even when $S$ is affine and Noetherian.
  \item[(ii)] Let $A$ be a Noetherian ring, $\mathfrak{J}$ an ideal of $A$ such that $A$ is \emph{separated and complete} for the $\mathfrak{J}$-preadic topology.
Then the couple $(\operatorname{Spec}(A),\operatorname{Spec}(A/\mathfrak{J}))$ is Henselian:
indeed, for every finite $A$-algebra $B$, $B$ is separated and complete for the $\mathfrak{J}$-preadic topology (\sref[0]{0.7.3.6}).
Replacing $B$ by $A$ and $\mathfrak{J}B$ by $\mathfrak{J}$, everything amounts to seeing that the map that to every idempotent of $A$ associates its class mod $\mathfrak{J}$, is \emph{bijective}.
Now $A=\varprojlim A/\mathfrak{J}^n$.
Write $\operatorname{Idem}(A)$ for the set of idempotents of $A$, and, for every ring homomorphism $\varphi:A\to B$, let $\operatorname{Idem}(\varphi)$ denote the map $\operatorname{Idem}(A)\to\operatorname{Idem}(B)$ induced by $\varphi$.
It follows from the definition of the projective limit that
\[
\operatorname{Idem}(A)=\varprojlim_n\operatorname{Idem}(A/\mathfrak{J}^n),
\]
for the transition maps $\psi_{nm}:\operatorname{Idem}(A/\mathfrak{J}^m)\to\operatorname{Idem}(A/\mathfrak{J}^n)$ induced by the canonical maps $A/\mathfrak{J}^m\to A/\mathfrak{J}^n$.
But since $\operatorname{Spec}(A/\mathfrak{J}^n)\to\operatorname{Spec}(A/\mathfrak{J}^m)$ is a homeomorphism, the $\psi_{nm}$ are bijections (as seen in the proof of \sref{IV.18.5.3});
this proves our assertion.
\end{enumerate}
\end{remarks}

\begin{theorem}[18.5.17]
\label{IV.18.5.17}
Let $A$ be a Henselian local ring, $S=\operatorname{Spec}(A)$, $s$ the closed point.
For every smooth morphism $f:X\to S$, the canonical map
\[
\Gamma(X/S)\to\Gamma(X_s/\kappa(s))
\]
(where $X_s=f^{-1}(s)$) is surjective.
\end{theorem}

\begin{proof}
The datum of a $\kappa(s)$-section of $X_s$ is equivalent to that of a point $x\in X$ above $s$ whose residue field $\kappa(x)$ is isomorphic to $\kappa(s)$, and it is a matter of proving that there exists an $S$-section $u:S\to X$ such that $u(s)=x$.
Taking into account \sref{IV.17.16.3}, (i), we can suppose that $f$ is \'etale.
Then the conclusion follows from the criterion \sref{IV.18.5.11}, $b)$.
(The reader will note that by virtue of this criterion, the validity of \sref{IV.18.5.17} for a local ring is necessary and sufficient for $A$ to be Henselian.)
\end{proof}

\begin{remark}[18.5.18]
\label{IV.18.5.18}
The conditions of \sref{IV.18.5.11} are also equivalent to the following:
\begin{enumerate}
  \item[d)] For every morphism locally of finite type $f:X\to S$ and every point $x\in X$ such that $f(x)$ be equal to the closed point $s$ of $S$ and that $\sh{O}_{X,x}/\mathfrak{m}_s\sh{O}_{X,x}$ be a field canonically isomorphic to $k=\kappa(s)$, there exists an open neighbourhood $U$ of $x$ in $X$ such that $f|U$ is a closed immersion.
\end{enumerate}
\end{remark}

\begin{proof}
\oldpage[IV]{135}
It is immediate that $d)$ implies condition $b)$ of \sref{IV.18.5.11} since a closed immersion \'etale is an open immersion \sref{IV.18.9.1}.
Conversely, suppose the conditions of \sref{IV.18.5.11} satisfied, and let us prove $d)$.
The hypothesis of $d)$ implies that $f$ is quasi-finite at the point $x$ (\sref[III]{III.Err.20}), by virtue of \sref{IV.13.1.4};
hence, by condition $c)$ of \sref{IV.18.5.11}, $\sh{O}_{X,x}$ is a finite $A$-algebra and $\operatorname{Spec}(\sh{O}_{X,x})$ is an open neighbourhood $U$ of $x$ in $X$.
Moreover, since $\sh{O}_{X,x}/\mathfrak{m}_s\sh{O}_{X,x}$ is isomorphic to $k=A/\mathfrak{m}_s$, the Nakayama lemma proves that the homomorphism $A\to\sh{O}_{X,x}$ is surjective, hence $f|U$ is a closed immersion.
\end{proof}

\begin{proposition}[18.5.19]
\label{IV.18.5.19}
Let $A$ be a Noetherian Henselian local ring, $S=\operatorname{Spec}(A)$, $s$ the closed point of $S$, and $f:X\to S$ a proper morphism;
set $X_s=f^{-1}(s)$.
Then the map $Y\to Y\cap X_s$ is a bijection from the set of connected components of $X$ to the set of connected components of $X_s$.
\end{proposition}

\begin{proof}
By considering a closed sub-prescheme of $X$ having for underlying space a connected component of $X$, we are reduced to showing that if $X$ is \emph{connected and non-empty}, then $X_s$ is \emph{connected and non-empty}.
The fact that $X_s$ is non-empty follows from (\sref[II]{II.7.2.1});
to show that $X_s$ is connected, we argue by contradiction, considering the Stein factorisation $f:X\xrightarrow{f'}S'\xrightarrow{g}S$ of the proper morphism $f$ (\sref[III]{III.4.3.3});
by hypothesis the finite discrete set $g^{-1}(s)$ would contain at least two points.
Since $A$ is Henselian and $g$ separated and of finite type, it would then follow from \sref{IV.18.5.11}, $c)$ that $S'$ would be the sum of two non-empty preschemes $S_1'$, $S_2'$ (since the intersection of any one of them with $g^{-1}(s)$ is reduced to a single point).
Since $f'$ is surjective (\sref[III]{III.4.3.1}), we would conclude that $X$ is a sum of two non-empty preschemes, contradicting the hypothesis.
\end{proof}
\subsection{Henselization.}
\label{subsection:IV.18.6}

\begin{env}[18.6.1]
\label{IV.18.6.1}
Given a local ring $A$, we shall say that a local $A$-algebra $B$ is \emph{essentially \'etale} if there exist an \'etale $A$-algebra $C$ and a prime ideal $\mathfrak{n}$ of $C$ such that $B$ is $A$-isomorphic to $C_{\mathfrak{n}}$ and the composite homomorphism $A\to C\to C_{\mathfrak{n}}$ is local (in other words, $\mathfrak{n}$ lies above the maximal ideal $\mathfrak{m}$ of $A$).
This implies \sref{IV.17.6.1} that $\mathfrak{n}$ is the only prime ideal of $B$ lying above $\mathfrak{m}$, since $B$ is a local ring;
if $B$, $B'$ are two local essentially \'etale $A$-algebras, every $A$-homomorphism $B\to B'$ is therefore local.

If $A'$ is a local essentially \'etale $A$-algebra and $A''$ a local essentially \'etale $A'$-algebra, then $A''$ is a local essentially \'etale $A$-algebra.
Indeed, by hypothesis $A'=B_{\mathfrak{n}}$, where $B$ is an \'etale $A$-algebra and $\mathfrak{n}$ a prime ideal of $B$ above the maximal ideal of $A$, and $A''=B'_{\mathfrak{n}'}$, where $B'$ is an \'etale $A'$-algebra and $\mathfrak{n}'$ a prime ideal of $B'$ above $\mathfrak{n}B_{\mathfrak{n}}$.
If we set $S=B-\mathfrak{n}$, so that $A'=S^{-1}B$, then $B'$ is of the form $S^{-1}C$, where $C$ is a $B$-algebra of finite presentation.
Consequently $C$ is an $A$-algebra of finite presentation and $A''$ is of the form $C_{\mathfrak{r}}$, where $\mathfrak{r}$ is a prime ideal of $C$ above the maximal ideal of $A$.
Since $A'$ is formally \'etale over $A$ and $A''$ is formally \'etale over $A'$, $A''$ is formally \oldpage[IV]{136}\'etale over $A$ (for the discrete topologies \sref{IV.17.1.3});
the morphism $\operatorname{Spec}(C)\to\operatorname{Spec}(A)$ is therefore \'etale at the point $\mathfrak{r}$ \sref{IV.17.6.1}, and hence there exists an element $g\in C$ such that $C_g$ is an \'etale $A$-algebra, which proves that $A''$ is an essentially \'etale $A$-algebra.

Given a local ring $A$, there exists a \emph{set} $\mathfrak{E}$ of local essentially \'etale $A$-algebras such that every local essentially \'etale $A$-algebra is $A$-isomorphic to an algebra belonging to $\mathfrak{E}$.
It suffices indeed to observe that there exists a set $\mathfrak{F}$ of $A$-algebras of finite type such that every $A$-algebra of finite type is isomorphic to an algebra belonging to $\mathfrak{F}$;
one may take $\mathfrak{F}$ to be the set of quotients of the polynomial algebras $A[\mathrm{T}_1,\ldots,\mathrm{T}_n]$ ($n\in\mathbf{N}$).

In what follows, if $A$ and $B$ are two local rings, we denote by $\Hom.\operatorname{loc}(A,B)$ the set of \emph{local} homomorphisms from $A$ to $B$.
\end{env}

\begin{lemma}[18.6.2]
\label{IV.18.6.2}
Let $\mathrm{A}$, $\mathrm{A}'$ be two local rings, $k$, $k'$ their residue fields,
$\vphi:\mathrm{A}\to\mathrm{A}'$ a local homomorphism making $\mathrm{A}'$ an $A$-algebra essentially \'{e}tale \sref{IV.18.6.1}
and such that the corresponding homomorphism $k\to k'$ is bijective.
Then, for every local Henselian ring $\mathrm{B}$, the canonical map
\[
  \Hom(\vphi,\,1_{\mathrm{B}}):\Hom.\operatorname{loc}(\mathrm{A}',\mathrm{B})\to\Hom.\operatorname{loc}(\mathrm{A},\mathrm{B})
\]
is bijective.
\end{lemma}

\begin{proof}
Let $\mathrm{S}=\Spec(\mathrm{A})$, $\mathrm{Y}=\Spec(\mathrm{B})$, and let $s$ and $y$ be the closed points of S and Y respectively.
By hypothesis, $\mathrm{A}'$ is isomorphic to a local ring $\sh{O}_{\mathrm{X},x}$ of an S-scheme X \'{e}tale over S, at a point $x\in\mathrm{X}$ above $s$.
Suppose given a local homomorphism $\psi:\mathrm{A}\to\mathrm{B}$, making Y an S-scheme;
it amounts to showing that there exists \emph{one and only one} S-morphism $f:\mathrm{Y}\to\mathrm{X}$ such that $f(y)=x$.
Let $\mathrm{X}'=\mathrm{X}\times_{\mathrm{S}}\mathrm{Y}$, and note that since $k(x)=k(s)$, there exists a unique point $x'\in\mathrm{X}'$ above $x$ and $y$, and that $k(x')=k(y)$.
It remains to show that there exists a unique Y-section $f'$ of $\mathrm{X}'$ such that $f'(y)=x'$.
Now, the morphism $g:\mathrm{X}'\to\mathrm{Y}$ is \'{e}tale and separated, and the fibre $\mathrm{X}'_0=g^{-1}(y)$ has at the point $x'$ the local ring $k(x')=k(y)$.
If we set $\mathrm{Y}_0=\Spec(k(y))$, there exists a unique $\mathrm{Y}_0$-section $f'_0$ of $\mathrm{X}'_0$ such that $f'_0(y)=x'$, and the conclusion follows from the fact that B is assumed Henselian and from \sref{IV.18.5.11},~$b$).
\end{proof}

We say that a local A-algebra $\mathrm{A}'$ satisfying the conditions of \sref{IV.18.6.2} is \emph{strictly essentially \'{e}tale}.
Note that the criterion \sref{IV.18.5.11},~$b$) means that, \emph{for A to be Henselian, it is necessary and sufficient that every strictly essentially \'{e}tale A-algebra be A-isomorphic to} A.

\begin{lemma}[18.6.3]
\label{IV.18.6.3}
Let $\mathrm{A}$ be a local ring, $\mathrm{A}_1$, $\mathrm{A}_2$ two A-algebras that are local and strictly essentially \'{e}tale.
\begin{enumerate}[label=\emph{\roman*)}]
  \item There exists at most one A-homomorphism (necessarily local) from $\mathrm{A}_1$ to $\mathrm{A}_2$.
  \item There exists a local A-algebra $\mathrm{A}_3$ that is strictly essentially \'{e}tale, and two A-homomorphisms $\mathrm{A}_1\to\mathrm{A}_3$, $\mathrm{A}_2\to\mathrm{A}_3$.
\end{enumerate}
\end{lemma}

\begin{proof}
Let $\mathrm{S}=\Spec(\mathrm{A})$;
by hypothesis there are two S-schemes \'{e}tale over S, $\mathrm{X}_1$, $\mathrm{X}_2$, and two points $x_1\in\mathrm{X}_1$, $x_2\in\mathrm{X}_2$ above the closed point $s$ of S and such that $\mathrm{A}_1=\sh{O}_{\mathrm{X}_1,x_1}$, $\mathrm{A}_2=\sh{O}_{\mathrm{X}_2,x_2}$;
Let $\mathrm{X}_3=\mathrm{X}_1\times_{\mathrm{S}}\mathrm{X}_2$;
the hypotheses $k(x_1)=k(x_2)=k(s)$ \sref[I]{I.3.4.9} imply that there exists a unique point $x_3\in\mathrm{X}_3$ above $x_1$ and $x_2$, and that $k(x_3)=k(s)$.
Furthermore, $\mathrm{X}_3$ is \'{e}tale over S \sref{IV.17.3.3}, so $\mathrm{A}_3=\sh{O}_{\mathrm{X}_3,x_3}$ satisfies the conditions of (ii).
On the other hand, \oldpage[IV]{137}we have seen that every $A$-homomorphism from $A_1$ to $A_2$ is necessarily local;
it corresponds to an $S$-morphism $f$ from $X_2'=\operatorname{Spec}(A_2)$ to $X_1$ such that $f(x_2)=x_1$, or, equivalently, on setting $X_3'=X_1\times_S X_2'$, to an $X_2'$-section $f'$ of $X_3'$ such that $f'(x_2)=x_3$.
Since $k(x_3)=k(x_2)$ and $X_2'$ is connected, the uniqueness of $f'$ follows from \sref{IV.17.4.9}, whence (i).
\end{proof}

\begin{env}[18.6.4]
\label{IV.18.6.4}
Let us denote by $\mathfrak{S}$ the subset of the set $\mathfrak{E}$ defined in \sref{IV.18.6.1} formed by the A-algebras that are strictly essentially \'{e}tale belonging to $\mathfrak{E}$.
It follows from \sref{IV.18.6.3} that the relation ``there exists an A-homomorphism from $\mathrm{A}_1$ to $\mathrm{A}_2$'' is a \emph{preorder} relation on $\mathfrak{S}$, and makes $\mathfrak{S}$ a \emph{filtered increasing} set.
Indexing $\mathfrak{S}$ by itself by means of the identity map, it follows from \sref{IV.18.6.3},~(i) that if $\lambda\leq\mu$ in $\mathfrak{S}$, there exists a \emph{unique} $\vphi_{\mu\lambda}:\mathrm{A}_\lambda\to\mathrm{A}_\mu$ and $(\mathrm{A}_\lambda,\vphi_{\mu\lambda})$ is clearly an \emph{inductive system} of local A-algebras, the $\vphi_{\mu\lambda}$ being local homomorphisms.
Furthermore, it follows from \sref{IV.17.3.5} that for $\lambda\leq\mu$, $\mathrm{A}_\mu$ is an $\mathrm{A}_\lambda$-algebra that is strictly essentially \'{e}tale.
\end{env}

\begin{definition}[18.6.5]
\label{IV.18.6.5}
The \emph{Henselization} of a local ring $\mathrm{A}$ is the inductive limit of the inductive system $(\mathrm{A}_\lambda,\vphi_{\mu\lambda})$, denoted ${}^h\!\mathrm{A}$, defined in \sref{IV.18.6.4}.
\end{definition}

This definition does not depend on the choice of $\mathfrak{E}$; if $\mathfrak{E}'$ is another set having the same property as $\mathfrak{E}$, and $\mathfrak{S}'$ the subset of $\mathfrak{E}'$ consisting of the strictly essentially \'{e}tale A-algebras, it follows from \sref{IV.18.6.3},~(ii) that, for the preorder relations considered, $\mathfrak{S}$ and $\mathfrak{S}'$ are \emph{cofinal} in $\mathfrak{S}''=\mathfrak{S}\cup\mathfrak{S}'$, so they give the same inductive limit up to an isomorphism.
We shall moreover see below \sref{IV.18.6.6},~(i) and~(ii) that ${}^h\!\mathrm{A}$ and the canonical homomorphism $\mathrm{A}\to{}^h\!\mathrm{A}$ are the solution of a universal problem, and are therefore determined up to a unique isomorphism.

One will observe that if A is a field, then evidently ${}^h\!\mathrm{A}=\mathrm{A}$.
In the general case, ${}^h\!\mathrm{A}$ is also the Henselization of all the A-algebras that are strictly essentially \'{e}tale $\mathrm{A}_\lambda$ ($\lambda\in\mathfrak{S}$), since if B is an $\mathrm{A}_\lambda$-algebra that is strictly essentially \'{e}tale, B is also an A-algebra that is strictly essentially \'{e}tale \sref{IV.18.6.1}, so (up to an A-isomorphism) one of the $\mathrm{A}_\mu$ for $\mu\geq\lambda$, and the $\mathrm{A}_\mu$ such that $\mu\geq\lambda$ and that $\mathrm{A}_\mu$ is an $\mathrm{A}_\lambda$-algebra that is strictly essentially \'{e}tale form a cofinal set in the preordered set of the $\mathrm{A}_\lambda$, by virtue of \sref{IV.18.6.1} and \sref{IV.18.6.3}.

\begin{theorem}[18.6.6]
\label{IV.18.6.6}
Let $\mathrm{A}$ be a local ring, ${}^h\!\mathrm{A}$ its Henselization.
\begin{enumerate}[label=\emph{\roman*)}]
  \item ${}^h\!\mathrm{A}$ is a Henselian local ring and the structural homomorphism $\mathrm{A}\to{}^h\!\mathrm{A}$ is local.
  \item For every Henselian local ring $\mathrm{B}$, the canonical map
  \[
    \Hom.\operatorname{loc}({}^h\!\mathrm{A},\mathrm{B})\to\Hom.\operatorname{loc}(\mathrm{A},\mathrm{B})
  \]
  is bijective.
  \item ${}^h\!\mathrm{A}$ is a faithfully flat $A$-module, and if $\mathfrak{m}$ is the maximal ideal of $\mathrm{A}$, $\mathfrak{m}.{}^h\!\mathrm{A}$ is the maximal ideal of ${}^h\!\mathrm{A}$, and the homomorphism $\mathrm{A}/\mathfrak{m}\to{}^h\!\mathrm{A}/\mathfrak{m}.{}^h\!\mathrm{A}$ of residue fields is bijective.
  \item If $\hat{\mathrm{A}}$ and $({}^h\!\mathrm{A})\hat{\phantom{x}}$ are the separated completions of the local rings $\mathrm{A}$ and ${}^h\!\mathrm{A}$, the homomorphism $\hat{\mathrm{A}}\to({}^h\!\mathrm{A})\hat{\phantom{x}}$ deduced from the structural homomorphism $\mathrm{A}\to{}^h\!\mathrm{A}$ by completion is bijective.
  \item For ${}^h\!\mathrm{A}$ to be Noetherian, it is necessary and sufficient that A be so.
  \item If A is Henselian, the canonical homomorphism $\mathrm{A}\to{}^h\!\mathrm{A}$ is bijective.
\end{enumerate}
\end{theorem}

\begin{proof}
\oldpage[IV]{138}
Let $\mathfrak{m}_\lambda$ be the maximal ideal of $\mathrm{A}_\lambda$;
it follows from the fact that, in \sref{IV.17.6.1}, $a)$ implies $c')$, that for $\lambda\leq\mu$ one has $\mathfrak{m}_\mu=\mathfrak{m}_\lambda\mathrm{A}_\mu$, that the homomorphism $\mathrm{A}_\lambda/\mathfrak{m}_\lambda\to\mathrm{A}_\mu/\mathfrak{m}_\mu$ is bijective, and that $\mathrm{A}_\mu$ is a flat $\mathrm{A}_\lambda$-module.
The fact that ${}^h\!\mathrm{A}$ is local, that the homomorphism $\mathrm{A}\to{}^h\!\mathrm{A}$ is local, assertion~(iii) and the sufficiency of~(v) follow therefore from \sref[0]{0.10.3.1.3};
the necessity of~(v) follows from the fact that ${}^h\!\mathrm{A}$ is a faithfully flat A-module \sref[0]{0.6.5.2}.

To prove that ${}^h\!\mathrm{A}$ is Henselian, let us apply the criterion \sref{IV.18.5.11},~$b$).
Let $\mathrm{S}=\Spec({}^h\!\mathrm{A})$, $\mathrm{S}_0=\Spec({}^h\!\mathrm{A}/\mathfrak{m}.{}^h\!\mathrm{A})$, let $g:\mathrm{S}'\to\mathrm{S}$ be an \'{e}tale morphism, set $\mathrm{S}_0'=g^{-1}(\mathrm{S}_0)$, and let $f_0:\mathrm{S}_0\to\mathrm{S}_0'$ be an $\mathrm{S}_0$-section.
reasoning as in \sref{IV.18.5.4}, one can suppose that $g$ is of finite presentation.
It follows then from \sref{IV.8.8.2} and \sref{IV.17.7.5} that there exist an index $\lambda\in\mathfrak{S}$, an \'etale morphism $g_\lambda:\mathrm{S}'^{(\lambda)}\to\mathrm{S}^{(\lambda)}=\Spec(\mathrm{A}_\lambda)$ and, on setting $\mathrm{S}_0^{(\lambda)}=\Spec(\mathrm{A}_\lambda/\mathfrak{m}\mathrm{A}_\lambda)$ and $\mathrm{S}_0'^{(\lambda)}=g_\lambda^{-1}(\mathrm{S}_0^{(\lambda)})$, an $\mathrm{S}_0^{(\lambda)}$-section $f_0^{(\lambda)}:\mathrm{S}_0^{(\lambda)}\to\mathrm{S}_0'^{(\lambda)}$, such that
\[
\mathrm{S}'=\mathrm{S}'^{(\lambda)}\times_{\mathrm{S}^{(\lambda)}}\mathrm{S},\qquad
g=g_\lambda\times 1,\qquad
f_0=f_0^{(\lambda)}\times 1.
\]
Let $s$ be the closed point of $S$, $x=f_0(s)$, let $x_\lambda$ be the projection of $x$ in $\mathrm{S}'^{(\lambda)}$, and let $s_\lambda$ be the closed point of $\mathrm{S}^{(\lambda)}$.
Since $x_\lambda$ lies above $s_\lambda$, the local ring $\mathrm{C}_\lambda$ of $\mathrm{S}'^{(\lambda)}$ at $x_\lambda$ is an essentially \'etale $\mathrm{A}_\lambda$-algebra;
moreover, since $f_0^{(\lambda)}(s_\lambda)=x_\lambda$, we have $k(x_\lambda)=k(s_\lambda)$, in other words $\mathrm{C}_\lambda$ is a local strictly essentially \'etale $\mathrm{A}_\lambda$-algebra, and consequently, by \sref{IV.18.6.1}, it is $\mathrm{A}_\lambda$-isomorphic to an $A$-algebra $\mathrm{A}_\mu$ with $\mu\geq\lambda$.
There is therefore an $\mathrm{A}_\lambda$-homomorphism $\mathrm{C}_\lambda\to{}^h\!\mathrm{A}$, that is, an $\mathrm{S}^{(\lambda)}$-morphism $h:\mathrm{S}\to\mathrm{S}'^{(\lambda)}$ such that $h(s)=x_\lambda$, and hence an $S$-section $f$ of $\mathrm{S}'$ such that $f(s)=x$;
this completes the proof of (i).

To prove~(ii), it suffices to note that $\Hom.\operatorname{loc}({}^h\!\mathrm{A},\mathrm{B})\cong\varprojlim\Hom.\operatorname{loc}(\mathrm{A}_\lambda,\mathrm{B})$ and that by \sref{IV.18.6.2} the canonical homomorphisms $\Hom.\operatorname{loc}(\mathrm{A}_\lambda,\mathrm{B})\leftarrow\Hom.\operatorname{loc}(\mathrm{A},\mathrm{B})$ are bijective.

To prove~(iv), note that for every $\lambda$ and every integer $n>0$, $\mathfrak{m}_\lambda^n=\mathfrak{m}^n\mathrm{A}_\lambda$, and
\[
(\mathfrak{m}.{}^h\!\mathrm{A})^n=\mathfrak{m}^n.{}^h\!\mathrm{A}=\varinjlim_\lambda\mathfrak{m}_\lambda^n.
\]
By exactness of the functor $\varinjlim$ in the category of $A$-modules,
\[
{}^h\!\mathrm{A}/(\mathfrak{m}.{}^h\!\mathrm{A})^n=\varinjlim_\lambda(\mathrm{A}_\lambda/\mathfrak{m}_\lambda^n),
\]
and it therefore suffices to show that, for every $n$ and every index $\lambda$, the homomorphism $\mathrm{A}/\mathfrak{m}^n\to\mathrm{A}_\lambda/\mathfrak{m}_\lambda^n$ is bijective.
On the other hand, since $\mathrm{A}_\lambda$ is a flat A-module, one has
\[
  \mathfrak{m}_\lambda^n/\mathfrak{m}_\lambda^{n+1}=(\mathfrak{m}^n/\mathfrak{m}^{n+1})\otimes_{\mathrm{A}}\mathrm{A}_\lambda
  =(\mathfrak{m}^n/\mathfrak{m}^{n+1})\otimes_{\mathrm{A}/\mathfrak{m}}(\mathrm{A}_\lambda/\mathfrak{m}\mathrm{A}_\lambda)
\]
and as $\mathrm{A}/\mathfrak{m}\to\mathrm{A}_\lambda/\mathfrak{m}\mathrm{A}_\lambda$ is bijective, the homomorphism $\mathfrak{m}^n/\mathfrak{m}^{n+1}\to\mathfrak{m}_\lambda^n/\mathfrak{m}_\lambda^{n+1}$ is also bijective;
the conclusion then follows from Bourbaki, \emph{Alg. comm.}, chap.~III, \S~2, no~8, cor.~3 du th.~1.

Finally, if A is Henselian, it follows from the remark preceding \sref{IV.18.6.3} that the canonical homomorphisms $\mathrm{A}\to\mathrm{A}_\lambda$ are bijective, which proves~(vi) by definition of ${}^h\!\mathrm{A}$.
\end{proof}

\begin{env}[18.6.7]
\label{IV.18.6.7}
Let A be a \emph{semi-local} ring, $\mathfrak{m}_i$ ($1\leq i\leq r$) its maximal ideals.
The \emph{Henselization} of A is the ring denoted again ${}^h\!\mathrm{A}$, defined as the \emph{product} $\prod_{i=1}^r{}^h(\mathrm{A}_{\mathfrak{m}_i})$ of the Henselizations of the local rings $\mathrm{A}_{\mathfrak{m}_i}$.
It is a faithfully flat A-module and a semi-local A-algebra, whose maximal ideals are the $\mathfrak{m}_i.{}^h\!\mathrm{A}$ by virtue of \sref{IV.18.6.6},~(iii);
furthermore, if $\mathfrak{r}=\bigcap\mathfrak{m}_i$ is the radical of A, it follows from the above that $\mathfrak{r}.{}^h\!\mathrm{A}$
\oldpage[IV]{139}
is the radical of ${}^h\!\mathrm{A}$ and that the canonical map $\mathrm{A}/\mathfrak{r}\to{}^h\!\mathrm{A}/\mathfrak{r}.{}^h\!\mathrm{A}$ is bijective.
As the separated completion $\hat{\mathrm{A}}$ of A for the $\mathfrak{r}$-adic topology is the product of the separated completions $\hat{\mathrm{A}}_{\mathfrak{m}_i}$ (Bourbaki, \emph{Alg. comm.}, chap.~III, \S~2, no~13, prop.~18), the canonical homomorphism $\hat{\mathrm{A}}\to({}^h\!\mathrm{A})\hat{\phantom{x}}$ is bijective by \sref{IV.18.6.6},~(iv), and it is clear by \sref{IV.18.6.6},~(v) that for ${}^h\!\mathrm{A}$ to be Noetherian, it is necessary and sufficient that A be so.

To obtain the analogue of the universal property \sref{IV.18.6.6},~(ii), when A and B are two \emph{semi-local} rings of radicals $\mathfrak{r}$, $\mathfrak{s}$, let us agree to call a \emph{semi-local homomorphism} $\vphi$ from A to B any homomorphism such that $\vphi(\mathfrak{r})\subset\mathfrak{s}$.
In the particular case where B is a \emph{product} of local rings $\mathrm{B}_j$ ($1\leq j\leq n$), the data of such a homomorphism amounts to the data of its projections $\vphi_j:\mathrm{A}\to\mathrm{B}_j$, which are homomorphisms subject only to the condition that $\vphi_j(\mathfrak{r})\subset\mathfrak{n}_j$, where $\mathfrak{n}_j$ is the maximal ideal of $\mathrm{B}_j$.
Moreover, since $\mathfrak{m}_i$ ($1\leq i\leq m$) are the maximal ideals of A, the prime ideal $\vphi_j^{-1}(\mathfrak{n}_j)$ must contain one of the $\mathfrak{m}_i$, since it must contain their intersection $\mathfrak{r}$, hence it is equal to one of the $\mathfrak{m}_i$ and $\vphi_j$ factorises as $\mathrm{A}\to\mathrm{A}_{\mathfrak{m}_i}\xrightarrow{\psi_{ij}}\mathrm{B}_j$, where $\psi_{ij}$ is a \emph{local} homomorphism.
If $\Hom.\operatorname{sloc}(\mathrm{A},\mathrm{B})$ denotes the set of semi-local homomorphisms from A to B, one can thus identify this set with $\prod_j\left(\bigcup_i\Hom.\operatorname{loc}(\mathrm{A}_{\mathfrak{m}_i},\mathrm{B}_j)\right)$;
in the case considered, since a semi-local Henselian ring is a direct product of Henselian local rings, one sees that the universal property \sref{IV.18.6.6},~(ii) is still valid for a semi-local ring A, provided one replaces ``local homomorphisms'' by ``semi-local homomorphisms'' \sref{IV.18.6.7}.
\end{env}

\begin{proposition}[18.6.8]
\label{IV.18.6.8}
Let $\mathrm{A}$ be a semi-local ring, and let $\mathrm{B}$ be a semi-local A-algebra integral over $\mathrm{A}$.
Then $\mathrm{B}\otimes_{\mathrm{A}}({}^h\!\mathrm{A})$ is a semi-local B-algebra isomorphic to ${}^h\!\mathrm{B}$.
\end{proposition}

\begin{proof}
Since $\mathrm{B}\otimes_{\mathrm{A}}({}^h\!\mathrm{A})$ is integral over ${}^h\!\mathrm{A}$, each of its maximal ideals lies above a maximal ideal of ${}^h\!\mathrm{A}$, hence above a maximal ideal of A.  Since B is integral over A, the prime ideals of B lying above a maximal ideal of A are maximal ideals of B.  Thus the projections in $\Spec({}^h\!\mathrm{A})$ and $\Spec(\mathrm{B})$ of a maximal ideal of $\mathrm{B}\otimes_{\mathrm{A}}({}^h\!\mathrm{A})$ are both maximal ideals.  Taking account of \sref{IV.18.6.6},~(iii), and \sref[I]{I.3.4.9}, one concludes that $\mathrm{B}\otimes_{\mathrm{A}}({}^h\!\mathrm{A})$ is semi-local.

Furthermore, for every Henselian local ring C, $\Hom.\operatorname{sloc}(\mathrm{B}\otimes_{\mathrm{A}}({}^h\!\mathrm{A}),\mathrm{C})$ is in bijection with the set of pairs $(\vphi,\psi)$ of semi-local homomorphisms $\vphi:\mathrm{B}\to\mathrm{C}$ and $\psi:{}^h\!\mathrm{A}\to\mathrm{C}$ such that the composites
\[
  \mathrm{A}\longrightarrow\mathrm{B}\xrightarrow{\vphi}\mathrm{C},
  \qquad
  \mathrm{A}\longrightarrow{}^h\!\mathrm{A}\xrightarrow{\psi}\mathrm{C}
\]
are equal.  By the bijection between $\Hom.\operatorname{sloc}(\mathrm{A},\mathrm{C})$ and $\Hom.\operatorname{sloc}({}^h\!\mathrm{A},\mathrm{C})$, for every $\vphi\in\Hom.\operatorname{sloc}(\mathrm{B},\mathrm{C})$ there exists a unique $\psi$ having this property.  Hence the map
\[
  \Hom.\operatorname{sloc}(\mathrm{B}\otimes_{\mathrm{A}}({}^h\!\mathrm{A}),\mathrm{C})
  \longrightarrow \Hom.\operatorname{sloc}(\mathrm{B},\mathrm{C})
\]
is bijective, which proves the proposition by the uniqueness of the solution of a universal problem.
\end{proof}

\begin{theorem}[18.6.9]
\label{IV.18.6.9}
Let $\mathrm{A}$ be a semi-local ring.
\begin{enumerate}[label=\emph{\roman*)}]
  \item For ${}^h\!\mathrm{A}$ to be reduced (resp. normal), it is necessary and sufficient that A be so.
  \item Suppose A Noetherian.
  Then, for every prime ideal $\mathfrak{p}$ of A, the ring $({}^h\!\mathrm{A})_{\mathfrak{p}}/\mathfrak{p}.({}^h\!\mathrm{A})_{\mathfrak{p}}$ is a finite direct product of separable algebraic extensions of $k(\mathfrak{p})$ (which implies that the fibres of the canonical morphism $\Spec({}^h\!\mathrm{A})\to\Spec(\mathrm{A})$ are \emph{geometrically regular}).
\end{enumerate}
\end{theorem}

\begin{proof}
\oldpage[IV]{140}
(i) As ${}^h\!\mathrm{A}$ is a faithfully flat A-module, it follows from \sref{IV.2.1.13} that if ${}^h\!\mathrm{A}$ is reduced (resp. normal), so is A.
To prove the converse, one can restrict to the case where A is local, so that, with the notations of \sref{IV.18.6.4}, one has ${}^h\!\mathrm{A}=\varinjlim\mathrm{A}_\alpha$.
Now, it follows from the definition of the $\mathrm{A}_\alpha$ and from \sref{IV.17.5.7} that if A is reduced (resp. integral and integrally closed), so is each of the $\mathrm{A}_\alpha$;
furthermore, the homomorphisms $\mathrm{A}_\alpha\to\mathrm{A}_\beta$ for $\alpha\leq\beta$ are injective by virtue of \sref[0]{0.6.5.1} since $\mathrm{A}_\beta$ is a faithfully flat $\mathrm{A}_\alpha$-module;
hence the morphisms $\Spec(\mathrm{A}_\beta)\to\Spec(\mathrm{A}_\alpha)$ are dominant, and one concludes from \sref{IV.5.13.2} (resp. \sref{IV.5.13.4}) that ${}^h\!\mathrm{A}$ is reduced (resp. integral and integrally closed).

(ii) One can restrict to the case where A is local.
By virtue of the fact that the functor $\varinjlim$ commutes with tensor products, the fibre of the morphism $\Spec({}^h\!\mathrm{A})\to\Spec(\mathrm{A})$ at a point $\mathfrak{p}$ is the projective limit of the fibres of the morphisms $\Spec(\mathrm{A}_\alpha)\to\Spec(\mathrm{A})$ at this point.
As ${}^h\!\mathrm{A}$ is Noetherian, one is reduced to proving the following lemma:
\end{proof}

\begin{lemma}[18.6.9.1]
\label{IV.18.6.9.1}
Let $(\mathrm{B}_\alpha)$ be an inductive system of Artinian rings, $\mathrm{B}=\varinjlim\mathrm{B}_\alpha$ its inductive limit.
If $\mathrm{B}$ is Noetherian, then $\mathrm{B}$ is Artinian;
if moreover, for $\alpha\leq\beta$, the homomorphisms $\vphi_{\beta\alpha}:\mathrm{B}_\alpha\to\mathrm{B}_\beta$ are injective, then there exists $\lambda$ such that, for $\alpha\geq\lambda$, the number of local components $\mathrm{B}_\alpha^{(i)}$ of $\mathrm{B}_\alpha$ is constant, that $\vphi_{\beta\alpha}(\mathrm{B}_\alpha^{(i)})\subset\mathrm{B}_\beta^{(i)}$ for every $i$ and that the $\mathrm{B}^{(i)}=\varinjlim\mathrm{B}_\alpha^{(i)}$ are the local components of $\mathrm{B}$.
\end{lemma}

\begin{proof}
Indeed $\mathrm{Y}_\alpha=\Spec(\mathrm{B}_\alpha)$, $\mathrm{Y}=\Spec(\mathrm{B})$, which, as a topological space, is equal to $\varprojlim\mathrm{Y}_\alpha$ \sref{IV.8.2.10}.
The spaces $\mathrm{Y}_\alpha$ are finite and discrete;
furthermore, if the $\vphi_{\beta\alpha}$ are injective, the canonical morphisms $\mathrm{Y}_\beta\to\mathrm{Y}_\alpha$ are dominant, hence surjective.
The lemma then follows from the following purely topological lemma:
\end{proof}

\begin{lemma}[18.6.9.2]
\label{IV.18.6.9.2}
Let $(\mathrm{Y}_\alpha,\psi_{\alpha\beta})$ be a projective system of finite discrete topological spaces.
If $\mathrm{Y}=\varprojlim\mathrm{Y}_\alpha$ is Noetherian, Y is finite and discrete.
If moreover the $\psi_{\alpha\beta}$ are surjective, then there exists $\lambda$ such that for $\alpha\geq\lambda$, the number of elements of $\mathrm{Y}_\alpha$ is constant and equal to the number of elements of $\mathrm{Y}$.
\end{lemma}

\begin{proof}
Indeed, since Y is compact and Noetherian, every subset of Y is compact, hence closed; thus Y is discrete, and therefore finite since it is compact.
If moreover the $\psi_{\alpha\beta}$ are surjective, the same holds for $\psi_\alpha:\mathrm{Y}\to\mathrm{Y}_\alpha$, so $\operatorname{Card}(\mathrm{Y})\geq\operatorname{Card}(\mathrm{Y}_\alpha)$, and as $\operatorname{Card}(\mathrm{Y}_\alpha)$ is an increasing function of $\alpha$, there exists $\lambda$ such that for $\alpha\geq\lambda$, $\operatorname{Card}(\mathrm{Y}_\alpha)=\operatorname{Card}(\mathrm{Y}_\lambda)$;
the $\psi_{\alpha\beta}$ are then bijective for $\alpha\geq\lambda$ and one therefore has $\operatorname{Card}(\mathrm{Y})=\operatorname{Card}(\mathrm{Y}_\lambda)$.
\end{proof}

\begin{corollary}[18.6.10]
\label{IV.18.6.10}
Let $\mathrm{A}$ be a local Noetherian ring.
For ${}^h\!\mathrm{A}$ to possess one of the following properties:
\begin{enumerate}[label=\emph{\alph*)}]
  \item being a Cohen-Macaulay ring \sref[0]{0.16.5.3};
  \item satisfying the property $(\mathrm{S}_n)$ \sref{IV.5.7.3};
  \item being regular;
  \item satisfying the property $(\mathrm{R}_n)$ \sref{IV.5.8.2};
\end{enumerate}
it is necessary and sufficient that A possess the same property.
\end{corollary}

\begin{proof}
\oldpage[IV]{141}
It follows immediately from the fact that the fibres of the morphism $\Spec({}^h\!\mathrm{A})\to\Spec(\mathrm{A})$ are geometrically regular \sref{IV.18.6.9} and from \sref{IV.6.4.1},~\sref{IV.6.5.1} and \sref{IV.6.5.3}.
\end{proof}

\begin{corollary}[18.6.11]
\label{IV.18.6.11}
Let $\mathrm{A}$ be a semi-local Noetherian ring.
For a prime ideal $\mathfrak{p}$ of ${}^h\!\mathrm{A}$ to belong to $\operatorname{Ass}({}^h\!\mathrm{A})$, it is necessary and sufficient that $\mathfrak{p}\cap\mathrm{A}$ belong to $\operatorname{Ass}(\mathrm{A})$.
\end{corollary}

\begin{proof}
Taking account of the description of the fibres of the morphism $\Spec({}^h\!\mathrm{A})\to\Spec(\mathrm{A})$, this also follows immediately from \sref{IV.3.3.1}.
\end{proof}

\begin{proposition}[18.6.12]
\label{IV.18.6.12}
Let $\mathrm{A}$ be a local ring;
the following properties are equivalent:
\begin{enumerate}[label=\emph{\alph*)}]
  \item A is unibranch \sref[0]{0.23.2.1} (resp. reduced and unibranch).
  \item For every A-algebra $\mathrm{A}'$ that is local and strictly essentially \'{e}tale, $\Spec(\mathrm{A}')$ is irreducible (resp. integral).
  \item $\Spec({}^h\!\mathrm{A})$ is irreducible (resp. integral).
\end{enumerate}
\end{proposition}

\begin{proof}
First note that ${}^h(\mathrm{A}_{\mathrm{red}})=({}^h\!\mathrm{A})_{\mathrm{red}}$: indeed, ${}^h(\mathrm{A}_{\mathrm{red}})=({}^h\!\mathrm{A})\otimes_{\mathrm{A}}\mathrm{A}_{\mathrm{red}}$ by \sref{IV.18.6.8}, and since ${}^h(\mathrm{A}_{\mathrm{red}})$ is reduced by \sref{IV.18.6.9}, it is equal to $({}^h\!\mathrm{A})_{\mathrm{red}}$.
This allows us, in the proof, to consider only the case where A is reduced; then so are $\mathrm{A}'$ \sref{IV.17.5.7} and ${}^h\!\mathrm{A}$ \sref{IV.18.6.9}.

With the notations of \sref{IV.18.6.4}, condition $b)$ means that all the $\mathrm{A}_\alpha$ are integral; one therefore concludes that ${}^h\!\mathrm{A}=\varinjlim\mathrm{A}_\alpha$ is integral \sref{IV.5.13.3}, so $b)$ implies $c)$.

It is clear that if ${}^h\!\mathrm{A}$ is integral, so are each of the $\mathrm{A}_\alpha\subset{}^h\!\mathrm{A}$, hence $c)$ implies $b)$.

Let us show that $c)$ implies $a)$.
Note first that $\mathrm{A}\subset{}^h\!\mathrm{A}$ and that ${}^h\!\mathrm{A}$ is an integral domain;
let K and L be the fields of fractions of A and ${}^h\!\mathrm{A}$ respectively.
It suffices to show that for every \emph{finite} A-algebra B contained in K, the ring B is local.  Now B is semi-local, and its Henselization is $\mathrm{B}\otimes_{\mathrm{A}}({}^h\!\mathrm{A})$ by \sref{IV.18.6.8}.  By flatness, $\mathrm{B}\otimes_{\mathrm{A}}({}^h\!\mathrm{A})$ identifies with a subring of $\mathrm{K}\otimes_{\mathrm{A}}({}^h\!\mathrm{A})$, and the latter with a subring of L; hence ${}^h\!\mathrm{B}$ is integral.  But since ${}^h\!\mathrm{B}$ is a semi-local Henselian ring, it is a finite direct product of local rings \sref{IV.18.5.9}; it can therefore be integral only if it is local, which indeed implies that B is local.

Let us finally prove that $a)$ implies $c)$.
Let C be the integral closure of the integral ring A.
As C is a local ring by hypothesis, $\mathrm{C}\otimes_{\mathrm{A}}({}^h\!\mathrm{A})$ is the Henselization of C \sref{IV.18.6.8}, hence is integral and integrally closed \sref{IV.18.6.9}.
Since A is a subring of C, flatness identifies ${}^h\!\mathrm{A}$ with a subring of ${}^h\!\mathrm{C}$; it is therefore integral as well, which completes the proof.
\end{proof}

\begin{corollary}[18.6.13]
\label{IV.18.6.13}
Let $\mathrm{A}$ be a local Henselian ring.
Then, for A to be unibranch, it is necessary and sufficient that $\mathrm{A}_{\mathrm{red}}$ be integral.
\end{corollary}

\begin{proposition}[18.6.14]
\label{IV.18.6.14}
\begin{enumerate}[label=\emph{\roman*)}]
  \item Every filtered inductive limit of Henselian local rings (where the transition homomorphisms are local) is a Henselian local ring.
  \item The functor $\mathrm{A}\leadsto{}^h\!\mathrm{A}$ in the category of local rings (where the morphisms are local homomorphisms) commutes with filtered inductive limits.
  \item Every Henselian local ring A is an inductive limit of an inductive filtered system of Henselian and Noetherian local rings $\mathrm{A}_\lambda$ (the transition homomorphisms $\mathrm{A}_\lambda\to\mathrm{A}_\mu$ being local).
\end{enumerate}
\end{proposition}

\begin{proof}
(i) Let $(\mathrm{A}_\lambda)$ be a filtered inductive system of Henselian local rings, where the transition homomorphisms are local.
Let us show that $\mathrm{A}'=\varinjlim\mathrm{A}_\lambda$, which is local \sref[0]{0.10.3.1.3}, is Henselian.
Now, let C be an $\mathrm{A}'$-algebra \'{e}tale, $\mathfrak{n}$ a prime ideal
\oldpage[IV]{142}
of C above the maximal ideal $\mathfrak{m}'$ of $\mathrm{A}'$, such that $\mathrm{C}_{\mathfrak{n}}$ is strictly essentially \'{e}tale over $\mathrm{A}'$ \sref{IV.18.6.1}, and in virtue of \sref{IV.8.8.2},~(ii) and \sref{IV.17.7.8}, there exist an index $\lambda$ and an $\mathrm{A}_\lambda$-algebra $\mathrm{C}_\lambda$ \'{e}tale such that $\mathrm{C}=\mathrm{C}_\lambda\otimes_{\mathrm{A}_\lambda}\mathrm{A}'$ up to an isomorphism.
Furthermore, if $\mathfrak{n}_\lambda$ is the inverse image of $\mathfrak{n}$ in $\mathrm{C}_\lambda$, and $\mathfrak{m}_\lambda$ the maximal ideal of $\mathrm{A}_\lambda$, one can, by virtue of \sref{IV.8.8.2.4} and of the transitivity of fibres \sref[I]{I.3.6.4}, suppose that $(\mathrm{C}_\lambda)_{\mathfrak{n}_\lambda}$ is strictly essentially \'{e}tale over $\mathrm{A}_\lambda$;
since $\mathrm{A}_\lambda$ is Henselian, $(\mathrm{C}_\lambda)_{\mathfrak{n}_\lambda}$ is necessarily isomorphic to $\mathrm{A}_\lambda$ \sref{IV.18.5.11},~$b$);
there is therefore a neighbourhood of $\mathfrak{n}_\lambda$ in $\Spec(\mathrm{C}_\lambda)$ isomorphic to $\Spec(\mathrm{A}_\lambda)$ \sref[I]{I.6.5.4};
one concludes that there is a neighbourhood of $\mathfrak{n}$ in $\Spec(\mathrm{C})$ isomorphic to $\Spec(\mathrm{A}')$, hence that $\mathrm{C}_\mathfrak{n}$ is isomorphic to $\mathrm{A}'$.  This proves that $\mathrm{A}'$ is Henselian \sref{IV.18.6.2} and completes the proof.

(ii) Suppose that $\mathrm{B}=\varinjlim\mathrm{B}_\lambda$, where the $\mathrm{B}_\lambda$ are local rings, the transition morphisms $\mathrm{B}_\lambda\to\mathrm{B}_\mu$ ($\lambda\leq\mu$) being local.
Set $\mathrm{A}_\lambda={}^h\!\mathrm{B}_\lambda$;
the $\mathrm{A}_\lambda$ are Henselian local rings \sref{IV.18.6.6},~(vi), and by virtue of \sref{IV.18.6.6},~(ii), the transition morphisms $\mathrm{B}_\lambda\to\mathrm{B}_\mu$ uniquely determine local homomorphisms $\mathrm{A}_\lambda\to\mathrm{A}_\mu$, so that $(\mathrm{A}_\lambda)$ is an inductive system.
It amounts to showing that $\mathrm{A}'=\varinjlim\mathrm{A}_\lambda$ is canonically isomorphic to $\mathrm{A}={}^h\!\mathrm{B}$.
By virtue of \sref{IV.18.6.6},~(ii) and of the definition of inductive limits, one has, for every Henselian local ring E,
\[
  \Hom.\operatorname{loc}(\mathrm{A}',\mathrm{E})=\varprojlim\Hom.\operatorname{loc}(\mathrm{A}_\lambda,\mathrm{E})=\varprojlim\Hom.\operatorname{loc}(\mathrm{B}_\lambda,\mathrm{E})=\Hom.\operatorname{loc}(\mathrm{B},\mathrm{E}).
\]
But as $\mathrm{A}'$ is Henselian by~(i), this proves our assertion.

(iii) The ring A is a filtered inductive limit of its Noetherian local subrings $\mathrm{B}_\lambda$, the transition homomorphisms being local \sref{IV.5.13.3},~(iii).
As the ${}^h\!\mathrm{B}_\lambda$ are Henselian and Noetherian local rings \sref{IV.18.6.6},~(v) and that $\mathrm{A}={}^h\!\mathrm{A}$ \sref{IV.18.6.6},~(vi), it suffices to apply~(ii) to the inductive system $(\mathrm{B}_\lambda)$.
\end{proof}

\begin{corollary}[18.6.15]
\label{IV.18.6.15}
Let $\mathrm{A}$ be a Henselian local ring, $\mathrm{X}$ an A-prescheme of finite presentation over $\mathrm{A}$.
Then there exists a Noetherian Henselian local ring $\mathrm{A}_0$, a local homomorphism $\mathrm{A}_0\to\mathrm{A}$, an $\mathrm{A}_0$-prescheme $\mathrm{X}_0$ of finite type over $\mathrm{A}_0$ and an A-isomorphism $\mathrm{X}_0\otimes_{\mathrm{A}_0}\mathrm{A}\simto\mathrm{X}$.
\end{corollary}

\begin{proof}
This follows from \sref{IV.18.6.14} and \sref{IV.8.8.2},~(ii).
\end{proof}


\subsection{Henselization and excellent rings.}
\label{subsection:IV.18.7}

\begin{env}[18.7.1]
\label{IV.18.7.1}
In this paragraph we denote by $\boldsymbol{P}(\mathrm{Z},k)$ a property of the form considered in \sref{IV.7.3.1}, where we further suppose that the property $\boldsymbol{Q}(\mathrm{A},k)$ satisfies the following condition:

For every separable algebraic extension $k'$ of $k$, and every local Noetherian $k'$-algebra A, the property $\boldsymbol{Q}(\mathrm{A},k)$ is \emph{equivalent} to $\boldsymbol{Q}(\mathrm{A},k')$.

It is immediate that all the properties $\boldsymbol{P}$ considered in \sref{IV.7.3.8} satisfy the preceding condition, since if K is a finite extension of $k$, $k'\otimes_k\mathrm{K}$ is a finite direct product of separable algebraic extensions of K.
\end{env}

\oldpage[IV]{143}

\begin{proposition}[18.7.2]
\label{IV.18.7.2}
Let $\mathrm{A}$ be a local Noetherian ring.
For ${}^h\!\mathrm{A}$ to be a $\boldsymbol{P}$-ring \sref{IV.7.3.13}, it is necessary and sufficient that A be so.
\end{proposition}

\begin{proof}
Indeed, by virtue of \sref{IV.18.6.6},~(iv), the completion $\hat{\mathrm{A}}$ of A is also the completion of ${}^h\!\mathrm{A}$ and one therefore has $\mathrm{A}\subset{}^h\!\mathrm{A}\subset\hat{\mathrm{A}}$.
By \sref{IV.18.6.9}, for every $x\in\Spec(\mathrm{A})$, the fibre at $x$ of the morphism $\Spec({}^h\!\mathrm{A})\to\Spec(\mathrm{A})$ is discrete and finite, and at each point $y_i$ of this fibre, $k(y_i)$ is a separable algebraic extension of $k(x)$.
The formal fibre of A at $x$ is therefore the prescheme \emph{sum} of the formal fibres of ${}^h\!\mathrm{A}$ at the points $y_i$.
The conclusion then follows from the hypothesis on $\boldsymbol{P}$ in \sref{IV.18.7.1}.
\end{proof}

\begin{corollary}[18.7.3]
\label{IV.18.7.3}
Let $\mathrm{A}$ be a local Noetherian ring.
For ${}^h\!\mathrm{A}$ to be universally Japanese, it is necessary and sufficient that A be so.
\end{corollary}

\begin{proof}
Indeed, for a Noetherian local ring B, it amounts to the same to say that B is universally Japanese, or that the formal fibres of B are geometrically reduced \sref{IV.7.6.4} and \sref{IV.7.7.1};
the conclusion therefore follows from \sref{IV.18.7.2}.
\end{proof}

\begin{corollary}[18.7.4]
\label{IV.18.7.4}
Let $\mathrm{A}$ be a local Noetherian ring.
For the formal fibres of ${}^h\!\mathrm{A}$ to be geometrically regular, it is necessary and sufficient that those of A be so.
\end{corollary}

\begin{proof}
This is a particular case of \sref{IV.18.7.2}.
\end{proof}

\begin{proposition}[18.7.5]
\label{IV.18.7.5}
Let $\mathrm{A}$ be a local Noetherian ring.
If A is universally catenary \sref{IV.5.6.2} (resp. formally catenary \sref{IV.7.1.9}, resp. strictly formally catenary \sref{IV.7.2.6}), so is ${}^h\!\mathrm{A}$.
\end{proposition}

\begin{proof}
With the notations of \sref{IV.18.6.4}, if A is universally catenary, the same holds for the $\mathrm{A}_\alpha$, which are A-algebras essentially of finite type \sref{IV.5.6.3}.
To prove that in this case ${}^h\!\mathrm{A}$ is universally catenary, it will therefore suffice to establish the following lemma:
\end{proof}

\begin{lemma}[18.7.5.1]
\label{IV.18.7.5.1}
Let $(\mathrm{B}_\alpha,\vphi_{\beta\alpha})$ be an inductive system of Noetherian rings and suppose that, for $\alpha\leq\beta$, $\vphi_{\beta\alpha}$ makes $\mathrm{B}_\beta$ a faithfully flat $\mathrm{B}_\alpha$-module and that $\mathrm{B}=\varinjlim\mathrm{B}_\alpha$ is Noetherian.
Then, if the $\mathrm{B}_\alpha$ are catenary (resp. universally catenary), the same holds for $\mathrm{B}$.
\end{lemma}

\begin{proof}
To say that B is universally catenary is equivalent to saying that $\mathrm{B}[\mathrm{T}_1,\ldots,\mathrm{T}_n]$ is catenary for all $n$ \sref{IV.5.6.2}, and it is immediate that the inductive system $(\mathrm{B}_\alpha[\mathrm{T}_1,\ldots,\mathrm{T}_n])$ satisfies the same conditions as $(\mathrm{B}_\alpha)$;
it therefore suffices to prove that if the $\mathrm{B}_\alpha$ are catenary, then B is catenary.
If $\mathfrak{p}_0\supset\mathfrak{p}_1\supset\cdots\supset\mathfrak{p}_r$ is a saturated chain of prime ideals, the inverse images $\mathfrak{p}_{0\alpha}\supset\mathfrak{p}_{1\alpha}\supset\cdots\supset\mathfrak{p}_{r\alpha}$ of these ideals in $\mathrm{B}_\alpha$ form a chain of prime ideals, and it suffices to show that for $\alpha$ sufficiently large, this chain is saturated.
Now, each $\mathfrak{p}_i$ is an ideal of finite type, hence there exists $\lambda$ such that, for $\alpha\geq\lambda$, one has $\mathfrak{p}_i=\mathfrak{p}_{i\alpha}.\mathrm{B}$.
One therefore has $1=\codim(\mathrm{V}(\mathfrak{p}_i),\mathrm{V}(\mathfrak{p}_{i+1}))=\codim(\mathrm{V}(\mathfrak{p}_{i\alpha}),\mathrm{V}(\mathfrak{p}_{i+1,\alpha}))$ since B is a flat $\mathrm{B}_\alpha$-module \sref{IV.6.1.4} and \sref[0]{0.6.2.3}, which completes the proof of the lemma.

Suppose now A formally catenary.
Let $\mathfrak{p}$ be a prime ideal of ${}^h\!\mathrm{A}$, and let $\mathfrak{q}$ be its inverse image in A, so that $\mathfrak{q}.{}^h\!\mathrm{A}\subset\mathfrak{p}$ and consequently ${}^h\!\mathrm{A}/\mathfrak{p}$ is a quotient ring of ${}^h\!\mathrm{A}/\mathfrak{q}.{}^h\!\mathrm{A}$.  It therefore suffices to prove that ${}^h\!\mathrm{A}/\mathfrak{q}.{}^h\!\mathrm{A}$ is formally catenary for every prime ideal $\mathfrak{q}$ of A \sref{IV.7.1.9}.  Since ${}^h\!\mathrm{A}/\mathfrak{q}.{}^h\!\mathrm{A}={}^h(\mathrm{A}/\mathfrak{q})$ by \sref{IV.18.6.8}, it suffices \sref{IV.7.1.11} to show that if A is \emph{integral} and formally catenary, ${}^h\!\mathrm{A}$ is formally \emph{equidimensional}.
But as the completion of ${}^h\!\mathrm{A}$ is equal to $\hat{\mathrm{A}}$, this follows from the hypothesis that A is formally catenary.

\oldpage[IV]{144}

Finally, suppose that A is strictly formally catenary.  We have just seen that ${}^h\!\mathrm{A}$ is formally catenary, and it remains to prove that the fibres of the morphism $\Spec(\hat{\mathrm{A}})\to\Spec({}^h\!\mathrm{A})$ satisfy property $(\mathrm{S}_1)$ \sref{IV.7.2.5},~$b)$.  This follows from the hypothesis on A and from \sref{IV.18.7.2}, applied to the case where $\boldsymbol{P}(\mathrm{Z},k)$ is property $(\mathrm{S}_1)$ for Z.
\end{proof}

\begin{corollary}[18.7.6]
\label{IV.18.7.6}
If a local Noetherian ring A is excellent \sref{IV.7.8.2}, so is ${}^h\!\mathrm{A}$.
\end{corollary}

\begin{remark}[18.7.7]
\label{IV.18.7.7}
It can happen that ${}^h\!\mathrm{A}$ is an excellent local ring without A being universally catenary (nor \emph{a fortiori} excellent).
To see this, let us take the example of A from \sref{IV.5.6.11}, with the same notations.
The ring E constructed in \sref{IV.5.6.11} is excellent when $k_0$ is of characteristic 0 \sref{IV.7.8.3},~(ii) and~(iii);
consider two rings $\mathrm{E}_1$, $\mathrm{E}_2$ isomorphic to E, and ``glue'' $\Spec(\mathrm{E}_1)$ and $\Spec(\mathrm{E}_2)$ such that if $\mathfrak{m}_i$ and $\mathfrak{m}'_i$ are the maximal ideals of $\mathrm{E}_i$ ($i=1,2$), corresponding to $\mathfrak{m}$ and $\mathfrak{m}'$, the point $\mathfrak{m}_1$ is ``glued'' to $\mathfrak{m}'_2$ and the point $\mathfrak{m}_2$ to $\mathfrak{m}'_1$;
to be precise, if $\varepsilon_i$ and $\varepsilon'_i$ are the canonical homomorphisms from $\mathrm{E}_i$ onto $k(\mathfrak{m}_i)$ and $k(\mathfrak{m}'_i)$ ($i=1,2$), the scheme obtained is $\Spec(\mathrm{R})$, where R is the subring of $\mathrm{E}_1\times\mathrm{E}_2$ formed by the pairs $(x_1,x_2)$ such that $\varepsilon'_2(x_2)=\sigma(\varepsilon_1(x_1))$ and $\varepsilon'_1(x_1)=\sigma(\varepsilon_2(x_2))$.
One verifies easily (for example with the help of \sref{IV.17.6.3}) that R is a finite \'{e}tale C-algebra, and that there are two maximal ideals $\mathfrak{r}_1$, $\mathfrak{r}_2$ of R above the maximal ideal $\mathfrak{n}$ of C, the completions of $\mathrm{R}_{\mathfrak{r}_1}$ and $\mathrm{R}_{\mathfrak{r}_2}$ being canonically isomorphic to that of $\mathrm{A}=\mathrm{C}_{\mathfrak{n}}$.  Thus $\mathrm{R}_{\mathfrak{r}_1}$ is an A-algebra that is strictly essentially \'{e}tale, and consequently its Henselization is equal to ${}^h\!\mathrm{A}$.
Furthermore, we saw in \sref{IV.7.8.4},~(ii), that the formal fibres of A are geometrically regular; the same therefore holds for those of ${}^h\!\mathrm{A}$ by \sref{IV.18.7.2}.  Finally, R is universally catenary, since its quotients by its two minimal prime ideals are isomorphic to E, the conclusion following from \sref{IV.5.6.3},~(iii).  Hence $\mathrm{R}_{\mathfrak{r}_1}$ is excellent, and consequently so is ${}^h\!\mathrm{A}$ by \sref{IV.18.7.6}, whereas A is not universally catenary.
\end{remark}


\subsection{Strictly local rings and strict Henselization.}
\label{subsection:IV.18.8}

\begin{proposition}[18.8.1]
\label{IV.18.8.1}
Let $\mathrm{A}$ be a local ring.
The following conditions are equivalent:
\begin{enumerate}[label=\emph{\alph*)}]
  \item A is a Henselian ring and its residue field $k$ is separably closed.
  \item A is a Henselian ring and every \'{e}tale covering of $\mathrm{S}=\Spec(\mathrm{A})$ is trivial.
  \item For every \'{e}tale morphism $f:\mathrm{X}\to\mathrm{S}$ and every point $x\in\mathrm{X}$ such that $f(x)$ is the closed point $s$ of S, there exists an S-section $u:\mathrm{S}\to\mathrm{X}$ of X such that $u(s)=x$.
\end{enumerate}
\end{proposition}

\begin{proof}
The hypothesis $c)$ implies that for every \'{e}tale morphism $f:\mathrm{X}\to\mathrm{S}$, the residue fields at the points of $f^{-1}(s)$ are all isomorphic to $k$, and moreover the condition $b)$ of \sref{IV.18.5.11} is satisfied.
Hence A is Henselian, and the fact that every \'{e}tale covering of S is trivial follows from \sref{IV.18.2.10},~(ii);
hence $c)$ implies $b)$.
As \'{e}tale coverings of $\Spec(k)$ are trivial if and only if $k$ is separably closed, $b)$ implies $a)$ by virtue of \sref{IV.18.5.11},~$b$).
Finally, it follows also from \sref{IV.18.5.11},~$b$) that $a)$ implies $c)$.
\end{proof}

\begin{definition}[18.8.2]
\label{IV.18.8.2}
A local ring is said to be \emph{strictly local} if it satisfies the equivalent conditions of \sref{IV.18.8.1}.
A scheme isomorphic to the spectrum of a strictly local ring is called a \emph{strictly local scheme}.
\end{definition}

\begin{remark}[18.8.3]
\label{IV.18.8.3}
The conditions of \sref{IV.18.8.1} are also equivalent to the following:
\begin{enumerate}
  \item[d)] For every morphism locally of finite type $f:\mathrm{X}\to\mathrm{S}$ and every point $x\in\mathrm{X}$ such that $f(x)$ is the closed point $s$ of S and that $\sh{O}_{\mathrm{X},x}/\mathfrak{m}_s\sh{O}_{\mathrm{X},x}$ is a field, a finite separable extension of $k(s)=k$, there exists an open neighbourhood U of $x$ in X such that $f|U$ is a closed immersion.
\end{enumerate}
\end{remark}

\oldpage[IV]{145}

(One will note that the hypothesis of $d)$ is satisfied if $f$ is \emph{formally unramified} at the point $x$ \sref{IV.17.4.1.2}.)

We leave the reader the proof, which is substantially the same as that of \sref{IV.18.5.18}.

\begin{lemma}[18.8.4]
\label{IV.18.8.4}
Let $\mathrm{A}$, $\mathrm{A}'$ be two local rings, $\vphi:\mathrm{A}\to\mathrm{A}'$ a local homomorphism making $\mathrm{A}'$ an A-algebra essentially \'{e}tale \sref{IV.18.6.1}.
Then, for every strictly local ring $\mathrm{B}$, every local homomorphism $\psi:\mathrm{A}\to\mathrm{B}$ and every homomorphism of $k(\mathrm{A})$-algebras $\alpha:k(\mathrm{A}')\to k(\mathrm{B})$, there exists one and only one local homomorphism $\psi':\mathrm{A}'\to\mathrm{B}$ such that $\psi=\psi'\circ\vphi$ and that $\alpha$ is equal to the homomorphism $\overline{\psi'}$ deduced from $\psi'$ by passage to quotients.
\end{lemma}

\begin{proof}
With the notations of \sref{IV.18.6.2}, the homomorphism $\alpha$ corresponds to a $k(s)$-morphism $\Spec(k(y))\to\Spec(k(x))$, hence to a point $x'$ well determined in $\mathrm{X}'$ above $y$.
The existence of a Y-section of $f':\mathrm{X}'\to\mathrm{Y}$ taking $y$ to $x'$ follows from \sref{IV.18.8.1},~$c$), since $f'$ is \'{e}tale and B is strictly local; its uniqueness follows from the fact that $f'$ is separated and from \sref{IV.17.4.9}.
\end{proof}

\begin{lemma}[18.8.5]
\label{IV.18.8.5}
Let $\mathrm{A}$ be a local ring, $\mathrm{A}_1$, $\mathrm{A}_2$ two A-algebras that are local and essentially \'{e}tale, K a field that is an extension of $k(\mathrm{A})$.
\begin{enumerate}[label=\emph{\roman*)}]
  \item For every $k(\mathrm{A})$-homomorphism $\gamma:k(\mathrm{A}_1)\to k(\mathrm{A}_2)$, there exists at most one A-homomorphism (necessarily local) $\psi:\mathrm{A}_1\to\mathrm{A}_2$ such that $\gamma$ is the homomorphism $\overline{\psi}$ deduced from $\psi$ by passage to quotients.
  \item For every pair of local A-homomorphisms $\beta_1:\mathrm{A}_1\to\mathrm{K}$, $\beta_2:\mathrm{A}_2\to\mathrm{K}$ there exists a local A-algebra $\mathrm{A}_3$ that is essentially \'{e}tale, two A-homomorphisms $\vphi_1:\mathrm{A}_1\to\mathrm{A}_3$, $\vphi_2:\mathrm{A}_2\to\mathrm{A}_3$ and an A-homomorphism $\beta:\mathrm{A}_3\to\mathrm{K}$ such that $\beta_1=\beta\circ\vphi_1$ and $\beta_2=\beta\circ\vphi_2$.
\end{enumerate}
\end{lemma}

\begin{proof}
With the notations of \sref{IV.18.6.3}, the homomorphisms $\beta_1$, $\beta_2$ in~(ii) correspond to two S-morphisms $\Spec(\mathrm{K})\to\mathrm{X}_1$, $\Spec(\mathrm{K})\to\mathrm{X}_2$, with respective images $x_1$, $x_2$;
one thus deduces a well-determined S-morphism $\Spec(\mathrm{K})\to\mathrm{X}_3$ \sref[I]{I.3.2.1} whose image $x_3$ lies above $x_1$ and $x_2$;
the A-algebra $\mathrm{A}_3=\sh{O}_{\mathrm{X}_3,x_3}$ and the homomorphism $\mathrm{A}_3\to\mathrm{K}$ corresponding to these data then satisfy the conditions of~(ii).
On the other hand, the data of $\gamma$ in~(i) corresponds to an S-morphism $\Spec(k(x_2))\to\Spec(k(x_1))$ or also to a point $x'_3$ of $\mathrm{X}'_3$ above $x_2$;
the uniqueness of $\psi$ follows from the uniqueness of an $\mathrm{X}_2$-section of $\mathrm{X}'_3$ passing through $x'_3$ \sref{IV.17.4.9}.
\end{proof}

\begin{env}[18.8.6]
\label{IV.18.8.6}
Let A be a local ring, $\mathfrak{m}$ its maximal ideal, $k$ its residue field, $\Omega$ a separable closure of $k$.
Consider the set $\mathfrak{E}$ of essentially \'{e}tale A-algebras defined in \sref{IV.18.6.1}, and denote by L the set of \emph{local A-homomorphisms} $\mathrm{A}'\to\Omega$, where $\mathrm{A}'$ ranges over $\mathfrak{E}$;
one will note that such a homomorphism factorises as $\lambda:\mathrm{A}'\to k(\mathrm{A}')\xrightarrow{\omega_\lambda}\Omega$, where $\omega_\lambda:k(\mathrm{A}')\to\Omega$ is a $k$-homomorphism (since $\mathfrak{m}\mathrm{A}'$ is the maximal ideal of $\mathrm{A}'$);
conversely the data of a homomorphism $\omega_\lambda$ uniquely determines a homomorphism $\lambda\in\mathrm{L}$.
For every $\lambda\in\mathrm{L}$, we denote by $\mathrm{A}_\lambda$ the A-algebra essentially \'{e}tale belonging to $\mathfrak{E}$ on which $\lambda$ is defined.
The set L is \emph{preordered} by the relation ``there exists an A-homomorphism $\vphi_{\mu\lambda}:\mathrm{A}_\lambda\to\mathrm{A}_\mu$ such that $\lambda=\mu\circ\vphi_{\mu\lambda}$'', and it follows from \sref{IV.18.8.5},~(i) that this homomorphism $\vphi_{\mu\lambda}$ is \emph{unique}.
Furthermore, L is \emph{filtered increasing} for the preceding preorder relation (which one also denotes $\lambda\leq\mu$), by virtue
\oldpage[IV]{146}
of \sref{IV.18.8.5},~(ii).
It is clear that $(\mathrm{A}_\lambda,\vphi_{\mu\lambda})$ is an \emph{inductive system} of local A-algebras, the $\vphi_{\mu\lambda}$ being local homomorphisms;
furthermore, it follows from \sref{IV.17.3.5} that for $\lambda\leq\mu$, $\mathrm{A}_\mu$ is an $\mathrm{A}_\lambda$-algebra that is essentially \'{e}tale.
If $\overline{\vphi}_{\mu\lambda}:k(\mathrm{A}_\lambda)\to k(\mathrm{A}_\mu)$ is the $k$-homomorphism deduced from $\vphi_{\mu\lambda}$ by passage to quotients, $(k(\mathrm{A}_\lambda),\overline{\vphi}_{\mu\lambda})$ is an inductive system of separable algebraic extensions of $k$, and the $\omega_\lambda:k(\mathrm{A}_\lambda)\to\Omega$ form an inductive system of homomorphisms of $k$.
Furthermore, the inductive limit
\begin{equation}
  \omega:\varinjlim k(\mathrm{A}_\lambda)\to\Omega
  \label{IV.18.8.6.1}
\tag{18.8.6.1}
\end{equation}
is a $k$\emph{-isomorphism}.
It suffices indeed to prove that for every finite separable extension $k'$ of $k$, there exists an A-algebra $\mathrm{A}'$ that is essentially \'{e}tale, such that $k'$ is the residue field of $\mathrm{A}'$ and that the homomorphism $k\to k'$ is deduced from $\mathrm{A}\to\mathrm{A}'$ by passage to quotients.
But this follows from \sref{IV.18.1.1} applied to \'{e}tale morphisms, since $\Spec(\mathrm{A})$ is the unique neighbourhood of its closed point.

Note now that if $\mathfrak{m}_\lambda$ is the maximal ideal of $\mathrm{A}_\lambda$, it follows from \sref{IV.17.6.1} that, for $\lambda\leq\mu$, $\mathfrak{m}_\mu=\mathfrak{m}_\lambda\mathrm{A}_\mu$ and that $\mathrm{A}_\mu$ is a flat $\mathrm{A}_\lambda$-module.
Hence it follows from \sref[0]{0.10.3.1.3} that the ring $\varinjlim\mathrm{A}_\lambda$ is \emph{local}, that the homomorphism $w:\mathrm{A}\to\varinjlim\mathrm{A}_\lambda$ is local, and that one has a canonical $k$-isomorphism
\[
  \varinjlim k(\mathrm{A}_\lambda)\xrightarrow{\sim}k(\varinjlim\mathrm{A}_\lambda);
\]
identifying these two fields, one deduces a canonical $k$-isomorphism
\begin{equation}
  \omega^{-1}:\Omega\xrightarrow{\sim}k(\varinjlim\mathrm{A}_\lambda).
  \label{IV.18.8.6.2}
\tag{18.8.6.2}
\end{equation}
\end{env}

\begin{definition}[18.8.7]
\label{IV.18.8.7}
Let $\mathrm{A}$ be a local ring, $k=k(\mathrm{A})$ its residue field, $i:k\to\Omega$ a homomorphism from $k$ into a separable closure $\Omega$ of $k$;
the \emph{strict Henselization} of $\mathrm{A}$ relative to $i$, denoted ${}^{hs}\!\mathrm{A}_{(i)}$ (or ${}^{hs}\!\mathrm{A}$ when this causes no confusion), is the inductive limit of the inductive system $(\mathrm{A}_\lambda,\vphi_{\mu\lambda})$ defined in \sref{IV.18.8.6}.
\end{definition}

As one has a $k$-isomorphism $\omega^{-1}:\Omega\xrightarrow{\sim}k({}^{hs}\!\mathrm{A}_{(i)})$, $\Omega$ is endowed canonically with the structure of a ${}^{hs}\!\mathrm{A}_{(i)}$-algebra by the local homomorphism
\[
  {}^{hs}\!\mathrm{A}_{(i)}\to k({}^{hs}\!\mathrm{A}_{(i)})\xrightarrow{\omega}\Omega.
\]

As in \sref{IV.18.6.5} one sees that the definition given in \sref{IV.18.8.7} does not depend on the choice of $\mathfrak{E}$;
we shall moreover see below \sref{IV.18.8.8},~(i) and~(ii) that ${}^{hs}\!\mathrm{A}_{(i)}$ is an object representing a functor entirely defined by the data of A and $i$ \sref[0]{0.8.1.8}, and is hence defined as such up to a unique isomorphism.

\begin{proposition}[18.8.8]
\label{IV.18.8.8}
Let $\mathrm{A}$ be a local ring, $k$ its residue field, $i:k\to\Omega$ a homomorphism from $k$ into a separable closure $\Omega$ of $k$, ${}^{hs}\!\mathrm{A}_{(i)}$ the strict Henselization of $\mathrm{A}$ corresponding to $i$.
\begin{enumerate}[label=\emph{\roman*)}]
  \item ${}^{hs}\!\mathrm{A}_{(i)}$ is a strictly local ring \sref{IV.18.8.2} and the structural homomorphism $w:\mathrm{A}\to{}^{hs}\!\mathrm{A}_{(i)}$ is local.
  \item For every local homomorphism $u:\mathrm{A}\to\mathrm{B}$, where $\mathrm{B}$ is a strictly local ring, and every $k$-homomorphism $\psi:\Omega\to k(\mathrm{B})$, there exists one and only one A-homomorphism $v:{}^{hs}\!\mathrm{A}_{(i)}\to\mathrm{B}$ such that $\psi$ factorises as $\Omega\xrightarrow{\omega^{-1}}k({}^{hs}\!\mathrm{A}_{(i)})\xrightarrow{\overline{v}}k(\mathrm{B})$, where $\overline{v}$ is deduced from $v$ by passage to quotients.
  \oldpage[IV]{147}
  \item ${}^{hs}\!\mathrm{A}_{(i)}$ is a faithfully flat A-module, and if $\mathfrak{m}$ is the maximal ideal of A, $\mathfrak{m}.{}^{hs}\!\mathrm{A}_{(i)}$ is the maximal ideal of ${}^{hs}\!\mathrm{A}_{(i)}$, and $k({}^{hs}\!\mathrm{A}_{(i)})$ is a separable closure of $k$.
  \item For ${}^{hs}\!\mathrm{A}_{(i)}$ to be Noetherian, it is necessary and sufficient that A be so.
  \item If A is strictly local (hence $\Omega=k$), one has ${}^{hs}\!\mathrm{A}\cong\mathrm{A}$.
  \item If $i':k\to\Omega'$ is a second homomorphism from $k$ into a separable closure of $k$, for every $k$-isomorphism $\sigma:\Omega\xrightarrow{\sim}\Omega'$, the corresponding A-homomorphism (by~(ii)) $v_\sigma:{}^{hs}\!\mathrm{A}_{(i)}\to{}^{hs}\!\mathrm{A}_{(i')}$ is an isomorphism.
\end{enumerate}
\end{proposition}

\begin{proof}
We have already seen above that ${}^{hs}\!\mathrm{A}_{(i)}$ is a local ring and that $w$ is a local homomorphism;
the fact that ${}^{hs}\!\mathrm{A}_{(i)}$ is Henselian (and hence strictly local) is demonstrated as in \sref{IV.18.6.6}.
The assertions of~(iii), as well as the sufficiency of condition~(iv), also follow from \sref[0]{0.10.3.1.3}; the necessity of condition~(iv) follows from the fact that ${}^{hs}\!\mathrm{A}_{(i)}$ is a faithfully flat A-module \sref[0]{0.6.5.2}.

To prove~(ii), remark, with the notations of \sref{IV.18.8.6}, that for every $\lambda\in\mathrm{L}$, there exists one and only one local homomorphism $v_\lambda:\mathrm{A}_\lambda\to\mathrm{B}$, such that the composite $k(\mathrm{A}_\lambda)\xrightarrow{\omega_\lambda}\Omega\xrightarrow{\psi}k(\mathrm{B})$ is deduced from $v_\lambda$ by passage to quotients, by virtue of \sref{IV.18.8.4} and the hypothesis that B is strictly local.
The uniqueness of $v_\lambda$ implies moreover that the $v_\lambda$ form an inductive system, whence the existence of the homomorphism $v$ satisfying the conditions of~(ii); its uniqueness follows from \sref{IV.17.4.9}.
The uniqueness of $v$ immediately implies assertion~(vi), by considering the composites $v_\sigma\circ v_{\sigma^{-1}}$ and $v_{\sigma^{-1}}\circ v_\sigma$.
Finally, if A is strictly local, and $\mathrm{A}'=\mathrm{B}_{\mathfrak{n}}$ is an A-algebra local and essentially \'{e}tale, where B is an A-algebra \'{e}tale, it follows from \sref{IV.18.8.1} that there exists a $\Spec(\mathrm{A})$-section of $\Spec(\mathrm{B})$ taking the value $\mathfrak{n}$ at the closed point of $\Spec(\mathrm{A})$, hence by \sref{IV.17.4.1} that the homomorphism $\mathrm{A}\to\mathrm{A}'$ is bijective, which proves~(v).
\end{proof}

\begin{env}[18.8.8.1]
\label{IV.18.8.8.1}
Let $\mathcal{C}$ be the category of strictly local rings, with local homomorphisms as morphisms;
for every $\mathrm{B}\in\mathcal{C}$, denote by $\mathrm{F}(\mathrm{B})$ the set of pairs $(u,\psi)$ consisting of a local homomorphism $u:\mathrm{A}\to\mathrm{B}$ and a $k$-homomorphism $\psi:\Omega\to k(\mathrm{B})$ such that the composite $k(\mathrm{A})=k\xrightarrow{i}\Omega\xrightarrow{\psi}k(\mathrm{B})$ is deduced from $u$ by passage to quotients.
One can say that the object ${}^{hs}\!\mathrm{A}_{(i)}$ \emph{represents} the covariant functor $\mathrm{F}:\mathcal{C}\to\cat{Ens}$ \sref[0]{0.8.1.11}.

One will remark that by virtue of \sref{IV.18.8.8},~(vi), the strict Henselizations ${}^{hs}\!\mathrm{A}_{(i)}$ are all A-isomorphic (for the various $k$-homomorphisms $i$ of $k$ into separable closures of $k$), but for given $i$ and $i'$, there is in general an infinity of A-isomorphisms ${}^{hs}\!\mathrm{A}_{(i)}\xrightarrow{\sim}{}^{hs}\!\mathrm{A}_{(i')}$.
To be precise, the group of A-automorphisms of ${}^{hs}\!\mathrm{A}_{(i)}$ is isomorphic to the Galois group of $\Omega$ over $k$.
\end{env}

\begin{env}[18.8.9]
\label{IV.18.8.9}
Let now A be a \emph{semi-local} ring, $\mathfrak{m}_j$ ($1\leq j\leq r$) its maximal ideals, and for every $j$, consider a homomorphism $i_j:k(\mathrm{A}_{\mathfrak{m}_j})\to\Omega_j$ of the residue field into a separable closure.
The \emph{strict Henselization} of A relative to the $i_j$ is the product $\prod_{j=1}^{r}({}^{hs}(\mathrm{A}_{\mathfrak{m}_j}))$, denoted ${}^{hs}\!\mathrm{A}$ (when there is no risk of confusion).
\oldpage[IV]{148}
It is therefore a faithfully flat A-module and a semi-local A-algebra, whose maximal ideals are the $\mathfrak{m}_j.{}^{hs}\!\mathrm{A}$, while its radical is $\mathfrak{r}.{}^{hs}\!\mathrm{A}$ (if $\mathfrak{r}$ denotes the radical of A).
Finally, the universal property \sref{IV.18.8.8},~(ii) still holds when replacing ``local homomorphism'' by ``semi-local homomorphism'' \sref{IV.18.6.7}, and replacing $\Omega$ by $\prod_j\Omega_j$.
\end{env}

\begin{proposition}[18.8.10]
\label{IV.18.8.10}
Let $\mathrm{A}$ be a semi-local ring, $\mathrm{B}$ a finite A-algebra.
Let $\mathfrak{n}_j$ ($1\leq j\leq r$) be the maximal ideals of $\mathrm{B}$.
Then there is, for every $j$, a maximal ideal $\mathfrak{q}_j$ of $\mathrm{B}\otimes_{\mathrm{A}}({}^{hs}\!\mathrm{A})$ above $\mathfrak{n}_j$, such that ${}^{hs}\!\mathrm{B}$ is isomorphic to the direct product of the local rings $(\mathrm{B}\otimes_{\mathrm{A}}({}^{hs}\!\mathrm{A}))_{\mathfrak{q}_j}$ ($1\leq j\leq r$).
\end{proposition}

\begin{proof}
One can evidently restrict to the case where A is local, replacing A by $\mathrm{A}_{\mathfrak{m}_j}$, where $\mathfrak{m}_j$ is the maximal ideal of A that is the inverse image of $\mathfrak{n}_j$.
On the other hand, let $\mathrm{C}=\mathrm{B}_{\mathfrak{n}_j}$;
by virtue of the definition of ${}^{hs}\!\mathrm{B}$ \sref{IV.18.8.9}, everything amounts to seeing that ${}^{hs}\!\mathrm{C}$ is C-isomorphic to $(\mathrm{C}\otimes_{\mathrm{A}}({}^{hs}\!\mathrm{A}))_{\mathfrak{q}}$, where $\mathfrak{q}$ is one of the maximal ideals of $\mathrm{C}\otimes_{\mathrm{A}}({}^{hs}\!\mathrm{A})$.
Let $k$ and K be the residue fields of A and C, $\Omega$ the separable closure of $k$, and let us suppose K and $\Omega$ contained in the same closure of $k$;
as K is a finite extension of $k$, it is immediate that $\mathrm{K}\otimes_k\Omega$ is a finite direct product of fields, all isomorphic to the separable closure $\mathrm{K}\Omega$ of K.
On the other hand, since ${}^{hs}\!\mathrm{A}$ is Henselian, $\mathrm{C}\otimes_{\mathrm{A}}({}^{hs}\!\mathrm{A})$ is a direct product of the local rings $(\mathrm{C}\otimes_{\mathrm{A}}({}^{hs}\!\mathrm{A}))_{\mathfrak{q}_i}$, where $\mathfrak{q}_i$ ($1\leq i\leq s$) ranges over the maximal ideals of $\mathrm{C}\otimes_{\mathrm{A}}({}^{hs}\!\mathrm{A})$; these local rings are Henselian \sref{IV.18.5.10} and have $\mathrm{K}\Omega$ as residue field, hence they are strictly local \sref{IV.18.8.1}.
Let $\mathrm{C}'=\mathrm{C}\otimes_{\mathrm{A}}({}^{hs}\!\mathrm{A})$ and $\mathrm{C}'_i=\mathrm{C}'_{\mathfrak{q}_i}$, which is therefore a quotient ring of $\mathrm{C}'$;
let $f:\mathrm{C}\to\mathrm{C}'$, $g:{}^{hs}\!\mathrm{A}\to\mathrm{C}'$, $\vphi_i:\mathrm{C}'\to\mathrm{C}'_i$ be the canonical maps.
\[
  \xymatrix{
    &&& {}^{hs}\!\mathrm{C} \ar@<0.8ex>[dl]_{w'} \\
    && \mathrm{C}'_i \ar@<0.8ex>[ur]_{w} \\
    \mathrm{C} \ar[uurrr]^{v} \ar[r]_{f} & \mathrm{C}' \ar[ur] \\
    \mathrm{A} \ar[u] \ar[r] & {}^{hs}\!\mathrm{A} \ar[u]^{g} \ar[uuurr]_{u}
  }
\]
Let $v:\mathrm{C}\to{}^{hs}\!\mathrm{C}$ be the canonical local homomorphism;
there exists a unique local homomorphism $u:{}^{hs}\!\mathrm{A}\to{}^{hs}\!\mathrm{C}$ which, by passage to quotients, gives the canonical injection $\Omega\to\mathrm{K}\Omega$ and such that the composite $\mathrm{A}\to\mathrm{C}\xrightarrow{v}{}^{hs}\!\mathrm{C}$ is equal to the composite $\mathrm{A}\to{}^{hs}\!\mathrm{A}\xrightarrow{u}{}^{hs}\!\mathrm{C}$ \sref{IV.18.8.8};
hence there exists a unique local homomorphism $w_0:\mathrm{C}'\to{}^{hs}\!\mathrm{C}$ such that $w_0\circ g=u$, $w_0\circ f=v$, and since ${}^{hs}\!\mathrm{C}$ is a local ring, this homomorphism factorises as $\mathrm{C}'\to\mathrm{C}'_i\xrightarrow{w}{}^{hs}\!\mathrm{C}$ for a well-determined index $i$, $w$ being a local homomorphism.
On the other hand, the ring $\mathrm{C}'_i$ being strictly local, there exists a local homomorphism $w':{}^{hs}\!\mathrm{C}\to\mathrm{C}'_i$ which, by passage to quotients, gives the identity automorphism of $\mathrm{K}\Omega$ and which is such that $w'\circ v=\vphi_i\circ f$ \sref{IV.18.8.8}.
One deduces first that $w\circ w'$ is an endomorphism of ${}^{hs}\!\mathrm{C}$ which, by passage to quotients, gives the identity automorphism of $\mathrm{K}\Omega$, hence
\oldpage[IV]{149}
is the identity \sref{IV.18.8.8}.
On the other hand, $w'\circ w\circ\vphi_i\circ f=w'\circ v=\vphi_i\circ f$ and $w'\circ w\circ\vphi_i\circ g=w'\circ u=\vphi_i\circ g$;
whence $w'\circ w\circ\vphi_i=\vphi_i$, which is to say $w'\circ w$ is the identity automorphism of $\mathrm{C}'_i$.
\end{proof}

\begin{remark}[18.8.11]
\label{IV.18.8.11}
The proof of \sref{IV.18.8.10} uses the fact that B is integral over A, but uses the fact that B is a \emph{finite} A-algebra only to establish that $\mathrm{K}\otimes_k\Omega$ is a finite direct product of fields.
This latter property still holds when B is assumed to be an \emph{integral semi-local} A-algebra whose maximal ideals $\mathfrak{n}_j$ have finite separable residue-field degrees $[k(\mathfrak{n}_j):k]_s$.
\end{remark}

\begin{proposition}[18.8.12]
\label{IV.18.8.12}
Let $\mathrm{A}$ be a semi-local ring.
\begin{enumerate}[label=\emph{\roman*)}]
  \item For ${}^{hs}\!\mathrm{A}$ to be reduced (resp. normal), it is necessary and sufficient that A be so.
  \item Suppose A Noetherian.
  Then, for every prime ideal $\mathfrak{p}$ of A, the ring $({}^{hs}\!\mathrm{A})_{\mathfrak{p}}/\mathfrak{p}.({}^{hs}\!\mathrm{A})_{\mathfrak{p}}$ is a finite direct product of separable algebraic extensions of $k(\mathfrak{p})$ (which implies that the fibres of the canonical morphism $\Spec({}^{hs}\!\mathrm{A})\to\Spec(\mathrm{A})$ are \emph{geometrically regular} and are \emph{discrete} spaces).
\end{enumerate}
\end{proposition}

\begin{corollary}[18.8.13]
\label{IV.18.8.13}
Let $\mathrm{A}$ be a local Noetherian ring.
For ${}^{hs}\!\mathrm{A}$ to possess one of the following properties:
\begin{enumerate}[label=\emph{\alph*)}]
  \item being a Cohen-Macaulay ring \sref[0]{0.16.5.3};
  \item satisfying the condition $(\mathrm{S}_n)$ \sref{IV.5.7.3};
  \item being regular;
  \item satisfying the condition $(\mathrm{R}_n)$ \sref{IV.5.8.2};
\end{enumerate}
it is necessary and sufficient that A possess the same property.
\end{corollary}

\begin{corollary}[18.8.14]
\label{IV.18.8.14}
Let $\mathrm{A}$ be a semi-local Noetherian ring.
For a prime ideal $\mathfrak{p}$ of ${}^{hs}\!\mathrm{A}$ to belong to $\operatorname{Ass}({}^{hs}\!\mathrm{A})$, it is necessary and sufficient that $\mathfrak{p}\cap\mathrm{A}$ belong to $\operatorname{Ass}(\mathrm{A})$.
\end{corollary}

\begin{proof}
The proofs are the same as those of \sref{IV.18.6.9},~\sref{IV.18.6.10} and \sref{IV.18.6.11} respectively.
\end{proof}

\begin{proposition}[18.8.15]
\label{IV.18.8.15}
Let $\mathrm{A}$ be a local ring;
the following properties are equivalent:
\begin{enumerate}[label=\emph{\alph*)}]
  \item A is geometrically unibranch (resp. reduced and geometrically unibranch).
  \item For every A-algebra $\mathrm{A}'$ that is local and essentially \'{e}tale, $\Spec(\mathrm{A}')$ is irreducible (resp. integral).
  \item $\Spec({}^{hs}\!\mathrm{A})$ is irreducible (resp. integral).
\end{enumerate}
\end{proposition}

\begin{proof}
As in \sref{IV.18.6.12}, one reduces to the case where A is reduced, using the relation ${}^{hs}(\mathrm{A}_{\mathrm{red}})=({}^{hs}\!\mathrm{A})\otimes_{\mathrm{A}}\mathrm{A}_{\mathrm{red}}$ \sref{IV.18.8.11}; the equivalence of $b)$ and $c)$ is proved as in \sref{IV.18.6.12}.  The same is true of the implication $a)\Rightarrow c)$, taking account of \sref{IV.18.8.11} and of the definition of a geometrically unibranch integral local ring.
Finally, to prove that $c)$ implies $a)$, using the same notations as in \sref{IV.18.6.12}, it suffices to see that ${}^{hs}\!\mathrm{B}$ is integral.
But, by \sref{IV.18.8.10}, ${}^{hs}\!\mathrm{B}$ is a localization of the semi-local ring $\mathrm{B}\otimes_{\mathrm{A}}({}^{hs}\!\mathrm{A})$.  By flatness, the latter again identifies with a subring of the field L and is therefore integral; consequently ${}^{hs}\!\mathrm{B}$ is integral as well, which completes the proof.
\end{proof}

\begin{corollary}[18.8.16]
\label{IV.18.8.16}
Let $\mathrm{A}$ be a local Henselian ring (resp. strictly local).
For A to be unibranch (resp. geometrically unibranch), it is necessary and sufficient that $\Spec(\mathrm{A})$ be irreducible.
\end{corollary}

\oldpage[IV]{150}

\begin{proposition}[18.8.17]
\label{IV.18.8.17}
Let $\mathrm{A}$ be a local Noetherian ring.
If A is universally catenary, the same holds for ${}^{hs}\!\mathrm{A}$.
\end{proposition}

\begin{proof}
The proof is the same as that of \sref{IV.18.7.5}.
\end{proof}

\begin{proposition}[18.8.18]
\label{IV.18.8.18}
\begin{enumerate}[label=\emph{\roman*)}]
  \item Every filtered inductive limit of strictly local rings (where the transition homomorphisms are local) is a strictly local ring.
  \item The functor $\mathrm{A}\leadsto{}^{hs}\!\mathrm{A}$ in the category of local rings (where the morphisms are local homomorphisms) commutes with filtered inductive limits.
  \item Every strictly local ring A is an inductive limit of a filtered inductive system of Noetherian strictly local rings (the transition homomorphisms being local).
\end{enumerate}
\end{proposition}

\begin{proof}
The proofs are modelled on those of \sref{IV.18.6.14}.
\end{proof}


\subsection{Formal fibres of Noetherian Henselian rings.}
\label{subsection:IV.18.9}

\begin{theorem}[18.9.1]
\label{IV.18.9.1}
Let $\mathrm{A}$ be a local Noetherian Henselian ring, whose formal fibres are geometrically normal.
Then the formal fibres of $\mathrm{A}$ are geometrically integral (hence geometrically connected).
\end{theorem}

\begin{proof}
The two assertions stated are in fact equivalent, a prescheme that is locally Noetherian, connected and normal being integral.
As every finite integral A-algebra is a Henselian local ring \sref{IV.18.5.9} and \sref{IV.18.5.10}, it follows from \sref{IV.7.3.16.2} that one is reduced to proving that if one further supposes A \emph{integral}, then the fibre of the morphism $\Spec(\hat{\mathrm{A}})\to\Spec(\mathrm{A})$ at the generic point of $\Spec(\mathrm{A})$ is \emph{integral}, which follows from the following corollary.
\end{proof}

\begin{corollary}[18.9.2]
\label{IV.18.9.2}
Let $\mathrm{A}$ be a local Noetherian Henselian integral ring, whose formal fibres are geometrically normal.
Then $\hat{\mathrm{A}}$ is integral.
\end{corollary}

\begin{proof}
Indeed, it follows from \sref{IV.18.8.16} that A is unibranch, hence the theorem follows from the hypothesis made on the formal fibres of A and from \sref{IV.7.6.3}.
\end{proof}

\begin{corollary}[18.9.3]
\label{IV.18.9.3}
Under the hypotheses of \sref{IV.18.9.2}, the field of fractions L of $\hat{\mathrm{A}}$ is a separable extension of the field of fractions K of A, and K is algebraically closed in L.
\end{corollary}

\begin{proof}
This follows from the fact that the fibre of the morphism $\Spec(\hat{\mathrm{A}})\to\Spec(\mathrm{A})$ at the generic point is geometrically integral and from \sref{IV.4.3.2} and \sref{IV.4.3.5}.
\end{proof}

\begin{corollary}[18.9.4]
\label{IV.18.9.4}
Let $\mathrm{A}$ be a local Noetherian ring, Henselian, and whose formal fibres are geometrically normal (for example a Noetherian Henselian local ring that is excellent), and let $\mathrm{A}'=\hat{\mathrm{A}}$.
Let $\mathrm{X}$ be an A-prescheme, $\mathrm{X}'=\mathrm{X}\otimes_{\mathrm{A}}\mathrm{A}'$, $g:\mathrm{X}'\to\mathrm{X}$ the canonical projection.
Then the map $\mathrm{U}\mapsto g^{-1}(\mathrm{U})$ is a bijection from the set of subsets of $\mathrm{X}$ that are both open and closed to the set of subsets of $\mathrm{X}'$ that are both open and closed.
\end{corollary}

\begin{proof}
As $g$ is a faithfully flat and quasi-compact morphism \sref{IV.2.3.12}, one knows that the topology of X is the \emph{quotient} of that of $\mathrm{X}'$ by the equivalence relation defined by $g$.
It therefore amounts to seeing that an open and closed subset $\mathrm{U}'$ of $\mathrm{X}'$ is \emph{saturated} for this relation.
Now, as the morphism $\Spec(\mathrm{A}')\to\Spec(\mathrm{A})$ has its fibres geometrically connected \sref{IV.18.9.1}, the fibres of $g$ are connected, whence our assertion immediately follows.

In particular, if X is locally connected (which will be the case if X is locally Noetherian), then, for every connected component U of X (which is open and closed
\oldpage[IV]{151}
in X), $g^{-1}(\mathrm{U})$ is connected (Bourbaki, \emph{Top. g\'{e}n.}, chap.~I\textsuperscript{er}, 3\textsuperscript{e} \'{e}d., \S~11, no~3, prop.~7).
If one denotes by $\pi_0(\mathrm{X})$ the set of connected components of X, the map $\pi_0(\mathrm{X}')\to\pi_0(\mathrm{X})$ canonically deduced from $g$ is therefore \emph{bijective}.
\end{proof}

\begin{corollary}[18.9.5]
\label{IV.18.9.5}
Under the hypotheses of \sref{IV.18.9.4}, the functor
\begin{equation}
  \mathrm{Z}\mapsto\mathrm{Z}\otimes_{\mathrm{A}}\mathrm{A}'
  \tag{18.9.5.1}
\label{IV.18.9.5.1}
\end{equation}
from the category of \'{e}tale preschemes over $\mathrm{X}$, to that of \'{e}tale preschemes over $\mathrm{X}'$, is fully faithful.
\end{corollary}

\begin{proof}
Let $\mathrm{Z}_1$, $\mathrm{Z}_2$ be two \'{e}tale preschemes over X, and set $\mathrm{Z}'_i=\mathrm{Z}_i\otimes_{\mathrm{A}}\mathrm{A}'$ ($i=1,2$);
it is a matter of proving that every $\mathrm{X}'$-morphism $\mathrm{Z}'_1\to\mathrm{Z}'_2$ comes, by base change, from a unique X-morphism $\mathrm{Z}_1\to\mathrm{Z}_2$.
Suppose first that $\mathrm{Z}_2$ is \emph{separated} over X;
then, as $\Hom_{\mathrm{X}}(\mathrm{Z}_1,\mathrm{Z}_2)$ identifies with $\Gamma((\mathrm{Z}_1\times_{\mathrm{X}}\mathrm{Z}_2)/\mathrm{Z}_1)$, and as $\mathrm{Z}_1\times_{\mathrm{X}}\mathrm{Z}_2$ is \'{e}tale and separated over $\mathrm{Z}_1$, $\Hom_{\mathrm{X}}(\mathrm{Z}_1,\mathrm{Z}_2)$ identifies functorially, by virtue of \sref{IV.17.9.3}, with the set of open and closed subsets U of $\mathrm{Z}_1\times_{\mathrm{X}}\mathrm{Z}_2$ such that the restriction to U of the projection $p:\mathrm{Z}_1\times_{\mathrm{X}}\mathrm{Z}_2\to\mathrm{Z}_1$ is a surjective and radicial morphism.
The assertion thus follows from \sref{IV.18.9.4} and the fact that $\mathrm{Z}'_1\times_{\mathrm{X}'}\mathrm{Z}'_2=(\mathrm{Z}_1\times_{\mathrm{X}}\mathrm{Z}_2)\otimes_{\mathrm{A}}\mathrm{A}'$.

Let us pass to the general case.
It will suffice to prove that the graph $\Gamma'$ of an $\mathrm{X}'$-morphism $\mathrm{Z}'_1\to\mathrm{Z}'_2$, which is open in $\mathrm{Z}'_1\times_{\mathrm{X}'}\mathrm{Z}'_2$ \sref{IV.17.9.3}, is of the form $\Gamma\otimes_{\mathrm{A}}\mathrm{A}'$, where $\Gamma$ is induced by an \emph{open} subset of $\mathrm{Z}=\mathrm{Z}_1\times_{\mathrm{X}}\mathrm{Z}_2$.
Indeed, the restriction to $\Gamma$ of the projection $p:\mathrm{Z}\to\mathrm{Z}_1$ is then an isomorphism: after the base change $\mathrm{A}\to\mathrm{A}'$, it becomes the restriction of $p':\mathrm{Z}'\to\mathrm{Z}'_1$ to $\Gamma'$, which is an isomorphism, and it suffices to apply \sref{IV.2.7.1},~(viii).
The conclusion then follows from the characterisation of graphs \sref[I]{I.5.3}.

If $q:\mathrm{Z}'=\mathrm{Z}\otimes_{\mathrm{A}}\mathrm{A}'\to\mathrm{Z}$ is the canonical projection, it remains to prove that $\Gamma'=q^{-1}(\Gamma)$ for an open subset $\Gamma$ of Z.
Since $q$ is faithfully flat and quasi-compact, it suffices \sref{IV.2.3.12} to prove that $\Gamma'=q^{-1}(\mathrm{U})$ for a subset $\mathrm{U}\subset\mathrm{Z}$.
Let $\mathrm{Z}''=\mathrm{Z}'\times_{\mathrm{Z}}\mathrm{Z}'=\mathrm{Z}\times_{\mathrm{X}}\mathrm{X}''$ where $\mathrm{X}''=\mathrm{X}'\times_{\mathrm{X}}\mathrm{X}'$, and let $q_1:\mathrm{Z}''\to\mathrm{Z}'$, $q_2:\mathrm{Z}''\to\mathrm{Z}'$ be the canonical projections.
Applying \sref{IV.4.5.19.1}, it suffices to show that $q_1^{-1}(\Gamma')=q_2^{-1}(\Gamma')$.
But this is a property which is true if and only if it is true after each base change $\Spec(k(x))\to\mathrm{X}$, where $x$ ranges over X.
In other words, it suffices to prove the corollary when X is the \emph{spectrum of a field};
but as every X-prescheme that is \'{e}tale is automatically \emph{separated} over X \sref{IV.17.6.2},~$c'$), one is reduced to the case considered at the beginning of the proof.
\end{proof}

\begin{remarks}[18.9.6]
\label{IV.18.9.6}
\begin{enumerate}[label=\emph{\roman*)}]
  \item It is possible that, if the residue field of A is of characteristic 0, the functor \sref{IV.18.9.5.1} induces even an \emph{equivalence} of the category of \'{e}tale coverings of X and the category of \'{e}tale coverings of $\mathrm{X}'$.
  One can indeed prove this when A is \emph{excellent}, using the resolution of singularities of Hironaka (M.~Artin).
  It is plausible that an analogous statement is still true without restriction on the characteristic of the residue field, provided one restricts to the ``principal Galois'' coverings whose group has order prime to the residual characteristic (cf.~[41]).

  \oldpage[IV]{152}
  \item One does not know whether the formal fibres of A are geometrically connected (or even geometrically irreducible) when A is any Henselian local ring.
  It would suffice that, for every Noetherian Henselian local \emph{integral} ring A, $\Spec(\hat{\mathrm{A}})$ be irreducible.
  One does not know if this is always the case.

  \item One does not know whether, when A is an excellent local ring, its strict Henselization ${}^{hs}\!\mathrm{A}$ is excellent.
  To elucidate this question, one can reduce to the case where $\mathrm{A}=k[[\mathrm{T}_1,\ldots,\mathrm{T}_n]]$, $k$ being a field of characteristic $p>0$.
  Taking account of \sref{IV.18.8.17}, the question is whether the formal fibres of ${}^{hs}\!\mathrm{A}$ are geometrically regular.
  The answer is affirmative when $n=1$, but is not known for $n=2$.
  One can show that the answer is affirmative whenever, for every scheme $\mathrm{Y}_1$ finite over $\mathrm{Y}=\Spec(\hat{\mathrm{A}})$, the singularities of $\mathrm{Y}_1$ can be resolved \sref{IV.7.9.1}.
\end{enumerate}
\end{remarks}

The connectedness assertion made in \sref{IV.18.9.1} generalises in the following way:

\begin{theorem}[18.9.7]
\label{IV.18.9.7}
Let $\mathrm{A}$, $\mathrm{B}$ be two local Noetherian rings, $\rho:\mathrm{A}\to\mathrm{B}$ a local homomorphism, $f:\Spec(\mathrm{B})=\mathrm{X}\to\Spec(\mathrm{A})=\mathrm{Y}$ the corresponding morphism.
Suppose that the following conditions are satisfied:
\begin{enumerate}[label=\emph{\roman*)}]
  \item $f$ is flat and all the fibres $\mathrm{X}_y=f^{-1}(y)$ ($y\in\mathrm{Y}$) are geometrically reduced \sref{IV.4.6.2}.
  \item Either the ring $\mathrm{A}$ is geometrically unibranch \sref[0]{0.23.2.1}, or else the ring A is unibranch \sref[0]{0.23.2.1} and $k(\mathrm{B})$ is a primary extension \sref{IV.4.3.1} of $k(\mathrm{A})$.
\end{enumerate}
Then, if $\eta$ is the generic point of $\mathrm{Y}$, $\mathrm{X}_\eta$ is connected, and for every closed subset $\mathrm{T}$ of X such that $f(\mathrm{T})$ is rare in Y, $\mathrm{X}-\mathrm{T}$ is connected.
\end{theorem}

\begin{proof}
The two conclusions stated are in fact equivalent;
indeed, the hypothesis that $f(\mathrm{T})$ is rare implies that $f(\mathrm{T})$ does not contain $\eta$.
On the other hand, since $\{\eta\}$ is proconstructible \sref{I.9.6} and $f$ is flat, $\mathrm{X}_\eta$ is dense in X \sref{IV.2.3.10}; consequently, if $\mathrm{X}_\eta$ is connected, then so is $\mathrm{X}-\mathrm{T}$, which contains it.
Conversely, as U ranges over the affine open neighbourhoods of $\eta$ in Y, $\mathrm{X}_\eta$ is the projective limit of the subschemes induced on the open subsets $f^{-1}(\mathrm{U})$ of X \sref{IV.8.1.2},~a).
If one supposes that $\mathrm{X}-\mathrm{T}$ is connected for every closed subset T of X such that $f(\mathrm{T})$ is rare in Y, the $f^{-1}(\mathrm{U})$ are connected, and therefore so is $\mathrm{X}_\eta$ \sref{IV.8.4.1}.

Let us first show that one can reduce to the case where A is integral and $\mathrm{Y}=\Spec(\mathrm{A})$ is geometrically unibranch \sref{IV.6.15.1} (which implies that A is geometrically unibranch, but is not equivalent to this condition \sref{IV.6.15.2}).
It is clear that one can first suppose Y reduced, by considering the morphism $\mathrm{X}\times_{\mathrm{Y}}\mathrm{Y}_{\mathrm{red}}\to\mathrm{Y}_{\mathrm{red}}$ obtained from $f$ by base change, which is flat and has the same fibres as $f$.
One can therefore suppose A integral and unibranch.
If K is the field of fractions of A, there then exists a finite A-subalgebra $\mathrm{A}''$ of K such that, if $\mathrm{A}'$ is the integral closure of A, the morphism $\Spec(\mathrm{A}')\to\Spec(\mathrm{A}'')$ is radicial \sref[0]{0.23.2.5}; it follows \sref{IV.6.15.5} that $\Spec(\mathrm{A}'')$ is geometrically unibranch.
Since A is unibranch, $\mathrm{A}'$ is a local ring, and hence so is $\mathrm{A}''$.
Let us show that the ring $\mathrm{B}''=\mathrm{B}\otimes_{\mathrm{A}}\mathrm{A}''$ is also local.
Indeed, $\mathrm{B}''$ is a finite B-algebra, hence a Noetherian semilocal ring.
If $k''$ is the residue field of $\mathrm{A}''$, the maximal ideals of $\mathrm{B}''$ are the points of $\Spec(\mathrm{B}'')$ lying over the closed point of $\Spec(\mathrm{B})$, hence the points of $\Spec(k(\mathrm{B})\otimes_{k(\mathrm{A})}k'')$ \sref[I]{I.3.4.9},
\oldpage[IV]{153}
since the closed point of $\Spec(\mathrm{A}'')$ is the only point lying over the closed point of $\Spec(\mathrm{A})$.
Now, when A is geometrically unibranch, $k''$ is a radicial extension of $k(\mathrm{A})$, so $\Spec(k(\mathrm{B})\otimes_{k(\mathrm{A})}k'')$ is radicial over $\Spec(k(\mathrm{B}))$ and consequently consists of a single point.
When A is unibranch and $k(\mathrm{B})$ is a primary extension of $k(\mathrm{A})$, $\Spec(k(\mathrm{B})\otimes_{k(\mathrm{A})}k'')$ is irreducible \sref{IV.4.3.2}; since it is a finite discrete space \sref[I]{I.6.4.4}, it again consists of a single point.
Thus $\mathrm{B}''$ is local in both cases considered.
If $\mathrm{Y}''=\Spec(\mathrm{A}'')$, $\mathrm{X}''=\Spec(\mathrm{B}'')=\mathrm{X}\times_{\mathrm{Y}}\mathrm{Y}''$, the morphism $f''=f_{(\mathrm{Y}'')}:\mathrm{X}''\to\mathrm{Y}''$ is flat and has geometrically reduced fibres.
Moreover, if $\eta''$ is the generic point of $\mathrm{Y}''$, then $k(\eta'')=k(\eta)=\mathrm{K}$ by definition, so the fibres $\mathrm{X}_\eta$ and $\mathrm{X}''_{\eta''}$ are isomorphic.
It therefore suffices to prove that $\mathrm{X}''_{\eta''}$ is connected.
\end{proof}

\begin{remark}[18.9.7.1]
\label{IV.18.9.7.1}
When the ring A is \emph{Japanese}, one can, in the case preceding \sref{IV.18.9.7.2}, take $\mathrm{A}''=\mathrm{A}'$ by definition \sref[0]{0.23.1.1};
one sees therefore that in this case one would even be able to prove the theorem when A is \emph{integrally closed}.
\end{remark}

\begin{env}[18.9.7.2]
\label{IV.18.9.7.2}
Suppose therefore now that Y is integral and geometrically unibranch.
Note that if T is a closed subset of X such that $f(\mathrm{T})$ is rare in Y, T contains no maximal point of X since $f$ is flat \sref{IV.2.3.4}, and is therefore rare in X.
Since X is a local scheme, hence connected, to prove that $\mathrm{X}-\mathrm{T}$ is connected it suffices, by virtue of \sref{IV.15.5.6.1}, to show that for every $x\in\mathrm{X}$ such that $f(x)\neq\eta$, $\Spec(\sh{O}_{\mathrm{X},x})-\{x\}$ is connected.
Set $y=f(x)$, $\mathrm{A}_1=\sh{O}_{\mathrm{Y},y}$, $\mathrm{B}_1=\sh{O}_{\mathrm{X},x}$, $\mathrm{Y}_1=\Spec(\mathrm{A}_1)$, $\mathrm{X}_1=\Spec(\mathrm{B}_1)$, and let $f_1:\mathrm{X}_1\to\mathrm{Y}_1$ be the morphism corresponding to the local homomorphism $\mathrm{A}_1\to\mathrm{B}_1$ induced by $\rho$.
It is clear that $f_1$ is flat, and it follows from \sref{IV.4.6.1} that its fibres are geometrically reduced.
Moreover, $\mathrm{A}_1$ is integral and geometrically unibranch \sref{IV.6.15.1}, and $\dim(\mathrm{A}_1)\geq 1$.
One is thus reduced to proving the following lemma.
\end{env}

\begin{lemma}[18.9.7.3]
\label{IV.18.9.7.3}
Let $\mathrm{A}$, $\mathrm{B}$ be two local Noetherian rings, $\mathrm{Y}=\Spec(\mathrm{A})$, $\mathrm{X}=\Spec(\mathrm{B})$, $\rho:\mathrm{A}\to\mathrm{B}$ a local homomorphism, $f:\mathrm{X}\to\mathrm{Y}$ the corresponding morphism.
Suppose that the following conditions are satisfied:
\begin{enumerate}[label=\emph{\roman*)}]
  \item $f$ is flat and all the fibres $\mathrm{X}_y$ ($y\in\mathrm{Y}$) are geometrically reduced.
  \item A is integral, geometrically unibranch and $\dim(\mathrm{A})\geq 1$.
\end{enumerate}
Then, if $b$ is the closed point of $\mathrm{X}$, $\mathrm{X}-\{b\}$ is connected.
\end{lemma}

\begin{proof}
Note first that by virtue of the theorem of Hartshorne \sref{IV.5.10.7}, $\mathrm{X}-\{b\}$ is connected if $\operatorname{prof}(\mathrm{B})\geq 2$.
Now one has \sref{IV.6.3.1}
\begin{equation}
  \operatorname{prof}(\mathrm{B})=\operatorname{prof}(\mathrm{A})+\operatorname{prof}(\mathrm{B}\otimes_{\mathrm{A}}k)
  \label{IV.18.9.7.4}
\tag{18.9.7.4}
\end{equation}
where $k$ denotes the residue field of A.
On the other hand, since A is integral and $\dim(\mathrm{A})\geq 1$, one has $\operatorname{prof}(\mathrm{A})\geq 1$, hence $\operatorname{prof}(\mathrm{B})\geq 2$ \emph{unless} $\operatorname{prof}(\mathrm{B}\otimes_{\mathrm{A}}k)=0$;
moreover, $\mathrm{B}\otimes_{\mathrm{A}}k$ is reduced by~(i), hence the relation $\operatorname{prof}(\mathrm{B}\otimes_{\mathrm{A}}k)=0$ means that $\mathrm{B}\otimes_{\mathrm{A}}k$ is a \emph{field} \sref[0]{0.16.4.7}.
We shall henceforth suppose that this last condition is satisfied.
We shall begin by treating a case where the proof is very simple.

A) \emph{Case where A is integrally closed.}
Then, if $\dim(\mathrm{A})\geq 2$, one has $\operatorname{prof}(\mathrm{A})\geq 2$ \sref[0]{0.16.5.1}, hence $\operatorname{prof}(\mathrm{B})\geq 2$.
One therefore only has to consider the case where $\dim(\mathrm{A})=1$ and where $\mathrm{B}\otimes_{\mathrm{A}}k$ is a field;
A is then a discrete valuation ring, hence regular, and since $\mathrm{B}\otimes_{\mathrm{A}}k$
\oldpage[IV]{154}
is a field and $f$ is flat, B is regular \sref[0]{0.17.3.3}.
Moreover \sref{IV.6.1.1.1}, one has $\dim(\mathrm{B})=\dim(\mathrm{A})+\dim(\mathrm{B}\otimes_{\mathrm{A}}k)=\dim(\mathrm{A})=1$, so B is also a discrete valuation ring.
Thus $\mathrm{X}-\{b\}$ consists of a single point, which proves the lemma \sref{IV.18.9.7.3} in this case.

One will observe that this proves the statement of the theorem \sref{IV.18.9.7} when one supposes in addition that, in that statement, the ring A is \emph{Japanese}, for by virtue of the remark \sref{IV.18.9.7.1}, one is reduced to the case where A is \emph{integrally closed}, and then, in the preceding reduction \sref{IV.18.9.7.3}, the ring $\mathrm{A}_1=\sh{O}_{\mathrm{Y},y}$ is also integrally closed and one can apply the result just proved.

B) \emph{Case where} $\dim(\mathrm{A})\geq 2$.
The result of \sref{IV.18.9.7.3} will follow from the two following lemmas.
\end{proof}

\begin{lemma}[18.9.7.5]
\label{IV.18.9.7.5}
Let $\mathrm{A}$ be a semi-local Noetherian ring, reduced, $\mathrm{A}'$ its integral closure in its total ring of fractions, $\mathrm{X}=\Spec(\mathrm{A})$, $\mathrm{X}'=\Spec(\mathrm{A}')$;
let S be the set of closed points of $\mathrm{X}$ and $\mathrm{U}=\mathrm{X}-\mathrm{S}$.
Suppose that for every point $x'\in\mathrm{X}'$ above a point of S, $\dim(\sh{O}_{\mathrm{X}',x'})\geq 2$;
then $\mathrm{A}^{(1)}=\Gamma(\mathrm{U},\sh{O}_{\mathrm{X}})$ is integral over $\mathrm{A}$ and is an A-algebra isomorphic to a subalgebra of $\mathrm{A}'$.
\end{lemma}

\begin{proof}
Recall that if the $\mathfrak{p}_i$ ($1\leq i\leq m$) are the minimal prime ideals of A, then $\mathrm{A}'$ is the direct product of the integral closures $\mathrm{A}'_i$ of the integral rings $\mathrm{A}_i=\mathrm{A}/\mathfrak{p}_i$, so that $\mathrm{X}'$ is the sum of the schemes $\mathrm{X}'_i=\Spec(\mathrm{A}'_i)$ ($1\leq i\leq m$).
If $\mathrm{X}_0$ is the sum of the schemes $\mathrm{X}_i=\Spec(\mathrm{A}_i)$ ($1\leq i\leq m$) and $\mathrm{U}_0$ the inverse image of U in $\mathrm{X}_0$, the canonical homomorphism $\Gamma(\mathrm{U},\sh{O}_{\mathrm{X}})\to\Gamma(\mathrm{U}_0,\sh{O}_{\mathrm{X}_0})$ is injective since X is reduced and the morphism $\mathrm{X}_0\to\mathrm{X}$ (and hence also $\mathrm{U}_0\to\mathrm{U}$) is surjective.
As $\Gamma(\mathrm{U}_0,\sh{O}_{\mathrm{X}_0})$ is the direct product of the $\Gamma(\mathrm{U}_i,\sh{O}_{\mathrm{X}_i})$ (where $\mathrm{U}_i$ is the inverse image of U in $\mathrm{X}_i$), it suffices to prove that for each $i$, $\Gamma(\mathrm{U}_i,\sh{O}_{\mathrm{X}_i})$ is isomorphic to a sub-$\mathrm{A}_i$-algebra of $\mathrm{A}'_i$.
In other words, one is reduced to the case where A is a Noetherian integral semilocal ring.
Since X is reduced and the morphism $\mathrm{X}'\to\mathrm{X}$ is surjective, the homomorphism $\Gamma(\mathrm{U},\sh{O}_{\mathrm{X}})\to\Gamma(\mathrm{U}',\sh{O}_{\mathrm{X}'})$ (where $\mathrm{U}'$ is the inverse image of U in $\mathrm{X}'$) is injective, and it suffices to verify that the canonical homomorphism $\mathrm{A}'=\Gamma(\mathrm{X}',\sh{O}_{\mathrm{X}'})\to\Gamma(\mathrm{U}',\sh{O}_{\mathrm{X}'})$ is \emph{bijective}.
Now \sref[I]{I.8.2.1.1}, $\Gamma(\mathrm{U}',\sh{O}_{\mathrm{X}'})$ is the intersection of the local rings $\mathrm{A}'_{\mathfrak{p}'}$, where $\mathfrak{p}'$ ranges over U', and by virtue of the hypothesis, among these local rings are included all those of height 1.
But the reasoning of \sref[0]{0.23.2.7} applies also to a Noetherian integral semilocal ring, so $\mathrm{A}'$ is a \emph{semilocal Krull ring}, and is therefore the intersection of the local rings $\mathrm{A}'_{\mathfrak{p}'}$, where $\mathfrak{p}'$ ranges over the set of prime ideals of height 1 of $\mathrm{A}'$ (Bourbaki, \emph{Alg. comm.}, chap.~VII, \S~1, no~6, th.~4);
\emph{a fortiori}, $\mathrm{A}'$ is the intersection of the $\mathrm{A}'_{\mathfrak{p}'}$ for $\mathfrak{p}'\in\mathrm{U}'$, which completes the proof of the lemma.
\end{proof}

\begin{remark}[18.9.7.6]
\label{IV.18.9.7.6}
We know \sref[0]{0.23.2.5} that there exists a finite sub-A-algebra $\mathrm{A}''$ of K, the field of fractions of A, such that the morphism $\Spec(\mathrm{A}')\to\Spec(\mathrm{A}'')$ is \emph{radicial};
as this morphism is also integral and dominant, hence surjective \sref[II]{II.6.1.10}, it is a \emph{homeomorphism} \sref{IV.2.4.5}.
By virtue of the geometric interpretation of $\dim(\sh{O}_{\mathrm{X}',x'})$ \sref{IV.5.1.2}, one has therefore, for every point $x''$ of $\mathrm{X}''=\Spec(\mathrm{A}'')$ above a point $x\in\mathrm{X}$, $\dim(\sh{O}_{\mathrm{X}'',x''})=\dim(\sh{O}_{\mathrm{X}',x'})\geq 2$, where $x'$ is the unique point of $\mathrm{X}'$ above $x$.
\oldpage[IV]{155}
This being so, it follows first from \sref[0]{0.16.1.5} and \sref[0]{0.16.1.4.1} that the relation $\dim(\sh{O}_{\mathrm{X}',x'})\geq 2$ implies $\dim(\sh{O}_{\mathrm{X},x})\geq 2$.
Conversely, suppose that $\dim(\sh{O}_{\mathrm{X},x})\geq 2$ and that A is universally catenary \sref{IV.5.6.2}; then so is $\sh{O}_{\mathrm{X},x}$ \sref{IV.5.6.3}, hence $\dim(\sh{O}_{\mathrm{X}'',x''})=\dim(\sh{O}_{\mathrm{X},x})\geq 2$ by virtue of \sref{IV.5.6.10}, and consequently $\dim(\sh{O}_{\mathrm{X}',x'})\geq 2$.
The same conclusion holds if $\sh{O}_{\mathrm{X},x}$ is \emph{geometrically unibranch}, for the fibre of the morphism $\mathrm{X}'\to\mathrm{X}$ at the point $x$ is then reduced to a single point, and $\sh{O}_{\mathrm{X}',x'}$ is an \emph{integral} $\sh{O}_{\mathrm{X},x}$-algebra;
it suffices then to apply \sref[0]{0.16.1.5}.
\end{remark}

\begin{lemma}[18.9.7.7]
\label{IV.18.9.7.7}
Let $\mathrm{A}$ be a local Noetherian ring, of dimension $\geq 2$, integral and geometrically unibranch, $\mathrm{Y}=\Spec(\mathrm{A})$, $y$ the closed point of Y.
Let $\mathrm{X}$ be a locally Noetherian and connected prescheme and $f:\mathrm{X}\to\mathrm{Y}$ a flat morphism.
Then, for every closed subset $\mathrm{T}\subset f^{-1}(y)$, $\mathrm{X}-\mathrm{T}$ is connected.
\end{lemma}

\begin{proof}
The same reasoning as at the beginning of \sref{IV.18.9.7.2} reduces to proving that, for every point $x\in f^{-1}(y)$, $\Spec(\sh{O}_{\mathrm{X},x})-\{x\}$ is connected.
One can therefore restrict to the case where $\mathrm{X}=\Spec(\mathrm{B})$, B being a local Noetherian ring and the homomorphism $\mathrm{A}\to\mathrm{B}$ being local;
let $\mathrm{U}=\mathrm{Y}-\{y\}$, and note that $\mathrm{V}=f^{-1}(\mathrm{U})$ is dense in X \sref{IV.2.3.10};
it suffices therefore to prove that V is \emph{connected}.
Keeping the notations of \sref{IV.18.9.7.5}, and setting $\mathrm{A}^{(1)}=\Gamma(\mathrm{U},\sh{O}_{\mathrm{Y}})$, $\mathrm{Y}_1=\Spec(\mathrm{A}^{(1)})$, $\mathrm{X}_1=\mathrm{X}\times_{\mathrm{Y}}\mathrm{Y}_1=\Spec(\mathrm{B}\otimes_{\mathrm{A}}\mathrm{A}^{(1)})$, we have the commutative diagram
\[
  \xymatrix{
    \mathrm{X} \ar[d]_{f} & \mathrm{X}_1 \ar[l]^{g_1} \ar[d]^{f_1} \\
    \mathrm{Y} & \mathrm{Y}_1 \ar[l]_{g}
  }
\]
By virtue of \sref{IV.2.3.1} and of the definition of $\mathrm{A}^{(1)}$, one has a canonical isomorphism $\Gamma(\mathrm{V},\sh{O}_{\mathrm{X}})\xrightarrow{\sim}\mathrm{B}\otimes_{\mathrm{A}}\mathrm{A}^{(1)}$, and it suffices therefore to prove that this last ring is \emph{local}, since a local ring cannot be a product of two nonzero rings \sref[III]{III.7.8.6.1}.
But we have seen \sref{IV.18.9.7.5} that $\mathrm{A}^{(1)}$ is a sub-A-algebra of the integral closure $\mathrm{A}'$ of A, since $\dim(\mathrm{A})\geq 2$.
The hypothesis that A is geometrically unibranch implies that the morphism $\Spec(\mathrm{A}')\to\Spec(\mathrm{A})$ is radicial at the point $y$ \sref{IV.6.15.3};
\emph{a fortiori} $g:\mathrm{Y}_1\to\mathrm{Y}$ is radicial at the point $y$, and consequently $g_1$ is radicial at the closed point $x$ of X;
furthermore $g_1$ is an integral morphism, hence the closed points of $\mathrm{X}_1$ are above $x$, hence there is only one closed point of $\mathrm{X}_1$, which completes the proof of \sref{IV.18.9.7.7}.

It is clear that \sref{IV.18.9.7.7} proves \sref{IV.18.9.7.3} when $\dim(\mathrm{A})\geq 2$ (even without supposing the fibres geometrically reduced).

C) \emph{Case where} $\dim(\mathrm{A})=1$.
One has seen at the beginning of the proof of \sref{IV.18.9.7.3} that one can suppose $\dim(\mathrm{B}\otimes_{\mathrm{A}}k)=0$, and by flatness \sref{IV.6.1.1.1}, one therefore has
\[
  \dim(\mathrm{B})=\dim(\mathrm{A})+\dim(\mathrm{B}\otimes_{\mathrm{A}}k)=1.
\]
Since A is integral, Y consists of two points, the generic point $\eta$ and the closed point $a$.
Furthermore, since $f$ is flat, every irreducible component of X dominates Y \sref{IV.2.3.4}, hence $\mathrm{X}-\{b\}$, which is the set of maximal points $\xi_i$ of X, is equal to
\oldpage[IV]{156}
the underlying set of $f^{-1}(\eta)$, and the fibre $f^{-1}(\eta)$ is the sum of the $\Spec(k(\xi_i))$.
The hypothesis on the $\mathrm{X}_y$ and the fact that these fibres are of dimension 0, show that they are \emph{geometrically regular} \sref{IV.6.7.6}, and \emph{a fortiori} geometrically normal, in other words $f$ is a \emph{normal morphism}.
But since A is integral and geometrically unibranch, it follows then from \sref{IV.6.15.10} that B is also integral and geometrically unibranch, hence $f^{-1}(\eta)=\mathrm{X}-\{b\}$ is reduced to a single point, and \emph{a fortiori} connected;
this concludes the proof of \sref{IV.18.9.7.3} and of \sref{IV.18.9.7}.
\end{proof}

\begin{remark}[18.9.7.8]
\label{IV.18.9.7.8}
In case C) of the proof of \sref{IV.18.9.7.3}, one can avoid making appeal to the delicate result \sref{IV.6.15.10} by reasoning as follows:
since A is integral and of dimension 1, it follows from the theorem of Krull-Akizuki (Bourbaki, \emph{Alg. comm.}, chap.~VII, \S~2, no~5, prop.~5) that its integral closure $\mathrm{A}'$ is a \emph{Noetherian} ring.
The same reasoning as at the beginning of the proof of \sref{IV.18.9.7} shows that $\mathrm{B}'=\mathrm{B}\otimes_{\mathrm{A}}\mathrm{A}'$ is a local ring;
furthermore, by flatness, $\mathrm{B}'$ is contained in the total ring of fractions R of the reduced $\mathrm{B}$ \sref{IV.3.3.5}.
Now, the theorem of Krull-Akizuki (Bourbaki, \emph{loc. cit.}) also applies to a local Noetherian \emph{reduced} ring of dimension 1, and shows that for every such ring B, everything contained between B and its total ring of fractions is \emph{Noetherian}.
The ring $\mathrm{B}'$ being a local Noetherian ring and the morphism $\Spec(\mathrm{B}')\to\Spec(\mathrm{A}')$ being normal, $\mathrm{B}'$ is an integral ring and integrally closed \sref{IV.6.5.4}, and one concludes as at the beginning of the proof of \sref{IV.18.9.7}.
\end{remark}

\begin{corollary}[18.9.8]
\label{IV.18.9.8}
With the notations of \sref{IV.18.9.7}, suppose that condition~(i) of \sref{IV.18.9.7} and one of the following two conditions hold:
\begin{enumerate}
  \item[\rm(ii$'$)] A is a strictly local ring \sref{IV.18.8.2}.
  \item[\rm(ii$''$)] A is a Henselian local ring and $k(\mathrm{B})$ is a primary extension of $k(\mathrm{A})$.
\end{enumerate}
Then all the fibres $\mathrm{X}_y$ ($y\in\mathrm{Y}$) are geometrically connected.
\end{corollary}

\begin{proof}
Let $\mathfrak{p}$ be the prime ideal of A corresponding to the point $y$, and set $\mathrm{A}'=\mathrm{A}/\mathfrak{p}$ and $\mathrm{B}'=\mathrm{B}\otimes_{\mathrm{A}}\mathrm{A}'=\mathrm{B}/\mathfrak{p}\mathrm{B}$, which are evidently local Noetherian rings;
it is clear that the morphism $f':\Spec(\mathrm{B}')\to\Spec(\mathrm{A}')$ satisfies condition~(i) of \sref{IV.18.9.7};
and as $y$ is the generic point of $\Spec(\mathrm{A}')$ and $k(\mathrm{A}')=k(\mathrm{A})$, $k(\mathrm{B}')=k(\mathrm{B})$, one sees that it suffices to verify that, in the case~(ii$'$) (resp.~(ii$''$)), $\mathrm{A}'$ is geometrically unibranch (resp. unibranch).
But $\mathrm{A}'$ is integral, and in the case~(ii$'$) (resp.~(ii$''$)) it is strictly local (resp. Henselian) by virtue of \sref{IV.18.5.10}.
One therefore concludes that $\mathrm{A}'$ is geometrically unibranch (resp. unibranch) by virtue of \sref{IV.18.8.15} (resp. \sref{IV.18.6.12}).
\end{proof}

\begin{corollary}[18.9.9]
\label{IV.18.9.9}
Let $\mathrm{A}$ be a local Noetherian ring, Henselian, universally Japanese (i.e.~\sref{IV.7.6.4} and \sref{IV.7.7.2}), whose formal fibres \sref{IV.7.3.13} are geometrically reduced;
then all the formal fibres of A are geometrically connected.
\end{corollary}

\begin{proof}
It suffices to apply \sref{IV.18.9.8} in the case where $\mathrm{B}=\hat{\mathrm{A}}$, since B is a flat A-module and $k(\mathrm{B})=k(\mathrm{A})$.
\end{proof}

\begin{remark}[18.9.10]
\label{IV.18.9.10}
The proofs of \sref{IV.18.9.4} and \sref{IV.18.9.5} use only the fact that the formal fibres of A are geometrically connected.
The conclusions
\oldpage[IV]{157}
of these two propositions are therefore still valid when one supposes the local ring A Noetherian, Henselian and universally Japanese.
\end{remark}

\begin{corollary}[18.9.11]
\label{IV.18.9.11}
Let $\mathrm{X}$, $\mathrm{Y}$ be two locally Noetherian preschemes, $\mathrm{Y}$ being supposed geometrically unibranch and X connected.
Let $f:\mathrm{X}\to\mathrm{Y}$ be a reduced morphism (i.e. flat and with geometrically reduced fibres; \sref{IV.6.8.1}).
Then, for every closed subset $\mathrm{T}$ of $\mathrm{X}$ such that $f(\mathrm{T})$ is rare in $\mathrm{Y}$, $\mathrm{X}-\mathrm{T}$ is connected.
\end{corollary}

\begin{proof}
As $f$ is flat and $f(\mathrm{T})$, being rare, cannot contain a maximal point of Y, T cannot contain a maximal point of X \sref{IV.2.3.4}, and hence is rare in X.
Utilising \sref{IV.15.5.6.1}, it suffices to show that, for every $x\in\mathrm{T}$, $\Spec(\sh{O}_{\mathrm{X},x})-\{x\}$ is connected.
Set $y=f(x)$, which is not a maximal point of Y.
Since the ring $\sh{O}_{\mathrm{Y},y}$ is geometrically unibranch, one can apply \sref{IV.18.9.7} to the morphism
\[
  f_1:\Spec(\sh{O}_{\mathrm{X},x})\to\Spec(\sh{O}_{\mathrm{Y},y}),
\]
which satisfies conditions~(i) and~(ii) of this theorem, and to the closed subset $\{x\}$ of $\Spec(\sh{O}_{\mathrm{X},x})$, since $f_1(\{x\})=\{y\}$ is rare in $\Spec(\sh{O}_{\mathrm{Y},y})$.
\end{proof}
\subsection{\'Etale preschemes over a geometrically unibranch or normal prescheme}
\label{subsection:IV.18.10}

\begin{theorem}[18.10.1]
\label{IV.18.10.1}
Let $f:X\to Y$ be a morphism locally of finite presentation, $x$ a point of $X$, $y=f(x)$.
Suppose that $Y$ is integral and geometrically unibranch at the point $y$.
Then, for $f$ to be \'etale at the point $x$, it is necessary and sufficient that the following two conditions be satisfied:
\begin{enumerate}[label=\emph{\roman*)}]
  \item $f$ is unramified at the point $x$;
  \item the homomorphism $\sh{O}_{Y,y}\to\sh{O}_{X,x}$ is injective.
\end{enumerate}

Moreover, if this is the case, $X$ is integral and geometrically unibranch at the point $x$.

Finally, if $\sh{O}_{Y,y}$ is Noetherian, one can, in the preceding criterion, replace the condition \emph{(ii)} by the condition

\emph{(ii bis)} $\dim\sh{O}_{Y,y}\leq\dim\sh{O}_{X,x}$.
\end{theorem}

\begin{proof}
The fact that, if $f$ is \'etale, $X$ is integral and geometrically unibranch at the point $x$ is a particular case of \sref{IV.17.5.7}.
It is clear on the other hand that the conditions (i) and (ii) are necessary for $f$ to be \'etale at the point $x$, since $\sh{O}_{X,x}$ is then an $\sh{O}_{Y,y}$-module faithfully flat \sref{IV.17.6.1}.
Let us prove that these conditions are sufficient.

Set $\mathrm{A}=\sh{O}_{Y,y}$, $\mathrm{A}'={}^{\mathrm{hs}}\mathrm{A}$ \sref{IV.18.8.7}, $\mathrm{Y}'=\operatorname{Spec}(\mathrm{A}')$, $\mathrm{X}'=\mathrm{X}\times_{\mathrm{Y}}\mathrm{Y}'$, $f'=f_{(\mathrm{Y}')}:\mathrm{X}'\to\mathrm{Y}'$.
Let $y'$ be the closed point of $\mathrm{Y}'$, $x'$ a point of $\mathrm{X}'$ above $x$ and $y'$; since the morphism $g:\mathrm{Y}'\to\mathrm{Y}$ is \emph{flat} \sref{IV.18.8.8}, it amounts to showing that $f'$ is \'etale at the point $x'$ \sref{IV.17.7.1},~(iii).
The hypothesis (ii) implies that the morphism $\operatorname{Spec}(\sh{O}_{X,x})\to\operatorname{Spec}(\sh{O}_{Y,y})$ is dominant \sref[I]{I.1.2.7}; as $\sh{O}_{Y,y}$ is integral, $\operatorname{Spec}(\sh{O}_{Y,y})$ has a single generic point $y_1$, and there is a generization $x_1$ of $x$ in $\mathrm{X}$ above $y_1$.
The projection morphism $\mathrm{X}'\to\mathrm{X}$ being flat, there exists a generization $x_1'$ of $x'$ above $x_1$ \sref{IV.2.3.4}; set $y_1'=f'(x_1')$. We then have $g(y_1')=y_1$.
Now, the hypothesis on $y$ implies that $\mathrm{A}'$ is \emph{integral} \sref{IV.18.8.15}, hence $\mathrm{Y}'$ has a single generic point $\eta$, necessarily above $y_1$ since $g$
\oldpage[IV]{158}
is flat \sref{IV.2.3.4}; on the other hand, $\eta$ is also the generic point of $g^{-1}(y_1)$ \sref[0]{0.2.1.8}, and one knows on the other hand that the fibres of $g$ are discrete spaces \sref{IV.18.8.12}, hence $g^{-1}(y_1)=\{\eta\}$, and consequently $y_1'=\eta$.
The morphism $\operatorname{Spec}(\sh{O}_{X',x'})\to\operatorname{Spec}(\sh{O}_{Y',y'})$, restriction of $f'$, is therefore dominant, and since $\sh{O}_{Y',y'}=\mathrm{A}'$ is integral, the corresponding homomorphism $\sh{O}_{Y',y'}\to\sh{O}_{X',x'}$ is injective \sref[I]{I.1.2.7}.
On the other hand, $f'$ is unramified at the point $x'$; hence \sref{IV.18.8.3},~d), there exists a neighbourhood $\mathrm{U}$ of $x'$ in $\mathrm{X}'$ such that $f'|\mathrm{U}$ is a \emph{closed immersion} of $\mathrm{U}$ into $\mathrm{Y}'$.
\emph{A fortiori} the homomorphism $\sh{O}_{Y',y'}\to\sh{O}_{X',x'}$ is surjective, and since it is injective, it is therefore \emph{bijective}.
But then $\sh{O}_{X',x'}$ is a flat $\sh{O}_{Y',y'}$-module, and $f'$ is therefore \'etale at the point $x'$ \sref{IV.17.6.1}.

Finally, when $\sh{O}_{Y,y}$ is Noetherian, we know that if $f$ is \'etale at the point $x$, we have $\dim\sh{O}_{Y,y}=\dim\sh{O}_{X,x}$ \sref{IV.17.6.4}.
Conversely, suppose the conditions (i) and (ii~bis) satisfied, and let us show that they imply (ii).
Let $\mathfrak{J}$ be the kernel of the canonical homomorphism $\sh{O}_{Y,y}\to\sh{O}_{X,x}$; restricting Y as needed to a neighbourhood of $y$, we may suppose that $\mathfrak{J}=\sh{J}_y$, where $\sh{J}$ is an ideal of $\sh{O}_Y$.
If $\mathrm{Y}_1$ is the closed subscheme of $\mathrm{Y}$ defined by $\sh{J}$, one can, restricting further $\mathrm{X}$ and $\mathrm{Y}$ to neighbourhoods of $x$ and $y$ respectively, suppose that $f$ factorises as $\mathrm{X}\xrightarrow{f_1}\mathrm{Y}_1\to\mathrm{Y}$ \sref[I]{I.6.5.1}.
It is clear that $f_1$ is still unramified at the point $x$, hence quasi-finite at this point \sref{IV.17.4.1}; the proof of \sref{IV.5.4.1},~(i) then shows that $\dim(\sh{O}_{X,x})\leq\dim(\sh{O}_{Y,y}/\mathfrak{J})$.
But if one had $\mathfrak{J}\neq 0$, one would conclude that $\dim(\sh{O}_{Y,y}/\mathfrak{J})<\dim(\sh{O}_{Y,y})$ since $\mathrm{Y}$ is integral \sref[0]{0.16.1.2.2}; one would then have $\dim\sh{O}_{X,x}<\dim\sh{O}_{Y,y}$, contradicting the hypothesis, which proves (ii).
\end{proof}

\begin{remark}[18.10.2]
\label{IV.18.10.2}
\medskip\noindent
\begin{enumerate}[label=(\roman*)]
  \item When $\mathrm{A}=\sh{O}_{Y,y}$ is Noetherian, and one knows that its completion $\hat{\mathrm{A}}$ is integral (for example if $\mathrm{A}$ is regular), one can, in the preceding proof, replace $\mathrm{A}'$ by $\hat{\mathrm{A}}$.
  \item When $\mathrm{Y}$ is locally Noetherian, one can give a faster proof of \sref{IV.18.10.1}, without using strict Henselization, but invoking the rather delicate results of \textsection\textsection~14 and 15.
As $f$ is quasi-finite at the point $x$ by hypothesis \sref{IV.17.4.1}, one can, replacing $\mathrm{X}$ by an open neighbourhood of $x$, suppose that $f^{-1}(y)=\{x\}$.
The hypothesis (ii) implies, as we have seen, the existence of an irreducible component $\mathrm{Z}$ of $\mathrm{X}$ containing $x$ and dominating the unique irreducible component $\mathrm{Y}_0$ of $\mathrm{Y}$ containing $y$.
The morphism $f|\mathrm{Z}$ being quasi-finite at the point $x$, hence equidimensional at this point \sref{IV.13.2.2}, it follows from the hypothesis on $\mathrm{Y}$ and from the Chevalley criterion \sref{IV.14.4.4} that $g=f|\mathrm{Z}$ is universally open at the point $x$.
Furthermore, since $f$ (and hence also $g$) is unramified at the point $x$, $g^{-1}(y)$ is geometrically reduced over $k(y)$ \sref{IV.17.4.1}; as $\sh{O}_{Y,y}$ is integral, one concludes from \sref{IV.15.2.3} that $g$ is \emph{flat} at the point $x$, hence $f$ is \emph{\'etale} at this point \sref{IV.18.4.9}.
  \item With the same notations and hypotheses on $\mathrm{Y}$ as in \sref{IV.18.10.1}, suppose only that $f$ is \emph{locally of finite type} (and not necessarily locally of finite presentation).
Then, for $\sh{O}_{X,x}$ to be an $\sh{O}_{Y,y}$-algebra \emph{essentially \'etale} \sref{IV.18.6.1}, it is necessary and sufficient that $f$ satisfy the two following conditions:
\oldpage[IV]{159}
\begin{enumerate}[label=\arabic*\textsuperscript{o}]
  \item $f$ is formally unramified at the point $x$;
  \item the homomorphism $\sh{O}_{Y,y}\to\sh{O}_{X,x}$ is injective.
\end{enumerate}

There is nothing more than proving the sufficiency of these conditions.
Taking into account \sref{IV.18.4.12}, it suffices to show that $f$ is \emph{flat} at the point $x$.
Now, by going back to the proof of \sref{IV.18.10.1} and the notations used in this proof, it suffices to show that $f'$ is flat at the point $x'$ \sref{IV.2.5.1}; by \sref{IV.17.1.3} furthermore $f'$ is formally unramified at the point $x'$.
One can therefore restrict to carrying out the proof when $\mathrm{Y}=\operatorname{Spec}(\mathrm{A})$, where $\mathrm{A}$ is a \emph{strictly local} ring and $y=f(x)$ the closed point of $\mathrm{Y}$.
Using then \sref{IV.18.8.3} we see that $\sh{O}_{Y,y}\to\sh{O}_{X,x}$ is surjective, hence \emph{bijective} by virtue of the hypothesis, and this suffices to show the flatness of $f$ at the point $x$.

Suppose furthermore that $\mathrm{Y}$ is \emph{locally integral} at the point $y$ \sref[I]{I.2.1.8}.
Then the conditions $1^{\circ}$ and $2^{\circ}$ above (together with the fact that $\mathrm{Y}$ is geometrically unibranch at the point $y$ and $f$ locally of finite type) already imply that $f$ is \emph{\'etale at the point} $x$.
Indeed, this follows from what precedes and from \sref{IV.18.4.13}.
\end{enumerate}
\end{remark}

\begin{corollary}[18.10.3]
\label{IV.18.10.3}
Let $\mathrm{Y}$ be an integral and geometrically unibranch prescheme, $\eta$ its generic point, $\mathrm{X}$ a connected prescheme, and $f:\mathrm{X}\to\mathrm{Y}$ a morphism locally of finite type such that $f^{-1}(\eta)$ is non-empty.
Then, if $f$ is formally unramified, $f$ is \'etale, and $\mathrm{X}$ is integral and geometrically unibranch.
\end{corollary}

\begin{proof}
It follows from the hypothesis and from \sref{IV.18.4.13} that $f$ is \'etale at all the points where it is flat, and in particular at the points of $f^{-1}(\eta)$, since $\sh{O}_{Y,\eta}=k(\eta)$ is a field.
The open set $\mathrm{U}$ of points of $\mathrm{X}$ where $f$ is \'etale is therefore \emph{non-empty}.
Let us show on the other hand that for all $x\in\overline{\mathrm{U}}$, the homomorphism $\sh{O}_{Y,f(x)}\to\sh{O}_{X,x}$ is \emph{injective}.
Indeed, let $s$ be a section of $\sh{O}_Y$ over an open neighbourhood of $f(x)$, which one may always suppose to be $\mathrm{Y}$ itself (the question being local on $\mathrm{X}$ and $\mathrm{Y}$), and suppose that its image $t\in\Gamma(\mathrm{Y},f_*(\sh{O}_X))=\Gamma(\mathrm{X},\sh{O}_X)$ has zero germ at the point $x$, hence vanishes in a neighbourhood $\mathrm{V}$ of $x$; then the restriction of $t$ to the non-empty open set $\mathrm{U}\cap\mathrm{V}$ is zero. But the restriction $f|(\mathrm{U}\cap\mathrm{V})$ is \'etale, hence $\mathrm{W}=f(\mathrm{U}\cap\mathrm{V})$ is open in $\mathrm{Y}$ and consequently contains the generic point $\eta$. Since the homomorphism $\sh{O}_{Y,\eta}\to\sh{O}_{X,z}$ is injective for every point $z\in f^{-1}(\eta)$, the hypothesis that $t|(\mathrm{U}\cap\mathrm{V})=0$ implies that $s_\eta=0$; and as $\mathrm{Y}$ is integral, this implies $s=0$, proving our assertion.
Applying now \sref{IV.18.10.2},~(iii), we conclude that $f$ is \emph{flat} at the point $x$, in other words $x\in\mathrm{U}$, hence $\mathrm{U}$ is both open and closed in $\mathrm{X}$; as it is non-empty and $\mathrm{X}$ is connected, we have $\mathrm{U}=\mathrm{X}$.

Note now that since $f$ is \'etale, hence locally quasi-finite and $\mathrm{Y}$ is irreducible, the maximal points of $\mathrm{X}$ belong to $f^{-1}(\eta)$ \sref{IV.2.3.4} and for all $x\in\mathrm{X}$, there is a neighbourhood of $x$ that contains only a finite number of these maximal points; in other words, the set of irreducible components of $\mathrm{X}$ is locally finite.
But by \sref{IV.18.10.1}, $\mathrm{X}$ is integral and geometrically unibranch at every point, hence every point of $\mathrm{X}$ belongs to only one irreducible component; and since the irreducible components form a locally finite family, they are open (and closed) in $\mathrm{X}$. Since $\mathrm{X}$ is connected, it is integral.
\end{proof}

\oldpage[IV]{160}
\begin{remark}[18.10.3.1]
\label{IV.18.10.3.1}
If $f:\mathrm{X}\to\mathrm{Y}$ is locally of finite presentation, dominant and quasi-compact, the fibre $f^{-1}(\eta)$ is non-empty by virtue of \sref[I]{I.1.1.5}.
On the other hand, if one does not suppose $f$ quasi-compact, $f$ \emph{can be unramified and dominant without being \'etale} (still assuming that the hypotheses of \sref{IV.18.10.3} are satisfied for $\mathrm{X}$ and $\mathrm{Y}$).
This is what the following example shows: one takes for $\mathrm{Y}$ the affine plane $\operatorname{Spec}(\mathbf{C}[\mathrm{T},\mathrm{U}])$; one considers on the other hand a family $(\mathrm{X}_j)_{j\in\mathbf{Z}}$ of preschemes isomorphic to the affine line $\operatorname{Spec}(\mathbf{C}[\mathrm{T}])$, and one forms the prescheme obtained by ``gluing together'' $\mathrm{X}_{2j}$ and $\mathrm{X}_{2j+1}$ at the point $-1$ and $\mathrm{X}_{2j+1}$ and $\mathrm{X}_{2j+2}$ at the point $+1$.
Let us define on the other hand $f:\mathrm{X}\to\mathrm{Y}$ as being equal on $\mathrm{X}_{2j}$ to the closed immersion whose image is the line of equation $y=2j$, transforming the point $-1$ to $(2j-1,2j)$ and the point $1$ to $(2j+1,2j)$; on $\mathrm{X}_{2j+1}$, $f_j$ is the closed immersion whose image is the line of equation $x=2j+1$, transforming the point $-1$ to $(2j+1,2j)$ and the point $1$ to $(2j+1,2j+2)$.
It is clear that $f$ is a local immersion (hence unramified) and dominant, but the fibre over the generic point of $\mathrm{Y}$ is empty and $f$ is not \'etale.
\end{remark}

\begin{corollary}[18.10.4]
\label{IV.18.10.4}
Let $g:\mathrm{Y}\to\mathrm{S}$ and $h:\mathrm{X}\to\mathrm{S}$ be two morphisms locally of finite presentation, $f:\mathrm{X}\to\mathrm{Y}$ an $\mathrm{S}$-morphism, $x$ a point of $\mathrm{X}$, $y=f(x)$, and $s=g(y)=h(x)$.
Suppose that $h$ is flat at the point $x$ and $g^{-1}(s)$ normal at the point $y$.
Then, for $f$ to be \'etale at the point $x$, it is necessary and sufficient that $f$ be unramified at the point $x$ and that $\dim_xh^{-1}(s)=\dim_yg^{-1}(s)$.
\end{corollary}

\begin{proof}
The conditions are necessary because if $f$ is \'etale at the point $x$, the same is true of $f_s:h^{-1}(s)\to g^{-1}(s)$, and the conclusion then follows from the fact that $f_s$ is flat and quasi-finite at the point $x$ \sref{IV.6.1.2}.
Conversely, if the conditions of the statement are satisfied, in order to see that $f$ is \'etale at the point $x$, it suffices, by virtue of \sref{IV.17.8.3}, to see that $f_s$ is so.
One can therefore restrict to the case where $\mathrm{S}$ is the spectrum of a field.
It then follows from the hypothesis $\dim_x\mathrm{X}=\dim_y\mathrm{Y}$ and from the fact that $f$ is unramified at the point $x$ that one has, by \sref{IV.5.2.3} and \sref{IV.17.4.1}, $\dim\sh{O}_{X,x}=\dim\sh{O}_{Y,y}$. Since $f$ is quasi-finite at the point $x$ and $\sh{O}_{Y,y}$ is integral, we deduce from \sref[0]{0.7.4.4} and \sref[0]{0.16.3.10} that the homomorphism $\sh{O}_{Y,y}\to\sh{O}_{X,x}$ is injective; it then suffices to apply \sref{IV.18.10.1}.
\end{proof}

\begin{corollary}[18.10.5]
\label{IV.18.10.5}
With the notations of \sref{IV.18.10.4}, suppose that $h$ is flat and that $g^{-1}(s)$ is normal for all $s\in\mathrm{S}$ (which will be the case for example if $g$ is smooth).
For $f$ to be an open immersion, it is necessary and sufficient that $f$ be a monomorphism and that, for all $x\in\mathrm{X}$, $\dim_xh^{-1}(s)=\dim_yg^{-1}(s)$, where $y=f(x)$ and $s=g(y)=h(x)$.
\end{corollary}

\begin{proof}
It remains only to prove the sufficiency of the conditions. It follows from \sref{IV.17.1.3} that a monomorphism locally of finite presentation is unramified; the hypotheses therefore imply that $f$ is \'etale by \sref{IV.18.10.4}, and one concludes by \sref{IV.17.9.1}.
\end{proof}

\begin{lemma}[18.10.6]
\label{IV.18.10.6}
Let $f:\mathrm{X}\to\mathrm{Y}$ be a flat and locally quasi-finite morphism $(\mathrm{Err}_{\mathrm{III}}, 20)$.
\begin{enumerate}[label=(\roman*)]
  \item The set $\mathrm{Max}(\mathrm{X})$ of maximal points of $\mathrm{X}$ is in canonical bijection with the set $\coprod_{y\in\mathrm{Max}(\mathrm{Y})}f^{-1}(y)$.
  \item If $\mathrm{Y}$ is irreducible with generic point $\eta$, the set of irreducible components of $\mathrm{X}$ is in canonical bijection with $f^{-1}(\eta)$, and is in particular finite if and only if $f^{-1}(\eta)$ is finite.
\oldpage[IV]{161}
  \item If the set of irreducible components of $\mathrm{Y}$ is locally finite, so is the set of irreducible components of $\mathrm{X}$, and in particular $\mathrm{X}$ is locally connected.
\end{enumerate}
\end{lemma}

\begin{proof}
The assertion (i) follows immediately from \sref{IV.2.3.4} and \sref[0]{0.2.1.6} and from the fact that the fibres $f^{-1}(y)$ are discrete sets.
It is clear that (ii) is a particular case of (i).
Finally, to prove (iii), one can, by virtue of (i), restrict to the case where $f$ is quasi-finite and where $\mathrm{Y}$ is irreducible \sref[0]{0.2.1.6} and the conclusion follows from (ii) and from the fact that the fibres of $f$ are finite sets \sref[II]{II.6.2.2}.
\end{proof}

\begin{proposition}[18.10.7]
\label{IV.18.10.7}
Let $\mathrm{Y}$ be a geometrically unibranch and irreducible prescheme (resp.\ integral), $f:\mathrm{X}\to\mathrm{Y}$ an \'etale morphism.
Then $\mathrm{X}$ is isomorphic to a disjoint union of irreducible (resp.\ integral) preschemes $\mathrm{X}_\lambda$, where $\lambda$ ranges over the fibre $f^{-1}(\eta)$ of the generic point $\eta$ of $\mathrm{Y}$.
The $\mathrm{X}_\lambda$ are geometrically unibranch; if $\mathrm{Y}$ is normal, the $\mathrm{X}_\lambda$ are normal.
\end{proposition}

\begin{proof}
Indeed, $\mathrm{X}$ is geometrically unibranch \sref{IV.17.5.7}, hence every point of $\mathrm{X}$ belongs to only one irreducible component. Moreover, by \sref{IV.18.10.6} and \sref{IV.17.6.1}, the family of irreducible components of $\mathrm{X}$ is locally finite; hence \sref[0]{0.2.1.6} these are the connected components $\mathrm{X}_\lambda$ of $\mathrm{X}$ ($\lambda\in f^{-1}(\eta)$), and $\mathrm{X}$ is their disjoint union. By \sref{IV.17.5.7}, if $\mathrm{Y}$ is integral (resp.\ normal), the $\mathrm{X}_\lambda$ are integral (resp.\ normal).
\end{proof}

\begin{proposition}[18.10.8]
\label{IV.18.10.8}
Let $\mathrm{Y}$ be a normal and integral prescheme, of generic point $\eta$, $\mathrm{K}=k(\eta)$ its field of rational functions, $f:\mathrm{X}\to\mathrm{Y}$ an \'etale morphism.
Suppose that $f^{-1}(\eta)$ is finite and non-empty, so that the $\mathrm{K}$-scheme $f^{-1}(\eta)$ is the spectrum of a finite separable $\mathrm{K}$-algebra $\mathrm{L}$, which is a direct product of finitely many fields $\mathrm{K}_i$ ($1\leq i\leq n$), finite separable extensions of $\mathrm{K}$ \sref{IV.17.6.2}.
Let $\mathrm{Y}'$ be the integral closure of $\mathrm{Y}$ in $\mathrm{L}$, isomorphic to the sum of the integral closures $\mathrm{Y}_i$ of $\mathrm{Y}$ in $\mathrm{K}_i$ \sref[II]{II.6.3.6}.
Then the morphism $f$ factorises uniquely as $\mathrm{X}\xrightarrow{f'}\mathrm{Y}'\xrightarrow{g}\mathrm{Y}$, such that the restriction $f_\eta':f^{-1}(\eta)\to g^{-1}(\eta)$ is the canonical isomorphism.
Moreover $f'$ is a local isomorphism; for $f'$ to be an open immersion, it is necessary and sufficient that $f$ be separated.
\end{proposition}

\begin{proof}
By virtue of \sref{IV.18.10.7}, one can restrict to the case where $\mathrm{X}$ is integral and normal, with $\mathrm{L}$ as its field of rational functions and $f$ dominant.
The existence and uniqueness of the factorisation considered for $f$ follow from \sref[II]{II.6.3.9}.
As $f'$ is birational and locally quasi-finite, the final assertions follow from \sref{IV.8.12.10}.
\end{proof}

\begin{corollary}[18.10.9]
\label{IV.18.10.9}
Under the hypotheses of \sref{IV.18.10.8}, for $f$ to be a finite morphism (in other words, for $\mathrm{X}$ to be an \'etale covering of $\mathrm{Y}$), it is necessary and sufficient that $\mathrm{X}$ be isomorphic to the integral closure $\mathrm{Y}'$ of $\mathrm{Y}$ in $\mathrm{L}$.
\end{corollary}

\begin{proof}
If $f$ is finite, the canonical injection $j:\mathrm{X}\to\mathrm{Y}'$ (which is an open immersion) is also a finite morphism \sref[II]{II.6.1.5},~(v), hence a closed morphism \sref[II]{II.6.1.10}, and as $j(\mathrm{X})$ is dense in $\mathrm{Y}'$ by \sref{IV.18.10.8}, we have $j(\mathrm{X})=\mathrm{Y}'$.
Conversely, as $\mathrm{Y}$ is normal and $\mathrm{L}$ is a direct product of finite separable extensions of $\mathrm{K}$, $\mathrm{Y}'$ is finite over $\mathrm{Y}$ (Bourbaki, \emph{Alg.\ comm.}, chap.~V, \textsection~1, n\textsuperscript{o}~6, cor.~1 of prop.~18), whence the corollary.
\end{proof}

\begin{env}[18.10.10]
\label{IV.18.10.10}
Let $\mathrm{Y}$ be a normal and integral prescheme, $\mathrm{K}=\mathrm{R}(\mathrm{Y})$ its field of rational functions.
We say that a finite-dimensional $\mathrm{K}$-algebra $\mathrm{L}$ is \emph{unramified over $\mathrm{Y}$} if: $1^{\circ}$ $\mathrm{L}$ is a separable $\mathrm{K}$-algebra, hence a direct product of finite separable extensions $\mathrm{K}_i$ of $\mathrm{K}$ ($1\leq i\leq n$); $2^{\circ}$ the integral closure $\mathrm{Y}'$ of $\mathrm{Y}$ in $\mathrm{L}$ (the disjoint union of the integral closures of $\mathrm{Y}$ in the $\mathrm{K}_i$) is \emph{unramified over} $\mathrm{Y}$ (which, by \sref{IV.18.10.3},
\oldpage[IV]{162}
is equivalent to saying that $\mathrm{Y}'$ is \'etale over $\mathrm{Y}$).
When $\mathrm{Y}=\operatorname{Spec}(\mathrm{A})$, where $\mathrm{A}$ is an integral, integrally closed ring and $\mathrm{K}$ its field of fractions, one also says (by a slight abuse of language, and when no confusion with the terminology of \sref{IV.17.3.2},~(ii) is to be feared) that $\mathrm{L}$ is unramified over $\mathrm{A}$ instead of ``unramified over $\mathrm{Y}$''.
\end{env}

\begin{remark}[18.10.11]
\label{IV.18.10.11}
Instead of saying that $\mathrm{L}$ is unramified over $\mathrm{Y}$, some authors also say that $\mathrm{L}$ is ``unramified over $\mathrm{K}$''; we shall not follow this potentially confusing usage.
\end{remark}

One can then express \sref{IV.18.10.9} in the following form:

\begin{corollary}[18.10.12]
\label{IV.18.10.12}
Let $\mathrm{Y}$ be a normal and integral prescheme, $\mathrm{K}=\mathrm{R}(\mathrm{Y})$ its field of rational functions.
Then the functor $\mathrm{X}\mapsto\mathrm{R}(\mathrm{X})$ is an equivalence between the category of \'etale coverings of $\mathrm{Y}$ and the category of finite \'etale $\mathrm{K}$-algebras unramified over $\mathrm{Y}$.
A quasi-inverse functor sends each finite \'etale $\mathrm{K}$-algebra $\mathrm{L}$ unramified over $\mathrm{Y}$ to the integral closure $\mathrm{Y}'$ of $\mathrm{Y}$ in $\mathrm{L}$.
\end{corollary}

\begin{proposition}[18.10.13]
\label{IV.18.10.13}
Let $\mathrm{Y}$ be a normal and integral prescheme, $\mathrm{K}=\mathrm{R}(\mathrm{Y})$ its field of rational functions.
\begin{enumerate}[label=(\roman*)]
  \item The field $\mathrm{K}$ is a $\mathrm{K}$-algebra unramified over $\mathrm{Y}$.
  \item Let $\mathrm{L}$ be a finite extension of $\mathrm{K}$ unramified over $\mathrm{Y}$, and let $\mathrm{Z}$ be the integral closure of $\mathrm{Y}$ in $\mathrm{L}$.
Then, if $\mathrm{M}$ is an $\mathrm{L}$-algebra unramified over $\mathrm{Z}$, $\mathrm{M}$ is also a $\mathrm{K}$-algebra unramified over $\mathrm{Y}$ (``transitivity'' of unramifiedness).
  \item Let $\mathrm{Y}'$ be a normal and integral prescheme, $\mathrm{K}'$ its field of rational functions, $g:\mathrm{Y}'\to\mathrm{Y}$ a dominant morphism.
If $\mathrm{L}$ is a $\mathrm{K}$-algebra unramified over $\mathrm{Y}$, then $\mathrm{L}\otimes_{\mathrm{K}}\mathrm{K}'$ is a $\mathrm{K}'$-algebra unramified over $\mathrm{Y}'$ (``base-change'' property).
Furthermore if $\mathrm{Y}=\operatorname{Spec}(\mathrm{A})$ and $\mathrm{Y}'=\operatorname{Spec}(\mathrm{A}')$ are affine, and if $\mathrm{C}$ is the integral closure of $\mathrm{A}$ in $\mathrm{L}$, then $\mathrm{C}'=\mathrm{C}\otimes_{\mathrm{A}}\mathrm{A}'$ is the integral closure of $\mathrm{A}'$ in $\mathrm{L}\otimes_{\mathrm{K}}\mathrm{K}'$.
\end{enumerate}
\end{proposition}

\begin{proof}
The assertion (i) is immediate, $\mathrm{Y}$ being its own normalisation in $\mathrm{K}$.
To prove (ii), let $\mathrm{Z}'$ be the integral closure of $\mathrm{Z}$ in $\mathrm{M}$, which is by hypothesis an \'etale covering of $\mathrm{Z}$.
As $\mathrm{Z}$ is \'etale and separated over $\mathrm{Y}$, so is $\mathrm{Z}'$ \sref{IV.17.3.3}, and it is clearly the integral closure of $\mathrm{Y}$ in $\mathrm{M}$.
As $\mathrm{M}$ is a finite separable $\mathrm{K}$-algebra, it follows from \sref{IV.18.10.10} that $\mathrm{M}$ is unramified over $\mathrm{Y}$.
Finally, to prove (iii), observe that if $\mathrm{Z}$ is the integral closure of $\mathrm{Y}$ in $\mathrm{L}$, then $\mathrm{Z}\times_{\mathrm{Y}}\mathrm{Y}'$ is \'etale and separated over $\mathrm{Y}'$ \sref{IV.17.3.3}, finite over $\mathrm{Y}'$ since $\mathrm{Z}$ is finite over $\mathrm{Y}$, and has $\mathrm{L}\otimes_{\mathrm{K}}\mathrm{K}'$ as its ring of rational functions \sref[I]{I.3.4.9}. Since $\mathrm{Y}'$ is normal, so is $\mathrm{Z}\times_{\mathrm{Y}}\mathrm{Y}'$ \sref{IV.17.5.7}; it is therefore the integral closure of $\mathrm{Y}'$ in $\mathrm{L}\otimes_{\mathrm{K}}\mathrm{K}'$, which completes the proof.
\end{proof}

\begin{corollary}[18.10.14]
\label{IV.18.10.14}
Let $\mathrm{Y}$ and $\mathrm{Y}'$ be two normal and integral preschemes, $\mathrm{K}$ and $\mathrm{K}'$ their respective fields of rational functions, and $g:\mathrm{Y}'\to\mathrm{Y}$ a dominant morphism, so that $\mathrm{K}'$ is an extension of $\mathrm{K}$.
Let $\mathrm{L}$ be a finite extension of $\mathrm{K}$ unramified over $\mathrm{Y}$, and let $\mathrm{L}_1$ be a composite extension (Bourbaki, \emph{Alg.}, chap.~VIII, \textsection~8, def.~1) of $\mathrm{L}$ and $\mathrm{K}'$.
Then $\mathrm{L}_1$ is an extension of $\mathrm{K}'$ unramified over $\mathrm{Y}'$.
Furthermore, if $\mathrm{Y}=\operatorname{Spec}(\mathrm{A})$ and $\mathrm{Y}'=\operatorname{Spec}(\mathrm{A}')$ are affine, and if $\mathrm{C}$ is the integral closure of $\mathrm{A}$ in $\mathrm{L}$, then the subring $\mathrm{C}_1=\mathrm{A}[\mathrm{C},\mathrm{A}']$ of $\mathrm{L}_1$ is the integral closure of $\mathrm{A}'$ in $\mathrm{L}_1$.
\end{corollary}

\begin{proof}
Indeed, $\mathrm{L}\otimes_{\mathrm{K}}\mathrm{K}'$ is a direct product of finite separable extensions of $\mathrm{K}'$, and $\mathrm{L}_1$
\oldpage[IV]{163}
is one of these extensions; hence the integral closure of $\mathrm{Y}'$ in $\mathrm{L}\otimes_{\mathrm{K}}\mathrm{K}'$ is a disjoint union of preschemes, one of which is the integral closure of $\mathrm{Y}'$ in $\mathrm{L}_1$.
The corollary therefore follows immediately from \sref{IV.18.10.13}.

When, for example, $\mathrm{A}$ is the ring of integers of an algebraic number field $\mathrm{K}$, and $\mathrm{K}'$ and $\mathrm{L}$ are algebraic extensions of $\mathrm{K}$, there are classical examples showing that the relation $\mathrm{C}_1=\mathrm{A}[\mathrm{C},\mathrm{A}']$ fails when $\mathrm{L}$ is not unramified over $\mathrm{A}$.
\end{proof}

\begin{env}[18.10.15]
\label{IV.18.10.15}
Suppose that $\mathrm{Y}=\operatorname{Spec}(\mathrm{A})$ is affine, $\mathrm{A}$ being integral and integrally closed; let $\mathrm{K}$ be its field of fractions and $\mathrm{L}$ a finite separable extension of $\mathrm{K}$, and suppose that the integral closure $\mathrm{C}$ of $\mathrm{A}$ in $\mathrm{L}$ is a \emph{projective} $\mathrm{A}$-module of finite type (which will be the case for example if $\mathrm{A}$ is a Dedekind ring (Bourbaki, \emph{Alg.\ comm.}, chap.~VII, \textsection~4, n\textsuperscript{o}~10, prop.~22)).
To say that $\mathrm{L}$ is unramified over $\mathrm{A}$ means that $\mathrm{C}$ is \'etale \sref{IV.18.10.10}; by \sref{IV.18.2.7},~(ii), this is equivalent to saying that the \emph{discriminant} $d_{\mathrm{C}/\mathrm{A}}$ of $\operatorname{Spec}(\mathrm{C})$ over $\operatorname{Spec}(\mathrm{A})$ is \emph{invertible} in $\mathrm{A}$.
In the particular case where $\mathrm{C}$ is a \emph{free} $\mathrm{A}$-module and $(x_i)_{1\leq i\leq n}$ a basis of $\mathrm{C}$ over $\mathrm{A}$, this means that $\det(\operatorname{Tr}_{\mathrm{C}/\mathrm{A}}(x_ix_j))$ is invertible in $\mathrm{A}$.
\end{env}

\begin{theorem}[18.10.16]
\label{IV.18.10.16}
Let $\mathrm{Y}$ be an integral prescheme, $\eta$ its generic point, and $f:\mathrm{X}\to\mathrm{Y}$ a separated morphism locally of finite type.
Suppose that every irreducible component of $\mathrm{X}$ dominates $\mathrm{Y}$ and that the generic fibre $\mathrm{X}_\eta=f^{-1}(\eta)$ is a finite prescheme over $\mathrm{K}=k(\eta)$; then $(\mathrm{X}_\eta)_{\mathrm{red}}=\operatorname{Spec}(\mathrm{L})$, where $\mathrm{L}=\prod_i\mathrm{L}_i$ is a product of finite extensions $\mathrm{L}_i$ of $\mathrm{K}$.
Set
\[
  n=[\mathrm{L}:\mathrm{K}]=\sum_i[\mathrm{L}_i:\mathrm{K}],\quad n_s=\sum_i[\mathrm{L}_i:\mathrm{K}]_s\quad\text{(sum of separable degrees of the $\mathrm{L}_i$)}.
\]
Let $y$ be a geometrically unibranch point of $\mathrm{Y}$, and denote by $n(y)$ the sum of the separable degrees over $k(y)$ of the residue fields of the isolated points of $f^{-1}(y)$.
Then:
\begin{enumerate}[label=(\roman*)]
  \item One has $n(y)\leq n_s\leq n$.
  \item Suppose $\mathrm{X}$ is reduced and the point $y$ is normal in $\mathrm{Y}$.
For $n(y)=n$, it is necessary and sufficient that there exists a neighbourhood $\mathrm{U}$ of $y$ in $\mathrm{Y}$ such that the restriction $f^{-1}(\mathrm{U})\to\mathrm{U}$ of $f$ is an \'etale and finite morphism.
\end{enumerate}
\end{theorem}

\begin{proof}
\emph{A) Reduction to the case where $f$ is quasi-finite.} ---
One knows \sref{IV.13.1.4} that the set $\mathrm{Z}$ of isolated points $x\in\mathrm{X}$ of $f^{-1}(f(x))$ is open in $\mathrm{X}$ and by hypothesis $\mathrm{Z}$ contains $f^{-1}(\eta)$.
Since $f^{-1}(\eta)$ is finite, $\mathrm{Z}$ is an increasing filtered union of quasi-compact open sets $\mathrm{V}_\lambda$ containing $f^{-1}(\eta)$ (dense in $\mathrm{X}$ by hypothesis).
If $n_\lambda(y)$ is the sum of the separable degrees over $k(y)$ of the residue fields of the points of $\mathrm{V}_\lambda\cap f^{-1}(y)$, we have $n(y)=\sup n_\lambda(y)$; to prove (i), it will therefore suffice to show that $n_\lambda(y)\leq n_s$.
On the other hand, if $n(y)=n$, there is an index $\lambda$ such that $n_\lambda(y)=n$. Suppose that we have shown that there then exists a neighbourhood $\mathrm{U}$ of $y$ in $\mathrm{Y}$ such that the restriction $\mathrm{V}_\lambda\cap f^{-1}(\mathrm{U})\to\mathrm{U}$ of $f$ is an \'etale and finite morphism; since $f$ is separated, the canonical injection $\mathrm{V}_\lambda\cap f^{-1}(\mathrm{U})\to f^{-1}(\mathrm{U})$ is then also a finite morphism \sref[II]{II.6.1.5},~(v), hence $\mathrm{V}_\lambda\cap f^{-1}(\mathrm{U})$ is closed in $f^{-1}(\mathrm{U})$; but being everywhere dense in $f^{-1}(\mathrm{U})$, it is necessarily equal to it.

\emph{B) Reduction to the case where $\mathrm{X}$ is reduced, $f$ is finite, and $\mathrm{Y}=\operatorname{Spec}(\sh{O}_{Y,y})$.} ---
We note that the hypothesis on $f$ implies that it is also satisfied for $f_{\mathrm{red}}$, and that the conclusion of (i) concerns only $f_{\mathrm{red}}$.
One can therefore replace $f$ everywhere by $f_{\mathrm{red}}$ and thus suppose $\mathrm{X}$ reduced.

\oldpage[IV]{164}
On the other hand, by virtue of \sref{IV.18.12.13} (${}^1$), one can (after having replaced $\mathrm{Y}$ by an open affine neighbourhood $\mathrm{V}$ of $y$ and $f$ by its restriction $f^{-1}(\mathrm{V})\to\mathrm{V}$) factorise $f$ as $\mathrm{X}\xrightarrow{j}\mathrm{X}'\xrightarrow{g}\mathrm{Y}$, where $g$ is \emph{finite} and $j$ an open immersion.
Since $\mathrm{X}$ is reduced, one can replace $\mathrm{X}'$ by the reduced closed sub-prescheme with underlying space $\overline{j(\mathrm{X})}$ \sref[I]{I.5.2.2}, hence suppose $\mathrm{X}'$ reduced and such that every irreducible component of $\mathrm{X}'$ dominates $\mathrm{Y}$.
Moreover, $j(\mathrm{X})$ is an open dense set in $\mathrm{X}'$ and has only a finite number of irreducible components by virtue of the hypothesis; hence \sref[0]{0.1.2.7} the irreducible components of $\mathrm{X}'$ are the closures of those of $j(\mathrm{X})$, and by construction the fibre $g^{-1}(\eta)$ identifies with $f^{-1}(\eta)$.
If one supposes the theorem proved when $f$ is finite, one can apply it to $g$, and as $f^{-1}(y)$ identifies with a sub-prescheme of $g^{-1}(y)$, the relation (i) for $g$ implies the same relation for $f$.
Suppose, on the other hand, that $n(y)=n$; then, by what precedes, one can apply the conclusion of (ii) to the morphism $g$, and after replacing $\mathrm{Y}$ if necessary by a neighbourhood of $y$, one may suppose that $g$ is \'etale, and consequently so is $f$; moreover, one then has $f^{-1}(y)=g^{-1}(y)$.
As the morphism $g$ is closed and $j(\mathrm{X})$ is open in $\mathrm{X}'$, there exists a neighbourhood $\mathrm{U}$ of $y$ in $\mathrm{Y}$ such that $j(\mathrm{X})\cap g^{-1}(\mathrm{U})=j(f^{-1}(\mathrm{U}))$, which proves that (ii) is also true for $f$.

We are thus reduced to the case where $f$ is moreover a \emph{finite} morphism.
It is immediate that, to prove (i), one can restrict to the case where $\mathrm{Y}=\operatorname{Spec}(\sh{O}_{Y,y})$.
Let us show that the same is true for proving (ii).
Indeed, suppose it has been proved that the condition $n(y)=n$ implies that $\mathrm{X}\times_{\mathrm{Y}}\operatorname{Spec}(\sh{O}_{Y,y})$ is \'etale and finite over $\operatorname{Spec}(\sh{O}_{Y,y})$. Since this implies that $f$ is flat and formally unramified at the points of $f^{-1}(y)$, and since $\mathrm{Y}$ is \emph{integral}, it follows from \sref{IV.18.4.13} that $f$ is \emph{\'etale at the points} of $f^{-1}(y)$, hence also on a neighbourhood $\mathrm{V}$ of $f^{-1}(y)$ in $\mathrm{X}$; but since $f$ is finite, hence closed, $\mathrm{V}$ contains an open set of the form $f^{-1}(\mathrm{U})$, where $\mathrm{U}$ is a neighbourhood of $y$ in $\mathrm{Y}$.

\emph{C) Reduction to the case where $\mathrm{Y}=\operatorname{Spec}(\sh{O}_{Y,y})$ and $\sh{O}_{Y,y}$ is strictly local.} ---
Set $\mathrm{A}=\sh{O}_{Y,y}$, and let $\mathrm{A}'={}^{\mathrm{hs}}\mathrm{A}$ be the strict henselisation of $\mathrm{A}$ \sref{IV.18.8.7}, which is an integral, geometrically unibranch local ring \sref{IV.18.8.16}.
Let $\mathrm{Y}'=\operatorname{Spec}(\mathrm{A}')$, and let $y'$ and $\eta'$ be the closed point and the generic point of $\mathrm{Y}'$; $y'$ is above $y$ and since the morphism $g:\mathrm{Y}'\to\mathrm{Y}$ is flat \sref{IV.18.8.8}, $\eta'$ is above $\eta$ \sref{IV.2.3.4}.
Set $\mathrm{X}'=\mathrm{X}\times_{\mathrm{Y}}\mathrm{Y}'$, $f'=f_{(\mathrm{Y}')}:\mathrm{X}'\to\mathrm{Y}'$; $f'$ is finite and since $g$ is flat, every irreducible component of $\mathrm{X}'$ dominates $\mathrm{Y}'$ \sref{IV.2.3.7}.
If $n'(y')$, $n'$ and $n_s'$ are the numbers defined for $f'$ in the same way as $n(y)$, $n$ and $n_s$ for $f$, we have $n_s'=n_s$ and $n'(y')=n(y)$ \sref[I]{I.6.4.8}, and it therefore amounts to proving the assertion (i) for $f$ or for $f'$.
Moreover, if $\mathrm{X}$ is reduced and $y$ is a normal point of $\mathrm{Y}$ (hence $\mathrm{A}$ is integrally closed), then $\mathrm{A}'$ is integrally closed \sref{IV.18.8.12}.
On the other hand, it follows from the definition \sref{IV.18.8.7}, from the remark \sref{IV.8.1.2},~$a$) and from the theorem of the double inductive limit, that $\mathrm{A}'$ is the inductive limit of rings $\mathrm{A}_\alpha$ such that the morphisms $\operatorname{Spec}(\mathrm{A}_\alpha)\to\operatorname{Spec}(\mathrm{A})=\mathrm{Y}$ are \'etale; consequently $\mathrm{X}'$ is the projective limit of the $\mathrm{X}_\alpha=\mathrm{X}\otimes_{\mathrm{A}}\mathrm{A}_\alpha$, which are \'etale over $\mathrm{X}$, hence reduced since $\mathrm{X}$ is \sref{IV.17.5.7}; by passage to the limit,

\footnotetext{(${}^1$) The reader will verify that \sref{IV.18.10.16} is not used in the proof of \sref{IV.18.12.13}. If $f$ is supposed locally of finite presentation (resp.\ quasi-projective), one can replace \sref{IV.18.12.13} by \sref{IV.8.12.8} (resp.\ \sref{IV.8.12.6}).}

\oldpage[IV]{165}
we deduce that $\mathrm{X}'$ is reduced \sref{IV.8.7.1}.
This being so, if $n(y)=n$, one has $n_s=n$, and consequently the residue fields at the points of $\mathrm{X}_\eta$ are necessarily separable over $k(\eta)$; the residue fields at the points of $\mathrm{X}_{\eta'}'$ are then algebraic and separable over $k(\eta')$ \sref{IV.4.6.1}, in other words $n_s'=n'=n$. It follows from what precedes that $n'(y')=n'$.
The morphism $f'$ therefore has the same properties as $f$, and if one proves that $f'$ is \'etale, it will follow from \sref{IV.17.7.1} that $f$ is \'etale.

\emph{D) End of the proof.} ---
Suppose therefore that $\mathrm{A}$ is strictly local.
Since $\mathrm{A}$ is henselian and the morphism $f$ is finite, it follows from \sref{IV.18.5.11} that if $x_j$ ($1\leq j\leq m$) are the distinct points of $f^{-1}(y)$, $\mathrm{X}$ is the \emph{sum} of open sub-preschemes $\operatorname{Spec}(\sh{O}_{X,x_j})=\mathrm{X}_j$, which are the connected components of $\mathrm{X}$.
By hypothesis, each of them dominates $\mathrm{Y}$, hence meets $f^{-1}(\eta)$, and so the number of points of $f^{-1}(\eta)$ is $\geq m$; \emph{a fortiori} we have $n_s\geq m$.
But since $k(y)$ is separably closed and the $k(x_j)$ are algebraic over $k(y)$, they are necessarily radicial extensions of $k(y)$; hence $m=n(y)$, and this proves (i).
Suppose now the hypotheses of (ii) satisfied; the relations $n(y)=n_s=n$ imply first of all, since every irreducible component of $\mathrm{X}$ dominates $\mathrm{Y}$, that each of the $\mathrm{X}_j$ is \emph{irreducible}; moreover, as the $\mathrm{X}_j$ are \emph{reduced}, they are \emph{integral}; finally the relations $m=n_s=n$ imply that for the generic point $\xi_j$ of $\mathrm{X}_j$, $k(\xi_j)$ is a separable extension of $k(\eta)$ of degree $1$, hence isomorphic to $k(\eta)$; in other words, the restriction $f|\mathrm{X}_j:\mathrm{X}_j\to\mathrm{Y}$ of $f$ is \emph{finite and birational}; moreover, $\mathrm{X}_j$ is \emph{integral} and by hypothesis $\mathrm{Y}$ is \emph{normal}; hence \sref{IV.8.12.10.1} $f|\mathrm{X}_j$ is an \emph{isomorphism}, which proves that $f$ is \'etale, and finishes the proof of \sref{IV.18.10.16}.
\end{proof}

\begin{remark}[18.10.17]
\label{IV.18.10.17}
The method of ``\'etale localisation'' used in the proof of \sref{IV.18.10.16} also allows one to improve the results of \sref{IV.15.5.1} by eliminating the Noetherian hypothesis.
One has first of all the following result which generalises \sref{IV.15.5.2}:
\end{remark}

\begin{env}[18.10.17.1]
\label{IV.18.10.17.1}
Let $f:\mathrm{X}\to\mathrm{Y}$ be a separated and quasi-finite morphism, $y$ a point of $\mathrm{Y}$ such that $f$ is universally open at the points of $f^{-1}(y)$.
Then, for every generisation $y'$ of $y$, we have $n(y)\leq n(y')$.
\end{env}

\begin{proof}
One can replace $\mathrm{Y}$ by the reduced closed sub-prescheme $\mathrm{Z}$ with underlying space $\overline{\{y'\}}$ and $\mathrm{X}$ by $f^{-1}(\mathrm{Z})$, hence suppose $\mathrm{Y}$ integral with generic point $y'$.
Using \sref{IV.18.12.13} (as in \sref{IV.18.10.16},~B)), one can suppose that $\mathrm{X}$ is an open dense subset of a prescheme $\mathrm{X}'$ and that $f$ is the restriction to $\mathrm{X}$ of a finite morphism $g:\mathrm{X}'\to\mathrm{Y}$.
Let us show that $g^{-1}(y')=f^{-1}(y')$.
Indeed, let $i:g^{-1}(y')\to\mathrm{X}'$ be the canonical injection, which is a flat morphism since $y'$ is the generic point of $\mathrm{Y}$ \sref[I]{I.3.6.5}; one can write $f^{-1}(y')=i^{-1}(\mathrm{X})$.
On the other hand, the canonical injection $j:\mathrm{X}\to\mathrm{X}'$ is a morphism of finite type since the composite $g\circ j=f$ is of finite type and $g$ is separated \sref{IV.1.5.4}; hence $\mathrm{X}$ is a retrocompact open subset of $\mathrm{X}'$, and consequently pro-constructible in $\mathrm{X}'$ \sref{IV.1.9.5},~(v).
One concludes therefore from \sref{IV.2.3.10} that $i^{-1}(\overline{\mathrm{X}})=\overline{i^{-1}(\mathrm{X})}$, in other words $f^{-1}(y')$ is dense in $g^{-1}(y')$; but as $g^{-1}(y')$ is discrete, we necessarily have $f^{-1}(y')=g^{-1}(y')$.

We are thus reduced to proving that the sum of the numbers $[k(x):k(y)]_s$ where $x$ ranges over the points of $g^{-1}(y)$ where $g$ is universally open, is at most equal to
\oldpage[IV]{166}
$n(y')$ (for the morphism $g$).
Using then the fact that the property of being universally open at a point is preserved by base change, one shows, as in \sref{IV.18.10.16},~C)), that one can reduce to the case where $\mathrm{Y}=\operatorname{Spec}(\sh{O}_{Y,y})$, with $\sh{O}_{Y,y}$ \emph{strictly local}.
Applying \sref{IV.18.5.11},~$c$)), we see that $\mathrm{X}$ is the disjoint union of the open subschemes $\mathrm{X}_j=\operatorname{Spec}(\sh{O}_{X,x_j})$, where $x_j$ ranges over $g^{-1}(y)$. The hypothesis that $g$ is open at one of these points $x_j$ implies that the corresponding subscheme $\mathrm{X}_j$ dominates $\mathrm{Y}$ \sref{IV.1.10.3}; one concludes as in \sref{IV.18.10.16},~D)).

We then deduce from \sref{IV.18.10.17.1} that \emph{the conclusions of \sref{IV.15.5.1} are still valid when one no longer supposes $\mathrm{Y}$ locally Noetherian, but only that $f$ is separated, quasi-finite, and of finite presentation}.
\end{proof}

Indeed, to prove assertion (i) of \sref{IV.15.5.1}, one remarks that, since $f$ is of finite presentation, the set $\mathrm{E}$ of points $z\in\mathrm{Y}$ such that $n(z)\geq n(y)$ is locally constructible \sref{IV.9.7.9}; to show that $y$ is interior to $\mathrm{E}$, it suffices, by virtue of \sref{IV.1.10.1}, to see that every generisation $y'$ of $y$ belongs to $\mathrm{E}$, which is nothing other than \sref{IV.18.10.17.1}.

To prove assertion (ii) of \sref{IV.15.5.1}, one may, since $f$ is of finite presentation, use \sref{IV.8.10.5},~(xii) according to the method of \sref{IV.8.1.2},~$a$)), and thus first suppose that $\mathrm{Y}=\operatorname{Spec}(\sh{O}_{Y,y})$. Applying next \sref{IV.2.7.1},~(vii), one sees that after the faithfully flat base change $\mathrm{Y}'\to\mathrm{Y}$, where $\mathrm{Y}'=\operatorname{Spec}({}^{\mathrm{hs}}\sh{O}_{Y,y})$, one may suppose that the ring $\sh{O}_{Y,y}$ is \emph{strictly local}.
Applying \sref{IV.18.5.11},~$c$)), we see then that if $x_j$ ($1\leq j\leq n$) are the points of $f^{-1}(y)$, $\mathrm{X}$ is the \emph{sum} of $n$ open sub-preschemes $\mathrm{X}_j=\operatorname{Spec}(\sh{O}_{X,x_j})$ which are \emph{finite} over $\mathrm{Y}$, and of an open prescheme $\mathrm{X}''$.
If one proves that $\mathrm{X}''=\varnothing$, one will have shown that $f$ is finite, hence proper. Now the fields $k(x_j)$, being algebraic extensions of the separably closed field $k(y)$, are radicial extensions, and consequently $n(y)=n$; since $f$ is open at each point $x_j$, $\mathrm{X}_j$ dominates $\mathrm{Y}$ \sref{IV.1.10.3}, and the restriction of $f$ to $\mathrm{X}_j$, being finite, is surjective \sref[II]{II.6.1.10}. The hypothesis that $z\mapsto n(z)$ is constant on $\mathrm{Y}$ therefore implies that $\mathrm{X}''\cap f^{-1}(z)=\varnothing$ for every $z\in\mathrm{Y}$, that is, $\mathrm{X}''=\varnothing$.

Finally, the proof of (iii) reproduces without change that given in \sref{IV.15.5.1}.

\begin{env}[18.10.17.2]
\label{IV.18.10.17.2}
By analogous methods using strict henselisation, one can, in most of the results of \textsection\textsection~14 and 15, eliminate the Noetherian hypotheses.
\end{env}

\begin{lemma}[18.10.18]
\label{IV.18.10.18}
Let $\mathrm{X}$, $\mathrm{Y}$ be two preschemes, $f:\mathrm{X}\to\mathrm{Y}$ a birational morphism (${}^1$), $x$ a point of $\mathrm{X}$.
For $f$ to be an isomorphism local at the point $x$, it is necessary and sufficient that $f$ be \'etale at the point $x$.
For $f$ to be an open immersion, it is necessary and sufficient that $f$ be \'etale and separated.
\end{lemma}

\footnotetext{(${}^1$) We understand by this the notion defined in \sref{IV.6.15.4}, but where we do not suppose $\mathrm{X}$ and $\mathrm{Y}$ reduced.}

\begin{proof}
The conditions stated being trivially necessary, all that is needed is to see that they are sufficient.
For the first assertion, the question is local on $\mathrm{X}$ and $\mathrm{Y}$; hence one may suppose that $f$ is \'etale and that $\mathrm{X}$ and $\mathrm{Y}$ are affine, so that $f$ is \emph{separated}, and one is reduced to proving the second assertion.
By virtue of \sref{IV.17.9.1}, it is a question of proving that $f$ is \emph{radicial}.
Let $y\in\mathrm{Y}$, and let us prove that $f^{-1}(y)$ is radicial on $k(y)$.
The base change $\operatorname{Spec}(\sh{O}_{Y,y})\to\mathrm{Y}$ not changing the fact that $f$ is \'etale, separated, and birational, one can restrict to the case where $\mathrm{Y}=\operatorname{Spec}(\mathrm{A})$, with $\mathrm{A}=\sh{O}_{Y,y}$.
\oldpage[IV]{167}
Let $\mathrm{A}'={}^{\mathrm{hs}}\mathrm{A}$, which is a local ring and a faithfully flat $\mathrm{A}$-module \sref{IV.18.8.8}, the homomorphism $\mathrm{A}\to\mathrm{A}'$ being local.
Set $\mathrm{Y}'=\operatorname{Spec}(\mathrm{A}')$, $f'=f_{(\mathrm{Y}')}:\mathrm{X}'\to\mathrm{Y}'$, so that $f'$ is \'etale and separated; moreover, as the morphism $\mathrm{Y}'\to\mathrm{Y}$ is flat, the reasoning of \sref{IV.6.15.4.1} shows that $f'$ is also birational.
If $y'$ is the closed point of $\mathrm{Y}'$, it suffices, by \sref{IV.2.6.1},~(v), to prove that $f'^{-1}(y')$ is radicial over $k(y')$; since $f'$ is \'etale, it is enough to prove that $f'^{-1}(y')$ contains at most one point \sref{IV.17.6.1},~$c'$).

Suppose therefore $\mathrm{Y}=\operatorname{Spec}(\mathrm{A})$, $\mathrm{A}$ strictly local, $f$ separated, \'etale and birational, and let us show that $f^{-1}(y)$ can contain only a single point.
Indeed, if $f^{-1}(y)$ contained two distinct points $x_1$, $x_2$, there would exist two $\mathrm{Y}$-sections $u_1$, $u_2$ of $f$ such that $u_1(y)=x_1$ and $u_2(y)=x_2$ \sref{IV.18.8.1}; but since $f$ is \'etale and separated, $u_1$ and $u_2$ would be isomorphisms of $\mathrm{Y}$ onto two connected components (open) of $\mathrm{X}$ \sref{IV.17.9.4} and there would then be two maximal points of $\mathrm{X}$ at least above every maximal point of $\mathrm{Y}$, which contradicts the hypothesis that $f$ is \emph{birational}.
\end{proof}

\begin{proposition}[18.10.19]
\label{IV.18.10.19}
Let $\mathrm{Y}$ be a reduced prescheme whose set of irreducible components is locally finite, $f:\mathrm{X}\to\mathrm{Y}$ a separated morphism, locally of finite type and formally unramified, $g$ a $\mathrm{Y}$-rational section of $f$ (i.e.\ a $\mathrm{Y}$-rational map from $\mathrm{Y}$ to $\mathrm{X}$), $\mathrm{U}$ the domain of definition of $g$ \sref[I]{I.7.2.1}, and let $\mathrm{Z}$ be the closure of $g(\mathrm{U})$ in $\mathrm{X}$.
Then, for all points $y\in\mathrm{Y}-\mathrm{U}$ such that $y$ is geometrically unibranch, $\mathrm{Z}_y=\mathrm{Z}\cap f^{-1}(y)=\varnothing$.
In particular, if $\mathrm{Y}$ is geometrically unibranch, $g(\mathrm{U})$ is a subset of $\mathrm{X}$ that is both open and closed.
If, for every closed irreducible $\mathrm{T}$ of $\mathrm{X}$, $f(\mathrm{T})$ is closed in $\mathrm{Y}$, $g$ is defined at every geometrically unibranch point of $\mathrm{Y}$.
\end{proposition}

\begin{proof}
Since $\mathrm{Y}$ is reduced and $f$ separated, it follows from \sref[I]{I.7.2.2} that there exists a $\mathrm{U}$-section $u$ of $f^{-1}(\mathrm{U})$ belonging to the class $g$; moreover, since $f$ is formally unramified, \sref{IV.17.4.1.2} shows that $u$ is an isomorphism of $\mathrm{U}$ onto the subscheme induced on the open set $u(\mathrm{U})$ of $\mathrm{X}$.
If one denotes also by $\mathrm{Z}$ the reduced sub-prescheme of $\mathrm{X}$ having $\mathrm{Z}$ for underlying space, $g(\mathrm{U})=u(\mathrm{U})$, being reduced, is induced also by $\mathrm{Z}$, and by construction $f'=f|\mathrm{Z}$ is a \emph{birational} morphism of $\mathrm{Z}$ to $\mathrm{Y}$, moreover formally unramified since $f$ is \sref{IV.17.1.3}.
The question being local on $\mathrm{Y}$, and since $\mathrm{Y}$ is reduced and geometrically unibranch (hence integral) at the point $y$, one may, after replacing $\mathrm{Y}$ by a neighbourhood of $y$, restrict to the case where $\mathrm{Y}$ is \emph{integral}; then $u(\mathrm{U})$ is irreducible, hence so is $\mathrm{Z}$, and consequently $\mathrm{Z}$ is integral as well.
Since $f'$ is dominant, \sref[I]{I.8.2.7} shows that the homomorphism $\sh{O}_{Y,y}\to\sh{O}_{Z,x}$ is injective for every $x\in\mathrm{Z}_y$, and \sref{IV.18.10.2},~(iii) implies that $f'$ is \'etale. Being separated and birational, it is a local isomorphism at $x$ by \sref{IV.18.10.18}; but this would imply that the $\mathrm{U}$-section $u$ extends to an open set strictly larger than $\mathrm{U}$, contrary to the definition of $\mathrm{U}$.
We therefore necessarily have $\mathrm{Z}_y=\varnothing$, which establishes the first assertion; the second is an obvious consequence.
Finally, under the last hypotheses of the statement, as $\mathrm{Z}$ is closed and irreducible, $f(\mathrm{Z})$ is closed in $\mathrm{Y}$, hence equal to $\mathrm{Y}$, i.e.\ $\mathrm{Z}_y\neq\varnothing$ for all $y\in\mathrm{Y}$, which finishes the proof.
\end{proof}

\begin{remark}[18.10.20]
\label{IV.18.10.20}
Let $\mathrm{Y}$ be a locally Noetherian prescheme.
One says that a morphism $f:\mathrm{X}\to\mathrm{Y}$ is \emph{essentially proper} if it is locally of finite type and if, for every $\mathrm{Y}$-scheme of the form $\mathrm{Y}'=\operatorname{Spec}(\mathrm{A})$, where $\mathrm{A}$ is a discrete valuation ring, the canonical map \sref[II]{II.7.3.2.2} relative to the $\mathrm{Y}'$-prescheme $\mathrm{X}\times_{\mathrm{Y}}\mathrm{Y}'$ is bijective.
\oldpage[IV]{168}
The reasoning of \sref[II]{II.7.3.8} proves that this is the same as saying that $f$ (supposed locally of finite type) is \emph{separated} and that for every base change $\mathrm{Y}''\to\mathrm{Y}$, where $\mathrm{Y}''$ is \emph{locally Noetherian}, the image by $f_{(\mathrm{Y}'')}: \mathrm{X}_{(\mathrm{Y}'')}\to\mathrm{Y}''$ of every closed irreducible subset of $\mathrm{X}''$ is \emph{closed}.
To say that $f$ is \emph{proper} (for $\mathrm{Y}$ locally Noetherian) means therefore that $f$ is \emph{essentially proper and of finite type} \sref[II]{II.7.3.8}; but one encounters important examples of essentially proper morphisms that are not of finite type (for example certain ``Picard preschemes'' or certain ``preschemes of N\'eron-Severi'' (chap.~VI)).
It follows immediately from \sref{IV.18.10.19} that if $\mathrm{Y}$ is \emph{locally Noetherian, reduced, and geometrically unibranch}, and if $f:\mathrm{X}\to\mathrm{Y}$ is \emph{essentially proper and unramified}, then every $\mathrm{Y}$-rational section of $f$ is \emph{defined everywhere}.
\end{remark}

The following proposition generalises the smoothness criterion \sref{IV.17.15.5} for a morphism $\mathrm{X}\to\mathrm{Y}$ when $\mathrm{Y}$ is not the spectrum of a field:

\begin{proposition}[18.10.21]
\label{IV.18.10.21}
Let $f:\mathrm{X}\to\mathrm{Y}$ be a morphism locally of finite presentation, $x$ a point of $\mathrm{X}$, $y=f(x)$, and suppose that the ring $\sh{O}_{Y,y}$ is integral and geometrically unibranch.
Let $r$ be the minimum number of generators of the $\sh{O}_{X,x}$-module $(\Omega^1_{\mathrm{X}/\mathrm{Y}})_x$ (equal to the rank of the $k(x)$-vector space $(\Omega^1_{\mathrm{X}/\mathrm{Y}})_x\otimes_{\sh{O}_{X,x}}k(x)$).
Then, for $f$ to be smooth at the point $x$, it is necessary and sufficient that, if $y'$ is the generic point of the unique irreducible component of $\mathrm{Y}$ containing $y$, there exists a generisation $x'$ of $x$ such that $f(x')=y'$ and $\dim_{x'}f^{-1}(y')\geq r$.
\end{proposition}

\begin{proof}
If $f$ is smooth at the point $x$, it is also so in an open neighbourhood $\mathrm{U}$ of $x$, hence in every generisation $x'$ of $x$, and $\dim_{x'}(f^{-1}(f(x')))=r$ in each of these points, by virtue of \sref{IV.17.10.2}.
As $f$ is moreover flat in $\mathrm{U}$, there exists a generisation of $x$ above every generisation of $y$ \sref{IV.2.3.4}, which proves that the condition is necessary.
To show that it is sufficient, let us first note that by \sref{IV.17.5.1}, it suffices to prove that the morphism deduced from $f$ by the base change $\operatorname{Spec}(\sh{O}_{Y,y})\to\mathrm{Y}$ is smooth at the point $x$; one can therefore suppose that $\mathrm{Y}=\operatorname{Spec}(\sh{O}_{Y,y})$.

The question is local on $\mathrm{X}$; hence, by \sref{IV.16.4.22} and \sref[0]{0.5.2.2}, one may suppose that $\mathrm{X}$ is affine and that there exist $r$ sections $s_i$ ($1\leq i\leq r$) of $\sh{O}_X$ over $\mathrm{X}$ such that the sections $d_{\mathrm{X}/\mathrm{Y}}(s_i)$ generate $\Omega^1_{\mathrm{X}/\mathrm{Y}}$.
Let
\[
  g:\mathrm{X}\to\mathrm{Z}=\mathrm{Y}[\mathrm{T}_1,\ldots,\mathrm{T}_r]=\mathbf{V}_{\mathrm{Y}}^r
\]
be the $\mathrm{Y}$-morphism corresponding to the homomorphism $\sh{O}_{\mathrm{Y}}^r\to f_*(\sh{O}_X)$ of $\sh{O}_Y$-modules defined by the $s_i$ (considered as sections of $f_*(\sh{O}_X)$ over $\mathrm{Y}$) \sref[II]{II.1.2.7}.
We have the exact sequence \sref{IV.16.4.19}
\[
  g^*(\Omega^1_{\mathrm{Z}/\mathrm{Y}})\to\Omega^1_{\mathrm{X}/\mathrm{Y}}\to\Omega^1_{\mathrm{X}/\mathrm{Z}}\to 0.
\]
As the $d\mathrm{T}_i$ ($1\leq i\leq r$) generate $\Omega^1_{\mathrm{Z}/\mathrm{Y}}$, it follows from the definition of $g$ that the homomorphism $g^*(\Omega^1_{\mathrm{Z}/\mathrm{Y}})\to\Omega^1_{\mathrm{X}/\mathrm{Y}}$ is surjective, hence $\Omega^1_{\mathrm{X}/\mathrm{Z}}=0$, and $g$ is consequently \emph{unramified} \sref{IV.17.4.1}.
We shall now see that $g$ is \emph{\'etale} at the point $x$, and taking into account \sref{IV.17.3.8}, this will prove that $f$ is smooth at the point $x$.
As $\sh{O}_{Y,y}$ is assumed to be integral and geometrically
\oldpage[IV]{169}
unibranch, so is $\sh{O}_{Z,z}$ by \sref{IV.17.5.7}, where $z=g(x)$, since the structural morphism $\mathrm{Z}\to\mathrm{Y}$ is smooth \sref{IV.17.3.8}.
By virtue of \sref{IV.18.10.1}, it suffices to show that the homomorphism $\sh{O}_{Z,z}\to\sh{O}_{X,x}$ is injective; it amounts to the same thing to see that the morphism $\operatorname{Spec}(\sh{O}_{X,x})\to\operatorname{Spec}(\sh{O}_{Z,z})$ corresponding is \emph{dominant}, since $\sh{O}_{Z,z}$ is integral \sref[I]{I.1.2.7}.
Let $y'$ be the generic point of $\mathrm{Y}=\operatorname{Spec}(\sh{O}_{Y,y})$, and let $z'$ be the generic point of the integral scheme $\mathrm{Z}=\operatorname{Spec}(\sh{O}_{Y,y}[\mathrm{T}_1,\ldots,\mathrm{T}_r])$; if $h:\mathrm{Z}\to\mathrm{Y}$ is the structural morphism, then $h(z')=y'$.
It suffices to prove that there exists in $f^{-1}(y')$ a generisation $x'$ of $x$ such that the image of $x'$ by the morphism $g_{y'}:f^{-1}(y')\to h^{-1}(y')$ equals $z'$.
Moreover, $z'$ is also the generic point of $h^{-1}(y')=\operatorname{Spec}(k(y')[\mathrm{T}_1,\ldots,\mathrm{T}_r])$; since the morphism $g_{y'}$ is unramified, hence quasi-finite \sref{IV.17.4.2}, it follows from \sref{IV.5.4.1} that the restriction of $g_{y'}$ to every irreducible component of \emph{dimension} $\geq r$ of $f^{-1}(y')$ is dominant.
Now, there exists by hypothesis a component containing a generisation $x'$ of $x$; \emph{a fortiori} its generic point is a generisation of $x$, whence the conclusion.
\end{proof}

\subsection{Application to Noetherian local algebras complete over a field}
\label{subsection:IV.18.11}

The following lemma generalizes \sref[0]{0.21.9.1} and \sref[0]{0.21.9.2}:

\begin{lemma}[18.11.1]
\label{IV.18.11.1}
Let $k$ be a field, let $\mathrm{A}$ be a complete Noetherian local ring that is a $k$-algebra, and let $\mathrm{K}$ be the residue field of $\mathrm{A}$.
\begin{enumerate}[label=(\roman*)]
  \item If the $\mathrm{K}$-vector space $\Omega^1_{\mathrm{K}/k}$ is of finite rank, the $\mathrm{A}$-module $\hat{\Omega}^1_{\mathrm{A}/k}$ is of finite type.
  \item If, moreover, $\mathrm{A}$ is formally smooth as a $k$-algebra (for the adic topology), $\hat{\Omega}^1_{\mathrm{A}/k}$ is a free $\mathrm{A}$-module of rank
\[
  \dim(\mathrm{A})+\operatorname{rg}_{\mathrm{K}}(\Omega^1_{\mathrm{K}/k})-\operatorname{rg}_{\mathrm{K}}(\Upsilon_{\mathrm{K}/k}).
\]
\end{enumerate}
Furthermore, for every subfield $k_0$ of $k$ such that $\Omega^1_{k/k_0}$ is a finite-dimensional $k$-vector space, $\hat{\Omega}^1_{\mathrm{A}/k_0}$ is a free $\mathrm{A}$-module of rank
\[
  \dim(\mathrm{A})+\operatorname{rg}_{\mathrm{K}}(\Omega^1_{\mathrm{K}/k_0})-\operatorname{rg}_{\mathrm{K}}(\Upsilon_{\mathrm{K}/k_0})+\operatorname{rg}_{k}(\Omega^1_{k/k_0}).
\]
\end{lemma}

\begin{proof}
Assertion (i) is recalled here, having been proved in \sref[0]{0.20.7.15}.
To prove (ii), note that this assertion was in fact established by the proof of \sref[0]{0.21.9.2}; the statement of \sref[0]{0.21.9.2} uses the finiteness of the transcendence degree of $\mathrm{K}$ over $k$ only through Cartier's equality \sref[0]{0.21.7.1}.
\end{proof}

\begin{lemma}[18.11.2]
\label{IV.18.11.2}
Let $k$ be a field of characteristic $p>0$, and let $\mathrm{C}$ be an integral complete Noetherian local $k$-algebra whose residue field is a finite extension of $k$; let $n$ be its dimension and $\mathrm{L}$ its field of fractions.
Then $\Omega^1_{\mathrm{L}/k}$ and $\Upsilon_{\mathrm{L}/k}$ are $\mathrm{L}$-vector spaces of finite rank, and one has
\begin{equation}
\label{IV.18.11.2.1}
\tag{18.11.2.1}
  \operatorname{rg}_{\mathrm{L}}(\Omega^1_{\mathrm{L}/k})-\operatorname{rg}_{\mathrm{L}}(\Upsilon_{\mathrm{L}/k})\geq n.
\end{equation}
\oldpage[IV]{170}
Suppose furthermore that $[k:k^p]<+\infty$.
Then one has the equality
\begin{equation}
\label{IV.18.11.2.2}
\tag{18.11.2.2}
  \operatorname{rg}_{\mathrm{L}}(\Omega^1_{\mathrm{L}/k})-\operatorname{rg}_{\mathrm{L}}(\Upsilon_{\mathrm{L}/k})=n.
\end{equation}
Furthermore, $\Omega^1_{\mathrm{C}/k}$ is then isomorphic to $\hat{\Omega}^1_{\mathrm{C}/k}$, and hence, by \sref{IV.18.11.1}, is a $\mathrm{C}$-module of finite type.
\end{lemma}

\begin{proof}
Indeed, by \sref[0]{0.19.8.9}, there exists a $k$-subalgebra $\mathrm{C}_0$ of $\mathrm{C}$ isomorphic to the formal power series algebra $k[[\mathrm{T}_1,\ldots,\mathrm{T}_n]]$ and such that $\mathrm{C}$ is finite over $\mathrm{C}_0$.
Consequently, $\mathrm{L}$ is a \emph{finite} extension of the field of fractions $\mathrm{L}_0=k((\mathrm{T}_1,\ldots,\mathrm{T}_n))$ of $\mathrm{C}_0$.
Applying \sref[0]{0.20.6.17.1} to the prime subfield of $k$ and to the three fields $k\subset\mathrm{L}_0\subset\mathrm{L}$ (and taking \sref[0]{0.20.6.21},~(i), into account) gives the exact sequence of $\mathrm{L}$-vector spaces
\[
  0\to\Upsilon_{\mathrm{L}_0/k}\otimes_{\mathrm{L}_0}\mathrm{L}\to\Upsilon_{\mathrm{L}/k}\to\Upsilon_{\mathrm{L}/\mathrm{L}_0}\to\Omega^1_{\mathrm{L}_0/k}\otimes_{\mathrm{L}_0}\mathrm{L}\to\Omega^1_{\mathrm{L}/k}\to\Omega^1_{\mathrm{L}/\mathrm{L}_0}\to 0.
\]
Since $\mathrm{L}$ is a finite extension of $\mathrm{L}_0$, the $\mathrm{L}$-vector spaces $\Omega^1_{\mathrm{L}/\mathrm{L}_0}$ and $\Upsilon_{\mathrm{L}/\mathrm{L}_0}$ have finite and equal ranks, by Cartier's equality \sref[0]{0.21.7.1}. Since $\mathrm{L}_0$ is separable over $k$ \sref[0]{0.21.9.6.4}, one has $\Upsilon_{\mathrm{L}_0/k}=0$ \sref[0]{0.20.6.3}. The preceding exact sequence therefore shows that $\Upsilon_{\mathrm{L}/k}$ has finite rank and, in every case,
\[
  \operatorname{rg}_{\mathrm{L}}(\Omega^1_{\mathrm{L}/k})-\operatorname{rg}_{\mathrm{L}}(\Upsilon_{\mathrm{L}/k})=\operatorname{rg}_{\mathrm{L}_0}(\Omega^1_{\mathrm{L}_0/k})-\operatorname{rg}_{\mathrm{L}_0}(\Upsilon_{\mathrm{L}_0/k}).
\]
It therefore suffices to prove \sref{IV.18.11.2.1} and \sref{IV.18.11.2.2} after replacing $\mathrm{C}$ and $\mathrm{L}$ by $\mathrm{C}_0$ and $\mathrm{L}_0$. Since $\mathrm{L}_0$ is separable over $k$ \sref[0]{0.21.9.6.4}, one has $\Upsilon_{\mathrm{L}_0/k}=0$ \sref[0]{0.20.6.3}; on the other hand,
\[
  \Omega^1_{\mathrm{L}_0/k}=\Omega^1_{\mathrm{C}_0/k}\otimes_{\mathrm{C}_0}\mathrm{L}_0
  \qquad\text{\sref[0]{0.20.5.9}.}
\]
We know from \sref[0]{0.21.9.4} that, if $\mathrm{C}_1=k[[\mathrm{T}_1^p,\ldots,\mathrm{T}_n^p]]$, then $\hat{\Omega}^1_{\mathrm{C}_0/k}$ identifies with $\Omega^1_{\mathrm{C}_0/\mathrm{C}_1}$. Moreover,
$\Omega^1_{\mathrm{C}_0/k}=\Omega^1_{\mathrm{C}_0/k[\mathrm{C}_0^p]}$ \sref[0]{0.21.1.5}, and, since $\mathrm{C}_0^p=k^p[[\mathrm{T}_1^p,\ldots,\mathrm{T}_n^p]]$, one has $k[\mathrm{C}_0^p]\subset\mathrm{C}_1$, with equality when $[k:k^p]<+\infty$. The exact sequence
\[
  \Omega^1_{\mathrm{C}_0/k[\mathrm{C}_0^p]}\longrightarrow\Omega^1_{\mathrm{C}_0/\mathrm{C}_1}\longrightarrow 0
  \qquad\text{\sref[0]{0.20.5.7}}
\]
therefore yields an exact sequence $\Omega^1_{\mathrm{L}_0/k}\to\hat{\Omega}^1_{\mathrm{C}_0/k}\otimes_{\mathrm{C}_0}\mathrm{L}_0\to0$. Since $\hat{\Omega}^1_{\mathrm{C}_0/k}$ is a free $\mathrm{C}_0$-module of rank $n$ \sref[0]{0.21.9.3}, it follows that $\operatorname{rg}_{\mathrm{L}_0}(\Omega^1_{\mathrm{L}_0/k})\geq n$ in every case, with equality when $[k:k^p]<+\infty$. This proves \sref{IV.18.11.2.1} and \sref{IV.18.11.2.2}.

Finally, when $[k:k^p]<+\infty$, to see that $\Omega^1_{\mathrm{C}/k}\to\hat{\Omega}^1_{\mathrm{C}/k}$ is an isomorphism, it suffices to prove that $\Omega^1_{\mathrm{C}/k}$ is a finite-type $\mathrm{C}$-module, since $\mathrm{C}$ is a complete Noetherian local ring \sref[0]{0.7.3.3}. Since $\Omega^1_{\mathrm{C}/k}\cong\Omega^1_{\mathrm{C}/k[\mathrm{C}^p]}$ and $k[\mathrm{C}_0^p]\subset k[\mathrm{C}^p]$, this reduces to proving that $\mathrm{C}$ is finite over $k[\mathrm{C}_0^p]$ \sref[0]{0.20.4.7}; this follows because $\mathrm{C}$ is finite over $\mathrm{C}_0$, while $\mathrm{C}_0$ is finite over $k[\mathrm{C}_0^p]$ under the hypothesis $[k:k^p]<+\infty$.
\end{proof}

\begin{lemma}[18.11.3]
\label{IV.18.11.3}
Let $k$ be a field, let $p$ be its characteristic exponent, and suppose that $[k:k^p]<+\infty$. Let $\mathrm{A}$ be a complete Noetherian local $k$-algebra whose residue field is a finite extension of $k$. Let $\mathfrak{p}$ be a prime ideal of $\mathrm{A}$ such that $\mathrm{A}_{\mathfrak{p}}$ is geometrically regular over $k$ \sref{IV.6.7.6}, so that there is a unique minimal prime ideal $\mathfrak{q}$ of $\mathrm{A}$ contained in $\mathfrak{p}$. Then $(\hat{\Omega}^1_{\mathrm{A}/k})_{\mathfrak{p}}$ is a free $\mathrm{A}_{\mathfrak{p}}$-module of rank $\dim(\mathrm{A}/\mathfrak{q})$.
\end{lemma}

\begin{proof}
Since $\mathrm{A}$ is Noetherian and $\operatorname{Spec}(\mathrm{A})$ is regular (and hence, a fortiori, integral) at $\mathfrak{p}$, the point $\mathfrak{p}$ lies on only one irreducible component of $\operatorname{Spec}(\mathrm{A})$. Thus it contains

\oldpage[IV]{171}

only one minimal ideal $\mathfrak{q}$ of $\mathrm{A}$, and moreover $\mathfrak{q}_{\mathfrak{p}}=0$. Put $\mathrm{B}=\mathrm{A}/\mathfrak{q}$. The sequence \sref[0]{0.20.7.20}
\[
  \mathfrak{q}/\mathfrak{q}^2\longrightarrow
  \hat{\Omega}^1_{\mathrm{A}/k}\widehat{\otimes}_{\mathrm{A}}\mathrm{B}
  \longrightarrow\hat{\Omega}^1_{\mathrm{B}/k}\longrightarrow0
\]
is exact. Indeed, $\hat{\Omega}^1_{\mathrm{A}/k}$ is a finite-type $\mathrm{A}$-module by \sref{IV.18.11.1}, and hence
\[
  \hat{\Omega}^1_{\mathrm{A}/k}\widehat{\otimes}_{\mathrm{A}}\mathrm{B}
  =\hat{\Omega}^1_{\mathrm{A}/k}\otimes_{\mathrm{A}}\mathrm{B}
  =\hat{\Omega}^1_{\mathrm{A}/k}/\mathfrak{q}\hat{\Omega}^1_{\mathrm{A}/k}.
\]
As this $\mathrm{B}$-module is of finite type, every one of its submodules is closed \sref[0]{0.7.3.5}, taking \sref[0]{0.20.4.5} into account. The image of $\mathfrak{q}/\mathfrak{q}^2$ is dense in the kernel of the displayed homomorphism \sref[0]{0.20.7.20}, and is therefore equal to that kernel. Localizing at $\mathfrak{p}$ also shows that the canonical homomorphism $(\hat{\Omega}^1_{\mathrm{A}/k})_{\mathfrak{p}}\to(\hat{\Omega}^1_{\mathrm{B}/k})_{\mathfrak{p}/\mathfrak{q}}$ is bijective. We may consequently reduce to the case $\mathfrak{q}=0$, that is, to the case where $\mathrm{A}$ is integral. We distinguish two cases.

\emph{I) $p>1$.} Apply \sref{IV.18.11.2} to the quotient algebra $\mathrm{C}=\mathrm{A}/\mathfrak{p}$, whose field of fractions $\mathrm{K}$ is the residue field of $\mathrm{A}_{\mathfrak{p}}$. Thus
\begin{equation}
\label{IV.18.11.3.1}
\tag{18.11.3.1}
  \operatorname{rg}_{\mathrm{K}}(\Omega^1_{\mathrm{K}/k})-
  \operatorname{rg}_{\mathrm{K}}(\Upsilon_{\mathrm{K}/k})=\dim(\mathrm{A}/\mathfrak{p}).
\end{equation}
Now $\mathrm{A}_{\mathfrak{p}}$ is formally smooth as a $k$-algebra for its $\mathfrak{p}\mathrm{A}_{\mathfrak{p}}$-adic topology \sref[0]{0.19.6.6}; hence $\Omega^1_{\mathrm{A}_{\mathfrak{p}}/k}$ is formally projective \sref[0]{0.20.4.9} for that topology \sref[0]{0.20.4.5}. Moreover,
$\Omega^1_{\mathrm{A}_{\mathfrak{p}}/k}=(\Omega^1_{\mathrm{A}/k})_{\mathfrak{p}}$ \sref[0]{0.20.5.9} is of finite type by \sref{IV.18.11.2} applied to $\mathrm{A}$. For every integer $j$, the module
\[
  \Omega^1_{\mathrm{A}_{\mathfrak{p}}/k}/\mathfrak{p}^{j+1}\Omega^1_{\mathrm{A}_{\mathfrak{p}}/k}
\]
is therefore projective over $\mathrm{A}_{\mathfrak{p}}/\mathfrak{p}^{j+1}\mathrm{A}_{\mathfrak{p}}$, of rank
$m=\operatorname{rg}_{\mathrm{K}}(\Omega^1_{\mathrm{A}_{\mathfrak{p}}/k}\otimes_{\mathrm{A}_{\mathfrak{p}}}\mathrm{K})$ \sref[0]{0.19.2.4}. It follows from \sref[0]{0.10.2.1} and \sref[0]{0.10.1.3} that $\hat{\Omega}^1_{\mathrm{A}_{\mathfrak{p}}/k}$ is a free $\mathrm{A}_{\mathfrak{p}}$-module of rank $m$. Let $\mathrm{A}'=(\mathrm{A}_{\mathfrak{p}})^{\wedge}$ be the completion of $\mathrm{A}_{\mathfrak{p}}$ for its $\mathfrak{p}\mathrm{A}_{\mathfrak{p}}$-adic topology. The algebra $\mathrm{A}'$ is again formally smooth over $k$ for its adic topology \sref[0]{0.19.3.6}, and \sref[0]{0.20.7.14} and \sref[0]{0.20.4.5} give $\hat{\Omega}^1_{\mathrm{A}'/k}=\hat{\Omega}^1_{\mathrm{A}_{\mathfrak{p}}/k}$. Thus $\hat{\Omega}^1_{\mathrm{A}'/k}$ is a free $\mathrm{A}'$-module of rank $m$. By \sref{IV.18.11.1} and $\dim(\mathrm{A}')=\dim(\mathrm{A}_{\mathfrak{p}})$ \sref[0]{0.16.2.4},
\begin{equation}
\label{IV.18.11.3.2}
\tag{18.11.3.2}
  m=\dim(\mathrm{A}_{\mathfrak{p}})+
  \bigl(\operatorname{rg}_{\mathrm{K}}(\Omega^1_{\mathrm{K}/k})-
  \operatorname{rg}_{\mathrm{K}}(\Upsilon_{\mathrm{K}/k})\bigr).
\end{equation}
By \sref{IV.18.11.3.1}, this is $m=\dim(\mathrm{A}/\mathfrak{p})+\dim(\mathrm{A}_{\mathfrak{p}})$. Since the complete Noetherian local ring $\mathrm{A}$ is a quotient of a regular ring \sref[0]{0.19.8.8}, one has
$\dim(\mathrm{A})=\dim(\mathrm{A}/\mathfrak{p})+\dim(\mathrm{A}_{\mathfrak{p}})$ by \sref[0]{0.16.5.12}, which proves the assertion in this case.

\emph{II) $p=1$.} As above, $\dim(\mathrm{A})=\dim(\mathrm{A}/\mathfrak{p})+\dim(\mathrm{A}_{\mathfrak{p}})$. Put $n=\dim(\mathrm{A})$, $r=\dim(\mathrm{A}/\mathfrak{p})$, and $s=\dim(\mathrm{A}_{\mathfrak{p}})$. We shall prove that $(\hat{\Omega}^1_{\mathrm{A}/k})_{\mathfrak{p}}$ is a free $\mathrm{A}_{\mathfrak{p}}$-module of rank $n$. We first prove the following lemma.

\begin{lemma}[18.11.3.3]
\label{IV.18.11.3.3}
Let $\mathrm{A}$ be a Noetherian local ring and $\mathfrak{p}$ a prime ideal such that $\mathrm{A}_{\mathfrak{p}}$ is Cohen--Macaulay. For every system of parameters $(t_i)_{1\leq i\leq s}$ of $\mathrm{A}_{\mathfrak{p}}$, there exists a sequence $(x_i)_{1\leq i\leq s}$ of elements of $\mathfrak{p}$ such that the $x_i$ form part of a system of parameters of $\mathrm{A}$
\oldpage[IV]{172}
and, for every $i$, $x_i/1$ is congruent to $t_i$ modulo $t_1\mathrm{A}_{\mathfrak{p}}+\cdots+t_{i-1}\mathrm{A}_{\mathfrak{p}}$. In particular, if $\mathrm{A}_{\mathfrak{p}}$ is regular and the $t_i$ form a regular system of parameters of $\mathrm{A}_{\mathfrak{p}}$, the same is true of the $x_i/1$.
\end{lemma}

\begin{proof}[Proof of \sref{IV.18.11.3.3}]
This lemma is itself a consequence of the following one.

\begin{lemma}[18.11.3.4]
\label{IV.18.11.3.4}
Let $\mathrm{A}$ be a Noetherian local ring, let $\mathfrak{p}$ be a prime ideal of $\mathrm{A}$, and let $x\in\mathrm{A}$. Then there exists $x'\in\mathrm{A}$ such that $x'/1=x/1$ in $\mathrm{A}_{\mathfrak{p}}$ and $x'$ belongs to none of the minimal prime ideals of $\mathrm{A}$ not contained in $\mathfrak{p}$.
\end{lemma}
We first show how \sref{IV.18.11.3.4} implies \sref{IV.18.11.3.3}. Multiplying the $t_i$, if necessary, by units of $\mathrm{A}_{\mathfrak{p}}$, we may assume that $t_i=x_i/1$ with $x_i\in\mathfrak{p}$ for $1\leq i\leq s$. We argue by induction on $s$. Since $t_1$ forms part of a system of parameters of $\mathrm{A}_{\mathfrak{p}}$, it belongs to no minimal prime ideal of $\mathrm{A}_{\mathfrak{p}}$ \sref[0]{0.16.5.5}; hence no $x_1'\in\mathrm{A}$ with $x_1'/1=t_1$ can belong to a minimal prime ideal of $\mathrm{A}$ contained in $\mathfrak{p}$. By \sref{IV.18.11.3.4}, we may moreover choose $x_1'$ such that $x_1'/1=t_1$ (so $x_1'\in\mathfrak{p}$) and $x_1'$ belongs to no minimal prime ideal of $\mathrm{A}$. Thus $x_1'$ forms part of a system of parameters of $\mathrm{A}$ \sref[0]{0.16.3.4} and \sref{IV.16.3.7}. We then apply induction to $\mathrm{A}'=\mathrm{A}/x_1'\mathrm{A}$ and $\mathfrak{p}'=\mathfrak{p}/x_1'\mathrm{A}$. Since $\mathrm{A}'_{\mathfrak{p}'}=\mathrm{A}_{\mathfrak{p}}/t_1\mathrm{A}_{\mathfrak{p}}$, this ring is Cohen--Macaulay \sref[0]{0.16.5.5}, and the images $t_i'$ ($2\leq i\leq s$) form a system of parameters of it. The induction hypothesis applied to $\mathrm{A}'$ and the $t_i'$ completes the construction. The last assertion follows because, in $\mathrm{A}_{\mathfrak{p}}$, a system of parameters that generates the maximal ideal is a regular system of parameters \sref[0]{0.17.1.1}.

It remains to prove \sref{IV.18.11.3.4}. Let $(\mathfrak{p}'_k)_{1\leq k\leq r}$ be the minimal prime ideals of $\mathrm{A}$ not contained in $\mathfrak{p}$ and containing $x$, and let $(\mathfrak{p}_j)_{1\leq j\leq n}$ be the other minimal prime ideals of $\mathrm{A}$; we may assume $r\geq1$. No $\mathfrak{p}'_k$ contains $\bigcap_{1\leq j\leq n}\mathfrak{p}_j$, so there exists
$y\in\bigcap_{1\leq j\leq n}\mathfrak{p}_j$ belonging to none of the $\mathfrak{p}'_k$ (Bourbaki, \emph{Alg.\ comm.}, chap.~II, \textsection~1, no.~1, prop.~2). Moreover, $y/1$ belongs to every minimal prime ideal of $\mathrm{A}_{\mathfrak{p}}$, and is therefore nilpotent. If $(y/1)^h=0$, the element $x'=x+y^h$ has the required properties: $x'/1=x/1$; since $y^h\notin\mathfrak{p}'_k$ and $x\in\mathfrak{p}'_k$, one has $x'\notin\mathfrak{p}'_k$ for $1\leq k\leq r$; and if $\mathfrak{p}_j$ is a minimal prime ideal not contained in $\mathfrak{p}$ with $x\notin\mathfrak{p}_j$, then $x'\notin\mathfrak{p}_j$ because $y\in\mathfrak{p}_j$.
\end{proof}

We return to the proof of \sref{IV.18.11.3} when $p=1$. By \sref{IV.18.11.3.3}, there exists a regular system of parameters $(t_i)_{1\leq i\leq s}$ of $\mathrm{A}_{\mathfrak{p}}$ such that $t_i=x_i/1$, where the $x_i$ ($1\leq i\leq s$) belong to $\mathfrak{p}$ and form part of a system of parameters $(x_j)_{1\leq j\leq n}$ of $\mathrm{A}$. Put $\mathrm{A}_0=k[[\mathrm{T}_1,\ldots,\mathrm{T}_n]]$. Since $\mathrm{A}$ is a complete $k$-algebra and the $x_j$ belong to its maximal ideal $\mathfrak{m}$, there is a local $k$-homomorphism $u:\mathrm{A}_0\to\mathrm{A}$ such that $u(\mathrm{T}_j)=x_j$ for $1\leq j\leq n$ (Bourbaki, \emph{Alg.\ comm.}, chap.~III, \textsection~4, no.~5, prop.~6). If $\mathfrak{n}$ is the ideal generated by the $x_j$, it is an ideal of definition of $\mathrm{A}$; hence \sref[0]{0.7.4.4} and \sref{IV.7.4.3} show that $u$ makes $\mathrm{A}$ finite over $\mathrm{A}_0$.

Put $\mathfrak{p}_0=\sum_{j=1}^{s}\mathrm{A}_0\mathrm{T}_j$, $\mathrm{B}_0=(\mathrm{A}_0)_{\mathfrak{p}_0}$, and $\mathrm{B}=\mathrm{A}\otimes_{\mathrm{A}_0}\mathrm{B}_0$, so that $u_{\mathfrak{p}_0}:\mathrm{B}_0\to\mathrm{B}$ makes
\oldpage[IV]{173}
$\mathrm{B}$ finite over $\mathrm{B}_0$. Moreover, if $\mathfrak{p}'$ is the ideal of $\mathrm{B}$ generated by $\mathfrak{p}$, then $\mathrm{B}_{\mathfrak{p}'}=\mathrm{A}_{\mathfrak{p}}$; since $\mathfrak{p}'$ contains $\mathfrak{p}_0\mathrm{B}$ by construction, it lies over the maximal ideal $\mathfrak{p}_0\mathrm{B}_0$ of $\mathrm{B}_0$. The morphism $\operatorname{Spec}(\mathrm{B})\to\operatorname{Spec}(\mathrm{B}_0)$ is unramified at $\mathfrak{p}'$: indeed, $k(\mathfrak{p}')$ is a finite extension of the characteristic-zero field $k(\mathfrak{p}_0)$, hence is necessarily separable, and $\mathrm{B}_{\mathfrak{p}'}/\mathfrak{p}_0\mathrm{B}_{\mathfrak{p}'}=k(\mathfrak{p}')$ by the choice of the $x_j$ for $1\leq j\leq s$ \sref{IV.17.4.1}.

We now record the following lemma.

\begin{lemma}[18.11.3.5]
\label{IV.18.11.3.5}
Let $k$ be a field, and let $\mathrm{R}$ and $\mathrm{S}$ be complete Noetherian local $k$-algebras such that, if $\mathrm{K}$ is the residue field of $\mathrm{S}$, $\Omega^1_{\mathrm{K}/k}$ has finite rank over $\mathrm{K}$. Let $u:\mathrm{R}\to\mathrm{S}$ be a $k$-homomorphism making $\mathrm{S}$ finite over $\mathrm{R}$. Then $\hat{\Omega}^1_{\mathrm{S}/\mathrm{R}}=\Omega^1_{\mathrm{S}/\mathrm{R}}$, and the sequence
\begin{equation}
\label{IV.18.11.3.6}
\tag{18.11.3.6}
  \hat{\Omega}^1_{\mathrm{R}/k}\otimes_{\mathrm{R}}\mathrm{S}\xrightarrow{v}\hat{\Omega}^1_{\mathrm{S}/k}\xrightarrow{w}\Omega^1_{\mathrm{S}/\mathrm{R}}\to 0
\end{equation}
(cf.\ \sref[0]{0.20.7.17.3}) is exact.
\end{lemma}

\begin{proof}
The first assertion follows because $\Omega^1_{\mathrm{S}/\mathrm{R}}$ is a finite-type $\mathrm{S}$-module \sref[0]{0.20.4.7} and by \sref[0]{0.7.3.6}. Moreover, \sref[0]{0.20.7.17.3} says that the image of $v$ is dense in $\ker(w)$ and that $w$ is surjective. By \sref{IV.18.11.1}, however, $\hat{\Omega}^1_{\mathrm{S}/k}$ is a finite-type $\mathrm{S}$-module, and therefore every one of its submodules is closed \sref[0]{0.7.3.5}. The lemma follows.
\end{proof}

Apply the lemma with $\mathrm{R}=\mathrm{A}_0$ and $\mathrm{S}=\mathrm{A}$. The module $\hat{\Omega}^1_{\mathrm{A}_0/k}$ is free of rank $n$ over $\mathrm{A}_0$ \sref[0]{0.21.9.3}; hence
$\hat{\Omega}^1_{\mathrm{A}_0/k}\widehat{\otimes}_{\mathrm{A}_0}\mathrm{A}=\hat{\Omega}^1_{\mathrm{A}_0/k}\otimes_{\mathrm{A}_0}\mathrm{A}$, and this is a free $\mathrm{A}$-module of rank $n$. Since $\operatorname{Spec}(\mathrm{B})$ is unramified over $\operatorname{Spec}(\mathrm{B}_0)$ at $\mathfrak{p}'$, one has
$(\Omega^1_{\mathrm{A}/\mathrm{A}_0})_{\mathfrak{p}}=\Omega^1_{\mathrm{B}_{\mathfrak{p}'}/\mathrm{B}_0}=0$ \sref{IV.16.4.15} and \sref{IV.17.4.1}. Localizing \sref{IV.18.11.3.6}, applied to $\mathrm{A}_0$ and $\mathrm{A}$, at $\mathfrak{p}$ therefore gives a surjective homomorphism
\[
  (\hat{\Omega}^1_{\mathrm{A}_0/k}\otimes_{\mathrm{A}_0}\mathrm{A})_{\mathfrak{p}}
  \longrightarrow(\hat{\Omega}^1_{\mathrm{A}/k})_{\mathfrak{p}}.
\]
Thus $(\hat{\Omega}^1_{\mathrm{A}/k})_{\mathfrak{p}}$ has a system of $n$ generators. On the other hand, $\hat{\Omega}^1_{\mathrm{A}/k}$ has rank $n$ by \sref[0]{0.21.9.5}, which applies to the complete integral ring $\mathrm{A}$ by Cohen's theorem \sref[0]{0.19.8.8},~(ii), and because the field of fractions of $\mathrm{A}$ has characteristic zero. Hence $(\hat{\Omega}^1_{\mathrm{A}/k})_{\mathfrak{p}}$ also has rank $n$; its quotient by its torsion submodule, being generated by $n$ elements, is necessarily free. It follows at once that the $n$ generators obtained above form a free system, which proves the assertion.
\end{proof}

\begin{lemma}[18.11.4]
\label{IV.18.11.4}
Let $k$ be a field, $\mathrm{A}$ a $k$-algebra, local Noetherian and complete, whose residue field is a finite extension of $k$.
Let $k'$ be an arbitrary extension of $k$, and set $\mathrm{A}'=\mathrm{A}\widehat{\otimes}_kk'$.
Then:
\begin{enumerate}[label=(\roman*)]
  \item $\mathrm{A}'$ is a semi-local Noetherian ring, complete, composed as a direct product of complete local rings $\mathrm{A}_i'$ ($1\leq i\leq r$) which are $\mathrm{A}$-modules faithfully flat, and whose residue fields are finite extensions of $k'$; if $\mathfrak{m}$ is the maximal ideal of $\mathrm{A}$, $\mathfrak{m}\mathrm{A}'$ is an ideal of definition of $\mathrm{A}'$; and $\dim(\mathrm{A}_i')=\dim(\mathrm{A})$ for all $i$.
  \item For every $i$, $\hat{\Omega}^1_{\mathrm{A}_i'/k'}$ is canonically isomorphic to $\hat{\Omega}^1_{\mathrm{A}/k}\otimes_{\mathrm{A}}\mathrm{A}_i'$.
\end{enumerate}
\end{lemma}

\begin{proof}
\oldpage[IV]{174}
The first assertions follow immediately from \sref{IV.7.5.5} and \sref[0]{0.6.6.2}; the fact that $\mathfrak{m}\mathrm{A}'$ is an ideal of definition of $\mathrm{A}'$ also follows from \sref{IV.7.5.5}, since $\mathrm{K}$ is the residue field of $\mathrm{A}$, $\mathrm{K}\otimes_k k'$ is a finite $k'$-algebra.
Finally, $\mathrm{A}_i'/\mathfrak{m}\mathrm{A}_i'$ is one of the Artinian local rings occurring as direct factors of $\mathrm{K}\otimes_k k'$ \sref{IV.7.5.5}, and is therefore zero-dimensional. Since $\mathrm{A}_i'$ is flat over $\mathrm{A}$, the equality $\dim(\mathrm{A}_i')=\dim(\mathrm{A})$ follows from \sref{IV.6.1.2}.

(ii) By the hypothesis on $\mathrm{K}$, $\hat{\Omega}^1_{\mathrm{A}/k}$ is an $\mathrm{A}$-module of finite type \sref{IV.18.11.1}; hence $\hat{\Omega}^1_{\mathrm{A}/k}\otimes_{\mathrm{A}}\mathrm{A}'$ is complete and identifies with the completed tensor product $\hat{\Omega}^1_{\mathrm{A}/k}\widehat{\otimes}_{\mathrm{A}}\mathrm{A}'$. By associativity of completed tensor products \sref[0]{0.7.7.4},
\[
  \hat{\Omega}^1_{\mathrm{A}/k}\otimes_{\mathrm{A}}\mathrm{A}'
  =\hat{\Omega}^1_{\mathrm{A}/k}\otimes_{\mathrm{A}}(\mathrm{A}\widehat{\otimes}_kk')
  \simeq\hat{\Omega}^1_{\mathrm{A}/k}\widehat{\otimes}_kk'.
\]
Now $\Omega^1_{\mathrm{A}/k}\otimes_kk'$ identifies with $\Omega^1_{\mathrm{A}''/k'}$, where $\mathrm{A}''=\mathrm{A}\otimes_kk'$ \sref[0]{0.20.5.5}; and since $\mathrm{A}'$ is by definition the separated completion of $\mathrm{A}''$, the separated completion of $\Omega^1_{\mathrm{A}''/k'}$ identifies by construction with that of $\Omega^1_{\mathrm{A}'/k'}$ \sref[0]{0.20.7.4}. In other words, there is a canonical isomorphism
\[
  \hat{\Omega}^1_{\mathrm{A}'/k'}\xrightarrow{\sim}
  \hat{\Omega}^1_{\mathrm{A}/k}\otimes_{\mathrm{A}}\mathrm{A}'.
\]
The conclusion of (ii) now follows because $\Omega^1_{\mathrm{A}'/k'}$ is the direct sum of the $\Omega^1_{\mathrm{A}_i'/k'}$ \sref[0]{0.20.4.13}. Indeed, if $\mathfrak r$ is the radical of $\mathrm{A}'$, its $\mathfrak r$-adic topology on $\Omega^1_{\mathrm{A}'/k'}$ is the product of the $\mathfrak m_i'$-adic topologies on the $\Omega^1_{\mathrm{A}_i'/k'}$, where $\mathfrak m_i'$ is the maximal ideal of $\mathrm{A}_i'$. Thus the separated completion $\hat{\Omega}^1_{\mathrm{A}'/k'}$ is the product of the separated completions $\hat{\Omega}^1_{\mathrm{A}_i'/k'}$, and it remains only to apply \sref[0]{0.20.4.5}.
\end{proof}

\begin{proposition}[18.11.5]
\label{IV.18.11.5}
Let $k$ be a field, $\mathrm{A}$ a $k$-algebra, local Noetherian and complete, whose residue field is a finite extension of $k$, $\mathfrak{p}$ a prime ideal of $\mathrm{A}$ distinct from the maximal ideal $\mathfrak{m}$, such that there exists a minimal prime ideal $\mathfrak{q}\subset\mathfrak{p}$ for which $\dim(\mathrm{A}/\mathfrak{q})=\dim(\mathrm{A})$ (which will in particular hold when $\mathrm{A}$ is equidimensional), $m$ an integer $\geq 0$.
The following conditions are equivalent:
\begin{enumerate}[label=\alph*)]
  \item The $\mathrm{A}_\mathfrak{p}$-module $(\hat{\Omega}^1_{\mathrm{A}/k})_\mathfrak{p}$ admits a system of $m$ generators.
  \item There exists a local $k$-homomorphism $u:\mathrm{B}\to\mathrm{A}$, where $\mathrm{B}=k[[\mathrm{T}_1,\ldots,\mathrm{T}_m]]$, making $\mathrm{A}$ a finite $\mathrm{B}$-algebra, and such that the corresponding morphism $\operatorname{Spec}(\mathrm{A})\to\operatorname{Spec}(\mathrm{B})$ is unramified at $\mathfrak{p}$.
\end{enumerate}
\end{proposition}

\begin{proof}
To prove that $b)$ implies $a)$, note that by Lemma \sref{IV.18.11.3.5} one has the exact sequence
\[
  \hat{\Omega}^1_{\mathrm{B}/k}\otimes_\mathrm{B}\mathrm{A}\to\hat{\Omega}^1_{\mathrm{A}/k}\to\Omega^1_{\mathrm{A}/\mathrm{B}}\to 0
\]
since $\mathrm{A}$ is a $\mathrm{B}$-algebra finite; localising at $\mathfrak{p}$ and noting that by hypothesis $(\Omega^1_{\mathrm{A}/\mathrm{B}})_\mathfrak{p}=0$ \sref{IV.17.4.1}, one obtains a surjective homomorphism $(\hat{\Omega}^1_{\mathrm{B}/k}\otimes_\mathrm{B}\mathrm{A})_\mathfrak{p}\to(\hat{\Omega}^1_{\mathrm{A}/k})_\mathfrak{p}$, and the conclusion follows because $\hat{\Omega}^1_{\mathrm{B}/k}$ is a free $\mathrm{B}$-module of rank $m$ \sref[0]{0.21.9.3}.

To prove that $a)$ implies $b)$, let us first prove the following lemmas.
\end{proof}

\begin{lemma}[18.11.5.1]
\label{IV.18.11.5.1}
Let $k$ be a field, $\mathrm{A}$ a $k$-algebra, local Noetherian and complete, whose residue field is a finite extension of $k$, $\mathfrak{p}$ a prime ideal of $\mathrm{A}$, $\mathfrak{q}$ a minimal prime ideal of $\mathrm{A}$ contained in $\mathfrak{p}$.
Then the minimum number of generators of the $\mathrm{A}_\mathfrak{p}$-module $(\hat{\Omega}^1_{\mathrm{A}/k})_\mathfrak{p}$ is at least equal to $\dim(\mathrm{A}/\mathfrak{q})$.
\end{lemma}

\oldpage[IV]{175}

\begin{proof}
Set $n=\dim(\mathrm{A}/\mathfrak{q})$, and let $m$ be the minimum number of generators of the $\mathrm{A}_\mathfrak{p}$-module $(\hat{\Omega}^1_{\mathrm{A}/k})_\mathfrak{p}$, which is equal to $\operatorname{rg}_{k(\mathfrak{p})}((\hat{\Omega}^1_{\mathrm{A}/k})_\mathfrak{p}\otimes_{\mathrm{A}_\mathfrak{p}}k(\mathfrak{p}))$ (Bourbaki, \emph{Alg.\ comm.}, chap.~II, \textsection~3, n\textsuperscript{o}~2, cor.~2 of prop.~4).
Note that $(\hat{\Omega}^1_{\mathrm{A}/k})_\mathfrak{q}=(\hat{\Omega}^1_{\mathrm{A}/k})_\mathfrak{p}\otimes_{\mathrm{A}_\mathfrak{p}}\mathrm{A}_\mathfrak{q}$, hence the minimum number of generators of the $\mathrm{A}_\mathfrak{q}$-module $(\hat{\Omega}^1_{\mathrm{A}/k})_\mathfrak{q}$ is \emph{at most} equal to $m$.
It suffices therefore to consider the case where $\mathfrak{p}=\mathfrak{q}$ is minimal.
In second place, let us show one can restrict to the case where $k$ is \emph{algebraically closed}.
Indeed, let $k'$ be an algebraic closure of $k$, and set $\mathrm{A}'=\mathrm{A}\widehat{\otimes}_kk'$, which is composed as a direct product of $k'$-local algebras $\mathrm{A}_i'$ ($1\leq i\leq r$), whose residue fields are isomorphic to $k'$, and which are $\mathrm{A}$-modules faithfully flat \sref{IV.18.11.4}.
Applying \sref{IV.2.3.4} and \sref{IV.6.1.1}, one sees that in a given $\mathrm{A}_i'$, there exists a minimal prime ideal $\mathfrak{q}_i'$ above $\mathfrak{q}$ such that $n'=\dim(\mathrm{A}_i'/\mathfrak{q}_i')\geq\dim(\mathrm{A}/\mathfrak{q})=n$.
On the other hand by \sref{IV.18.11.4}, the minimum $m'$ of generators of the $(\mathrm{A}_i')_{\mathfrak{q}_i'}$-module $(\hat{\Omega}^1_{\mathrm{A}_i'/k'})_{\mathfrak{q}_i'}$ is \emph{at most} equal to $m$; whence our assertion.

Suppose therefore that $k$ is algebraically closed, and let us show that we may further reduce to the case $\mathfrak q=0$. Indeed, $\Omega^1_{(\mathrm{A}/\mathfrak q)/\mathrm{A}}=0$ \sref[0]{0.20.4.12}, so the exact sequence \sref{IV.18.11.3.6} gives a surjection
\[
  \hat{\Omega}^1_{\mathrm{A}/k}\otimes_{\mathrm{A}}(\mathrm{A}/\mathfrak q)
  \longrightarrow \hat{\Omega}^1_{(\mathrm{A}/\mathfrak q)/k}.
\]
It follows at once that the minimum number of generators of the $\mathrm{K}$-vector space $\hat{\Omega}^1_{(\mathrm{A}/\mathfrak q)/k}\otimes_{\mathrm{A}/\mathfrak q}\mathrm{K}$, where $\mathrm{K}$ is the fraction field of $\mathrm{A}/\mathfrak q$, is at most $m$.

We may thus assume that $\mathrm{A}$ is integral. Its fraction field $\mathrm{K}$ is separable over the algebraically closed field $k$, and hence geometrically regular over $k$ \sref{IV.6.7.6}. We may therefore apply \sref{IV.18.11.3} with $\mathfrak p=0$; since $k=k^p$, this gives $m=n$.
\end{proof}

\begin{lemma}[18.11.5.2]
\label{IV.18.11.5.2}
Let $k$ be a field, $\mathrm{A}$ a $k$-algebra, local Noetherian and complete, whose residue field is a finite extension of $k$, $\mathfrak{m}$ its maximal ideal, $\mathfrak{p}\subset\mathfrak{m}$ a prime ideal of $\mathrm{A}$.
Let $m$ be the minimum number of generators of the $\mathrm{A}_\mathfrak{p}$-module $(\hat{\Omega}^1_{\mathrm{A}/k})_\mathfrak{p}$, and let $(x_i)_{1\leq i\leq m}$ be a family of elements of $\mathfrak{m}$ not belonging to $\mathfrak{p}$.
Then there exist $m$ invertible elements $u_i$ ($1\leq i\leq m$) of $\mathrm{A}$ such that if $y_i=u_ix_i$, the canonical images of the $dy_i$ in $(\hat{\Omega}^1_{\mathrm{A}/k})_\mathfrak{p}$ generate this $\mathrm{A}_\mathfrak{p}$-module.
\end{lemma}

\begin{proof}
For all $x\in\mathrm{A}$, denote by $\delta(x)$ the canonical image of $dx$ $(=d_{\mathrm{A}/k}x)$ in $(\hat{\Omega}^1_{\mathrm{A}/k})_\mathfrak{p}\otimes_{\mathrm{A}_\mathfrak{p}}k(\mathfrak{p})=\hat{\Omega}^1_{\mathrm{A}/k}\otimes_{\mathrm{A}}k(\mathfrak{p})$; as this is a $k(\mathfrak{p})$-vector space of dimension $m$ by hypothesis, it suffices (by virtue of the Nakayama lemma) to prove that one can determine the $u_i$ such that the $\delta(u_ix_i)$ form a free system.
We reason by induction and suppose for an integer $r<m$, we have determined the $u_i$ ($1\leq i\leq r$) such that the $\delta(u_ix_i)$ for $1\leq i\leq r$ are linearly independent over $k(\mathfrak{p})$; if $\delta(x_{r+1})$ is not a linear combination of the $\delta(u_ix_i)$ for $1\leq i\leq r$, one takes $u_{r+1}=1$ to continue the induction.
In the contrary case, for $u\in\mathrm{A}$ the element $\delta(ux_{r+1})$ is the canonical image of $(du)x_{r+1}+u(dx_{r+1})$. The second term has image in the span of the $\delta(u_ix_i)$ for $1\leq i\leq r$, while the image of $x_{r+1}$ in $k(\mathfrak p)$ is nonzero. It is therefore enough to find a unit $u\in\mathrm{A}$ such that $\delta(u)$ does not belong to that span.
\oldpage[IV]{176}
Now \sref[0]{0.20.7.15} and \sref[0]{0.7.2.9} show that the $\delta(x)$, for $x\in\mathrm{A}$, generate the $k(\mathfrak p)$-vector space $\hat{\Omega}^1_{\mathrm{A}/k}\otimes_{\mathrm{A}}k(\mathfrak p)$. Since $r<m$, there is therefore a $z\in\mathrm{A}$ for which $\delta(z)$ is not in the span of the $\delta(u_ix_i)$. If $z\notin\mathfrak m$, take $u=z$; otherwise $1+z$ is invertible and $\delta(1+z)=\delta(z)$ because $d(1+z)=dz$, so take $u=1+z$. This completes the proof of \sref{IV.18.11.5.2}.
\end{proof}

\begin{proof}[Proof of {\sref{IV.18.11.5}}, continued]
We now return to the proof of the implication $a)\Rightarrow b)$ in \sref{IV.18.11.5}.
Let us first note that since $\mathfrak{p}\neq\mathfrak{m}$, there exists a system of parameters $(x_i)_{1\leq i\leq n}$ of $\mathrm{A}$ (with $n=\dim(\mathrm{A})$) such that $x_i\notin\mathfrak{p}$ for all $i$, without which $\mathrm{A}/\mathfrak{p}$ would be of finite length, and $\mathfrak{p}$ would be maximal, contrary to the hypothesis.
But if for example $x_1\notin\mathfrak{p}$, it suffices, for each index $i$ such that $x_i\in\mathfrak{p}$, to replace $x_i$ by $x_1+x_i$, to have a system of parameters none of whose elements belong to $\mathfrak{p}$.
The relation $\dim(\mathrm{A}/\mathfrak{q})=n$ implies, by virtue of \sref{IV.18.11.5.1}, that $m\geq n$; one can therefore consider a family $(x_i)_{1\leq i\leq m}$ of elements $x_i\notin\mathfrak{p}$, whose $n$ first form a system of parameters of $\mathrm{A}$.
Multiplying furthermore the $x_i$ by invertible elements $u_i$ of $\mathrm{A}$, one can, by \sref{IV.18.11.5.2}, suppose that the images of $dx_i$ ($1\leq i\leq m$) in $(\hat{\Omega}^1_{\mathrm{A}/k})_\mathfrak{p}$ generate this $\mathrm{A}_\mathfrak{p}$-module, and the multiplication by the $u_i$ has not altered the fact that the $x_i$ for $1\leq i\leq n$ form a system of parameters.
We shall consider the local $k$-homomorphism $u:\mathrm{B}\to\mathrm{A}$ such that $u(\mathrm{T}_i)=x_i$ for $1\leq i\leq m$ (Bourbaki, \emph{Alg.\ comm.}, chap.~III, \textsection~4, n\textsuperscript{o}~5, prop.~6); as the $x_i$ generate an ideal of definition of $\mathrm{A}$, it follows from \sref[0]{0.7.4.4} and \sref{IV.7.4.3} that $u$ makes $\mathrm{A}$ a $\mathrm{B}$-algebra \emph{finite}.

But the $d_\mathrm{A}x_i$ are the canonical images by $v$ of the elements $d_\mathrm{B}\mathrm{T}_i\otimes 1$ \sref[0]{0.20.5.2.6}.
Localising the preceding exact sequence at $\mathfrak p$, we see that
\[
  v_\mathfrak p:(\hat{\Omega}^1_{\mathrm{B}/k}\otimes_\mathrm{B}\mathrm{A})_\mathfrak p
  \longrightarrow(\hat{\Omega}^1_{\mathrm{A}/k})_\mathfrak p
\]
is surjective. Consequently
\[
  0=(\Omega^1_{\mathrm{A}/\mathrm{B}})_\mathfrak p
   =\Omega^1_{\mathrm{A}_\mathfrak p/\mathrm{B}_\mathfrak r},
\]
where $\mathfrak r=u^{-1}(\mathfrak p)$ is the inverse image of $\mathfrak p$ in $\mathrm{B}$; cf.\ \sref{IV.16.4.15}. By \sref{IV.17.4.1}, the morphism $\operatorname{Spec}(\mathrm{A})\to\operatorname{Spec}(\mathrm{B})$ is therefore unramified at $\mathfrak p$, which proves $b)$.
\end{proof}

\begin{remark}[18.11.6]
\label{IV.18.11.6}
In the statement of \sref{IV.18.11.5}, the hypothesis $\mathfrak{p}\neq\mathfrak{m}$ cannot be omitted.
Indeed, for every local $k$-homomorphism $u:\mathrm{B}\to\mathrm{A}$, where $\mathrm{B}=k[[\mathrm{T}_1,\ldots,\mathrm{T}_r]]$ ($r$ any integer) making $\mathrm{A}$ a $\mathrm{B}$-algebra finite, $\mathfrak{m}$ is the only point of $\operatorname{Spec}(\mathrm{A})$ above the maximal ideal $\mathfrak{n}$ of $\mathrm{B}$; the morphism $\operatorname{Spec}(\mathrm{A})\to\operatorname{Spec}(\mathrm{B})$ cannot be unramified at $\mathfrak{m}$ unless $\mathrm{A}/\mathfrak{n}\mathrm{A}$ is a field, an extension separable of $k$ \sref{IV.17.4.1}, which implies that the residue field $\mathrm{K}$ of $\mathrm{A}$ is a separable extension of $k$.
If this condition is not satisfied, the conclusion of \sref{IV.18.11.5} can never be satisfied for $\mathfrak{p}=\mathfrak{m}$, whatever the integer $m$.
\end{remark}

\begin{corollary}[18.11.7]
\label{IV.18.11.7}
Let $k$ be a field, $\mathrm{A}$ a $k$-algebra, local Noetherian and complete, equidimensional, whose residue field is a finite extension of $k$, $\mathfrak{p}$ a prime ideal of $\mathrm{A}$.
Then the minimum number of generators of the $\mathrm{A}_\mathfrak{p}$-module $(\hat{\Omega}^1_{\mathrm{A}/k})_\mathfrak{p}$ is at least equal to $\dim(\mathrm{A})$.
\end{corollary}

\begin{proof}
This is a special case of \sref{IV.18.11.5.1}.
\end{proof}

\oldpage[IV]{177}

\begin{corollary}[18.11.8]
\label{IV.18.11.8}
Let $k$ be a field, $\mathrm{A}$ a $k$-algebra, local Noetherian and complete, integral, whose residue field is a finite extension of $k$.
If $\mathrm{K}$ is the field of fractions of $\mathrm{A}$, one has $\operatorname{rg}_{\mathrm{K}}((\hat{\Omega}^1_{\mathrm{A}/k})\otimes_{\mathrm{A}}\mathrm{K})\geq\dim(\mathrm{A})$.
\end{corollary}

\begin{proof}
It suffices to set $\mathfrak{p}=0$ in \sref{IV.18.11.7}.
\end{proof}

\begin{corollary}[18.11.9]
\label{IV.18.11.9}
Let $k$ be a field, $\mathrm{A}$ a $k$-algebra, local Noetherian and complete, of dimension $n$, whose residue field is a finite extension of $k$.
Let $\mathfrak{p}$ be a prime ideal of $\mathrm{A}$ distinct from the maximal ideal, and containing a minimal prime ideal $\mathfrak{q}$ such that $\dim(\mathrm{A}/\mathfrak{q})=n$.
Suppose that the $\mathrm{A}_\mathfrak{p}$-module $(\hat{\Omega}^1_{\mathrm{A}/k})_\mathfrak{p}$ admits $n$ generators.
Then:
\begin{enumerate}[label=(\roman*)]
  \item The $\mathrm{A}_\mathfrak{p}$-module $(\hat{\Omega}^1_{\mathrm{A}/k})_\mathfrak{p}$ is free of rank $n$.
  \item There exists a local $k$-homomorphism $u:\mathrm{B}\to\mathrm{A}$, where $\mathrm{B}=k[[\mathrm{T}_1,\ldots,\mathrm{T}_n]]$, making $\mathrm{A}$ a $\mathrm{B}$-algebra finite, and such that the corresponding morphism $\operatorname{Spec}(\mathrm{A})\to\operatorname{Spec}(\mathrm{B})$ is \'etale at the point $\mathfrak{p}$.
  \item The $k$-algebra $\mathrm{A}_\mathfrak{p}$ is geometrically regular over $k$.
\end{enumerate}
\end{corollary}

\begin{proof}
We first prove (ii). By \sref{IV.18.11.5}, there is a local $k$-homomorphism
\[
  u:\mathrm{B}=k[[\mathrm{T}_1,\ldots,\mathrm{T}_n]]\longrightarrow\mathrm{A}
\]
making $\mathrm{A}$ a finite $\mathrm{B}$-algebra and such that $\operatorname{Spec}(\mathrm{A})\to\operatorname{Spec}(\mathrm{B})$ is unramified at $\mathfrak p$; put $\mathfrak r=u^{-1}(\mathfrak p)$. The hypothesis on $\mathfrak q\subset\mathfrak p$, together with the fact that $\mathrm{A}$ is a quotient of a regular ring \sref[0]{0.19.8.8}, gives
\[
  \dim(\mathrm{A}_\mathfrak p)=n-\dim(\mathrm{A}/\mathfrak p)
\]
by \sref[0]{0.16.5.12}; likewise $\dim(\mathrm{B}_\mathfrak r)=n-\dim(\mathrm{B}/\mathfrak r)$. Since the morphism is unramified at $\mathfrak p$, its fibre over $\mathfrak r$ has dimension zero, so \sref[0]{0.16.3.9} gives $\dim(\mathrm{A}/\mathfrak p)\leq\dim(\mathrm{B}/\mathfrak r)$ and hence $\dim(\mathrm{B}_\mathfrak r)\leq\dim(\mathrm{A}_\mathfrak p)$. It follows from \sref{IV.18.10.1} that the morphism is \'etale at $\mathfrak p$.

Assertion (iii) follows because $\mathrm{A}_\mathfrak p$ is formally smooth over $\mathrm{B}_\mathfrak r$, while $\mathrm{B}_\mathfrak r$ is formally smooth over $k$, for their preadic topologies \sref[0]{0.19.3.4} and \sref[0]{0.19.3.5}. Thus $\mathrm{A}_\mathfrak p$ is formally smooth over $k$ for its preadic topology, and consequently geometrically regular over $k$ \sref[0]{0.22.5.8}.

Finally, we prove (i). Let $k'$ be an algebraically closed extension of $k$, and consider the semilocal $k'$-algebra $\mathrm{A}'=\mathrm{A}\widehat{\otimes}_kk'$. Regarding $\mathrm{A}$ as a finite $\mathrm{B}$-module via $u$, there is a canonical isomorphism
\[
  \mathrm{A}'\simeq\mathrm{A}\otimes_{\mathrm{B}}(\mathrm{B}\widehat{\otimes}_kk')
\]
by \sref{IV.7.5.7.1}; that result does not require the residue field of $\mathrm{B}$ to be finite over $k$. Put $\mathrm{B}'=\mathrm{B}\widehat{\otimes}_kk'$, which identifies canonically with $k'[[\mathrm{T}_1,\ldots,\mathrm{T}_n]]$.

Since $\operatorname{Spec}(\mathrm{A})\to\operatorname{Spec}(\mathrm{B})$ is finite and \'etale at $\mathfrak p$, the morphism $\operatorname{Spec}(\mathrm{A}')\to\operatorname{Spec}(\mathrm{B}')$ is finite and \'etale at every point $\mathfrak p'$ over $\mathfrak p$ \sref{IV.17.3.3}. Moreover, $\mathrm{A}'$ is the direct product of local rings $\mathrm{A}_i'$ of dimension $n$ \sref{IV.18.11.4}, and $\mathfrak p'$ identifies with a prime ideal $\mathfrak p_i'$ of one of them. The preceding argument shows that $(\mathrm{A}_i')_{\mathfrak p_i'}$ is geometrically regular over $k'$. Since $k'$ is perfect, \sref{IV.18.11.3} shows that $(\hat{\Omega}^1_{\mathrm{A}_i'/k'})_{\mathfrak p_i'}$ is a free $(\mathrm{A}_i')_{\mathfrak p_i'}$-module of finite type.

But \sref{IV.18.11.4} gives
\[
  (\hat{\Omega}^1_{\mathrm{A}_i'/k'})_{\mathfrak p_i'}
  \simeq (\hat{\Omega}^1_{\mathrm{A}/k})_\mathfrak p
  \otimes_{\mathrm{A}_\mathfrak p}(\mathrm{A}_i')_{\mathfrak p_i'},
\]
and $(\mathrm{A}_i')_{\mathfrak p_i'}$ is faithfully flat over $\mathrm{A}_\mathfrak p$. Therefore $(\hat{\Omega}^1_{\mathrm{A}/k})_\mathfrak p$ is a free $\mathrm{A}_\mathfrak p$-module of finite type \sref{IV.2.5.2};

\oldpage[IV]{178}

and \sref{IV.18.11.5.1}, together with $\dim(\mathrm{A}/\mathfrak q)=n$, shows that its rank is $n$.
\end{proof}

\begin{theorem}[18.11.10]
\label{IV.18.11.10}
Let $k$ be a field of characteristic exponent $p$, $\mathrm{A}$ a $k$-algebra, local Noetherian and complete, whose residue field is a finite extension of $k$, $\mathfrak{p}$ a prime ideal of $\mathrm{A}$ distinct from the maximal ideal.
The following conditions are equivalent:
\begin{enumerate}[label=\alph*)]
  \item For every extension $k'$ of $k$, and every prime ideal $\mathfrak{p}'$ of $\mathrm{A}'=\mathrm{A}\widehat{\otimes}_kk'$ above $\mathfrak{p}$, $\mathrm{A}'_{\mathfrak{p}'}$ is a regular ring.
\item[a$'$)] There exists a perfect extension $k'$ of $k$ and a prime ideal $\mathfrak{p}'$ of $\mathrm{A}'=\mathrm{A}\widehat{\otimes}_kk'$ above $\mathfrak{p}$ such that $\mathrm{A}'_{\mathfrak{p}'}$ is regular.
  \item Let $n$ be the largest of the dimensions of the irreducible components of $\operatorname{Spec}(\mathrm{A})$ to which $\mathfrak{p}$ belongs; then $(\hat{\Omega}^1_{\mathrm{A}/k})_\mathfrak{p}$ is an $\mathrm{A}_\mathfrak{p}$-module free of rank $n$.
\end{enumerate}
Suppose furthermore that $n=\dim(\mathrm{A})$ (which will be the case if $\mathrm{A}$ is equidimensional); then the preceding conditions are also equivalent to:
\begin{enumerate}[label=\alph*)]
  \setcounter{enumi}{2}
  \item There exists a local $k$-homomorphism $u:\mathrm{B}\to\mathrm{A}$, where $\mathrm{B}=k[[\mathrm{T}_1,\ldots,\mathrm{T}_n]]$, making $\mathrm{A}$ a $\mathrm{B}$-algebra finite, and such that the corresponding morphism $\operatorname{Spec}(\mathrm{A})\to\operatorname{Spec}(\mathrm{B})$ is \'etale at the point $\mathfrak{p}$.
\end{enumerate}
Each of the conditions $a)$, $a')$, and $b)$ implies the following:
\begin{enumerate}[label=\alph*)]
  \setcounter{enumi}{3}
  \item The ring $\mathrm{A}_\mathfrak{p}$ is geometrically regular over $k$.
\end{enumerate}
Furthermore, if $[k:k^p]<+\infty$, condition $d)$ is equivalent to $a)$, $a')$, and $b)$.
\end{theorem}

\begin{proof}
That $d)$ implies $b)$ when $[k:k^p]<+\infty$ is precisely \sref{IV.18.11.3}. It is clear that $a)$ implies $a')$. Let us show that $c)$ implies $a)$ when $\dim(\mathrm{A})=n$. With the notation of $a)$, one then has
\[
  \mathrm{A}'=\mathrm{A}\otimes_{\mathrm{B}}\mathrm{B}',
  \qquad
  \mathrm{B}'=\mathrm{B}\widehat{\otimes}_kk'=k'[[\mathrm{T}_1,\ldots,\mathrm{T}_n]]
\]
by \sref{IV.7.5.7.1}. The morphism $\operatorname{Spec}(\mathrm{A}')\to\operatorname{Spec}(\mathrm{B}')$ is finite and \'etale at $\mathfrak p'$, and the argument proving \sref{IV.18.11.9},~(iii), shows that $\mathrm{A}'_{\mathfrak p'}$ is regular.

That $b)$ implies $c)$ when $\dim(\mathrm{A})=n$ follows from \sref{IV.18.11.9}. Let us prove that $a')$ implies $b)$ under the same hypothesis. The ring $\mathrm{A}'$ is flat over $\mathrm{A}$ \sref{IV.18.11.4}, and $\dim(\mathrm{A}')=\dim(\mathrm{A})=n$. By \sref{IV.2.3.4} and \sref{IV.6.1.1}, there is a minimal prime $\mathfrak q'$ of $\mathrm{A}'$, contained in $\mathfrak p'$ and lying over a minimal prime $\mathfrak q$ of $\mathrm{A}$, such that
\[
  \dim(\mathrm{A}'/\mathfrak q')=\dim(\mathrm{A}/\mathfrak q)=n.
\]
Moreover, \sref{IV.18.11.4} gives
\[
  (\hat{\Omega}^1_{\mathrm{A}'/k'})_{\mathfrak p'}
  \simeq(\hat{\Omega}^1_{\mathrm{A}/k})_\mathfrak p
  \otimes_{\mathrm{A}_\mathfrak p}\mathrm{A}'_{\mathfrak p'}.
\]
Since $k'$ is perfect, the regularity of $\mathrm{A}'_{\mathfrak p'}$ implies that it is geometrically regular over $k'$ \sref{IV.6.7.7}. As $k'^p=k'$, we may apply the implication $d)\Rightarrow b)$ to $\mathrm{A}'$ and $\mathfrak p'$. Thus $(\hat{\Omega}^1_{\mathrm{A}'/k'})_{\mathfrak p'}$ is a free $\mathrm{A}'_{\mathfrak p'}$-module of rank $n$; faithful flatness \sref{IV.18.11.4} and \sref{IV.2.5.2} then shows that $(\hat{\Omega}^1_{\mathrm{A}/k})_\mathfrak p$ is a free $\mathrm{A}_\mathfrak p$-module of rank $n$.

This proves the equivalence of $a)$, $a')$, $b)$, and $c)$ when $\dim(\mathrm{A})=n$.
It remains to prove that $a)$, $a')$, and $b)$ are still equivalent in the general case.
Let $\mathfrak J$ be the kernel of the canonical homomorphism $\mathrm{A}\to\mathrm{A}_\mathfrak p$, and put $\mathrm{A}_1=\mathrm{A}/\mathfrak J$. Necessarily $\mathfrak J\subset\mathfrak p$; and if $\mathfrak p_1=\mathfrak p/\mathfrak J$, the canonical homomorphism $\mathrm{A}_\mathfrak p\to(\mathrm{A}_1)_{\mathfrak p_1}$ is bijective. It follows from \sref[I]{I.6.5.4} that the canonical injection $\operatorname{Spec}(\mathrm{A}_1)\to\operatorname{Spec}(\mathrm{A})$ is a local isomorphism at $\mathfrak p_1$, and $\mathfrak J_\mathfrak p=0$. As in \sref{IV.18.11.3}, there is an exact sequence
\[
  \mathfrak J/\mathfrak J^2
  \longrightarrow \hat{\Omega}^1_{\mathrm{A}/k}\otimes_{\mathrm{A}}\mathrm{A}_1
  \longrightarrow \hat{\Omega}^1_{\mathrm{A}_1/k}\longrightarrow0.
\]
\oldpage[IV]{179}
After localising at $\mathfrak p$, this yields an isomorphism
\[
  (\hat{\Omega}^1_{\mathrm{A}/k})_\mathfrak p
  \xrightarrow{\sim}
  (\hat{\Omega}^1_{\mathrm{A}_1/k})_{\mathfrak p_1}.
\]
Thus condition $b)$ for $\mathrm{A}$ and $\mathfrak p$ is equivalent to condition $b)$ for $\mathrm{A}_1$ and $\mathfrak p_1$.

On the other hand, with the notation of $a)$, there is a canonical isomorphism
\[
  \mathrm{A}_1'=\mathrm{A}_1\widehat{\otimes}_kk'
  \simeq \mathrm{A}'/\mathfrak J\mathrm{A}'
\]
by \sref{IV.7.5.7.1}, where the hypothesis on the residue field of $\mathrm{B}$ is superfluous. If $\mathfrak p'$ is a prime of $\mathrm{A}'$ above $\mathfrak p$, every element of $\mathfrak J\mathrm{A}'$ annihilates an element of $\mathrm{A}'-\mathfrak p'$, and hence $\mathfrak J\mathrm{A}'\subset\mathfrak p'$. Putting $\mathfrak p_1'=\mathfrak p'/\mathfrak J\mathrm{A}'$, the prime $\mathfrak p_1'$ lies above $\mathfrak p_1$ and $(\mathrm{A}_1')_{\mathfrak p_1'}$ identifies canonically with $\mathrm{A}'_{\mathfrak p'}$. This proves that condition $a)$ (respectively $a')$) for $\mathrm{A}$ and $\mathfrak p$ is equivalent to the same condition for $\mathrm{A}_1$ and $\mathfrak p_1$.

All minimal prime ideals of $\mathrm{A}$ contained in $\mathfrak p$ contain $\mathfrak J$, since $\operatorname{Spec}(\mathrm{A}_1)\to\operatorname{Spec}(\mathrm{A})$ is a local isomorphism at $\mathfrak p_1$. Moreover, the primes in $\operatorname{Ass}_{\mathrm{A}}(\mathrm{A}/\mathfrak J)$ are precisely those in $\operatorname{Ass}(\mathrm{A})$ which are contained in $\mathfrak p$ (Bourbaki, \emph{Alg.\ comm.}, chap.~IV, \textsection~1, n\textsuperscript{o}~2, prop.~6). Hence all minimal primes of $\mathrm{A}_1$ are contained in $\mathfrak p_1$, and consequently $\dim(\mathrm{A}_1)=n$. It remains only to apply to $\mathrm{A}_1$ and $\mathfrak p_1$ what was proved above.
\end{proof}

\begin{remark}[18.11.11]
\label{IV.18.11.11}
\begin{enumerate}[label=(\roman*)]
  \item The equivalence of conditions $d)$ and $b)$ in \sref{IV.18.11.10} is no longer valid when one no longer supposes that $[k:k^p]<+\infty$.
Indeed, by the example of \sref[0]{0.22.7.7},~(ii)), the ring $\mathrm{B}=\mathrm{A}/\mathfrak{q}$ is integral and of dimension $1$; on the other hand, the sequence
\[
  (\mathfrak{q}/\mathfrak{q}^2)\otimes_{\mathrm{A}}\mathrm{L}\xrightarrow{j}\hat{\Omega}^1_{\mathrm{A}/k}\otimes_{\mathrm{A}}\mathrm{L}\to\hat{\Omega}^1_{\mathrm{B}/k}\otimes_{\mathrm{B}}\mathrm{L}\to 0
\]
is exact (same demonstration as in \sref{IV.18.11.3}), and we know that $j$ \emph{is not injective}, hence $j=0$ since $(\mathfrak{q}/\mathfrak{q}^2)\otimes_\mathrm{A}\mathrm{L}$ is of rank $1$.
We conclude that $\operatorname{rg}_{\mathrm{L}}(\hat{\Omega}^1_{\mathrm{B}/k}\otimes_{\mathrm{B}}\mathrm{L})=2$; as $\mathrm{B}_\mathfrak{p}$ is a $k$-algebra geometrically regular, one sees that the condition $d)$ does not imply here the condition $b)$.
  \item With the notations of \sref{IV.18.11.10} and supposing that $n=\dim(\mathrm{A})$ and $(x_i)_{1\leq i\leq n}$ a system of parameters of $\mathrm{A}$ not belonging to $\mathfrak{p}$; one deduces a local $k$-homomorphism $u:\mathrm{B}\to\mathrm{A}$ such that $u(\mathrm{T}_i)=x_i$ for $1\leq i\leq n$, making $\mathrm{A}$ a $\mathrm{B}$-algebra finite.
For the morphism $\operatorname{Spec}(\mathrm{A})\to\operatorname{Spec}(\mathrm{B})$ to be \emph{\'etale} at the point $\mathfrak{p}$, it is necessary and sufficient that the images of $d_\mathrm{A}x_i$ in $(\hat{\Omega}^1_{\mathrm{A}/k})_\mathfrak{p}$ form a \emph{system of generators} of this $\mathrm{A}_\mathfrak{p}$-module.
Indeed, as one has seen in the proof of \sref{IV.18.11.5} that if the morphism $\operatorname{Spec}(\mathrm{A})\to\operatorname{Spec}(\mathrm{B})$ is not ramified at $\mathfrak{p}$, the canonical homomorphism $(\hat{\Omega}^1_{\mathrm{B}/k}\otimes_\mathrm{B}\mathrm{A})_\mathfrak{p}\to(\hat{\Omega}^1_{\mathrm{A}/k})_\mathfrak{p}$ is surjective and the images of the elements $d_\mathrm{B}\mathrm{T}_i\otimes 1$, which generate $(\hat{\Omega}^1_{\mathrm{B}/k}\otimes_\mathrm{B}\mathrm{A})_\mathfrak{p}$, are the $d_\mathrm{A}x_i$; whence the \emph{necessity} of the condition.
Conversely, the same reasoning shows that if this condition is satisfied, the morphism $\operatorname{Spec}(\mathrm{A})\to\operatorname{Spec}(\mathrm{B})$ is not ramified at $\mathfrak{p}$, and the reasoning of \sref{IV.18.11.9} proves in fact that this morphism is \'etale at $\mathfrak{p}$.
\end{enumerate}
\end{remark}

\begin{corollary}[18.11.12]
\label{IV.18.11.12}
Let $k$ be a field of characteristic exponent $p$ such that $[k:k^p]<+\infty$, $\mathrm{A}$ a $k$-algebra, local Noetherian and complete and integral, which is not a field, and whose residue field is a finite extension of $k$.
Then there exists a finite radicial extension $k'$ of $k$ such that, setting $\mathrm{A}'=\mathrm{A}\otimes_kk'$, the $k'$-algebra $\mathrm{A}'_{\mathrm{red}}$ and the prime ideal $\mathfrak{o}$ of this algebra satisfy the equivalent conditions $a)$, $a')$, $b)$, $c)$, and $d)$ of \sref{IV.18.11.10}.
In particular, if $n=\dim(\mathrm{A})=\dim(\mathrm{A}')$,
\oldpage[IV]{180}
there exists a $k'$-homomorphism local $\mathrm{B}'=k'[[\mathrm{T}_1,\ldots,\mathrm{T}_n]]\to\mathrm{A}'_\mathrm{red}$ making $\mathrm{A}'_\mathrm{red}$ a $\mathrm{B}'$-algebra finite, and such that the field of fractions $\mathrm{K}'$ of $\mathrm{A}'_\mathrm{red}$ is a finite separable extension of the field of fractions $k'((\mathrm{T}_1,\ldots,\mathrm{T}_n))$ of $\mathrm{B}'$.
\end{corollary}

\begin{proof}
Since $\operatorname{Spec}(\mathrm{A}')\to\operatorname{Spec}(\mathrm{A})$ is radicial and finite, $\mathrm{A}'$ is a complete local ring and its nilradical is the unique prime above the zero ideal of $\mathrm{A}$. By flatness, $\mathrm{A}'$ identifies with a subring of $\mathrm{K}\otimes_kk'$, which is also the total ring of fractions of $\mathrm{A}'$; hence
\[
  \mathrm{K}'=(\mathrm{K}\otimes_kk')_{\mathrm{red}}.
\]
In view of the equivalence $d)\Leftrightarrow c)$ in \sref{IV.18.11.10}, it is therefore enough to find a finite radicial extension $k'$ of $k$ such that $(\mathrm{K}\otimes_kk')_{\mathrm{red}}$ is separable over $k'$ \sref{IV.6.7.6}.

By \sref[0]{0.19.8.9}, under the present hypotheses there is a sub-$k$-algebra $\mathrm{C}=k[[\mathrm{T}_1,\ldots,\mathrm{T}_m]]$ of $\mathrm{A}$ such that $\mathrm{A}$ is finite over $\mathrm{C}$. Thus $\mathrm{K}$ is a finite extension of the fraction field $\mathrm{K}_1=k((\mathrm{T}_1,\ldots,\mathrm{T}_m))$ of $\mathrm{C}$, and $\mathrm{K}_1$ is separable over $k$ \sref[0]{0.21.9.6.4}. We may write
\[
  \mathrm{K}\otimes_kk^{p^{-\infty}}
  =\mathrm{K}\otimes_{\mathrm{K}_1}(\mathrm{K}_1\otimes_kk^{p^{-\infty}}),
\]
and $\mathrm{K}_1\otimes_kk^{p^{-\infty}}$ is a field, radicial over $\mathrm{K}_1$. Consequently $\mathrm{K}\otimes_kk^{p^{-\infty}}$ is finite over a field and is therefore Artinian. The conclusion follows from the more general lemma below.
\end{proof}

\begin{lemma}[18.11.12.1]
\label{IV.18.11.12.1}
Let $k$ be a field of characteristic $p>0$, $\mathrm{K}$ an extension of $k$.
The following conditions are equivalent:
\begin{enumerate}[label=\alph*)]
  \item The ring $\mathrm{K}\otimes_kk^{p^{-\infty}}$ is Artinian.
  \item There exists an extension $k'$ of $k$ of finite type such that $(\mathrm{K}\otimes_kk')_{\mathrm{red}}$ is a $k'$-algebra separable.
  \item There exists a finite radicial extension $k'$ of $k$ such that $(\mathrm{K}\otimes_kk')_\mathrm{red}$ is a separable extension of $k'$.
\end{enumerate}
\end{lemma}

\begin{proof}
Let us first show that $c)$ implies $a)$.
The ring $\mathrm{A}=\mathrm{K}\otimes_kk'$ is then local Artinian. If $\mathfrak N$ is its nilradical, its residue field $\mathrm{L}=\mathrm{A}/\mathfrak N$ is separable over $k'$, and hence $\mathrm{L}\otimes_{k'}k^{p^{-\infty}}$ is a field. Since this field is
\[
  (\mathrm{A}\otimes_{k'}k^{p^{-\infty}})/(\mathfrak N\otimes_{k'}k^{p^{-\infty}}),
\]
the ideal $\mathfrak N\otimes_{k'}k^{p^{-\infty}}$ is the nilradical of $\mathrm{K}\otimes_kk^{p^{-\infty}}$. The ideal $\mathfrak N$ is finitely generated and nilpotent, so this nilradical is finitely generated; it follows that $\mathrm{K}\otimes_kk^{p^{-\infty}}$ is Artinian.

Conversely, let us prove that $a)$ implies $c)$.
Let $\mathfrak N$ be the nilradical of the local Artinian ring $\mathrm{B}=\mathrm{K}\otimes_kk^{p^{-\infty}}$. It is generated by finitely many elements
\[
  y_i=\sum_j x_{ij}\otimes\xi_{ij},
  \qquad x_{ij}\in\mathrm{K},\quad \xi_{ij}\in k^{p^{-\infty}}.
\]
Let $k'$ be the finite radicial extension of $k$ generated by the $\xi_{ij}$, and let $\mathfrak N_0$ be the ideal of $\mathrm{B}_0=\mathrm{K}\otimes_kk'$ generated by the $y_i$. The $y_i$ are nilpotent in $\mathrm{B}_0$; moreover
\[
  \mathfrak N_0\otimes_{k'}k^{p^{-\infty}}=\mathfrak N,
  \qquad
  \mathfrak N\cap\mathrm{B}_0=\mathfrak N_0.
\]
Thus $\mathfrak N_0$ is exactly the nilradical of $\mathrm{B}_0$. Since
\[
  \mathrm{B}/\mathfrak N=(\mathrm{B}_0/\mathfrak N_0)\otimes_{k'}k^{p^{-\infty}}
\]
is reduced, \sref{IV.4.6.1} shows that $\mathrm{B}_0/\mathfrak N_0=(\mathrm{K}\otimes_kk')_{\mathrm{red}}$ is separable over $k'$.

It is clear that $c)$ implies $b)$.
Conversely, suppose $b)$ holds. There is a separable extension $k_1$ of $k$ such that $k'$ is finite radicial over $k_1$; put $\mathrm{K}_1=\mathrm{K}\otimes_kk_1$, which is a field. Applying the equivalence of $a)$ and $c)$ to the extension $\mathrm{K}_1/k_1$, we see that $\mathrm{K}_1\otimes_{k_1}k_1^{p^{-\infty}}$ is Artinian. But
\oldpage[IV]{181}
\[
  \mathrm{K}\otimes_kk_1^{p^{-\infty}}
  =(\mathrm{K}\otimes_kk^{p^{-\infty}})
   \otimes_{k^{p^{-\infty}}}k_1^{p^{-\infty}},
\]
so $\mathrm{K}\otimes_kk^{p^{-\infty}}$ is Artinian as well (Bourbaki, \emph{Alg.\ comm.}, chap.~I, \textsection~3, n\textsuperscript{o}~5, cor.\ of prop.~8). Thus $b)$ implies $a)$, completing the proofs of \sref{IV.18.11.12.1} and \sref{IV.18.11.12}.
\end{proof}


\subsection{Applications of \'etale localisation to quasi-finite morphisms (generalisations of earlier results)}
\label{subsection:IV.18.12}

The results of this section were communicated by P.~Deligne.

\begin{theorem}[18.12.1]
\label{IV.18.12.1}
Let $f:\mathrm{X}\to\mathrm{Y}$ be a morphism locally of finite type, let $x$ be a point of $\mathrm{X}$, and put $y=f(x)$.
Suppose that $x$ is an isolated point of the space $f^{-1}(y)$.
Then there exist an \'etale morphism $\mathrm{Y}'\to\mathrm{Y}$, a point $x'$ of $\mathrm{X}'=\mathrm{X}\times_{\mathrm{Y}}\mathrm{Y}'$ above $x$, and an open neighbourhood $\mathrm{V}'$ of $x'$ in $\mathrm{X}'$ such that, if $f'=f_{(\mathrm{Y}')}:\mathrm{X}'\to\mathrm{Y}'$, the restriction $f'|\mathrm{V}'$ is finite.
If, moreover, $f$ is separated, then $\mathrm{V}'$ is both open and closed, and $\mathrm{X}'$ is therefore the disjoint union of two subschemes induced on open subsets of $\mathrm{X}'$, one of which is finite over $\mathrm{Y}'$ and contains $x'$.
\end{theorem}

\begin{proof}
The last assertion follows because, if $f$ is separated, so is $f'$. Thus, if $j':\mathrm{V}'\to\mathrm{X}'$ denotes the canonical injection, the finiteness of $f'|\mathrm{V}'=f'\circ j'$ implies that $j'$ is finite \sref[II]{II.6.1.5}; consequently $\mathrm{V}'=j'(\mathrm{V}')$ is closed in $\mathrm{X}'$ \sref[II]{II.6.1.10}.

The question is local on $\mathrm{X}$, so we may suppose that $\mathrm{Y}=\operatorname{Spec}(\mathrm{A})$ and $\mathrm{X}=\operatorname{Spec}(\mathrm{B})$ are affine, with $\mathrm{B}$ an $\mathrm{A}$-algebra of finite type. Write $\mathrm{B}=\mathrm{C}/\mathfrak{J}$, where $\mathrm{C}=\mathrm{A}[\mathrm{T}_1,\ldots,\mathrm{T}_r]$ and $\mathfrak{J}$ is an ideal of $\mathrm{C}$. Since $f$ is then separated, we may also suppose that $f^{-1}(y)=\{x\}$.

Let $(\mathfrak{J}_\lambda)$ be the filtered family of finitely generated ideals of $\mathrm{C}$ contained in $\mathfrak{J}$; then $\mathfrak{J}$ is their union. Regard $\mathrm{X}$ and the schemes $\mathrm{X}_\lambda=\operatorname{Spec}(\mathrm{C}/\mathfrak{J}_\lambda)$ as closed subschemes of $\mathrm{Z}=\operatorname{Spec}(\mathrm{C})$. On the underlying spaces, $\mathrm{X}=\bigcap_\lambda\mathrm{X}_\lambda$. If $f_\lambda:\mathrm{X}_\lambda\to\mathrm{Y}$ is the structural morphism, then
\[
  f^{-1}(y)=\bigcap_\lambda f_\lambda^{-1}(y).
\]
The sets $f_\lambda^{-1}(y)$ are closed in the Noetherian space $\mathrm{Z}_y=\operatorname{Spec}(k(y)[\mathrm{T}_1,\ldots,\mathrm{T}_r])$; hence, for some $\lambda$, $f_\lambda^{-1}(y)=f^{-1}(y)=\{x\}$.

We may replace $\mathrm{X}$ by such an $\mathrm{X}_\lambda$. Indeed, if the assertion is proved for $\mathrm{X}_\lambda$, let $x'$ lie above $x$ in $\mathrm{Z}'=\mathrm{Z}\times_{\mathrm{Y}}\mathrm{Y}'$ and put $\mathrm{X}_\lambda'=p^{-1}(\mathrm{X}_\lambda)$, where $p:\mathrm{Z}'\to\mathrm{Z}$ is the projection (so also $\mathrm{X}'=p^{-1}(\mathrm{X})$). If $\mathrm{V}_\lambda'$ is an open neighbourhood of $x'$ in $\mathrm{X}_\lambda'$ finite over $\mathrm{Y}'$, then the restriction to the closed subscheme $\mathrm{X}'\cap\mathrm{V}_\lambda'$ is finite as well. We have therefore reduced to the case where $f$ is of finite presentation.

The set of points of $\mathrm{X}$ which are isolated in their fibres is open in $\mathrm{X}$ \sref{IV.13.1.4}; after replacing $\mathrm{X}$ by an open neighbourhood of $x$, we may therefore suppose that $f$ is quasi-finite. Let $\mathrm{A}''$ be the Henselization of $\sh{O}_{Y,y}$ and put $\mathrm{Y}''=\operatorname{Spec}(\mathrm{A}'')$. If $\mathrm{X}''=\mathrm{X}\times_{\mathrm{Y}}\mathrm{Y}''$ and $f''=f_{(\mathrm{Y}'')}:\mathrm{X}''\to\mathrm{Y}''$, then $f''$ is separated, quasi-finite, and of finite presentation. Since $\mathrm{A}''$ is henselian, every point $x''\in\mathrm{X}''$ above $x$ has an open-and-closed neighbourhood $\mathrm{V}''$ in $\mathrm{X}''$ which is finite over $\mathrm{Y}''$ \sref{IV.18.5.11},~c). The immersion $\mathrm{V}''\to\mathrm{X}''$ is open and closed, hence quasi-compact and of finite presentation \sref{IV.1.6.2}; consequently $f''|\mathrm{V}''$ is of finite presentation.

By definition, $\mathrm{A}''$ is the filtered direct limit of strictly essentially \'etale $\sh{O}_{Y,y}$-algebras $\mathrm{A}_\lambda$ \sref{IV.18.6.5}. Each $\mathrm{A}_\lambda$ is itself the filtered direct limit of \'etale $\sh{O}_{Y,y}$-algebras $\mathrm{B}_{\lambda\mu}$ \sref{IV.18.6.1}, \sref{IV.8.1.2},~a). Since the $\mathrm{B}_{\lambda\mu}$ may be assumed finitely presented over $\sh{O}_{Y,y}$, it follows again from \sref{IV.8.1.2},~a) and \sref{IV.17.7.8} that they are filtered direct limits of \'etale $\mathrm{A}$-algebras. The double direct-limit theorem therefore shows that $\mathrm{Y}''$ is the projective limit of a filtered family $(\mathrm{Y}_\alpha')$ of affine schemes \'etale over $\mathrm{Y}$.

Let $x_\alpha'$ be the projection of $x''$ in $\mathrm{X}_\alpha'=\mathrm{X}\times_{\mathrm{Y}}\mathrm{Y}_\alpha'$. Since $\mathrm{V}''$ is quasi-compact, after increasing $\alpha$ we may write
\[
  \mathrm{V}''=\mathrm{V}_\alpha'\times_{\mathrm{Y}_\alpha'}\mathrm{Y}'',
\]
where $\mathrm{V}_\alpha'$ is an open neighbourhood of $x_\alpha'$ in $\mathrm{X}_\alpha'$ \sref{IV.8.2.11}. Finally, since $\mathrm{V}''$ is of finite presentation over $\mathrm{Y}''$, \sref{IV.8.10.5},~(x) shows that, for a suitable $\alpha$, $\mathrm{V}_\alpha'$ is finite over $\mathrm{Y}_\alpha'$, as required.
\end{proof}

\begin{remark}[18.12.2]
\label{IV.18.12.2}
The preceding proof shows (taking \sref{IV.18.6.2} into account) that $\mathrm{Y}'$ may moreover be chosen so that, if $y'=f'(x')$, the homomorphism $k(y)\to k(y')$ is bijective.
\end{remark}

\begin{corollary}[18.12.3]
\label{IV.18.12.3}
Let $f:\mathrm{X}\to\mathrm{Y}$ be a separated morphism locally of finite type, and let $y$ be a point of $\mathrm{Y}$ such that the space $f^{-1}(y)$ is finite and discrete.
Then there exist an \'etale morphism $\mathrm{Y}'\to\mathrm{Y}$, a point $y'\in\mathrm{Y}'$ above $y$ such that $k(y')=k(y)$, and a decomposition of $\mathrm{X}'=\mathrm{X}\times_{\mathrm{Y}}\mathrm{Y}'$ as a disjoint union of two subschemes $\mathrm{X}_1'$ and $\mathrm{X}_2'$ induced on open subsets of $\mathrm{X}'$, such that the restriction of $f'=f_{(\mathrm{Y}')}:\mathrm{X}'\to\mathrm{Y}'$ to $\mathrm{X}_1'$ is finite and $\mathrm{X}_2'\cap f'^{-1}(y')=\varnothing$.
\end{corollary}

\begin{proof}
Let $n$ be the number of points of $f^{-1}(y)$. We argue by induction on $n$, the assertion being trivial for $n=0$. Choose $x\in f^{-1}(y)$. By \sref{IV.18.12.1} and \sref{IV.18.12.2}, there are an \'etale morphism $\mathrm{Y}_1\to\mathrm{Y}$, a point $y_1\in\mathrm{Y}_1$ above $y$ with $k(y_1)=k(y)$, and, on putting $\mathrm{S}_1=\mathrm{X}\times_{\mathrm{Y}}\mathrm{Y}_1$ and $f_1=f_{(\mathrm{Y}_1)}$, a decomposition
\[
  \mathrm{S}_1=\mathrm{V}_1\amalg\mathrm{X}_1
\]
into open subschemes, where $\mathrm{V}_1$ is finite over $\mathrm{Y}_1$ and contains a point $x_1$ above $x$.

\oldpage[IV]{182}
Because $k(y_1)=k(y)$, the fibre $f_1^{-1}(y_1)$ is isomorphic to $f^{-1}(y)$; hence $\mathrm{X}_1\cap f_1^{-1}(y_1)$ is finite, discrete, and has $n-1$ points. Apply the induction hypothesis to $f_1|\mathrm{X}_1$. There are an \'etale morphism $\mathrm{Y}'\to\mathrm{Y}_1$, a point $y'\in\mathrm{Y}'$ above $y_1$ with $k(y')=k(y_1)$, and, if $p:\mathrm{X}'\to\mathrm{S}_1$ is the canonical projection, a decomposition
\[
  p^{-1}(\mathrm{X}_1)=\mathrm{U}'\amalg\mathrm{X}_2'
\]
such that $\mathrm{U}'$ is finite over $\mathrm{Y}'$ and $\mathrm{X}_2'\cap f'^{-1}(y')=\varnothing$. The open subscheme $\mathrm{V}'=p^{-1}(\mathrm{V}_1)$ is finite over $\mathrm{Y}'$, and
\[
  \mathrm{X}'=\mathrm{U}'\amalg\mathrm{V}'\amalg\mathrm{X}_2'.
\]
Taking $\mathrm{X}_1'=\mathrm{U}'\amalg\mathrm{V}'$ proves the assertion.
\end{proof}

The following corollary improves \sref{IV.8.11.1}.

\begin{corollary}[18.12.4]
\label{IV.18.12.4}
For a morphism $f:\mathrm{X}\to\mathrm{Y}$ to be finite, it is necessary and sufficient that it be separated, universally closed, locally of finite type, and that $f^{-1}(y)$ be a finite discrete space for every $y\in\mathrm{Y}$. In particular, every proper quasi-finite morphism is finite.
\end{corollary}

\begin{proof}
The necessity is clear. Conversely, apply \sref{IV.18.12.3} at an arbitrary point $y\in\mathrm{Y}$. With the notation of that corollary, $f'$ is closed. Since $\mathrm{X}_2'$ is closed in $\mathrm{X}'$, the set $f'(\mathrm{X}_2')$ is closed in $\mathrm{Y}'$ and does not contain $y'$. There is therefore an open neighbourhood $\mathrm{U}'$ of $y'$ in $\mathrm{Y}'$ such that $\mathrm{U}'\to\mathrm{Y}$ is of finite type and
\[
  f'^{-1}(\mathrm{U}')=\mathrm{X}\times_{\mathrm{Y}}\mathrm{U}'
\]
is finite over $\mathrm{U}'$.

Let $\mathrm{U}$ be the image of $\mathrm{U}'$ in $\mathrm{Y}$. It is open by \sref[II]{II.3.1}, since $\mathrm{Y}'\to\mathrm{Y}$ is \'etale, hence flat and locally of finite presentation. Moreover
\[
  \mathrm{X}\times_{\mathrm{Y}}\mathrm{U}'
  =f^{-1}(\mathrm{U})\times_{\mathrm{U}}\mathrm{U}'.
\]
The morphism $\mathrm{U}'\to\mathrm{U}$ is faithfully flat and quasi-compact; hence \sref{IV.2.7.1},~(xv) implies that $f^{-1}(\mathrm{U})\to\mathrm{U}$ is finite. Since $y$ was arbitrary, $f$ is finite.
\end{proof}

\begin{remark}[18.12.5]
\label{IV.18.12.5}
Observe that the proof of \sref{IV.18.12.4} does not use universal closedness in its full form, but only the fact that $f_{(\mathrm{Y}')}$ is closed for every \'etale morphism $\mathrm{Y}'\to\mathrm{Y}$.
\end{remark}

\begin{corollary}[18.12.6]
\label{IV.18.12.6}
Let $f:\mathrm{X}\to\mathrm{Y}$ be a morphism of preschemes.
The following conditions are equivalent:
\begin{enumerate}[label=\alph*)]
  \item $f$ is a closed immersion.
  \item $f$ is a proper monomorphism.
  \item $f$ is proper and, for all $y\in\mathrm{Y}$, the $k(y)$-prescheme $\mathrm{X}_y=f^{-1}(y)$ is radicial and geometrically reduced on $k(y)$ (i.e.\ empty or $k(y)$-isomorphic to $\operatorname{Spec}(k(y))$).
\end{enumerate}
\end{corollary}

\begin{proof}
This follows from \sref{IV.18.12.4} in the same way that \sref{IV.8.11.5} follows from \sref{IV.8.11.1}.
\end{proof}

\begin{proposition}[18.12.7]
\label{IV.18.12.7}
Let $f:\mathrm{X}\to\mathrm{Y}$ be a morphism locally of finite type, and let $y$ be a point of $\mathrm{Y}$.
There exists an open neighbourhood $\mathrm{U}$ of $y$ such that the restriction $f^{-1}(\mathrm{U})\to\mathrm{U}$ is a closed immersion if and only if $\mathrm{X}_y$ is a radicial and geometrically reduced $k(y)$-scheme (that is, empty or $k(y)$-isomorphic to $\operatorname{Spec}(k(y))$), and there exists an open neighbourhood $\mathrm{V}$ of $y$ such that $f^{-1}(\mathrm{V})\to\mathrm{V}$ is universally closed.
\end{proposition}

The following lemma will be used in the proof.

\begin{lemma}[18.12.7.1]
\label{IV.18.12.7.1}
Let $f:\mathrm{X}\to\mathrm{Y}$ be a closed morphism, and let $y\in\mathrm{Y}$ be a point such that every neighbourhood of $\mathrm{X}_y$ contains an affine neighbourhood of $\mathrm{X}_y$ (a condition which is always satisfied when $\mathrm{X}_y$ is empty or consists of a single point).
Then there exists an affine open neighbourhood $\mathrm{U}$ of $y$ such that $f^{-1}(\mathrm{U})\to\mathrm{U}$ is affine.
\end{lemma}

\begin{proof}
Let $\mathrm{U}_0$ be an affine open neighbourhood of $y$ in $\mathrm{Y}$, and let $\mathrm{V}$ be an affine open neighbourhood of $\mathrm{X}_y$ contained in $f^{-1}(\mathrm{U}_0)$. Since $f$ is closed, there is an affine open neighbourhood $\mathrm{U}\subset\mathrm{U}_0$ of $y$ such that $f^{-1}(\mathrm{U})\subset\mathrm{V}$. The restriction $g:\mathrm{V}\to\mathrm{U}_0$ of $f$ is affine, and hence so is its restriction $f^{-1}(\mathrm{U})\to\mathrm{U}$ \sref[II]{II.1.2.5}.
\end{proof}

\begin{proof}[Proof of {\hyperlink{ega4.c04.prop.18-12-7}{Proposition 18.12.7}}]
The conditions are plainly necessary. To prove their sufficiency, we may suppose first that $\mathrm{X}_y\ne\varnothing$, and, by the second condition, that $f$ is universally closed. By {\hyperlink{ega4.c04.lem.18-12-7-1}{Lemma~18.12.7.1}}, after shrinking $\mathrm{Y}$ about $y$, we may also suppose that $f$ is affine, and hence separated.

We may now suppose that $\mathrm{X}$ and $\mathrm{Y}$ are affine. Since $f$ is of finite type, it is proper. It remains, by \sref{IV.18.12.4} and \sref{IV.18.12.6}, to make $f$ quasi-finite over a neighbourhood of $y$. There is an open neighbourhood $\mathrm{V}$ of the one-point set $\mathrm{X}_y$ such that $f|\mathrm{V}$ is quasi-finite \sref{IV.13.1.4}. Since $f$ is closed, there is an open neighbourhood $\mathrm{U}$ of $y$ such that $f^{-1}(\mathrm{U})\subset\mathrm{V}$. Thus $f^{-1}(\mathrm{U})\to\mathrm{U}$ is proper and quasi-finite, hence finite by \sref{IV.18.12.4}; its fibres satisfy condition~c) of \sref{IV.18.12.6}, so it is a closed immersion. The case $\mathrm{X}_y=\varnothing$ is immediate by the same closedness argument.
\end{proof}

\begin{proposition}[18.12.8]
\label{IV.18.12.8}
Let $f:\mathrm{X}\to\mathrm{Y}$ be a morphism of preschemes.
For $f$ to be integral, it is necessary and sufficient that $f$ be affine and universally closed.
\end{proposition}

\begin{proof}
The necessity follows from \sref[II]{II.6.1.10}. Conversely, suppose that $f$ is affine and universally closed. We may take $\mathrm{Y}=\operatorname{Spec}(\mathrm{A})$ and $\mathrm{X}=\operatorname{Spec}(\mathrm{B})$; it is enough to prove that every $b\in\mathrm{B}$ is integral over $\mathrm{A}$ \sref[II]{II.6.1.1}.

Let $\mathrm{B}'=\mathrm{A}[b]\subset\mathrm{B}$ and set $\mathrm{X}'=\operatorname{Spec}(\mathrm{B}')$. Then $f$ factors as
\[
  \mathrm{X}\xrightarrow{g}\mathrm{X}'\xrightarrow{h}\mathrm{Y},
\]
where $h$ is of finite type and $g$ is dominant because $\mathrm{B}'\to\mathrm{B}$ is injective \sref[I]{I.1.2.7}. Since $h$ is separated and $h\circ g$ is universally closed, $g$ is universally closed \sref[II]{II.5.4.3}, \sref[II]{II.5.4.9}; being dominant, it is therefore surjective. It follows from the same results that $h$ is universally closed as well.

\oldpage[IV]{183}
Thus $h$ is proper, since it is of finite type and separated. For every $y\in\mathrm{Y}$, the fibre $h^{-1}(y)$ is proper and affine over $k(y)$, hence finite \sref[III]{III.4.4.2}. Thus $h$ is quasi-finite, and \sref{IV.18.12.4} shows that $h$ is finite. Consequently $\mathrm{B}'$ is a finite $\mathrm{A}$-algebra and $b$ is integral over $\mathrm{A}$.
\end{proof}

\begin{remark}[18.12.9]
\label{IV.18.12.9}
It may be that a morphism $f:\mathrm{X}\to\mathrm{Y}$ is integral whenever it is separated and universally closed and, for every $y\in\mathrm{Y}$, the induced morphism $f^{-1}(y)\to\operatorname{Spec}(k(y))$ is integral.
To establish this it would be enough to prove that these conditions imply that $f$ is affine, or equivalently that every fibre $\mathrm{X}_y$ is contained in an affine open neighbourhood.
\end{remark}

\begin{corollary}[18.12.10]
\label{IV.18.12.10}
Every injective universally closed morphism $f:\mathrm{X}\to\mathrm{Y}$ is integral.
\end{corollary}

\begin{proof}
It suffices, by \sref{IV.18.12.8}, to prove that $f$ is affine, which follows from the lemma \sref{IV.18.12.7.1} and the hypothesis.
\end{proof}

\begin{corollary}[18.12.11]
\label{IV.18.12.11}
Let $f:\mathrm{X}\to\mathrm{Y}$ be a morphism of preschemes (resp.\ a locally of finite type morphism).
For $f$ to be a universal homeomorphism \sref{IV.2.4.2}, it is necessary and sufficient that $f$ be integral (resp.\ finite), radicial, and surjective.
\end{corollary}

\begin{proof}
One knows that the conditions are sufficient \sref{IV.2.4.5}; they are necessary by virtue of \sref{IV.18.12.10} and \sref{IV.18.12.4}.
\end{proof}

The following proposition improves \sref{IV.8.11.2}.

\begin{proposition}[18.12.12]
\label{IV.18.12.12}
Every morphism $f:\mathrm{X}\to\mathrm{Y}$ which is quasi-finite and separated is quasi-affine.
\end{proposition}

\begin{proof}
Put $\mathscr{A}=f_*(\sh{O}_X)$. This is a quasi-coherent $\sh{O}_Y$-algebra because $f$ is quasi-compact and separated \sref{IV.1.7.4}. Let $\mathrm{Z}=\operatorname{Spec}(\mathscr{A})$ \sref[II]{II.1.3.1}; then $f$ has the canonical factorisation
\[
  \mathrm{X}\xrightarrow{g}\mathrm{Z}\xrightarrow{h}\mathrm{Y},
\]
where $h$ is affine and $g$ corresponds to the identity automorphism of $\mathscr{A}$ \sref[II]{II.1.2.7}. It is enough to prove that $g$ is an open immersion, or equivalently \sref{IV.17.9.1} that $g$ is \'etale and radicial. It is plainly enough to establish this at every point of $f^{-1}(y)$, for each $y\in\mathrm{Y}$.

Fix $y\in\mathrm{Y}$ and apply \sref{IV.18.12.3} to $f$, retaining its notation. The morphism $\mathrm{Y}'\to\mathrm{Y}$ is flat, and $f$ is quasi-compact and separated; hence there is a canonical isomorphism
\[
  f'_*(\sh{O}_{X'})
  =\mathscr{A}\otimes_{\sh{O}_Y}\sh{O}_{Y'}
  =:\mathscr{A}'
\]
by \sref{IV.2.3.1}. Thus $\mathrm{Z}'=\mathrm{Z}\times_{\mathrm{Y}}\mathrm{Y}'$ identifies with $\operatorname{Spec}(f'_*(\sh{O}_{X'}))$, and the canonical factorisation
\[
  \mathrm{X}'\xrightarrow{g'}\mathrm{Z}'\xrightarrow{h'}\mathrm{Y}'
\]
is obtained from the preceding one by base change \sref[II]{II.1.2.7}.

The decomposition $\mathrm{X}'=\mathrm{X}_1'\amalg\mathrm{X}_2'$ yields a direct-product decomposition $\mathscr{A}'=\mathscr{A}_1'\times\mathscr{A}_2'$ and hence $\mathrm{Z}'=\mathrm{Z}_1'\amalg\mathrm{Z}_2'$, where $\mathrm{Z}_i'=\operatorname{Spec}(\mathscr{A}_i')$ and $g'(\mathrm{X}_i')\subset\mathrm{Z}_i'$. Since $\mathrm{X}_1'$ is finite over $\mathrm{Y}'$, it is affine over $\mathrm{Y}'$, and therefore $g'|\mathrm{X}_1':\mathrm{X}_1'\to\mathrm{Z}_1'$ is an isomorphism. Moreover $\mathrm{X}_2'\cap f'^{-1}(y')=\varnothing$. It follows that $g'$ is \'etale and radicial at every point of $f'^{-1}(y')$.

Since $\mathrm{Y}'\to\mathrm{Y}$ is flat and locally of finite presentation, \sref{IV.17.7.4} shows that $g$ is \'etale at the projections of the points of $f'^{-1}(y')$, that is, at all points of $f^{-1}(y)$ \sref[I]{I.3.5.2}. On the other hand, the induced map
\[
  g'_{y'}:f'^{-1}(y')\longrightarrow h'^{-1}(y')
\]
is radicial; since $k(y')=k(y)$, it is the scalar extension of $g_y:f^{-1}(y)\to h^{-1}(y)$. Hence $g_y$ is radicial as well. Thus $g$ is \'etale and radicial at every point of $\mathrm{X}$, which proves the proposition.
\end{proof}

The following assertion likewise improves \sref{IV.8.12.6}.

\begin{corollary}[18.12.13]
\label{IV.18.12.13}
(``Main Theorem'' of Zariski).
Let $\mathrm{Y}$ be a quasi-compact and quasi-separated scheme, and let $f:\mathrm{X}\to\mathrm{Y}$ be a quasi-finite separated morphism.
Then $f$ admits a factorisation
\[
  \mathrm{X}\xrightarrow{g}\mathrm{Z}\xrightarrow{u}\mathrm{Y}
\]
where $g$ is an open immersion (necessarily quasi-compact) and $u$ is finite.
\end{corollary}

\begin{proof}
Indeed, it follows from \sref{IV.18.12.12} that $f$ is quasi-affine, hence quasi-projective.
One can therefore apply \sref{IV.8.12.8}, where the hypothesis on $\sh{O}_Y$-Module ample can in fact be replaced by the sole hypothesis that $\mathrm{Y}$ is quasi-compact and quasi-separated: indeed, it follows from \sref{IV.8.12.3} that the existence of the factorisation announced is a \emph{local} property on $\mathrm{Y}$, and it suffices to prove it when $\mathrm{Y}$ is affine.
\end{proof}

\begin{remark}[18.12.14]
\label{IV.18.12.14}
There is also a proof of \sref{IV.18.12.13} analogous to that of \sref{IV.18.12.12}. It does not use \sref{IV.8.12.8}, but it does use \sref{IV.18.12.3}, and hence \sref{IV.18.5.11}, which itself uses the ``Main Theorem'' in its local form \sref{IV.8.12.9}.

Retain the notation of the proof of \sref{IV.18.12.12}, and let $\mathscr{B}$ be the quasi-coherent $\sh{O}_Y$-algebra which is the integral closure of $\sh{O}_Y$ in $\mathscr{A}$ \sref[II]{II.6.3.2}. Put $\mathrm{T}=\operatorname{Spec}(\mathscr{B})$. By \sref{IV.8.12.3}, it is enough to prove that the $\mathrm{Y}$-morphism $g:\mathrm{X}\to\mathrm{T}$ corresponding to the canonical injection $\mathscr{B}\to\mathscr{A}$ is an immersion; for this it is enough, by \sref{IV.17.9.1}, to show that $g$ is \'etale and radicial.

With the notation of \sref{IV.18.12.12}, we may suppose that $\mathrm{Y}=\operatorname{Spec}(\mathrm{C})$ and $\mathrm{Y}'=\operatorname{Spec}(\mathrm{C}')$ are affine, with $\mathrm{C}'$ an \'etale $\mathrm{C}$-algebra \sref{IV.17.3.2}. Write $\mathscr{A}=\widetilde{\mathrm{A}}$, where $\mathrm{A}$ is a $\mathrm{C}$-algebra, put $\mathrm{A}'=\mathrm{A}\otimes_\mathrm{C}\mathrm{C}'$, and write $\mathscr{B}=\widetilde{\mathrm{B}}$, where $\mathrm{B}$ is the integral closure of $\mathrm{C}$ in $\mathrm{A}$. Recall that $\mathrm{A}'\simeq\mathrm{A}_1'\times\mathrm{A}_2'$, with $\mathrm{A}_1'$ a finite $\mathrm{C}'$-algebra. It is enough to prove that
\[
  \mathrm{B}'=\mathrm{B}\otimes_\mathrm{C}\mathrm{C}',
\]
\oldpage[IV]{184}
which identifies by flatness with a $\mathrm{C}'$-subalgebra of $\mathrm{A}'$, contains $\mathrm{A}_1'$. Indeed, $\mathrm{B}'$ would then decompose as $\mathrm{A}_1'\times\mathrm{A}_2''$, where $\mathrm{A}_2''$ is a $\mathrm{C}'$-subalgebra of $\mathrm{A}_2'$. Thus $\mathrm{T}'=\mathrm{T}\times_\mathrm{Y}\mathrm{Y}'$ would be the disjoint union of $\mathrm{Z}_1'=\operatorname{Spec}(\mathrm{A}_1')$ and $\mathrm{T}_2''=\operatorname{Spec}(\mathrm{A}_2'')$, and $g'|\mathrm{X}_1'$ would be an isomorphism from $\mathrm{X}_1'$ onto $\mathrm{Z}_1'$. One would then conclude as in \sref{IV.18.12.12}.

To prove that $\mathrm{B}'$ contains $\mathrm{A}_1'$, it plainly suffices to prove the following proposition, which partially extends \sref{IV.6.14.4} when the Noetherian hypotheses are removed.
\end{remark}

\begin{proposition}[18.12.15]
\label{IV.18.12.15}
Let $\mathrm{C}$ be a ring, let $\mathrm{C}'$ be an \'etale $\mathrm{C}$-algebra, let $\mathrm{A}$ be a $\mathrm{C}$-algebra, and let $\mathrm{B}$ be the integral closure of $\mathrm{C}$ in $\mathrm{A}$. Put $\mathrm{A}'=\mathrm{A}\otimes_\mathrm{C}\mathrm{C}'$ and $\mathrm{B}'=\mathrm{B}\otimes_\mathrm{C}\mathrm{C}'$, where $\mathrm{B}'$ is identified with a subalgebra of $\mathrm{A}'$. Then $\mathrm{B}'$ is the integral closure of $\mathrm{C}'$ in $\mathrm{A}'$.
\end{proposition}

\begin{proof}
Consider the filtered increasing family of finite-type $\mathrm{C}$-subalgebras of $\mathrm{A}$ and argue as in \sref{IV.6.14.4},~II. We may first suppose that $\mathrm{A}$ is a finite-type $\mathrm{C}$-algebra, say $\mathrm{A}=\mathrm{E}/\mathfrak{J}$, where $\mathrm{E}=\mathrm{C}[\mathrm{T}_1,\ldots,\mathrm{T}_r]$. Then $\mathrm{A}'=\mathrm{E}'/\mathfrak{J}\mathrm{E}'$, where $\mathrm{E}'=\mathrm{C}'[\mathrm{T}_1,\ldots,\mathrm{T}_r]$.

Let $(\mathfrak{J}_\lambda)$ be the filtered increasing family of finitely generated ideals of $\mathrm{E}$ contained in $\mathfrak{J}$, and let $\mathrm{B}_\lambda$ be the integral closure of $\mathrm{C}$ in $\mathrm{E}/\mathfrak{J}_\lambda$. The reasoning of \sref{IV.5.13.4} shows that $\mathrm{B}=\varinjlim_\lambda\mathrm{B}_\lambda$. Likewise, the integral closure of $\mathrm{C}'$ in $\mathrm{E}'/\mathfrak{J}\mathrm{E}'$ is the direct limit of the integral closures of $\mathrm{C}'$ in $\mathrm{E}'/\mathfrak{J}_\lambda\mathrm{E}'$. If the latter closure equals $\mathrm{B}_\lambda\otimes_\mathrm{C}\mathrm{C}'$ for every $\lambda$, then
\[
  \mathrm{B}'=\varinjlim_\lambda
  (\mathrm{B}_\lambda\otimes_\mathrm{C}\mathrm{C}')
\]
is the integral closure of $\mathrm{C}'$ in $\mathrm{A}'$. We are therefore reduced to the case where $\mathrm{A}$ is finitely presented over $\mathrm{C}$.

We next reduce to the case where $\mathrm{C}$ is Noetherian. Write $\mathrm{C}$ as the filtered union of its finite-type $\mathbb{Z}$-subalgebras $\mathrm{C}_\alpha$. By \sref{IV.17.7.8}, for some $\alpha$ there is an \'etale $\mathrm{C}_\alpha$-algebra $\mathrm{C}_\alpha'$ such that
\[
  \mathrm{C}'=\mathrm{C}_\alpha'\otimes_{\mathrm{C}_\alpha}\mathrm{C}.
\]
For $\beta\geq\alpha$, put $\mathrm{C}_\beta'=\mathrm{C}_\alpha'\otimes_{\mathrm{C}_\alpha}\mathrm{C}_\beta$; then $\mathrm{C}'=\varinjlim_{\beta\geq\alpha}\mathrm{C}_\beta'$. Since $\mathrm{A}$ is finitely presented over $\mathrm{C}$, after increasing $\alpha$ we may write
\[
  \mathrm{A}=\mathrm{A}_\alpha\otimes_{\mathrm{C}_\alpha}\mathrm{C},
\]
with $\mathrm{A}_\alpha$ a finite-type $\mathrm{C}_\alpha$-algebra. Put $\mathrm{A}_\beta=\mathrm{A}_\alpha\otimes_{\mathrm{C}_\alpha}\mathrm{C}_\beta$; then $\mathrm{A}=\varinjlim_{\beta\geq\alpha}\mathrm{A}_\beta$ \sref{IV.8.9.1}. If $\mathrm{B}_\beta$ is the integral closure of $\mathrm{C}_\beta$ in $\mathrm{A}_\beta$, the reasoning of \sref{IV.5.13.4} gives $\mathrm{B}=\varinjlim\mathrm{B}_\beta$. Similarly,
\[
  \mathrm{A}'=\varinjlim_{\beta\geq\alpha}\mathrm{A}_\beta',
  \qquad
  \mathrm{A}_\beta'
   =\mathrm{A}_\beta\otimes_{\mathrm{C}_\beta}\mathrm{C}_\beta'
   =\mathrm{A}_\alpha\otimes_{\mathrm{C}_\alpha}\mathrm{C}_\beta'.
\]
The same argument shows that it is enough to prove that $\mathrm{B}_\alpha'=\mathrm{B}_\alpha\otimes_{\mathrm{C}_\alpha}\mathrm{C}_\alpha'$ is the integral closure of $\mathrm{C}_\alpha'$ in $\mathrm{A}_\alpha'$. Thus we may suppose that $\mathrm{C}$ is Noetherian.

In that case the proposition is a special case of \sref{IV.6.14.4}. One should note, however, that when $\mathrm{C}'$ is \'etale over the Noetherian ring $\mathrm{C}$, the proof of \sref{IV.6.14.4} no longer requires the difficult theorem \sref{IV.6.14.1}. Indeed, the successive reductions in that proof lead, without using \sref{IV.6.14.1}, to the case in which $\mathrm{C}$ is integral and $\mathrm{A}$ is its field of fractions. Then $\mathrm{B}$ is normal, and since $\operatorname{Spec}(\mathrm{B}')\to\operatorname{Spec}(\mathrm{B})$ is \'etale, $\mathrm{B}'$ is normal by \sref{IV.17.5.7}. The argument then uses only the elementary lemma \sref{IV.6.14.1.1}, not the difficult part of \sref{IV.6.14.1}.
\end{proof}

\begin{proposition}[18.12.16]
\label{IV.18.12.16}
Let $g:\mathrm{Y}\to\mathrm{S}$ be a quasi-compact morphism, and let $f:\mathrm{X}\to\mathrm{Y}$ be separated and quasi-finite.
For every invertible $\sh{O}_Y$-module $\mathscr{L}$ which is ample relative to $\mathrm{S}$ \sref[II]{II.4.6.1}, the invertible $\sh{O}_X$-module $f^*(\mathscr{L})$ is ample relative to $\mathrm{S}$.
\end{proposition}

\begin{proof}
By \sref{IV.18.12.12}, the morphism $f$ is quasi-affine; hence $\sh{O}_X$ is ample for $f$ by \sref[II]{II.5.1.6}. It follows from \sref[II]{II.4.6.13},~(ii) that
\[
  \sh{O}_X\otimes_{\sh{O}_X}f^*(\mathscr{L}^{\otimes n})
  =f^*(\mathscr{L}^{\otimes n})
\]
is ample relative to $\mathrm{S}$ for all sufficiently large $n$. Since
\[
  f^*(\mathscr{L}^{\otimes n})=(f^*\mathscr{L})^{\otimes n},
\]
this implies that $f^*(\mathscr{L})$ is ample relative to $\mathrm{S}$ \sref[II]{II.4.6.9},~(i).
\end{proof}

\begin{corollary}[18.12.17]
\label{IV.18.12.17}
Let $\mathrm{Z}$ be an affine scheme, let $h:\mathrm{X}\to\mathrm{Z}$ be a morphism of finite type, and let $\mathscr{L}$ be an invertible $\sh{O}_X$-module.
With the notation of \sref[II]{II.4.5.2}, the conditions in \sref[II]{II.4.5.2} (equivalent to the ampleness of $\mathscr{L}$) are also equivalent to each of the following:

\begin{enumerate}
  \item[\rm b$''$)] $h$ is separated, $\mathrm{G}(\varepsilon)=\mathrm{X}$, and the canonical morphism $u:\mathrm{G}(\varepsilon)\to\operatorname{Proj}(\mathrm{S})$ has finite discrete fibres.
  \item[\rm b$'''$)] $\mathrm{G}(\varepsilon)=\mathrm{X}$ and the canonical morphism $u$ is radicial.
\end{enumerate}

If $\mathrm{Z}=\operatorname{Spec}(\mathrm{A})$, then $\mathrm{S}$ is canonically a graded $\mathrm{A}$-algebra; hence there is a structural morphism $g:\operatorname{Proj}(\mathrm{S})\to\mathrm{Z}$, which is separated \sref[II]{II.2.4.2}, and $h=g\circ u$. Since every radicial morphism is separated \sref[I]{I.1.8.7.1}, condition~b$'''$) implies condition~b$''$). Condition~b) of \sref[II]{II.4.5.2} plainly implies b$'''$), so it remains to prove that b$''$) implies that $\mathscr{L}$ is ample.

Since $h=g\circ u$ is of finite type and $g$ is separated, $u$ is of finite type \sref[I]{I.6.3.4}; condition~b$''$) therefore implies that $u$ is quasi-finite and separated \sref[I]{I.5.5.1}. Put $\mathrm{Y}=\operatorname{Proj}(\mathrm{S})$. As $\mathrm{X}$ is quasi-compact, there exist an integer $n>0$ and finitely many elements $f_j\in\mathrm{S}_n$ such that the inverse images $u^{-1}(\mathrm{D}_+(f_j))$ cover $\mathrm{X}$. We may thus regard $u$ as a quasi-finite separated morphism from $\mathrm{X}$ to the open subscheme
\[
  \mathrm{Y}'=\bigcup_j\mathrm{D}_+(f_j)
\]
of $\mathrm{Y}$.

The invertible $\sh{O}_{Y'}$-module
\[
  \mathscr{L}'=\sh{O}_Y(n)|\mathrm{Y}'
\]
\sref[II]{II.2.5.8}; it is very ample relative to $\mathrm{Z}$ \sref[II]{II.4.4.3}, and by definition
\[
  u^*(\mathscr{L}')=\mathscr{L}^{\otimes n}
\]
\sref[II]{II.3.7.9},~(ii). By \sref{IV.18.12.16}, $\mathscr{L}^{\otimes n}$ is therefore ample relative to $\mathrm{Z}$ (equivalently, ample \sref[II]{II.4.6.6}), and consequently $\mathscr{L}$ itself is ample \sref[II]{II.4.5.6},~(i).

\oldpage[IV]{185}
We leave it to the reader to state, in the same way, conditions equivalent to those of \sref[II]{II.4.6.3} for relatively ample $\sh{O}_X$-modules.
\end{corollary}
% END NATIVE INPUT source_aligned/ega4-18
% BEGIN NATIVE INPUT source_aligned/ega4-19
\section{Regular immersions and normal flatness}
\label{section:IV.19}

This section is devoted, on the one hand, to a study of regular immersions \sref{IV.16.9.2} $Y\to X$ of preschemes, especially when $X$ and $Y$ are flat over the same prescheme $S$; on the other hand, it gives supplements concerning $M$-regular sequences and normal flatness, notably generalizing flatness results established by H.~Hironaka in his theory of resolution of singularities~[35].


\subsection{Properties of regular immersions}
\label{subsection:IV.19.1}


\begin{proposition}[19.1.1]
\label{IV.19.1.1}
Let $X$ be a locally Noetherian prescheme, $j:Y\to X$ an immersion, and $y$ a regular point of $Y$. For $j$ to be a regular immersion at the point $y$, it is necessary and sufficient that $X$ be regular at the point $j(y)$.
\end{proposition}

Taking \sref{IV.16.9.10} into account, we are reduced to proving the following corollary.

\begin{corollary}[19.1.2]
\label{IV.19.1.2}
Let $A$ be a Noetherian local ring and $\mathfrak J$ an ideal of $A$; suppose that the quotient ring $B=A/\mathfrak J$ is regular. For $A$ to be regular, it is necessary and sufficient that the ideal $\mathfrak J$ be regular.
\end{corollary}

\begin{proof}
Indeed, if $A$ is regular, then $\mathfrak J$ is generated by part of a regular system of parameters of $A$ \sref[0]{0.17.1.9}, and hence $\mathfrak J$ is regular. Conversely, if $\mathfrak J$ is regular and $(x_i)_{1\leq i\leq r}$ is an $A$-regular sequence generating $\mathfrak J$, then $(x_i)$ is part of a system of parameters of $A$ \sref[0]{0.16.4.1}; it therefore follows from \sref[0]{0.17.1.7} that $A$ is regular.
\end{proof}

\begin{definition}[19.1.3]
\label{IV.19.1.3}
Let $X$ be a prescheme, $Y$ a subprescheme of $X$, and $U$ an open subset of $X$ containing $Y$ such that $Y$ is defined by an ideal $\sh J$ of $\sh O_U$; suppose that $\sh J/\sh J^2$ is a locally free $(\sh O_U/\sh J)$-module (as is the case when $Y$ is quasi-regularly immersed in $X$). For every $y\in Y$, the \emph{transverse codimension} of $Y$ in $X$ at $y$, denoted by $\operatorname{codim}_y^*(Y,X)$, is the rank of the free $(\sh O_{U,y}/\sh J_y)$-module $\sh J_y/\sh J_y^2$. If $f:Z\to X$ is an immersion with image $Y$, the transverse codimension of $f$ at $z\in Z$ is likewise the transverse codimension in $X$, at $f(z)$, of the subprescheme $Y$.
\end{definition}

\begin{proposition}[19.1.4]
\label{IV.19.1.4}
Let $X$ be a locally Noetherian prescheme and $Y$ a subprescheme regularly immersed in $X$. For every $y\in Y$, one has
\begin{equation}
  \operatorname{codim}_y^*(Y,X)=\operatorname{codim}_y(Y,X).
  \tag{19.1.4.1}\label{IV.19.1.4.1}
\end{equation}
\end{proposition}

\begin{proof}
By virtue of \sref{IV.5.1.3.2}, $\operatorname{codim}_y(Y,X)$ is equal to the least value of the dimensions of the local rings $A=\sh O_{X,z}$, where $z$ runs through the generic points of the irreducible components of $Y$ containing $y$. Such a point $z$ belongs to every neighbourhood of $y$ in $X$, and $\sh J_z/\sh J_z^2$ is a free $(A/\sh J_z)$-module of rank
\[
  n=\operatorname{codim}_y^*(Y,X).
\]
It therefore remains to prove that $\dim(A)=n$.

\oldpage[IV-4]{186}

Since $z$ is a maximal point of the subprescheme $Y$ defined by $\sh J$, one has $\dim(A/\sh J_z)=0$; hence $\sh J_z$ is an ideal of definition of the Noetherian local ring $A$. Moreover, the quasi-regularity of $\sh J$ implies that $\sh J_z^m/\sh J_z^{m+1}$ is a free $(A/\sh J_z)$-module of rank $\binom{m+n-1}{n-1}$. If $r$ denotes the length of the Artinian ring $A/\sh J_z$, the length of $\sh J_z^m/\sh J_z^{m+1}$ is therefore
\[
  r\binom{m+n-1}{n-1}.
\]
Consequently, the Hilbert--Samuel polynomial of $A$ for the $\sh J_z$-adic filtration has degree $n$, which proves the proposition by \sref[0]{0.16.2.3}.
\end{proof}

By virtue of \sref{IV.19.1.4}, henceforth we shall say ``codimension'' instead of ``transverse codimension'' when speaking of a subprescheme $Y$ regularly immersed in $X$, even when $X$ is not locally Noetherian.

\begin{proposition}[19.1.5]
\label{IV.19.1.5}
\medskip\noindent
\begin{enumerate}[label=(\roman*)]
  \item For an immersion $j:Y\to X$ to be open, it is necessary and sufficient that it be regular and of codimension $0$ at every point.
  \item Let $f:Y\to X$ be a regular (resp.\ quasi-regular) immersion, $g:X'\to X$ a flat morphism, $Y'=Y\times_X X'$; then the morphism $f'=f_{(X')}:Y'\to X'$ is a regular (resp.\ quasi-regular) immersion, and for every $y'\in Y'$, if $y$ is the projection of $y'$ in $Y$, the codimension of $f'$ at the point $y'$ is equal to the codimension of $f$ at the point $y$.
Conversely, if $f$ is an immersion and $g:X'\to X$ a faithfully flat and quasi-compact morphism, then, if $f'$ is a quasi-regular immersion, so is $f$.
  \item If $f:Y\to X$ and $g:Z\to Y$ are regular immersions, then $f\circ g:Z\to X$ is a regular immersion, and for every $z\in Z$, the codimension of $f\circ g$ at the point $z$ is the sum of the codimension of $g$ at the point $z$ and of the codimension of $f$ at the point $g(z)$.
Moreover, the sequence of canonical homomorphisms \sref{IV.16.2.7.1}
\[
  0\to g^*(\sh{N}_{Y/X})\to\sh{N}_{Z/X}\to\sh{N}_{Z/Y}\to 0
\]
{\hyperlink{ega4.c04.item.19-1-5-1}{(\textbf{19.1.5.1})}}
\label{IV.19.1.5.1}
is exact, and for every $z\in Z$, there exists a neighbourhood of $z$ in $Z$ in which the restrictions of the homomorphisms of this sequence form a split exact sequence.
  \item Let $X$ be a locally Noetherian prescheme, $f:Y\to X$ and $g:Z\to Y$ two immersions, $z$ a point of $Z$.
The following conditions are equivalent:
  \begin{enumerate}[label=\emph{\alph*)}]
    \item $g$ is a regular immersion at the point $z$ and $f$ is a regular immersion at the point $g(z)$.
    \item $g$ and $f\circ g$ are regular immersions at the point $z$.
    \item $f\circ g$ is a regular immersion at the point $z$, $f$ is a regular immersion at the point $g(z)$, and the sequence of canonical homomorphisms \sref{IV.19.1.5.1}, restricted to a sufficiently small open neighbourhood of $z$ in $Z$, is exact and split.
  \end{enumerate}
  \item Let $X$ be a locally Noetherian prescheme, $f:Y\to X$, $g:Z\to Y$ two immersions, $z$ a point of $Z$; suppose that $y=g(z)$ is a regular point of $Y$ and $x=f(g(z))$ a regular point
\oldpage[IV-4]{187}
of $X$ (so that, by \sref{IV.19.1.1}, $f$ is a regular immersion at the point $g(z)$); then,
for $f\circ g$ to be a regular immersion at the point $z$, it is necessary and sufficient that $g$ be a regular immersion at the point $z$.
\end{enumerate}
\end{proposition}

\begin{proof}
Assertion~(i) is another way of expressing \sref{IV.16.1.10}.
To prove the first assertion of~(ii), we can restrict ourselves to the case where $Y$ is a closed sub-prescheme of $X$ defined by a regular (resp.\ quasi-regular) ideal $\sh{J}$; then $Y'$ is the closed sub-prescheme of $X'$ defined by the ideal $g^*(\sh{J})\sh{O}_{X'}$, which is identified here with $g^*(\sh{J})$ since $g$ is flat; by replacing $X$ by a sufficiently small affine open, we can suppose that $\sh{J}$ admits a regular (resp.\ quasi-regular) sequence of generators $f_i$ (sections of $\sh{O}_X$ over $X$).
If the sequence $(f_i)$ is regular, so is the sequence $(f_i\otimes 1)$ (which generates $g^*(\sh{J})\sh{O}_{X'}$), by virtue of \sref[0]{0.15.2.5}; the analogous result for quasi-regular sequences follows immediately from the flatness of $g$ and from the criterion \sref{IV.16.9.4}.
Similarly, if $g$ is faithfully flat and quasi-compact, and if $g^*(\sh{J})$ is quasi-regular, $\sh{J}$ is of finite type by virtue of \sref{IV.2.5.2}; by flatness, we have $g^*(\sh{J})/g^*(\sh{J})^2=g^*(\sh{J}/\sh{J}^2)$, hence, since $g^*(\sh{J}/\sh{J}^2)$ is locally free by hypothesis \sref{IV.16.9.4}, it follows from \sref{IV.2.5.2} that $\sh{J}/\sh{J}^2$ is locally free.
Finally, condition~(iii) of \sref{IV.16.9.4} relative to $g^*(\sh{J})$ implies the same condition for $\sh{J}$ by virtue of the faithful flatness of $g$ \sref{IV.2.2.7}.
We conclude from \sref{IV.16.9.4} that $\sh{J}$ is quasi-regular.

To prove~(iii), we can likewise suppose that $Y$ and $Z$ are closed sub-preschemes of $X$, defined respectively by ideals $\sh{J}$, $\sh{K}$ of $\sh{O}_X$ such that $\sh{J}\subset\sh{K}$; moreover, we can suppose that there exist sections $f_i$ ($1\leq i\leq m$) of $\sh{J}$ over $X$ generating $\sh{J}$ and such that for $1\leq i\leq m$, the endomorphism of $\sh{O}_X/(\sum_{j=1}^{i-1}f_j\sh{O}_X)$ defined by $f_i$ is injective, and sections $\bar{f}_i$ ($m+1\leq i\leq n$) of $\sh{K}/\sh{J}$ over $X$ generating $\sh{K}/\sh{J}$ and such that for $m+1\leq i\leq n$, the endomorphism of
\[
  (\sh{O}_X/\sh{J})\big/\Bigl(\sum_{j=m+1}^{i-1}\bar{f}_j(\sh{O}_X/\sh{J})\Bigr)
\]
defined by $\bar{f}_i$ is injective.
For every $z\in Z$, let $U$ be an open neighbourhood of $z$ in $X$ such that there exist sections $f_i$ of $\sh{K}$ over $U$, whose classes in $\Gamma(U,\sh{K}/\sh{J})$ are $\bar{f}_i|U$ for $m+1\leq i\leq n$; then the restriction to $U$ of
\[
  (\sh{O}_X/\sh{J})\big/\Bigl(\sum_{j=m+1}^{i-1}\bar{f}_j(\sh{O}_X/\sh{J})\Bigr)
\]
is isomorphic to
\[
  \sh{O}_U\big/\Bigl((\sh{J}|U)+\sum_{j=m+1}^{i-1}f_j\sh{O}_U\Bigr)
  =\sh{O}_U\big/\Bigl(\sum_{j=1}^{i-1}f_j\sh{O}_U\Bigr),
\]
and we see therefore that the sequence $(f_i)_{1\leq i\leq n}$ is regular in $\sh{O}_U$; as it generates $\sh{K}|U$, this proves the first assertion of~(iii).
To prove the last assertion of~(iii), we can restrict ourselves (with the same notation) to the case where the images $t_i$ of the $f_i$ in $\sh{K}/\sh{K}^2$ (for $1\leq i\leq n$) form a basis of this $(\sh{O}_X/\sh{K})$-module \sref{IV.16.9.5}; for the same reason, we can suppose that the $t_i$ for $1\leq i\leq m$ are the canonical images of elements of a basis of the $(\sh{O}_X/\sh{K})$-module $g^*(\sh{J}/\sh{J}^2)$, and that the canonical images $\bar{t}_i$ of the $t_i$ in $(\sh{K}/\sh{J})/(\sh{K}/\sh{J})^2$ for $m+1\leq i\leq n$ form a basis of this $(\sh{O}_X/\sh{K})$-module.
This completes the proof of~(iii).

Let us pass to the proof of~(iv).
The fact that \emph{a)} implies \emph{c)} follows from~(iii).
Let us prove that \emph{c)} implies \emph{a)}.
Set again $\mathrm{A}=\sh{O}_{X,f(g(z))}$, so that $\sh{O}_{Y,g(z)}=\mathrm{A}/\mathfrak{J}$,
\oldpage[IV-4]{188}
$\sh{O}_{Z,z}=\mathrm{A}/\mathfrak{R}$ with $\mathfrak{J}\subset\mathfrak{R}$.
By hypothesis, $\mathfrak{J}$ and $\mathfrak{R}$ are regular ideals of $\mathrm{A}$ and we have a split exact sequence
\[
  0\to\mathfrak{J}/\mathfrak{J}\mathfrak{R}\to\mathfrak{R}/\mathfrak{R}^2\to(\mathfrak{R}/\mathfrak{J})/(\mathfrak{R}/\mathfrak{J})^2\to 0
\]
(cf.\ \sref{IV.16.2.7}).
This implies that there exist elements $f_j$ ($1\leq j\leq r+s$) of $\mathfrak{R}$ whose images $f_j'$ in $\mathfrak{R}/\mathfrak{R}^2$ form a basis of this $(\mathrm{A}/\mathfrak{R})$-module, and such that, in addition, the $f_j$ of index $j\leq r$ belong to $\mathfrak{J}$ and are such that their images in $\mathfrak{J}/\mathfrak{J}\mathfrak{R}$ form a basis of this $(\mathrm{A}/\mathfrak{R})$-module.
The $f_j$ for $1\leq j\leq r+s$ generate $\mathfrak{R}$, and the $f_j$ for $1\leq j\leq r$ generate $\mathfrak{J}$ since $\mathrm{A}$ is Noetherian (Bourbaki, \emph{Alg.\ comm.}, chap.~II, \S~3, n\textsuperscript{o}~2, prop.~5); as $\mathfrak{R}$ is a regular ideal, it follows from \sref{IV.16.9.5} that the sequence $(f_j)_{1\leq j\leq r+s}$ is regular.
By definition, the images of $f_{r+1},\ldots,f_{r+s}$ in $\mathfrak{R}/\mathfrak{J}$ form a regular sequence in this $(\mathrm{A}/\mathfrak{J})$-module, which completes the proof that \emph{c)} implies \emph{a)} in~(iv).

Finally, taking into account \sref{IV.16.9.10}, it is immediate that~(v), as well as the equivalence of \emph{a)} and \emph{b)} in~(iv), follow from the following corollary.
\end{proof}

\begin{corollary}[19.1.6]
\label{IV.19.1.6}
Let $\mathrm{A}$ be a local Noetherian ring, $\mathfrak{J}$, $\mathfrak{R}$ two ideals of $\mathrm{A}$ contained in its maximal ideal $\mathfrak{m}$ and such that $\mathfrak{J}\subset\mathfrak{R}$, $\mathrm{A}'=\mathrm{A}/\mathfrak{J}$, $\mathfrak{R}'=\mathfrak{R}/\mathfrak{J}$.
Then:
\begin{enumerate}[label=(\roman*)]
  \item If $\mathfrak{J}$ and $\mathfrak{R}'$ are regular ideals, so is $\mathfrak{R}$.
  \item If $\mathfrak{R}$ and $\mathfrak{R}'$ are regular ideals, so is $\mathfrak{J}$.
  \item Suppose the rings $\mathrm{A}$ and $\mathrm{A}'=\mathrm{A}/\mathfrak{J}$ are regular.
For the ideal $\mathfrak{R}$ to be regular, it is necessary and sufficient that $\mathfrak{R}'$ be so.
\end{enumerate}
\end{corollary}

\begin{proof}
Assertion~(i) is a particular case of \sref{IV.19.1.5},~(iii).
To prove~(ii), consider the canonical surjective homomorphism $u:\mathfrak{R}/\mathfrak{R}^2\to\mathfrak{R}'/\mathfrak{R}'^2=\mathfrak{R}/(\mathfrak{R}^2+\mathfrak{J})$.
Since by hypothesis $\mathfrak{R}/\mathfrak{R}^2$ and $\mathfrak{R}'/\mathfrak{R}'^2$ are free $(\mathrm{A}/\mathfrak{R})$-modules, whose ranks we denote respectively by $p+q$ and $q$, the kernel of $u$ is an $(\mathrm{A}/\mathfrak{R})$-module \emph{projective} (Bourbaki, \emph{Alg.}, chap.~II, 3\textsuperscript{e} ed., \S~2, n\textsuperscript{o}~2, prop.~4), hence \emph{free} of rank $p$ since $\mathrm{A}/\mathfrak{R}$ is a local Noetherian ring (\sref[0]{0.10.1.3}); in other words, there exists a basis of $\mathfrak{R}/\mathfrak{R}^2$ of which the first $p$ elements form a basis of $(\mathfrak{R}^2+\mathfrak{J})/\mathfrak{R}^2$.
This means (by virtue of Nakayama's lemma) that there exists a system of generators $(f_i)_{1\leq i\leq p+q}$ of $\mathfrak{R}$, forming a regular sequence in $\mathrm{A}$, such that the sequence $(f_i)_{1\leq i\leq p}$ consists of elements of $\mathfrak{J}$; it remains to show that this last sequence generates $\mathfrak{J}$.
For this, denote by $\mathfrak{J}'\subset\mathfrak{J}$ the ideal it generates; considering the ring $\mathrm{A}/\mathfrak{J}'$, we see that we are reduced to proving the following lemma:
\end{proof}

\begin{lemma}[19.1.6.1]
\label{IV.19.1.6.1}
Let $\mathrm{B}$ be a Noetherian local ring with maximal ideal $\mathfrak{n}$, $(g_j)_{1\leq j\leq q}$ a regular sequence in $\mathrm{B}$ consisting of elements of $\mathfrak{n}$, $\mathfrak{b}$ the ideal it generates, and $\mathfrak{c}\subset\mathfrak{b}$ an ideal such that the canonical images of the $g_j$ in $\mathrm{B}/\mathfrak{c}$ still form a regular sequence in this ring.
Then $\mathfrak{c}=0$.
\end{lemma}

\begin{proof}
The lemma being trivial for $q=0$, we reason by induction on $q$.
The induction hypothesis, applied to the ring $\mathrm{B}/g_1\mathrm{B}$ and to the canonical images of the $g_j$ ($2\leq j\leq q$)
\oldpage[IV-4]{189}
and of $\mathfrak{c}$ in this quotient ring, shows that $\mathfrak{c}\subset g_1\mathrm{B}$.
Let $\mathfrak{d}$ be the ideal of $\mathrm{B}$ consisting of the $x\in\mathrm{B}$ such that $g_1x\in\mathfrak{c}$; we thus have $\mathfrak{c}=g_1\mathfrak{d}$; but since by hypothesis $g_1$ is regular in $\mathrm{B}/\mathfrak{c}$, we necessarily have $\mathfrak{d}=\mathfrak{c}$, hence $\mathfrak{c}=g_1\mathfrak{c}$.
Since $\mathfrak{c}$ is of finite type and $g_1$ is contained in the radical of $\mathrm{B}$, Nakayama's lemma shows that $\mathfrak{c}=0$.

Let us now prove assertion~(iii) of \sref{IV.19.1.6}.
By virtue of~(i) and of \sref{IV.19.1.2}, it suffices to show that if $\mathfrak{R}$ is regular, so is $\mathfrak{R}'$.
Note that $\mathfrak{J}$ is regular \sref{IV.19.1.2}, and reason by induction on the rank $q$ of the $(\mathrm{A}/\mathfrak{J})$-module $\mathfrak{J}/\mathfrak{J}^2$.
If $q=0$, we have $\mathfrak{J}=0$ by Nakayama's lemma and the assertion is trivial.
Since $\mathrm{A}$ and $\mathrm{A}/\mathfrak{J}$ are regular, we know (\sref[0]{0.17.1.9}) that $\mathfrak{J}$ is generated by a subset of $q$ elements of a regular system of parameters of $\mathrm{A}$, hence, if $f$ is one of the elements of this system of generators of $\mathfrak{J}$, $\mathrm{A}_1=\mathrm{A}/f\mathrm{A}$ is regular and $f\notin\mathfrak{m}^2$ (\sref[0]{0.17.1.8}).
Let $\mathfrak{J}_1,\mathfrak{R}_1$ be the images of $\mathfrak{J},\mathfrak{R}$ in $\mathrm{A}_1$; we have $\mathfrak{J}_1\subset\mathfrak{R}_1$, $\mathrm{A}_1/\mathfrak{J}_1=\mathrm{A}/\mathfrak{J}$, hence $\mathrm{A}_1/\mathfrak{J}_1$ is regular and consequently $\mathfrak{J}_1$ is a regular ideal \sref{IV.19.1.2}.
Let us show that $\mathfrak{R}_1$ is a regular ideal in $\mathrm{A}_1$; since $\mathfrak{R}$ is a regular ideal in $\mathrm{A}$, it suffices to show that $f$ is part of a regular system of generators of $\mathfrak{R}$, and for this \sref{IV.16.9.5} it suffices to show that the image of $f$ in $\mathfrak{R}/(\mathfrak{R}^2+\mathfrak{m}\mathfrak{R})=\mathfrak{R}/\mathfrak{m}\mathfrak{R}$ is part of a basis of this $(\mathrm{A}/\mathfrak{m})$-vector space, in other words that $f\notin\mathfrak{m}\mathfrak{R}$, which indeed follows from $f\notin\mathfrak{m}^2$.
Then, as $\mathfrak{J}_1/\mathfrak{J}_1^2$ is a free $(\mathrm{A}_1/\mathfrak{J}_1)$-module of rank $q-1$, the induction hypothesis shows that $\mathfrak{R}_1'=\mathfrak{R}_1/\mathfrak{J}_1$ is a regular ideal of $\mathrm{A}_1/\mathfrak{J}_1$, and since $\mathfrak{R}'$ is identified with $\mathfrak{R}_1'$ in the canonical isomorphism between $\mathrm{A}/\mathfrak{J}$ and $\mathrm{A}_1/\mathfrak{J}_1$, $\mathfrak{R}'$ is regular.
\end{proof}

\begin{remarks}[19.1.7]
\label{IV.19.1.7}
\medskip\noindent
\begin{enumerate}[label=(\roman*)]
  \item We do not know whether the composite of two quasi-regular immersions is quasi-regular.
  \item In a ringed space with local rings $X$, for every ideal $\sh{J}$ of $\sh{O}_X$, $\sh{O}_X/\sh{J}$ is of finite type and therefore has closed support in $X$ (\sref[0]{0.5.2.2}).
We can therefore define, as in \sref[I]{I.4.1.3} and \sref{IV.4.2.1}, the notions of a (locally closed) ringed subspace of $X$, locally defined by an ideal $\sh{J}$ of a sheaf $\sh{O}_U$, where $U$ is an open subset of $X$, and of an immersion.
The definitions of quasi-regular and regular immersions and of transverse codimension then apply without change, and the results of \sref{IV.19.1.5},~(i) to~(iv) are still valid,
\footnote{The printed text reads ``{\hyperlink{ega4.c04.prop.19-1-4}{(19.1.4)}}, (i) to (iv)''; the reference must be to {\hyperlink{ega4.c04.prop.19-1-5}{Proposition~19.1.5}}, the itemized proposition immediately above.}
provided that, in the second assertion of~(ii) and in~(iv), the sheaf $\sh{O}_X$ is coherent and the local rings $\sh{O}_{X,x}$ are Noetherian; as for~(ii), it is necessary to \emph{define} $Y'$ as the ringed subspace defined by the $\sh{O}_{X'}$-ideal $g^*(\sh{J})$; the proofs are unchanged.
\end{enumerate}
\end{remarks}

\begin{env}[19.1.8]
\label{IV.19.1.8}
We have already noted (in particular \sref{IV.16.9.6},~(ii)) that in non-Noetherian rings, the notion of ``regular sequence'' of elements of the ring does not possess the properties that would make it usable.
Recent research seems to show that in this case the notion that should be substituted for that of regular sequence is that of the sequence $\mathbf{f}=(f_1,\ldots,f_n)$ having the property that the complex of the exterior algebra $\mathrm{K}_0(\mathbf{f},\mathrm{M})$ (\sref[III]{III.1.1.3}) has \emph{zero homology in positive degrees} (cf.\ \sref[III]{III.1.1.4}); see in particular A.~Grothendieck, S\'eminaire de g\'eom\'etrie alg\'ebrique, 1966, to appear.
\end{env}

\subsection{Transversally regular immersions}
\label{subsection:IV.19.2}

\oldpage[IV-4]{190}

\begin{definition}[19.2.1]
\label{IV.19.2.1}
Let $u:X\to S$ be a morphism of preschemes and $\sh{F}$ a quasi-coherent $\sh{O}_X$-module.
We say that a sequence $(f_i)_{1\leq i\leq n}$ of sections of $\sh{O}_X$ over $X$ is \emph{transversally $\sh{F}$-regular relative to $S$} (or relative to $u$) if the sequence $(f_i)$ is $\sh{F}$-regular and the $\sh{O}_X$-modules
\[
  \sh{F}\big/\Bigl(\sum_{1\leq j\leq i}f_j\sh{F}\Bigr)
\]
are flat over $S$ for $0\leq i\leq n$.
We say that a quasi-coherent ideal $\sh{J}$ of $\sh{O}_X$ is \emph{transversally regular relative to $S$} (or relative to $u$) at a point $x\in\operatorname{Supp}(\sh{O}_X/\sh{J})$ if there exists an open neighbourhood $U$ of $x$ and a finite sequence $(f_i)$ of sections of $\sh{J}$ over $U$ that is transversally $\sh{O}_U$-regular relative to $S$ and generates $\sh{J}|U$.

In an $A$-algebra $B$, we say that an ideal $\mathfrak{J}$ is \emph{transversally regular relative to $A$} if $\widetilde{\mathfrak{J}}$ is transversally regular relative to $S=\operatorname{Spec}(A)$ at every point of $V(\mathfrak{J})$ in $\operatorname{Spec}(B)$.
\end{definition}

\begin{definition}[19.2.2]
\label{IV.19.2.2}
Let $S$ be a prescheme, $X$ and $Y$ two $S$-preschemes, and $u:Y\to X$ an $S$-morphism which is an immersion.
We say that $u$ is a \emph{transversally regular immersion relative to $S$} at a point $y\in Y$ if there exists a neighbourhood $V$ of $u(y)$ in $X$ such that the sub-prescheme induced on $u(Y)\cap V$ by the sub-prescheme of $X$ associated to $u$ is defined by an ideal of $\sh{O}_V$ transversally regular relative to $S$ at the point $u(y)$.

It is clear that the set of points of $Y$ where an immersion is transversally regular relative to $S$ is \emph{open} in $Y$; if it is equal to $Y$, we say simply that $u$ is a \emph{transversally regular immersion relative to $S$}.

If $S\to T$ is a flat morphism and $\sh{J}$ is an ideal of $\sh{O}_X$ transversally regular relative to $S$ at $x\in X$, then it is also transversally regular relative to $T$ (\sref[0]{0.6.2.1}).
An $S$-immersion $Y\to X$ transversally regular relative to $S$ at a point $y\in Y$ is therefore also a $T$-immersion transversally regular relative to $T$ at this point.
\end{definition}

\begin{proposition}[19.2.3]
\label{IV.19.2.3}
Let $S$ be a prescheme, $X$ and $Y$ two $S$-preschemes, and $u:Y\to X$ an $S$-morphism that is a transversally regular immersion relative to $S$ at a point $y\in Y$.
Then, for every base change $S'\to S$, if we set $X'=X\times_S S'$ and $Y'=Y\times_S S'$, the immersion $u'=u_{(S')}:Y'\to X'$ is transversally regular relative to $S'$ at every point $y'\in Y'$ above $y$.
\end{proposition}

\begin{proof}
This follows immediately from the definition and from \sref[0]{0.15.1.15}.
\end{proof}

\begin{proposition}[19.2.4]
\label{IV.19.2.4}
Let $S$ be a prescheme, $X$ and $Y$ two $S$-preschemes, and $u:Y\to X$ an $S$-morphism that is an immersion.
Suppose either that $S$ and $X$ are locally Noetherian, or that the structure morphisms $g:X\to S$ and $h:Y\to S$ are locally of finite presentation.
Let $y$ be a point of $Y$, $s$ its image in $S$.
The following conditions are equivalent:
\begin{enumerate}[label=\emph{\alph*)}]
  \item $u$ is transversally regular relative to $S$ at the point $y$.
  \item The morphisms $g$ and $h$ are flat in a neighbourhood of the points $u(y)$ and $y$ respectively, and if we set $X_s=g^{-1}(s)$, $Y_s=h^{-1}(s)$, the immersion $u_s:Y_s\to X_s$ is regular at the point $y$.
  \item[\rm b$'$)] The morphisms $g$ and $h$ are flat in a neighbourhood of the points $u(y)$ and $y$ respectively, and, for every change of base $S'\to S$, if we set $X'=X\times_S S'$, $Y'=Y\times_S S'$, the immersion $u':Y'\to X'$ is regular at every point of $Y'$ above $y$.
  \item The morphism $h$ is flat in a neighbourhood of the point $y$, and the immersion $u$ is regular in a neighbourhood of the point $y$.
  \oldpage[IV-4]{191}
  \item[\rm c$'$)] The morphism $h$ is flat in a neighbourhood of the point $y$, and the immersion $u$ is quasi-regular in a neighbourhood of the point $y$.
\end{enumerate}
\end{proposition}

\begin{proof}
We have already proved in \sref{IV.19.2.3} that \emph{a)} implies \emph{b$'$)} without the hypothesis of finiteness, and it is trivial that \emph{b$'$)} implies \emph{b)}.
Let us show that \emph{b)} implies \emph{c)}.
For this, since the question is local on $X$ and on $S$, we may suppose that $Y$ is a closed sub-prescheme of $X$ defined by a quasi-coherent ideal $\sh{J}$ of $\sh{O}_X$, and that $u$ is the canonical injection.
The hypothesis that $\sh{O}_X/\sh{J}$ is $g$-flat implies that the sequence
\[
  0\to\sh{J}\otimes_{\sh{O}_S}\mathbf{k}(s)\to\sh{O}_X\otimes_{\sh{O}_S}\mathbf{k}(s)\to(\sh{O}_X/\sh{J})\otimes_{\sh{O}_S}\mathbf{k}(s)\to 0
\]
is exact (\sref[0]{0.6.1.2}); hence $Y_s$ is the closed sub-prescheme of $X_s$ defined by the ideal $\sh{J}_s$ of $\sh{O}_{X_s}=\sh{O}_X\otimes_{\sh{O}_S}\mathbf{k}(s)$ that is identified with $\sh{J}\otimes_{\sh{O}_S}\mathbf{k}(s)=\sh{J}/\mathfrak{m}_s\sh{J}$.
Consider a finite sequence $(f_i)$ of sections of $\sh{J}$ over a neighbourhood of $y$ in $X$ whose images in
\[
  \sh{J}_y/(\mathfrak{m}_s\sh{O}_{X,y}\sh{J}_y+\sh{J}_y^2)
\]
form a basis of this $\bigl(\sh{O}_{X,y}/(\mathfrak{m}_s\sh{O}_{X,y}+\sh{J}_y)\bigr)$-module, which is free by hypothesis~\emph{b)}.
Since $X_s$ is locally Noetherian, it follows from \sref{IV.16.9.11},
\sref{IV.16.9.5}, and hypothesis~\emph{b)} that the images $(\overline f_i)_y=(f_i)_y\otimes 1$, sections of $\sh{J}_s=\sh{J}\otimes_{\sh{O}_S}\mathbf{k}(s)$ over a neighbourhood of $y$ in $X_s$, form a regular sequence generating the restriction of $\sh{J}_s$ to this neighbourhood.
Since $\mathfrak{m}_s\sh{O}_{X,y}$ is contained in the radical of $\sh{O}_{X,y}$, and since in both cases under consideration $\sh{J}_y$ is an ideal of finite type of $\sh{O}_{X,y}$, Nakayama's lemma shows that the $(f_i)_y$ also generate $\sh{J}_y$; since $\sh{J}$ is an ideal of finite type of $\sh{O}_X$ in both cases, there is a neighbourhood $U$ of $y$ in $X$ such that the $f_i|U$ generate $\sh{J}|U$ (\sref[0]{0.5.2.2}).
The implication \emph{b)}$\Rightarrow$\emph{c)} now follows from \sref[0]{0.15.1.16} when $S$ and $X$ are locally Noetherian, and from \sref{IV.11.3.8} when $g$ and $h$ are locally of finite presentation; in the latter case, \sref{IV.11.3.8} also proves the equivalence of \emph{c)} and \emph{c$'$)}, which is immediate when $X$ and $S$ are locally Noetherian (\sref[0]{0.15.1.11}).
Finally, if \emph{c)} holds, then to prove \emph{a)} we may again reduce to the case where $Y$ is a closed sub-prescheme of $X$ defined by $\sh{J}$. By hypothesis there is a regular sequence $(f_i)$ of sections of $\sh{J}$ over a neighbourhood $U$ of $y$ in $X$ such that the $f_i$ generate $\sh{J}|U$ and $h$ is flat on $U$; to see that \emph{c)} implies \emph{a)}, it again suffices to apply \sref[0]{0.15.1.16} when $S$ and $X$ are locally Noetherian, and \sref{IV.11.3.8} when $g$ and $h$ are locally of finite presentation.
\end{proof}

\begin{remarks}[19.2.5]
\label{IV.19.2.5}
\medskip\noindent
\begin{enumerate}[label=(\roman*)]
  \item If the hypothesis that $h$ is flat in a neighbourhood of $y$ is removed from condition~\emph{b)}, the resulting condition no longer implies \emph{a)}.
It suffices to take for $S$ the spectrum of a local ring that is not a field, for $s$ the closed point of $S$, and $Y=X_s$.
  \item Suppose that $Y$ is a closed sub-prescheme of $X$ defined by a quasi-coherent ideal $\sh{J}$ of $\sh{O}_X$.
The proof of \sref{IV.19.2.4} shows that, when the equivalent conditions of that proposition hold, then for every system $(f_i)$ of sections of $\sh{J}$ over a neighbourhood of $y$ in $X$ whose images in $\sh{J}_y/(\mathfrak{m}_s\sh{O}_{X,y}\sh{J}_y+\sh{J}_y^2)$ form a basis of this $\bigl(\sh{O}_{X,y}/(\mathfrak{m}_s\sh{O}_{X,y}+\sh{J}_y)\bigr)$-module, there is an open neighbourhood $U$ of $y$ in $X$ such that the $f_i|U$ satisfy the conditions of Definition~\sref{IV.19.2.1}.
\end{enumerate}
\end{remarks}

\begin{corollary}[19.2.6]
\label{IV.19.2.6}
\oldpage[IV-4]{192}
Suppose either that $S$ and $X$ are locally Noetherian, or that $g$ and $h$ are locally of finite presentation.
Suppose moreover that the fibres $X_s$ and $Y_s$ are regular at the points $u(y)$ and $y$ respectively, and that the morphisms $g$ and $h$ are flat in a neighbourhood of $u(y)$ and $y$ respectively.
Then the immersion $Y\to X$ is transversally regular at the point $y$.\footnote{The printed text has ``$x$'' here and again in the proof; the point introduced in {\hyperlink{ega4.c04.prop.19-2-4}{Proposition~19.2.4}} is $y$.}
\end{corollary}

\begin{proof}
Indeed, \sref{IV.19.1.1} shows that the immersion $Y_s\to X_s$ is regular at the point $y$, and it suffices to apply \sref{IV.19.2.4},~\emph{b)}.
\end{proof}

\begin{proposition}[19.2.7]
\label{IV.19.2.7}
\medskip\noindent
\begin{enumerate}[label=(\roman*)]
  \item Every open $S$-immersion whose source is flat over $S$ is transversally regular relative to $S$.
  \item Let $f:Y\to X$ be an S-immersion, $g:S'\to S$ a morphism and set $X'=X_{(S')}$, $Y'=Y_{(S')}$, $f'=f_{(S')}:Y'\to X'$.
If $f$ is transversally regular relative to $S$, then $f'$ is transversally regular relative to $S'$. The converse holds if $g$ is faithfully flat and, moreover, $X$ and $Y$ are locally of finite presentation over $S$.
  \item If $f:Y\to X$ and $g:Z\to Y$ are transversally regular immersions relative to $S$, so is $f\circ g:Z\to X$.
  \item Let $X$ be an S-prescheme, $f:Y\to X$, $g:Z\to Y$ two S-immersions.
Suppose either that $S$ and $X$ are locally Noetherian, or that $X$, $Y$, and $Z$ are locally of finite presentation over $S$.
Then, if $g$ and $f\circ g$ are transversally regular at the point $z\in Z$ relative to $S$, $f$ is transversally regular at the point $g(z)$ relative to $S$.
  \item Under the general conditions of~(iv), suppose that $X$ and $Y$ are $S$-flat, $Y_s$ is regular at the point $g(z)$, and $X_s$ is regular at the point $f(g(z))$; then, if $f\circ g$ is transversally regular at the point $z$, so is $g$.
\end{enumerate}
\end{proposition}

\begin{proof}
Assertion~(i) is immediate, and the first assertion of~(ii) follows immediately from \sref{IV.19.2.3}.
To prove the second assertion of~(ii), apply \sref{IV.19.2.4}. Since $X'$ and $Y'$ are locally of finite presentation over $S'$, this criterion first shows that $X'$ and $Y'$ are flat over $S'$; hence $X$ and $Y$ are flat over $S$ (\sref{IV.2.5.1}). On the other hand, for every $s\in S$, if $s'\in S'$ lies above $s$, then $Y'_{s'}\to X'_{s'}$ is a regular immersion by hypothesis; \sref{IV.19.1.5},~(ii), therefore shows that the same is true of $Y_s\to X_s$, since these preschemes are locally Noetherian and, in that case, regularity and quasi-regularity of the immersion $Y_s\to X_s$ are equivalent.
To prove~(iii), reduce to the case where $Z$ is a closed sub-prescheme of $Y$ and $Y$ a closed sub-prescheme of $X$, and form, as in \sref{IV.19.1.5},~(iii), a system of generators of the ideal of $\sh{O}_X$ defining $Z$.
For~(iv), the hypothesis together with \sref{IV.19.2.4} first shows that $X$, $Y$, and $Z$ are $S$-flat in a neighbourhood of $z$. If $s$ is the image of $z$ in $S$, it remains only to prove that $f_s:Y_s\to X_s$ is regular at the point $g(z)$, knowing that $g_s:Z_s\to Y_s$ and $f_s\circ g_s$ are regular at $z$; this follows from \sref{IV.19.1.5},~(iv).
Finally, for~(v), argue in the same way as for~(iv). The hypothesis on $f\circ g$ already implies by \sref{IV.19.2.4},~\emph{b)}, that $Z$ is $S$-flat in a neighbourhood of $z$, and the same is true of $X$ and $Y$ by hypothesis. We are therefore reduced to proving that $g_s$ is regular at $z$, knowing that $f_s\circ g_s$ is regular there; this follows from \sref{IV.19.1.5},~(v), under the stated hypotheses on $X_s$ and $Y_s$.
\end{proof}

\oldpage[IV-4]{193}

\begin{proposition}[19.2.8]
\label{IV.19.2.8}
Let $X$ be an $S$-prescheme, let $Y$ be a closed sub-prescheme of $X$ defined by a quasi-coherent ideal $\sh{J}$ of $\sh{O}_X$, and let $u:Y\to X$ be the canonical injection.
If $u$ is transversally regular relative to $S$, the $\sh{O}_X$-modules $\sh{J}^n$ and $\sh{J}^n/\sh{J}^m$ ($0\leq n<m$) are $S$-flat at the points of $Y$.
For every base change $S'\to S$, if $X'=X_{(S')}$, $Y'=Y_{(S')}$, and $u'=u_{(S')}$, then, up to canonical isomorphism,
\[
  \sh{G}r_n(u')=\sh{G}r_n(u)\otimes_{\sh{O}_S}\sh{O}_{S'}
\]
for every $n>0$.
\end{proposition}

\begin{proof}
Indeed, $\sh{O}_X$ and $\sh{O}_X/\sh{J}$ are, by hypothesis, $S$-flat at the points of $Y$, while $\sh{J}^n/\sh{J}^{n+1}$ is a locally free $(\sh{O}_X/\sh{J})$-module (\sref{IV.16.9.3}), and hence is $S$-flat at those points.
The first assertion then follows from \sref[0]{0.6.1.2}, and the second from \sref{IV.11.2.9.2}.
\end{proof}

The following proposition generalizes and sharpens \sref{IV.17.16.1}.

\begin{proposition}[19.2.9]
\label{IV.19.2.9}
Let $f:X\to S$ be a morphism, $x$ a point of $X$, $s=f(x)$.
Suppose that $f$ is of finite presentation at $x$ and is a Cohen--Macaulay morphism at $x$ (\sref{IV.6.8.1}).
Then there exists a sub-prescheme $Y$ of $X$, containing $x$, with the following properties:
\begin{enumerate}[label=(\roman*)]
  \item The canonical injection $j:Y\to X$ is transversally regular relative to $S$.
  \item The morphism $g=f\circ j:Y\to S$ is Cohen--Macaulay; equivalently, in view of~(i), all the fibres of $g$ are Cohen--Macaulay preschemes.
  \item The point $x$ is a maximal point of $Y_s=g^{-1}(s)$.
\end{enumerate}
\end{proposition}

\begin{proof}
Set $X_s=f^{-1}(s)$. By hypothesis, the ring $\sh{O}_{X_s,x}$ is Cohen--Macaulay; let $(t_i)_{1\leq i\leq n}$ be a system of parameters of this ring, and hence a regular sequence (\sref[0]{0.16.5.7}).
There is an open neighbourhood $U$ of $x$ in $X$ and sections $h_i$ ($1\leq i\leq n$) of $\sh{O}_X$ over $U$ such that the $t_i$ are the canonical images of the germs $(h_i)_x$ in the quotient ring $\sh{O}_{X_s,x}$ of $\sh{O}_{X,x}$.
Let $\sh{J}$ be the ideal of $\sh{O}_U$ generated by the $h_i$, and let $Y$ be the closed sub-prescheme of $U$ defined by $\sh{J}$.
We may suppose that $f|U$ is locally of finite presentation; by the definition of $Y$, the same is then true of $g$.
The equivalence of \emph{c)} and \emph{a)} in \sref{IV.11.3.8} shows that the sequence $((h_i)_x)$ is regular in $\sh{O}_{X,x}$ and that $\sh{O}_{Y,x}$ is flat over $\sh{O}_{S,s}$. By \sref{IV.11.3.1}, $f$ and $g$ are therefore flat in a neighbourhood of $x$, and \sref{IV.19.2.4} shows that $j$ is a transversally regular immersion relative to $S$, after replacing $Y$, if necessary, by $Y\cap V$ for an open neighbourhood $V$ of $x$ in $X$.
Thus~(i) is proved, and we may suppose that $g:Y\to S$ is flat.
Since $\sh{O}_{Y_s,x}$, the quotient of $\sh{O}_{X_s,x}$ by the ideal generated by the $t_i$, is Artinian by construction (\sref[0]{0.16.3.6}), $x$ is a maximal point of $Y_s$.
Finally, every Artinian ring is Cohen--Macaulay (\sref[0]{0.16.5.1}); replacing $Y$, if necessary, by its intersection with an open neighbourhood of $x$ in $X$, we may therefore suppose, by \sref{IV.12.1.7},~(iii), that $g$ is a Cohen--Macaulay morphism.

In particular, we recover the existence of flat ``quasi-sections'' $Y$ of a morphism $f:X\to S$ at a point $x\in X_s$ that is closed in $X_s$, under the hypotheses that $f$ is flat at $x$ and $\sh{O}_{X_s,x}$ is Cohen--Macaulay (\sref{IV.17.16.1}); moreover, the immersion $Y\to X$ is transversally regular relative to $S$.
This last refinement also applies to the \emph{\'etale} quasi-section defined in \sref{IV.17.16.3},~(i).
\end{proof}

\subsection{Relative complete intersections (flat case)}
\label{subsection:IV.19.3}

\oldpage[IV-4]{194}

\begin{definition}[19.3.1]
\label{IV.19.3.1}
Let $\mathrm{A}$ be a Noetherian local ring.
We say that $\mathrm{A}$ is a \emph{complete-intersection ring} (or an \emph{absolute complete-intersection ring}, when there is danger of confusion) if its completion $\widehat{\mathrm{A}}$ is isomorphic to the quotient of a complete regular Noetherian local ring $\mathrm{B}$ by a regular ideal, that is, by an ideal generated by a regular sequence of elements of $\mathrm{B}$ (\sref{IV.16.9.7}).

Every regular local ring is a complete intersection ring.

We say that a locally Noetherian prescheme $X$ is an \emph{(absolute) complete intersection} at a point $x\in X$ if the ring $\sh{O}_{X,x}$ is a complete-intersection ring.
\end{definition}

\begin{proposition}[19.3.2]
\label{IV.19.3.2}
Let $\mathrm{B}$ be a regular Noetherian local ring and $\mathfrak{J}$ an ideal of $\mathrm{B}$.
For $\mathrm{A}=\mathrm{B}/\mathfrak{J}$ to be a complete-intersection ring, it is necessary and sufficient that $\mathfrak{J}$ be a regular ideal, that is, be generated by a regular sequence of elements of $\mathrm{B}$ (\sref{IV.16.9.7}).
\end{proposition}

\begin{proof}
We may write $\widehat{\mathrm{A}}=\widehat{\mathrm{B}}/\mathfrak{J}\widehat{\mathrm{B}}$. Since $\widehat{\mathrm{B}}$ is faithfully flat over $\mathrm{B}$ and is Noetherian, $\mathfrak{J}\widehat{\mathrm{B}}$ is regular if and only if $\mathfrak{J}$ is regular (\sref{IV.19.1.5},~(ii)).
This shows, by \sref[0]{0.17.1.5}, both that the condition in the statement is sufficient and that, in proving its necessity, we may reduce to the case where $\mathrm{A}$ and $\mathrm{B}$ are complete.
By hypothesis, there is then a complete regular Noetherian local ring $\mathrm{B}'$ such that $\mathrm{A}$ is isomorphic to a quotient $\mathrm{B}'/\mathfrak{J}'$, where $\mathfrak{J}'$ is a regular ideal of $\mathrm{B}'$.
Consider the fibre product $\mathrm{B}''=\mathrm{B}\times_{\mathrm{A}}\mathrm{B}'$ for the surjective homomorphisms $\mathrm{B}\to\mathrm{A}$ and $\mathrm{B}'\to\mathrm{A}$ (\sref[0]{0.18.1.2}). We first prove the following lemma.

\begin{lemma}[19.3.2.1]
\label{IV.19.3.2.1}
Let $\mathrm{A}$, $\mathrm{E}$, $\mathrm{F}$ be three rings, $f:\mathrm{E}\to\mathrm{A}$, $g:\mathrm{F}\to\mathrm{A}$ two homomorphisms, and set $\mathrm{G}=\mathrm{E}\times_{\mathrm{A}}\mathrm{F}$ (\sref[0]{0.18.1.2}).
\begin{enumerate}[label=(\roman*)]
  \item If $f$ is surjective and if $\mathrm{E}$ and $\mathrm{F}$ are Noetherian rings, $\mathrm{G}$ is Noetherian.
  \item If $\mathrm{A}$, $\mathrm{E}$, and $\mathrm{F}$ are local rings and $f$ and $g$ are local homomorphisms, then $\mathrm{G}$ is a local ring and the canonical homomorphisms $\mathrm{G}\to\mathrm{E}$ and $\mathrm{G}\to\mathrm{F}$ are local.
  \item If the conditions of~(i) and~(ii) are satisfied, if $g$ is surjective, if $\mathrm{E}$ and $\mathrm{F}$ are complete Noetherian local rings, $\mathrm{G}$ is complete.
\end{enumerate}
\end{lemma}

\begin{proof}
\begin{enumerate}[label=(\roman*)]
  \item In view of \sref[0]{0.18.1.3} and \sref{0.18.1.5}, the projection $\mathrm{G}\to\mathrm{F}$ is surjective. Its kernel $\mathfrak{J}'$ is canonically identified with the kernel $\mathfrak{J}$ of $f$, and the $\mathrm{G}$-module structure on $\mathfrak{J}'$ corresponds to the $\mathrm{E}$-module structure on $\mathfrak{J}$ via the canonical homomorphism $\mathrm{G}\to\mathrm{E}$.
If $\mathfrak{a}$ is any ideal of $\mathrm{G}$, then $\mathfrak{a}/(\mathfrak{a}\cap\mathfrak{J}')$ is isomorphic to $(\mathfrak{a}+\mathfrak{J}')/\mathfrak{J}'$, and hence to an ideal of $\mathrm{F}$, so it is finitely generated. On the other hand, $\mathfrak{a}\cap\mathfrak{J}'$ is isomorphic to an ideal of $\mathrm{E}$ contained in $\mathfrak{J}$, and is therefore also finitely generated. Thus $\mathfrak{a}$ is finitely generated and $\mathrm{G}$ is Noetherian.
  \item Let $\mathfrak{r}$ and $\mathfrak{s}$ be the maximal ideals of $\mathrm{E}$ and $\mathrm{F}$ respectively. The set $\mathfrak{m}=\mathfrak{r}\times_{\mathrm{A}}\mathfrak{s}$ is an ideal of $\mathrm{G}$, the inverse image of $\mathfrak{r}$ under $\mathrm{G}\to\mathrm{E}$ and of $\mathfrak{s}$ under $\mathrm{G}\to\mathrm{F}$. Its elements are therefore non-invertible in $\mathrm{G}$.
On the other hand, if $(x,y)\notin\mathfrak{m}$ and, for example, $x\notin\mathfrak{r}$, then $f(x)$ does not belong to the maximal ideal of $\mathrm{A}$. Since $g(y)=f(x)$ and $g$ is local, we also have $y\notin\mathfrak{s}$; we
\oldpage[IV-4]{195}
conclude that $x$ and $y$ are invertible, hence $(x,y)$ is invertible, which completes the proof of~(ii).
  \item The image of $\mathfrak{m}$ under the surjective homomorphism $\mathrm{G}\to\mathrm{E}$ is the maximal ideal $\mathfrak{r}$ of $\mathrm{E}$; hence the image of $\mathfrak{m}^n\cap\mathfrak{J}'$ is $\mathfrak{r}^n\cap\mathfrak{J}$. Since $\mathfrak{J}$ is closed, and therefore complete, for the topology induced by the $\mathfrak{r}$-adic topology of $\mathrm{E}$ (\sref[0]{0.7.3.5}), $\mathfrak{J}'$ is complete for the topology induced by the $\mathfrak{m}$-preadic topology of $\mathrm{G}$.
On the other hand, $\mathrm{G}/\mathfrak{J}'$, endowed with the $\mathfrak{m}$-preadic topology, is isomorphic to $\mathrm{F}$ endowed with the $\mathfrak{s}$-adic topology, and is therefore complete for the quotient topology of the $\mathfrak{m}$-preadic topology of $\mathrm{G}$.
It follows that $\mathrm{G}$ is complete for the $\mathfrak{m}$-preadic topology.
\end{enumerate}
\end{proof}

The lemma established, Cohen's theorem (\sref[0]{0.19.8.8}) shows that $\mathrm{B}''$ is isomorphic to a quotient of a complete regular Noetherian local ring $\mathrm{C}$.
It follows from \sref{IV.19.1.2} that the immersion $\operatorname{Spec}(\mathrm{B}')\to\operatorname{Spec}(\mathrm{C})$ is regular. Since, by hypothesis, the immersion $\operatorname{Spec}(\mathrm{A})\to\operatorname{Spec}(\mathrm{B}')$ is regular, \sref{IV.19.1.5},~(iii), shows that $\operatorname{Spec}(\mathrm{A})\to\operatorname{Spec}(\mathrm{C})$ is regular. But this immersion is also the composite
\[
  \operatorname{Spec}(\mathrm{A})\longrightarrow\operatorname{Spec}(\mathrm{B})\longrightarrow\operatorname{Spec}(\mathrm{C}).
\]
Since $\mathrm{B}$ is regular, the immersion $\operatorname{Spec}(\mathrm{B})\to\operatorname{Spec}(\mathrm{C})$ is regular (\sref{IV.19.1.2}); hence $\operatorname{Spec}(\mathrm{A})\to\operatorname{Spec}(\mathrm{B})$ is regular by \sref{IV.19.1.5},~(iv).
\end{proof}

\begin{corollary}[19.3.3]
\label{IV.19.3.3}
Let $X$ be a locally Noetherian prescheme locally immersible in a regular prescheme.
Then the set of $x\in X$ where $X$ is a complete intersection is open in $X$ and contains the set of regular points of $X$.
\end{corollary}

\begin{proof}
We can indeed restrict ourselves to the case where $X=\operatorname{Spec}(\mathrm{B}/\mathfrak{J})$, where $\mathrm{B}$ is a regular ring and $\mathfrak{J}$ an ideal of $\mathrm{B}$.
The set of $\mathfrak{p}\supset\mathfrak{J}$ in $\operatorname{Spec}(\mathrm{B})$ such that $\mathfrak{J}_\mathfrak{p}$ is regular in $\mathrm{B}_\mathfrak{p}$ is open (\sref{IV.16.9.5}), and, by \sref{IV.19.3.2}, it is the set of points of $X$ at which $X$ is a complete intersection.
\end{proof}

\begin{corollary}[19.3.4]
\label{IV.19.3.4}
Let $k$ be a field, $X$ a prescheme locally of finite type over $k$, $k'$ an extension of $k$, $X'=X\otimes_k k'$.
Let $x'$ be a point of $X'$, $x$ its projection in $X$.
For $X$ to be a complete intersection at $x$, it is necessary and sufficient that $X'$ be a complete intersection at $x'$.
\end{corollary}

\begin{proof}
The question being local on $X$, we can suppose that $X=\operatorname{Spec}(\mathrm{A})$, where $\mathrm{A}$ is a $k$-algebra of finite type, hence a quotient of a polynomial algebra $k[\mathrm{T}_1,\ldots,\mathrm{T}_n]$, so $X$ is a closed sub-prescheme of the regular scheme $Y=\operatorname{Spec}(k[\mathrm{T}_1,\ldots,\mathrm{T}_n])$ (\sref[0]{0.17.3.7}).
If we set $Y'=Y\otimes_k k'=\operatorname{Spec}(k'[\mathrm{T}_1,\ldots,\mathrm{T}_n])$, $Y'$ is also a regular scheme.
Now, since $\operatorname{Spec}(k')$ is faithfully flat over $\operatorname{Spec}(k)$, it follows from \sref{IV.19.1.5},~(ii), and the fact that these preschemes are Noetherian that the immersion $X\to Y$ is regular at $x$ if and only if the immersion $X'\to Y'$ is regular at $x'$.
The conclusion then follows from the criterion \sref{IV.19.3.2}.
\end{proof}

\begin{corollary}[19.3.5]
\label{IV.19.3.5}
Let $\mathrm{A}$ and $\mathrm{C}$ be Noetherian local rings, let $\rho:\mathrm{A}\to\mathrm{C}$ be a local homomorphism, let $\mathrm{B}=\mathrm{C}/\mathfrak{J}$ be a quotient ring of $\mathrm{C}$, and let $k$ be the residue field of $\mathrm{A}$.
Suppose that $\rho$ makes $\mathrm{C}$ a flat $\mathrm{A}$-module and that $\mathrm{C}\otimes_{\mathrm{A}}k$ is regular.
Then the following conditions are equivalent:
\begin{enumerate}[label=\emph{\alph*)}]
  \item The ideal $\mathfrak{J}$ is transversally regular relatively to $\mathrm{A}$.
  \item $\mathrm{B}$ is a flat $\mathrm{A}$-module, and $\mathrm{B}\otimes_{\mathrm{A}} k$ is a complete intersection ring.
\end{enumerate}
\end{corollary}

\begin{proof}
\oldpage[IV-4]{196}
By virtue of \sref{IV.19.2.4}, condition \emph{a)} amounts to saying that $\mathrm{B}$ is a flat $\mathrm{A}$-module and that $\mathfrak{J}\otimes_{\mathrm{A}} k$ (which is identified with an ideal of $\mathrm{C}\otimes_{\mathrm{A}} k$ by virtue of the flatness of $\mathrm{B}$ over $\mathrm{A}$ (\sref[0]{0.6.1.2})) is regular in this ring.
Since the ring $\mathrm{C}\otimes_{\mathrm{A}} k$ is regular, it amounts to the same, by virtue of \sref{IV.19.3.2}, to say (when $\mathrm{B}$ is a flat $\mathrm{A}$-module) that $\mathfrak{J}\otimes_{\mathrm{A}} k$ is a regular ideal of $\mathrm{C}\otimes_{\mathrm{A}} k$, or that $\mathrm{B}\otimes_{\mathrm{A}} k=(\mathrm{C}\otimes_{\mathrm{A}} k)/(\mathfrak{J}\otimes_{\mathrm{A}} k)$ is a complete intersection ring; whence the corollary.
\end{proof}

\begin{definition}[19.3.6]
\label{IV.19.3.6}
Let $f:X\to S$ be a flat morphism locally of finite presentation.
We say that $X$ is a \emph{complete intersection relatively to $S$} at a point $x$, if the fibre $f^{-1}(f(x))$ is a complete intersection (absolute) at the point $x$.
We say that $X$ is a \emph{complete intersection relatively to $S$}, or that the morphism $f$ is a \emph{morphism of complete intersection}, if $X$ is a complete intersection relatively to $S$ at each of its points.
\end{definition}

\begin{proposition}[19.3.7]
\label{IV.19.3.7}
Let $g:X\to S$, $h:Y\to S$ be two flat morphisms locally of finite presentation, $f:Y\to X$ an S-immersion, $y$ a point of $Y$, $x=f(y)$, $s=g(x)=h(y)$.
Suppose that the fibre $X_s=g^{-1}(s)$ is regular at the point $x$.
Then, for $f$ to be a transversally regular immersion at the point $y$, it is necessary and sufficient that $Y$ be a complete intersection relatively to $S$ at the point $y$.
\end{proposition}

\begin{proof}
To say that $f$ is transversally regular at the point $y$ amounts, by virtue of \sref{IV.19.2.4}, to saying that $f_s:Y_s\to X_s$ is a regular immersion at the point $y$.
But since $X_s$ is regular at the point $x=f(y)$, this is equivalent, by virtue of \sref{IV.19.3.2}, to saying that $Y_s$ is a complete intersection (absolute) at the point $y$, whence the proposition.
\end{proof}

\begin{corollary}[19.3.8]
\label{IV.19.3.8}
Let $f:X\to S$ be a flat morphism locally of finite presentation.
The set $U$ of $x\in X$ for which $X$ is a complete intersection relatively to $S$ is open in $X$.
If moreover $f$ is proper, the set $E$ of $s\in S$ such that $X_s=f^{-1}(s)$ is a complete intersection at each of its points is an open subset of $S$.
\end{corollary}

\begin{proof}
Since $E=S-f(X-U)$,\footnote{The printed text has $Y-f(X-U)$; the ambient set in which $E$ was just defined is $S$.} the second assertion follows from the first and from the fact that $f$ is a closed map.
The first assertion is local on $X$, hence we can suppose that $X$ is a closed sub-prescheme of a prescheme of the form $Z=S[\mathrm{T}_1,\ldots,\mathrm{T}_n]$ ($\mathrm{T}_i$ indeterminate).
Since $Z$ is smooth over $S$ (\sref{IV.17.3.8}), its fibres $Z_s$ are regular, and, by \sref{IV.19.3.7}, the assertion that $x\in U$ is equivalent to the assertion that the immersion $X\to Z$ is transversally regular at $x$; the result follows from \sref{IV.19.2.2}.
\end{proof}

\begin{proposition}[19.3.9]
\label{IV.19.3.9}
\medskip\noindent
\begin{enumerate}[label=(\roman*)]
  \item Every open immersion is a complete-intersection morphism.
  \item Let $f:X\to S$ be a flat morphism locally of finite presentation that is a complete-intersection morphism.
For every change of base $g:S'\to S$, $f'=f_{(S')}:X_{(S')}\to S'$ is a morphism of complete intersection.
The converse is true if $g$ is faithfully flat and quasi-compact.
  \item If $f:X\to Y$ and $g:Y\to Z$ are flat morphisms locally of finite presentation which are complete intersection morphisms, so is $g\circ f:X\to Z$.
\end{enumerate}
\end{proposition}

\begin{proof}
\oldpage[IV-4]{197}
Assertion~(i) follows immediately from the definition \sref{IV.19.3.1}.
To prove~(ii), note first that if $f$ is flat and locally of finite presentation, then so is $f'$ (\sref{IV.2.1.4}); conversely, these two properties descend when $g$ is faithfully flat and quasi-compact (\sref{IV.2.5.1} and \sref{IV.2.7.1}). Moreover, for every $s'\in S'$, if $s=g(s')$, the fibre $f^{-1}(s)$ is a complete intersection if and only if $f'^{-1}(s')$ is one (\sref{IV.19.3.4}). This proves~(ii), since a faithfully flat morphism is surjective.

Finally, under the hypotheses of~(iii), $g\circ f$ is flat and locally of finite presentation.
By definition \sref{IV.19.3.6}, we are reduced to the case where $Z$ is the spectrum of a field, and must then prove that the local rings $\sh{O}_{X,x}$ are complete-intersection rings. This is contained in the following more general result.
\end{proof}

\begin{corollary}[19.3.10]
\label{IV.19.3.10}
Let $f:X\to S$ be a flat morphism locally of finite presentation that is a complete-intersection morphism, let $x$ be a point of $X$, and set $s=f(x)$.
If $\sh{O}_{S,s}$ is a Noetherian complete-intersection ring, then so is $\sh{O}_{X,x}$.
\end{corollary}

\begin{proof}
We may plainly reduce to the case where $S=\operatorname{Spec}(\mathrm{A})$ with $\mathrm{A}=\sh{O}_{S,s}$. We first show that we may further reduce to the case where $\mathrm{A}$ is complete.
Set $\mathrm{A}'=\widehat{\mathrm{A}}$, $X'=X\otimes_{\mathrm{A}}\mathrm{A}'$, and let $s'$ be the unique closed point of $S'=\operatorname{Spec}(\mathrm{A}')$, hence the unique point of $S'$ above $s$. If $f'=f_{(S')}:X'\to S'$, then $f'^{-1}(s')$ is canonically isomorphic to $f^{-1}(s)$ because $\mathrm{A}$ and $\mathrm{A}'$ have the same residue field (\sref[I]{I.3.6.4}). Thus there is a unique point $x'\in X'$ above $x$ and $s'$, and $\sh{O}_{X',x'}=\sh{O}_{X,x}\otimes_{\mathrm{A}}\mathrm{A}'$.
The local rings $\sh{O}_{X,x}$ and $\sh{O}_{X',x'}$ have isomorphic completions. Indeed, if $\mathrm{E}$ is a local ring that is a local $\mathrm{A}$-algebra, then the separated completion $(\mathrm{E}\otimes_{\mathrm{A}}\widehat{\mathrm{A}})^\wedge$ of $\mathrm{E}\otimes_{\mathrm{A}}\widehat{\mathrm{A}}$, endowed with the tensor-product topology, is isomorphic to the separated completion of $\widehat{\mathrm{E}}\widehat{\otimes}_{\mathrm{A}}\widehat{\mathrm{A}}\cong\widehat{\mathrm{E}}$, and hence to $\widehat{\mathrm{E}}$ itself (\sref[0]{0.7.7.1}). By \sref{IV.19.3.1}, it is therefore equivalent for $\sh{O}_{X,x}$ or $\sh{O}_{X',x'}$ to be a complete-intersection ring.

Suppose, then, that $\mathrm{A}$ is complete. We may also suppose that $X=\operatorname{Spec}(\mathrm{B})$, where $\mathrm{B}$ is a quotient of the polynomial algebra $\mathrm{C}=\mathrm{A}[\mathrm{T}_1,\ldots,\mathrm{T}_r]$. Since polynomial rings over a field are regular (\sref[0]{0.17.3.7}), \sref{IV.19.3.7} shows that the immersion $X\to Z=\operatorname{Spec}(\mathrm{C})$ may be taken transversally regular relative to $S$. Thus $\mathrm{B}=\mathrm{C}/\mathfrak{J}$, where $\mathfrak{J}$ is generated by a regular sequence $(g_i)_{1\leq i\leq n}$ in $\mathrm{C}$.
On the other hand, since $\mathrm{A}$ is complete, Cohen's theorem allows us to write $\mathrm{A}=\mathrm{A}'/\mathfrak{R}$, where $\mathrm{A}'$ is a regular local ring and $\mathfrak{R}$ an ideal of $\mathrm{A}'$ (\sref[0]{0.19.8.8}). Since $\mathrm{A}$ is a complete-intersection ring, $\mathfrak{R}$ is generated by a regular sequence $(f_j)_{1\leq j\leq m}$ (\sref{IV.19.3.2}).
In $\mathrm{C}'=\mathrm{A}'[\mathrm{T}_1,\ldots,\mathrm{T}_r]$, the elements $f_j$ still form a regular sequence, generating $\mathfrak{R}'=\mathfrak{R}[\mathrm{T}_1,\ldots,\mathrm{T}_r]$ (\sref[0]{0.15.1.4}). Since $\mathrm{C}'/\mathfrak{R}'=\mathrm{C}$, choose, for each $i$, an element $g_i'\in\mathrm{C}'$ whose image in $\mathrm{C}$ is $g_i$. The sequence consisting of the $f_j$ and the $g_i'$ is regular in $\mathrm{C}'$. If $\mathfrak{J}'$ is the ideal it generates, then $\mathrm{B}=\mathrm{C}'/\mathfrak{J}'$; the conclusion follows from \sref{IV.19.3.2}, since $\mathrm{C}'$ is regular (\sref[0]{0.17.3.7}).
\end{proof}

\subsection{Application: criteria of regularity and smoothness for blown-up preschemes}
\label{subsection:IV.19.4}

\oldpage[IV-4]{198}

\begin{env}[19.4.1]
\label{IV.19.4.1}
Let $X$ be a prescheme and $Y$ a closed sub-prescheme of $X$, defined by a quasi-coherent ideal $\sh{J}$ of $\sh{O}_X$.
Recall (\sref[II]{II.8.1.3}) that the $X$-scheme $X'$ obtained by blowing up $Y$ is the prescheme
\[
  X'=\operatorname{Proj}(\sh{S}), \quad \text{where}\quad \sh{S}=\bigoplus_{n\geq 0}\sh{J}^n.
\]
If $\sh{J}$ is of finite type, the structural morphism $f:X'\to X$ is projective.
Without hypothesis on $\sh{J}$, the closed sub-prescheme
\[
  Y'=f^{-1}(Y)
\]
of $X'$ is defined by the quasi-coherent Ideal $\sh{J}\sh{O}_{X'}$ of $\sh{O}_{X'}$, canonically isomorphic to $\sh{O}_{X'}(1)$; more precisely, we have an exact sequence
\[
  0\to\sh{O}_{X'}(1)\xrightarrow{\widetilde{u}_0}\sh{O}_{X'}\to\sh{O}_{Y'}\to 0
\]
{\hyperlink{ega4.c04.subsec.19-4-1-1}{(\textbf{19.4.1.1})}}
\label{IV.19.4.1.1}
where $\sh{J}\sh{O}_{X'}$ is the image of $\widetilde{u}_0$ (\sref[II]{II.8.1.7} and \sref{II.8.1.8}).
Since $\sh{O}_{X'}(1)$ is free of rank $1$ in a neighbourhood of every point of $X'$, $\sh{J}\sh{O}_{X'}$ is locally generated by a non-zero-divisor of $\sh{O}_{X'}$, and $\sh{J}\sh{O}_{X'}/\sh{J}^2\sh{O}_{X'}$ is consequently a free $(\sh{O}_{X'}/\sh{J}\sh{O}_{X'})$-module of rank $1$.
This proves the first assertion of the following lemma:
\end{env}

\begin{lemma}[19.4.2]
\label{IV.19.4.2}
The canonical immersion $Y'\to X'$ is regular and of codimension $1$, and $Y'$ is canonically isomorphic to the Y-scheme $\operatorname{Proj}(\operatorname{gr}_\sh{J}^*(\sh{O}_X))$.
\end{lemma}

\begin{proof}
The second assertion follows from \sref[II]{II.3.5.3}, applied to the canonical injection $Y\to X$, because
\[
  \sh{S}\otimes_{\sh{O}_X}\sh{O}_Y=\operatorname{gr}_{\sh{J}}^*(\sh{O}_X)
\]
by definition: $\sh{O}_Y\cong\sh{O}_X/\sh{J}$ and $\sh{J}^n\otimes_{\sh{O}_X}(\sh{O}_X/\sh{J})\cong\sh{J}^n/\sh{J}^{n+1}$.
\end{proof}

\begin{corollary}[19.4.3]
\label{IV.19.4.3}
If the canonical immersion $Y\to X$ is quasi-regular, the restriction $g:Y'\to Y$ of the morphism $f$ is smooth.
\end{corollary}

\begin{proof}
Indeed, $\sh{N}_{Y/X}=\sh{J}/\sh{J}^2$ is then an $\sh{O}_Y$-Module locally free of finite rank and $\operatorname{gr}_\sh{J}^*(\sh{O}_X)$ is isomorphic to $\mathbf{S}_{\sh{O}_Y}^*(\sh{N}_{Y/X})$ \sref{IV.16.9.8}, hence $Y'$ is Y-isomorphic to $\mathbf{P}(\sh{N}_{Y/X})$ and is therefore smooth over $Y$ \sref{IV.17.3.9}.
\end{proof}

We shall see later (chap.~V) that $g$ is smooth also when $Y\to X$ ``has only ordinary singularities''.

\begin{proposition}[19.4.4]
\label{IV.19.4.4}
With the notations of \sref{IV.19.4.1}, suppose that $X$ is locally Noetherian and that the morphism $g:Y'\to Y$, restriction of $f$, is smooth.
Let $x'$ be a point of $Y'$, $x=g(x')$.
For $Y$ to be regular at the point $x$, it is necessary and sufficient that $Y'$ be regular at the point $x'$; when this is so, $X'$ is then regular at the point $x'$.
\end{proposition}

\begin{proof}
The first assertion follows from \sref{IV.17.5.8} and the second from \sref{IV.19.1.1} and \sref{IV.19.4.2}.
\end{proof}

\begin{corollary}[19.4.5]
\label{IV.19.4.5}
Under the general hypotheses of \sref{IV.19.4.4}, suppose that $Y$ is regular at the point $x$.
Then, for every generization $x_1$ of $x$ in $X$ not belonging to $Y$, $X$ is regular at the point $x_1$.
If $\operatorname{Reg}(X)$ is open ($(\sref{IV.6.12})$, for example if $X$ is excellent (\sref{IV.7.8.6},~(iii)))
\oldpage[IV-4]{199}
and if $Y$ is regular, then there exists an open neighbourhood $U$ of $Y$ in $X$ such that $X$ is regular at the points of $U-Y$.
\end{corollary}

\begin{proof}
To prove the first assertion, we may reduce to the case where $X$ is a local scheme with closed point $x$. The hypothesis then implies that $Y$ is regular (\sref[0]{0.17.3.2}), and hence so is $Y'$ by \sref{IV.19.4.4}.
The morphism $f:X'\to X$ is projective, hence proper. If $x_1'$ denotes the unique point of $X'$ above $x_1$ (\sref[II]{II.8.1.3}), the valuative specialization property (\sref[II]{II.7.2.1}) gives a specialization of $x_1'$ lying in $Y'$. Applying \sref{IV.19.4.4} and \sref[0]{0.17.3.2}, we conclude that $X'$ is regular at $x_1'$, and therefore that $X$ is regular at $x_1$, since $X'-Y'$ is isomorphic to $X-Y$ (\sref[II]{II.8.1.3}).
If $Y$ is regular, every point of $X-Y$ which is a generization of a point of $Y$ belongs to $\operatorname{Reg}(X)$; if $\operatorname{Reg}(X)$ is open, $U=Y\cup\operatorname{Reg}(X)$ is then locally constructible and stable under generization, hence open (\sref[0]{0.9.2.5}), and responds to the question.
\end{proof}

\begin{proposition}[19.4.6]
\label{IV.19.4.6}
Let $g:X\to S$ be a morphism, $Y$ a closed sub-prescheme of $X$ defined by a quasi-coherent Ideal $\sh{J}$, $X'$ the X-scheme obtained by blowing up $Y$, $Y'$ the inverse image of $Y$ in $X'$.
\begin{enumerate}[label=(\roman*)]
  \item Conditions \emph{a)} and \emph{b)} below are equivalent. They imply condition \emph{c)}; if $X$ is S-flat, condition \emph{c)} is also equivalent to them:
  \begin{enumerate}[label=\emph{\alph*)}]
    \item $\operatorname{gr}_\sh{J}^*(\sh{O}_X)$ is a $g$-flat $\sh{O}_X$-algebra.
    \item For every $n\geq 0$, the $\sh{O}_X$-Module $\sh{O}_X/\sh{J}^{n+1}$ is $g$-flat (in other words, the $n$-th infinitesimal neighbourhood $Y^{(n)}$ of $Y$ in $X$ \sref{IV.16.1.2} is S-flat).
    \item For every base change $S_1\to S$ and every integer $n\geq 1$, if $X_1=X\times_S S_1$, the canonical homomorphism $\sh{J}^n\otimes_{\sh{O}_S}\sh{O}_{S_1}\to\sh{J}^n\sh{O}_{X_1}$ is bijective.
  \end{enumerate}
When condition \emph{a)}, or equivalently \emph{b)}, holds, $Y'$ is S-flat, and for every change of base $S_1\to S$, if $Y_1$ is the inverse image of $Y$ in $X_1$, the prescheme $X_1'$ obtained by blowing up $Y_1$ in $X_1$ is canonically isomorphic to $X'\times_S S_1$.
  \item Suppose that conditions \emph{a)} and \emph{b)} of~(i) hold and moreover that the morphisms $X\to S$ and $Y\to S$ are locally of finite presentation.
Then $\sh{S}=\bigoplus_{n\geq 0}\sh{J}^n$ is a $\sh{O}_X$-Algebra of finite presentation, the morphism $X'\to X$ is of finite presentation, the canonical immersion $Y'\to X'$ is transversally regular of codimension $1$ relatively to $S$, and the set of points of $X$ (resp.\ $X'$) where $X$ (resp.\ $X'$) is S-flat, is a neighbourhood of $Y$ (resp.\ $Y'$).
\end{enumerate}
\end{proposition}

\begin{proof}
\oldpage[IV-4]{200}
\begin{enumerate}[label=(\roman*)]
  \item We may work affinely, with $S=\operatorname{Spec}(\mathrm{A})$, $X=\operatorname{Spec}(\mathrm{B})$, and $\sh{J}=\widetilde{\mathfrak{J}}$. Put $Q_m=\mathrm{B}/\mathfrak{J}^m$ for $m\geq1$ and $G_r=\mathfrak{J}^r/\mathfrak{J}^{r+1}$ for $r\geq0$. The graded algebra in~\emph{a)} is flat if and only if every $G_r$ is flat. Since $Q_1=G_0$, the exact sequences $0\to G_r\to Q_{r+1}\to Q_r\to0$ for $r\geq1$ prove the equivalence of~\emph{a)} and~\emph{b)}: extensions of flat modules are flat, and the kernel of a surjection between flat modules is flat \sref[0]{0.6.1.2}. The Tor exact sequence associated to $0\to\mathfrak{J}^m\to\mathrm{B}\to Q_m\to0$ shows that~\emph{b)} implies~\emph{c)}. Conversely, suppose that $\mathrm{B}$ is flat over $\mathrm{A}$ and that~\emph{c)} holds. For every $\mathrm{A}$-algebra $\mathrm{A}'$, the same exact sequence and the injectivity in~\emph{c)} give $\operatorname{Tor}_1^{\mathrm{A}}(Q_m,\mathrm{A}')=0$. Taking the square-zero algebra $\mathrm{A}'=\mathrm{A}\oplus\mathrm{M}$, with multiplication $(a,m)(a',m')=(aa',am'+a'm)$, for an arbitrary $\mathrm{A}$-module $\mathrm{M}$, gives $\operatorname{Tor}_1^{\mathrm{A}}(Q_m,\mathrm{M})=0$. Thus every $Q_m$ is flat, proving the conditional converse. The following conclusions are made under~\emph{a)}, or equivalently~\emph{b)}.
If $\mathfrak{J}_1=\mathfrak{J}\otimes_{\mathrm{A}}\mathrm{A}_1$, then
\[
  \bigoplus_{n\geq0}\mathfrak{J}_1^n=\left(\bigoplus_{n\geq0}\mathfrak{J}^n\right)\otimes_{\mathrm{A}}\mathrm{A}_1,
\]
and the assertion concerning the blown-up prescheme $X_1'$ follows immediately.
Finally, to prove that $Y'$ is $S$-flat, we may again suppose that $S$ and $X$ are affine. Put $\mathrm{C}=\operatorname{gr}_\mathfrak{J}^*(\mathrm{B})$. By \sref{IV.19.4.2} and \sref[II]{II.2.3.6}, every point of $Y'$ has an affine neighbourhood whose ring is $\mathrm{C}_{(t)}$, with $t$ homogeneous of degree $1$. Since $\mathrm{C}$ is flat over $\mathrm{A}$, so is the localization $\mathrm{C}_t$, and hence its degree-zero direct summand $\mathrm{C}_{(t)}$. Thus $Y'$ is $S$-flat.

  \item Again suppose that $S$ and $X$ are affine, so that $\mathrm{B}$ is a finitely presented $\mathrm{A}$-algebra and $\mathfrak{J}$ is a finitely generated ideal of $\mathrm{B}$. By \sref{IV.8.9.1}, \sref{IV.8.6.3}, and \sref{IV.11.2.9}, there exist a Noetherian subring $\mathrm{A}_0\subset\mathrm{A}$, a finite-type $\mathrm{A}_0$-algebra $\mathrm{B}_0$, and a finitely generated ideal $\mathfrak{J}_0\subset\mathrm{B}_0$ such that
\[
  \mathrm{B}=\mathrm{B}_0\otimes_{\mathrm{A}_0}\mathrm{A},\qquad
  \mathfrak{J}=\mathfrak{J}_0\mathrm{B},\qquad
  \operatorname{gr}_{\mathfrak{J}}^*(\mathrm{B})=\operatorname{gr}_{\mathfrak{J}_0}^*(\mathrm{B}_0)\otimes_{\mathrm{A}_0}\mathrm{A},
\]
with $\operatorname{gr}_{\mathfrak{J}_0}^*(\mathrm{B}_0)$ flat over $\mathrm{A}_0$. By~(i),
\[
  \bigoplus_{n\geq0}\mathfrak{J}^n=\left(\bigoplus_{n\geq0}\mathfrak{J}_0^n\right)\otimes_{\mathrm{A}_0}\mathrm{A}=\left(\bigoplus_{n\geq0}\mathfrak{J}_0^n\right)\otimes_{\mathrm{B}_0}\mathrm{B}.
\]
The algebra $\bigoplus_{n\geq0}\mathfrak{J}_0^n$ is generated over $\mathrm{B}_0$ by the finitely generated ideal $\mathfrak{J}_0$; it is therefore of finite type, and hence of finite presentation because $\mathrm{B}_0$ is Noetherian. Thus $\bigoplus_{n\geq0}\mathfrak{J}^n$ is a finitely presented $\mathrm{B}$-algebra. Likewise, $\operatorname{Proj}(\bigoplus_{n\geq0}\mathfrak{J}_0^n)\to\operatorname{Spec}(\mathrm{B}_0)$ is of finite type (\sref[II]{II.2.7.1}), hence of finite presentation, and its base change $X'\to X$ is of finite presentation.
By \sref{IV.19.4.2}, the canonical immersion $Y'\to X'$ is regular. Since $Y'$ is $S$-flat by~(i), \sref{IV.19.2.4} shows that $Y'\to X'$ is transversally regular relative to $S$ and that $X'$ is $S$-flat at the points of $Y'$, hence on an open neighbourhood of $Y'$ by \sref{IV.11.3.1}. Finally, since $\operatorname{gr}_\sh{J}^*(\sh{O}_X)$ is an $S$-flat $\sh{O}_X$-algebra, \sref{IV.11.3.4} shows that $X$ is $S$-flat at the points of $Y$, and hence on a neighbourhood of $Y$ by \sref{IV.11.3.1}.
\end{enumerate}
\end{proof}

\begin{corollary}[19.4.7]
\label{IV.19.4.7}
Let $g:X\to S$ be a morphism locally of finite presentation, $Y$ a closed sub-prescheme of $X$ such that the composite morphism $Y\to X\to S$ is locally of finite presentation; suppose moreover that $Y$ is S-flat and that $\sh{O}_X$ is normally flat along $Y$ (\sref{IV.11.3.4}).
Then (with the notations of \sref{IV.19.4.1}), the morphism $X'\to X$ is of finite presentation, the canonical immersion $Y'\to X'$ is transversally regular of codimension $1$ relative to $S$, $Y'$ is flat over $Y$ (and hence over $S$), and the set of points of $X$ (resp.\ $X'$) where $X$ (resp.\ $X'$) is $S$-flat is a neighbourhood of $Y$ (resp.\ $Y'$).
\end{corollary}

\begin{proof}
Indeed, by hypothesis $\operatorname{gr}_\sh{J}^*(\sh{O}_X)$ is a flat $(\sh{O}_X/\sh{J})$-module, and hence is $S$-flat because $Y$ is $S$-flat.
We can therefore apply the results of \sref{IV.19.4.6},~(i) and~(ii).
Since $Y'$ is $Y$-isomorphic to $\operatorname{Proj}(\operatorname{gr}_\sh{J}^*(\sh{O}_X))$ (\sref{IV.19.4.2}), the same argument as in \sref{IV.19.4.6} shows that $Y'$ is flat over $Y$.
\end{proof}

\begin{proposition}[19.4.8]
\label{IV.19.4.8}
Suppose the hypotheses of \sref{IV.19.4.7} satisfied and moreover that the morphism $Y'\to Y$ is smooth.
Let $x'$ be a point of $Y'$, $x$ its image in $Y$.
For $Y$ to be smooth over $S$ at the point $x$, it is necessary and sufficient that $Y'$ be smooth over $S$ at the point $x'$; when this is so, $X'$ is then smooth over $S$ at the point $x'$.
The preceding conclusions hold in particular when $X\to S$ is locally of finite presentation and the immersion $Y\to X$ is transversally regular relative to $S$.
\end{proposition}

\begin{proof}
Taking into account the compatibility of the formation of a blown-up prescheme with changes of base in the case envisaged \sref{IV.19.4.6},~(i) and of (\sref{IV.17.5.1}) and \sref{IV.17.7.1},
\oldpage[IV-4]{201}
it suffices to consider the case where $S$ is the spectrum of an algebraically closed field.
But then it amounts to the same to say that an S-prescheme locally of finite type is smooth or that it is regular at this point (\sref{IV.17.15.1}), hence the conclusions follow from \sref{IV.19.4.4}.

When $X\to S$ is locally of finite presentation and the canonical immersion $Y\to X$ is quasi-regular, that immersion is, by definition, of finite presentation, since the Ideal defining it is of finite type; hence the composite $Y\to X\to S$ is locally of finite presentation.
We have already seen in \sref{IV.19.4.3} that under these conditions $Y'\to Y$ is smooth; moreover, $\sh{N}_{Y/X}$ is a locally free $\sh{O}_Y$-Module and $\operatorname{gr}_\sh{J}^*(\sh{O}_X)$ is isomorphic to $\mathbf{S}^*_{\sh{O}_Y}(\sh{N}_{Y/X})$, hence is a flat $\sh{O}_Y$-Module.
The hypotheses of \sref{IV.19.4.7} are thus all satisfied and we can apply the conclusions of \sref{IV.19.4.8}.
\end{proof}

\begin{corollary}[19.4.9]
\label{IV.19.4.9}
Under the general hypotheses of \sref{IV.19.4.8}, suppose moreover that $Y$ is smooth over $S$ at the point $x$.
Then, for every generization $x_1$ of $x$ in $X$ not belonging to $Y$, $X$ is smooth over $S$ at the point $x_1$.
If $Y$ is smooth over $S$, then there exists an open neighbourhood $U$ of $Y$ in $X$ such that $X$ is smooth over $S$ at the points of $U-Y$.
\end{corollary}

\begin{proof}
To prove the first assertion, we can restrict ourselves to the case where $X$ is a local scheme with closed point $x$; the hypothesis then implies that $Y$ is smooth over $S$, since every neighbourhood in $Y$ of its closed point is necessarily all of $Y$. We conclude by the same reasoning as in \sref{IV.19.4.5}, using the fact that the set of points where $X'$ is smooth over $S$ is open in $X'$.
To prove the second assertion, we can reduce to the case where $S$ is Noetherian, and $X$ of finite type over $S$, using (\sref{IV.8.9.1}), (\sref{IV.8.6.3}), (\sref{IV.11.2.9}) and \sref{IV.17.7.6}; we then conclude as in \sref{IV.19.4.5}, using the fact that the set of points where $X$ is smooth over $S$ is open in $X$.
\end{proof}

\begin{remark}[19.4.10]
\label{IV.19.4.10}
In \sref{IV.19.4.6},~(ii), replace the hypothesis that $X$ and $Y$ are locally of finite presentation over $S$ by the hypothesis that $S$ and $X$ (hence also $Y$) are \emph{locally Noetherian}.
Then it is clear that $\bigoplus_{n\geq 0}\sh{J}^n$ is an $\sh{O}_X$-Algebra of finite type, that the morphism $X'\to X$ is of finite type, and it follows again from \sref{IV.19.2.4} applied to locally Noetherian preschemes, that the immersion $Y'\to X'$ is transversally regular relatively to $S$ and that $X'$ is S-flat at the points of $Y'$.
\end{remark}

\begin{proposition}[19.4.11]
\label{IV.19.4.11}
Let $X$ be a prescheme, $\sh{E}$ a quasi-coherent $\sh{O}_X$-Module, $u:\sh{E}\to\sh{O}_X$ a homomorphism of $\sh{O}_X$-Modules, $\sh{J}=u(\sh{E})$ the quasi-coherent Ideal of $\sh{O}_X$ which is the image of $u$, and $Y$ the closed sub-prescheme of $X$ defined by $\sh{J}$. Let $X'$ be the $X$-scheme obtained by blowing up $Y$.
On the other hand, let $\mathrm{P}=\mathbf{P}(\sh{E})$ be the projective bundle over $X$ defined by $\sh{E}$ (\sref[II]{II.4.1.1}), $p:\mathrm{P}\to X$ the structural morphism, $\alpha_1^{\#}:p^*(\sh{E})\to\sh{O}_{\mathrm{P}}(1)$ the canonical homomorphism (\sref[II]{II.4.1.5.1}) and $\sh{H}$ its kernel, so that we have an exact sequence
\[
  0\to\sh{H}\to p^*(\sh{E})\xrightarrow{\alpha_1^{\#}}\sh{O}_{\mathrm{P}}(1)\to 0.
\]
Finally, let $\sh{K}$ be the quasi-coherent Ideal of $\sh{O}_{\mathrm{P}}$ which is the image of the restriction $v:\sh{H}\to\sh{O}_{\mathrm{P}}$ of $p^*(u)$, and let $Z$ be the closed sub-prescheme of $\mathrm{P}$ defined by $\sh{K}$.
\oldpage[IV-4]{202}
\begin{enumerate}[label=(\roman*)]
  \item The image of $X'$ by the closed immersion $j:X'\to\mathbf{P}(\sh{E})$ corresponding to the surjective homomorphism of graded $\sh{O}_X$-Algebras
  \[
    \mathbf{S}_{\sh{O}_X}(\sh{E})\longrightarrow\mathbf{S}_{\sh{O}_X}(\sh{J})\longrightarrow\bigoplus_{n\geq 0}\sh{J}^n
  \]
  (\sref[II]{II.3.6.2}) is a closed sub-prescheme of $Z$.
  \item If $\sh{E}$ is an $\sh{O}_X$-Module locally free of rank $m$, $\sh{H}$ is an $\sh{O}_{\mathrm{P}}$-Module locally free of rank $m-1$.
  \item Suppose $X$ is locally Noetherian, that $\sh{E}$ is locally free of finite type, that the sub-prescheme $Y$ is regular, and that at every point $x\in Y$, the canonical immersion $i:Y\to X$ is regular and of codimension equal to $\operatorname{rg}_{\sh{O}_{X,x}}(\sh{E}_x)$ (which we shall express by saying that the homomorphism $u$ is regular at the point $x$ (cf.\ chap.~V)).
Then the image of $X'$ by the closed immersion $j:X'\to\mathrm{P}$ is the sub-prescheme $Z$; at all points of $X'-Y'$ ($Y'$ being the inverse image of $Y$ in $X'$) the immersion $j$ is transversally regular relatively to $X$; and finally, at all points $x'$ of $Y'$, $X'$ is regular, the immersion $j$ is regular and of codimension $\operatorname{rg}_{\sh{O}_{P,z}}(\sh{H}_z)$ where $z=j(x')$ (in other words the homomorphism $v:\sh{H}\to\sh{O}_{\mathrm{P}}$ is regular in $z$).
\end{enumerate}
\end{proposition}

\begin{proof}
All the questions being local on $X$, we can restrict ourselves to the case where $X=\operatorname{Spec}(\mathrm{A})$ is affine and $\sh{E}=\widetilde{\mathrm{E}}$, where $\mathrm{E}$ is an $\mathrm{A}$-module.
To verify~(i), we can in addition examine what happens in an affine open of $\mathrm{P}$ of the form $\mathrm{D}_+(t)$, where $t\in\mathrm{E}=\mathbf{S}_\mathrm{A}^1(\mathrm{E})$ (\sref[II]{II.2.3.14}).
If we set $\mathrm{S}=\mathbf{S}_\mathrm{A}(\mathrm{E})$, then
$\Gamma(\mathrm{D}_+(t),\sh{O}_{\mathrm{P}})=\mathrm{S}_{(t)}$ (\sref[II]{II.2.4.1}) and
$p^*(\sh{E})|_{\mathrm{D}_+(t)}=(\mathrm{E}\otimes_\mathrm{A}\mathrm{S}_{(t)})^\sim$.
Referring to the definition of $\alpha_1^{\#}$ (\sref[II]{II.2.6.2.2}), we see that to every element
\[
  a=\sum_i x^{(i)}\otimes\frac{x_1^{(i)}x_2^{(i)}\cdots x_n^{(i)}}{t^n}
\]
of $\mathrm{E}\otimes_\mathrm{A}\mathrm{S}_{(t)}$, where the $x^{(i)}$ and $x_k^{(i)}$ belong to $\mathrm{E}$, it assigns
\[
  \sum_i\frac{x^{(i)}x_1^{(i)}x_2^{(i)}\cdots x_n^{(i)}}{t^n}
  \quad\text{in }(\mathrm{S}(1))_{(t)}.
\]
To prove~(i), it is enough to show that, if this last element is zero, then the image of
\[
  a'=v(a)=\sum_i u(x^{(i)})\frac{x_1^{(i)}x_2^{(i)}\cdots x_n^{(i)}}{t^n}
\]
under the canonical homomorphism
$\mathrm{S}_{(t)}\to\left(\bigoplus_{n\geq0}u(\mathrm{E})^n\right)_{(u(t))}$ (\sref[II]{II.3.6.2}) is also zero.
Its image in $\mathrm{A}_{u(t)}$ is
\[
  \sum_i\frac{u(x^{(i)})u(x_1^{(i)})\cdots u(x_n^{(i)})}{u(t)^n},
\]
that is, the product of $u(t)$ and the canonical image of
$\sum_i x^{(i)}x_1^{(i)}\cdots x_n^{(i)}/t^{n+1}$ under the algebra homomorphism extending $u:\mathrm{E}\to\mathrm{A}$ from $\mathrm{S}_{(t)}$. This proves~(i).

To establish~(ii), we may suppose that $\mathrm{E}$ is a free $\mathrm{A}$-module of rank $m$. Since $\alpha_1^{\#}:\mathrm{E}\otimes_\mathrm{A}\mathrm{S}_{(t)}\to(\mathrm{S}(1))_{(t)}$ is surjective and $(\mathrm{S}(1))_{(t)}$ is a free $\mathrm{S}_{(t)}$-module of rank $1$, the exact sequence
\[
  0\to\mathrm{H}_{(t)}\to\mathrm{E}\otimes_\mathrm{A}\mathrm{S}_{(t)}\xrightarrow{\alpha_1^{\#}}(\mathrm{S}(1))_{(t)}\to 0
\]
splits, and $\mathrm{H}_{(t)}$ is therefore a projective $\mathrm{S}_{(t)}$-module of rank $m-1$, whence~(ii).

To prove~(iii), let us examine first what happens above a point $x\in X-Y$.
The homomorphism $u:\sh{E}\to\sh{O}_X$, restricted to a neighbourhood of $x$, is then surjective, so we may restrict to the case $\sh{J}=\sh{O}_X$, in which $X'$ is canonically identified with $X$.
To see that the closed sub-prescheme of $\mathrm{P}$ which is the image of $X'$ is then identical with $Z$, we may, as at the beginning, suppose $X=\operatorname{Spec}(\mathrm{A})$ and $\mathrm{E}=\mathrm{A}^m$, so that $\mathrm{S}=\mathrm{A}[\mathrm{T}_0,\ldots,\mathrm{T}_{m-1}]$.
With the notation above we may suppose $u(t)\ne0$; the elements $t\in\mathrm{E}$ having this property generate the symmetric algebra $\mathbf{S}_\mathrm{A}(\mathrm{E})$.
After changing the basis of $\mathrm{E}$ if necessary, we may therefore suppose $t=\mathrm{T}_0$ and $u(\mathrm{T}_i)=0$ for $1\leq i\leq m-1$, so that $\mathrm{S}_{(t)}$ is canonically identified with $\mathrm{A}[\mathrm{T}_1,\ldots,\mathrm{T}_{m-1}]$.
The elements of $\mathrm{H}$ are then those of the form
\[
  \sum_\alpha \mathrm{L}_\alpha(\mathbf{T})\otimes\mathbf{T}^\alpha
\]
in $\mathrm{E}\otimes_\mathrm{A}\mathrm{A}[\mathrm{T}_1,\ldots,\mathrm{T}_{m-1}]$, where
$\mathrm{L}_\alpha(\mathbf{T})=\lambda_{\alpha0}+\lambda_{\alpha1}\mathrm{T}_1+\cdots+\lambda_{\alpha,m-1}\mathrm{T}_{m-1}$ and
$\sum_\alpha\mathrm{L}_\alpha(\mathbf{T})\mathbf{T}^\alpha=0$.
One checks readily that, subject to $\lambda_{\mathbf{0},0}=0$, the values of $\lambda_{\alpha0}$ may be chosen arbitrarily for $|\alpha|\geq1$, after which the $\lambda_{\alpha i}$ for $i\geq1$ may be chosen so as to satisfy the displayed relation.
The corresponding elements of the Ideal $\sh{K}$ are therefore precisely the polynomials
$\sum_\alpha\lambda_{\alpha0}\mathbf{T}^\alpha$ without constant term.
Since $(\mathbf{S}_\mathrm{A}(\mathrm{A}))_{(t)}$ is identified with $\mathrm{A}$, these polynomials are exactly the kernel of
$(\mathbf{S}_\mathrm{A}(\mathrm{E}))_{(t)}\to(\mathbf{S}_\mathrm{A}(\mathrm{A}))_{(t)}$.
We have thus proved, by \sref{IV.19.2.1}, that $j:X'\to\mathrm{P}$ is transversally regular relative to $X$ at every point of $X'-Y'$ and that $j(X')=Z$ at those points.

\oldpage[IV-4]{203}
It remains to examine what happens above a point $x\in Y$.
The regularity hypothesis on $u$ at $x$ in~(iii) is equivalent, by \sref{IV.19.1.1} and (\sref[0]{0.17.1.7}), to saying that $X$ is regular at $x$ and that
\[
  u_x\otimes1:\sh{E}_x\otimes_{\sh{O}_{X,x}}\mathbf{k}(x)\longrightarrow\mathfrak{m}_x/\mathfrak{m}_x^2
\]
is injective.
Since $p:\mathrm{P}\to X$ is smooth (\sref{IV.17.3.9}), for every $z\in\mathrm{P}$ above $x$, $\mathrm{P}$ is regular at $z$ (\sref{IV.17.5.8}); moreover, by (\sref[0]{0.17.3.3}), the canonical homomorphism
\[
  (\mathfrak{m}_x/\mathfrak{m}_x^2)\otimes_{\mathbf{k}(x)}\mathbf{k}(z)\longrightarrow\mathfrak{m}_z/\mathfrak{m}_z^2
  \tag{19.4.11.1}\label{IV.19.4.11.1}
\]
is injective.
It follows that the homomorphism
\[
  (\sh{E}_x\otimes_{\sh{O}_{X,x}}\sh{O}_{\mathrm{P},z})\otimes_{\sh{O}_{\mathrm{P},z}}\mathbf{k}(z)
  \longrightarrow\mathfrak{m}_z/\mathfrak{m}_z^2
\]
induced by $p^*(u)$ is also injective, since it is the composite
\[
  (\sh{E}_x\otimes_{\sh{O}_{X,x}}\mathbf{k}(x))\otimes_{\mathbf{k}(x)}\mathbf{k}(z)
  \longrightarrow(\mathfrak{m}_x/\mathfrak{m}_x^2)\otimes_{\mathbf{k}(x)}\mathbf{k}(z)
  \longrightarrow\mathfrak{m}_z/\mathfrak{m}_z^2,
\]
where the first arrow is injective by flatness and the second is the injective homomorphism \sref{IV.19.4.11.1}.
Since $\sh{H}$ is, in a neighbourhood of $z$ in $\mathrm{P}$, a direct factor of $p^*(\sh{E})$, the homomorphism
$v_z\otimes1:\sh{H}_z\otimes_{\sh{O}_{\mathrm{P},z}}\mathbf{k}(z)\to\mathfrak{m}_z/\mathfrak{m}_z^2$ is also injective; this establishes the regularity of $v$ at $z$.
It remains to see that the image by $j$ of $X'$ is identical to the closed sub-prescheme $Z$.
It will suffice for this to show that the image by $j$ of the points of $q^{-1}(x)$ (where $q:X'\to X$ is the structural morphism) are exactly those of $Z\cap p^{-1}(x)$ and that on the other hand at each of these points $z=j(x')$, the surjective homomorphism $\sh{O}_{Z,z}\to\sh{O}_{X',x'}$ is bijective.
By the hypothesis and \sref{IV.19.4.3}, the restriction $Y'\to Y$ of $q$ is smooth; since $\sh{O}_{Y,x}$ is regular, so is $\sh{O}_{Y',x'}$ (\sref{IV.17.5.8}). Moreover, $Y'\to X'$ is a regular immersion (\sref{IV.19.4.2}), and hence $\sh{O}_{X',x'}$ is regular (\sref{IV.19.1.1}). The regularity of $v$ and of $\mathrm{P}$ at $z$ similarly shows that $\sh{O}_{Z,z}$ is regular.
To prove that $\sh{O}_{Z,z}\to\sh{O}_{X',x'}$ is bijective, it therefore suffices to show that the dimensions
\oldpage[IV-4]{204}
of these two regular rings are equal (\sref[0]{0.17.1.9}).
By (\sref[0]{0.17.3.3}), this is equivalent to equality of the dimensions of
$\sh{O}_{X',x'}\otimes_{\sh{O}_{X,x}}\mathbf{k}(x)$ and
$\sh{O}_{Z,z}\otimes_{\sh{O}_{X,x}}\mathbf{k}(x)$.
We may therefore replace $X'$ and $Z$ by the fibres $q^{-1}(x)$ and $Z\cap p^{-1}(x)$ respectively, and reduce to proving that the morphism
$q^{-1}(x)\to Z\cap p^{-1}(x)$ induced by $j$ is an isomorphism.
We may now suppose that $X=\operatorname{Spec}(\mathrm{B})$, where $\mathrm{B}=\sh{O}_{X,x}$, and that $\mathrm{E}=\mathrm{B}^m$, so that $\mathrm{S}=\mathrm{B}[\mathrm{T}_1,\ldots,\mathrm{T}_m]$.
The hypothesis means that there is a regular system of parameters $(t_i)_{1\leq i\leq n}$ of $\mathrm{B}$ such that $u(\mathrm{T}_i)=t_i$ for $1\leq i\leq m$ (\sref[0]{0.17.1.7}).
Let $\mathfrak{J}$ be the image of $u$, and let $\mathfrak{m}$ be the maximal ideal of $\mathrm{B}$. Then
\[
  \mathfrak{J}^n\otimes_\mathrm{B}(\mathrm{B}/\mathfrak{m})
  =\mathfrak{J}^n/\mathfrak{m}\mathfrak{J}^n
  =(\mathfrak{J}^n/\mathfrak{J}^{n+1})\otimes_\mathrm{B}(\mathrm{B}/\mathfrak{m}),
\]
and consequently, by \sref{IV.16.9.4.1},
$(\bigoplus_{n\geq0}\mathfrak{J}^n)\otimes_\mathrm{B}\mathbf{k}(x)$ is identified with
$\mathbf{k}(x)[\mathrm{T}_1,\ldots,\mathrm{T}_m]$.
It follows that the immersion $q^{-1}(x)\to p^{-1}(x)$ is an isomorphism; this implies \emph{a fortiori} that the closed sub-prescheme $Z\cap p^{-1}(x)$ is identical with $j(q^{-1}(x))$, and completes the proof of \sref{IV.19.4.11}.
\end{proof}

\subsection{Criteria for M-regularity}
\label{subsection:IV.19.5}

We resume here, in completed form, the criteria for a sequence of elements of a ring $\mathrm{A}$ to be M-regular or M-quasi-regular, where $\mathrm{M}$ is an $\mathrm{A}$-module, already studied in (\sref[0]{0.15.1}).

\begin{theorem}[19.5.1]
\label{IV.19.5.1}
Let $\mathrm{A}$ be a ring, $\mathbf{f}=(f_i)_{1\leq i\leq n}$ a sequence of elements of $\mathrm{A}$, and $\mathrm{M}$ an $\mathrm{A}$-module. Consider the following properties:
\begin{enumerate}
\item[\rm a)] The sequence $\mathbf{f}$ is M-regular.
\item[\rm b)] $\mathrm{H}_i(\mathbf{f},\mathrm{M})=0$ for every $i>0$ (\sref[III]{III.1.1.3}).
\item[\rm b')] $\mathrm{H}_1(\mathbf{f},\mathrm{M})=0$.
\item[\rm c)] The sequence $\mathbf{f}$ is M-quasi-regular.
\end{enumerate}
Then we have the implications
\[
  \text{a)}\Longrightarrow\text{b)}\Longrightarrow\text{b')}
  \qquad\text{and}\qquad
  \text{a)}\Longrightarrow\text{c)}.
\]
Put $\mathfrak{J}=\sum_{i=1}^n f_i\mathrm{A}$.
If the modules $\mathrm{M}_i=\mathrm{M}/(\sum_{j=1}^i f_j\mathrm{M})$ are separated for the $\mathfrak{J}$-preadic topology (respectively, if every quotient module of a submodule of $\mathrm{M}$ is separated for the $\mathfrak{J}$-preadic topology), then we also have the implication c)$\Rightarrow$a) (respectively, b')$\Rightarrow$a)).
\end{theorem}

\begin{proof}
The implications a)$\Rightarrow$c), and c)$\Rightarrow$a) when the $\mathrm{M}_i$ are separated for the $\mathfrak{J}$-preadic topology, were already proved in (\sref[0]{0.15.1.9}) and are recalled only for completeness.
We have also shown in (\sref[III]{III.1.1.4}) and (\sref[III]{III.1.1.3.3}) that a) implies b), and b) trivially implies b').
It therefore remains to prove that, when every quotient module of a submodule of $\mathrm{M}$ is separated for the $\mathfrak{J}$-preadic topology, b') implies a).
We argue by induction on $n$.
For $n=1$, the assertion follows immediately from the definitions, since $\mathrm{H}_1(\mathbf{f},\mathrm{M})$ is precisely the kernel in $\mathrm{M}$ of multiplication $z\mapsto f_1z$ (\sref[III]{III.1.1.1}, \sref[III]{III.1.1.2}).
Suppose then that $n\geq2$, and put $\mathbf{f}'=(f_i)_{1\leq i\leq n-1}$.
With the notation of (\sref[III]{III.1.1.2}),
$\mathrm{K}_\bullet(\mathbf{f},\mathrm{M})=\mathrm{K}_\bullet(f_n)\otimes_\mathrm{A}\mathrm{K}_\bullet(\mathbf{f}',\mathrm{M})$.
We therefore have (\sref[III]{III.1.1.4.1}) the exact sequence
\[
  0\longrightarrow\mathrm{H}_0(f_n,\mathrm{H}_1(\mathbf{f}',\mathrm{M}))
  \longrightarrow\mathrm{H}_1(\mathbf{f},\mathrm{M})
  \longrightarrow\mathrm{H}_1(f_n,\mathrm{H}_0(\mathbf{f}',\mathrm{M}))
  \longrightarrow0.
  \tag{19.5.1.1}\label{IV.19.5.1.1}
\]

\oldpage[IV-4]{205}

The relation $\mathrm{H}_1(\mathbf{f},\mathrm{M})=0$ implies therefore $\mathrm{H}_0(f_n,\mathrm{H}_1(\mathbf{f}',\mathrm{M}))=0$ and $\mathrm{H}_1(f_n,\mathrm{H}_0(\mathbf{f}',\mathrm{M}))=0$.
The first of these two relations means that we have $f_n\mathrm{H}_1(\mathbf{f}',\mathrm{M})=\mathrm{H}_1(\mathbf{f}',\mathrm{M})$ \quad (\textbf{III},
$1.1.3.5$). Now, by definition (\textbf{III}, $1.1.1$), $\mathrm{H}_1(\mathbf{f}',\mathrm{M})$ is isomorphic to a quotient of a submodule $\mathrm{N}$ of $\mathrm{M}^{n-1}$; taking on $\mathrm{M}^{n-1}$ the filtration of the $\mathrm{M}^j$ for $j\leqslant n-1$, on $\mathrm{N}$ the induced filtration and on $\mathrm{H}_1(\mathbf{f}',\mathrm{M})$ the quotient filtration from that of $\mathrm{N}$, we obtain on $\mathrm{H}_1(\mathbf{f}',\mathrm{M})$ a finite filtration whose quotients are isomorphic to quotients of submodules of $\mathrm{M}$, hence are separated for the $\mathfrak{J}$-preadic topology by hypothesis.
As the relation $f_n\mathrm{H}_1(\mathbf{f}',\mathrm{M})=\mathrm{H}_1(\mathbf{f}',\mathrm{M})$ implies \emph{a fortiori} $\mathfrak{J}\mathrm{H}_1(\mathbf{f}',\mathrm{M})=\mathrm{H}_1(\mathbf{f}',\mathrm{M})$,
we deduce from the hypothesis and the preceding remarks that $\mathrm{H}_1(\mathbf{f}',\mathrm{M})=0$. The induction hypothesis proves then that the sequence $\mathbf{f}'$ is M-regular. On the other hand, the relation
$\mathrm{H}_1(f_n,\mathrm{H}_0(\mathbf{f}',\mathrm{M}))=0$ means that $f_n$ is $(\mathrm{M}/(\sum_{i=1}^{n-1}f_i\mathrm{M}))$-regular, hence the sequence $\mathbf{f}$ is M-regular.
\end{proof}

\begin{corollary}[19.5.2]
\label{IV.19.5.2}
Suppose that $\mathrm{A}$ is a Noetherian ring, that the $f_i$ $(1\leqslant i\leqslant n)$ belong to the radical of $\mathrm{A}$, and that $\mathrm{M}$ is a finitely generated $\mathrm{A}$-module. Then the four conditions \emph{a)}, \emph{b)}, \emph{b')}, \emph{c)} of \emph{(\sref{IV.19.5.1})} are equivalent.
\end{corollary}

\begin{proof}
We know indeed from $(\mathbf{0}_\mathrm{I},7.3.5)$ that every finitely generated $\mathrm{A}$-module is separated for the $\mathfrak{J}$-preadic topology.
\end{proof}

\begin{proposition}[19.5.3]
\label{IV.19.5.3}
Let
\[
  0\to\mathrm{M}'\to\mathrm{M}\to\mathrm{M}''\to 0
\]
be an exact sequence of $\mathrm{A}$-modules, $\mathbf{f}=(f_i)_{1\leqslant i\leqslant n}$ an M-regular sequence, $\mathfrak{J}=\sum_{i=1}^{n}f_i\mathrm{A}$. Consider the following properties:
\begin{enumerate}
\item[\rm a)] The sequence $\mathbf{f}$ is $\mathrm{M}''$-regular.
\item[\rm b)] The $\mathfrak{J}$-preadic filtration of $\mathrm{M}'$ is induced on $\mathrm{M}'$ by the $\mathfrak{J}$-preadic filtration of $\mathrm{M}$ (in other words, the canonical homomorphism $\operatorname{gr}_\mathfrak{J}^*(\mathrm{M}')\to\operatorname{gr}_\mathfrak{J}^*(\mathrm{M})$ is injective).
\item[\rm c)] The canonical homomorphism $\mathrm{M}'/(\sum_{i=1}^{n}f_i\mathrm{M}')\to\mathrm{M}/(\sum_{i=1}^{n}f_i\mathrm{M})$ is injective.
\end{enumerate}
Then we have the implications \emph{a)} $\Rightarrow$ \emph{b)} $\Rightarrow$ \emph{c)}, and \emph{a)} implies moreover that the sequence $\mathbf{f}$ is $\mathrm{M}'$-regular.
If every quotient of a submodule of $\mathrm{M}''$ is separated for the $\mathfrak{J}$-preadic topology, the conditions \emph{a)}, \emph{b)} and \emph{c)} are equivalent.
\end{proposition}

\begin{proof}
It is clear that $c)$ is a consequence of $b)$. Let us prove that under all the hypotheses of the last assertion, $c)$ implies $a)$: since $\mathbf{f}$ is M-regular, by \sref{IV.19.5.1}, $\mathrm{H}_1(\mathbf{f},\mathrm{M})=0$, whence an exact sequence, portion of the exact homology sequence
\[
  0\to\mathrm{H}_1(\mathbf{f},\mathrm{M}'')\to\mathrm{H}_0(\mathbf{f},\mathrm{M}')\to\mathrm{H}_0(\mathbf{f},\mathrm{M}).
\]
Now, condition $c)$ expresses that the homomorphism $\mathrm{H}_0(\mathbf{f},\mathrm{M}')\to\mathrm{H}_0(\mathbf{f},\mathrm{M})$ is injective (\textbf{III}, $1.1.3.5$); it is therefore equivalent to $\mathrm{H}_1(\mathbf{f},\mathrm{M}'')=0$, whence, by virtue of the separation hypothesis and \sref{IV.19.5.1}, the fact that $\mathbf{f}$ is $\mathrm{M}''$-regular.

Let us show next that $a)$ implies $c)$, and also that the sequence $\mathbf{f}$ is $\mathrm{M}'$-regular, proceeding by induction on $n$. For $n=1$, $f_1$ being M-regular is also $\mathrm{M}'$-regular and property $c)$ is precisely the lemma (\sref{IV.3.4.1.4}). For $n>1$, the induction hypothesis shows that
\oldpage[IV-4]{206}
if we set $\mathbf{f}'=(f_1,\ldots,f_{n-1})$, the sequence $\mathbf{f}'$ is $\mathrm{M}'$-regular and, by virtue of $c)$, we have an exact sequence
\[
  0\to\mathrm{M}'/(\textstyle\sum_{i=1}^{n-1}f_i\mathrm{M}')\to\mathrm{M}/(\textstyle\sum_{i=1}^{n-1}f_i\mathrm{M})\to\mathrm{M}''/(\textstyle\sum_{i=1}^{n-1}f_i\mathrm{M}'')\to 0.
\]
By hypothesis, $f_n$ is $(\mathrm{M}/(\sum_{i=1}^{n-1}f_i\mathrm{M}))$-regular and $(\mathrm{M}''/(\sum_{i=1}^{n-1}f_i\mathrm{M}''))$-regular. The same reasoning shows, on the one hand, that $f_n$ is also $(\mathrm{M}'/(\sum_{i=1}^{n-1}f_i\mathrm{M}'))$-regular and, on the other hand, by (\sref{IV.3.4.1.4}), that the sequence
\[
  0\to\mathrm{M}'/(\textstyle\sum_{i=1}^{n}f_i\mathrm{M}')\to\mathrm{M}/(\textstyle\sum_{i=1}^{n}f_i\mathrm{M})\to\mathrm{M}''/(\textstyle\sum_{i=1}^{n}f_i\mathrm{M}'')\to 0
\]
is exact, whence $c)$.

Let us show finally that $a)$ implies $b)$. As the sequence $\mathbf{f}$ is then M-regular and $\mathrm{M}'$-regular, hence M-quasi-regular and $\mathrm{M}'$-quasi-regular, we have canonical isomorphisms
\[
  \operatorname{gr}_\mathfrak{J}^*(\mathrm{M}')\xrightarrow{\sim}\operatorname{gr}_\mathfrak{J}^0(\mathrm{M}')[\mathrm{T}_1,\ldots,\mathrm{T}_n],\qquad
  \operatorname{gr}_\mathfrak{J}^*(\mathrm{M})\xrightarrow{\sim}\operatorname{gr}_\mathfrak{J}^0(\mathrm{M})[\mathrm{T}_1,\ldots,\mathrm{T}_n]
\]
$(\mathbf{0},15.1.7)$, and as we saw above that $\operatorname{gr}_\mathfrak{J}^0(\mathrm{M}')\to\operatorname{gr}_\mathfrak{J}^0(\mathrm{M})$ is injective, the same holds for $\operatorname{gr}_\mathfrak{J}^*(\mathrm{M}')\to\operatorname{gr}_\mathfrak{J}^*(\mathrm{M})$.

We note moreover that the equivalence of conditions $a)$, $b)$, $c)$ of \sref{IV.19.5.3} holds in particular when $\mathrm{A}$ is Noetherian, $\mathrm{M}''$ is a finitely generated $\mathrm{A}$-module, and the $f_i$ belong to the radical of $\mathrm{A}$.
\end{proof}

\begin{env}[19.5.4]
\label{IV.19.5.4}
In the rest of this number, we preserve the notations $\mathrm{A}$, $\mathbf{f}$, $\mathrm{M}$, $\mathfrak{J}$, and we suppose moreover $\mathrm{M}$ endowed with a decreasing filtration $(\mathrm{M}_k)$ formed of sub-$\mathrm{A}$-modules, such that $\mathrm{M}_0=\mathrm{M}$. Recall $(\mathbf{0},15.1.5)$ that we then define on $\mathrm{M}$ a second decreasing filtration formed of submodules
\[
  \mathrm{M}_k'=\mathrm{M}_k+\mathfrak{J}\mathrm{M}_{k-1}+\ldots+\mathfrak{J}^{k-1}\mathrm{M}_1+\mathfrak{J}^k\mathrm{M}_0
\]
and that, if $\operatorname{gr}_\bullet(\mathrm{M})$ and $\operatorname{gr}_\bullet'(\mathrm{M})$ are the graded $\mathrm{A}$-modules associated to the filtrations $(\mathrm{M}_k)$ and $(\mathrm{M}_k')$ respectively, we define a surjective graded homomorphism of degree $0$ $(\mathbf{0},15.1.5.2)$
\[
  \psi_\mathbf{M}:(\operatorname{gr}_\bullet(\mathrm{M})\otimes_\mathrm{A}(\mathrm{A}/\mathfrak{J}))[\mathrm{T}_1,\ldots,\mathrm{T}_n]\to\operatorname{gr}_\bullet'(\mathrm{M})
\]
Recall also that we have the canonical surjective homomorphism $(\mathbf{0},15.1.1.1)$
\[
  \varphi_{\operatorname{gr}_\bullet(\mathrm{M})}:(\operatorname{gr}_\bullet(\mathrm{M})\otimes_\mathrm{A}(\mathrm{A}/\mathfrak{J}))[\mathrm{T}_1,\ldots,\mathrm{T}_n]\to\operatorname{gr}_\mathfrak{J}^*(\operatorname{gr}_\bullet(\mathrm{M}))
\]
\end{env}

\begin{theorem}[19.5.5]
\label{IV.19.5.5}
With the notations of \emph{(\sref{IV.19.5.4})}, consider the following properties:
\begin{enumerate}
\item[\rm a)] The sequence $\mathbf{f}$ is $\operatorname{gr}_\bullet(\mathrm{M})$-regular.
\item[\rm b)] The canonical homomorphism $\psi_\mathbf{M}$ is bijective.
\item[\rm c)] The canonical homomorphism $\varphi_{\operatorname{gr}_\bullet(\mathrm{M})}$ is bijective (in other words $(\mathbf{0},15.1.7)$, the sequence $\mathbf{f}$ is $\operatorname{gr}_\bullet(\mathrm{M})$-quasi-regular).
\end{enumerate}
\oldpage[IV-4]{207}
Then we have the implications\footnote{The implication $b)\Rightarrow c)$, without the separation hypothesis on the $\operatorname{gr}_k(\mathrm{M})$, is due to P.~Deligne, whose demonstration we reproduce below.}
\[
  \text{a)}\Rightarrow\text{b)}\Rightarrow\text{c)}.
\]
Moreover, if for every integer $k\geqslant 0$, the quotient modules of submodules of $\operatorname{gr}_k(\mathrm{M})$ are separated for the $\mathfrak{J}$-preadic topology (which is the case when $\mathrm{A}$ is Noetherian, the $f_i$ belong to the radical of $\mathrm{A}$, and the $\operatorname{gr}_k(\mathrm{M})$ are finitely generated $\mathrm{A}$-modules), \emph{then all conditions} \emph{a)}, \emph{b)} and \emph{c)} \emph{are equivalent}.
\end{theorem}

\begin{proof}
The implication $a)\Rightarrow b)$ has been demonstrated in $(\mathbf{0},15.1.8)$; the implication $c)\Rightarrow a)$ under the separation hypothesis is a particular case of $(\mathbf{0},15.1.9)$. It remains to prove that $b)$ implies $c)$, which will be done in several steps.

We denote by $\mathfrak{J}$ the ideal $\sum_{i=1}^{n}f_i\mathrm{A}$ generated by the $f_i$; for every sequence $\mathbf{q}=(q_1,\ldots,q_n)$ of $n$ integers $\geqslant 0$, we set $|\mathbf{q}|=\sum_{i=1}^{n}q_i$ and $\mathbf{f}^\mathbf{q}=f_1^{q_1}f_2^{q_2}\cdots f_n^{q_n}$.

\begin{env}[19.5.5.1]
\label{IV.19.5.5.1}
For the map $\psi_\mathbf{M}$ to be injective (and hence bijective), it is necessary and sufficient that the following condition be satisfied:

$(\mathrm{I}_{q,p})$ \emph{For every family $(x_\mathbf{q}^p)_{|\mathbf{q}|=q}$ of elements of $\mathrm{M}_p$, the relation}
\[
  \textstyle\sum_{|\mathbf{q}|=q}\mathbf{f}^\mathbf{q} x_\mathbf{q}^p\in\textstyle\sum_{i=1}^{q}\mathfrak{J}^{q-i}\mathrm{M}_{p+i}+\mathrm{M}_{p+q+1}'
\]
\emph{implies} $x_\mathbf{q}^p\in\mathrm{M}_{p+1}+\mathfrak{J}\mathrm{M}_p$ \emph{for every $\mathbf{q}$ such that $|\mathbf{q}|=q$.}

We proceed as in $(\mathbf{0},15.1.6)$ by considering the sub-$\mathrm{A}$-module $\mathrm{Q}_k$ (resp. $\mathrm{Q}_k'$) as a term of degree $k$ in the first (resp. second) member of $(\mathbf{0},15.1.5.2)$. We endow $\mathrm{Q}_k$ with the filtration
\[
  (\mathrm{Q}_k)_i=\textstyle\sum_{j\leqslant k-i}\bigl(\textstyle\sum_{|\mathbf{l}|=j}(\operatorname{gr}_{k-j}(\mathrm{M})\otimes_\mathrm{A}(\mathrm{A}/\mathfrak{J}))\bigr)\mathrm{T}^\mathbf{l}
\]
and $\mathrm{Q}_k'$ the image filtration formed by $\psi_\mathbf{M}((\mathrm{Q}_k)_i)=(\mathrm{Q}_k')_i$; it suffices then to prove that the homomorphisms $(\psi_\mathbf{M})_{ki}:\operatorname{gr}_i(\mathrm{Q}_k)\to\operatorname{gr}_i(\mathrm{Q}_k')$ are injective. Now, we have
\[
  \operatorname{gr}_i(\mathrm{Q}_k)=\textstyle\sum_{|\mathbf{l}|=k-i}\bigl((\mathrm{M}_i/\mathrm{M}_{i+1})\otimes_\mathrm{A}(\mathrm{A}/\mathfrak{J})\bigr)\mathrm{T}^\mathbf{l}
  =\textstyle\sum_{|\mathbf{l}|=k-i}\bigl(\mathrm{M}_i/(\mathfrak{J}\mathrm{M}_i+\mathrm{M}_{i+1})\bigr)\mathrm{T}^\mathbf{l}
\]
\oldpage[IV-4]{208}
on the other hand, $(\mathrm{Q}_k')_{i+1}$ is the image of $\sum_{h=0}^{k-i-1}\mathfrak{J}^{k-i-h-1}\mathrm{M}_{i+h+1}$ in $\mathrm{M}_k'/\mathrm{M}_{k+1}'$. To write that $(\psi_\mathbf{M})_{ki}$ is injective is therefore equivalent to writing the condition $(\mathrm{I}_{k-i,i})$; whence our assertion.
\end{env}

\begin{env}[19.5.5.2]
\label{IV.19.5.5.2}
\emph{For the sequence $\mathbf{f}$ to be $\operatorname{gr}_\bullet(\mathrm{M})$-quasi-regular, it suffices that, for $p\geqslant 0$ and $q\geqslant 0$, the following condition be satisfied:}

$(\mathrm{S}_{q,p})$ \emph{For every family $(x_\mathbf{q}^p)_{|\mathbf{q}|=q}$ of elements of $\mathrm{M}_p$, the relation}
\[
  \textstyle\sum_{|\mathbf{q}|=q}\mathbf{f}^\mathbf{q}x_\mathbf{q}^p\in\mathrm{M}_{p+1}
\]
\emph{implies $x_\mathbf{q}^p\in\mathrm{M}_{p+1}+\mathfrak{J}\mathrm{M}_p$ for every $\mathbf{q}$ such that $|\mathbf{q}|=q$.}

By definition, saying that $\mathbf{f}$ is $\operatorname{gr}_\bullet(\mathrm{M})$-quasi-regular means that, for every $p\geqslant 0$ and every $q\geqslant 0$, the relation
\[
  \textstyle\sum_{|\mathbf{q}|=q}\mathbf{f}^\mathbf{q} x_\mathbf{q}^p\in\mathrm{M}_{p+1}+\mathfrak{J}^{q+1}\mathrm{M}_p
\]
for a family $(x_\mathbf{q}^p)_{|\mathbf{q}|=q}$ of elements of $\mathrm{M}_p$ implies $x_\mathbf{q}^p\in\mathrm{M}_{p+1}+\mathfrak{J}\mathrm{M}_p$ for every $\mathbf{q}$ such that $|\mathbf{q}|=q$. The hypothesis on the family $(x_\mathbf{q}^p)_{|\mathbf{q}|=q}$ means that there exists a family $(y_\mathbf{q}^p)_{|\mathbf{q}|=q}$ of elements of $\mathfrak{J}\mathrm{M}_p$ such that $\sum_{|\mathbf{q}|=q}\mathbf{f}^\mathbf{q}(x_\mathbf{q}^p-y_\mathbf{q}^p)\in\mathrm{M}_{p+1}$; condition $(\mathrm{S}_{q,p})$ then implies $x_\mathbf{q}^p-y_\mathbf{q}^p\in\mathrm{M}_{p+1}+\mathfrak{J}\mathrm{M}_p$, whence $x_\mathbf{q}^p\in\mathrm{M}_{p+1}+\mathfrak{J}\mathrm{M}_p$; hence one sees that the conditions $(\mathrm{S}_{q,p})$ imply that the sequence $\mathbf{f}$ is $\operatorname{gr}_\bullet(\mathrm{M})$-quasi-regular.
\end{env}

It suffices therefore to prove the following proposition:

\begin{env}[19.5.5.3]
\label{IV.19.5.5.3}
\emph{For every $r\geqslant 0$, if the conditions $(\mathrm{I}_{q,p})$ are satisfied for $p+q<r$, then the conditions $(\mathrm{S}_{q,p})$ are satisfied for $p+q<r$.}

Let us introduce the conditions

$(\mathrm{A}_{q,p})$ \emph{For every family $(x_\mathbf{q}^p)_{|\mathbf{q}|=q}$ of elements of $\mathrm{M}_p$, the relation}
\[
  \textstyle\sum_{|\mathbf{q}|=q}\mathbf{f}^\mathbf{q}x_\mathbf{q}^p\in\mathrm{M}_{p+1}
\]
\emph{implies}
\[
  \textstyle\sum_{|\mathbf{q}|=q}\mathbf{f}^\mathbf{q}x_\mathbf{q}^p
  \in\textstyle\sum_{i=1}^{q}\mathfrak{J}^{q-i}\mathrm{M}_{p+i}.
\]

It is clear that the conjunction of $(\mathrm{A}_{q,p})$ and $(\mathrm{I}_{q,p})$ implies $(\mathrm{S}_{q,p})$. On the other hand, the conditions $(\mathrm{A}_{1,p})$ and $(\mathrm{S}_{0,p})$ are trivial. Consequently, \sref{IV.19.5.5.3} will follow from the following proposition:
\end{env}

\begin{env}[19.5.5.4]
\label{IV.19.5.5.4}
\emph{For all $p\geqslant 0$ and $q\geqslant 1$, the condition $(\mathrm{A}_{q+1,p})$ is implied by the conjunction of the conditions $(\mathrm{A}_{q,p})$, $(\mathrm{I}_{q-1,p+1})$, $(\mathrm{I}_{q-2,p+2})$, $\ldots$, $(\mathrm{I}_{0,p+q})$.}

Indeed, once this proposition is proved, the conditions $(\mathrm{I}_{q,p})$ supposed satisfied for $p+q<r$ will imply, for each $p\leqslant r$, the condition $(\mathrm{A}_{q,p})$ for $1\leqslant q<r-p$, by induction on $q$.

Let us prove \sref{IV.19.5.5.4}. Let us note that the condition $(\mathrm{A}_{q,p})$ is also equivalent to
\end{env}

\begin{env}[19.5.5.5]
\label{IV.19.5.5.5}
\[
  \mathfrak{J}^q\mathrm{M}_p\cap\mathrm{M}_{p+1}\subset\textstyle\sum_{i=1}^{q}\mathfrak{J}^{q-i}\mathrm{M}_{p+i}.
\]

Let therefore $m\in\mathfrak{J}^{q+1}\mathrm{M}_p\cap\mathrm{M}_{p+1}\subset\mathfrak{J}^q\mathrm{M}_p\cap\mathrm{M}_{p+1}$. Applying condition $(\mathrm{A}_{q,p})$, we see that for each $i$ such that $1\leqslant i\leqslant q$, there exists a family $(y_\mathbf{q}^{p+i})_{|\mathbf{q}|=q-i}$ of elements of $\mathrm{M}_{p+i}$ such that
\[
  m=\textstyle\sum_{i=1}^{q}\bigl(\textstyle\sum_{|\mathbf{q}|=q-i}\mathbf{f}^\mathbf{q} y_\mathbf{q}^{p+i}\bigr).
\]
Suppose that for $1\leqslant i<j$ ($j\geqslant 1$), we have proved that
\[
  y_\mathbf{q}^{p+i}\in\mathfrak{J}\mathrm{M}_{p+i}+\mathrm{M}_{p+i+1}\qquad(\text{for }|\mathbf{q}|=q-i).
\]
We deduce by definition \sref{IV.19.5.4}
\[
  m-\textstyle\sum_{i<j}\bigl(\textstyle\sum_{|\mathbf{q}|=q-i}\mathbf{f}^\mathbf{q} y_\mathbf{q}^{p+i}\bigr)\in\mathrm{M}_{p+q+1}',
\]
Since by hypothesis $(\mathrm{I}_{q-j,p+j})$ is satisfied, we deduce
\[
  y_\mathbf{q}^{p+j}\in\mathfrak{J}\mathrm{M}_{p+j}+\mathrm{M}_{p+j+1}
\]
\oldpage[IV-4]{209}
and, by induction on $j$, one sees therefore that this condition is satisfied for \emph{every} $j$ such that $1\leqslant j\leqslant q$. We have therefore
\[
  m\in\textstyle\sum_{i=1}^{q}\mathfrak{J}^{q-i}(\mathfrak{J}\mathrm{M}_{p+i}+\mathrm{M}_{p+i+1})=\textstyle\sum_{i=1}^{q+1}\mathfrak{J}^{q-i+1}\mathrm{M}_{p+i}
\]
and we have thus proved $(\mathrm{A}_{q+1,p})$ and completed the demonstration of \sref{IV.19.5.5}.
\end{env}
\end{proof}

\begin{remarks}[19.5.6]
\label{IV.19.5.6}
\medskip\noindent
\begin{enumerate}
\item[\rm(i)] The results of this number remain valid when, instead of an $\mathrm{A}$-module $\mathrm{M}$ and elements $f_i\in\mathrm{A}$, one takes an object $\mathrm{M}$ of any abelian category, and $n$ endomorphisms $f_i$ of $\mathrm{M}$ that mutually commute; the demonstrations adapt without difficulty to this more general case. The hypothesis of separation for an $\mathrm{A}$-module $\mathrm{N}$ relative to the $\mathfrak{J}$-preadic topology must here be replaced by the hypothesis that every sub-object of $\mathrm{N}$ contained in all the $\sum f_i^r(\mathrm{N})$ ($r$ an arbitrary integer) is necessarily zero. The same remark applies to the results of \sref{IV.19.7}.
\item[\rm(ii)] Suppose that $\mathrm{A}$ is a \emph{graded} ring in degrees $\geq0$. If $\mathrm{M}$ (resp. each $\operatorname{gr}_p(\mathrm{M})$) is a graded $\mathrm{A}$-module with degrees bounded below, and if the $f_i$ contain no homogeneous component of degree $0$, the separation hypothesis of $(\mathbf{0},15.1.9)$ is \emph{ipso facto} satisfied for $\mathrm{M}$ (resp. each $\operatorname{gr}_p(\mathrm{M})$), and one sees that in this case we have the implication $c)\Rightarrow a)$ in \sref{IV.19.5.1} (resp. \sref{IV.19.5.5}).
\item[\rm(iii)] Recall that, without the separation hypothesis, the implication $c)\Rightarrow b)$ is no longer valid even for $n=1$ $(\mathbf{0},15.1.12,\text{ (iii)})$.
\end{enumerate}
\end{remarks}

\subsection{Regular sequences relative to a filtered quotient module}
\label{subsection:IV.19.6}

\begin{env}[19.6.1]
\label{IV.19.6.1}
Let $\mathrm{A}$ be a ring, $\mathbf{f}=(f_1,\ldots,f_n)$ a finite sequence of elements of $\mathrm{A}$, $\mathfrak{J}=f_1\mathrm{A}+\ldots+f_n\mathrm{A}$ the ideal it generates, and consider an exact sequence of $\mathrm{A}$-modules
\[
  0\to\mathrm{R}\xrightarrow{j}\mathrm{N}\xrightarrow{p}\mathrm{M}\to 0
\]
where one supposes that $\mathrm{N}$ is endowed with a decreasing filtration $(\mathrm{N}_k)$ of sub-$\mathrm{A}$-modules, with $\mathrm{N}_0=\mathrm{N}$. We set $\mathrm{R}_k=\mathrm{R}\cap\mathrm{N}_k$ (filtration induced by $(\mathrm{N}_k)$), $\mathrm{M}_k=p(\mathrm{N}_k)$ (quotient filtration of $(\mathrm{N}_k)$) and we denote by $\operatorname{gr}_\bullet(\mathrm{N})$, $\operatorname{gr}_\bullet(\mathrm{R})$ and $\operatorname{gr}_\bullet(\mathrm{M})$ the graded $\mathrm{A}$-modules associated with these three filtrations. We set on the other hand for all $k$,
\[
  \mathrm{N}_k'=\mathrm{N}_k+\mathfrak{J}\mathrm{N}_{k-1}+\ldots+\mathfrak{J}^{k-1}\mathrm{N}_1+\mathfrak{J}^k\mathrm{N}_0
\]
so that $\mathrm{M}_k'=p(\mathrm{N}_k')=\mathrm{M}_k+\mathfrak{J}\mathrm{M}_{k-1}+\ldots+\mathfrak{J}^{k-1}\mathrm{M}_1+\mathfrak{J}^k\mathrm{M}_0$; we set on the other hand $\mathrm{R}_k'=\mathrm{R}\cap\mathrm{N}_k'$ (filtration \emph{induced} by $(\mathrm{N}_k')$), and we denote by $\operatorname{gr}_\bullet'(\mathrm{N})$, $\operatorname{gr}_\bullet'(\mathrm{M})$ and $\operatorname{gr}_\bullet'(\mathrm{R})$ the graded $\mathrm{A}$-modules associated with these three filtrations; we thus have a commutative diagram of exact sequences
\[
\xymatrix{
  0 \ar[r] & \operatorname{gr}_\bullet(\mathrm{R}) \ar[r]^-{\operatorname{gr}(j)} \ar[d] & \operatorname{gr}_\bullet(\mathrm{N}) \ar[r]^-{\operatorname{gr}(p)} \ar[d] & \operatorname{gr}_\bullet(\mathrm{M}) \ar[r] \ar[d] & 0 \\
  0 \ar[r] & \operatorname{gr}_\bullet'(\mathrm{R}) \ar[r]^-{\operatorname{gr}'(j)} & \operatorname{gr}_\bullet'(\mathrm{N}) \ar[r]^-{\operatorname{gr}'(p)} & \operatorname{gr}_\bullet'(\mathrm{M}) \ar[r] & 0
}
\tag{19.6.1.1}
\label{IV.19.6.1.1}
\]
\oldpage[IV-4]{210}
where the vertical arrows are the canonical homomorphisms from a graded module associated with a filtration to a graded module associated with a coarser filtration.

We note that if one sets
\[
  \mathrm{R}_k''=\mathrm{R}_k+\mathfrak{J}\mathrm{R}_{k-1}+\ldots+\mathfrak{J}^{k-1}\mathrm{R}_1+\mathfrak{J}^k\mathrm{R}_0
\]
we evidently have $\mathrm{R}_k''\subset\mathrm{R}_k'$, but the filtrations $(\mathrm{R}_k'')$ and $(\mathrm{R}_k')$ are in general distinct; we denote by $\operatorname{gr}_\bullet''(\mathrm{R})$ the graded $\mathrm{A}$-module associated with the filtration $(\mathrm{R}_k'')$.
\end{env}

\begin{env}[19.6.2]
\label{IV.19.6.2}
We defined in $(\mathbf{0},15.1.5.2)$ the surjective graded homomorphisms of degree $0$
\begin{align*}
  \psi_\mathrm{N} &: (\operatorname{gr}_\bullet(\mathrm{N})\otimes_\mathrm{A}(\mathrm{A}/\mathfrak{J}))[\mathrm{T}_1,\ldots,\mathrm{T}_n] \to \operatorname{gr}_\bullet'(\mathrm{N}) \\
  \psi_\mathrm{M} &: (\operatorname{gr}_\bullet(\mathrm{M})\otimes_\mathrm{A}(\mathrm{A}/\mathfrak{J}))[\mathrm{T}_1,\ldots,\mathrm{T}_n] \to \operatorname{gr}_\bullet'(\mathrm{M}) \\
  \psi_\mathrm{R} &: (\operatorname{gr}_\bullet(\mathrm{R})\otimes_\mathrm{A}(\mathrm{A}/\mathfrak{J}))[\mathrm{T}_1,\ldots,\mathrm{T}_n] \to \operatorname{gr}_\bullet''(\mathrm{R})
\end{align*}
the first two of which are deduced from the last two vertical arrows of \sref{IV.19.6.1.1}.

As the filtration $(\mathrm{R}_k'')$ is finer than $(\mathrm{R}_k')$, on $\mathrm{R}$ we have a canonical homomorphism $\gamma:\operatorname{gr}_\bullet''(\mathrm{R})\to\operatorname{gr}_\bullet'(\mathrm{R})$; we will denote by $\psi_\mathrm{R}'$ the composite homomorphism $\gamma\circ\psi_\mathrm{R}$. It then follows immediately from the definitions that we have a commutative diagram
\[
\xymatrix{
  & 0 \ar[d] \\
  (\operatorname{gr}_\bullet(\mathrm{R})\otimes_\mathrm{A}(\mathrm{A}/\mathfrak{J}))[\mathrm{T}_1,\ldots,\mathrm{T}_n] \ar[r]^-{\psi_\mathrm{R}'} \ar[d]^{j'} & \operatorname{gr}_\bullet'(\mathrm{R}) \ar[d]^{\operatorname{gr}'(j)} \\
  (\operatorname{gr}_\bullet(\mathrm{N})\otimes_\mathrm{A}(\mathrm{A}/\mathfrak{J}))[\mathrm{T}_1,\ldots,\mathrm{T}_n] \ar[r]^-{\psi_\mathrm{N}} \ar[d]^{p'} & \operatorname{gr}_\bullet'(\mathrm{N}) \ar[d]^{\operatorname{gr}'(p)} \\
  (\operatorname{gr}_\bullet(\mathrm{M})\otimes_\mathrm{A}(\mathrm{A}/\mathfrak{J}))[\mathrm{T}_1,\ldots,\mathrm{T}_n] \ar[r]^-{\psi_\mathrm{M}} \ar[d] & \operatorname{gr}_\bullet'(\mathrm{M}) \ar[d] \\
  0 & 0
}
\tag{19.6.2.1}
\label{IV.19.6.2.1}
\]
where the columns are exact.
\end{env}

\begin{proposition}[19.6.3]
\label{IV.19.6.3}
With the preceding notations, suppose that the sequence $\mathbf{f}$ is $\operatorname{gr}_\bullet(\mathrm{N})$-regular. Consider the following properties:
\begin{enumerate}
\item[\rm a)] $\mathbf{f}$ is $\operatorname{gr}_\bullet(\mathrm{M})$-regular.
\item[\rm b)] $\psi_\mathrm{M}$ is bijective.
\item[\rm c)] $\psi_\mathrm{R}'$ is surjective (in other words $\operatorname{gr}_\bullet'(\mathrm{R})$ is generated as $\operatorname{gr}_\mathfrak{J}^*(\mathrm{A})$-module by $\operatorname{Im}(\operatorname{gr}_\bullet(\mathrm{R})\to\operatorname{gr}_\bullet'(\mathrm{R})))$.
\item[\rm d)] $\psi_\mathrm{R}'$ is injective.
\item[\rm d$'$)] The homomorphism $(\psi_\mathrm{R}')_0:\operatorname{gr}_\bullet(\mathrm{R})\otimes_\mathrm{A}(\mathrm{A}/\mathfrak{J})\to\operatorname{gr}_\bullet'(\mathrm{R})$ obtained by restricting $\psi_\mathrm{R}'$ to polynomials of degree $0$, is injective.
\item[\rm e)] $\psi_\mathrm{R}'$ is bijective.
\item[\rm f)] $j'$ is injective.
\oldpage[IV-4]{211}
\item[\rm f$'$)] The homomorphism $\operatorname{gr}_\bullet(\mathrm{R})\otimes_\mathrm{A}(\mathrm{A}/\mathfrak{J})\to\operatorname{gr}_\bullet(\mathrm{N})\otimes_\mathrm{A}(\mathrm{A}/\mathfrak{J})$ obtained by restricting $j'$ to polynomials of degree $0$, is injective.
\item[\rm g)] The $\mathfrak{J}$-preadic filtration of $\operatorname{gr}_\bullet(\mathrm{R})$ is induced by that of $\operatorname{gr}_\bullet(\mathrm{N})$.
\item[\rm h)] The filtrations $(\mathrm{R}_k')$ and $(\mathrm{R}_k'')$ on $\mathrm{R}$ are identical.
\end{enumerate}
Then we have the following implications:
\[
\xymatrix@C=2.4em@R=1.8em{
  \text{a)} \ar@{=>}[r] \ar@{=>}[dr] &
  \text{e)} \ar@{=>}[r] \ar@{=>}[dr] &
  \text{b)} \ar@{<=>}[r] &
  \text{c)} \ar@{<=>}[r] &
  \text{h)} \\
  & \text{g)} \ar@{=>}[r] &
  \text{d)} \ar@{<=>}[r] &
  \text{d$'$)} \ar@{<=>}[r] &
  \text{f)} \ar@{<=>}[r] &
  \text{f$'$)}
}
\]

Moreover, when every quotient module of a submodule of $\operatorname{gr}_k(\mathrm{M})$ is separated for the $\mathfrak{J}$-preadic topology (which is the case when $\mathrm{A}$ is Noetherian, the $f_i$ belong to the radical of $\mathrm{A}$, and the $\operatorname{gr}_k(\mathrm{M})$ are finitely generated $\mathrm{A}$-modules), \emph{then all conditions} \emph{a)} to \emph{h)} \emph{are equivalent}.
\end{proposition}

\begin{proof}
Let us note that, from \sref{IV.19.5.5}, the hypothesis that $\mathbf{f}$ is $\operatorname{gr}_\bullet(\mathrm{N})$-regular implies that $\psi_\mathrm{N}$ is \emph{bijective}; as $\operatorname{gr}'(j)$ is injective, the equivalence of $d)$ and $f)$ results from the diagram \sref{IV.19.6.2.1}. Likewise $d')$ and $f')$ are equivalent, while $f)$ and $f')$ are trivially equivalent; this proves the equivalence of $d)$, $d')$, $f)$ and $f')$.

As one already knows that $\psi_\mathrm{M}$ is surjective and $\psi_\mathrm{N}$ bijective, the equivalence of $b)$ and $c)$ also follows from the diagram \sref{IV.19.6.2.1} by a particular case of the five lemma.

The relation $\psi_\mathrm{R}'=\gamma\circ\psi_\mathrm{R}$ \sref{IV.19.6.2} and the fact that $\psi_\mathrm{R}$ is surjective show that condition $c)$ is equivalent to saying that $\gamma$ is surjective, which is also equivalent to condition $h)$; hence we have proved the equivalence of $b)$, $c)$ and $h)$.

It is trivial that $e)$ is equivalent to the conjunction of $c)$ and $d)$, and hence also of $b)$ and $d)$. Let us note now that $a)$ implies $b)$ and that $b)$ implies $a)$ under the separation hypothesis \sref{IV.19.5.5}. On the other hand, the implications $a)\Rightarrow g)\Rightarrow f')$, and the implication $f')\Rightarrow a)$ under the separation hypothesis, result from \sref{IV.19.5.3}. Hence we have proved $a)$ implies $e)$ (equivalent to the conjunction of $b)$ and $f')$), and also that all conditions are equivalent under the separation hypothesis.
\end{proof}

\begin{corollary}[19.6.4]
\label{IV.19.6.4}
Let $\mathfrak{R}$ be an ideal of $\mathrm{A}$, $\mathfrak{L}=\mathfrak{J}+\mathfrak{R}$. Suppose that the filtration $(\mathrm{N}_k)$ (hence also $(\mathrm{M}_k)$) is the $\mathfrak{R}$-preadic filtration, so that the filtrations $(\mathrm{M}_k')$ and $(\mathrm{N}_k')$ are $\mathfrak{L}$-preadic filtrations. Consider $\operatorname{gr}_\bullet(\mathrm{R})$ (resp. $\operatorname{gr}_\bullet'(\mathrm{R})$) as a graded module over $\operatorname{gr}_\mathfrak{R}^*(\mathrm{A})$ (resp. $\operatorname{gr}_\mathfrak{L}^*(\mathrm{A})$).
\begin{enumerate}
\item[\rm(i)] Let $\mathrm{S}$ be a part of $\operatorname{gr}_\bullet(\mathrm{R})$ that generates $\operatorname{gr}_\bullet(\mathrm{R})$ as a $\operatorname{gr}_\mathfrak{R}^*(\mathrm{A})$-module. Then condition \emph{c)} of \emph{(\sref{IV.19.6.3})} is equivalent to the following condition:
\item[\rm c$'$)] The image of $\mathrm{S}$ in $\operatorname{gr}_\bullet'(\mathrm{R})$ by $\psi_\mathrm{R}'$ generates $\operatorname{gr}_\bullet'(\mathrm{R})$ as a $\operatorname{gr}_\mathfrak{L}^*(\mathrm{A})$-module.
\item[\rm(ii)] Suppose condition \emph{e)} of \emph{(\sref{IV.19.6.3})} is satisfied, and also that the quotients of $\operatorname{gr}_k(\mathrm{R})$ are separated for the $\mathfrak{J}$-preadic topology (which is the case when $\mathrm{A}$ is Noetherian, the $f_i$ belong to the radical of $\mathrm{A}$, and the $\operatorname{gr}_k(\mathrm{N})$ are finitely generated $\mathrm{A}$-modules). Then, \emph{for a part $\mathrm{S}$ of $\operatorname{gr}_\bullet(\mathrm{R})$ formed of homogeneous elements to generate $\operatorname{gr}_\bullet(\mathrm{R})$ as a $\operatorname{gr}_\mathfrak{R}^*(\mathrm{A})$-module, it is necessary and sufficient that its image by $\psi_\mathrm{R}'$ generates $\operatorname{gr}_\bullet'(\mathrm{R})$ as a $\operatorname{gr}_\mathfrak{L}^*(\mathrm{A})$-module.}
\end{enumerate}
\end{corollary}

\begin{proof}
\oldpage[IV-4]{212}
\begin{enumerate}
\item[\rm(i)] Consider the homomorphism of degree $0$
\[
  \psi_\mathrm{A}:(\operatorname{gr}_\mathfrak{R}^*(\mathrm{A})\otimes_\mathrm{A}(\mathrm{A}/\mathfrak{J}))[\mathrm{T}_1,\ldots,\mathrm{T}_n]\to\operatorname{gr}_\bullet'(\mathrm{A})=\operatorname{gr}_\mathfrak{L}^*(\mathrm{A})
\]
$(\mathbf{0},15.1.5.2)$. The fact that this homomorphism is \emph{surjective} implies that the sub-$\operatorname{gr}_\mathfrak{L}^*(\mathrm{A})$-module of $\operatorname{gr}_\bullet'(\mathrm{R})$ generated by $\mathrm{S}'=\psi_\mathrm{R}'(\mathrm{S})$ is none other than $\operatorname{Im}(\psi_\mathrm{R}')$, whence the equivalence of $c)$ and $c')$.

\item[\rm(ii)] Since condition $e)$ is satisfied, it follows by (i) that it remains to prove the sufficiency of the condition in the statement. Now, as $\psi_\mathrm{R}'$ is bijective, to say that $\mathrm{S}'$ generates $\operatorname{gr}_\bullet'(\mathrm{R})$ considered as a $\operatorname{gr}_\mathfrak{L}^*(\mathrm{A})$-module amounts to saying that $\mathrm{S}$ generates $(\operatorname{gr}_\bullet(\mathrm{R})\otimes_\mathrm{A}(\mathrm{A}/\mathfrak{J}))[\mathrm{T}_1,\ldots,\mathrm{T}_n]$ as a $\operatorname{gr}_\mathfrak{L}^*(\mathrm{A})$-module, and by the surjectivity of $\psi_\mathrm{A}$, we can in this statement replace the ring $\operatorname{gr}_\mathfrak{L}^*(\mathrm{A})$ by $(\operatorname{gr}_\mathfrak{R}^*(\mathrm{A})\otimes_\mathrm{A}(\mathrm{A}/\mathfrak{J}))[\mathrm{T}_1,\ldots,\mathrm{T}_n]$. It amounts to saying that $\mathrm{S}$ generates $\operatorname{gr}_\bullet(\mathrm{R})\otimes_\mathrm{A}(\mathrm{A}/\mathfrak{J})$ as a module over the ring $\operatorname{gr}_\mathfrak{R}^*(\mathrm{A})\otimes_\mathrm{A}(\mathrm{A}/\mathfrak{J})$, or also as a $\operatorname{gr}_\mathfrak{R}^*(\mathrm{A})$-module. Now, if $\mathrm{T}$ is the graded sub-$\operatorname{gr}_\mathfrak{R}^*(\mathrm{A})$-module of $\operatorname{gr}_\bullet(\mathrm{R})$ generated by $\mathrm{S}$, and if we set $\mathrm{T}'=\operatorname{gr}_\bullet(\mathrm{R})/\mathrm{T}$, then by hypothesis $\mathrm{T}'/\mathfrak{J}\mathrm{T}'=0$, or $\mathfrak{J}\mathrm{T}'=\mathrm{T}'$; but $\mathrm{T}'$ is a graded module, and each $\mathrm{T}_p'$ is a quotient of a submodule of $\operatorname{gr}_p(\mathrm{R})$, hence separated for the $\mathfrak{J}$-preadic topology. The hypothesis therefore implies $\mathrm{T}'=0$, which completes the proof of the corollary.
\end{enumerate}
\end{proof}

\subsection{Hironaka's normal flatness criterion}
\label{subsection:IV.19.7}
\label{IV.19.7}

\begin{theorem}[19.7.1]
\label{IV.19.7.1}
\textup{(Hironaka).}
Let $\mathrm{A}$ be a ring, $\mathfrak{R}$ an ideal of $\mathrm{A}$, $\mathbf{f}=(f_1,\ldots,f_n)$ a sequence of elements of $\mathrm{A}$ that is $(\mathrm{A}/\mathfrak{R})$-regular, $\mathfrak{J}=f_1\mathrm{A}+\ldots+f_n\mathrm{A}$ the ideal generated by $\mathbf{f}$, $\mathfrak{L}=\mathfrak{J}+\mathfrak{R}$, $\mathrm{M}$ an $\mathrm{A}$-module.
In this case $(\operatorname{gr}_\mathfrak{R}^*(\mathrm{M}))\otimes_\mathrm{A}(\mathrm{A}/\mathfrak{J})=(\operatorname{gr}_\mathfrak{R}^*(\mathrm{M}))\otimes_\mathrm{A}(\mathrm{A}/\mathfrak{L})=(\operatorname{gr}_\mathfrak{R}^*(\mathrm{M}))\otimes_{\mathrm{A}/\mathfrak{R}}(\mathrm{A}/\mathfrak{L})$, since $(\mathfrak{R}^n\mathrm{M}/\mathfrak{R}^{n+1}\mathrm{M})\otimes_\mathrm{A}(\mathrm{A}/\mathfrak{J})=\mathfrak{R}^n\mathrm{M}/(\mathfrak{R}^{n+1}+\mathfrak{J}\mathfrak{R}^n)\mathrm{M}=\mathfrak{R}^n\mathrm{M}/(\mathfrak{R}^{n+1}+\mathfrak{L}\mathfrak{R}^n)\mathrm{M}$, and the last equality is a result of Bourbaki, \emph{Alg.}, chap.~II, $3^e$~ed., \textsection~3, n\textsuperscript{o}~7, cor.~2 of prop.~6.

Consider the following conditions:
\begin{enumerate}
\item[\rm a)] $\operatorname{gr}_\mathfrak{R}^*(\mathrm{M})$ is a flat $(\mathrm{A}/\mathfrak{R})$-module.
\item[\rm b)] $\operatorname{gr}_\mathfrak{L}^*(\mathrm{M})$ is a flat $(\mathrm{A}/\mathfrak{L})$-module and the canonical homomorphism $(\mathbf{0},15.1.5.2)$
\[
  \psi_\mathrm{M}:((\operatorname{gr}_\mathfrak{R}^*(\mathrm{M}))\otimes_\mathrm{A}(\mathrm{A}/\mathfrak{L}))[\mathrm{T}_1,\ldots,\mathrm{T}_n]\to\operatorname{gr}_\mathfrak{L}^*(\mathrm{M})
\]
is bijective.
\item[\rm c)] $(\operatorname{gr}_\mathfrak{R}^*(\mathrm{M}))\otimes_{\mathrm{A}/\mathfrak{R}}(\mathrm{A}/\mathfrak{L})$ is a flat $(\mathrm{A}/\mathfrak{L})$-module and the sequence $\mathbf{f}$ is $(\operatorname{gr}_\mathfrak{R}^*(\mathrm{M}))$-regular.
\item[\rm d)] $\operatorname{gr}_\mathfrak{L}^*(\mathrm{M})$ is a flat $(\mathrm{A}/\mathfrak{L})$-module, and for every $\mathfrak{p}\in\operatorname{Ass}_{\mathrm{A}/\mathfrak{R}}(\mathrm{A}/\mathfrak{L})$, $(\operatorname{gr}_\mathfrak{R}^*(\mathrm{M}))_\mathfrak{p}$ is a flat $(\mathrm{A}/\mathfrak{R})_\mathfrak{p}$-module.
\end{enumerate}
Then we have the implications
\[
\xymatrix@C=3em@R=1.6em{
  \text{a)} \ar@{=>}[r] \ar@{=>}[dr] & \text{c)} \ar@{=>}[r] & \text{b)} \\
  & \text{d)}
}
\]

When every quotient of a submodule of $\operatorname{gr}_\mathfrak{R}^k(\mathrm{M})$ $(k\geqslant 0)$ is ideally separated for $\mathfrak{J}$ (Bourbaki, \emph{Alg.~comm.}, chap.~III, \textsection~5, n\textsuperscript{o}~1), the conditions \emph{a)}, \emph{b)}, \emph{c)} are equivalent.

\oldpage[IV-4]{213}
Finally, when $\mathrm{A}$ is Noetherian, the $f_i$ belong to the radical of $\mathrm{A}$, the sequence $\mathbf{f}$ is $\operatorname{gr}_\mathfrak{R}^*(\mathrm{A})$-regular, and $\mathrm{M}$ is a finitely generated $\mathrm{A}$-module, the conditions \emph{a)}, \emph{b)}, \emph{c)}, \emph{d)} are equivalent.
\end{theorem}

\begin{proof}
The implication $a)\Rightarrow c)$ is immediate $(\mathbf{0},15.1.13)$. If $c)$ is satisfied, it results from \sref{IV.19.5.5} that $\psi_\mathrm{M}$ is bijective; moreover, $(\operatorname{gr}_\mathfrak{R}^*(\mathrm{M})\otimes_{\mathrm{A}/\mathfrak{R}}(\mathrm{A}/\mathfrak{L}))[\mathrm{T}_1,\ldots,\mathrm{T}_n]$ is a flat $(\mathrm{A}/\mathfrak{L})$-module, and hence so is $\operatorname{gr}_\mathfrak{L}^*(\mathrm{M})$, since $\psi_\mathrm{M}$ is an $(\mathrm{A}/\mathfrak{L})$-isomorphism; thus $c)$ implies $b)$. The fact that $a)$ also implies $d)$ is immediate.

When every quotient of a submodule of $\operatorname{gr}_\mathfrak{R}^k(\mathrm{M})$ is ideally separated for $\mathfrak{J}$, it results from \sref{IV.19.5.5} that condition $b)$ implies that $\mathbf{f}$ is $(\operatorname{gr}_\mathfrak{R}^*(\mathrm{M}))$-regular; moreover, $(\operatorname{gr}_\mathfrak{R}^*(\mathrm{M})\otimes_{\mathrm{A}/\mathfrak{R}}(\mathrm{A}/\mathfrak{L}))[\mathrm{T}_1,\ldots,\mathrm{T}_n]$ is isomorphic to $\operatorname{gr}_\mathfrak{L}^*(\mathrm{M})$ and hence is flat over $\mathrm{A}/\mathfrak{L}$. Its degree-zero direct summand $\operatorname{gr}_\mathfrak{R}^*(\mathrm{M})\otimes_{\mathrm{A}/\mathfrak{R}}(\mathrm{A}/\mathfrak{L})$ is therefore flat; this proves that $b)$ implies $c)$. On the other hand, under the same hypothesis, $c)$ implies $a)$ by virtue of $(\mathbf{0},15.1.21)$ (where we can replace the Noetherian hypothesis by the hypothesis that $\mathrm{M}$ is ideally separated for $\mathfrak{J}$, by virtue of $(\mathbf{0}_\mathrm{III},10.2.1)$).

It remains therefore to prove that when $\mathrm{A}$ is Noetherian, $\mathrm{M}$ of finite type, and the $f_i$ belong to the radical of $\mathrm{A}$, $d)$ implies $a)$.

There exists a free $\mathrm{A}$-module $\mathrm{N}$ of \emph{finite type} and an exact sequence $0\to\mathrm{R}\to\mathrm{N}\to\mathrm{M}\to 0$. It suffices to prove that $d)$ implies $b)$, since every quotient of a sub-$\operatorname{gr}_\mathfrak{R}^k(\mathrm{M})$ is then ideally separated for $\mathfrak{J}$, by virtue of the fact that $\operatorname{gr}_\mathfrak{R}^k(\mathrm{M})$ is a finitely generated $\mathrm{A}$-module, that $\mathfrak{J}$ is contained in the radical of $\mathrm{A}$, and that $\mathrm{A}$ is Noetherian (Bourbaki, \emph{Alg.~comm.}, chap.~III, \textsection~5, n\textsuperscript{o}~1); in other words, it suffices to show that $\psi_\mathrm{M}$ is \emph{bijective}.

As $\mathrm{N}$ is free and the sequence $\mathbf{f}$ is $\operatorname{gr}_\mathfrak{R}^*(\mathrm{A})$-regular by hypothesis, it is also $\operatorname{gr}_\mathfrak{R}^*(\mathrm{N})$-regular. One can therefore apply \sref{IV.19.6.3}, since every quotient module of a sub-$\operatorname{gr}_k(\mathrm{M})$ is separated for the $\mathfrak{J}$-preadic topology, the $f_i$ belonging to the radical of $\mathrm{A}$ by hypothesis. The diagram \sref{IV.19.6.2.1} here reads
\[
\xymatrix{
  & 0 \ar[d] \\
  (\operatorname{gr}_\mathfrak{R}^*(\mathrm{R},\mathrm{N})\otimes_{\mathrm{A}/\mathfrak{R}}(\mathrm{A}/\mathfrak{L}))[\mathrm{T}_1,\ldots,\mathrm{T}_n] \ar[r]^-{\psi_\mathrm{R}'} \ar[d] & \operatorname{gr}_\mathfrak{L}^*(\mathrm{R},\mathrm{N}) \ar[d] \\
  (\operatorname{gr}_\mathfrak{R}^*(\mathrm{N})\otimes_{\mathrm{A}/\mathfrak{R}}(\mathrm{A}/\mathfrak{L}))[\mathrm{T}_1,\ldots,\mathrm{T}_n] \ar[r]^-{\psi_\mathrm{N}} \ar[d] & \operatorname{gr}_\mathfrak{L}^*(\mathrm{N}) \ar[d] \\
  (\operatorname{gr}_\mathfrak{R}^*(\mathrm{M})\otimes_{\mathrm{A}/\mathfrak{R}}(\mathrm{A}/\mathfrak{L}))[\mathrm{T}_1,\ldots,\mathrm{T}_n] \ar[r]^-{\psi_\mathrm{M}} \ar[d] & \operatorname{gr}_\mathfrak{L}^*(\mathrm{M}) \ar[d] \\
  0 & 0
}
\tag{19.7.1.1}
\label{IV.19.7.1.1}
\]
where $\operatorname{gr}_\mathfrak{R}^*(\mathrm{R},\mathrm{N})$ (resp. $\operatorname{gr}_\mathfrak{L}^*(\mathrm{R},\mathrm{N})$) is the graded associated to $\mathrm{R}$ for the filtration $\mathrm{R}\cap\mathfrak{R}^m\mathrm{N}$ (resp. $\mathrm{R}\cap\mathfrak{L}^m\mathrm{N}$); recall that in this diagram the columns are exact and that $\psi_\mathrm{N}$ is \emph{bijective}.
\oldpage[IV-4]{214}

\begin{env}[19.7.1.2]
\label{IV.19.7.1.2}
Set $\mathrm{B}=\operatorname{gr}_\mathfrak{R}^*\mathrm{A}\otimes_{\mathrm{A}/\mathfrak{R}}(\mathrm{A}/\mathfrak{L})$, $\mathrm{P}=\operatorname{gr}_\mathfrak{R}^*(\mathrm{N})\otimes_{\mathrm{A}/\mathfrak{R}}(\mathrm{A}/\mathfrak{L})$, $\mathrm{Q}=\operatorname{gr}_\mathfrak{L}^*(\mathrm{R},\mathrm{N})$; since $\psi_\mathrm{N}$ is bijective, $\mathrm{P}[\mathrm{T}_1,\ldots,\mathrm{T}_n]$ identifies with $\operatorname{gr}_\mathfrak{L}^*(\mathrm{N})$, and $\mathrm{Q}$ is a sub-$\mathrm{B}[\mathrm{T}_1,\ldots,\mathrm{T}_n]$-module of $\mathrm{P}[\mathrm{T}_1,\ldots,\mathrm{T}_n]$; we will first show that we have
\end{env}
\begin{env}[19.7.1.3]
\label{IV.19.7.1.3}
\[
  \mathrm{Q}=(\mathrm{Q}\cap\mathrm{P})[\mathrm{T}_1,\ldots,\mathrm{T}_n].
\]
\end{env}

For this, set
\[
  \mathrm{Z}=\mathrm{Q}/((\mathrm{Q}\cap\mathrm{P})[\mathrm{T}_1,\ldots,\mathrm{T}_n]).
\]
This is a sub-$\mathrm{B}[\mathrm{T}_1,\ldots,\mathrm{T}_n]$-module of
\[
  \begin{aligned}
  &\mathrm{P}[\mathrm{T}_1,\ldots,\mathrm{T}_n]
  /((\mathrm{Q}\cap\mathrm{P})[\mathrm{T}_1,\ldots,\mathrm{T}_n])\\
  &\qquad=(\mathrm{P}/(\mathrm{Q}\cap\mathrm{P}))[\mathrm{T}_1,\ldots,\mathrm{T}_n].
  \end{aligned}
\]
As an $(\mathrm{A}/\mathfrak{L})$-module, it is therefore isomorphic to a submodule of a direct sum of $(\mathrm{A}/\mathfrak{L})$-modules isomorphic to
\[
  \mathrm{P}/(\mathrm{Q}\cap\mathrm{P})=(\mathrm{P}+\mathrm{Q})/\mathrm{Q}.
\]
But $(\mathrm{P}+\mathrm{Q})/\mathrm{Q}$ is an $(\mathrm{A}/\mathfrak{L})$-module isomorphic to a submodule of $\mathrm{P}[\mathrm{T}_1,\ldots,\mathrm{T}_n]/\mathrm{Q}$, which is none other than $\operatorname{gr}_\mathfrak{L}^*(\mathrm{M})$ by the diagram \sref{IV.19.7.1.1}.

\begin{equation}
  \operatorname{Ass}_{\mathrm{A}/\mathfrak{L}}(\mathrm{Z})\subset\operatorname{Ass}_{\mathrm{A}/\mathfrak{L}}(\operatorname{gr}_\mathfrak{L}^*(\mathrm{M})).
  \tag{19.7.1.4}\label{IV.19.7.1.4}
\end{equation}

But each $\operatorname{gr}_\mathfrak{L}^m(\mathrm{M})$ is by hypothesis a \emph{flat} $(\mathrm{A}/\mathfrak{L})$-module of finite type, hence projective since $\mathrm{A}/\mathfrak{L}$ is Noetherian, and we therefore have $\operatorname{Ass}_{\mathrm{A}/\mathfrak{L}}(\operatorname{gr}_\mathfrak{L}^*(\mathrm{M}))\subset\operatorname{Ass}_{\mathrm{A}/\mathfrak{L}}(\mathrm{A}/\mathfrak{L})$. To see that $\mathrm{Z}=0$, it suffices by (Bourbaki, \emph{Alg.~comm.}, chap.~IV, \textsection~1, n\textsuperscript{o}~1, cor.~1 of prop.~2 and n\textsuperscript{o}~3, prop.~7) to see that for every $\mathfrak{q}\in\operatorname{Ass}_{\mathrm{A}/\mathfrak{L}}(\mathrm{A}/\mathfrak{L})$, $\mathrm{Z}_\mathfrak{q}=0$. But the ideals of $\operatorname{Ass}_{\mathrm{A}/\mathfrak{L}}(\mathrm{A}/\mathfrak{L})$ are of the form $\mathfrak{p}/(\mathfrak{L}/\mathfrak{R})$, where $\mathfrak{p}\in\operatorname{Ass}_{\mathrm{A}/\mathfrak{R}}(\mathrm{A}/\mathfrak{L})$, and it comes down to the same to see that $\mathrm{Z}_\mathfrak{p}=0$ for every $\mathfrak{p}\in\operatorname{Ass}_{\mathrm{A}/\mathfrak{R}}(\mathrm{A}/\mathfrak{L})$. Now, if $\mathfrak{p}=\mathfrak{r}/\mathfrak{R}$, where $\mathfrak{r}$ is a prime ideal of $\mathrm{A}$, then $(\operatorname{gr}_\mathfrak{R}^*(\mathrm{M}))_\mathfrak{p}=\operatorname{gr}_{\mathfrak{R}_\mathfrak{r}}^*(\mathrm{M}_\mathfrak{r})$ and $(\mathrm{A}/\mathfrak{L})_\mathfrak{p}=\mathrm{A}_\mathfrak{r}/\mathfrak{L}_\mathfrak{r}$, and the images of the $f_i$ in $\mathrm{A}_\mathfrak{r}$ form a $\operatorname{gr}_{\mathfrak{R}_\mathfrak{r}}^*(\mathrm{A}_\mathfrak{r})$-regular sequence $(\mathbf{0},15.1.14)$; applying to $\mathrm{A}_\mathfrak{r}$, $\mathfrak{R}_\mathfrak{r}$ and $\mathrm{M}_\mathfrak{r}$ the implication $a)\Rightarrow b)$ of the statement, we see from the hypothesis $d)$ that $\psi_{\mathrm{M}_\mathfrak{r}}$ is bijective, hence also $\psi_{\mathrm{R}_\mathfrak{r}}'$ by virtue of \sref{IV.19.6.3}; in other words, $\mathrm{Q}_\mathfrak{r}=\mathrm{Q}_\mathfrak{r}'[\mathrm{T}_1,\ldots,\mathrm{T}_n]$, where we have set $\mathrm{Q}'=\operatorname{gr}_\mathfrak{R}^*(\mathrm{R},\mathrm{N})\otimes_{\mathrm{A}/\mathfrak{R}}(\mathrm{A}/\mathfrak{L})$, which is here a sub-$\mathrm{B}$-module of $\mathrm{P}$; in particular $\mathrm{Q}_\mathfrak{r}'=\mathrm{Q}_\mathfrak{r}\cap\mathrm{P}_\mathfrak{r}$; hence finally $\mathrm{Z}_\mathfrak{r}=\mathrm{Z}_\mathfrak{p}=0$, which completes the demonstration of \sref{IV.19.7.1.3}.

\begin{env}[19.7.1.5]
\label{IV.19.7.1.5}
By virtue of \sref{IV.19.7.1.3}, to prove that $\psi_\mathrm{R}'$ is surjective, it suffices to show that every homogeneous element of degree $m$ of $\mathrm{Q}\cap\mathrm{P}$ is the image of an element of the same degree of $\mathrm{Q}'=\operatorname{gr}_\mathfrak{R}^*(\mathrm{R},\mathrm{N})\otimes_{\mathrm{A}/\mathfrak{R}}(\mathrm{A}/\mathfrak{L})$. Now, the elements of degree $m$ of $\mathrm{P}$ identify (by $\psi_\mathrm{N}$) with the elements of $(\mathfrak{R}^m\mathrm{N}+\mathfrak{L}^{m+1}\mathrm{N})/\mathfrak{L}^{m+1}\mathrm{N}$; those which belong to $\mathrm{Q}$ identify with the elements of $((\mathrm{R}\cap\mathfrak{L}^m\mathrm{N})+\mathfrak{L}^{m+1}\mathrm{N})/\mathfrak{L}^{m+1}\mathrm{N}$; and the images of the elements of degree $m$ of $\mathrm{Q}'$ identify with the elements of $((\mathrm{R}\cap\mathfrak{R}^m\mathrm{N})+\mathfrak{L}^{m+1}\mathrm{N})/\mathfrak{L}^{m+1}\mathrm{N}$. It will therefore suffice to prove, for $m>0$, the inclusion
\[
  ((\mathrm{R}\cap\mathfrak{L}^m\mathrm{N})+\mathfrak{L}^{m+1}\mathrm{N})\cap(\mathfrak{R}^m\mathrm{N}+\mathfrak{L}^{m+1}\mathrm{N})\subset(\mathrm{R}\cap\mathfrak{R}^m\mathrm{N})+\mathfrak{L}^{m+1}\mathrm{N}.
\]
As $\mathfrak{R}\subset\mathfrak{L}$, we see immediately that the first member of this relation is equal to $(\mathrm{R}\cap(\mathfrak{R}^m\mathrm{N}+\mathfrak{L}^{m+1}\mathrm{N}))+\mathfrak{L}^{m+1}\mathrm{N}$, so that everything reduces to proving
\end{env}
\begin{env}[19.7.1.6]
\label{IV.19.7.1.6}
\[
  \mathrm{R}\cap(\mathfrak{R}^m\mathrm{N}+\mathfrak{L}^{m+1}\mathrm{N})\subset(\mathrm{R}\cap\mathfrak{R}^m\mathrm{N})+\mathfrak{L}^{m+1}\mathrm{N}.
\]
\end{env}

\oldpage[IV-4]{215}
We will proceed by induction on $m$, assuming therefore \sref{IV.19.7.1.6} verified when $m$ is replaced by an integer $m'<m$. We will consider on the other hand, for $m$ fixed and for an integer $d\geqslant 0$, the relation
\begin{equation}
  \mathrm{R}\cap(\mathfrak{R}^m\mathrm{N}+\mathfrak{L}^{m+1}\mathrm{N})\subset(\mathrm{R}\cap\mathfrak{R}^d\mathrm{N}\cap\mathfrak{L}^m\mathrm{N})+\mathfrak{L}^{m+1}\mathrm{N}.
  \tag{$*_d$}
\end{equation}

It is clear that it suffices to demonstrate $(*_d)$ for $d=m$, and on the other hand $(*_0)$ is trivially true. We will prove $(*_d)$ by \emph{induction on $d$}; in other words, we suppose, for a $d\geqslant 0$ fixed (with $d<m$), that $(*_d)$ is true, and we wish to prove
\begin{equation}
  \mathrm{R}\cap(\mathfrak{R}^m\mathrm{N}+\mathfrak{L}^{m+1}\mathrm{N})\subset(\mathrm{R}\cap\mathfrak{R}^{d+1}\mathrm{N}\cap\mathfrak{L}^m\mathrm{N})+\mathfrak{L}^{m+1}\mathrm{N}.
  \tag{$*_{d+1}$}
\end{equation}

Consider for this, for an $h\geqslant 0$, the relation
\begin{equation}
  \mathrm{R}\cap(\mathfrak{R}^m\mathrm{N}+\mathfrak{R}^d\mathfrak{L}^h\mathrm{N})\subset(\mathrm{R}\cap(\mathfrak{R}^{d+1}\mathrm{N}+\mathfrak{R}^d\mathfrak{L}^h\mathrm{N})\cap\mathfrak{L}^m\mathrm{N})+\mathfrak{L}^{m+1}\mathrm{N}.
  \tag{$**_{d+1,h}$}
\end{equation}

The hypothesis $(*_d)$ implies that $(**_{d+1,0})$ is true. We will prove by \emph{induction on $h$} that $(**_{d+1,h})$ is true for every $h\geqslant 0$. But we have

\begin{lemma}[19.7.1.7]
\label{IV.19.7.1.7}
Let $\mathrm{A}$ be a Noetherian ring, $\mathrm{N}$ a finitely generated $\mathrm{A}$-module, $\mathfrak{L}$ an ideal of $\mathrm{A}$, $\mathrm{E}$, $\mathrm{F}$ two submodules of $\mathrm{N}$. For every $k>0$ there exists $h>0$ such that $\mathrm{E}\cap(\mathrm{F}+\mathfrak{L}^h\mathrm{N})\subset(\mathrm{E}\cap\mathrm{F})+\mathfrak{L}^k\mathrm{N}$.
\end{lemma}

\begin{proof}
Indeed, let $\varphi:\mathrm{N}\to\mathrm{N}/(\mathrm{E}\cap\mathrm{F})=\mathrm{N}_1$ be the canonical homomorphism, and set $\mathrm{E}_1=\varphi(\mathrm{E})$, $\mathrm{F}_1=\varphi(\mathrm{F})$, so that $\mathrm{E}_1\cap\mathrm{F}_1=0$ and $\varphi(\mathrm{E}\cap(\mathrm{F}+\mathfrak{L}^h\mathrm{N}))=\mathrm{E}_1\cap(\mathrm{F}_1+\mathfrak{L}^h\mathrm{N}_1)$.
We know $(\mathbf{0}_\mathrm{I},7.3.2)$ that there exists $h$ such that $(\mathrm{E}_1+\mathrm{F}_1)\cap\mathfrak{L}^h\mathrm{N}_1\subset\mathfrak{L}^k(\mathrm{E}_1+\mathrm{F}_1)$; if $h$ is thus chosen and if $x_1\in\mathrm{E}_1$ is such that there exists $y_1\in\mathrm{F}_1$ for which $x_1-y_1\in\mathfrak{L}^h\mathrm{N}_1$, we deduce $x_1-y_1\in\mathfrak{L}^k\mathrm{N}_1$, hence $\mathrm{E}_1\cap(\mathrm{F}_1+\mathfrak{L}^h\mathrm{N}_1)\subset\mathfrak{L}^k\mathrm{N}_1$; therefore $\mathrm{E}\cap(\mathrm{F}+\mathfrak{L}^h\mathrm{N})\subset\mathrm{E}\cap\mathrm{F}+\mathfrak{L}^k\mathrm{N}$, which proves the lemma.
\end{proof}

It suffices then to apply this lemma with $\mathrm{E}=\mathrm{R}\cap\mathfrak{L}^m\mathrm{N}$, $\mathrm{F}=\mathfrak{R}^{d+1}\mathrm{N}$ and to take $k=m+1$ to deduce that $(**_{d+1,h})$ implies $(*_{d+1})$.

We suppose therefore in all that follows that $(**_{d+1,h})$ \emph{is satisfied.}

\begin{env}[19.7.1.8]
\label{IV.19.7.1.8}
\emph{Demonstration of $(**_{d+1,h+1})$; First case: $h\leqslant m-d$.}

Start with an element
\end{env}
\begin{env}[19.7.1.9]
\label{IV.19.7.1.9}
\[
  g\in\mathrm{R}\cap(\mathfrak{R}^{d+1}\mathrm{N}+\mathfrak{R}^d\mathfrak{L}^h\mathrm{N})\cap\mathfrak{L}^m\mathrm{N}
\]
\end{env}

congruent modulo $\mathfrak{L}^{m+1}\mathrm{N}$ to an element $g_1\in\mathrm{R}\cap(\mathfrak{R}^m\mathrm{N}+\mathfrak{L}^{m+1}\mathrm{N})$; we therefore have

\begin{env}[19.7.1.10]
\label{IV.19.7.1.10}
\[
  g\in\mathfrak{R}^m\mathrm{N}+\mathfrak{L}^{m+1}\mathrm{N}.
\]
\end{env}

It will suffice to prove that $g$ is also congruent modulo $\mathfrak{L}^{m+1}\mathrm{N}$ to an element $g_2\in\mathrm{R}\cap(\mathfrak{R}^{d+1}\mathrm{N}+\mathfrak{R}^d\mathfrak{L}^{h+1}\mathrm{N})$, for this will imply $g_2\in\mathfrak{L}^m\mathrm{N}$, since $g\in\mathfrak{L}^m\mathrm{N}$ and $g_2-g\in\mathfrak{L}^{m+1}\mathrm{N}$. In the case $h\leqslant m-d$ under consideration, it is enough to prove

\begin{env}[19.7.1.11]
\label{IV.19.7.1.11}
\[
  (\mathfrak{R}^m\mathrm{N}+\mathfrak{L}^{m+1}\mathrm{N})\cap\mathfrak{R}^d\mathrm{N}
  \subset\mathfrak{L}^{m-d+1}\mathfrak{R}^d\mathrm{N}+\mathfrak{R}^{d+1}\mathrm{N}.
\]
\end{env}
\oldpage[IV-4]{216}

Indeed, relations \sref{IV.19.7.1.9} and \sref{IV.19.7.1.10} imply that $g$ belongs to the left-hand side of \sref{IV.19.7.1.11}; and since $h\leqslant m-d$, we have $\mathfrak{L}^{m-d+1}\subset\mathfrak{L}^{h+1}$, which will establish $(**_{d+1,h+1})$. Moreover $\mathfrak{R}^m\mathrm{N}\subset\mathfrak{R}^d\mathrm{N}$, and the relation

\begin{env}[19.7.1.12]
\label{IV.19.7.1.12}
\[
  \mathfrak{L}^{m+1}\mathrm{N}\cap\mathfrak{R}^d\mathrm{N}
  \subset\mathfrak{L}^{m-d+1}\mathfrak{R}^d\mathrm{N}
\]
\end{env}

will imply \sref{IV.19.7.1.11}. Since $\mathrm{N}$ is a free $\mathrm{A}$-module of finite type, it is therefore enough to prove

\begin{env}[19.7.1.13]
\label{IV.19.7.1.13}
\[
  \mathfrak{L}^{m+1}\cap\mathfrak{R}^d\subset\mathfrak{L}^{m-d+1}\mathfrak{R}^d.
\]
\end{env}

It will suffice to prove that, in general, for $d<m$, one has

\begin{env}[19.7.1.14]
\label{IV.19.7.1.14}
\[
  \mathfrak{L}^{m+1}\cap\mathfrak{R}^d\subset\mathfrak{L}^{m-d+1}\mathfrak{R}^d+\mathfrak{R}^{d+1}.
\]
\end{env}

Indeed, since $\mathfrak{R}\subset\mathfrak{L}$, hence $\mathfrak{L}^{m-d+1}\mathfrak{R}^d\subset\mathfrak{L}^{m+1}$, the preceding relation, applied successively with $d,d+1,\ldots,m$, implies
\[
  \mathfrak{L}^{m+1}\cap\mathfrak{R}^d
  \subset\mathfrak{L}^{m-d+1}\mathfrak{R}^d+\mathfrak{R}^{m+1},
\]
whence \sref{IV.19.7.1.13}, since $\mathfrak{R}^{m+1}\subset\mathfrak{L}^{m-d+1}\mathfrak{R}^d$.

\begin{env}[19.7.1.15]
\label{IV.19.7.1.15}
As $\mathfrak{L}=\mathfrak{J}+\mathfrak{R}$, the relation \sref{IV.19.7.1.14} also reads
\[
\begin{aligned}
  &(\mathfrak{J}^{m+1}+\mathfrak{J}^m\mathfrak{R}+\cdots+\mathfrak{J}\mathfrak{R}^m+\mathfrak{R}^{m+1})\cap\mathfrak{R}^d\\
  &\qquad\subset\mathfrak{J}^{m-d+1}\mathfrak{R}^d+\mathfrak{J}^{m-d}\mathfrak{R}^{d+1}+\cdots+\mathfrak{J}\mathfrak{R}^m+\mathfrak{R}^{m+1}+\mathfrak{R}^{d+1}\\
  &\qquad=\mathfrak{J}^{m-d+1}\mathfrak{R}^d+\mathfrak{R}^{d+1}.
\end{aligned}
\]

We will show that this inclusion is itself a consequence of
\end{env}
\begin{env}[19.7.1.16]
\label{IV.19.7.1.16}
\[
  \mathfrak{R}^a\mathfrak{J}^b\cap\mathfrak{R}^{a+1}\subset\mathfrak{R}^{a+1}\mathfrak{J}^{b-1}
\]
\end{env}
valid for $a\geqslant 0$, $b\geqslant 1$. Indeed, it suffices, to prove \sref{IV.19.7.1.15}, to show that for $0<q\leqslant d$,
\begin{env}[19.7.1.17]
\label{IV.19.7.1.17}
\[
  (\mathfrak{J}^{m+1}+\mathfrak{J}^m\mathfrak{R}+\ldots+\mathfrak{J}\mathfrak{R}^m+\mathfrak{R}^{m+1})\cap\mathfrak{R}^d\subset(\mathfrak{J}^{m+1-q}\mathfrak{R}^q+\mathfrak{J}^{m-q}\mathfrak{R}^{q+1}+\ldots+\mathfrak{R}^{m+1})\cap\mathfrak{R}^d
\]
\end{env}
since for $q=d$, this will imply \sref{IV.19.7.1.15}. To prove \sref{IV.19.7.1.17}, proceed by induction on $q$, supposing the relation true for $q<d$. An element of the right-hand side corresponding to $q$ can be written $y+z$, with $y\in\mathfrak{J}^{m+1-q}\mathfrak{R}^q$ and $z\in\mathfrak{J}^{m-q}\mathfrak{R}^{q+1}+\cdots+\mathfrak{R}^{m+1}$. Since $y+z\in\mathfrak{R}^d\subset\mathfrak{R}^{q+1}$ and $z\in\mathfrak{R}^{q+1}$, we have $y\in\mathfrak{R}^{q+1}$; by \sref{IV.19.7.1.16}, $y\in\mathfrak{J}^{m-q}\mathfrak{R}^{q+1}$, and hence
\[
  y+z\in\mathfrak{J}^{m-q}\mathfrak{R}^{q+1}+\ldots+\mathfrak{R}^{m+1},\quad\text{and}\quad y+z\in\mathfrak{R}^d
\]
by hypothesis, which proves \sref{IV.19.7.1.17} where $q$ has been replaced by $q+1$.

It remains to establish \sref{IV.19.7.1.16}. The case $b=1$ is trivial; suppose $b\geq2$. An element of the left-hand side is written $a_1f_1+\cdots+a_nf_n$, with $a_i\in\mathfrak R^a\mathfrak J^{b-1}$, and its image in $Q=\mathfrak R^a/\mathfrak R^{a+1}$ is zero. Since the sequence $\mathbf f$ is $Q$-regular, the Koszul relations allow us to choose $c_{ij}\in\mathfrak R^a\mathfrak J^{b-2}$, with $c_{ij}=-c_{ji}$ and $c_{ii}=0$, such that $a_i-\sum_j c_{ij}f_j\in\mathfrak R^{a+1}$.\footnote{Correction note. The required filtration bound is proved as follows. Put $I=\mathfrak J$, $s=b-1$ and $K_\bullet=K_\bullet(\mathbf f;Q)$, with differential $\delta$. For $u\in I^sK_1$ with $\delta_1u=0$, regularity gives $u=\delta_2C$. By $(\mathbf{0},15.1.7,15.1.9)$, $\operatorname{gr}_I(Q)\simeq(Q/IQ)[T_1,\ldots,T_n]$; the Koszul complex of the indeterminates on this module is exact in positive degrees. If $C\in I^dK_2$ with $d<s-1$, its degree-$d$ class is a cycle. For $d=0$, this class is zero, since a preimage under $\delta_3$ would have coefficient degree $-1$. For $d\geq1$, it is the boundary of a class of coefficient degree $d-1$; lifting that class to $D\in I^{d-1}K_3$ and replacing $C$ by $C-\delta_3D$ gives $C\in I^{d+1}K_2$, without changing $\delta_2C$. After at most $s-1$ steps, $C\in I^{s-1}K_2$. Lifting its coefficients to $\mathfrak R^a\mathfrak J^{b-2}$ gives the required $c_{ij}$. No passage to the completion is needed. The printed text incorrectly concludes that the originally chosen coefficients necessarily belong to $\mathfrak R^{a+1}$.} Replace the $a_i$ by these new coefficients: the antisymmetric terms cancel, so the sum $a_1f_1+\cdots+a_nf_n$ is unchanged, while all its coefficients belong to $\mathfrak R^a\mathfrak J^{b-1}\cap\mathfrak R^{a+1}$. It now suffices to argue by
\oldpage[IV-4]{217}
induction on $b$ for $b\geqslant 1$ (the case $b=1$ being trivial) to establish \sref{IV.19.7.1.16} and by extension also $(**_{d+1,h+1})$ for $h\leqslant m-d$.

\begin{env}[19.7.1.18]
\label{IV.19.7.1.18}
\emph{Demonstration of $(**_{d+1,h+1})$; Second case: $h>m-d$.}

We will first prove that, for every $h\geqslant 0$, we have
\end{env}
\begin{env}[19.7.1.19]
\label{IV.19.7.1.19}
\[
  \mathrm{R}\cap(\mathfrak{R}^{d+1}\mathrm{N}+\mathfrak{R}^d\mathfrak{L}^h\mathrm{N})\subset\mathfrak{J}^h(\mathrm{R}\cap\mathfrak{R}^d\mathrm{N})+\mathfrak{L}^{h+1}\mathfrak{R}^d\mathrm{N}+\mathfrak{R}^{d+1}\mathrm{N}.
\]
\end{env}

Consider an element $z$ of the left-hand side of \sref{IV.19.7.1.19}. We note that $\mathfrak{R}^{d+1}\mathrm{N}+\mathfrak{R}^d\mathfrak{L}^h\mathrm{N}=\mathfrak{R}^{d+1}\mathrm{N}+\mathfrak{R}^d\mathfrak{J}^h\mathrm{N}$, since $\mathfrak{L}=\mathfrak{J}+\mathfrak{R}$. We consider the class $\bar{z}$ of $z$ in $\operatorname{gr}_\mathfrak{J}^{h}(\operatorname{gr}_\mathfrak{R}^{d}(\mathrm{N}))$, which, by virtue of \sref{IV.19.5.5} and the hypothesis on $\mathbf{f}$, identifies with a sub-$(\mathrm{A}/\mathfrak{L})$-module of $\operatorname{gr}_\mathfrak{L}^{h+d}(\mathrm{N})$. Since $z\in\mathrm{R}$, we will show that, under this identification, one even has

\begin{env}[19.7.1.20]
\label{IV.19.7.1.20}
\[
  \bar{z}\in\operatorname{gr}_\mathfrak{L}^{h+d}(\mathrm{R},\mathrm{N}).
\]
\end{env}

Indeed, we noted in the proof of \sref{IV.19.7.1.3} that, for every prime ideal $\mathfrak{r}$ of $\mathrm{A}$ such that $\mathfrak{p}=\mathfrak{r}/\mathfrak{R}\in\operatorname{Ass}_{\mathrm{A}/\mathfrak{R}}(\mathrm{A}/\mathfrak{L})$, the map $\psi_{\mathrm{R}_\mathfrak{r}}'$ is bijective, and consequently relation \sref{IV.19.7.1.6}, with all modules replaced by their localisations at $\mathfrak{r}$, is true. Since $z$ belongs to $\mathrm{R}\cap(\mathfrak{R}^d\mathrm{N}+\mathfrak{L}^{d+1}\mathrm{N})$, its image $z/1$ in $\mathrm{R}_\mathfrak{r}$ belongs to $(\mathrm{R}_\mathfrak{r}\cap\mathfrak{R}_\mathfrak{r}^d\mathrm{N}_\mathfrak{r})+\mathfrak{L}_\mathfrak{r}^{d+1}\mathrm{N}_\mathfrak{r}$. Hence the image $\bar{z}/1$ in $\operatorname{gr}_{\mathfrak{L}_\mathfrak{r}}^{h+d}(\mathrm{N}_\mathfrak{r})$ belongs to
\[
  \operatorname{gr}_{\mathfrak{L}_\mathfrak{r}}^{h+d}(\mathrm{R}_\mathfrak{r},\mathrm{N}_\mathfrak{r})
  =(\operatorname{gr}_\mathfrak{L}^{h+d}(\mathrm{R},\mathrm{N}))_\mathfrak{r}
  =(\operatorname{gr}_\mathfrak{L}^{h+d}(\mathrm{R},\mathrm{N}))_\mathfrak{q},
\]
where $\mathfrak{q}=\mathfrak{r}/\mathfrak{L}$. In other words, the image of $\bar{z}$ by the canonical map
\[
  (\operatorname{gr}_\mathfrak{L}^{h+d}(\mathrm{N}))_\mathfrak{q}
  \longrightarrow
  (\operatorname{gr}_\mathfrak{L}^{h+d}(\mathrm{M}))_\mathfrak{q}
\]
is zero for every $\mathfrak{q}\in\operatorname{Ass}_{\mathrm{A}/\mathfrak{L}}(\mathrm{A}/\mathfrak{L})$. As we saw in the proof of \sref{IV.19.7.1.3},
\[
  \operatorname{Ass}_{\mathrm{A}/\mathfrak{L}}(\operatorname{gr}_\mathfrak{L}^{h+d}(\mathrm{M}))\subset\operatorname{Ass}_{\mathrm{A}/\mathfrak{L}}(\mathrm{A}/\mathfrak{L}).
\]
It follows (Bourbaki, \emph{Alg.~comm.}, chap.~IV, \textsection~1, n\textsuperscript{o}~1, cor.~1 of prop.~2 and n\textsuperscript{o}~3, prop.~7) that the image of $\bar{z}$ by the canonical map $\operatorname{gr}_\mathfrak{L}^{h+d}(\mathrm{N})\to\operatorname{gr}_\mathfrak{L}^{h+d}(\mathrm{M})$ is zero, which is precisely \sref{IV.19.7.1.20}.

\begin{env}[19.7.1.21]
\label{IV.19.7.1.21}
This being so, we can write by definition
\[
  z\equiv\textstyle\sum_{|\mathbf{p}|=h}c_\mathbf{p}\mathbf{f}^\mathbf{p}\bmod{\mathfrak{R}^{d+1}\mathrm{N}}
\]
\end{env}
where one sets as usual $\mathbf{p}=(p_1,\ldots,p_n)$, $\mathbf{f}^\mathbf{p}=f_1^{p_1}f_2^{p_2}\cdots f_n^{p_n}$, $|\mathbf{p}|=p_1+\ldots+p_n$ and $c_\mathbf{p}\in\mathfrak{R}^d\mathrm{N}$. Moreover, as $\psi_\mathrm{N}$ identifies the relevant graded module with $(\operatorname{gr}_\mathfrak{R}^d(\mathrm{N})\otimes_{\mathrm{A}/\mathfrak{R}}(\mathrm{A}/\mathfrak{L}))[\mathrm{T}_1,\ldots,\mathrm{T}_n]$, the $c_\mathbf{p}$ are determined modulo $\mathfrak{R}^d\mathfrak{L}\mathrm{N}$; if $\bar c_\mathbf{p}$ denotes the class of $c_\mathbf{p}$ modulo $\mathfrak{R}^d\mathfrak{L}\mathrm{N}$, then, under the preceding identification,
\begin{env}[19.7.1.22]
\label{IV.19.7.1.22}
\[
  \bar{z}=\textstyle\sum_{|\mathbf{p}|=h}\bar{c}_\mathbf{p}\mathrm{T}^\mathbf{p}.
\]
\end{env}
Using \sref{IV.19.7.1.3}, we deduce that necessarily, for every $\mathbf{p}$ such that $|\mathbf{p}|=h$ ($\bar{c}_\mathbf{p}$ being this time identified by $\psi_\mathrm{N}$ with an element of $\operatorname{gr}_\mathfrak{R}^d(\mathrm{N})$),
\begin{env}[19.7.1.23]
\label{IV.19.7.1.23}
\[
  \bar{c}_\mathbf{p}\in\operatorname{gr}_\mathfrak{R}^d(\mathrm{R},\mathrm{N}).
\]
\end{env}

But since $d<m$, the hypothesis that $(*_d)$ is true implies that $\bar{c}_\mathbf{p}$ belongs to the image of $\psi_\mathrm{R}'$, hence we can suppose that $c_\mathbf{p}\in(\mathrm{R}\cap\mathfrak{R}^d\mathrm{N})+\mathfrak{R}^d\mathfrak{L}\mathrm{N}$. The relation \sref{IV.19.7.1.21} then gives
\[
  z\in\mathfrak{J}^h(\mathrm{R}\cap\mathfrak{R}^d\mathrm{N})+\mathfrak{J}^h\mathfrak{R}^d\mathfrak{L}\mathrm{N}+\mathfrak{R}^{d+1}\mathrm{N}
\]
and as $\mathfrak{J}\subset\mathfrak{L}$, this proves \sref{IV.19.7.1.19}.

\oldpage[IV-4]{218}
In particular, we can therefore write
\[
  g=t+g_2
\]
with $t\in\mathfrak{J}^h(\mathrm{R}\cap\mathfrak{R}^d\mathrm{N})$ and $g_2\in\mathfrak{R}^{d+1}\mathrm{N}+\mathfrak{R}^d\mathfrak{L}^{h+1}\mathrm{N}$; moreover, as $t\in\mathrm{R}$ and $g\in\mathrm{R}$, we have $g_2\in\mathrm{R}$. On the other hand, using the hypothesis $h\geqslant m-d+1$, one sees that $t\in\mathfrak{L}^{m+1}\mathrm{N}$. The element $g_2$ hence satisfies all the conditions stated at the start of \sref{IV.19.7.1.8}, which completes the demonstration of \sref{IV.19.7.1}.
\end{proof}

\begin{remarks}[19.7.2]
\label{IV.19.7.2}
\medskip\noindent
\begin{enumerate}
\item[\rm(i)] When one supposes that $\operatorname{gr}_\mathfrak{R}^*(\mathrm{A})$ is an $(\mathrm{A}/\mathfrak{R})$-module \emph{flat}, the hypothesis that the sequence $\mathbf{f}$ is $(\mathrm{A}/\mathfrak{R})$-regular implies that it is also $\operatorname{gr}_\mathfrak{R}^*(\mathrm{A})$-regular $(\mathbf{0},15.1.14)$. We note that this will be the case when $\mathrm{A}$ and $\mathrm{A}/\mathfrak{R}$ \emph{are regular rings} $(\mathbf{0},17.3.6)$: indeed, $\mathfrak{R}$ is then a regular ideal of $\mathrm{A}$ \sref{IV.19.1.2}, hence quasi-regular, and localising at the maximal ideals of $\mathrm{A}$ containing $\mathfrak{R}$ and using $(\mathbf{0},15.1.7)$, one sees that $\operatorname{gr}_\mathfrak{R}^*(\mathrm{A})$ is a \emph{projective} $(\mathrm{A}/\mathfrak{R})$-module.
\item[\rm(ii)] It would be interesting to be able to prove the last conclusion of \sref{IV.19.7.1} under the sole hypothesis that there exists an $\mathrm{A}$-algebra $\mathrm{B}$ which is a Noetherian ring, for which $\mathfrak{J}\mathrm{B}$ is contained in the radical of $\mathrm{B}$, and that $\mathrm{M}$ is a finitely generated $\mathrm{B}$-module.
\end{enumerate}
\end{remarks}

\begin{corollary}[19.7.3]
\label{IV.19.7.3}
Let $\mathrm{A}$ be a local Noetherian ring, $\mathfrak{m}$ its maximal ideal, $\mathfrak{R}$ an ideal of $\mathrm{A}$ such that the ring $\mathrm{A}/\mathfrak{R}$ is regular and of dimension $n$, $\mathrm{M}$ a finitely generated $\mathrm{A}$-module. The following conditions are equivalent:
\begin{enumerate}
\item[\rm a)] $\operatorname{gr}_\mathfrak{R}^*(\mathrm{M})$ is a flat $(\mathrm{A}/\mathfrak{R})$-module.
\item[\rm b)] If $\mathrm{P}(\mathrm{T})$ and $\mathrm{Q}(\mathrm{T})$ are the Poincar\'e series of the $(\mathrm{A}/\mathfrak{m})$-modules graded $\operatorname{gr}_\mathfrak{m}^*(\mathrm{M})$ and $\operatorname{gr}_\mathfrak{R}^*(\mathrm{M})\otimes_{\mathrm{A}/\mathfrak{R}}(\mathrm{A}/\mathfrak{m})$, we have
\end{enumerate}
\begin{equation}
  \mathrm{Q}(\mathrm{T})=\mathrm{P}(\mathrm{T})(1-\mathrm{T})^n.
  \tag{19.7.3.1}
\label{IV.19.7.3.1}
\end{equation}

Let $\mathbf{f}=(f_1,\ldots,f_n)$ be a sequence of elements of $\mathrm{A}$ such that the images of the $f_i$ in $\mathrm{A}/\mathfrak{R}$ form a regular system of parameters of $\mathrm{A}/\mathfrak{R}$ $(\mathbf{0},17.1.6)$, so that $\mathbf{f}$ is an $(\mathrm{A}/\mathfrak{R})$-regular sequence; moreover, if $\mathfrak{J}$ is the ideal generated by $\mathbf{f}$, $\mathfrak{J}+\mathfrak{R}=\mathfrak{m}$, and as $\mathrm{A}/\mathfrak{m}$ is a field, $\operatorname{gr}_\mathfrak{m}^*(\mathrm{M})$ is a flat $(\mathrm{A}/\mathfrak{m})$-module. As $\mathrm{A}$ is Noetherian and $\mathrm{M}$ a finitely generated $\mathrm{A}$-module, one can apply the equivalence of $a)$ and $b)$ in \sref{IV.19.7.1}. Moreover, $\mathrm{A}/\mathfrak{m}$ being a field, the fact that $\psi_\mathrm{M}$ is bijective is equivalent to saying that for every $h\geqslant 0$, $\operatorname{gr}_\mathfrak{m}^h(\mathrm{M})$ and the submodule of elements of degree $h$ of $(\operatorname{gr}_\mathfrak{R}^*(\mathrm{M})\otimes_{\mathrm{A}/\mathfrak{R}}(\mathrm{A}/\mathfrak{m}))[\mathrm{T}_1,\ldots,\mathrm{T}_n]$ have the same rank over $\mathrm{A}/\mathfrak{m}$; but for the second of these modules the rank in question is clearly equal to
\[
  \textstyle\sum_{r=0}^{h}\binom{r+n-1}{n-1}\operatorname{rg}(\operatorname{gr}_\mathfrak{R}^{h-r}(\mathrm{M})\otimes_{\mathrm{A}/\mathfrak{R}}(\mathrm{A}/\mathfrak{m})).\quad\text{As}\quad(1-\mathrm{T})^{-n}=\textstyle\sum_{r=0}^{\infty}\binom{r+n-1}{n-1}\mathrm{T}^r,\quad\text{this}
\]
proves that the relation \sref{IV.19.7.3.1} is necessary and sufficient for $\psi_\mathrm{M}$ to be bijective.
\end{corollary}

\begin{corollary}[19.7.4]
\label{IV.19.7.4}
Let $\mathrm{X}$ be a locally Noetherian prescheme, $\mathrm{Y}$ a closed sub-prescheme of $\mathrm{X}$, $\mathrm{Y}$ regular and connected, $\sh{F}$ a coherent $\sh{O}_\mathrm{X}$-Module normally flat $(\sref{IV.6.10.1})$ along $\mathrm{Y}$. Then there exists a formal power series $\mathrm{R}\in\bb{Q}[[\mathrm{T}]]$ (independent of $x\in\mathrm{Y}$) such that, for every $x\in\mathrm{Y}$, the Poincar\'e series of $\operatorname{gr}_{\mathfrak{m}_x}^*(\sh{F}_x)$ is given by the formula
\begin{equation}
  \mathrm{P}_x(\sh{F})(\mathrm{T})=\mathrm{R}(\mathrm{T})(1-\mathrm{T})^{-n}\qquad\text{where }n=\dim(\sh{O}_{\mathrm{Y},x}).
  \label{IV.19.7.4.1}
\tag{19.7.4.1}
\end{equation}
\oldpage[IV-4]{219}

Let $\sh{J}$ be the coherent Ideal of $\sh{O}_\mathrm{X}$ defining $\mathrm{Y}$; the hypothesis on $\sh{F}$ means that $\operatorname{gr}_\sh{J}^*(\sh{F})$ is a flat $\sh{O}_\mathrm{Y}$-Module. We can therefore, for every $x\in\mathrm{Y}$, apply \sref{IV.19.7.3} taking $\mathrm{A}=\sh{O}_{\mathrm{X},x}$, $\mathfrak{R}=\sh{J}_x$, $\mathrm{M}=\sh{F}_x$; condition $a)$ being satisfied, the Poincar\'e series $\mathrm{P}_x(\sh{F})(\mathrm{T})$ is equal to $\mathrm{R}_x(\mathrm{T})(1-\mathrm{T})^{-n}$, where $\mathrm{R}_x(\mathrm{T})$ is the Poincar\'e series of $\operatorname{gr}_{\sh{J}_x}^*(\sh{F}_x)\otimes_{\sh{O}_{\mathrm{Y},x}}k(x)$. Each $\sh{O}_\mathrm{Y}$-Module $\sh{J}^h\sh{F}/\sh{J}^{h+1}\sh{F}$ is flat and coherent by hypothesis, hence locally free $(2.1.12)$; and since $\mathrm{Y}$ is connected, its rank is constant on $\mathrm{Y}$. This proves that $\mathrm{R}_x(\mathrm{T})$ is independent of $x\in\mathrm{Y}$.
\end{corollary}

\begin{corollary}[19.7.5]
\label{IV.19.7.5}
Let $\mathrm{X}$ be a locally Noetherian prescheme, $\mathrm{Y}$ a closed sub-prescheme of $\mathrm{X}$, $\mathrm{Z}$ a closed sub-prescheme of $\mathrm{Y}$; one supposes $\mathrm{Y}$ and $\mathrm{Z}$ regular. Let $\sh{F}$ be a coherent $\sh{O}_\mathrm{X}$-Module.
\begin{enumerate}
\item[\rm(i)] If $\sh{F}$ is normally flat along $\mathrm{Y}$ at the points of $\mathrm{Z}$ (cf. \sref{IV.11.3.4} for the definition of normal flatness at a point), then $\sh{F}$ is normally flat along $\mathrm{Z}$.
\item[\rm(ii)] Suppose moreover that $\mathrm{Z}$ is connected and that $\sh{O}_\mathrm{X}$ is normally flat along $\mathrm{Y}$ at the points of $\mathrm{Z}$ (which is the case in particular if $\mathrm{X}$ is regular at the points of $\mathrm{Z}$). If $\sh{F}$ is normally flat along $\mathrm{Z}$ and if there exists a point $z\in\mathrm{Z}$ at which $\sh{F}$ is normally flat along $\mathrm{Y}$, then $\sh{F}$ is normally flat along $\mathrm{Y}$ at all the points of $\mathrm{Z}$.
\end{enumerate}

Let $\sh{K}$ and $\sh{L}\supset\sh{K}$ be the coherent Ideals of $\sh{O}_\mathrm{X}$ defining respectively $\mathrm{Y}$ and $\mathrm{Z}$. For every point $z\in\mathrm{Z}$, the canonical immersion $\mathrm{Z}\to\mathrm{Y}$ is regular at $z$, since $\mathrm{Y}$ and $\mathrm{Z}$ are regular at that point \sref{IV.19.1.1}; in other words, there exists a sequence $(f_i)_{1\leqslant i\leqslant n}$ of elements of $\sh{O}_{\mathrm{X},z}$ which is $\sh{K}_z$-regular and such that $\sh{L}_z=\sh{K}_z+\sum_i f_i\sh{O}_{\mathrm{X},z}$.
\begin{enumerate}
\item[\rm(i)] The hypothesis means that $\operatorname{gr}_{\sh{K}_z}^*(\sh{F}_z)$ is a flat $\sh{O}_{\mathrm{Y},z}$-module; it therefore implies that $\operatorname{gr}_{\sh{L}_z}^*(\sh{F}_z)$ is a flat $\sh{O}_{\mathrm{Z},z}$-module, by the implication $a)\Rightarrow b)$ of \sref{IV.19.7.1}.
\item[\rm(ii)] The supplementary hypothesis on $\mathrm{X}$ implies that the sequence $(f_i)$ is $\operatorname{gr}_{\sh{K}_z}^*(\sh{O}_{\mathrm{X},z})$-regular at every point $z\in\mathrm{Z}$ \sref{IV.19.7.2}. On the other hand, $\mathrm{Z}$, being regular and connected, is integral; if $\operatorname{gr}_{\sh{K}_z}^*(\sh{F}_z)$ is flat over $\sh{O}_{\mathrm{Y},z}$ at one point $z\in\mathrm{Z}$, then, since the generic point $\zeta$ of $\mathrm{Z}$ is a generization of $z$, $\operatorname{gr}_{\sh{K}_\zeta}^*(\sh{F}_\zeta)$ is flat over $\sh{O}_{\mathrm{Y},\zeta}$ $(\mathbf{0}_\mathrm{I},6.3.1)$. If, at every point $z'$ of $\mathrm{Z}$, $\operatorname{gr}_{\sh{L}_{z'}}^*(\sh{F}_{z'})$ is flat over $\sh{O}_{\mathrm{Z},z'}$, we then conclude that $\operatorname{gr}_{\sh{K}_{z'}}^*(\sh{F}_{z'})$ is flat over $\sh{O}_{\mathrm{Y},z'}$, by applying the equivalence of $a)$ and $d)$ in \sref{IV.19.7.1}.
\end{enumerate}
\end{corollary}

\subsection{Properties of passage to the projective limit}
\label{subsection:IV.19.8}
\label{IV.19.8}

In this number, the notations and conventions on projective limits are those of $(8.5.1)$ and $(8.8.1)$.

\begin{proposition}[19.8.1]
\label{IV.19.8.1}
Suppose that the transition morphisms $\mathrm{S}_\mu\to\mathrm{S}_\lambda$ $(\lambda\leqslant\mu)$ are flat, and moreover that one of the two following hypotheses is satisfied:
\begin{enumerate}
\item[$1^\circ$] The preschemes $\mathrm{S}_\lambda$ are locally Noetherian.
\item[$2^\circ$] The transition morphisms $\mathrm{S}_\mu\to\mathrm{S}_\lambda$ are surjective (and hence faithfully flat).
\end{enumerate}
\oldpage[IV-4]{220}

\emph{Under these conditions:}
\begin{enumerate}
\item[\rm(i)] Suppose $\mathrm{S}_\alpha$ quasi-compact. Let $\sh{F}_\alpha$ be a quasi-coherent $\sh{O}_{\mathrm{S}_\alpha}$-Module, which one supposes moreover of finite type when hypothesis $1^\circ$ is satisfied. Let $(f_{i\alpha})_{1\leqslant i\leqslant n}$ be a finite sequence of sections of $\sh{F}_\alpha$ above $\mathrm{S}_\alpha$, the sequence of sections $f_i$ of the $\sh{O}_\mathrm{S}$-Module $\sh{F}$ above $\mathrm{S}$, corresponding to the $f_{i\alpha}$. For the sequence of sections $f_i$ $(1\leqslant i\leqslant n)$ to be $\sh{F}$-regular, it is necessary and sufficient that there exist $\lambda\geqslant\alpha$ such that the sequence of sections $f_{i\lambda}$ of $\sh{F}_\lambda$ above $\mathrm{S}_\lambda$ is $\sh{F}_\lambda$-regular.
\item[\rm(ii)] Suppose $\mathrm{X}_\alpha$ quasi-compact, and let $j_\alpha:\mathrm{X}_\alpha\to\mathrm{S}_\alpha$ be an immersion, which one supposes locally of finite presentation when hypothesis $2^\circ$ is satisfied. Then, for the corresponding immersion $j:\mathrm{X}\to\mathrm{S}$ to be regular, it is necessary and sufficient that there exist $\lambda\geqslant\alpha$ such that $j_\lambda:\mathrm{X}_\lambda\to\mathrm{S}_\lambda$ is regular.
\end{enumerate}
\end{proposition}

\begin{proof}
\begin{enumerate}
\item[\rm(i)] If $p_\lambda:\mathrm{S}\to\mathrm{S}_\lambda$ is the canonical projection, we have
\[
  \sh{F}/\Bigl(\sum_{j=1}^{i-1}f_j\sh{F}\Bigr)
  =p_\lambda^*\Bigl(\sh{F}_\lambda/\Bigl(\sum_{j=1}^{i-1}f_{j\lambda}\sh{F}_\lambda\Bigr)\Bigr),
\]
hence we are reduced to the case $n=1$, in which case we omit the index $i$. Sufficiency follows from the flatness of $p_\lambda$ $(8.3.8)$ and $(\mathbf{0},15.1.5)$. In case $2^\circ$, $p_\lambda$ is faithfully flat $(8.3.8)$ and necessity follows from $(\mathbf{0},15.2.5)$, with $\lambda=\alpha$. In case $1^\circ$, let $\sh{N}_\lambda$ (resp. $\sh{N}$) be the kernel of multiplication by $f_\lambda$ on $\sh{F}_\lambda$ (resp. by $f$ on $\sh{F}$); since $\sh{F}_\lambda$ is coherent, so is $\sh{N}_\lambda$, and the hypothesis $\sh{N}=0$ implies $\sh{N}_\lambda=0$ for some $\lambda\geqslant\alpha$ by $(8.5.8)$, (ii).
\item[\rm(ii)] As $p_\lambda$ is flat, sufficiency follows from \sref{IV.19.1.5}, (ii). For necessity, since $j_\alpha(\mathrm{X}_\alpha)$ is quasi-compact, it is contained in a quasi-compact open subset of $\mathrm{S}_\alpha$; we may therefore reduce to the case in which $\mathrm{S}_\alpha$ is quasi-compact and $j_\alpha$ is a closed immersion. Its image is then defined by a quasi-coherent Ideal $\sh{J}_\alpha$ of $\sh{O}_{\mathrm{S}_\alpha}$, which may moreover be supposed of finite type in cases $1^\circ$ and $2^\circ$. The image of $\mathrm{X}_\lambda$ (resp. $\mathrm{X}$) in $\mathrm{S}_\lambda$ (resp. $\mathrm{S}$) is defined by $\sh{J}_\lambda=\sh{J}_\alpha\sh{O}_{\mathrm{S}_\lambda}$ (resp. $\sh{J}=\sh{J}_\alpha\sh{O}_\mathrm{S}$), again of finite type. By $(8.2.11)$, we may suppose that $\sh{J}$ is generated by a regular sequence $(f_i)_{1\leqslant i\leqslant n}$ of sections over $\mathrm{S}$, and hence by a surjection $u:\sh{O}_\mathrm{S}^{,n}\to\sh{J}$. By $(8.5.2)$, (i), and $(8.5.7)$, there exist $\lambda\geqslant\alpha$ and a surjection $u_\lambda:\sh{O}_{\mathrm{S}_\lambda}^{,n}\to\sh{J}_\lambda$ such that $u=p_\lambda^*(u_\lambda)$; thus the $f_i$ are the canonical images of sections $f_{i\lambda}$ of $\sh{J}_\lambda$ generating this Ideal. By (i), there is then $\mu\geqslant\lambda$ for which $(f_{i\mu})$ is $\sh{O}_{\mathrm{S}_\mu}$-regular, and hence $j_\mu$ is a regular immersion.
\end{enumerate}
\end{proof}

\begin{proposition}[19.8.2]
\label{IV.19.8.2}
Let $\mathrm{X}_\alpha$ be an $\mathrm{S}_\alpha$-prescheme quasi-compact and locally of finite presentation.
\begin{enumerate}
\item[\rm(i)] Let $\sh{F}_\alpha$ be a quasi-coherent $\sh{O}_{\mathrm{X}_\alpha}$-Module of finite presentation, and $(f_{i\alpha})_{1\leqslant i\leqslant n}$ a finite sequence of sections of $\sh{O}_{\mathrm{X}_\alpha}$ over $\mathrm{X}_\alpha$, corresponding to sections $f_i$ of $\sh{O}_\mathrm{X}$ over $\mathrm{X}$. For the sequence $(f_i)$ to be $\sh{F}$-transversally regular relative to $\mathrm{S}$ \emph{(\sref{IV.19.2.1})}, it is necessary and sufficient that there exist $\lambda\geqslant\alpha$ such that the sequence $(f_{i\lambda})$ over $\mathrm{X}_\lambda$ is $\sh{F}_\lambda$-transversally regular relative to $\mathrm{S}_\lambda$.
\item[\rm(ii)] Let $\mathrm{Y}_\alpha$ be an $\mathrm{S}_\alpha$-prescheme quasi-compact, $j_\alpha:\mathrm{Y}_\alpha\to\mathrm{X}_\alpha$ an $\mathrm{S}_\alpha$-immersion locally of finite presentation. For the corresponding $\mathrm{S}$-immersion $j:\mathrm{Y}\to\mathrm{X}$ to be transversally regular relative to $\mathrm{S}$, it is necessary and sufficient that there exist $\lambda\geqslant\alpha$ such that $j_\lambda:\mathrm{Y}_\lambda\to\mathrm{X}_\lambda$ is transversally regular relative to $\mathrm{S}_\lambda$.
\end{enumerate}
\end{proposition}

\begin{proof}
\begin{enumerate}
\item[\rm(i)]
\oldpage[IV-4]{221}
The sufficiency of the condition follows from $(\mathbf{0},15.1.15)$. To prove necessity, it is enough to show that every point $x\in\mathrm{X}$ has an open neighbourhood $\mathrm{V}(x)$ for which there exist an index $\lambda=\lambda(x)\geqslant\alpha$ and an open neighbourhood $\mathrm{V}_\lambda$ of the image $x_\lambda$ of $x$ in $\mathrm{X}_\lambda$, such that $\mathrm{V}(x)$ is the inverse image of $\mathrm{V}_\lambda$ and the restrictions of the $f_{i\lambda}$ to $\mathrm{V}_\lambda$ form an $(\sh{F}_\lambda|\mathrm{V}_\lambda)$-transversally regular sequence relative to $\mathrm{S}_\lambda$. Indeed, one then applies $(8.3.4)$ to the open subset $\mathrm{U}_\lambda$ of $\mathrm{X}_\lambda$ defined as the largest open on which the restrictions of the $f_{i\lambda}$ form such a sequence.

Let $s$ be the image of $x$ in $\mathrm{S}$ and, for every $\lambda\geqslant\alpha$, let $s_\lambda$ be the image of $s$ in $\mathrm{S}_\lambda$. The fibre $\mathrm{X}_s$ is the projective limit of the fibres $(\mathrm{X}_\lambda)_{s_\lambda}$ and is locally Noetherian, since $\mathrm{X}\to\mathrm{S}$ is locally of finite presentation. Moreover $(\mathrm{X}_\alpha)_{s_\alpha}$ is quasi-compact and the transition morphisms $(\mathrm{X}_\mu)_{s_\mu}\to(\mathrm{X}_\lambda)_{s_\lambda}$ are flat, since this is so for $\operatorname{Spec}(k(s_\mu))\to\operatorname{Spec}(k(s_\lambda))$. We may therefore apply \sref{IV.19.8.1}, (i), to the sections $f_i\otimes1$ of $\sh{F}_s=\sh{F}\otimes_{\sh{O}_\mathrm{S}}k(s)$ over $\mathrm{X}_s$. Thus there exists $\lambda\geqslant\alpha$ such that the sections $f_{i\lambda}\otimes1$ of $(\sh{F}_\lambda)_{s_\lambda}=\sh{F}_\lambda\otimes_{\sh{O}_{\mathrm{S}_\lambda}}k(s_\lambda)$ form an $(\sh{F}_\lambda)_{s_\lambda}$-regular sequence. By \sref{IV.11.2.6}, one may moreover suppose that $\sh{F}_\lambda$ is flat over $\mathrm{S}_\lambda$ at $x_\lambda$. It then follows from \sref{IV.11.3.8} that there exists an open neighbourhood $\mathrm{V}_\lambda$ of $x_\lambda$ in $\mathrm{X}_\lambda$ satisfying the required conditions.

\item[\rm(ii)] Sufficiency follows from \sref{IV.19.2.7}, (ii). For necessity, it again suffices to show that every point $x\in j(\mathrm{Y})$ has an open neighbourhood $\mathrm{V}(x)$ for which there exist an index $\lambda=\lambda(x)\geqslant\alpha$ and an open neighbourhood $\mathrm{V}_\lambda$ of the image $x_\lambda$ of $x$ in $\mathrm{X}_\lambda$, such that $\mathrm{V}(x)$ is the inverse image of $\mathrm{V}_\lambda$ and the restriction $j_\lambda^{-1}(\mathrm{V}_\lambda)\to\mathrm{V}_\lambda$ is transversally regular relative to $\mathrm{S}_\lambda$.

Let $s$ be the image of $x$ in $\mathrm{S}$ and, for every $\lambda\geqslant\alpha$, let $s_\lambda$ be the image of $s$ in $\mathrm{S}_\lambda$. Since $\mathrm{X}_s$ (resp. $\mathrm{Y}_s$) is the projective limit of the fibres $(\mathrm{X}_\lambda)_{s_\lambda}$ (resp. $(\mathrm{Y}_\lambda)_{s_\lambda}$), we may use \sref{IV.19.8.1}, (ii), as in (i); and since the immersion $\mathrm{Y}_s\to\mathrm{X}_s$ is regular by hypothesis, there exists an index $\lambda(s)$ such that, for $\lambda\geqslant\lambda(s)$, the immersion $(\mathrm{Y}_\lambda)_{s_\lambda}\to(\mathrm{X}_\lambda)_{s_\lambda}$ is regular. Moreover, by \sref{IV.19.2.4}, there exists a quasi-compact open neighbourhood $\mathrm{W}(x)$ of $x$ in $\mathrm{X}$ such that the structural morphisms $\mathrm{W}(x)\to\mathrm{S}$ and $\mathrm{Y}\cap\mathrm{W}(x)\to\mathrm{S}$ are flat. Applying \sref{IV.11.2.6} and $(8.2.11)$, one may suppose that there exist $\lambda\geqslant\lambda(s)$ and a quasi-compact open neighbourhood $\mathrm{W}(x_\lambda)$ of $x_\lambda$ in $\mathrm{X}_\lambda$, such that $\mathrm{W}(x)$ is the inverse image of $\mathrm{W}(x_\lambda)$ and the restrictions to $\mathrm{W}(x_\lambda)$ and $\mathrm{Y}_\lambda\cap\mathrm{W}(x_\lambda)$ of the structural morphisms $\mathrm{X}_\lambda\to\mathrm{S}_\lambda$ and $\mathrm{Y}_\lambda\to\mathrm{S}_\lambda$ are flat. We conclude from \sref{IV.19.2.4} that there exists an open neighbourhood $\mathrm{V}_\lambda$ of $x_\lambda$ satisfying the required conditions.
\end{enumerate}
\end{proof}

\begin{remark}[19.8.3]
\label{IV.19.8.3}
The preceding demonstrations show that the statements of \sref{IV.19.8.1} and \sref{IV.19.8.2} subsist when one replaces the conditions concerning regularity or transversal regularity throughout $\mathrm{X}$ (resp. $\mathrm{X}_\lambda$) by these conditions in a neighbourhood of a point $x\in\mathrm{X}$ (resp. in a neighbourhood of $x_\lambda$, projection of $x$).
\end{remark}

\subsection{\texorpdfstring{$\sh{F}$}{F}-regular sequences and depth}
\label{subsection:IV.19.9}
\label{IV.19.9}

\oldpage[IV-4]{222}

\begin{env}[19.9.1]
\label{IV.19.9.1}
Recall that if $\mathrm{X}$ is a locally Noetherian prescheme, $\sh{F}$ a coherent $\sh{O}_\mathrm{X}$-Module, $\mathrm{T}$ a part of $\mathrm{X}$, one sets $(5.10.1)$
\begin{equation}
  \operatorname{prof}_\mathrm{T}(\sh{F})=\inf_{x\in\mathrm{T}}\operatorname{prof}(\sh{F}_x).
  \label{IV.19.9.1.1}
\tag{19.9.1.1}
\end{equation}
We set on the other hand, for every point $t\in\mathrm{T}$,
\begin{equation}
  \operatorname{prof}_{\mathrm{T},t}(\sh{F})=\inf_{x\in\mathrm{T}\cap\operatorname{Spec}(\sh{O}_{\mathrm{X},t})}\operatorname{prof}(\sh{F}_x),
  \label{IV.19.9.1.2}
\tag{19.9.1.2}
\end{equation}
and we say that $\operatorname{prof}_{\mathrm{T},t}(\sh{F})$ is the \emph{T-depth of $\sh{F}$ at the point $t$}; it is clear that we have
\begin{equation}
  \operatorname{prof}_\mathrm{T}(\sh{F})=\inf_{t\in\mathrm{T}}\operatorname{prof}_{\mathrm{T},t}(\sh{F}).
  \label{IV.19.9.1.3}
\tag{19.9.1.3}
\end{equation}
\end{env}

\begin{lemma}[19.9.2]
\label{IV.19.9.2}
Let $\mathrm{A}$ be a Noetherian ring, $\mathrm{M}$ a finitely generated $\mathrm{A}$-module, $\mathfrak{J}$ an ideal of $\mathrm{A}$. In order that, for every prime ideal $\mathfrak{p}\supset\mathfrak{J}$, we have $\operatorname{prof}_{\mathrm{A}_\mathfrak{p}}(\mathrm{M}_\mathfrak{p})\geqslant r$, it is necessary and sufficient that there exist an M-regular sequence of $r$ elements belonging to $\mathfrak{J}$.
\end{lemma}

\begin{proof}
The sufficiency of the condition follows immediately from $(\mathbf{0},16.4.5)$; let us prove its necessity. We argue by induction on $r$ (there is nothing to prove if $r=0$): since $\operatorname{prof}_{\mathrm{A}_\mathfrak{p}}(\mathrm{M}_\mathfrak{p})\geqslant r-1$ for every $\mathfrak{p}\supset\mathfrak{J}$, there exists by the induction hypothesis an M-regular sequence $(g_i)_{1\leqslant i\leqslant r-1}$ of elements of $\mathfrak{J}$. Set $\mathrm{N}=\mathrm{M}/(g_1\mathrm{M}+\ldots+g_{r-1}\mathrm{M})$; for every $\mathfrak{p}\supset\mathfrak{J}$, the images of the $g_i$ in $\mathrm{A}_\mathfrak{p}$ belong to the maximal ideal $\mathfrak{p}\mathrm{A}_\mathfrak{p}$ and form an $\mathrm{M}_\mathfrak{p}$-regular sequence $(\mathbf{0},15.1.14)$. It follows from $(\mathbf{0},16.4.6)$ that $\operatorname{prof}_{\mathrm{A}_\mathfrak{p}}(\mathrm{N}_\mathfrak{p})=\operatorname{prof}_{\mathrm{A}_\mathfrak{p}}(\mathrm{M}_\mathfrak{p})-(r-1)\geqslant 1$, so everything reduces to proving the lemma for $r=1$. The hypothesis then means that, for every $\mathfrak{p}\supset\mathfrak{J}$, $\mathfrak{p}\mathrm{A}_\mathfrak{p}$ is not associated with $\mathrm{M}_\mathfrak{p}$ $(\mathbf{0},16.4.6)$; hence (Bourbaki, \emph{Alg.~comm.}, chap.~IV, \textsection~1, n\textsuperscript{o}~2, prop.~5), $\mathfrak{p}$ is not associated with $\mathrm{M}$. In other words, $\mathfrak{J}$ is contained in none of the prime ideals of $\operatorname{Ass}(\mathrm{M})$, and therefore is not contained in their union (Bourbaki, \emph{Alg.~comm.}, chap.~II, \textsection~1, n\textsuperscript{o}~1, prop.~2). Since this union is the set of elements that are not M-regular (Bourbaki, \emph{Alg.~comm.}, chap.~IV, \textsection~1, cor.~2 of prop.~2), this proves the lemma.
\end{proof}

\begin{proposition}[19.9.3]
\label{IV.19.9.3}
Let $\mathrm{X}$ be a locally Noetherian prescheme, $\sh{F}$ a coherent $\sh{O}_\mathrm{X}$-Module, $\mathrm{Y}$ a closed subset of $\mathrm{X}$, $y$ a point of $\mathrm{Y}$, $n$ an integer $>0$. The following conditions are equivalent:
\begin{enumerate}
\item[\rm a)] There exists an open neighbourhood $\mathrm{U}$ of $y$ in $\mathrm{X}$ such that $\operatorname{prof}_{\mathrm{Y}\cap\mathrm{U}}(\sh{F}|\mathrm{U})\geqslant n$.
\item[\rm b)] There exists an open neighbourhood $\mathrm{U}$ of $y$ in $\mathrm{X}$ and an $(\sh{F}|\mathrm{U})$-regular sequence of $n$ sections $f_i$ of $\sh{O}_\mathrm{U}$ above $\mathrm{U}$, such that $f_i(x)=0$ $(\mathbf{0}_\mathrm{I},5.5.1)$ for $1\leqslant i\leqslant n$ at every point $x\in\mathrm{U}\cap\mathrm{Y}$.
\item[\rm c)] $\operatorname{prof}_{\mathrm{Y},y}(\sh{F})\geqslant n$.
\end{enumerate}

Moreover, the set of $y\in\mathrm{Y}$ satisfying these conditions is open in $\mathrm{Y}$.
\end{proposition}

\begin{proof}
The last assertion follows trivially from $a)$. The equivalence of $a)$ and $b)$ follows from \sref{IV.19.9.2}, applied to an affine open neighbourhood of $y$ in $\mathrm{X}$, and from $(\mathbf{0},15.1.14)$. It is clear that $a)$ implies $c)$, since every open neighbourhood of $y$ in $\mathrm{X}$ contains $\operatorname{Spec}(\sh{O}_{\mathrm{X},y})$; let us prove conversely that $c)$ implies $b)$. It follows from $c)$ and the definition \sref{IV.19.9.1} that
\oldpage[IV-4]{223}
we can apply \sref{IV.19.9.2} to the $\sh{O}_{\mathrm{X},y}$-module $\sh{F}_y$ and the ideal $\sh{J}_y$ of $\sh{O}_{\mathrm{X},y}$, where $\sh{J}$ is an ideal defining $\mathrm{Y}$ near $y$. Hence there exists an $\sh{F}_y$-regular sequence $(t_i)_{1\leqslant i\leqslant n}$ of $n$ elements of $\sh{J}_y$. There is therefore an open neighbourhood $\mathrm{U}$ of $y$ in $\mathrm{X}$ and a sequence $(f_i)_{1\leqslant i\leqslant n}$ of sections of $\sh{J}$ above $\mathrm{U}$ such that the $t_i$ are the germs of the $f_i$ at $y$; using $(\mathbf{0},15.2.4)$, one deduces that there exists an open neighbourhood $\mathrm{U}'\subset\mathrm{U}$ of $y$ in $\mathrm{X}$ such that the $f_i|\mathrm{U}'$ form an $(\sh{F}|\mathrm{U}')$-regular sequence.
\end{proof}

\begin{corollary}[19.9.4]
\label{IV.19.9.4}
With the notations of \emph{(\sref{IV.19.9.3})}, the function $y\mapsto\operatorname{prof}_{\mathrm{Y},y}(\sh{F})$ is lower semi-continuous on $\mathrm{Y}$.
\end{corollary}

\begin{proposition}[19.9.5]
\label{IV.19.9.5}
With the notations of \emph{(\sref{IV.19.9.3})}, let $\mathrm{X}'$ be a locally Noetherian prescheme, $g:\mathrm{X}'\to\mathrm{X}$ a flat morphism, $\mathrm{Y}'=g^{-1}(\mathrm{Y})$, $\sh{F}'=g^*(\sh{F})$. Then, for every point $y'\in\mathrm{Y}'$ above $y$, we have $\operatorname{prof}_{\mathrm{Y}',y'}(\sh{F}')=\operatorname{prof}_{\mathrm{Y},y}(\sh{F})$.
\end{proposition}

\begin{proof}
Indeed, for every $z'\in\operatorname{Spec}(\sh{O}_{\mathrm{X}',y'})$, it follows from $(6.3.1)$ that $\operatorname{prof}(\sh{F}'_{z'})\geqslant\operatorname{prof}(\sh{F}_{g(z')})$, and if $z''$ is a maximal point of the fibre $g^{-1}(g(z'))$ that is a generization of $z'$ (and hence also belongs to $\operatorname{Spec}(\sh{O}_{\mathrm{X}',y'})$), we have (\emph{loc.~cit.}) $\operatorname{prof}(\sh{F}'_{z''})=\operatorname{prof}(\sh{F}_{g(z')})$. The proposition then follows from the fact that $g(\operatorname{Spec}(\sh{O}_{\mathrm{X}',y'}))=\operatorname{Spec}(\sh{O}_{\mathrm{X},g(y')})$, since $g$ is flat $(2.3.4)$.
\end{proof}

\begin{proposition}[19.9.6]
\label{IV.19.9.6}
Let $f:\mathrm{X}\to\mathrm{S}$ be a morphism locally of finite presentation, $\mathrm{Y}$ a closed subset of $\mathrm{X}$, $\sh{F}$ a coherent $\sh{O}_\mathrm{X}$-Module of finite presentation and $f$-flat, $n$ an integer $>0$. For every point $y\in\mathrm{Y}$, denoting by $\mathrm{X}_{f(y)}$ the fibre of $f$ at the point $f(y)$, by $\sh{F}_{f(y)}$ the inverse image of $\sh{F}$ in $\mathrm{X}_{f(y)}$, and setting $\mathrm{Y}_{f(y)}=\mathrm{X}_{f(y)}\cap\mathrm{Y}$:
\begin{enumerate}
\item[\rm a)] $\operatorname{prof}_{\mathrm{Y}_{f(y)},y}(\sh{F}_{f(y)})\geqslant n$.
\item[\rm b)] There exists an open neighbourhood $\mathrm{U}$ of $y$ in $\mathrm{X}$ and a sequence $(g_i)_{1\leqslant i\leqslant n}$ of $n$ sections of $\sh{O}_\mathrm{X}$ above $\mathrm{U}$, that is transversally $\sh{F}$-regular relative to $\mathrm{S}$ (cf. \sref{IV.19.2.1} for the terminology) and such that $g_i(x)=0$ for $1\leqslant i\leqslant n$ at every point $x\in\mathrm{U}\cap\mathrm{Y}$.
\end{enumerate}

The set $\mathrm{V}_n$ of $y\in\mathrm{Y}$ satisfying these conditions is open in $\mathrm{Y}$; if moreover $\mathrm{Y}$ is locally constructible in $\mathrm{X}$, $\mathrm{V}_n$ is retrocompact in $\mathrm{X}$ $(\mathbf{0}_\mathrm{III},9.1.1)$.
\end{proposition}

\begin{proof}
By virtue of \sref{IV.19.9.3}, condition $a)$ is equivalent to the existence of an open neighbourhood $\mathrm{U}$ of $y$ in $\mathrm{X}$ and of an $(\sh{F}_{f(y)}|(\mathrm{U}\cap\mathrm{X}_{f(y)}))$-regular sequence $(h_i)_{1\leqslant i\leqslant n}$ of sections of $\sh{O}_{\mathrm{X}_{f(y)}}$ above $\mathrm{U}\cap\mathrm{X}_{f(y)}$, such that $h_i(z)=0$ for $1\leqslant i\leqslant n$ at every point $z\in\mathrm{U}\cap\mathrm{Y}_{f(y)}$. If we regard $\mathrm{Y}$ as a closed sub-prescheme of $\mathrm{X}$ having the given closed subset as its underlying space, and denote by $\sh{J}$ the quasi-coherent ideal of $\sh{O}_\mathrm{X}$ defining it, one can further require that $h_i\in\sh{J}_{f(y)}$ for $1\leqslant i\leqslant n$. Replacing $\mathrm{U}$ by a smaller affine open neighbourhood of $y$, one can suppose that the $h_i$ are the images of a sequence $(g_i)_{1\leqslant i\leqslant n}$ of sections of $\sh{J}$ above $\mathrm{U}$, so that the germs $(g_i)_y$ belong to the maximal ideal $\mathfrak{m}_y$ of $\sh{O}_{\mathrm{X},y}$. The equivalence of $a)$ and $b)$ then follows from \sref{IV.11.3.8}, which also proves that the set $\mathrm{V}_n$ is open in $\mathrm{Y}$.

It remains to prove the last assertion. Since this again concerns properties local on $\mathrm{X}$, we may suppose that $\mathrm{S}=\operatorname{Spec}(\mathrm{A})$ is affine and that $\mathrm{X}$ is of finite presentation over $\mathrm{S}$. There then exist a Noetherian subring $\mathrm{A}_0$ of $\mathrm{A}$, a prescheme $\mathrm{X}_0$ of finite type over $\mathrm{S}_0=\operatorname{Spec}(\mathrm{A}_0)$, and a coherent $\sh{O}_{\mathrm{X}_0}$-module $\sh{F}_0$ such that $\mathrm{X}=\mathrm{X}_0\times_{\mathrm{S}_0}\mathrm{S}$ and $\sh{F}=\sh{F}_0\otimes_{\sh{O}_{\mathrm{X}_0}}\sh{O}_\mathrm{X}$ $(8.9.1)$. We may moreover suppose that, if $f_0:\mathrm{X}_0\to\mathrm{S}_0$ is the structural morphism, $\sh{F}_0$ is $f_0$-flat $(\sref{IV.11.2.6})$. Since $\mathrm{Y}$ is constructible, we may further suppose $(8.3.11)$ that $\mathrm{Y}=p^{-1}(\mathrm{Y}_0)$, where $p:\mathrm{X}\to\mathrm{X}_0$ is the canonical projection. Then, for every $y\in\mathrm{Y}$, \sref{IV.19.9.5} and the transitivity of fibres give
\[
  \operatorname{prof}_{\mathrm{Y}_{f(y)},y}(\sh{F}_{f(y)})
  =\operatorname{prof}_{(\mathrm{Y}_0)_{f_0(y_0)},y_0}((\sh{F}_0)_{f_0(y_0)}),
\]
where $y_0=p(y)$. If $\mathrm{V}_n^0$ is the set of $y_0\in\mathrm{Y}_0$ for which the right-hand side is at least $n$, then $\mathrm{V}_n=p^{-1}(\mathrm{V}_n^0)$. Since $\mathrm{X}_0$ is Noetherian, $\mathrm{V}_n^0$ is a retrocompact open subset of $\mathrm{X}_0$, and consequently (cf. the proof of $(\mathrm{I},8.2)$) $\mathrm{V}_n$ is retrocompact in $\mathrm{X}$.
\end{proof}

\oldpage[IV-4]{224}

\begin{corollary}[19.9.7]
\label{IV.19.9.7}
Under the general hypotheses of \sref{IV.19.9.6}, the function $y\mapsto\operatorname{prof}_{Y_{f(y)},y}(\sh{F}_{f(y)})$ is lower semi-continuous in $Y$.
If moreover $Y$ is locally constructible, this function is locally constructible.
\end{corollary}

\begin{proposition}[19.9.8]
\label{IV.19.9.8}
Let $f:X\to S$ be a morphism locally of finite presentation, $Z$ a closed subset of $X$, $\sh{F}$ an $\sh{O}_X$-module of finite presentation and $f$-flat, and $\sh{G}$ a quasi-coherent $\sh{O}_S$-module.
Suppose that, for every $x\in Z$, we have $\operatorname{prof}((\sh{F}_{f(x)})_x)\geqslant 1$ (resp.\ $\operatorname{prof}((\sh{F}_{f(x)})_x)\geqslant 2$).
Then, if $j:X-Z\to X$ is the canonical injection and if we set $\sh{H}=\sh{F}\otimes_{\sh{O}_X}f^*(\sh{G})$, the canonical homomorphism $\rho_{\sh{H}}:\sh{H}\to j_*(j^*(\sh{H}))$ relative to $j$ is injective (resp.\ bijective).
\end{proposition}

\begin{proof}
The question is evidently local on $X$ and $S$, and it suffices to prove the proposition in a neighbourhood of a point $x\in Z$. In other words, we may restrict to the case where $X$ and $S$ are affine. By \sref{IV.19.9.6}, we may suppose that there exists a section $g_1$ of $\sh{O}_X$ over $X$ that is transversally $\sh{F}$-regular relative to $S$ (resp.\ a transversally $\sh{F}$-regular sequence $(g_1,g_2)$ of two sections of $\sh{O}_X$ over $X$ relative to $S$) such that, if $Z'$ is the set of $x'\in X$ for which $g_1(x')=0$ (resp.\ $g_1(x')=g_2(x')=0$), then $Z\subset Z'$.

This being so, return to the proof of \sref{IV.19.9.8}. By \sref{IV.19.9.6}, we still have $\operatorname{prof}((\sh{F}_{f(x)})_x)\geqslant 1$ (resp.\ $\operatorname{prof}((\sh{F}_{f(x)})_x)\geqslant 2$) for every $x\in Z'$. If the proposition is proved after replacing the pair $(X,Z)$ by each of $(X,Z')$ and $(X-Z,Z'\cap(X-Z))$, it is immediate that it also holds for $(X,Z)$; we may therefore restrict to proving it for $X$ and $Z'$. In other words, we may restrict to proving \sref{IV.19.9.8} when $S=\Spec(A)$ and $X=\Spec(B)$, where $B$ is an $A$-algebra of finite presentation, $Z=V(\mathfrak{J})$ for an ideal $\mathfrak{J}$ of finite type of $B$, $\sh{F}=\widetilde M$, and $\sh{G}=\widetilde N$, where $N$ is an $A$-module and $M$ is a finitely presented $B$-module that is flat as an $A$-module. We may moreover reduce to the case where $N$ is a finitely presented $A$-module. Indeed, $N$ is the filtered inductive limit of a system of finitely presented $A$-modules $N_\alpha$ (by taking a double inductive limit). If $\sh{G}_\alpha=\widetilde{N_\alpha}$, then $\sh{H}$ is the inductive limit of the $\sh{H}_\alpha=\sh{F}\otimes_{\sh{O}_X}f^*(\sh{G}_\alpha)$; since $j_*$ and $j^*$ commute with inductive limits and $\varinjlim$ is an exact functor in the category of quasi-coherent modules, the proposition for every $\sh{H}_\alpha$ implies it for $\sh{H}$.

It further suffices to prove that the canonical homomorphism
\[
  \Gamma(X,\sh{H})\longrightarrow\Gamma(X-Z,\sh{H})
\]
is injective (resp.\ bijective). There then exist a Noetherian subring $A_0$ of $A$, a finite-type $A_0$-algebra $B_0$, an ideal $\mathfrak{J}_0$ of $B_0$, a finitely generated $B_0$-module $M_0$ that is flat as an $A_0$-module, and an $A_0$-module $N_0$, such that
\oldpage[IV-4]{225}
$B=B_0\otimes_{A_0}A$, $\mathfrak{J}=\mathfrak{J}_0B$, $M=M_0\otimes_{A_0}A$, and $N=N_0\otimes_{A_0}A$ $(8.9.1)$, $(8.5.11)$, and $(\sref{IV.11.2.7})$. Let $(A_\lambda)$ be the family of finite-type $A_0$-subalgebras of $A$, so that $A=\varinjlim A_\lambda$; set
\[
\begin{gathered}
  B_\lambda=B_0\otimes_{A_0}A_\lambda,\qquad
  \mathfrak{J}_\lambda=\mathfrak{J}_0B_\lambda,\qquad
  M_\lambda=M_0\otimes_{A_0}A_\lambda=M_0\otimes_{B_0}B_\lambda,\\
  N_\lambda=N_0\otimes_{A_0}A_\lambda,\qquad
  S_\lambda=\Spec(A_\lambda),\quad X_\lambda=\Spec(B_\lambda),\quad
  Z_\lambda=V(\mathfrak{J}_\lambda).
\end{gathered}
\]
If $p_\lambda:X\to X_\lambda$ is the canonical projection, then $Z=p_\lambda^{-1}(Z_\lambda)$ and $X-Z=p_\lambda^{-1}(X_\lambda-Z_\lambda)$. Since $X_\lambda-Z_\lambda$ is quasi-compact and quasi-separated, it follows from $(8.5.2)$ that, on setting $\sh{F}_\lambda=\widetilde{M_\lambda}$ and $\sh{H}_\lambda=(M_\lambda\otimes_{A_\lambda}N_\lambda)^\sim$, the homomorphisms
\[
  \varinjlim\Gamma(X_\lambda,\sh{H}_\lambda)\longrightarrow\Gamma(X,\sh{H}),
  \qquad
  \varinjlim\Gamma(X_\lambda-Z_\lambda,\sh{H}_\lambda)\longrightarrow\Gamma(X-Z,\sh{H})
\]
are bijective. By exactness of $\varinjlim$, it therefore suffices to prove that, for every sufficiently large $\lambda$, the canonical homomorphism $\Gamma(X_\lambda,\sh{H}_\lambda)\to\Gamma(X_\lambda-Z_\lambda,\sh{H}_\lambda)$ is injective (resp.\ bijective).

For every $x\in Z$, if $x_\lambda$ is the projection of $x$ in $Z_\lambda$, then $(4.2.7)$ and $(6.7.1)$ give
\[
  \operatorname{prof}((\sh{F}_{f(x)})_x)
  =\operatorname{prof}(((\sh{F}_\lambda)_{f_\lambda(x_\lambda)})_{x_\lambda}).
\]
Thus, to reduce to proving \sref{IV.19.9.8} when $A$ is Noetherian, it remains to prove that, for every sufficiently large $\lambda$, $\operatorname{prof}(((\sh{F}_\lambda)_{f_\lambda(x_\lambda)})_{x_\lambda})\geqslant1$ (resp.\ $\geqslant2$) for every $x_\lambda\in Z_\lambda$. Let $Z_\lambda'$ be the set of points $x_\lambda\in Z_\lambda$ such that, for every generization $z_\lambda$ of $x_\lambda$ in $Z_\lambda\cap(X_\lambda)_{f_\lambda(x_\lambda)}$, one has $\operatorname{prof}(((\sh{F}_\lambda)_{f_\lambda(x_\lambda)})_{z_\lambda})\geqslant1$ (resp.\ $\geqslant2$). If $p_{\lambda\mu}:Z_\mu\to Z_\lambda$ is the canonical projection for $\lambda\leqslant\mu$, the hypothesis and \sref{IV.19.9.5} give $Z_\mu'=p_{\lambda\mu}^{-1}(Z_\lambda')$ and $Z=p_\lambda^{-1}(Z_\lambda')$. On the other hand, \sref{IV.19.9.6} shows that $Z_\lambda'$ is open in $Z_\lambda$; the assertion therefore follows from $(8.3.4)$.

Finally, when $A$ is Noetherian, the flatness hypothesis and $(6.3.1)$ give $\operatorname{prof}_Z(\sh{H})\geqslant1$ (resp.\ $\operatorname{prof}_Z(\sh{H})\geqslant2$); in this case the conclusion of \sref{IV.19.9.8} follows from $(5.10.2)$ (resp.\ $(5.10.4)$).
\end{proof}

\begin{remark}[19.9.9]
\label{IV.19.9.9}
The same proof, combined with the results of chap.~III, 3\textsuperscript{rd} part on depth and local cohomology, allows one to establish the following generalisation of \sref{IV.19.9.8}:

\emph{Under the general conditions of \sref{IV.19.9.8}, suppose that, for every $x\in Z$, we have $\operatorname{prof}((\sh{F}_{f(x)})_x)\geqslant k$; then the canonical homomorphism}
\[
  H^i(X,\sh{H})\to H^i(X-Z,\sh{H}|(X-Z))
\]
\emph{is bijective for $i\leqslant k-2$ and injective for $i=k-1$} (which, expressed using the general cohomological notion of depth introduced in chap.~III, is also written $\operatorname{prof}_Z(\sh{H})\geqslant k$).
\end{remark}
% END NATIVE INPUT source_aligned/ega4-19
% BEGIN NATIVE INPUT source_aligned/ega4-20
\section{Meromorphic functions and pseudo-morphisms}
\label{IV.20}

\setcounter{subsection}{-1}
\subsection{Introduction}
\label{IV.20.0}

Most of the concepts and results of \textsection\textsection20 and 21 relate directly to Chapter~I and hardly depend on Chapters~II to IV, except for the occasional use of the notion of depth and of a regular local ring (in \sref{IV.20.6}, \sref{IV.21.11}, \sref{IV.21.13}, and \sref{IV.21.15}), of Zariski's ``Main Theorem'' in \sref{IV.20.4} and \sref{IV.21.12}, and of the properties of transversally regular immersions in \sref{IV.20.6} and \sref{IV.21.15}.

\oldpage[IV-4]{226}
In \textsection20, we introduce several variants of the concept of a rational map, already studied in \sref[I]{section:I.7} from a point of view still fairly close to the classical point of view, and for this reason quite ill-suited to the case of not necessarily reduced preschemes.
The notions and results of \textsection20 are used in \textsection21 (n\textsuperscript{os}\sref{IV.21.1} and \sref{IV.21.7}) to develop the general notion of a divisor and its most elementary properties.
This notion is especially convenient when the local rings of the preschemes considered are Noetherian and integrally closed, and especially when they are also \emph{factorial} (\sref{IV.21.6} and \sref{IV.21.7}), because in the latter case it is identified with the notion of a \emph{codimension-$1$ cycle} (a linear combination of irreducible subpreschemes of codimension~$1$).
In \sref{IV.21.9},  we determine the divisors on a Noetherian prescheme of dimension~$1$ but not necessarily normal, which is useful for various applications.
Numbers \sref{IV.21.11} and \sref{IV.21.12} give two important theorems, due respectively to Auslander--Buchsbaum and Van der Waerden, and relating to the notion of a factorial ring (numbers \sref{IV.21.9}, \sref{IV.21.11}, and \sref{IV.21.12} are independent of one another).
In numbers \sref{IV.21.13} and \sref{IV.21.14}, which are likewise independent of the preceding three, we study a useful variant of the notion of a factorial local ring, namely that of a \emph{parafactorial local ring}; this notion is introduced in particular in~\cite{IV-41} in the development of comparison theorems for the Picard group of a projective prescheme $X$ over a field $k$ and of a ``hyperplane section''.
We will see in \sref{IV.21.14.1} (Ramanujam--Samuel theorem) that the local parafactorial rings are much more numerous than one would have expected \emph{a priori}.

In \sref{IV.20.5}, \sref{IV.20.6}, and \sref{IV.21.15}, we review the previous notions but from a point of view ``relative'' to a fixed base prescheme.
For the moment these notions are used only rather rarely; in particular, the notion of a relative divisor is hardly used except for positive divisors, and in that case it can be described advantageously without recourse to relative meromorphic functions, by using the notion of a transversally regular immersion of codimension~$1$.
It is therefore advisable to omit these sections on a first reading.

\subsection{Meromorphic functions}
\label{IV.20.1}

\begin{env}[20.1.1]
\label{IV.20.1.1}
Let $(X,\sh{O}_X)$ be a ringed space, and let $\sh{S}$ be a subsheaf \emph{of sets} of $\sh{O}_X$.
For every open $U$ of $X$, consider the \emph{ring of fractions} $\Gamma(U,\sh{O}_X)[\Gamma(U,\sh{S})^{-1}]$ (Bourbaki, \emph{Alg. comm.}, chap.~II, \textsection2, n\textsuperscript{o}1).
It is immediate that the map $U\mapsto\Gamma(U,\sh{O}_X)[\Gamma(U,\sh{S})^{-1}]$ is a \emph{presheaf of rings} (\sref[0]{0.1.5.1} and \sref[0]{0.1.5.7}).
We denote by $\sh{O}_X[\sh{S}^{-1}]$ the \emph{sheaf of rings} associated to this presheaf and we say that this is the \emph{sheaf of rings of fractions of $\sh{O}_X$ with denominators in $\sh{S}$};
this is a \emph{flat} $\sh{O}_X$-module.
It is immediate that for every $x\in X$, we have a canonical isomorphism
\[
\label{IV.20.1.1.1}
  (\sh{O}_X[\sh{S}^{-1}])_x\isoto\sh{O}_{X,x}[\sh{S}_x^{-1}],
  \tag{20.1.1.1}
\]
since the reasoning of \sref[0]{0.1.4.5} generalizes immediately in the case where we have an inductive system $(A_\alpha,\vphi_{\beta\alpha})$ of rings, and for each index $\alpha$ a subset $S_\alpha$ of $A_\alpha$ such that
\oldpage[IV-4]{227}
$\vphi_{\beta\alpha}(S_\alpha)\subset S_\beta$ for $\alpha\leq\beta$;
we then take $S$ to be the inductive limit, in $A=\varinjlim A_\alpha$, of the inductive system of subsets $(S_\alpha)$.
\end{env}

\begin{env}[20.1.2]
\label{IV.20.1.2}
Now let $\sh{F}$ be an $\sh{O}_X$-module.
We then set
\[
\label{IV.20.1.2.1}
  \sh{F}[\sh{S}^{-1}]=\sh{F}\otimes_{\sh{O}_X}\sh{O}_X[\sh{S}^{-1}],
  \tag{20.1.2.1}
\]
and we say that this is the \emph{sheaf of modules of fractions of $\sh{F}$ with denominators in $\sh{S}$};
it is immediate that it is associated to the presheaf of modules $U\mapsto\Gamma(U,\sh{F})[\Gamma(U,\sh{S})^{-1}]$, and that for every $x\in X$, we have a canonical isomorphism
\[
\label{IV.20.1.2.2}
  (\sh{F}[\sh{S}^{-1}])_x\isoto\sh{F}_x[\sh{S}_x^{-1}].
  \tag{20.1.2.2}
\]
\end{env}

\begin{env}[20.1.3]
\label{IV.20.1.3}
We will focus here on the case where $\sh{S}$ is the subsheaf $\sh{S}(\sh{O}_X)$ of $\sh{O}_X$ such that for every open $U$, $\Gamma(U,\sh{S})$ is the \emph{set of regular elements} of the ring $\Gamma(U,\sh{O}_X)$;
it is immediate that it is a sheaf (and not merely a presheaf), since regularity of a section of $\sh{O}_X$ over $U$ can be checked ``fibre by fibre'' (that is, its germ at $x$ is regular in $\sh{O}_{X,x}$ for every $x\in U$); in other words, $\sh{S}(\sh{O}_X)_x$ is precisely the set of regular elements of $\sh{O}_{X,x}$.
The corresponding sheaf of rings
\[
  \sh{M}_X=\sh{O}_X[\sh{S}^{-1}]
\]
is called the \emph{sheaf of germs of meromorphic functions on $X$}, and the sections of $\sh{M}_X$ over $X$ are called the \emph{meromorphic functions on $X$};
they form a ring which we denote by $M(X)$.
For every \emph{$\sh{O}_X$-module $\sh{F}$},
\[
  \sh{F}\otimes_{\sh{O}_X}\sh{M}_X=\sh{F}[\sh{S}^{-1}]
\]
is also denoted $\sh{M}_X(\sh{F})$ and called the \emph{sheaf of germs of meromorphic sections of $\sh{F}$};
its sections over $X$ form an $M(X)$-module denoted by $M(X,\sh{F})$, whose elements are called \emph{meromorphic sections of $\sh{F}$ over $X$}.
These definitions imply that for every open $U$ of $X$, we have a canonical isomorphism $\sh{M}_X(\sh{F})|U\isoto\sh{M}_U(\sh{F}|U)$, in particular $\sh{M}_X|U\isoto\sh{M}_U$.
\end{env}

\begin{env}[20.1.3.1]
\label{IV.20.1.3.1}
If $X$ is a \emph{reduced prescheme}, then we note that if an element $s\in\Gamma(U,\sh{O}_X)$ is such that $s_\xi\neq 0$ for every maximal point $\xi$ of $U$, then $s$ is \emph{regular}.
Indeed, if $st=0$ for a $t\in\Gamma(U,\sh{O}_X)$, then we have $s_\xi t_\xi=0$, so $t_\xi=0$ since $\sh{O}_{X,\xi}$ is a field, and say that $t_\xi=0$ for every maximal point $\xi$ of $X$ means that $t=0$:
we are immediately reduced to the case where $U$ is affine, and an element of a reduced ring which belongs to all the minimal prime ideals is zero by definition.
The converse is true if the set of irreducible components of $X$ is \emph{locally finite}.
We immediately reduce to the case where $X=\Spec(A)$ is affine;
if $\mathfrak{p}_i$ ($1\leq i\leq n$) are the minimal prime ideals of $A$ and if $s\in\mathfrak{p}_i$ for an index $i$, then there exists $t\in A$ such that $t\in\mathfrak{p}_j$ for $j\neq i$ and $t\not\in\mathfrak{p}_i$ (Bourbaki, \emph{Alg. comm.}, chap.~II, \textsection1, n\textsuperscript{o}1, Prop.~1);
so we have $st\in\mathfrak{p}_i$ for all $i$, and as a result $st=0$ since $A$ is reduced;
$s$ is therefore nonregular.
\end{env}

\begin{env}[20.1.4]
\label{IV.20.1.4}
For every open $U$ of $X$, the homomorphism $t\mapsto t/1$ from $\Gamma(U,\sh{O}_X)$ to $\Gamma(U,\sh{O}_X)[\Gamma(U,\sh{S})^{-1}]$ (which is none other than the \emph{total ring of fractions} of
\oldpage[IV-4]{228}
$\Gamma(U,\sh{O}_X)$) is injective;
these homomorphisms thus define a \emph{canonical injective homomorphism}
\[
\label{IV.20.1.4.1}
  i:\sh{O}_X\to\sh{M}_X
  \tag{20.1.4.1}
\]
which allows us to identify $\sh{O}_X$ with a subsheaf of $\sh{M}_X$.
Given a meromorphic function $\vphi\in M(X)$, we say that $\vphi$ is \emph{defined} on an open $U$ of $X$ if $\vphi|U$ is a \emph{section of $\sh{O}_U$} over $U$;
the axioms of sheaves show that there is, for a given section $\vphi$, a \emph{largest} open on which $\vphi$ is defined;
we call this the \emph{domain of definition} of $\vphi$ and denote it by $\operatorname{dom}(\vphi)$.
\end{env}

\begin{env}[20.1.5]
\label{IV.20.1.5}
For every $\sh{O}_X$-module $\sh{F}$, we obtain from \sref{IV.20.1.4.1} a di-homomorphism consisting of $i$ and the homomorphism of sheaves of additive groups
\[
\label{IV.20.1.5.1}
  1_\sh{F}\otimes i:\sh{F}\to\sh{M}_X(\sh{F})=\sh{F}\otimes_{\sh{O}_X}\sh{M}_X.
  \tag{20.1.5.1}
\]
We note that the latter is not injective in general;
when it is injective, we say that $\sh{F}$ is \emph{torsion-free}:
this means that for every open $U$ of $X$ and every section $s\in\Gamma(U,\sh{O}_X)$ which is a regular element in this ring the homothety $z\mapsto sz$ of $\Gamma(U,\sh{F})$ is injective;
this condition is evidently satisfied if $\sh{F}$ is \emph{locally free}.
\end{env}

\begin{proposition}[20.1.6]
\label{IV.20.1.6}
Let $X$ be a locally Noetherian prescheme, $\sh{F}$ a quasi-coherent $\sh{O}_X$-module.
For $\sh{F}$ to be torsion-free, it is necessary and sufficient that $\operatorname{Ass}(\sh{F})\subset\operatorname{Ass}(\sh{O}_X)$.
\end{proposition}
 
\begin{proof}
We immediately reduce to the case where $X=\Spec(A)$ is affine, $\sh{F}=\widetilde{M}$, and we know that the elements $s$ of $A$ belonging to an ideal of $\operatorname{Ass}(M)$ are exactly those for which the homothety $z\mapsto sz$ is not injective (Bourbaki, \emph{Alg. comm.}, chap.~IV, \textsection1, n\textsuperscript{o}1, cor.~2 of prop.~2).
\end{proof}

\begin{env}[20.1.7]
\label{IV.20.1.7}
If $u$ is a section of $\sh{M}_X(\sh{F})$ over $X$, then we say that $u$ is \emph{defined} at a point $x\in X$ if there exists an open neighbourhood $V$ of $x$ in $X$ such that $u|V$ is the image of a section of $\sh{F}$ over $V$ under the di-homomorphism \sref{IV.20.1.5.1}.
We say that $u$ is \emph{defined} on an open $U$ of $X$ if it is defined at every point in $U$;
there is still a larger open in which $u$ is defined, called the \emph{domain of definition} of $u$ and denoted $\operatorname{dom}(u)$.
When $\sh{F}$ is torsion-free, so that $\sh{F}$ is identified by \sref{IV.20.1.5.1} with a subsheaf of $\sh{M}_X(\sh{F})$, saying that $u$ is defined on $U$ means that $u|U$ is a \emph{section of $\sh{F}$ over $U$}.
\end{env}

\begin{env}[20.1.8]
\label{IV.20.1.8}
In accordance with the general notation of \sref[0\textsubscript{I}]{0.5.4.7}, we denote by $\sh{M}_X^*$ the sheaf of multiplicative groups such that $\Gamma(U,\sh{M}_X^*)$ is (for every open $U$ of $X$) the group of \emph{invertible elements} of $\Gamma(U,\sh{M}_X)$.
This sheaf is none other than the sheaf $\sh{S}(\sh{M}_X)$ defined in \sref{IV.20.1.3}:
indeed, if $s\in\Gamma(U,\sh{S}(\sh{M}_X))$, then for every $x\in U$, there exists an open neighbourhood $V\subset U$ of $x$ such that $s|V$ is a regular element in the \emph{total ring of fractions} of $\Gamma(V,\sh{O}_X)$, and we know that such an element is necessarily invertible in this ring of fractions.
We say that the sections of $\sh{M}_X^*$ over $X$ are the \emph{regular meromorphic functions} (note that we are deviating here from the terminology followed by certain authors, who call ``regular'' meromorphic functions those which are \emph{sections of $\sh{O}_X$}, identified with a subsheaf of $\sh{M}_X$).

Let $\sh{L}$ be an \emph{invertible $\sh{O}_X$-module} \sref[0\textsubscript{I}]{0.5.4.1};
then it is clear that $\sh{M}_X(\sh{L})=\sh{L}\otimes_{\sh{O}_X}\sh{M}_X$
\oldpage[IV-4]{229}
is an \emph{invertible $\sh{M}_X$-module}.
Let $U$ be an open such that $\sh{L}|U$ is isomorphic to $\sh{O}_U$;
as every automorphism of $\sh{M}_U$ is multiplication by an invertible element of $\Gamma(U,\sh{M}_X)$ \sref[0\textsubscript{I}]{0.5.4.7}, it amounts to the same thing to say that a section $s\in\Gamma(U,\sh{M}_X(\sh{L}))$ has an invertible image in $\Gamma(U,\sh{M}_X)$ under \emph{some} isomorphism or under \emph{every} isomorphism onto $\Gamma(U,\sh{M}_X)$;
we say in this case that $s$ is a \emph{regular meromorphic section of $\sh{L}$} over $U$;
a section $s$ of $\sh{L}$ over $X$ is called a \emph{regular meromorphic section of $\sh{L}$} if, for every open $U$ such that $\sh{L}|U$ is isomorphic to $\sh{O}_U$, $s|U$ is a regular meromorphic section of $\sh{L}$ over $U$.
We denote by $(\sh{M}_X(\sh{L}))^*$ the subsheaf of $\sh{M}_X(\sh{L})$ such that for every open $U$, $\Gamma(U,(\sh{M}_X(\sh{L}))^*)$ is the set of regular meromorphic sections of $\sh{L}$ over $U$.
Let $s$ be a meromorphic section of $\sh{L}$ over $X$ (i.e. a section of $\sh{M}_X(\sh{L})$);
it defines a homomorphism $h_s:\sh{M}_X\to\sh{M}_X(\sh{L})$ which sends every section $t$ of $\sh{M}_X$ over an open $U$ to $(s|U)t$.
It follows immediately from the above that for $s$ to be \emph{regular}, it is necessary and sufficient for $h_s$ to be \emph{injective}, and in fact $h_s$ is then a \emph{bijective} homomorphism from $\sh{M}_X$ to $\sh{M}_X(\sh{L})$, and its restriction to $\sh{M}_X^*$ is a bijection to $(\sh{M}_X(\sh{L}))^*$.
We conclude that the homothety $t\mapsto ts$ is an isomorphism from $M(X)$ to $M(X,\sh{L})$.
\end{env}

\begin{env}[20.1.9]
\label{IV.20.1.9}
Let $s$ be a regular meromorphic section of an invertible $\sh{O}_X$-module $\sh{L}$ over $X$;
then for every $\sh{O}_X$-module $\sh{F}$, $s$ similarly defines a homomorphism $h_s\otimes 1_\sh{F}:\sh{M}_X(\sh{F})\to\sh{M}_X(\sh{F}\otimes_{\sh{O}_X}\sh{L})$, which is again \emph{bijective}.
\end{env}

\begin{env}[20.1.10]
\label{IV.20.1.10}
Let $s$ be a meromorphic section of an invertible $\sh{O}_X$-module $\sh{L}$ over $X$;
for $s$ to be regular, it is necessary and sufficient for there to exist a meromorphic section $s'$ of $\sh{L}^{-1}$ over $X$ such that the canonical image of $s\otimes s'$ in $\sh{M}_X$ \sref[0\textsubscript{I}]{0.5.4.3} is the unit section, and this section $s'$ is then unique:
indeed, the necessity of the local existence of such a section is evident, and its local uniqueness implies its global (and unique) existence;
moreover, the existence of $s'$ is trivially sufficient for $s$ to be regular.
We will take $s'=s^{-1}$.

Finally, if $\sh{L}'$ is a second invertible $\sh{O}_X$-module, $s$ (resp. $s'$) a regular meromorphic section of $\sh{L}$ (resp. $\sh{L}'$) over $X$, then $s\otimes s'$ is evidently a regular meromorphic section of $\sh{L}\otimes\sh{L}'$ over $X$.
\end{env}

\begin{env}[20.1.11]
\label{IV.20.1.11}
If $f:X'\to X$ is a morphism of ringed spaces, then there is in general no natural map sending a meromorphic function on $X$ to a meromorphic function on $X'$.
For example, if $X$ is the spectrum of a local integral domain $A$ and $X'$ is the spectrum of its residue field $\kres$, then there is no natural homomorphism from the field of fractions $K$ of $A$ to $\kres$, and we can only send an element of $K$ to an element of $\kres$ if it is already in $A$.

In general, if $f=(\psi,\theta)$, then for every open $U$ of $X$, denote by $\sh{S}_f(U)$ the set of \emph{regular} sections $s\in\Gamma(U,\sh{O}_X)$ such that the image of $s$ under
\[
  \Gamma(\theta^\#):\Gamma(U,\sh{O}_X)\to\Gamma(f^{-1}(U),\sh{O}_{X'})
\]
is a \emph{regular} section.
It is immediate that $U\mapsto\sh{S}_f(U)$ is a \emph{subsheaf} of the sheaf of sets $\sh{S}(\sh{O}_X)$, which we denote by $\sh{S}_f$.
We set $\sh{M}_f=\sh{O}_X[\sh{S}_f^{-1}]$;
this is a subsheaf
\oldpage[IV-4]{230}
of rings of $\sh{M}_X$, and we canonically obtain from $\theta^\#:\psi^*(\sh{O}_X)\to\sh{O}_{X'}$ a homomorphism of sheaves of rings $\theta^{\prime\#}:\psi^*(\sh{M}_f)\to\sh{M}_{X'}$ extending $\theta^\#$ (Bourbaki, \emph{Alg. comm.}, chap.~II, \textsection2, n\textsuperscript{o}1, prop.~2);
hence, recalling that $f^*(\sh{M}_f)=\psi^*(\sh{M}_f)\otimes_{\psi^*(\sh{O}_X)}\sh{O}_{X'}$,
we get a canonical homomorphism of $\sh{O}_{X'}$-algebras
\[
\label{IV.20.1.11.1}
  f^*(\sh{M}_f)\to\sh{M}_{X'}.
  \tag{20.1.11.1}
\]
For every meromorphic function $\vphi$ on $X$ that \emph{is a section of $\sh{M}_f$}, $\Gamma(\theta^{\prime\#})(\vphi)$ is a meromorphic function on $X'$, called the \emph{inverse image of $\vphi$ under $f$}, and denoted by $\vphi\circ f$ if there is no cause for confusion.

Similarly, if $\sh{F}$ is an $\sh{O}_X$-module, then we set $\sh{M}_f(\sh{F})=\sh{F}\otimes_{\sh{O}_X}\sh{M}_f$, and we immediately obtain from $\theta^{\prime\#}$ a canonical homomorphism (which is also written as $u\mapsto u\circ f$)
\[
  \Gamma(X,\sh{M}_f(\sh{F}))\to\Gamma(X',\sh{M}_{X'}(f^*(\sh{F}))).
\]
In addition, if $u\in\Gamma(X,\sh{M}_f(\sh{F}))$ is defined \sref{IV.20.1.7} at a point $x$, then $u$ coincides, on a neighbourhood $U$ of $x$, with a section of the form $\sum_i h_i\otimes(t_i/s_i)$, where the $h_i$ belong to $\Gamma(U,\sh{F})$, the $t_i$ to $\Gamma(U,\sh{O}_X)$, and the $s_i$ to $\Gamma(U,\sh{S}_f)$.
As by hypothesis the images of the $s_i$ in $\Gamma(f^{-1}(U),\sh{O}_{X'})$ are regular, we see that $u\circ f$ is defined at every point of $f^{-1}(U)$;
in other words, we have
\[
\label{IV.20.1.11.2}
  f^{-1}(\operatorname{dom}(u))\subset\operatorname{dom}(u\circ f).
  \tag{20.1.11.2}
\]
We will see later \sref{IV.20.6.5}[(i)] examples (with $\sh{F}=\sh{O}_X$) where the two sides of \sref{IV.20.1.11.2} can be different.

Consider in particular the case where $\sh{M}_f=\sh{M}_X$;
then, if $\sh{L}$ is an invertible $\sh{O}_X$-module, the image in $\sh{M}_{X'}(f^*(\sh{L}))$, under $\Gamma(\theta^{\prime\#})$, of a \emph{regular} meromorphic section of $\sh{L}$ over $X$ \sref{IV.20.1.8} is a \emph{regular} meromorphic section of $f^*(\sh{L})$ over $X'$, as it follows immediately from the definition of its sections, and from the fact that a homomorphism of rings sends an invertible element to an invertible element.

Let $f':X''\to X'$ be a second morphism of ringed spaces, and suppose that $\sh{M}_f=\sh{M}_X$ and $\sh{M}_{f'}=\sh{M}_{X'}$;
then, if we set $f''=f\circ f'$, we also have $\sh{M}_{f''}=\sh{M}_X$, and we immediately see that for every meromorphic section $u$ of $\sh{F}$ over $X$, we have $u\circ f''=(u\circ f)\circ f'$.
\end{env}

\begin{proposition}[20.1.12]
\label{IV.20.1.12}
If the morphism $f:X'\to X$ is flat \sref[0\textsubscript{I}]{0.6.7.1}, then we have $\sh{M}_f=\sh{M}_X$, and the homomorphism $\vphi\mapsto\vphi\circ f$ is defined on all of $M(X)$.
In addition, if $f$ is a (flat) morphism of locally ringed spaces, then we have $\operatorname{dom}(\vphi\circ f)=f^{-1}(\operatorname{dom}(\vphi))$;
if in addition $f$ is surjective (thus faithfully flat), then the homomorphism $\vphi\mapsto\vphi\circ f$ is injective.
\end{proposition}

\begin{proof}
The first assertion follows from the fact that, if $B$ is an $A$-algebra which is a flat $A$-module, then every element of $A$ not a divisor of $0$ in $A$ is not a divisor of $0$ in $B$ \sref[0\textsubscript{I}]{0.6.3.4}.
To prove the other assertions, note that, for every $x'\in X'$, if $x=f(x')$, then $\sh{O}_{X',x'}$ is a flat $\sh{O}_{X,x}$-module, and as the homomorphism $\sh{O}_{X,x}\to\sh{O}_{X',x'}$ is \emph{local} by hypothesis, it is injective (\sref[0\textsubscript{I}]{0.6.5.1} and \sref[0\textsubscript{I}]{0.6.6.2});
if we set $A=\sh{O}_{X,x}$, $B=\sh{O}_{X',x'}$, such that $A$ identifies with a subring of $B$, then $(f^*(\sh{M}_X))_{x'}$ is equal to $S^{-1}A\otimes_A B=S^{-1}B$, where $S$ is the set of regular elements of $A$, $(\sh{M}_{X'})_{x'}$ is equal to $T^{-1}B$, where $T$ is the set
\oldpage[IV-4]{231}
of regular elements of $B$, and as we have seen that $S\subset T$, the homomorphism $S^{-1}B\to T^{-1}B$ is injective;
in other words, this proves that the homomorphism \sref{IV.20.1.11.1} $f^*(\sh{M}_X)\to\sh{M}_{X'}$ is \emph{injective} (hence the last assertion of the statement).
The quotient $f^*(\sh{M}_X)/\sh{O}_{X'}$ identifies with an $\sh{O}_{X'}$-submodule of $\sh{M}_{X'}/\sh{O}_{X'}$, and $(f^*(\sh{M}_X)/\sh{O}_{X'})_{x'}$ identifies with $(\sh{M}_X/\sh{O}_X)_x\otimes_{\sh{O}_{X,x}}\sh{O}_{X',x'}$.
Then suppose that $x\not\in\operatorname{dom}(\vphi)$;
the image of $\vphi_x$ in $(\sh{M}_X/\sh{O}_X)_x$ is therefore $\neq 0$;
by faithful flatness, we deduce that it is the same for the image of $(\vphi\circ f)_{x'}$, so $x'\not\in\operatorname{dom}(\vphi\circ f)$, which finishes the proof.
\end{proof}

\begin{remark}[20.1.13]
\label{IV.20.1.13}
Let $X$ be a \emph{reduced} complex analytic space;
then the notion of a meromorphic function on $X$ defined above coincides with the usual notion.
Consider on the other hand a prescheme $Y$, locally of finite type over the field $\bb{C}$;
we then know that we can associate to $Y$ an analytic space $Y^\mathrm{an}$ having the same underlying topological space, and the canonical morphism $f:Y^\mathrm{an}\to Y$ is \emph{flat}~\cite{IV-37};
by virtue of \sref{IV.20.1.12}, the canonical homomorphism $u\mapsto u\circ f$ from $M(Y)$ to $M(Y^\mathrm{an})$ is therefore always defined and is injective;
but it is not \emph{surjective} in general.
For example, when $Y=\bb{V}_{\bb{C}}^r$ (\textbf{Err}\textsubscript{\textbf{III}},~14) is affine $r$-space over $\bb{C}$, $M(Y)$ canonically identifies with the field $R(Y)$ of rational functions on $Y$ \sref{IV.20.2.13}[(i)], while $M(Y^\mathrm{an})$ is the field of usual meromorphic functions on $\bb{C}^r$.
Because of this fact, it is often preferable, in algebraic geometry, to abstain from the terminology introduced in this section, and to use the equivalent terminology of ``pseudo-function'' which will be defined below.
\end{remark}

\subsection{Pseudo-morphisms and pseudo-functions}
\label{IV.20.2}

The only ringed spaces discussed in this section are \emph{preschemes}.

\begin{env}[20.2.1]
\label{IV.20.2.1}
Recall \sref{IV.11.10.2} that in a prescheme $X$, we say that an open $U$ is \emph{schematically dense} if, for every open $V$ of $X$, the canonical homomorphism $\Gamma(V,\sh{O}_X)\to\Gamma(V\cap U,\sh{O}_X)$ is \emph{injective}.

Consider two preschemes $X$, $Y$, two \emph{schematically dense} opens $U$, $U'$ of $X$;
we say that two morphisms $u:U\to Y$, $u':U'\to Y$ are \emph{equivalent} if there exists an open $U''\subset U\cap U'$, \emph{schematically dense} in $X$, such that $u|U''=u'|U''$.
As it results immediately from the definition of schematically dense opens that the intersection of two such opens is again one, it is immediate that the preceding relation is indeed an equivalence relation.
An equivalence class for this relation is called a \emph{pseudo-morphism} of $X$ in $Y$, or a \emph{strict rational application} of $X$ in $Y$.

If $S$ is a prescheme and $X$, $Y$ two $S$-preschemes, we call \emph{pseudo-$S$-morphism} the equivalence class (for the preceding relation) of an $S$-\emph{morphism} of a schematically dense open of $X$ in $Y$.
A pseudo-morphism is therefore nothing other than a pseudo-$S$-morphism for $S=\Spec(\bb{Z})$.

We denote by $\operatorname{Ps.hom}(X,Y)$ (resp.\ $\operatorname{Ps.hom}_S(X,Y)$) the set of pseudo-morphisms (resp.\ pseudo-$S$-morphisms) of $X$ in $Y$.
\end{env}

\begin{env}[20.2.2]
\label{IV.20.2.2}
\oldpage[IV-4]{232}
It follows from the definition recalled above that if $U$ is a schematically dense open in $X$, then, for every open $V$ of $X$, $U\cap V$ is schematically dense in $V$.
If two morphisms $u:U\to Y$, $u':U'\to Y$ of schematically dense opens of $X$ in $Y$ are equivalent, it follows that their restrictions $u|(V\cap U)$ and $u'|(V\cap U')$ are also equivalent morphisms (for the prescheme induced on $V$);
their class is called the \emph{restriction} to $V$ of the pseudo-morphism $\omega$, the class of $u$, and we denote this pseudo-morphism $\omega|V$.
If $W\subset V$ is an open of $X$, it is clear that $(\omega|V)|W=\omega|W$.
This shows that the restriction maps define a \emph{presheaf} of sets $U\mapsto\operatorname{Ps.hom}(U,Y)$;
we note that this presheaf \emph{is not} in general a sheaf (cf.\ \sref{IV.20.2.16});
the associated sheaf is denoted $\mathscr{P}\!s.\!\operatorname{hom}(X,Y)$.
For pseudo-$S$-morphisms, we see that $U\mapsto\operatorname{Ps.hom}_S(U,Y)$ is a presheaf of sets, whose associated sheaf we denote $\mathscr{P}\!s.\!\operatorname{hom}_S(X,Y)$.

If $V$ is a \emph{schematically dense} open in $X$, then for every open $U$ schematically dense in $X$, $U\cap V$ is also schematically dense in $X$, so the map $\omega\mapsto\omega|V$ is a \emph{bijection} of $\operatorname{Ps.hom}(X,Y)$ (resp.\ $\operatorname{Ps.hom}_S(X,Y)$) onto $\operatorname{Ps.hom}(V,Y)$ (resp.\ $\operatorname{Ps.hom}_S(V,Y)$).
\end{env}

\begin{env}[20.2.3]
\label{IV.20.2.3}
Given a pseudo-$S$-morphism $\omega$ of $X$ in $Y$, we say that it is \emph{defined} at a point $x\in X$ if there exists an open neighbourhood $V$ of $x$ in $X$, an open $U\subset V$ \emph{containing $x$ and schematically dense in} $V$, and finally an $S$-morphism $u:U\to Y$ whose class in $\operatorname{Ps.hom}_S(V,Y)$ is equal to $\omega|V$ \sref{IV.20.2.2};
we call the \emph{domain of definition} of $\omega$ and denote $\operatorname{dom}_S(\omega)$ (or simply $\operatorname{dom}(\omega)$ if $S=\Spec(\bb{Z})$) the set of $x\in X$ where $\omega$ is defined;
this is evidently an \emph{open} of $X$.
Moreover, for every open $W$ of $X$, we have
\[
\label{IV.20.2.3.1}
  \operatorname{dom}_S(\omega|W)=(\operatorname{dom}_S(\omega))\cap W
  \tag{20.2.3.1}
\]
by virtue of the property of schematically dense opens recalled in \sref{IV.20.2.2}.
\end{env}

\begin{proposition}[20.2.4]
\label{IV.20.2.4}
Suppose $X$, $Y$ are $S$-preschemes such that the structural morphism $Y\to S$ is separated;
then, if $\omega$ is a pseudo-$S$-morphism of $X$ in $Y$, $\operatorname{dom}_S(\omega)$ is the largest of the schematically dense opens $U$ of $X$ such that there exists an $S$-morphism $u:U\to Y$ belonging to the class~$\omega$.
\end{proposition}

\begin{proof}
It suffices to prove that if $U$, $U'$ are two schematically dense opens in $X$ such that two $S$-morphisms $u:U\to Y$ and $u':U'\to Y$ are equivalent, then the restrictions of $u$ and $u'$ to $U\cap U'$ are equal.
But there exists by hypothesis an open $U''\subset U\cap U'$, schematically dense in $X$, in which $u$ and $u'$ coincide;
as $U''$ is also schematically dense in $U\cap U'$, our assertion results from \sref{IV.11.10.1}[(d)].
\end{proof}

\begin{corollary}[20.2.5]
\label{IV.20.2.5}
Let $S$ be an $S_0$-scheme, and let $X$, $Y$ be two $S$-preschemes such that the composite morphism $Y\to S\to S_0$ is separated (which implies, by \sref[I]{I.5.5.1}, that $Y\to S$ is also separated).
For every pseudo-$S$-morphism $\omega$ of $X$ in $Y$, one then has $\operatorname{dom}_S(\omega)=\operatorname{dom}_{S_0}(\omega)$.
In particular, if $Y$ is a prescheme, one has $\operatorname{dom}_S(\omega)=\operatorname{dom}(\omega)$.
\end{corollary}

\begin{proof}
Indeed, it is clear that $\operatorname{dom}_S(\omega)\subset\operatorname{dom}_{S_0}(\omega)$ without the hypothesis of separation;
by virtue of \sref{IV.20.2.4}, it suffices to prove that if $u_0:U_0\to Y$ is an $S_0$-morphism of a
\oldpage[IV-4]{233}
schematically dense open $U_0$ of $X$ in $Y$ such that the composed $v_0:U_0\xrightarrow{u_0}Y\to S$ coincides with the restriction of the structural morphism $w_0:U_0\to S$ in an open $U\subset U_0$ schematically dense in $X$, then one has $v_0=w_0$.
But by virtue of the hypothesis that the morphism $S\to S_0$ is separated, this again results from \sref{IV.11.10.1}[, d)].
\end{proof}

\begin{corollary}[20.2.6]
\label{IV.20.2.6}
Under the hypotheses of \sref{IV.20.2.4}, the presheaf $U\mapsto\operatorname{Ps.hom}_S(U,Y)$ is a sheaf.
\end{corollary}

\begin{proof}
Indeed, let $U$ be an open of $X$, $(U_\alpha)$ a covering of $U$ by opens of $X$, and suppose given for each $\alpha$ a pseudo-$S$-morphism $\omega_\alpha$ of $U_\alpha$ in $Y$, such that for every pair of indices $\alpha$, $\beta$, the restrictions \sref{IV.20.2.2} of $\omega_\alpha$ and $\omega_\beta$ to $U_\alpha\cap U_\beta$ are equal;
by virtue of \sref{IV.20.2.3.1}, this implies that $\operatorname{dom}_S(\omega_\alpha)\cap U_\beta=\operatorname{dom}_S(\omega_\beta)\cap U_\alpha$.
Let $W$ be the union of the opens $\operatorname{dom}_S(\omega_\alpha)$, and, for each $\alpha$, let $u_\alpha$ be the unique $S$-morphism $\operatorname{dom}_S(\omega_\alpha)\to Y$ belonging to the class $\omega_\alpha$ \sref{IV.20.2.4};
by virtue of \sref{IV.20.2.4}, the restrictions of $u_\alpha$ and $u_\beta$ to $\operatorname{dom}_S(\omega_\alpha)\cap\operatorname{dom}_S(\omega_\beta)$ are equal, so there exists a morphism $u:W\to Y$ whose restriction to each $\operatorname{dom}_S(\omega_\alpha)$ is equal to $u_\alpha$;
it is clear that $W$ is schematically dense in $U$, so it defines a pseudo-$S$-morphism $\omega$ of $U$ in $Y$ whose restrictions to the $U_\alpha$ are the $\omega_\alpha$.
\end{proof}

\begin{remark}[20.2.7]
\label{IV.20.2.7}
We know \sref{IV.11.10.4} that when $X$ is \emph{reduced}, it amounts to the same to say that an open of $X$ is dense in $X$ or that it is schematically dense in $X$;
the notion of pseudo-morphism (resp.\ pseudo-$S$-morphism) of $X$ in $Y$ then coincides with that of \emph{rational application} (resp.\ $S$-\emph{rational application}) of $X$ in $Y$ \sref[I]{I.7.1.2}.
In the general case, the notion of pseudo-morphism seems more useful than that of rational application.
\end{remark}

\begin{env}[20.2.8]
\label{IV.20.2.8}
We call \emph{pseudo-function} on $X$ a pseudo-morphism of $X$ in $\Spec(\bb{Z}[T])$ ($T$ indeterminate), or, what amounts to the same, an $X$-pseudo-morphism of $X$ in $X\otimes_{\bb{Z}}\bb{Z}[T]$;
this amounts also to the same \sref[I]{I.3.3.15} as saying that a pseudo-function on $X$ is an \emph{equivalence class of sections of $\sh{O}_X$ over schematically dense opens} of $X$, two sections $g\in\Gamma(U,\sh{O}_X)$, $g'\in\Gamma(U',\sh{O}_X)$ over such opens being equivalent if there exists an open $U''\subset U\cap U'$, schematically dense in $X$, where $g$ and $g'$ coincide.
We can apply \sref{IV.20.2.4} with $S=X$ and $Y=X\otimes_{\bb{Z}}\bb{Z}[T]$;
then, for every pseudo-function $\varphi$ on $X$, there is a \emph{largest} schematically dense open $\operatorname{dom}(\varphi)$ in $X$ such that there is a section of $\sh{O}_X$ over $\operatorname{dom}(\varphi)$ belonging to the class $\varphi$.
It is clear that the sheaf $\mathscr{P}\!s.\!\operatorname{hom}_X(X,X\otimes_{\bb{Z}}\bb{Z}[T])$ is here a \emph{sheaf of rings}, and even an $\sh{O}_X$-\emph{Algebra}, which we denote $\sh{M}'_X$;
it results from \sref{IV.20.2.6} that, for every open $U$ of $X$, $\Gamma(U,\sh{M}'_X)$ is equal to the set of pseudo-morphisms of $U$ in $\Spec(\bb{Z}[T])$;
we set $M'(X)=\Gamma(X,\sh{M}'_X)$.
To say that a section $\varphi$ of $\sh{M}'_X$ over $U$ is \emph{invertible} in the ring $\Gamma(U,\sh{M}'_X)$ means that there exists an open $U'$ schematically dense in $\operatorname{dom}(\varphi)$, thus in $U$, such that, if $g$ is the unique section of $\sh{O}_X$ over $\operatorname{dom}(\varphi)$ belonging to the class $\varphi$, $g|U'$ is \emph{invertible} in $\Gamma(U',\sh{O}_X)$.
It results from \sref[I]{I.3.3.15} that, in the canonical correspondence between $\Gamma(V,\sh{O}_X)$ and $\Hom(V,\bb{Z}[T])$ ($V$ open of $X$), the invertible elements of $\Gamma(V,\sh{O}_X)$ correspond to the morphisms which
\oldpage[IV-4]{234}
\emph{factor through} $V\to\Spec(\bb{Z}[T,T^{-1}])\to\Spec(\bb{Z}[T])$.
We conclude that the sheaf ${\sh{M}'}^*_X$ of invertible germs of sections of $\sh{M}'_X$ is canonically identified with the sheaf $\mathscr{P}\!s.\!\operatorname{hom}_X(X,X\otimes_{\bb{Z}}\bb{Z}[T,T^{-1}])$.
\end{env}

\begin{lemma}[20.2.9]
\label{IV.20.2.9}
Let $U$ be an open of $X$, $s$ a regular element of the ring $\Gamma(U,\sh{O}_X)$;
then the open $U_s$ of the $x\in U$ such that $s(x)\neq 0$ is schematically dense in $U$.
\end{lemma}

\begin{proof}
Indeed, let $V$ be an open of $U$, $t$ a section of $\sh{O}_X$ over $V$ such that $t|(V\cap U_s)=0$.
For every affine open $W\subset V$, there exists an integer $n$ such that $s^n t|W=0$ \sref[I]{I.1.4.1};
but since $s$ is a \emph{regular} element of $\Gamma(U,\sh{O}_X)$, this implies $t|W=0$, hence $t=0$.
\end{proof}

\begin{env}[20.2.10]
\label{IV.20.2.10}
Consider a meromorphic function $f\in M(X)$ \sref{IV.20.1.4};
then $\operatorname{dom}(f)$ is \emph{schematically dense} in $X$:
indeed, every point of $X$ admits an open neighbourhood $U$ such that there exists a section $s\in\Gamma(U,\sh{O}_X)$ which is a regular element of this ring, and for which $s(f|U)\in\Gamma(U,\sh{O}_X)$;
as $s|U_s$ is invertible, we conclude that $f|U_s\in\Gamma(U_s,\sh{O}_X)$, hence $U_s\subset\operatorname{dom}(f)$ by definition \sref{IV.20.1.4}, and our assertion results from the lemma \sref{IV.20.2.9}.
We can thus associate to $f$ the pseudo-function equal to the class of the section of $\sh{O}_X$ over $\operatorname{dom}(f)$, restriction of $f$;
operating in the same way replacing $X$ by an arbitrary open of $X$, we thus obtain a canonical \emph{homomorphism of $\sh{O}_X$-Algebras}
\[
\label{IV.20.2.10.1}
  \sh{M}_X\to\sh{M}'_X
  \tag{20.2.10.1}
\]
which, by restriction, gives evidently a homomorphism of sheaves of abelian groups
\[
\label{IV.20.2.10.2}
  \sh{M}_X^*\to{\sh{M}'}^*_X
  \tag{20.2.10.2}
\]
for the sheaves of invertible germs of sections of $\sh{M}_X$ and $\sh{M}'_X$.
\end{env}

\begin{proposition}[20.2.11]
\label{IV.20.2.11}
\medskip
\begin{enumerate}[label=\emph{(\roman*)}]
  \item The canonical homomorphism \emph{\sref{IV.20.2.10.1}} (and hence also \emph{\sref{IV.20.2.10.2}}) is \emph{injective}.
  \item Suppose either that $X$ is locally Noetherian, or that $X$ is reduced and the set of its irreducible components is locally finite.
  Then the canonical homomorphism \emph{\sref{IV.20.2.10.1}} (and hence also \emph{\sref{IV.20.2.10.2}}) is \emph{bijective}.
\end{enumerate}
\end{proposition}

\begin{proof}
The questions being local on $X$, we can restrict to the case where $X=\Spec(A)$ is affine, and show that the canonical homomorphism $M(X)\to M'(X)$ is injective, and that it is bijective in the cases considered in~(ii).
Taking into account the definition of the sheaf $\sh{M}_X$ \sref{IV.20.1.3}, we can in fact remark that \sref{IV.20.2.10.1} comes from a homomorphism of \emph{presheaves}
\[
  \Gamma(U,\sh{O}_X)[\Gamma(U,\sh{S})^{-1}]\to\Gamma(U,\sh{M}'_X)
\]
and it suffices to show that the latter is injective (resp.\ bijective under the hypotheses of~(ii)).
Denoting by $S$ the set of regular elements of $A$ (so that $S^{-1}A$ is the \emph{total ring of fractions} of $A$), we must therefore consider the canonical homomorphism
\[
\label{IV.20.2.11.1}
  S^{-1}A\to M'(X).
  \tag{20.2.11.1}
\]
\oldpage[IV-4]{235}
and show that it is injective (resp.\ bijective under the conditions of~(ii)).
We can write $S^{-1}A=\varinjlim A_t$, where $t$ runs through the set of \emph{regular} elements of $A$, ordered by the relation ``$t$ is a divisor of $t'$'' \sref[0]{0.1.4.5};
we can also write $A_t=\Gamma(\mathrm{D}(t),\sh{O}_X)$.
On the other hand, by definition $M'(X)=\varinjlim\Gamma(U,\sh{O}_X)$, where $U$ runs through the set of schematically dense opens of $X$ (ordered by the relation $\supset$), and the map \sref{IV.20.2.11.1} is none other than the canonical map coming from the fact that the $\mathrm{D}(t)$ constitute a part of the set of schematically dense opens in $X$ \sref{IV.20.2.9}.
Note that by definition of a schematically dense open, the homomorphism $\Gamma(U,\sh{O}_X)\to\Gamma(U',\sh{O}_X)$ for any two such opens $U'\subset U$ is always \emph{injective}, hence the homomorphisms $\Gamma(U,\sh{O}_X)\to M'(X)$ are also injective, and we conclude that \sref{IV.20.2.11.1} is injective.
To prove that \sref{IV.20.2.11.1} is bijective, it suffices to show that the $\mathrm{D}(t)$ are \emph{cofinal} in the set of schematically dense opens, i.e., that for such an open $U$, there exists $t$ regular in $A$ such that $\mathrm{D}(t)\subset U$.
Now, when $X$ is reduced and the irreducible components of $X$ form a locally finite set, this set is finite since $X$ has been supposed affine, so $A$ has only a finite number of minimal primes $\mathfrak{p}_i$;
as $A$ is reduced, the intersection of the $\mathfrak{p}_i$ is zero, and to say that $t$ is regular is equivalent to saying that $t$ does not belong to any of the $\mathfrak{p}_i$;
we conclude by the reasoning of \sref[I]{I.7.1.9.1}.
When $A$ is Noetherian, to say that $U=X-Y$ (where $Y=V(\mathfrak{J})$ is closed in $X$) is schematically dense means \sref{IV.5.10.2} that $Y$ \emph{does not meet} $\operatorname{Ass}(\sh{O}_X)$, and by virtue of Bourbaki, \emph{Alg.~comm.}, chap.~IV, \textsection1, n\textsuperscript{o}~4, prop.~8, this implies the existence of a $t\in\mathfrak{J}$ such that $t$ is $A$-regular, hence $U\supset\mathrm{D}(t)$.

We have moreover proved in the course of this demonstration the
\end{proof}

\begin{corollary}[20.2.12]
\label{IV.20.2.12}
If $X=\Spec(A)$, where $A$ is Noetherian, or reduced and having only a finite number of minimal prime ideals, the schematically dense opens in $X$ are those which contain an open of the form $\mathrm{D}(t)$, where $t$ is regular in $A$, and $M(X)=M'(X)$ is the total ring of fractions $S^{-1}A$, where $S$ is the set of regular elements of~$A$.
\end{corollary}

\begin{remarks}[20.2.13]
\label{IV.20.2.13}
\begin{enumerate}[label=(\roman*)]
  \item Let $\varphi$ be an element of $M(X)$, $\varphi'$ its image in $M'(X)$;
  by definition (\sref{IV.20.1.4} and \sref{IV.20.2.8}) $\operatorname{dom}(\varphi)\subset\operatorname{dom}(\varphi')$;
  but in fact one has the equality $\operatorname{dom}(\varphi)=\operatorname{dom}(\varphi')$, because there exists a section of $\sh{O}_X$ over $\operatorname{dom}(\varphi')$, belonging to the class $\varphi'$, and the corresponding meromorphic function equals $\varphi$ by \sref{IV.20.2.11}[(i)], hence $\operatorname{dom}(\varphi')\subset\operatorname{dom}(\varphi)$.

  \item We have already noted that when $X$ is \emph{reduced}, one has $\sh{M}'_X=\sh{R}(X)$ (sheaf of rational functions on $X$ \sref[I]{I.7.3.2});
  if moreover the irreducible components of $X$ form a locally finite set, one has $\sh{M}_X=\sh{M}'_X=\sh{R}(X)$.
  In general, since every schematically dense open is dense, one has a canonical homomorphism $\sh{M}'_X\to\sh{R}(X)$, but even when $X$ is locally Noetherian, this homomorphism is not necessarily injective.
  For example, if $X=\Spec(A)$, where $A$ is a Noetherian ring such that $\operatorname{Ass}(A)$ contains embedded prime ideals (which implies that $A$ is not reduced), then $M(X)$ and $M'(X)$ identify with the total ring of fractions $S^{-1}A$,
  \oldpage[IV-4]{236}
  where $S$ is the set of regular elements of $A$, the complement of the union of the ideals $\mathfrak{p}\in\operatorname{Ass}(A)$;
  on the other hand, $R(X)$ identifies with $Q^{-1}A$, where $Q$ is the complement of the union of the \emph{minimal} prime ideals of $A$ \sref[I]{I.7.1.9}, and the canonical homomorphism $A\to Q^{-1}A$ (and \emph{a fortiori} $S^{-1}A\to Q^{-1}A$) is therefore not injective, since there exist nonzero elements of $A-Q$ annihilated by elements of $Q$ (Bourbaki, \emph{Alg.~comm.}, chap.~IV, \textsection1, n\textsuperscript{o}~1, cor.~2 of prop.~1).

  \item We note also that when $X$ is locally Noetherian, the $\sh{O}_X$-Module $\sh{M}'_X$ is \emph{not necessarily quasi-coherent}.
  Consider for example a local Noetherian ring $A$ of dimension $\geq 2$, with maximal ideal $\mathfrak{m}\in\operatorname{Ass}(A)$, and if we set $X=\Spec(A)$, the scheme induced on the open $U$ of $X$, complementary of $\{\mathfrak{m}\}$, is \emph{integral}.
  The only regular elements of $A$ are the invertible elements, so $\Gamma(X,\sh{M}'_X)=M'(X)=A$;
  if $\sh{M}'_X$ were quasi-coherent, it would therefore be equal to $\sh{O}_X$;
  but as $U$ is an integral scheme, $\sh{M}'_X|U=R(U)$ is a simple sheaf \sref[I]{I.7.3.5}, whereas $\sh{O}_U$ is not a simple sheaf since $\dim(A)\geq 2$.

  It remains to give an example of a ring $A$ having the above properties.
  It suffices to consider an integral local ring $B$ of dimension $\geq 2$ and residue field $k$, and to take $A=B\oplus k$ with the law of multiplication $(b,x)(b',x')=(bb',bx'+b'x)$.

  \item If $X$ is locally Noetherian, it results from \sref{IV.5.10.2} that the schematically dense opens in $X$ are those which \emph{contain} $\operatorname{Ass}(\sh{O}_X)$.

  In particular, if $X$ is reduced and has only finitely many irreducible
  components, the notion of torsion-free $\sh{O}_X$-Module defined in
  \sref{IV.20.1.5} coincides with that of \sref[I]{I.7.4.1}.
\end{enumerate}
\end{remarks}

\begin{env}[20.2.14]
\label{IV.20.2.14}
Let $X$ be a prescheme, $\sh{F}$ an $\sh{O}_X$-Module quasi-coherent and \emph{torsion-free} \sref{IV.20.1.5}, so that $\sh{F}$ identifies via \sref{IV.20.1.5.1} with a sub-$\sh{O}_X$-Module of $\sh{M}_X(\sh{F})$.
For every meromorphic section $u$ of $\sh{F}$ over $X$, we call the \emph{ideal of denominators} of $u$ the Ideal $\sh{J}$ of the section $\overline{u}$, image of $u$ in $\sh{M}_X(\sh{F})/\sh{F}$.
The Ideal $\sh{J}$ is \emph{quasi-coherent}:
indeed, the question being local on $X$, we can restrict to the case where $X$ is affine, and there exists a section $s\in\Gamma(X,\sh{S}(\sh{O}_X))$ such that $v=su\in\Gamma(X,\sh{F})$;
to say that a section $f\in\Gamma(U,\sh{O}_X)$ belongs to $\Gamma(U,\sh{J})$ means that $f(u|U)\in\Gamma(U,\sh{F})$, and since $s|U$ is a regular element of $\Gamma(U,\sh{O}_X)$ and $\sh{F}$ is torsion-free, the preceding relation is equivalent to $f((su)|U)\in\Gamma(U,s\sh{F})$;
so if $\overline{v}$ is the section of $\sh{F}/s\sh{F}$, canonical image of $v$, we see that $\sh{J}$ is the kernel of the homomorphism $\sh{O}_X\to\sh{F}/s\sh{F}$ obtained by multiplication by the section $\overline{v}$.
As $\sh{F}/s\sh{F}$ is quasi-coherent, so is $\sh{J}$.

It results immediately from the preceding definition that $\operatorname{dom}(u)$ is the open complement of the closed sub-prescheme of $X$ defined by the Ideal of denominators of~$u$.
\end{env}

\begin{proposition}[20.2.15]
\label{IV.20.2.15}
Let $f:X'\to X$ be a morphism, $\sh{F}$ an $\sh{O}_X$-Module quasi-coherent, $\varphi$ a section of $\sh{M}_f(\sh{F})$ over $X$ \sref{IV.20.1.11}.
Then $f^{-1}(\operatorname{dom}(\varphi))$ is schematically dense in $X'$.
\end{proposition}

\begin{proof}
The question being evidently local on $X$ and $X'$, we can suppose $X=\Spec(A)$, $X'=\Spec(A')$ affines, $\sh{F}=\widetilde{M}$, and $\varphi=h\otimes(1/s)$, where $h\in M$ and $s$ is a regular element of $A$ whose image $s'$ in $A'$ is a regular element.
We know then \sref{IV.20.2.9} that $\mathrm{D}(s')$ is a schematically dense open in $X'$, and it is the inverse image by $f$ of $\mathrm{D}(s)$, which is contained in $\operatorname{dom}(\varphi)$.
\end{proof}

\begin{remark}[20.2.16]
\label{IV.20.2.16}
When $Y$ is not separated, the presheaf $U\mapsto\operatorname{Ps.hom}_S(U,Y)$ on $X$ is not necessarily a sheaf, even when $X$ is Noetherian, as the following example shows.
We take for $X$ the spectrum of a semi-local Noetherian ring $A$ of dimension $\geq 2$, having exactly two
\oldpage[IV-4]{237}
maximal ideals $\mathfrak{m}'$, $\mathfrak{m}''$ (so that $X$ has exactly two closed points $x'$, $x''$), such that $\mathfrak{m}'$ and $\mathfrak{m}''$ belong to $\operatorname{Ass}(A)$ and that the open $U=X-\{x',x''\}$ is \emph{integral}.
Note that the only schematically dense opens in $X$ are those containing $x'$ and $x''$ \sref{IV.20.2.13}[(iv)], and hence they are necessarily equal to $X$.
Let $U'=X-\{x'\}$, $U''=X-\{x''\}$, so that $U=U'\cap U''$.
To have a counterexample, it suffices to define two $S$-morphisms $u':U'\to Y$, $u'':U''\to Y$ (with $S=X$) whose restrictions to $U$ belong to the same pseudo-morphism, but which for every neighbourhood $V'$ of $x'$ in $U'$ and every neighbourhood $V''$ of $x''$ in $U''$, the restrictions of $u'$ and $u''$ to $V'\cap V''$ are distinct.
For this, let us consider a closed irreducible part $Z$ of $X$, distinct from $X$;
let $Y$ be the $X$-prescheme obtained by gluing two copies $Y'$, $Y''$ of $X$ along the open $X-Z$ \sref[I]{I.2.3.1}.
We take for $u'$ and $u''$ the restrictions to $U'$ and $U''$ respectively of the structural isomorphisms $Y'\xrightarrow{\sim}X$, $Y''\xrightarrow{\sim}X$.
Since $V'$ and $V''$ contain the generic point of $Z$, the restrictions of $u'$ and $u''$ to $V'\cap V''$ are distinct; but the restrictions of $u'$ and $u''$ to $U-(U\cap Z)$ are equal, and this is a dense open in $U$, hence schematically dense since $U$ is reduced.

It remains therefore only to define $A$ and $Z$ so as to have the preceding properties.
Let $X_0=\Spec(B)$ be an integral affine scheme (for example the affine plane over a field $k$), $Y$ an irreducible closed subvariety of $X_0$ (for example a line) containing two distinct closed points $x'$, $x''$, corresponding to maximal ideals $\mathfrak{n}'$, $\mathfrak{n}''$ of $B$.
Consider the ring $C=B\oplus(B/\mathfrak{n}'\oplus B/\mathfrak{n}'')$ with the multiplication $(b,z)(b',z')=(bb',bz'+b'z)$.
If $X_1=\Spec(C)$, we have $X_0=(X_1)_\mathrm{red}$ and $X_1$ is reduced except at the points $x'$, $x''$.
It then suffices to take $A=R^{-1}C$, where $R$ is the complement of the union of the maximal ideals of $C$ at the points $x'$, $x''$, and for $Z$ the trace of $Y$ on $X=\Spec(A)$.
\end{remark}

\subsection{Composition of pseudo-morphisms}
\label{IV.20.3}

\begin{env}[20.3.1]
\label{IV.20.3.1}
Let $X$, $Y$, $Z$ be three preschemes, $\omega$ a pseudo-morphism of $X$ in $Y$, $f:Y\to Z$ a morphism.
It is clear that if $U'$, $U''$ are two schematically dense opens in $X$, $u':U'\to Y$, $u'':U''\to Y$ two morphisms of the class $\omega$, the morphisms $fu'$ and $fu''$ are equivalent (for the relation defined in \sref{IV.20.2.1}), so, for all morphisms $u$ of the class $\omega$, the morphisms $fu$ belong to a single equivalence class, which we call the pseudo-morphism of $X$ in $Z$ \emph{composed} of $f$ and $\omega$.
One has $\operatorname{dom}(f\omega)=\operatorname{dom}(\omega)$.
If $g:Z\to T$ is a morphism, it is clear that $g\circ(f\omega)=(g\circ f)\omega$.
\end{env}

\begin{env}[20.3.2]
\label{IV.20.3.2}
Let us now suppose given a pseudo-$S$-morphism $\omega$ of $X$ in $Y$, where $Y$ is \emph{separated over $S$}, so that there exists an $S$-morphism $u:\operatorname{dom}_S(\omega)\to Y$ of the class $\omega$ \sref{IV.20.2.4}.
Let $f:X'\to X$ be an $S$-morphism such that the open $V'=f^{-1}(\operatorname{dom}_S(\omega))$ is schematically dense in $X'$;
then we say that the class (for the equivalence relation of \sref{IV.20.2.1}) of the $S$-morphism $u\circ(f|V')$\footnote{The printed text has $f|U'$ here, but the open introduced in the preceding sentence is $V'=f^{-1}(\operatorname{dom}_S(\omega))$.} (a class defined by virtue of the preceding hypothesis) is the pseudo-$S$-morphism \emph{composed} of $\omega$ and $f$, and we denote it $\omega\circ f$;
it is clear that one has
\[
\label{IV.20.3.2.1}
  f^{-1}(\operatorname{dom}_S(\omega))\subset\operatorname{dom}_S(\omega\circ f).
  \tag{20.3.2.1}
\]

For $f:X'\to X$ given, we denote by $\operatorname{Ps.hom}_S(X,Y)^f$ the set of pseudo-$S$-morphisms $\omega$ of $X$ in $Y$ which verify the preceding condition.
If $\omega$ is such a pseudo-$S$-morphism, it is clear that for every open $V$ of $X$,
\[
  f^{-1}(\operatorname{dom}_S(\omega|V))=f^{-1}(V\cap\operatorname{dom}_S(\omega))=f^{-1}(V)\cap f^{-1}(\operatorname{dom}_S(\omega))
\]
is schematically dense in $f^{-1}(V)$, so, if $f^V:f^{-1}(V)\to V$ is the restriction of $f$, the composed $(\omega|V)\circ f^V$ is defined and equal to $(\omega\circ f)|f^{-1}(V)$.
One thus defines a
\oldpage[IV-4]{238}
canonical restriction map $\operatorname{Ps.hom}_S(X,Y)^f\to\operatorname{Ps.hom}_S(V,Y)^{f^V}$, which
allows to define a \emph{presheaf} of sets $V\mapsto\operatorname{Ps.hom}_S(V,Y)^{f^V}$ on $X$, sub-presheaf of $V\mapsto\operatorname{Ps.hom}_S(V,Y)$, which gives a sheaf of sets that we denote $\mathscr{P}\!s.\!\operatorname{hom}_S(X,Y)^f$.
Moreover, for every open $V$ of $X$, we have a map $\omega\mapsto\omega\circ f^V$ from $\operatorname{Ps.hom}_S(V,Y)^{f^V}$ to $\operatorname{Ps.hom}_S(f^{-1}(V),Y)$, which defines an $f$-morphism of sheaves of sets
\[
  \mathscr{P}\!s.\!\operatorname{hom}_S(X,Y)^f\to\mathscr{P}\!s.\!\operatorname{hom}_S(X',Y).
\]
\end{env}

\begin{env}[20.3.3]
\label{IV.20.3.3}
Let now $f':X''\to X'$ be an $S$-morphism such that the open $f'^{-1}(f^{-1}(\operatorname{dom}_S(\omega)))$ is schematically dense in $X''$;
then $\omega\circ(f\circ f')$ is defined and $u\circ f\circ f'$ belongs to this pseudo-$S$-morphism;
moreover, by virtue of \sref{IV.20.3.2.1}, $f'^{-1}(\operatorname{dom}_S(\omega\circ f))$ is \emph{a fortiori} schematically dense in $X''$, so $(\omega\circ f)\circ f'$ is also defined and $u\circ f\circ f'$ belongs to this pseudo-$S$-morphism, hence $(\omega\circ f)\circ f'=\omega\circ(f\circ f')$.
On the other hand, for every $S$-morphism $g:Y\to Z$, one has $\operatorname{dom}_S(g\omega)=\operatorname{dom}_S(\omega)$ \sref{IV.20.3.1}, so $(g\omega)\circ f$ is defined, and $g\circ u\circ f$ belongs to this pseudo-$S$-morphism, which shows that $(g\omega)\circ f=g\circ(\omega\circ f)$.
\end{env}

\begin{env}[20.3.4]
\label{IV.20.3.4}
The most important case in the definition \sref{IV.20.3.2} is that where one has $\mathscr{P}\!s.\!\operatorname{hom}_S(X,Y)^f=\mathscr{P}\!s.\!\operatorname{hom}_S(X,Y)$:
it suffices for this that for every open $U$ of $X$ and every open $V$ schematically dense in $U$, $f^{-1}(V)$ is schematically dense in $f^{-1}(U)$;
when this holds, then, for \emph{every} open $U$ of $X$ and \emph{every} pseudo-$S$-morphism $\omega:U\to Y$, one can define the composed $\omega\circ f^U$, \emph{even if $Y$ is not separated over $S$}.
Indeed, if $u':U'\to Y$ and $u'':U''\to Y$ are two morphisms of the class $\omega$, they coincide in an open $U_0$ schematically dense in $U$, so the composed morphisms $f^{-1}(U')\to U'\xrightarrow{u'}Y$ and $f^{-1}(U'')\to U''\xrightarrow{u''}Y$ coincide in $f^{-1}(U_0)$, which is by hypothesis schematically dense in $f^{-1}(U)$;
one can therefore define $\omega\circ f^U$ as the class of any of the morphisms $f^{-1}(U')\to U'\xrightarrow{u'}Y$, where $u'$ belongs to $\omega$.
\end{env}

\begin{proposition}[20.3.5]
\label{IV.20.3.5}
Let $X$, $X'$ be two preschemes, $f:X'\to X$ a morphism.
Then, in each of the three following cases, for every open $U$ of $X$ and every open $V$ schematically dense in $U$, $f^{-1}(V)$ is schematically dense in $f^{-1}(U)$:
\begin{enumerate}[label=\emph{(\roman*)}]
  \item $f$ is flat and locally of finite presentation.
  \item $f$ is flat and the underlying space of $X$ is locally Noetherian.
  \item $X'$ is reduced, the set of irreducible components of $X$ is locally finite, and every irreducible component of $X'$ dominates an irreducible component of $X$.
\end{enumerate}
\end{proposition}

\begin{proof}
Assertion~(i) results from \sref{IV.11.10.5}[(ii), b)];
assertion~(ii) results from \sref{IV.11.10.5}[(ii), a)], since then every open $V$ in $U$ is retrocompact, i.e.\ the canonical injection $j:V\to U$ is a quasi-compact morphism.
Finally, to prove~(iii), note that since $X'$ is reduced, it suffices to show that for every open $U$ of $X$ and every open $V$ dense in $U$, $f^{-1}(V)$ is \emph{dense} in $f^{-1}(U)$.
Now, for $f^{-1}(V)$ to be dense in $f^{-1}(U)$, it suffices that $f^{-1}(V)$ contain all the maximal points of $X'$ belonging to $f^{-1}(U)$;
  the conclusion therefore results from the hypothesis that the image under $f$ of every
\oldpage[IV-4]{239}
maximal point of $X'$ belonging to $f^{-1}(U)$ is a maximal point of $X$ belonging to $U$, hence belongs to $V$ since the set of irreducible components of $X$ is locally finite.
\end{proof}

\begin{env}[20.3.6]
\label{IV.20.3.6}
Let $X$, $Y$ be two $S$-preschemes;
suppose that $X$ satisfies one of the two following hypotheses:
\begin{enumerate}[label=\emph{\alph*)}]
  \item $X$ is locally Noetherian;
  \item $X$ is reduced and the set of its irreducible components is locally finite.
\end{enumerate}
Then, for every $x\in X$, the canonical $S$-morphism $j:\Spec(\sh{O}_{X,x})\to X$ is flat and satisfies condition~(ii) of \sref{IV.20.3.5} in case~a), condition~(iii) of \sref{IV.20.3.5} in case~b).
For every pseudo-$S$-morphism $\omega$ of $X$ in $Y$, the composite $\omega\circ j$ is therefore defined and is a pseudo-$S$-morphism of $\Spec(\sh{O}_{X,x})$ in $Y$, called the \emph{restriction} of $\omega$ to $\Spec(\sh{O}_{X,x})$.
Note now that if $X$ satisfies condition~a) (resp.~b)) of \sref{IV.20.3.6}, the same holds for every prescheme induced on an open $U$ of $X$ containing $x$.
By passage to the inductive limit, we deduce from the canonical maps $\operatorname{Ps.hom}_S(U,Y)\to\operatorname{Ps.hom}_S(\Spec(\sh{O}_{X,x}),Y)$ thus obtained, a canonical map
\[
\label{IV.20.3.6.1}
  (\mathscr{P}\!s.\!\operatorname{hom}_S(X,Y))_x\to\operatorname{Ps.hom}_S(\Spec(\sh{O}_{X,x}),Y)
  \tag{20.3.6.1}
\]
where the left-hand side is the stalk at the point $x$ of the sheaf $\mathscr{P}\!s.\!\operatorname{hom}_S(X,Y)$, the set of germs at $x$ of pseudo-$S$-morphisms from open neighbourhoods of $x$ into $Y$.
\end{env}

\begin{proposition}[20.3.7]
\label{IV.20.3.7}
Under the hypotheses of \sref{IV.20.3.6}, the canonical map \sref{IV.20.3.6.1} is injective.
If moreover $Y$ is locally of finite presentation over $S$, the map \sref{IV.20.3.6.1} is bijective.
\end{proposition}

\begin{proof}
By application of the method of \sref{IV.8.1.2}[(a)], this proposition will be a particular case of the following:
\end{proof}

\begin{proposition}[20.3.8]
\label{IV.20.3.8}
The notations being those of \sref{IV.8.8.1}, suppose $X_\alpha$ quasi-compact (resp.\ quasi-compact and quasi-separated) and $Y_\alpha$ locally of finite type (resp.\ locally of finite presentation) over $S_\alpha$.
Suppose moreover that one of the following conditions is satisfied:
\begin{enumerate}[label=\emph{(\roman*)}]
  \item The transition morphisms $S_\mu\to S_\lambda$ (for $\lambda\leq\mu$) are flat, and the $X_\lambda$ and $X$ are Noetherian.
  \item The $X_\lambda$ are reduced, the sets of irreducible components of $X$ and of each $X_\lambda$ are finite, and, for $\lambda\leq\mu$, every irreducible component of $X_\mu$ dominates an irreducible component of $X_\lambda$.
\end{enumerate}
Then, the canonical map
\[
\label{IV.20.3.8.1}
  \varinjlim\operatorname{Ps.hom}_{S_\lambda}(X_\lambda,Y_\lambda)\to\operatorname{Ps.hom}_S(X,Y)
  \tag{20.3.8.1}
\]
is injective (resp.\ bijective).
\end{proposition}

\begin{proof}
Note first that, in case~(i), the morphisms $X_\mu\to X_\lambda$ (for $\lambda\leq\mu$) and $X\to X_\lambda$ are flat, so by \sref{IV.20.3.4} and \sref{IV.20.3.5} the canonical maps
\[
  \operatorname{Ps.hom}_{S_\lambda}(X_\lambda,Y_\lambda)\to\operatorname{Ps.hom}_{S_\mu}(X_\mu,Y_\mu)
\]
\oldpage[IV-4]{240}
for $\lambda\leq\mu$, as well as the maps
\[
  \operatorname{Ps.hom}_{S_\lambda}(X_\lambda,Y_\lambda)\to\operatorname{Ps.hom}_S(X,Y),
\]
are defined (without any separation hypothesis on the $Y_\lambda$ or on $Y$);
the same is consequently true of the map \sref{IV.20.3.8.1}.
The proposition will result from the following lemma.

\begin{lemma}[20.3.8.2]
\label{IV.20.3.8.2}
The notations being those of \sref{IV.8.8.1}, suppose $X_\alpha$ quasi-compact, and let $U_\alpha$ be an open in $X_\alpha$;
let $U_\lambda$ and $U$ be its inverse images in $X_\lambda$ and $X$.
Suppose that one of the conditions~(i),~(ii) of \sref{IV.20.3.8} is satisfied.
Then, for $U$ to be schematically dense in $X$, it is necessary and sufficient that there exist $\lambda\geq\alpha$ such that $U_\lambda$ is schematically dense in $X_\lambda$, and in this case $U_\mu$ is schematically dense in $X_\mu$ for $\mu\geq\lambda$.
\end{lemma}

Suppose first that condition~(i) is satisfied.
Denote by $j_\lambda:U_\lambda\to X_\lambda$ and $j:U\to X$ the canonical injections, and by $\sh{J}_\lambda$ the kernel of the canonical homomorphism $\sh{O}_{X_\lambda}\to(j_\lambda)_*(\sh{O}_{U_\lambda})$.
Since the immersion $j_\lambda$ is quasi-compact and quasi-separated, $(j_\lambda)_*(\sh{O}_{U_\lambda})$ is a quasi-coherent $\sh{O}_{X_\lambda}$-module; hence $\sh{J}_\lambda$ is a quasi-coherent ideal of $\sh{O}_{X_\lambda}$, and, since $X_\lambda$ is Noetherian, $\sh{J}_\lambda$ is coherent and therefore of finite type.
On the other hand, since the transition morphism $p_{\lambda\mu}:X_\mu\to X_\lambda$ (resp.\ $p_\lambda:X\to X_\lambda$) is flat, it follows from \sref{IV.2.3.1} that one may identify $\sh{J}_\mu$ with $p_{\lambda\mu}^*(\sh{J}_\lambda)$ (resp.\ $\sh{J}$ with $p_\lambda^*(\sh{J}_\lambda)$).
The assertion results then from the definition of a schematically dense open, and of \sref{IV.8.5.8}[(ii)].

To establish \sref{IV.20.3.8.2} when condition~(ii) is satisfied, we first prove two lemmas.

\begin{lemma}[20.3.8.3]
\label{IV.20.3.8.3}
Under the hypotheses of \sref{IV.8.2.2}, let $M_\lambda$ (resp.\ $M$) be the set of maximal points of $S_\lambda$ (resp.\ $S$).
Suppose that for every $\lambda$, the set $M_\lambda$ is finite, and that the $M_\lambda$ form a projective system of sets.
Then $M$ is the projective limit of the $M_\lambda$.
\end{lemma}

Let us first show that a point $s\in\varprojlim M_\lambda$ is maximal in $S$: indeed, if $s'$ is a generalisation of $s$, the image $s'_\lambda$ of $s'$ in $S_\lambda$ is a generalisation of the image $s_\lambda$ of $s$, hence equal to $s_\lambda$ for all $\lambda$, which implies $s'=s$, since the set of points underlying $S$ is the projective limit of the sets underlying the $S_\lambda$ \sref{IV.8.2.9}.
Inversely, let $s$ be a maximal point of $S$ and let us prove that it belongs to $\varprojlim M_\lambda$.
Let $s_\lambda$ be the image of $s$ in $S_\lambda$, $M'_\lambda$ the set of maximal points of $S_\lambda$ that are generalisations of $s_\lambda$; the $M'_\lambda$ are non-empty finite sets forming a projective system, hence $M'=\varprojlim M'_\lambda$ is non-empty and the map $M'\to M'_\lambda$ is surjective (Bourbaki, \emph{Ens.}, chap.~III, 2\textsuperscript{e}~\'ed., \textsection~7, n\textsuperscript{o}~4, Example~I).
On the other hand, $\Spec(\sh{O}_{S,s})=\varprojlim\Spec(\sh{O}_{S_\lambda,s_\lambda})$ by virtue of \sref{IV.8.2.12} and \sref{IV.8.2.9}, hence the points of $\varprojlim M'_\lambda$ are also maximal points of $\Spec(\sh{O}_{S,s})$ by the first part of the reasoning.
Thus $M'=\varprojlim M'_\lambda$ necessarily reduces to the point $s$; one concludes that $M'_\lambda$ reduces to the point $s_\lambda$ and consequently the $s_\lambda$ are maximal in the $S_\lambda$, which completes the proof of the lemma.

\begin{lemma}[20.3.8.4]
\label{IV.20.3.8.4}
The hypotheses being those of \sref{IV.20.3.8.3}, suppose moreover that $S_\alpha$ is quasi-compact, and let $U_\alpha$ be an open of $S_\alpha$, with inverse images $U_\lambda$ in $S_\lambda$ for $\lambda\geq\alpha$ and $U$ in $S$.\footnote{The printed statement introduces $V_\alpha,V_\lambda,V$ but immediately and consistently uses $U_\alpha,U_\lambda,U$; the latter notation is used here.}
If $U_\alpha$ is dense in $S_\alpha$, then $U_\lambda$ is dense in $S_\lambda$ for $\lambda\geq\alpha$ and $U$ is dense in $S$.
Inversely, if $U$ is dense in $S$ and if moreover the set $M$ of maximal points of $S$ is finite, there exists $\lambda\geq\alpha$ such that $U_\lambda$ is dense in $S_\lambda$.
\end{lemma}

Indeed, since $M_\alpha$ is finite, the hypothesis that $U_\alpha$ is dense in $S_\alpha$ implies that
\oldpage[IV-4]{241}
$M_\alpha\subset U_\alpha$, hence $M_\lambda\subset U_\lambda$ for $\lambda\geq\alpha$ and $M\subset U$ by virtue of \sref{IV.20.3.8.3}, which proves the first assertion.
Inversely, suppose $M$ finite and $U$ dense in $S$, hence $M\subset U$; note that the $U_\lambda$ are open, hence ind-constructible, and the $M_\lambda$ finite, hence pro-constructible \sref{IV.1.9.6}.
The second assertion then follows from \sref{IV.8.3.2}.

\phantomsection\label{IV.20.3.8.5}
\emph{20.3.8.5. End of the proof of \sref{IV.20.3.8.2}.}
Suppose condition~(ii) is verified, and note that $X$ is then \emph{reduced} by virtue of \sref{IV.8.7.1};
it amounts to the same to say that $U_\lambda$ (resp.\ $U$) is schematically dense in $X_\lambda$ (resp.\ $X$) or that it is dense, and the conclusion results from \sref{IV.20.3.8.4} applied to $X_\lambda$ and~$X$.

\phantomsection\label{IV.20.3.8.6}
\emph{20.3.8.6. End of the proof of \sref{IV.20.3.8}.}
To show that the map \sref{IV.20.3.8.1} is injective, consider two morphisms $u'_\lambda$, $u''_\lambda$ of a schematically dense open $U_\lambda\subset X_\lambda$ in $Y_\lambda$, and suppose that their images $u'$, $u''$, morphisms of $U$ in $Y$, coincide in a schematically dense open $V$ of $U$.
Under either hypothesis~(i) or~(ii), one can suppose $V$ \emph{quasi-compact};
otherwise, since $X$ has only a finite number of maximal points and is reduced, it suffices to replace $V$ by a union of a finite number of affine open neighbourhoods of these maximal points (contained in $V$ by hypothesis).
Then there exists $\lambda\geq\alpha$ such that $V$ is the inverse image of a quasi-compact open $V_\lambda$ of $X_\lambda$ \sref{IV.8.2.11}, and by \sref{IV.20.3.8.2} we can suppose that $V_\lambda$ is schematically dense in $X_\lambda$.
It results from \sref{IV.8.8.2}[(i)] that, on taking $\lambda$ sufficiently large, $u'_\lambda$ and $u''_\lambda$ coincide on $V_\lambda$, hence belong to the same pseudo-morphism.

Let us prove finally that the map \sref{IV.20.3.8.1} is surjective.
Consider now a morphism $u$ of a schematically dense open $U$ of $X$ in $Y$;
as above, we can suppose $U$ quasi-compact and quasi-separated ($U$ being replaceable, in case~(ii), by a union of a finite number of disjoint affine opens).
We can furthermore suppose there exists $\lambda\geq\alpha$ such that $U$ is the inverse image of a quasi-compact open $U_\lambda$ of $X_\lambda$ \sref{IV.8.2.11} which is automatically quasi-separated, and by \sref{IV.20.3.8.2} we can suppose that $U_\lambda$ is schematically dense in $X_\lambda$.
Since the $Y_\lambda$ are supposed locally of finite presentation, it results from \sref{IV.8.8.2}[(i)] that there exists $\lambda$ such that $u$ is the image of a morphism $u_\lambda:U_\lambda\to Y_\lambda$; whence the conclusion.

\begin{remarks}[20.3.8.7]
\label{IV.20.3.8.7}
\begin{enumerate}[label=(\roman*)]
  \item To show that the map \sref{IV.20.3.8.1} is \emph{injective}, it is not necessary, in the hypothesis~(i) of \sref{IV.20.3.8}, to suppose $X$ Noetherian.
  Indeed, the lemma \sref{IV.20.3.8.2} does not use this hypothesis.
  With the notations of \sref{IV.20.3.8.6}, let $Z_\lambda$ be the sub-prescheme of coincidences of $u'_\lambda$ and $u''_\lambda$, and let $Z$ be the sub-prescheme of coincidences of $u'$ and $u''$;
  it results from the definition \sref{IV.17.4.5} and from \sref[I]{I.3.3.10.1} that $Z$ is the projective limit of the $Z_\mu$ for $\mu\geq\lambda$.
  Now, by hypothesis, $Z$ majorizes a schematically dense open of $X$;
  it follows that $Z$ is itself induced on an open of $X$ by virtue of the following lemma:

\begin{lemma}[20.3.8.8]
\label{IV.20.3.8.8}
Let $T$ be a prescheme.
Then every sub-prescheme $W$ of $T$ which majorizes a schematically dense open $V$ of $T$ is induced on an open (schematically dense) of $T$.
\end{lemma}

  Indeed, the subspace of $T$ underlying $W$ is locally closed, hence open
  \oldpage[IV-4]{242}
  in its closure, which already proves that the space underlying $W$ is open in $T$;
  the conclusion results from \sref{IV.11.10.1}[(c)].

  This lemma being established, we conclude that for $\mu\geq\lambda$ sufficiently large, $Z_\mu$ is induced on an open of $X_\mu$ by virtue of \sref{IV.8.6.3}; indeed, $Z_\lambda$, being a sub-prescheme of a Noetherian prescheme, is of finite presentation over $X_\lambda$ \sref{IV.1.6.2}, and the same is therefore true of $Z_\mu$ over $X_\mu$ for $\mu\geq\lambda$ and of $Z$ over $X$.
  One can now apply \sref{IV.20.3.8.2}, which shows that for $\mu\geq\lambda$ sufficiently large, $Z_\mu$ is schematically dense in $X_\mu$, whence the conclusion.

  \item If, in the hypothesis~(i) of \sref{IV.20.3.8}, we drop the condition that $X$ is Noetherian, the reasoning of \sref{IV.20.3.8.6} shows yet that the image of \sref{IV.20.3.8.1} is formed of the pseudo-$S$-morphisms having a representative which is an $S$-morphism $U\to Y$, where $U$ is schematically dense in $X$ and \emph{quasi-compact and quasi-separated}.
\end{enumerate}
\end{remarks}
\end{proof}

\begin{corollary}[20.3.9]
\label{IV.20.3.9}
Suppose that either hypothesis~a) or~b) of \sref{IV.20.3.6} is satisfied by $X$, and let $Y$ be locally of finite presentation over $S$.
Then, for a pseudo-$S$-morphism $\omega$ of $X$ in $Y$ to be defined at a point $x$ \sref{IV.20.2.3}, it is necessary and sufficient that its restriction to $\Spec(\sh{O}_{X,x})$ be everywhere defined (in other words, that it be an $S$-morphism of $\Spec(\sh{O}_{X,x})$ in $Y$).
\end{corollary}

The following result, which will be used in the proof of \sref{IV.20.3.11}, uses the theory of faithfully flat descent from Chapter~VI.
The reader may verify that the results of \sref{IV.20.3} will not be used in that theory.

\begin{lemma}[20.3.10]
\label{IV.20.3.10}
Let $f:X'\to X$ be an $S$-morphism faithfully flat and quasi-compact, $X''=X'\times_X X'$, $p_1$ and $p_2$ the canonical projections of $X''$ to $X'$, and let $Y$ be an $S$-prescheme separated over $S$.
Let $U$ be an open of $X$, $U'=f^{-1}(U)$, $U''=p_1^{-1}(U')=p_2^{-1}(U')$, and suppose that $U''$ is schematically dense in $X''$.
Let $g:U\to Y$ be an $S$-morphism;
then, if $g\circ(f|U')$ extends to an $S$-morphism $X'\to Y$, $g$ extends to an $S$-morphism $X\to Y$.
\end{lemma}

The hypotheses imply that $U$ (resp.\ $U'$) is schematically dense in $X$ (resp.\ $X'$) \sref{IV.11.10.5}[(i)].
Thus, if $\omega$ denotes the pseudo-$S$-morphism given by the class of $g$, the assertion of \sref{IV.20.3.10} can equivalently be stated as follows: if $\omega\circ f$ is everywhere defined, then so is $\omega$.

\begin{proof}
To prove \sref{IV.20.3.10}, denote by $g'$ an $S$-morphism $X'\to Y$ which extends $g\circ(f|U')$, and set $g''_i=g'\circ p_i:X''\to Y$ ($i=1,2$).
If one sets $f''=f\circ p_1=f\circ p_2:X''\to X$, it is clear that $g''_1$ and $g''_2$ coincide on $U''$ with $g\circ(f''|U'')$.
But as $Y$ is separated over $S$ and $U''$ is schematically dense in $X''$, we have $g''_1=g''_2$ \sref{IV.11.10.1}[(d)].
As $f$ is faithfully flat and quasi-compact, it results from descent theory (chap.~VI) that there exists a unique $S$-morphism $h:X\to Y$ such that $h\circ f=g'$;
as the restriction $U'\to U$ of $f$ is a faithfully flat and quasi-compact morphism and $U''=U'\times_U U'$, the preceding uniqueness result, applied to $U$ instead of $X$, shows that $h|U=g$, which proves the lemma.
\end{proof}

\begin{proposition}[20.3.11]
\label{IV.20.3.11}
Let $Y$ be an $S$-prescheme separated over $S$, $\omega$ a pseudo-$S$-morphism of $X$ in $Y$, $f:X'\to X$ an $S$-morphism.
Suppose that $f$ is flat, and that one of the following conditions is satisfied:
\begin{enumerate}[label=\emph{(\roman*)}]
\oldpage[IV-4]{243}
  \item $f$ is an open morphism, or surjective and quasi-compact, and $\operatorname{dom}_S(\omega)$ contains an open $U$ schematically dense in $X$ and retrocompact in~$X$.
  \item $f$ is locally of finite presentation.
  \item $Y$ is locally of finite presentation over $S$, and $X$ satisfies one of the conditions~a),~b) of \sref{IV.20.3.6}.
\end{enumerate}
Then, $f^{-1}(\operatorname{dom}_S(\omega))$ is schematically dense in $X'$, so that $\omega\circ f$ is defined \sref{IV.20.3.2} and one has
\[
\label{IV.20.3.11.1}
  \operatorname{dom}_S(\omega\circ f)=f^{-1}(\operatorname{dom}_S(\omega)).
  \tag{20.3.11.1}
\]
\end{proposition}

\begin{proof}
Let us first prove that $f^{-1}(\operatorname{dom}_S(\omega))$ is schematically dense in $X'$.
The question being local on $X'$, we can suppose $X$ and $X'$ affines and since $f$ is flat, it suffices to verify, by \sref{IV.11.10.5}[(ii), a)], that $\operatorname{dom}_S(\omega)$ contains a retrocompact open $U$ schematically dense in $X$.
This is the hypothesis in case~(i), and follows from \sref{IV.20.2.12} in case~(iii), since in an affine scheme every open of the form $\mathrm{D}(t)$ is retrocompact;
in case~(ii), we can directly apply \sref{IV.20.3.5}[(i)].

Let us now prove \sref{IV.20.3.11.1}, i.e., that, for every point $x'\in\operatorname{dom}_S(\omega\circ f)$, one has $x=f(x')\in\operatorname{dom}_S(\omega)$.
Note first that we can restrict to the case where $f$ is \emph{faithfully flat and quasi-compact}.
This is already the hypothesis in the second case of~(i).
In the other cases the question is local on $X'$, so we may suppose $X$ and $X'$ affine, and hence $f$ quasi-compact.
In the first case of~(i), and in case~(ii), $f$ is an open morphism \sref{IV.11.3.1}, so, replacing $X$ by the open $f(X')$, we may suppose $f$ surjective, hence faithfully flat.
In case~(iii), by \sref{IV.20.3.9}, it suffices to prove that the restriction of $\omega$ to $\Spec(\sh{O}_{X,x})$ is everywhere defined, and we can therefore replace $X$ by $\Spec(\sh{O}_{X,x})$, $X'$ by $X'\times_X\Spec(\sh{O}_{X,x})$ and $f$ by its restriction to this last prescheme, which is a surjective morphism \sref{IV.2.3.4}, hence faithfully flat.

Suppose therefore $f$ faithfully flat and quasi-compact;
with the notations of the lemma \sref{IV.20.3.10}, it suffices to see that $\mathrm{U}''$ is schematically dense in $X''$, by taking for $U$ an open schematically dense in $X$, contained in $\operatorname{dom}_S(\omega)$;
this will be the case, by virtue of \sref{IV.11.10.5}[(ii), a)], if $U$ is taken retrocompact in $X$ (since the morphism $f'':X''\to X$ is flat).
The existence of such an open $U$ is part of the hypothesis in case~(i);
in case~(iii) it results from \sref{IV.20.2.12} and the fact that in an affine scheme $\Spec(A)$, every open of the form $\mathrm{D}(t)$ is retrocompact.
Finally, in case~(ii), we take $U=\operatorname{dom}_S(\omega)$ and show directly that $\mathrm{U}''$ is schematically dense in $X''$:
it suffices for this to remark that $f'':X''\to X$ is flat and locally of finite presentation and to apply \sref{IV.20.3.5}[(i)].
\end{proof}

\begin{corollary}[20.3.12]
\label{IV.20.3.12}
Let $\varphi$ be a pseudo-function on a prescheme $X$.
Then, for every flat morphism locally of finite presentation $f:X'\to X$, the pseudo-function $\varphi\circ f$ is defined and one has $\operatorname{dom}(\varphi\circ f)=f^{-1}(\operatorname{dom}(\varphi))$.
\end{corollary}

\begin{remark}[20.3.13]
\label{IV.20.3.13}
\oldpage[IV-4]{244}
When $X$ satisfies one of the conditions~a),~b) of \sref{IV.20.3.6}, then we have seen \sref{IV.20.2.11} that the pseudo-functions on $X$ are identified with the meromorphic functions on $X$.
By virtue of \sref{IV.20.1.12} and of \sref{IV.20.2.13}[(i)], if one only supposes that the morphism $f:X'\to X$ is \emph{flat}, then, for every pseudo-function $\varphi$ on $X$, $\varphi\circ f$ is defined and one has $\operatorname{dom}(\varphi\circ f)=f^{-1}(\operatorname{dom}(\varphi))$.
\end{remark}

\subsection{Properties of domains of definition of rational functions}
\label{IV.20.4}

\begin{env}[20.4.1]
\label{IV.20.4.1}
Let $X$, $Y$ be two $S$-preschemes, $\omega$ a pseudo-$S$-morphism of $X$ in $Y$.
Let $u$ be an $S$-morphism $U\to Y$ belonging to the class $\omega$, where $U$ is schematically dense in $X$, and consider the \emph{graph} $\Gamma_u$, sub-prescheme of $U\times_S Y$ \sref[I]{I.5.3.11}, and therefore sub-prescheme of $X\times_S Y$.
Suppose that this sub-prescheme admits an \emph{adherence} \sref[I]{I.9.5.11} in $X\times_S Y$;
if $j:\Gamma_u\to X\times_S Y$ is the canonical injection, this will be the case when the $\sh{O}_{X\times_S Y}$-Module $j_*(\sh{O}_{\Gamma_u})$ is quasi-coherent;
it results from the definition of the class of equivalence \sref{IV.20.2.1} that $j_*(\sh{O}_{\Gamma_u})$ does not depend on the representative $u$ considered, hence the adherence of this sub-prescheme is well determined by $\omega$;
we say that this closed sub-prescheme is the \emph{graph} of the pseudo-$S$-morphism $\omega$ and we denote it $\Gamma_\omega$.
We note that $\Gamma_\omega$ is defined if there exists a morphism $u:U\to Y$ in the class $\omega$ such that $U$ is retrocompact in $X$ and if, moreover, $Y$ is quasi-separated over $S$;
indeed, the injection $j$ is then a quasi-compact and quasi-separated morphism \sref[I]{I.2.2} and \sref[I]{I.7.4}.
The first hypothesis is always satisfied when $X$ is locally Noetherian.
We note also that when $X$ is \emph{reduced}, so is $\Gamma_u$, which is isomorphic to $U$ \sref[I]{I.5.3.11};
then $\Gamma_\omega$ exists and is none other than the \emph{reduced} closed sub-prescheme of $X\times_S Y$ having for underlying space the closure in $X\times_S Y$ of the underlying space of $\Gamma_u$ \sref[I]{I.5.2.1} and \sref[I]{I.5.2.2}.
Note finally that if $Y$ is \emph{separated} over $S$, then $\Gamma_u$ is \emph{closed} in $U\times_S Y$ \sref[I]{I.5.4.3}, and hence is induced by $\Gamma_\omega$ (whenever the latter exists) on the open $\Gamma_\omega\cap(U\times_S Y)$ \sref[I]{I.9.5.10};
by contrast, the prescheme induced by $\Gamma_\omega$ is in general different from $\Gamma_u$ when $Y$ is not separated.
In particular, if $v:X\to Y$ is an $S$-\emph{morphism}, the graph $\Gamma_\omega$ of the class $\omega$ of $v$ may be distinct from the graph $\Gamma_v$.
In particular also in what follows, when the question is about the graph of a pseudo-$S$-morphism, we will restrict to the case where \emph{$Y$ is separated over $S$}.
\end{env}

\begin{env}[20.4.2]
\label{IV.20.4.2}
Suppose therefore $\Gamma_\omega$ defined and $Y$ separated over $S$;
denote by $p$ and $q$ the restrictions to $\Gamma_\omega$ of the canonical projections.
\[
  \xymatrix{
  & \Gamma_\omega \ar[dl]_p \ar[dr]^q & \\
  X & & Y
  }
\]
Then, if $U\subset\operatorname{dom}_S(\omega)$, the restriction $p^{-1}(U)\to U$ of $p$ is an isomorphism \sref[I]{I.5.3.11};
reciprocally, if $U$ is an open of $X$ having this property, and if $u$ is the isomorphism reciprocal of the restriction $p^{-1}(U)\to U$ of $p$, $q\circ u$ is an $S$-morphism of $U$ in $Y$ which coincides with a morphism of the class $\omega$ in $U\cap\operatorname{dom}_S(\omega)$.
We conclude that $\operatorname{dom}_S(\omega)$ is \emph{the largest open} $U$ of $X$ such that the restriction $p^{-1}(U)\to U$ of $p$ is an
\oldpage[IV-4]{245}
isomorphism;
let $v:\operatorname{dom}_S(\omega)\to\Gamma_\omega$ be the $S$-morphism inverse to the restriction $p^{-1}(\operatorname{dom}_S(\omega))\to\operatorname{dom}_S(\omega)$ of $p$.
The pseudo-$S$-morphism from $X$ to $\Gamma_\omega$ represented by $v$ is sometimes denoted $p^{-1}$; its associated rational map \sref{IV.20.2.13}[(ii)] is then \emph{birational}.
Since $p^{-1}(\operatorname{dom}_S(\omega))$ is the graph of the $S$-morphism $u=q\circ v:\operatorname{dom}_S(\omega)\to Y$, it is schematically dense in $\Gamma_\omega$ \sref{IV.11.10.3}[(iv)], so $\omega$ can be considered as the \emph{composite} $q\circ p^{-1}$ \sref{IV.20.3.2}.
For every subset $M$ of the underlying space of $X$, one sometimes sets $\omega(M)=q(p^{-1}(M))$, which amounts to considering $\omega$ as a map from $X$ to the set of subsets of $Y$ (or, as some authors say, a ``multiform function'').
We note that when $x\in\operatorname{dom}_S(\omega)$, $\omega(\{x\})$ is the singleton $\{u(x)\}$;
in general, for any $x\in X$, if one denotes by $s$ the image of $x$ in $S$, via the structural morphism, $\omega(\{x\})$ is a part of the prescheme $Y_s$, the fibre at the point $s$ of the structural morphism.
\end{env}

\begin{env}[20.4.3]
\label{IV.20.4.3}
In all the rest of this n\textsuperscript{o}, we will restrict to the case where $X$ is \emph{reduced}, so that there is then an identity between pseudo-$S$-morphisms and $S$-\emph{rational applications} \sref{IV.20.2.7}.
Moreover, the graph of every $S$-rational application of $X$ in $Y$ is then defined \sref{IV.20.4.1}.
\end{env}

\begin{proposition}[20.4.4]
\label{IV.20.4.4}
Let $X$ be an $S$-prescheme locally integral, $Y$ an $S$-prescheme locally of finite type and separated over $S$, $\omega$ an $S$-rational application of $X$ in $Y$, $p:\Gamma_\omega\to X$ the canonical projection.
Then, if $x\in X$ is a normal point such that $p^{-1}(x)$ contains an isolated point, $\omega$ is defined at the point $x$.
\end{proposition}

\begin{proof}
Indeed, the first projection $p_1:X\times_S Y\to X$ is then a separated and locally of finite type morphism, and so is its restriction $p:\Gamma_\omega\to X$, which is moreover \emph{birational}, and as $\Gamma_\omega$ is reduced and $X$ integral, $\Gamma_\omega$ is \emph{integral};
it results then from (\textbf{Err}\textsubscript{\textbf{IV}}, 30) that the hypothesis on $x$ implies the existence of an open neighbourhood $U$ of $x$ such that the restriction $p^{-1}(U)\to U$ of $p$ is an isomorphism; whence the conclusion by \sref{IV.20.4.2}.

We note that \sref{IV.20.4.4} is the original formulation given by Zariski of his ``Main theorem'' (up to the fact that it was originally restricted to algebraic schemes over a base field).
\end{proof}

\begin{proposition}[20.4.5]
\label{IV.20.4.5}
The hypotheses and notations being those of \sref{IV.20.4.4}, suppose moreover $X$ normal, and let $X'$ be a reduced prescheme, $f:X'\to X$ a morphism locally of finite type and universally open.
Then $\omega\circ f$ is defined, and one has $\operatorname{dom}_S(\omega\circ f)=f^{-1}(\operatorname{dom}_S(\omega))$ (in other words, if $x'\in X'$ and $x=f(x')$, then for $\omega$ to be defined at $x$ it is necessary and sufficient that $\omega\circ f$ be defined at $x'$).
\end{proposition}

\begin{proof}
As $X'$ is reduced and $f^{-1}(\operatorname{dom}_S(\omega))$ is everywhere dense in $X'$ by virtue of \sref{IV.2.4.11}, the composite $\omega\circ f$ is defined.
To prove that, whenever $\omega\circ f$ is defined at $x'$, $\omega$ is defined at $x$, we may plainly replace $X'$ by an open neighbourhood of $x'$ and hence suppose $\omega\circ f$ everywhere defined;
moreover, since $f(X')$ is open in $X$, we may replace $X$ by $f(X')$ and thus suppose $f$ surjective.

By the hypothesis on $f$, the morphism
\[
  f_{(Y)}:X'\times_S Y\longrightarrow X\times_S Y
\]
is then open, and consequently
\[
  \overline{f_{(Y)}^{-1}(M)}=f_{(Y)}^{-1}(\overline{M})
\]
for every subset $M$ of $X\times_S Y$.
Taking for $M$ the graph $\Gamma_u$, where $u:\operatorname{dom}_S(\omega)\to Y$ is the restriction of $\omega$ to its domain, it follows from the preceding relation and from \sref[I]{I.3.3.12} that the underlying space of $\Gamma_{\omega\circ f}$ is $f_{(Y)}^{-1}(\Gamma_\omega)$.
Since
\oldpage[IV-4]{246}
$\Gamma_{\omega\circ f}$ is a reduced prescheme \sref{IV.20.4.1}, the $X'$-prescheme $\Gamma_{\omega\circ f}$ is therefore equal to
\[
  (\Gamma_\omega\times_X X')_{\mathrm{red}}.
\]
But by hypothesis the composite morphism
\[
  \Gamma_{\omega\circ f}\longrightarrow\Gamma_\omega\times_X X'
  \xrightarrow{p^{(X')}}X'
\]
is an isomorphism, so $p^{(X')}$ is necessarily radicial.
Since $f$ is surjective, the same is true of $p:\Gamma_\omega\to X$ \sref[I]{I.3.4.8}; hence $p^{-1}(x)$ consists of a single point \sref[I]{I.3.6.4}, and \sref{IV.20.4.4} shows that $\omega$ is defined at $x$.
\end{proof}

\begin{proposition}[20.4.6]
\label{IV.20.4.6}
Let $S$ be a prescheme, $X$, $Y$ two $S$-preschemes;
suppose $X$ locally Noetherian, $Y$ locally of finite type on $S$.
Let $U$ be a dense open in $X$, $h:U\to Y$ an $S$-morphism, $x\in X-U$ a normal point of $X$, $h'_x:\Spec(k(x))\to Y$ an $S$-morphism.
In order that $h$ can be extended to an $S$-morphism $h':U'\to Y$, where $U'$ is an open of $X$ containing $U$ and $x$, and where the composed morphism $\Spec(k(x))\to U'\xrightarrow{h'}Y$ is the $S$-morphism given by $h'_x$, it is necessary and sufficient that the following condition be satisfied:
\begin{quote}
\emph{(P) For every spectrum $S_1$ of a discrete valuation ring, of closed point $a$ and generic point $b$, and for every morphism $u:S_1\to X$ such that $u(a)=x$, $u(b)\in U$, the composed morphism $h_1=h\circ(u|\{b\}):\{b\}=u^{-1}(U)\to Y$ extends to a morphism $h'_1:S_1\to Y$ such that the diagram}
\[
  \xymatrix{
  \Spec(k(a)) \ar[r] & S_1 \ar[d]^{h'_1} \\
  \Spec(k(x)) \ar[r]_{h'_x} & Y
  }
\]
\emph{is commutative.}
\end{quote}
Moreover, if this condition is satisfied, and if $h'':U''\to Y$ is a morphism satisfying the same conditions as $h'$, then there exists an open set containing $U\cup\{x\}$, in which $h'$ and $h''$ coincide.
\end{proposition}

\begin{proof}
Let us first prove the last assertion;
we can suppose $h'$ and $h''$ defined in the same open $U_0\supset U\cup\{x\}$.
The sub-prescheme $Z$ of coincidences of $h'$ and $h''$ \sref{IV.17.4.5} contains $U$ and $x$, so there is an open neighbourhood $V$ of $x$ in $U_0$ such that the sub-prescheme induced by $Z$ on the open $Z\cap V$ is a \emph{closed} sub-prescheme of the prescheme induced by $V$ on $V$;
as this sub-prescheme $Z\cap V$ majorizes the prescheme induced by $V$ on the open $U\cap V$, and the latter is schematically dense in $V$, $Z\cap V$ is necessarily equal to $V$ \sref{IV.20.3.8.8}, which proves that $h'$ and $h''$ coincide in $U\cup V$.

The necessity of the condition is evident, let us prove that it is sufficient.
By the bijective correspondence between $S$-morphisms of $X$ in $Y$ and $X$-sections of $X\times_S Y$ \sref[I]{I.3.3.15}, we can restrict to the case where $S=X$ and $h$ is therefore a $U$-section of $Y$;
we note that then $Y$ is locally Noetherian, and we may plainly (since $X$ is locally integral and locally Noetherian) suppose that $X=S$ is irreducible.
Moreover, by virtue of \sref{IV.20.3.7}, we can suppose that $X=\Spec(A)$, where $A$ is a local Noetherian ring, integral and integrally closed.

\oldpage[IV-4]{247}
Let $y=h'_x(x)$.
For every $x'\in U$, the point $y$ is a specialization of $h(x')$.
Indeed, there exists the spectrum $S_1$ of a discrete valuation ring and a morphism $u:S_1\to X$ such that $u(a)=x$ and $u(b)=x'$ \sref[II]{II.7.1.9}.
Applying condition~(P), we obtain the assertion at once, since $h(x')=h_1(b)=h'_1(b)$ and $y=h'_1(a)$.
If $Y'$ is an affine open neighbourhood of $y$, it follows that $h(X\cap U)\subset Y'$;
we may therefore replace $Y$ by $Y'$ and hence suppose that $Y$ is affine, and thus separated over $X$.
Let $\omega$ be the rational $X$-section of $Y$ to which $h$ belongs, so that its graph has as underlying space the closure of $h(U)$ in $Y$.
Since $Y$ is separated over $X$, we may apply \sref{IV.20.4.4} to $\omega$;
it is enough to prove that, if $p:\Gamma_\omega\to X$ is the canonical projection, then $p^{-1}(x)$ consists of the single point $y$ and $h'_x(x)=y$.
Indeed, \sref{IV.20.4.4} then extends $h$ to a section $h'$ over an open $U'$ containing $U\cup\{x\}$, with $h'(x)=y$;
and since there is only one $X$-morphism $\Spec(k(x))\to Y$ carrying $x$ to $y$, this proves that $h'_x$ is the composite of $h'$ with $\Spec(k(x))\to U'$.

Since, for every $x'\in U$, $y$ is a specialization of $h(x')$, we have $y\in p^{-1}(x)$.
Now let $z$ be any point of $p^{-1}(x)$.
Since $\Gamma_\omega$ is the closure of $h(U)$, and every point of $h(U)$ is adherent to $h(\xi)$, where $\xi$ is the generic point of $X$, $\Gamma_\omega$ is the closure of $h(\xi)$ in $Y$.
There therefore exist the spectrum $S_1$ of a discrete valuation ring and a morphism $v:S_1\to Y$ such that $v(a)=z$ and $v(b)=h(\xi)$ \sref[II]{II.7.1.9};
set $u=p\circ v$, so that $u(a)=x$ and $u(b)=\xi$.
Applying condition~(P) to $u$, we obtain a morphism $h'_1:S_1\to Y$ such that $h'_1(a)=y$ and $h'_1(b)=h(\xi)$.
But then $z=y$ by \sref[II]{II.7.2.3}, since $Y$ is separated over $X$ and $v$ and $h'_1$ must coincide, being equal at $b$.
C.Q.F.D.
\end{proof}

\begin{corollary}[20.4.7]
\label{IV.20.4.7}
The hypotheses on $S$, $X$, $Y$, $U$ and $h$ being those of \sref{IV.20.4.6}, let $E$ be a part of $X-U$ such that $X$ is normal at every point of $E$, and for each $x\in E$, let $h'_x:\Spec(k(x))\to Y$ be an $S$-morphism, such that the condition~(P) of \sref{IV.20.4.6} is satisfied.
Suppose moreover that the union $F$ of the graphs of the $h'_x$ (identified with parts of $X\times_S Y$) for $x\in E$, is contained in a union of a finite number of opens $V_i$ of $X\times_S Y$ which are \emph{separated} on $X$ (a condition automatically satisfied if $Y$ is \emph{separated} on $S$, or if $X$ is \emph{quasi-compact} and $Y$ of \emph{finite type} on $S$).
Then there exists an $S$-morphism $h':U'\to Y$, where $U'$ is an open of $X$ containing $U\cup E$, such that, for every $x\in E$, the composed
\[
  \Spec(k(x))\to U'\xrightarrow{h'}Y
\]
is equal to $h'_x$.
\end{corollary}

\begin{proof}
Note first that, in \sref{IV.20.4.6}, if $Y$ is supposed \emph{separated}, there is a largest open $U_0\supset U$, equal to the domain of the $S$-rational application corresponding to $h$, in which $h$ can be extended, and this extension is unique $(\mathbf{I}$, 7.2.2);
whence the conclusion in this case, by virtue of \sref{IV.20.4.6}.
In the general case, let $E_i$ be the set of $x\in E$ such that $(x,h'_x(x))\in V_i$.
By virtue of \sref{IV.20.4.6}, for every $x\in E$, there is an extension $h^{(x)}$ of $h$ to $U\cup W^{(x)}$, where $W^{(x)}$ is an open neighbourhood of $x$ in $X$ such that the graph of $h^{(x)}|W^{(x)}$ is contained in all the $V_i$ such that $x\in E_i$.
Since $V_i$ is separated over $X$ and $X$ is reduced, for two points $x',x''\in E_i$, the morphisms $h^{(x')}$ and $h^{(x'')}$ coincide on $W^{(x')}\cap W^{(x'')}$, because they
\oldpage[IV-4]{248}
coincide on the everywhere dense open $U\cap W^{(x')}\cap W^{(x'')}$.
There is therefore a morphism $h_i:U\cup U_i\to Y$ extending $h$, where $U_i$ is an open of $X$ containing $E_i$.
Moreover, for every pair of indices $i,j$, the graphs of the restrictions $h_i|(U_i\cap U_j)$ and $h_j|(U_i\cap U_j)$ are contained in $V_i\cap V_j$;
since $V_i\cap V_j$ is separated over $X$ and the preceding morphisms coincide on the everywhere dense open $U\cap U_i\cap U_j$ of $U_i\cap U_j$, they are equal.
The morphism $h'$, equal to $h$ on $U$ and to $h_i$ on each $U_i$, answers the question.
\end{proof}

\begin{remark}[20.4.8]
\label{IV.20.4.8}
When $E=X-U$, it is clear that if, for every affine open $T$ of $X$, there exists an $S$-morphism $h'_T:T\to Y$ extending $h|(U\cap T)$ and such that, for every $x\in T-(U\cap T)$, the composite $\Spec(k(x))\to T\xrightarrow{h'_T}Y$ is equal to $h'_x$, then the $h'_T$ are the restrictions of an everywhere-defined $S$-morphism $h':X\to Y$ extending $h$, by the uniqueness assertion in \sref{IV.20.4.6}.
To prove the existence of $h'$, one is therefore reduced to the case where $X$ is \emph{affine}, and it then suffices that the union $F$ of the graphs of the $h'_x$, is \emph{quasi-compact} in $X\times_S Y$ so that one can apply \sref{IV.20.4.7}.
The latter will hold in particular when the $h'_x$ are of the form $\Spec(k(x))\to Z\xrightarrow{h_0}Y$, where $Z$ is a closed sub-prescheme of $X$ having $X-U$ for underlying space, and $h_0$ an $S$-morphism.
\end{remark}

\begin{corollary}[20.4.9]
\label{IV.20.4.9}
Let $S$ be a locally Noetherian prescheme, $X$ a locally Noetherian prescheme, $f:X\to S$ a flat morphism, and $g:Y\to S$ a morphism locally of finite type.
Let $U$ be a dense open of $S$, $h:f^{-1}(U)\to Y$ an $S$-morphism, $T=S-U$, $Z$ the reduced closed sub-prescheme of $X$ having $f^{-1}(T)$ as underlying space, and $h_0:Z\to Y$ an $S$-morphism.
Suppose $X$ normal at all points of $Z$.
In order that there exist a (necessarily unique) $S$-morphism $h':X\to Y$ extending $h$ and $h_0$, it is necessary and sufficient that the following condition be satisfied:

\emph{For every spectrum $S_1$ of a discrete valuation ring, of closed point $a$ and generic point $b$, and for every morphism $u:S_1\to S$ such that $u(a)\in T$ and $u(b)\in U$, there exists an $S_1$-morphism $h'_1:X_{(S_1)}\to Y_{(S_1)}$ which extends $h_{(S_1)}:f_{(S_1)}^{-1}(b)\to Y_{(S_1)}$ and is such that the diagram}
\[
  \xymatrix{
  \Spec(k(z_1)) \ar[r] \ar[d] & Z_{(S_1)} \ar[d]^{(h_0)_{(S_1)}} \\
  X_{(S_1)} \ar[r]_{h'_1} & Y_{(S_1)}
  }
\]
\emph{is commutative for every $z_1\in Z_{(S_1)}$.}

The hypothesis that $f$ is flat implies indeed that $f^{-1}(U)$ is dense in $X$ \sref{IV.2.3.10}, and it then suffices to apply \sref{IV.20.4.8}.
\end{corollary}

\begin{corollary}[20.4.10]
\label{IV.20.4.10}
Under the hypotheses of \sref{IV.20.4.6}, suppose moreover $Y$ separated and locally quasi-finite on $S$.
Let $U$ be a dense open in $X$, $h:U\to Y$ an $S$-morphism, $\omega$ the $S$-rational application corresponding to it, $x\in X-U$ a normal point of $X$.
In order that $\omega$ be defined at the point $x$, it is necessary and sufficient that the following condition be satisfied:
\emph{there exists a spectrum $S_1$ of a discrete valuation ring, of closed point $a$ and generic point $b$, and a morphism $u:S_1\to X$}
\oldpage[IV-4]{249}
\emph{such that $u(a)=x$, $u(b)\in U$, and the $S$-rational application $\omega\circ u$ is everywhere defined.}

Indeed, by hypothesis all the fibres of the projection morphism $X\times_S Y\to X$ consist of isolated points, and to apply \sref{IV.20.4.4} it suffices to know that the fibre $p^{-1}(x)$ is nonempty.
Now, if $h_1$ is the unique morphism in the class $\omega\circ u$, the unique point of $X\times_S Y$ lying over $x$ and over $h_1(a)$ belongs to $\Gamma_\omega$, whence the conclusion.
\end{corollary}

\begin{proposition}[20.4.11]
\label{IV.20.4.11}
Let $X$ be a locally Noetherian prescheme, $Y$ an $S$-prescheme affine over $S$, $U$ an open of $X$, and $Z=X-U$.
Suppose that $\operatorname{prof}_Z(\sh{O}_X)\geq 2$ \sref{IV.5.10.1};
then every $S$-morphism $f:U\to Y$ extends uniquely to an $S$-morphism of $X$ in $Y$.
\end{proposition}

\begin{proof}
We can restrict to the case where $S$ and $X$ are affines, and (by virtue of \sref[I]{I.3.3.14}) to the case $S=X$;
so one has $X=\Spec(A)$, $Y=\Spec(B)$, $B$ being a finite type $A$-algebra.
As $B$ is a quotient of a polynomial algebra $A[(T_\lambda)]_{\lambda\in L}$, $Y$ is a closed sub-prescheme of $Y'=X[(T_\lambda)]_{\lambda\in L}$.
On the other hand, the hypothesis on $Z$ implies that $U$ is schematically dense in $X$ by virtue of \sref{IV.20.2.13}[(iv)] and \sref{IV.5.10.2}.
If we prove that every $X$-morphism $u:U\to Y$ extends uniquely to an $X$-morphism $v':X\to Y'$, it will result that $v'$ will factorise as $X\xrightarrow{v}Y\to Y'$:
indeed, the sub-prescheme $v'^{-1}(Y)$ is closed and contains $U$ \sref[I]{I.4.4.1}, and hence is equal to $X$ \sref{IV.20.3.8.8}.
In these conditions, $v$ will be the unique $S$-morphism of $X$ in $Y$ extending $u$.
We can therefore restrict to the case where $Y=Y'$.
But then there is a bijective correspondence between the $X$-morphisms of an open $U\subset X$ in $Y'$ and the families $(s_\lambda)_{\lambda\in L}$ of sections of $\sh{O}_X$ over $U$ \sref[II]{II.1.7.9};
the conclusion results therefore from \sref{IV.5.10.5}.
\end{proof}

\begin{corollary}[20.4.12]
\label{IV.20.4.12}
Let $X$ be an $S$-prescheme locally Noetherian reduced and satisfying the condition~$(S_2)$ \sref{IV.5.7.2} (for example a locally Noetherian and \emph{normal} $S$-prescheme \sref{IV.5.8.6}), $Y$ an $S$-prescheme affine on $S$, $f$ an $S$-rational application of $X$ in $Y$;
then every irreducible component of $X-\operatorname{dom}(f)$ is of codimension $1$ in $X$.
\end{corollary}

\begin{proof}
This amounts to the same as saying that if $Z^{(2)}$ is the set of $x\in X$ such that $\dim(\sh{O}_{X,x})\geq 2$, then, for every closed part $Z\subset Z^{(2)}$ of $X$, every $S$-morphism of $X-Z$ in $Y$ extends to an $S$-morphism of $X$ in $Y$;
but this results from the hypothesis on $X$ \sref{IV.5.7.2} and from \sref{IV.20.4.11}.
\end{proof}

\subsection{Relative pseudo-morphisms}
\label{IV.20.5}

\begin{env}[20.5.1]
\label{IV.20.5.1}
Let $X$, $Y$ be two $S$-preschemes.
It results from the definitions \sref{IV.11.10.8} that the intersection of two opens $U$, $U'$ of $X$, \emph{universally schematically dense with respect to $S$}, possesses this same property.
We therefore define an equivalence relation between $S$-morphisms $u:U\to Y$ by replacing in \sref{IV.20.2.1} ``schematically dense'' by ``universally schematically dense with respect to $S$''.
An equivalence class for this relation is called a \emph{pseudo-morphism of $X$ in $Y$ relatively to $S$}, and the set of these classes is denoted $\operatorname{Ps.hom}_{X/S}(X,Y)$.
\end{env}

\begin{env}[20.5.2]
\label{IV.20.5.2}
As every open of $X$ universally schematically dense with respect to $S$ is in particular schematically dense, the elements of a pseudo-morphism
\oldpage[IV-4]{250}
of $X$ in $Y$ relatively to $S$ are equivalent in the sense of \sref{IV.20.2.1}, whence a canonical map
\[
\label{IV.20.5.2.1}
  \operatorname{Ps.hom}_{X/S}(X,Y)\to\operatorname{Ps.hom}_S(X,Y).
  \tag{20.5.2.1}
\]
\end{env}

\begin{proposition}[20.5.3]
\label{IV.20.5.3}
Suppose $Y$ separated over $S$.
Then:
\begin{enumerate}[label=\emph{(\roman*)}]
  \item The map \sref{IV.20.5.2.1} is injective and identifies $\operatorname{Ps.hom}_{X/S}(X,Y)$ with the subset of $\operatorname{Ps.hom}_S(X,Y)$ formed of the pseudo-$S$-morphisms $\omega$ such that $\operatorname{dom}_S(\omega)$ is universally schematically dense relatively to $S$.
  \item The presheaf $U\mapsto\operatorname{Ps.hom}_{U/S}(U,Y)$ on $X$ is a sheaf.
\end{enumerate}
\end{proposition}

\begin{proof}
Assertion~(i) is immediate, for two $S$-morphisms $u:U\to Y$, $u':U'\to Y$ are equivalent in the sense of \sref{IV.20.2.1} means that they are the restrictions of a single morphism $\operatorname{dom}_S(\omega)\to Y$ \sref{IV.20.2.4}, and if $U$ and $U'$ are universally schematically dense with respect to $S$, the same holds \emph{a fortiori} for $\operatorname{dom}_S(\omega)$.
To prove~(ii), note that $U\mapsto\operatorname{Ps.hom}_S(U,Y)$ is then a sheaf \sref{IV.20.2.6};
on the other hand, given an open covering $(U_\alpha)$ of $U$ and a pseudo-$S$-morphism $\omega$ from $U$ to $Y$, it follows immediately from the definitions (cf.\ \sref{IV.20.2.1}) that $\operatorname{dom}_S(\omega)$ is universally schematically dense in $U$ relative to $S$ if and only if each $\operatorname{dom}_S(\omega)\cap U_\alpha=\operatorname{dom}_S(\omega|U_\alpha)$ is universally schematically dense in $U_\alpha$ relative to $S$;
whence~(ii).
\end{proof}

When $Y$ is separated, we denote by $\mathscr{P}\!s.\!\operatorname{hom}_{X/S}(X,Y)$ the sub-sheaf
\[
  U\mapsto\operatorname{Ps.hom}_{U/S}(U,Y)
\]
of $\mathscr{P}\!s.\!\operatorname{hom}_S(X,Y)$.

When $X$ is \emph{flat and locally of finite presentation} over $S$ and $Y$ separated over $S$, the pseudo-morphisms of $X$ in $Y$ relatively to $S$ are identified with the pseudo-$S$-morphisms $\omega$ such that, for every fibre $X_s$ of the morphism $X\to S$, $\operatorname{dom}_S(\omega)\cap X_s$ is \emph{schematically dense} in $X_s$ \sref{IV.11.10.10}.

\begin{env}[20.5.4]
\label{IV.20.5.4}
A particularly important case is that where $Y$ is separated over $S$ and $Y=S[T]=S\otimes_{\bb{Z}}\Spec(\bb{Z}[T])$ ($T$ indeterminate).
There is then a bijective correspondence between the pseudo-$S$-morphisms of $X$ in $Y$ and the \emph{pseudo-functions} on $X$ \sref{IV.20.2.8} by the definition of a product of $\bb{Z}$-preschemes.
The pseudo-morphisms of $X$ in $Y$ relatively to $S$ are identified then, by virtue of \sref{IV.20.5.3}, with the \emph{pseudo-functions $\varphi$ on $X$ such that $\operatorname{dom}(\varphi)$ is universally schematically dense relatively to $S$}.
The sheaf $\mathscr{P}\!s.\!\operatorname{hom}_{X/S}(X,Y)$ is a \emph{sub-sheaf of rings} of $\sh{M}'_X$, that we denote $\sh{M}'_{X/S}$.

We then set $\operatorname{Ps.hom}_{X/S}(X,Y)=M'(X/S)$ and we say that its elements are the \emph{pseudo-functions on $X$ relative to $S$}.
\end{env}

\begin{env}[20.5.5]
\label{IV.20.5.5}
Let $X$, $Y$, $Z$ be three $S$-preschemes, $f:Y\to Z$ an $S$-morphism.
We can, in the reasoning of \sref{IV.20.3.1}, replace everywhere ``schematically dense'' by ``universally schematically dense relatively to $S$'';
for every pseudo-morphism $\omega\in\operatorname{Ps.hom}_{X/S}(X,Y)$, the morphisms $f\circ u$, where $u$ runs through $\omega$, belong therefore to a single equivalence class (for the relation defined in \sref{IV.20.5.1}), which we
\oldpage[IV-4]{251}
call the pseudo-morphism from $X$ to $Z$, relative to $S$, \emph{composite} of $f$ and $\omega$, and denote $f\circ\omega$.
If $g:Z\to T$ is a morphism, it is immediate that $g\circ(f\circ\omega)=(g\circ f)\circ\omega$.
\end{env}

\begin{env}[20.5.6]
\label{IV.20.5.6}
Suppose $Y$ \emph{separated} on $S$, and let $\omega$ be a pseudo-morphism of $X$ in $Y$ relatively to $S$.
If $f:X'\to X$ is an $S$-morphism such that $f^{-1}(\operatorname{dom}_S(\omega))$ is universally schematically dense in $X'$ \emph{relative to $S$}, it follows from \sref{IV.20.3.2} that the pseudo-$S$-morphism $\omega\circ f$ has a domain universally schematically dense relative to $S$, and hence, by \sref{IV.20.5.3}, may be regarded as a pseudo-morphism \emph{relative to $S$}.
When $X'$ is flat and locally of finite presentation over $S$, the condition that $f^{-1}(\operatorname{dom}_S(\omega))$ be universally schematically dense in $X'$ relative to $S$ is equivalent to saying that for every $s\in S$, $(f^{-1}(\operatorname{dom}_S(\omega)))_s$ (notation of \sref{IV.11.10.10}) is schematically dense in $X'_s$.
Equivalently, writing $f_s:X'_s\to X_s$ for the morphism deduced from $f$ by base change, the inverse image under $f_s$ of $(\operatorname{dom}_S(\omega))_s$ is schematically dense in $X'_s$.
This latter condition will in particular be satisfied if, for every $s\in S$, $X_s$, $X'_s$ and $f_s$ satisfy one of the three conditions~(i),~(ii),~(iii) of \sref{IV.20.3.5}.
\end{env}

\begin{env}[20.5.7]
\label{IV.20.5.7}
Suppose now that $X$ and $X'$ are both $S$-preschemes \emph{flat and locally of finite presentation} over $S$, and that $f:X'\to X$ is an $S$-morphism \emph{flat} (or, what amounts to the same \sref{IV.11.3.10}, that for every $s\in S$, $f_s:X'_s\to X_s$ is a flat morphism).
Then, for every open $U$ of $X$ and every open $V\subset U$ universally schematically dense in $U$ relatively to $S$, it results from \sref{IV.11.10.5} and \sref{IV.11.10.9} that $f^{-1}(V)$ is universally schematically dense in $f^{-1}(U)$ relatively to $S$.
For every pseudo-$S$-morphism $\omega$ of $X$ in $Y$ relatively to $S$, it results from \sref{IV.20.3.4} that the pseudo-$S$-morphism $\omega\circ f$ is defined and is a pseudo-morphism of $X'$ in $Y$ relatively to $S$, \emph{even when $Y$ is not supposed separated over $S$}.
We deduce from this that for every $S$-morphism $g:Y\to Z$, $(g\circ\omega)\circ f$ is still defined and equal to $g\circ(\omega\circ f)$ (with the definitions of \sref{IV.20.5.1}), and that it is therefore a pseudo-morphism \emph{relative} to $S$.
\end{env}

\begin{env}[20.5.8]
\label{IV.20.5.8}
Let $X$ be an $S$-prescheme, $S'\to S$ a morphism, $g:X'\to X$ the canonical projection, $X'=X\times_S S'$.
Then, for every open $U$ of $X$ and every open $V\subset U$ universally schematically dense in $U$ relatively to $S$, $V'=g^{-1}(V)$ is universally schematically dense in $U'=g^{-1}(U)$ relatively to $S'$ by definition \sref{IV.11.10.8}.
Let $\omega$ be a pseudo-$S$-morphism of $X$ in an $S$-prescheme $Y$ relatively to $S$;
if $u_1$, $u_2$ are $S$-morphisms of the class $\omega$, defined respectively on opens $U_1$, $U_2$ of $X$, universally schematically dense in $X$ relative to $S$, it follows from what precedes that $u'_1=u_1\circ(g|g^{-1}(U_1))$ and $u'_2=u_2\circ(g|g^{-1}(U_2))$ coincide on an open $U'_3$ universally schematically dense relative to $S'$.
Now, if $Y'=Y_{(S')}$ and $h:Y'\to Y$ is the canonical projection, $u'_1$ and $u'_2$ factor canonically as $u'_1=h\circ v_1$ and $u'_2=h\circ v_2$, where $v_1$ and $v_2$ are $S'$-morphisms into $Y'$ which coincide on $U'_3$.
Thus, as $u_1$ runs through the class $\omega$, the corresponding $S'$-morphisms $v_1$ belong to a single pseudo-morphism relative to $S'$, called the \emph{inverse image} of $\omega$ under the base change $S'\to S$ and denoted $\omega_{(S')}$.
It is clear that if $S''\to S'$ is a second morphism, $(\omega_{(S')})_{(S'')}=\omega_{(S'')}$ (for the composed morphism $S''\to S'\to S$).
\end{env}

\oldpage[IV-4]{252}
\subsection{Relative meromorphic functions}
\label{IV.20.6}

\begin{env}[20.6.1]
\label{IV.20.6.1}
Let $S$ be a prescheme, $X$ an $S$-prescheme which is \emph{flat and locally of finite presentation} over $S$;
for every $s\in S$, we denote by $X_s$ the fibre at the point $s$ of the structural morphism $X\to S$.
In general, if $\varphi$ is a meromorphic function on $X$, it is not possible, for each $s\in S$, to associate to it in a ``natural'' way a meromorphic function ``induced'' on $X_s$ \sref{IV.20.1.11}.
For every open $U$ of $X$, denote by $T_{X/S}(U)$ the set of sections $t\in\Gamma(U,\sh{O}_X)$ such that, for every $s\in S$, the image $t_s$ of $t$ by the canonical homomorphism $\Gamma(U,\sh{O}_X)\to\Gamma(U\cap X_s,\sh{O}_{X_s})$ is a \emph{regular} section;
this also implies, by the equivalence of a) and b) in \sref{IV.11.3.7}, that $t$ is itself a \emph{regular} section.
It is clear that $U\mapsto T_{X/S}(U)$ is a \emph{sub-sheaf} of the sheaf of sets $\sh{S}=\sh{S}(\sh{O}_X)$ (notation of \sref{IV.20.1.3}), that we denote $\sh{T}_{X/S}$;
we set
\[
\label{IV.20.6.1.1}
  \sh{M}_{X/S}=\sh{O}_X[\sh{T}_{X/S}^{-1}]
  \tag{20.6.1.1}
\]
(notation of \sref{IV.20.1.1}) and we say that this sheaf of rings is the \emph{sheaf of germs of meromorphic functions on $X$ relative to $S$};
its sections over $X$ are called \emph{meromorphic functions on $X$ relative to $S$} and their set is denoted $M(X/S)$.
It is clear that $\sh{M}_{X/S}$ is a \emph{sub-sheaf} of $\sh{M}_X=\sh{O}_X[\sh{S}^{-1}]$;
for every meromorphic function $\varphi\in M(X/S)$ and every $s\in S$, the inverse image of $\varphi$ by the canonical injection morphism $j_s:X_s\to X$ is then defined \sref{IV.20.1.11}, and denoted $\varphi_s$.
\end{env}

\begin{env}[20.6.2]
\label{IV.20.6.2}
Let now $\sh{F}$ be an $\sh{O}_X$-Module quasi-coherent;
we set
\[
\label{IV.20.6.2.1}
  \sh{M}_{X/S}(\sh{F})=\sh{F}[\sh{T}_{X/S}^{-1}]=\sh{F}\otimes_{\sh{O}_X}\sh{M}_{X/S};
  \tag{20.6.2.1}
\]
the sections of $\sh{M}_{X/S}(\sh{F})$ are called \emph{meromorphic sections of $\sh{F}$ over $X$, relative to $S$} and their set is denoted $M(X/S,F)$.
The canonical homomorphism $\sh{F}\to\sh{M}_{X/S}(\sh{F})$ is not necessarily injective;
when it is, we say that $\sh{F}$ is \emph{without torsion relatively to $S$}:
this means that for every open $U$ of $X$ and every section $t\in T_{X/S}(U)$, $t$ is $(\sh{F}|U)$-regular;
this condition is satisfied \emph{a fortiori} when $\sh{F}$ is torsion-free \sref{IV.20.1.5}.
In this last case, it results immediately from the definitions \sref{IV.20.1.2} that the canonical homomorphism of $\sh{O}_X$-Modules
\[
\label{IV.20.6.2.2}
  \sh{M}_{X/S}(\sh{F})\to\sh{M}_X(\sh{F})
  \tag{20.6.2.2}
\]
is \emph{injective}, so that the meromorphic sections of $\sh{F}$ relative to $S$ are meromorphic sections of $\sh{F}$ in the sense of \sref{IV.20.1.3}.
\end{env}

\begin{proposition}[20.6.3]
\label{IV.20.6.3}
The image by the injective homomorphism \sref{IV.20.2.10.1} of the sub-sheaf $\sh{M}_{X/S}$ of $\sh{M}_X$ is the sub-sheaf $\sh{M}'_{X/S}$ of pseudo-functions on $X$ relative to $S$ (i.e.\ of the pseudo-functions whose domain of definition is universally schematically dense relatively to $S$ \sref{IV.20.5.4}).
\end{proposition}

\begin{proof}
We can evidently restrict to proving that the image of $M(X/S)$ by the canonical homomorphism $M(X)\to M'(X)$ is equal to $M'(X/S)$;
the proposition is then a consequence of the following more general proposition:
\end{proof}

\oldpage[IV-4]{253}
\begin{proposition}[20.6.4]
\label{IV.20.6.4}
Let $\sh{F}$ be an $\sh{O}_X$-Module quasi-coherent of finite presentation and torsion-free.
Then, for a meromorphic section $\varphi$ of $\sh{F}$ over $X$ to be a meromorphic section relative to $S$, it is necessary and sufficient that $\operatorname{dom}(\varphi)$ be universally schematically dense relatively to $S$.
\end{proposition}

\begin{proof}
The necessity of the condition results from \sref{IV.20.2.15} applied to each canonical injection $X_s\to X$ ($s\in S$), taking into account \sref{IV.11.10.9}.
To see that the condition is sufficient, it suffices to prove that for every $x\in X$, there exists a neighbourhood $U$ of $x$ and a section of $\sh{M}_{X/S}(\sh{F})$ over $U$ whose restriction to $U\cap\operatorname{dom}(\varphi)$ coincides with $\varphi$ in an open schematically dense of $U\cap\operatorname{dom}(\varphi)$.
Consider the ideal of denominators $\sh{J}$ of $\varphi$ \sref{IV.20.2.14}, which is quasi-coherent and defines the closed sub-prescheme of $X$ whose underlying space is $X-\operatorname{dom}(\varphi)$.
By hypothesis, if $s$ is the image of $x$ in $S$, $\operatorname{dom}(\varphi)\cap X_s$ is schematically dense in the locally Noetherian prescheme $X_s$, and hence contains $\operatorname{Ass}(\sh{O}_{X_s})$ \sref{IV.20.2.13}[(iv)];
this implies that the Ideal $\sh{J}_x$ has an image in $\sh{O}_{X_s,x}=\sh{O}_{X,x}\otimes_{\sh{O}_{S,s}}k(s)$ which is not contained in any of the prime ideals $\mathfrak{p}_i\in\operatorname{Ass}(\sh{O}_{X_s,x})$ (finitely many);
hence (Bourbaki, \emph{Alg.~comm.}, chap.~II, \textsection1, n\textsuperscript{o}~1, prop.~2) there exists an element $t_x\in\sh{J}_x$ whose image in $\sh{O}_{X_s,x}$ does not belong to any of the $\mathfrak{p}_i$, and is therefore \emph{regular} in this Noetherian ring.
Let $t$ be a section of $\sh{J}$ over an affine neighbourhood $U$ of $x$, whose germ at $x$ is $t_x$;
since $X$ is flat and locally of finite presentation over $S$, we can suppose \sref{IV.11.3.8} that $t$ is a regular section of $\sh{O}_X$ over $U$ and that, for every $s'\in S$, the image of $t$ in $\Gamma(U\cap X_{s'},\sh{O}_{X_{s'}})$ is also regular;
in other words, $t\in T_{X/S}(U)$.
But then, since $\sh{F}$ is torsion-free, $t(\varphi|(U\cap\operatorname{dom}(\varphi)))$ is a section $u$ of $\sh{F}$ over $U\cap\operatorname{dom}(\varphi)$;
on the other hand, $U\cap\operatorname{dom}(\varphi)$ contains the open $U_t$ of $x'\in U$ where $t(x')\neq 0$, and this last is schematically dense in $U$ \sref{IV.20.2.9}.
We see that in $U_t$, $\varphi$ coincides with the restriction to $U_t$ of the section $u/t$ of $\sh{M}_{X/S}(\sh{F})$ over $U$.
C.Q.F.D.
\end{proof}

\begin{remarks}[20.6.5]
\label{IV.20.6.5}
\begin{enumerate}[label=(\roman*)]
  \item Let $\varphi$ be a meromorphic function on $X$ relative to $S$, so that for every $s\in S$, $\varphi_s$ is a meromorphic function on $X_s$ \sref{IV.20.6.1};
  by virtue of \sref{IV.20.1.11.1}, one has
  \[
  \label{IV.20.6.5.1}
    \operatorname{dom}(\varphi)\cap X_s\subset\operatorname{dom}(\varphi_s).
    \tag{20.6.5.1}
  \]
  But it is worth noting that even when $S$ is the spectrum of a discrete valuation ring $A$ and $X=S[T]$ ($T$ indeterminate), the two sides of \sref{IV.20.6.5.1} are not necessarily equal:
  for example, if $\pi$ is a uniformizer of $A$, it is immediate that $\varphi=\pi/T$ is a meromorphic function on $X$ relative to $S$, since if $a$ and $b$ are the closed and generic points of $S$, $T$ is regular in $\Gamma(X_a,\sh{O}_{X_a})=k[T]$ and in $\Gamma(X_b,\sh{O}_{X_b})=K[T]$, $k$ being the residue field and $K$ the field of fractions of $A$.
  One has $\operatorname{dom}(\varphi)=\mathrm{D}(T)$ in $X$, but $\operatorname{dom}(\varphi_a)=X_a$ since $\varphi_a=0$.

  \item For a meromorphic function $\varphi$ relative to $S$ to be \emph{invertible in the ring} $M(X/S)$, it is necessary and sufficient that for every $s\in S$, $\varphi_s$ be invertible in $M(X_s)$ (in other words, that $\varphi_s$ be a regular meromorphic function on $X_s$ \sref{IV.20.1.8}).
  The condition is
  \oldpage[IV-4]{254}
  trivially necessary.
  Conversely, if it is satisfied, and if $x$ is any point of $X$, $s$ its image in $S$, there exists by hypothesis a neighbourhood $U$ of $x$ and two sections $u$, $t$ of $\sh{O}_X$ over $U$ such that $t\in T_{X/S}(U)$ and $\varphi|U=u/t$;
  the hypothesis implies that $u_s$ is regular at $x$ in $\Gamma(U\cap X_s,\sh{O}_{X_s})$, so $u_s$ is regular at point $x$.
  Restricting $U$, we can therefore suppose, by virtue of \sref{IV.11.3.8}, that $u\in T_{X/S}(U)$, whence the conclusion.

  When $\varphi$ is invertible in $M(X/S)$, we also say that $\varphi$ is a \emph{regular meromorphic function relative to $S$}.
  We note that $\varphi\in M(X/S)$ can be invertible in $M(X)$ (in other words, be \emph{regular} in the sense of \sref{IV.20.1.8}) without being invertible in $M(X/S)$, as the above example of~(i) shows.

  \item Let $\sh{L}$ be an invertible $\sh{O}_X$-Module, and let $\varphi$ be a regular meromorphic section of $\sh{L}$ over $X$ \sref{IV.20.1.8};
  we say that $\varphi$ is \emph{regular relatively to $S$} if, for every open $U$ of $X$ such that $\sh{L}|U$ is isomorphic to $\sh{O}_U$, $\varphi|U$ corresponds to an element of $\Gamma(U,\sh{M}_X)$ which is \emph{regular relatively to $S$} \sref{IV.20.1.8};
  by virtue of~(ii), it is immediate that it is necessary and sufficient that, for every $s\in S$, $\varphi_s$ be a regular meromorphic section \sref{IV.20.1.8} of the invertible $\sh{O}_{X_s}$-Module $\sh{L}_s=\sh{L}\otimes_{\sh{O}_X}k(s)$.
  If $\varphi'$ is the inverse of $\varphi$ in $\sh{L}^{-1}$ \sref{IV.20.1.10}, $\varphi'$ is then also regular relatively to $S$, and if $\sh{L}_1$ is a second invertible $\sh{O}_X$-module, $\varphi_1$ a meromorphic section of $\sh{L}_1$ over $X$, regular relatively to $S$, then $\varphi\otimes\varphi_1$ is a meromorphic section of $\sh{L}\otimes\sh{L}_1$ over $X$, regular relatively to $S$.
\end{enumerate}
\end{remarks}

\begin{proposition}[20.6.6]
\label{IV.20.6.6}
Let $X$ be an $S$-prescheme flat and locally of finite presentation over $S$, and let $\sh{F}$ be an $\sh{O}_X$-Module locally free of finite type;
for every $s\in S$, denote by $X_s$ the fibre at the point $s$ of the structural morphism $f:X\to S$.
Let $\varphi$ be a meromorphic section of $\sh{F}$ over $X$, \emph{relatively to $S$}, and suppose that $\varphi$ is defined at all points $x\in X$ such that $\operatorname{prof}(\sh{O}_{X_{f(x)},x})=1$.
Then $\varphi$ is everywhere defined.
\end{proposition}

\begin{proof}
By hypothesis, $\operatorname{dom}(\varphi)\cap X_s$ is schematically dense in $X_s$ for every $s\in S$, hence contains the points $x$ of $X_s$ such that $\operatorname{prof}(\sh{O}_{X_s,x})=0$ \sref{IV.5.10.2};
the hypothesis therefore signifies that if one sets $Z=X-\operatorname{dom}(\varphi)$, one has $\operatorname{prof}(\sh{O}_{X_{f(x)},x})\geq 2$ for all points $x$ of $Z$.
It suffices then to apply \sref{IV.19.9.8}.
\end{proof}

\begin{env}[20.6.7]
\label{IV.20.6.7}
Let $X$, $X'$ be two $S$-preschemes flat and locally of finite presentation, $f=(\psi,\theta):X'\to X$ an $S$-morphism.
For every open $U$ of $X$, denote by $T_f(U)$ the set of sections $t\in T_{X/S}(U)$ whose image in $\Gamma(f^{-1}(U),\sh{O}_{X'})$ belongs to $T_{X'/S}(f^{-1}(U))$;
it is immediate that $U\mapsto T_f(U)$ is a sub-sheaf of the sheaf of sets $\sh{T}_{X/S}$, that we denote $\sh{T}_f$.
We set $\sh{M}_{X/S,f}=\sh{O}_X[\sh{T}_f^{-1}]$;
this is a sub-sheaf of rings of $\sh{M}_{X/S}$, and we canonically obtain from $\theta^\#:\psi^*(\sh{O}_X)\to\sh{O}_{X'}$ a homomorphism of sheaves of rings $\theta'^{\#}:\psi^*(\sh{M}_{X/S,f})\to\sh{M}_{X'/S}$ extending $\theta^\#$.
A meromorphic function $\varphi$ on $X$, relative to $S$, that \emph{is a section of $\sh{M}_{X/S,f}$}, then has an inverse image $\Gamma(\theta'^{\#})(\varphi)$ which is a meromorphic function on $X'$ relative to $S$, called the \emph{inverse image of $\varphi$ by $f$}, and denoted $\varphi\circ f$ if this causes no confusion.
We extend immediately in the same way the definitions of \sref{IV.20.1.11} relative to $\sh{O}_X$-Modules.
\end{env}

\begin{proposition}[20.6.8]
\label{IV.20.6.8}
With the notations of \sref{IV.20.6.7}, if the $S$-morphism $f:X'\to X$ is \emph{flat}, one has $\sh{M}_{X/S,f}=\sh{M}_{X/S}$, and the homomorphism $\varphi\mapsto\varphi\circ f$ is defined on all of $M(X/S)$.
\end{proposition}

\begin{proof}
\oldpage[IV-4]{255}
Indeed, the hypothesis implies, by virtue of \sref{IV.11.3.10}, that for every $s\in S$, $f_s:X'_s\to X_s$ is flat;
hence, for every section $t\in T_{X/S}(U)$, if $t'$ is its inverse image, $t'_s$ is a \emph{regular} section of $\sh{O}_{X'_s}$ over $f^{-1}(U)\cap X'_s$;
by virtue of \sref{IV.20.1.12}, we conclude that by definition $t'\in T_{X'/S}(f^{-1}(U))$, whence the proposition.
\end{proof}

We deduce from this, as in \sref{IV.20.1.12}, a canonical homomorphism of $\sh{O}_{X'}$-algebras
\[
\label{IV.20.6.8.1}
  f^*(\sh{M}_{X/S})\longrightarrow\sh{M}_{X'/S}.
  \tag{20.6.8.1}
\]

\begin{env}[20.6.9]
\label{IV.20.6.9}
Consider finally a morphism $S'\to S$, and set $X'=X\times_S S'$, which is flat and locally of finite presentation over $S'$;
let $p:X'\to X$ be the canonical projection.
Let $U$ be an open of $X$, $t$ a section belonging to $T_{X/S}(U)$, $t'$ its image in $\Gamma(p^{-1}(U),\sh{O}_{X'})$;
for every $s'\in S'$, if $s$ is the image of $s'$, we have $X'_{s'}=X_s\otimes_{k(s)}k(s')$, hence the morphism $X'_{s'}\to X_s$ is flat, and by virtue of \sref{IV.20.1.12} the inverse image $t'_{s'}$ of $t_s$ is regular;
this proves that $t'\in T_{X'/S'}(p^{-1}(U))$.
We deduce from this, as in \sref{IV.20.1.12}, a canonical homomorphism of $\sh{O}_{X'}$-Algebras
\[
\label{IV.20.6.9.1}
  p^*(\sh{M}_{X/S})\to\sh{M}_{X'/S'}.
  \tag{20.6.9.1}
\]

Using the identification permitted by \sref{IV.20.6.3}, this notion of base change for relative meromorphic functions is a particular case of the analogous notion for relative pseudo-morphisms \sref{IV.20.5.8}.
\end{env}
% END NATIVE INPUT source_aligned/ega4-20
% BEGIN NATIVE INPUT source_aligned/ega4-21
\section{Divisors}
\label{section:IV.21}

For the contents of the present section, see the comments in the introduction to §20.
For the \emph{global} properties of divisors, the reader is referred to the section devoted to them in Chapter~V.

\subsection{Divisors on a ringed space}
\label{subsection:IV.21.1}
\label{IV.21.1}

\begin{env}[21.1.1]
\label{IV.21.1.1}
Let $(X,\sh{O}_X)$ be a ringed space.
Let $\sh{M}_X$ be the sheaf of germs of meromorphic functions on $X$ \sref{IV.20.1.3}, and $\sh{M}_X^*$ the sheaf (of multiplicative groups) of germs of regular meromorphic functions on $X$ \sref{IV.20.1.8}.
It is clear that the sheaf $\sh{O}_X^*$ (of multiplicative groups) of germs of invertible sections of $\sh{O}_X$ is canonically identified with a subsheaf (of multiplicative groups) of $\sh{M}_X^*$.
\end{env}

\begin{definition}[21.1.2]
\label{IV.21.1.2}
We call the quotient sheaf (of commutative groups) $\sh{M}_X^*/\sh{O}_X^*$ the \emph{sheaf of divisors on $X$}, and denote it by $\sh{D}iv_X$.
Its sections over $X$ are called the \emph{divisors on $X$}; they form a commutative group denoted $\operatorname{Div}(X)$.
For every section $f$ of $\sh{M}_X^*$ over $X$ (that is, every regular meromorphic function on $X$ \sref{IV.20.1.8}), the \emph{divisor of $f$}, denoted $\operatorname{div}(f)$ (or $\operatorname{div}_X(f)$), is the divisor on $X$ which is the image of $f$ under the canonical homomorphism
\[
  \Gamma(X,\sh{M}_X^*)\longrightarrow\Gamma(X,\sh{D}iv_X).
\]
\end{definition}

\oldpage[IV-4]{256}

The \emph{support} of a divisor $D$ is the closed set of $x\in X$ such that $D_x\neq 0$.
We denote it $\operatorname{Supp}(D)$.

For every open $U$ of $X$, we evidently have $\sh{M}_X^*|U=\sh{M}_U^*$, $\sh{O}_X^*|U=\sh{O}_U^*$, hence
$\sh{D}iv_X|U=\sh{D}iv_U$, and consequently the sheaf $\sh{D}iv_X$ is equal to the presheaf $U\mapsto\operatorname{Div}(U)$.

When $X=\operatorname{Spec}(A)$ is affine, we write $\operatorname{Div}(A)$ instead of $\operatorname{Div}(\operatorname{Spec}(A))$.

\begin{env}[21.1.3]
\label{IV.21.1.3}
We will always denote \emph{additively} the group $\operatorname{Div}(X)$ of divisors on $X$.
For two regular meromorphic functions $f$, $g$ on $X$, we thus have
\[
\label{IV.21.1.3.1}
  \operatorname{div}(fg)=\operatorname{div}(f)+\operatorname{div}(g),
  \tag{21.1.3.1}
\]
\[
\label{IV.21.1.3.2}
  \operatorname{div}(f^{-1})=-\operatorname{div}(f).
  \tag{21.1.3.2}
\]
By definition, for every regular meromorphic function $f$ on $X$, we have the equivalence
\[
\label{IV.21.1.3.3}
  \operatorname{div}(f)=0 \iff f\in\Gamma(X,\sh{O}_X^*),
  \tag{21.1.3.3}
\]
whence, for two regular meromorphic functions $f$, $g$ on $X$, we have
\[
\label{IV.21.1.3.4}
  \operatorname{div}(f)=\operatorname{div}(g) \iff fg^{-1}\in\Gamma(X,\sh{O}_X^*).
  \tag{21.1.3.4}
\]
\end{env}

\begin{env}[21.1.4]
\label{IV.21.1.4}
Now let $\sh{L}$ be an invertible $\sh{O}_X$-Module, and let $s$ be a \emph{regular meromorphic section} \sref{IV.20.1.8} of $\sh{L}$ over $X$.
Every $x\in X$ admits a neighbourhood $U$ such that $\sh{L}|U$ is isomorphic to $\sh{O}_U$, hence $\sh{M}_X(\sh{L})|U$ is isomorphic to $\sh{M}_X|U$;
by one of these isomorphisms, $s|U$ corresponds to a section $f\in\Gamma(U,\sh{M}_X^*)$, and since two isomorphisms differ only by multiplication by an element of $\Gamma(U,\sh{O}_X^*)$ (\hyperref[0.5.4.3]{$\mathbf{0_I}$, 5.4.3}), the element $\operatorname{div}_U(f)$ of $\Gamma(U,\sh{D}iv_X)$ is \emph{independent} of the isomorphism chosen;
it is clear that these elements (for $U$ variable) are the restrictions of a section of $\sh{D}iv_X$ over $X$, which we call the \emph{divisor of $s$} and denote $\operatorname{div}(s)$ (such a divisor is not necessarily of the form $\operatorname{div}(g)$ for a regular meromorphic function $g$ on $X$;
see \sref{IV.21.2.9}).
For $\sh{L}=\sh{O}_X$, the definition of $\operatorname{div}(s)$ coincides with that of \sref{IV.21.1.2}.
If $\sh{L}$, $\sh{L}'$ are two invertible $\sh{O}_X$-Modules, $s$ (resp.\ $s'$) a regular meromorphic section of $\sh{L}$ (resp.\ $\sh{L}'$) over $X$, we immediately have
\[
\label{IV.21.1.4.1}
  \operatorname{div}(s\otimes s')=\operatorname{div}(s)+\operatorname{div}(s')
  \tag{21.1.4.1}
\]
\[
\label{IV.21.1.4.2}
  \operatorname{div}(s^{\otimes n})=n.\operatorname{div}(s) \quad\text{for all}\quad n\in\mathbf{Z}
  \tag{21.1.4.2}
\]
($s^{-1}$ being the regular meromorphic section of $\sh{L}^{-1}$ over $X$ defined in \sref{IV.20.1.10})
and, for two regular meromorphic sections $s$, $s'$ of $\sh{L}$ over $X$, we have the relation
\[
\label{IV.21.1.4.3}
  \operatorname{div}(s)=\operatorname{div}(s') \iff s'=ts \quad\text{with}\quad t\in\Gamma(X,\sh{O}_X^*).
  \tag{21.1.4.3}
\]
\end{env}

\begin{env}[21.1.5]
\label{IV.21.1.5}
The sheaf $\sh{S}(\sh{O}_X)$ \sref{IV.20.1.3} whose sections over an open $U$ of $X$ are the regular elements of $\Gamma(U,\sh{O}_X)$ is a \emph{sub-sheaf of monoids} of $\sh{M}_X^*$;
we can write
\[
\label{IV.21.1.5.1}
  \sh{S}(\sh{O}_X)=\sh{O}_X\cap\sh{M}_X^*.
  \tag{21.1.5.1}
\]
\end{env}

\oldpage[IV-4]{257}

If we denote by $\sh{S}(\sh{O}_X)^{-1}$ the sheaf whose sections over $U$ are the inverses in $\Gamma(U,\sh{M}_X^*)$ of the elements of $\Gamma(U,\sh{S}(\sh{O}_X))$, it is clear that $\Gamma(U,\sh{S}(\sh{O}_X))\cap\Gamma(U,\sh{S}(\sh{O}_X)^{-1})=\Gamma(U,\sh{O}_X^*)$, hence
\[
\label{IV.21.1.5.2}
  \sh{S}(\sh{O}_X)\cap\sh{S}(\sh{O}_X)^{-1}=\sh{O}_X^*.
  \tag{21.1.5.2}
\]

\begin{definition}[21.1.6]
\label{IV.21.1.6}
The sub-sheaf of sets of $\sh{D}iv_X$, the canonical image of the sub-sheaf $\sh{S}(\sh{O}_X)$ of $\sh{M}_X^*$, is denoted $\sh{D}iv_X^+$;
its sections over $X$ are called \emph{positive divisors} and their set is denoted $\operatorname{Div}^+(X)$.
\end{definition}

Since $\sh{S}(\sh{O}_X)$ is a sheaf of monoids (multiplicative), we have
\[
\label{IV.21.1.6.1}
  \operatorname{Div}^+(X)+\operatorname{Div}^+(X)\subset\operatorname{Div}^+(X)
  \tag{21.1.6.1}
\]
and on the other hand, by virtue of \sref{IV.21.1.5.2} and \sref{IV.21.1.3.3}
\[
\label{IV.21.1.6.2}
  \operatorname{Div}^+(X)\cap(-\operatorname{Div}^+(X))=\{0\}.
  \tag{21.1.6.2}
\]
These two relations show that $\operatorname{Div}^+(X)$ is the set of positive elements for an \emph{order structure} on the group $\operatorname{Div}(X)$, compatible with this group structure;
we denote this order relation by $D\leq D'$, in other words, we have
\[
\label{IV.21.1.6.3}
  D\geq 0 \iff D\in\operatorname{Div}^+(X).
  \tag{21.1.6.3}
\]

We will always suppose in what follows that $\operatorname{Div}(X)$ is equipped with this order structure.

It is clear that $\sh{D}iv_X^+|U=\sh{D}iv_U^+$ for every open $U$ of $X$, and hence that
$\Gamma(U,\sh{D}iv_X^+)=\operatorname{Div}^+(U)$;
thus $\sh{D}iv_X^+$ defines on $\sh{D}iv_X$ a structure of sheaf of ordered groups.
The stalk $(\sh{D}iv_X^+)_x$ is a submonoid of $(\sh{D}iv_X)_x$, consisting of its positive elements.
Consequently, for a divisor $D$ on $X$, to say that $D\geq0$ is equivalent to saying that $D_x\geq0$ for every $x\in X$.

By definition, for every regular meromorphic function $f$ on $X$, we have the relation
\[
\label{IV.21.1.6.4}
  \operatorname{div}(f)\geq 0 \iff f\in\Gamma(X,\sh{S}(\sh{O}_X))
  \tag{21.1.6.4}
\]
in other words, $\operatorname{div}(f)\geq 0$ means that $f$ is a \emph{regular section} of $\sh{O}_X$, or equivalently a section of $\sh{O}_X$ \emph{invertible in} $M(X)$.

More generally, given a divisor $D$ on $X$, the relation $\operatorname{div}(f)\geq D$ is equivalent to the following: for every open $U$ of $X$ such that $D|U=\operatorname{div}(g)$, where $g\in\Gamma(U,\sh{M}_X^*)$, there exists a regular element $t$ of $\Gamma(U,\sh{O}_X)$ such that $f|U=tg$.

\begin{env}[21.1.7]
\label{IV.21.1.7}
Let $\sh{L}$ be an invertible $\sh{O}_X$-Module, $s$ a regular meromorphic section of $\sh{L}$ over $X$;
we have the relation
\[
\label{IV.21.1.7.1}
  \operatorname{div}(s)\geq 0 \iff s\in\Gamma(X,\sh{L})\cap\Gamma(X,(\sh{M}_X(\sh{L}))^*)
  \tag{21.1.7.1}
\]
as follows immediately from the definitions \sref{IV.21.1.4} and \sref{IV.21.1.6}.
\end{env}

\begin{proposition}[21.1.8]
\label{IV.21.1.8}
Let $X$ be a locally Noetherian prescheme, $D$ a divisor on $X$.
Suppose that for every $x\in X$ such that $\operatorname{prof}(\sh{O}_{X,x})=1$ we have $D_x\geq 0$ (resp.\ $D_x=0$).
Then $D\geq 0$ (resp.\ $D=0$).
\end{proposition}

\begin{proof}
The question being local on $X$, we can suppose that $D=\operatorname{div}(f)$, $f$ being a
\oldpage[IV-4]{258}
regular meromorphic function on $X$;
the relation $D_x\geq 0$ is equivalent to $x\in\operatorname{dom}(f)$, so the hypothesis means that if $T=X-\operatorname{dom}(f)$, then $\operatorname{prof}(\sh{O}_{X,x})\geq 2$ for every $x\in T$ (since $\operatorname{dom}(f)$ contains the maximal points of $X$).
By \sref{IV.5.10.5}, the restriction homomorphism $\Gamma(X,\sh{O}_X)\to\Gamma(X-T,\sh{O}_X)$ is \emph{bijective}, which shows that there exists a section $s$ of $\sh{O}_X$ over $X$ such that $f=s|(X-T)$.
But by definition of $T$, this implies $T=\varnothing$, hence $f=s$ and $D\geq 0$.
The assertion concerning the relation $D_x=0$ follows immediately by applying what precedes to $-D$, by virtue of \sref{IV.21.1.6.2}.
\end{proof}

\begin{corollary}[21.1.9]
\label{IV.21.1.9}
Let $X$ be a locally Noetherian prescheme, $D$ a divisor on $X$.
Let $S$ be the support of $D$.
Then, for every maximal point $x$ of $S$, $\operatorname{prof}(\sh{O}_{X,x})=1$.
\end{corollary}

\begin{proof}
Set $X_1=\operatorname{Spec}(\sh{O}_{X,x})$;
taking into account \sref{IV.20.2.11} and \sref{IV.20.3.6}, the sheaf $\sh{M}_{X_1}^*$ is induced on $X_1$ by $\sh{M}_X^*$, hence we can restrict to the case where $X=X_1$, in which case $x\in S$ and $x$ is the maximal point of $S$, so necessarily $S=\{x\}$.
If $\operatorname{prof}(\sh{O}_{X,x})\neq 1$, we would conclude, by virtue of \sref{IV.21.1.8}, that $D=0$, which contradicts the definition of $S$.
\end{proof}

\begin{proposition}[21.1.10]
\label{IV.21.1.10}
Let $A$ be a local Noetherian ring;
in order that $\operatorname{Div}(A)=0$, it is necessary and sufficient that $\operatorname{prof}(A)=0$, in other words, that the maximal ideal $\mathfrak{m}$ of $A$ be associated to $A$
(\hyperref[0.16.4.6]{$\mathbf{0}$, 16.4.6}).
\end{proposition}

\begin{proof}
Indeed, to say that $\operatorname{Div}(A)=0$ means that in $A$ every regular element is invertible, or equivalently that all the elements of $\mathfrak{m}$ are divisors of zero, which means that $\mathfrak{m}\in\operatorname{Ass}(A)$ (Bourbaki, \emph{Alg.\ comm.}, chap.~IV, \textsection~1, n\textsuperscript{o}~1, cor.~3 of prop.~2).
\end{proof}

\subsection{Divisors and invertible fractionary ideals}
\label{subsection:IV.21.2}

\begin{env}[21.2.1]
\label{IV.21.2.1}
Let $(X,\sh{O}_X)$ be a ringed space.
We call \emph{fractionary Ideal} on $X$ a sub-$\sh{O}_X$-Module of the $\sh{O}_X$-Module $\sh{M}_X$ of germs of meromorphic functions on $X$.
A fractionary Ideal $\sh{J}$ on $X$ which is an invertible $\sh{O}_X$-Module is called an \emph{invertible fractionary Ideal}.
\end{env}

\begin{proposition}[21.2.2]
\label{IV.21.2.2}
For a fractionary Ideal $\sh{J}$ on $X$ to be invertible, it is necessary and sufficient that for every $x\in X$, there exists a neighbourhood $U$ of $x$ and a section $f\in\Gamma(U,\sh{M}_X^*)$ such that $\sh{J}|U=\sh{O}_U.f$.
\end{proposition}

\begin{proof}
The condition is evidently sufficient, since the map $s\mapsto s(f|V)$ of $\Gamma(V,\sh{O}_X)$ is evidently bijective for every open $V\subset U$.
To see that it is necessary, let us note that there exists by hypothesis a neighbourhood $U$ of $x$ and an isomorphism of $\sh{O}_X$-Modules $\sh{O}_U\xrightarrow{\sim}\sh{J}|U$.
If $f$ is the image of the section $1\in\Gamma(U,\sh{O}_X)$ by this isomorphism, we may suppose, after restricting $U$, that $f=u/s$, where $u\in\Gamma(U,\sh{O}_X)$ and $s\in\Gamma(U,\sh{S}(\sh{O}_X))$, and the isomorphism considered then associates, to every section $v\in\Gamma(V,\sh{O}_X)$ (where $V$ is an open contained in $U$), the section $v(u|V)/(s|V)$;
to say that this map is bijective means that $u|V$ is a regular element of $\Gamma(V,\sh{O}_X)$, hence $f\in\Gamma(U,\sh{M}_X^*)$.
\end{proof}

We note that for every open $U$ of $X$ such that $\sh{J}|U=\sh{O}_U.f$ with $f\in\Gamma(U,\sh{M}_X^*)$, the section $f$ is determined \emph{up to multiplication} by an element of $\Gamma(U,\sh{O}_X^*)$, since the
\oldpage[IV-4]{259}
multiplication by these elements gives all the automorphisms of the $\sh{O}_U$-Module $\sh{O}_U$ (\hyperref[0.5.4.3]{$\mathbf{0_I}$, 5.4.3}).

\begin{corollary}[21.2.3]
\label{IV.21.2.3}
\begin{enumerate}[label=(\roman*)]
  \item Let $\sh{J}$ be an invertible fractionary Ideal;
  then the invertible $\sh{O}_X$-Module $\sh{J}^{-1}$ is canonically identified with the fractionary Ideal $\sh{J}'$ (the \emph{transporter} of $\sh{J}$ into $\sh{O}_X$) defined as follows:
  for every open $U$ of $X$ such that $\sh{J}|U=\sh{O}_U.f$, where $f\in\Gamma(U,\sh{M}_X^*)$, we have $\sh{J}'|U=\sh{O}_U.f^{-1}$.
  \item If $\sh{J}_1$ and $\sh{J}_2$ are two invertible fractionary Ideals, the canonical map $\sh{J}_1\otimes\sh{J}_2\to\sh{J}_1\sh{J}_2$ is an isomorphism of $\sh{O}_X$-Modules.
\end{enumerate}
\end{corollary}

\begin{proof}
The assertion (ii) follows immediately from \sref{IV.21.2.2}.
On the other hand, the remark made at the end of \sref{IV.21.2.2} shows that there is indeed a unique fractionary Ideal $\sh{J}'$ defined by the condition of the statement;
the canonical isomorphism from $\sh{J}'$ onto $\sh{J}^{-1}=\shHom_{\sh{O}_X}(\sh{J},\sh{O}_X)$ is obtained by associating, to every section $s(f^{-1}|V)$ of $\Gamma(V,\sh{J}')$ (where $V$ is an open contained in $U$ and $s\in\Gamma(V,\sh{O}_X)$), the homomorphism $t(f|V)\mapsto st$ from $\Gamma(V,\sh{J})$ to $\Gamma(V,\sh{O}_X)$.
\end{proof}

By virtue of \sref{IV.21.2.3},~(i), we will identify the invertible fractionary ideals $\sh{J}'$ and $\sh{J}^{-1}$ considered as sub-$\sh{O}_X$-Modules of $\sh{M}_X$.

\begin{env}[21.2.4]
\label{IV.21.2.4}
It follows from \sref{IV.21.2.3} that the set $\operatorname{Id.inv}(X)$ of invertible fractionary Ideals on $X$ is equipped with a structure of \emph{commutative group} for the law of composition $(\sh{J}_1,\sh{J}_2)\mapsto\sh{J}_1\sh{J}_2$, the neutral element of this group being $\sh{O}_X$.
It is clear that for every open $U$ of $X$, we have $(\sh{J}_1\sh{J}_2)|U=(\sh{J}_1|U)(\sh{J}_2|U)$, hence the restriction map $\sh{J}\mapsto\sh{J}|U$ is a homomorphism of groups $\operatorname{Id.inv}(X)\to\operatorname{Id.inv}(U)$;
we thus define a \emph{presheaf of commutative groups} $U\mapsto\operatorname{Id.inv}(U)$;
it is immediate that in fact this presheaf is a \emph{sheaf of commutative groups}, which we denote $\sh{J}d.inv_X$.
\end{env}

\begin{env}[21.2.5]
\label{IV.21.2.5}
For every regular meromorphic function $f\in\Gamma(X,\sh{M}_X^*)$, it follows from \sref{IV.21.2.2} that $\sh{J}(f)=\sh{O}_X.f$ is an invertible fractionary Ideal, and we evidently have $\sh{J}(f_1f_2)=\sh{J}(f_1)\sh{J}(f_2)$, in other words the map $f\mapsto\sh{J}(f)$ is a homomorphism of the commutative group $\Gamma(X,\sh{M}_X^*)$ to the commutative group $\operatorname{Id.inv}(X)$.
Replacing $X$ by an arbitrary open $U$ of $X$ and noting that the homomorphisms thus obtained are compatible with the restriction operations, we obtain a canonical homomorphism of sheaves of commutative groups:
\[
\label{IV.21.2.5.1}
  I_0:\sh{M}_X^*\to\sh{J}d.inv_X.
  \tag{21.2.5.1}
\]
Let us note that if $f\in\Gamma(X,\sh{O}_X^*)$, then $\sh{J}(f)=\sh{O}_X$;
we deduce immediately that the homomorphism $I_0$ factorises as
\[
\label{IV.21.2.5.2}
  \sh{M}_X^*\to\sh{M}_X^*/\sh{O}_X^*=\sh{D}iv_X\xrightarrow{I}\sh{J}d.inv_X
  \tag{21.2.5.2}
\]
where $I$ is a homomorphism of the sheaf of additive groups $\sh{D}iv_X$ to the sheaf of multiplicative groups $\sh{J}d.inv_X$;
we thus have for every open $U$ of $X$ a homomorphism $\sh{J}_U:\operatorname{Div}(U)\to\operatorname{Id.inv}(U)$ of commutative groups, such that, for every section $f\in\Gamma(U,\sh{M}_X^*)$, we have
\[
\label{IV.21.2.5.3}
  \sh{J}_U(\operatorname{div}_U(f))=\sh{O}_U.f.
  \tag{21.2.5.3}
\]
\oldpage[IV-4]{260}

We conclude that, for every divisor $D\in\operatorname{Div}(X)$, $\sh{J}_X(D)$ is the invertible fractionary ideal defined as follows:
for every open $U$ of $X$ such that $D|U=\operatorname{div}_U(f)$, where $f\in\Gamma(U,\sh{M}_U^*)$, one has $\sh{J}_X(D)|U=\sh{O}_U.f$.
Hence, by virtue of \sref{IV.21.1.6}, for every regular meromorphic function $f\in\Gamma(X,\sh{M}_X^*)$, we have the relation
\[
\label{IV.21.2.5.4}
  f\in\Gamma(X,\sh{J}_X(D)) \iff \operatorname{div}(f)\geq D.
  \tag{21.2.5.4}
\]
\end{env}

\begin{proposition}[21.2.6]
\label{IV.21.2.6}
The homomorphism $I:\sh{D}iv_X\to\sh{J}d.inv_X$ is bijective.
\end{proposition}

\begin{proof}
We define indeed a homomorphism $I_X'$ of $\operatorname{Id.inv}(X)$ to $\operatorname{Div}(X)$ by making correspond to every invertible fractionary Ideal $\sh{J}$ on $X$ the divisor $I_X'(\sh{J})$ defined as follows:
for every open $U$ of $X$ such that $\sh{J}|U=\sh{O}_U.f$, where $f\in\Gamma(U,\sh{M}_X^*)$ \sref{IV.21.2.2}, we take $I_X'(\sh{J})|U=\operatorname{div}_U(f)$;
by virtue of the remark following \sref{IV.21.2.2}, this definition is independent of the generator $f$ chosen for $\sh{J}|U$, and determines a divisor on $X$.
Moreover, this definition immediately shows that the homomorphisms $\sh{J}_X$ and $I_X'$ are inverses of one another.
Replacing $X$ by an arbitrary open $U$, we deduce the definition of the isomorphism $I':\sh{J}d.inv_X\to\sh{D}iv_X$, inverse of $I$, whence the proposition.
We set $I_X'(\sh{J})=\operatorname{div}(\sh{J})$, so that for every regular meromorphic function $f$ on $X$
\[
\label{IV.21.2.6.1}
  \operatorname{div}(\sh{O}_X.f)=\operatorname{div}(f).
  \tag{21.2.6.1}
\]
\end{proof}

\begin{env}[21.2.7]
\label{IV.21.2.7}
We will often identify the sheaves $\sh{D}iv_X$ and $\sh{J}d.inv_X$ (resp.\ the groups $\operatorname{Div}(X)$ and $\operatorname{Id.inv}(X)$) by means of the isomorphisms $I$ and $I'$ (resp.\ $\sh{J}_X$ and $I_X'$) above.
We note that we have the relation
\[
\label{IV.21.2.7.1}
  D\geq 0 \iff \sh{J}_X(D)\subset\sh{O}_X \quad\text{for}\quad D\in\operatorname{Div}(X)
  \tag{21.2.7.1}
\]
as follows immediately from the definitions \sref{IV.21.1.6} and \sref{IV.21.1.5.1};
in other terms, the image $\sh{J}_X(\operatorname{Div}^+(X))$ is the set of \emph{Ideals} of $\sh{O}_X$ (also sometimes called \emph{integral Ideals}) which are invertible $\sh{O}_X$-Modules:
such an Ideal $\sh{J}$ is further characterised by the fact that for every $x\in X$, there is a neighbourhood $U$ of $x$ in $X$ such that $\sh{J}|U=\sh{O}_U.f$, where $f$ is a regular element of $\Gamma(U,\sh{O}_X)$.
The set $\sh{J}_X(\operatorname{Div}^+(X))$ of these Ideals is thus a sub-monoid of $\operatorname{Id.inv}(X)$, equal to the set of positive elements for an order relation on $\operatorname{Id.inv}(X)$, and it is immediate that this relation is none other than the relation \emph{opposite} to the inclusion;
in other words, we have
\[
\label{IV.21.2.7.2}
  \sh{J}_1\subset\sh{J}_2 \iff \operatorname{div}(\sh{J}_1)\geq\operatorname{div}(\sh{J}_2)
  \tag{21.2.7.2}
\]
for $\sh{J}_1$, $\sh{J}_2$ in $\operatorname{Id.inv}(X)$.
\end{env}

\begin{env}[21.2.8]
\label{IV.21.2.8}
For every divisor $D$ on $X$, we set
\[
\label{IV.21.2.8.1}
  \sh{O}_X(D)=(\sh{J}_X(D))^{-1};
  \tag{21.2.8.1}
\]
$\sh{O}_X(D)$ is thus an invertible fractionary Ideal, defined as follows:
for every open $U$ of $X$ such that $D|U=\operatorname{div}_U(f)$, where $f\in\Gamma(U,\sh{M}_X^*)$, $\sh{O}_X(D)|U$ is the invertible fractionary Ideal
\oldpage[IV-4]{261}
$\sh{O}_U.f^{-1}$;
by virtue of \sref{IV.21.1.6}, for every regular meromorphic function $f$, we have the relation
\[
\label{IV.21.2.8.2}
  f\in\Gamma(X,\sh{O}_X(D)) \iff \operatorname{div}(f)\geq -D.
  \tag{21.2.8.2}
\]
Moreover, it is clear that we have canonical isomorphisms \sref{IV.21.2.3}
\[
\label{IV.21.2.8.3}
  \left\{\begin{aligned}
    &\sh{O}_X(0)=\sh{O}_X,\quad \sh{O}_X(D+D')=\sh{O}_X(D)\sh{O}_X(D')\xrightarrow{\sim}\sh{O}_X(D)\otimes\sh{O}_X(D')\\
    &\sh{O}_X(nD)=(\sh{O}_X(D))^n\xrightarrow{\sim}(\sh{O}_X(D))^{\otimes n}
  \end{aligned}\right.
  \tag{21.2.8.3}
\]
for every integer $n\in\mathbf{Z}$, and for any two divisors $D$, $D'$ on $X$.
\end{env}

\begin{env}[21.2.9]
\label{IV.21.2.9}
Let $\sh{J}$ be an invertible fractionary Ideal on $X$.
The canonical injection $\sh{J}\hookrightarrow\sh{M}_X$ defines by tensorisation a homomorphism of $\sh{O}_X$-Modules
\[
\label{IV.21.2.9.1}
  \sh{M}_X(\sh{J})=\sh{J}\otimes_{\sh{O}_X}\sh{M}_X\to\sh{M}_X\otimes_{\sh{O}_X}\sh{M}_X=\sh{M}_X
  \tag{21.2.9.1}
\]
which is an \emph{isomorphism}:
indeed, if $U$ is an open of $X$ such that $\sh{J}|U=\sh{O}_U.f$, where $f\in\Gamma(U,\sh{M}_X^*)$, the homomorphism \sref{IV.21.2.9.1} restricted to $U$ is none other than the isomorphism which, for every section $t$ of $\sh{M}_X(\sh{J})|U=\sh{M}_X|U$ over $V\subset U$, makes correspond the section $t(f|V)^{-1}$ of the same sheaf.
In the isomorphism \sref{IV.21.2.9.1}, the regular meromorphic sections of $\sh{J}$ over $X$ correspond to the regular meromorphic functions on $X$.

Consider in particular the case where $\sh{J}=\sh{O}_X(D)$, where $D$ is a divisor on $X$;
we then have a canonical isomorphism
\[
\label{IV.21.2.9.2}
  \sh{M}_X(\sh{O}_X(D))\xrightarrow{\sim}\sh{M}_X
  \tag{21.2.9.2}
\]
and we denote by $s_D$ the regular meromorphic section of $\sh{O}_X(D)$ over $X$ which corresponds by this isomorphism to the section $1$ of $\sh{M}_X$.
If $U$ is an open of $X$ such that $\sh{O}_X(D)|U=\sh{O}_U.f^{-1}$, where $f\in\Gamma(U,\sh{M}_X^*)$, then $s_D|U=1$ in $\Gamma(U,\sh{M}_X)$;
as $D|U=\operatorname{div}_U(f)$, we deduce from \sref{IV.21.1.4} that
\[
\label{IV.21.2.9.3}
  \operatorname{div}(s_D)=D.
  \tag{21.2.9.3}
\]
On the other hand, we also deduce from the canonical isomorphisms \sref{IV.21.2.8.3} the formulas
\[
\label{IV.21.2.9.4}
  s_0=1,\quad s_{D+D'}=s_D\otimes s_{D'},\quad s_{nD}=s_D^{\otimes n}\quad (n\in\mathbf{Z}).
  \tag{21.2.9.4}
\]
\end{env}

\begin{env}[21.2.10]
\label{IV.21.2.10}
Consider, between two pairs $(\sh{L},s)$, $(\sh{L}',s')$, where $\sh{L}$ and $\sh{L}'$ are two invertible $\sh{O}_X$-Modules, $s$ (resp.\ $s'$) a regular meromorphic section of $\sh{L}$ (resp.\ $\sh{L}'$) over $X$, the relation:
``there exists an isomorphism $u:\sh{L}\xrightarrow{\sim}\sh{L}'$ such that $\overline{u}(s)=s'$'',
where $\overline{u}:\Gamma(X,\sh{M}_X(\sh{L}))\xrightarrow{\sim}\Gamma(X,\sh{M}_X(\sh{L}'))$ is the isomorphism derived from $u$ (we note that the isomorphism $u$ satisfying this condition is then determined in a unique way).
It is clear that this is an equivalence relation, and since there exists a set of invertible $\sh{O}_X$-Modules such that every invertible $\sh{O}_X$-Module is isomorphic to an element of this set (\hyperref[0.5.4.7]{$\mathbf{0_I}$, 5.4.7}), we can speak of the \emph{set} $D(X)$ of \emph{equivalence classes} of pairs $(\sh{L},s)$ for the above relation.
For every invertible $\sh{O}_X$-Module $\sh{L}$ and every regular meromorphic section
\oldpage[IV-4]{262}
$s$ of $\sh{L}$ over $X$, we denote by $\operatorname{cl}(\sh{L},s)$ the element of $D(X)$ corresponding to the pair $(\sh{L},s)$.
It follows from (\hyperref[0.5.4.3]{$\mathbf{0_I}$, 5.4.3}) that if $s$, $s'$ are two regular meromorphic sections of $\sh{L}$ over $X$, the relation $\operatorname{cl}(\sh{L},s)=\operatorname{cl}(\sh{L},s')$ is equivalent to the existence of a section $t\in\Gamma(X,\sh{O}_X^*)$ such that $s'=ts$.

It is immediate that if $(\sh{L},s)$ is equivalent to $(\sh{L}_1,s_1)$ and $(\sh{L}',s')$ to $(\sh{L}_1',s_1')$, the pairs $(\sh{L}\otimes\sh{L}',s\otimes s')$ and $(\sh{L}_1\otimes\sh{L}_1',s_1\otimes s_1')$ are equivalent;
we thus define in $D(X)$ a law of composition by setting
\[
\label{IV.21.2.10.1}
  \operatorname{cl}(\sh{L},s)\operatorname{cl}(\sh{L}',s')=\operatorname{cl}(\sh{L}\otimes\sh{L}',s\otimes s');
  \tag{21.2.10.1}
\]
it is immediate that this is the law of a \emph{commutative group}, whose neutral element is $\operatorname{cl}(\sh{O}_X,1)$ and where the inverse of $\operatorname{cl}(\sh{L},s)$ is $\operatorname{cl}(\sh{L}^{-1},s^{\otimes(-1)})$.
\end{env}

\begin{proposition}[21.2.11]
\label{IV.21.2.11}
The maps
\[
\label{IV.21.2.11.1}
  D\mapsto\operatorname{cl}(\sh{O}_X(D),s_D),\qquad \operatorname{cl}(\sh{L},s)\mapsto\operatorname{div}(s)
  \tag{21.2.11.1}
\]
are inverse isomorphisms of $\operatorname{Div}(X)$ to $D(X)$ and of $D(X)$ to $\operatorname{Div}(X)$ respectively.
\end{proposition}

\begin{proof}
Taking into account \sref{IV.21.2.8.3}, \sref{IV.21.2.9.4}, and \sref{IV.21.2.10.1}, it suffices to verify that the composites of these two maps are the identities in $\operatorname{Div}(X)$ and $D(X)$ respectively.
The first assertion is none other than \sref{IV.21.2.9.3}.
On the other hand, let $D=\operatorname{div}(s)$, where $s$ is a regular meromorphic section of $\sh{L}$ over $X$, and let $U$ be an open of $X$ such that there exists an isomorphism of $\sh{L}|U$ onto $\sh{O}_U$, transforming $s|U$ to $f\in\Gamma(U,\sh{M}_X^*)$, so that $D|U=\operatorname{div}_U(f)$, $\sh{O}_X(D)|U=\sh{O}_U.f^{-1}$ and $s_D|U$ is the unit element of $\Gamma(U,\sh{M}_X)$.
There is thus an isomorphism $v_U:\sh{L}|U\to\sh{O}_X(D)|U$ such that $\overline{v}_U$ (notation of \sref{IV.21.2.10}) transforms $s|U$ to $f.f^{-1}=1$;
one sees immediately that these isomorphisms are compatible with the restriction operations, and thus define an isomorphism $v:\sh{L}\to\sh{O}_X(D)$ such that $\overline{v}(s)=s_D$.
\end{proof}

We can \emph{transport} by the first of the isomorphisms \sref{IV.21.2.11.1} the structure of ordered group of $\operatorname{Div}(X)$ to $D(X)$;
the elements $\geq 0$ of $D(X)$ are thus the classes $\operatorname{cl}(\sh{L},s)$ such that $\operatorname{div}(s)\geq 0$, that is to say \sref{IV.21.1.7.1} such that
\[
  s\in\Gamma(X,\sh{L})\cap\Gamma(X,(\sh{M}_X(\sh{L}))^*).
\]

\begin{env}[21.2.12]
\label{IV.21.2.12}
Let $D$ be a \emph{positive} divisor on a \emph{prescheme} $X$;
the fractionary Ideal $\sh{J}_X(D)$ is an Ideal of $\sh{O}_X$, which is an invertible $\sh{O}_X$-Module;
let $Y(D)$ be the closed sub-prescheme of $X$ that it defines.
For every $x\in Y(D)$, there is by hypothesis a neighbourhood $U$ of $x$ in $X$ and a regular section $t\in\Gamma(U,\sh{O}_X)$ such that $\sh{J}_X(D)|U=\sh{O}_U.t$ \sref{IV.21.2.7};
in other terms, the canonical immersion $Y(D)\to X$ is \emph{regular and of codimension $1$ at every point of $Y(D)$} \sref{IV.19.1.4}.
Conversely, if $Y$ is a closed sub-prescheme of $X$, regularly immersed in $X$ and of codimension $1$ at every point of $Y$, there exists \emph{a unique positive divisor $D$} such that $Y(D)=Y$, since for every $x\in Y$ there is a neighbourhood $U$ of $x$ in $X$ such that $Y\cap U$ is defined by an Ideal of $\sh{O}_U$ of the form $\sh{O}_U.t$, where $t$ is a regular element of $\Gamma(U,\sh{O}_X)$.
We note that $\operatorname{Supp}(D)=Y(D)$, since to say that $D_x\neq 0$ means that $t_x$ (with the above notation) is not invertible, that is to say $x\in Y(D)$.
\end{env}

\oldpage[IV-4]{263}

\subsection{Linear equivalence of divisors}
\label{subsection:IV.21.3}

\begin{env}[21.3.1]
\label{IV.21.3.1}
We say that a divisor $D$ on $X$ is \emph{principal} if it is of the form $\operatorname{div}(f)$, where $f$ is a regular meromorphic function on $X$;
the meromorphic functions $f'$ such that $\operatorname{div}(f')=D$ are then all of the form $tf$, where $t\in\Gamma(X,\sh{O}_X^*)$ \sref{IV.21.1.3.4}.
The set of principal divisors is a subgroup of $\operatorname{Div}(X)$, denoted $\operatorname{Div.princ}(X)$, isomorphic to $\Gamma(X,\sh{M}_X^*)/\Gamma(X,\sh{O}_X^*)$.
Two divisors $D$, $D'$ are said to be \emph{linearly equivalent} if $D-D'$ is a principal divisor;
the principal divisors are thus the divisors linearly equivalent to $0$.
\end{env}

\begin{env}[21.3.2]
\label{IV.21.3.2}
Recall (\hyperref[0.5.4.7]{$\mathbf{0_I}$, 5.4.7}) that we can speak of the set of equivalence classes of invertible $\sh{O}_X$-Modules for the isomorphism relation;
we denote this set by $\operatorname{Pic}(X)$, and for every invertible $\sh{O}_X$-Module $\sh{L}$, we denote by $\operatorname{cl}(\sh{L})$ the equivalence class of $\sh{O}_X$-Modules isomorphic to $\sh{L}$;
moreover, $\operatorname{Pic}(X)$ is a commutative group for the multiplication defined by $\operatorname{cl}(\sh{L})\operatorname{cl}(\sh{L}')=\operatorname{cl}(\sh{L}\otimes\sh{L}')$.
It is clear that the map
\[
\label{IV.21.3.2.1}
  r:\operatorname{cl}(\sh{L},s)\mapsto\operatorname{cl}(\sh{L})
  \tag{21.3.2.1}
\]
is a homomorphism of the group $D(X)$ \sref{IV.21.2.10} to the group $\operatorname{Pic}(X)$.
By composition, we thus deduce a homomorphism
\[
\label{IV.21.3.2.2}
  l:\operatorname{Div}(X)\xrightarrow{\sim}D(X)\xrightarrow{r}\operatorname{Pic}(X)
  \tag{21.3.2.2}
\]
(also denoted $l_X$) such that, for every divisor $D$, we have
\[
\label{IV.21.3.2.3}
  l(D)=\operatorname{cl}(\sh{O}_X(D)).
  \tag{21.3.2.3}
\]

Let us note finally that, if $u:X'\to X$ is a morphism of ringed spaces, $\sh{L}_1$, $\sh{L}_2$ two invertible isomorphic $\sh{O}_X$-Modules, the invertible $\sh{O}_{X'}$-Modules $u^*(\sh{L}_1)$ and $u^*(\sh{L}_2)$ (\hyperref[0.5.4.5]{$\mathbf{0_I}$, 5.4.5}) are isomorphic;
as moreover, for any two invertible $\sh{O}_X$-Modules $\sh{L}_1$, $\sh{L}_2$, we have $u^*(\sh{L}_1\otimes\sh{L}_2)=u^*(\sh{L}_1)\otimes u^*(\sh{L}_2)$ up to canonical isomorphism (\hyperref[0.4.3.3]{$\mathbf{0_I}$, 4.3.3}), we see that the morphism $u$ canonically defines a homomorphism of commutative groups
\[
\label{IV.21.3.2.4}
  \operatorname{Pic}(u):\operatorname{Pic}(X)\to\operatorname{Pic}(X').
  \tag{21.3.2.4}
\]
\end{env}

\begin{proposition}[21.3.3]
\label{IV.21.3.3}
\begin{enumerate}[label=(\roman*)]
  \item The kernel of the canonical homomorphism $l:\operatorname{Div}(X)\to\operatorname{Pic}(X)$ is the subgroup $\operatorname{Div.princ}(X)$, in other words, for $\sh{O}_X(D)$ and $\sh{O}_X(D')$ to be isomorphic, it is necessary and sufficient that $D$ and $D'$ be linearly equivalent.
  There is thus a canonical injective homomorphism
\[
\label{IV.21.3.3.1}
  \operatorname{Div}(X)/\operatorname{Div.princ}(X)\to\operatorname{Pic}(X)
  \tag{21.3.3.1}
\]
  derived from $l$.
  \item For an invertible $\sh{O}_X$-Module $\sh{L}$ to be such that $\operatorname{cl}(\sh{L})$ is of the form $l(D)$, or equivalently that $\sh{L}$ is isomorphic to an $\sh{O}_X$-Module of the form $\sh{O}_X(D)$, it is necessary and sufficient that there exists a regular meromorphic section of $\sh{L}$.
\end{enumerate}
\end{proposition}

\begin{proof}
\oldpage[IV-4]{264}
The proposition follows immediately from the definitions and from \sref{IV.21.2.10}.
\end{proof}

\begin{proposition}[21.3.4]
\label{IV.21.3.4}
Let $X$ be a prescheme satisfying one of the two following hypotheses:
\begin{enumerate}[label=\emph{\alph*)}]
  \item $X$ is locally Noetherian and $\operatorname{Ass}(\sh{O}_X)$ is contained in an affine open of $X$.
  \item $X$ is reduced and the set of its irreducible components is locally finite.
\end{enumerate}
Then the canonical homomorphism $l:\operatorname{Div}(X)\to\operatorname{Pic}(X)$ is surjective, and gives by passage to the quotient an isomorphism
\[
  \operatorname{Div}(X)/\operatorname{Div.princ}(X)\xrightarrow{\sim}\operatorname{Pic}(X).
\]
\end{proposition}

\begin{proof}
It suffices to show that every invertible $\sh{O}_X$-Module $\sh{L}$ admits a regular meromorphic section over $X$ \sref{IV.21.3.3}.
In both cases, it will suffice, thanks to \sref{IV.20.2.11},~(ii), to define a section $s$ of $(\sh{M}_X(\sh{L}))^*$ over a \emph{schematically dense} open $U$ of $X$, or else, in case $a)$, over an open $U$ containing $\operatorname{Ass}(\sh{O}_X)$ \sref{IV.20.2.13},~(iv).
Indeed, let $V$ be an arbitrary open of $X$ such that $\sh{L}|V$ is isomorphic to $\sh{O}_V$, so that there exists an isomorphism of $\sh{M}_X(\sh{L})|V$ onto $\sh{M}_X|V$, transforming $s|(U\cap V)$ to a section $f_V$ of $\sh{M}_X^*$ over $U\cap V$.
As $U\cap V$ is schematically dense in $V$, it follows from \sref{IV.20.2.11} that there exists a unique regular meromorphic function $g_V$ such that $g_V|(U\cap V)=f_V$, and this section corresponds by the isomorphism considered to a regular meromorphic section $u_V$ of $\sh{L}$ over $V$ such that $u_V|(U\cap V)=s|(U\cap V)$.
Moreover, if $V'$ is a second open of $X$ such that $\sh{L}|V'$ is isomorphic to $\sh{O}_{V'}$, the restrictions of $u_V$ and $u_{V'}$ to $V\cap V'$ are equal, since they correspond by isomorphism to two meromorphic functions which coincide on an open schematically dense $U\cap V\cap V'$, and we conclude again by \sref{IV.20.2.11};
the $u_V$ are thus well the restrictions of a single regular meromorphic section of $\sh{L}$ over $X$.

In case $b)$, we take for each of the maximal points $x_\lambda$ of $X$ an open $U_\lambda$ containing $x_\lambda$, meeting no irreducible component of $X$ distinct from $\overline{\{x_\lambda\}}$, and such that $\sh{L}|U_\lambda$ is isomorphic to $\sh{O}_{U_\lambda}$;
we take for $s$ the section such that $s|U_\lambda$ is the section of $\sh{L}|U_\lambda$ which corresponds by the preceding isomorphism to the unit section of $\sh{O}_{U_\lambda}$.

In case $a)$, we can take $U$ to be an affine open (hence Noetherian), in other words, suppose that $X=\operatorname{Spec}(A)$, where $A$ is Noetherian, and $\sh{L}=\widetilde{P}$, where $P$ is a projective $A$-module of rank $1$.
If $S$ is the set of regular elements of $A$, we have $\Gamma(X,\sh{M}_X)=S^{-1}A$ \sref{IV.20.2.12} and $\Gamma(X,\sh{M}_X(\sh{L}))=S^{-1}P$;
now $S$ is the set of elements not belonging to any of the ideals associated to $A$, hence $S^{-1}A$ is a semi-local ring whose maximal ideals come from the maximal elements of $\operatorname{Ass}(A)$, and $S^{-1}P$ is a projective $S^{-1}A$-module of rank $1$, hence a \emph{free} $S^{-1}A$-module of rank $1$ (Bourbaki, \emph{Alg.\ comm.}, chap.~II, \textsection~5, n\textsuperscript{o}~3, prop.~5);
an element forming a base of this $S^{-1}A$-module is thus \sref{IV.20.1.8} a regular meromorphic section of $\sh{L}$ over $X$.
\end{proof}

\begin{corollary}[21.3.5]
\label{IV.21.3.5}
If $X$ is a Noetherian prescheme such that there exists an ample invertible $\sh{O}_X$-Module (\textbf{II}, 4.5.3) (for example, a quasi-projective prescheme over the spectrum of a Noetherian ring (\textbf{II}, 5.3.1 and 4.6.6)), the canonical homomorphism $l:\operatorname{Div}(X)\to\operatorname{Pic}(X)$ is surjective.
\end{corollary}

\begin{proof}
\oldpage[IV-4]{265}

Indeed (\textbf{II}, 4.5.4), there then exists an affine open neighbourhood of the finite set $\operatorname{Ass}(\sh{O}_X)$.
\end{proof}

\begin{remark}[21.3.6]
\label{IV.21.3.6}
Recall (\hyperref[0.5.4.7]{$\mathbf{0_I}$, 5.4.7}) that there is a canonical isomorphism $\pi:\mathrm{H}^1(X,\sh{O}_X^*)\xrightarrow{\sim}\operatorname{Pic}(X)$ defined as follows.
We start from a $1$-cocycle $(c_{\alpha\beta})$ with values in $\sh{O}_X^*$ corresponding to a covering $(U_\alpha)$ of $X$, $c_{\alpha\beta}$ being an element of $\Gamma(U_\alpha\cap U_\beta,\sh{O}_X^*)$, and we associate to it the class of the invertible $\sh{O}_X$-Module obtained by gluing the $\sh{O}_{U_\alpha}$ along the isomorphisms $\sh{O}_{U_\alpha}|(U_\alpha\cap U_\beta)\xrightarrow{\sim}\sh{O}_{U_\beta}|(U_\alpha\cap U_\beta)$ defined by multiplication by $c_{\alpha\beta}$.
On the other hand, we deduce from the exact sequence of sheaves of commutative groups
\[
  1\to\sh{O}_X^*\to\sh{M}_X^*\to\sh{D}iv_X\to 0
\]
the boundary homomorphism of the exact cohomology sequence
\[
\label{IV.21.3.6.1}
  \partial:\operatorname{Div}(X)\to\mathrm{H}^1(X,\sh{O}_X^*).
  \tag{21.3.6.1}
\]
Let us show that the composite homomorphism
\[
  \operatorname{Div}(X)\xrightarrow{\partial}\mathrm{H}^1(X,\sh{O}_X^*)\xrightarrow{\pi}\operatorname{Pic}(X)
\]
is the homomorphism $l$ defined in \sref{IV.21.3.2.2}.
Indeed, we must start from a covering $(U_\alpha)$ of $X$ and a divisor $D$ such that $D|U_\alpha=\operatorname{div}_{U_\alpha}(g_\alpha)$, where $g_\alpha$ is a regular meromorphic function over $U_\alpha$;
$\partial(D)$ is the cohomology class of the cocycle $(c_{\alpha\beta})$, where $c_{\alpha\beta}=g_{\alpha|\beta}g_{\beta|\alpha}^{-1}$, $g_{\alpha|\beta}$ denoting the restriction of $g_\alpha$ to $U_\alpha\cap U_\beta$.
It is clear that the image by $\pi$ of this cohomology class is the class of the invertible fractionary Ideal $\sh{L}$ such that for every $\alpha$, $\sh{L}|U_\alpha=\sh{O}_{U_\alpha}.g_\alpha^{-1}$, which is none other by definition than $\sh{O}_X(D)$ \sref{IV.21.2.8}.
\end{remark}

\subsection{Inverse images of divisors}
\label{subsection:IV.21.4}
\label{IV.21.4}

\begin{env}[21.4.1]
\label{IV.21.4.1}
Let $f:X'\to X$ be a morphism of ringed spaces;
we propose to give conditions allowing us to associate to a divisor $D$ on $X$ a divisor $D'$ on $X'$, the \emph{inverse image} of $D$ by $f$.
Let us first note that for every section $t\in\Gamma(X,\sh{O}_X^*)$, the image of $t$ by the canonical homomorphism $\Gamma(X,\sh{O}_X)\to\Gamma(X',\sh{O}_{X'})$ is still invertible, in other words belongs to $\Gamma(X',\sh{O}_{X'}^*)$.
Consider then $D$ as given by the class of equivalence of a pair $(\sh{L},s)$, where $\sh{L}$ is an invertible $\sh{O}_X$-Module and $s$ a regular meromorphic section of $\sh{L}$ over $X$ \sref{IV.21.2.11}.
Form the invertible $\sh{O}_{X'}$-Module $f^*(\sh{L})=\sh{L}'$;
to say that the inverse image $s\circ f$ of $s$ by $f$ exists \sref{IV.20.1.11} and is a regular meromorphic section of $\sh{L}'$ over $X'$ amounts to saying that the inverse images by $f$ of $s$ and $s^{\otimes(-1)}$ exist, in other words that $s\in\Gamma(X,\sh{M}_f(\sh{L}))$ and $s^{\otimes(-1)}\in\Gamma(X,\sh{M}_f(\sh{L}^{-1}))$;
the remark made above shows that if $(\sh{L}_1,s_1)$ is a pair equivalent to $(\sh{L},s)$ in the sense of \sref{IV.21.2.10}, the inverse image $s_1\circ f$ exists and is a regular meromorphic section of $\sh{L}_1'=f^*(\sh{L}_1)$ over $X'$, and the pairs $(\sh{L}',s\circ f)$ and $(\sh{L}_1',s_1\circ f)$ are equivalent.
We can therefore pose the following definition:
\end{env}

\begin{definition}[21.4.2]
\label{IV.21.4.2}
Given a morphism $f:X'\to X$ of ringed spaces, we say that the \emph{inverse image by $f$} of a divisor $D$ on $X$ exists if $s_D\in\Gamma(X,\sh{M}_f(\sh{O}_X(D)))$ and $s_{-D}\in\Gamma(X,\sh{M}_f(\sh{O}_X(-D)))$ (cf.\ \sref{IV.20.1.11}).
We call then the \emph{inverse image} of $D$ by $f$ and denote $f^*(D)$ the divisor on $X'$ which corresponds canonically \sref{IV.21.2.11} to the class of the pair $(f^*(\sh{O}_X(D)),s_D\circ f)$.
\end{definition}

It follows immediately from this definition that if $D$ and $D'$ have inverse images by $f$, so does $-D$ and $D+D'$ (taking into account \sref{IV.21.2.9.4}) and we have
\oldpage[IV-4]{266}
$f^*(D+D')=f^*(D)+f^*(D')$.
In other words, the set $\operatorname{Div}^f(X)$ of divisors on $X$ whose inverse image by $f$ exists is a \emph{subgroup} of $\operatorname{Div}(X)$, and the map $D\mapsto f^*(D)$ is an \emph{increasing homomorphism} of the ordered subgroup $\operatorname{Div}^f(X)$ to the ordered group $\operatorname{Div}(X')$, making commutative the diagram
\[
\label{IV.21.4.2.1}
  \xymatrix{
    \operatorname{Div}^f(X) \ar[r]^{l_X} \ar[d]_{f^*} & \operatorname{Pic}(X) \ar[d]^{\operatorname{Pic}(f)} \\
    \operatorname{Div}(X') \ar[r]^{l_{X'}} & \operatorname{Pic}(X')
  }
  \tag{21.4.2.1}
\]

\begin{env}[21.4.3]
\label{IV.21.4.3}
The definition \sref{IV.21.4.2} also shows immediately that, for $f^*(D)$ to exist, it is necessary and sufficient that, for every open $U$ of $X$, the inverse image of $D|U$ by $f^U:f^{-1}(U)\to U$ (the restriction of $f$) exist.
Now, if $D=\operatorname{div}(g)$, where $g$ is a regular meromorphic function on $X$, to say that the inverse image of $s_D$ by $f$ exists and is a regular meromorphic section of $f^*(\sh{O}_X(D))$ means \sref{IV.21.2.9} that the inverse image of $g$ by $f$ exists and is a regular meromorphic function on $X'$.
We deduce also immediately a second description of $\operatorname{Div}^f(X)$ and of $f^*(D)$: we consider the sub-sheaf of groups of $\sh{M}_X^*$, denoted $\sh{M}_f^*$, formed of germs of meromorphic functions on an open of $X$ whose inverse image exists and is regular on the inverse image open \sref{IV.20.1.11}.
Then if $f=(\psi,\theta)$, the canonical homomorphism \sref{IV.20.1.11} $\psi^*(\sh{M}_f)\to\sh{M}_{X'}$ gives by restriction a homomorphism of sheaves of groups $\psi^*(\sh{M}_f^*)\to\sh{M}_{X'}^*$.
Setting $\sh{D}iv_X'=\sh{M}_f^*/\sh{O}_X^*$, we have $\operatorname{Div}^f(X)=\Gamma(X,\sh{D}iv_X')$, and the map $D\mapsto f^*(D)$ corresponds to the homomorphism of sheaves of groups $\psi^*(\sh{M}_f^*)/\psi^*(\sh{O}_X^*)\to\sh{M}_{X'}^*/\sh{O}_{X'}^*=\sh{D}iv_{X'}$, derived from the preceding by passage to quotients.
\end{env}

\begin{env}[21.4.4]
\label{IV.21.4.4}
It follows also from the preceding definitions that if $f':X''\to X'$ is a second morphism of ringed spaces, $D$ a divisor on $X$ such that the inverse images $f^*(D)$ and $f'^*(f^*(D))$ exist, then $(f\circ f')^*(D)$ exists and is equal to $f'^*(f^*(D))$.
\end{env}

\begin{proposition}[21.4.5]
\label{IV.21.4.5}
Let $f:X'\to X$ be a morphism of ringed spaces.
In each of the three following cases, the inverse image by $f$ of every divisor on $X$ is defined:
\begin{enumerate}[label=(\roman*)]
  \item $f$ is flat.
  \item $X$ and $X'$ are locally Noetherian preschemes and $f(\operatorname{Ass}(\sh{O}_{X'}))\subset\operatorname{Ass}(\sh{O}_X)$.
  \item $X$ and $X'$ are preschemes, the set of irreducible components of $X$ is locally finite, $X'$ is reduced and every irreducible component of $X'$ dominates an irreducible component of $X$.
\end{enumerate}
\end{proposition}

\begin{proof}
It suffices indeed to show that in all three cases we have $\sh{M}_f=\sh{M}_X$;
in case (i), this follows from \sref{IV.20.1.12}.
In case (iii), we can restrict to the case where $X=\operatorname{Spec}(A)$ and $X'=\operatorname{Spec}(A')$ are affine;
if $s\in A$ is regular, it does not belong to any minimal prime ideal of $A$ \sref{IV.20.1.3.1}, hence the hypothesis implies that its image in $A'$ does not belong to any minimal prime ideal of $A'$, and is therefore a regular element of $A'$ \sref{IV.20.1.3.1}.
In case (ii) the meromorphic functions on $X'$ are identified with the pseudo-functions on $X'$ \sref{IV.20.2.11}, and the hypothesis, combined with \sref{IV.20.2.13},~(iv), ensures that the inverse image
\oldpage[IV-4]{267}
by $f$ of every schematically dense open in $X$ is a schematically dense open in $X'$;
we conclude therefore by \sref{IV.20.3.12}.
\end{proof}

\begin{corollary}[21.4.6]
\label{IV.21.4.6}
Let $X$ be a prescheme having one of the following properties:
\begin{enumerate}[label=(\roman*)]
  \item $X$ is locally Noetherian.
  \item $X$ is reduced and the set of its irreducible components is locally finite.
\end{enumerate}
Then, for every $x\in X$, we have a canonical isomorphism
\[
\label{IV.21.4.6.1}
  (\sh{D}iv_X)_x\xrightarrow{\sim}\operatorname{Div}(\sh{O}_{X,x}).
  \tag{21.4.6.1}
\]
\end{corollary}

\begin{proof}
This follows indeed from \sref{IV.20.2.11}, \sref{IV.20.3.7}, and from the fact that $(\sh{O}_X^*)_x$ is identified with the group of invertible elements of the ring $\sh{O}_{X,x}$.
\end{proof}

\begin{env}[21.4.7]
\label{IV.21.4.7}
Let $X$, $X'$ be two preschemes, $f:X'\to X$ a morphism.
If $D$ is a \emph{positive} divisor on $X$ such that the inverse image $f^*(D)$ is defined \sref{IV.21.4.2}, then the closed sub-prescheme $Y(f^*(D))$ of $X'$ is none other than the inverse image $f^{-1}(Y(D))$;
this follows immediately from the definitions \sref{IV.21.4.2} and \sref{IV.21.2.12}.
\end{env}

\begin{proposition}[21.4.8]
\label{IV.21.4.8}
Let $X$, $Y$ be two preschemes, $f:X\to Y$ a faithfully flat morphism.
Then, if $D$ is a divisor on $Y$ such that $f^*(D)\geq 0$ (the existence of $f^*(D)$ resulting from \sref{IV.21.4.5}), we have $D\geq 0$.
In particular, the map $D\mapsto f^*(D)$ of $\operatorname{Div}(Y)$ to $\operatorname{Div}(X)$ is injective.
\end{proposition}

\begin{proof}
The question being local on $Y$, we can restrict to the case where $D=\operatorname{div}(w)$, with $w=uv^{-1}$, $u$ and $v$ being two regular sections of $\sh{O}_Y$ over $Y$.
By hypothesis we have $u\sh{O}_X\subset v\sh{O}_X$, hence, for every $x\in X$, setting $y=f(x)$, we have $u_y\sh{O}_{X,x}\subset v_y\sh{O}_{X,x}$.
We conclude that $u_y\sh{O}_{Y,y}\subset v_y\sh{O}_{Y,y}$ because $\sh{O}_{X,x}$ is a faithfully flat $\sh{O}_{Y,y}$-module (Bourbaki, \emph{Alg. comm.}, chap.~I, \textsection~3, n\textsuperscript{o}~5, prop.~10,~(ii));
since $f$ is surjective, we have $u\sh{O}_Y\subset v\sh{O}_Y$, and consequently $D\geq 0$.
\end{proof}

% BEGIN NATIVE INPUT source_aligned/ega4-21-erratum53-insertion
% EGA IV, Errata and Addenda, List 3, no. 53.
% This component is inserted after 21.4.8 in the corrected reading text
% and repeated verbatim under Err_IV, 53 in the printed errata appendix.

The following proposition and its corollaries were pointed out to us by
M.~Raynaud.

\begin{proposition}[21.4.9]
\ifdefined\EGAIVSuppressErratumLabels\else\label{IV.21.4.9}\fi
Let $X$ and $Y$ be two locally Noetherian preschemes, and let
$f:X\to Y$ be a faithfully flat morphism.  Suppose in addition either
that $f$ is quasi-compact or that $f$ is open.  A divisor $D'$ on $X$
is of the form $f^*(D)$, with $D$ a divisor on $Y$, if and only if the
following condition holds: for every $y\in Y$ such that
$\operatorname{prof}(\sh{O}_{Y,y})\leq 1$, put
\[
 Y_1=\operatorname{Spec}(\sh{O}_{Y,y}),\qquad
 X_1=X\times_Y Y_1,
 \qquad f_1=f_{(Y_1)}:X_1\to Y_1;
\]
then the inverse image $D_1'$ of $D'$ on $X_1$ is of the form
$f_1^*(D_1)$ for some $D_1\in\operatorname{Div}(Y_1)$.
\end{proposition}

\oldpage[IV]{359}
\begin{proof}
The necessity of the condition is immediate from the transitivity of
inverse images of divisors, all the morphisms in question being flat
{\hyperlink{ega4.c04.statement.21-4-4}{(21.4.4)}}.  For every open subset $U$ of $Y$, write $f_U$ for the
restriction $f^{-1}(U)\to U$ of $f$.  By {\hyperlink{ega4.c04.prop.21-4-8}{(21.4.8)}}, the map
\[
 D_U\longmapsto f_U^*(D_U)
\]
from $\operatorname{Div}(U)$ to
$\operatorname{Div}(f^{-1}(U))$ is injective.  Thus, if $U_1,U_2$ are
two open subsets of $Y$ and
\[
 D'|f^{-1}(U_1)=f_{U_1}^*(D_1),\qquad
 D'|f^{-1}(U_2)=f_{U_2}^*(D_2),
\]
then necessarily
$D_1|U_1\cap U_2=D_2|U_1\cap U_2$.  It follows at once that there is a
largest open subset $U\subset Y$ for which
\[
 D'|f^{-1}(U)=f_U^*(D_U)
 \quad\text{with}\quad D_U\in\operatorname{Div}(U).
\]
We must show that $U=Y$.  Suppose otherwise, and let $y$ be a maximal
point of $Y\setminus U$.  We shall find an open neighbourhood $V$ of
$y$ in $Y$ such that
$D'|f^{-1}(V)=f_V^*(D_V)$ for some
$D_V\in\operatorname{Div}(V)$, contradicting the maximality of $U$.

\textup{I) The case where $f$ is quasi-compact.}
Using the quasi-compactness of $f$ and {\hyperlink{ega4.c04.prop.20-3-8}{(20.3.8)}}, by the method of
({\hyperlink{ega4.c04.statement.8-1-2}{8.1.2}}, a), we are reduced to proving---in the notation of the
statement, but without initially imposing a bound on
$\operatorname{prof}(\sh{O}_{Y,y})$---that $D_1'$ is of the form
$f_1^*(D_1)$.  If $\operatorname{prof}(\sh{O}_{Y,y})\leq 1$, this is
precisely the hypothesis.  Suppose therefore that
$\operatorname{prof}(\sh{O}_{Y,y})\geq 2$.  Since $y$ is a maximal
point of $Y\setminus U$, the open subset $W=Y_1\cap U$ of $Y_1$ is
$Y_1\setminus\{y\}$.  We may therefore restrict to the case
$Y=Y_1$ and $U=Y\setminus\{y\}$.

By hypothesis there is a divisor $D_U\in\operatorname{Div}(U)$ such
that $f_U^*(D_U)=D'|f^{-1}(U)$.  The divisor $D_U$ is the equivalence
class of a pair $(\sh{L}_U,s)$, where $\sh{L}_U$ is an invertible
$\sh{O}_U$-Module and $s$ a regular meromorphic section of $\sh{L}_U$
over $U$ {\hyperlink{ega4.c04.prop.21-2-11}{(21.2.11)}}.  Hence, by {\hyperlink{ega4.c04.statement.21-4-1}{(21.4.1)}}, $D'|f^{-1}(U)$ is the
equivalence class of $(\sh{L}'_U,s')$, where
\[
 \sh{L}'_U=f_U^*(\sh{L}_U),\qquad s'=s\circ f_U.
\]
Let $i:U\to Y$ and $j:f^{-1}(U)\to X$ be the canonical injections,
and put
\[
 \sh{L}=i_*(\sh{L}_U),\qquad
 \sh{L}'=j_*(\sh{L}'_U).
\]
Since $f$ is flat, $f^*(\sh{L})$ is canonically isomorphic to
$\sh{L}'$ {\hyperlink{ega4.c04.lem.2-3-1}{(2.3.1)}}.  Moreover, flatness and the inequality
$\operatorname{prof}(\sh{O}_{Y,y})\geq 2$ imply
$\operatorname{prof}(\sh{O}_{X,x})\geq 2$ for every
$x\in f^{-1}(y)$ {\hyperlink{ega4.c04.prop.6-3-1}{(6.3.1)}}.

The divisor $D'$ is the equivalence class of a pair
$(\sh{L}'',s'')$, where $\sh{L}''$ is an invertible $\sh{O}_X$-Module
and $s''$ a regular meromorphic section of $\sh{L}''$ over $X$.
Consequently $\sh{L}'_U$ is isomorphic to
$\sh{L}''|f^{-1}(U)$, and the preceding observations together with
{\hyperlink{ega4.c04.prop.5-9-8}{(5.9.8)}} and {\hyperlink{ega4.c04.thm.5-10-5}{(5.10.5)}} show that $\sh{L}'$ is an invertible
$\sh{O}_X$-Module isomorphic to $\sh{L}''$.  Since $f$ is faithfully
flat and quasi-compact, it follows that $\sh{L}$ is an invertible
$\sh{O}_Y$-Module {\hyperlink{ega4.c04.prop.2-5-2}{(2.5.2)}}, and hence isomorphic to $\sh{O}_Y$, because
$Y$ is a local scheme.

Furthermore, since $\operatorname{prof}(\sh{O}_{Y,y})\geq 1$, the open
subset $U=Y\setminus\{y\}$ is schematically dense in $Y$
({\hyperlink{ega4.c04.statement.20-2-13}{20.2.13}}, (iv)).  The regular meromorphic section $s$ of $\sh{L}_U$
over $U$ therefore extends to a regular meromorphic section $s_0$ of
$\sh{L}$ over $Y$ {\hyperlink{ega4.c04.prop.20-2-11}{(20.2.11)}}.  The equivalence class of
$(\sh{L},s_0)$ is thus a divisor $D$ on $Y$.  Plainly $f^*(D)$ and
$D'$ induce the same divisor on
$f^{-1}(U)=X\setminus f^{-1}(y)$; and since
$\operatorname{prof}(\sh{O}_{X,x})\geq 2$ for every
$x\in f^{-1}(y)$, {\hyperlink{ega4.c04.prop.21-1-8}{(21.1.8)}} gives $D'=f^*(D)$.  This proves case I.

\textup{II) The case where $f$ is open.}
The problem is to find an open neighbourhood $V$ of $y$ in $Y$ such
that $D'|f^{-1}(V)=f_V^*(D_V)$ for some
$D_V\in\operatorname{Div}(V)$, so we may first suppose $Y$ affine.
Because $f$ is open, the images $f(W)$, with $W$ running through the
affine open subsets of $X$, form an open covering of $Y$.  Since $Y$
is quasi-compact, there is a quasi-compact open subset $T\subset X$
such that $f(T)=Y$, and $f|T$ is quasi-compact (I, 1.1.1).  A fortiori,
for every quasi-compact open $T_\alpha\supset T$ in $X$,
\oldpage[IV]{360}
one has
$f(T_\alpha)=Y$ and $f|T_\alpha$ is quasi-compact.

Applying case I to $f|T_\alpha$, we obtain a divisor $D_\alpha$ on
$Y$ such that
\[
 (f|T_\alpha)^*(D_\alpha)=D'|T_\alpha.
\]
By transitivity of inverse images of divisors under flat morphisms
{\hyperlink{ega4.c04.statement.21-4-4}{(21.4.4)}}, one also has
$(f|T)^*(D_\alpha)=D'|T$.  Hence for any $\alpha,\beta$,
\[
 (f|T)^*(D_\alpha)=(f|T)^*(D_\beta),
\]
and {\hyperlink{ega4.c04.prop.21-4-8}{(21.4.8)}} gives $D_\alpha=D_\beta$.  If $D$ denotes their common
value, then $f^*(D)|T_\alpha=D'|T_\alpha$ for every $\alpha$, and
therefore $f^*(D)=D'$.
\end{proof}

\begin{corollary}[21.4.10]
\ifdefined\EGAIVSuppressErratumLabels\else\label{IV.21.4.10}\fi
Assume the general hypotheses of {\hyperlink{ega4.c04.prop.21-4-9}{(21.4.9)}}, and suppose in addition that
$Y$ is normal and that, for every $y\in Y$ with
$\dim(\sh{O}_{Y,y})=1$, the fibre $X_y=f^{-1}(y)$ has no embedded
associated prime cycle.  A divisor $D'$ on $X$ is of the form $f^*(D)$,
with $D\in\operatorname{Div}(Y)$, if and only if the following two
conditions hold:
\begin{enumerate}[label=\textup{\arabic*\textdegree}]
\item for every maximal point $y$ of $Y$, the inverse image of $D'$ in
$\operatorname{Div}(X_y)$ is zero;
\item for every $y\in Y$ such that $\dim(\sh{O}_{Y,y})=1$, there is an
integer $n_y$ such that, for every maximal point $x$ of $X_y$, the
inverse image of $D'$ in
$\operatorname{Div}(\operatorname{Spec}(\sh{O}_{X,x}))$ equals
$\operatorname{div}(t_y^{n_y})$, where $t_y$ is the image under
$\sh{O}_{Y,y}\to\sh{O}_{X,x}$ of a uniformizer of the discrete
valuation ring $\sh{O}_{Y,y}$.
\end{enumerate}
\end{corollary}

\begin{proof}
Apply the criterion {\hyperlink{ega4.c04.prop.21-4-9}{(21.4.9)}}.  At a maximal point $y$ of $Y$,
$\operatorname{Div}(Y_1)=0$ because $\sh{O}_{Y,y}$ is Artinian;
condition 1\textdegree{} is therefore equivalent to the criterion of
{\hyperlink{ega4.c04.prop.21-4-9}{(21.4.9)}} at such a point.  It remains to show that conditions
1\textdegree{} and 2\textdegree{} likewise imply the criterion at a
point $y\in Y$ with $\operatorname{prof}(\sh{O}_{Y,y})=1$.  Since $Y$
is normal, these are precisely the points for which
$\dim(\sh{O}_{Y,y})=1$, equivalently those for which
$\sh{O}_{Y,y}$ is a discrete valuation ring {\hyperlink{ega4.c04.thm.5-8-6}{(5.8.6)}}.

We may thus restrict to $Y=\operatorname{Spec}(\sh{O}_{Y,y})$.  If $t$
is a uniformizer of $\sh{O}_{Y,y}$, the map
$n\mapsto\operatorname{div}(t^n)$ is an isomorphism from $\mathbb Z$
onto $\operatorname{Div}(Y)$.  The condition of {\hyperlink{ega4.c04.prop.21-4-9}{(21.4.9)}} says that
$D'$ is of the form $\operatorname{div}(t^n\circ f)$ for some
$n\in\mathbb Z$, and plainly implies condition 2\textdegree{}.
Conversely, suppose 2\textdegree{} holds, as does 1\textdegree{}, and
write $n=n_y$.  It remains to prove
$D'=\operatorname{div}(t^n\circ f)$.  If $\eta$ is the generic point of
$Y$, the germs of these two divisors at every point of $X_\eta$ are
zero by condition 1\textdegree{}.  Their germs at every maximal point
of $X_y$ are equal by condition 2\textdegree{}.  At a nonmaximal point
$x\in X_y$, one has
$\operatorname{prof}(\sh{O}_{X_y,x})\geq 1$, because $X_y$ has no
embedded associated prime cycle, and hence
$\operatorname{prof}(\sh{O}_{X,x})\geq 2$ by {\hyperlink{ega4.c04.prop.6-3-1}{(6.3.1)}}.  The desired
equality follows from {\hyperlink{ega4.c04.prop.21-1-8}{(21.1.8)}}.
\end{proof}

\begin{corollary}[21.4.11]
\ifdefined\EGAIVSuppressErratumLabels\else\label{IV.21.4.11}\fi
Assume the general hypotheses of {\hyperlink{ega4.c04.prop.21-4-9}{(21.4.9)}}, and suppose in addition that
$Y$ is normal and that, for every $y\in Y$ with
$\dim(\sh{O}_{Y,y})=1$, the fibre $X_y$ is integral.  A divisor $D'$ is
of the form $f^*(D)$, with $D\in\operatorname{Div}(Y)$, if and only if,
for every maximal point $y$ of $Y$, the inverse image of $D'$ in
$\operatorname{Div}(X_y)$ is zero.
\end{corollary}

\begin{proof}
It is enough to check that condition 2\textdegree{} of {\hyperlink{ega4.c04.cor.21-4-10}{(21.4.10)}} is then
automatic.  If $\dim(\sh{O}_{Y,y})=1$ and $X_y$ is integral, then at
the unique maximal point $x$ of $X_y$ one has
$\dim(\sh{O}_{X,x})=1$ by {\hyperlink{ega4.c04.cor.6-1-3}{(6.1.3)}}.  Moreover,
$\sh{O}_{Y,y}$ is regular and $\sh{O}_{X_y,x}$ is a field, so
$\sh{O}_{X,x}$ is regular {\hyperlink{ega4.c04.prop.6-5-1}{(6.5.1)}}, hence a discrete valuation ring.
If $t$ is a uniformizer of $\sh{O}_{Y,y}$, its image $t_y$ in
$\sh{O}_{X,x}$ is a uniformizer of $\sh{O}_{X,x}$
($0$, 17.3.3).  The inverse image of $D'$ in
$\operatorname{Div}(\operatorname{Spec}(\sh{O}_{X,x}))$ is therefore
\oldpage[IV]{361}
of the form $\operatorname{div}(t_y^n)$, which is exactly condition
2\textdegree{} of {\hyperlink{ega4.c04.cor.21-4-10}{(21.4.10)}}, since $X_y$ has only one maximal point.
\end{proof}

\begin{remark}[21.4.12]
\ifdefined\EGAIVSuppressErratumLabels\else\label{IV.21.4.12}\fi
If in {\hyperlink{ega4.c04.cor.21-4-11}{(21.4.11)}} one assumes only that the $X_y$ are reduced, one must
add to the condition in the statement that the inverse images of $D'$
in each group
$\operatorname{Div}(\operatorname{Spec}(\sh{O}_{X,x}))$, with $x$
running over the maximal points of $X_y$, are all of the form
$\operatorname{div}((t_x')^n)$, where $t_x'$ is a uniformizer of
$\sh{O}_{X,x}$ and the integer $n$ is the same for all these points.
\end{remark}

\begin{corollary}[21.4.13]
\ifdefined\EGAIVSuppressErratumLabels\else\label{IV.21.4.13}\fi
Assume the hypotheses of {\hyperlink{ega4.c04.cor.21-4-11}{(21.4.11)}}.  For every invertible
$\sh{O}_X$-Module $\sh{L}$ such that, for every maximal point $\eta$ of
$Y$, the inverse image
\[
 \sh{L}_\eta
 =\sh{L}\otimes_{\sh{O}_X}\sh{O}_{X_\eta}
\]
is isomorphic to $\sh{O}_{X_\eta}$, there exist an invertible
$\sh{O}_Y$-Module $\sh{M}$ and an isomorphism
$f^*(\sh{M})\simeq\sh{L}$.
\end{corollary}

\begin{proof}
It is enough to find a regular meromorphic section $s$ of $\sh{L}$ over
$X$ such that the divisor $D'$ corresponding to $(\sh{L},s)$ has
$D'_x=0$ at every point $x$ of a fibre $X_\eta$ over a maximal point
$\eta$ of $Y$.  Applying {\hyperlink{ega4.c04.cor.21-4-11}{(21.4.11)}} to $D'$ then gives the result by
{\hyperlink{ega4.c04.def.21-4-2}{(21.4.2)}}.

The points of $\operatorname{Ass}(\sh{O}_X)$ lie over the maximal
points of $Y$ {\hyperlink{ega4.c04.prop.3-3-1}{(3.3.1)}}.  Since $Y$ is reduced and locally Noetherian,
there is a dense open subset $V$ of $Y$ which is the union of a
disjoint family of affine open subsets, each containing exactly one
maximal point of $Y$; and $f^{-1}(V)$ is schematically dense in $X$
{\hyperlink{ega4.c04.prop.5-10-2}{(5.10.2)}}.  We may therefore restrict to the case in which $Y$ is
affine and integral.

For the unique maximal point $\eta$ of $Y$, the fibre $X_\eta$ is then
integral, so $\operatorname{Ass}(\sh{O}_X)$ consists of a single point
$z$, the generic point of $X$.  Let $s_\eta$ be the section of
$\sh{L}_\eta$ over $X_\eta$ which corresponds, under an isomorphism,
to the unit section of $\sh{O}_{X_\eta}$.  Let $U$ be an affine open
neighbourhood in $X$ of a point $x\in X_\eta$, and let $s$ be a section
of $\sh{L}$ over $U$ with the same germ at $x$ as $s_\eta$.  Since $U$
is schematically dense in $X$ ({\hyperlink{ega4.c04.statement.20-2-13}{20.2.13}}, (iv)), $s$ identifies with a
uniquely determined regular meromorphic section of $\sh{L}$ over $X$.
It is clear from {\hyperlink{ega4.c04.prop.20-2-11}{(20.2.11)}} and {\hyperlink{ega4.c04.prop.20-3-5}{(20.3.5)}} that this section induces on
$X_\eta$ a regular meromorphic section of $\sh{L}_\eta$ equal to
$s_\eta$, and hence has the required property.
\end{proof}
% END NATIVE INPUT source_aligned/ega4-21-erratum53-insertion

\subsection{Direct images of divisors}
\label{subsection:IV.21.5}
\label{IV.21.5}

\begin{env}[21.5.1]
\label{IV.21.5.1}
Let $X$, $X'$ be two preschemes, $f:X'\to X$ a morphism.
In this section, we shall give sufficient conditions for one to associate to a divisor $D'$ on $X'$ a divisor $D$ on $X$, the \emph{direct image} of $D'$ by $f$.
We will restrict ourselves to the case where $f$ is a \emph{finite} morphism (for more general conditions, see the chapter of this Treatise devoted to the theory of intersections).
\end{env}

\begin{lemma}[21.5.2]
\label{IV.21.5.2}
Let $A$ be a ring, $E$ an $A$-module free of finite rank.
For an endomorphism $u$ of $E$ to be injective, it is necessary and sufficient that $\det(u)$ be a regular element of $A$.
\end{lemma}

\begin{proof}
This is proved in Bourbaki, \emph{Alg.}, chap.~III, 3\textsuperscript{e}~\'{e}d., \textsection~8, n\textsuperscript{o}~2, prop.~3.
\end{proof}

\begin{env}[21.5.3]
\label{IV.21.5.3}
Suppose now that the morphism $f:X'\to X$ is \emph{finite}, and in addition that $f$ satisfies one of the two following properties:
\begin{enumerate}[label=\Roman*)]
  \item $f$ is a finite morphism \emph{locally free}, in other words \sref{IV.18.2.7} the quasi-coherent $\sh{O}_X$-Module of finite type $f_*(\sh{O}_{X'})$ is \emph{locally free}.
  \item $X$ is a locally Noetherian reduced prescheme, the $\sh{R}(X)$-Module $f_*(\sh{O}_{X'})\otimes_{\sh{O}_X}\sh{R}(X)$ is locally free, and for every section $s'\in\Gamma(f^{-1}(U),\sh{O}_{X'})$ ($U$ open
\oldpage[IV-4]{268}
in $X$), $\mathrm{N}_{f_*(\sh{O}_{X'})|\sh{O}_X}(s')$ is a section of $\sh{O}_X$ over $U$ (cf.\ (\textbf{II}, 6.5.1)).
  (We recall that condition (II) is satisfied for every \emph{finite morphism} $f$ when $X$ is a locally Noetherian normal prescheme (\emph{loc.\ cit.}).)
\end{enumerate}

We know then (\textbf{II}, 6.5.5) that for every invertible $\sh{O}_{X'}$-Module $\sh{L}'$ we define the \emph{norm} $\sh{L}=\mathrm{N}_{X'/X}(\sh{L}')$ (which we also write $\mathrm{N}(\sh{L}')$), which is an invertible $\sh{O}_X$-Module.
Moreover, for every open $U$ of $X$ and every \emph{regular} section $s'\in\Gamma(f^{-1}(U),\sh{O}_{X'})$, the norm $\mathrm{N}_{X'/X}(s')=\mathrm{N}_{f_*(\sh{O}_{X'})|\sh{O}_X}(s')$ (which we also write $\mathrm{N}(s')$) is a regular element of $\Gamma(U,\sh{O}_X)$;
we are indeed immediately reduced to the case where $U=X$ is affine, and then the conclusion follows from \sref{IV.21.5.2} in hypothesis (I);
on the other hand, in hypothesis (II), the fact that $\sh{R}(X)$ is a flat $\sh{O}_X$-Module implies that the section $s'\otimes 1$ of $\Gamma(U,f_*(\sh{O}_{X'})\otimes_{\sh{O}_X}\sh{R}(X))$ is also regular (\hyperref[0.6.1.2]{$\mathbf{0_I}$, 6.1.2}), and the conclusion follows again from \sref{IV.21.5.2} applied to the ring $\Gamma(U,\sh{R}(X))$, taking into account the definition of the norm (\textbf{II}, 6.5.3).
This being so, let $u'$ be a meromorphic section of $\sh{L}'$ over $X'$;
the morphism $f$ being affine, every point $x\in X$ admits a neighbourhood $U$ such that $u'|f^{-1}(U)$ is of the form $t'/s'$, where $t'\in\Gamma(f^{-1}(U),\sh{L}')$ and $s'$ is a regular section of $\sh{O}_{X'}$ over $f^{-1}(U)$;
the element $\mathrm{N}(t')/\mathrm{N}(s')$ (where $\mathrm{N}(t')$ is the section of $\sh{L}$ defined in (\textbf{II}, 6.5.3)) is then a meromorphic section of $\sh{L}$ over $U$, and from what precedes and the multiplicative properties of the norm (\textbf{II}, 6.5.3.1) this section depends only on $u'|f^{-1}(U)$ and not on its representation in the form $t'/s'$;
for the same reason, the sections of $\sh{L}$ over $U$ thus defined are the restrictions of a single meromorphic section of $\sh{L}$ over $X$, which we call the \emph{norm} of $u'$ and denote $\mathrm{N}_{X'/X}(u')$ (or simply $\mathrm{N}(u')$).
The map
\[
\label{IV.21.5.3.1}
  \mathrm{N}_{X'/X}:\Gamma(X',\sh{M}_{X'}(\sh{L}'))\to\Gamma(X,\sh{M}_X(\mathrm{N}_{X'/X}(\sh{L}')))
  \tag{21.5.3.1}
\]
extends the norm defined in (\textbf{II}, 6.5.5.3);
if $u'$ is a \emph{regular} meromorphic section of $\sh{L}'$ over $X'$, it follows immediately from what precedes that $\mathrm{N}(u')$ is a regular meromorphic section of $\sh{L}$ over $X$, since (with the same notation) $\mathrm{N}(t')$ is regular if $t'$ is.
Finally, if $\sh{L}_1'$, $\sh{L}_2'$ are two invertible $\sh{O}_{X'}$-Modules, $s_1'$ (resp.\ $s_2'$) a regular meromorphic section of $\sh{L}_1'$ (resp.\ $\sh{L}_2'$) over $X'$, we have, by virtue of what precedes and of (\textbf{II}, 6.5.3.1)
\[
\label{IV.21.5.3.2}
  \mathrm{N}(s_1'\otimes s_2')=\mathrm{N}(s_1')\otimes\mathrm{N}(s_2').
  \tag{21.5.3.2}
\]
\end{env}

\begin{env}[21.5.4]
\label{IV.21.5.4}
Suppose always that $f$ satisfies one of hypotheses I), II) of \sref{IV.21.5.3}.
If $(\sh{L}_1',s_1')$, $(\sh{L}_2',s_2')$ are two pairs each formed of an invertible $\sh{O}_{X'}$-Module and a regular meromorphic section of this Module over $X'$, which are moreover \emph{equivalent} in the sense of \sref{IV.21.2.10}, then the pairs $(\mathrm{N}(\sh{L}_1'),\mathrm{N}(s_1'))$ and $(\mathrm{N}(\sh{L}_2'),\mathrm{N}(s_2'))$ are also equivalent, since an isomorphism of invertible $\sh{O}_{X'}$-Modules has for norm an isomorphism of their norms (\textbf{II}, 6.5.3), and we have already seen above that $\mathrm{N}_{X'/X}$ transforms the sections of $\Gamma(X',\sh{O}_{X'}^*)$ to sections of $\Gamma(X,\sh{O}_X^*)$.
We can therefore pose the following definition:
\end{env}

\oldpage[IV-4]{269}

\begin{definition}[21.5.5]
\label{IV.21.5.5}
Given a finite morphism $f:X'\to X$ of preschemes, satisfying one of conditions I), II) of \sref{IV.21.5.3}, we call the \emph{direct image} (or \emph{norm}) of a divisor $D'$ on $X'$ by $f$, and denote $f_*(D')$ (or $\mathrm{N}_{X'/X}(D')$), the divisor on $X$ which corresponds canonically \sref{IV.21.2.11} to the class of the pair $(\mathrm{N}_{X'/X}(\sh{O}_{X'}(D')),\mathrm{N}_{X'/X}(s_{D'}))$.
\end{definition}

It follows immediately from this definition, taking into account \sref{IV.21.2.9.4} and \sref{IV.21.5.3.2}, that if $D_1'$, $D_2'$, $D'$ are divisors on $X'$, we have
\[
\label{IV.21.5.5.1}
  f_*(D_1'+D_2')=f_*(D_1')+f_*(D_2')
  \tag{21.5.5.1}
\]
and $D'\geq 0$ implies $f_*(D')\geq 0$, in other words, $D'\mapsto f_*(D')$ is an \emph{increasing homomorphism} of the ordered group $\operatorname{Div}(X')$ to the ordered group $\operatorname{Div}(X)$.
The definition \sref{IV.21.5.5} also shows immediately that for every open $U$ of $X$, we have $f_*^U(D'|f^{-1}(U))=f_*(D')|U$ ($f^U$ being the restriction $f^{-1}(U)\to U$ of $f$), and the homomorphisms $f_*^U$, for $U$ variable, thus define a \emph{homomorphism of sheaves of ordered groups}
\[
\label{IV.21.5.5.2}
  \mathrm{N}_{X'/X}:f_*(\sh{D}iv_{X'})\to\sh{D}iv_X.
  \tag{21.5.5.2}
\]
Moreover, for every invertible $\sh{O}_{X'}$-Module $\sh{L}'$ and every regular meromorphic section $s'$ of $\sh{L}'$ over $f^{-1}(U)$, we have, from the preceding definitions and \sref{IV.21.1.4}
\[
\label{IV.21.5.5.3}
  \operatorname{div}_U(\mathrm{N}(s'))=f_*^U(\operatorname{div}_{f^{-1}(U)}(s')).
  \tag{21.5.5.3}
\]

\begin{proposition}[21.5.6]
\label{IV.21.5.6}
Let $f:X'\to X$ be a finite morphism locally free and suppose that $f_*(\sh{O}_{X'})$ is of constant rank $n$.
Then, for every divisor $D$ on $X$, $f^*(D)$ is defined and we have
\[
\label{IV.21.5.6.1}
  f_*(f^*(D))=nD.
  \tag{21.5.6.1}
\]
\end{proposition}

\begin{proof}
The first assertion follows from the fact that $f$ is flat \sref{IV.21.4.5}, and the second is an immediate consequence of the definitions and of (\textbf{II}, 6.5.3.2).
\end{proof}

\begin{proposition}[21.5.7]
\label{IV.21.5.7}
Let $f:X'\to X$ be a finite morphism satisfying one of hypotheses I), II) of \sref{IV.21.5.3}, $f':X''\to X'$ a finite morphism locally free of constant rank $n$.
Then $f''=f\circ f':X''\to X$ satisfies the same hypothesis as $f$, and for every divisor $D''$ on $X''$, we have
\[
\label{IV.21.5.7.1}
  f_*''(D'')=f_*(f_*'(D'')).
  \tag{21.5.7.1}
\]
\end{proposition}

Taking into account the definition \sref{IV.21.5.5}, it suffices to prove the following result:

\begin{lemma}[21.5.7.2]
\label{IV.21.5.7.2}
Under the hypotheses of \sref{IV.21.5.7}, we have a functorial isomorphism
\[
\label{IV.21.5.7.3}
  \mathrm{N}_{X''/X}(\sh{L}'')\xrightarrow{\sim}\mathrm{N}_{X'/X}(\mathrm{N}_{X''/X'}(\sh{L}''))
  \tag{21.5.7.3}
\]
in the category of invertible $\sh{O}_X$-Modules.
\end{lemma}

Indeed, taking into account the definition of the norm of a section of $\sh{L}''$ (\textbf{II}, 6.5.3) and definition \sref{IV.21.5.5}, formula \sref{IV.21.5.7.1} follows at once from the lemma.
To prove the lemma, it suffices, by the definitions of (\textbf{II}, 6.5.2 and 6.5.3), to prove that for every section $s$ of $f_*''(\sh{O}_{X''})=f_*(f_*'(\sh{O}_{X''}))$ over an open $U$ of $X$, we have
\[
\label{IV.21.5.7.4}
  \mathrm{N}_{f_*''(\sh{O}_{X''})|\sh{O}_X}(s)=\mathrm{N}_{f_*(\sh{O}_{X'})|\sh{O}_X}(\mathrm{N}_{f_*'(\sh{O}_{X''})|\sh{O}_{X'}}(s))
  \tag{21.5.7.4}
\]
\oldpage[IV-4]{270}
The question is evidently local on $X$, and we may therefore restrict to the case where $X=\operatorname{Spec}(A)$ is affine;
then $X'=\operatorname{Spec}(A')$ and $X''=\operatorname{Spec}(A'')$, and we may suppose that $A''$ is a projective $A'$-module of rank $n$.
When $f$ is locally free, we may suppose that $A'$ is a free $A$-module of rank $m$;
then $A''$ is a projective $A$-module of rank $mn$, and after restricting $X$ to a suitable open, we may suppose that $A''$ is a free $A$-module of rank $mn$.
Formula \sref{IV.21.5.7.4} then follows from the transitivity of the norm (Bourbaki, \emph{Alg.}, chap.~VIII, \textsection~12, n\textsuperscript{o}~2, prop.~7).
When $f$ satisfies hypothesis~II), $A$ is Noetherian and reduced, and if $R$ is its total ring of fractions, $A'\otimes_A R$ is a free $R$-module of rank $m$;
hence $A''\otimes_A R=A''\otimes_{A'}(A'\otimes_A R)$ is a projective $R$-module of rank $mn$, and since $R$ is a semi-local ring, this module is free.
The proposition then follows again from the transitivity of norms.

\begin{proposition}[21.5.8]
\label{IV.21.5.8}
Let $f:X'\to X$ be a finite morphism, $g:Y\to X$ a morphism;
set $Y'=X'\times_XY$, $f'=f_{(Y)}:Y'\to Y$, $g'=g_{(X')}:Y'\to X'$.
Suppose that one of the following conditions holds:
\begin{enumerate}[label=(\roman*)]
  \item $f$ is locally free and $g$ is flat.
  \item $f$ is locally free, $X$ and $Y$ are locally Noetherian, $g(\operatorname{Ass}(\sh{O}_Y))\subset\operatorname{Ass}(\sh{O}_X)$ and $g'(\operatorname{Ass}(\sh{O}_{Y'}))\subset\operatorname{Ass}(\sh{O}_{X'})$.
  \item $f$ satisfies hypothesis \emph{II)} of \sref{IV.21.5.3}, $Y$ is locally Noetherian, $Y$ and $Y'$ are reduced and every irreducible component of $Y$ (resp.\ $Y'$) dominates an irreducible component of $X$ (resp.\ $X'$).
\end{enumerate}
Then, for every divisor $D'$ on $X'$, $g'^*(D')$ is defined, $g^*(f_*(D'))$ is defined and we have
\[
\label{IV.21.5.8.1}
  g^*(f_*(D'))=f_*'(g'^*(D')).
  \tag{21.5.8.1}
\]
\end{proposition}

\begin{proof}
Indeed, in all cases, it follows from (\textbf{II}, 6.5.8) that there is a functorial isomorphism
\[
  g^*(\mathrm{N}_{X'/X}(\sh{L}'))\xrightarrow{\sim}\mathrm{N}_{Y'/Y}(g'^*(\sh{L}'))
\]
in the category of invertible $\sh{O}_Y$-Modules;
moreover (\textbf{II}, 6.5.4), if $s'$ is a section of $\sh{O}_{X'}$ over $f^{-1}(U)$, and $s''$ the corresponding section of $\sh{O}_{Y'}$ over $g'^{-1}(f^{-1}(U))$ (where $U$ is open in $X$), then $\mathrm{N}_{Y'/Y}(s'')$ is the section of $\sh{O}_Y$ over $g^{-1}(U)$ corresponding to the section $\mathrm{N}_{X'/X}(s')$ of $\sh{O}_X$ over $U$.
Formula \sref{IV.21.5.8.1} therefore follows from the definitions once it is proved that $g'^*(D')$ and $g^*(D)$ are defined for arbitrary divisors $D'$ on $X'$ and $D$ on $X$.
For $D$, this follows from the hypotheses and \sref{IV.21.4.5}.
For $D'$, $g'$ is flat in case~(i), while in cases~(ii) and~(iii) the corresponding hypotheses of \sref{IV.21.4.5} hold;
thus $g'^*(D')$ is defined in every case.
\end{proof}
\subsection{1-codimensional cycle associated to a divisor}
\label{subsection:IV.21.6}
\label{IV.21.6}


\begin{env}[21.6.1]
\label{IV.21.6.1}
Let $X$ be a locally Noetherian prescheme, and let $\mathfrak{I}(X)$ be the set of \emph{irreducible closed subsets} of $X$ (which is in bijective correspondence with $X$ by the map $x\mapsto\overline{\{x\}}$).
In the product group $\bb{Z}^X$, consider the subgroup $\mathfrak{Z}(X)$ of elements $(n_x)_{x\in X}$ such that the set of $\overline{\{x\}}\in\mathfrak{I}(X)$ for which $n_x\neq0$ (or, equivalently, the set of such $x\in X$) is locally finite.
It is clear that $\mathfrak{Z}(X)$ is a subgroup of $\bb{Z}^X$ containing the direct sum $\bb{Z}^{(X)}$ (the free group with basis $\mathfrak{I}(X)$), and equal to it when $X$ is Noetherian.
The elements of $\mathfrak{Z}(X)$ are called \emph{cycles} on $X$, and those of $\mathfrak{I}(X)$ \emph{prime cycles} (they do not in general form a basis of $\mathfrak{Z}(X)$ when $X$ is not Noetherian).
We always regard $\mathfrak{Z}(X)$ as ordered by the order induced from the product order on $\bb{Z}^X$, and denote by $\mathfrak{Z}^+(X)$ the set of cycles $\geq0$.

\oldpage[IV-4]{271}

For every cycle $Z\in\mathfrak{Z}(X)$, equal to $(n_x)_{x\in X}$, we write, by abuse of language,
\[
  Z=\sum_{x\in X}n_x\cdot\overline{\{x\}}.
\]
We call \emph{multiplicity of $Z$ at the point $x$} and denote by $\operatorname{mult}_x(Z)$ the integer $n_x$ (positive or negative).
To say that $Z\geq 0$ means that $\operatorname{mult}_x(Z)\geq 0$ for all $x\in X$.
We call the \emph{support} of $Z$ and denote by $\Supp(Z)$ the union of the $\overline{\{x\}}$ such that $\operatorname{mult}_x(Z)\neq 0$;
by definition of $\mathfrak{Z}(X)$, this is a closed subset of $X$, as the union of a locally finite family of closed subsets.
We call the \emph{dimension} (resp.\ \emph{codimension}) of $Z$ and denote by $\dim(Z)$ (resp.\ $\codim(Z)$) the dimension (resp.\ codimension in $X$) of $\Supp(Z)$.
\end{env}

\begin{env}[21.6.2]
\label{IV.21.6.2}
We say that a closed subspace $\mathrm{Y}$ of $\mathrm{X}$ is \emph{purely of codimension $p$} (in $\mathrm{X}$) if \emph{each} of its irreducible components is of codimension $p$ in $\mathrm{X}$.
We say that a cycle is \emph{purely of codimension $p$}, or (by abuse of language) is a \emph{$p$-codimensional cycle}, if its support is purely of codimension $p$.
We denote by $\mathrm{X}^{(p)}$ the set of $x\in\mathrm{X}$ such that $\codim(\overline{\{x\}},\mathrm{X})=p$, or, what amounts to the same \sref{IV.5.1.2},{} $\dim(\sh{O}_{\mathrm{X},x})=p$:
the cycles purely of codimension $p$ form a subgroup $\mathfrak{Z}^p(\mathrm{X})$ of $\mathfrak{Z}(\mathrm{X})$, isomorphic to the sub-group of $\bb{Z}^{\mathrm{X}^{(p)}}$ consisting of the $(n_x)_{x\in\mathrm{X}^{(p)}}$ such that the set of $\overline{\{x\}}$ (or the set of $x$) for which $n_x\neq 0$ is locally finite;
this sub-group contains the free group $\bb{Z}^{(\mathrm{X}^{(p)})}$ and is identical to it when $\mathrm{X}$ is Noetherian.
We denote by $\mathfrak{Z}^{p+}(\mathrm{X})$ the set of elements $\geq 0$ of $\mathfrak{Z}^p(\mathrm{X})$.
It is clear that the ordered group $\mathfrak{Z}(\mathrm{X})$ contains the direct sum of the ordered sub-groups $\mathfrak{Z}^p(\mathrm{X})$, and is identical to this direct sum when $\mathrm{X}$ is Noetherian;
in the latter case, we consider $\mathfrak{Z}(\mathrm{X})$ as \emph{graded} by the $\mathfrak{Z}^p(\mathrm{X})$ for $p\in\bb{N}$.
\end{env}

\begin{env}[21.6.3]
\label{IV.21.6.3}
Let $\mathrm{Z}=\sum_{x\in\mathrm{X}}n_x.\overline{\{x\}}$ be a cycle on $\mathrm{X}$, $\mathrm{U}$ an open subset of $\mathrm{X}$;
we call the \emph{restriction} of $\mathrm{Z}$ to $\mathrm{U}$ and denote by $\mathrm{Z}|\mathrm{U}$ the cycle $\sum_{x\in\mathrm{U}}n_x.(\mathrm{U}\cap\overline{\{x\}})$ on $\mathrm{U}$;
we have $\Supp(\mathrm{Z}|\mathrm{U})=\Supp(\mathrm{Z})\cap\mathrm{U}$.
It is clear that $\mathrm{Z}\mapsto\mathrm{Z}|\mathrm{U}$ is a homomorphism of ordered groups from $\mathfrak{Z}(\mathrm{X})$ to $\mathfrak{Z}(\mathrm{U})$ (resp.\ from $\mathfrak{Z}^p(\mathrm{X})$ to $\mathfrak{Z}^p(\mathrm{U})$), and that if $\mathrm{V}$ is an open subset contained in $\mathrm{U}$, we have $\mathrm{Z}|\mathrm{V}=(\mathrm{Z}|\mathrm{U})|\mathrm{V}$;
the map $\mathrm{U}\mapsto\mathfrak{Z}(\mathrm{U})$ (resp.\ $\mathrm{U}\mapsto\mathfrak{Z}^{p}(\mathrm{U})$) is therefore a \emph{presheaf} on $\mathrm{X}$ of ordered commutative groups.
In fact, this presheaf is a \emph{sheaf}, as a direct sum of sheaves $(i_x)_*(\bb{Z}_{\overline{\{x\}}})$, where $x$ ranges over $\mathrm{X}$ (resp.\ $\mathrm{X}^{(p)}$) and for each $x\in\mathrm{X}$, $i_x$ is the canonical injection $\overline{\{x\}}\to\mathrm{X}$ and $\bb{Z}_{\overline{\{x\}}}$ is the simple sheaf associated to the constant sheaf $\bb{Z}$ on the space $\overline{\{x\}}$:
this follows at once from the description of the sections of a direct sum of sheaves of commutative groups (G,~II,~2.7).
We denote this sheaf by $\mathscr{Z}_\mathrm{X}$ (resp.\ $\mathscr{Z}_\mathrm{X}^p$).
We denote by $\mathscr{Z}_\mathrm{X}^+$ (resp.\ $\mathscr{Z}_\mathrm{X}^{p+}$) the sub-sheaf
\oldpage[IV-4]{272}
of monoids $\mathrm{U}\mapsto\mathfrak{Z}^+(\mathrm{U})$ (resp.\ $\mathrm{U}\mapsto\mathfrak{Z}^{p+}(\mathrm{U})$) of $\mathscr{Z}_\mathrm{X}$ (resp.\ $\mathscr{Z}_\mathrm{X}^p$).
The sheaf $\mathscr{Z}_\mathrm{X}$ is evidently a \emph{direct sum} of sheaves $\mathscr{Z}_\mathrm{X}^p$.

Let us note finally that the sheaves $\mathscr{Z}_\mathrm{X}^p$ (and therefore also $\mathscr{Z}_\mathrm{X}$) are \emph{flasque}.
Indeed, suppose given a section $\mathrm{Z}$ of $\mathscr{Z}_\mathrm{X}^p$ over an open $\mathrm{U}$, so that the set $\mathrm{S}$ of $x\in\mathrm{U}$ such that $\operatorname{mult}_x(\mathrm{Z})\neq 0$ is locally finite \emph{in} $\mathrm{U}$;
this set is also locally finite \emph{in} $\mathrm{X}$, since for all $z\in\mathrm{X}$ and every open Noetherian neighbourhood $\mathrm{V}$ of $z$, $\mathrm{U}\cap\mathrm{V}$ is Noetherian, and thus contains only a finite number of points of $\mathrm{S}$.
We thus define a section $\mathrm{Z}'$ of $\mathscr{Z}_\mathrm{X}^p$ over $\mathrm{X}$ by setting $\mathrm{Z}'=\sum_{x\in\mathrm{U}}\operatorname{mult}_x(\mathrm{Z}).\overline{\{x\}}$ and since the closure of $x$ in $\mathrm{U}$ is $\overline{\{x\}}\cap\mathrm{U}$, we have $\mathrm{Z}'|\mathrm{U}=\mathrm{Z}$, whence our assertion.
\end{env}

\begin{env}[21.6.4]
\label{IV.21.6.4}
We now propose to define a canonical homomorphism
\[
\label{IV.21.6.4.1}
  c:\sh{D}iv_\mathrm{X}\to\mathscr{Z}_\mathrm{X}^1
  \tag{21.6.4.1}
\]
of sheaves of commutative groups.
It clearly suffices to define a homomorphism of $\sh{M}_\mathrm{X}^*$ to $\mathscr{Z}_\mathrm{X}^1$, which vanishes on $\sh{O}_\mathrm{X}^*$ \sref{IV.21.1.2};
or, by definition, $\sh{M}_\mathrm{X}^*$ is the \emph{symmetrised} monoid sheaf of $\sh{O}_\mathrm{X}\cap\sh{M}_\mathrm{X}^*$, and a homomorphism $\sh{M}_\mathrm{X}^*\to\mathscr{Z}_\mathrm{X}^1$ is uniquely determined by the data of its restriction $\sh{S}(\sh{O}_\mathrm{X})\to\mathscr{Z}_\mathrm{X}^1$, which is a homomorphism of sheaves of \emph{monoids} subject to the sole condition of being zero on $\sh{O}_\mathrm{X}^*$;
it therefore amounts to the same thing to say that, in order to define $c$, it suffices to define a homomorphism of \emph{sheaves of monoids}
\[
\label{IV.21.6.4.2}
  c:\sh{D}iv_\mathrm{X}^+\to\mathscr{Z}_\mathrm{X}^{1+}.
  \tag{21.6.4.2}
\]
\end{env}

\begin{env}[21.6.5]
\label{IV.21.6.5}
Now, we have seen that every positive divisor $\mathrm{D}$ on $\mathrm{X}$ gives rise to the closed sub-prescheme $\mathrm{Y}(\mathrm{D})$ of $\mathrm{X}$, defined by the Ideal $\sh{J}_\mathrm{X}(\mathrm{D})\subset\sh{O}_\mathrm{X}$, and which is regularly immersed and of codimension $1$.
At each of the maximal points $x$ of $\mathrm{Y}(\mathrm{D})$, we have (\sref{IV.19.1.4}{} and \sref{IV.5.1.2}{}) $\codim(\overline{\{x\}},\mathrm{X})=\dim(\sh{O}_{\mathrm{X},x})=1$, in other words $x\in\mathrm{X}^{(1)}$;
moreover the set of these maximal points is locally finite in $\mathrm{X}$ \sref{IV.3.1.6}, and $\sh{O}_{\mathrm{Y}(\mathrm{D}),x}$ is an \emph{Artinian} ring.
At every point $x\in\mathrm{X}^{(1)}$ which is not a maximal point of $\mathrm{Y}(\mathrm{D})$, we necessarily have $x\notin\mathrm{Y}(\mathrm{D})$, so $\sh{O}_{\mathrm{Y}(\mathrm{D}),x}=0$.
Set
\[
\label{IV.21.6.5.1}
  \operatorname{cyc}(\mathrm{D})=\sum_{x\in\mathrm{X}^{(1)}}\operatorname{length}(\sh{O}_{\mathrm{Y}(\mathrm{D}),x}).\overline{\{x\}}
  \tag{21.6.5.1}
\]
which is thus an element of $\mathfrak{Z}^1(\mathrm{X})$.
\end{env}

\begin{proposition}[21.6.6]
\label{IV.21.6.6}
The map $\mathrm{D}\mapsto\operatorname{cyc}(\mathrm{D})$ is a homomorphism of the monoid $\operatorname{Div}^+(\mathrm{X})$ to $\mathfrak{Z}^{1+}(\mathrm{X})$.
\end{proposition}

\begin{proof}
It amounts to showing that for two positive divisors $\mathrm{D}$, $\mathrm{D}'$, we have
\[
  \operatorname{cyc}(\mathrm{D}+\mathrm{D}')=\operatorname{cyc}(\mathrm{D})+\operatorname{cyc}(\mathrm{D}').
\]
Now, by \sref{IV.21.2.5}{} $\sh{J}_\mathrm{X}(\mathrm{D}+\mathrm{D}')=\sh{J}_\mathrm{X}(\mathrm{D})\sh{J}_\mathrm{X}(\mathrm{D}')$;
it thus amounts to showing that if $x\in\mathrm{X}^{(1)}$, setting $\mathrm{A}=\sh{O}_{\mathrm{X},x}$, and if $t$ and $t'$ are two regular elements of $\mathrm{A}$, then $\operatorname{length}(\mathrm{A}/tt'\mathrm{A})=\operatorname{length}(\mathrm{A}/t\mathrm{A})+\operatorname{length}(\mathrm{A}/t'\mathrm{A})$;
but since $t$ is regular, $t\mathrm{A}/tt'\mathrm{A}$ is isomorphic to $\mathrm{A}/t'\mathrm{A}$, whence the proposition at once.
\end{proof}

It follows at once from the definitions that for every open $\mathrm{U}\subset\mathrm{X}$, we have
\[
  \mathrm{Y}(\mathrm{D}|\mathrm{U})=\mathrm{Y}(\mathrm{D})\cap\mathrm{U},
\]
\oldpage[IV-4]{273}
whence $\operatorname{cyc}(\mathrm{D}|\mathrm{U})=\operatorname{cyc}(\mathrm{D})|\mathrm{U}$, and the homomorphisms $\operatorname{cyc}:\operatorname{Div}^+(\mathrm{U})\to\mathfrak{Z}^1(\mathrm{U})$ define a homomorphism of sheaves of monoids \sref{IV.21.6.4.2}, whence the sought-after homomorphism \sref{IV.21.6.4.1}{} of sheaves of groups.

We have $\Supp(\operatorname{cyc}(\mathrm{D}))\subset\Supp(\mathrm{D})$ for every divisor $\mathrm{D}$ and
\begin{env}[21.6.6.2]
\label{IV.21.6.6.2}
\[
  \Supp(\operatorname{cyc}(\mathrm{D}))=\Supp(\mathrm{D})\quad\text{when}\quad\mathrm{D}\geq 0;
\]
\end{env}
we have indeed already seen the second relation \sref{IV.21.2.12};
when $\mathrm{D}$ is arbitrary, the relation $\mathrm{D}_x=0$ implies the existence of an open neighbourhood $\mathrm{U}$ of $x$ such that $\mathrm{D}|\mathrm{U}=0$, whence $\operatorname{cyc}(\mathrm{D})|\mathrm{U}=0$, which proves our assertion.

\begin{env}[21.6.7]
\label{IV.21.6.7}
The homomorphism \sref{IV.21.6.4.1}{} gives, by passage to groups of sections over $\mathrm{X}$, a \emph{homomorphism of ordered groups} $\operatorname{Div}(\mathrm{X})\to\mathfrak{Z}^1(\mathrm{X})$, which we will also denote $\mathrm{D}\mapsto\operatorname{cyc}(\mathrm{D})$;
the number $\operatorname{mult}_x(\operatorname{cyc}(\mathrm{D}))$ for $x\in\mathrm{X}^{(1)}$ is also denoted $\operatorname{mult}_x(\mathrm{D})$ or $\operatorname{mult}_{\overline{\{x\}}}(\mathrm{D})$ and is called the \emph{multiplicity of the divisor $\mathrm{D}$ at the point $x$}, or the \emph{multiplicity of the prime cycle $\overline{\{x\}}$ in $\mathrm{D}$};
it is a positive or negative integer, equal as we have seen to $\operatorname{length}(\sh{O}_{\mathrm{Y}(\mathrm{D}),x})$ when $\mathrm{D}$ is a \emph{positive} divisor;
we thus have by definition
\[
\label{IV.21.6.7.1}
  \operatorname{cyc}(\mathrm{D})=\sum_{x\in\mathrm{X}^{(1)}}\operatorname{mult}_x(\mathrm{D}).\overline{\{x\}},
  \tag{21.6.7.1}
\]
and by virtue of the fact that $\mathrm{D}\mapsto\operatorname{cyc}(\mathrm{D})$ is a homomorphism, we have
\[
\label{IV.21.6.7.2}
  \operatorname{mult}_x(-\mathrm{D})=-\operatorname{mult}_x(\mathrm{D}),\qquad\operatorname{mult}_x(\mathrm{D}+\mathrm{D}')=\operatorname{mult}_x(\mathrm{D})+\operatorname{mult}_x(\mathrm{D}')
  \tag{21.6.7.2}
\]
for any two divisors $\mathrm{D}$, $\mathrm{D}'$.

For every regular meromorphic function $f$ on $\mathrm{X}$, we set in particular, for $x\in\mathrm{X}^{(1)}$
\[
\label{IV.21.6.7.3}
  o_x(f)=\operatorname{mult}_x(\operatorname{div}(f))
  \tag{21.6.7.3}
\]
and we say that $o_x(f)$ (a positive or negative integer) is the \emph{order of $f$ at the point} $x\in\mathrm{X}^{(1)}$.
If $f_x\in\sh{O}_{\mathrm{X},x}$ (a regular element of this local ring by hypothesis), then we have
\[
\label{IV.21.6.7.4}
  o_x(f)=\operatorname{length}(\sh{O}_{\mathrm{X},x}/f_x\sh{O}_{\mathrm{X},x});
  \tag{21.6.7.4}
\]
for two regular meromorphic functions $f$, $f'$ on $\mathrm{X}$, we have
\[
\label{IV.21.6.7.5}
  o_x(ff')=o_x(f)+o_x(f'),\qquad o_x(f^{-1})=-o_x(f)
  \tag{21.6.7.5}
\]
for all $x\in\mathrm{X}^{(1)}$.
The $1$-codimensional cycles
\[
  \mathrm{Z}^+(f)=\sum_{x\in\mathrm{X}^{(1)}}(o_x(f))^+.\overline{\{x\}},\qquad
  \mathrm{Z}^-(f)=\sum_{x\in\mathrm{X}^{(1)}}(o_x(f))^-.\overline{\{x\}}
\]
are called respectively the \emph{cycle of zeros} and the \emph{cycle of poles} (or \emph{polar cycle}) of $f$, and we evidently have $\operatorname{cyc}(\operatorname{div}(f))=\mathrm{Z}^+(f)-\mathrm{Z}^-(f)$.
The $1$-codimensional cycles of the form $\operatorname{cyc}(\operatorname{div}(f))$ are called \emph{principal} (or \emph{linearly equivalent to zero}) and form a subgroup $\mathfrak{Z}^1_\mathrm{princ}(\mathrm{X})$ of $\mathfrak{Z}^1(\mathrm{X})$.
The sections over $\mathrm{X}$ of the image $c(\sh{D}iv_\mathrm{X})\subset\mathscr{Z}_\mathrm{X}^1$ are
\oldpage[IV-4]{274}
called \emph{locally principal cycles};
such a cycle $\mathrm{Z}$ is therefore characterised by the fact that, for all $y\in\mathrm{X}$, there is an open neighbourhood $\mathrm{U}$ of $y$ in $\mathrm{X}$ and a regular meromorphic function $f$ on $\mathrm{U}$ such that $\mathrm{Z}|\mathrm{U}=\sum_{x\in\mathrm{U}\cap\mathrm{X}^{(1)}}o_x(f)(\overline{\{x\}}\cap\mathrm{U})$.
If $\mathrm{Z}=\sum_{x\in\mathrm{X}^{(1)}}n_x.\overline{\{x\}}$, it is equivalent to saying that, for every $y\in\mathrm{X}$, setting $\mathrm{T}_y=\Spec(\sh{O}_{\mathrm{X},y})$ and $\mathrm{Z}_y=\sum_{x\in\mathrm{X}^{(1)}\cap\mathrm{T}_y}n_x(\overline{\{x\}}\cap\mathrm{T}_y)$, $\mathrm{Z}_y$ is \emph{principal}, otherwise said there exists a regular meromorphic function $g$ on $\mathrm{T}_y$ such that $\mathrm{Z}_y=\operatorname{cyc}(\operatorname{div}(g))$.
This follows at once from the definition above and from \sref{IV.20.3.7}{} which establishes a bijective correspondence between the regular meromorphic functions on $\mathrm{T}_y$ and the germs of regular meromorphic functions on the open neighbourhoods of $y$ in $\mathrm{X}$ when $\mathrm{X}$ is locally Noetherian.
We also express the fact that $\mathrm{Z}_y$ is principal by saying that $\mathrm{Z}$ is \emph{principal at the point $y$}.
The set of points where $\mathrm{Z}$ is principal is evidently \emph{open}, by virtue of what precedes.

We deduce from the canonical homomorphism $\operatorname{cyc}:\operatorname{Div}(\mathrm{X})\to\mathfrak{Z}^1(\mathrm{X})$ a homomorphism $\operatorname{Div}(\mathrm{X})/\operatorname{Div.princ}(\mathrm{X})\to\mathfrak{Z}^1(\mathrm{X})/\mathfrak{Z}^1_\mathrm{princ}(\mathrm{X})$, by virtue of the definition of $\mathfrak{Z}^1_\mathrm{princ}(\mathrm{X})$.
We say that the group $\mathfrak{Z}^1(\mathrm{X})/\mathfrak{Z}^1_\mathrm{princ}(\mathrm{X})$ is the \emph{group of classes of $1$-codimensional cycles} of $\mathrm{X}$ and we denote it $\mathfrak{Cl}(\mathrm{X})$.
Two elements of $\mathfrak{Z}^1(\mathrm{X})$ having the same image in $\mathfrak{Cl}(\mathrm{X})$ are called \emph{linearly equivalent}.
\end{env}

\begin{env}[21.6.8]
\label{IV.21.6.8}
Let us consider in particular the case where $\mathrm{X}=\Spec(\mathrm{A})$, where $\mathrm{A}$ is a \emph{Noetherian integral and integrally closed} ring.
Then $\mathrm{X}^{(1)}$ is the set of \emph{prime ideals of height $1$} of $\mathrm{A}$, and $\mathfrak{Z}^1(\mathrm{X})$ is identified with the \emph{group of divisors} (in the sense of N.~Bourbaki) of the Krull ring $\mathrm{A}$ (Bourbaki, \emph{Alg.\ comm.}, chap.~VII, \textsection~1, n\textsuperscript{o}~3, cor.\ du th.~2 et n\textsuperscript{o}~6, th.~3).

On the other hand, the regular meromorphic functions on $\mathrm{X}$ are identified with the nonzero elements of the field of fractions $\mathrm{K}$ of $\mathrm{A}$, and the map $f\mapsto\operatorname{cyc}(\operatorname{div}(f))$ from $\mathrm{M}(\mathrm{X})$ to $\mathfrak{Z}^1(\mathrm{X})$ is identified with the map denoted $f\mapsto\operatorname{div}(f)$ in Bourbaki (\emph{loc.\ cit.}, \textsection~1, n\textsuperscript{o}~1);
$\mathfrak{Z}^1_\mathrm{princ}(\mathrm{X})$ is identified with the \emph{group of principal divisors} of $\mathrm{A}$ in the sense of Bourbaki, and $\mathfrak{Cl}(\mathrm{X})$ with the \emph{group of classes of divisors} of $\mathrm{A}$ in the sense of Bourbaki (\emph{loc.\ cit.}, \textsection~1, n\textsuperscript{o}~2 et n\textsuperscript{o}~10).
\end{env}

\begin{theorem}[21.6.9]
\label{IV.21.6.9}
Let $\mathrm{X}$ be a locally Noetherian normal prescheme.
\begin{enumerate}[label=(\roman*)]
  \item The canonical homomorphism $\operatorname{cyc}:\operatorname{Div}(\mathrm{X})\to\mathfrak{Z}^1(\mathrm{X})$ is injective and its image is formed of locally principal cycles.
  \item The following conditions are equivalent:
  \begin{enumerate}[label=\emph{\alph*)}]
    \item The homomorphism $\operatorname{cyc}:\operatorname{Div}(\mathrm{X})\to\mathfrak{Z}^1(\mathrm{X})$ is bijective.
    \item Every $1$-codimensional cycle is locally principal.
    \item For all $x\in\mathrm{X}$, the local ring $\sh{O}_{\mathrm{X},x}$ is factorial.
  \end{enumerate}
  (We then say that $\mathrm{X}$ is a \emph{locally factorial} prescheme.)
\end{enumerate}
\end{theorem}

\begin{proof}
\begin{enumerate}[label=(\roman*)]
  \item It suffices to show that
\[
\label{IV.21.6.9.1}
  \operatorname{cyc}^{-1}(\mathfrak{Z}^{1+}(\mathrm{X}))=\operatorname{Div}^+(\mathrm{X}),
  \tag{21.6.9.1}
\]
since we have $\operatorname{Div}^+(\mathrm{X})\cap(-\operatorname{Div}^+(\mathrm{X}))=0$ and $\mathfrak{Z}^{1+}(\mathrm{X})\cap(-\mathfrak{Z}^{1+}(\mathrm{X}))=0$.
We are thus reduced to proving that if $\mathrm{D}$ is a divisor on $\mathrm{X}$ such that $\operatorname{mult}_x(\mathrm{D})\geq 0$ for all $x\in\mathrm{X}^{(1)}$, then $\mathrm{D}\geq 0$.
Now, for all $x\notin\mathrm{X}^{(1)}$, the local ring $\sh{O}_{\mathrm{X},x}$ is integral and integrally closed and
\oldpage[IV-4]{275}
of dimension $0$ or $\geq 2$, therefore on the one hand $\operatorname{prof}(\sh{O}_{\mathrm{X},x})=0$ or $\operatorname{prof}(\sh{O}_{\mathrm{X},x})\geq 2$ \sref[0]{0.16.5.1}.
On the other hand, for points $x\in\mathrm{X}^{(1)}$, the ring $\sh{O}_{\mathrm{X},x}$ is a discrete valuation ring (\sref[II]{II.7.1.6}),
whence $\operatorname{prof}(\sh{O}_{\mathrm{X},x})=1$ \sref[0]{0.16.5.1};
the only points of $\mathrm{X}$ such that $\operatorname{prof}(\sh{O}_{\mathrm{X},x})=1$ are thus the points of $\mathrm{X}^{(1)}$, and the conclusion follows from \sref{IV.21.1.8}.

  \item The equivalence of \emph{a)} and \emph{b)} is clear by virtue of (i).
From the characterisation of locally principal cycles given in \sref{IV.21.6.7}, and the relation between $1$-codimensional cycles on $\Spec(\mathrm{A})$ and divisors (in the sense of Bourbaki) of the ring $\mathrm{A}$ when $\mathrm{A}$ is a Noetherian integrally closed ring \sref{IV.21.6.8}, condition \emph{b)} is equivalent to saying that for all $x\in\mathrm{X}$, every divisor of the ring $\sh{O}_{\mathrm{X},x}$ is principal, in other words that the ring $\sh{O}_{\mathrm{X},x}$ is factorial (Bourbaki, \emph{Alg.\ comm.}, chap.~VII, \textsection~3, n\textsuperscript{o}~1), whence the equivalence of \emph{b)} and \emph{c)}.
\end{enumerate}
\end{proof}

\begin{env}[21.6.9.2]
\label{IV.21.6.9.2}
When $\mathrm{A}$ is a Noetherian \emph{factorial} ring, it is clear that $\mathrm{X}=\Spec(\mathrm{A})$ is locally factorial (Bourbaki, \emph{Alg.\ comm.}, chap.~VII, \textsection~3, n\textsuperscript{o}~4, prop.~3).
In this case, if one writes an element $f\neq 0$ of the field of fractions $\mathrm{K}$ of $\mathrm{A}$ in the form $r/s$, where $r$ and $s$ are two \emph{coprime} elements of $\mathrm{A}$, whose divisors are well-determined (\emph{loc.\ cit.}, \textsection~3, n\textsuperscript{o}~3), these divisors are identified respectively with the \emph{divisor of zeros} and the \emph{divisor of poles} of $f$ \sref{IV.21.6.7}.
\end{env}

\begin{corollary}[21.6.10]
\label{IV.21.6.10}
Let $\mathrm{X}$ be a locally Noetherian normal prescheme.
\begin{enumerate}[label=(\roman*)]
  \item There exists a canonical injective homomorphism
\[
\label{IV.21.6.10.1}
  \operatorname{Pic}(\mathrm{X})\to\mathfrak{Cl}(\mathrm{X}).
  \tag{21.6.10.1}
\]

  \item If $\mathrm{X}$ is locally factorial, the homomorphism \sref{IV.21.6.10.1}{} is bijective, and conversely.
\end{enumerate}
\end{corollary}

\begin{proof}
We have seen \sref{IV.21.6.7}{} that the image of $\operatorname{Div.princ}(\mathrm{X})$ by the homomorphism $\operatorname{cyc}:\operatorname{Div}(\mathrm{X})\to\mathfrak{Z}^1(\mathrm{X})$ is $\mathfrak{Z}^1_\mathrm{princ}(\mathrm{X})$;
we thus deduce from \sref{IV.21.6.9}{} that the homomorphism $\operatorname{Div}(\mathrm{X})/\operatorname{Div.princ}(\mathrm{X})\to\mathfrak{Z}^1(\mathrm{X})/\mathfrak{Z}^1_\mathrm{princ}(\mathrm{X})=\mathfrak{Cl}(\mathrm{X})$ deduced from $\operatorname{cyc}$ is injective, and that it is bijective if and only if $\mathrm{X}$ is locally factorial.
The conclusion follows from the fact that, when $\mathrm{X}$ is locally Noetherian and reduced, the canonical homomorphism $\operatorname{Div}(\mathrm{X})/\operatorname{Div.princ}(\mathrm{X})\to\operatorname{Pic}(\mathrm{X})$ \sref{IV.21.3.3.1}{} is bijective \sref{IV.21.3.4},{}~(ii).
\end{proof}

\begin{corollary}[21.6.11]
\label{IV.21.6.11}
Let $\mathrm{X}$ be a locally Noetherian prescheme which is locally factorial.
Then the sheaf $\sh{D}iv_\mathrm{X}$ is flasque, and for every open $\mathrm{U}$ of $\mathrm{X}$, the canonical homomorphism $\operatorname{Pic}(\mathrm{X})\to\operatorname{Pic}(\mathrm{U})$ is surjective.
\end{corollary}

\begin{proof}
Taking into account \sref{IV.21.6.9},{}~(ii), the first assertion is equivalent to saying that the sheaf $\mathscr{Z}_\mathrm{X}^1$ is flasque, which was seen above \sref{IV.21.6.3}.
For every open $\mathrm{U}$ of $\mathrm{X}$, the canonical homomorphism $\mathfrak{Z}^1(\mathrm{X})\to\mathfrak{Z}^1(\mathrm{U})$ is thus surjective;
as in \sref{IV.21.6.10}, the homomorphism $\operatorname{Pic}(\mathrm{X})\to\operatorname{Pic}(\mathrm{U})$ is identified canonically with $\mathfrak{Z}^1(\mathrm{X})/\mathfrak{Z}^1_\mathrm{princ}(\mathrm{X})\to\mathfrak{Z}^1(\mathrm{U})/\mathfrak{Z}^1_\mathrm{princ}(\mathrm{U})$, it is also surjective.
\end{proof}

\begin{proposition}[21.6.12]
\label{IV.21.6.12}
Let $\mathrm{X}$ be a Noetherian reduced prescheme.
Let $(\mathrm{U}_\lambda)_{\lambda\in\mathrm{L}}$ be a decreasing filtered family of opens of $\mathrm{X}$ satisfying the following conditions:
\begin{enumerate}
  \item[\rm 1\textsuperscript{o}] If $\mathrm{Y}_\lambda=\mathrm{X}-\mathrm{U}_\lambda$, then $\codim(\mathrm{Y}_\lambda,\mathrm{X})\geq 2$ for all $\lambda\in\mathrm{L}$.
  \item[\rm 2\textsuperscript{o}] For all $x\in\bigcap_{\lambda\in\mathrm{L}}\mathrm{U}_\lambda$, the ring $\sh{O}_{\mathrm{X},x}$ is factorial.
\end{enumerate}
Then there exist canonical isomorphisms
\[
\label{IV.21.6.12.1}
  \varinjlim\operatorname{Div}(\mathrm{U}_\lambda)\simto\mathfrak{Z}^1(\mathrm{X}),\qquad
  \varinjlim\operatorname{Pic}(\mathrm{U}_\lambda)\simto\mathfrak{Cl}(\mathrm{X})
  \tag{21.6.12.1}
\]
such that, for every open $\mathrm{V}$ of $\mathrm{X}$, the diagrams
\[
\label{IV.21.6.12.2}
  \xymatrix{
    \varinjlim\operatorname{Div}(\mathrm{U}_\lambda) \ar[r]^-\sim \ar[d] & \mathfrak{Z}^1(\mathrm{X}) \ar[d] & & \varinjlim\operatorname{Pic}(\mathrm{U}_\lambda) \ar[r]^-\sim \ar[d] & \mathfrak{Cl}(\mathrm{X}) \ar[d]\\
    \varinjlim\operatorname{Div}(\mathrm{U}_\lambda\cap\mathrm{V}) \ar[r]^-\sim & \mathfrak{Z}^1(\mathrm{V}) & & \varinjlim\operatorname{Pic}(\mathrm{U}_\lambda\cap\mathrm{V}) \ar[r]^-\sim & \mathfrak{Cl}(\mathrm{V})
  }
  \tag{21.6.12.2}
\]
are commutative.
\end{proposition}

\begin{proof}
The hypothesis 1\textsuperscript{o} on $\mathrm{U}_\lambda$ implies that for all $\lambda\in\mathrm{L}$, we have $\mathrm{X}^{(1)}\subset\mathrm{U}_\lambda$ \sref{IV.5.1.3.1}{} and therefore the restriction homomorphism $\mathfrak{Z}^1(\mathrm{X})\to\mathfrak{Z}^1(\mathrm{U}_\lambda)$ is bijective and consequently we have an isomorphism $\varinjlim\mathfrak{Z}^1(\mathrm{U}_\lambda)\simto\mathfrak{Z}^1(\mathrm{X})$.
The canonical homomorphisms $\operatorname{cyc}:\operatorname{Div}(\mathrm{U}_\lambda)\to\mathfrak{Z}^1(\mathrm{U}_\lambda)$ thus define, by passage to the inductive limit, the first of the canonical homomorphisms \sref{IV.21.6.12.1}, and the second is deduced by passage to quotients.
Furthermore, it follows from condition 1\textsuperscript{o} that the $\mathrm{U}_\lambda$ are schematically dense in $\mathrm{X}$ since $\mathrm{X}$ is reduced \sref{IV.11.10.4}, and therefore every meromorphic function on $\mathrm{X}$ is entirely determined by its restriction to each $\mathrm{U}_\lambda$;
we deduce from this that in the isomorphism $\mathfrak{Z}^1(\mathrm{X})\simto\mathfrak{Z}^1(\mathrm{U}_\lambda)$, the image of $\mathfrak{Z}^1_\mathrm{princ}(\mathrm{X})$ is $\mathfrak{Z}^1_\mathrm{princ}(\mathrm{U}_\lambda)$, and therefore we also have a canonical isomorphism $\varinjlim\mathfrak{Cl}(\mathrm{U}_\lambda)\simto\mathfrak{Cl}(\mathrm{X})$.
The second of the canonical homomorphisms \sref{IV.21.6.12.1}{} will therefore be an isomorphism if the first is, and it remains to prove this last point, the commutativity of the diagrams \sref{IV.21.6.12.2}{} being trivial.

\oldpage[IV-4]{276}
Let us first show that the homomorphism $\varinjlim\operatorname{Div}(\mathrm{U}_\lambda)\to\mathfrak{Z}^1(\mathrm{X})$ is \emph{injective}.
Set $\mathrm{T}=\bigcap_\lambda\mathrm{U}_\lambda$, and note that the $\mathrm{U}_\lambda$ form a \emph{fundamental system of neighbourhoods} of $\mathrm{T}$.
Indeed, for every open $\mathrm{V}\supset\mathrm{T}$, $\mathrm{X}-\mathrm{V}$ is a closed subspace of $\mathrm{X}$, thus a Noetherian space whose every irreducible closed subspace admits a generic point;
as the sets $(\mathrm{X}-\mathrm{V})\cap(\mathrm{X}-\mathrm{U}_\lambda)$ are closed in $\mathrm{X}-\mathrm{V}$, they form an increasing filtered family and their union is $\mathrm{X}-\mathrm{V}$;
our assertion follows from \sref[0]{0.9.2.4}.
This being so, it will suffice to show that if $\mathrm{D}\in\operatorname{Div}(\mathrm{U}_\lambda)$ is such that $\operatorname{cyc}(\mathrm{D})=0$ in $\mathfrak{Z}^1(\mathrm{U}_\lambda)$, then there exists $\mu\geq\lambda$ such that $\mathrm{D}|\mathrm{U}_\mu=0$.
By virtue of what precedes, it will suffice to show that for all $x\in\mathrm{T}$, $\mathrm{D}_x=0$;
indeed, if this holds, there will be an open neighbourhood $\mathrm{W}(x)$ of $x$ in $\mathrm{X}$ such that $\mathrm{D}|\mathrm{W}(x)=0$, and the union of the $\mathrm{W}(x)$ for $x\in\mathrm{T}$ contains a $\mathrm{U}_\mu$.
Taking into account \sref{IV.21.4.6}, we are thus reduced to the case where $\mathrm{X}=\Spec(\sh{O}_{\mathrm{X},x})$ with $x\in\mathrm{T}$;
but by hypothesis $\sh{O}_{\mathrm{X},x}$ is factorial, thus integral and integrally closed, and $\Spec(\sh{O}_{\mathrm{X},x})$ is therefore normal;
whence the conclusion follows from \sref{IV.21.6.9},{}~(i).

\oldpage[IV-4]{277}
To prove that the homomorphism $\varinjlim\operatorname{Div}(\mathrm{U}_\lambda)\to\mathfrak{Z}^1(\mathrm{X})$ is bijective, it will suffice to prove that for all $\mathrm{Z}\in\mathfrak{Z}^1(\mathrm{X})$ and all $x\in\mathrm{T}$, there is a neighbourhood $\mathrm{W}(x)$ of $x$ in $\mathrm{X}$ and a divisor $\mathrm{D}^{(x)}$ on $\mathrm{W}(x)$ such that $\operatorname{cyc}(\mathrm{D}^{(x)})=\mathrm{Z}|\mathrm{W}(x)$.
Indeed, we can then cover the quasi-compact set $\mathrm{T}$ by a finite number of the $\mathrm{W}(x_i)$, and by the first part of the proof, since the pairs $(i,j)$ are finite in number, there will exist an index $\lambda$ such that in each of the $\mathrm{W}(x_i)\cap\mathrm{W}(x_j)\cap\mathrm{U}_\lambda$, the restrictions of $\mathrm{D}^{(x_i)}$ and $\mathrm{D}^{(x_j)}$ coincide;
since we can moreover assume that $\mathrm{U}_\lambda$ is contained in the union of the $\mathrm{W}(x_i)$, we see that the restrictions $\mathrm{D}^{(x_i)}|(\mathrm{W}(x_i)\cap\mathrm{U}_\lambda)$ are the restrictions of a single divisor $\mathrm{D}\in\operatorname{Div}(\mathrm{U}_\lambda)$ which will be such that $\operatorname{cyc}(\mathrm{D})=\mathrm{Z}|\mathrm{U}_\lambda$.
We are thus reduced again to the case where $\mathrm{X}=\Spec(\sh{O}_{\mathrm{X},x})$ with $x\in\mathrm{T}$;
but as $\sh{O}_{\mathrm{X},x}$ is factorial, the same is true of its localizations (Bourbaki, \emph{Alg.\ comm.}, chap.~VII, \textsection~3, n\textsuperscript{o}~4, prop.~3), and it suffices to apply \sref{IV.21.6.9},{}~(ii).
\end{proof}

\begin{corollary}[21.6.13]
\label{IV.21.6.13}
Let $\mathrm{A}$ be a Noetherian local ring which is integral and integrally closed and of dimension $\geq 2$.
Set $\mathrm{X}=\Spec(\mathrm{A})$, and let $a$ be the closed point of $\mathrm{X}$, $\mathrm{U}=\mathrm{X}-\{a\}$.
For $\mathrm{A}$ to be factorial, it is necessary and sufficient that $\mathrm{U}$ is locally factorial and that $\operatorname{Pic}(\mathrm{U})=0$.
\end{corollary}

\begin{proof}
Indeed, to say that $\mathrm{A}$ is factorial is equivalent to saying that $\mathfrak{Cl}(\mathrm{X})=0$ (Bourbaki, \emph{Alg.\ comm.}, chap.~VII, \textsection~1, n\textsuperscript{o}~4, cor.\ du th.~2 et \textsection~3, n\textsuperscript{o}~2, th.~1);
it therefore suffices to use the existence of the second isomorphism \sref{IV.21.6.12.1}, taking the family $(\mathrm{U}_\lambda)$ restricted to the single open $\mathrm{U}$.
\end{proof}

\begin{corollary}[21.6.14]
\label{IV.21.6.14}
Let $\mathrm{A}$ be a Noetherian local ring of dimension $\geq 2$, $\mathrm{X}=\Spec(\mathrm{A})$, $a$ the closed point of $\mathrm{X}$, $\mathrm{U}=\mathrm{X}-\{a\}$.
The following conditions are equivalent:
\begin{enumerate}[label=\emph{\alph*)}]
  \item $\mathrm{A}$ is factorial.
  \item $\operatorname{Pic}(\mathrm{U})=0$ and $\operatorname{prof}(\mathrm{A})\geq 2$ (conditions which we will express later \sref{IV.21.13}{} by saying that the ring $\mathrm{A}$ is \emph{parafactorial}), and $\mathrm{U}$ is locally factorial.
\end{enumerate}
\end{corollary}

\begin{proof}
It is clear that if $\mathrm{A}$ is factorial, it is normal, hence $\operatorname{prof}(\mathrm{A})\geq 2$ since $\dim(\mathrm{A})\geq 2$ \sref[0]{0.16.5.1}, and \sref{IV.21.6.13}{} shows that \emph{a)} implies \emph{b)}.
Conversely, if \emph{b)} is satisfied, it suffices to see that $\mathrm{A}$ is integrally closed to conclude by \sref{IV.21.6.13}{} that \emph{b)} implies \emph{a)}.
By virtue of the Serre criterion \sref{IV.5.8.6}, it suffices to verify for $\mathrm{X}$ the conditions $(\mathrm{R}_1)$ and $(\mathrm{S}_2)$.
Now, $\mathrm{U}$ being locally factorial satisfies these conditions, and the hypothesis $\operatorname{prof}(\mathrm{A})\geq 2$ implies that $\mathrm{X}$ also satisfies them.
\end{proof}


\subsection{Interpretation of positive 1-codimensional cycles in terms of sub-preschemes}
\label{subsection:IV.21.7}
\label{IV.21.7}

\begin{env}[21.7.1]
\label{IV.21.7.1}
Let $\mathrm{X}$ be a locally Noetherian prescheme, $\mathrm{C}=\sum_{x\in\mathrm{X}^{(1)}}n_x.\overline{\{x\}}$ a positive $1$-codimensional cycle (so that $n_x\geq 0$ for all $x\in\mathrm{X}^{(1)}$, and $n_x=0$ except at a locally finite set of points).
Let $\mathrm{Y}(\mathrm{C})$ denote the closed sub-prescheme of $\mathrm{X}$, \emph{closed image} \sref[I]{I.9.5.3}{} and \sref{IV.9.5.1}{} of the sub-prescheme \emph{sum} $\mathrm{Y}'(\mathrm{C})=\coprod_{x\in\mathrm{X}^{(1)}}\Spec(\sh{O}_{\mathrm{X},x}/\mathfrak{m}_x^{n_x})$ by the canonical morphism, and let $\sh{J}_\mathrm{X}(\mathrm{C})$ (or $\sh{J}(\mathrm{C})$) denote the Ideal of $\sh{O}_\mathrm{X}$ defining $\mathrm{Y}(\mathrm{C})$.
\end{env}

\begin{proposition}[21.7.2]
\label{IV.21.7.2}
Let $\mathrm{X}$ be a locally Noetherian prescheme satisfying the condition $(\mathrm{R}_1)$ \sref{IV.5.8.2}.
For a closed sub-prescheme $\mathrm{Y}$ of $\mathrm{X}$ to be of the form $\mathrm{Y}(\mathrm{C})$, where $\mathrm{C}$ is a positive $1$-codimensional cycle, it is necessary and sufficient that it satisfies the two following conditions:
\begin{enumerate}[label=(\roman*)]
  \item $\mathrm{Y}$ is purely of codimension $1$.
\oldpage[IV-4]{278}
  \item $\mathrm{Y}$ satisfies the condition $(\mathrm{S}_1)$, in other words \sref{IV.5.7.5}{} it has no embedded associated prime cycle.
\end{enumerate}
We then have
\[
\label{IV.21.7.2.1}
  \operatorname{mult}_x(\mathrm{C})=\operatorname{length}\,\sh{O}_{\mathrm{Y},x}.
  \tag{21.7.2.1}
\]
The map $\mathrm{C}\mapsto\mathrm{Y}(\mathrm{C})$ is a bijection from $\mathfrak{Z}^{1+}(\mathrm{X})$ onto the set of closed sub-preschemes of $\mathrm{X}$ satisfying conditions (i) and (ii).
\end{proposition}

\begin{proof}
The conditions are necessary (without assuming that $\mathrm{X}$ satisfies $(\mathrm{R}_1)$).
Indeed, since the question is local on $\mathrm{X}$, we can assume that $\mathrm{X}=\Spec(\mathrm{A})$, where $\mathrm{A}$ is a Noetherian ring;
we then have $\mathrm{Y}(\mathrm{C})=\Spec(\mathrm{A}/\mathfrak{J})$, where $\mathfrak{J}$ is, by definition \sref[I]{I.9.5.1}, the kernel of the canonical homomorphism $\mathrm{A}\to\bigoplus_i\mathrm{A}_{\mathfrak{p}_i}/\mathfrak{p}_i^{n_i}\mathrm{A}_{\mathfrak{p}_i}$, where the $\mathfrak{p}_i$ are the prime ideals corresponding to points $x\in\mathrm{X}^{(1)}$ for which $n_x>0$, and we have set $n_i=n_x$.
The inverse image $\mathfrak{q}_i$ of $\mathfrak{p}_i^{n_i}\mathrm{A}_{\mathfrak{p}_i}$ in $\mathrm{A}$ is a $\mathfrak{p}_i$-primary ideal and we have $\mathfrak{J}=\bigcap_i\mathfrak{q}_i$.
Moreover, as the $x\in\mathrm{X}^{(1)}$ are such that $\overline{\{x\}}$ is of codimension $1$, no point of $\mathrm{X}^{(1)}$ can be adherent to another point of $\mathrm{X}^{(1)}$, so the $\mathfrak{p}_i$ are minimal primes of $\mathrm{A}/\mathfrak{J}$ and the set of $\mathfrak{p}_i$ is equal to $\operatorname{Ass}(\mathrm{A}/\mathfrak{J})$, which proves conditions (i) and (ii).

The conditions are sufficient.
Let $\sh{J}$ denote the Ideal of $\sh{O}_\mathrm{X}$ defining $\mathrm{Y}$.
Hypothesis (ii) implies that $\operatorname{Ass}(\sh{O}_\mathrm{X}/\sh{J})$ is identical to the set $\mathrm{F}$ of maximal points of $\mathrm{Y}$, and hypothesis (i) implies that $\mathrm{F}$ is contained in $\mathrm{X}^{(1)}$;
hence \sref{IV.3.2.6}{} $\sh{O}_\mathrm{X}/\sh{J}$ is identified with a sub-$\sh{O}_\mathrm{X}$-Module of $\bigoplus_{x\in\mathrm{F}}\mathscr{G}(x)$, where, for every $x\in\mathrm{F}$, $\mathscr{G}(x)$ is an irredundant $\sh{O}_\mathrm{X}$-Module such that $\operatorname{Ass}(\mathscr{G}(x))=\{x\}$.
Now, since $\mathrm{X}$ satisfies $(\mathrm{R}_1)$, for every $x\in\mathrm{F}$ the ring $\sh{O}_{\mathrm{X},x}$ is a discrete valuation ring, and consequently the nonzero primary ideals of this ring are the powers $\mathfrak{m}_x^{n_x}$ of its maximal ideal.
Supposing again that $\mathrm{X}=\Spec(\mathrm{A})$, we conclude that $\mathfrak{J}=\widetilde{\mathfrak{J}}$, where $\widetilde{\mathfrak{J}}=\bigcap_i\mathfrak{q}_i$ and the $\mathfrak{q}_i$ are the inverse images in $\mathrm{A}$ of the ideals $\mathfrak{p}_i^{n_i}\mathrm{A}_{\mathfrak{p}_i}$, the $\mathfrak{p}_i$ corresponding to the points of $\mathrm{F}$.
Thus $\mathrm{Y}$ is indeed of the form $\mathrm{Y}(\mathrm{C})$.
\end{proof}

\begin{corollary}[21.7.3]
\label{IV.21.7.3}
Let $\mathrm{X}$ be a locally Noetherian prescheme.
For every positive divisor $\mathrm{D}$ on $\mathrm{X}$ such that $\mathrm{X}$ satisfies $(\mathrm{R}_1)$ at the maximal points of the closed sub-prescheme $\mathrm{Y}(\mathrm{D})$ \sref{IV.21.2.12}, $\mathrm{Y}(\mathrm{D})$ majorises the closed sub-prescheme $\mathrm{Y}(\operatorname{cyc}(\mathrm{D}))$ \sref{IV.21.7.1}, and the two have the same underlying space;
for these two sub-preschemes to be equal, it is necessary and sufficient that $\mathrm{Y}(\mathrm{D})$ satisfies the condition $(\mathrm{S}_1)$, or else, for every $z\in\mathrm{Y}(\mathrm{D})$ distinct from a maximal point of $\mathrm{Y}(\mathrm{D})$, that $\operatorname{prof}(\sh{O}_{\mathrm{X},z})\geq 2$ (a condition always satisfied when $\mathrm{X}$ is \emph{normal}).
\end{corollary}

\begin{proof}
The question being local, we may assume that $\sh{J}_\mathrm{X}(\mathrm{D})=t\sh{O}_\mathrm{X}$, where $t$ is a regular section of $\sh{O}_\mathrm{X}$ over $\mathrm{X}$.
At every maximal point $x$ of $\mathrm{Y}(\mathrm{D})$, which necessarily belongs to $\mathrm{X}^{(1)}$, the ring $\sh{O}_{\mathrm{X},x}$ is by hypothesis a discrete valuation ring, and therefore $t_x\sh{O}_{\mathrm{X},x}=\mathfrak{m}_x^{n_x}$, where $n_x=\operatorname{mult}_x(\mathrm{D})$.
We may suppose that $\mathrm{X}=\Spec(\mathrm{A})$; if $\sh{J}_\mathrm{X}(\operatorname{cyc}(\mathrm{D}))=\widetilde{\mathfrak{J}}$, it follows from what precedes and from \sref{IV.21.7.1}{} that $\mathfrak{J}=t\mathrm{A}$ and $\widetilde{\mathfrak{J}}$ have the same \emph{non-embedded primary ideals}, whence $\widetilde{\mathfrak{J}}\subset\mathfrak{J}$, since $\widetilde{\mathfrak{J}}$ is the intersection of
\oldpage[IV-4]{279}
these primary ideals \sref{IV.21.7.2}.
This proves that $\mathrm{Y}(\mathrm{D})$ majorises $\mathrm{Y}(\operatorname{cyc}(\mathrm{D}))$ and that these two sub-preschemes are equal if and only if $\mathrm{Y}(\mathrm{D})$ has no associated immersed prime cycle.
For every $x\in\mathrm{Y}(\mathrm{D})$, there is by hypothesis an open neighbourhood $\mathrm{U}$ of $x$ in $\mathrm{X}$ and a regular element $t\in\Gamma(\mathrm{U},\sh{O}_\mathrm{X})$ such that $\sh{O}_{\mathrm{Y}(\mathrm{D})}|(\mathrm{U}\cap\mathrm{Y}(\mathrm{D}))$ is the restriction to $\mathrm{Y}(\mathrm{D})\cap\mathrm{U}$ of $\sh{O}_\mathrm{U}/t\sh{O}_\mathrm{U}$.
Thus, to say that $\mathrm{Y}(\mathrm{D})$ satisfies $(\mathrm{S}_1)$ means that at every non-maximal point $z\in\mathrm{Y}(\mathrm{D})$ one has $\operatorname{prof}(\sh{O}_{\mathrm{X},z}/t_z\sh{O}_{\mathrm{X},z})\geq 1$, or equivalently \sref[0]{0.16.4.6}{} that $\operatorname{prof}(\sh{O}_{\mathrm{X},z})\geq 2$.
The assertion for normal $\mathrm{X}$ follows immediately, since $\mathrm{X}$ satisfies $(\mathrm{S}_2)$ and, at every non-maximal point $z$ of $\mathrm{Y}(\mathrm{D})$, one has $\dim(\sh{O}_{\mathrm{X},z})\geq 2$ \sref[0]{0.16.3.4}.
\end{proof}

\begin{env}[21.7.3.1]
\label{IV.21.7.3.1}
We see therefore that when $\mathrm{X}$ is \emph{normal}, $\operatorname{Div}^+(\mathrm{X})$ is identified canonically with the set of closed sub-preschemes $\mathrm{Y}$ of $\mathrm{X}$ satisfying conditions (i) and (ii) of \sref{IV.21.7.2}{} and \emph{regularly immersed} in $\mathrm{X}$.
\end{env}

\begin{proposition}[21.7.4]
\label{IV.21.7.4}
Let $\mathrm{X}$ be a locally Noetherian prescheme which is reduced at each of its isolated points.
The following conditions are equivalent:
\begin{enumerate}
  \item[\emph{a)}] The canonical homomorphism $c:\sh{D}iv_\mathrm{X}\to\mathscr{Z}_\mathrm{X}^1$ \sref{IV.21.6.4}{} is an isomorphism of sheaves of ordered groups.
  \item[\emph{a$'$)}] Every prime $1$-codimensional cycle on $\mathrm{X}$ is the image under $\operatorname{cyc}$ of a positive divisor, and the canonical homomorphism $c:\sh{D}iv_\mathrm{X}\to\mathscr{Z}_\mathrm{X}^1$ is injective.
  \item[\emph{a$''$)}] For every integral closed sub-prescheme $\mathrm{Y}$ of $\mathrm{X}$ of codimension $1$, the canonical immersion $\mathrm{Y}\to\mathrm{X}$ is regular.
  \item[\emph{b)}] $\mathrm{X}$ is normal and the homomorphism $\operatorname{cyc}:\operatorname{Div}(\mathrm{X})\to\mathfrak{Z}^1(\mathrm{X})$ is bijective.
  \item[\emph{c)}] $\mathrm{X}$ is locally factorial.
\end{enumerate}
\end{proposition}

\begin{proof}
The equivalence of \emph{b)} and \emph{c)}, as well as the fact that \emph{c)} implies \emph{a)}, was proved in \sref{IV.21.6.9}.
Moreover, \emph{b)} implies \emph{a$''$)} by \sref{IV.21.7.3.1}, and \emph{a)} trivially implies \emph{a$'$)}.
It remains to show that either \emph{a$'$)} or \emph{a$''$)} implies \emph{c)}.

Suppose first that \emph{a$'$)} holds, and let us show first that $\mathrm{X}$ is normal.
The condition \emph{a$'$)} is local, so it also holds for every local scheme $\Spec(\sh{O}_{\mathrm{X},x})$.
Let $x\in\mathrm{X}^{(1)}$, so that $\mathrm{A}=\sh{O}_{\mathrm{X},x}$ is a Noetherian local ring of dimension $1$.
Applying \emph{a$'$)} to $\Spec(\mathrm{A})$ and to the prime $1$-codimensional cycle formed by its closed point, we see that the maximal ideal of $\mathrm{A}$ is generated by a single regular element; hence \sref[0]{0.17.1.1}{} $\mathrm{A}$ is a regular ring.
Since the localisations $\mathrm{A}_{\mathfrak p}$ are then also regular \sref[0]{0.17.3.2}, $\mathrm{X}$ is regular at all of its non-isolated maximal points; since it is reduced, and therefore regular, at its isolated points by hypothesis, $\mathrm{X}$ satisfies $(\mathrm{R}_1)$.
Let us also show that $\mathrm{X}$ satisfies $(\mathrm{S}_1)$, that is, for every $x\in\mathrm{X}$ such that $\dim(\sh{O}_{\mathrm{X},x})\geq 1$, one has $\mathfrak m_x\notin\operatorname{Ass}(\sh{O}_{\mathrm{X},x})$ \sref[0]{0.16.4.6}.
Indeed, condition \emph{a$'$)} applied to $\mathrm{X}_1=\Spec(\sh{O}_{\mathrm{X},x})$ shows that this scheme carries at least one nonzero divisor, in other words $\sh{M}_{\mathrm{X}_1}^*\neq\sh{O}_{\mathrm{X}_1}^*$, and it suffices to apply \sref{IV.21.1.10}.
These results already show that $\mathrm{X}$ is reduced \sref{IV.5.8.5}.
Let us next prove that it satisfies $(\mathrm{S}_2)$, which will establish that $\mathrm{X}$ is normal by Serre's criterion \sref{IV.5.8.6}.
Argue by contradiction: suppose that the set $\mathrm{E}$ of points $x\in\mathrm{X}$ at which $\mathrm{X}$ does not satisfy $(\mathrm{S}_2)$ is nonempty, and choose $x\in\mathrm{E}$ for which $\dim(\sh{O}_{\mathrm{X},x})$ is as small as possible.
Since $\mathrm{X}$ satisfies $(\mathrm{S}_1)$, one has $\dim(\sh{O}_{\mathrm{X},x})\geq 2$.
In $\mathrm{X}_1=\Spec(\sh{O}_{\mathrm{X},x})$, the open $\mathrm{U}=\mathrm{X}_1-\{x\}$ satisfies $(\mathrm{S}_2)$; let us show that $\mathrm{X}_1$ also satisfies $(\mathrm{S}_2)$, which gives a contradiction.
By \sref{IV.5.10.5}, it suffices to show that every section $f$ of $\sh{O}_{\mathrm{X}_1}$ over $\mathrm{U}$ extends to a section over $\mathrm{X}_1$.
Since $\mathrm{X}_1$ is reduced and $\mathrm{U}$ is dense in $\mathrm{X}_1$, $\mathrm{U}$ is schematically dense in $\mathrm{X}_1$, and $f$ is therefore the restriction to $\mathrm{U}$ of a regular meromorphic section $g\in\mathrm{M}(\mathrm{X}_1)$.
Moreover, since $\dim(\sh{O}_{\mathrm{X},x})\geq 2$, one has $\operatorname{cyc}(\operatorname{div}(g))\geq 0$ because $g|\mathrm{U}=f$.
As $\operatorname{cyc}$ is injective, necessarily $\operatorname{div}(g)\geq 0$, so \sref{IV.21.6.4}{} $g$ is a section of $\sh{O}_{\mathrm{X}_1}$ over $\mathrm{X}_1$, as required.

To show that \emph{a$'$)} implies \emph{c)}, observe now that \emph{a$'$)} implies the surjectivity of the canonical homomorphism $\operatorname{cyc}:\operatorname{Div}(\mathrm{X})\to\mathfrak{Z}^1(\mathrm{X})$; since $\mathrm{X}$ is normal, it suffices to apply \sref{IV.21.6.9}.

Finally, suppose that \emph{a$''$)} holds.
It follows from \sref{IV.21.6.5}{} that every prime $1$-codimensional cycle on $\mathrm{X}$ is the image under $\operatorname{cyc}$ of a positive divisor; the first part of the preceding argument again proves that $\mathrm{X}$ satisfies $(\mathrm{R}_1)$ and $(\mathrm{S}_1)$.
It remains only to prove that $\mathrm{X}$ is normal, the rest of the argument then being the same; equivalently, if $x\in\mathrm{X}$ is such that $\dim(\sh{O}_{\mathrm{X},x})\geq 2$, we must prove that $\operatorname{prof}(\sh{O}_{\mathrm{X},x})\geq 2$.
If $y$ is a generalisation of $x$ such that $\overline{\{y\}}$ has codimension $1$ in $\mathrm{X}$, let $\mathrm{Y}$ be the reduced closed sub-prescheme of $\mathrm{X}$ having $\overline{\{y\}}$ as its
\oldpage[IV-4]{280}
underlying space.
Then $\mathrm{Y}$ is integral and therefore, by hypothesis, regularly immersed in $\mathrm{X}$.
It follows that there is a regular element $t$ of $\sh{O}_{\mathrm{X},x}$ such that the ideal $t\sh{O}_{\mathrm{X},x}$ is the prime ideal defining the inverse image $\mathrm{Y}_1$ of $\mathrm{Y}$ in $\mathrm{X}_1=\Spec(\sh{O}_{\mathrm{X},x})$.
Since $\sh{O}_{\mathrm{X},x}/t\sh{O}_{\mathrm{X},x}$ is integral, this proves that $\operatorname{prof}(\sh{O}_{\mathrm{X},x})\geq 2$ and completes the proof of \sref{IV.21.7.4}.
\end{proof}

\begin{remark}[21.7.5]
\label{IV.21.7.5}
We cannot replace in \sref{IV.21.7.4}{} the condition \emph{a$'$)} by the weaker condition that every prime $1$-codimensional cycle on $\mathrm{X}$ is the image by $\operatorname{cyc}$ of a positive divisor.
This is shown by the example where $\mathrm{X}$ is the affine scheme defined in \sref{IV.14.1.5}{} (``the union of a plane and a line that intersects it'') \emph{which is not normal};
with the notation introduced in \sref{IV.14.1.5}, the prime $1$-codimensional cycles of $\mathrm{X}$ are those of the plane $\mathrm{X}_1$ and the closed points of the line $\mathrm{X}_2$ distinct from the point common to $\mathrm{X}_1$ and $\mathrm{X}_2$;
if $t_1$, $t_2$, $t_3$ are the images of $\mathrm{T}_1$, $\mathrm{T}_2$, $\mathrm{T}_3$ in $\mathrm{A}/\mathfrak{a}$, the prime $1$-codimensional cycles of $\mathrm{X}$ are defined by the monogenic prime ideals generated by $\mathrm{P}(t_1,t_2)$ (with $\mathrm{P}$ irreducible in $\mathrm{K}[\mathrm{T}_1,\mathrm{T}_2]$) or by $t_3-a$, with $a\neq 0$ in $\mathrm{K}$;
these elements are indeed regular elements in $\mathrm{A}/\mathfrak{a}$, so every $1$-codimensional cycle is indeed the canonical image of a positive divisor.
\end{remark}

\begin{corollary}[21.7.6]
\label{IV.21.7.6}
Let $\mathrm{X}$ be a locally Noetherian prescheme which is reduced at each of its isolated points.
The following conditions are equivalent:
\begin{enumerate}
  \item[\emph{a)}] The canonical homomorphism $c:\sh{D}iv_\mathrm{X}\to\mathscr{Z}_\mathrm{X}^1$ is an isomorphism of sheaves of ordered groups, and $\operatorname{Div}(\mathrm{X})=\operatorname{Div.princ}(\mathrm{X})$.
  \item[\emph{a$'$)}] The canonical homomorphism $c:\sh{D}iv_\mathrm{X}\to\mathscr{Z}_\mathrm{X}^1$ is injective, and every prime $1$-codimensional cycle is the image under $\operatorname{cyc}$ of a positive principal divisor, in other words is of the form $\operatorname{cyc}(\operatorname{div}(f))$, where $f$ is a regular section of $\sh{O}_\mathrm{X}$ over $\mathrm{X}$.
  \item[\emph{a$''$)}] For every integral closed sub-prescheme $\mathrm{Y}$ of $\mathrm{X}$ of codimension $1$, there exists a regular section $f$ of $\sh{O}_\mathrm{X}$ over $\mathrm{X}$ such that $\mathrm{Y}$ is defined by the Ideal $f\sh{O}_\mathrm{X}$.
  \item[\emph{b)}] $\mathrm{X}$ is normal and every prime $1$-codimensional cycle on $\mathrm{X}$ is principal.
  \item[\emph{c)}] $\mathrm{X}$ is locally factorial and $\operatorname{Pic}(\mathrm{X})=0$.
  \item[\emph{d)}] $\mathrm{X}$ is normal, and for every open $\mathrm{U}$ of $\mathrm{X}$, one has $\operatorname{Pic}(\mathrm{U})=0$.
\end{enumerate}
\end{corollary}

\begin{proof}
The equivalence of \emph{a)}, \emph{a$'$)}, \emph{a$''$)}, \emph{b)} and \emph{c)} follows at once from \sref{IV.21.7.4}{} and \sref{IV.21.6.11}.
Moreover, it follows at once from \emph{a$''$)} that every nonempty open $\mathrm{U}$ of $\mathrm{X}$ again satisfies the same conditions; in other words, these conditions imply \emph{d)}.
It remains to show that \emph{d)} implies \emph{b)}.
Let $\mathrm{U}_\lambda$ be the open complements of finite unions of sets of the form $\overline{\{x\}}$, where $\dim(\sh{O}_{\mathrm{X},x})\geq 2$.
It is clear that the $\mathrm{U}_\lambda$ form a decreasing filtered family and that, for every $x\in\bigcap_\lambda\mathrm{U}_\lambda$, one has $\dim(\sh{O}_{\mathrm{X},x})\leq 1$, so $\sh{O}_{\mathrm{X},x}$ is a principal ring by hypothesis.
We may therefore apply \sref{IV.21.6.12}{} to the family $(\mathrm{U}_\lambda)$; it shows that $\mathfrak{Cl}(\mathrm{X})$ is isomorphic to $\varinjlim\operatorname{Pic}(\mathrm{U}_\lambda)$, hence that $\mathfrak{Cl}(\mathrm{X})=0$ by hypothesis, which proves \emph{b)} by definition.
\end{proof}

\begin{remark}[21.7.7]
\label{IV.21.7.7}
When $\mathrm{X}=\Spec(\mathrm{A})$, where $\mathrm{A}$ is a Noetherian integral ring, the conditions of \sref{IV.21.7.6}{} are equivalent to saying that $\mathrm{A}$ is a \emph{factorial} ring.
\end{remark}


\subsection{Divisors and normalization}
\label{subsection:IV.21.8}

\begin{lemma}[21.8.1]
\label{IV.21.8.1}
Let $f:\mathrm{X}'\to\mathrm{X}$ be an integral morphism of preschemes.
For every locally free $\sh{O}_{\mathrm{X}'}$-Module $\sh{E}'$ of constant rank $n$ and every $x\in\mathrm{X}$, there exists an open neighbourhood $\mathrm{U}$ of $x$ in $\mathrm{X}$ such that, if we set $\mathrm{U}'=f^{-1}(\mathrm{U})$, $\sh{E}'|\mathrm{U}'$ is isomorphic to $\sh{O}_{\mathrm{U}'}^n$.
\end{lemma}

\begin{proof}
The question being local on $\mathrm{X}$, we can restrict to the case where $\mathrm{X}=\Spec(\mathrm{A})$ is affine, in which case $\mathrm{X}'=\Spec(\mathrm{A}')$, where $\mathrm{A}'$ is an $\mathrm{A}$-algebra that is integral.
Then $\mathrm{A}'$ is the inductive limit of its \emph{finite} sub-$\mathrm{A}$-algebras $\mathrm{A}'_\lambda$.
Setting $\mathrm{X}'_\lambda=\Spec(\mathrm{A}'_\lambda)$ and denoting by $p_\lambda:\mathrm{X}'\to\mathrm{X}'_\lambda$ the structural morphism, it follows from \sref{IV.8.5.2}{} and \sref{IV.8.5.5}{} that there exists an index $\lambda$ and an $\sh{O}_{\mathrm{X}'_\lambda}$-Module locally free $\sh{E}'_\lambda$ of rank $n$ such that $\sh{E}'=p_\lambda^*(\sh{E}'_\lambda)$, and it will clearly suffice to prove the lemma for $\mathrm{X}'_\lambda$ and $\sh{E}'_\lambda$;
in other words, we may assume that the morphism $f$ is \emph{finite}.
Set $\mathrm{B}=\sh{O}_{\mathrm{X},x}$ and let $\mathrm{B}'$ be the ring of the affine scheme $\mathrm{X}'_0=\mathrm{X}'\times_\mathrm{X}\Spec(\sh{O}_{\mathrm{X},x})$;
as $\mathrm{B}$ is a local ring and $\mathrm{B}'$ is a $\mathrm{B}$-algebra which is finite, $\mathrm{B}'$ is a semi-local ring (Bourbaki, \emph{Alg.\ comm.}, chap.~V, \textsection~2, n\textsuperscript{o}~1, prop.~3);
we conclude that the locally free $\sh{O}_{\mathrm{X}'_0}$-Module $\sh{E}'_0=\sh{E}'\otimes_{\sh{O}_{\mathrm{X}'}}\sh{O}_{\mathrm{X}'_0}$ is isomorphic to $\sh{O}_{\mathrm{X}'_0}^n$ (Bourbaki, \emph{Alg.\ comm.}, chap.~II, \textsection~5, n\textsuperscript{o}~3, prop.~5).
Considering $\mathrm{X}'_0$ as the
\oldpage[IV-4]{281}
projective limit of the preschemes induced by $\mathrm{X}'$ on the opens $f^{-1}(\mathrm{U})$, where $\mathrm{U}$ ranges over the set of affine open neighbourhoods of $x$ in $\mathrm{X}$, following the method of \sref{IV.8.1.2},{}~\emph{a)}), and applying \sref{IV.8.5.2.5}, we obtain the sought conclusion.
\end{proof}

\begin{corollary}[21.8.2]
\label{IV.21.8.2}
Let $f:\mathrm{X}'\to\mathrm{X}$ be an integral morphism of preschemes.
Then:
\begin{enumerate}[label=(\roman*)]
  \item We have $\RR^1 f_*(\sh{O}^*_{\mathrm{X}'})=0$ \sref[0]{0.12.2.1}.
  \item The group $\operatorname{Pic}(\mathrm{X}')$ is canonically isomorphic to $\HH^1(\mathrm{X},f_*(\sh{O}^*_{\mathrm{X}'}))$.
\end{enumerate}
\end{corollary}

\begin{proof}
\begin{enumerate}[label=(\roman*)]
  \item $\RR^1f_*(\sh{O}^*_{\mathrm{X}'})$ is the sheaf associated to the presheaf of commutative groups $\mathrm{U}\mapsto\HH^1(f^{-1}(\mathrm{U}),\sh{O}^*_{\mathrm{X}'})$ on $\mathrm{X}$ (\emph{loc.\ cit.});
  the fibre of this sheaf at a point $x$ can therefore be identified \sref[0]{0.5.4.7}{} with the commutative group $\varinjlim\operatorname{Pic}(f^{-1}(\mathrm{U}))$, where $\mathrm{U}$ ranges over the open neighbourhoods of $x$, the transition homomorphisms $\operatorname{Pic}(f^{-1}(\mathrm{U}'))\to\operatorname{Pic}(f^{-1}(\mathrm{U}))$ for $\mathrm{U}\subset\mathrm{U}'$ being defined by \sref{IV.21.3.2.4}.
Now, for every $x\in\mathrm{X}$, every open neighbourhood $\mathrm{U}'$ of $x$ in $\mathrm{X}$, and every element $\zeta\in\operatorname{Pic}(f^{-1}(\mathrm{U}'))$, it follows from \sref{IV.21.8.1}{} that there is an open neighbourhood $\mathrm{U}\subset\mathrm{U}'$ of $x$ such that the image of $\zeta$ in $\operatorname{Pic}(f^{-1}(\mathrm{U}))$ is zero.
Therefore the fibre of $\RR^1 f_*(\sh{O}^*_{\mathrm{X}'})$ at the point $x$ is zero.

  \item The Leray spectral sequence for the composed functor $\sh{F}'\mapsto\Gamma(\mathrm{X},f_*(\sh{F}'))$ of sheaves of commutative groups on $\mathrm{X}'$ (T,~{\hyperlink{ega4.c04.subsec.2-4}{2.4}}) gives the exact sequence of low degree
\[
  0\to\HH^1(\mathrm{X},f_*(\sh{O}^*_{\mathrm{X}'}))\to\HH^1(\mathrm{X}',\sh{O}^*_{\mathrm{X}'})\to\HH^0(\mathrm{X},\RR^1 f_*(\sh{O}^*_{\mathrm{X}'}))
\]
and the conclusion follows from (i) and from the isomorphism of $\operatorname{Pic}(\mathrm{X}')$ and $\HH^1(\mathrm{X}',\sh{O}^*_{\mathrm{X}'})$ \sref[0]{0.5.4.7}.
\end{enumerate}
\end{proof}

\begin{proposition}[21.8.3]
\label{IV.21.8.3}
Let $f=(\psi,\theta):\mathrm{X}'\to\mathrm{X}$ be an integral morphism of preschemes;
suppose that, for every open $\mathrm{U}$ of $\mathrm{X}$, the homomorphism $\Gamma(\theta):\Gamma(\mathrm{U},\sh{O}_\mathrm{X})\to\Gamma(\mathrm{U},f_*(\sh{O}_{\mathrm{X}'}))$ sends regular elements to regular elements, which allows us to canonically extend the homomorphism $\theta:\sh{O}_\mathrm{X}\to f_*(\sh{O}_{\mathrm{X}'})$ to a homomorphism of sheaves of rings $\theta':\sh{M}_\mathrm{X}\to f_*(\sh{M}_{\mathrm{X}'})$, whence homomorphisms of sheaves of multiplicative groups $\theta^*:\sh{O}^*_\mathrm{X}\to f_*(\sh{O}^*_{\mathrm{X}'})$ and $\theta'^*:\sh{M}^*_\mathrm{X}\to f_*(\sh{M}^*_{\mathrm{X}'})$, giving by passage to quotients a homomorphism $\theta''^*:\sh{D}iv_\mathrm{X}\to f_*(\sh{D}iv_{\mathrm{X}'})$.
We then have a commutative diagram
\[
\label{IV.21.8.3.1}
  \xymatrix{
    1 \ar[r] & \sh{O}^*_\mathrm{X} \ar[r] \ar[d]^{\theta^*} & \sh{M}^*_\mathrm{X} \ar[r] \ar[d]^{\theta'^*} & \sh{D}iv_\mathrm{X} \ar[r] \ar[d]^{\theta''^*} & 0\\
    1 \ar[r] & f_*(\sh{O}^*_{\mathrm{X}'}) \ar[r] & f_*(\sh{M}^*_{\mathrm{X}'}) \ar[r] & f_*(\sh{D}iv_{\mathrm{X}'}) \ar[r] & 0
  }
  \tag{21.8.3.1}
\]
in which the two rows are exact.
\end{proposition}

\begin{proof}
The only point to verify is the exactness of the second row of the diagram, which follows from the application of the exact sequence of cohomology for the functor $f_*$ to the exact sequence of sheaves of commutative groups on $\mathrm{X}'$
\[
  1\to\sh{O}^*_{\mathrm{X}'}\to\sh{M}^*_{\mathrm{X}'}\to\sh{D}iv_{\mathrm{X}'}\to 0
\]
and from \sref{IV.21.8.2}.
\end{proof}

\begin{corollary}[21.8.4]
\label{IV.21.8.4}
If, in addition to the hypotheses of \sref{IV.21.8.3}, the homomorphism $\theta'$ is an isomorphism of sheaves of rings, then $\theta^*:\sh{O}^*_\mathrm{X}\to f_*(\sh{O}^*_{\mathrm{X}'})$ is injective, $\theta''^*:\sh{D}iv_\mathrm{X}\to f_*(\sh{D}iv_{\mathrm{X}'})$ is surjective, and $\operatorname{Ker}(\theta''^*)$ is isomorphic to $\operatorname{Coker}(\theta^*)$.
\end{corollary}

\begin{proof}
This is an immediate consequence of the snake diagram (Bourbaki, \emph{Alg.\ comm.}, chap.~I, \textsection~2, n\textsuperscript{o}~4, prop.~2) applied to the diagram \sref{IV.21.8.3.1}.
\end{proof}

\oldpage[IV-4]{282}

\begin{proposition}[21.8.5]
\label{IV.21.8.5}
Let $\mathrm{X}$, $\mathrm{X}'$ be two locally Noetherian preschemes, $f:\mathrm{X}'\to\mathrm{X}$ an integral morphism;
suppose that there exists a schematically dense open $\mathrm{U}$ in $\mathrm{X}$ such that $f^{-1}(\mathrm{U})$ is schematically dense in $\mathrm{X}'$ (cf.\ \sref{IV.20.3.5}), and that the morphism $f^{-1}(\mathrm{U})\to\mathrm{U}$, restriction of $f$, is an isomorphism.
Then:
\begin{enumerate}[label=(\roman*)]
  \item The homomorphism $\theta':\sh{M}_\mathrm{X}\to f_*(\sh{M}_{\mathrm{X}'})$ is defined and bijective, the homomorphism
\[
  \theta^*:\sh{O}^*_\mathrm{X}\to f_*(\sh{O}^*_{\mathrm{X}'})
\]
is injective, the homomorphism $\theta''^*:\sh{D}iv_\mathrm{X}\to f_*(\sh{D}iv_{\mathrm{X}'})$ is surjective, $\operatorname{Ker}(\theta''^*)$ is isomorphic to $\sh{N}=\operatorname{Coker}(\theta^*)$, and the support of the sheaf of multiplicative groups $\sh{N}$ is contained in $\mathrm{S}=\mathrm{X}-\mathrm{U}$.
  \item Suppose moreover that the set $\mathrm{S}$ is discrete, and put $\sh{O}'_\mathrm{X}=f_*(\sh{O}_{\mathrm{X}'})$ for short.
Then there is a commutative diagram
\[
\label{IV.21.8.5.1}
  \xymatrix{
    1 \ar[r] & \prod_{s\in\mathrm{S}}\sh{O}'^{*}_{\mathrm{X},s}/\sh{O}^*_{\mathrm{X},s} \ar[r]^-{j} \ar[d] & \operatorname{Div}(\mathrm{X}) \ar[r] \ar[d]^{k_\mathrm{X}} & \operatorname{Div}(\mathrm{X}') \ar[r] \ar[d]^{k_{\mathrm{X}'}} & 0\\
    1 \ar[r] & \bigl(\prod_{s\in\mathrm{S}}\sh{O}'^{*}_{\mathrm{X},s}/\sh{O}^*_{\mathrm{X},s}\bigr)/\operatorname{Im}\Gamma(\mathrm{X}',\sh{O}^*_{\mathrm{X}'}) \ar[r]^-{j'} & \operatorname{Pic}(\mathrm{X}) \ar[r] & \operatorname{Pic}(\mathrm{X}') \ar[r] & 0
  }
  \tag{21.8.5.1}
\]
whose rows are exact and whose left vertical arrow is the canonical homomorphism.
\end{enumerate}
\end{proposition}

\begin{proof}
\begin{enumerate}[label=(\roman*)]
  \item The hypothesis gives a canonical isomorphism $\sh{M}''_\mathrm{X}\simto f_*(\sh{M}''_{\mathrm{X}'})$ for the sheaves of germs of pseudo-functions \sref{IV.20.2.2}; the assertion concerning $\theta'$ follows, in view of the canonical isomorphisms $\sh{M}_\mathrm{X}\simto\sh{M}''_\mathrm{X}$ and $\sh{M}_{\mathrm{X}'}\simto\sh{M}''_{\mathrm{X}'}$ \sref{IV.20.2.11}.
The rest of assertion (i) follows from \sref{IV.21.8.4}, except what concerns the support of $\sh{N}$, which follows directly from the hypothesis on $\mathrm{U}$.

  \item Applying the cohomology exact sequence to the two exact sequences of sheaves of commutative groups
\[
\label{IV.21.8.5.2}
\begin{gathered}
  1\to\sh{N}\to\sh{D}iv_\mathrm{X}\to f_*(\sh{D}iv_{\mathrm{X}'})\to 0,\\
  1\to\sh{O}^*_\mathrm{X}\to f_*(\sh{O}^*_{\mathrm{X}'})\to\sh{N}\to 1,
\end{gathered}
  \tag{21.8.5.2}
\]
we obtain respectively the exact sequences
\[
\label{IV.21.8.5.3}
\begin{gathered}
  1\to\Gamma(\mathrm{X},\sh{N})\xrightarrow{j}\operatorname{Div}(\mathrm{X})\to\operatorname{Div}(\mathrm{X}')\to\HH^1(\mathrm{X},\sh{N}),\\
  \Gamma(\mathrm{X}',\sh{O}^*_{\mathrm{X}'})\to\Gamma(\mathrm{X},\sh{N})\xrightarrow{\partial}\operatorname{Pic}(\mathrm{X})\to\operatorname{Pic}(\mathrm{X}')\to\HH^1(\mathrm{X},\sh{N}),
\end{gathered}
  \tag{21.8.5.3}
\]
\oldpage[IV-4]{283}
taking into account the canonical isomorphism $\operatorname{Pic}(\mathrm{X}')\simto\HH^1(\mathrm{X},f_*(\sh{O}^*_{\mathrm{X}'}))$ \sref{IV.21.8.2}.
Now, since the support of $\sh{N}$ is contained in the discrete set $\mathrm{S}$, which is closed in $\mathrm{X}$, we have
\[
  \HH^1(\mathrm{X},\sh{N})=\HH^1(\mathrm{S},\sh{N}|\mathrm{S}) \quad\text{(G, II, 4.9.2)},
  \qquad
  \HH^1(\mathrm{S},\sh{N}|\mathrm{S})=\prod_{s\in\mathrm{S}}\HH^1(\{s\},\sh{N}_s)=0
\]
by the definition of cohomology (G, II, 4.4).
Similarly, $\Gamma(\mathrm{X},\sh{N})=\prod_{s\in\mathrm{S}}\sh{N}_s$ and $\sh{N}_s=\sh{O}'^{*}_{\mathrm{X},s}/\sh{O}^*_{\mathrm{X},s}$, by the second exact sequence \sref{IV.21.8.5.2}.
It remains to describe the injections $j$ and $j'$.
A section $t$ of $\sh{N}$ over $\mathrm{X}$ may be described using a covering of $\mathrm{X}$ consisting of $\mathrm{U}$ and open sets $\mathrm{U}_i$ such that $\mathrm{U}_i\cap\mathrm{S}$ consists of a single point $s_i$ and $\mathrm{U}_i\cap\mathrm{U}_j\subset\mathrm{U}$ for $i\neq j$.
For each $i$, take a section $t_i$ of $\sh{N}$ over $\mathrm{U}_i$, necessarily with $(t_i)_{s_i}\in\sh{O}'^{*}_{\mathrm{X},s_i}/\sh{O}^*_{\mathrm{X},s_i}$ and $(t_i)_x=0$ for $x\in\mathrm{U}_i-\{s_i\}$.
The germ $(t_i)_{s_i}$ comes from a section $u_i$ of $\sh{O}^*_{\mathrm{X}'}$ over $f^{-1}(\mathrm{U}_i)$; this may also be regarded as a section of $\sh{M}^*_{\mathrm{X}'}$ over $f^{-1}(\mathrm{U}_i)$, hence, by the hypothesis, as a section of $\sh{M}^*_{\mathrm{X}}$ over $\mathrm{U}_i$.
It canonically determines a section $d_i\in\Gamma(\mathrm{U}_i,\sh{D}iv_\mathrm{X})$.
On $f^{-1}(\mathrm{U}\cap\mathrm{U}_i)$ the restriction of $u_i$ identifies, by the hypothesis, with a section of $\sh{O}^*_{\mathrm{X}}$ over $\mathrm{U}\cap\mathrm{U}_i$; consequently the restrictions of the $d_i$ to $\mathrm{U}\cap\mathrm{U}_i$ are all zero.
Thus the $d_i$ are the restrictions of a single divisor $d\in\operatorname{Div}(\mathrm{X})$, which is the image of $t$ under $j$.
Likewise, the image of $t$ under $\partial:\Gamma(\mathrm{X},\sh{N})\to\HH^1(\mathrm{X},\sh{O}^*_{\mathrm{X}})$ is represented by the cocycle equal to the restriction of $u_i$ on $\mathrm{U}\cap\mathrm{U}_i$, and to
$(u_i|_{\mathrm{U}_i\cap\mathrm{U}_j})(u_j|_{\mathrm{U}_i\cap\mathrm{U}_j})^{-1}$ on $\mathrm{U}_i\cap\mathrm{U}_j$.
This gives the expression for $j'(t)$ and proves the commutativity of \sref{IV.21.8.5.1}.
\end{enumerate}
\end{proof}

\begin{remark}[21.8.6]
\label{IV.21.8.6}
\begin{enumerate}[label=(\roman*)]
  \item The hypotheses of \sref{IV.21.8.5},{}~(i) are in particular satisfied when $\mathrm{X}$ is a reduced locally Noetherian prescheme with only finitely many irreducible components, $\mathrm{X}'$ is its normalization, and $\mathrm{X}'$ is also locally Noetherian (\sref[II]{II.6.3.8});
the $\sh{O}_\mathrm{X}$-module $\sh{D}iv_\mathrm{X}$ is then an \emph{extension} of $f_*(\sh{D}iv_{\mathrm{X}'})$ by the cokernel $\sh{N}$ of $\sh{O}^{*}_\mathrm{X}\to f_*(\sh{O}^*_{\mathrm{X}'})$, which may be regarded as known.
If moreover $\mathrm{X}$ is an algebraic curve (reduced) over a field $k$, we are in the conditions of application of \sref{IV.21.8.5},{}~(ii).

  \item When the hypotheses of \sref{IV.21.8.5},{}~(ii) hold and, in addition, the canonical homomorphism on global sections $\Gamma(\mathrm{X},\sh{O}_\mathrm{X})\to\Gamma(\mathrm{X}',\sh{O}_{\mathrm{X}'})$ is \emph{bijective}, the kernels of the surjective homomorphisms $\operatorname{Div}(\mathrm{X})\to\operatorname{Div}(\mathrm{X}')$ and $\operatorname{Pic}(\mathrm{X})\to\operatorname{Pic}(\mathrm{X}')$ are \emph{isomorphic}.

  \item When the hypotheses of \sref{IV.21.8.5},{}~(ii) hold, the homomorphism $\operatorname{Div}(\mathrm{X})\to\operatorname{Div}(\mathrm{X}')$ can be injective (and hence bijective) only if $\sh{O}'^{*}_{\mathrm{X},s}=\sh{O}^*_{\mathrm{X},s}$ for every $s\in\mathrm{S}$.
For an $s\in\mathrm{S}$ such that $\sh{O}'_{\mathrm{X},s}$ is \emph{local}---which, in view of $\mathrm{X}'=\Spec(\sh{O}'_\mathrm{X})$ \sref[II]{II.1.3}, means that there is a single point $s'\in\mathrm{X}'$ over $s$---this necessarily entails that the residue fields $\boldsymbol{k}(s)$ and $\boldsymbol{k}(s')$ are equal and that $1+\mathfrak{m}_s=1+\mathfrak{m}_{s'}$, and hence finally that $\sh{O}'_{\mathrm{X},s}=\sh{O}_{\mathrm{X},s}$.
Equivalently, in view of the hypothesis, there is an open neighbourhood $\mathrm{V}$ of $s$ in $\mathrm{X}$ such that $f^{-1}(\mathrm{V})\to\mathrm{V}$ is an isomorphism.
In general, the ring $\sh{O}'_{\mathrm{X},s}$ is semilocal (this is evident when $f$ is finite, and in general follows from \sref[0]{0.23.2.6}); the canonical homomorphism $\sh{O}^{*}_{\mathrm{X},s}\to\sh{O}'^{*}_{\mathrm{X},s}$ gives, on passing to quotients, a homomorphism from the multiplicative group $\boldsymbol{k}(s)^*$ to the product of the multiplicative groups $\boldsymbol{k}(s'_j)^*$, where the $s'_j$ are the points (in
\oldpage[IV-4]{284}
  number $>1$) of $\mathrm{X}'$ above $s$.
It is immediate that this homomorphism can be bijective only if $\boldsymbol{k}(s)$ and all the $\boldsymbol{k}(s'_j)$ are equal to the field $\bb{F}_2$ with two elements;
moreover, if this condition holds, the multiplicative group $1+\mathfrak{m}_s$ must have image $1+\mathfrak{r}$, where $\mathfrak{r}$ is the radical of $\sh{O}'_{\mathrm{X},s}$, or equivalently $\mathfrak{m}_s=\mathfrak{r}$.
This entails that $\sh{O}'_{\mathrm{X},s}\otimes_{\sh{O}_{\mathrm{X},s}}\boldsymbol{k}(s)$ is a direct product of fields isomorphic to $\boldsymbol{k}(s)$ (and, when $f$ is finite, that $f$ is unramified at the points $s'_j$ \sref{IV.17.4.1}).
Thus, if for example \emph{no} residue field of $\mathrm{X}$ is isomorphic to $\bb{F}_2$, the canonical homomorphism $\operatorname{Div}(\mathrm{X})\to\operatorname{Div}(\mathrm{X}')$ can be bijective only if $f$ is an \emph{isomorphism}.
In the case where the canonical homomorphism $\Gamma(\mathrm{X},\sh{O}_\mathrm{X})\to\Gamma(\mathrm{X}',\sh{O}_{\mathrm{X}'})$ is bijective, we conclude from this and from (ii) that in the preceding considerations, one can replace the homomorphism $\operatorname{Div}(\mathrm{X})\to\operatorname{Div}(\mathrm{X}')$ by the homomorphism $\operatorname{Pic}(\mathrm{X})\to\operatorname{Pic}(\mathrm{X}')$.
\end{enumerate}
\end{remark}


\subsection{Divisors on preschemes of dimension 1}
\label{subsection:IV.21.9}
\label{IV.21.9}

\begin{env}[21.9.1]
\label{IV.21.9.1}
Let $\mathrm{X}$ be a topological space, $x$ a point of $\mathrm{X}$, $i_x:\{x\}\to\mathrm{X}$ the canonical injection.
If $\mathrm{A}(x)$ is a commutative group, we may regard it as a sheaf of commutative groups on the one-point space $\{x\}$ and take its direct image $(i_x)_*(\mathrm{A}(x))$, a sheaf of commutative groups on $\mathrm{X}$ \sref[0]{0.3.4.1};
it follows at once from the definitions that for every open $\mathrm{U}$ of $\mathrm{X}$, $\Gamma(\mathrm{U},(i_x)_*(\mathrm{A}(x)))=\mathrm{A}(x)$ if $x\in\mathrm{U}$, and is reduced to $0$ in the contrary case, and for $y\in\overline{\{x\}}$, $((i_x)_*(\mathrm{A}(x)))_y=\mathrm{A}(x)$, and for $y\notin\overline{\{x\}}$, $((i_x)_*(\mathrm{A}(x)))_y=0$.

If now $\sh{F}$ is a sheaf of commutative groups on $\mathrm{X}$ and $\mathrm{Y}$ a part of $\mathrm{X}$, for every open $\mathrm{U}$ of $\mathrm{X}$, we have a canonical homomorphism
\[
\label{IV.21.9.1.1}
  \Gamma(\mathrm{U},\sh{F})\to\prod_{x\in\mathrm{U}\cap\mathrm{Y}}\sh{F}_x=\prod_{x\in\mathrm{Y}}\Gamma(\mathrm{U},(i_x)_*(\sh{F}_x))
  \tag{21.9.1.1}
\]
and since these homomorphisms commute with the restriction operators, they define a canonical homomorphism of \emph{sheaves}
\[
\label{IV.21.9.1.2}
  j_\mathrm{Y}:\sh{F}\to\prod_{x\in\mathrm{Y}}(i_x)_*(\sh{F}_x).
  \tag{21.9.1.2}
\]
\end{env}

\begin{lemma}[21.9.2]
\label{IV.21.9.2}
Let $\mathrm{X}$ be a locally Noetherian topological space, $\mathrm{X}_0$ the set of its closed points, $\sh{F}$ a sheaf of commutative groups on $\mathrm{X}$.
The following conditions are equivalent:
\begin{enumerate}[label=\emph{\alph*)}]
  \item The canonical homomorphism $j_{\mathrm{X}_0}:\sh{F}\to\prod_{x\in\mathrm{X}_0}(i_x)_*(\sh{F}_x)$ is injective and has for image $\bigoplus_{x\in\mathrm{X}_0}(i_x)_*(\sh{F}_x)$.
  \item[\rm a$'$)] There exists a family of commutative groups $(\mathrm{A}(x))_{x\in\mathrm{X}_0}$ such that $\sh{F}$ is isomorphic to $\bigoplus_{x\in\mathrm{X}_0}(i_x)_*(\mathrm{A}(x))$.
  \item Every section of $\sh{F}$ over an open $\mathrm{U}$ has a discrete support contained in $\mathrm{X}_0$.
\end{enumerate}
When these conditions hold, for every $x\in\mathrm{X}_0$ the group $\mathrm{A}(x)$ is determined up to unique isomorphism and is isomorphic to $\sh{F}_x$.
Furthermore, the sheaf $\sh{F}$ is then flasque.
\end{lemma}

\begin{proof}
\oldpage[IV-4]{285}
As the points of $\mathrm{X}_0$ are closed in $\mathrm{X}$, we have $((i_x)_*(\mathrm{A}(x)))_y=0$ for $y\neq x\in\mathrm{X}_0$;
this remark proves the uniqueness assertion relative to the groups $\mathrm{A}(x)$, and it is clear moreover that \emph{a)} implies \emph{a$'$)}.
To see that \emph{a$'$)} implies \emph{b)}, we may restrict to the case where $\mathrm{U}$ is Noetherian;
then (G,~II,~3.10) we have
\[
  \Gamma(\mathrm{U},\bigoplus_{x\in\mathrm{X}_0}(i_x)_*(\mathrm{A}(x)))=\bigoplus_{x\in\mathrm{X}_0}\Gamma(\mathrm{U},(i_x)_*(\mathrm{A}(x))),
\]
and thus every section $s$ of $\sh{F}$ over $\mathrm{U}$ is a direct sum of a \emph{finite} number of sections $s_k\in\Gamma(\mathrm{U},(i_{x_k})_*(\mathrm{A}(x_k)))$ ($1\leq k\leq n$) with $x_k\in\mathrm{X}_0$.
Since $x_k$ is closed in $\mathrm{X}$, the support of $s_k$ is reduced to $\{x_k\}$, so the support of $s$ is contained in the finite set of the $x_k$, which is evidently discrete.
Let us show that \emph{b)} implies \emph{a)}.
For every Noetherian open $\mathrm{U}$, the support of every section of $\sh{F}$ over $\mathrm{U}$, being discrete and quasi-compact, is \emph{finite}.
It follows from hypothesis \emph{b)} that, for every Noetherian open $\mathrm{U}$, the image of the homomorphism \sref{IV.21.9.1.1}{} for $\mathrm{Y}=\mathrm{X}_0$ is contained in $\bigoplus_{x\in\mathrm{X}_0}\Gamma(\mathrm{U},(i_x)_*(\sh{F}_x))$ and that this homomorphism is injective; this establishes \emph{a)}.

Finally, to show that $\sh{F}$ is flasque, consider a section $s$ of $\sh{F}$ over an open $\mathrm{U}$ of $\mathrm{X}$;
as the support of $s$ is a discrete and closed subspace in $\mathrm{X}$, we extend $s$ to a section $s'$ of $\sh{F}$ over $\mathrm{X}$ by setting $s'_x=0$ in $\mathrm{X}-\mathrm{U}$.
\end{proof}

\begin{remark}[21.9.3]
\label{IV.21.9.3}
\begin{enumerate}[label=(\roman*)]
  \item In the condition \emph{b)} of \sref{IV.21.9.2}, it does not suffice to assume that the support of every section of $\sh{F}$ over any open of $\mathrm{X}$ is \emph{discrete};
this is shown by the example where we take for $\mathrm{X}$ the spectrum of a discrete valuation ring, with a closed point $b$ and a generic point $a$, and for $\sh{F}$ the sheaf of commutative groups such that $\sh{F}_b=0$ and $\sh{F}_a=\bb{Z}$.

  \item The conditions of \sref{IV.21.9.2}{} being supposed satisfied, let $\mathrm{E}$ be a discrete part of $\mathrm{X}_0$, and for each $x\in\mathrm{E}$, let $a(x)$ be an element of $\mathrm{A}(x)$.
There then exists a unique section $t$ of $\bigoplus_{x\in\mathrm{X}_0}(i_x)_*(\mathrm{A}(x))$ over $\mathrm{X}$ such that $t_x=a(x)$ for all $x\in\mathrm{E}$ and whose support is contained in $\mathrm{E}$.
Indeed, for every Noetherian open $\mathrm{U}$, $\mathrm{E}\cap\mathrm{U}$ is finite, and it suffices to take for $t$ the section $\bigoplus_{x\in\mathrm{E}}(i_x)_*(\mathrm{A}(x))$ whose restriction to each Noetherian open $\mathrm{U}$ is the sum of the $a(x)$ for $x\in\mathrm{E}\cap\mathrm{U}$.
\end{enumerate}
\end{remark}

\begin{proposition}[21.9.4]
\label{IV.21.9.4}
Let $\mathrm{X}$ be a locally Noetherian prescheme of dimension $\leq 1$, and let $\mathrm{X}^{(1)}$ be the set of points $x\in\mathrm{X}$ such that $\dim(\sh{O}_{\mathrm{X},x})=1$.
We then have a canonical isomorphism
\[
\label{IV.21.9.4.1}
  \sh{D}iv_\mathrm{X}\simto\bigoplus_{x\in\mathrm{X}^{(1)}}(i_x)_*(\operatorname{Div}(\sh{O}_{\mathrm{X},x}))
  \tag{21.9.4.1}
\]
and $\sh{D}iv_\mathrm{X}$ is flasque.
\end{proposition}

\begin{proof}
Taking into account the isomorphism \sref{IV.21.4.6.1}, the homomorphism \sref{IV.21.9.4.1}{} is defined by \sref{IV.21.9.1.2}:
let us prove that it is a bijection.
As $\dim(\mathrm{X})\leq 1$, the points of $\mathrm{X}^{(1)}$ are the \emph{non-isolated closed points} of $\mathrm{X}$.
We are thus reduced to proving that:
1\textsuperscript{o} $\sh{D}iv_\mathrm{X}$ satisfies the condition \emph{b)} of \sref{IV.21.9.2};
2\textsuperscript{o} for every isolated point $x\in\mathrm{X}$, we have $(\sh{D}iv_\mathrm{X})_x=0$.
The second point follows from the fact that $\sh{O}_{\mathrm{X},x}$ is an Artinian ring and that every regular element of $\sh{O}_{\mathrm{X},x}$ is therefore invertible.
\oldpage[IV-4]{286}
For 1\textsuperscript{o}, it suffices to note that, for every open $\mathrm{U}$ of $\mathrm{X}$ and every divisor $\mathrm{D}\in\operatorname{Div}(\mathrm{U})$, every maximal point $x$ of the support $\mathrm{S}$ of $\mathrm{D}$ satisfies $\operatorname{prof}(\sh{O}_{\mathrm{X},x})=1$ \sref{IV.21.1.9}, and hence belongs a fortiori to $\mathrm{X}^{(1)}\cap\mathrm{U}$.
Since $\dim(\mathrm{X})\leq1$, the set of these points is equal to $\mathrm{S}$; thus $\mathrm{S}$ is discrete, the set of irreducible components of $\mathrm{S}$ being locally finite.
\end{proof}

\begin{corollary}[21.9.5]
\label{IV.21.9.5}
Let $\mathrm{X}$ be a locally Noetherian prescheme of dimension $\leq 1$.
For every discrete part $\mathrm{E}\subset\mathrm{X}^{(1)}$ and every family $(\mathrm{D}(x))_{x\in\mathrm{E}}$ with $\mathrm{D}(x)\in\operatorname{Div}(\sh{O}_{\mathrm{X},x})$, there exists a unique divisor $\mathrm{D}$ on $\mathrm{X}$ such that the support of $\mathrm{D}$ is contained in $\mathrm{E}$ and that $\mathrm{D}_x=\mathrm{D}(x)$ for all $x\in\mathrm{E}$.
\end{corollary}

\begin{proof}
This follows from \sref{IV.21.9.4}{} and from \sref{IV.21.9.3},{}~(ii).
\end{proof}

\begin{corollary}[21.9.6]
\label{IV.21.9.6}
Let $\mathrm{X}$ be a locally Noetherian prescheme of dimension $\leq 1$.
Every divisor $\mathrm{D}$ on $\mathrm{X}$ is of the form $\mathrm{D}'-\mathrm{D}''$, where $\mathrm{D}'$, $\mathrm{D}''$ are two divisors $\geq 0$ with supports contained in that of $\mathrm{D}$.
\end{corollary}

\begin{proof}
By virtue of \sref{IV.21.9.5}, it is at once reduced to the case where $\mathrm{X}=\Spec(\mathrm{A})$, where $\mathrm{A}$ is a Noetherian local ring;
it then suffices to observe that $\mathrm{M}(\mathrm{X})$ is the total ring of fractions of $\mathrm{A}$, and that a section of $\sh{M}^*_\mathrm{X}$ over $\mathrm{X}$ is written by definition as the quotient $b/a$ of two regular elements of $\mathrm{A}$.
\end{proof}

\begin{corollary}[21.9.7]
\label{IV.21.9.7}
Let $\mathrm{X}$ be a locally Noetherian prescheme of dimension $\leq 1$, without isolated point, and $\mathrm{U}$ a dense open in $\mathrm{X}$.
There then exists a divisor $\mathrm{D}\geq 0$ on $\mathrm{X}$, with support contained in $\mathrm{U}$ and meeting each of the irreducible components of $\mathrm{X}$.
\end{corollary}

\begin{proof}
We can assume that $\mathrm{U}$ is the union of disjoint non-empty opens $\mathrm{U}_\alpha$, each of which is contained in a single irreducible component of $\mathrm{X}$;
it suffices then to take in each $\mathrm{U}_\alpha$ a point $x_\alpha$ closed in $\mathrm{X}$ (there exists one since the unique generic point of $\mathrm{U}_\alpha$ cannot be isolated by hypothesis), and to apply \sref{IV.21.9.5}{} to the discrete set of the $x_\alpha$.
\end{proof}

The interest of this corollary is that it will allow us to prove that a separated algebraic curve $\mathrm{X}$ over a field $k$ is quasi-projective, the divisor $\mathrm{D}\geq 0$ defined in \sref{IV.21.9.7}{} being then necessarily \emph{ample} by virtue of the Riemann--Roch theorem for curves (chap.~V).

\begin{env}[21.9.8]
\label{IV.21.9.8}
For locally Noetherian preschemes $\mathrm{X}$ of dimension $\leq 1$ the proposition \sref{IV.21.9.4}{} brings back the determination of $\sh{D}iv_\mathrm{X}$ to the case where $\mathrm{X}=\Spec(\mathrm{A})$, $\mathrm{A}$ being a Noetherian local ring of dimension $1$.

When $\mathrm{A}$ is a regular local ring of dimension $1$ (in other words, a discrete valuation ring), the group $\operatorname{Div}(\mathrm{A})$ is canonically identified with $\bb{Z}$ \sref{IV.21.6.8};
by \sref{IV.21.9.2}:
\end{env}

\begin{proposition}[21.9.9]
\label{IV.21.9.9}
Let $\mathrm{X}$ be a locally Noetherian regular prescheme of dimension $\leq 1$.
Then the sheaf $\sh{D}iv_\mathrm{X}$ is canonically isomorphic to $\bigoplus_{x\in\mathrm{X}^{(1)}}(i_x)_*(\bb{Z}(x))$, where $\bb{Z}(x)$ is the additive group $\bb{Z}$ considered as a sheaf of groups on the space $\{x\}$.
\end{proposition}

\begin{env}[21.9.10]
\label{IV.21.9.10}
Suppose only that the Noetherian local ring $\mathrm{A}$ of dimension $1$ is \emph{reduced};
then, if $\mathrm{A}'$ is the \emph{normalization} of $\mathrm{A}$ (the integral closure of $\mathrm{A}$ in its total ring of fractions), $\mathrm{A}'$ is Noetherian by virtue of the Krull--Akizuki theorem (Bourbaki, \emph{Alg.\ comm.}, chap.~VII, \textsection~2, n\textsuperscript{o}~5, prop.~5), and we have seen in \sref{IV.21.8.6}{} how
\oldpage[IV-4]{287}
$\operatorname{Div}(\mathrm{A})$ can be obtained as an \emph{extension} of $\operatorname{Div}(\mathrm{A}')$; the latter group is of the form $\bb{Z}^r$.
\end{env}

\begin{proposition}[21.9.11]
\label{IV.21.9.11}
Let $\mathrm{X}$ be a Noetherian prescheme, $\mathrm{X}_0$ a closed sub-prescheme of $\mathrm{X}$ having the following properties:
\begin{enumerate}
  \item[\rm 1\textsuperscript{o}] $\dim(\mathrm{X}_0)\leq 1$.
  \item[\rm 2\textsuperscript{o}] For every locally closed part $\mathrm{Y}$ of $\mathrm{X}$ such that $\mathrm{Y}\cap\mathrm{X}_0$ is discrete, there exists a part $\mathrm{Y}'$ of $\mathrm{Y}$, closed in $\mathrm{X}$ and open in $\mathrm{Y}$, containing $\mathrm{Y}\cap\mathrm{X}_0$.
\end{enumerate}
Under these conditions:
\begin{enumerate}[label=(\roman*)]
  \item Let $\mathrm{Z}_0$ be the union of the sets $\overline{\{x\}}\cap\mathrm{X}_0$ when $x$ ranges over the part of $\operatorname{Ass}(\sh{O}_\mathrm{X})$ formed of points such that $\overline{\{x\}}\cap\mathrm{X}_0$ is finite.
Then, for every divisor $\mathrm{D}_0$ on $\mathrm{X}_0$, whose support does not meet $\mathrm{Z}_0$, there exists a divisor $\mathrm{D}$ on $\mathrm{X}$ whose inverse image by the injection $\mathrm{X}_0\to\mathrm{X}$ exists \sref{IV.21.4}{} and is equal to $\mathrm{D}_0$;
if moreover $\mathrm{D}_0\geq 0$, we can assume $\mathrm{D}\geq 0$.

  \item Suppose in addition that there exists in $\mathrm{X}_0$ an affine open $\mathrm{U}_0$ containing $(\operatorname{Ass}(\sh{O}_{\mathrm{X}_0}))\cup\mathrm{Z}_0$ (a condition automatically satisfied when an ample $\sh{O}_{\mathrm{X}_0}$-module exists \sref[II]{II.4.5.4}).
Then the canonical homomorphism \sref{IV.21.3.2.4}{} $\operatorname{Pic}(\mathrm{X})\to\operatorname{Pic}(\mathrm{X}_0)$ is surjective.
\end{enumerate}
\end{proposition}

\begin{proof}
\begin{enumerate}[label=(\roman*)]
  \item In view of \sref{IV.21.9.6}, it is enough to prove the assertion for divisors $\mathrm{D}_0\geq0$.
The support $\mathrm{T}$ of $\mathrm{D}_0$ is a finite discrete set, closed in $\mathrm{X}_0$, by \sref{IV.21.9.4}.
We may suppose $\mathrm{D}_0\neq0$, that is, $\mathrm{T}\neq\emp$; otherwise there is nothing to prove.
For each $x\in\mathrm{T}$ there is an element $s_x\in\sh{O}_{\mathrm{X}_0,x}$ which is not a zero divisor, belongs to the maximal ideal of that ring, and whose image in $(\sh{D}iv_{\mathrm{X}_0})_x$ is $(\mathrm{D}_0)_x$.
There is an affine open neighbourhood $\mathrm{U}^{(x)}$ of $x$ in $\mathrm{X}$ and a section $g^{(x)}$ of $\sh{O}_\mathrm{X}$ over $\mathrm{U}^{(x)}$ such that $s_x$ is the image in $\sh{O}_{\mathrm{X}_0,x}$ of the germ $(g^{(x)})_x\in\sh{O}_{\mathrm{X},x}$.
After shrinking $\mathrm{U}^{(x)}$, we may arrange that $g^{(x)}$ is not a zero divisor in $\Gamma(\mathrm{U}^{(x)},\sh{O}_\mathrm{X})$; equivalently \sref{IV.3.1.9}, if $\mathrm{V}^{(x)}$ denotes the closed subset of $\mathrm{U}^{(x)}$ defined by $g^{(x)}=0$, then $\mathrm{V}^{(x)}\cap\operatorname{Ass}(\sh{O}_\mathrm{X})=\emp$.

Indeed, if $z\in\operatorname{Ass}(\sh{O}_\mathrm{X})$, the hypothesis $x\notin\mathrm{Z}_0$ implies either that $x\notin\overline{\{z\}}$, or that $x$ is not isolated in $\overline{\{z\}}\cap\mathrm{X}_0$.
By shrinking $\mathrm{U}^{(x)}$, we may suppose that the second alternative holds for every $z\in\mathrm{V}^{(x)}\cap\operatorname{Ass}(\sh{O}_\mathrm{X})$.
But then $\mathrm{V}^{(x)}$ would contain an irreducible component of dimension $1$ of $\mathrm{X}_0\cap\mathrm{U}^{(x)}$ containing $x$; this would mean that the image $\overline{g}^{(x)}$ of $g^{(x)}$ in $\Gamma(\mathrm{U}^{(x)}\cap\mathrm{X}_0,\sh{O}_{\mathrm{X}_0})$ belongs to the nilradical of that ring, and hence that $s_x$ belongs to the nilradical of $\sh{O}_{\mathrm{X}_0,x}$, a contradiction.
Thus $x\notin\overline{\{z\}}$ for every $z\in\mathrm{V}^{(x)}\cap\operatorname{Ass}(\sh{O}_\mathrm{X})$; since this set is finite, a further shrinking of $\mathrm{U}^{(x)}$ gives $\mathrm{V}^{(x)}\cap\operatorname{Ass}(\sh{O}_\mathrm{X})=\emp$.

The choice of $\mathrm{U}^{(x)}$ therefore defines a divisor $\mathrm{D}'^{(x)}=\operatorname{div}(g^{(x)})$ on $\mathrm{U}^{(x)}$.
We have also seen that $x$ is necessarily isolated in $\mathrm{V}^{(x)}\cap\mathrm{X}_0$; shrinking once more, we may suppose $\mathrm{V}^{(x)}\cap\mathrm{X}_0=\{x\}$.
By condition $2^{\mathrm{o}}$, there is then a subset $\mathrm{W}^{(x)}$ of $\mathrm{V}^{(x)}$, closed in $\mathrm{X}$ and open in $\mathrm{V}^{(x)}$, such that $\mathrm{W}^{(x)}\cap\mathrm{X}_0=\{x\}$.
Replacing $\mathrm{U}^{(x)}$ once more by a smaller open neighbourhood of $x$, we may suppose
\oldpage[IV-4]{288}
$\mathrm{W}^{(x)}=\mathrm{V}^{(x)}$, or in other words that $\mathrm{V}^{(x)}$ is closed in $\mathrm{X}$.
We may now define a divisor $\mathrm{D}^{(x)}$ on $\mathrm{X}$ by requiring $\mathrm{D}^{(x)}|\mathrm{U}^{(x)}=\mathrm{D}'^{(x)}$ and $\mathrm{D}^{(x)}|(\mathrm{X}-\mathrm{V}^{(x)})=0$; this is well defined because on $\mathrm{U}^{(x)}\cap(\mathrm{X}-\mathrm{V}^{(x)})$ the section $g^{(x)}$ is invertible.
It is then enough to take $\mathrm{D}=\sum_{x\in\mathrm{T}}\mathrm{D}^{(x)}$, a meaningful sum since $\mathrm{T}$ is finite.

  \item In view of the commutative diagram \sref{IV.21.4.2.1}, the conclusion follows from (i) once we prove that every invertible $\sh{O}_{\mathrm{X}_0}$-module $\sh{L}_0$ is isomorphic to one of the form $\sh{O}_{\mathrm{X}_0}(\mathrm{D}_0)$ \sref{IV.21.2.8}, where $\mathrm{D}_0$ is a divisor on $\mathrm{X}_0$ whose support does not meet $\mathrm{Z}_0$.
Since $\mathrm{U}_0$ is schematically dense in $\mathrm{X}_0$ \sref{IV.20.2.13},{}~(iv), it is enough \sref{IV.3.1.9} to find a section $t\in\Gamma(\mathrm{U}_0,\sh{L}_0)$ such that $t(z_0)\neq0$ at every point of $\operatorname{Ass}(\sh{O}_{\mathrm{X}_0})\cup\mathrm{Z}_0$.
We may therefore reduce to the case where $\mathrm{X}_0=\Spec(\mathrm{A}_0)$ is affine; then $\sh{L}_0$ is ample \sref[II]{II.5.1.4}, and since $\operatorname{Ass}(\sh{O}_{\mathrm{X}_0})\cup\mathrm{Z}_0$ is finite, the conclusion follows from \sref[II]{II.4.5.4}.
\end{enumerate}
\end{proof}

\begin{corollary}[21.9.12]
\label{IV.21.9.12}
Let $\mathrm{A}$ be a Henselian local ring \sref{IV.18.5.8}, $\mathrm{S}=\Spec(\mathrm{A})$, $s_0$ the closed point of $\mathrm{S}$, $f:\mathrm{X}\to\mathrm{S}$ a separated morphism of finite presentation, and suppose that the set $\mathrm{X}_0=f^{-1}(s_0)$ is of dimension $\leq 1$.
Then, for every closed sub-scheme $\mathrm{X}_0'$ of $\mathrm{X}$, having $\mathrm{X}_0$ as underlying set and of finite presentation over $\mathrm{S}$, the canonical homomorphism \sref{IV.21.3.2.4}{} $\operatorname{Pic}(\mathrm{X})\to\operatorname{Pic}(\mathrm{X}_0')$ is surjective.
\end{corollary}

\begin{proof}
Let us first show that it suffices to prove the corollary when the local ring $\mathrm{A}$ is moreover Noetherian.
Indeed \sref{IV.18.6.15}{} we can write $\mathrm{A}=\varinjlim\mathrm{A}_\lambda$, where the $\mathrm{A}_\lambda$ are local Noetherian and Henselian rings, the homomorphisms $\mathrm{A}_\lambda\to\mathrm{A}$ being local;
there exists in addition an index $\alpha$ and a separated morphism of finite presentation $f_\alpha:\mathrm{X}_\alpha\to\mathrm{S}_\alpha=\Spec(\mathrm{A}_\alpha)$ such that $\mathrm{X}$ and $f$ are deduced from $\mathrm{X}_\alpha$ and $f_\alpha$ by base change $\mathrm{S}\to\mathrm{S}_\alpha$ \sref{IV.8.10.5},{}~(v);
with the usual notation \sref{IV.8.8.1}, the morphisms $f_\lambda:\mathrm{X}_\lambda\to\mathrm{S}_\lambda$ will be separated for $\lambda\geq\alpha$ and we will have $\mathrm{X}=\mathrm{X}_\lambda\times_{\mathrm{S}_\lambda}\mathrm{S}$ \sref{IV.8.6.3};
Moreover, we can assume that if $s_{0\alpha}$ is the unique closed point of $\mathrm{S}_\alpha$, there is a closed sub-scheme $\mathrm{X}'_{0\alpha}$ of $\mathrm{X}_\alpha$, such that $\mathrm{X}_0'=\mathrm{X}'_{0\alpha}\times_{\mathrm{S}_\alpha}\mathrm{S}$ \sref{IV.8.6.3};
this $\mathrm{X}'_{0\alpha}$ has the same underlying set as $f_\alpha^{-1}(s_{0\alpha})$.
For $\lambda\geq\alpha$, let $\mathrm{X}'_{0\lambda}$ denote its base change to $\mathrm{S}_\lambda$.
By transitivity of fibres and \sref{IV.4.1.4}, we have $\dim(\mathrm{X}'_{0\lambda})=\dim(\mathrm{X}'_0)\leq1$.
It remains to see that, if $\operatorname{Pic}(\mathrm{X}_\lambda)\to\operatorname{Pic}(\mathrm{X}'_{0\lambda})$ is surjective for every $\lambda\geq\alpha$, then the same is true for $\mathrm{X}$ and $\mathrm{X}'_0$.
We evidently have the commutative diagram
\[
  \xymatrix{
    \operatorname{Pic}(\mathrm{X}_\lambda) \ar[r] \ar[d] & \operatorname{Pic}(\mathrm{X}) \ar[d]\\
    \operatorname{Pic}(\mathrm{X}'_{0\lambda}) \ar[r] & \operatorname{Pic}(\mathrm{X}_0')
  }
\]
For every invertible $\sh{O}_{\mathrm{X}'_0}$-module $\sh{L}'_0$, there exist $\lambda\geq\alpha$ and an invertible $\sh{O}_{\mathrm{X}'_{0\lambda}}$-module $\sh{L}'_{0\lambda}$ from which $\sh{L}'_0$ is obtained by the base change $\mathrm{S}\to\mathrm{S}_\lambda$ \sref{IV.8.5.2}{} and \sref{IV.8.5.5}.
Choose an invertible $\sh{O}_{\mathrm{X}_\lambda}$-module $\sh{L}_\lambda$ whose class maps to $\operatorname{cl}(\sh{L}'_{0\lambda})$, and let $\sh{L}$ be its base change to $\mathrm{X}$; by the commutativity of the preceding diagram, $\operatorname{cl}(\sh{L}'_0)$ is the image of $\operatorname{cl}(\sh{L})$.

Suppose therefore that $\mathrm{A}$ is Noetherian, and verify that $\mathrm{X}$ and $\mathrm{X}'_0$ satisfy the conditions of \sref{IV.21.9.11},{}~(ii).
By hypothesis, $\dim(\mathrm{X}'_0)\leq1$.
To verify condition $2^{\mathrm{o}}$ of \sref{IV.21.9.11}, let $\mathrm{Y}$ be the reduced subprescheme of $\mathrm{X}$ having a locally closed subset $Y$ as underlying space.
The morphism $g:\mathrm{Y}\to\mathrm{S}$ induced by $f$ is quasi-finite at each point $x_i$ of $Y\cap\mathrm{X}_0$ ($1\leq i\leq n$).
By \sref{IV.18.5.11},{}~c), $\mathrm{Y}$ is the sum of the open subpreschemes $\mathrm{Y}_i=\Spec(\sh{O}_{\mathrm{Y},x_i})$ ($1\leq i\leq n$) and a prescheme $\mathrm{Y}''$ such that $\mathrm{Y}''\cap\mathrm{X}_0=\emp$; moreover, the canonical injections $j_i:\mathrm{Y}_i\to\mathrm{X}$ are such that $f\circ j_i$ is finite.
Since $f$ is separated, $j_i$ is also finite, and hence $\mathrm{Y}_i$ is closed in $\mathrm{X}$.
The required subset is therefore $Y'=\bigcup_{1\leq i\leq n}Y_i$.

It remains to verify the supplementary hypothesis of \sref{IV.21.9.11},{}~(ii).
Now $\mathrm{X}'_0$, being a separated curve over the field $\boldsymbol{k}(s_0)$, is a quasi-projective $\boldsymbol{k}(s_0)$-scheme (chap.~V)\footnote{The reader will verify that \sref{IV.21.9.12} is not used to prove this property in Chapter~V.}; hence an ample $\sh{O}_{\mathrm{X}'_0}$-module exists \sref[II]{II.5.3.1}, and this completes the proof.
\end{proof}

\begin{remark}[21.9.13]
\label{IV.21.9.13}
Under the conditions of \sref{IV.21.9.12}, supposing in addition $f$ proper, the morphism $f$ is even \emph{projective}:
indeed, if $\sh{L}_0$ is an ample $\sh{O}_{\mathrm{X}'_0}$-module, there exists, by virtue of \sref{IV.21.9.12}{}, an invertible $\sh{O}_\mathrm{X}$-module $\sh{L}$ whose inverse image in $\mathrm{X}_0'$ is isomorphic to $\sh{L}_0$, hence is ample and, as every neighbourhood of $s_0$ in $\mathrm{S}$ is necessarily the whole of $\mathrm{S}$, we deduce then from \sref{IV.9.6.4}{} that $\sh{L}$ is an ample $\sh{O}_\mathrm{X}$-module, whence the conclusion (\sref[II]{II.5.3.1}{} and \sref[II]{II.5.5.3}).
\end{remark}
\subsection{Inverse images and direct images of 1-codimensional cycles}
\label{subsection:IV.21.10}

In a later chapter devoted to intersection theory, the notions of inverse image and direct image of cycles will be developed systematically.
In the present section we shall confine ourselves to defining these notions in certain useful special cases; for $1$-codimensional cycles, the definitions will be chosen so as to be compatible with the corresponding definitions for divisors \sref{IV.21.4}{} and \sref{IV.21.5}, in view of the homomorphism $\operatorname{cyc}$ defined in \sref{IV.21.6}.

\begin{env}[21.10.1]
\label{IV.21.10.1}
\oldpage[IV-4]{289}
\enlargethispage{2pt}
Let $\mathrm{X}$ and $\mathrm{X}'$ be two locally Noetherian preschemes, $f:\mathrm{X}'\to\mathrm{X}$ a morphism, and $\mathrm{T}$ a subset of $\mathrm{X}$.
Suppose that the image under $f$ of every maximal point of $\mathrm{X}'$ is a maximal point of $\mathrm{X}$, and that, for every $x'\in\mathrm{X}'^{(1)}$ (that is, $\dim(\sh{O}_{\mathrm{X}',x'})=1$), one of the following three conditions holds for $x=f(x')$:
\begin{enumerate}[label=(\roman*)]
  \item $x\notin\mathrm{T}$;
  \item $x\in\mathrm{X}^{(1)}$ and $\sh{O}_{\mathrm{X}',x'}$ is a flat $\sh{O}_{\mathrm{X},x}$-module;
  \item the ring $\sh{O}_{\mathrm{X},x}$ is factorial and $\mathfrak{m}_{x'}\notin\operatorname{Ass}(\sh{O}_{\mathrm{X}',x'})$.
\end{enumerate}
Under these conditions, for every $1$-codimensional cycle $\mathrm{Z}$ on $\mathrm{X}$ whose support is contained in $\mathrm{T}$, we shall define a $1$-codimensional cycle $\mathrm{Z}'=f^*(\mathrm{Z})$ so that $\mathrm{Z}\mapsto f^*(\mathrm{Z})$ is a homomorphism of ordered groups from the subgroup of $\mathfrak{Z}^1(\mathrm{X})$ formed by cycles supported in $\mathrm{T}$ to the ordered group $\mathfrak{Z}^1(\mathrm{X}')$.
Write
\[
\label{IV.21.10.1.1}
  \mathrm{Z} = \sum_{x\in\mathrm{T}\cap\mathrm{X}^{(1)}} n_x\,\overline{\{x\}}
  \tag{21.10.1.1}
\]
where the family of the $x\in\mathrm{T}\cap\mathrm{X}^{(1)}$ such that $n_x\neq 0$ is locally finite.
For every $x'\in\mathrm{X}'^{(1)}$, define an integer $n_{x'}$ as follows, putting $x=f(x')$:
\begin{enumerate}[label=\arabic*\textsuperscript{o}]
  \item if $x\notin\mathrm{T}$, we set $n_{x'}=0$;
  \item if $x\in\mathrm{X}^{(1)}$ and if $\sh{O}_{\mathrm{X}',x'}$ is a flat $\sh{O}_{\mathrm{X},x}$-module, we know \sref{IV.6.1.1} that $\dim(\sh{O}_{\mathrm{X}',x'}/\mathfrak{m}_x\sh{O}_{\mathrm{X}',x'})=0$, in other words $\sh{O}_{\mathrm{X}',x'}/\mathfrak{m}_x\sh{O}_{\mathrm{X}',x'}$ is an $\sh{O}_{\mathrm{X}',x'}$-module of finite length $\lambda_{x'}$, and we set $n_{x'}=\lambda_{x'}n_x$;
  \item if $\sh{O}_{\mathrm{X},x}$ is factorial and $\mathfrak{m}_{x'}\notin\operatorname{Ass}(\sh{O}_{\mathrm{X}',x'})$, then the canonical homomorphism $\operatorname{cyc}:\operatorname{Div}(\sh{O}_{\mathrm{X},x})\to\mathfrak{Z}^1(\Spec(\sh{O}_{\mathrm{X},x}))$ is bijective \sref{IV.21.6.9}.
Furthermore, since $\dim(\sh{O}_{\mathrm{X}',x'})=1$ and $\mathfrak{m}_{x'}\notin\operatorname{Ass}(\sh{O}_{\mathrm{X}',x'})$, $\operatorname{Ass}(\sh{O}_{\mathrm{X}',x'})$ consists precisely of the maximal points of $\Spec(\sh{O}_{\mathrm{X}',x'})$; the hypothesis on $f$ therefore implies
$f(\operatorname{Ass}(\sh{O}_{\mathrm{X}',x'}))\subset\operatorname{Ass}(\sh{O}_{\mathrm{X},x})$.
It follows from \sref{IV.21.4.5},{}~(ii) that $f^*:\operatorname{Div}(\sh{O}_{\mathrm{X},x})\to\operatorname{Div}(\sh{O}_{\mathrm{X}',x'})$ is defined.
Finally, since $x'$ is the unique closed point of $\Spec(\sh{O}_{\mathrm{X}',x'})$, $\mathfrak{Z}^1(\Spec(\sh{O}_{\mathrm{X}',x'}))$ is canonically isomorphic to $\bb{Z}$.
\end{enumerate}
We have thus defined a canonical composite homomorphism
\[
\label{IV.21.10.1.2}
  \mathfrak{Z}^1(\Spec(\sh{O}_{\mathrm{X},x})) \xrightarrow{\operatorname{cyc}^{-1}} \operatorname{Div}(\sh{O}_{\mathrm{X},x}) \xrightarrow{f^*} \operatorname{Div}(\sh{O}_{\mathrm{X}',x'}) \xrightarrow{\operatorname{cyc}} \mathfrak{Z}^1(\Spec(\sh{O}_{\mathrm{X}',x'})) \xrightarrow{\sim} \bb{Z}
  \tag{21.10.1.2}
\]
If $\mathrm{Z}_x$ is the cycle
$\sum_{y\in\Spec(\sh{O}_{\mathrm{X},x})\cap\mathrm{X}^{(1)}}n_y\bigl(\overline{\{y\}}\cap\Spec(\sh{O}_{\mathrm{X},x})\bigr)$ on $\Spec(\sh{O}_{\mathrm{X},x})$, we take $n_{x'}$ to be the \emph{image} of $\mathrm{Z}_x$ under \sref{IV.21.10.1.2}.
\end{env}

\oldpage[IV-4]{290}
\begin{env}[21.10.2]
\label{IV.21.10.2}
We propose to show that:
\begin{enumerate}[label=\emph{\Alph*)}]
  \item When two of the conditions $1^{\mathrm{o}}$, $2^{\mathrm{o}}$, $3^{\mathrm{o}}$ of \sref{IV.21.10.1} are \emph{simultaneously} satisfied, the corresponding values of $n_{x'}$ \emph{coincide}.
  \item The set of $x'\in\mathrm{X}'^{(1)}$ such that $n_{x'}\neq 0$ is \emph{locally finite} in $\mathrm{X}'$.
\end{enumerate}

To prove A), let us first suppose that $x\notin\mathrm{T}$ and that $x$ satisfies one of the conditions $2^{\mathrm{o}}$ or $3^{\mathrm{o}}$ of \sref{IV.21.10.1}; then $n_x=0$ and if one is in case $2^{\mathrm{o}}$, we have $n_{x'}=0$; if one is in case $3^{\mathrm{o}}$, we have $\mathrm{Z}_x=0$ since $\operatorname{Supp}(\mathrm{Z})\subset\mathrm{T}$, so again $n_{x'}=0$.
It remains to consider the case where one is simultaneously in cases $2^{\mathrm{o}}$ and $3^{\mathrm{o}}$; then, since $\dim(\sh{O}_{\mathrm{X},x})=1$, $\sh{O}_{\mathrm{X},x}$ is a discrete valuation ring; if $t$ is a uniformiser of this ring, the divisor corresponding to $\mathrm{Z}_x$ is $\operatorname{div}(t^{n_x})$ in $\operatorname{Div}(\sh{O}_{\mathrm{X},x})$, and its image in $\operatorname{Div}(\sh{O}_{\mathrm{X}',x'})$ is the divisor of $t'^{n_x}$, where $t'$ is the (regular) element of $\sh{O}_{\mathrm{X}',x'}$ that is the image of $t$.
We may clearly reduce to the case $n_x=1$; then the definition \sref{IV.21.6.5.1} shows that the image of $\mathrm{Z}_x$ under \sref{IV.21.10.1.2} is $\lambda_{x'}$, completing the proof of A).

Let us now prove B).
Set $\mathrm{T}_0=\operatorname{Supp}(\mathrm{Z})\subset\mathrm{T}$, $\mathrm{T}_0'=f^{-1}(\mathrm{T}_0)$; it suffices to prove that the relation $n_{x'}\neq 0$ implies that $x'$ belongs to the set of maximal points of $\mathrm{T}_0'$, this last being locally finite in $\mathrm{X}'$.
It is immediate that $x'\in\mathrm{T}_0'$ necessarily; if $x'$ was not maximal in the closed set $\mathrm{T}_0'$, there would exist
\oldpage[IV-4]{291}
a generisation $y'$ of $x'$ in $\mathrm{T}_0'$, distinct from $x'$, and since $x'\in\mathrm{X}'^{(1)}$, $y'$ would necessarily be a maximal point of $\mathrm{X}'$; as a result $y=f(y')$ would be a maximal point of $\mathrm{X}$ by hypothesis; but this is absurd since $y\in\mathrm{T}_0$ and $\mathrm{T}_0$ is purely of codimension $1$ in $\mathrm{X}$.
\end{env}

\begin{env}[21.10.3]
\label{IV.21.10.3}
We can now set
\[
  f^*(Z)=\sum_{x'\in\mathrm{X}'^{(1)}} n_{x'}\overline{\{x'\}},
\]
the sum of the second member having a meaning by virtue of what we proved in \sref{IV.21.10.2}; we say that the $1$-codimensional cycle $f^*(Z)$ is the \emph{inverse image} of $\mathrm{Z}$ by $f$.
It is immediate that the map $f^*$ thus defined is a homomorphism of ordered groups.
Furthermore, if $\mathrm{U}$ is an open of $\mathrm{X}$, $\mathrm{V}$ an open of $\mathrm{X}'$ such that $f(\mathrm{V})\subset\mathrm{U}$ and $f':\mathrm{V}\to\mathrm{U}$ the restriction of $f$, it follows immediately from the definitions that one has
\[
\label{IV.21.10.3.1}
  f'^*(Z|\mathrm{U})=f^*(Z)|\mathrm{V}.
  \tag{21.10.3.1}
\]

Let us denote by $\Gamma_{\mathrm{T}}(\mathscr{Z}^1_{\mathrm{X}})$ the largest sub-sheaf of commutative groups of $\mathscr{Z}^1_{\mathrm{X}}$ with support contained in $\mathrm{T}$; it follows from the relation \sref{IV.21.10.3.1} that the maps $\Gamma(\mathrm{U},\Gamma_{\mathrm{T}}(\mathscr{Z}^1_{\mathrm{X}}))\to\Gamma(\mathrm{V},\mathscr{Z}^1_{\mathrm{X}'})$ that we have just defined define a homomorphism of sheaves of ordered commutative groups on $\mathrm{X}'$
\[
  \psi^*(\Gamma_{\mathrm{T}}(\mathscr{Z}^1_{\mathrm{X}}))\to\mathscr{Z}^1_{\mathrm{X}'}
\]
where $\psi$ is the continuous map underlying the morphism $f$.
\end{env}

\begin{proposition}[21.10.4]
\label{IV.21.10.4}
Suppose satisfied the conditions of \sref{IV.21.10.1}.
Then, for every divisor $\mathrm{D}$ on $\mathrm{X}$ such that $\operatorname{Supp}(\mathrm{D})\subset\mathrm{T}$ and such that $f^*(\mathrm{D})$ is defined \sref{IV.21.4.2}, we have
\[
\label{IV.21.10.4.1}
  \operatorname{cyc}(f^*(\mathrm{D}))=f^*(\operatorname{cyc}(\mathrm{D})).
  \tag{21.10.4.1}
\]
\end{proposition}

\begin{proof}
The question being local on $\mathrm{X}$, we can restrict to the case where $\mathrm{X}=\operatorname{Spec}(\mathrm{A})$ is affine, where $\mathrm{D}=\operatorname{div}(t)$, where $t$ is a regular non-invertible element of $\mathrm{A}$, and where the sub-prescheme $\mathrm{Y}(\mathrm{D})$ \sref{IV.21.2.12} has only a single maximal point $y$, so that $\operatorname{cyc}(\mathrm{D})=n_y.\overline{\{y\}}$, where $n_y$ is the length of $\sh{O}_{\mathrm{X},y}/t_y\sh{O}_{\mathrm{X},y}$ \sref{IV.21.6.5.1}.
We saw in the proof of \sref{IV.21.10.2},~B) that the points $x'\in\mathrm{X}'$ such that $\operatorname{mult}_{x'}(f^*(\operatorname{cyc}(\mathrm{D})))\neq 0$ are maximal points of $f^{-1}(\overline{\{y\}})$.
If the point $x'$ is in case $3^{\mathrm{o}}$ of \sref{IV.21.10.1}, equality of the multiplicities at $x'$ of the two sides of \sref{IV.21.10.4.1} follows from the definition of $n_{x'}$ by means of the homomorphism \sref{IV.21.10.1.2}.
Suppose on the contrary that $x'$ is in case $2^{\mathrm{o}}$ of \sref{IV.21.10.1}, and let $x=f(x')$; since $x\in\mathrm{X}^{(1)}\cap\overline{\{y\}}$, we have necessarily $x=y$.
Now $f^*(\mathrm{D})=\operatorname{div}(t')$, where $t'$ is the image of $t$ in $\Gamma(\mathrm{X}',\sh{O}_{\mathrm{X}'})$, and $\mathrm{Y}(f^*(\mathrm{D}))=f^{-1}(\mathrm{Y}(\mathrm{D}))$.
The multiplicity of $f^*(\mathrm{D})$ at $x'$ is therefore the length $n'_{x'}$ of the $\sh{O}_{\mathrm{X}',x'}$-module
$\sh{O}_{\mathrm{X}',x'}/t'_{x'}\sh{O}_{\mathrm{X}',x'}=(\sh{O}_{\mathrm{X},y}/t_y\sh{O}_{\mathrm{X},y})\otimes_{\sh{O}_{\mathrm{X},y}}\sh{O}_{\mathrm{X}',x'}$.
It follows from \sref{IV.4.7.1} that $n'_{x'}=\lambda_{x'}n_y$, so the multiplicities at $x'$ of the two sides of \sref{IV.21.10.4.1} are again equal, completing the proof.
\end{proof}

\begin{env}[21.10.5]
\label{IV.21.10.5}
\oldpage[IV-4]{292}
Suppose now that $f:\mathrm{X}'\to\mathrm{X}$ is a morphism of locally Noetherian preschemes, transforming every maximal point of $\mathrm{X}'$ into a maximal point of $\mathrm{X}$, and suppose furthermore that for every \emph{closed rare} subset $\mathrm{T}$ of $\mathrm{X}$, $f$ satisfies the conditions of \sref{IV.21.10.1}; this means that for all $x'\in\mathrm{X}'^{(1)}$, either $x=f(x')$ is a maximal point of $\mathrm{X}$, or $x'$ satisfies one of the conditions (ii), (iii) of \sref{IV.21.10.1}.
Since every $1$-codimensional cycle has rare support in $\mathrm{X}$, $f^*(\mathrm{Z})$ is therefore defined for \emph{every} $1$-codimensional cycle $\mathrm{Z}$ on $\mathrm{X}$; by virtue of \sref{IV.21.10.3.1}, we have thus defined a homomorphism of sheaves of ordered commutative groups
\[
  \psi^*(\mathscr{Z}^1_{\mathrm{X}})\to\mathscr{Z}^1_{\mathrm{X}'}.
\]

If furthermore, for every divisor $\mathrm{D}$ on $\mathrm{X}$, $f^*(\mathrm{D})$ is defined \sref{IV.21.4.5}, the fact that the support of $\mathrm{D}$ is rare in $\mathrm{X}$ \sref{IV.21.6.6} implies that \sref{IV.21.10.4} is applicable, and we therefore have the formula \sref{IV.21.10.4.1} for every divisor $\mathrm{D}$ on $\mathrm{X}$.
In particular:
\end{env}

\begin{proposition}[21.10.6]
\label{IV.21.10.6}
Let $\mathrm{X}$, $\mathrm{X}'$ be two locally Noetherian preschemes, $f:\mathrm{X}'\to\mathrm{X}$ a flat morphism.
Then $f^*(\mathrm{Z})$ is defined for every $1$-codimensional cycle $\mathrm{Z}$ on $\mathrm{X}$, $f^*(\mathrm{D})$ is defined for every divisor $\mathrm{D}$ on $\mathrm{X}$ and we have the relation \sref{IV.21.10.4.1}.
\end{proposition}

\begin{proof}
Indeed, if $x'\in\mathrm{X}'^{(1)}$ is such that $x=f(x')$ is not maximal, it follows from \sref{IV.6.1.1} that we have necessarily $x\in\mathrm{X}^{(1)}$, so one is in case (ii) of \sref{IV.21.10.1}.
We can therefore apply \sref{IV.21.10.5}, taking into account \sref{IV.21.4.5} and \sref{IV.2.3.4}.
\end{proof}

\begin{remark}[21.10.7]
\label{IV.21.10.7}
The existence of $f^*(\mathrm{Z})$ for every $1$-codimensional cycle $\mathrm{Z}$ on $\mathrm{X}$ already follows from the hypothesis that $f$ is \emph{flat at all points $x'$ of $\mathrm{X}'$ of codimension $\leq 1$} (i.e.\ such that $\dim\sh{O}_{\mathrm{X}',x'}\leq 1$); indeed, for every maximal $z'\in\mathrm{X}'$, it follows from \sref{IV.6.1.1} that $z=f(z')$ is a maximal point of $\mathrm{X}$ since $\dim\sh{O}_{\mathrm{X},z}\leq\dim\sh{O}_{\mathrm{X}',z'}=0$.
Moreover, if $x'\in\mathrm{X}'^{(1)}$, $x=f(x')$ is maximal or belongs to $\mathrm{X}^{(1)}$ by \sref{IV.6.1.1}; one can again apply \sref{IV.21.10.5}.
\end{remark}

\begin{proposition}[21.10.8]
\label{IV.21.10.8}
Let $\mathrm{X}$, $\mathrm{X}'$, $\mathrm{X}''$ be three locally Noetherian preschemes, $f:\mathrm{X}'\to\mathrm{X}$, $g:\mathrm{X}''\to\mathrm{X}'$ two morphisms; suppose that $f$ (resp.\ $g$) is flat at every point of codimension $\leq 1$ in $\mathrm{X}'$ (resp.\ $\mathrm{X}''$).
Then $f\circ g$ is flat at every point of codimension $\leq 1$ in $\mathrm{X}''$ and for every $1$-codimensional cycle $\mathrm{Z}$ on $\mathrm{X}$, we have $(f\circ g)^*(\mathrm{Z})=g^*(f^*(\mathrm{Z}))$.
\end{proposition}

\begin{proof}
The first assertion follows from \sref{IV.6.1.1} and \sref{IV.2.1.6}.
The second follows from the fact that $f$ (resp.\ $g$, resp.\ $f\circ g$) sends maximal points of $\mathrm{X}'$ (resp.\ $\mathrm{X}''$, resp.\ $\mathrm{X}''$) to maximal points of $\mathrm{X}$ (resp.\ $\mathrm{X}'$, resp.\ $\mathrm{X}$), and from the fact that, if $x''\in\mathrm{X}''^{(1)}$ is such that $x'=g(x'')\in\mathrm{X}'^{(1)}$ and $x=f(x')\in\mathrm{X}^{(1)}$, then
\[
  \operatorname{length}(\sh{O}_{\mathrm{X}'',x''}/\mathfrak{m}_x\sh{O}_{\mathrm{X}'',x''})=\operatorname{length}(\sh{O}_{\mathrm{X}'',x''}/\mathfrak{m}_{x'}\sh{O}_{\mathrm{X}'',x''})\cdot\operatorname{length}(\sh{O}_{\mathrm{X}',x'}/\mathfrak{m}_x\sh{O}_{\mathrm{X}',x'})
\]
Indeed, if we set $\mathrm{A}=\sh{O}_{\mathrm{X},x}$, $\mathrm{A}'=\sh{O}_{\mathrm{X}',x'}$, $\mathrm{A}''=\sh{O}_{\mathrm{X}'',x''}$, $\mathfrak{m}=\mathfrak{m}_x$, $\mathfrak{m}'=\mathfrak{m}_{x'}$, we have $\mathrm{A}''/\mathfrak{m}\mathrm{A}''=(\mathrm{A}'/\mathfrak{m}\mathrm{A}')\otimes_{\mathrm{A}'}\mathrm{A}''$, and the formula
\[
  \operatorname{length}_{\mathrm{A}''}(\mathrm{A}''/\mathfrak{m}\mathrm{A}'')=\operatorname{length}_{\mathrm{A}'}(\mathrm{A}'/\mathfrak{m}\mathrm{A}')\cdot\operatorname{length}_{\mathrm{A}''}(\mathrm{A}''/\mathfrak{m}'\mathrm{A}'')
\]
follows from \sref{IV.4.7.1} and the hypothesis of flatness of $\mathrm{A}''$ over $\mathrm{A}'$.
\end{proof}

\begin{env}[21.10.9]
\label{IV.21.10.9}
Let $\mathrm{X}$ be a locally Noetherian prescheme, $\Lambda$ a commutative group, written additively, and considered as a $\bb{Z}$-module; we will again denote by $\bb{Z}$ and $\Lambda$ the simple sheaves on $\mathrm{X}$ associated to the constant presheaves respectively equal to $\bb{Z}$ and $\Lambda$ \sref[0]{0.3.6}.
We call \emph{sheaf of germs of $1$-codimensional cycles with coefficients in $\Lambda$} the sheaf of commutative groups $\mathscr{Z}^1_{\mathrm{X}}\otimes_{\bb{Z}}\Lambda$.
If $\Lambda=\bb{Q}$, we say that the sections
\oldpage[IV-4]{293}
of $\mathscr{Z}^1_{\mathrm{X}}\otimes_{\bb{Z}}\bb{Q}$ over $\mathrm{X}$ are the \emph{$1$-codimensional cycles with rational coefficients}.
Since the fibres of $\mathscr{Z}^1_{\mathrm{X}}$ are torsion-free $\bb{Z}$-modules \sref{IV.21.6.3}, the canonical homomorphism $\mathscr{Z}^1_{\mathrm{X}}\to\mathscr{Z}^1_{\mathrm{X}}\otimes_{\bb{Z}}\bb{Q}$ is \emph{injective}, so that the $1$-codimensional cycles are identified with $1$-codimensional cycles with rational coefficients.
\end{env}

\begin{env}[21.10.10]
\label{IV.21.10.10}
We will now see that under certain conditions, we can enlarge the definition of $f^*(Z)$ given in \sref{IV.21.10.3} for a $1$-codimensional cycle $\mathrm{Z}$ on $\mathrm{X}$, but to take for $f^*(Z)$ a $1$-codimensional cycle \emph{with rational coefficients} on $\mathrm{X}'$.
The most general case where we place ourselves is the one where $f$ transforms every maximal point of $\mathrm{X}'$ into a maximal point of $\mathrm{X}$, and where, at every point $x'\in\mathrm{X}'^{(1)}$, we have one of the conditions (i), (ii), (iii) of \sref{IV.21.10.1} \emph{or a fourth condition} (setting $x=f(x')$):
\begin{enumerate}
  \item[(iv)] $x\in\mathrm{X}^{(1)}$, $\mathfrak{m}_{x}\notin\operatorname{Ass}(\sh{O}_{\mathrm{X},x})$, and, in addition, if $\mathrm{A}=\widehat{\sh{O}}_{\mathrm{X},x}$, $\mathrm{A}'=\widehat{\sh{O}}_{\mathrm{X}',x'}$, and $\mathrm{K}$ denotes the total ring of fractions of $\mathrm{A}$, then $\mathrm{A}'$ is a finite $\mathrm{A}$-algebra and $\mathrm{K}'=\mathrm{A}'\otimes_{\mathrm{A}}\mathrm{K}$ is a free $\mathrm{K}$-module.
\end{enumerate}

Let us note then $r_{x'}$ the rank of the free $\mathrm{K}$-module $\mathrm{K}'$, $q_{x'}$ the degree of $\boldsymbol{k}(x')=\boldsymbol{k}(\mathrm{A}')$ over $\boldsymbol{k}(x)=\boldsymbol{k}(\mathrm{A})$, and set $\mu_{x'}=r_{x'}/q_{x'}$.
For a $1$-codimensional cycle with support contained in $\mathrm{T}$, given by \sref{IV.21.10.1.1} and for all $x'\in\mathrm{X}'^{(1)}$, we define $c_{x'}\in\bb{Q}$ to equal $n_{x'}$ when one is in one of the cases $1^{\mathrm{o}}$, $2^{\mathrm{o}}$, $3^{\mathrm{o}}$ of \sref{IV.21.10.1}; but there remains here a fourth possibility:
\begin{enumerate}
  \item[$4^{\mathrm{o}}$] if $x'$ satisfies condition (iv) above, we set $c_{x'}=\mu_{x'}n_x\in\bb{Q}$.
\end{enumerate}
\end{env}

\begin{env}[21.10.11]
\label{IV.21.10.11}
We must still prove that when condition $4^{\mathrm{o}}$ of \sref{IV.21.10.10} is satisfied \emph{at the same time} as one of the conditions $1^{\mathrm{o}}$, $2^{\mathrm{o}}$, $3^{\mathrm{o}}$ of \sref{IV.21.10.1}, we have $n_{x'}=c_{x'}$.
This is evident if $x\notin\mathrm{T}$ since then $n_x=0$.
To study the other two cases, note that $\mathfrak{m}_x\sh{O}_{\mathrm{X}',x'}$ is closed for the $\mathfrak{m}_{x'}$-adic topology; hence the completion of $\sh{O}_{\mathrm{X}',x'}/\mathfrak{m}_x\sh{O}_{\mathrm{X}',x'}$ for that topology is $\mathrm{A}'/\mathfrak{m}_x\mathrm{A}'$.
If cases $2^{\mathrm{o}}$ and $4^{\mathrm{o}}$ hold simultaneously, $\sh{O}_{\mathrm{X}',x'}/\mathfrak{m}_x\sh{O}_{\mathrm{X}',x'}$ is discrete for the $\mathfrak{m}_{x'}$-adic topology and is therefore isomorphic to $\mathrm{A}'/\mathfrak{m}_x\mathrm{A}'$.
Since $\mathrm{A}'$ is a finite and flat $\mathrm{A}$-algebra $(\mathbf{0}_{\mathrm{III}},10.2.3)$, it is a free $\mathrm{A}$-module $(\mathbf{0}_{\mathrm{III}},10.1.3)$.
The rank of $\mathrm{A}'/\mathfrak{m}_x\mathrm{A}'$ over $\mathrm{A}/\mathfrak{m}_x\mathrm{A}=\boldsymbol{k}(x)$ is thus equal to the rank $r_{x'}$ of $\mathrm{K}'$ over $\mathrm{K}$.
On the other hand, this rank is the product of the length $\lambda_{x'}$ of the $\sh{O}_{\mathrm{X}',x'}$-module $\sh{O}_{\mathrm{X}',x'}/\mathfrak{m}_x\sh{O}_{\mathrm{X}',x'}$ by $[\boldsymbol{k}(x'): \boldsymbol{k}(x)]=q_{x'}$.
Consequently $\mu_{x'}=r_{x'}/q_{x'}=\lambda_{x'}$ in this case.

Suppose finally that one is simultaneously in cases $3^{\mathrm{o}}$ and $4^{\mathrm{o}}$.
Then, since $\dim(\sh{O}_{\mathrm{X},x})=1$, $\sh{O}_{\mathrm{X},x}$ is a discrete valuation ring, hence regular.
On the other hand, $\sh{O}_{\mathrm{X}',x'}$ is of dimension $1$, and since $\mathfrak{m}_{x'}\notin\operatorname{Ass}(\sh{O}_{\mathrm{X}',x'})$, $\sh{O}_{\mathrm{X}',x'}$ is a Cohen-Macaulay ring $(\mathbf{0}$, $16.4.6)$; finally, since $\mathrm{A}'$ is an A-module of finite type, $\mathrm{A}'/\mathfrak{m}_x\mathrm{A}'$ is a $\boldsymbol{k}(x)$-vector space of finite rank; since $\sh{O}_{\mathrm{X}',x'}/\mathfrak{m}_x\sh{O}_{\mathrm{X}',x'}$ is contained in $\mathrm{A}'/\mathfrak{m}_x\mathrm{A}'$, it is also a $\boldsymbol{k}(x)$-vector space of finite rank, hence an Artinian ring.
The application of \sref{IV.6.1.5} then shows that $\sh{O}_{\mathrm{X}',x'}$ is a flat $\sh{O}_{\mathrm{X},x}$-module, so one is also in case $2^{\mathrm{o}}$, and we conclude by what was seen above.

Accordingly, in the situation under consideration, set
\[
  f^*(Z)=\sum_{x'\in\mathrm{X}'^{(1)}} c_{x'}.\overline{\{x'\}}
\]
\oldpage[IV-4]{294}
which is thus a $1$-codimensional cycle on $\mathrm{X}'$ with \emph{rational} coefficients.
We have again defined a homomorphism $f^*$ of ordered groups satisfying \sref{IV.21.10.3.1}, and hence a homomorphism of sheaves of ordered commutative groups
\[
  \psi^*(\Gamma_{\mathrm{T}}(\mathscr{Z}^1_{\mathrm{X}}))\to\mathscr{Z}^1_{\mathrm{X}'}\otimes_{\bb{Z}}\bb{Q}.
\]
When $f$ satisfies the preceding conditions for every \emph{closed rare} subset $\mathrm{T}$ in $\mathrm{X}$, i.e.\ that for all $x'\in\mathrm{X}'^{(1)}$, either $x=f(x')$ is a maximal point of $\mathrm{X}$, or $x'$ satisfies one of the conditions (ii), (iii) or (iv), $f^*(Z)$ is then defined for \emph{every} $1$-codimensional cycle $\mathrm{Z}$ on $\mathrm{X}$, and we have defined a homomorphism of sheaves of commutative ordered groups
\[
  \psi^*(\mathscr{Z}^1_{\mathrm{X}})\to\mathscr{Z}^1_{\mathrm{X}'}\otimes_{\bb{Z}}\bb{Q}
\]
and hence, by tensoring, a homomorphism of sheaves of ordered $\bb{Q}$-vector spaces
\[
\label{IV.21.10.11.1}
  \psi^*(\mathscr{Z}^1_{\mathrm{X}}\otimes_{\bb{Z}}\bb{Q})\to\mathscr{Z}^1_{\mathrm{X}'}\otimes_{\bb{Z}}\bb{Q}.
  \tag{21.10.11.1}
\]
\end{env}

\begin{remark}[21.10.12]
\label{IV.21.10.12}
When one is in the situation of \sref{IV.21.10.10}, one can effectively, for $1$-codimensional cycles $\mathrm{Z}$ on $\mathrm{X}$, have for $f^*(Z)$ cycles $1$-codimensional with \emph{non-integer} coefficients; in other words, the numbers $\mu_{x'}$ can be non-integers.
An example is obtained by taking the complete integral domain $\mathrm{A}$ and its integral closure $\mathrm{A}'$ from $(\mathbf{0}_{\mathrm{IV}},6.15.11)$,~(ii): the closed point $x'$ of $\Spec(\mathrm{A}')$ satisfies condition (iv) of \sref{IV.21.10.10}, and $\mu_{x'}=\frac12$.
\end{remark}

\begin{lemma}[21.10.13]
\label{IV.21.10.13}
Let $\mathrm{A}$ be a local Noetherian ring of dimension $1$, $t$ a regular element of $\mathrm{A}$ belonging to the maximal ideal $\mathfrak{m}$ (which implies that $\mathfrak{m}\notin\operatorname{Ass}(\mathrm{A})$).
\begin{enumerate}[label=(\roman*)]
  \item For every A-module $\mathrm{M}$ of finite type, the kernel $\mathrm{N}_t(\mathrm{M})$ and the cokernel $\mathrm{P}_t(\mathrm{M})$ of the homothety $t_{\mathrm{M}}:\mathrm{M}\to\mathrm{M}$ of ratio $t$ are of finite length. We set $d_t(\mathrm{M})=\operatorname{length}\mathrm{P}_t(\mathrm{M})-\operatorname{length}\mathrm{N}_t(\mathrm{M})$.
  \item If $0\to\mathrm{M}'\to\mathrm{M}\to\mathrm{M}''\to 0$ is an exact sequence of A-modules of finite type, we have $d_t(\mathrm{M})=d_t(\mathrm{M}')+d_t(\mathrm{M}'')$.
  \item We have $d_t(\mathrm{M})\geq 0$ for every A-module $\mathrm{M}$ of finite type; for $d_t(\mathrm{M})=0$, it is necessary and sufficient that $\mathrm{M}$ be of finite length.
  \item If $\mathrm{K}$ is the total ring of fractions of $\mathrm{A}$ and if $\mathrm{M}\otimes_{\mathrm{A}}\mathrm{K}$ is a free $\mathrm{K}$-module of rank $n$, we have $d_t(\mathrm{M})=n\cdot d_t(\mathrm{A})=n\cdot\operatorname{length}(\mathrm{A}/t\mathrm{A})$.
  \item If $\mathrm{M}$ satisfies the hypotheses of (iv) and if furthermore $t$ is $\mathrm{M}$-regular, we have $\operatorname{length}(\mathrm{M}/t\mathrm{M})=n\cdot\operatorname{length}(\mathrm{A}/t\mathrm{A})$.
\end{enumerate}
\end{lemma}

\begin{proof}
(i) $\operatorname{Spec}(\mathrm{A})$ is formed of the point $\mathfrak{m}$ and the minimal prime ideals $\mathfrak{p}_i$; since by hypothesis $t\notin\mathfrak{p}_i$ for all $i$ (Bourbaki, \emph{Alg.\ comm.}, chap.~IV, \textsection1, n$^{\mathrm{o}}$~1, cor.~3 of prop.~2), the image of $t$ in each of the $\mathrm{A}_{\mathfrak{p}_i}$ is invertible, and the supports of the A-modules of finite type $\mathrm{N}_t(\mathrm{M})$ and $\mathrm{P}_t(\mathrm{M})$ are therefore empty or reduced to $\mathfrak{m}$; we conclude from this $(\mathbf{0}$, $16.1.10)$ that these modules are of finite length.

(ii) Since $t$ is regular, we have an exact sequence
\[
  0\to\mathrm{A}\xrightarrow{t_{\mathrm{A}}}\mathrm{A}\to\mathrm{A}/t\mathrm{A}\to 0
\]
\oldpage[IV-4]{295}
and since $\operatorname{Tor}_i^{\mathrm{A}}(\mathrm{M},\mathrm{A})=0$ for $i\geq 1$, the Tor exact sequence gives on one hand the exact sequence
\[
  0\to\operatorname{Tor}_1^{\mathrm{A}}(\mathrm{M},\mathrm{A}/t\mathrm{A})\to\mathrm{M}\xrightarrow{t_{\mathrm{M}}}\mathrm{M}\to\mathrm{M}/t\mathrm{M}\to 0
\]
and on the other hand, for $i\geq 2$,
\[
  0=\operatorname{Tor}_i^{\mathrm{A}}(\mathrm{M},\mathrm{A})\to\operatorname{Tor}_i^{\mathrm{A}}(\mathrm{M},\mathrm{A}/t\mathrm{A})\to\operatorname{Tor}_{i-1}^{\mathrm{A}}(\mathrm{M},\mathrm{A})=0
\]
hence $\mathrm{N}_t(\mathrm{M})=\operatorname{Tor}_1^{\mathrm{A}}(\mathrm{M},\mathrm{A}/t\mathrm{A})$ and $\operatorname{Tor}_i^{\mathrm{A}}(\mathrm{M},\mathrm{A}/t\mathrm{A})=0$ for $i\geq 2$; the Tor exact sequence gives the exact sequence
\[
  0\to\operatorname{Tor}_1^{\mathrm{A}}(\mathrm{M}',\mathrm{A}/t\mathrm{A})\to\operatorname{Tor}_1^{\mathrm{A}}(\mathrm{M},\mathrm{A}/t\mathrm{A})\to\operatorname{Tor}_1^{\mathrm{A}}(\mathrm{M}'',\mathrm{A}/t\mathrm{A})\to
\]
\[
  \mathrm{M}'\otimes_{\mathrm{A}}(\mathrm{A}/t\mathrm{A})\to\mathrm{M}\otimes_{\mathrm{A}}(\mathrm{A}/t\mathrm{A})\to\mathrm{M}''\otimes_{\mathrm{A}}(\mathrm{A}/t\mathrm{A})\to 0
\]
and this sequence can be written, in view of what precedes
\[
  0\to\mathrm{N}_t(\mathrm{M}')\to\mathrm{N}_t(\mathrm{M})\to\mathrm{N}_t(\mathrm{M}'')\to\mathrm{P}_t(\mathrm{M}')\to\mathrm{P}_t(\mathrm{M})\to\mathrm{P}_t(\mathrm{M}'')\to 0
\]
which proves (ii).

(iii) To prove (iii), we note that for $\mathrm{M}$ there is a composition series $(\mathrm{M}_h)$ of $\mathrm{M}$ such that the quotients $\mathrm{M}_h/\mathrm{M}_{h+1}$ are isomorphic to $\mathrm{A}/\mathfrak{m}$ or to one of the $\mathrm{A}/\mathfrak{p}_i$ (Bourbaki, \emph{Alg.\ comm.}, chap.~IV, \textsection1, n$^{\mathrm{o}}$~4, th.~1), and that for $\mathrm{M}$ to be of finite length, it is necessary and sufficient that all these quotients be isomorphic to $\mathrm{A}/\mathfrak{m}$.
It all comes down therefore (by virtue of (ii)) to proving that $d_t(\mathrm{A}/\mathfrak{m})=0$ and $d_t(\mathrm{A}/\mathfrak{p}_i)>0$.
Now, the image of $t$ in $\mathrm{A}/\mathfrak{m}$ being $0$, we have $\mathrm{N}_t(\mathrm{A}/\mathfrak{m})=\mathrm{A}/\mathfrak{m}$ and $\mathrm{P}_t(\mathrm{A}/\mathfrak{m})=\mathrm{A}/\mathfrak{m}$, whence the first assertion; on the other hand, the image of $t$ in $\mathrm{A}/\mathfrak{p}_i$ is regular, so $\mathrm{N}_t(\mathrm{A}/\mathfrak{p}_i)=0$ and $\mathrm{P}_t(\mathrm{A}/\mathfrak{p}_i)=\mathrm{A}/(t\mathrm{A}+\mathfrak{p}_i)$ which is not reduced to $0$, whence the second assertion.

(iv) There is a basis of $\mathrm{M}\otimes_{\mathrm{A}}\mathrm{K}$ of the form $(x_j/s)_{1\leq j\leq n}$, where $s$ is a regular element of $\mathrm{A}$ and $x_j\in\mathrm{M}$.
Consider the homomorphism $u:\mathrm{A}^n\to\mathrm{M}$ which transforms the element $e_j$ ($1\leq j\leq n$) of the canonical basis of $\mathrm{A}^n$ into $x_j$ and let us show that $\operatorname{Ker}(u)$ and $\operatorname{Coker}(u)$ are of \emph{finite length}; indeed, for each $i$, the image of $s$ in $\mathrm{A}_{\mathfrak{p}_i}$ is invertible, and since $\mathrm{K}_{\mathfrak{p}_i}=\mathrm{A}_{\mathfrak{p}_i}$, the images of the $x_j/s$ in $\mathrm{M}_{\mathfrak{p}_i}$ form a basis of this $\mathrm{A}_{\mathfrak{p}_i}$-module; hence $u_{\mathfrak{p}_i}:\mathrm{A}_{\mathfrak{p}_i}^n\to\mathrm{M}_{\mathfrak{p}_i}$ is bijective.
As in (i), we conclude that the supports of $\operatorname{Ker}(u)$ and of $\operatorname{Coker}(u)$ are contained in $\{\mathfrak{m}\}$, so these modules are of finite length.
This being so, it follows from (ii) and (iii) that $d_t(\mathrm{M})=d_t(\mathrm{A}^n)=n\cdot d_t(\mathrm{A})=n\cdot\operatorname{length}(\mathrm{A}/t\mathrm{A})$.

Finally, it is clear that (v) follows immediately from (iv), since then $\mathrm{N}_t(\mathrm{M})=0$.
\end{proof}

This lemma allows us to generalise \sref{IV.21.10.4}:
\begin{proposition}[21.10.13]
\label{IV.21.10.13bis}
Suppose that $f$ satisfies the conditions of \sref{IV.21.10.10}.
Then, for every divisor $\mathrm{D}$ on $\mathrm{X}$ such that $\operatorname{Supp}(\mathrm{D})\subset\mathrm{T}$ and such that $f^*(\mathrm{D})$ is defined, we have
\[
\label{IV.21.10.13.1}
  \operatorname{cyc}(f^*(\mathrm{D}))=f^*(\operatorname{cyc}(\mathrm{D})).
  \tag{21.10.13.1}
\]
\end{proposition}

\begin{proof}
Reasoning as in \sref{IV.21.10.4}, it all amounts to seeing (with the same notation) that if $x'$ is in case $4^{\mathrm{o}}$ of \sref{IV.21.10.10} and if $n_y$ is the length of $\sh{O}_{\mathrm{X},y}/t_y\sh{O}_{\mathrm{X},y}$, then the
\oldpage[IV-4]{296}
length $n_{x'}$ of $\sh{O}_{\mathrm{X}',x'}/t'_{x'}\sh{O}_{\mathrm{X}',x'}$ is equal to $\mu_{x'}n_y$, $\mu_{x'}$ being the rational number defined in \sref{IV.21.10.10}.
Since $\sh{O}_{\mathrm{X}',x'}/t'_{x'}\sh{O}_{\mathrm{X}',x'}$ is of finite length, it has the same length as its $\mathfrak{m}_{x'}$-preadic completion $\mathrm{A}'/t'_{x'}\mathrm{A}'$, which can also be written $\mathrm{A}'/t_y\mathrm{A}'$; moreover, since $t'_{x'}$ is regular by hypothesis in $\sh{O}_{\mathrm{X}',x'}$, it is also regular in $\mathrm{A}'$ by flatness $(\mathbf{0}_{\mathrm{I}}$, $6.3.4)$, and when $\mathrm{A}'$ is considered as an A-module, we can also say that $t_y$ is $\mathrm{A}'$-regular.
Since $\mathrm{A}$ has dimension $1$ and $\mathrm{A}'$ is a finite $\mathrm{A}$-module such that $\mathrm{A}'\otimes_{\mathrm{A}}\mathrm{K}$ is a free $\mathrm{K}$-module of rank $r_{x'}$, we may apply \sref{IV.21.10.13},{}~(v) to $\mathrm{M}=\mathrm{A}'$ and $t_y$; the length of $\mathrm{A}'/t_y\mathrm{A}'$ as an $\mathrm{A}$-module is therefore $r_{x'}n_y$.
Since $\boldsymbol{k}(x')$ is a $\boldsymbol{k}(x)$-vector space of rank $q_{x'}$, the length of $\mathrm{A}'/t_y\mathrm{A}'$ as an $\mathrm{A}'$-\emph{module} is therefore $r_{x'}n_y/q_{x'}=\mu_{x'}n_y$, which finishes the proof.
\end{proof}

\begin{env}[21.10.14]
\label{IV.21.10.14}
Let now $\mathrm{X}$, $\mathrm{X}'$ be two locally Noetherian preschemes, $f:\mathrm{X}'\to\mathrm{X}$ a morphism having the \emph{two} following properties:
\begin{enumerate}[label=\emph{\alph*)}]
  \item $f$ is finite;
  \item the image by $f$ of every maximal point of $\mathrm{X}'$ is a maximal point of $\mathrm{X}$.
\end{enumerate}

For all $x\in\mathrm{X}^{(1)}$, the points $x'\in f^{-1}(x)$ all belong to $\mathrm{X}'^{(1)}$, as a consequence of hypothesis $b)$ and of the inequality $(\mathbf{0}$, $16.3.9.1)$, the fibre $f^{-1}(x)$ being discrete.
Let then
\[
  Z'=\sum_{x'\in\mathrm{X}'^{(1)}} n_{x'}.\overline{\{x'\}}
\]
be a $1$-codimensional cycle on $\mathrm{X}'$.
For all $x\in\mathrm{X}^{(1)}$, we define an integer $n_x$ by the formula
\[
  n_x=\sum_{x'\in f^{-1}(x)} n_{x'}\cdot[\boldsymbol{k}(x'):\boldsymbol{k}(x)]
\]
which has a meaning, the points of $f^{-1}(x)$ being in finite number and $\boldsymbol{k}(x')$ being a field of finite degree over $\boldsymbol{k}(x)$ $(\mathbf{I}$, $6.4.4)$.
Furthermore, the set of $x\in\mathrm{X}^{(1)}$ such that $n_x\neq 0$ is \emph{locally finite} in $\mathrm{X}$, since it is contained in the image by $f$ of the set of $x'\in\mathrm{X}'^{(1)}$ such that $n_{x'}\neq 0$, and the conclusion follows from the fact that the morphism $f$ is quasi-compact.
We can therefore define a $1$-codimensional cycle on $\mathrm{X}$ by setting
\[
\label{IV.21.10.14.1}
  f_*(Z')=\sum_{x\in\mathrm{X}^{(1)}} n_x.\overline{\{x\}}
  \tag{21.10.14.1}
\]
and we say that $f_*(Z')$ is the \emph{direct image} of $Z'$ by $f$.
It is clear that the map $f_*:\mathfrak{Z}^1(\mathrm{X}')\to\mathfrak{Z}^1(\mathrm{X})$ thus defined is a homomorphism of ordered groups.
Furthermore, if $\mathrm{U}$ is an open of $\mathrm{X}$, $f_{\mathrm{U}}:f^{-1}(\mathrm{U})\to\mathrm{U}$ the restriction of $f$, it follows immediately from the definitions that we have
\[
\label{IV.21.10.14.2}
  (f_{\mathrm{U}})_*(Z'|f^{-1}(\mathrm{U}))=f_*(Z')|\mathrm{U}
  \tag{21.10.14.2}
\]
hence, denoting by $\psi$ the continuous map underlying the morphism $f$, the maps $\Gamma(f^{-1}(\mathrm{U}),\mathscr{Z}^1_{\mathrm{X}'})\to\Gamma(\mathrm{U},\mathscr{Z}^1_{\mathrm{X}})$ that we have just defined define a homomorphism of sheaves of ordered commutative groups on $\mathrm{X}$
\[
  \psi_*(\mathscr{Z}^1_{\mathrm{X}'})\to\mathscr{Z}^1_{\mathrm{X}}.
\]
\end{env}

\begin{proposition}[21.10.15]
\label{IV.21.10.15}
\oldpage[IV-4]{297}
Let $\mathrm{X}$, $\mathrm{X}'$, $\mathrm{X}''$ be three locally Noetherian preschemes, $f:\mathrm{X}'\to\mathrm{X}$, $g:\mathrm{X}''\to\mathrm{X}'$ two morphisms satisfying conditions \emph{a)} and \emph{b)} of \sref{IV.21.10.14}.
Then $f\circ g$ satisfies the same conditions, and for every $1$-codimensional cycle $\mathrm{Z}''$ on $\mathrm{X}''$, we have $(f\circ g)_*(Z'')=f_*(g_*(Z''))$.
\end{proposition}

\begin{proof}
This follows immediately from the definitions.
\end{proof}

\begin{proposition}[21.10.16]
\label{IV.21.10.16}
Let $\mathrm{X}$, $\mathrm{X}'$, $\mathrm{X}_1$ be three locally Noetherian preschemes, $f:\mathrm{X}'\to\mathrm{X}$ a morphism satisfying conditions \emph{a)} and \emph{b)} of \sref{IV.21.10.14}, $g:\mathrm{X}_1\to\mathrm{X}$ a flat morphism.
Set $\mathrm{X}_1'=\mathrm{X}'\times_{\mathrm{X}}\mathrm{X}_1$ (so that $\mathrm{X}_1'$ is locally Noetherian) and let $f_1:\mathrm{X}_1'\to\mathrm{X}_1$ and $g_1:\mathrm{X}_1'\to\mathrm{X}'$ be the canonical projections.
Then $f_1$ satisfies conditions \emph{a)} and \emph{b)} of \sref{IV.21.10.14}, and for every $1$-codimensional cycle $\mathrm{Z}'$ on $\mathrm{X}'$, we have
\[
\label{IV.21.10.16.1}
  g^*(f_*(Z'))=(f_1)_*(g_1^*(Z')).
  \tag{21.10.16.1}
\]
\end{proposition}

\begin{proof}
It is clear that $f_1$ is finite, and it satisfies condition $b)$ of \sref{IV.21.10.14} by virtue of \sref{IV.2.3.7}.
To prove \sref{IV.21.10.16.1}, we are immediately reduced to the case where $\mathrm{X}$, $\mathrm{X}'$ and $\mathrm{X}_1$ are spectra of local Noetherian rings of dimension $1$, $\mathrm{A}$, $\mathrm{A}'$ and $\mathrm{A}_1$, with $\mathrm{A}'$ a finite A-module and $\mathrm{A}_1$ a flat A-module.
Denoting by $x$, $x'$ and $x_1$ the closed points of $\mathrm{X}$, $\mathrm{X}'$ and $\mathrm{X}_1$ respectively, it is a matter of seeing that we have
\[
\label{IV.21.10.16.2}
  \sum_{x_1'}\lambda_{x_1'}[\boldsymbol{k}(x_1'):\boldsymbol{k}(x_1)]=\lambda_{x_1}[\boldsymbol{k}(x'):\boldsymbol{k}(x)]
  \tag{21.10.16.2}
\]
where $x_1'$ ranges over the closed points of $\mathrm{X}_1'$ (i.e.\ the points lying above both $x'$ and $x_1$) and where $\lambda_{x_1'}=\operatorname{length}(\sh{O}_{\mathrm{X}_1',x_1'}/\mathfrak{m}_{x'}\sh{O}_{\mathrm{X}_1',x_1'})$ and $\lambda_{x_1}=\operatorname{length}(\mathrm{A}_1/\mathfrak{m}_x\mathrm{A}_1)$.
The left-hand side of \sref{IV.21.10.16.2} may also be written
\[
  \operatorname{length}_{\mathrm{A}_1}(\mathrm{A}_1'/\mathfrak{m}_{x'}\mathrm{A}_1'),
\]
where $\mathrm{A}_1'=\mathrm{A}'\otimes_{\mathrm{A}}\mathrm{A}_1$.
Now $\mathrm{A}_1'/\mathfrak{m}_{x'}\mathrm{A}_1'=(\mathrm{A}'/\mathfrak{m}_{x'}\mathrm{A}')\otimes_{\mathrm{A}}\mathrm{A}_1$, and since $\mathrm{A}_1$ is a flat A-module, \sref{IV.4.7.1} gives
\[
  \operatorname{length}_{\mathrm{A}_1}(\mathrm{A}_1'/\mathfrak{m}_{x'}\mathrm{A}_1')
  =\operatorname{length}_{\mathrm{A}}(\mathrm{A}'/\mathfrak{m}_{x'}\mathrm{A}')
   \operatorname{length}_{\mathrm{A}_1}(\mathrm{A}_1/\mathfrak{m}_x\mathrm{A}_1)
  =[\boldsymbol{k}(x'):\boldsymbol{k}(x)]\lambda_{x_1},
\]
which completes the proof.
\end{proof}

\begin{proposition}[21.10.17]
\label{IV.21.10.17}
Let $\mathrm{X}$, $\mathrm{X}'$ be two locally Noetherian preschemes, $f:\mathrm{X}'\to\mathrm{X}$ a finite locally free morphism.
Then $f$ satisfies condition \emph{b)} of \sref{IV.21.10.14}, and for every divisor $\mathrm{D}'$ on $\mathrm{X}'$, we have
\[
\label{IV.21.10.17.1}
  f_*(\operatorname{cyc}(\mathrm{D}'))=\operatorname{cyc}(f_*(\mathrm{D}')).
  \tag{21.10.17.1}
\]
\end{proposition}

\begin{proof}
Since $f$ is flat and finite, condition $b)$ of \sref{IV.21.10.14} and the relation $f(\mathrm{X}'^{(1)})\subset\mathrm{X}^{(1)}$ follow from \sref{IV.6.1.1}.
The definition \sref{IV.21.10.14.1} shows that we can reduce to the case where $\mathrm{X}=\operatorname{Spec}(\mathrm{A})=\operatorname{Spec}(\sh{O}_{\mathrm{X},x})$ for an $x\in\mathrm{X}^{(1)}$; then $\mathrm{X}'=\operatorname{Spec}(\mathrm{A}')$, where $\mathrm{A}'$ is an $\mathrm{A}$-algebra which is a
\oldpage[IV-4]{298}
free $\mathrm{A}$-module of finite rank; furthermore, we can suppose that $\mathrm{D}'=\operatorname{div}(t')$, where $t'$ is a regular element of $\mathrm{A}'$; we then have $f_*(\mathrm{D}')=\operatorname{div}(t)$, where $t=\mathrm{N}_{\mathrm{A}'/\mathrm{A}}(t')$ is a regular element of $\mathrm{A}$ \sref{IV.21.5.2}.
We can restrict to the case where $t'$ is not invertible in $\mathrm{A}'$, which is equivalent to $t$ not being invertible in $\mathrm{A}$; the ring $\mathrm{A}/t\mathrm{A}$ is then nonzero and of finite length, and $\mathrm{A}'/t\mathrm{A}'$ is a finite direct product of Artinian local rings $\mathrm{A}_i'$ ($1\leq i\leq r$), whose residue fields are $\boldsymbol{k}(x_i')$, where the $x_i'$ are the points of $\mathrm{X}'$ lying above $x$.
If $t_i'$ is the image of $t'$ in $\mathrm{A}_i'$, $\mathrm{A}'/t'\mathrm{A}'$ is the direct product of the $\mathrm{A}_i'/t_i'\mathrm{A}_i'$; since the product $\operatorname{length}_{\mathrm{A}_i'}(\mathrm{A}_i'/t_i'\mathrm{A}_i')\cdot[\boldsymbol{k}(x_i'):\boldsymbol{k}(x)]$ is equal to $\operatorname{length}_{\mathrm{A}}(\mathrm{A}_i'/t_i'\mathrm{A}_i')$, we see that the multiplicity at point $x$ of the left-hand side of \sref{IV.21.10.17.1} is $\operatorname{length}_{\mathrm{A}}(\mathrm{A}'/t'\mathrm{A}')$, so the formula to prove reduces to
\[
\label{IV.21.10.17.2}
  \operatorname{length}_{\mathrm{A}}(\mathrm{A}'/t'\mathrm{A}')=\operatorname{length}_{\mathrm{A}}(\mathrm{A}/t\mathrm{A}).
  \tag{21.10.17.2}
\]

This relation follows from the following more general lemma:
\end{proof}

\begin{lemma}[21.10.17.3]
\label{IV.21.10.17.3}
Let $\mathrm{A}$ be a local Noetherian ring of dimension $1$, $\mathrm{M}$ a free A-module of finite rank, $u$ an injective endomorphism of $\mathrm{M}$.
Then we have
\[
\label{IV.21.10.17.4}
  \operatorname{length}_{\mathrm{A}}(\operatorname{Coker} u)=\operatorname{length}_{\mathrm{A}}(\mathrm{A}/(\det u)\mathrm{A}).
  \tag{21.10.17.4}
\]
\end{lemma}

\begin{proof}
We distinguish several cases.

I) \emph{$\mathrm{A}$ is a discrete valuation ring}.
Let $\pi$ be a uniformiser of $\mathrm{A}$, and note $\operatorname{length}_{\mathrm{A}}(\mathrm{A}/\pi^k\mathrm{A})=k$; if $n$ is the rank of $\mathrm{M}$ and if $\pi^{m_i}$ ($1\leq i\leq n$) are the invariant factors of $u$, $\operatorname{Coker}(u)$ is a direct sum of A-modules $\mathrm{A}/\pi^{m_i}\mathrm{A}$, so has length $m=\sum_{i=1}^n m_i$, and $\det(u)$ is the product of an invertible element and of $\pi^m$, whence the conclusion in this case (Bourbaki, \emph{Alg.}, chap.~VII, \textsection4, n$^{\mathrm{o}}$~5, cor.~1 of prop.~4).

II) \emph{$\mathrm{A}$ is a complete integral ring (of dimension $1$)}.
We know then $(\mathbf{0}$, $19.8.8$,~(ii)) that there is a subring $\mathrm{B}$ of $\mathrm{A}$, which is a discrete valuation ring, such that $\mathrm{B}\to\mathrm{A}$ is a local homomorphism making $\mathrm{A}$ a $\mathrm{B}$-module of finite type; since this $\mathrm{B}$-module is evidently torsion-free, it is \emph{free} (Bourbaki, \emph{Alg.\ comm.}, chap.~VI, \textsection3, n$^{\mathrm{o}}$~6, lemma~1).
Let $\mathrm{M}'$ denote the set $\mathrm{M}$ equipped with its structure of $\mathrm{B}$-module (free), by $u'$ the endomorphism $u$ considered as a $\mathrm{B}$-endomorphism.
It follows from I) that
\[
\label{IV.21.10.17.5}
  \operatorname{length}_{\mathrm{B}}(\operatorname{Coker} u')=\operatorname{length}_{\mathrm{B}}(\mathrm{B}/(\det u')\mathrm{B}).
  \tag{21.10.17.5}
\]
But $\det u'=\mathrm{N}_{\mathrm{A}/\mathrm{B}}(\det u)$ (Bourbaki, \emph{Alg.}, chap.~VIII, \textsection~12, n\textsuperscript{o}~2, prop.~7), therefore, by applying \sref{IV.21.10.17.4} to the homothety $x\mapsto(\det u)x$ of the free $\mathrm{B}$-module $\mathrm{A}$, we get
\[
  \operatorname{length}_{\mathrm{B}}(\mathrm{B}/(\det u')\mathrm{B})=\operatorname{length}_{\mathrm{B}}(\mathrm{A}/(\det u)\mathrm{A}),
\]
whence, substituting in \sref{IV.21.10.17.5} and dividing by $[\boldsymbol{k}(\mathrm{A}):\boldsymbol{k}(\mathrm{B})]$, the length of the $\mathrm{B}$-module $\boldsymbol{k}(\mathrm{A})$, we obtain \sref{IV.21.10.17.4} in the case considered.

\oldpage[IV-4]{299}
III) \emph{$\mathrm{A}$ is a complete ring}.
Let us note that we can further suppose that $\mathfrak{m}\notin\operatorname{Ass}(\mathrm{A})$; indeed, since $\det(u)$ is a regular element of $\mathrm{A}$, if we had $\mathfrak{m}\in\operatorname{Ass}(\mathrm{A})$, then $\det(u)$ would be invertible, $u$ an automorphism of $\mathrm{M}$, and the formula \sref{IV.21.10.17.4} would then be trivial, both members being zero.

In what follows, for an endomorphism $v$ of a module $\mathrm{N}$ over a ring $\mathrm{R}$, such that $\operatorname{Ker} v$ and $\operatorname{Coker} v$ are of \emph{finite length}, we set
\[
  \chi(\mathrm{N},v)=\operatorname{length}_{\mathrm{R}}(\operatorname{Ker}(v))-\operatorname{length}_{\mathrm{R}}(\operatorname{Coker}(v)).
\]
We note that the hypothesis on $v$ amounts to saying that the complex
\[
  \mathrm{K}^0:0\to\mathrm{N}\xrightarrow{v}\mathrm{N}\to 0
\]
has its cohomology modules of finite length and that $\chi(\mathrm{N},v)=\chi(\mathrm{H}^0(\mathrm{K}^0))$ $(\mathbf{0}_{\mathrm{III}}$, $11.10)$.
We deduce that if $\mathrm{N}'$, $\mathrm{N}$, $\mathrm{N}''$ are three $\mathrm{R}$-modules, and if there is a commutative diagram
\[
  \xymatrix{
    0 \ar[r] & \mathrm{N}' \ar[r]^r \ar[d]^{v'} & \mathrm{N} \ar[r]^s \ar[d]^v & \mathrm{N}'' \ar[r] \ar[d]^{v''} & 0\\
    0 \ar[r] & \mathrm{N}' \ar[r]_r & \mathrm{N} \ar[r]_s & \mathrm{N}'' \ar[r] & 0
  }
\]
whose rows are exact, and if two of the three numbers $\chi(\mathrm{N},v)$, $\chi(\mathrm{N}',v')$ and $\chi(\mathrm{N}'',v'')$ are defined, then so is the third and one has
\[
\label{IV.21.10.17.6}
  \chi(\mathrm{N},v)=\chi(\mathrm{N}',v')+\chi(\mathrm{N}'',v'').
  \tag{21.10.17.6}
\]
This follows indeed immediately from the cohomology exact sequence.

Finally, if $\mathrm{N}$ is an $\mathrm{R}$-module of \emph{finite length}, we have $\chi(\mathrm{N},v)=0$.

With this notation, we have the following lemma:
\end{proof}

\begin{lemma}[21.10.17.7]
\label{IV.21.10.17.7}
Let $\mathrm{A}$ be a local Noetherian ring of dimension $1$ whose maximal ideal $\mathfrak{m}$ is such that $\mathfrak{m}\notin\operatorname{Ass}(\mathrm{A})$; let $\mathfrak{p}_i$ ($1\leq i\leq n$) be the minimal prime ideals of $\mathrm{A}$.
Let $\mathrm{M}$ be a free A-module of finite type, $u$ an endomorphism of $\mathrm{M}$ such that $\chi(\mathrm{M},u)$ is defined; for each $i$, set $\mathrm{M}_i=\mathrm{M}\otimes_{\mathrm{A}}(\mathrm{A}/\mathfrak{p}_i)$ and let $u_i$ be the endomorphism $u\otimes 1_{\mathrm{A}/\mathfrak{p}_i}$ of $\mathrm{M}_i$; then if $\chi(\mathrm{M}_i,u_i)$ is defined for each $i$, we have
\[
\label{IV.21.10.17.8}
  \chi(\mathrm{M},u)=\sum_{i=1}^n\operatorname{length}(\mathrm{A}_{\mathfrak{p}_i})\cdot\chi(\mathrm{M}_i,u_i).
  \tag{21.10.17.8}
\]
\end{lemma}

\begin{proof}
Since $\mathfrak{m}\notin\operatorname{Ass}(\mathrm{A})$, we have a unique reduced primary decomposition $(0)=\bigcap_{1\leq i\leq n}\mathfrak{q}_i$, where the ideal $\mathfrak{q}_i$ is $\mathfrak{p}_i$-primary for $1\leq i\leq n$.
If we set $\mathrm{M}_i'=\mathrm{M}\otimes_{\mathrm{A}}(\mathrm{A}/\mathfrak{q}_i)$, we have therefore an exact sequence
\[
  0\to\mathrm{M}\to\bigoplus_i\mathrm{M}_i'\to\mathrm{M}''\to 0
\]
of A-modules, where $\mathrm{M}''$ is of \emph{finite length}: indeed by localising the preceding exact sequence at each of the $\mathfrak{p}_i$, we obtain $\mathrm{M}''_{\mathfrak{p}_i}=0$, since $(\mathfrak{q}_i)_{\mathfrak{p}_i}=0$ and $(\mathfrak{q}_j)_{\mathfrak{p}_i}=\mathrm{A}_{\mathfrak{p}_i}$ for $j\neq i$ (Bourbaki, \emph{Alg.\ comm.}, chap.~IV, \textsection2, n$^{\mathrm{o}}$~4, prop.~6), so $(\mathrm{M}_i')_{\mathfrak{p}_i}=\mathrm{M}_{\mathfrak{p}_i}$ and $(\mathrm{M}_j')_{\mathfrak{p}_i}=0$;
\oldpage[IV-4]{300}
the support of $\mathrm{M}''$ is therefore reduced to $\mathfrak{m}$, hence $\mathrm{M}''$ is of finite length $(\mathbf{0}$, $16.1.10)$.
If we set $u_i'=u\otimes 1_{\mathrm{A}/\mathfrak{q}_i}$, and if $u''$ is the endomorphism of $\mathrm{M}''$ deduced from $\bigoplus u_i'$ by passing to quotients, we deduce from \sref{IV.21.10.17.6} that $\chi(\mathrm{M},u)+\chi(\mathrm{M}'',u'')=\sum_i\chi(\mathrm{M}_i',u_i')$, and since $\mathrm{M}''$ is of finite length, $\chi(\mathrm{M}'',u'')=0$.
To prove \sref{IV.21.10.17.8}, we can therefore reduce to the case where $n=1$.
We will then denote by $\mathfrak{p}$ the unique minimal prime ideal of $\mathrm{A}$, which is the nilradical of $\mathrm{A}$; if $\mathrm{A}_0=\mathrm{A}/\mathfrak{p}$, $\mathrm{M}_0=\mathrm{M}\otimes_{\mathrm{A}}\mathrm{A}_0$, $u_0=u\otimes 1_{\mathrm{A}_0}$, it is a matter of proving that if $\chi(\mathrm{M}_0,u_0)$ is defined, then
\[
\label{IV.21.10.17.9}
  \chi(\mathrm{M},u)=\operatorname{length}\mathrm{A}_{\mathfrak{p}}\cdot\chi(\mathrm{M}_0,u_0).
  \tag{21.10.17.9}
\]

Let $\mathfrak{n}_j$ ($0\leq j\leq r$) be the ``$j$-th symbolic power'' of $\mathfrak{p}$, the inverse image of $(\mathfrak{p}\mathrm{A}_{\mathfrak{p}})^j$ in $\mathrm{A}$ ($1\leq j\leq r$), with $\mathfrak{n}_0=\mathrm{A}$ and $\mathfrak{n}_r=0$; set
\[
  \mathrm{M}_j=\mathfrak{n}_j\mathrm{M}/\mathfrak{n}_{j+1}\mathrm{M}\qquad(0\leq j\leq r-1),
\]
and denote by $v_j$ the endomorphism of $\mathrm{M}_j$ deduced from $u$ by restriction to $\mathfrak{n}_j\mathrm{M}$ and passage to quotients.
We shall first show that each of the numbers $\chi(\mathrm{M}_j,v_j)$ is defined and that
\[
\label{IV.21.10.17.10}
  \chi(\mathrm{M},u)=\sum_j\chi(\mathrm{M}_j,v_j).
  \tag{21.10.17.10}
\]
The first assertion implies the second, by applying \sref{IV.21.10.17.6} to each of the exact sequences
\[
  0\to\mathfrak{n}_j\mathrm{M}/\mathfrak{n}_{j+1}\mathrm{M}\to\mathrm{M}/\mathfrak{n}_{j+1}\mathrm{M}\to\mathrm{M}/\mathfrak{n}_j\mathrm{M}\to 0.
\]

To prove the first assertion, we note that if $m$ is the rank of the free A-module $\mathrm{M}$, $\mathrm{M}_j$ is isomorphic to $(\mathfrak{n}_j/\mathfrak{n}_{j+1})^m$, or since $\mathfrak{p}$ annihilates each of the quotients (since $\mathfrak{p}\cdot\mathfrak{n}_j\subset\mathfrak{n}_{j+1}$), $\mathrm{M}_j$ is an $\mathrm{A}_0$-module isomorphic to $\mathrm{M}_0\otimes_{\mathrm{A}_0}(\mathfrak{n}_j/\mathfrak{n}_{j+1})$.
Let us denote by $l_j$ the rank of the $\mathrm{A}_0$-module $\mathfrak{n}_j/\mathfrak{n}_{j+1}$; since the field of fractions $\mathrm{K}_0$ of $\mathrm{A}_0$ is the residue field of $\mathrm{A}_{\mathfrak{p}}$, $l_j$ is also the length of the $\mathrm{A}_{\mathfrak{p}}$-module $(\mathfrak{p}\mathrm{A}_{\mathfrak{p}})^j/(\mathfrak{p}\mathrm{A}_{\mathfrak{p}})^{j+1}$.
There is a system of generators of the $\mathrm{A}_0$-module $\mathrm{M}_j$ which contains a basis of $\mathrm{M}_j\otimes_{\mathrm{A}_0}\mathrm{K}_0$; hence there is an $\mathrm{A}_0$-homomorphism
\[
  w_j:\mathrm{M}_0^{l_j}\to\mathrm{M}_j
\]
whose localisation at the ideal $(0)$ of $\mathrm{A}_0$ is an isomorphism, so that the support of $\operatorname{Ker}(w_j)$ and of $\operatorname{Coker}(w_j)$ is reduced to the maximal ideal $\mathfrak{m}/\mathfrak{p}$ of $\mathrm{A}_0$; $\operatorname{Ker}(w_j)$ and $\operatorname{Coker}(w_j)$ are therefore $\mathrm{A}_0$-modules of finite length $(\mathbf{0}$, $16.1.10)$.
Since by hypothesis $\chi(\mathrm{M}_0,u_0)$ is defined, it follows from the same of $\chi(\mathrm{M}_0^{l_j},u_0^{l_j})=l_j\chi(\mathrm{M}_0,u_0)$ and by virtue of \sref{IV.21.10.17.6} and the fact that $\operatorname{Ker}(w_j)$ and $\operatorname{Coker}(w_j)$ are of finite length, we see that $\chi(\mathrm{M}_j,v_j)$ is defined and equal to $l_j\chi(\mathrm{M}_0,u_0)$; the relation \sref{IV.21.10.17.10} then gives
\[
  \chi(\mathrm{M},u)=(\sum_j l_j)\chi(\mathrm{M}_0,u_0),
\]
and by virtue of a preceding remark, $\sum_j l_j$ is none other than the length of $\mathrm{A}_{\mathfrak{p}}$, which finishes the proof of the lemma \sref{IV.21.10.17.7}.
\end{proof}

\oldpage[IV-4]{301}
To apply this lemma when $\mathrm{A}$ is a complete ring and $\mathfrak{m}\notin\operatorname{Ass}(\mathrm{A})$, we observe that if $u$ is injective, so is each of the $u_i$ (with the notation of the lemma): indeed, $\det u$ is then a regular element of $\mathrm{A}$, so is not contained in any of the $\mathfrak{p}_i$, and its image $\det u_i$ in $\mathrm{A}/\mathfrak{p}_i$ is therefore a nonzero element of this integral ring, which proves that $u_i$ is injective.
Since $\operatorname{Coker}(u_i)$ is an image of $\operatorname{Coker}(u)$, it too has finite length, so $\chi(\mathrm{M}_i,u_i)$ is defined for every $i$; formula \sref{IV.21.10.17.8} therefore applies.
On the other hand, since $\det(u)$ is a regular element of $\mathrm{A}$, it is not contained in any of the $\mathfrak{p}_i$; the ideal $(\det u)\mathrm{A}$ is therefore $\mathfrak{m}$-primary and the quotient $\mathrm{A}/(\det u)\mathrm{A}$ is of finite length.
Applying the same reasoning above to the injective homothety $t:\xi\mapsto(\det u)\xi$ of $\mathrm{A}$ and to its images $t_i=\det u_i$ in the $\mathrm{A}/\mathfrak{p}_i$, we get
\[
  \chi(\mathrm{A},t)=\sum_i\operatorname{length}(\mathrm{A}_{\mathfrak{p}_i})\cdot\chi(\mathrm{A}/\mathfrak{p}_i,t_i).
\]
But all the rings $\mathrm{A}/\mathfrak{p}_i$ are integral and complete, and by applying the result of II) to each of them, one obtains again the relation \sref{IV.21.10.17.4} for $\mathrm{M}$ and $u$.

IV) \emph{General case}.
Set $\mathrm{A}'=\widehat{\mathrm{A}}$, $\mathrm{M}'=\mathrm{M}\otimes_{\mathrm{A}}\mathrm{A}'$, $u'=u\otimes 1_{\mathrm{A}'}$; we have $\det(u')=\det(u)$, and by flatness, $\operatorname{Coker}(u')=(\operatorname{Coker}(u))_{(\mathrm{A}')}$ and $\mathrm{A}'/(\det u)\mathrm{A}'=(\mathrm{A}/(\det u)\mathrm{A})_{(\mathrm{A}')}$; since the formula \sref{IV.21.10.17.4} is true for $\mathrm{A}'$ and $u'$ by virtue of III), it is also true for $\mathrm{A}$ and $u$.

This finishes the proof of \sref{IV.21.10.17}.

\begin{proposition}[21.10.18]
\label{IV.21.10.18}
Let $\mathrm{X}$, $\mathrm{X}'$ be two locally Noetherian preschemes, $f:\mathrm{X}'\to\mathrm{X}$ a finite morphism, transforming every maximal point of $\mathrm{X}'$ into a maximal point of $\mathrm{X}$, and satisfying for all $x'\in\mathrm{X}'^{(1)}$ one of the conditions (ii), (iii), (iv) of \sref{IV.21.10.10}.
Suppose furthermore that there exists an integer $n$ such that, for every maximal point $x$ of $\mathrm{X}$, $(f_*(\sh{O}_{\mathrm{X}'}))_x$ is a free $\sh{O}_{\mathrm{X},x}$-module of rank $n$.
Then, for every $1$-codimensional cycle $\mathrm{Z}$ on $\mathrm{X}$, we have
\[
\label{IV.21.10.18.1}
  f_*(f^*(\mathrm{Z}))=n\cdot\mathrm{Z}
  \tag{21.10.18.1}
\]
(``projection formula'').
\end{proposition}

\begin{proof}
By virtue of the definitions, we are immediately reduced to the case where $\mathrm{X}$ is the spectrum of a local Noetherian ring $\mathrm{A}$ of dimension $1$, with closed point $x$, and $\mathrm{Z}=\overline{\{x\}}$; write $\mathrm{X}'=\operatorname{Spec}(\mathrm{A}')$, where $\mathrm{A}'$ is a finite A-algebra and, for every minimal prime ideal $\mathfrak{p}_i$ of $\mathrm{A}$, $\mathrm{A}'_{\mathfrak{p}_i}$ is a free $\mathrm{A}_{\mathfrak{p}_i}$-module of rank $n$.
Let us show that we may moreover reduce to the case where $\mathrm{A}$ is \emph{complete}.  Make the base change $h:\mathrm{Y}\to\mathrm{X}$, where $\mathrm{Y}=\operatorname{Spec}(\mathrm{B})$ and $\mathrm{B}=\widehat{\mathrm{A}}$, and set
\[
  \mathrm{Y}'=\mathrm{X}'\times_{\mathrm{X}}\mathrm{Y}
  =\operatorname{Spec}(\mathrm{B}\otimes_{\mathrm{A}}\mathrm{A}'),
  \qquad g=f_{\{\mathrm{Y}\}}:\mathrm{Y}'\to\mathrm{Y}.
\]
The morphism $g$ is finite.  Since $h$ is flat, the maximal points of $\mathrm{Y}$ lie above those of $\mathrm{X}$; above each $\mathfrak{p}_i$ there are only finitely many minimal prime ideals $\mathfrak{p}_{ij}$ of $\mathrm{B}$, and $(\mathrm{B}\otimes_{\mathrm{A}}\mathrm{A}')_{\mathfrak{p}_{ij}}$ is a free $\mathrm{B}_{\mathfrak{p}_{ij}}$-module of rank $n$.
Finally, if $f$ satisfies one of conditions (ii), (iii), (iv) of \sref{IV.21.10.10} at every point $x'$ of $f^{-1}(x)$, then $g$ satisfies the corresponding condition at the unique point $y'$ of $g^{-1}(y)$ lying above $x'$, where $y$ is the closed point of $\mathrm{Y}$.  This is immediate for (ii) and (iv); for (iii), it follows because $\mathrm{A}$ is then a discrete valuation ring, hence so is $\mathrm{B}$, while
\oldpage[IV-4]{302}
$\mathfrak{m}_{y'}\notin\operatorname{Ass}(\sh{O}_{\mathrm{Y}',y'})$ follows by flatness from $\mathfrak{m}_{x'}\notin\operatorname{Ass}(\sh{O}_{\mathrm{X}',x'})$ $(\mathbf{0}$, $3.3.1)$.
Thus $g$ satisfies the same hypotheses as $f$; if \sref{IV.21.10.18.1} holds for $g$ and $\overline{\{y\}}$, it holds for $f$ and $\overline{\{x\}}$ by \sref{IV.21.10.16.1}.

We may therefore suppose that $\mathrm{A}$ is complete.  Then $\mathrm{A}'$ is complete as well and is a finite direct product of complete local rings, so we may further reduce to the case where $\mathrm{A}'$ is local, and verify \sref{IV.21.10.18.1} separately in cases (ii), (iii), and (iv), for the closed point $x'$ of $\mathrm{X}'$.

In case (ii), $\mathrm{A}'$ is a finite flat A-module, hence a free A-module $(\mathbf{0}_{\mathrm{III}}$, $10.1.3)$, of rank $n$ by hypothesis.  By definitions \sref{IV.21.10.1} and \sref{IV.21.10.3},
\[
  f^*(\mathrm{Z})=\lambda_{x'}\cdot\overline{\{x'\}},
\]
where $\lambda_{x'}$ is the length of the $\mathrm{A}'$-module $\mathrm{A}'/\mathfrak{m}_x\mathrm{A}'$.  Hence
\[
  f_*(f^*(\mathrm{Z}))
  =\lambda_{x'}[\boldsymbol{k}(x'):\boldsymbol{k}(x)]\cdot\overline{\{x\}}.
\]
But the coefficient is the length of $\mathrm{A}'/\mathfrak{m}_x\mathrm{A}'$ as an A-module, equivalently its dimension as a $\boldsymbol{k}(x)$-vector space, and is therefore equal to $n$.

In case (iii), $\mathrm{A}$ is a discrete valuation ring, hence regular, and the hypothesis $\mathfrak{m}_{x'}\notin\operatorname{Ass}(\sh{O}_{\mathrm{X}',x'})$ implies that $\mathrm{A}'$ is Cohen--Macaulay $(\mathbf{0}$, $16.4.6)$.  Since $\dim(\mathrm{A}')=\dim(\mathrm{A})=1$ and $\mathrm{A}'/\mathfrak{m}_x\mathrm{A}'$ is Artinian, \sref{IV.6.1.5} shows that $\mathrm{A}'$ is flat over $\mathrm{A}$, reducing this to case (ii).

In case (iv), if $\mathrm{K}$ is the total ring of fractions of $\mathrm{A}$, then $\mathrm{K}'=\mathrm{A}'\otimes_{\mathrm{A}}\mathrm{K}$ is by hypothesis a free $\mathrm{K}$-module of rank $n$ and, by definition \sref{IV.21.10.10},
\[
  f^*(\mathrm{Z})=\frac{n}{[\boldsymbol{k}(x'):\boldsymbol{k}(x)]}\cdot\overline{\{x'\}},
\]
whence $f_*(f^*(\mathrm{Z}))=n\cdot\overline{\{x\}}$.
\end{proof}

We note that the formula \sref{IV.21.10.18.1} is applicable in particular when the morphism $f$ is \emph{finite} and \emph{flat} and such that for every maximal point $x$ of $\mathrm{X}$, $(f_*(\sh{O}_{\mathrm{X}'}))_x$ is a free $\sh{O}_{\mathrm{X},x}$-module of rank $n$.

\begin{corollary}[21.10.19]
\label{IV.21.10.19}
Under the hypotheses of \sref{IV.21.10.18}, let $\mathrm{D}$ be a divisor on $\mathrm{X}$ such that $f^*(\mathrm{D})$ is defined \sref{IV.21.4.5}; then we have
\[
\label{IV.21.10.19.1}
  f_*(\operatorname{cyc}(f^*(\mathrm{D})))=n\cdot\operatorname{cyc}(\mathrm{D}).
  \tag{21.10.19.1}
\]
\end{corollary}

\begin{proof}
This follows from \sref{IV.21.10.18} and \sref{IV.21.10.13bis}.
\end{proof}
\subsection{Factoriality of regular local rings}
\label{subsection:IV.21.11}
\label{IV.21.11}

\begin{theorem}[21.11.1]
\label{IV.21.11.1}
\textup{(Auslander-Buchsbaum)} ---
A regular local Noetherian ring is factorial.
\end{theorem}

\begin{proof}
The proof that follows is due to I.~Kaplansky.

Let $\mathrm{A}$ be a regular local ring of dimension $n$; we will reason by induction on $n$.
For $n=0$, $\mathrm{A}$ is a field, and for $n=1$, a discrete valuation ring, hence principal (and \emph{a fortiori} factorial).
Suppose therefore $n\geq 2$ and the theorem proved for regular rings of dimension $<n$.
Set $\mathrm{X}=\operatorname{Spec}(\mathrm{A})$, denote by $a$ the closed point of $\mathrm{X}$ and set $\mathrm{U}=\mathrm{X}-\{a\}$.
At every point $y\in\mathrm{U}$, we have $\dim(\sh{O}_{\mathrm{X},y})\leq n-1$, and since $\mathrm{A}$ is regular, the rings $\sh{O}_{\mathrm{X},y}$ are also regular $(\mathbf{0}$, $17.3.2)$, hence by the induction hypothesis they are factorial.
Moreover, since $\operatorname{prof}(\mathrm{A})=\dim(\mathrm{A})\geq 2$ since $\mathrm{A}$ is regular, hence Cohen-Macaulay $(\mathbf{0}$, $17.1.3)$.
\oldpage[IV-4]{303}
Using \sref{IV.21.6.14}, we are reduced to proving that $\operatorname{Pic}(\mathrm{U})=0$.
Consider therefore an invertible $\sh{O}_{\mathrm{U}}$-Module $\mathscr{L}$, and let us prove that it is isomorphic to $\sh{O}_{\mathrm{U}}$.
It follows from $(\mathbf{I}$, $9.4.5)$ that there exists a \emph{coherent} $\sh{O}_{\mathrm{X}}$-Module $\mathscr{F}$ such that $\mathscr{F}|\mathrm{U}=\mathscr{L}$.
Since $\mathrm{A}$ is regular, hence of finite cohomological dimension $(\mathbf{0}$, $17.3.1)$, there exists a \emph{finite} free resolution of $\mathscr{F}$:
\[
  0\leftarrow\mathscr{F}\leftarrow\sh{O}_{\mathrm{X}}^{n_1}\leftarrow\sh{O}_{\mathrm{X}}^{n_2}\leftarrow\ldots\leftarrow\sh{O}_{\mathrm{X}}^{n_h}\leftarrow 0
\]
$(\mathbf{0}$, $17.2.8$ and $\mathbf{0}$, $17.2.2$,~(iii)).
By restriction to $\mathrm{U}$, we have therefore a finite resolution
\[
\label{IV.21.11.1.1}
  0\leftarrow\mathscr{L}\leftarrow\sh{O}_{\mathrm{U}}^{n_1}\leftarrow\sh{O}_{\mathrm{U}}^{n_2}\leftarrow\ldots\leftarrow\sh{O}_{\mathrm{U}}^{n_h}\leftarrow 0.
  \tag{21.11.1.1}
\]
\end{proof}

The theorem will follow from the following general considerations.
On a ringed space $\mathrm{X}$, let $\mathscr{E}$ be a locally free $\sh{O}_{\mathrm{X}}$-module of finite rank; we will denote by $\Lambda^{\max}\mathscr{E}$ the invertible $\sh{O}_{\mathrm{X}}$-module which, in a neighbourhood of each point of $\mathrm{X}$, is equal to $\Lambda^p\mathscr{E}$, denoting by $p$ the rank of $\mathscr{E}$ in this neighbourhood (which can vary with the connected component).
With this notation, we have the

\begin{lemma}[21.11.1.2]
\label{IV.21.11.1.2}
Let $\mathrm{X}$ be a ringed space in local rings, and
\[
  0\leftarrow\mathscr{E}_0\xleftarrow{u_0}\mathscr{E}_1\leftarrow\ldots\leftarrow\mathscr{E}_h\leftarrow 0
\]
an exact sequence of $\sh{O}_{\mathrm{X}}$-Modules locally free of finite rank; then the invertible $\sh{O}_{\mathrm{X}}$-Module
\[
  \bigotimes_{0\leq i\leq h}(\Lambda^{\max}\mathscr{E}_i)^{\otimes(-1)^i}
\]
is isomorphic to $\sh{O}_{\mathrm{X}}$.
\end{lemma}

\begin{proof}
Let us show how this lemma will allow us to conclude the proof of \sref{IV.21.11.1}.
It suffices to note that, for every integer $n$, $\Lambda^{\max}(\sh{O}_{\mathrm{U}}^n)=\Lambda^n\sh{O}_{\mathrm{U}}^n$ is isomorphic to $\sh{O}_{\mathrm{U}}$, and on the other hand $\Lambda^{\max}\mathscr{L}=\mathscr{L}$ for every invertible $\sh{O}_{\mathrm{U}}$-Module $\mathscr{L}$; the lemma \sref{IV.21.11.1.2}, applied to the exact sequence \sref{IV.21.11.1.1}, shows that $\mathscr{L}$ is isomorphic to $\sh{O}_{\mathrm{U}}$.

It remains to prove \sref{IV.21.11.1.2}; we proceed by induction on $h$, the lemma being trivial for $h=1$.
For $h>1$, $\mathscr{N}=\operatorname{Ker}(u_0)$ is an $\sh{O}_{\mathrm{X}}$-Module locally free of finite rank $(\mathbf{0}_{\mathrm{I}}$, $5.5.5)$, and we have the two exact sequences
\[
  0\leftarrow\mathscr{E}_0\leftarrow\mathscr{E}_1\leftarrow\mathscr{N}\leftarrow 0
\]
\[
  0\leftarrow\mathscr{N}\leftarrow\mathscr{E}_2\leftarrow\ldots\leftarrow\mathscr{E}_h\leftarrow 0
\]
By virtue of the induction hypothesis, $(\Lambda^{\max}\mathscr{N})\otimes(\bigotimes_{2\leq i\leq h}(\Lambda^{\max}\mathscr{E}_i)^{\otimes(-1)^{i-1}})$ is isomorphic to $\sh{O}_{\mathrm{X}}$; it will therefore suffice to define a canonical isomorphism $(\Lambda^{\max}\mathscr{N})\otimes(\Lambda^{\max}\mathscr{E}_0)\xrightarrow{\sim}\Lambda^{\max}\mathscr{E}_1$ to finish the proof.
Now, there exists an open covering $(\mathrm{U}_\alpha)$ of $\mathrm{X}$ such that in each $\mathrm{U}_\alpha$, $\mathscr{E}_1|\mathrm{U}_\alpha$ is a direct sum of $\mathscr{N}|\mathrm{U}_\alpha$ and of a free $\sh{O}_{\mathrm{U}_\alpha}$-Module $\mathscr{M}_\alpha$ $(\mathbf{0}_{\mathrm{I}}$, $5.5.5)$, whence a canonical isomorphism $v_\alpha:\mathscr{M}_\alpha\xrightarrow{\sim}\mathscr{E}_0|\mathrm{U}_\alpha$.
Since we have a canonical isomorphism
\[
  r_\alpha:(\Lambda^{\max}\mathscr{N}|\mathrm{U}_\alpha)\otimes(\Lambda^{\max}\mathscr{M}_\alpha)\xrightarrow{\sim}(\Lambda^{\max}\mathscr{E}_1)|\mathrm{U}_\alpha
\]
we deduce, by means of $v_\alpha$, an isomorphism
\[
  u_\alpha:(\Lambda^{\max}\mathscr{N}|\mathrm{U}_\alpha)\otimes(\Lambda^{\max}\mathscr{E}_0|\mathrm{U}_\alpha)\xrightarrow{\sim}\Lambda^{\max}\mathscr{E}_1|\mathrm{U}_\alpha
\]
\oldpage[IV-4]{304}
and it is a matter of showing that $u_\alpha$ and $u_\beta$ coincide in $\mathrm{U}_\alpha\cap\mathrm{U}_\beta$ for any two indices $\alpha$, $\beta$.
Now, if $v_\alpha'$ and $v_\beta'$ are the restrictions to $\mathrm{U}_\alpha\cap\mathrm{U}_\beta$ of $v_\alpha$ and $v_\beta$ respectively, we have $v_\alpha'=v_\beta'\circ w_{\beta\alpha}$, where $w_{\beta\alpha}:\mathscr{M}_\alpha|(\mathrm{U}_\alpha\cap\mathrm{U}_\beta)\xrightarrow{\sim}\mathscr{M}_\beta|(\mathrm{U}_\alpha\cap\mathrm{U}_\beta)$ is the ``parallel projection'' such that for every section $s\in\Gamma(\mathrm{U}_\alpha\cap\mathrm{U}_\beta,\mathscr{M}_\alpha)$, $w_{\beta\alpha}(s)=s+t$ with $t\in\Gamma(\mathrm{U}_\alpha\cap\mathrm{U}_\beta,\mathscr{N})$; the identity of $u_\alpha$ and $u_\beta$ follows immediately from this fact and from the definition of the canonical isomorphism $r_\alpha$ (Bourbaki, \emph{Alg.}, chap.~III, 3$^{\text{e}}$ ed.).
\end{proof}

\subsection{The Van der Waerden purity theorem for the ramification locus of a birational morphism}
\label{subsection:IV.21.12}
\label{IV.21.12}

\begin{env}[21.12.1]
\label{IV.21.12.1}
Let $\mathrm{X}$ and $\mathrm{U}$ be two preschemes, $f:\mathrm{U}\to\mathrm{X}$ a quasi-compact and quasi-separated morphism, so that $f_*(\sh{O}_{\mathrm{U}})$ is a quasi-coherent $\sh{O}_{\mathrm{X}}$-Algebra \sref{IV.1.7.4}.
We call \emph{affine envelope} of the $\mathrm{X}$-prescheme $\mathrm{U}$ the $\mathrm{X}$-prescheme affine over $\mathrm{X}$
\[
  \mathrm{U}^0=\operatorname{Aff}(\mathrm{U}/\mathrm{X})=\operatorname{Spec}(f_*(\sh{O}_{\mathrm{U}}))=\operatorname{Spec}(\mathscr{A}(\mathrm{U}))\quad(\mathbf{II},~1.1.1).
\]
If $f^0:\mathrm{U}^0\to\mathrm{X}$ is the structural morphism, we have therefore by definition
\[
  \mathscr{A}(\mathrm{U}^0)=f_*^0(\sh{O}_{\mathrm{U}^0})=f_*(\sh{O}_{\mathrm{U}}),
\]
and at the identity isomorphism of $f_*(\sh{O}_{\mathrm{U}})$, there corresponds by $(\mathbf{II}$, $1.2.7)$ a canonical $\mathrm{X}$-morphism
\[
\label{IV.21.12.1.1}
  i_{\mathrm{U}}:\mathrm{U}\to\mathrm{U}^0.
  \tag{21.12.1.1}
\]

For $i_{\mathrm{U}}$ to be an \emph{isomorphism}, it is necessary and sufficient that the morphism $f:\mathrm{U}\to\mathrm{X}$ be \emph{affine}.
For every affine $\mathrm{X}$-prescheme $\mathrm{V}$, the map $u\mapsto u\circ i_{\mathrm{U}}$ is a \emph{bijection}
\[
  \operatorname{Hom}_{\mathrm{X}}(\mathrm{U}^0,\mathrm{V})\xrightarrow{\sim}\operatorname{Hom}_{\mathrm{X}}(\mathrm{U},\mathrm{V})
\]
functorial in $\mathrm{V}$: this follows from the existence of canonical bijections
\[
  \operatorname{Hom}_{\mathrm{X}}(\mathrm{U},\mathrm{V})\xrightarrow{\sim}\operatorname{Hom}_{\sh{O}_{\mathrm{X}}}(\mathscr{A}(\mathrm{V}),\mathscr{A}(\mathrm{U}))
\]
and $\operatorname{Hom}_{\mathrm{X}}(\mathrm{U}^0,\mathrm{V})\xrightarrow{\sim}\operatorname{Hom}_{\sh{O}_{\mathrm{X}}}(\mathscr{A}(\mathrm{V}),\mathscr{A}(\mathrm{U}^0))$ $(\mathbf{II}$, $1.2.7)$.
We can therefore say that $\mathrm{U}^0$ \emph{represents} the covariant functor $\mathrm{V}\mapsto\operatorname{Hom}_{\mathrm{X}}(\mathrm{U},\mathrm{V})$ in the category of X-preschemes affine over $\mathrm{X}$.
We deduce from this $(\mathbf{0}_{\mathrm{III}}$, $8.1.7)$ that for a fixed $\mathrm{X}$, $\mathrm{U}\mapsto\operatorname{Aff}(\mathrm{U}/\mathrm{X})$ is a covariant functor from the category of $\mathrm{X}$-preschemes quasi-compact and quasi-separated over $\mathrm{X}$, into the category of $\mathrm{X}$-preschemes affine over $\mathrm{X}$; more precisely, if $\mathrm{U}_1$, $\mathrm{U}_2$ are two $\mathrm{X}$-preschemes quasi-compact and quasi-separated over $\mathrm{X}$, to every $\mathrm{X}$-morphism $g:\mathrm{U}_1\to\mathrm{U}_2$ corresponds the unique $\mathrm{X}$-morphism $g^0:\mathrm{U}_1^0\to\mathrm{U}_2^0$ making commutative the diagram
\[
  \xymatrix{
    \mathrm{U}_1 \ar[r]^g \ar[d]_{i_{\mathrm{U}_1}} & \mathrm{U}_2 \ar[d]^{i_{\mathrm{U}_2}} \\
    \mathrm{U}_1^0 \ar[r]_{g^0} & \mathrm{U}_2^0
  }
\]
\end{env}

\oldpage[IV-4]{305}
More generally, consider a commutative diagram
\[
  \xymatrix{
    \mathrm{U} \ar[r]^v \ar[d]_f & \mathrm{U}' \ar[d]^{f'} \\
    \mathrm{X} \ar[r]_u & \mathrm{X}'
  }
\]
where the morphisms $f$, $f'$ are quasi-compact and quasi-separated and the morphism $u$ \emph{affine}.
If $h=u\circ f$, we have $h_*(\sh{O}_{\mathrm{U}})=u_*(f_*(\sh{O}_{\mathrm{U}}))=u_*(f_*^{0}(\sh{O}_{\mathrm{U}^0}))$ and $u\circ f^0$ is an affine morphism; hence $\mathrm{U}^0=\operatorname{Aff}(\mathrm{U}/\mathrm{X}')$ (relative to the morphism $h$), whence a unique $\mathrm{X}'$-morphism $v^0:\mathrm{U}^0\to\mathrm{U}'^0$ making commutative the diagram
\[
  \xymatrix{
    \mathrm{U} \ar[r]^v \ar[d]_{i_{\mathrm{U}}} & \mathrm{U}' \ar[d]^{i_{\mathrm{U}'}} \\
    \mathrm{U}^0 \ar[r]_{v^0} & \mathrm{U}'^0
  }
\]

\begin{proposition}[21.12.2]
\label{IV.21.12.2}
Let $f:\mathrm{U}\to\mathrm{X}$ be a quasi-compact and quasi-separated morphism, $h:\mathrm{X}'\to\mathrm{X}$ a flat morphism; set $\mathrm{U}'=\mathrm{U}\times_{\mathrm{X}}\mathrm{X}'$, $f'=f_{\{X'\}}:\mathrm{U}'\to\mathrm{X}'$.
Then there is a canonical $\mathrm{X}'$-isomorphism
\[
\label{IV.21.12.2.1}
  \operatorname{Aff}(\mathrm{U}'/\mathrm{X}')\xrightarrow{\sim}\operatorname{Aff}(\mathrm{U}/\mathrm{X})\times_{\mathrm{X}}\mathrm{X}'.
  \tag{21.12.2.1}
\]
\end{proposition}

\begin{proof}
Indeed, we have $\operatorname{Aff}(\mathrm{U}'/\mathrm{X}')=\operatorname{Spec}(f'_*(\sh{O}_{\mathrm{U}'}))$, and $\operatorname{Aff}(\mathrm{U}/\mathrm{X})\times_{\mathrm{X}}\mathrm{X}'=\operatorname{Spec}(h^*(f_*(\sh{O}_{\mathrm{U}})))$; the isomorphism of the assertion comes from the canonical isomorphism $h^*(f_*(\sh{O}_{\mathrm{U}}))\xrightarrow{\sim}f'_*(\sh{O}_{\mathrm{U}'})$ \sref{IV.2.3.1}.
\end{proof}

\begin{corollary}[21.12.3]
\label{IV.21.12.3}
For every quasi-compact and quasi-separated morphism $f:\mathrm{U}\to\mathrm{X}$ and every $x\in\mathrm{X}$, there is, up to canonical isomorphism, an equality
\[
  \mathrm{U}^0\times_{\mathrm{X}}\operatorname{Spec}(\sh{O}_{\mathrm{X},x})=(\mathrm{U}\times_{\mathrm{X}}\operatorname{Spec}(\sh{O}_{\mathrm{X},x}))^0.
\]
\end{corollary}

It also follows from \sref{IV.21.12.2} that, up to canonical isomorphism, for every open $\mathrm{V}$ of $\mathrm{X}$
\[
\label{IV.21.12.4}
  (f^{-1}(\mathrm{V}))^0=(f^0)^{-1}(\mathrm{V}).
  \tag{21.12.4}
\]

\begin{env}[21.12.5]
\label{IV.21.12.5}
We will consider in particular the case where $f:\mathrm{U}\to\mathrm{X}$ is an \emph{open immersion}, $\mathrm{U}$ thus being identified with an open of $\mathrm{X}$.
Since the morphism $f^{-1}(\mathrm{U})\to\mathrm{U}$ is the identity, it follows from \sref{IV.21.12.4} that the morphism $(f^0)^{-1}(\mathrm{U})\to\mathrm{U}$ restriction of $f^0$ is an isomorphism, $i_{\mathrm{U}}:\mathrm{U}\to\mathrm{U}^0$ being thus an \emph{open immersion} allowing us to identify $\mathrm{U}$ with an open of $\mathrm{U}^0$.
More generally, for an open $\mathrm{V}$ of $\mathrm{X}$, the restriction $(f^0)^{-1}(\mathrm{V})\to\mathrm{V}$ of $f^0$ is an isomorphism of $(f^0)^{-1}(\mathrm{V})$ onto the prescheme induced on the open $\mathrm{V}\cap\mathrm{U}$ if and only if the open immersion $\mathrm{U}\cap\mathrm{V}\to\mathrm{V}$ is an \emph{affine} morphism.
It is clear $(\mathbf{II}$, $1.2.1)$
\oldpage[IV-4]{306}
that the union of these opens $\mathrm{V}$ is the \emph{largest} of them, $\mathrm{U}_1$, which contains $\mathrm{U}$; by virtue of what precedes, $\mathrm{U}_1$ is also the largest open not meeting the set
\[
  \operatorname{Daf}(\mathrm{U}/\mathrm{X})=f^0(\mathrm{U}^0)-\mathrm{U}
\]
(``affinity defect'' of the open $\mathrm{U}$ relative to $\mathrm{X}$, which is empty if and only if $\mathrm{U}$ is affine over $\mathrm{X}$); in other words, the closed set $\mathrm{Z}=\mathrm{X}-\mathrm{U}_1$ is the \emph{closure} of $\operatorname{Daf}(\mathrm{U}/\mathrm{X})$.

We note that for every \emph{flat} morphism $h:\mathrm{X}'\to\mathrm{X}$, if we set $\mathrm{U}'=h^{-1}(\mathrm{U})$, we have
\[
\label{IV.21.12.5.1}
  \operatorname{Daf}(\mathrm{U}'/\mathrm{X}')=h^{-1}(\operatorname{Daf}(\mathrm{U}/\mathrm{X}))
  \tag{21.12.5.1}
\]
as follows immediately from \sref{IV.21.12.2.1} and from $(\mathbf{I}$, $3.4.8)$.
In particular, for every open $\mathrm{V}$ of $\mathrm{X}$, we have
\[
\label{IV.21.12.5.2}
  \operatorname{Daf}((\mathrm{U}\cap\mathrm{V})/\mathrm{V})=\operatorname{Daf}(\mathrm{U}/\mathrm{X})\cap\mathrm{V}
  \tag{21.12.5.2}
\]
and for all $x\in\mathrm{X}$,
\[
\label{IV.21.12.5.3}
  \operatorname{Daf}((\mathrm{U}\cap\operatorname{Spec}(\sh{O}_{\mathrm{X},x}))/\operatorname{Spec}(\sh{O}_{\mathrm{X},x}))=\operatorname{Daf}(\mathrm{U}/\mathrm{X})\cap\operatorname{Spec}(\sh{O}_{\mathrm{X},x}).
  \tag{21.12.5.3}
\]
\end{env}

When $\mathrm{X}$ is locally Noetherian, we shall give more precise information on the nature of $\operatorname{Daf}(\mathrm{U}/\mathrm{X})$; it will show, for example, that when $\mathrm{U}$ is everywhere dense in $\mathrm{X}$, $\operatorname{Daf}(\mathrm{U}/\mathrm{X})$ is not just an arbitrary rare closed subset:

\begin{theorem}[21.12.6]
\label{IV.21.12.6}
Let $\mathrm{X}$ be a locally Noetherian prescheme, $\mathrm{U}$ a non-empty open, $f:\mathrm{U}\to\mathrm{X}$ the canonical injection.
Then:
\begin{enumerate}[label=(\roman*)]
  \item The closure $\mathrm{Z}=\mathrm{X}-\mathrm{U}_1$ of $\operatorname{Daf}(\mathrm{U}/\mathrm{X})$ is of codimension $\geq 2$ in $\mathrm{X}$.
  \item If $\mathrm{T}=\mathrm{X}-\mathrm{U}\supset\mathrm{Z}$ is of codimension $\geq 2$, the morphism $f^0:\mathrm{U}^0\to\mathrm{X}$ is surjective.
\end{enumerate}
\end{theorem}

\begin{proof}
(i) Let us first show that for every point $x\in\operatorname{Daf}(\mathrm{U}/\mathrm{X})$ we have necessarily $\dim(\sh{O}_{\mathrm{X},x})>1$.
Indeed, $x$ plainly cannot be a maximal point of $\mathrm{X}$, since it lies in $\overline{\mathrm{U}}-\mathrm{U}$; it remains to see that $\dim(\sh{O}_{\mathrm{X},x})$ cannot equal $1$.
That equality would imply, by \sref{IV.21.12.5.3}, that $x\in\operatorname{Daf}(\mathrm{U}^{(x)}/\operatorname{Spec}(\sh{O}_{\mathrm{X},x}))$, where $\mathrm{U}^{(x)}=\mathrm{U}\cap\operatorname{Spec}(\sh{O}_{\mathrm{X},x})$.
But the only open subsets of $\operatorname{Spec}(\sh{O}_{\mathrm{X},x})$ are $\operatorname{Spec}(\sh{O}_{\mathrm{X},x})$ itself and the parts of the (finite) set of maximal points of $\operatorname{Spec}(\sh{O}_{\mathrm{X},x})$.
Now, an open set reduced to a single maximal point of $\operatorname{Spec}(\sh{O}_{\mathrm{X},x})$ is affine by the definition of preschemes; we conclude that \emph{all} open subsets of $\operatorname{Spec}(\sh{O}_{\mathrm{X},x})$ are affine, hence $\operatorname{Daf}(\mathrm{U}^{(x)}/\operatorname{Spec}(\sh{O}_{\mathrm{X},x}))=\varnothing$, contradicting the hypothesis.

To prove (i), it remains to prove the stronger assertion that if $x\in\mathrm{X}$ satisfies $\dim(\sh{O}_{\mathrm{X},x})=1$, then there exists an open neighbourhood $\mathrm{W}$ of $x$ in $\mathrm{X}$ which does not meet $\operatorname{Daf}(\mathrm{U}/\mathrm{X})$, i.e.\ for which the canonical injection $f_{\mathrm{W}}:\mathrm{U}\cap\mathrm{W}\to\mathrm{W}$ is affine.
But, with the same notation as above, the injection $f^{(x)}:\mathrm{U}^{(x)}\to\operatorname{Spec}(\sh{O}_{\mathrm{X},x})$ is affine.
We may plainly reduce to the case where $\mathrm{X}$ is Noetherian; since $f$ is then of finite presentation, the conclusion follows from $(\mathbf{IV}$, $8.10.5$,~(viii)), by the method of $(\mathbf{IV}$, $8.1.2$,~a)).

\oldpage[IV-4]{307}

(ii) It remains to prove that $x\in f^0(\mathrm{U}^0)$ for every $x\in\mathrm{T}$.
We first reduce to the case where $\mathrm{X}=\operatorname{Spec}(\mathrm{A})$, with $\mathrm{A}$ a complete local Noetherian ring and $x$ the closed point of $\mathrm{X}$.
Indeed, make the base change $h:\mathrm{X}'=\operatorname{Spec}(\widehat{\sh{O}}_{\mathrm{X},x})\to\mathrm{X}$, set $\mathrm{U}'=h^{-1}(\mathrm{U})$, and let $f'=f_{\{\mathrm{X}'\}}:\mathrm{U}'\to\mathrm{X}'$ be the canonical injection.
Since $h$ is flat, \sref{IV.21.12.2} shows that if the closed point belongs to $f^{\prime 0}(\mathrm{U}^{\prime 0})$, then $x\in f^0(\mathrm{U}^0)$; by \sref{IV.6.1.1}, this gives the desired reduction.

Let $\mathrm{X}_1$ be the reduced closed sub-prescheme of $\mathrm{X}$ whose underlying space is an irreducible component of $\mathrm{X}$ of \emph{maximal} dimension among those containing an irreducible component of $\mathrm{T}$, and set
\[
  \mathrm{U}_1=\mathrm{U}\cap\mathrm{X}_1,
  \qquad
  \mathrm{T}_1=\mathrm{T}\cap\mathrm{X}_1=\mathrm{X}_1-\mathrm{U}_1.
\]
Then $\operatorname{codim}(\mathrm{T}_1,\mathrm{X}_1)\geq 2$ $(\mathbf{0}$, $14.2.1)$, and the pair $(\mathrm{U}_1,\mathrm{X}_1)$ satisfies the same hypotheses as $(\mathrm{U},\mathrm{X})$; indeed, $\mathrm{X}_1$ is the spectrum of a quotient of $\mathrm{A}$ and is therefore local Noetherian and complete.
We also have a commutative diagram
\[
  \xymatrix{
    \mathrm{U}_1 \ar[r]^j \ar[d]_{f_1} & \mathrm{U} \ar[d]^f \\
    \mathrm{X}_1 \ar[r]_i & \mathrm{X},
  }
\]
where $i$ and $j$ are the canonical injections, and $i$ is therefore \emph{affine}.  From \sref{IV.21.12.1} we obtain a morphism $j^0:\mathrm{U}_1^0\to\mathrm{U}^0$ making commutative the diagram
\[
  \xymatrix{
    \mathrm{U}_1^0 \ar[r]^{j^0} \ar[d]_{f_1^0} & \mathrm{U}^0 \ar[d]^{f^0} \\
    \mathrm{X}_1 \ar[r]_i & \mathrm{X}.
  }
\]
Consequently, to prove that $x\in f^0(\mathrm{U}^0)$, it suffices to prove that $x\in f_1^0(\mathrm{U}_1^0)$.
We may therefore replace $\mathrm{X}$ by $\mathrm{X}_1$ and suppose in addition that $\mathrm{A}$ is \emph{integral}.
By the Cohen structure theorem $(\mathbf{0}$, $19.8.8$,~(i)), $\mathrm{A}$ is a quotient of a regular Noetherian ring; since it is integral, we may apply $(\mathbf{IV}$, $5.11.1)$ with the family $(x_\alpha)$ reduced to the single maximal point of $\mathrm{X}$.
As $\operatorname{codim}(\mathrm{T},\mathrm{X})\geq 2$, it follows that $f_*(\sh{O}_{\mathrm{U}})$ is a \emph{coherent} $\sh{O}_{\mathrm{X}}$-Module, so $f^0:\mathrm{U}^0\to\mathrm{X}$ is \emph{finite}.  Since it is dominant ($\mathrm{U}$ being everywhere dense), it is surjective $(\mathbf{II}$, $6.1.10)$, and therefore $x\in f^0(\mathrm{U}^0)$.
\end{proof}

\begin{corollary}[21.12.7]
\label{IV.21.12.7}
Let $\mathrm{X}$ be a locally Noetherian prescheme, $\mathrm{U}$ a sub-prescheme induced on an open of $\mathrm{X}$, $j:\mathrm{U}\to\mathrm{X}$ the canonical immersion; suppose that $j$ is an affine morphism.
Then every irreducible component of $\mathrm{T}=\mathrm{X}-\mathrm{U}$ is of codimension $\leq 1$ (and hence of codimension $1$ if $\mathrm{U}$ is everywhere dense).
\end{corollary}

\begin{proof}
Suppose indeed that one of the irreducible components $\mathrm{T}_1$ of $\mathrm{T}$ is of codimension $\geq 2$ in $\mathrm{X}$.
Replacing $\mathrm{X}$ if needed by an open neighbourhood of the generic point of $\mathrm{T}_1$, we can suppose $\mathrm{T}$ irreducible and of codimension $\geq 2$.
But then the hypothesis
\oldpage[IV-4]{308}
that $j:\mathrm{U}\to\mathrm{X}$ is affine implies that $\mathrm{U}^0$ is identified with $\mathrm{U}$ and $j^0$ with $j$, in other words $j^0(\mathrm{U}^0)=\mathrm{U}$; but this contradicts the conclusion of \sref{IV.21.12.6} which, under the hypothesis of dimension, says that $j^0$ must be surjective.
\end{proof}

\begin{env}[21.12.8]
\label{IV.21.12.8}
Let $\mathrm{A}$ be a local Noetherian ring, $\mathrm{Y}=\operatorname{Spec}(\mathrm{A})$, $y$ the unique closed point of $\mathrm{Y}$, $\mathrm{Y}'=\mathrm{Y}-\{y\}$.
Consider the following condition:
\begin{itemize}
  \item[\textup{(W)}] \emph{For every open $\mathrm{U}$ of $\mathrm{Y}$ contained in $\mathrm{Y}'$, containing no irreducible component of $\mathrm{Y}'$ and such that the canonical injection $\mathrm{U}\to\mathrm{Y}'$ is affine, $\mathrm{U}$ is itself an affine open.}
\end{itemize}

This condition is implied by the following:
\begin{itemize}
  \item[\textup{(W')}] \emph{For every closed subset $\mathrm{T}$ of $\mathrm{Y}$ whose every irreducible component is of codimension $1$ and such that for every irreducible component $\mathrm{Y}_i$ of $\mathrm{Y}$, $\mathrm{Y}_i\cap\mathrm{T}\neq\{y\}$, the open $\mathrm{U}=\mathrm{Y}-\mathrm{T}$ is affine.}
\end{itemize}

Indeed, if (W') holds and $\mathrm{U}$ satisfies the hypotheses of (W), \sref{IV.21.12.7} shows that no irreducible component of $\mathrm{T}=\mathrm{Y}-\mathrm{U}$ can have codimension $\geq 2$; moreover, $\mathrm{U}$ contains none of the irreducible components $\mathrm{Y}_i-\{y\}$ of $\mathrm{Y}'$, so (W') shows that $\mathrm{U}$ is affine.

We note that condition (W') simplifies when $\mathrm{Y}$ is \emph{irreducible} and is then equivalent to the following:
\begin{itemize}
  \item[\textup{(W'')}] \emph{For every closed irreducible subset $\mathrm{T}$ of $\mathrm{Y}$, of codimension $1$ in $\mathrm{Y}$, $\mathrm{Y}-\mathrm{T}$ is an affine open.}
\end{itemize}

Indeed, it is clear that (W') implies (W'') when $\mathrm{Y}$ is irreducible, and the converse is true by considering the irreducible components of $\mathrm{T}$ and of $\mathrm{Y}$ and noting that the intersection of a finite number of affine open subsets is an affine open $(\mathbf{I}$, $5.5.6)$.
\end{env}

\begin{env}[21.12.9]
\label{IV.21.12.9}
\emph{Examples}. ---
If $\mathrm{A}$ is a local Noetherian \emph{factorial} ring, it satisfies (W') (and \emph{a fortiori} (W)), since every prime ideal of height $1$ is principal $(\mathbf{I}$, $1.3.6)$.
There are nevertheless nonfactorial local Noetherian rings satisfying (W''), for example in dimension $\leq 1$: indeed, as remarked in the proof of \sref{IV.21.12.6},~(i), \emph{all} open subsets of $\mathrm{Y}$ are then affine.
On the other hand, the theory of local duality (chap.~III, Part~3) shows that every Noetherian local ring of dimension $2$ satisfies (W').
\end{env}

The interest of condition (W) resides in the following result:
\begin{proposition}[21.12.10]
\label{IV.21.12.10}
Let $\mathrm{X}$, $\mathrm{Y}$ be two locally Noetherian preschemes, $g:\mathrm{X}\to\mathrm{Y}$ a morphism, $y$ a closed point of $\mathrm{Y}$, $\mathrm{Y}'=\mathrm{Y}-\{y\}$, and $\mathrm{X}'=g^{-1}(\mathrm{Y}')$; suppose that the restriction $g':\mathrm{X}'\to\mathrm{Y}'$ is an open immersion, and that the local ring $\sh{O}_{\mathrm{Y},y}$ satisfies condition \textup{(W)} of \sref{IV.21.12.8}.
Then, for every irreducible component $\mathrm{Z}$ of $\mathrm{X}_y=g^{-1}(y)$, either $\mathrm{Z}$ has codimension $\leq 1$ in $\mathrm{X}$, or its generic point is isolated in $\mathrm{Z}$.
If $\mathrm{Z}$ is locally of finite type over $\boldsymbol{k}(y)$, the second alternative implies that $\mathrm{Z}$ is reduced to a single point.
\end{proposition}

\begin{proof}
The last assertion of the statement results from the fact that $\mathrm{Z}$ is then a prescheme of Jacobson $(\mathbf{I}$, $10.4.7)$, and since the set of closed points of $\mathrm{Z}$ is then dense in $\mathrm{Z}$, the generic point of $\mathrm{Z}$ cannot be isolated in $\mathrm{Z}$ unless $\mathrm{Z}$ is reduced to a single point.

\oldpage[IV-4]{309}
Suppose that $\mathrm{X}_y$ has an irreducible component $\mathrm{Z}$ whose generic point $z$ is not isolated in $\mathrm{Z}$ and for which $\operatorname{codim}(\mathrm{Z},\mathrm{X})\geq 2$.
The question is local on $\mathrm{X}$ and $\mathrm{Y}$; since $z$ is not isolated in $\mathrm{Z}$, after replacing $\mathrm{X}$ by an open neighbourhood of $z$ we may suppose that $\mathrm{X}$ and $\mathrm{Y}$ are affine, that $\mathrm{X}_y=\mathrm{Z}$ is irreducible and not reduced to a single point, and that $\operatorname{codim}(\mathrm{X}_y,\mathrm{X})\geq 2$.
The image $g(\mathrm{X}-\mathrm{X}_y)=\mathrm{U}$ is an open subset of $\mathrm{Y}'$ isomorphic to $\mathrm{X}-\mathrm{X}_y$ by hypothesis.  We claim that, after shrinking $\mathrm{X}$ around $z$ if necessary, $\mathrm{U}$ contains none of the irreducible components of $\mathrm{Y}'$ whose closure contains $y$.

We may first suppose that every maximal point of $\mathrm{X}$ is a generisation of $z$; there are only finitely many such points.  The maximal points of $\mathrm{U}$ are then the images $y_i=g(x_i)$ of the maximal points $x_i$ of $\mathrm{X}$ (no maximal point of $\mathrm{X}$ can lie in $\mathrm{X}_y$, since $\dim(\sh{O}_{\mathrm{X},z})\geq 2$).  Put $\mathrm{X}_i=\overline{\{x_i\}}$.  By hypothesis $z\in\mathrm{X}_i$, and because $z$ is not isolated in $\mathrm{X}_y$, there is a point $x\in\mathrm{X}_y$ with $x\neq z$.  Since $\mathrm{X}_i\neq\mathrm{Z}$, there exists a point $t_i\in\mathrm{X}_i$ such that, if $\mathrm{T}_i=\overline{\{t_i\}}$, then $\dim(\sh{O}_{\mathrm{T}_i,x})=1$ $(\mathbf{I}$, $10.5.9)$.  Hence $t_i\notin\mathrm{Z}$, so $t_i$ is not a generisation of $z$.
Replacing $\mathrm{X}$ by $\mathrm{X}'=\mathrm{X}-\bigcup_i\mathrm{T}_i$, the image under $g$ of $\mathrm{X}'-\mathrm{X}_y$ contains none of the points $g(t_i)$, and consequently none of the irreducible components of $\mathrm{Y}'$ whose closure contains $y$.

Since $\mathrm{X}$ and $\mathrm{Y}$ are affine, $g:\mathrm{X}\to\mathrm{Y}$ is affine, and so is its restriction $g':\mathrm{X}'\to\mathrm{Y}'$.  Because $g'$ is also an open immersion, the canonical immersion $\mathrm{U}\to\mathrm{Y}'$ is affine.
Set $\mathrm{Y}_1=\operatorname{Spec}(\sh{O}_{\mathrm{Y},y})$, $\mathrm{Y}_1'=\mathrm{Y}_1-\{y\}=\mathrm{Y}'\cap\mathrm{Y}_1$, and $\mathrm{U}_1=\mathrm{U}\cap\mathrm{Y}_1$.  The preceding paragraph shows that $\mathrm{U}_1\to\mathrm{Y}_1'$ is affine, hence (W) says that $\mathrm{U}_1$ is an affine open of $\mathrm{Y}_1$, or equivalently that $\mathrm{U}_1\to\mathrm{Y}_1$ is affine.
This immersion is of finite presentation $(\mathbf{I}$, $6.2)$; applying $(\mathbf{IV}$, $8.10.5$,~(viii)) by the method of $(\mathbf{IV}$, $8.1.2$,~a)), and shrinking $\mathrm{Y}$ to an affine open neighbourhood of $y$ if necessary, we may suppose that $\mathrm{U}\to\mathrm{Y}$ is affine.  It would follow that the open $\mathrm{X}-\mathrm{X}_y$ of $\mathrm{X}$, isomorphic to $\mathrm{U}$, is affine, contradicting \sref{IV.21.12.7}.  This proves the proposition.
\end{proof}

\begin{corollary}[21.12.11]
\label{IV.21.12.11}
Let $\mathrm{X}$, $\mathrm{Y}$ be two locally Noetherian preschemes, with $\mathrm{X}$ irreducible; let $g:\mathrm{X}\to\mathrm{Y}$ be a morphism locally of finite type, and let $\mathrm{V}$ be the largest open subset of $\mathrm{X}$ such that the restriction $g|\mathrm{V}:\mathrm{V}\to\mathrm{Y}$ is a local isomorphism.
Suppose that for every point $y\in\mathrm{Y}$, the local ring $\sh{O}_{\mathrm{Y},y}$ satisfies condition \textup{(W)} of \sref{IV.21.12.8}.
Then every irreducible component of $\mathrm{T}=\mathrm{X}-\mathrm{V}$ either has codimension $\leq1$, or has a generic point $z$ isolated in $g^{-1}(g(z))$.
\end{corollary}

\begin{proof}
Put $g(z)=y$. Since the question depends only on the fibre $g^{-1}(y)$ and on the ring $\sh{O}_{\mathrm{Y},y}$, base change reduces us to the case $\mathrm{Y}=\operatorname{Spec}(\sh{O}_{\mathrm{Y},y})$ with $\mathrm{X}$ affine $((\mathbf I,3.6.5)$ and $(\mathbf I,4.5.5))$; after replacing $\mathrm{X}$ by an open neighbourhood of $z$, we may suppose that $\mathrm{T}$ is irreducible, and we may also restrict to the case where $\mathrm{V}$ is non-empty.
The morphism $g$ may therefore be supposed separated and $y$ closed in $\mathrm{Y}$; since $g|\mathrm{V}:\mathrm{V}\to\mathrm{Y}$ is a local isomorphism and the non-empty open $\mathrm{V}$ of $\mathrm{X}$ is irreducible, $(\mathbf I,8.2.8)$ shows that $g|\mathrm{V}$ is an open immersion.
The restriction $g|(\mathrm{V}\cap\mathrm{X}_y):\mathrm{V}\cap\mathrm{X}_y\to\{y\}=\operatorname{Spec}(\boldsymbol{k}(y))$ is therefore also an open immersion, which shows that $\mathrm{V}\cap\mathrm{X}_y$ is either empty or reduced to one point rational
\oldpage[IV-4]{310}
over $\boldsymbol{k}(y)$, hence closed in $\mathrm{X}_y$ $(\mathbf I,6.4.2)$.
After replacing $\mathrm{X}$ once more by a smaller open neighbourhood of $z$, we may thus reduce to the case $\mathrm{V}\cap\mathrm{X}_y=\varnothing$, i.e.\ $\mathrm{T}\supset\mathrm{X}_y$; since moreover $z\in\mathrm{X}_y$ is the generic point of $\mathrm{T}$, we have $g(\mathrm{T})\subset\overline{\{y\}}=\{y\}$, hence $\mathrm{T}=\mathrm{X}_y$. We are then exactly under the hypotheses of \sref{IV.21.12.10}, which gives the conclusion.
\end{proof}

\begin{theorem}[21.12.12]
\label{IV.21.12.12}
\textup{(van der Waerden)} ---
Let $\mathrm{X}$, $\mathrm{Y}$ be two locally Noetherian integral preschemes, $g:\mathrm{X}\to\mathrm{Y}$ a birational morphism and locally of finite type.
Suppose furthermore that $\mathrm{Y}$ is normal and that for all $y\in\mathrm{Y}$, the local ring $\sh{O}_{\mathrm{Y},y}$ satisfies the condition \textup{(W)} of \sref{IV.21.12.8} (conditions satisfied in particular when $\mathrm{Y}$ is \emph{locally factorial} \sref{IV.21.12.9}).
If $\mathrm{V}$ is the largest open of $\mathrm{X}$ such that the restriction $g|\mathrm{V}:\mathrm{V}\to\mathrm{Y}$ is a local isomorphism, all the irreducible components of $\mathrm{T}=\mathrm{X}-\mathrm{V}$ are of codimension $1$ in $\mathrm{X}$.
\end{theorem}

\begin{proof}
We note that since $g$ is birational, the open $\mathrm{V}$ is not empty; it suffices therefore to prove that a maximal point $z$ of $\mathrm{T}$ cannot be isolated in $\mathrm{X}_y$, where $y=g(z)$.
But, after restricting to an open neighbourhood of $z$, we may reduce to the case $\mathrm{T}=\mathrm{X}_y$; then every fibre $g^{-1}(y')$ ($y'\in\mathrm{Y}$) would be empty or reduced to a single point, and the ``Main theorem'' \sref{IV.8.12.10} would imply that $g$ is a local isomorphism, contradicting $z\in\mathrm{T}$.
\end{proof}

\begin{corollary}[21.12.13]
\label{IV.21.12.13}
Suppose the hypotheses of \sref{IV.21.12.12} hold and, in addition, that $g$ is quasi-finite at every point of $\mathrm{X}^{(1)}$ (recall that these are the points $x$ for which $\dim(\sh{O}_{\mathrm{X},x})=1$).  Then $g$ is a local isomorphism; if moreover $g$ is separated, it is an open immersion.
\end{corollary}

\begin{proof}
It suffices to prove the first assertion, the second being a consequence of the first and of $(\mathbf{I}$, $8.2.8)$.
It all comes down to proving, with the notations of \sref{IV.21.12.12}, that $\mathrm{T}=\varnothing$; in the contrary case, a generic point $z$ of an irreducible component of $\mathrm{T}$ would belong to $\mathrm{X}^{(1)}$ by virtue of \sref{IV.21.12.12}, and by hypothesis it would be isolated in $\mathrm{X}_y$ with $y=g(z)$; but as we saw in the proof of \sref{IV.21.12.12} this is not possible.
\end{proof}

\begin{remarks}[21.12.14]
\label{IV.21.12.14}
(i)
The conclusion of \sref{IV.21.12.12} applies when $\mathrm{Y}$ is regular and integral and $\mathrm{X}$ is integral; by virtue of the Auslander-Buchsbaum theorem \sref{IV.21.11.1}, $\mathrm{Y}$ is then locally factorial.
On the contrary, one knows examples where $\mathrm{X}$ and $\mathrm{Y}$ are normal algebraic schemes of dimension $3$, over an algebraically closed field of any characteristic, and where the conclusion of \sref{IV.21.12.12} is no longer valid.

(ii)
The set $\mathrm{T}$ of \sref{IV.21.12.12} is also the set of points where $g$ is \emph{ramified}.
Indeed, if $g$ is unramified at a point $x\in\mathrm{X}$, it is also unramified in an open neighbourhood of $x$ \sref{IV.17.3.7}, hence $g|\mathrm{U}$ is a separated, quasi-finite and birational morphism.
Since $\mathrm{Y}$ is normal, we conclude from the ``Main theorem'' $(\mathbf{III}$, $4.4.9)$ that $g$ is an isomorphism local at the point $x$, hence $x\in\mathrm{X}-\mathrm{T}$.
Conversely, $g$ is evidently unramified at every point where it is a local isomorphism.
This justifies the title given to this section.

\oldpage[IV-4]{311}
(iii)
Without assuming that $\mathrm{Y}$ is locally factorial, but assuming instead that $\mathrm{X}$ is a complete intersection over $\mathrm{Y}$ (chap.~V), we shall see in chap.~V a result more precise than \sref{IV.21.12.12}, expressing $\mathrm{T}$ as the $1$-codimensional cycle of a section of an invertible $\sh{O}_{\mathrm{X}}$-Module canonically associated with $g$.
This applies notably when $\mathrm{X}$ and $\mathrm{Y}$ are both regular, or when $g$ is a S-morphism, $\mathrm{X}$ and $\mathrm{Y}$ being smooth preschemes over $\mathrm{S}$.

(iv) One may ask whether \sref{IV.21.12.10} has a converse.  In other words, let $\mathrm{A}$ be a local Noetherian ring, set $\mathrm{Y}=\operatorname{Spec}(\mathrm{A})$, and let $y$ be its closed point.  Suppose that, for every locally Noetherian prescheme $\mathrm{X}$ and every morphism $g:\mathrm{X}\to\mathrm{Y}$ for which $g':\mathrm{X}'\to\mathrm{Y}'$ is an open immersion, every irreducible component of $\mathrm{X}_y$ has codimension $\leq 1$ in $\mathrm{X}$ or is reduced to a single point.  Must $\mathrm{A}$ then satisfy condition (W) of \sref{IV.21.12.8}?

(v) Let $\mathrm{Y}$ be an integral \emph{regular} prescheme, $\mathrm{X}$ a locally Noetherian \emph{normal} prescheme, $g:\mathrm{X}\to\mathrm{Y}$ a morphism locally of finite type; suppose furthermore that, if $\eta$ is the generic point of $\mathrm{Y}$, the fibre $\mathrm{X}_\eta$ is \emph{\'etale} over $\boldsymbol{k}(\eta)$, and let $\mathrm{T}'$ be the complement of the largest open $\mathrm{V}'$ of $\mathrm{X}$ such that $g|\mathrm{V}':\mathrm{V}'\to\mathrm{Y}$ is \emph{\'etale}; is it true that all the irreducible components of $\mathrm{T}'$ are of codimension $1$?
This is indeed the case as well when one moreover supposes that $g$ is \emph{locally quasi-finite} (``purity theorem'' of Zariski-Nagata).
One can show that the preceding conjecture would follow from the following statement (and would even be equivalent to it if the conjecture in (iv) held): if a regular local ring $\mathrm{A}$ is contained in a local integral and integrally closed ring $\mathrm{B}$ which is finite as an A-module, then the open subset $\mathrm{U}$ of $\operatorname{Spec}(\mathrm{B})$ at whose points $\operatorname{Spec}(\mathrm{B})$ is unramified over $\operatorname{Spec}(\mathrm{A})$---or, equivalently by $(\mathbf{IV}$, $18.10.1)$, \'etale over $\operatorname{Spec}(\mathrm{A})$---is affine.
\end{remarks}

\begin{proposition}[21.12.15]
\label{IV.21.12.15}
Let $\mathrm{S}$ be a prescheme, $g:\mathrm{X}\to\mathrm{S}$ a flat morphism locally of finite presentation, whose fibres $\mathrm{X}_s=g^{-1}(s)$ are geometrically irreducible \sref{IV.4.5.2}, $h:\mathrm{Y}\to\mathrm{S}$ a smooth morphism; denote by $\mathrm{Y}_s=h^{-1}(s)$ the fibres of this morphism.
Let $f:\mathrm{X}\to\mathrm{Y}$ be a proper S-morphism such that, for every maximal point $\eta$ of $\mathrm{S}$, the morphism $f_\eta:\mathrm{X}_\eta\to\mathrm{Y}_\eta$ is an isomorphism.
Then $f$ is an isomorphism.
\end{proposition}

\begin{proof}
\oldpage[IV-4]{312}
Since $h$ is flat, every maximal point of $\mathrm{Y}$ lies above a maximal point of $\mathrm{S}$ \sref{IV.2.3.4}, hence the union of the $\mathrm{Y}_\eta$, as $\eta$ ranges over the set of maximal points of $\mathrm{S}$, is dense in $\mathrm{Y}$; since the $f_\eta$ are isomorphisms, $f$ is dominant and hence \emph{surjective} since it is proper.
We must therefore show that $f$ is an \emph{open immersion}.
In view of \sref{IV.17.9.5}, it suffices to prove that $f_s:\mathrm{X}_s\to\mathrm{Y}_s$ is an open immersion for every $s\in\mathrm{S}$.

Every $s\in\mathrm{S}$ is a specialization of a maximal point $\eta$.  Making the base change $\operatorname{Spec}(\sh{O}_{\mathrm{S}',s})\to\mathrm{S}$, where $\mathrm{S}'$ is the reduced sub-prescheme of $\mathrm{S}$ with underlying space $\overline{\{\eta\}}$, we may first reduce to the case where $\mathrm{S}$ is a local integral scheme with closed point $s$ and generic point $\eta$; the fibres at $s$ and $\eta$ and the relevant properties of $g$, $h$, and $f$ are unchanged by this base change.
The question is local on $\mathrm{Y}$.  Replacing $\mathrm{Y}$ by an affine open neighbourhood of a point of $\mathrm{Y}_s$, and $\mathrm{X}$ by its inverse image (which is quasi-compact because $f$ is proper), we may also suppose that $\mathrm{X}$ and $\mathrm{Y}$ are of finite presentation over $\mathrm{S}$.
There then exist a local Noetherian subring $\mathrm{A}_0$ of $\mathrm{A}=\sh{O}_{\mathrm{S},s}$, with $\mathrm{A}_0\to\mathrm{A}$ local, finite-type morphisms
\[
  g_0:\mathrm{X}_0\to\mathrm{S}_0=\operatorname{Spec}(\mathrm{A}_0),
  \qquad h_0:\mathrm{Y}_0\to\mathrm{S}_0,
\]
and an $\mathrm{S}_0$-morphism $f_0:\mathrm{X}_0\to\mathrm{Y}_0$ from which $g$, $h$, and $f$ are obtained by base change $(\mathbf{IV}$, $8.9.1$ and $5.13.3)$.  We may moreover suppose that $g_0$ is flat $(\mathbf{IV}$, $11.2.7)$, $h_0$ smooth $(\mathbf{IV}$, $17.7.9)$, and $f_0$ proper $(\mathbf{IV}$, $8.10.5$,~(xii)).
The image in $\mathrm{S}_0$ of the generic point $\eta$ of $\mathrm{S}$ is the generic point $\eta_0$ of $\mathrm{S}_0$, and $(f_0)_{\eta_0}:(\mathrm{X}_0)_{\eta_0}\to(\mathrm{Y}_0)_{\eta_0}$ is an isomorphism $(\mathbf{IV}$, $2.7.1$,~(viii)).  The closed point $s$ of $\mathrm{S}$ lies above the closed point $s_0$ of $\mathrm{S}_0$, and $(\mathrm{X}_0)_{s_0}$ is geometrically irreducible $(\mathbf{IV}$, $4.5.6)$.
Thus we may reduce to the case where $\mathrm{S}$ is the spectrum of a Noetherian local integral ring, while weakening the fibre hypothesis to require only that $\mathrm{X}_s$ and $\mathrm{X}_\eta$ be geometrically irreducible; under these hypotheses we must prove that $f_s$ is an open immersion.

There exists a discrete valuation ring $\mathrm{A}'$ and a morphism $\mathrm{S}'=\operatorname{Spec}(\mathrm{A}')\to\mathrm{S}$ carrying the closed point $a$ of $\mathrm{S}'$ to $s$ and the generic point $b$ to $\eta$ $(\mathbf{II}$, $7.1.9)$.  Let $g'$, $h'$, and $f'$ be obtained from $g$, $h$, and $f$ by this base change.  They again satisfy the hypotheses of \sref{IV.21.12.15}; if $f'_a:\mathrm{X}'_a\to\mathrm{Y}'_a$ is an open immersion, then so is $f_s:\mathrm{X}_s\to\mathrm{Y}_s$ by $(\mathbf{IV}$, $2.7.1$,~(x)).  We may therefore finally suppose that $\mathrm{S}$ is the spectrum of a discrete valuation ring.

Since $h:\mathrm{Y}\to\mathrm{S}$ is smooth, $\mathrm{Y}$ is regular $(\mathbf{IV}$, $17.5.8)$; as the question is local on $\mathrm{Y}$, we may suppose that $\mathrm{Y}$ is integral.
Because $g:\mathrm{X}\to\mathrm{S}$ is flat, the maximal points of $\mathrm{X}$ lie in $\mathrm{X}_\eta$ $(\mathbf{IV}$, $2.3.4)$; since $f_\eta$ is an isomorphism, $\mathrm{X}$ is irreducible and the local ring at its generic point is a field.  Moreover, flatness $(\mathbf{IV}$, $3.3.2)$ shows that $\mathrm{X}$ has no embedded prime cycles, so $\mathrm{X}$ is reduced $(\mathbf{IV}$, $3.2.1)$ and hence integral, while $f$ is birational and separated.
We may therefore apply \sref{IV.21.12.13}; it remains only to check that $f$ is quasi-finite at the points of $\mathrm{X}^{(1)}$.  This is clear for those lying in $\mathrm{X}_\eta$.  By \sref{IV.6.1.1}, the only point of $\mathrm{X}^{(1)}$ lying in $\mathrm{X}_s$ is the generic point $x$ of $\mathrm{X}_s$.
By \sref{IV.2.4.6} and \sref{IV.14.2.2}, $g$ and $h$ are equidimensional; since $\mathrm{X}_\eta$ and $\mathrm{Y}_\eta$ are isomorphic, the irreducible components of $\mathrm{X}_s$ and $\mathrm{Y}_s$ all have the same dimension.  The scheme $\mathrm{X}_s$ is irreducible by hypothesis, and $f_s$ is surjective because $f$ is; hence $\mathrm{Y}_s$ is also irreducible.  If $y$ is its generic point, then $f(x)=y$, and $\dim(\mathrm{X}_s)=\dim(\mathrm{Y}_s)$ gives $\dim(f_s^{-1}(y))=0$ by \sref{IV.5.6.6}.  Thus $f_s$, and hence $f$, is quasi-finite at $x$, completing the proof.
\end{proof}

\begin{corollary}[21.12.16]
\label{IV.21.12.16}
Let $g:\mathrm{X}\to\mathrm{S}$ be a proper, flat morphism of finite presentation, $h:\mathrm{Y}\to\mathrm{S}$ a smooth, proper morphism of finite presentation, $f:\mathrm{X}\to\mathrm{Y}$ an S-morphism.
Suppose that the fibres $\mathrm{X}_s=g^{-1}(s)$ of $g$ are geometrically irreducible.
Let $\mathrm{U}$ be the set of $s\in\mathrm{S}$ such that $f_s:\mathrm{X}_s\to\mathrm{Y}_s$ is an isomorphism.
Then $\mathrm{U}$ is open and closed in $\mathrm{S}$, and the morphism $g^{-1}(\mathrm{U})\to h^{-1}(\mathrm{U})$, restriction of $f$, is an isomorphism.
\end{corollary}

\begin{proof}
The last assertion follows from the first and from \sref{IV.17.9.5}.
We already know \sref{IV.9.6.1},~(xi) that $\mathrm{U}$ is locally constructible in $\mathrm{S}$.
By virtue of \sref{IV.1.10.1}, it suffices to prove that $\mathrm{U}$ is stable by \emph{specialisation} and by \emph{generisation}.
To prove these assertions, one can, by a change of base $\operatorname{Spec}(\sh{O}_{\mathrm{S},s})\to\mathrm{S}$, reduce to the case where $\mathrm{S}$ is a local integral scheme of closed point $s$ and generic point $\eta$.

\oldpage[IV-4]{313}
Stability under specialisation follows from \sref{IV.21.12.15}.

To prove stability under generisation, we argue as in the proof of \sref{IV.21.12.15} (noting, with the same notation, that the image of the closed point of $\mathrm{S}$ is the closed point of $\mathrm{S}_0$) and reduce to the case where $\mathrm{S}$ is moreover \emph{Noetherian}; the conclusion then follows from $(\mathbf{III}$, $4.6.7$,~(ii)).
\end{proof}

\begin{remarks}[21.12.17]
\label{IV.21.12.17}
(i)
The conclusion of the proposition \sref{IV.21.12.15} is no longer valid if we eliminate the hypothesis that the fibres $\mathrm{X}_s$ are irreducible.
Take indeed $\mathrm{S}=\operatorname{Spec}(\mathrm{A})$, where $\mathrm{A}$ is a discrete valuation ring, $\mathrm{Y}=\mathbb{P}^1_{\mathrm{S}}$, which is proper and smooth over $\mathrm{S}$ \sref{IV.17.3.9}.
Let $s$ and $\eta$ denote the closed and generic points of $\mathrm{S}$; let $z$ be a closed point of $\mathrm{Y}_s$, for example a $k$-rational point where $k=\boldsymbol{k}(s)$, and let $\mathrm{X}$ be the $\mathrm{Y}$-prescheme obtained by blowing up $z$.
Since $\mathrm{A}[\mathrm{T}]$ is regular of dimension $2$, the argument of the proof of \sref{IV.15.1.1.6} shows that, if $f:\mathrm{X}\to\mathrm{Y}$ is the structural morphism, $f^{-1}(z)$ is isomorphic to $\mathbb{P}^1_k$; moreover, the complement of $f^{-1}(z)$ in $\mathrm{X}_s$ is isomorphic to $\mathrm{Y}_s-\{z\}$.  Thus $\mathrm{Z}_1=f^{-1}(z)$ and the closure $\mathrm{Z}_2$ in $\mathrm{X}_s$ of its complement are the two irreducible components of $\mathrm{X}_s$.
The morphism $f$ is plainly proper, and $g=h\circ f$ is flat: $\mathrm{X}$ is integral $(\mathbf{II}$, $8.1.4)$ and, for every affine open $\mathrm{U}$ of $\mathrm{X}$, the homomorphism $\mathrm{A}\to\Gamma(\mathrm{U},\sh{O}_{\mathrm{X}})$ is injective $(\mathbf{I}$, $1.2.7)$, so the latter is a torsion-free A-module and hence flat because $\mathrm{A}$ is a discrete valuation ring $(\mathbf{0}_{\mathrm{I}}$, $6.3.4)$.
It is clear that $f_\eta:\mathrm{X}_\eta\to\mathrm{Y}_\eta$ is an isomorphism, whereas $f_s$ is not.

(ii)
The conclusion of \sref{IV.21.12.15} also fails if, under the hypotheses, we suppress the hypothesis that $f$ is \emph{proper}.
Take $\mathrm{S}$ and $\mathrm{Y}$ as in (i), but replace $\mathrm{X}$ by the complement $\mathrm{X}'$ of $\mathrm{Z}_2$ in the prescheme constructed there, and replace $f$ by its restriction $f':\mathrm{X}'\to\mathrm{Y}$.  The restriction $g':\mathrm{X}'\to\mathrm{S}$ of $g$ is still flat, and now $\mathrm{X}'_s$ is geometrically irreducible.  Moreover, $\mathrm{X}'_s$ is an affine scheme isomorphic to the complement of a closed point in $\mathbb{P}^1_k$, hence to $\mathbb{A}^1_k$; since its image under $f'_s$ is reduced to the closed point $z$, $f'$ is not proper and $f'_s$ is not an isomorphism, although $f'_\eta$ is.

(iii)
It is possible that the statement of the proposition \sref{IV.21.12.15} remains valid when we replace the word ``isomorphism'' by ``\'etale morphism'' (cf.\ \sref{IV.21.12.14},~(v)).
It will then also be the same of \sref{IV.21.12.16}.
\end{remarks}

\subsection{Parafactorial pairs. Parafactorial local rings}
\label{subsection:IV.21.13}
\label{IV.21.13}

\begin{definition}[21.13.1]
\label{IV.21.13.1}
Let $\mathrm{X}$ be a ringed space, $\mathrm{Y}$ a closed subset of $\mathrm{X}$, $\mathrm{U}=\mathrm{X}-\mathrm{Y}$.
We say that the pair $(\mathrm{X},\mathrm{Y})$ is \emph{parafactorial} if, for every open $\mathrm{V}$ of $\mathrm{X}$, the restriction functor $\mathscr{L}\mapsto\mathscr{L}|(\mathrm{U}\cap\mathrm{V})$, from the category of $\sh{O}_{\mathrm{V}}$-Modules invertibles into that of $\sh{O}_{\mathrm{U}\cap\mathrm{V}}$-Modules invertibles, is an \emph{equivalence of categories}.
\end{definition}

To say that the pair $(\mathrm{X},\mathrm{Y})$ is parafactorial means therefore that, for every open $\mathrm{V}$ of $\mathrm{X}$, the following conditions are satisfied:
\begin{enumerate}[label=\arabic*\textsuperscript{o}]
  \item the functor $\mathscr{L}\mapsto\mathscr{L}|(\mathrm{U}\cap\mathrm{V})$ is \emph{fully faithful}, in other words, for two $\sh{O}_{\mathrm{V}}$-Modules invertibles $\mathscr{L}$, $\mathscr{L}'$, the map of restriction
  \[
    \operatorname{Hom}_{\sh{O}_{\mathrm{V}}}(\mathscr{L},\mathscr{L}')\to\operatorname{Hom}_{\sh{O}_{\mathrm{U}\cap\mathrm{V}}}(\mathscr{L}|(\mathrm{U}\cap\mathrm{V}),\mathscr{L}'|(\mathrm{U}\cap\mathrm{V}))
  \]
  is \emph{bijective};
  \item the functor $\mathscr{L}\mapsto\mathscr{L}|(\mathrm{U}\cap\mathrm{V})$ is \emph{essentially surjective}, in other words, for every invertible $\sh{O}_{\mathrm{U}\cap\mathrm{V}}$-module $\mathscr{L}_0$, there exists an invertible $\sh{O}_{\mathrm{V}}$-module $\mathscr{L}$ such that $\mathscr{L}_0$ is
\oldpage[IV-4]{314}
isomorphic to $\mathscr{L}|(\mathrm{U}\cap\mathrm{V})$; this is again expressed by saying that the canonical homomorphism \sref{IV.21.3.2.4}
  \[
    \operatorname{Pic}(\mathrm{V})\to\operatorname{Pic}(\mathrm{U}\cap\mathrm{V})
  \]
  is \emph{surjective}.
\end{enumerate}

\begin{lemma}[21.13.2]
\label{IV.21.13.2}
Let $f:\mathrm{X}'\to\mathrm{X}$ be a morphism of ringed spaces; for every open $\mathrm{V}$ of $\mathrm{X}$, denote $f_{\mathrm{V}}:f^{-1}(\mathrm{V})\to\mathrm{V}$ the restriction of $f$.
The following conditions are equivalent:
\begin{enumerate}[label=\emph{\alph*)}]
  \item For every open $\mathrm{V}$ of $\mathrm{X}$, the functor $\mathscr{E}\mapsto f_{\mathrm{V}}^*(\mathscr{E})$ from the category of $\sh{O}_{\mathrm{V}}$-Modules locally free of finite rank, into the category of $\sh{O}_{f^{-1}(\mathrm{V})}$-Modules locally free of finite rank, is faithful (resp.\ fully faithful).
  \item[\emph{a')}] For every open $\mathrm{V}$ of $\mathrm{X}$, the functor $\mathscr{L}\mapsto f_{\mathrm{V}}^*(\mathscr{L})$ from the category of $\sh{O}_{\mathrm{V}}$-Modules locally free of rank $1$ into the category of $\sh{O}_{f^{-1}(\mathrm{V})}$-Modules locally free of rank $1$ is faithful (resp.\ fully faithful).
  \item For every open $\mathrm{V}$ of $\mathrm{X}$, the canonical homomorphism $(\mathbf{0}_{\mathrm{I}}$, $4.4.3.2)$ $\mathscr{E}\to(f_{\mathrm{V}})_*(f_{\mathrm{V}}^*(\mathscr{E}))$ is injective (resp.\ an isomorphism) for every $\sh{O}_{\mathrm{V}}$-Module locally free of finite rank $\mathscr{E}$.
  \item[\emph{b')}] The canonical homomorphism $\sh{O}_{\mathrm{X}}\to f_*(\sh{O}_{\mathrm{X}'})$ is a monomorphism (resp.\ an isomorphism).
\end{enumerate}

Suppose that the canonical homomorphism $\sh{O}_{\mathrm{X}}\to f_*(\sh{O}_{\mathrm{X}'})$ is bijective.
Then, for every $\sh{O}_{\mathrm{X}'}$-Module locally free of finite rank $\mathscr{F}$, $\mathscr{F}$ is of the form $f^*(\mathscr{E})$, where $\mathscr{E}$ is an $\sh{O}_{\mathrm{X}}$-Module locally free of finite rank, if and only if the two following conditions are satisfied:
\begin{enumerate}[label=(\roman*)]
  \item $f_*(\mathscr{F})$ is an $\sh{O}_{\mathrm{X}}$-Module locally free;
  \item the canonical homomorphism $(\mathbf{0}_{\mathrm{I}}$, $4.4.3.3)$ $f^*(f_*(\mathscr{F}))\to\mathscr{F}$ is an isomorphism.
\end{enumerate}
When these two conditions are satisfied, $\mathscr{E}$ is isomorphic to $f_*(\mathscr{F})$.
\end{lemma}

\begin{proof}
\oldpage[IV-4]{315}
To say that the functor $\mathscr{E}\mapsto f_{\mathrm{V}}^*(\mathscr{E})$ is faithful (resp.\ fully faithful) means that, for two $\sh{O}_{\mathrm{V}}$-Modules $\mathscr{E}_1$, $\mathscr{E}_2$ locally free of finite rank, the map
\[
  \operatorname{Hom}_{\sh{O}_{\mathrm{V}}}(\mathscr{E}_1,\mathscr{E}_2)
  \longrightarrow
  \operatorname{Hom}_{\sh{O}_{f^{-1}(\mathrm{V})}}(f_{\mathrm{V}}^*(\mathscr{E}_1),f_{\mathrm{V}}^*(\mathscr{E}_2))
\]
is injective (resp.\ bijective). Since this must hold after replacing $\mathrm{V}$ by an open $\mathrm{W}\subset\mathrm{V}$ and $\mathscr{E}_1$, $\mathscr{E}_2$ by $\mathscr{E}_1|\mathrm{W}$, $\mathscr{E}_2|\mathrm{W}$, and since
\[
  \operatorname{Hom}_{\sh{O}_{\mathrm{W}}}(\mathscr{E}_1|\mathrm{W},\mathscr{E}_2|\mathrm{W})
  =\Gamma(\mathrm{W},\shHom_{\sh{O}_{\mathrm{V}}}(\mathscr{E}_1,\mathscr{E}_2)),
\]
the condition also means that the canonical homomorphism of sheaves
\[
\label{IV.21.13.2.1}
  \shHom_{\sh{O}_{\mathrm{V}}}(\mathscr{E}_1,\mathscr{E}_2)\to(f_{\mathrm{V}})_*(f_{\mathrm{V}}^*(\shHom_{\sh{O}_{\mathrm{V}}}(\mathscr{E}_1,\mathscr{E}_2)))
  \tag{21.13.2.1}
\]
is injective (resp.\ bijective).
But since $\mathscr{E}_1$ and $\mathscr{E}_2$ are locally free of finite rank, $\shHom_{\sh{O}_{\mathrm{V}}}(\mathscr{E}_1,\mathscr{E}_2)$, isomorphic to $\mathscr{E}_1^\vee\otimes_{\sh{O}_{\mathrm{V}}}\mathscr{E}_2$ $(\mathbf{0}_{\mathrm{I}}$, $5.4.2)$, is also locally free of finite rank. This already shows that $b)$ implies $a)$; conversely, $b)$ is a particular case of $a)$, in view of the isomorphism $\mathscr{E}\xrightarrow{\sim}\shHom_{\sh{O}_{\mathrm{V}}}(\sh{O}_{\mathrm{V}},\mathscr{E})$.
It is clear that $a')$ is a particular case of $a)$ and that $b')$ is a particular case of $a')$.
Conversely, since $b)$ is a local property, one may, in order to verify it, restrict to the case $\mathscr{E}=\sh{O}_{\mathrm{V}}^n$; this shows that $b')$ implies $b)$.

If the canonical homomorphism $\sh{O}_{\mathrm{X}}\to f_*(\sh{O}_{\mathrm{X}'})$ is bijective, and if conditions (i) and (ii) are satisfied, it is clear that $\mathscr{F}=f^*(\mathscr{E})$ with $\mathscr{E}=f_*(\mathscr{F})$, up to canonical isomorphism.
Conversely, suppose that $\mathscr{F}=f^*(\mathscr{E})$, with $\mathscr{E}$ locally free. Since the question is local on $\mathrm{X}$, we may suppose that $\mathscr{E}=\sh{O}_{\mathrm{X}}^n$, whence $\mathscr{F}=\sh{O}_{\mathrm{X}'}^n$; conditions (i) and (ii) then follow from the hypothesis that $\sh{O}_{\mathrm{X}}\to f_*(\sh{O}_{\mathrm{X}'})$ is an isomorphism.
\end{proof}

In this paragraph, we will apply the preceding lemma to a canonical injection $j:\mathrm{U}\to\mathrm{X}$, where $\mathrm{U}$ is the ringed space induced by $\mathrm{X}$ on an \emph{open} of $\mathrm{X}$.
With these notations, \sref{IV.21.13.2} translates as:

\begin{corollary}[21.13.3]
\label{IV.21.13.3}
Let $\mathrm{X}$ be a ringed space, $\mathrm{Y}$ a closed subset of $\mathrm{X}$, $\mathrm{U}=\mathrm{X}-\mathrm{Y}$, $j:\mathrm{U}\to\mathrm{X}$ the canonical injection.
The following conditions are equivalent:
\begin{enumerate}[label=\emph{\alph*)}]
  \item For every open $\mathrm{V}$ of $\mathrm{X}$, the functor restriction $\mathscr{E}\mapsto\mathscr{E}|(\mathrm{U}\cap\mathrm{V})$ from the category of $\sh{O}_{\mathrm{V}}$-Modules locally free of finite rank into the category of $\sh{O}_{\mathrm{U}\cap\mathrm{V}}$-Modules locally free of finite rank is faithful (resp.\ fully faithful).
  \item[\emph{a')}] For every open $\mathrm{V}$ of $\mathrm{X}$, the restriction functor $\mathscr{L}\mapsto\mathscr{L}|(\mathrm{U}\cap\mathrm{V})$ from the category of $\sh{O}_{\mathrm{V}}$-Modules locally free of rank $1$ into the category of $\sh{O}_{\mathrm{U}\cap\mathrm{V}}$-Modules locally free of rank $1$ is faithful (resp.\ fully faithful).
  \item For every open $\mathrm{V}$ of $\mathrm{X}$, the canonical homomorphism $\mathscr{E}\to(j_{\mathrm{V}})_*(\mathscr{E}|(\mathrm{U}\cap\mathrm{V}))$ is injective (resp.\ bijective) for every $\sh{O}_{\mathrm{V}}$-Module locally free of finite rank $\mathscr{E}$.
  \item[\emph{b')}] The canonical homomorphism $\sh{O}_{\mathrm{X}}\to j_*(\sh{O}_{\mathrm{U}})$ is injective (resp.\ bijective).
\end{enumerate}

Suppose that the canonical homomorphism $\sh{O}_{\mathrm{X}}\to j_*(\sh{O}_{\mathrm{U}})$ is bijective.
Then, for an $\sh{O}_{\mathrm{U}}$-Module locally free of finite rank $\mathscr{F}$ to be of the form $\mathscr{E}|\mathrm{U}$, where $\mathscr{E}$ is an $\sh{O}_{\mathrm{X}}$-Module locally free of finite rank, it is necessary and sufficient that $j_*(\mathscr{F})$ is an $\sh{O}_{\mathrm{X}}$-Module locally free of finite rank, and in this case, we can take $\mathscr{E}=j_*(\mathscr{F})$.
\end{corollary}

\begin{lemma}[21.13.4]
\label{IV.21.13.4}
Let $\mathrm{X}$ be a locally Noetherian prescheme, $\mathrm{Y}$ a closed subset of $\mathrm{X}$, $\mathrm{U}=\mathrm{X}-\mathrm{Y}$, and $j:\mathrm{U}\to\mathrm{X}$ the canonical injection.
For the canonical homomorphism $\sh{O}_{\mathrm{X}}\to j_*(\sh{O}_{\mathrm{U}})$ to be injective (resp.\ bijective), it is necessary and sufficient that $\operatorname{prof}_{\mathrm{Y}}(\sh{O}_{\mathrm{X}})\geq 1$ (resp.\ $\operatorname{prof}_{\mathrm{Y}}(\sh{O}_{\mathrm{X}})\geq 2$) (in other words, that $\operatorname{prof}(\sh{O}_{\mathrm{X},x})\geq 1$ (resp.\ $\operatorname{prof}(\sh{O}_{\mathrm{X},x})\geq 2$) for every $x\in\mathrm{Y}$).
\end{lemma}

\begin{proof}
These assertions are indeed particular cases of \sref{IV.5.10.2} and \sref{IV.5.10.5}.
\end{proof}

\begin{proposition}[21.13.5]
\label{IV.21.13.5}
Let $\mathrm{X}$ be a ringed space, $\mathrm{Y}$ a closed subset of $\mathrm{X}$, $\mathrm{U}=\mathrm{X}-\mathrm{Y}$, $j:\mathrm{U}\to\mathrm{X}$ the canonical injection.
For the pair $(\mathrm{X},\mathrm{Y})$ to be parafactorial, it is necessary and sufficient that the canonical homomorphism $\sh{O}_{\mathrm{X}}\to j_*(\sh{O}_{\mathrm{U}})$ is bijective, and that for every open $\mathrm{V}$ of $\mathrm{X}$ and every invertible $\sh{O}_{\mathrm{U}\cap\mathrm{V}}$-module $\mathscr{L}$, $(j_{\mathrm{V}})_*(\mathscr{L})$ is an invertible $\sh{O}_{\mathrm{V}}$-module (notation of \sref{IV.21.13.3}).
\end{proposition}

\begin{proof}
This is an immediate consequence of the definition \sref{IV.21.13.1} and of \sref{IV.21.13.3}.
\end{proof}

\begin{corollary}[21.13.6]
\label{IV.21.13.6}
Let $\mathrm{X}$ be a ringed space, $\mathrm{Y}$ a closed subset of $\mathrm{X}$.
\begin{enumerate}[label=(\roman*)]
  \item If the pair $(\mathrm{X},\mathrm{Y})$ is parafactorial, then so is $(\mathrm{W},\mathrm{Y}\cap\mathrm{W})$ for every open $\mathrm{W}$ of $\mathrm{X}$.
  Conversely, if $(\mathrm{W}_\alpha)$ is an open covering of $\mathrm{X}$ such that each of the pairs $(\mathrm{W}_\alpha,\mathrm{Y}\cap\mathrm{W}_\alpha)$ is parafactorial, then the pair $(\mathrm{X},\mathrm{Y})$ is parafactorial.
  \item If the pair $(\mathrm{X},\mathrm{Y})$ is parafactorial, then so is $(\mathrm{X},\mathrm{Y}')$ for every closed subset $\mathrm{Y}'$ of $\mathrm{Y}$.
  \item Suppose that $\mathrm{X}$ is a prescheme, and let $\mathrm{X}'$ be a prescheme, $f:\mathrm{X}'\to\mathrm{X}$ a faithfully flat and quasi-compact morphism; set $\mathrm{Y}'=f^{-1}(\mathrm{Y})$.
  Suppose that the pair $(\mathrm{X}',\mathrm{Y}')$ is parafactorial and that the open $\mathrm{U}=\mathrm{X}-\mathrm{Y}$ is retrocompact in $\mathrm{X}$ $(\mathbf{0}_{\mathrm{III}}$, $9.1.1)$.
  Then the pair $(\mathrm{X},\mathrm{Y})$ is parafactorial.
\end{enumerate}
\end{corollary}

\begin{proof}
\oldpage[IV-4]{316}
(i) Since the bijectivity of the canonical homomorphism $\sh{O}_{\mathrm{X}}\to j_*(\sh{O}_{\mathrm{U}})$ is a local property on $\mathrm{X}$, it remains only to see, by virtue of \sref{IV.21.13.5}, that the following property holds. Denoting by $j_\alpha:\mathrm{U}\cap\mathrm{W}_\alpha\to\mathrm{W}_\alpha$ the canonical injection, if, for every $\alpha$ and every invertible $\sh{O}_{\mathrm{U}}$-Module $\mathscr{L}$, $(j_\alpha)_*(\mathscr{L}|(\mathrm{U}\cap\mathrm{W}_\alpha))$ is an invertible $\sh{O}_{\mathrm{W}_\alpha}$-Module, then $j_*(\mathscr{L})$ is an invertible $\sh{O}_{\mathrm{X}}$-Module.
But being an invertible $\sh{O}_{\mathrm{X}}$-Module is local on $\mathrm{X}$, and $(\mathrm{W}_\alpha)$ is an open covering of $\mathrm{X}$; since
\[
  j_*(\mathscr{L})|\mathrm{W}_\alpha=(j_\alpha)_*(\mathscr{L}|(\mathrm{U}\cap\mathrm{W}_\alpha)),
\]
this proves the assertion.

(ii) Set $\mathrm{U}'=\mathrm{X}-\mathrm{Y}'$ and let $j':\mathrm{U}'\to\mathrm{X}$ be the canonical injection.
To say that the canonical homomorphism $\sh{O}_{\mathrm{X}}\to j'_*(\sh{O}_{\mathrm{U}'})$ is bijective amounts to saying that, for every open $\mathrm{V}$ of $\mathrm{X}$, the homomorphism $\Gamma(\mathrm{V},\sh{O}_{\mathrm{X}})\to\Gamma(\mathrm{V}\cap\mathrm{U}',\sh{O}_{\mathrm{X}})$ is bijective. But the composite
\[
  \Gamma(\mathrm{V},\sh{O}_{\mathrm{X}})\to\Gamma(\mathrm{V}\cap\mathrm{U}',\sh{O}_{\mathrm{X}})\to\Gamma(\mathrm{V}\cap\mathrm{U},\sh{O}_{\mathrm{X}})
\]
is bijective by hypothesis, and, on replacing $\mathrm{V}$ by $\mathrm{V}\cap\mathrm{U}'$, one sees that the second arrow is bijective; hence the first arrow is bijective.

Next let $j'':\mathrm{U}\to\mathrm{U}'$ be the canonical injection, and let $\mathscr{L}'$ be an invertible $\sh{O}_{\mathrm{U}'}$-Module. Then $\mathscr{L}'|\mathrm{U}$ is an invertible $\sh{O}_{\mathrm{U}}$-Module, so by hypothesis
\[
  j_*(\mathscr{L}'|\mathrm{U})=j'_*(j''_*(\mathscr{L}'|\mathrm{U}))
\]
is an invertible $\sh{O}_{\mathrm{X}}$-Module. Since the pair $(\mathrm{U}',\mathrm{U}'\cap\mathrm{Y})$ is parafactorial by (i), one has $j''_*(\mathscr{L}'|\mathrm{U})\cong\mathscr{L}'$, and therefore $j'_*(\mathscr{L}')$ is an invertible $\sh{O}_{\mathrm{X}}$-Module. To conclude, it suffices to replace $\mathrm{X}$, $\mathrm{U}$, and $\mathrm{U}'$ in the preceding argument by $\mathrm{V}$, $\mathrm{V}\cap\mathrm{U}$, and $\mathrm{V}\cap\mathrm{U}'$, where $\mathrm{V}$ is an open subset of $\mathrm{X}$.

(iii) Set $\mathrm{U}'=\mathrm{X}'-\mathrm{Y}'=f^{-1}(\mathrm{U})$ and note that, since we are dealing with preschemes, one may write $\mathrm{U}'=\mathrm{U}\times_{\mathrm{X}}\mathrm{X}'$. Let $f_{\mathrm{U}}:f^{-1}(\mathrm{U})\to\mathrm{U}$ be the restriction of $f$, and let $j':\mathrm{U}'\to\mathrm{X}'$ be the canonical injection, also denoted $j_{(\mathrm{X}')}$.
Let us first show that the canonical homomorphism $\rho:\sh{O}_{\mathrm{X}}\to j_*(\sh{O}_{\mathrm{U}})$ is bijective. By hypothesis, the canonical homomorphism $\sh{O}_{\mathrm{X}'}\to j'_*(\sh{O}_{\mathrm{U}'})$ is bijective. Since $j$ is quasi-compact and separated by hypothesis, and $f$ is flat, $j'_*(\sh{O}_{\mathrm{U}'})=j'_*(f_{\mathrm{U}}^*(\sh{O}_{\mathrm{U}}))$ identifies canonically with $f^*(j_*(\sh{O}_{\mathrm{U}}))$ by \sref{IV.2.3.1}. Since $\sh{O}_{\mathrm{X}'}=f^*(\sh{O}_{\mathrm{X}})$, the homomorphism $f^*(\rho):f^*(\sh{O}_{\mathrm{X}})\to f^*(j_*(\sh{O}_{\mathrm{U}}))$ is therefore bijective; hence $\rho$ is bijective because $f$ is faithfully flat \sref{IV.2.2.7}.

Now let $\mathscr{L}$ be an invertible $\sh{O}_{\mathrm{U}}$-Module, and set $\mathscr{L}'=f_{\mathrm{U}}^*(\mathscr{L})$, which is an invertible $\sh{O}_{\mathrm{U}'}$-Module. By hypothesis, $j'_*(\mathscr{L}')=j'_*(f_{\mathrm{U}}^*(\mathscr{L}))$ is an invertible $\sh{O}_{\mathrm{X}'}$-Module; as above, it is isomorphic by \sref{IV.2.3.1} to $f^*(j_*(\mathscr{L}))$. Since $f$ is faithfully flat and quasi-compact, it follows that $j_*(\mathscr{L})$ is a locally free, hence invertible, $\sh{O}_{\mathrm{X}}$-Module \sref{IV.2.5.2}. To complete the proof, it suffices to replace $\mathrm{X}$ by an open $\mathrm{V}$ and $\mathrm{X}'$ by $f^{-1}(\mathrm{V})$ in the preceding argument.
\end{proof}

\begin{definition}[21.13.7]
\label{IV.21.13.7}
We say that a \emph{local ring} $\mathrm{A}$ is parafactorial if, setting $\mathrm{X}=\operatorname{Spec}(\mathrm{A})$ and denoting by $a$ the closed point of $\mathrm{X}$, the pair $(\mathrm{X},\{a\})$ is parafactorial.
\end{definition}

\begin{proposition}[21.13.8]
\label{IV.21.13.8}
The notations being those of \sref{IV.21.13.7}, set $\mathrm{U}=\mathrm{X}-\{a\}$.
For $\mathrm{A}$ to be parafactorial, it is necessary and sufficient that it satisfies the following conditions:
\begin{enumerate}[label=(\roman*)]
  \item The canonical homomorphism $\mathrm{A}=\Gamma(\mathrm{X},\sh{O}_{\mathrm{X}})\to\Gamma(\mathrm{U},\sh{O}_{\mathrm{X}})$ is bijective.
  \item $\operatorname{Pic}(\mathrm{U})=0$.
\end{enumerate}

\oldpage[IV-4]{317}
If furthermore $\mathrm{A}$ is Noetherian, condition (i) is equivalent to
\begin{enumerate}
  \item[\textup{(i bis)}] $\operatorname{prof}(\mathrm{A})\geq 2$.
\end{enumerate}
\end{proposition}

\begin{proof}
Indeed, the only open of $\mathrm{X}$ containing $a$ is $\mathrm{X}$ itself, and hence every invertible $\sh{O}_{\mathrm{X}}$-module is isomorphic to $\sh{O}_{\mathrm{X}}$, i.e.\ $\operatorname{Pic}(\mathrm{X})=0$; the first assertion follows therefore from the definition \sref{IV.21.13.1} and from \sref{IV.21.13.3}.
The second assertion is nothing but a particular case of \sref{IV.21.13.4}.
\end{proof}

\begin{env}[21.13.9]
\label{IV.21.13.9}
\emph{Examples}. ---
(i)
A local Noetherian parafactorial ring is necessarily of dimension $\geq 2$ by virtue of \sref{IV.21.13.8}; in other words a local Noetherian ring of dimension $\leq 1$ is not parafactorial.

(ii)
A local Noetherian \emph{factorial} ring $\mathrm{A}$ of dimension $\geq 2$ is parafactorial, as follows from \sref{IV.21.13.8} and \sref{IV.21.6.14}.

(iii)
If $\mathrm{A}$ is a local Noetherian ring of dimension $\geq 3$ and parafactorial, it is not necessarily normal nor even reduced.
Consider a regular local ring $\mathrm{B}$ of dimension $\geq 3$, and let $\mathrm{A}=\mathrm{D}_{\mathrm{B}}(\mathrm{B})$ $(\mathbf{0}$, $18.2.3)$, isomorphic to $\mathrm{B}[\mathrm{T}]/(\mathrm{T}^2)$; we see immediately that $\operatorname{prof}(\mathrm{A})=\operatorname{prof}(\mathrm{B})\geq 3$.
To show that $\mathrm{A}$ is parafactorial, it suffices therefore, with the notations of \sref{IV.21.13.8}, to prove that $\operatorname{Pic}(\mathrm{U})=0$.
Let $\mathfrak{J}$ be the ideal of the augmentation $\mathrm{A}\to\mathrm{B}$, which is such that $\mathfrak{J}^2=0$ and which, as a $\mathrm{B}$-module, is isomorphic to $\mathrm{B}$.
Since $\mathrm{B}=\mathrm{A}/\mathfrak{J}$, $\mathrm{X}=\operatorname{Spec}(\mathrm{A})$ and $\mathrm{X}_0=\operatorname{Spec}(\mathrm{B})$ have the same underlying space; if $\mathscr{J}=\widetilde{\mathfrak{J}}$, $\mathrm{X}_0$ is the sub-prescheme defined by $\mathscr{J}$; we will denote $\mathrm{U}_0$ the prescheme induced by $\mathrm{X}_0$ on the open $\mathrm{U}$.
For every $z\in\mathfrak{J}$, set $\varphi(z)=1+z$; since $(1+z)(1-z)=1$, $1+z$ is invertible in $\mathrm{A}$, and $\varphi(\mathfrak{J})$ is the kernel of the canonical surjective homomorphism of multiplicative groups $\mathrm{A}^*\to\mathrm{B}^*$; in other words, we have an exact sequence of commutative groups
\[
  0\to\mathfrak{J}\xrightarrow{\varphi}\mathrm{A}^*\to\mathrm{B}^*\to 1
\]
(the three last groups being multiplicative, the two first being additive).
For the same reason, for every $t\in\mathrm{A}$, we have the exact sequence
\[
  0\to\mathfrak{J}_t\xrightarrow{\varphi_t}(\mathrm{A}_t)^*\to(\mathrm{B}_{t_0})^*\to 1
\]
denoting by $t_0$ the image of $t$ in $\mathrm{B}$, since $\mathrm{B}_{t_0}=\mathrm{A}_t/\mathfrak{J}_t$; in other words, we have an exact sequence of sheaves of commutative groups on the topological space $\mathrm{X}$
\[
  0\to\mathscr{J}\xrightarrow{u}\sh{O}_{\mathrm{X}}^*\to\sh{O}_{\mathrm{X}_0}^*\to 1
\]
whence by restriction to the open $\mathrm{U}$, an exact sequence
\[
  0\to\mathscr{J}|\mathrm{U}\to\sh{O}_{\mathrm{U}}^*\to\sh{O}_{\mathrm{U}_0}^*\to 1.
\]

By the cohomology exact sequence, we deduce an exact sequence
\[
\label{IV.21.13.9.1}
  \mathrm{H}^1(\mathrm{U},\mathscr{J})\to\mathrm{H}^1(\mathrm{U},\sh{O}_{\mathrm{U}}^*)\xrightarrow{u_1}\mathrm{H}^1(\mathrm{U},\sh{O}_{\mathrm{U}_0}^*).
  \tag{21.13.9.1}
\]
But since $\mathfrak{J}$ is a $\mathrm{B}$-module isomorphic to $\mathrm{B}$, we have $\mathrm{H}^1(\mathrm{U},\mathscr{J})=\mathrm{H}^1(\mathrm{U}_0,\sh{O}_{\mathrm{U}_0})$, and it follows from chap.~III, Part~3 (see also \cite[III, Example~III-1]{IV-41}) that
\oldpage[IV-4]{318}
$\mathrm{H}^1(\mathrm{U}_0,\sh{O}_{\mathrm{U}_0})=0$ by reason of the relation $\operatorname{prof}(\mathrm{B})\geq 3$.
On the other hand, since $\mathrm{B}$ is factorial, $\mathrm{H}^1(\mathrm{U},\sh{O}_{\mathrm{U}_0}^*)=\operatorname{Pic}(\mathrm{U}_0)=0$; we obtain therefore from the preceding exact sequence $\operatorname{Pic}(\mathrm{U})=\mathrm{H}^1(\mathrm{U},\sh{O}_{\mathrm{U}}^*)=0$.

(iv) There also exist local Noetherian rings of dimension $3$ which are integral, integrally closed and parafactorial, but not factorial.
Let indeed $\mathrm{B}$ be a local Noetherian ring, complete, integral, integrally closed, of dimension $\geq 2$ and not factorial (for example the completion of the localised polynomial algebra $k[\mathrm{U},\mathrm{V},\mathrm{W}]/(\mathrm{W}^2-\mathrm{U}\mathrm{V})$ at the ideal image of $(\mathrm{U})+(\mathrm{V})+(\mathrm{W})$).
Then it follows from what we will see below \sref{IV.21.14.2} that the local ring $\mathrm{A}=\mathrm{B}[[\mathrm{T}]]$ of formal power series in $\mathrm{T}$ over $\mathrm{B}$ is parafactorial, but it is not factorial, without which $\mathrm{B}$ would be so by virtue of \sref{IV.21.13.12} below.

(v) One can show that an absolute complete-intersection local ring \sref{IV.19.3.1} of dimension $\geq 4$ is parafactorial (cf.\ \cite[XI, $3.13$~(i)]{IV-41}).

(vi) We have seen (Remark~(ii)) that every local Noetherian factorial ring $\mathrm{A}$ of dimension $2$ is parafactorial.
But there are local Noetherian rings of dimension $2$ which are parafactorial and not factorial.
One can show that, for a local Noetherian ring $\mathrm{A}$ of dimension $2$ to be parafactorial and not factorial, it is necessary and sufficient that it satisfy the following three conditions:
\begin{enumerate}[label=\arabic*\textsuperscript{o}]
  \item $\mathrm{A}$ is a Cohen-Macaulay ring (in other words $\operatorname{prof}(\mathrm{A})=2$).
  \item $\mathrm{A}$ is integral and if $\mathrm{A}'$ is its integral closure, $\mathrm{A}'$ is factorial and is a finite A-algebra.
  \item Let $\mathfrak{J}$ be the conductor of $\mathrm{A}$ in $\mathrm{A}'$ (the annihilator of the $\mathrm{A}$-module $\mathrm{A}'/\mathrm{A}$, or equivalently the largest ideal of $\mathrm{A}'$ contained in $\mathrm{A}$); set $\mathrm{B}=\mathrm{A}/\mathfrak{J}$, $\mathrm{B}'=\mathrm{A}'/\mathfrak{J}$. Then $\dim(\mathrm{B})=1$ (which implies $\mathrm{A}'\neq\mathrm{A}$, i.e.\ $\mathrm{A}$ is not integrally closed), and the canonical map $\operatorname{Div}(\mathrm{B})\to\operatorname{Div}(\mathrm{B}')$ \sref{IV.21.4.5} is \emph{surjective}.
\end{enumerate}
One can furthermore show that these conditions imply the following property:

$4^{\mathrm{o}}$ The ring $\mathrm{B}'$ (and, a fortiori, $\mathrm{B}\subset\mathrm{B}'$) is reduced, and the morphism $\operatorname{Spec}(\mathrm{B}')\to\operatorname{Spec}(\mathrm{B})$ is bijective.

If we set $\mathrm{X}=\operatorname{Spec}(\mathrm{A})$, $\mathrm{X}'=\operatorname{Spec}(\mathrm{A}')$, $\mathrm{Y}=\operatorname{Spec}(\mathrm{B})$, and $\mathrm{Y}'=\operatorname{Spec}(\mathrm{B}')$, so that $\mathrm{Y}$ (resp.\ $\mathrm{Y}'$) is defined by the ideal $\mathfrak{J}$ of $\mathrm{A}$ (resp.\ $\mathrm{A}'$), the structural morphism $f:\mathrm{X}'\to\mathrm{X}$ is an isomorphism of $\mathrm{X}'-\mathrm{Y}'$ onto $\mathrm{X}-\mathrm{Y}$ (Bourbaki, \emph{Alg. comm.}, chap.~V, \S1, no.~5, Cor.~5 of Prop.~16). By $4^{\mathrm{o}}$, $f$ is a bijective morphism; in other words, $\mathrm{X}$ is a Noetherian local unibranch prescheme \sref{IV.6.15.1} (and in particular $\mathrm{A}'$ is a Noetherian local ring). In general, $\mathrm{X}$ is not geometrically unibranch.

The space $\mathrm{Y}$, of dimension $1$, consists of the closed point $x$ of $\mathrm{X}$ and the maximal points $y_i$ ($1\leq i\leq r$) of $\mathrm{Y}$; the unique point $y_i'$ of $\mathrm{Y}'$ above $y_i$ is likewise a maximal point of $\mathrm{Y}'$. Since $\mathrm{B}$ and $\mathrm{B}'$ are reduced, it follows that
\[
  \mathfrak{m}_{y_i}\sh{O}_{\mathrm{X}',y_i'}=\mathfrak{m}_{\sh{O}_{\mathrm{X}',y_i'}}.
\]
Thus the relation $\mathfrak{m}_z\sh{O}_{\mathrm{X}',z'}=\mathfrak{m}_{z'}$ holds for $z=y_i$ and $z'=y_i'$, and it is plainly also true for $z\notin\mathrm{Y}$ and the unique point $z'$ of $f^{-1}(z)$; hence it holds for every $z\in\mathrm{U}=\mathrm{X}-\{x\}$ and the unique point $z'$ of $f^{-1}(z)$. In particular, if $\mathrm{A}$ has characteristic $0$ \sref{IV.21.1.1}, then $\mathrm{U}'$ is \emph{unramified} over $\mathrm{U}$ (but not \'etale in general).

We shall confine ourselves here to proving that conditions $1^{\mathrm{o}}$, $2^{\mathrm{o}}$, and $3^{\mathrm{o}}$ above are sufficient for $\mathrm{A}$ to be parafactorial. By condition $1^{\mathrm{o}}$ and \sref{IV.21.13.8}, it suffices to show that $\operatorname{Pic}(\mathrm{U})=0$. Since $\mathrm{A}'$ is factorial by $2^{\mathrm{o}}$, one has $\operatorname{Pic}(\mathrm{U}')=0$, and \sref{IV.21.8.5},~(ii) gives an exact sequence
\[
  1\to
  \left(\prod_{i=1}^r \sh{O}_{\mathrm{U}',y_i'}^*/\sh{O}_{\mathrm{U},y_i}^*\right)
  \big/\operatorname{Im}\Gamma(\mathrm{U}',\sh{O}_{\mathrm{U}'}^*)
  \longrightarrow\operatorname{Pic}(\mathrm{U})\longrightarrow\operatorname{Pic}(\mathrm{U}')\to 0.
\]
It therefore remains to show that the second term of this sequence is reduced to the neutral element. Here we take $\mathrm{S}=\{y_i\mid 1\leq i\leq r\}$ and note that $f_*(\sh{O}_{\mathrm{X}'})$ is the $\sh{O}_{\mathrm{X}}$-algebra $\widetilde{\mathrm{A}'}$. Let $\mathfrak{p}_i$, $\mathfrak{p}_i'$ be the prime ideals of $\mathrm{A}$ and $\mathrm{A}'$ corresponding to $y_i$, $y_i'$, and note that $\mathfrak{p}_i'\cap\mathrm{A}=\mathfrak{p}_i$. For each $i$, choose an invertible element $a_i'/s_i$ of $\mathrm{A}'_{\mathfrak{p}_i'}$, with $a_i'\in\mathrm{A}'$ and $s_i\notin\mathfrak{p}_i'$. Since only the quotient groups $(\mathrm{A}'_{\mathfrak{p}_i'})^*/(\mathrm{A}_{\mathfrak{p}_i})^*$ are at issue, we may suppose that $s_i=1$ for every $i$. Let $b_i'\in\mathrm{B}'=\mathrm{A}'/\mathfrak{J}$ be the canonical image of $a_i'$, so that $b_i'/1$ is a nonzero element of $\kappa(y_i')=\sh{O}_{\mathrm{Y}',y_i'}$.

The open $\mathrm{U}\cap\mathrm{Y}$, consisting of the points $y_i$, is an affine open of the form $\mathrm{D}(t)$ in $\mathrm{Y}$, with $t\in\mathfrak{m}$, the maximal ideal of $\mathrm{A}$. There is therefore an invertible element $b'\in\mathrm{B}'_t$ whose canonical images are the $b_i'/1$. Since $t$ is regular in the integral ring $\mathrm{A}'$, we may write $b'=b''/t^m$, where $b''\in\mathrm{B}'$ is invertible. Let $a''\in\mathrm{A}'$ represent the class $b''$; it is necessarily invertible. Then $a'=a''/t^m$ is invertible in $\mathrm{A}'_t$, and for every $i$ one has
\[
  u_i=a''-a_i't^m\in\mathfrak{J},
  \qquad
  t^m a_i'=a''-u_i=a''(1-a''^{-1}u_i).
\]
But $a''^{-1}u_i\in\mathfrak{J}\subset\mathfrak{m}$, so $1-a''^{-1}u_i$ is invertible in $\mathrm{A}$; hence the classes of $a'$ and $a_i'$ in $(\mathrm{A}'_{\mathfrak{p}_i'})^*/(\mathrm{A}_{\mathfrak{p}_i})^*$ are the same, which completes the proof.

\oldpage[IV-4]{319}
For an explicit example of a parafactorial ring of dimension $2$ obtained in this way and not factorial, consider
\[
  \mathrm{E}=\mathbb{R}[[\mathrm{U},\mathrm{V}]]/(\mathrm{U}^2+\mathrm{V}^2),
\]
whose integral closure $\mathrm{E}'$ identifies with $\mathbb{C}[[\mathrm{U}]]$ \sref{IV.6.15.11}. If $\mathrm{A}=\mathrm{E}[[\mathrm{T}]]$, its integral closure is $\mathrm{A}'=\mathrm{E}'[[\mathrm{T}]]$ (Bourbaki, \emph{Alg. comm.}, chap.~V, \S1, no.~4, Prop.~14). One verifies at once that the conductor of $\mathrm{E}$ in $\mathrm{E}'$ is the maximal ideal $\mathfrak{n}$ of $\mathrm{E}$, hence that the conductor $\mathfrak{J}$ of $\mathrm{A}$ in $\mathrm{A}'$ is $\mathfrak{n}[[\mathrm{T}]]$, and one has $\mathrm{B}=\mathrm{A}/\mathfrak{J}\cong\mathbb{R}[[\mathrm{T}]]$, $\mathrm{B}'=\mathrm{A}'/\mathfrak{J}\cong\mathbb{C}[[\mathrm{T}]]$. Conditions $1^{\mathrm{o}}$, $2^{\mathrm{o}}$, and $3^{\mathrm{o}}$ above are therefore satisfied, although $\mathrm{A}$ is not even integrally closed.

One may vary this example: if $k$ is an algebraically closed field, the localisation at the maximal ideal generated by the images of $\mathrm{U}$, $\mathrm{V}$, and $\mathrm{W}$ of
\[
  k[\mathrm{U},\mathrm{V},\mathrm{W}]/(\mathrm{U}^2-\mathrm{W}\mathrm{V}^2)
\]
also satisfies conditions $1^{\mathrm{o}}$, $2^{\mathrm{o}}$, and $3^{\mathrm{o}}$ above.

It may be that these three conditions even imply that each of the rings $\mathrm{B}'/\mathfrak{p}_i'$ is a discrete valuation ring, which would imply that $\mathrm{A}'$ itself is a regular ring.
\end{env}

\begin{proposition}[21.13.10]
\label{IV.21.13.10}
Let $\mathrm{X}$ be a locally Noetherian prescheme, $\mathrm{Y}$ a closed subset of $\mathrm{X}$.
For the pair $(\mathrm{X},\mathrm{Y})$ to be parafactorial, it is necessary and sufficient that, for all $y\in\mathrm{Y}$, the local ring $\sh{O}_{\mathrm{X},y}$ is parafactorial.
\end{proposition}

\begin{proof}
\oldpage[IV-4]{320}
Set $\mathrm{U}=\mathrm{X}-\mathrm{Y}$ and let $j:\mathrm{U}\to\mathrm{X}$ be the canonical injection.

By \sref{IV.21.13.4}, to say that the canonical homomorphism $\sh{O}_{\mathrm{X}}\to j_*(\sh{O}_{\mathrm{U}})$ is bijective amounts to saying that each of the local rings $\sh{O}_{\mathrm{X},y}$, for $y\in\mathrm{Y}$, satisfies condition (i bis) of \sref{IV.21.13.8}.

Let us first show that if the pair $(\mathrm{X},\mathrm{Y})$ is parafactorial, each of the rings $\sh{O}_{\mathrm{X},y}$ ($y\in\mathrm{Y}$) is parafactorial.
Taking into account the preceding remark, it all amounts to seeing that if $\mathrm{T}_y=\operatorname{Spec}(\sh{O}_{\mathrm{X},y})$ and $\mathrm{U}_y=\mathrm{T}_y-\{y\}$, every invertible $\sh{O}_{\mathrm{U}_y}$-module $\mathscr{L}_0$ is isomorphic to $\sh{O}_{\mathrm{U}_y}$.
Now, when $\mathrm{V}$ ranges over the open neighbourhoods of $y$ in $\mathrm{X}$, it follows from \sref{IV.8.2.13} and \sref[I]{I.2.4.2} that the prescheme $\mathrm{U}_y$ is projective limit of the preschemes induced by $\mathrm{X}$ on the opens $\mathrm{V}-(\mathrm{V}\cap\overline{\{y\}})$, which are quasi-compact since $\mathrm{X}$ is Noetherian.
Since the opens $\mathrm{V}-(\mathrm{V}\cap\overline{\{y\}})$ are separated, it follows from \sref{IV.8.5.2},~(ii) and \sref{IV.8.5.5} that there exist an affine open neighbourhood $\mathrm{V}$ of $y$ in $\mathrm{X}$ and an invertible module $\mathscr{L}_{\mathrm{V}}'$ on $\mathrm{V}-(\mathrm{V}\cap\overline{\{y\}})$ which induces $\mathscr{L}_0$ on $\mathrm{U}_y$.
But the hypothesis implies, by \sref{IV.21.13.6},~(i) and (ii), that the pair $(\mathrm{V},\mathrm{V}\cap\overline{\{y\}})$ is parafactorial. Hence there exists an invertible $\sh{O}_{\mathrm{V}}$-Module $\mathscr{L}_{\mathrm{V}}$ inducing $\mathscr{L}_{\mathrm{V}}'$ on $\mathrm{V}-(\mathrm{V}\cap\overline{\{y\}})$. Replacing $\mathrm{V}$ if necessary by a smaller open neighbourhood of $y$, we may suppose that $\mathscr{L}_{\mathrm{V}}\cong\sh{O}_{\mathrm{V}}$; it follows that $\mathscr{L}_0\cong\sh{O}_{\mathrm{U}_y}$.

Conversely, let us prove that if all the $\sh{O}_{\mathrm{X},y}$ ($y\in\mathrm{Y}$) are parafactorial, the pair $(\mathrm{X},\mathrm{Y})$ is parafactorial.
This is evidently reduced, taking into account the remark at the beginning, to showing that for every invertible $\sh{O}_{\mathrm{U}}$-module $\mathscr{L}$, $j_*(\mathscr{L})$ is an invertible $\sh{O}_{\mathrm{X}}$-module.
The set of $x\in\mathrm{X}$ such that the restriction of $j_*(\mathscr{L})$ to an open neighbourhood of $x$ in $\mathrm{X}$ is invertible is evidently open in $\mathrm{X}$; it is a matter of showing that $\mathrm{V}=\mathrm{X}$.
For this, suppose the contrary and set $\mathrm{Z}=\mathrm{X}-\mathrm{V}\neq\varnothing$.
Let $y$ be a \emph{maximal} point of $\mathrm{Z}$; since $\mathrm{Z}\subset\mathrm{Y}$, $\sh{O}_{\mathrm{X},y}$ is by hypothesis a parafactorial ring.
Replacing $\mathrm{X}$ if necessary by an open neighbourhood of $y$, we may suppose that the restriction of $j_*(\mathscr{L})$ to $\mathrm{X}-\overline{\{y\}}$ is invertible. With the notation introduced above, the invertible $\sh{O}_{\mathrm{U}_y}$-Module $\mathscr{L}_0$ induced by $j_*(\mathscr{L})$ is therefore invertible. Since $\sh{O}_{\mathrm{X},y}$ is parafactorial, $\mathscr{L}_0$ is induced by an invertible $\sh{O}_{\mathrm{T}_y}$-Module $\mathscr{L}'$.

As $\mathrm{T}_y$ is the projective limit of the open neighbourhoods $\mathrm{W}$ of $y$ in $\mathrm{X}$, another application of \sref{IV.8.5.2},~(ii) and \sref{IV.8.5.5} yields such a neighbourhood $\mathrm{W}$ and an invertible $\sh{O}_{\mathrm{W}}$-Module $\mathscr{L}''$ inducing $\mathscr{L}'$ on $\mathrm{T}_y$, and therefore inducing $\mathscr{L}_0$ on $\mathrm{U}_y$. By \sref{IV.8.5.2.5} and \sref{IV.8.5.2},~(i), after replacing $\mathrm{W}$ by a smaller open neighbourhood of $y$, we may suppose that the restrictions of $j_*(\mathscr{L})$ and $\mathscr{L}''$ to $\mathrm{W}-(\mathrm{W}\cap\overline{\{y\}})$ are equal.

Set $\mathrm{V}_1=\mathrm{V}\cup\mathrm{W}$, and let $\mathscr{L}_1$ be the invertible $\sh{O}_{\mathrm{V}_1}$-Module equal to $\mathscr{L}''$ on $\mathrm{W}$ and to $j_*(\mathscr{L})|\mathrm{V}$ on $\mathrm{V}$. Then $\mathrm{U}\subset\mathrm{V}_1$ and $\mathscr{L}_1|\mathrm{U}=\mathscr{L}$. Since $\sh{O}_{\mathrm{X}}\to j_*(\sh{O}_{\mathrm{U}})$ is bijective, \sref{IV.21.13.2},~$b)$ shows that $j_*(\mathscr{L})|\mathrm{V}_1\cong\mathscr{L}_1$, hence is invertible.
But this contradicts the definition of $\mathrm{V}$, and concludes the proof of \sref{IV.21.13.10}.
\end{proof}

\begin{corollary}[21.13.11]
\label{IV.21.13.11}
Let $\mathrm{X}$ be a locally Noetherian prescheme, $\mathrm{Y}$ a closed subset of $\mathrm{X}$, $\mathrm{U}=\mathrm{X}-\mathrm{Y}$.
Suppose that the pair $(\mathrm{X},\mathrm{Y})$ is parafactorial and that $\mathrm{U}$ is locally factorial; in other terms \sref{IV.21.13.10}, for all $x\in\mathrm{X}$, the local ring $\sh{O}_{\mathrm{X},x}$ is:
$1^{\mathrm{o}}$~parafactorial if $x\in\mathrm{Y}$;
$2^{\mathrm{o}}$~factorial if $x\notin\mathrm{Y}$.
Then $\mathrm{X}$ is locally factorial (in other words, $\sh{O}_{\mathrm{X},x}$ is in fact factorial for \emph{all} $x\in\mathrm{X}$).
\end{corollary}

\begin{proof}
Suppose indeed the contrary, and let $y$ be a point of $\mathrm{Y}$ such that $\sh{O}_{\mathrm{X},y}$ is not factorial and has a minimal dimension among all the points of $\mathrm{Y}$ having this property.
Then, if we set $\mathrm{T}_y=\operatorname{Spec}(\sh{O}_{\mathrm{X},y})$, $\mathrm{U}_y=\mathrm{T}_y-\{y\}$, the hypothesis that $\sh{O}_{\mathrm{X},y}$ is parafactorial implies $\operatorname{Pic}(\mathrm{U}_y)=0$ and $\operatorname{prof}(\sh{O}_{\mathrm{X},y})\geq 2$ \sref{IV.21.13.8}; moreover, the choice of $y$ implies that $\mathrm{U}_y$ is locally factorial.
But then \sref{IV.21.6.14} proves that $\sh{O}_{\mathrm{X},y}$ is factorial, contrary to the hypothesis.
\end{proof}

\begin{proposition}[21.13.12]
\label{IV.21.13.12}
Let $\mathrm{A}$, $\mathrm{B}$ be two local Noetherian rings, $\rho:\mathrm{A}\to\mathrm{B}$ a local homomorphism making $\mathrm{B}$ a flat A-module.
Then, if $\mathrm{B}$ is factorial, so is $\mathrm{A}$.
\end{proposition}

\begin{proof}
We already know that $\mathrm{A}$ is integral and integrally closed \sref{IV.6.5.4}; we may therefore restrict to the case $\dim(\mathrm{A})\geq 2$ and argue by induction on $n=\dim(\mathrm{A})$.
Let $\mathrm{X}=\operatorname{Spec}(\mathrm{A})$, $a$ the closed point of $\mathrm{X}$, $\mathrm{X}'=\operatorname{Spec}(\mathrm{B})$, $f:\mathrm{X}'\to\mathrm{X}$ the structural morphism, $\mathrm{U}=\mathrm{X}-\{a\}$, $\mathrm{U}'=f^{-1}(\mathrm{U})$.
The local rings $\sh{O}_{\mathrm{X}',x'}$ at points $x'\in f^{-1}(a)$ are factorial and of dimension $\geq 2$ \sref{IV.6.1.2}, hence parafactorial \sref{IV.21.13.9},~(ii); we conclude from \sref{IV.21.13.10} that the pair $(\mathrm{X}',f^{-1}(a))$ is parafactorial.
By virtue of \sref{IV.21.13.6},~(iii), this shows that the ring $\mathrm{A}$ is parafactorial, hence $\operatorname{prof}(\mathrm{A})\geq 2$ and $\operatorname{Pic}(\mathrm{U})=0$ \sref{IV.21.13.8}.
Finally, the local rings of $\mathrm{U}$ being of dimension $<n$, the induction hypothesis proves that $\mathrm{U}$ is \emph{locally factorial}; the conclusion follows then from \sref{IV.21.6.14}.
\end{proof}

\begin{env}[21.13.12.1]
\label{IV.21.13.12.1}
\oldpage[IV-4]{321}
We note that the statement of \sref{IV.21.13.12} where one replaces ``factorial'' by ``parafactorial'' is no longer exact, as shown by the example where $\mathrm{A}=k$ is a field and $\mathrm{B}$ the local factorial ring $k[[\mathrm{T}_1,\mathrm{T}_2]]$.
However, it follows from \sref{IV.21.13.6},~(iii) that, under the conditions of \sref{IV.21.13.12}, if $\mathfrak{m}$ denotes the maximal ideal of $\mathrm{A}$ and if $\mathfrak{m}\mathrm{B}$ is an \emph{ideal of definition} of $\mathrm{B}$ (i.e.\ $\dim(\mathrm{B}/\mathfrak{m}\mathrm{B})=0$), then, when $\mathrm{B}$ is parafactorial, so is $\mathrm{A}$.
In particular, if the completion $\widehat{\mathrm{A}}$ of a local Noetherian ring is parafactorial, so is $\mathrm{A}$.
\end{env}

\begin{remark}[21.13.13]
\label{IV.21.13.13}
Some of the preceding definitions and results belong to more general considerations on the cohomology of sheaves of commutative groups, developed in chap.~III, Part~3, for which we have given an independent account here for the convenience of the reader.
Indeed, if $\mathrm{X}$ is a ringed space, there is a canonical \emph{equivalence} between the category of invertible $\sh{O}_{\mathrm{X}}$-modules and that of \emph{principal homogeneous sheaves under the sheaf of groups} $\sh{O}_{\mathrm{X}}^*$ \sref{IV.16.5.15}.
This equivalence associates with every invertible $\sh{O}_{\mathrm{X}}$-module $\mathscr{L}$ the sheaf $\mathscr{L}^*=\mathscr{I}som_{\sh{O}_{\mathrm{X}}}(\sh{O}_{\mathrm{X}},\mathscr{L})$ of germs of isomorphisms of $\sh{O}_{\mathrm{X}}$ onto $\mathscr{L}$; it is immediate that $\mathscr{L}^*$ is canonically equipped with a structure of principal homogeneous sheaf under $\sh{O}_{\mathrm{X}}^*\cong\mathscr{I}som_{\sh{O}_{\mathrm{X}}}(\sh{O}_{\mathrm{X}},\sh{O}_{\mathrm{X}})$.
The verification of this equivalence of categories is immediate.
This being said, let $\mathrm{X}$ be a topological space, $\mathscr{G}$ a sheaf of groups on $\mathrm{X}$, $\mathrm{Y}$ a closed subset of $\mathrm{X}$, and $i$ an integer such that $0\leq i\leq 2$. We say that the pair $(\mathrm{X},\mathrm{Y})$ is \emph{$i$-pure} for $\mathscr{G}$ if, for every open $\mathrm{V}$ of $\mathrm{X}$, the restriction functor $\mathscr{P}\mapsto\mathscr{P}|(\mathrm{U}\cap\mathrm{V})$, from the category of principal homogeneous sheaves under the sheaf of groups $\mathscr{G}|\mathrm{V}$ into the analogous category under $\mathscr{G}|(\mathrm{U}\cap\mathrm{V})$, is \emph{faithful} ($i=0$), resp.\ \emph{fully faithful} ($i=1$), resp.\ an \emph{equivalence of categories} ($i=2$).
In the case where $\mathrm{X}$ is a ringed space and $\mathscr{G}=\sh{O}_{\mathrm{X}}^*$, to say that $(\mathrm{X},\mathrm{Y})$ is $i$-pure therefore means, for $i=0$, that the canonical homomorphism $\sh{O}_{\mathrm{X}}\to j_*(\sh{O}_{\mathrm{U}})$ is \emph{injective}; for $i=1$, that this homomorphism is \emph{bijective}; and finally, for $i=2$, that the pair $(\mathrm{X},\mathrm{Y})$ is \emph{parafactorial} \sref{IV.21.13.3}.

Returning to the general case, recall that the morphisms of the category of principal sheaves are by definition isomorphisms.
To say that the pair $(\mathrm{X},\mathrm{Y})$ is $0$-pure means therefore that for every open $\mathrm{V}$ of $\mathrm{X}$, the canonical homomorphism $\mathrm{H}^0(\mathrm{V},\mathscr{G})\to\mathrm{H}^0(\mathrm{U}\cap\mathrm{V},\mathscr{G})$ is injective; to say that the pair $(\mathrm{X},\mathrm{Y})$ is $1$-pure means that the canonical homomorphism $\mathrm{H}^0(\mathrm{V},\mathscr{G})\to\mathrm{H}^0(\mathrm{U}\cap\mathrm{V},\mathscr{G})$ is bijective (and so also the canonical homomorphism $\mathrm{H}^1(\mathrm{V},\mathscr{G})\to\mathrm{H}^1(\mathrm{U}\cap\mathrm{V},\mathscr{G})$ is injective); finally, one can show that for the pair $(\mathrm{X},\mathrm{Y})$ to be $2$-pure, it is necessary and sufficient that the homomorphisms $\mathrm{H}^i(\mathrm{V},\mathscr{G})\to\mathrm{H}^i(\mathrm{U}\cap\mathrm{V},\mathscr{G})$ be bijective for $i=0$ and $i=1$.
When $\mathscr{G}$ is a sheaf of commutative groups, introduce the cohomology sheaves $\mathscr{H}^p_{\mathrm{Y}}(\mathscr{G})$ defined in chap.~III, Part~3. To say that $(\mathrm{X},\mathrm{Y})$ is $i$-pure for $i\leq2$ then means that $\mathscr{H}^p_{\mathrm{Y}}(\mathscr{G})=0$ for $p\leq i$; in this form, the notion generalises immediately to every integer $i$.
\end{remark}

\begin{proposition}[21.13.14]
\label{IV.21.13.14}
Let $\mathrm{X}$ be a locally Noetherian \emph{normal} prescheme, $(\mathrm{U}_\lambda)_{\lambda\in\mathrm{L}}$ a decreasing filtered family of opens of $\mathrm{X}$.
The following conditions are equivalent:
\begin{enumerate}[label=\emph{\alph*)}]
  \item Every $1$-codimensional cycle $\mathrm{Z}$ on $\mathrm{X}$ such that there exists an index $\lambda$ for which $\mathrm{Z}|\mathrm{U}_\lambda$ is locally principal \sref{IV.21.6.7}, is locally principal.
  \item For all $x\in\mathrm{X}$ such that $\dim(\sh{O}_{\mathrm{X},x})\geq 2$ and such that $x$ does not belong to $\bigcap_{\lambda\in\mathrm{L}}\mathrm{U}_\lambda$, the local ring $\sh{O}_{\mathrm{X},x}$ is parafactorial.
\end{enumerate}
\end{proposition}

\begin{proof}
\oldpage[IV-4]{322}
Condition $b)$ may also be expressed as follows:

\emph{$b')$ For every closed subset $\mathrm{Y}$ of $\mathrm{X}$ such that $\operatorname{codim}(\mathrm{Y},\mathrm{X})\geq 2$ and $\mathrm{Y}\subset\mathrm{X}-\mathrm{U}_\lambda$ for some $\lambda$, the pair $(\mathrm{X},\mathrm{Y})$ is parafactorial.}

Property $a)$ may in turn be expressed as follows. If the closed set $\mathrm{N}$ of points $x\in\mathrm{X}$ at which $\mathrm{Z}$ is not principal is contained in one of the sets $\mathrm{X}-\mathrm{U}_\lambda$, then $\mathrm{Z}$ is locally principal. Since $\mathrm{X}$ is normal, the condition $\dim(\sh{O}_{\mathrm{X},x})\leq 1$ implies that $\sh{O}_{\mathrm{X},x}$ is a field or a discrete valuation ring, and therefore that $\mathrm{Z}_x$ is principal; thus necessarily $\operatorname{codim}(\mathrm{N},\mathrm{X})\geq 2$ \sref{IV.5.1.3}. Consequently, $a)$ is equivalent to:

\emph{$a')$ If there exists a closed subset $\mathrm{Y}$ contained in one of the sets $\mathrm{X}-\mathrm{U}_\lambda$, such that $\operatorname{codim}(\mathrm{Y},\mathrm{X})\geq 2$ and $\mathrm{Z}|(\mathrm{X}-\mathrm{Y})$ is locally principal, then $\mathrm{Z}$ is locally principal.}

Condition $b)$ implies $b')$ by \sref{IV.21.13.10}. Conversely, if $b')$ holds and $x\notin\mathrm{U}_\lambda$, then $\mathrm{Y}=\overline{\{x\}}$ is contained in $\mathrm{X}-\mathrm{U}_\lambda$ and $\operatorname{codim}(\mathrm{Y},\mathrm{X})\geq 2$; it follows again from \sref{IV.21.13.10} that $\sh{O}_{\mathrm{X},x}$ is parafactorial. This proves the equivalence of $b)$ and $b')$.

We are thus reduced to proving the equivalence of $a')$ and $b')$. Since $\mathrm{X}$ is normal, for every closed subset $\mathrm{Y}$ of $\mathrm{X}$ such that $\operatorname{codim}(\mathrm{Y},\mathrm{X})\geq 2$, one has $\operatorname{prof}_{\mathrm{Y}}(\sh{O}_{\mathrm{X}})\geq 2$ \sref{IV.5.8.6}. If $\mathrm{U}=\mathrm{X}-\mathrm{Y}$ and $j:\mathrm{U}\to\mathrm{X}$ is the canonical injection, the homomorphism $\sh{O}_{\mathrm{X}}\to j_*(\sh{O}_{\mathrm{U}})$ is therefore bijective \sref{IV.21.13.4}. Since $a')$ and $b')$ are local conditions, it suffices to prove, for $\mathrm{X}$ Noetherian and normal and $\mathrm{Y}$ closed with $\operatorname{codim}(\mathrm{Y},\mathrm{X})\geq 2$, the equivalence of:

\emph{$a'')$ For every $1$-codimensional cycle $\mathrm{Z}$ on $\mathrm{X}$, if $\mathrm{Z}|\mathrm{U}$ is locally principal, then $\mathrm{Z}$ is locally principal.}

\emph{$b'')$ The canonical homomorphism $\operatorname{Pic}(\mathrm{X})\to\operatorname{Pic}(\mathrm{U})$ is surjective.}

First suppose $a'')$, and let $\zeta_0\in\operatorname{Pic}(\mathrm{U})$. Its image in $\mathfrak{Cl}(\mathrm{U})$ \sref{IV.21.6.11} is the class of a locally principal $1$-codimensional cycle $\mathrm{Z}_0$ on $\mathrm{U}$. The hypothesis $\operatorname{codim}(\mathrm{Y},\mathrm{X})\geq 2$ implies that the restriction homomorphism $\mathscr{Z}^1(\mathrm{X})\to\mathscr{Z}^1(\mathrm{U})$ is bijective; hence $\mathrm{Z}_0=\mathrm{Z}|\mathrm{U}$ for a $1$-codimensional cycle $\mathrm{Z}$ on $\mathrm{X}$. By $a'')$, $\mathrm{Z}$ is locally principal. It follows from \sref{IV.21.6.11} that the image of $\mathrm{Z}$ in $\mathfrak{Cl}(\mathrm{X})$ is the image of a unique element $\zeta\in\operatorname{Pic}(\mathrm{X})$, and it is then clear that $\zeta_0$ is the image of $\zeta$.

Conversely, suppose $b'')$, and let $\mathrm{Z}$ be a $1$-codimensional cycle on $\mathrm{X}$ such that $\mathrm{Z}|\mathrm{U}$ is locally principal. The image of $\mathrm{Z}|\mathrm{U}$ in $\mathfrak{Cl}(\mathrm{U})$ is the image of a unique element $\zeta_0\in\operatorname{Pic}(\mathrm{U})$ \sref{IV.21.6.11}. By hypothesis, there exists $\zeta\in\operatorname{Pic}(\mathrm{X})$ whose image in $\operatorname{Pic}(\mathrm{U})$ is $\zeta_0$; the image of $\zeta$ in $\mathfrak{Cl}(\mathrm{X})$ is therefore the class of a locally principal $1$-codimensional cycle $\mathrm{Z}'$ on $\mathrm{X}$ such that $\mathrm{Z}|\mathrm{U}$ and $\mathrm{Z}'|\mathrm{U}$ are linearly equivalent. But $\mathrm{U}$ is schematically dense in $\mathrm{X}$, and the isomorphism $\mathscr{Z}^1(\mathrm{X})\to\mathscr{Z}^1(\mathrm{U})$ carries $\mathscr{Z}_{\mathrm{princ}}^1(\mathrm{X})$ onto $\mathscr{Z}_{\mathrm{princ}}^1(\mathrm{U})$; hence $\mathrm{Z}$ and $\mathrm{Z}'$ are linearly equivalent. Since $\mathrm{Z}'$ is locally principal, so is $\mathrm{Z}$.
\end{proof}

\begin{corollary}[21.13.15]
\label{IV.21.13.15}
Let $\mathrm{X}$ be a locally Noetherian normal prescheme, $\mathrm{S}$ a subset of $\mathrm{X}$.
The following conditions are equivalent:
\begin{enumerate}[label=\emph{\alph*)}]
  \item Every $1$-codimensional cycle on $\mathrm{X}$ which is principal at the points of $\mathrm{S}$ is locally principal.
  \oldpage[IV-4]{323}
  \item For all $x\in\mathrm{X}$ which is not a generisation of a point of $\mathrm{S}$ (in other words, such that $\overline{\{x\}}\cap\mathrm{S}=\varnothing$) and such that $\dim(\sh{O}_{\mathrm{X},x})\geq 2$, the local ring $\sh{O}_{\mathrm{X},x}$ is parafactorial.
\end{enumerate}
\end{corollary}

\begin{proof}
It is clear that the set $\mathrm{S}'$ of generisations in $\mathrm{X}$ of points of $\mathrm{S}$ is the intersection of the open neighbourhoods of $\mathrm{S}$. We may restrict to the case where $\mathrm{X}$ is Noetherian, since properties $a)$ and $b)$ are local on $\mathrm{X}$. Every $1$-codimensional cycle on $\mathrm{X}$ which is principal at the points of $\mathrm{S}$ is also principal at the points of an open neighbourhood of $\mathrm{S}$; condition $a)$ of \sref{IV.21.13.15} is therefore condition $a)$ of \sref{IV.21.13.14} applied to the family $(\mathrm{U}_\lambda)$ of open neighbourhoods of $\mathrm{S}$. The corollary now follows from the equivalence of $a)$ and $b)$ in \sref{IV.21.13.14}.
\end{proof}

\begin{remark}[21.13.16]
\label{IV.21.13.16}
Under the general conditions of \sref{IV.21.13.14}, suppose moreover $\varinjlim\operatorname{Pic}(\mathrm{U}_\lambda)=0$.
Then condition $a)$ of \sref{IV.21.13.14} is also equivalent to the following:
\begin{enumerate}
  \item[c)] \emph{Every $1$-codimensional cycle on $\mathrm{X}$ whose support is contained in one of the sets $\mathrm{X}-\mathrm{U}_\lambda$ is locally principal.}
\end{enumerate}

We can restrict to the case where $\mathrm{X}$ is irreducible and all the $\mathrm{U}_\lambda$ are non-empty.
It is clear that $a)$ implies $c)$, since if the $1$-codimensional cycle $\mathrm{Z}$ on $\mathrm{X}$ has its support in $\mathrm{X}-\mathrm{U}_\lambda$, $\mathrm{Z}|\mathrm{U}_\lambda=0$, hence $\mathrm{Z}|\mathrm{U}_\lambda$ is locally principal.
Inversely $c)$ implies $a)$: let $\mathrm{Z}$ be a $1$-codimensional cycle on $\mathrm{X}$ such that $\mathrm{Z}|\mathrm{U}_\lambda$ is locally principal; since $\mathrm{X}$ is normal, $\mathrm{Z}|\mathrm{U}_\lambda=\operatorname{cyc}(\mathrm{D}_\lambda)$, where $\mathrm{D}_\lambda$ is a divisor on $\mathrm{U}_\lambda$.
The hypothesis $\varinjlim\operatorname{Pic}(\mathrm{U}_\lambda)=0$ implies, by \sref{IV.21.3.4}, that there is an index $\mu$ with $\mathrm{U}_\mu\subset\mathrm{U}_\lambda$ such that $\mathrm{D}_\mu=\mathrm{D}_\lambda|\mathrm{U}_\mu$ is equivalent to $0$.
If $\mathrm{D}_\mu=\operatorname{div}(f_\mu)$, $f_\mu$ may be considered as a regular rational function on $\mathrm{X}$, regular on $\mathrm{U}_\mu$; hence $\mathrm{Z}'=\operatorname{cyc}(\operatorname{div}(f_\mu))$, we see therefore that $\mathrm{Z}''=\mathrm{Z}-\mathrm{Z}'$ has its support in $\mathrm{X}-\mathrm{U}_\mu$; by hypothesis, $\mathrm{Z}''$ is locally principal, whence the conclusion.

The condition $\varinjlim\operatorname{Pic}(\mathrm{U}_\lambda)=0$ is satisfied in particular when $(\mathrm{U}_\lambda)$ is the family of open neighbourhoods of a point $z\in\mathrm{X}$: for every invertible $\sh{O}_{\mathrm{U}_\lambda}$-Module $\mathscr{L}_\lambda$, there exists by definition an open neighbourhood $\mathrm{U}_\mu\subset\mathrm{U}_\lambda$ of $z$ such that $\mathscr{L}_\lambda|\mathrm{U}_\mu\cong\sh{O}_{\mathrm{U}_\mu}$.
We see therefore that, in \sref{IV.21.13.15}, if $\mathrm{S}=\{z\}$, conditions $a)$ and $b)$ are equivalent also to the following:
\begin{enumerate}
  \item[c)] \emph{Every $1$-codimensional cycle on $\mathrm{X}$ whose support does not contain $z$ is locally principal.}
\end{enumerate}

If moreover $\operatorname{Pic}(\mathrm{X})=0$, it follows that this condition implies that every $1$-codimensional cycle on $\mathrm{X}$ whose support does not contain $z$ is principal.
\end{remark}

\subsection{The Ramanujam-Samuel theorem}
\label{subsection:IV.21.14}
\label{IV.21.14}

\begin{theorem}[21.14.1]
\label{IV.21.14.1}
\textup{(Ramanujam-Samuel)} ---
Let $\mathrm{A}$ be a local Noetherian ring with maximal ideal $\mathfrak{m}$, such that its completion $\widehat{\mathrm{A}}$ is integral and integrally closed.
Let $\mathrm{B}$ be a local Noetherian ring such that $\dim(\mathrm{B})\geq\dim(\mathrm{A})$, and let $\rho:\mathrm{A}\to\mathrm{B}$ be a local homomorphism making $\mathrm{B}$ a formally smooth $\mathrm{A}$-algebra (for the preadic topologies $(\mathbf{0}$, $19.3.1)$) and such that the residue field of $\mathrm{B}$ is finite over that of $\mathrm{A}$.
Then every $1$-codimensional cycle on $\operatorname{Spec}(\mathrm{B})$ which is principal at the point $\mathfrak{p}=\mathfrak{m}\mathrm{B}$, is a principal $1$-codimensional cycle.
\end{theorem}

\begin{proof}
\oldpage[IV-4]{324}
If $k=\mathrm{A}/\mathfrak{m}$ is the residue field of $\mathrm{A}$, then $\mathrm{B}/\mathfrak{m}\mathrm{B}$ is a formally smooth $k$-algebra for its preadic topology $(\mathbf{0}$, $19.3.5)$, hence regular and in particular integral. Thus $\mathfrak{p}=\mathfrak{m}\mathrm{B}$ is indeed a prime ideal of $\mathrm{B}$, which justifies the statement.
It evidently all amounts to proving that every prime ideal $\mathfrak{q}$ of height one of $\mathrm{B}$ not contained in $\mathfrak{m}\mathrm{B}$ is \emph{principal}.

Let $\widehat{\mathrm{A}}$, $\widehat{\mathrm{B}}$ be the completions of $\mathrm{A}$ and $\mathrm{B}$ respectively, so that the maximal ideal of $\widehat{\mathrm{A}}$ is $\mathfrak{m}\widehat{\mathrm{A}}$; we know $(\mathbf{0}$, $19.3.6)$ that $\widehat{\mathrm{B}}$ is a formally smooth $\widehat{\mathrm{A}}$-algebra for the adic topologies.
Let $k'$ be the residue field of $\widehat{\mathrm{B}}$, a finite extension of $k$. There exists a local homomorphism $\widehat{\mathrm{A}}\to\mathrm{C}$, where $\mathrm{C}$ is a local Noetherian ring which is finite and flat (hence free) as an $\widehat{\mathrm{A}}$-module and such that $\mathrm{C}/\mathfrak{m}\mathrm{C}\cong k'$ $(\mathbf{0}_{\mathrm{III}}$, $10.3.1)$. It follows that $\mathrm{C}$ is complete, and then from \sref{IV.7.5.1}, \sref{IV.7.5.3}, and \sref{IV.6.5.4},~(ii) that $\mathrm{C}$ is \emph{integral} and \emph{integrally closed}.

Moreover,
\[
  \mathrm{D}=\widehat{\mathrm{B}}\widehat{\otimes}_{\widehat{\mathrm{A}}}\mathrm{C}
\]
is a complete semilocal ring, a direct product of complete local rings, one of which, $\mathrm{D}_0$, has residue field $k'$ (since $k'\otimes_k k'$ is a direct product of local rings, one of them isomorphic to $k'$). Since $\mathrm{D}$ is formally smooth over $\mathrm{C}$, so is $\mathrm{D}_0$; consequently $\mathrm{D}_0/\mathfrak{m}\mathrm{D}_0$ is a formally smooth $k'$-algebra with residue field $k'$, and is therefore $k'$-isomorphic to a formal power-series algebra $k'[[\mathrm{T}_1,\ldots,\mathrm{T}_n]]$ $(\mathbf{0}$, $19.6.4)$. It follows from $(\mathbf{0}$, $19.7.1.5)$ that $\mathrm{D}_0$ is $\mathrm{C}$-isomorphic to $\mathrm{C}[[\mathrm{T}_1,\ldots,\mathrm{T}_n]]$, and hence is integral and integrally closed (Bourbaki, \emph{Alg.\ comm.}, chap.~V, \S1, no.~4, Prop.~14). Since the morphisms $\operatorname{Spec}(\mathrm{D}_0)\to\operatorname{Spec}(\widehat{\mathrm{B}})$ and $\operatorname{Spec}(\mathrm{D}_0)\to\operatorname{Spec}(\mathrm{B})$ are faithfully flat, it follows that $\mathrm{B}$ and $\widehat{\mathrm{B}}$ are likewise integral and integrally closed \sref{IV.2.1.13}.

The ideal $\mathfrak{q}\mathrm{D}_0$ is therefore divisorial (Bourbaki, \emph{Alg.\ comm.}, chap.~VII, \S1, no.~10, Prop.~15), and it is not contained in $\mathfrak{m}\mathrm{D}_0$; otherwise faithful flatness would give
\[
  \mathfrak{q}=(\mathfrak{q}\mathrm{D}_0)\cap\mathrm{B}
  \subset(\mathfrak{m}\mathrm{D}_0)\cap\mathrm{B}
  =\mathfrak{m}\mathrm{B}=\mathfrak{p}
\]
$(\mathbf{0}_{\mathrm{I}}$, $6.5.1)$. Hence none of the height-one prime ideals $\mathfrak{r}_h$ of $\mathrm{D}_0$ containing $\mathfrak{q}\mathrm{D}_0$ is contained in $\mathfrak{m}\mathrm{D}_0$. If these ideals are principal, then $\mathfrak{q}\mathrm{D}_0$ is principal, since the divisors (in Bourbaki's sense) of $\mathfrak{q}\mathrm{D}_0$ and of a product of powers of the $\mathfrak{r}_h$ are equal (Bourbaki, \emph{Alg.\ comm.}, chap.~VII, \S1, no.~4, Prop.~5). As $\mathfrak{q}\mathrm{D}_0\cap\mathrm{B}=\mathfrak{q}$, faithful flatness then shows that $\mathfrak{q}$ is principal \sref{IV.2.5.2}, using the fact that $\mathrm{B}$ is local.

We are thus reduced to proving the theorem in the particular case where $\mathrm{A}$ is complete and $\mathrm{B}=\mathrm{A}[[\mathrm{T}_1,\ldots,\mathrm{T}_n]]$. Let us first show that we may reduce to $n=1$. We argue by induction on $n$, and choose $f\in\mathfrak{q}$ with $f\notin\mathfrak{p}=\mathfrak{m}\mathrm{B}$. It is enough to show that there exists an $\mathrm{A}$-automorphism $\sigma$ of $\mathrm{B}$ such that $\sigma(f)$ does not belong to
\[
  \mathfrak{p}'=\mathfrak{p}+\mathrm{B}\mathrm{T}_1+\cdots+\mathrm{B}\mathrm{T}_{n-1}.
\]
Indeed, if $\mathrm{B}'=\mathrm{A}[[\mathrm{T}_1,\ldots,\mathrm{T}_{n-1}]]$, then $\mathrm{B}'$ is complete and has maximal ideal $\mathfrak{m}\mathrm{B}'+\mathrm{B}'\mathrm{T}_1+\cdots+\mathrm{B}'\mathrm{T}_{n-1}=\mathfrak{m}'$. Since $\mathrm{B}=\mathrm{B}'[[\mathrm{T}_n]]$ and $\mathfrak{p}'=\mathfrak{m}'\mathrm{B}$, the induction hypothesis, taking into account that $\mathrm{B}'$ is integral and integrally closed, shows that $\sigma(\mathfrak{q})$ is principal in $\mathrm{B}$, and hence so is $\mathfrak{q}$. Now, if $\overline{f}\in k[[\mathrm{T}_1,\ldots,\mathrm{T}_n]]$ is the image of $f$, then $\overline{f}\neq0$; one knows (Bourbaki, \emph{Alg.\ comm.}, chap.~VII, \S3, no.~7, Lemma~3) that there is an $\mathrm{A}$-automorphism $\sigma$ of $\mathrm{B}$ such that $(\overline{\sigma(f)})(0,\ldots,0,\mathrm{T}_n)\neq0$, which plainly implies $\sigma(f)\notin\mathfrak{p}'$.

Suppose therefore $n=1$, and write $\mathrm{T}$ instead of $\mathrm{T}_1$, so that $\mathrm{B}=\mathrm{A}[[\mathrm{T}]]$. It is enough to show that the ideal $\mathfrak{q}\cap\mathrm{A}[\mathrm{T}]$ of the polynomial ring $\mathrm{A}[\mathrm{T}]$ is principal.
\oldpage[IV-4]{325}
Choose $f_0\in\mathfrak{q}$ with $f_0\notin\mathfrak{m}\mathrm{B}$; it is a formal power series whose reduced series $\overline{f}_0$ is nonzero. By the preparation theorem (Bourbaki, \emph{Alg.\ comm.}, chap.~VII, \S3, no.~8, Prop.~5), for every $f\in\mathfrak{q}$ there exist $g\in\mathrm{B}$ and $r\in\mathrm{A}[\mathrm{T}]$ such that $f=gf_0+r$, so $r\in\mathfrak{q}\cap\mathrm{A}[\mathrm{T}]$. On the other hand (\emph{loc.\ cit.}, Prop.~6), there exist a nonconstant distinguished polynomial $\mathrm{F}_0\in\mathrm{A}[\mathrm{T}]$ and a unit $u\in\mathrm{B}$ such that $f_0=u\mathrm{F}_0$; hence $\mathrm{F}_0\in\mathfrak{q}\cap\mathrm{A}[\mathrm{T}]$. Thus, putting $\mathfrak{q}_1=\mathfrak{q}\cap\mathrm{A}[\mathrm{T}]$, one has $\mathfrak{q}=\mathfrak{q}_1\mathrm{B}$. Since $\mathrm{B}$ is flat over $\mathrm{A}[\mathrm{T}]$ $(\mathbf{0}_{\mathrm{I}}$, $7.3.3)$, it follows from Bourbaki, \emph{Alg.\ comm.}, chap.~VII, \S1, no.~10, Prop.~15 that $\mathfrak{q}_1$ is a height-one prime ideal of $\mathrm{A}[\mathrm{T}]$.

Necessarily $\mathfrak{q}_1\cap\mathrm{A}=0$. Otherwise $\mathfrak{q}_1\cap\mathrm{A}$ would have height at least $1$, and \sref{IV.5.5.3} would give $\mathfrak{q}_1=(\mathfrak{q}_1\cap\mathrm{A})\mathrm{A}[\mathrm{T}]$; since $\mathfrak{q}_1\cap\mathrm{A}\subset\mathfrak{m}$, this would imply $\mathfrak{q}_1\subset\mathfrak{m}\mathrm{A}[\mathrm{T}]$, contrary to the hypothesis on $\mathfrak{q}$.

Let $\mathrm{K}$ be the field of fractions of $\mathrm{A}$. Then $\mathfrak{q}_1\mathrm{K}[\mathrm{T}]$ is a prime ideal distinct from both $0$ and $\mathrm{K}[\mathrm{T}]$, hence is of the form $h\mathrm{K}[\mathrm{T}]$, where
\[
  h(\mathrm{T})=\mathrm{T}^m+a_1\mathrm{T}^{m-1}+\cdots+a_m,
  \qquad m\geq1,
\]
the $a_i$ belonging to $\mathrm{K}$ and $h$ being irreducible in $\mathrm{K}[\mathrm{T}]$. We saw above that $\mathfrak{q}_1$ contains a nonconstant distinguished polynomial $\mathrm{F}_0\in\mathrm{A}[\mathrm{T}]$. If $t$ is the class of $\mathrm{T}$ in $\mathrm{K}[\mathrm{T}]/\mathfrak{q}_1\mathrm{K}[\mathrm{T}]$, then $t$ is a root of $\mathrm{F}_0$ in an extension of $\mathrm{K}$, and therefore $h$ divides $\mathrm{F}_0$ in $\mathrm{K}[\mathrm{T}]$. Since $h$ and $\mathrm{F}_0$ are monic, the coefficients $a_i$ of $h$ are integral over $\mathrm{A}$ (Bourbaki, \emph{Alg.\ comm.}, chap.~V, \S1, no.~3, Prop.~11), and thus belong to $\mathrm{A}$ because $\mathrm{A}$ is integrally closed. In other words,
\[
  h\in\mathrm{A}[\mathrm{T}]\cap\mathfrak{q}_1\mathrm{K}[\mathrm{T}]=\mathfrak{q}_1
\]
(Bourbaki, \emph{Alg.\ comm.}, chap.~II, \S1, no.~5, Prop.~11). Since every polynomial $g\in\mathfrak{q}_1$ is divisible by $h$ in $\mathrm{K}[\mathrm{T}]$ and $h$ is monic, the coefficients of $g/h$ belong to $\mathrm{A}$; hence $\mathfrak{q}_1=h\mathrm{A}[\mathrm{T}]$. This completes the proof.
\end{proof}

The statement \sref{IV.21.14.1} is \emph{equivalent} to the following:

\begin{corollary}[21.14.2]
\label{IV.21.14.2}
Under the hypotheses of \sref{IV.21.14.1} on $\mathrm{A}$, $\mathrm{B}$ and $\rho$, for every prime ideal $\mathfrak{q}$ of $\mathrm{B}$ not contained in $\mathfrak{p}=\mathfrak{m}\mathrm{B}$ and such that $\dim(\mathrm{B}_{\mathfrak{q}})\geq 2$, the ring $\mathrm{B}_{\mathfrak{q}}$ is parafactorial.
In particular, if $\dim(\mathrm{B}\otimes_{\mathrm{A}}k)>0$ (i.e.\ $\dim\mathrm{B}>\dim\mathrm{A}$) $((\mathbf{0}$, $19.7.1)$ and $(\mathbf{0}$, $6.1.1)))$, the ring $\mathrm{B}$ is parafactorial.
\end{corollary}

\begin{proof}
Indeed, we saw in the proof of \sref{IV.21.14.1} that the conditions of the statement imply that $\mathrm{B}$ is integral and integrally closed; and the equivalence of the statements \sref{IV.21.14.1} and \sref{IV.21.14.2} follows then from \sref{IV.21.13.15} applied to $\mathrm{X}=\operatorname{Spec}(\mathrm{B})$ and $\mathrm{S}=\{\mathfrak{p}\}$.
\end{proof}

\begin{proposition}[21.14.3]
\label{IV.21.14.3}
Let $\mathrm{S}$ be a normal prescheme, $f:\mathrm{X}\to\mathrm{S}$ a smooth morphism.
\begin{enumerate}[label=(\roman*)]
  \item If $\mathrm{S}$ is locally Noetherian (hence also $\mathrm{X}$), then, for all $x\in\mathrm{X}$ such that $\dim(\sh{O}_{\mathrm{X},x})\geq 2$ and such that $x$ is not a maximal point of its fibre $f^{-1}(f(x))$, the ring $\sh{O}_{\mathrm{X},x}$ is parafactorial.
  \oldpage[IV-4]{326}
  Every $1$-codimensional cycle $\mathrm{Z}$ on $\mathrm{X}$ such that $\operatorname{Supp}(\mathrm{Z})$ contains no irreducible component of any fibre $\mathrm{X}_s$, $f^{-1}(s)=\mathrm{X}_s$, is locally principal.
  \item Let $\mathrm{Y}$ be a closed subset of $\mathrm{X}$ containing no irreducible component of any fibre $\mathrm{X}_s$, and such that, for every maximal point $\eta$ of $\mathrm{S}$, $\mathrm{Y}_\eta=\mathrm{Y}\cap\mathrm{X}_\eta$ has codimension $\geq 2$ in $\mathrm{X}_\eta$.
  Then the pair $(\mathrm{X},\mathrm{Y})$ is parafactorial.
\end{enumerate}
\end{proposition}

\begin{proof}
We note that the conditions of (ii) are satisfied if $\mathrm{Y}$ is a closed subset of $\mathrm{X}$ such that, for all $s\in\mathrm{S}$, $\mathrm{Y}_s=\mathrm{Y}\cap\mathrm{X}_s$ is of codimension $\geq 2$ in $\mathrm{X}_s$.

To prove (i), suppose first that $\mathrm{Y}=\overline{\{x\}}$. There is an open neighbourhood $\mathrm{V}$ of $x$ in $\mathrm{X}$ such that $\mathrm{Y}\cap\mathrm{V}$ and $\mathrm{V}$ satisfy the conditions of (ii). Indeed, the hypothesis and \sref{IV.9.5.3} allow us to choose $\mathrm{V}$ so that $\mathrm{Y}\cap\mathrm{V}$ contains no irreducible component of a fibre $\mathrm{X}_s$. Moreover, for every maximal point $\eta$ of $\mathrm{S}$ for which $\mathrm{Y}\cap\mathrm{X}_\eta\neq\varnothing$, the points $z$ of $\mathrm{Y}\cap\mathrm{X}_\eta$ are specializations of $x$; hence $\dim(\sh{O}_{\mathrm{X},z})\geq2$. Since $\sh{O}_{\mathrm{X},z}=\sh{O}_{\mathrm{X}_\eta,z}$ (the prescheme $\mathrm{S}$ being reduced at $\eta$, by \sref{IV.2.1.13}), $\mathrm{Y}\cap\mathrm{X}_\eta$ has codimension $\geq2$ in $\mathrm{X}_\eta$. Thus (ii) makes the pair $(\mathrm{V},\mathrm{Y}\cap\mathrm{V})$ parafactorial, and \sref{IV.21.13.10} shows that $\sh{O}_{\mathrm{X},x}$ is parafactorial. This proves the first assertion of (i).

For the second assertion, it is enough to consider the case $\mathrm{Z}=\overline{\{z\}}$, where $z\in\mathrm{X}^{(1)}$. Since $\sh{O}_{\mathrm{X},z}$ is an integrally closed Noetherian local ring of dimension $1$, it is a discrete valuation ring, and $\mathrm{Z}$ is therefore principal at $z$. At every other point $x\in\operatorname{Supp}(\mathrm{Z})$, one has $\dim(\sh{O}_{\mathrm{X},x})\geq2$, and by hypothesis $x$ is not a maximal point of its fibre; hence $\sh{O}_{\mathrm{X},x}$ is parafactorial. Apply \sref{IV.21.13.15}, taking for $\mathrm{S}$ the set formed by $z$ and by the maximal points of the fibres of $f$. The cycle $\mathrm{Z}$ is principal at the points of this set, since $z$ is the only point of it belonging to $\operatorname{Supp}(\mathrm{Z})$; and condition~(b) of \sref{IV.21.13.15} is plainly satisfied. It follows that $\mathrm{Z}$ is locally principal.

It remains to prove (ii). Put $\mathrm{U}=\mathrm{X}-\mathrm{Y}$, and let $j:\mathrm{U}\to\mathrm{X}$ be the canonical injection. Since the hypotheses and the conclusion are local on $\mathrm{S}$ and on $\mathrm{X}$ by \sref{IV.21.13.6},~(i), we may assume that $\mathrm{S}$ and $\mathrm{X}$ are affine; then $f$ is of finite presentation. Since the fibres of $f$ are regular, the hypothesis on $\mathrm{Y}$ implies that, for every $x\in\mathrm{Y}$,
\[
  \operatorname{prof}(\sh{O}_{\mathrm{X}_{f(x)},x})
  =\dim(\sh{O}_{\mathrm{X}_{f(x)},x})\geq1.
\]
It follows from \sref{IV.19.9.8} that the canonical homomorphism $\sh{O}_{\mathrm{X}}\to j_*(\sh{O}_{\mathrm{U}})$ is injective. Moreover, replacing $\mathrm{Y}$ by a larger closed subset by the method used in the proof of \sref{IV.19.9.8}, and using \sref{IV.21.13.6},~(ii), we may suppose that $\mathrm{Y}$ is defined by an ideal of finite type in the ring of $\mathrm{X}$. Consequently $\mathrm{U}=\mathrm{X}-\mathrm{Y}$ is quasi-compact and quasi-separated, and $\mathrm{Y}$ is constructible.

For every invertible $\sh{O}_{\mathrm{U}}$-module $\mathscr{L}$ and every $x\in\mathrm{Y}$, it is enough to establish the existence of an open neighbourhood $\mathrm{V}$ of $x$ in $\mathrm{X}$ such that: $1^{\mathrm{o}}$ the canonical homomorphism $\sh{O}_{\mathrm{V}}\to(j_{\mathrm{V}})_*(\sh{O}_{\mathrm{U}\cap\mathrm{V}})$ is surjective; and $2^{\mathrm{o}}$ there is an invertible $\sh{O}_{\mathrm{V}}$-module $\mathscr{L}'_{\mathrm{V}}$ whose restriction to $\mathrm{U}\cap\mathrm{V}$ is isomorphic to $\mathscr{L}|(\mathrm{U}\cap\mathrm{V})$ (\sref{IV.21.13.5} and \sref{IV.21.13.3}).

Set $s=f(x)$. We first reduce to the case $\mathrm{S}=\mathrm{S}_0=\operatorname{Spec}(\sh{O}_{\mathrm{S},s})$. Let $(\mathrm{S}_\nu)$ be a fundamental system of affine open neighbourhoods of $s$ in $\mathrm{S}$, and put $\mathrm{X}_\nu=f^{-1}(\mathrm{S}_\nu)$, $\mathrm{Y}_\nu=\mathrm{Y}\cap\mathrm{X}_\nu$, $\mathrm{U}_\nu=\mathrm{X}_\nu-\mathrm{Y}_\nu$, with canonical injection $j_\nu:\mathrm{U}_\nu\to\mathrm{X}_\nu$. Then $\mathrm{X}_0=\mathrm{X}\times_{\mathrm{S}}\mathrm{S}_0$ is the projective limit of the $\mathrm{X}_\nu$ $(\mathbf{8.1.2},\,a))$. Put $\mathrm{Y}_0=\mathrm{Y}\cap\mathrm{X}_0$, $\mathrm{U}_0=\mathrm{X}_0-\mathrm{Y}_0$, and let $j_0:\mathrm{U}_0\to\mathrm{X}_0$ be the canonical injection. If $p_\nu:\mathrm{X}_0\to\mathrm{X}_\nu$ denotes the affine projection, then
\[
  (j_0)_*(\sh{O}_{\mathrm{U}_0})
  =p_\nu^*((j_\nu)_*(\sh{O}_{\mathrm{U}_\nu}))
  \qquad (\mathbf{II},\ 1.5.2),
\]
and the canonical homomorphism $\rho_0:\sh{O}_{\mathrm{X}_0}\to(j_0)_*(\sh{O}_{\mathrm{U}_0})$ is $p_\nu^*(\rho_\nu)$. Since $\rho_0$ is surjective by hypothesis, $\rho_\nu$ is surjective for sufficiently large $\nu$ \sref{IV.8.5.7}.

Let $\mathscr{L}_0$ be the restriction of $\mathscr{L}$ to $\mathrm{U}_0$. The $\mathrm{X}_\nu$ are affine, hence quasi-compact and quasi-separated. By hypothesis there exists an invertible $\sh{O}_{\mathrm{X}_0}$-module $\mathscr{L}'_0$ extending $\mathscr{L}_0$. For sufficiently large $\nu$ there is a quasi-coherent $\sh{O}_{\mathrm{X}_\nu}$-module $\mathscr{L}'_\nu$ of finite presentation such that $\mathscr{L}'_0\simeq p_\nu^*(\mathscr{L}'_\nu)$ \sref{IV.8.5.2}; by \sref{IV.8.5.5} we may suppose that $\mathscr{L}'_\nu$ is invertible. Finally, since the $\mathrm{U}_\nu$ are quasi-compact and quasi-separated and $\mathscr{L}'_0|\mathrm{U}_0\simeq\mathscr{L}_0$, for sufficiently large $\nu$ the modules $\mathscr{L}'_\nu|\mathrm{U}_\nu$ and $\mathscr{L}|\mathrm{U}_\nu$ are isomorphic \sref{IV.8.5.2.5}. This proves the reduction.

\oldpage[IV-4]{327}
We are thus reduced to the case $\mathrm{S}=\operatorname{Spec}(\mathrm{A})$, $\mathrm{X}=\operatorname{Spec}(\mathrm{B})$, where $\mathrm{A}$ is local. Since $\mathrm{X}$ is normal and $f$ smooth, $\mathrm{S}$ is normal \sref{IV.17.5.7}; hence $\mathrm{A}$ is integral and integrally closed. Regard $\mathrm{A}$ as the filtered direct limit of its $\bb{Z}$-subalgebras of finite type. The integral closure of such a subalgebra is a subring of $\mathrm{A}$ and is again a $\bb{Z}$-algebra of finite type $(\mathbf{0},\ 7.8.3,\ \text{(ii), (iii), and (vi)})$. Thus $\mathrm{A}$ is the filtered direct limit of integrally closed $\bb{Z}$-subalgebras $\mathrm{A}_\lambda$ of finite type. By \sref{IV.1.8.4.2}, for some $\lambda$ there is an $\mathrm{A}_\lambda$-algebra $\mathrm{B}_\lambda$ of finite type such that, up to an $\mathrm{A}$-isomorphism,
\[
  \mathrm{B}=\mathrm{B}_\lambda\otimes_{\mathrm{A}_\lambda}\mathrm{A}.
\]
Put $\mathrm{S}_\lambda=\operatorname{Spec}(\mathrm{A}_\lambda)$ and $\mathrm{X}_\lambda=\operatorname{Spec}(\mathrm{B}_\lambda)$. If $p_\lambda:\mathrm{X}\to\mathrm{X}_\lambda$ is the canonical projection, we may also suppose, since $\mathrm{Y}$ is constructible, that $\mathrm{Y}=p_\lambda^{-1}(\mathrm{Y}_\lambda)$ for a closed subset $\mathrm{Y}_\lambda$ of $\mathrm{X}_\lambda$ \sref{IV.8.3.11}. Write $f_\lambda:\mathrm{X}_\lambda\to\mathrm{S}_\lambda$ for the structural morphism. By transitivity of fibres and \sref{IV.4.2.6}, after increasing $\lambda$ we may suppose that $\mathrm{Y}_\lambda$ contains no irreducible component of any fibre of $f_\lambda$. Since $f$ is smooth, we may likewise suppose that $f_\lambda$ is smooth \sref{IV.17.7.8}. Finally, the image $\eta_\lambda$ in $\mathrm{S}_\lambda$ of the generic point $\eta$ of $\mathrm{S}$ is the generic point of $\mathrm{S}_\lambda$, and
\[
  \operatorname{codim}((\mathrm{Y}_\lambda)_{\eta_\lambda},
  (\mathrm{X}_\lambda)_{\eta_\lambda})
  =\operatorname{codim}(\mathrm{Y}_\eta,\mathrm{X}_\eta)\geq2
\]
by \sref{IV.6.1.4}. Thus $\mathrm{S}_\lambda$, $\mathrm{X}_\lambda$, and $\mathrm{Y}_\lambda$ satisfy all the hypotheses of (ii), as do their base changes $\mathrm{S}_\mu$, $\mathrm{X}_\mu$, and $\mathrm{Y}_\mu$ for $\lambda\leq\mu$.

If $(\mathrm{X}_\mu,\mathrm{Y}_\mu)$ is parafactorial for all sufficiently large $\mu$, then so is $(\mathrm{X},\mathrm{Y})$. Indeed, put $\mathrm{U}=\mathrm{X}-\mathrm{Y}$, $\mathrm{U}_\mu=\mathrm{X}_\mu-\mathrm{Y}_\mu$, and let $j$ and $j_\mu$ be the corresponding open immersions. Since $p_\mu:\mathrm{X}\to\mathrm{X}_\mu$ is affine,
\[
  j_*(\sh{O}_{\mathrm{U}})=p_\mu^*((j_\mu)_*(\sh{O}_{\mathrm{U}_\mu}))
  \qquad (\mathbf{II},\ 1.5.2),
\]
so the canonical map $\sh{O}_{\mathrm{X}}\to j_*(\sh{O}_{\mathrm{U}})$ is the pullback of the corresponding map on $\mathrm{X}_\mu$, and is therefore bijective. Moreover, for every invertible $\sh{O}_{\mathrm{U}}$-module $\mathscr{L}$, there is, for sufficiently large $\mu$, an invertible $\sh{O}_{\mathrm{U}_\mu}$-module $\mathscr{L}_\mu$ such that $\mathscr{L}\simeq(p'_\mu)^*(\mathscr{L}_\mu)$ \sref{IV.8.5.2} and \sref{IV.8.5.5}; here we use that $\mathrm{U}_\lambda$ is quasi-compact because $\mathrm{X}_\lambda$ is Noetherian. By parafactoriality, $\mathscr{L}_\mu$ extends to an invertible $\sh{O}_{\mathrm{X}_\mu}$-module $\mathscr{L}'_\mu$, whose pullback extends $\mathscr{L}$. Hence it remains only to prove (ii) when $\mathrm{A}$ is an integrally closed $\bb{Z}$-algebra of finite type.

The local rings $\sh{O}_{\mathrm{S},s}$ are then excellent integrally closed rings, and their completions are integrally closed $(\mathbf{0},\ 7.8.3,\ \text{(ii), (iii), and (vii)})$. By \sref{IV.21.13.10}, it is enough to show that $\sh{O}_{\mathrm{X},x}$ is parafactorial for every $x\in\mathrm{Y}$. Choose a closed point $y$ of $\mathrm{Y}_{f(x)}$ which is a specialization of $x$ \sref{IV.5.1.11}, and put $s=f(x)=f(y)$. The completion of $\sh{O}_{\mathrm{S},s}$ is integrally closed, and $\sh{O}_{\mathrm{X},y}$ is a formally smooth $\sh{O}_{\mathrm{S},s}$-algebra for the preadic topologies \sref{IV.17.5.3}. Since $\mathrm{Y}$ contains no irreducible component of $\mathrm{X}_s$,
\[
  \dim(\sh{O}_{\mathrm{X},y})>\dim(\sh{O}_{\mathrm{S},s})
  \qquad\text{by \sref{IV.17.5.8}}.
\]
If $s\neq\eta$, then $\dim(\sh{O}_{\mathrm{S},s})\geq1$; if $s=\eta$, the hypothesis gives $\dim(\sh{O}_{\mathrm{X},x})\geq2$. In either case $\dim(\sh{O}_{\mathrm{X},y})\geq2$. Furthermore, the prime ideal of $\sh{O}_{\mathrm{X},y}$ corresponding to $x$ is not contained in $\mathfrak{m}_s\sh{O}_{\mathrm{X},y}$, because $\mathrm{Y}$ contains no irreducible component of $\mathrm{X}_s$. Finally $y$ is closed in $\mathrm{X}_s$, so $k(y)$ is a finite extension of $k(s)$ $(\mathbf{I},\ 6.4.2)$. We may therefore apply \sref{IV.21.14.2} with $\mathrm{A}=\sh{O}_{\mathrm{S},s}$, $\mathrm{B}=\sh{O}_{\mathrm{X},y}$, and $\mathrm{B}_{\mathfrak{q}}=\sh{O}_{\mathrm{X},x}$. This completes the proof.
\end{proof}

\begin{env}[21.14.4]
\label{IV.21.14.4}
\oldpage[IV-4]{328}
\emph{Remarks}. ---
(i) It is possible that the statements \sref{IV.21.14.1} and \sref{IV.21.14.2} remain valid \emph{without hypothesis} on the residue field of $\mathrm{B}$.
The statement with this generality would be equivalent, by virtue of \sref{IV.21.13.15} and \sref{IV.21.13.12.1}, to the following: let $\mathrm{A}$ be a complete, integral, integrally closed Noetherian local ring, and let $\mathrm{B}$ be a Noetherian local ring which is a formally smooth $\mathrm{A}$-algebra, such that $\dim(\mathrm{B})>\dim(\mathrm{A})$; then $\mathrm{B}$ is parafactorial.

(ii) We will see in chap.~VI, using a technique of ``finite descent'' applied to the morphism $\mathrm{Y}'\to\mathrm{Y}=\operatorname{Spec}(\mathrm{A})$, where $\mathrm{Y}'$ is the normalization of $\mathrm{Y}$, that the conclusion of \sref{IV.21.14.1} (or of \sref{IV.21.14.2}) remains valid when the hypothesis that $\widehat{\mathrm{A}}$ is integrally closed is replaced by the hypothesis that $\widehat{\mathrm{A}}$ is reduced, provided that $\dim\mathrm{B}\geq\dim\mathrm{A}+2$.
If this last condition is replaced by $\dim\mathrm{B}\geq\dim\mathrm{A}+3$, the hypothesis that $\widehat{\mathrm{A}}$ is reduced may even be omitted.
Similarly, the conclusion of \sref{IV.21.14.3} remains valid by replacing the hypothesis that $\mathrm{X}$ is normal by the hypothesis that $\mathrm{X}$ is reduced, provided that $\dim(\sh{O}_{\mathrm{X},x})\geq\dim(\sh{O}_{\mathrm{S},f(x)})+2$.

(iii) In chap.~III, Part~3, we prove that, if $f:\mathrm{X}\to\mathrm{S}$ is a smooth morphism and $\mathrm{Y}$ is a closed locally constructible subset of $\mathrm{X}$ such that, for every $s\in\mathrm{S}$, $\operatorname{codim}(\mathrm{Y}_s,\mathrm{X}_s)\geq3$ (with the notation of \sref{IV.21.14.3}), then the pair $(\mathrm{X},\mathrm{Y})$ is parafactorial.
This conclusion is no longer valid when we only suppose $\operatorname{codim}(\mathrm{Y}_s,\mathrm{X}_s)\geq 2$ for all $s\in\mathrm{S}$ and when we do not suppose $\mathrm{X}$ reduced.
For example, let $k$ be a field, let $\mathrm{A}=k[\mathrm{T}]/(\mathrm{T}^2)$ be the algebra of dual numbers over $k$, and put $\mathrm{S}=\operatorname{Spec}(\mathrm{A})$ and $\mathrm{X}=\operatorname{Spec}(\mathrm{A}[\mathrm{T}_1,\mathrm{T}_2])=\mathbb{V}_{\mathrm{S}}^2$, where $\mathrm{T}_1$ and $\mathrm{T}_2$ are indeterminates; let $\mathrm{Y}$ be the ``zero section'' of this fibration. If $s$ is the unique closed point of $\mathrm{S}$, then $\mathrm{X}_s=\mathbb{V}_k^2$ and $\mathrm{Y}_s$ is the ``zero section'' of $\mathrm{X}_s$.
To see that the pair $(\mathrm{X},\mathrm{Y})$ is not parafactorial, it suffices to show that the ring $\mathrm{B}=\mathrm{A}[\mathrm{T}_1,\mathrm{T}_2]_\mathfrak{m}$, where $\mathfrak{m}$ is the ideal $(\mathrm{T}_1)+(\mathrm{T}_2)$ which defines $\mathrm{Y}$, is not parafactorial \sref{IV.21.13.10}.
Let $\mathrm{B}_0=\mathrm{B}_{\mathrm{red}}$, and let $\mathrm{U}$ and $\mathrm{U}_0$ be the complementary open subsets of the closed point in $\operatorname{Spec}(\mathrm{B})$ and $\operatorname{Spec}(\mathrm{B}_0)$; reasoning as in \sref{IV.21.13.9}, we have the exact sequence, extension of \sref{IV.21.13.9.1}
\[
  \Gamma(\mathrm{U},\sh{O}_{\mathrm{U}})^*\to\Gamma(\mathrm{U}_0,\sh{O}_{\mathrm{U}_0})^*\to\mathrm{H}^1(\mathrm{U}_0,\sh{O}_{\mathrm{U}_0})\to\operatorname{Pic}(\mathrm{U})\to\operatorname{Pic}(\mathrm{U}_0).
\]
Now, one has $\Gamma(\mathrm{U},\sh{O}_{\mathrm{U}})=\mathrm{B}$, $\Gamma(\mathrm{U}_0,\sh{O}_{\mathrm{U}_0})=\mathrm{B}_0$ since $\operatorname{prof}(\mathrm{B})=\operatorname{prof}(\mathrm{B}_0)=2$ \sref{IV.5.10.5}.
On the other hand $\operatorname{Pic}(\mathrm{U}_0)=0$ since $\mathrm{B}_0$ is factorial, and $\mathrm{H}^1(\mathrm{U}_0,\sh{O}_{\mathrm{U}_0})\neq 0$ \cite[3, Example~III-1]{IV-41}; since the homomorphism $\mathrm{B}^*\to\mathrm{B}_0^*$ is surjective, we conclude that $\operatorname{Pic}(\mathrm{U})\neq 0$, hence that $\mathrm{B}$ is not parafactorial.
\oldpage[IV-4]{329}
(iv) The result of \sref{IV.21.14.3} was first proved by C.~Seshadri \cite{IV-44} in the particular case where $\mathrm{S}$ is an algebraic normal scheme over an algebraically closed field $k$ and $\mathrm{X}=\mathrm{S}\times_k\mathrm{T}$, where $\mathrm{T}$ is a $k$-algebraic prescheme smooth over $k$.
The proof of Seshadri \cite[pp.~188--189]{IV-44} is of global nature and uses the theory of Picard schemes.
It gives moreover (\emph{loc.\ cit.}) results such as the following (for which one does not currently possess a proof by local methods).
Let $\mathrm{S}$ and $\mathrm{T}$ be two preschemes locally of finite type over a field $k$, put $\mathrm{X}=\mathrm{S}\times_k\mathrm{T}$, and let $\mathrm{Z}$ be a $1$-codimensional cycle on $\mathrm{X}$ (regarded as an $\mathrm{S}$-prescheme). Suppose that the following conditions hold:
\begin{enumerate}[label=\arabic*\textsuperscript{o}]
  \item $\mathrm{S}$ and $\mathrm{T}$ are geometrically normal over $k$ \sref{IV.6.7.6};
  \item For every maximal point $\eta$ of $\mathrm{S}$, the $1$-codimensional cycle $\mathrm{Z}_\eta$ on the fibre $\mathrm{X}_\eta$, having the same multiplicity as $\mathrm{Z}$ at every point of $\mathrm{X}^{(1)}\cap\mathrm{X}_\eta$, is locally principal (in other words, it is the image of a divisor of $\mathrm{X}_\eta$, since $\mathrm{X}_\eta$ is normal);
  \item For all $s\in\mathrm{S}$, $\mathrm{Z}$ is principal at the maximal points of the fibre $\mathrm{X}_s$.
\end{enumerate}
Then $\mathrm{Z}$ is locally principal.
In other words, since $\mathrm{X}$ is normal \sref{IV.6.14.1}, for every $x\in\mathrm{X}$ which is not maximal in its fibre $\mathrm{X}_s$ and which belongs to none of the ``generic fibres'' $\mathrm{X}_\eta$ (which implies $\dim(\sh{O}_{\mathrm{X},x})\geq2$ by \sref{IV.6.1.1}), the local ring $\sh{O}_{\mathrm{X},x}$ is parafactorial, by \sref{IV.21.13.15},~(b).\footnote{In the cited paper Seshadri in fact assumes that $k$ is algebraically closed and that $\mathrm{T}$ is separated and ``semi-complete'' (that is, $\Gamma(\mathrm{T},\sh{O}_{\mathrm{T}})\simeq k$), and replaces condition $3^{\mathrm{o}}$ by the stronger hypothesis that $\operatorname{Supp}(\mathrm{Z})$ contains no fibre $\mathrm{X}_s$, $s\in\mathrm{S}$. Since the statement is local on $\mathrm{S}$, however, condition $3^{\mathrm{o}}$ suffices. The resulting conclusion, interpreted as above in terms of parafactoriality of the rings $\sh{O}_{\mathrm{X},x}$, is local on both $\mathrm{S}$ and $\mathrm{T}$; this removes altogether the assumptions that $\mathrm{T}$ is semi-complete and that $k$ is algebraically closed. Indeed, after first passing to an algebraic closure of $k$, one may assume that $\mathrm{T}$ is affine, embed it as an open subscheme of a normal projective prescheme over $k$, and apply Seshadri's result. This reduction also shows that it is enough to prove Seshadri's theorem when $\mathrm{T}$ is projective (rather than merely semi-complete), the case in which the Picard-scheme theory used by Seshadri is contained in chap.~VI of this Treatise.}
\end{env}

\subsection{Relative divisors}
\label{subsection:IV.21.15}
\label{IV.21.15}

\begin{env}[21.15.1]
\label{IV.21.15.1}
Let $\mathrm{S}$ be a prescheme, $f:\mathrm{X}\to\mathrm{S}$ a \emph{flat} morphism and locally of finite presentation.
We defined in \sref{IV.20.6.1} the sheaf of rings $\mathscr{M}_{\mathrm{X}/\mathrm{S}}$ of the \emph{germs of meromorphic functions on $\mathrm{X}$ relative to $\mathrm{S}$}, sub-sheaf of $\mathscr{M}_{\mathrm{X}}$; it is clear that the canonical injection $\sh{O}_{\mathrm{X}}\hookrightarrow\mathscr{M}_{\mathrm{X}}$ \sref{IV.20.1.4.1} maps $\sh{O}_{\mathrm{X}}$ into a sub-sheaf of $\mathscr{M}_{\mathrm{X}/\mathrm{S}}$, with which we identify it.
Let $\mathscr{M}_{\mathrm{X}/\mathrm{S}}^*$ be the sheaf (of multiplicative groups) of germs of invertible sections of $\mathscr{M}_{\mathrm{X}/\mathrm{S}}$; it is therefore a sub-sheaf of $\mathscr{M}_{\mathrm{X}}^*$ and contains $\sh{O}_{\mathrm{X}}^*$ as a sub-sheaf.
\end{env}

\begin{definition}[21.15.2]
\label{IV.21.15.2}
We call \emph{sheaf of divisors on $\mathrm{X}$ relative to $\mathrm{S}$}, or \emph{sheaf of divisors on $\mathrm{X}$ transversal to $f$}, and denote $\mathscr{D}iv_{\mathrm{X}/\mathrm{S}}$ the quotient sheaf (of commutative groups) $\mathscr{M}_{\mathrm{X}/\mathrm{S}}^*/\sh{O}_{\mathrm{X}}^*$; the sections of this sheaf over $\mathrm{X}$ are called \emph{divisors relative to $\mathrm{S}$}, or \emph{divisors on $\mathrm{X}$ transversal to $f$}; they form a commutative group denoted $\operatorname{Div}(\mathrm{X}/\mathrm{S})$.
\end{definition}

It is clear that $\mathscr{D}iv_{\mathrm{X}/\mathrm{S}}$ identifies canonically with a sub-sheaf of $\mathscr{D}iv_{\mathrm{X}}$, and hence $\operatorname{Div}(\mathrm{X}/\mathrm{S})$ with a subgroup of $\operatorname{Div}(\mathrm{X})$, which we denote also \emph{additively}.

\oldpage[IV-4]{330}
For every open $\mathrm{U}$ of $\mathrm{X}$, we have $\mathscr{M}_{\mathrm{X}/\mathrm{S}}^*|\mathrm{U}=\mathscr{M}_{\mathrm{U}/\mathrm{S}}^*$, hence $\mathscr{D}iv_{\mathrm{X}/\mathrm{S}}|\mathrm{U}=\mathscr{D}iv_{\mathrm{U}/\mathrm{S}}$, and the sheaf $\mathscr{D}iv_{\mathrm{X}/\mathrm{S}}$ is therefore equal to the presheaf $\mathrm{U}\mapsto\operatorname{Div}(\mathrm{U}/\mathrm{S})$.

Since $\mathscr{M}_{\mathrm{X}/\mathrm{S}}^*$ is a sub-sheaf of $\mathscr{M}_{\mathrm{X}}^*$, the definitions, notations and formulas relative to divisors of sections of $\mathscr{M}_{\mathrm{X}}$ over $\mathrm{X}$ \sref{IV.21.1.3} apply without change to sections of $\mathscr{M}_{\mathrm{X}/\mathrm{S}}$ over $\mathrm{X}$.

\begin{env}[21.15.3]
\label{IV.21.15.3}
The structure of sheaf of ordered groups on $\mathscr{D}iv_{\mathrm{X}}$ \sref{IV.21.1.6} induces on the sub-sheaf $\mathscr{D}iv_{\mathrm{X}/\mathrm{S}}$ a structure of sheaf of ordered groups, for which the sheaf of monoids of germs of positive sections is $\mathscr{D}iv_{\mathrm{X}/\mathrm{S}}^+=\mathscr{D}iv_{\mathrm{X}}^+\cap\mathscr{D}iv_{\mathrm{X}/\mathrm{S}}$, which we denote $\mathscr{D}iv_{\mathrm{X}/\mathrm{S}}^+$.
We have $\Gamma(\mathrm{X},\mathscr{D}iv_{\mathrm{X}/\mathrm{S}}^+)=\operatorname{Div}^+(\mathrm{X})\cap\operatorname{Div}(\mathrm{X}/\mathrm{S})$; we denote this sub-monoid of $\operatorname{Div}(\mathrm{X}/\mathrm{S})$ by $\operatorname{Div}^+(\mathrm{X}/\mathrm{S})$, and it consists of elements $\geq 0$ with respect to the ordered group structure on $\operatorname{Div}(\mathrm{X}/\mathrm{S})$.
It follows from \sref{IV.21.5.1} that $\mathscr{D}iv_{\mathrm{X}/\mathrm{S}}^+$ is the image in $\mathscr{D}iv_{\mathrm{X}/\mathrm{S}}$ of the sub-sheaf of monoids
\[
\label{IV.21.15.3.1}
  \sh{O}_{\mathrm{X}/\mathrm{S}}^{\operatorname{reg}}=\sh{O}_{\mathrm{X}}\cap\mathscr{M}_{\mathrm{X}/\mathrm{S}}^*
  \tag{21.15.3.1}
\]
\end{env}

For every open $\mathrm{U}$ of $\mathrm{X}$, $\Gamma(\mathrm{U},\sh{O}_{\mathrm{X}/\mathrm{S}}^{\operatorname{reg}})$ is the set of sections $t$ of $\sh{O}_{\mathrm{X}}$ over $\mathrm{U}$ such that $t$ is regular and that $1/t$ belongs to $\Gamma(\mathrm{U},\mathscr{M}_{\mathrm{X}/\mathrm{S}})$, which means, with the notations of \sref{IV.20.6.1}, that $t\in\mathrm{T}_{\mathrm{X}/\mathrm{S}}(\mathrm{U})$, so that the sheaf $\sh{O}_{\mathrm{X}/\mathrm{S}}^{\operatorname{reg}}$ is none other than the sheaf denoted $\mathscr{T}_{\mathrm{X}/\mathrm{S}}$ in \sref{IV.20.6.1}.
We can therefore write
\[
\label{IV.21.15.3.2}
  \mathscr{D}iv_{\mathrm{X}/\mathrm{S}}^+=\sh{O}_{\mathrm{X}/\mathrm{S}}^{\operatorname{reg}}/\sh{O}_{\mathrm{X}}^*
  \tag{21.15.3.2}
\]
up to canonical isomorphism.

\begin{env}[21.15.3.3]
\label{IV.21.15.3.3}
Let $\mathrm{D}\in\operatorname{Div}^+(\mathrm{X}/\mathrm{S})$ and consider the closed sub-prescheme $\mathrm{Y}(\mathrm{D})$ defined by the ideal $\mathscr{J}_{\mathrm{X}}(\mathrm{D})$ of $\sh{O}_{\mathrm{X}}$ \sref{IV.21.6.5}; by virtue of what precedes, for all $x\in\mathrm{Y}(\mathrm{D})$, there is an open neighbourhood $\mathrm{U}$ of $x$ in $\mathrm{X}$ and a section $t\in\mathrm{T}_{\mathrm{X}/\mathrm{S}}(\mathrm{U})$ such that $\mathscr{J}_{\mathrm{X}}(\mathrm{D})|\mathrm{U}=\sh{O}_{\mathrm{U}}\cdot t$; since $x\in\mathrm{Y}(\mathrm{D})$, the image $t_{f(x)}$ of $t$ in $\Gamma(\mathrm{U}\cap\mathrm{X}_{f(x)},\sh{O}_{\mathrm{X}_{f(x)}})$ belongs to the maximal ideal of $\sh{O}_{\mathrm{X}_{f(x)},x}$, and furthermore, by definition, for all $s\in\mathrm{S}$, the image $t_s$ of $t$ in $\Gamma(\mathrm{U}\cap\mathrm{X}_s,\sh{O}_{\mathrm{X}_s})$ is a regular section.
We deduce therefore from \sref{IV.11.3.8} and \sref{IV.19.2.4} that the canonical injection $\mathrm{Y}(\mathrm{D})\to\mathrm{X}$ is \emph{transversely regular relative to $\mathrm{S}$ and of codimension $1$}.
The converse being immediate, we see that one can \emph{identify} canonically the positive divisors on $\mathrm{X}$ relative to $\mathrm{S}$ with the closed sub-preschemes $\mathrm{Y}$ of $\mathrm{X}$ such that the canonical injection $\mathrm{Y}\to\mathrm{X}$ is a transversely regular immersion relative to $\mathrm{S}$ and of codimension $1$.
We will ordinarily make this identification.
\end{env}

\begin{proposition}[21.15.4]
\label{IV.21.15.4}
Let $\mathrm{D}$ be a divisor on $\mathrm{X}$, and let $\mathscr{J}_{\mathrm{X}}(\mathrm{D})$ be the corresponding invertible fractional ideal \sref{IV.21.2.5}.
For $\mathrm{D}\in\operatorname{Div}(\mathrm{X}/\mathrm{S})$, it is necessary and sufficient that, for all $x\in\mathrm{X}$, $(\mathscr{J}_{\mathrm{X}}(\mathrm{D})\cap\mathscr{M}_{\mathrm{X}/\mathrm{S}}^*)_x\neq 0$ (or equivalently, that $\mathscr{J}_{\mathrm{X}}(\mathrm{D})\subset\mathscr{M}_{\mathrm{X}/\mathrm{S}}$ and $\mathscr{J}_{\mathrm{X}}(\mathrm{D})^{-1}\subset\mathscr{M}_{\mathrm{X}/\mathrm{S}}$).
\end{proposition}

\begin{proof}
Indeed, to say that $\mathrm{D}\in\operatorname{Div}(\mathrm{X}/\mathrm{S})$ means that for all $x\in\mathrm{X}$, there exists an open neighbourhood $\mathrm{U}$ of $x$ and a section $f\in\Gamma(\mathrm{U},\mathscr{M}_{\mathrm{X}/\mathrm{S}}^*)$ such that $\mathrm{D}|\mathrm{U}=\operatorname{div}_{\mathrm{U}}(f)$, where $f\in\Gamma(\mathrm{U},\mathscr{M}_{\mathrm{X}}^*)$; since $\mathscr{J}_{\mathrm{X}}(\mathrm{D})|\mathrm{U}$ is the fractional ideal $\sh{O}_{\mathrm{U}}\cdot f$, the proposition follows immediately.
\end{proof}

\oldpage[IV-4]{331}
\begin{proposition}[21.15.5]
\label{IV.21.15.5}
Let $\mathrm{D}$ be a divisor on $\mathrm{X}$, let $\sh{O}_{\mathrm{X}}(\mathrm{D})$ be the invertible fractional ideal, and let $s_{\mathrm{D}}$ be the regular meromorphic section of $\sh{O}_{\mathrm{X}}(\mathrm{D})$ canonically defined by $\mathrm{D}$ \sref{IV.21.2.8} and \sref{IV.21.2.9}.
For $\mathrm{D}\in\operatorname{Div}(\mathrm{X}/\mathrm{S})$, it is necessary and sufficient that $s_{\mathrm{D}}\in\Gamma(\mathrm{X},\mathscr{M}_{\mathrm{X}/\mathrm{S}}(\sh{O}_{\mathrm{X}}(\mathrm{D})))$.
\end{proposition}

\begin{proof}
Indeed, if $\mathrm{U}$ is an open of $\mathrm{X}$ such that $\mathrm{D}|\mathrm{U}=\operatorname{div}_{\mathrm{U}}(f)$, where $f\in\Gamma(\mathrm{U},\mathscr{M}_{\mathrm{X}}^*)$, to say that $s_{\mathrm{D}}\in\Gamma(\mathrm{U},\mathscr{M}_{\mathrm{X}/\mathrm{S}}(\sh{O}_{\mathrm{X}}(\mathrm{D})))$ means, by virtue of the definitions \sref{IV.20.6.2}, that $f^{-1}\in\Gamma(\mathrm{U},\mathscr{M}_{\mathrm{X}/\mathrm{S}})$, whence the proposition.
\end{proof}

The interpretation of the divisors on $\mathrm{X}$ by means of the classes $\operatorname{cl}(\mathscr{L},s)$ \sref{IV.21.2.11} allows us therefore to interpret the elements of $\operatorname{Div}(\mathrm{X}/\mathrm{S})$ as the pairs $(\mathscr{L},s)$ (up to isomorphism) such that $\mathscr{L}$ is an invertible $\sh{O}_{\mathrm{X}}$-module and $s$ is a meromorphic section of $\mathscr{L}$ regular relative to $\mathrm{S}$ \sref{IV.20.6.5},~(iii)).

\begin{proposition}[21.15.6]
\label{IV.21.15.6}
Let $\mathrm{D}$ be a divisor on $\mathrm{X}$ relative to $\mathrm{S}$, and suppose that for all $x\in\mathrm{X}$ such that $\operatorname{prof}(\sh{O}_{\mathrm{X}_{f(x)},x})=1$, we have $\mathrm{D}_x\geq 0$ (resp.\ $\mathrm{D}_x=0$).
Then $\mathrm{D}\geq 0$ (resp.\ $\mathrm{D}=0$).
\end{proposition}

\begin{proof}
As in \sref{IV.21.1.8}, we may restrict to the case where $\mathrm{D}=\operatorname{div}(\varphi)$, with $\varphi$ a meromorphic function regular relative to $\mathrm{S}$. It is then enough to show that, if $\mathrm{D}_x\geq0$ at every point $x\in\mathrm{X}$ for which $\operatorname{prof}(\sh{O}_{\mathrm{X}_{f(x)},x})=1$, then $\varphi$ is defined everywhere on $\mathrm{X}$.
This hypothesis means that, if $\mathrm{T}=\mathrm{X}-\operatorname{dom}(\varphi)$, then $\operatorname{prof}(\sh{O}_{\mathrm{X}_{f(x)},x})\geq2$ for every $x\in\mathrm{T}$; the conclusion follows from \sref{IV.20.6.6}.
\end{proof}

\begin{env}[21.15.7]
\label{IV.21.15.7}
Let $\mathrm{X}'$ be a second S-prescheme flat and locally of finite presentation over $\mathrm{S}$, and let $g:\mathrm{X}'\to\mathrm{X}$ be an S-morphism.
If the S-morphism $g$ is \emph{flat}, we know \sref{IV.21.4.5} that the inverse image by $g$ of every divisor on $\mathrm{X}$ is defined; if furthermore $\mathrm{D}\in\operatorname{Div}(\mathrm{X}/\mathrm{S})$, it follows from the definition \sref{IV.21.15.2} and from \sref{IV.20.6.8} that $g^*(\mathrm{D})\in\operatorname{Div}(\mathrm{X}'/\mathrm{S})$.
\end{env}

\begin{env}[21.15.8]
\label{IV.21.15.8}
Let $\mathrm{X}'$ be a second S-prescheme flat and locally of finite presentation over $\mathrm{S}$, and let $g:\mathrm{X}'\to\mathrm{X}$ be an S-morphism \emph{finite} and \emph{flat}.
We note that $g$ is also necessarily of finite presentation \sref{IV.1.4.3},~\sref{IV.1.4.6} and \sref{IV.1.6.3}, hence $g_*(\sh{O}_{\mathrm{X}'})$ is a $\sh{O}_{\mathrm{X}}$-Module flat and of finite presentation, and therefore \emph{locally free} \sref{IV.2.1.12}, in other words $g$ is a locally free morphism \sref{IV.18.2.7}; for all $s\in\mathrm{S}$, the corresponding morphism $g_s:\mathrm{X}_s'\to\mathrm{X}_s$ is also finite and locally free.
We deduce from \sref{IV.21.5.2} and \sref{IV.20.6.1} that, for every open $\mathrm{U}$ of $\mathrm{X}$ and every section $t'\in\Gamma(g^{-1}(\mathrm{U}),\mathscr{T}_{\mathrm{X}'/\mathrm{S}})$, the norm $\mathrm{N}_{\mathrm{X}'/\mathrm{X}}(t')$ belongs to $\Gamma(\mathrm{U},\mathscr{T}_{\mathrm{X}/\mathrm{S}})$. The argument of \sref{IV.21.5.3} then shows that, for every invertible $\sh{O}_{\mathrm{X}'}$-module $\mathscr{L}'$ and every meromorphic section $u'$ of $\mathscr{L}'$ over $\mathrm{X}'$ which is regular relative to $\mathrm{S}$, the norm $\mathrm{N}_{\mathrm{X}'/\mathrm{X}}(u')$ is a meromorphic section of $\mathscr{L}=\mathrm{N}_{\mathrm{X}'/\mathrm{X}}(\mathscr{L}')$ over $\mathrm{X}$, regular relative to $\mathrm{S}$.
The interpretation of divisors relative to $\mathrm{S}$ given in \sref{IV.21.15.5}, together with the definition of the direct image of a divisor \sref{IV.21.5.5}, therefore proves that \emph{for every divisor} $\mathrm{D}'\in\operatorname{Div}(\mathrm{X}'/\mathrm{S})$, one has $g_*(\mathrm{D}')\in\operatorname{Div}(\mathrm{X}/\mathrm{S})$.
\end{env}

\begin{env}[21.15.9]
\label{IV.21.15.9}
\oldpage[IV-4]{332}
Consider finally an arbitrary morphism $\mathrm{S}'\to\mathrm{S}$, and (under the hypotheses of \sref{IV.21.15.1}) set $\mathrm{X}'=\mathrm{X}\times_{\mathrm{S}}\mathrm{S}'$, which is flat and locally of finite presentation over $\mathrm{S}'$; if $p:\mathrm{X}'\to\mathrm{X}$ is the canonical projection, we have seen \sref{IV.20.6.9} that there is a canonical homomorphism $p^*(\mathscr{M}_{\mathrm{X}/\mathrm{S}})\to\mathscr{M}_{\mathrm{X}'/\mathrm{S}'}$, which transforms every section of $\mathscr{M}_{\mathrm{X}/\mathrm{S}}$ over $\mathrm{U}$, \emph{regular relative to $\mathrm{S}$}, into a section of $\mathscr{M}_{\mathrm{X}'/\mathrm{S}'}$ regular relative to $\mathrm{S}'$ \sref{IV.20.6.5},~(iii); we conclude then from the definition \sref{IV.21.15.2} and the right exactness of the functor $p^*$, that the preceding homomorphism defines by passage to quotients a canonical homomorphism
\[
\label{IV.21.15.9.1}
  \operatorname{Div}(\mathrm{X}/\mathrm{S})\to\operatorname{Div}(\mathrm{X}'/\mathrm{S}')
  \tag{21.15.9.1}
\]
which evidently transforms elements of $\operatorname{Div}^+(\mathrm{X}/\mathrm{S})$ into elements of $\operatorname{Div}^+(\mathrm{X}'/\mathrm{S}')$.
We also set $\operatorname{Div}(\mathrm{X}'/\mathrm{S}')=\operatorname{Div}_{\mathrm{X}/\mathrm{S}}(\mathrm{S}')$ (resp.\ $\operatorname{Div}^+(\mathrm{X}'/\mathrm{S}')=\operatorname{Div}_{\mathrm{X}/\mathrm{S}}^+(\mathrm{S}')$), and we see immediately that we have thus defined two contravariant functors
\[
  \operatorname{Div}_{\mathrm{X}/\mathrm{S}}:\mathbf{Sch}_{/\mathrm{S}}\to\mathbf{Ab},\qquad\operatorname{Div}_{\mathrm{X}/\mathrm{S}}^+:\mathbf{Sch}_{/\mathrm{S}}\to\mathbf{Ens}
\]
from the category of S-preschemes into that of commutative groups (resp.\ sets).
We will see further (chap.~VI) the important cases where the functor $\operatorname{Div}_{\mathrm{X}/\mathrm{S}}^+$ is \emph{representable} $(\mathbf{0}_{\mathrm{III}}$, $8.1.8)$.

For every divisor $\mathrm{D}\in\operatorname{Div}(\mathrm{X}/\mathrm{S})$, the image of $\mathrm{D}$ by the homomorphism \sref{IV.21.15.9.1} is none other than the \emph{inverse image} $p^*(\mathrm{D})$ (in the sense of \sref{IV.21.4.2}): the existence of this inverse image and the preceding assertions are immediate consequences of \sref{IV.20.6.5},~(iii) and \sref{IV.20.6.9}.

The elements of $\operatorname{Div}(\mathrm{X}'/\mathrm{S}')$ are often called, by abuse of language, ``\emph{families of divisors on $\mathrm{X}$ relative to $\mathrm{S}$, parametrised by the S-prescheme $\mathrm{S}'$}''; this terminology is used especially when it concerns positive divisors.
\end{env}
% END NATIVE INPUT source_aligned/ega4-21

% BEGIN NATIVE INPUT source_aligned/ega4-backmatter-bibliography
% Source-aligned English projection of EGA IV, Part 4, p. 333.
% Entries [32], [34], and [37] are carried forward from Parts 1 and 2
% because they are cited in the cumulative reader. The final two entries
% close explicit cross-volume citations to EGA I and SGA 1.
\clearpage
\oldpage[IV]{333}
\phantomsection
\addcontentsline{toc}{section}{Bibliography}
\section*{Bibliography}

\begin{thebibliography}{SGA 1}

\bibitem[32]{IV4-print-IV-32}
A.~N\'eron, \emph{Mod\`eles minimaux des vari\'et\'es ab\'eliennes sur les
corps locaux et globaux}, \emph{Publications Math\'ematiques de l'IH\'ES}
\textbf{21} (1964).

\bibitem[34]{IV4-print-IV-34}
H.~Hironaka, forthcoming article on \emph{Projective morphisms}.

\bibitem[37]{IV4-print-IV-37}
H.~Cartan, \emph{S\'eminaire de l'\'Ecole Normale Sup\'erieure},
13th year (1960--1961).

\bibitem[41]{IV4-print-IV-41}
A.~Grothendieck, \emph{S\'eminaire de G\'eom\'etrie alg\'ebrique, 1962:
Cohomologie locale des faisceaux coh\'erents, et th\'eor\`emes de Lefschetz
locaux et globaux}, Institut des Hautes \'Etudes Scientifiques,
Bures-sur-Yvette.

\bibitem[42]{IV4-print-IV-42}
M.~Demazure and A.~Grothendieck, \emph{S\'eminaire de G\'eom\'etrie
alg\'ebrique, 1963--1964: Sch\'emas en groupes}, Institut des Hautes
\'Etudes Scientifiques, Bures-sur-Yvette.

\bibitem[43]{IV4-print-IV-43}
M.~Artin and A.~Grothendieck, \emph{S\'eminaire de G\'eom\'etrie
alg\'ebrique, 1963--1964: Cohomologie \'etale des sch\'emas}, Institut des
Hautes \'Etudes Scientifiques, Bures-sur-Yvette.

\bibitem[44]{IV4-print-IV-44}
C.~S.~Seshadri, ``Quotient space by an abelian variety,''
\emph{Mathematische Annalen} \textbf{152} (1962), 185--194.

\bibitem[I-4]{IV4-print-I-4}
R.~Godement, \emph{Th\'eorie des faisceaux}, Actualit\'es Scientifiques et
Industrielles, no.~1252, Paris (Hermann), 1958.

\bibitem[SGA 1]{IV4-print-SGA1}
A.~Grothendieck (ed.), \emph{Rev\^etements \'etales et groupe fondamental
(SGA~1)}, Lecture Notes in Mathematics 224, Springer-Verlag, 1971.

\end{thebibliography}

\begin{quote}\footnotesize
Entries [32] and [34] reproduce the continued bibliographies printed in
Parts~1 and~2. The EGA~I and SGA~1 entries are included so that the two
explicit cross-volume citations resolve in this standalone cumulative reader.
\end{quote}
% END NATIVE INPUT source_aligned/ega4-backmatter-bibliography
% BEGIN NATIVE INPUT source_aligned/ega4-backmatter-notation
% Source-aligned English projection of EGA IV, Part 4, pp. 334--335.
\clearpage
\oldpage[IV]{334}
\phantomsection
\addcontentsline{toc}{section}{Index of notation}
\section*{Index of notation}

\begingroup
\small
\setlength{\parindent}{0pt}
\setlength{\parskip}{1.5pt}
\newcommand{\EGAIVidxref}[1]{\hyperref[IV.#1]{#1}}

$Y_f^{(n)}$, $Y^{(n)}$, $\mathcal{G}r_{\bullet}(f)$,
$\mathcal{G}r_n(f)$, $\mathcal{N}_{Y/X}$
($f:Y\to X$ a morphism of ringed spaces): {\hyperlink{ega4.c04.def.16-1-2}{16.1.2}}.

$\mathcal{O}_{Y^{(\infty)}}$: {\hyperlink{ega4.c04.rem.16-1-11}{16.1.11}}.

$\mathcal{P}^{n}_{f}$, $\mathcal{P}^{n}_{X/S}$,
$\mathcal{P}^{\infty}_{f}$, $\mathcal{P}^{\infty}_{X/S}$,
$\Omega^{1}_{f}$, $\Omega^{1}_{X/S}$
($f:X\to S$ a morphism of preschemes): {\hyperlink{ega4.c04.def.16-3-1}{16.3.1}}.

$d^{n}_{f}$, $d^{n}_{X/S}$, $d^{\infty}_{f}$,
$d^{\infty}_{X/S}$, $d^{n}$, $d^{\infty}$, $dt$, $d_{X/S}(t)$
($f:X\to S$ a morphism of preschemes): {\hyperlink{ega4.c04.def.16-3-6}{16.3.6}}.

$P^{n}(u)$, $\Gr_{\bullet}(u)$, $\Gr_{1}(u)$
($u:X'\to X$ a morphism of preschemes): {\hyperlink{ega4.c04.statement.16-4-3}{16.4.3}}.

$\operatorname{Der}_{S}(\mathcal{O}_{X},\mathcal{F})$:
{\hyperlink{ega4.c04.statement.16-5-1}{16.5.1}}.

$\mathcal{D}er_{S}(\mathcal{O}_{X},\mathcal{F})$:
{\hyperlink{ega4.c04.statement.16-5-4}{16.5.4}}.

$\mathcal{T}_{X/S}$: {\hyperlink{ega4.c04.statement.16-5-7}{16.5.7}}.

$D_i$, $\partial/\partial s_i$: {\hyperlink{ega4.c04.statement.16-5-8}{16.5.8}}.

$T_{X/S}$: {\hyperlink{ega4.c04.statement.16-5-12}{16.5.12}}.

$T_{X/S}(x)$, $T_Y(s)$: {\hyperlink{ega4.c04.statement.16-5-13}{16.5.13}}.

$\Omega^p_{X/S}$, $\Omega^\bullet_{X/S}$, $\Omega^p_{B/A}$:
{\hyperlink{ega4.c04.statement.16-6-1}{16.6.1}}.

$d$, $d_{X/S}$: {\hyperlink{ega4.c04.def.16-6-3}{16.6.3}}.

$\mathcal{P}^n_{X/S}(\mathcal{F})$: {\hyperlink{ega4.c04.statement.16-7-2}{16.7.2}}.

$d^n_{X/S,\mathcal{F}}$, $d^n_{X/S}$: {\hyperlink{ega4.c04.statement.16-7-5}{16.7.5}}.

$\mathcal{P}^{\infty}_{X/S}(\mathcal{F})$,
$d^{\infty}_{X/S,\mathcal{F}}$: {\hyperlink{ega4.c04.statement.16-7-7}{16.7.7}}.

$\operatorname{Diff}^{n}_{X/S}(\mathcal{F},\mathcal{G})$,
$\operatorname{Diff}^{n}_{X/S}$,
$\mathcal{D}iff^{n}_{X/S}(\mathcal{F},\mathcal{G})$:
{\hyperlink{ega4.c04.statement.16-8-3}{16.8.3}}.

$\mathcal{D}iff^{n}_{X/S}$, $\langle t,D\rangle$:
{\hyperlink{ega4.c04.prop.16-8-4}{16.8.4}}.

$aD$ ($D$ a differential operator, $a\in\Gamma(U,\mathcal{O}_X)$):
{\hyperlink{ega4.c04.statement.16-8-5}{16.8.5}}.

$\operatorname{Diff}_{X/S}(\mathcal{F},\mathcal{G})$,
$\mathcal{D}iff_{X/S}(\mathcal{F},\mathcal{G})$:
{\hyperlink{ega4.c04.statement.16-8-7}{16.8.7}}.

$\mathcal{D}iff_{X/S}$: {\hyperlink{ega4.c04.cor.16-8-10}{16.8.10}}.

$|\mathbf p|$, $\mathbf p!$, $\binom{\mathbf p}{\mathbf q}$,
$\mathbf z^{\mathbf p}$, $\boldsymbol\zeta^{\mathbf p}$
($\mathbf p,\mathbf q$ multi-indices): {\hyperlink{ega4.c04.statement.16-11-1}{16.11.1}}.

$\dim_x(f)$ ($f:X\to Y$ locally of finite type, $x\in X$):
{\hyperlink{ega4.c04.def.17-10-1}{17.10.1}}.

$\operatorname{Tr}_{B/A}$, $\operatorname{astr}_{B/A}$
($B$ a finite free $A$-algebra): {\hyperlink{ega4.c04.statement.18-2-1}{18.2.1}}.

$\operatorname{Tr}_{\mathcal B/\mathcal O_X,U}$,
$\operatorname{astr}_{\mathcal B/\mathcal O_X,U}$,
$\operatorname{Tr}_{\mathcal B/\mathcal O_X}$,
$\operatorname{astr}_{\mathcal B/\mathcal O_X}$
($X$ a ringed space, $\mathcal B$ an $\mathcal O_X$-algebra locally free
of finite rank): {\hyperlink{ega4.c04.statement.18-2-2}{18.2.2}}.

$\operatorname{Tr}_f$ ($f$ finite and \'etale): {\hyperlink{ega4.c04.cor.18-2-6}{18.2.6}}.

$d_{X/Y}$, $\mathcal D_{X/Y}$: {\hyperlink{ega4.c04.item.18-2-7}{18.2.7}}.

${}^{t}V_{\mathcal E}$: {\hyperlink{ega4.c04.statement.18-5-1}{18.5.1}}.

$Al(X,\mathcal E,\mathcal J)$, $Al(\mathcal E,\mathcal J)$:
{\hyperlink{ega4.c04.statement.18-5-2}{18.5.2}}.

$Of(S)$, $Of(X)$, $Id(\mathcal B')$: {\hyperlink{ega4.c04.lem.18-5-3}{18.5.3}}.

$\operatorname{Hom.loc}(A,B)$: {\hyperlink{ega4.c04.statement.18-6-1}{18.6.1}}.

${}^{h}A$: {\hyperlink{ega4.c04.def.18-6-5}{18.6.5}} and {\hyperlink{ega4.c04.statement.18-6-7}{18.6.7}}.

$\operatorname{Hom.sloc}(A,B)$: {\hyperlink{ega4.c04.statement.18-6-7}{18.6.7}}.

${}^{hs}A_{(i)}$, ${}^{hs}A$: {\hyperlink{ega4.c04.def.18-8-7}{18.8.7}} and
{\hyperlink{ega4.c04.statement.18-8-9}{18.8.9}}.

$\operatorname{codim}^{t}_{y}(Y,X)$: {\hyperlink{ega4.c04.def.19-1-3}{19.1.3}}.

$\operatorname{prof}_{T,t}(\mathcal F)$: {\hyperlink{ega4.c04.statement.19-9-1}{19.9.1}}.

$\mathcal O_X[\mathcal S^{-1}]$, $\mathcal F[\mathcal S^{-1}]$:
{\hyperlink{ega4.c04.statement.20-1-1}{20.1.1}}.

$\mathcal S(\mathcal O_X)$, $\mathcal M_X$, $M(X)$,
$\mathcal M_X(\mathcal F)$, $M(X,\mathcal F)$:
{\hyperlink{ega4.c04.statement.20-1-3}{20.1.3}}.

$\operatorname{dom}(\varphi)$ ($\varphi$ a meromorphic function):
{\hyperlink{ega4.c04.statement.20-1-4}{20.1.4}}.

$\operatorname{dom}(u)$ ($u$ a meromorphic section of $\mathcal F$):
{\hyperlink{ega4.c04.statement.20-1-7}{20.1.7}}.

$(\mathcal M_X(\mathcal L))^{*}$
($\mathcal L$ an invertible $\mathcal O_X$-Module):
{\hyperlink{ega4.c04.statement.20-1-8}{20.1.8}}.

$s^{-1}$ ($s$ a meromorphic section of an invertible
$\mathcal O_X$-Module): {\hyperlink{ega4.c04.statement.20-1-10}{20.1.10}}.

$\mathcal S_f(U)$, $\mathcal S_f$, $\mathcal M_f$
($f$ a morphism of ringed spaces), $\varphi\circ f$ ($\varphi$ a
meromorphic function), $u\circ f$ ($u$ a meromorphic section of
$\mathcal F$), $\mathcal M_f(\mathcal F)$: {\hyperlink{ega4.c04.statement.20-1-11}{20.1.11}}.

$\operatorname{Ps.hom}(X,Y)$, $\operatorname{Ps.hom}_S(X,Y)$:
{\hyperlink{ega4.c04.statement.20-2-1}{20.2.1}}.

$\omega|U$ ($\omega$ a pseudo-morphism),
$\mathcal P\!s.\operatorname{hom}(X,Y)$,
$\mathcal P\!s.\operatorname{hom}_S(X,Y)$:
{\hyperlink{ega4.c04.statement.20-2-2}{20.2.2}}.

$\operatorname{dom}_S(\omega)$, $\operatorname{dom}(\omega)$
($\omega$ a pseudo-morphism): {\hyperlink{ega4.c04.statement.20-2-3}{20.2.3}}.

$\mathcal M'_X$, $M'(X)$: {\hyperlink{ega4.c04.statement.20-2-8}{20.2.8}}.

$f\circ\omega$: {\hyperlink{ega4.c04.statement.20-3-1}{20.3.1}}.

$\omega\circ f$: {\hyperlink{ega4.c04.statement.20-3-2}{20.3.2}}.

$\operatorname{Ps.hom}_S(X,Y)^f$,
$\mathcal P\!s.\operatorname{hom}_S(X,Y)^f$: {\hyperlink{ega4.c04.statement.20-3-2}{20.3.2}}.

$\Gamma_\omega$: {\hyperlink{ega4.c04.statement.20-4-1}{20.4.1}}.

$p^{-1}$ ($p:\Gamma_\omega\to X$ the restriction of the first
projection): {\hyperlink{ega4.c04.statement.20-4-2}{20.4.2}}.

$\omega(M)$: {\hyperlink{ega4.c04.statement.20-4-2}{20.4.2}}.

$\operatorname{Ps.hom}_{X/S}(X,Y)$: {\hyperlink{ega4.c04.statement.20-5-1}{20.5.1}}.

$\mathcal P\!s.\operatorname{hom}_{X/S}(X,Y)$: {\hyperlink{ega4.c04.prop.20-5-3}{20.5.3}}.

$\mathcal M'_{X/S}$, $M'(X/S)$: {\hyperlink{ega4.c04.statement.20-5-4}{20.5.4}}.

$\mathcal S_{X/S}(U)$, $\mathcal S_{X/S}$, $\mathcal M_{X/S}$,
$M(X/S)$: {\hyperlink{ega4.c04.statement.20-6-1}{20.6.1}}.

$\mathcal M_{X/S}(\mathcal F)$, $M(X/S,\mathcal F)$:
{\hyperlink{ega4.c04.statement.20-6-2}{20.6.2}}.

$\mathcal S_f(U)$, $\mathcal S_f$, $\mathcal M_{X/S,f}$:
{\hyperlink{ega4.c04.statement.20-6-7}{20.6.7}}.

\oldpage[IV]{335}
$\mathcal{D}iv_X$, $\operatorname{Div}(X)$, $\operatorname{div}(f)$,
$\operatorname{div}_X(f)$ ($X$ a ringed space): {\hyperlink{ega4.c04.statement.21-1-2}{21.1.2}}.

$\operatorname{Supp}(D)$: {\hyperlink{ega4.c04.statement.21-1-2}{21.1.2}}.

$\operatorname{Div}(A)$: {\hyperlink{ega4.c04.statement.21-2-1}{21.2.1}}.

$\operatorname{div}(s)$ ($s$ a section of an invertible Module):
{\hyperlink{ega4.c04.statement.21-1-4}{21.1.4}}.

$\mathcal{D}iv_X^+$, $\operatorname{Div}^+(X)$, $D\leq D'$:
{\hyperlink{ega4.c04.def.21-1-6}{21.1.6}}.

$\operatorname{Id.inv}(X)$, $\sh{J}d.inv_X$:
{\hyperlink{ega4.c04.statement.21-2-4}{21.2.4}}.

$\mathcal J_U$, $\mathcal J_X(D)$: {\hyperlink{ega4.c04.statement.21-2-5}{21.2.5}}.

$\operatorname{div}(\mathcal I)$: {\hyperlink{ega4.c04.prop.21-2-6}{21.2.6}}.

$\mathcal O_X(D)$: {\hyperlink{ega4.c04.statement.21-2-8}{21.2.8}}.

$s_D$: {\hyperlink{ega4.c04.statement.21-2-9}{21.2.9}}.

$D(X)$, $\operatorname{cl}(\mathcal L,s)$: {\hyperlink{ega4.c04.statement.21-2-10}{21.2.10}}.

$Y(D)$: {\hyperlink{ega4.c04.statement.21-2-12}{21.2.12}}.

$\operatorname{Div.princ}(X)$: {\hyperlink{ega4.c04.statement.21-3-1}{21.3.1}}.

$\operatorname{Pic}(X)$, $\operatorname{cl}(\mathcal L)$, $l$, $l_X$,
$\operatorname{Pic}(u)$: {\hyperlink{ega4.c04.statement.21-3-2}{21.3.2}}.

$f^*(D)$, $\operatorname{Div}^f(X)$: {\hyperlink{ega4.c04.def.21-4-2}{21.4.2}}.

$\mathcal M_f^{**}$, $\mathcal{D}iv_X^f$: {\hyperlink{ega4.c04.statement.21-4-3}{21.4.3}}.

$N_{X'/X}(u')$, $N(u')$: {\hyperlink{ega4.c04.statement.21-5-3}{21.5.3}}.

$f_*(D')$, $N_{X'/X}(D')$: {\hyperlink{ega4.c04.def.21-5-5}{21.5.5}}.

$\mathfrak I(X)$, $\mathfrak Z(X)$, $\mathfrak Z^+(X)$
($X$ locally Noetherian): {\hyperlink{ega4.c04.statement.21-6-1}{21.6.1}}.

$\operatorname{mult}_x(Z)$, $\operatorname{Supp}(Z)$,
$\dim(Z)$, $\operatorname{codim}(Z)$: {\hyperlink{ega4.c04.statement.21-6-1}{21.6.1}}.

$X^{(p)}$, $\mathfrak Z^p(X)$, $\mathfrak Z^{p+}(X)$:
{\hyperlink{ega4.c04.statement.21-6-2}{21.6.2}}.

$Z|U$, $\mathscr Z_X$, $\mathscr Z_X^p$, $\mathscr Z_X^+$,
$\mathscr Z_X^{p+}$: {\hyperlink{ega4.c04.statement.21-6-3}{21.6.3}}.

$c$: {\hyperlink{ega4.c04.statement.21-6-4}{21.6.4}}.

$\operatorname{cyc}(D)$: {\hyperlink{ega4.c04.statement.21-6-5}{21.6.5}} and
{\hyperlink{ega4.c04.statement.21-6-7}{21.6.7}}.

$\operatorname{mult}_x(D)$, $\operatorname{mult}_{\{x\}}(D)$,
$o_x(f)$: {\hyperlink{ega4.c04.statement.21-6-7}{21.6.7}}.

$\mathrm Z^+(f)$, $\mathrm Z^-(f)$, $\mathfrak Z^1_{\mathrm{princ}}(X)$,
$\mathfrak{Cl}(X)$: {\hyperlink{ega4.c04.statement.21-6-7}{21.6.7}}.

$Y(C)$, $Y'(C)$, $\sh{J}_X(C)$, $\sh{J}(C)$:
{\hyperlink{ega4.c04.statement.21-7-1}{21.7.1}}.

$f^*(Z)$: {\hyperlink{ega4.c04.statement.21-10-3}{21.10.3}} and {\hyperlink{ega4.c04.statement.21-10-11}{21.10.11}}.

$\Gamma_T(\mathscr Z_X^1)$: {\hyperlink{ega4.c04.statement.21-10-3}{21.10.3}}.

$f_*(Z')$: {\hyperlink{ega4.c04.statement.21-10-14}{21.10.14}}.

$\Lambda^{\max}\mathcal E$: {\hyperlink{ega4.c04.thm.21-11-1}{21.11.1}}.

$U'$, $\operatorname{Aff}(U/X)$: {\hyperlink{ega4.c04.statement.21-12-1}{21.12.1}}.

$\operatorname{Daf}(U/X)$: {\hyperlink{ega4.c04.statement.21-12-5}{21.12.5}}.

$\mathcal{D}iv_{X/S}$, $\operatorname{Div}(X/S)$:
{\hyperlink{ega4.c04.def.21-15-2}{21.15.2}}.

$\mathcal{D}iv_{X/S}^+$, $\operatorname{Div}^+(X/S)$:
{\hyperlink{ega4.c04.statement.21-15-3}{21.15.3}}.

$\mathcal O_{X/S}^{\mathrm{reg}}$: {\hyperlink{ega4.c04.statement.21-15-3}{21.15.3}}.

$\operatorname{Div}_{X/S}(S')$, $\operatorname{Div}_{X/S}^+(S')$:
{\hyperlink{ega4.c04.statement.21-15-9}{21.15.9}}.

\endgroup
% END NATIVE INPUT source_aligned/ega4-backmatter-notation
% BEGIN NATIVE INPUT source_aligned/ega4-backmatter-terminology
% Source-aligned English projection of EGA IV, Part 4, pp. 336--339.
% Entries remain in the source index order; only the head terms are translated.
\clearpage
\oldpage[IV]{336}
\phantomsection
\addcontentsline{toc}{section}{Terminological index}
\section*{Terminological index}

\begingroup
\small
\setlength{\parindent}{0pt}
\setlength{\parskip}{2pt}
\newcommand{\EGAIVterm}[2]{\textbf{#1}: #2.\par}
\newcommand{\EGAIVidxref}[1]{\hyperref[IV.#1]{#1}}

\EGAIVterm{\'Etale algebra}{{\hyperlink{ega4.c04.subsec.17-3-2}{17.3.2}}}
\EGAIVterm{Essentially \'etale algebra}{{\hyperlink{ega4.c04.statement.18-6-1}{18.6.1}}}
\EGAIVterm{Finite and \'etale $\sh{O}_X$-algebra}{{\hyperlink{ega4.c04.prop.18-2-3}{18.2.3}}}
\EGAIVterm{Smooth algebra}{{\hyperlink{ega4.c04.subsec.17-3-2}{17.3.2}}}
\EGAIVterm{Unramified algebra}{{\hyperlink{ega4.c04.subsec.17-3-2}{17.3.2}}}
\EGAIVterm{$K$-algebra unramified over $Y$, over $A$}{{\hyperlink{ega4.c04.rem.18-10-11}{18.10.11}}}
\EGAIVterm{Strictly essentially \'etale algebra}{{\hyperlink{ega4.c04.lem.18-6-2}{18.6.2}}}
\EGAIVterm{Complete-intersection ring; absolute complete-intersection ring}{{\hyperlink{ega4.c04.def.19-3-1}{19.3.1}}}
\EGAIVterm{Henselian ring}{{\hyperlink{ega4.c04.def.18-5-8}{18.5.8}}}
\EGAIVterm{Parafactorial local ring}{{\hyperlink{ega4.c04.def.21-13-7}{21.13.7}}}
\EGAIVterm{Strictly local ring}{{\hyperlink{ega4.c04.def.18-8-2}{18.8.2}}}
\EGAIVterm{Tangent linear map}{{\hyperlink{ega4.c04.statement.16-5-13}{16.5.13}}}
\EGAIVterm{Strict rational map}{{\hyperlink{ega4.c04.statement.20-2-1}{20.2.1}}}
\EGAIVterm{Codimension of a cycle}{{\hyperlink{ega4.c04.statement.21-6-1}{21.6.1}}}
\EGAIVterm{Transversal codimension at a point of a sub-prescheme}{{\hyperlink{ega4.c04.def.19-1-3}{19.1.3}}}
\EGAIVterm{Transversal codimension of an immersion morphism at a point}{{\hyperlink{ega4.c04.def.19-1-3}{19.1.3}}}
\EGAIVterm{Pair of morphisms transversal at a point}{{\hyperlink{ega4.c04.cor.17-13-6}{17.13.6}}}
\EGAIVterm{Henselian pair}{{\hyperlink{ega4.c04.def.18-5-5}{18.5.5}}}
\EGAIVterm{Parafactorial pair}{{\hyperlink{ega4.c04.def.21-13-1}{21.13.1}}}
\EGAIVterm{Cycle}{{\hyperlink{ega4.c04.statement.21-6-1}{21.6.1}}}
\EGAIVterm{Codimension-one cycle with rational coefficients}{{\hyperlink{ega4.c04.statement.21-10-9}{21.10.9}}}
\EGAIVterm{Codimension-$p$ cycle}{{\hyperlink{ega4.c04.statement.21-6-2}{21.6.2}}}
\EGAIVterm{Cycle of poles; cycle of zeros of a meromorphic function}{{\hyperlink{ega4.c04.statement.21-6-7}{21.6.7}}}
\EGAIVterm{Cycle linearly equivalent to zero}{{\hyperlink{ega4.c04.statement.21-6-7}{21.6.7}}}
\EGAIVterm{Locally principal cycle}{{\hyperlink{ega4.c04.statement.21-6-7}{21.6.7}}}
\EGAIVterm{Polar cycle of a meromorphic function}{{\hyperlink{ega4.c04.statement.21-6-7}{21.6.7}}}
\EGAIVterm{Prime cycle}{{\hyperlink{ega4.c04.statement.21-6-1}{21.6.1}}}
\EGAIVterm{Principal cycle}{{\hyperlink{ega4.c04.statement.21-6-7}{21.6.7}}}
\EGAIVterm{Cycle principal at a point}{{\hyperlink{ega4.c04.statement.21-6-7}{21.6.7}}}
\EGAIVterm{Cycle purely of codimension $p$}{{\hyperlink{ega4.c04.statement.21-6-2}{21.6.2}}}
\EGAIVterm{Affineness defect of an open in a prescheme}{{\hyperlink{ega4.c04.statement.21-12-5}{21.12.5}}}
\EGAIVterm{$S$-derivation, $(X/S)$-derivation, $f$-derivation}{{\hyperlink{ega4.c04.statement.16-5-1}{16.5.1}}}
\EGAIVterm{Differential of a section of $\sh{O}_X$}{{\hyperlink{ega4.c04.def.16-3-6}{16.3.6}}}
\EGAIVterm{Exterior differential}{{\hyperlink{ega4.c04.def.16-6-3}{16.6.3}}}
\EGAIVterm{Dimension of a cycle}{{\hyperlink{ega4.c04.statement.21-6-1}{21.6.1}}}
\EGAIVterm{Relative dimension of a morphism at a point; of a $Y$-prescheme at a point}{{\hyperlink{ega4.c04.def.17-10-1}{17.10.1}}}
\EGAIVterm{Discriminant of a finite locally free $Y$-prescheme}{{\hyperlink{ega4.c04.item.18-2-7}{18.2.7}}}
\EGAIVterm{Divisor on a ringed space}{{\hyperlink{ega4.c04.statement.21-1-2}{21.1.2}}}
\EGAIVterm{Divisor of a meromorphic function}{{\hyperlink{ega4.c04.statement.21-1-2}{21.1.2}}}
\EGAIVterm{Divisor of a meromorphic section of an invertible Module}{{\hyperlink{ega4.c04.statement.21-1-4}{21.1.4}}}
\EGAIVterm{Positive divisor}{{\hyperlink{ega4.c04.def.21-1-6}{21.1.6}}}
\EGAIVterm{Principal divisor}{{\hyperlink{ega4.c04.statement.21-3-1}{21.3.1}}}
\EGAIVterm{Divisor relative to $S$; divisor transversal to $f$}{{\hyperlink{ega4.c04.def.21-15-2}{21.15.2}}}
\EGAIVterm{Linearly equivalent divisors}{{\hyperlink{ega4.c04.statement.21-3-1}{21.3.1}}}

\oldpage[IV]{337}
\EGAIVterm{Domain of definition of a meromorphic function}{{\hyperlink{ega4.c04.statement.20-1-4}{20.1.4}}}
\EGAIVterm{Domain of definition of a pseudo-morphism}{{\hyperlink{ega4.c04.statement.20-2-3}{20.2.3}}}
\EGAIVterm{Domain of definition of a meromorphic section of a Module}{{\hyperlink{ega4.c04.statement.20-1-7}{20.1.7}}}
\EGAIVterm{Affine envelope of an $X$-prescheme}{{\hyperlink{ega4.c04.statement.21-12-1}{21.12.1}}}
\EGAIVterm{Tangent space to a prescheme at a point}{{\hyperlink{ega4.c04.statement.16-5-13}{16.5.13}}}
\EGAIVterm{Conormal sheaf}{{\hyperlink{ega4.c04.def.16-1-2}{16.1.2}}}
\EGAIVterm{Sheaf of rings of fractions with denominators in $\sh{S}$}{{\hyperlink{ega4.c04.statement.20-1-1}{20.1.1}}}
\EGAIVterm{Sheaf of graded rings associated with a morphism}{{\hyperlink{ega4.c04.def.16-1-2}{16.1.2}}}
\EGAIVterm{Sheaf of modules of fractions with denominators in $\sh{S}$}{{\hyperlink{ega4.c04.statement.20-1-2}{20.1.2}}}
\EGAIVterm{Sheaf of $S$-derivations of $\sh{O}_X$}{{\hyperlink{ega4.c04.statement.16-5-7}{16.5.7}}}
\EGAIVterm{Sheaf of $S$-derivations of $\sh{O}_X$ into $\sh{F}$}{{\hyperlink{ega4.c04.statement.16-5-4}{16.5.4}}}
\EGAIVterm{Sheaf of $p$-differentials of $X$ relative to $S$}{{\hyperlink{ega4.c04.statement.16-6-1}{16.6.1}}}
\EGAIVterm{Sheaf of divisors on a ringed space}{{\hyperlink{ega4.c04.statement.21-1-2}{21.1.2}}}
\EGAIVterm{Sheaf of divisors relative to $S$, or transversal to $f$}{{\hyperlink{ega4.c04.def.21-15-2}{21.15.2}}}
\EGAIVterm{Sheaf of germs of codimension-one cycles with coefficients in $\Lambda$}{{\hyperlink{ega4.c04.statement.21-10-9}{21.10.9}}}
\EGAIVterm{Sheaf of germs of meromorphic functions on $X$}{{\hyperlink{ega4.c04.statement.20-1-4}{20.1.4}}}
\EGAIVterm{Sheaf of germs of meromorphic functions relative to $S$}{{\hyperlink{ega4.c04.statement.20-6-1}{20.6.1}}}
\EGAIVterm{Sheaf of germs of meromorphic sections of a Module}{{\hyperlink{ega4.c04.statement.20-1-3}{20.1.3}}}
\EGAIVterm{Sheaf of principal parts of order $n$ of the $S$-prescheme $X$}{{\hyperlink{ega4.c04.def.16-3-1}{16.3.1}}}
\EGAIVterm{Formally principal homogeneous sheaf (or pseudo-torsor) under a sheaf of groups}{{\hyperlink{ega4.c04.statement.16-5-15}{16.5.15}}}
\EGAIVterm{Principal homogeneous sheaf (or torsor) under a sheaf of groups}{{\hyperlink{ega4.c04.statement.16-5-15}{16.5.15}}}
\EGAIVterm{Trivial principal homogeneous sheaf (or torsor) under a sheaf of groups}{{\hyperlink{ega4.c04.statement.16-5-15}{16.5.15}}}
\EGAIVterm{Tangent sheaf of $X$ relative to $S$}{{\hyperlink{ega4.c04.statement.16-5-7}{16.5.7}}}
\EGAIVterm{Family of divisors on $X$ relative to $S$, parametrized by $S'$}{{\hyperlink{ega4.c04.statement.21-15-9}{21.15.9}}}
\EGAIVterm{Family of morphisms transversal at a point}{{\hyperlink{ega4.c04.cor.17-13-10}{17.13.10}}}
\EGAIVterm{Family of sub-preschemes intersecting transversally at a point}{{\hyperlink{ega4.c04.statement.17-13-11}{17.13.11}}}
\EGAIVterm{Tangent bundle of $X$ relative to $S$}{{\hyperlink{ega4.c04.statement.16-5-12}{16.5.12}}}
\EGAIVterm{Meromorphic function on a ringed space}{{\hyperlink{ega4.c04.statement.20-1-3}{20.1.3}}}
\EGAIVterm{Meromorphic function defined at a point}{{\hyperlink{ega4.c04.statement.20-1-4}{20.1.4}}}
\EGAIVterm{Regular meromorphic function}{{\hyperlink{ega4.c04.statement.20-1-8}{20.1.8}}}
\EGAIVterm{Meromorphic function regular relative to $S$}{{\hyperlink{ega4.c04.statement.20-6-5}{20.6.5}}}
\EGAIVterm{Meromorphic function relative to $S$}{{\hyperlink{ega4.c04.statement.20-6-1}{20.6.1}}}
\EGAIVterm{Trace form}{{\hyperlink{ega4.c04.statement.18-2-1}{18.2.1}}}
\EGAIVterm{Graph of a pseudo-$S$-morphism}{{\hyperlink{ega4.c04.statement.20-4-1}{20.4.1}}}
\EGAIVterm{Group of classes of codimension-one cycles}{{\hyperlink{ega4.c04.statement.21-6-7}{21.6.7}}}
\EGAIVterm{Henselization of a local ring}{{\hyperlink{ega4.c04.def.18-6-5}{18.6.5}}}
\EGAIVterm{Henselization of a semilocal ring}{{\hyperlink{ega4.c04.statement.18-6-7}{18.6.7}}}
\EGAIVterm{Strict Henselization of a local ring}{{\hyperlink{ega4.c04.def.18-8-7}{18.8.7}}}
\EGAIVterm{Strict Henselization of a semilocal ring}{{\hyperlink{ega4.c04.statement.18-8-9}{18.8.9}}}
\EGAIVterm{Homomorphism associated with the trace}{{\hyperlink{ega4.c04.statement.18-2-2}{18.2.2}}}
\EGAIVterm{Regular homomorphism of $\sh{O}_X$-Modules}{{\hyperlink{ega4.c04.prop.19-4-11}{19.4.11}}}
\EGAIVterm{Semilocal homomorphism}{{\hyperlink{ega4.c04.statement.18-6-7}{18.6.7}}}
\EGAIVterm{Trace homomorphism}{{\hyperlink{ega4.c04.statement.18-2-2}{18.2.2}}}
\EGAIVterm{Ideal of denominators of a meromorphic section}{{\hyperlink{ega4.c04.statement.20-2-14}{20.2.14}}}
\EGAIVterm{Discriminant ideal of a finite locally free $Y$-prescheme}{{\hyperlink{ega4.c04.item.18-2-7}{18.2.7}}}
\EGAIVterm{Integral ideal}{{\hyperlink{ega4.c04.statement.21-2-7}{21.2.7}}}
\EGAIVterm{Fractional ideal}{{\hyperlink{ega4.c04.statement.21-2-1}{21.2.1}}}
\EGAIVterm{Invertible fractional ideal}{{\hyperlink{ega4.c04.statement.21-2-1}{21.2.1}}}
\EGAIVterm{Quasi-regular ideal of a sheaf of rings}{{\hyperlink{ega4.c04.def.16-9-1}{16.9.1}}}
\EGAIVterm{Quasi-regular ideal of a ring}{{\hyperlink{ega4.c04.statement.16-9-7}{16.9.7}}}
\EGAIVterm{Regular ideal of a sheaf of rings}{{\hyperlink{ega4.c04.def.16-9-1}{16.9.1}}}
\EGAIVterm{Regular ideal of a ring}{{\hyperlink{ega4.c04.statement.16-9-7}{16.9.7}}}
\EGAIVterm{Transversally regular ideal at a point of a prescheme}{{\hyperlink{ega4.c04.def.19-2-1}{19.2.1}}}

\oldpage[IV]{338}
\EGAIVterm{Transversally regular ideal of an algebra}{{\hyperlink{ega4.c04.def.19-2-1}{19.2.1}}}
\EGAIVterm{Direct image of a codimension-one cycle}{{\hyperlink{ega4.c04.statement.21-10-14}{21.10.14}}}
\EGAIVterm{Direct image of a divisor}{{\hyperlink{ega4.c04.def.21-5-5}{21.5.5}}}
\EGAIVterm{Inverse image of a codimension-one cycle}{{\hyperlink{ega4.c04.statement.21-10-3}{21.10.3}} and {\hyperlink{ega4.c04.statement.21-10-11}{21.10.11}}}
\EGAIVterm{Inverse image of a divisor}{{\hyperlink{ega4.c04.def.21-4-2}{21.4.2}}}
\EGAIVterm{Inverse image of a meromorphic function}{{\hyperlink{ega4.c04.statement.20-1-11}{20.1.11}}}
\EGAIVterm{Inverse image of a meromorphic function relative to $S$}{{\hyperlink{ega4.c04.statement.20-6-7}{20.6.7}}}
\EGAIVterm{Quasi-regular immersion}{{\hyperlink{ega4.c04.def.16-9-2}{16.9.2}}}
\EGAIVterm{Regular immersion}{{\hyperlink{ega4.c04.def.16-9-2}{16.9.2}}}
\EGAIVterm{Regular immersion at a point}{{\hyperlink{ega4.c04.prop.16-9-10}{16.9.10}}}
\EGAIVterm{Transversally regular immersion at a point relative to $S$; transversally regular immersion relative to $S$}{{\hyperlink{ega4.c04.def.19-2-2}{19.2.2}}}
\EGAIVterm{Complete intersection at a point relative to $S$; complete intersection relative to $S$ (for a prescheme flat and locally of finite presentation)}{{\hyperlink{ega4.c04.def.19-3-6}{19.3.6}}}
\EGAIVterm{$n$th normal invariant of a morphism}{{\hyperlink{ega4.c04.def.16-1-2}{16.1.2}}}
\EGAIVterm{Normal invariant of infinite order of a morphism}{{\hyperlink{ega4.c04.rem.16-1-11}{16.1.11}}}
\EGAIVterm{$\sh{O}_X$-Module of first differentials of $f$ (or of $X$ relative to $S$, or of the $S$-prescheme $X$)}{{\hyperlink{ega4.c04.def.16-3-1}{16.3.1}}}
\EGAIVterm{$\sh{O}_X$-Module torsion-free relative to $S$}{{\hyperlink{ega4.c04.statement.20-6-2}{20.6.2}}}
\EGAIVterm{Torsion-free $\sh{O}_X$-Module}{{\hyperlink{ega4.c04.statement.20-1-5}{20.1.5}}}
\EGAIVterm{Differentially smooth morphism}{{\hyperlink{ega4.c04.def.16-10-1}{16.10.1}}}
\EGAIVterm{Differentially smooth morphism up to order $r$}{{\hyperlink{ega4.c04.statement.16-11-3}{16.11.3}}}
\EGAIVterm{(Flat) complete-intersection morphism}{{\hyperlink{ega4.c04.def.19-3-6}{19.3.6}}}
\EGAIVterm{Essentially proper morphism}{{\hyperlink{ega4.c04.rem.18-10-20}{18.10.20}}}
\EGAIVterm{\'Etale morphism}{{\hyperlink{ega4.c04.def.17-3-1}{17.3.1}}}
\EGAIVterm{\'Etale morphism at a point}{{\hyperlink{ega4.c04.def.17-3-7}{17.3.7}}}
\EGAIVterm{Finite locally free morphism of rank $n$; finite locally free morphism}{{\hyperlink{ega4.c04.item.18-2-7}{18.2.7}}}
\EGAIVterm{Formally \'etale, formally smooth, formally net, formally unramified morphism}{{\hyperlink{ega4.c04.def.17-1-1}{17.1.1}}}
\EGAIVterm{Smooth morphism}{{\hyperlink{ega4.c04.def.17-3-1}{17.3.1}}}
\EGAIVterm{Smooth morphism at a point}{{\hyperlink{ega4.c04.def.17-3-7}{17.3.7}}}
\EGAIVterm{Unramified morphism; net morphism}{{\hyperlink{ega4.c04.def.17-3-1}{17.3.1}}}
\EGAIVterm{Unramified morphism at a point}{{\hyperlink{ega4.c04.def.17-3-7}{17.3.7}}}
\EGAIVterm{Morphism transversal to $Y$ at $x$ relative to $S$}{{\hyperlink{ega4.c04.def.17-13-3}{17.13.3}}}
\EGAIVterm{Equivalent morphisms}{{\hyperlink{ega4.c04.statement.20-2-1}{20.2.1}}}
\EGAIVterm{Multiplicity of a cycle at a point}{{\hyperlink{ega4.c04.statement.21-6-1}{21.6.1}}}
\EGAIVterm{Norm of a divisor}{{\hyperlink{ega4.c04.def.21-5-5}{21.5.5}}}
\EGAIVterm{Norm of a meromorphic section}{{\hyperlink{ega4.c04.statement.21-5-3}{21.5.3}}}
\EGAIVterm{Differential operator relative to $S$}{{\hyperlink{ega4.c04.def.16-8-1}{16.8.1}}}
\EGAIVterm{Order of a differential operator}{{\hyperlink{ega4.c04.def.16-8-1}{16.8.1}}}
\EGAIVterm{Order of a meromorphic function at a point $x\in X^{(1)}$}{{\hyperlink{ega4.c04.statement.21-6-7}{21.6.7}}}
\EGAIVterm{Principal part of order $n$, or of infinite order, of a section of $\sh{O}_X$}{{\hyperlink{ega4.c04.def.16-3-6}{16.3.6}}}
\EGAIVterm{Closed part purely of codimension $p$}{{\hyperlink{ega4.c04.statement.21-6-2}{21.6.2}}}
\EGAIVterm{Separable polynomial over a ring}{{\hyperlink{ega4.c04.cor.18-4-3}{18.4.3}}}
\EGAIVterm{Prescheme of characteristic $p$}{{\hyperlink{ega4.c04.statement.16-12-1}{16.12.1}}}
\EGAIVterm{Prescheme of coincidences of two morphisms}{{\hyperlink{ega4.c04.statement.17-4-5}{17.4.5}}}
\EGAIVterm{Differentially smooth prescheme over $S$}{{\hyperlink{ega4.c04.def.16-10-1}{16.10.1}}}
\EGAIVterm{Prescheme \'etale over $S$}{{\hyperlink{ega4.c04.def.17-3-1}{17.3.1}}}
\EGAIVterm{Prescheme \'etale over $S$ at a point}{{\hyperlink{ega4.c04.def.17-3-7}{17.3.7}}}
\EGAIVterm{Prescheme formally \'etale, formally smooth, or formally unramified over $S$}{{\hyperlink{ega4.c04.def.17-1-1}{17.1.1}}}
\EGAIVterm{Absolute complete-intersection prescheme at a point}{{\hyperlink{ega4.c04.def.19-3-1}{19.3.1}}}
\EGAIVterm{Prescheme smooth over $S$}{{\hyperlink{ega4.c04.def.17-3-1}{17.3.1}}}
\EGAIVterm{Prescheme smooth over $S$ at a point}{{\hyperlink{ega4.c04.def.17-3-7}{17.3.7}}}
\EGAIVterm{Locally factorial prescheme}{{\hyperlink{ega4.c04.thm.21-6-9}{21.6.9}}}
\EGAIVterm{Prescheme net over $S$; prescheme unramified over $S$}{{\hyperlink{ega4.c04.def.17-3-1}{17.3.1}}}

\oldpage[IV]{339}
\EGAIVterm{Prescheme net over $S$ at a point; prescheme unramified over $S$ at a point}{{\hyperlink{ega4.c04.def.17-3-7}{17.3.7}}}
\EGAIVterm{$T$-depth of a coherent $\sh{O}_X$-Module at a point}{{\hyperlink{ega4.c04.statement.19-9-1}{19.9.1}}}
\EGAIVterm{Pseudo-function}{{\hyperlink{ega4.c04.statement.20-2-8}{20.2.8}}}
\EGAIVterm{Pseudo-function relative to $S$}{{\hyperlink{ega4.c04.statement.20-5-4}{20.5.4}}}
\EGAIVterm{Pseudo-morphism; pseudo-$S$-morphism}{{\hyperlink{ega4.c04.statement.20-2-1}{20.2.1}}}
\EGAIVterm{Pseudo-morphism composed with a morphism}{{\hyperlink{ega4.c04.statement.20-3-1}{20.3.1}}}
\EGAIVterm{Pseudo-morphism composed with a morphism on the other side}{{\hyperlink{ega4.c04.statement.20-3-2}{20.3.2}}}
\EGAIVterm{Pseudo-morphism defined at a point}{{\hyperlink{ega4.c04.statement.20-2-3}{20.2.3}}}
\EGAIVterm{Pseudo-morphism relative to $S$}{{\hyperlink{ega4.c04.statement.20-5-1}{20.5.1}}}
\EGAIVterm{Pseudo-torsor}{{\hyperlink{ega4.c04.statement.16-5-15}{16.5.15}}}
\EGAIVterm{Restriction of a cycle to an open}{{\hyperlink{ega4.c04.statement.21-6-3}{21.6.3}}}
\EGAIVterm{Restriction of a pseudo-morphism to an open}{{\hyperlink{ega4.c04.statement.20-2-2}{20.2.2}}}
\EGAIVterm{Restriction of a pseudo-morphism to a local scheme}{{\hyperlink{ega4.c04.statement.20-3-6}{20.3.6}}}
\EGAIVterm{\'Etale covering; locally trivial \'etale covering; trivial \'etale covering}{{\hyperlink{ega4.c04.item.18-2-7}{18.2.7}}}
\EGAIVterm{Henselian local scheme}{{\hyperlink{ega4.c04.def.18-5-8}{18.5.8}}}
\EGAIVterm{Strictly local scheme}{{\hyperlink{ega4.c04.def.18-8-2}{18.8.2}}}
\EGAIVterm{Meromorphic section of an $\sh{O}_X$-Module}{{\hyperlink{ega4.c04.statement.20-1-3}{20.1.3}}}
\EGAIVterm{Meromorphic section of an $\sh{O}_X$-Module defined at a point}{{\hyperlink{ega4.c04.statement.20-1-7}{20.1.7}}}
\EGAIVterm{Regular meromorphic section of an invertible $\sh{O}_X$-Module}{{\hyperlink{ega4.c04.statement.20-1-8}{20.1.8}}}
\EGAIVterm{Meromorphic section of an $\sh{O}_X$-Module relative to $S$}{{\hyperlink{ega4.c04.statement.20-6-2}{20.6.2}}}
\EGAIVterm{Meromorphic section of an invertible $\sh{O}_X$-Module regular relative to $S$}{{\hyperlink{ega4.c04.statement.20-6-5}{20.6.5}}}
\EGAIVterm{Quasi-regularly immersed or regularly immersed sub-prescheme}{{\hyperlink{ega4.c04.def.16-9-2}{16.9.2}}}
\EGAIVterm{Sub-preschemes intersecting transversally at a point relative to $S$}{{\hyperlink{ega4.c04.statement.17-13-7}{17.13.7}}}
\EGAIVterm{Support of a cycle}{{\hyperlink{ega4.c04.statement.21-6-1}{21.6.1}}}
\EGAIVterm{Support of a divisor}{{\hyperlink{ega4.c04.statement.21-1-2}{21.1.2}}}
\EGAIVterm{Sequence of sections transversally regular relative to $S$}{{\hyperlink{ega4.c04.cor.19-1-2}{19.1.2}}}
\EGAIVterm{Canonical symmetry of $X\times_S X$}{{\hyperlink{ega4.c04.statement.16-3-3}{16.3.3}}}
\EGAIVterm{Canonical symmetry of $\sh{P}^{n}_{X/S}$, of $\sh{P}^{\infty}_{X/S}$}{{\hyperlink{ega4.c04.prop.16-3-4}{16.3.4}}}
\EGAIVterm{Quasi-regular (regular) system of generators of an Ideal}{{\hyperlink{ega4.c04.def.16-9-1}{16.9.1}}}
\EGAIVterm{Torsor}{{\hyperlink{ega4.c04.statement.16-5-15}{16.5.15}}}
\EGAIVterm{Formal neighbourhood of a sub-prescheme}{{\hyperlink{ega4.c04.rem.16-1-11}{16.1.11}}}
\EGAIVterm{$n$th infinitesimal neighbourhood}{{\hyperlink{ega4.c04.def.16-1-2}{16.1.2}}}

\endgroup
% END NATIVE INPUT source_aligned/ega4-backmatter-terminology
% BEGIN NATIVE INPUT source_aligned/ega4-backmatter-contents
% Source-aligned English projection of EGA IV, Part 4, pp. 341--343.
\clearpage
\oldpage[IV]{341}
\phantomsection
\addcontentsline{toc}{section}{Original table of contents}
\section*{Original table of contents}

\begingroup
\small
\newcommand{\EGAIVcontentsline}[3]{%
  \textbf{#1}\quad #2\dotfill #3\par}

\begin{center}\textbf{Chapter IV. --- Local study of schemes and morphisms of schemes (end)}\end{center}

\EGAIVcontentsline{\S16}{Differential invariants. Differentially smooth morphisms}{5}
\EGAIVcontentsline{{\texorpdfstring{\protect\hyperlink{ega4.c04.subsec.16-1}{16.1}}{16.1}}}{Normal invariants of an immersion}{5}
\EGAIVcontentsline{{\texorpdfstring{\protect\hyperlink{ega4.c04.subsec.16-2}{16.2}}{16.2}}}{Functorial properties of the normal invariants of an immersion}{9}
\EGAIVcontentsline{{\texorpdfstring{\protect\hyperlink{ega4.c04.subsec.16-3}{16.3}}{16.3}}}{Fundamental differential invariants of a morphism of preschemes}{14}
\EGAIVcontentsline{{\texorpdfstring{\protect\hyperlink{ega4.c04.subsec.16-4}{16.4}}{16.4}}}{Functorial properties of differential invariants}{16}
\EGAIVcontentsline{{\texorpdfstring{\protect\hyperlink{ega4.c04.subsec.16-5}{16.5}}{16.5}}}{Relative tangent sheaves and bundles; derivations}{27}
\EGAIVcontentsline{{\texorpdfstring{\protect\hyperlink{ega4.c04.subsec.16-6}{16.6}}{16.6}}}{Sheaves of $p$-differentials and exterior differential}{34}
\EGAIVcontentsline{{\texorpdfstring{\protect\hyperlink{ega4.c04.subsec.16-7}{16.7}}{16.7}}}{The $\sh{P}^{n}_{X/S}(\sh{F})$}{36}
\EGAIVcontentsline{{\texorpdfstring{\protect\hyperlink{ega4.c04.subsec.16-8}{16.8}}{16.8}}}{Differential operators}{39}
\EGAIVcontentsline{{\texorpdfstring{\protect\hyperlink{ega4.c04.subsec.16-9}{16.9}}{16.9}}}{Regular and quasi-regular immersions}{46}
\EGAIVcontentsline{{\texorpdfstring{\protect\hyperlink{ega4.c04.subsec.16-10}{16.10}}{16.10}}}{Differentially smooth morphisms}{51}
\EGAIVcontentsline{{\texorpdfstring{\protect\hyperlink{ega4.c04.subsec.16-11}{16.11}}{16.11}}}{Differential operators on a differentially smooth $S$-prescheme}{53}
\EGAIVcontentsline{{\texorpdfstring{\protect\hyperlink{ega4.c04.subsec.16-12}{16.12}}{16.12}}}{Characteristic zero: the Jacobian criterion for differentially smooth morphisms}{55}

\EGAIVcontentsline{\S17}{Smooth, unramified, and \'etale morphisms}{56}
\EGAIVcontentsline{{\texorpdfstring{\protect\hyperlink{ega4.c04.subsec.17-1}{17.1}}{17.1}}}{Formally smooth, formally unramified, and formally \'etale morphisms}{56}
\EGAIVcontentsline{{\texorpdfstring{\protect\hyperlink{ega4.c04.subsec.17-2}{17.2}}{17.2}}}{General differential properties}{59}
\EGAIVcontentsline{{\texorpdfstring{\protect\hyperlink{ega4.c04.subsec.17-3}{17.3}}{17.3}}}{Smooth, unramified, and \'etale morphisms}{61}
\EGAIVcontentsline{{\texorpdfstring{\protect\hyperlink{ega4.c04.subsec.17-4}{17.4}}{17.4}}}{Characterizations of unramified morphisms}{65}
\EGAIVcontentsline{{\texorpdfstring{\protect\hyperlink{ega4.c04.subsec.17-5}{17.5}}{17.5}}}{Characterizations of smooth morphisms}{67}
\EGAIVcontentsline{{\texorpdfstring{\protect\hyperlink{ega4.c04.subsec.17-6}{17.6}}{17.6}}}{Characterizations of \'etale morphisms}{70}
\EGAIVcontentsline{{\texorpdfstring{\protect\hyperlink{ega4.c04.subsec.17-7}{17.7}}{17.7}}}{Descent and passage-to-the-limit properties}{72}
\EGAIVcontentsline{{\texorpdfstring{\protect\hyperlink{ega4.c04.subsec.17-8}{17.8}}{17.8}}}{Fibrewise criteria for smoothness and unramifiedness}{79}
\EGAIVcontentsline{{\texorpdfstring{\protect\hyperlink{ega4.c04.subsec.17-9}{17.9}}{17.9}}}{\'Etale morphisms and open immersions}{79}
\EGAIVcontentsline{{\texorpdfstring{\protect\hyperlink{ega4.c04.subsec.17-10}{17.10}}{17.10}}}{Relative dimension of a prescheme smooth over another}{81}
\EGAIVcontentsline{{\texorpdfstring{\protect\hyperlink{ega4.c04.subsec.17-11}{17.11}}{17.11}}}{Smooth morphisms of smooth preschemes}{82}
\EGAIVcontentsline{{\texorpdfstring{\protect\hyperlink{ega4.c04.subsec.17-12}{17.12}}{17.12}}}{Smooth sub-preschemes of a smooth prescheme; smooth and differentially smooth morphisms}{85}

\oldpage[IV]{342}
\EGAIVcontentsline{{\texorpdfstring{\protect\hyperlink{ega4.c04.subsec.17-13}{17.13}}{17.13}}}{Transversal morphisms}{89}
\EGAIVcontentsline{{\texorpdfstring{\protect\hyperlink{ega4.c04.subsec.17-14}{17.14}}{17.14}}}{Local and infinitesimal characterizations of smooth, unramified, and \'etale morphisms}{98}
\EGAIVcontentsline{{\texorpdfstring{\protect\hyperlink{ega4.c04.subsec.17-15}{17.15}}{17.15}}}{Preschemes over a base field}{99}
\EGAIVcontentsline{{\texorpdfstring{\protect\hyperlink{ega4.c04.subsec.17-16}{17.16}}{17.16}}}{Quasi-sections of flat or smooth morphisms}{105}

\EGAIVcontentsline{\S18}{Complements on \'etale morphisms; Henselian local rings and strictly local rings}{109}
\EGAIVcontentsline{{\texorpdfstring{\protect\hyperlink{ega4.c04.subsec.18-1}{18.1}}{18.1}}}{A remarkable equivalence of categories}{109}
\EGAIVcontentsline{{\texorpdfstring{\protect\hyperlink{ega4.c04.subsec.18-2}{18.2}}{18.2}}}{\'Etale coverings}{111}
\EGAIVcontentsline{{\texorpdfstring{\protect\hyperlink{ega4.c04.subsec.18-3}{18.3}}{18.3}}}{Finite and \'etale algebras}{114}
\EGAIVcontentsline{{\texorpdfstring{\protect\hyperlink{ega4.c04.subsec.18-4}{18.4}}{18.4}}}{Local structure of unramified and \'etale morphisms}{118}
\EGAIVcontentsline{{\texorpdfstring{\protect\hyperlink{ega4.c04.subsec.18-5}{18.5}}{18.5}}}{Henselian local rings}{125}
\EGAIVcontentsline{{\texorpdfstring{\protect\hyperlink{ega4.c04.subsec.18-6}{18.6}}{18.6}}}{Henselization}{135}
\EGAIVcontentsline{{\texorpdfstring{\protect\hyperlink{ega4.c04.subsec.18-7}{18.7}}{18.7}}}{Henselization and excellent rings}{142}
\EGAIVcontentsline{{\texorpdfstring{\protect\hyperlink{ega4.c04.subsec.18-8}{18.8}}{18.8}}}{Strictly local rings and strict Henselization}{144}
\EGAIVcontentsline{{\texorpdfstring{\protect\hyperlink{ega4.c04.subsec.18-9}{18.9}}{18.9}}}{Formal fibres of Noetherian Henselian rings}{150}
\EGAIVcontentsline{{\texorpdfstring{\protect\hyperlink{ega4.c04.subsec.18-10}{18.10}}{18.10}}}{Preschemes \'etale over a geometrically unibranch or normal prescheme}{157}
\EGAIVcontentsline{{\texorpdfstring{\protect\hyperlink{ega4.c04.subsec.18-11}{18.11}}{18.11}}}{Application to complete Noetherian local algebras over a field}{169}
\EGAIVcontentsline{{\texorpdfstring{\protect\hyperlink{ega4.c04.subsec.18-12}{18.12}}{18.12}}}{Applications of \'etale localization to quasi-finite morphisms (generalizations of earlier results)}{181}

\EGAIVcontentsline{\S19}{Regular immersions and normal flatness}{185}
\EGAIVcontentsline{{\texorpdfstring{\protect\hyperlink{ega4.c04.subsec.19-1}{19.1}}{19.1}}}{Properties of regular immersions}{185}
\EGAIVcontentsline{{\texorpdfstring{\protect\hyperlink{ega4.c04.subsec.19-2}{19.2}}{19.2}}}{Transversally regular immersions}{190}
\EGAIVcontentsline{{\texorpdfstring{\protect\hyperlink{ega4.c04.subsec.19-3}{19.3}}{19.3}}}{Relative complete intersections (flat case)}{194}
\EGAIVcontentsline{{\texorpdfstring{\protect\hyperlink{ega4.c04.subsec.19-4}{19.4}}{19.4}}}{Application: criteria of regularity and smoothness for blown-up preschemes}{198}
\EGAIVcontentsline{{\texorpdfstring{\protect\hyperlink{ega4.c04.subsec.19-5}{19.5}}{19.5}}}{Criteria of $M$-regularity}{204}
\EGAIVcontentsline{{\texorpdfstring{\protect\hyperlink{ega4.c04.subsec.19-6}{19.6}}{19.6}}}{Regular sequences relative to a filtered quotient module}{209}
\EGAIVcontentsline{{\texorpdfstring{\protect\hyperlink{ega4.c04.subsec.19-7}{19.7}}{19.7}}}{Hironaka's normal-flatness criterion}{212}
\EGAIVcontentsline{{\texorpdfstring{\protect\hyperlink{ega4.c04.subsec.19-8}{19.8}}{19.8}}}{Passage-to-the-projective-limit properties}{219}
\EGAIVcontentsline{{\texorpdfstring{\protect\hyperlink{ega4.c04.subsec.19-9}{19.9}}{19.9}}}{$F$-regular sequences and depth}{222}

\EGAIVcontentsline{\S20}{Meromorphic functions and pseudo-morphisms}{223}
\EGAIVcontentsline{{\texorpdfstring{\protect\hyperlink{ega4.c04.subsec.20-0}{20.0}}{20.0}}}{Introduction}{223}
\EGAIVcontentsline{{\texorpdfstring{\protect\hyperlink{ega4.c04.subsec.20-1}{20.1}}{20.1}}}{Meromorphic functions}{224}
\EGAIVcontentsline{{\texorpdfstring{\protect\hyperlink{ega4.c04.subsec.20-2}{20.2}}{20.2}}}{Pseudo-morphisms and pseudo-functions}{231}

\oldpage[IV]{343}
\EGAIVcontentsline{{\texorpdfstring{\protect\hyperlink{ega4.c04.subsec.20-3}{20.3}}{20.3}}}{Composition of pseudo-morphisms}{237}
\EGAIVcontentsline{{\texorpdfstring{\protect\hyperlink{ega4.c04.subsec.20-4}{20.4}}{20.4}}}{Properties of domains of definition of rational functions}{244}
\EGAIVcontentsline{{\texorpdfstring{\protect\hyperlink{ega4.c04.subsec.20-5}{20.5}}{20.5}}}{Relative pseudo-morphisms}{249}
\EGAIVcontentsline{{\texorpdfstring{\protect\hyperlink{ega4.c04.subsec.20-6}{20.6}}{20.6}}}{Relative meromorphic functions}{252}

\EGAIVcontentsline{\S21}{Divisors}{255}
\EGAIVcontentsline{{\texorpdfstring{\protect\hyperlink{ega4.c04.subsec.21-1}{21.1}}{21.1}}}{Divisors on a ringed space}{255}
\EGAIVcontentsline{{\texorpdfstring{\protect\hyperlink{ega4.c04.subsec.21-2}{21.2}}{21.2}}}{Divisors and invertible fractional Ideals}{258}
\EGAIVcontentsline{{\texorpdfstring{\protect\hyperlink{ega4.c04.subsec.21-3}{21.3}}{21.3}}}{Linear equivalence of divisors}{263}
\EGAIVcontentsline{{\texorpdfstring{\protect\hyperlink{ega4.c04.subsec.21-4}{21.4}}{21.4}}}{Inverse images of divisors}{265}
\EGAIVcontentsline{{\texorpdfstring{\protect\hyperlink{ega4.c04.subsec.21-5}{21.5}}{21.5}}}{Direct images of divisors}{267}
\EGAIVcontentsline{{\texorpdfstring{\protect\hyperlink{ega4.c04.subsec.21-6}{21.6}}{21.6}}}{Codimension-one cycle associated with a divisor}{270}
\EGAIVcontentsline{{\texorpdfstring{\protect\hyperlink{ega4.c04.subsec.21-7}{21.7}}{21.7}}}{Interpretation of positive codimension-one cycles in terms of sub-preschemes}{277}
\EGAIVcontentsline{{\texorpdfstring{\protect\hyperlink{ega4.c04.subsec.21-8}{21.8}}{21.8}}}{Divisors and normalization}{280}
\EGAIVcontentsline{{\texorpdfstring{\protect\hyperlink{ega4.c04.subsec.21-9}{21.9}}{21.9}}}{Divisors on one-dimensional preschemes}{284}
\EGAIVcontentsline{{\texorpdfstring{\protect\hyperlink{ega4.c04.subsec.21-10}{21.10}}{21.10}}}{Inverse and direct images of codimension-one cycles}{289}
\EGAIVcontentsline{{\texorpdfstring{\protect\hyperlink{ega4.c04.subsec.21-11}{21.11}}{21.11}}}{Factoriality of regular rings}{302}
\EGAIVcontentsline{{\texorpdfstring{\protect\hyperlink{ega4.c04.subsec.21-12}{21.12}}{21.12}}}{Van der Waerden's purity theorem for the ramification locus of a birational morphism}{304}
\EGAIVcontentsline{{\texorpdfstring{\protect\hyperlink{ega4.c04.subsec.21-13}{21.13}}{21.13}}}{Parafactorial pairs; parafactorial local rings}{313}
\EGAIVcontentsline{{\texorpdfstring{\protect\hyperlink{ega4.c04.subsec.21-14}{21.14}}{21.14}}}{The Ramanujam--Samuel theorem}{323}
\EGAIVcontentsline{{\texorpdfstring{\protect\hyperlink{ega4.c04.subsec.21-15}{21.15}}{21.15}}}{Relative divisors}{329}
\EGAIVcontentsline{}{Bibliography}{333}
\EGAIVcontentsline{}{Index of notation}{334}
\EGAIVcontentsline{}{Terminological index}{336}

\endgroup
% END NATIVE INPUT source_aligned/ega4-backmatter-contents
% BEGIN NATIVE INPUT source_aligned/ega4-errata-addenda-list3
% Source-aligned English projection of EGA IV, Part 4, pp. 345--361.
\clearpage
\oldpage[IV]{345}
\phantomsection
\addcontentsline{toc}{section}{Errata and Addenda (List 3)}
\section*{Errata and Addenda}
\begin{center}\large (List 3)\end{center}

\subsection*{A) Typographical errors}

\paragraph{(II, 3.3.2).}
On line 1 of p.~58, replace $\sh{S}$ by $\sh{M}$.  On line 6 of p.~58,
replace $S(n)$ by $M(n)$.

\paragraph{(II, 4.1.1).}
On line 8 of p.~71, replace $\sh{O}$-$P$-Module by
$\sh{O}_P$-Module.

\paragraph{(II, 5.4.5).}
On line 6 of p.~102, replace $u$ by $f\circ u$.

\paragraph{(II, 5.4.8).}
Before the fourth line from the bottom of p.~102, replace
$f\times 1_{S'}$ by $f_{(Y')}$ in the diagram.  On the fourth line from
the bottom of p.~102, replace $f\times 1_{S'}$ by $f_{(Y')}$.

\paragraph{(II, 6.4.4).}
On the first line at the top of p.~122, replace ``normal'' by
``integrally closed.''

\paragraph{(II, 6.4.7).}
On the fourteenth line from the bottom of p.~122, replace ``dimension''
by ``rank.''

\paragraph{(II, 6.5.2).}
On line 5 of p.~127, replace $(\sh{B}|U_\lambda)$ by $(B|U_\lambda)$.
On line 15 of p.~127, delete the first occurrence of
$(N_{\sh{B}/\sh{A}}\omega_{\lambda_\mu})$.

\paragraph{($0_{\mathrm{III}}$, 8.1.11).}
On the eighth line from the bottom of p.~8, replace $h_X(u)$ by
$h'_X(u)$.

\paragraph{($0_{\mathrm{III}}$, 9.1.2).}
On the tenth line from the bottom of p.~12, replace $\sh{F}$ by $X$.

\paragraph{($0_{\mathrm{III}}$, 10.2.1).}
On the sixth line from the bottom of p.~18, replace ``in addition if''
by ``if in addition.''

\paragraph{(III, 4.4.12).}
On line 10 of p.~138, replace ``chap.~V'' by ``chap.~IV (8.11.2).''

\paragraph{(III, 5.2.5).}
On line 12 of p.~153, replace (I, 5.5.4) by (II, 5.5.4).

\paragraph{(III, 7.4.7).}
On line 15 of p.~57, replace ``every $A$-module'' by
``every $A$-module $M$.''

\paragraph{(III, 7.8.5).}
On the eighteenth line from the bottom of p.~74, replace $F$ by $\sh{F}$.

\paragraph{(III, 7.8.9).}
On the ninth line from the bottom of p.~75, replace $M$ by $\sh{M}$.

\paragraph{(III, 7.9.11).}
On the seventeenth line from the bottom of p.~79, replace chap.~V by
chap.~VI.

\paragraph{($0_{\mathrm{IV}}$, 16.2.2.2).}
On line 13 of p.~26, replace $M'\cap\mathfrak q^nM'$ by
$M'\cap\mathfrak q^nM$.

\paragraph{($0_{\mathrm{IV}}$, 17.3.7).}
On the sixteenth line from the bottom of p.~50, replace (17.1.10) by
(17.3.3).

\paragraph{($0_{\mathrm{IV}}$, 18.1.3).}
In diagram (18.1.3.1), $p_2$ must be beside the arrow $G\to F$.

\paragraph{($0_{\mathrm{IV}}$, 19.8.7).}
On the seventh, eighth, and ninth lines from the bottom of p.~111,
replace (21.5.5) by (21.5.3).

\paragraph{($0_{\mathrm{IV}}$, 19.8.8).}
On the second line from the bottom of p.~112, replace $A$-module by
$B$-module.

\paragraph{($0_{\mathrm{IV}}$, 20.5.4).}
On the fifth line from the bottom of p.~130, replace
$\Omega_{B'/A'}$ by $\Omega^1_{B'/A'}$.

\paragraph{($0_{\mathrm{IV}}$, 20.7.8).}
On the ninth line from the bottom of p.~149, replace (18.4.3) by
(18.5.4).  On the fourth line from the bottom, replace (18.4.1) by
(18.5.1).

\paragraph{($0_{\mathrm{IV}}$, 21.5.3).}
On the last line of p.~166, replace the feminine form of ``ramified'' by
the masculine form.

\paragraph{($0_{\mathrm{IV}}$, 21.7.1).}
On line 4 of p.~170, replace (20.6.15.1) by (20.6.17.1).

\paragraph{($0_{\mathrm{IV}}$, 22.1.1).}
On line 12 of p.~183, replace (20.6.20) by (20.6.10).

\oldpage[IV]{346}
\paragraph{($0_{\mathrm{IV}}$, 22.3.3).}
On the thirteenth line from the bottom of p.~192, replace
$\widehat B_{\mathfrak p}$ by $\widehat B_{\mathfrak q}$.

\paragraph{($0_{\mathrm{IV}}$, 22.7.3).}
On the twenty-first line from the bottom of p.~212, replace
$\operatorname{Ker}(i)$ by $\operatorname{Ker}(i_0)$.  On the twelfth
line from the bottom of p.~212, replace
$(\mathfrak p/\mathfrak p^2)\otimes_L L$ by
$(\mathfrak p/\mathfrak p^2)\otimes_A L$.

\paragraph{($0_{\mathrm{IV}}$, 23.1.3).}
On line 4 of p.~215, replace $x'^q\in K'$ by $x'\in K'$.  On line 9,
replace $x'^q\in A'$ by $x'^q\in A$.

\paragraph{({\texorpdfstring{\protect\hyperlink{ega4.c04.statement.2-1-1}{IV, 2.1.1}}{IV, 2.1.1}}).}
On the fourth line from the bottom of p.~5, replace (III, 1.4.15.1) by
(III, 1.4.15.5).

\paragraph{({\texorpdfstring{\protect\hyperlink{ega4.c04.prop.2-4-3}{IV, 2.4.3}}{IV, 2.4.3}}).}
On line 12 of p.~20, replace
$(X_{\mathrm{red}}\times_{Y_{\mathrm{red}}}Y_{\mathrm{red}})_{\mathrm{red}}$
by
$(X_{\mathrm{red}}\times_{Y_{\mathrm{red}}}Y'_{\mathrm{red}})_{\mathrm{red}}$.

\paragraph{({\texorpdfstring{\protect\hyperlink{ega4.c04.prop.2-6-1}{IV, 2.6.1}}{IV, 2.6.1}}).}
On line 10 of p.~27, replace $X'\to X_{(S')}$ by $X'=X_{(S')}$.

\paragraph{({\texorpdfstring{\protect\hyperlink{ega4.c04.prop.4-2-1}{IV, 4.2.1}}{IV, 4.2.1}}).}
On the twentieth line from the bottom of p.~55, replace
$\operatorname{tr.deg}_kE$ by $\operatorname{tr.deg}_LE$ and
$\operatorname{tr.deg}_kL(t)$ by $\operatorname{tr.deg}_LL(t)$; on the
nineteenth line from the bottom, interchange ``first'' and ``second.''

\paragraph{({\texorpdfstring{\protect\hyperlink{ega4.c04.lem.4-4-2}{IV, 4.4.2}}{IV, 4.4.2}}).}
On the last line of p.~59, replace $X$ by $X'$.

\paragraph{({\texorpdfstring{\protect\hyperlink{ega4.c04.prop.4-7-3}{IV, 4.7.3}}{IV, 4.7.3}}).}
On the tenth line from the bottom of p.~75, replace ``prescheme'' by
``$k$-prescheme.''

\paragraph{({\texorpdfstring{\protect\hyperlink{ega4.c04.prop.5-6-1}{IV, 5.6.1}}{IV, 5.6.1}}).}
On line 16 of p.~98, replace the subscript $k(\mathfrak p($ by
$k(\mathfrak p)$.

\paragraph{({\texorpdfstring{\protect\hyperlink{ega4.c04.prop.5-6-5}{IV, 5.6.5}}{IV, 5.6.5}}).}
On line 20 of p.~99, delete the second closing parenthesis in $k(\xi))$.

\paragraph{({\texorpdfstring{\protect\hyperlink{ega4.c04.lem.5-6-11-1}{IV, 5.6.11.1}}{IV, 5.6.11.1}}).}
On line 16 of p.~102, replace $K$ by $\mathfrak K$.

\paragraph{({\texorpdfstring{\protect\hyperlink{ega4.c04.cor.5-11-2}{IV, 5.11.2}}{IV, 5.11.2}}).}
In the footnote on p.~123, replace ``5.11.2)'' by ``(5.11.2).''

\paragraph{({\texorpdfstring{\protect\hyperlink{ega4.c04.prop.6-1-5}{IV, 6.1.5}}{IV, 6.1.5}}).}
On the second line from the bottom of p.~136, replace (5.5.1.2) by
($0$, 16.3.9.2).

\paragraph{({\texorpdfstring{\protect\hyperlink{ega4.c04.prop.6-3-1}{IV, 6.3.1}}{IV, 6.3.1}}).}
On line 9 of p.~139, replace $y\in\mathfrak m$ by $y\in\mathfrak n$.

\paragraph{({\texorpdfstring{\protect\hyperlink{ega4.c04.cor.6-7-2}{IV, 6.7.2}}{IV, 6.7.2}}).}
On line 18 of p.~146, replace $\sh{O}'_{x'}$ by $\sh{O}_{x'}$.

\paragraph{({\texorpdfstring{\protect\hyperlink{ega4.c04.statement.6-10-1}{IV, 6.10.1}}{IV, 6.10.1}}).}
On line 13 of p.~155, replace
$\sh{F}\otimes_{\sh{O}_Y}\sh{O}_Y$ by
$\sh{F}\otimes_{\sh{O}_X}\sh{O}_Y$.

\paragraph{({\texorpdfstring{\protect\hyperlink{ega4.c04.prop.6-14-1}{IV, 6.14.1}}{IV, 6.14.1}}).}
On line 21 of p.~170, replace ``no.~3, prop.~12'' by
``no.~2, prop.~8.''

\paragraph{({\texorpdfstring{\protect\hyperlink{ega4.c04.statement.6-15-2}{IV, 6.15.2}}{IV, 6.15.2}}).}
On line 13 of p.~177, replace $t(t-v)$ by $t^2(t-v)^2$.

\paragraph{({\texorpdfstring{\protect\hyperlink{ega4.c04.prop.6-15-7}{IV, 6.15.7}}{IV, 6.15.7}}).}
On the sixth line from the bottom of p.~179, replace
$X\otimes_Ak'$ by $X\otimes_kk'$.

\paragraph{({\texorpdfstring{\protect\hyperlink{ega4.c04.item.7-5-4}{IV, 7.5.4}}{IV, 7.5.4}}).}
On the fourteenth line from the bottom of p.~206, replace
$\operatorname{Spec}(\widehat B\otimes k)$ by
$\operatorname{Spec}(\widehat B\otimes_{\widehat A}k)$.

\paragraph{({\texorpdfstring{\protect\hyperlink{ega4.c04.item.7-8-4}{IV, 7.8.4}}{IV, 7.8.4}}).}
On the tenth line from the bottom of p.~216, move the opening parenthesis
so that the phrase reads ``satisfying ({\hyperlink{ega4.c04.def.7-8-2}{7.8.2}}, (ii) and (iii)).''

\paragraph{(IV, Part 2, Bibliography).}
On line 6 of p.~224, replace \emph{Amer. J. of Math.} by
\emph{Annals of Math.}

\paragraph{({\texorpdfstring{\protect\hyperlink{ega4.c04.cor.8-5-8}{IV, 8.5.8}}{IV, 8.5.8}}).}
On line 6 of p.~24, replace $\operatorname{Ker}g_\mu$ at the end of the
line by $\operatorname{Ker}g_\lambda$.

\paragraph{({\texorpdfstring{\protect\hyperlink{ega4.c04.prop.8-10-1}{IV, 8.10.1}}{IV, 8.10.1}}).}
On the thirteenth line from the bottom of p.~36, replace (2.3.11) by
{\hyperlink{ega4.c04.cor.8-2-11}{(8.2.11)}}.

\paragraph{({\texorpdfstring{\protect\hyperlink{ega4.c04.prop.10-2-1}{IV, 10.2.1}}{IV, 10.2.1}}).}
On the fourteenth line from the bottom of p.~98, replace $y$ by $Y$.

\paragraph{({\texorpdfstring{\protect\hyperlink{ega4.c04.statement.10-4-11}{IV, 10.4.11}}{IV, 10.4.11}}).}
On the tenth line from the bottom of p.~104, replace (17.9.7) by
{\hyperlink{ega4.c04.prop.17-9-6}{(17.9.6)}}.

\paragraph{({\texorpdfstring{\protect\hyperlink{ega4.c04.thm.11-2-9}{IV, 11.2.9}}{IV, 11.2.9}}).}
On the third line from the bottom of p.~126, replace $A_\lambda$ by
$A_\mu$.

\paragraph{({\texorpdfstring{\protect\hyperlink{ega4.c04.prop.11-3-4}{IV, 11.3.4}}{IV, 11.3.4}}).}
On line 16 of p.~134, replace $F$ by $\sh{F}$.

\paragraph{({\texorpdfstring{\protect\hyperlink{ega4.c04.thm.11-3-10}{IV, 11.3.10}}{IV, 11.3.10}}).}
On the ninth line from the bottom of p.~139, replace (11.3.7) by
{\hyperlink{ega4.c04.thm.11-3-10}{(11.3.10)}}.

\paragraph{({\texorpdfstring{\protect\hyperlink{ega4.c04.statement.11-7-5}{IV, 11.7.5}}{IV, 11.7.5}}).}
On line 11 of p.~159, replace (6.5.11) by (6.15.11).

\paragraph{({\texorpdfstring{\protect\hyperlink{ega4.c04.thm.12-1-1}{IV, 12.1.1}}{IV, 12.1.1}}).}
On line 13 of p.~174, add at the end: ``for $k\geq 1$.''

\paragraph{({\texorpdfstring{\protect\hyperlink{ega4.c04.thm.14-2-1}{IV, 14.2.1}}{IV, 14.2.1}}).}
On line 11 of p.~202, replace (6.1.2) by (6.1.1).

\oldpage[IV]{347}
\paragraph{({\texorpdfstring{\protect\hyperlink{ega4.c04.prop.14-3-8}{IV, 14.3.8}}{IV, 14.3.8}}).}
On the tenth line from the bottom of p.~206, replace $(f^{-1}y)$ by
$f^{-1}(y)$.

\paragraph{({\texorpdfstring{\protect\hyperlink{ega4.c04.rem.14-5-7}{IV, 14.5.7}}{IV, 14.5.7}}).}
On line 16 of p.~218, replace (14.5.5) by (14.5.6).

\paragraph{({\texorpdfstring{\protect\hyperlink{ega4.c04.cor.14-5-11}{IV, 14.5.11}}{IV, 14.5.11}}).}
In the diagram at the top of p.~223, every horizontal arrow must point
from right to left.

\paragraph{({\texorpdfstring{\protect\hyperlink{ega4.c04.prop.15-4-2}{IV, 15.4.2}}{IV, 15.4.2}}).}
On line 23 of p.~230, replace (6.1.4) by (6.1.5).

\paragraph{({\texorpdfstring{\protect\hyperlink{ega4.c04.cor.15-5-2}{IV, 15.5.2}}{IV, 15.5.2}}).}
On line 16 of p.~233, replace $f^{-1}(y)$ by $f^{-1}(y')$.

\paragraph{({\texorpdfstring{\protect\hyperlink{ega4.c04.item.15-6-10}{IV, 15.6.10}}{IV, 15.6.10}}).}
On line 6 of p.~242, replace ${}^{-1}(y')$ by $f^{-1}(y')$.  On line
11, replace ``$\alpha$) and $\beta$)'' by ``$\alpha$) or $\beta$).''

\subsection*{B) Textual modifications}

\paragraph{(Err\textsubscript{IV}, 1).}
In ($0_{\mathrm{III}}$, 11.1.6), the statement (which reproduces a
statement of (G, I, 4.4.1)) is false when $q(n)$ is an arbitrary integer;
it is correct only when $q(n)=0$ or $q(n)=n$.  Indeed, in the first case
the kernel of $d_r^{p,0}$ is $E_r^{p,0}$, since the image of
$d_r^{p,0}$ lies in $E_r^{p+r,-r+1}=0$ by hypothesis; likewise the
image of $d_r^{p-r,r-1}$ is zero since $E_r^{p-r,r-1}=0$.  Hence
$E_{r+1}^{p,0}=E_r^{p,0}$.  The case $q(n)=n$ for every $n$ is handled
in the same way.

\paragraph{(Err\textsubscript{IV}, 2).}
In (III, 3.4), delete lines 2 and 3 of p.~119.

\oldpage[IV]{348}
\paragraph{(Err\textsubscript{IV}, 3).}
In (III, 4.4.9), on the fifteenth line from the bottom of p.~137,
replace ``finite'' by ``finite as sets.''

\paragraph{(Err\textsubscript{IV}, 4).}
In (III, 7.8.10), on line 9 of p.~76, replace ``simple'' by ``smooth.''

\paragraph{(Err\textsubscript{IV}, 5).}
In ($0_{\mathrm{IV}}$, 16.1.4), on the eighth line from the bottom of
p.~23, replace ``To say that $A$ is catenary'' by ``To say that an
\emph{integral} ring $A$ is catenary.''

\paragraph{(Err\textsubscript{IV}, 6).}
In ($0_{\mathrm{IV}}$, 16.1.9), the hypothesis that $A$ is Noetherian
may be deleted.  Indeed, for every $u\in B_M$, $A[u]$ is an
$A$-subalgebra of $C$, so $u$ is integral over $A$ (Bourbaki,
\emph{Alg. comm.}, chap.~V, \S1, no.~1, corollary to prop.~1), and
hence also over $A_M$; one concludes again by {\hyperlink{ega4.c04.prop.16-1-5}{(16.1.5)}}.

\paragraph{(Err\textsubscript{IV}, 7).}
In ($0_{\mathrm{IV}}$, 16.5.12), note that relation (16.5.11.1) is
equivalent to the same relation after replacing $A$ by the quotient rings
$A/\mathfrak r$, where $\mathfrak r$ ranges over the minimal prime ideals
of $A$.  Thus in the proof of {\hyperlink{ega4.c04.statement.16-5-12}{(16.5.12)}} one may reduce to the case in
which $A$ is integral and apply (16.1.4.2).

\paragraph{(Err\textsubscript{IV}, 8).}
Delete ($0_{\mathrm{IV}}$, 19.3.12), whose proof is erroneous (the sets
$W_\lambda$, $U_\lambda$, and $V_\lambda$ are not linear varieties).

\paragraph{(Err\textsubscript{IV}, 9).}
In ($0_{\mathrm{IV}}$, 19.8.3), on the sixth and eighth lines from the
bottom of p.~110, replace ``homomorphism'' by ``local homomorphism'';
make the same replacement in ($0_{\mathrm{IV}}$, 19.8.4), line 1 of
p.~111.

\paragraph{(Err\textsubscript{IV}, 10).}
In ($0_{\mathrm{IV}}$, 20.1.2), on the second line from the bottom of
p.~117, replace $a\in A$ by $b\in B$ and $aD$ by $bD$; on the last line,
replace $A$-module by $B$-module.

\paragraph{(Err\textsubscript{IV}, 11).}
In ($0_{\mathrm{IV}}$, 20.2.1), on line 8 of p.~119, replace
$A$-modules by $B$-modules.  On lines 21 and 22, replace
``di-homomorphisms of modules (relative to $u$)'' by ``homomorphism of
$B$-modules.''

\paragraph{(Err\textsubscript{IV}, 12).}
In ($0_{\mathrm{IV}}$, 20.4.6.1), replace $p_2-p_1$ by $p_1-p_2$; on
the tenth line from the bottom of p.~125, replace
$x\otimes1-1\otimes x$ by $1\otimes x-x\otimes1$.

\paragraph{(Err\textsubscript{IV}, 13).}
In ($0_{\mathrm{IV}}$, 20.4.8), on line 14 of p.~126, replace
$p_2=p_1-d_{B/A}$ by $p_1=p_2-d_{B/A}$.

\paragraph{(Err\textsubscript{IV}, 14).}
In ($0_{\mathrm{IV}}$, 20.4.13), replace the sixth through eleventh
lines from the bottom of p.~127 by the following text:

\medskip
\noindent
(iv) Let $A$ and $B$ be integral rings such that $A\subset B$, with
$B$ integral over $A$, and suppose that the fraction field of $B$ is a
separable extension of that of $A$.  Then the $B$-module
$\Omega^1_{B/A}$ is a torsion module.  Indeed, for every $x\in B$, let
$f$ be the minimal polynomial of $x$ over the fraction field $K$ of
$A$; there is a nonzero $a\in A$ such that $af[T]\in A[T]$.  Since $x$
is separable over $K$, $f'(x)\neq0$, while the relation $af(x)=0$
gives $af'(x)d_{B/A}x=0$; the assertion follows from {\hyperlink{ega4.c04.cor.20-4-7}{(20.4.7)}}.

\paragraph{(Err\textsubscript{IV}, 15).}
In ($0_{\mathrm{IV}}$, 22.2.7), on line 5 of p.~189, replace
``formally smooth'' by ``formally smooth for the preadic topologies.''

\paragraph{(Err\textsubscript{IV}, 16).}
In the Summary of Chapter IV, on the fourth line from the bottom of
p.~222, add ``and strictly local rings.''  Replace the last three lines
by:
\begin{quote}
\S19. Regular immersions and normal flatness.\\
\S20. Meromorphic functions and pseudo-morphisms.\\
\S21. Divisors.
\end{quote}
On line 16 of p.~223, replace ``of finite presentation'' by ``locally of
finite presentation''; on line 19, replace ``as'' by ``as being.''

\paragraph{(Err\textsubscript{IV}, 17).}
In (IV, 3.1.10), add after line 16 of p.~39:

\begin{remark}[3.1.10.2]
\label{IV.3.1.10.2}
Let $f:X\to Y$ be a quasi-compact and quasi-separated morphism, and let
$\sh{F}$ be a quasi-coherent $\sh{O}_X$-Module, so that $f_*(\sh{F})$
is a quasi-coherent $\sh{O}_Y$-Module {\hyperlink{ega4.c04.statement.1-7-4}{(1.7.4)}}.  Then
\begin{align}
 \operatorname{Ass}(f_*(\sh{F}))&\subset f(X), \label{IV.3.1.10.3}
\tag{3.1.10.3}\\
 \operatorname{Ass}(f_*(\sh{F}))&\subset f(\operatorname{Ass}(\sh{F}))
       \quad\text{if $X$ is locally Noetherian}. \label{IV.3.1.10.4}
\tag{3.1.10.4}
\end{align}

Indeed, let $y\in\operatorname{Ass}(f_*(\sh{F}))$.  By {\hyperlink{ega4.c04.lem.2-3-1}{(2.3.1)}} and the
flatness of $\operatorname{Spec}(\sh{O}_{Y,y})\to Y$,
$(f_*(\sh{F}))_y$ equals $(f'_*(\sh{F}'))_y$, where
\[
 f':X\times_Y\operatorname{Spec}(\sh{O}_{Y,y})
       \longrightarrow\operatorname{Spec}(\sh{O}_{Y,y})
\]
is obtained from $f$ by base change and
$\sh{F}'=\sh{F}\otimes_{\sh{O}_Y}\sh{O}_{Y,y}$.  In view of (I, 3.6.5)
and the definition of $\operatorname{Ass}(\sh{F}')$, we may therefore
replace $Y$ by $\operatorname{Spec}(\sh{O}_{Y,y})$, $f$ by $f'$, and
$\sh{F}$ by $\sh{F}'$.

By definition there is a nonzero element
$t_y\in(f_*(\sh{F}))_y$ annihilated by $\mathfrak m_y$.  Since $Y$ is
local, $(f_*(\sh{F}))_y$ identifies with
$\Gamma(Y,f_*(\sh{F}))$, and hence $t_y$ identifies with a section
$s\in\Gamma(X,\sh{F})$.  The annihilator Ideal of $s$ in $\sh{O}_X$
contains $\mathfrak m_y\sh{O}_X$, so the support of $s$ is contained in
$f^{-1}(y)$; if this inverse image were empty then $s=0$, a
contradiction.  If moreover $X$ is locally Noetherian, the coherent
sub-$\sh{O}_X$-Module $\sh{O}_Xs$ of $\sh{F}$ is nonzero, and therefore
$\operatorname{Ass}(\sh{O}_Xs)\neq\varnothing$ {\hyperlink{ega4.c04.cor.3-1-5}{(3.1.5)}}.  But
\[
 \operatorname{Ass}(\sh{O}_Xs)\subset\operatorname{Supp}(\sh{O}_Xs)
 \subset f^{-1}(y),
 \qquad
 \operatorname{Ass}(\sh{O}_Xs)\subset\operatorname{Ass}(\sh{F})
\]
by {\hyperlink{ega4.c04.prop.3-1-7}{(3.1.7)}}, and hence $y\in f(\operatorname{Ass}(\sh{F}))$.
\end{remark}

\long\def\EGAIVErrataThirtyTwoThroughFiftyThree{%
\paragraph{(Err\textsubscript{IV}, 32).}
In (IV, 11.3.14), on line 7 of p.~141, insert ``(i)'' before ``Let.''
After line 9 of p.~141, add:

\smallskip
\noindent
\textup{(ii)} Conversely, let $f:X\to Y$ be a flat morphism.  If $X$
is unibranch (respectively geometrically unibranch) at a point $x$, then
$Y$ is unibranch (respectively geometrically unibranch) at the point
$y=f(x)$.

At the beginning of line 10 of p.~141, insert ``(i).''  After line 20
of p.~141, add:

\smallskip
\noindent
\textup{(ii)} Put $Y_1=Y_{\mathrm{red}}$ and
$X_1=X\times_Y Y_1$.  The underlying space of $X_1$ is that of $X$,
and the morphism $f_1:X_1\to Y_1$ obtained from $f$ by base change is
flat.  We may therefore restrict to the case
$Y=\operatorname{Spec}(A)$, with $A$ a reduced local ring and $y$ the
closed point of $Y$.  Put $B=\sh{O}_{X,x}$.  By hypothesis
$B_{\mathrm{red}}$ is unibranch, so $\operatorname{Spec}(B)$ has only
one maximal point; flatness {\hyperlink{ega4.c04.prop.2-3-4}{(2.3.4)}} then gives the same conclusion for
\oldpage[IV]{356}
$\operatorname{Spec}(A)$.  Since $A$ is reduced, it is therefore
integral.

Let $K$ be the fraction field of $A$, and let $A'$ be its integral
closure.  By flatness, $B\otimes_A A'$ identifies with a subring of
$B\otimes_A K$ containing $B$.  Since the nonzero elements of $A$ are
$B$-regular by flatness, $B\otimes_A K$ identifies with a subring of the
total ring of fractions $R$ of $B$.  To show that $A$ is unibranch, it
is enough to show that $A'$ cannot have two distinct maximal ideals
(necessarily lying over $y$; Bourbaki, \emph{Alg. comm.}, chap.~V,
\S2, no.~1, prop.~1).  If it did, then
$B'=B\otimes_A A'$ would have two distinct prime ideals over $x$
(I, 3.4.7), and we need only show that this is impossible.

Let $\mathfrak N$ be the nilradical of $B$.  If $S$ is the set of
regular elements of $B$, then $R=S^{-1}B$, and
$S^{-1}\mathfrak N=R\mathfrak N$ is the nilradical of $R$.  Hence
\[
 R/R\mathfrak N=S^{-1}B/S^{-1}\mathfrak N
 =S^{-1}(B/\mathfrak N)
\]
is contained in the fraction field $L$ of the integral local ring
$B/\mathfrak N=B_{\mathrm{red}}$.  The same is therefore true of
\[
 B'_{\mathrm{red}}=B'/(B'\cap R\mathfrak N),
\]
which is integral over $B_{\mathrm{red}}$.  Since
$B'_{\mathrm{red}}$ is contained in the integral closure of
$B_{\mathrm{red}}$, that ring is local by hypothesis; hence
$B'_{\mathrm{red}}$, and therefore $B'$, is local.  This proves the
assertion.

If $B$ is moreover geometrically unibranch, the residue field of
$B'_{\mathrm{red}}$, being contained in the residue field of the
integral closure of $B_{\mathrm{red}}$, is a radicial extension of
$k(x)$.  But the residue field of $B'_{\mathrm{red}}$ is also that of
$B'=B\otimes_A A'$, hence is
$k(x)\otimes_{k(y)}k'$, where $k'$ is the residue field of $A'$.  It
follows from {\hyperlink{ega4.c04.prop.2-6-1}{(2.6.1)}} that $k'/k(y)$ is radicial, so $A$ is geometrically
unibranch as well.

\paragraph{(Err\textsubscript{IV}, 33).}
The example in ({\hyperlink{ega4.c04.item.12-1-2}{IV, 12.1.2}}, (i)) is incorrect: the morphism considered
there is open.  An example of a flat morphism $f:X\to Y$ whose
restriction to an irreducible component of $X$ is not open may be
constructed as follows.  Let
$Z=\operatorname{Spec}(\mathbb C[U,V])$ be the affine plane over the
complex numbers, and let $Y$ be the scheme obtained by identifying two
distinct closed points $a,b$ of $Z$.  Thus $Y=\operatorname{Spec}(B)$,
where $B$ is the subring of $\mathbb C[U,V]$ formed by the polynomials
taking the same value at $a$ and $b$ (the set of closed points of $Z$
being identified with $\mathbb C^2$).  Plainly $B$ is a complex vector
subspace of codimension one in $\mathbb C[U,V]$, so
$\mathbb C[U,V]$ is a finite $B$-module.  Consequently $Z\to Y$ is
finite, and the inverse image of the point $c$ of $Y$ obtained by
identifying $a$ and $b$ is $\{a,b\}$.

Now form $X$ by gluing two copies $Z_1,Z_2$ of $Z$ so that, if
$a_i,b_i$ ($i=1,2$) are the points corresponding to $a,b$, then $a_1$
is glued to $b_2$ and $b_1$ to $a_2$ (cf. ({\hyperlink{ega4.c04.item.6-5-5}{6.5.5}}, (ii))).  It is easy
to see (\emph{loc. cit.}) that the resulting morphism $X\to Y$ is flat
and finite (indeed, \'{e}tale), and that $Z_1$ and $Z_2$ identify with
the two irreducible components of $X$.  If $D$ is a line in $Z_1$
through $a_1$ but not through $b_1$, the image in $Y$ of the open set
$Z_1\setminus D$ contains the image of $b_1$, namely $c$, but does not
contain the image of the generic point of $D$.  That image is therefore
not a neighbourhood of $c$.

\paragraph{(Err\textsubscript{IV}, 34).}
In the title of (IV, 13.2), on the fifth line from the bottom of p.~190,
replace ``equidimensional'' by ``locally equidimensional.''  Likewise,
in (IV, 13.2.2), on line 4 of p.~191, replace both occurrences of
``equidimensional'' by ``locally equidimensional.''
\oldpage[IV]{357}
On line 5 of
p.~191, add:
\begin{quote}
If, in addition, the fibres $f^{-1}(f(x))$ of $f$ are equidimensional
($0_{\mathrm{IV}}$, 14.1.3), then $f$ is said to be
\emph{equidimensional}.
\end{quote}
On the sixth line from the bottom of p.~194, replace
``equidimensional'' by ``locally equidimensional.''

\paragraph{(Err\textsubscript{IV}, 35).}
In the title of (IV, 13.3), on the fourth line from the bottom of p.~194,
replace ``equidimensional'' by ``locally equidimensional.''  In
(IV, 13.3.2), on the sixth line from the bottom of p.~196, make the same
replacement twice.  On the fifth line from the bottom of p.~196, add:
\begin{quote}
If, in addition, the fibres $f^{-1}(f(x))$ of $f$ are equidimensional
($0_{\mathrm{IV}}$, 14.1.3), then $f$ is said to be
\emph{equidimensional}.
\end{quote}
In (IV, 13.3.7), on line 14 of p.~198, replace ``equidimensional'' by
``locally equidimensional.''  In ({\hyperlink{ega4.c04.statement.13-3-9}{IV, 13.3.9}}), on the third line from
the bottom of p.~198, replace ``equidimensionality'' by ``local
equidimensionality.''  On lines 10, 12, 20, and 24 of p.~199, replace
``equidimensional'' by ``locally equidimensional.''

\paragraph{(Err\textsubscript{IV}, 36).}
In the introduction to \S14, on the sixth line from the bottom of
p.~199, replace ``equidimensional'' by ``locally equidimensional.''

\paragraph{(Err\textsubscript{IV}, 37).}
In (IV, 14.4.1), on the fifth line from the bottom of p.~209, replace
``locally of finite type'' by ``locally of finite presentation.''  The
present statement of ({\hyperlink{ega4.c04.thm.14-4-1}{IV, 14.4.1}}) is false even when $f:X\to Y$ is a
closed immersion, $X$ consists of a single closed point $y$ of $Y$, and
$\sh{O}_{Y,y}$ is a field.  Indeed, there exist rings $A$ such that
$Y=\operatorname{Spec}(A)$ is compact, has no isolated points, and is
totally disconnected, while every prime ideal $\mathfrak p$ of $A$ is
maximal and $A_\mathfrak p$ is a field (Bourbaki,
\emph{Alg. comm.}, chap.~II, \S4, exercises 16 and 17).  Clearly the
closed immersion $\{y\}\to Y$ is not open, although it satisfies
conditions b) and b') of ({\hyperlink{ega4.c04.thm.14-4-1}{IV, 14.4.1}}).

\paragraph{(Err\textsubscript{IV}, 38).}
In (IV, 14.4.1.3), on line 14 of p.~211, replace ``locally quasi-finite
and dominant'' by ``locally quasi-finite, locally of finite presentation,
and dominant.''

\paragraph{(Err\textsubscript{IV}, 39).}
In (IV, 14.4.2), on the tenth line from the bottom of p.~211, replace
``locally of finite type'' by ``locally of finite presentation.''

\paragraph{(Err\textsubscript{IV}, 40).}
In (IV, 14.4.3), on line 18 of p.~212, replace ``locally of finite type''
by ``locally of finite presentation.''

\paragraph{(Err\textsubscript{IV}, 41).}
In ({\hyperlink{ega4.c04.cor.14-4-4}{IV, 14.4.4}}), on the fifteenth and sixteenth lines from the bottom of
p.~212, replace ``locally of finite type'' by ``locally of finite
presentation.''  On the tenth and eleventh lines from the bottom,
replace ``equidimensional'' by ``locally equidimensional.''

\paragraph{(Err\textsubscript{IV}, 42).}
In (IV, 14.4.8), after line 1 of p.~214, add:
\begin{quote}
If, in addition, the set of irreducible components of $Y$ is locally
finite, these two conditions are equivalent to the following ones:
\end{quote}

\paragraph{(Err\textsubscript{IV}, 43).}
In (IV, 14.4.10), on line 15 of p.~215, replace ``arbitrary'' by ``whose
set of irreducible components is locally finite.''  On line 16 of
p.~215, replace ``of finite type'' by ``of finite presentation.''

\paragraph{(Err\textsubscript{IV}, 44).}
In (IV, 15.5.6.1), on the fifteenth line from the bottom of p.~234,
after the first sentence, add:

\begin{quote}
It is enough to consider the case in which every irreducible closed
subset of $X$ has a unique generic point.  Otherwise, consider the
space $Y=X/R$,
\oldpage[IV]{358}
the quotient of $X$ by the equivalence relation
\[
 \overline{\{x\}}=\overline{\{x'\}}.
\]
It follows immediately from this definition that every open subset of
$X$ is saturated for $R$; hence the topology of $X$ is the inverse
image of that of $Y$ under the canonical map $\pi:X\to X/R$.  A dense
open subset of $X$ is of the form $\pi^{-1}(U)$, where $U$ is dense and
open in $Y$; consequently a nowhere dense closed subset $Z$ of $X$ is
of the form $\pi^{-1}(S)$, where $S=\pi(Z)$ is nowhere dense and closed
in $Y$.  Thus $X\setminus Z$ is connected if and only if
$Y\setminus S$ is connected.

Similarly, $\pi(T_x)$ is the set of generalizations of $\pi(x)$, and
$T_x\setminus\{x\}$ is connected if and only if
$\pi(T_x)\setminus\{\pi(x)\}$ is connected.  Finally, by construction,
$Y$ is a locally Noetherian Kolmogorov space, so every irreducible
closed subset of $Y$ has a unique generic point.  Proving the lemma for
$X$ is therefore equivalent to proving it for $Y$, and we may suppose
that every irreducible closed subset of $X$ has a unique generic point.
\end{quote}

\paragraph{(Err\textsubscript{IV}, 45).}
In (IV, 15.6.1), on the eighteenth line from the bottom of p.~236,
delete ``locally Noetherian'' (this hypothesis is not used in the
proof).

\paragraph{(Err\textsubscript{IV}, 46).}
In (IV, 15.6.3), on line 11 of p.~237, replace ``suppose that for every
$y\in Y$'' by ``suppose that $f$ is locally of finite presentation and
that for every $y\in Y$.''  Replace the proof (lines 14--24 of p.~237)
by the following:

\begin{quote}
By {\hyperlink{ega4.c04.prop.1-10-3}{(1.10.3)}}, it is enough to prove that for every $x\in X_y^0$ and
every generalization $y'$ of $y$, there exists a generalization $x'$ of
$x$ such that $f(x')=y'$.  By {\hyperlink{ega4.c04.cor.15-6-2}{(15.6.2)}}, it is enough to take for $x'$
the generic point of $X_{y'}^0$.
\end{quote}

\paragraph{(Err\textsubscript{IV}, 47).}
In (IV, 15.6.7), on line 2 of p.~241, after ``Noetherian,'' add ``or if
condition $\beta$) of ({\hyperlink{ega4.c04.prop.15-6-6}{15.6.6}}, (iii)) is satisfied.''  After line 9 of
p.~241, add:
\begin{quote}
When condition $\beta$) of ({\hyperlink{ega4.c04.prop.15-6-6}{15.6.6}}, (iii)) is satisfied, the implication
a)$\Rightarrow$b') follows from ({\hyperlink{ega4.c04.prop.15-6-6}{15.6.6}}, (ii)), and the implication
b')$\Rightarrow$a) follows from ({\hyperlink{ega4.c04.prop.15-6-6}{15.6.6}}, (iii)).
\end{quote}

\paragraph{(Err\textsubscript{IV}, 48).}
In (IV, 20.1.5), line 13 of p.~228; (IV, 20.1.6), line 17 of p.~228;
and (IV, 20.1.7), the thirteenth line from the bottom of p.~228, delete
``strictly.''

\paragraph{(Err\textsubscript{IV}, 49).}
In (IV, 20.2.13), after line 6 of p.~236, add:
\begin{quote}
In particular, if $X$ is reduced and has only finitely many irreducible
components, the notion of torsion-free $\sh{O}_X$-Module defined in
({\hyperlink{ega4.c04.statement.20-1-5}{IV, 20.1.5}}) coincides with that of (I, 7.4.1).
\end{quote}

\paragraph{(Err\textsubscript{IV}, 50).}
In (IV, 20.2.14), on lines 18 and 25 of p.~236, delete ``strictly.''

\paragraph{(Err\textsubscript{IV}, 51).}
In (IV, 20.6.2), on the thirteenth line from the bottom of p.~252,
delete ``strictly.''

\paragraph{(Err\textsubscript{IV}, 52).}
In (IV, 20.6.4), on lines 23 and 22 of p.~253, delete ``strictly.''

\paragraph{(Err\textsubscript{IV}, 53).}
In (IV, 21.4), after line 23 of p.~267, add:

\def\EGAIVSuppressErratumLabels{1}
% BEGIN NATIVE INPUT source_aligned/ega4-21-erratum53-insertion
% EGA IV, Errata and Addenda, List 3, no. 53.
% This component is inserted after 21.4.8 in the corrected reading text
% and repeated verbatim under Err_IV, 53 in the printed errata appendix.

The following proposition and its corollaries were pointed out to us by
M.~Raynaud.

\begin{proposition}[21.4.9]
\ifdefined\EGAIVSuppressErratumLabels\else\label{IV.21.4.9}\fi
Let $X$ and $Y$ be two locally Noetherian preschemes, and let
$f:X\to Y$ be a faithfully flat morphism.  Suppose in addition either
that $f$ is quasi-compact or that $f$ is open.  A divisor $D'$ on $X$
is of the form $f^*(D)$, with $D$ a divisor on $Y$, if and only if the
following condition holds: for every $y\in Y$ such that
$\operatorname{prof}(\sh{O}_{Y,y})\leq 1$, put
\[
 Y_1=\operatorname{Spec}(\sh{O}_{Y,y}),\qquad
 X_1=X\times_Y Y_1,
 \qquad f_1=f_{(Y_1)}:X_1\to Y_1;
\]
then the inverse image $D_1'$ of $D'$ on $X_1$ is of the form
$f_1^*(D_1)$ for some $D_1\in\operatorname{Div}(Y_1)$.
\end{proposition}

\oldpage[IV]{359}
\begin{proof}
The necessity of the condition is immediate from the transitivity of
inverse images of divisors, all the morphisms in question being flat
{\hyperlink{ega4.c04.statement.21-4-4}{(21.4.4)}}.  For every open subset $U$ of $Y$, write $f_U$ for the
restriction $f^{-1}(U)\to U$ of $f$.  By {\hyperlink{ega4.c04.prop.21-4-8}{(21.4.8)}}, the map
\[
 D_U\longmapsto f_U^*(D_U)
\]
from $\operatorname{Div}(U)$ to
$\operatorname{Div}(f^{-1}(U))$ is injective.  Thus, if $U_1,U_2$ are
two open subsets of $Y$ and
\[
 D'|f^{-1}(U_1)=f_{U_1}^*(D_1),\qquad
 D'|f^{-1}(U_2)=f_{U_2}^*(D_2),
\]
then necessarily
$D_1|U_1\cap U_2=D_2|U_1\cap U_2$.  It follows at once that there is a
largest open subset $U\subset Y$ for which
\[
 D'|f^{-1}(U)=f_U^*(D_U)
 \quad\text{with}\quad D_U\in\operatorname{Div}(U).
\]
We must show that $U=Y$.  Suppose otherwise, and let $y$ be a maximal
point of $Y\setminus U$.  We shall find an open neighbourhood $V$ of
$y$ in $Y$ such that
$D'|f^{-1}(V)=f_V^*(D_V)$ for some
$D_V\in\operatorname{Div}(V)$, contradicting the maximality of $U$.

\textup{I) The case where $f$ is quasi-compact.}
Using the quasi-compactness of $f$ and {\hyperlink{ega4.c04.prop.20-3-8}{(20.3.8)}}, by the method of
({\hyperlink{ega4.c04.statement.8-1-2}{8.1.2}}, a), we are reduced to proving---in the notation of the
statement, but without initially imposing a bound on
$\operatorname{prof}(\sh{O}_{Y,y})$---that $D_1'$ is of the form
$f_1^*(D_1)$.  If $\operatorname{prof}(\sh{O}_{Y,y})\leq 1$, this is
precisely the hypothesis.  Suppose therefore that
$\operatorname{prof}(\sh{O}_{Y,y})\geq 2$.  Since $y$ is a maximal
point of $Y\setminus U$, the open subset $W=Y_1\cap U$ of $Y_1$ is
$Y_1\setminus\{y\}$.  We may therefore restrict to the case
$Y=Y_1$ and $U=Y\setminus\{y\}$.

By hypothesis there is a divisor $D_U\in\operatorname{Div}(U)$ such
that $f_U^*(D_U)=D'|f^{-1}(U)$.  The divisor $D_U$ is the equivalence
class of a pair $(\sh{L}_U,s)$, where $\sh{L}_U$ is an invertible
$\sh{O}_U$-Module and $s$ a regular meromorphic section of $\sh{L}_U$
over $U$ {\hyperlink{ega4.c04.prop.21-2-11}{(21.2.11)}}.  Hence, by {\hyperlink{ega4.c04.statement.21-4-1}{(21.4.1)}}, $D'|f^{-1}(U)$ is the
equivalence class of $(\sh{L}'_U,s')$, where
\[
 \sh{L}'_U=f_U^*(\sh{L}_U),\qquad s'=s\circ f_U.
\]
Let $i:U\to Y$ and $j:f^{-1}(U)\to X$ be the canonical injections,
and put
\[
 \sh{L}=i_*(\sh{L}_U),\qquad
 \sh{L}'=j_*(\sh{L}'_U).
\]
Since $f$ is flat, $f^*(\sh{L})$ is canonically isomorphic to
$\sh{L}'$ {\hyperlink{ega4.c04.lem.2-3-1}{(2.3.1)}}.  Moreover, flatness and the inequality
$\operatorname{prof}(\sh{O}_{Y,y})\geq 2$ imply
$\operatorname{prof}(\sh{O}_{X,x})\geq 2$ for every
$x\in f^{-1}(y)$ {\hyperlink{ega4.c04.prop.6-3-1}{(6.3.1)}}.

The divisor $D'$ is the equivalence class of a pair
$(\sh{L}'',s'')$, where $\sh{L}''$ is an invertible $\sh{O}_X$-Module
and $s''$ a regular meromorphic section of $\sh{L}''$ over $X$.
Consequently $\sh{L}'_U$ is isomorphic to
$\sh{L}''|f^{-1}(U)$, and the preceding observations together with
{\hyperlink{ega4.c04.prop.5-9-8}{(5.9.8)}} and {\hyperlink{ega4.c04.thm.5-10-5}{(5.10.5)}} show that $\sh{L}'$ is an invertible
$\sh{O}_X$-Module isomorphic to $\sh{L}''$.  Since $f$ is faithfully
flat and quasi-compact, it follows that $\sh{L}$ is an invertible
$\sh{O}_Y$-Module {\hyperlink{ega4.c04.prop.2-5-2}{(2.5.2)}}, and hence isomorphic to $\sh{O}_Y$, because
$Y$ is a local scheme.

Furthermore, since $\operatorname{prof}(\sh{O}_{Y,y})\geq 1$, the open
subset $U=Y\setminus\{y\}$ is schematically dense in $Y$
({\hyperlink{ega4.c04.statement.20-2-13}{20.2.13}}, (iv)).  The regular meromorphic section $s$ of $\sh{L}_U$
over $U$ therefore extends to a regular meromorphic section $s_0$ of
$\sh{L}$ over $Y$ {\hyperlink{ega4.c04.prop.20-2-11}{(20.2.11)}}.  The equivalence class of
$(\sh{L},s_0)$ is thus a divisor $D$ on $Y$.  Plainly $f^*(D)$ and
$D'$ induce the same divisor on
$f^{-1}(U)=X\setminus f^{-1}(y)$; and since
$\operatorname{prof}(\sh{O}_{X,x})\geq 2$ for every
$x\in f^{-1}(y)$, {\hyperlink{ega4.c04.prop.21-1-8}{(21.1.8)}} gives $D'=f^*(D)$.  This proves case I.

\textup{II) The case where $f$ is open.}
The problem is to find an open neighbourhood $V$ of $y$ in $Y$ such
that $D'|f^{-1}(V)=f_V^*(D_V)$ for some
$D_V\in\operatorname{Div}(V)$, so we may first suppose $Y$ affine.
Because $f$ is open, the images $f(W)$, with $W$ running through the
affine open subsets of $X$, form an open covering of $Y$.  Since $Y$
is quasi-compact, there is a quasi-compact open subset $T\subset X$
such that $f(T)=Y$, and $f|T$ is quasi-compact (I, 1.1.1).  A fortiori,
for every quasi-compact open $T_\alpha\supset T$ in $X$,
\oldpage[IV]{360}
one has
$f(T_\alpha)=Y$ and $f|T_\alpha$ is quasi-compact.

Applying case I to $f|T_\alpha$, we obtain a divisor $D_\alpha$ on
$Y$ such that
\[
 (f|T_\alpha)^*(D_\alpha)=D'|T_\alpha.
\]
By transitivity of inverse images of divisors under flat morphisms
{\hyperlink{ega4.c04.statement.21-4-4}{(21.4.4)}}, one also has
$(f|T)^*(D_\alpha)=D'|T$.  Hence for any $\alpha,\beta$,
\[
 (f|T)^*(D_\alpha)=(f|T)^*(D_\beta),
\]
and {\hyperlink{ega4.c04.prop.21-4-8}{(21.4.8)}} gives $D_\alpha=D_\beta$.  If $D$ denotes their common
value, then $f^*(D)|T_\alpha=D'|T_\alpha$ for every $\alpha$, and
therefore $f^*(D)=D'$.
\end{proof}

\begin{corollary}[21.4.10]
\ifdefined\EGAIVSuppressErratumLabels\else\label{IV.21.4.10}\fi
Assume the general hypotheses of {\hyperlink{ega4.c04.prop.21-4-9}{(21.4.9)}}, and suppose in addition that
$Y$ is normal and that, for every $y\in Y$ with
$\dim(\sh{O}_{Y,y})=1$, the fibre $X_y=f^{-1}(y)$ has no embedded
associated prime cycle.  A divisor $D'$ on $X$ is of the form $f^*(D)$,
with $D\in\operatorname{Div}(Y)$, if and only if the following two
conditions hold:
\begin{enumerate}[label=\textup{\arabic*\textdegree}]
\item for every maximal point $y$ of $Y$, the inverse image of $D'$ in
$\operatorname{Div}(X_y)$ is zero;
\item for every $y\in Y$ such that $\dim(\sh{O}_{Y,y})=1$, there is an
integer $n_y$ such that, for every maximal point $x$ of $X_y$, the
inverse image of $D'$ in
$\operatorname{Div}(\operatorname{Spec}(\sh{O}_{X,x}))$ equals
$\operatorname{div}(t_y^{n_y})$, where $t_y$ is the image under
$\sh{O}_{Y,y}\to\sh{O}_{X,x}$ of a uniformizer of the discrete
valuation ring $\sh{O}_{Y,y}$.
\end{enumerate}
\end{corollary}

\begin{proof}
Apply the criterion {\hyperlink{ega4.c04.prop.21-4-9}{(21.4.9)}}.  At a maximal point $y$ of $Y$,
$\operatorname{Div}(Y_1)=0$ because $\sh{O}_{Y,y}$ is Artinian;
condition 1\textdegree{} is therefore equivalent to the criterion of
{\hyperlink{ega4.c04.prop.21-4-9}{(21.4.9)}} at such a point.  It remains to show that conditions
1\textdegree{} and 2\textdegree{} likewise imply the criterion at a
point $y\in Y$ with $\operatorname{prof}(\sh{O}_{Y,y})=1$.  Since $Y$
is normal, these are precisely the points for which
$\dim(\sh{O}_{Y,y})=1$, equivalently those for which
$\sh{O}_{Y,y}$ is a discrete valuation ring {\hyperlink{ega4.c04.thm.5-8-6}{(5.8.6)}}.

We may thus restrict to $Y=\operatorname{Spec}(\sh{O}_{Y,y})$.  If $t$
is a uniformizer of $\sh{O}_{Y,y}$, the map
$n\mapsto\operatorname{div}(t^n)$ is an isomorphism from $\mathbb Z$
onto $\operatorname{Div}(Y)$.  The condition of {\hyperlink{ega4.c04.prop.21-4-9}{(21.4.9)}} says that
$D'$ is of the form $\operatorname{div}(t^n\circ f)$ for some
$n\in\mathbb Z$, and plainly implies condition 2\textdegree{}.
Conversely, suppose 2\textdegree{} holds, as does 1\textdegree{}, and
write $n=n_y$.  It remains to prove
$D'=\operatorname{div}(t^n\circ f)$.  If $\eta$ is the generic point of
$Y$, the germs of these two divisors at every point of $X_\eta$ are
zero by condition 1\textdegree{}.  Their germs at every maximal point
of $X_y$ are equal by condition 2\textdegree{}.  At a nonmaximal point
$x\in X_y$, one has
$\operatorname{prof}(\sh{O}_{X_y,x})\geq 1$, because $X_y$ has no
embedded associated prime cycle, and hence
$\operatorname{prof}(\sh{O}_{X,x})\geq 2$ by {\hyperlink{ega4.c04.prop.6-3-1}{(6.3.1)}}.  The desired
equality follows from {\hyperlink{ega4.c04.prop.21-1-8}{(21.1.8)}}.
\end{proof}

\begin{corollary}[21.4.11]
\ifdefined\EGAIVSuppressErratumLabels\else\label{IV.21.4.11}\fi
Assume the general hypotheses of {\hyperlink{ega4.c04.prop.21-4-9}{(21.4.9)}}, and suppose in addition that
$Y$ is normal and that, for every $y\in Y$ with
$\dim(\sh{O}_{Y,y})=1$, the fibre $X_y$ is integral.  A divisor $D'$ is
of the form $f^*(D)$, with $D\in\operatorname{Div}(Y)$, if and only if,
for every maximal point $y$ of $Y$, the inverse image of $D'$ in
$\operatorname{Div}(X_y)$ is zero.
\end{corollary}

\begin{proof}
It is enough to check that condition 2\textdegree{} of {\hyperlink{ega4.c04.cor.21-4-10}{(21.4.10)}} is then
automatic.  If $\dim(\sh{O}_{Y,y})=1$ and $X_y$ is integral, then at
the unique maximal point $x$ of $X_y$ one has
$\dim(\sh{O}_{X,x})=1$ by {\hyperlink{ega4.c04.cor.6-1-3}{(6.1.3)}}.  Moreover,
$\sh{O}_{Y,y}$ is regular and $\sh{O}_{X_y,x}$ is a field, so
$\sh{O}_{X,x}$ is regular {\hyperlink{ega4.c04.prop.6-5-1}{(6.5.1)}}, hence a discrete valuation ring.
If $t$ is a uniformizer of $\sh{O}_{Y,y}$, its image $t_y$ in
$\sh{O}_{X,x}$ is a uniformizer of $\sh{O}_{X,x}$
($0$, 17.3.3).  The inverse image of $D'$ in
$\operatorname{Div}(\operatorname{Spec}(\sh{O}_{X,x}))$ is therefore
\oldpage[IV]{361}
of the form $\operatorname{div}(t_y^n)$, which is exactly condition
2\textdegree{} of {\hyperlink{ega4.c04.cor.21-4-10}{(21.4.10)}}, since $X_y$ has only one maximal point.
\end{proof}

\begin{remark}[21.4.12]
\ifdefined\EGAIVSuppressErratumLabels\else\label{IV.21.4.12}\fi
If in {\hyperlink{ega4.c04.cor.21-4-11}{(21.4.11)}} one assumes only that the $X_y$ are reduced, one must
add to the condition in the statement that the inverse images of $D'$
in each group
$\operatorname{Div}(\operatorname{Spec}(\sh{O}_{X,x}))$, with $x$
running over the maximal points of $X_y$, are all of the form
$\operatorname{div}((t_x')^n)$, where $t_x'$ is a uniformizer of
$\sh{O}_{X,x}$ and the integer $n$ is the same for all these points.
\end{remark}

\begin{corollary}[21.4.13]
\ifdefined\EGAIVSuppressErratumLabels\else\label{IV.21.4.13}\fi
Assume the hypotheses of {\hyperlink{ega4.c04.cor.21-4-11}{(21.4.11)}}.  For every invertible
$\sh{O}_X$-Module $\sh{L}$ such that, for every maximal point $\eta$ of
$Y$, the inverse image
\[
 \sh{L}_\eta
 =\sh{L}\otimes_{\sh{O}_X}\sh{O}_{X_\eta}
\]
is isomorphic to $\sh{O}_{X_\eta}$, there exist an invertible
$\sh{O}_Y$-Module $\sh{M}$ and an isomorphism
$f^*(\sh{M})\simeq\sh{L}$.
\end{corollary}

\begin{proof}
It is enough to find a regular meromorphic section $s$ of $\sh{L}$ over
$X$ such that the divisor $D'$ corresponding to $(\sh{L},s)$ has
$D'_x=0$ at every point $x$ of a fibre $X_\eta$ over a maximal point
$\eta$ of $Y$.  Applying {\hyperlink{ega4.c04.cor.21-4-11}{(21.4.11)}} to $D'$ then gives the result by
{\hyperlink{ega4.c04.def.21-4-2}{(21.4.2)}}.

The points of $\operatorname{Ass}(\sh{O}_X)$ lie over the maximal
points of $Y$ {\hyperlink{ega4.c04.prop.3-3-1}{(3.3.1)}}.  Since $Y$ is reduced and locally Noetherian,
there is a dense open subset $V$ of $Y$ which is the union of a
disjoint family of affine open subsets, each containing exactly one
maximal point of $Y$; and $f^{-1}(V)$ is schematically dense in $X$
{\hyperlink{ega4.c04.prop.5-10-2}{(5.10.2)}}.  We may therefore restrict to the case in which $Y$ is
affine and integral.

For the unique maximal point $\eta$ of $Y$, the fibre $X_\eta$ is then
integral, so $\operatorname{Ass}(\sh{O}_X)$ consists of a single point
$z$, the generic point of $X$.  Let $s_\eta$ be the section of
$\sh{L}_\eta$ over $X_\eta$ which corresponds, under an isomorphism,
to the unit section of $\sh{O}_{X_\eta}$.  Let $U$ be an affine open
neighbourhood in $X$ of a point $x\in X_\eta$, and let $s$ be a section
of $\sh{L}$ over $U$ with the same germ at $x$ as $s_\eta$.  Since $U$
is schematically dense in $X$ ({\hyperlink{ega4.c04.statement.20-2-13}{20.2.13}}, (iv)), $s$ identifies with a
uniquely determined regular meromorphic section of $\sh{L}$ over $X$.
It is clear from {\hyperlink{ega4.c04.prop.20-2-11}{(20.2.11)}} and {\hyperlink{ega4.c04.prop.20-3-5}{(20.3.5)}} that this section induces on
$X_\eta$ a regular meromorphic section of $\sh{L}_\eta$ equal to
$s_\eta$, and hence has the required property.
\end{proof}
% END NATIVE INPUT source_aligned/ega4-21-erratum53-insertion
\let\EGAIVSuppressErratumLabels\undefined
}

\begin{remark}[3.1.10.5]
\label{IV.3.1.10.5}
Even when $f$ is proper, one need not have
$f(\operatorname{Ass}(\sh{F}))\subset
\operatorname{Ass}(f_*(\sh{F}))$.  Take $Y$ to be the spectrum of a
field and $X$ a projective space over $k$, for example $\mathbf P^1_k$;
then $\sh{F}=\sh{O}_X(-1)$ is a nonzero coherent $\sh{O}_X$-Module with
$\Gamma(X,\sh{F})=0$ (III, 2.1.13), so
$\operatorname{Ass}(f_*(\sh{F}))=\varnothing$ whereas
$\operatorname{Ass}(\sh{F})\neq\varnothing$.
\end{remark}

\oldpage[IV]{349}
\paragraph{(Err\textsubscript{IV}, 18).}
After (IV, 3.2.8), insert after line 17 of p.~43:

\begin{remark}[3.2.9]
\label{IV.3.2.9}
For every point $s$ of a prescheme $X$, write
$i_s:\operatorname{Spec}(k(s))\to X$ for the canonical morphism, which
is a quasi-compact monomorphism.  For every $k(s)$-module $G(s)$,
regarded as a Module on $\operatorname{Spec}(k(s))$,
$(i_s)_*(G(s))$ is therefore a quasi-coherent $\sh{O}_X$-Module
{\hyperlink{ega4.c04.statement.1-7-4}{(1.7.4)}}.  Its fibre is zero at every point
$x\notin\overline{\{s\}}$ and is isomorphic to $G(s)$ otherwise.  If
$X$ is locally Noetherian and $G(s)$ has finite length, then
$\operatorname{Ass}((i_s)_*(G(s)))$ consists of the single point $s$;
in other words, $(i_s)_*(G(s))$ is irredundant.

Now suppose that $X$ is locally Noetherian and that $\sh{F}$ is a
coherent $\sh{O}_X$-Module.  Let $S$ be a discrete subset of $X$, and
for every $s\in S$ put $\sh{F}(s)=i_s^*(\sh{F})$, a $k(s)$-module.
For each $s\in S$, choose a finite-length quotient $G(s)$ of
$\sh{F}(s)$.  There is a canonical homomorphism
\[
 u:\sh{F}\longrightarrow
 \bigoplus_{s\in S}(i_s)_*(\sh{F}(s))\longrightarrow
 \bigoplus_{s\in S}(i_s)_*(G(s))
\]
of coherent $\sh{O}_X$-Modules, since a discrete subset of $X$ is
locally finite.  If $\sh{H}=\operatorname{Ker}u$ and
$\sh{G}=\sh{F}/\sh{H}$, then $\sh{G}$ is coherent and $u$ induces an
injection
\[
 v:\sh{G}\longrightarrow\bigoplus_{s\in S}(i_s)_*(G(s)).
\]
If $\sh{G}^{(s)}$ denotes the image of $\sh{G}$ in
$(i_s)_*(G(s))$, the family $(\sh{G}^{(s)})_{s\in S}$ is an
irredundant decomposition, reduced when all the $G(s)$ are nonzero, and
$\operatorname{Ass}(\sh{G})=S$; thus $\sh{G}$ has no embedded
associated prime cycle.

This gives a bijection $(G(s))_{s\in S}\leftrightarrow\sh{G}$, where
$S$ ranges over the discrete subsets of $X$ for which
$\sh{F}(s)\neq0$ for every $s\in S$, and, for each $S$, the $G(s)$
range over the nonzero finite-length quotients of $\sh{F}(s)$; on the
other side, $\sh{G}$ ranges over the coherent quotient
$\sh{O}_X$-Modules of $\sh{F}$ without embedded associated prime
cycles.  The inverse bijection associates to $\sh{G}$ its reduced
irredundant decomposition.
\end{remark}

\paragraph{(Err\textsubscript{IV}, 19).}
In (IV, 4.2.1), before the seventeenth line from the bottom of p.~55,
insert:

\begin{remark}[4.2.1.4]
\label{IV.4.2.1.4}
Under the hypotheses of {\hyperlink{ega4.c04.prop.4-2-1}{(4.2.1)}}, one also has
\begin{equation}
 \dim(K\otimes_kL)=inf(\operatorname{tr.deg}_kK,
                           \operatorname{tr.deg}_kL). \label{IV.4.2.1.5}
\tag{4.2.1.5}
\end{equation}
If both fields have infinite transcendence degree (which may be an
arbitrary cardinal), the left-hand side is to be understood as
$+\infty$; the dimension here is that of ($0$, 16.1.1).

Suppose, for example, that
$\operatorname{tr.deg}_kK\leq\operatorname{tr.deg}_kL$, and first that
$\operatorname{tr.deg}_kK=n<+\infty$.  In the proof of {\hyperlink{ega4.c04.prop.4-2-1}{(4.2.1)}} we saw
that $K\otimes_kL$ is an algebraic extension of $K'\otimes_kL$, and
therefore has the same dimension ($0$, 16.1.5).  We may reduce to
$K=k(T_1,\ldots,T_n)$; then $K\otimes_kL$ is a ring of fractions of
\[
 k[T_1,\ldots,T_n]\otimes_kL=L[T_1,\ldots,T_n],
\]
which has dimension $n$ {\hyperlink{ega4.c04.cor.5-5-4}{(5.5.4)}}, so
$\dim(K\otimes_kL)\leq n$ ($0$, 16.1.3.1).  To prove equality, it is
enough to find in $\operatorname{Spec}(L[T_1,\ldots,T_n])$ an
$L$-rational point belonging to $\operatorname{Spec}(K\otimes_kL)$
{\hyperlink{ega4.c04.cor.5-2-3}{(5.2.3)}}, or equivalently (I, 1.7.4) an $L$-homomorphism
$u:K\otimes_kL\to L$.  By the definition of the tensor product this is
equivalent to a $k$-homomorphism $v:K\to L$, whose existence follows
from the inequality of transcendence degrees.

If the two transcendence degrees are infinite, the same argument gives
a rational point of $\operatorname{Spec}(L[(t_\alpha)_{\alpha\in I}])$
belonging to $\operatorname{Spec}(K\otimes_kL)$.  The kernel of an
$L$-homomorphism from this polynomial ring to $L$ is a maximal ideal
generated by a family $(t_\alpha-a_\alpha)_{\alpha\in I}$ with
$a_\alpha\in L$.  Such a maximal ideal contains an infinite descending
chain of distinct prime ideals; the corresponding local ring is
non-Noetherian and infinite-dimensional.  Since it is a local ring of
$K\otimes_kL$, the latter is also non-Noetherian and
infinite-dimensional.
\end{remark}

\oldpage[IV]{350}
\paragraph{(Err\textsubscript{IV}, 20).}
After (IV, 4.5.19), before the twelfth line from the bottom of p.~67,
insert:

\begin{corollary}[4.5.19.3]
\label{IV.4.5.19.3}
Let $k$ be a field and $X$ a finite-type $k$-prescheme which is
irreducible but not geometrically irreducible.  There is a closed subset
$Y$ of $X$, nowhere dense in $X$, containing every $x\in X$ for which
$k(x)$ is radicial over $k$ (and in particular every $k$-rational point
of $X$).
\end{corollary}

Let $k'$ be an algebraic closure of $k$, let $X'=X\otimes_kk'$, and let
$p:X'\to X$ be the canonical projection.  If $k(x)$ is radicial over
$k$, there is only one point $x'\in X'$ over $x$ (I, 3.4.9); this point
must belong to two distinct irreducible components of $X'$, since
otherwise $X$ would be geometrically irreducible by {\hyperlink{ega4.c04.prop.4-5-19}{(4.5.19)}}.  Let
$X'_i$ ($1\leq i\leq n$) be the irreducible components of the Noetherian
prescheme $X'$, and for $1\leq i<j\leq n$ put
$X'_{ij}=X'_i\cap X'_j$.  Then
\[
 Y=p\left(\bigcup_{i<j}X'_{ij}\right)
\]
has the required properties.  Since $p$ is integral and the $X'_{ij}$
are closed in $X'$, $Y$ is closed in $X$.  If $\xi$ is the generic point
of $X$, $p^{-1}(\xi)$ is discrete and consists of the generic points of
the $X'_i$ {\hyperlink{ega4.c04.cor.4-5-11}{(4.5.11)}}, so it meets none of the $X'_{ij}$; hence $Y$ is
nowhere dense.

\paragraph{(Err\textsubscript{IV}, 21).}
In (IV, 5.6.11.1), on the seventh line from the bottom of p.~102, add
``of dimension $2$'' after ``integral.''

\paragraph{(Err\textsubscript{IV}, 22).}
After (IV, 6.2.3), after line 18 of p.~138, insert:

\begin{proposition}[6.2.4]
\label{IV.6.2.4}
Under the hypotheses of {\hyperlink{ega4.c04.prop.6-2-3}{(6.2.3)}}, let in addition $M$ be a finite
$A$-module, and suppose that $M\neq0$ and $N\neq0$.  Then
\begin{equation}
 \operatorname{pd}_B(M\otimes_AN)=
 \operatorname{pd}_A(M)+
 \operatorname{pd}_{B\otimes_Ak}(N\otimes_Ak). \label{IV.6.2.4.1}
\tag{6.2.4.1}
\end{equation}
\end{proposition}

Consider the left resolution $L_\bullet$ of $N$ introduced in {\hyperlink{ega4.c04.prop.6-2-3}{(6.2.3)}},
and a left resolution $K_\bullet$ of $M$ by finite free $A$-modules:
\[
 \cdots\longrightarrow K_i\longrightarrow K_{i-1}\longrightarrow
 \cdots\longrightarrow K_0\longrightarrow M\longrightarrow0.
\]
We saw that the $L_i$, $N$, and the $Z_i(L_\bullet)$ are flat
$A$-modules.  The same is true of the $B_i(L_\bullet)$, which equal
$Z_i(L_\bullet)$ except for $i=0$; in that case $B_0(L_\bullet)$ is flat
\oldpage[IV]{351}
by ($0_{\mathrm I}$, 6.1.2).  Finally the $H_i(L_\bullet)$ vanish except
for $H_0(L_\bullet)=N$, which is flat over $A$.  Under these conditions
one knows that
\[
 H_n(K_\bullet\otimes_AL_\bullet)\simeq
 \bigoplus_{p+q=n}H_p(K_\bullet)\otimes_AH_q(L_\bullet)
\]
(G, I, 5.5).  Thus $K_\bullet\otimes_AL_\bullet$, with the usual total
grading, is a resolution of $M\otimes_AN$ by free $B$-modules.

Put
\[
 r=\operatorname{pd}_A(M),\qquad
 s=\operatorname{pd}_B(N)=
   \operatorname{pd}_{B\otimes_Ak}(N\otimes_Ak).
\]
There are free resolutions $K_\bullet$ and $L_\bullet$ of lengths $r$
and $s$, so their tensor product is a free resolution of length $r+s$;
this proves
$\operatorname{pd}_B(M\otimes_AN)\leq r+s$.  It remains, by
($0$, 17.2.6), to show
\[
 \operatorname{Tor}^{B}_{r+s}(M\otimes_AN,k(B))\neq0.
\]
Using the preceding resolution, the left-hand side is
\[
 H_{r+s}((K_\bullet\otimes_AL_\bullet)\otimes_Bk(B)).
\]
But
\[
 (K_\bullet\otimes_AL_\bullet)\otimes_Bk(B)
 = (K_\bullet\otimes_Ak)\otimes_k
   (L_\bullet\otimes_Bk(B)).
\]
This is a tensor product of complexes of vector spaces over $k$, so the
K\"unneth formula gives
\[
 H_{r+s}\simeq H_r(K_\bullet\otimes_Ak)\otimes_k
 H_s(L_\bullet\otimes_Bk(B)),
\]
since the homology of the first, respectively second, complex vanishes
above $r$, respectively $s$.  Hence, up to canonical isomorphism,
\[
 \operatorname{Tor}^{B}_{r+s}(M\otimes_AN,k(B))=
 \operatorname{Tor}^{A}_{r}(M,k)\otimes_k
 \operatorname{Tor}^{B}_{s}(N,k(B)),
\]
and both factors are nonzero by hypothesis; so is their tensor product.

\oldpage[IV]{352}
\paragraph{(Err\textsubscript{IV}, 23).}
In (IV, 6.10.1), on line 8 of p.~155, add after ``canonical'':
``$\sh{F}$ a coherent $\sh{O}_X$-Module.''

\paragraph{(Err\textsubscript{IV}, 24).}
In ({\hyperlink{ega4.c04.prop.6-14-4}{IV, 6.14.4}}), the reference on the twelfth line from the bottom of
p.~175 is insufficient; argue as follows.  First, if $M$ is a
transcendence basis of $K'$ over $A$, then $A(M)$ is algebraically closed
in $B(A(M))$ (Bourbaki, \emph{Alg.}, chap.~V, \S6, exercise 8(c)).  We
may therefore reduce to the case in which $K'$ is an algebraic separable
extension of $A$.

Let $p$ be the characteristic exponent of $A$.  It follows from
Bourbaki, \emph{loc. cit.}, \S9, exercise 2, that the algebraic closure
of $K'$ in $B(K')$ is radicial over $K'$.  We show that every
$x\in B(K')$ such that $x^{p^r}\in K'$ already belongs to $K'$.
Choose a basis $(e_\lambda)$ of $K'$ over $A$.  Since $K'/A$ is
algebraic separable, the $e_\lambda^{p^r}$ also form a basis of $K'$ over
$A$ (Bourbaki, \emph{loc. cit.}, \S8, no.~3, corollary to prop.~4).
Write $x=\sum_\lambda e_\lambda b_\lambda$ with $b_\lambda\in B$.  By
hypothesis there are $a_\lambda\in A$ such that
\[
 \sum_\lambda e_\lambda^{p^r}b_\lambda^{p^r}
 =\sum_\lambda e_\lambda^{p^r}a_\lambda.
\]
Therefore $b_\lambda^{p^r}=a_\lambda$ for every $\lambda$.  Since $A$
is algebraically closed in $B$, each $b_\lambda$ lies in $A$, and hence
$x\in K'$.

\paragraph{(Err\textsubscript{IV}, 25).}
In (IV, 7.4.8, B), p.~203, the question was answered affirmatively by
R.~Kiehl, ``Ausgezeichnete Ringe der nichtarchimedischen analytischen
Geometrie,'' forthcoming in \emph{Inventiones Mathematicae}.

\paragraph{(Err\textsubscript{IV}, 26).}
In (IV, 7.8.4), on the ninth line from the bottom of p.~217, replace
``We shall see later (\S18) that if $k$ is a complete nondiscrete valued
field\ldots'' by:
\begin{quote}
R.~Kiehl proved (in the article cited in
(Err\textsubscript{IV}, 25)) that if $k$ is a complete nondiscrete
valued field\ldots
\end{quote}

\paragraph{(Err\textsubscript{IV}, 27).}
In (IV, 7.8.4), insert before the sixth line from the bottom of p.~217:

\smallskip
\noindent
\textup{(vi)} Let $A$ be a one-dimensional Noetherian local integral
ring.  It is then universally catenary {\hyperlink{ega4.c04.cor.7-2-9}{(7.2.9)}}, and hence, by {\hyperlink{ega4.c04.statement.7-3-19}{(7.3.19)}},
it is excellent if and only if its formal fibres are geometrically
reduced; equivalently, if and only if its completion $\widehat A$ is
reduced and the component fields of the total ring of fractions of
$\widehat A$ are separable extensions of the fraction field $K$ of
$A$.  By {\hyperlink{ega4.c04.thm.7-6-4}{(7.6.4)}} and {\hyperlink{ega4.c04.thm.7-7-2}{(7.7.2)}}, this is also equivalent to saying that
$A$ is universally Japanese.  When $A$ is a discrete valuation ring,
so is $\widehat A$; to say that $A$ is excellent then means that the
fraction field of $\widehat A$ is a separable extension of $K$, a
condition which is always satisfied when $K$ has characteristic zero.

\paragraph{(Err\textsubscript{IV}, 28).}
After (IV, 8.10.3), following line 14 of p.~37, insert:

\begin{remark}[8.10.3.1]
\label{IV.8.10.3.1}
The hypothesis in {\hyperlink{ega4.c04.prop.8-10-3}{(8.10.3)}} that the sets
$\overline{f_\alpha(X_\alpha)}$ be constructible cannot be omitted.
For example, take $S=\operatorname{Spec}(A)$, where $A$ is a ring such
that $S$ is compact, totally disconnected, and has no isolated point,
every prime ideal $\mathfrak p$ of $A$ is maximal, and $A_\mathfrak p$
is a field (Bourbaki, \emph{Alg. comm.}, chap.~II, \S4,
exercises 16 and 17).  Let $s\in S$, and take for
$Y_\alpha=S_\alpha$ the affine open neighbourhoods of $s$ in $S$.
For every $\alpha$, take $X_\alpha=\operatorname{Spec}(k(s))$, with
$f_\alpha$ and $f$ the canonical morphisms.  It is clear that $f$,
which is the identity, is dominant, whereas none of the $f_\alpha$ is.
\end{remark}

\paragraph{(Err\textsubscript{IV}, 29).}
In (IV, 8.12.8), on line 7 of p.~47, replace ``such that there exists an
ample $\sh{O}_Y$-Module'' by ``and quasi-separated.''  On line 12 of
p.~47, replace ``The hypothesis'' by:
\begin{quote}
It follows from {\hyperlink{ega4.c04.prop.8-12-3}{(8.12.3)}} that the existence of the factorisation in the
statement is a property local on $Y$; we may therefore suppose $Y$
affine, which
\end{quote}

\paragraph{(Err\textsubscript{IV}, 30).}
After (IV, 8.12.11), add at the end of p.~48:

\begin{corollary}[8.12.12]
\label{IV.8.12.12}
Let $X$ and $Y$ be two integral preschemes and let $f:X\to Y$ be a
separated, birational morphism locally of finite type.  Let $y\in Y$ be
a normal point of $Y$, and suppose that a point $x\in f^{-1}(y)$ is
isolated in $f^{-1}(y)$.  Then there exists an open neighbourhood $U$
of $y$ in $Y$ such that the restriction $f^{-1}(U)\to U$ of $f$ is an
isomorphism.
\end{corollary}

\begin{proof}
It is enough to prove that there exist an open neighbourhood $U$ of
$y$ in $Y$ and a $U$-section $g:U\to f^{-1}(U)$ of $f^{-1}(U)$ such
that $g(y)=x$.  Indeed, $U$ contains the generic point $\eta$ of $Y$
and $f^{-1}(U)$ contains the generic point $\xi$ of $X$; necessarily
$g(\eta)=\xi$, since $f$ is birational {\hyperlink{ega4.c04.statement.6-15-4}{(6.15.4)}}.  As $g$ is an
immersion, $g(U)$ is locally closed in $f^{-1}(U)$, hence open in its
closure there; the preceding observation shows that this closure is
$f^{-1}(U)$.  The subprescheme of $f^{-1}(U)$ associated with $g$ is
isomorphic to $U$, hence reduced, and $f^{-1}(U)$ is reduced, so (I,
5.2.1) shows that $g$ is an open immersion.  On the other hand, since
$f$ is separated, $g$ is a closed immersion (I, 5.4.6).  Thus $g(U)$
is both open and closed in $f^{-1}(U)$.  Since $f^{-1}(U)$ is integral,
this gives $g(U)=f^{-1}(U)$, as required.

We consider two cases.

\textup{I)} First suppose $Y$ normal, and let $V$ be the set of points
of $X$ at which $f$ is quasi-finite; this is open in $X$ {\hyperlink{ega4.c04.cor.13-1-4}{(13.1.4)}}.\footnote{The
reader will verify that neither {\hyperlink{ega4.c04.cor.8-12-12}{(8.12.12)}} nor any of its consequences is
used in the proof of {\hyperlink{ega4.c04.cor.13-1-4}{(13.1.4)}}.}  By {\hyperlink{ega4.c04.cor.8-12-10}{(8.12.10)}},
\oldpage[IV]{353}
the restriction
$f|V:V\to Y$ is an open immersion.  Hence there is an open
neighbourhood $U$ of $y$ in $Y$ such that $f|V$ is an isomorphism from
$V$ onto the subprescheme induced by $Y$ on $U$; its inverse $g$ is a
$U$-section of $f^{-1}(U)$ satisfying $g(y)=x$.

\textup{II)} The construction of a $U$-section $g$ with $g(y)=x$ is
local on $X$ and $Y$.  We may therefore suppose that $X$ and $Y$ are
affine and that $f$ is of finite type, so that $X$ is a closed
subprescheme of a $Y$-prescheme $Z$ of finite presentation over $Y$
(for example one of the form $\mathbb A_Y^r$).  Put
\[
 Y'=\operatorname{Spec}(\sh{O}_{Y,y}),\qquad
 X'=X_{(Y')},\qquad Z'=Z_{(Y')};
\]
thus $X'$ is a closed subprescheme of $Z'$.  If $f'=f_{(Y')}$, then
$f'^{-1}(y)=f^{-1}(y)$ (I, 3.6.5), so $x$ is isolated in
$f'^{-1}(y)$.  Since $Y'$ is normal, case I applies to $f'$; and since
$Y'$ is the only neighbourhood of $y$ in $Y'$, there is a $Y'$-section
$g'$ of $X'$ with $g'(y)=x$, which may also be regarded as a
$Y'$-section of $Z'$.

The structural morphism $h:Z\to Y$ is of finite presentation.  Applying
({\hyperlink{ega4.c04.thm.8-8-2}{8.8.2}}, (i) and (ii)) by the method of ({\hyperlink{ega4.c04.statement.8-1-2}{8.1.2}}, a), we obtain an open
neighbourhood $U$ of $y$ in $Y$ and a $U$-section $g$ of $h^{-1}(U)$
such that $g(y)=x$.  It remains to see that $g$ is in fact a $U$-section
of $f^{-1}(U)$.  Since $f^{-1}(U)$ is a closed subprescheme of
$h^{-1}(U)$ and $U$ is reduced, it is enough by (I, 5.2.2) to show that
$g(U)\subset f^{-1}(U)$.  As $U$ is integral and $f^{-1}(U)$ is closed
in $h^{-1}(U)$, it is enough to check this at the generic point $\eta$
of $Y$ (and of $U$).  But the restriction of $g$ to $Y'$ is a
$Y'$-section of $X'\subset f^{-1}(U)$, and $\eta\in Y'$.
\end{proof}

\paragraph{(Err\textsubscript{IV}, 31).}
In ({\hyperlink{ega4.c04.subsec.10-4}{IV, 10.4}}), after Proposition {\hyperlink{ega4.c04.statement.10-4-11}{(10.4.11)}}, add:

\begin{corollary}[10.4.12]
\label{IV.10.4.12}
Let $S$ be a prescheme, let $X$ be an $S$-prescheme of finite
presentation, and let $f$ be a radicial $S$-endomorphism of $X$.
\begin{enumerate}[label=\textup{(\roman*)}]
\item The endomorphism $f$ is bijective and finite (and hence is a
universal homeomorphism {\hyperlink{ega4.c04.prop.2-4-5}{(2.4.5)}}).
\item Suppose in addition that $X$ is reduced, that the local rings of
$S$ are Noetherian and universally Japanese, and that $f$ is
birational (a condition automatically satisfied if $k(x)$ is perfect
for every maximal point $x$ of $X$).  Then $f$ is an
$S$-automorphism.
\end{enumerate}
\end{corollary}

\begin{proof}
\textup{(i)} The question is local on $S$, so we may restrict to the
case $S=\operatorname{Spec}(A)$ affine.  There is a finitely generated
$\mathbb Z$-subalgebra $A_0$ of $A$ such that, on putting
$S_0=\operatorname{Spec}(A_0)$, one has
$X=X_0\times_{S_0}S$, where $X_0$ is a prescheme of finite type over
$S_0$, and such that $f$ is obtained from an $S_0$-endomorphism $f_0$
of $X_0$ by the base change $S\to S_0$ {\hyperlink{ega4.c04.prop.8-9-1}{(8.9.1)}}.  We may moreover
suppose $f_0$ radicial ({\hyperlink{ega4.c04.thm.8-10-5}{8.10.5}}, (vii)).  It is enough to prove that
$f_0$ is bijective and finite, so from now on we may suppose that $A$
is a finitely generated $\mathbb Z$-algebra.  It follows that $S$ and
$X$ are excellent preschemes ({\hyperlink{ega4.c04.statement.7-8-3}{(7.8.3)}} and {\hyperlink{ega4.c04.prop.7-8-6}{(7.8.6)}}), and in particular
that $X$ is universally Japanese ({\hyperlink{ega4.c04.cor.5-11-4}{(5.11.4)}} and ({\hyperlink{ega4.c04.statement.7-8-3}{7.8.3}}, (vi))).

We may also suppose $X$ reduced.  Indeed,
$f_{\mathrm{red}}:X_{\mathrm{red}}\to X_{\mathrm{red}}$ is a radicial
$S_{\mathrm{red}}$-endomorphism (I, 5.1.6).  If
$f_{\mathrm{red}}$ is finite, then so is $f$: because $f$ is radicial,
it is separated {\hyperlink{ega4.c04.lem.1-8-7-1}{(1.8.7.1)}}; since $f_{\mathrm{red}}$ is proper, $f$ is
proper as well (II, 5.4.6), and since $f$ is plainly quasi-finite, it is
finite (III, 4.4.2).

\oldpage[IV]{354}
Let $X'$ be the normalization of $X$ (II, 6.3.8).  If
$X_i$ ($1\leq i\leq n$) are the reduced subpreschemes of $X$ whose
underlying spaces are its irreducible components, then $X'$ is the
disjoint sum of the normalizations $X_i'$ of the integral preschemes
$X_i$.  Write $x_i$ (respectively $x_i'$) for the generic point of
$X_i$ (respectively $X_i'$).  The canonical morphism $p:X'\to X$ is
finite because $X$ is universally Japanese, and
$p^{-1}(x_i)=\{x_i'\}$, with $k(x_i)\to k(x_i')$ bijective.

Write $f=(\psi,\theta)$.  We already know from {\hyperlink{ega4.c04.statement.10-4-11}{(10.4.11)}} that $\psi$ is
a continuous bijection of the space $X$ onto itself.  Put
$y_i=\psi(x_i)$.  Then $y_i$ is again a maximal point of $X$: otherwise
a generalization $y_i'\neq y_i$ of $y_i$ would make $x_i$ a
generalization, different from itself, of $\psi^{-1}(y_i')$.  The set
$M$ of maximal points of $X$ is finite, so $\psi|M$ is a permutation;
write $\psi(x_i)=x_{\sigma(i)}$ for a permutation $\sigma$ of
$\{1,\ldots,n\}$.  For each $i$, the restriction of $f$ to $X_i$
therefore factors as
\[
 X_i\xrightarrow{f_i}X_{\sigma(i)}\longrightarrow X,
\]
the second arrow being the canonical injection (I, 5.2.2).  Since the
$X_i$ are integral, there is by (II, 6.3.9) a unique morphism
$f_i':X_i'\to X_{\sigma(i)}'$ making the diagram
\[
\begin{array}{ccc}
X_i'&\xrightarrow{f_i'}&X_{\sigma(i)}'\\
\downarrow p_i&&\downarrow p_{\sigma(i)}\\
X_i&\xrightarrow{f_i}&X_{\sigma(i)}
\end{array}
\]
commute.  Taking $f'$ to agree with $f_i'$ on every $X_i'$ gives a
commutative diagram
\phantomsection\label{IV.10.4.12.1}
\begin{equation}
\label{IV.10.4.12.1.errata}
\begin{array}{ccc}
X'&\xrightarrow{f'}&X'\\
\downarrow p&&\downarrow p\\
X&\xrightarrow{f}&X.
\end{array}
\tag{10.4.12.1}
\end{equation}

Since $f$ is separated, so is $f\circ p=p\circ f'$, and hence $f'$ is
separated (I, 5.5.1).  Since $p$ is finite and $f$ bijective, the
diagram {\hyperlink{ega4.c04.sec.10-4-12-1}{(10.4.12.1)}} also shows that $f'$ is quasi-finite.  The
$X_i'$ being integral and normal, {\hyperlink{ega4.c04.cor.8-12-11}{(8.12.11)}} gives a factorisation
\[
 X_i'\xrightarrow{f_i''}X_{\sigma(i)}''
 \xrightarrow{u_{\sigma(i)}}X_{\sigma(i)}',
\]
where $X_{\sigma(i)}''$ is the integral closure of
$X_{\sigma(i)}'$ relative to the extension $k(x_i')$ of
$k(x_{\sigma(i)}')$, and $f_i''$ is an open immersion.  As
$k(x_i')/k(x_{\sigma(i)}')$ is radicial, $u_{\sigma(i)}$ is radicial
(Bourbaki, \emph{Alg. comm.}, chap.~V, \S2, no.~3, lemma 4).
Consequently $f'$ is radicial, and {\hyperlink{ega4.c04.statement.10-4-11}{(10.4.11)}} shows that it is
surjective.  Since the $u_{\sigma(i)}$ are injective, each $f_i''$ is
bijective and hence an isomorphism.  It follows that $f'$ is finite,
because the $u_{\sigma(i)}$ are finite, $X'$ being universally
Japanese {\hyperlink{ega4.c04.statement.7-8-3}{(7.8.3)}}.  Thus $p\circ f'=f\circ p$ is finite; as $p$ is
surjective and $f$ separated, $f$ is proper (II, 5.4.3), and being
quasi-finite it is finite {\hyperlink{ega4.c04.thm.8-11-1}{(8.11.1)}}.

\oldpage[IV]{355}
\textup{(ii)} By ({\hyperlink{ega4.c04.thm.8-10-5}{8.10.5}}, (i)), applied by the method of ({\hyperlink{ega4.c04.statement.8-1-2}{8.1.2}}, a), we
may restrict to the case in which $S$ is a local scheme.  The preliminary
reductions made in the proof of (i) are therefore unnecessary here.  In
the same notation, one obtains the commutative diagram {\hyperlink{ega4.c04.sec.10-4-12-1}{(10.4.12.1)}}, in
which $X$ is an $S$-prescheme of finite type, $p$ is finite, $f'$ is an
isomorphism, and the canonical homomorphism
$\sh{O}_X\to p_*(\sh{O}_{X'})$ is injective.  In Chapter VI, in the
theory of quotients of preschemes by equivalence relations, we shall
prove that under these conditions $f$ itself is an isomorphism (the
reader will verify that {\hyperlink{ega4.c04.cor.10-4-12}{(10.4.12)}} is not used in that theory).
\end{proof}

\begin{corollary}[10.4.13]
\label{IV.10.4.13}
Let $k$ be a field and $X$ an algebraic prescheme over $k$.  Every
radicial $k$-endomorphism $f$ of $X$ is finite and bijective, hence a
universal homeomorphism.  If, in addition, $X$ is reduced and $f$ is
birational (for example if $k$ has characteristic zero), then $f$ is a
$k$-automorphism.
\end{corollary}

\begin{remark}[10.4.14]
\label{IV.10.4.14}
\begin{enumerate}[label=\textup{(\roman*)}]
\item In {\hyperlink{ega4.c04.cor.10-4-12}{(10.4.12)}} and {\hyperlink{ega4.c04.cor.10-4-13}{(10.4.13)}}, the hypothesis that $X$ is reduced
cannot be omitted from the second conclusion.  This is shown by the
example of the $B$-endomorphism corresponding to the augmentation
$D_B(L)\to B$ of the ring $D_B(L)=A$ defined in (0, 18.2.3), where $B$
is identified with a subring of $A$.

\item Let $X$ be a Noetherian prescheme.  We do not know whether every
radicial endomorphism $f$ of $X$ is necessarily surjective and finite.
If $f$ is assumed surjective and $X$ is Japanese, the argument of
{\hyperlink{ega4.c04.cor.10-4-12}{(10.4.12)}} shows that $f$ is finite.  If moreover $X$ is reduced and
$f$ is birational---for example, if $k(x)$ is perfect for every maximal
point $x\in X$---the same argument shows that $f$ is an automorphism.

\item There are examples of nonsurjective radicial endomorphisms of
non-Noetherian integral affine schemes.  Let $k$ be a ring, let
$A=k[(T_\lambda)_{\lambda\in L}]$ be a polynomial ring with $L$
countably infinite, and let $\mathfrak J$ be the ideal generated by a
subfamily $(T_\lambda)_{\lambda\in H}$, where $L\setminus H$ is infinite
and $H\neq\varnothing$.  Then $A/\mathfrak J$ is isomorphic to
$k[(T_\lambda)_{\lambda\in L\setminus H}]$, hence to $A$.  The
endomorphism of $\operatorname{Spec}(A)$ corresponding to the composite
\[
 A\xrightarrow{p}A/\mathfrak J\xrightarrow{u}A,
\]
where $p$ is the canonical projection and $u$ an isomorphism, is such
an example: its image is the closed subprescheme of
$\operatorname{Spec}(A)$ defined by $\mathfrak J$.
\end{enumerate}
\end{remark}

\EGAIVErrataThirtyTwoThroughFiftyThree
% END NATIVE INPUT source_aligned/ega4-errata-addenda-list3

\clearpage
\phantomsection
\addcontentsline{toc}{section}{Consolidated bibliography}
\section*{Consolidated bibliography}
\nocite{*}
\bibliography{the}
\bibliographystyle{amsalpha}

\end{document}
